\documentclass[UTF8,12pt,a4paper,oneside]{ctexbook}
\usepackage[margin=25mm]{geometry}
\usepackage{fontspec}
\setmainfont{Times New Roman}
\setsansfont{Times New Roman}
\setmonofont{DejaVu Sans Mono}
\setCJKmainfont[
  BoldFont=FandolSong-Bold.otf,
  ItalicFont=FandolKai-Regular.otf
]{FandolSong-Regular.otf}
\setCJKsansfont[
  BoldFont=FandolHei-Bold.otf
]{FandolHei-Regular.otf}
\setCJKmonofont[
  BoldFont=FandolHei-Bold.otf
]{FandolHei-Regular.otf}
\usepackage{amsmath,amssymb,mathtools}
\usepackage{unicode-math}
\setmathfont{TeX Gyre Termes Math}
\usepackage{newunicodechar}
\newunicodechar{∅}{\ensuremath{\varnothing}}
\newunicodechar{ℰ}{\ensuremath{\mathcal{E}}}
\newunicodechar{ℜ}{\ensuremath{\mathfrak{R}}}
\newunicodechar{⊗}{\ensuremath{\otimes}}
\newunicodechar{∈}{\ensuremath{\in}}
\newunicodechar{∘}{\ensuremath{\circ}}
\newunicodechar{↦}{\ensuremath{\mapsto}}
\newunicodechar{∇}{\ensuremath{\nabla}}
\newunicodechar{⇒}{\ensuremath{\Rightarrow}}
\newunicodechar{∨}{\ensuremath{\vee}}
\newunicodechar{∪}{\ensuremath{\cup}}
\newunicodechar{⊕}{\ensuremath{\oplus}}
\newunicodechar{⇔}{\ensuremath{\Leftrightarrow}}
\newunicodechar{⊂}{\ensuremath{\subset}}
\newunicodechar{∧}{\ensuremath{\wedge}}
\newunicodechar{⊤}{\ensuremath{\top}}
\usepackage{graphicx}
\usepackage{longtable,booktabs,array}
\usepackage{xcolor,framed,fancyvrb}
\usepackage{setspace}
\usepackage{enumitem}
\usepackage{microtype}
\usepackage[hidelinks,unicode,pdfusetitle,bookmarksnumbered]{hyperref}
\usepackage[open,openlevel=2]{bookmark}
\doublespacing
\setlength{\parindent}{1.5em}
\setlength{\parskip}{0pt}
\setcounter{tocdepth}{2}
\setcounter{secnumdepth}{2}
\setlist{itemsep=0.25\baselineskip,topsep=0.25\baselineskip}
\providecommand{\tightlist}{\setlength{\itemsep}{0pt}\setlength{\parskip}{0pt}}
\providecommand{\pandocbounded}[1]{#1}
\newcommand{\specialchapter}[1]{%
  \chapter*{#1}%
  \phantomsection
  \addcontentsline{toc}{chapter}{#1}%
}
% 印刷的“第 k 节习题”分组：由计数器生成，不手写序号。
\newcounter{exercisepar}
\newcommand{\exercisesection}[1]{%
  \refstepcounter{exercisepar}%
  \section*{第\theexercisepar 节习题}%
  \phantomsection
  \addcontentsline{toc}{section}{第\theexercisepar 节习题}%
}
% 历史注记的编号条目（该部分没有 § 层级）。
\newcounter{histitem}
\newcommand{\histitem}[1]{%
  \refstepcounter{histitem}%
  \section*{\thehistitem.\quad #1}%
  \phantomsection
  \addcontentsline{toc}{section}{#1}%
}
% 术语索引中的字母分组：仅排版用途，不进入目录。
\newcommand{\lettergroup}[1]{\par\bigskip\noindent\textbf{#1}\par\nobreak\medskip}
\renewcommand{\thechapter}{\Roman{chapter}}
\definecolor{shadecolor}{RGB}{248,248,248}
\newenvironment{Shaded}{\begin{snugshade}}{\end{snugshade}}
\newenvironment{Highlighting}[1][]{}{}
\newcommand{\NormalTok}[1]{#1\par}
\makeatletter
\def\maxwidth{\ifdim\Gin@nat@width>\linewidth\linewidth\else\Gin@nat@width\fi}
\def\maxheight{\ifdim\Gin@nat@height>0.8\textheight 0.8\textheight\else\Gin@nat@height\fi}
\makeatother
\setkeys{Gin}{width=\maxwidth,height=\maxheight,keepaspectratio}
\graphicspath{{images/}}
\emergencystretch=3em
\pagestyle{plain}
\title{谱理论（第3--5章）}
\author{尼古拉·布尔巴基}
\date{}
\begin{document}
\pagenumbering{gobble}
\maketitle
\clearpage
\frontmatter
\tableofcontents
\mainmatter
\setcounter{chapter}{2}
\specialchapter{引言}

在本书第 III 章至第 V 章中，我们详细考察实或复拓扑向量空间的自同态代数。在有限维空间的情形，这一代数的元素逐个来看的主要性质都可由若尔当分解（A, VII, p. 42, § 5, 第 9 目）导出。我们在此研究它们向无穷维空间的推广，然而这推广要更为精细。

在第 III 章中，我们在非常一般的框架下考察拓扑向量空间之间的紧线性映射。它们的谱性质与有限维情形最为接近。例如，在巴拿赫空间的特殊情形，紧自同态的谱中的非零元素都是具有有限谱重数的广义特征值。

在希尔伯特空间的框架下，会出现一种易于处理的分类和结构。这正是第 IV 章的对象。该章从希尔伯特空间中紧算子的研究开始：它们借助奇异值的描述，以及正规自同态情形下的对角化，都完全类似于有限维空间上已知的性质。更一般地，谱定理给出了希尔伯特空间正规自同态的简单模型。

某些重要应用，例如涉及微分算子的应用，但也涉及量子力学的应用，要求一种更精微的理论：即定义在希尔伯特空间的稠密子空间上的线性映射---未必连续---的理论，我们把这类映射称为部分算子。在这一框架下，自伴算子或正规算子的定义本身就必须极其精确，否则便会产生重大的谬误。一旦掌握了这些定义，就会出现一种与连续自同态情形相类似的分类；而对有界或无界的普遍可测函数的函数演算，则在连续正规自同态与正规部分算子之间建立起紧密的联系。

顺带地，我们还会谈到带迹的正规自同态，特别是希尔伯特空间 \(L^{2}(X,\mu)\) 中带积分核的积分自同态。人们熟知，在自守形式与表示的理论中，断言这类自同态（当其存在时）的迹等于核沿对角线的积分的那条公式具有至关重要的意义。我们还用一节的篇幅叙述 \(R^{n}\) 中广义函数与缓增广义函数的基本定义和性质。这些概念不仅给出了欧几里得傅里叶理论诸基本原理（参见 II, § 1）的最方便的表述，而且还使人们能够确定许多重要微分算子的定义域。

最后，在第 V 章中，我们把希尔伯特空间自同态的研究与群的酉表示联系起来。在若干初等的一般性考虑之后---其中舒尔引理是最重要的命题---我们研究局部紧群的情形。然后我们考察正定函数，并给出它们的最初应用。这些应用之一是盖尔范德--奈马克--塞加尔定理，它断言每个星代数都同构于某个希尔伯特空间的自同态代数的一个闭子代数。另一些应用则涉及局部紧群的不可约酉表示所成的空间的基本性质，特别是盖尔范德--赖科夫定理，它确立了这些表示分离点。在本章最后一节中，我们证明关于紧群表示的基本结果，它们曾在《群与李代数》第 IX 章中为我们所用。

A. BOREL, Representations of locally compact groups, Lecture Notes in Mathematics 276, Springer, 1972.

M. REED 与 B. SIMON, Methods of Mathematical Physics, Academic Press, Vol. 1, 1980 与 Vol. 2, 1975.

K. SCHMÜDGEN, Unbounded Self-adjoint Operators on Hilbert Space, Graduate Texts in Mathematics 265, Springer, 2012.

在前面的各章与各书中，对“TS, to appear”的下列引用应作如下修改：

EVT, V, p. 35, 第 1 行。应将“TS, IV, § 6”改为“(TS, IV, p. 160, remark)”。

INT, IX, p. 94, § 6, 第 12 目。应将“One can prove a converse \ldots{} bounded”改为“The converse of prop. 11 is a particular case of TS, V, p. 455, th. 5.”

LIE, IX, p. 61, § 6, 第 4 目, 第 6 行。应将“TS, to appear”改为“TS, III, p. 36, prop. 9”。

LIE, IX, p. 63, § 6, 第 4 目, 第 6 行。应将“TS, to appear”改为“TS, III, p. 36, rem.”。

LIE, IX, p. 66, § 7, 注 1。应将“to a chapter to appear”改为“to paragraph 4 of Chapter V”。

LIE, IX, p. 69, § 7, 第 2 目。应将“then according to TS”改为“then (TS, V, p. 464, prop. 6)”。

LIE, IX, p. 71, § 7, 第 3 目。应将“which, according to TS, is an isomorphism”改为“which, according to TS, V, p. 469, cor. 2, is an isomorphism.”

LIE, IX, p. 72, § 7, 第 4 目。应将“(TS)”改为“(TS, V, p. 459, cor. 3)”。

LIE, IX, p. 74, § 7, 第 4 目。在“of the space of central representative functions on G”之后加上“(TS, V, p. 469, cor. 2, a)”；在“of the space ZL²(G) of central square-integrable functions on G”之后加上“(TS, V, p. 468, prop. 9).”

LIE, IX, p. 78, § 8, 第 1 目。应将“In this number, one recalls definitions and results from TS”改为“One recalls definitions and results from TS, V, p. 386, no. 8; p. 392, no. 10; p. 470, no. 7.”

LIE, IX, p. 84, § 8, 第 3 目。应将“Recall (TS) the formulas”改为“One has (TS, V, p. 467, formulas (1) and (2)) the formulas”。

LIE, IX, p. 89, § 9, 第 1 目。应将“TS, III, §2, no. 7, prop. 16”改为“TS, III, p. 67, prop. 10”。

LIE, IX, p. 90, § 9, 第 1 目。应将“TS, III, §2, no. 6”改为“TS, III, p. 61, no. 4”。

LIE, IX, p. 91, § 9, 第 2 目。应将“TS, to appear”改为“TS, III, p. 36, prop. 9 and TS, V, p. 464, prop. 7”。

LIE, IX, p. 92, 注 1。应将“(TS, to appear)”改为“(TS, V, p. 379, def. 3)”。

LIE, IX, p. 99, 附录 1, 第 1 目。应将“(TS, to appear)”改为“(TS, V, p. 464, remark 4)”。

LIE, IX, p. 128, § 9 习题 3 (b)。应将“(TS, to appear)”改为“(Use the Peter--Weyl theorem, TS, V, p. 462, th. 1)”。

TS, I, p. 103, 例 6。应将“V, to appear”改为“V, p. 442, th. 2.”

TS, I, p. 173, 习题 5, d)。应将“III, to appear”改为“III, p. 82, § 6.”

TS, II, p. 237, 第 9 目。应将“IV, to appear”改为“IV, p. 196, § 3”。

TS, II, p. 251, 第 1 目。应将“V, to appear”改为“V, p. 425, prop. 10”。

\chapter{紧线性映射与扰动}
\setcounter{exercisepar}{0}

在本章中，所考虑的向量空间都是等于 R 或 C 的域 K 上的向量空间。本书中所作的关于 EVT, II 的引用一般针对 K = R 的情形；关于 K = C 的情形，参见 EVT, II, p. 64 至 68。

半赋范空间是指赋有一个半范数 p 以及由 p 定义的拓扑的向量空间 E（参见 EVT, II, p. 2）。E 的单位球，或 p 的单位球，是 E 中满足 \(p(x) \leqslant 1\) 的元素 x 所成的集。

对每个向量空间 E，我们用 \(1_{E}\) 表示 E 的恒等映射。若 E, F 和 G 是向量空间，且 u: \(E \rightarrow F\)、\(v: F \rightarrow G\) 是线性映射，我们有时用 vu 表示从 E 到 G 的线性映射 \(v \circ u\)。

设 E 和 F 是拓扑向量空间。从 E 到 F 的连续线性映射也称为态射，而当 F = E 时称为自同态。从 E 到 F 的双射线性映射，若它连续且其逆也连续，则称为同构，而当 F = E 时称为自同构。当 E 和 F 是半赋范空间时，自同态、同构和自同构这几个概念都是就底层的拓扑向量空间结构而言的。

给定拓扑向量空间 E 和 F，我们用 \(\mathcal{L}^{\mathrm{f}}(\mathrm{E};\mathrm{F})\) 表示从 E 到 F 的有限秩连续线性映射所成的向量空间。我们也用 \(\mathcal{L}^{\mathrm{f}}(\mathrm{E})\) 表示向量空间 \(\mathcal{L}^{\mathrm{f}}(\mathrm{E};\mathrm{E})\)。

我们回顾，从 E 到 F 的连续线性映射 u 称为严格的，如果它通过过渡到商空间诱导出从 E/ Ker(u) 到 u(E) 上的一个同构（TG, III, p. 16, def. 1）；这等价于说，E 中 0 的每个邻域的像是 u(E) 中 0 的一个邻域（TG, III, p. 16, prop. 24）。

\section{紧线性映射}

\subsection{紧线性映射}

定义 1. --- 设 E 是拓扑向量空间，F 是分离的拓扑向量空间。从 E 到 F 的线性映射 u 称为紧的，如果存在 E 中 0 的一个邻域 V，使得 \(u(V)\) 是 F 的一个相对紧子集。

我们用 \(\mathcal{L}^{\mathrm{c}}(\mathrm{E};\mathrm{F})\) 表示从 E 到 F 的紧线性映射所成的集；我们也用 \(\mathcal{L}^{\mathrm{c}}(\mathrm{F})\) 表示集 \(\mathcal{L}^{\mathrm{c}}(\mathrm{F};\mathrm{F})\)。

注记. --- 1) E 中 0 的一个邻域吸收 E 的每个有界子集；因此，E 的有界子集在从 E 到 F 的紧线性映射下的像是 F 的一个相对紧子集。

\begin{enumerate}
\def\labelenumi{\arabic{enumi})}
\setcounter{enumi}{1}
\item
  设 E 是半赋范空间，B 是 E 的单位球。为使从 E 到一个分离拓扑向量空间 F 的线性映射 u 是紧的，必须且只需 \(u(\mathrm{B})\) 是 F 的一个相对紧子集。
\item
  设 F 是分离的拓扑向量空间。为使 F 的向量子空间 E 是有限维的，必须且只需 E 到 F 中的典范单射是紧的（EVT, I, p. 15, th. 3）。
\item
  对域 R 上的每个拓扑向量空间 E，我们用 \(\mathrm{E}_{(\mathbf{C})}\) 表示 E 的复化拓扑向量空间（EVT, II, p. 65）。我们把 E 等同于 \(\mathrm{E}_{(\mathbf{C})}\) 的一个实拓扑向量子空间（经由映射 \(x \mapsto 1 \otimes x\)）；于是 \(\mathrm{E}_{(\mathbf{C})}\) 的底层实拓扑向量空间是 E 与 iE 的拓扑直和。
\end{enumerate}

设 \(u\) 是从 E 到 \(\mathbf{R}\) 上的一个分离拓扑向量空间 F 的线性映射。用 \(u_{(\mathbf{C})}\) 表示 \(\mathbf{C}\)-线性映射（它从 \(\mathrm{E}_{(\mathbf{C})}\) 到 \(\mathrm{F}_{(\mathbf{C})}\) 中，并且拓广 \(u\)）。形如 \(V + iV\) 的集合（其中 \(V\) 是 E 中 0 的一个邻域）构成 \(\mathrm{E}_{(\mathbf{C})}\) 中 0 的一个基本邻域系。为使集合 \(u_{(\mathbf{C})}(V + iV) = u(V) + iu(V)\) 在 \(\mathrm{F}_{(\mathbf{C})}\) 中相对紧，必须且只需 \(u(V)\) 在 F 中相对紧。因此，为使 \(\mathbf{C}\)-线性映射 \(u_{(\mathbf{C})}\) 是紧的，必须且只需 \(\mathbf{R}\)-线性映射 \(u\) 是紧的。

出于同样的思路，当 \(v\) 是从 \(\mathrm{E}_{(\mathbf{C})}\) 到一个 C-拓扑向量空间 G 的连续 C-线性映射时，为使 \(v\) 是紧的，必须且只需它在 E 上的限制是紧的。

\begin{enumerate}
\def\labelenumi{\arabic{enumi})}
\setcounter{enumi}{4}
\item
  设 E 是拓扑向量空间，F 是分离的拓扑向量空间。设 \(F_{1}\) 是 F 的一个闭向量子空间，j 是 \(F_{1}\) 到 F 中的典范单射。为使从 E 到 \(F_{1}\) 的线性映射 u 是紧的，必须且只需从 E 到 F 的线性映射 \(j \circ u\) 是紧的。
\item
  设 \((\mathrm{G}_{i})_{i\in\mathrm{I}}\) 是一族分离的拓扑向量空间，且对每个 \(i\in I\)，设 \(u_{i}\) 是从拓扑向量空间 E 到 \(G_{i}\) 的一个紧线性映射。若集合 I 有限，或者 E 是半赋范空间，则线性映射 \(x\mapsto(u_{i}(x))_{i\in\mathrm{I}}\)（从 E 到 \(\prod_{i\in I}G_{i}\) 中）是紧的（注记 2 及 TG, I, p. 63, Th. 3）。
\end{enumerate}

命题 1. --- 设 E 是拓扑向量空间，F 是分离的拓扑向量空间。

\begin{enumerate}
\def\labelenumi{\alph{enumi})}
\item
  从 E 到 F 的每个紧线性映射都是连续的。
\item
  从 E 到 F 的每个有限秩连续线性映射都是紧的；
\item
  从 E 到 F 的紧线性映射所成的集 \(\mathcal{L}^{\mathrm{c}}(\mathrm{E};\mathrm{F})\) 是 \(\mathcal{L}(\mathrm{E};\mathrm{F})\) 的一个向量子空间。
\end{enumerate}

我们来证明 a)。设 \(v: E \rightarrow F\) 是紧线性映射，U 是 F 中 0 的一个邻域。取 E 中 0 的一个邻域 V，使得 \(v(V)\) 在 F 中相对紧。于是 \(v(V)\) 有界，因此存在实数 \(\lambda > 0\)，使得集合 \(v(\lambda V)\) 包含于 U；从而 v 连续。

设 \(u \colon \mathrm{E} \to \mathrm{F}\) 是有限秩连续线性映射。由于 \(u(\mathrm{E})\)（即 \(u\) 的像）是分离的且是有限维的，故存在 \(u(\mathrm{E})\) 中 0 的一个紧邻域 A。集合 \(\mathrm{V} = u^{-1}(\mathrm{A})\) 是 E 中 0 的一个邻域，并且 \(u(\mathrm{V}) \subset \mathrm{A}\)，所以线性映射 \(u\) 是紧的。这就证明了断言 \(b\)）。

设 \(u_{1}\) 和 \(u_{2}\) 是从 E 到 F 的紧线性映射。取 E 中 0 的邻域 \(V_{1}\) 和 \(V_{2}\)，使得集合 \(u_{1}(V_{1})\) 与 \(u_{2}(V_{2})\) 都在 F 中相对紧。令 \(V = V_{1} \cap V_{2}\)，\(u = u_{1} + u_{2}\)。于是 \(u(V)\) 包含于 \(u_{1}(V_{1}) + u_{2}(V_{2})\)，从而在 F 中相对紧，故线性映射 u 是紧的。\(\mathcal{L}^{c}(E;F)\) 在位似变换下稳定是显然的。这就证明了 \(\mathcal{L}^{c}(E;F)\) 是 \(\mathcal{L}(E;F)\) 的一个向量子空间，从而得到断言 c)。

命题 2. --- 设 E 是半赋范空间，F 是分离的拓扑向量空间。假设 F 的每个有界闭子集都是完备的。那么向量子空间 \(\mathcal{L}^{\mathrm{c}}(\mathrm{E};\mathrm{F})\)（\(\mathcal{L}(\mathrm{E};\mathrm{F})\) 的）对有界收敛拓扑是闭的。

设 \(u\) 是 \(\mathcal{L}(\mathrm{E};\mathrm{F})\) 中对有界收敛拓扑附着于 \(\mathcal{L}^{\mathrm{c}}(\mathrm{E};\mathrm{F})\) 的一个元素。用 B 表示 E 的单位球，V 表示 F 中 0 的一个邻域。取 F 中 0 的一个邻域 W，使得 \(\mathrm{W} + \mathrm{W} \subset \mathrm{V}\)。按假设，存在元素 \(v\)（它属于 \(\mathcal{L}^{\mathrm{c}}(\mathrm{E};\mathrm{F})\)），使得 \((u - v)(\mathrm{B}) \subset \mathrm{W}\)。集合 \(v(\mathrm{B})\) 在 F 中相对紧，故存在 F 的一个有限子集 M，使得 \(v(\mathrm{B}) \subset \mathrm{M} + \mathrm{W}\)，从而 \(u(\mathrm{B}) \subset \mathrm{M} + \mathrm{V}\)。由于这对每个 V 都成立，集合 \(u(\mathrm{B})\) 在 F 中预紧，因而其闭包也预紧（TG, II, prop. 1, p. 30）。后者是闭的且有界，故由对 F 所作的假设，它是完备的，从而是紧的（TG, II, cor., p. 30）。由此推出 \(u(\mathrm{B})\) 相对紧。这就证明了 \(u\) 属于 \(\mathcal{L}^{\mathrm{c}}(\mathrm{E};\mathrm{F})\)。

若 F 是局部凸的，则命题的假设意味着 F 是拟完备的（EVT, III, p. 8, def. 6）。

推论. --- 设 E 和 F 是巴拿赫空间。集合 \(\mathcal{L}^{\mathrm{c}}(\mathrm{E};\mathrm{F})\) 是巴拿赫空间 \(\mathcal{L}(\mathrm{E};\mathrm{F})\) 的一个闭向量子空间。它包含从 E 到 F 的有限秩连续线性映射所成的空间 \(\mathcal{L}^{\mathrm{f}}(\mathrm{E};\mathrm{F})\) 的闭包。

这由命题 1 和命题 2 得出。

\pandocbounded{\includegraphics[keepaspectratio]{images/part001_746f5095.jpg}}

存在巴拿赫空间 E 和 F，使得 \(\mathcal{L}^{\mathrm{f}}(\mathrm{E};\mathrm{F})\) 在 \(\mathcal{L}^{\mathrm{c}}(\mathrm{E};\mathrm{F})\) 中不稠密（参见 III, p. 16, 注记 6 及 III, p. 19, 定理 4, b)）。

命题 3. --- 设 \(\mathrm{E}_{1}\)、E、F、\(\mathrm{F}_{1}\) 是拓扑向量空间，其中 F 和 \(\mathrm{F}_{1}\) 假设为分离的。设 v: \(\mathrm{E}_{1}\) → E、u: E → F、w: F → \(\mathrm{F}_{1}\) 是线性映射。若 v 和 w 连续而 u 紧，则 w ∘ u ∘ v 紧。

按假设，存在 E 中 0 的一个邻域 V，使得 \(u(\mathrm{V})\) 在 F 中相对紧。令 \(\mathrm{U} = v^{-1}(\mathrm{V})\)。于是 U 是 \(\mathrm{E}_{1}\) 中 0 的一个邻域，且它在 \(w \circ u \circ v\) 下的像包含于 \(w(u(\mathrm{V}))\)，从而在 \(\mathrm{F}_{1}\) 中相对紧。因此，线性映射 \(w \circ u \circ v\) 是紧的。

设 E 是分离的拓扑向量空间；由命题 1 和命题 3，\(\mathcal{L}^{\mathrm{c}}(\mathrm{E})\) 是代数 \(\mathcal{L}(\mathrm{E})\) 的一个双边理想。当 E 是巴拿赫空间时，\(\mathcal{L}(\mathrm{E})\) 是巴拿赫代数，而 \(\mathcal{L}^{\mathrm{c}}(\mathrm{E})\) 是 \(\mathcal{L}(\mathrm{E})\) 的一个闭双边理想（命题 2 的推论）。若 E 是无穷维的，则它是一个真理想（参见 p. 2, 注记 3）。

推论. — 设 E 是希尔伯特空间。空间 \(\mathcal{L}^{\mathrm{c}}(\mathrm{E})\) 是 \(\mathcal{L}(\mathrm{E})\) 的一个闭的自伴双边理想。特别地，它是一个星代数。

事实上，空间 \(\mathcal{L}^{c}(E)\) 是 \(\mathcal{L}(E)\) 的一个闭双边理想，从而是自伴的（I, p. 122, 引理 15）。

命题 4. — 设 E 是拓扑向量空间，\(\widehat{\mathrm{E}}\) 是它的分离完备化，\(j\) 是 E 到 \(\widehat{\mathrm{E}}\) 中的典范映射。设 \(u\) 是从 E 到一个分离拓扑向量空间 F 的紧线性映射。则存在唯一的紧线性映射 \(v\)（从 \(\widehat{\mathrm{E}}\) 到 F 中），使得 \(u = v \circ j\)。

把 F 等同于 \(\widehat{\mathrm{F}}\) 的一个拓扑向量子空间，并用 \(\widehat{u}\colon \widehat{\mathrm{E}}\to \widehat{\mathrm{F}}\) 表示使 \(\widehat{u}\circ j\) 在 E 上与 \(u\) 重合的唯一连续线性映射。由于 \(u\) 紧，存在 E 中 0 的一个邻域 V 和 F 的一个紧子集 A，使得 \(u(\mathrm{V})\subset \mathrm{A}\)。我们有 \(\widehat{u}(j(\mathrm{V}))\subset \mathrm{A}\)，从而 \(\widehat{u} (\overline{j(\mathrm{V})})\subset \mathrm{A}\)。而 \(\overline{j(\mathrm{V})}\) 是 \(\widehat{\mathrm{E}}\) 中 0 的一个邻域（TG, III, p. 24, prop. 7）。因此，\(\widehat{u}\) 的像包含于 F，且连续线性映射 \(v\colon \widehat{\mathrm{E}}\rightarrow \mathrm{F}\)（由 \(\widehat{u}\) 过渡到子空间得到的）是紧的。由于 \(j(\mathrm{E})\) 在 \(\widehat{\mathrm{E}}\) 中稠密，映射 \(v\) 是从 \(\widehat{\mathrm{E}}\) 到 F 的满足 \(u = v\circ j\) 的唯一连续线性映射。

命题 5. --- 设 E 是局部凸空间，F 是分离的拓扑向量空间，u 是从 E 到 F 的紧线性映射。则存在一个巴拿赫空间 G、一个从 E 到 G 的连续线性映射 v 和一个从 G 到 F 的紧线性映射 w，使得 \(u = w \circ v\)。

取 E 中 0 的一个邻域 V，使得 \(u(\mathrm{V})\) 在 F 中相对紧。设 \(p\) 是 E 上的一个连续半范数，使得 V 包含 \(p\) 的单位球（EVT, II, p. 26, cor.）。用 \(\mathrm{E}_p\) 表示给 E 赋以半范数 \(p\) 所得到的半赋范空间。映射 \(u\) 是从 \(\mathrm{E}_p\) 到 F 的紧线性映射。\(\mathrm{E}_p\) 的分离完备化 G 是一个巴拿赫空间。用 \(v\) 表示从 \(\mathrm{E}_p\) 到 G 的典范线性映射。由命题 4，存在一个紧线性映射 \(w: \mathrm{G} \to \mathrm{F}\)，使得 \(u = w \circ v\)。另一方面，v 是从 E 到 G 的连续映射，因为 E 的拓扑比 \(\mathrm{E}_p\) 的拓扑更细。

\subsection{紧线性映射与弱拓扑}

设 E 是局部凸空间，\(E'\) 是它的对偶。回顾，E 上的拓扑 \(\sigma(\mathrm{E},\mathrm{E}')\) 称为 E 上的弱拓扑（TVS, IV, p. 4）。在本节中，\(E_{\sigma}\) 表示赋有弱拓扑的空间 E。

命题 6. --- 设 E 是局部凸空间，F 是分离的局部凸空间，\(u: \mathrm{E} \to \mathrm{F}\) 是紧线性映射，B 是 E 的一个有界子集。\(u\) 在 B 上的限制是从赋有 \(\sigma(\mathrm{E}, \mathrm{E}')\) 诱导拓扑的集合 B 到空间 F 的连续映射。

映射 \(u\) 是从 E 到 F 的连续线性映射，也是从 \(\mathrm{E}_{\sigma}\) 到 \(\mathrm{F}_{\sigma}\) 的连续线性映射。因此，它在 B 上的限制是从赋有 \(\sigma(\mathrm{E}, \mathrm{E}')\) 诱导拓扑的集合 B 到空间 \(\mathrm{F}_{\sigma}\) 的连续映射。而集合 \(u(\mathrm{B})\) 包含于 F 的一个紧子集 C（III, p. 2, 注记 1），且空间 \(\mathrm{F}_{\sigma}\) 是豪斯多夫的，故由 F 与由 \(\mathrm{F}_{\sigma}\) 在 C 上诱导的拓扑重合。命题得证。

推论. --- 设 \((x_{n})_{n\in\mathbf{N}}\) 是 E 的点的一个序列，它对弱拓扑收敛于 E 的一个点 x。则序列 \((u(x_{n}))_{n\in\mathbf{N}}\) 在 F 中收敛于 \(u(x)\)。

事实上，由点 \(x\) 与序列 \((x_{n})_{n\in \mathbf{N}}\) 的诸点所组成的集是 E 的一个有界子集（EVT, III, p. 3, 命题 2 之推论）。

命题 7. --- 设 E 是半赋范空间，E' 是它的对偶，B 是 E' 的单位球，u 是从 E' 到一个分离局部凸空间 F 的线性映射。若 u 在赋有 \(\sigma(\mathrm{E}',\mathrm{E})\) 诱导拓扑的 B 上的限制是连续的，则 \(u(\mathrm{B})\) 是 F 的一个紧子集，且 u 是从 E 的强对偶 \(\mathrm{E}_b'\) 到 F 的紧线性映射。

事实上，赋有 \(\sigma(\mathrm{E}^{\prime},\mathrm{E})\) 诱导拓扑的集合 B 是紧的（EVT, III, p. 17, cor. 2）。

推论 1. --- 设 E 是可数型的半赋范空间，E' 是它的对偶，B 是 E' 的单位球。设 u 是从 E' 到一个分离局部凸空间 F 的线性映射。假设对 B 中元素的每个序列 \((x_{n})_{n\in\mathbf{N}}\)（它对 \(\sigma(\mathrm{E}',\mathrm{E})\) 收敛于 0），序列 \((u(x_{n}))_{n\in\mathbf{N}}\) 都在 F 中收敛于 0。那么，\(u(\mathrm{B})\) 是 F 的一个紧子集，且 u 是从 E 的强对偶 \(E_{b}'\) 到 F 的紧线性映射。

事实上，\(\sigma(\mathrm{E}^{\prime},\mathrm{E})\) 在 B 上诱导的拓扑是可度量化的（EVT, III, p. 19, cor. 2）。因此，推论的假设蕴含：\(u\) 在 B 上的限制（当 B 赋有拓扑 \(\sigma(\mathrm{E}^{\prime},\mathrm{E})\) 时）是从 B 到 F 的连续映射，于是推论由前一命题得出。

推论 2. --- 设 E 是半赋范空间，S 是 E 的预紧子集所成的一个集，其并在 E 中稠密。用 \(E'_{b}\) 表示 E 的强对偶，用 \(E'_{S}\) 表示赋有 S-拓扑的 E 的对偶。空间 \(E'_{S}\) 是分离的，且从 \(E'_{b}\) 到 \(E'_{S}\) 的恒等映射是紧的。

依照 TG, X, p. 8, prop. 7，空间 \(E_{S}^{\prime}\) 是分离的；在 \(E^{\prime}\) 的单位球上，S-拓扑与弱拓扑 \(\sigma(\mathrm{E}^{\prime},\mathrm{E})\) 重合（EVT, III, p. 17, prop. 5），由此即得推论。

例. --- 设 E 是半赋范空间。从 \(E_{b}^{\prime}\) 到赋有弱拓扑、紧收敛拓扑或预紧收敛拓扑的空间 \(E^{\prime}\) 的恒等映射是紧的。

命题 8. --- 设 E 是自反的巴拿赫空间，B 是它的单位球，F 是分离的局部凸空间，u: E → F 是线性映射。下列条件等价：

\begin{enumerate}
\def\labelenumi{(\roman{enumi})}
\item
  线性映射 u 是紧的；
\item
  集合 \(u(\mathrm{B})\) 是 F 的一个紧子集；
\item
  u 在 B 上的限制是从赋有 \(\sigma(\mathrm{E},\mathrm{E}')\) 诱导拓扑的集合 B 到 F 的连续映射；
\item
  对 B 中点的任何序列 \((x_{n})_{n\in \mathbf{N}}\)（它对拓扑 \(\sigma (\mathrm{E},\mathrm{E}^{\prime})\) 收敛于 0），序列 \((u(x_n))_{n\in \mathbf{N}}\) 在 F 中收敛于 0；
\item
  从 \(u(\mathrm{B})\) 的点的任何无穷序列中都能取出一个在 F 中收敛的子序列；
\item
  \(u(\mathrm{B})\) 的点的任何无穷序列在 F 中都有聚点。
\end{enumerate}

由于 E 是自反空间，B 是 \(E_{\sigma}\) 的一个紧子集（EVT, IV, p. 17, prop. 6），故 (iii) 蕴含 (ii)。按定义，条件 (ii) 蕴含 (i)，而由命题 6，条件 (i) 蕴含 (iii)。同样初等的是 (iii) 蕴含 (iv) 且 (v) 蕴含 (vi)。

我们现在证明 (iv) 蕴含 (v)。设 \((x_{n})\) 是 B 的点的一个无穷序列。由于 B 是 \(E_{\sigma}\) 的紧子集，由什穆利安定理（EVT, IV, p. 36, th. 2），存在一个序列 \((y_{n})\)（从序列 \((x_{n})\) 中取出），它在 \(E_{\sigma}\) 中收敛于一个极限 y。序列 \((y_{n}-y)\) 在 \(E_{\sigma}\) 中有界，从而在 E 中有界（EVT, IV, p. 1, prop. 1）；因此它包含于 B 的一个位似集合。若条件 (iv) 成立，则序列 \((u(y_{n}-y))\) 在 F 中收敛于 0，而 \((u(y_{n}))\) 作为 \((u(x_{n}))\) 的一个子序列收敛于 \(u(y)\)。这就证明了 (iv) 蕴含 (v)。

最后，我们证明 (vi) 蕴含 (ii)。设 \(j: \mathrm{F} \to \widehat{\mathrm{F}}\) 是 F 到其完备化中的典范单射。在假设 (vi) 之下，\(u(\mathrm{B})\) 是 F 的一个预紧子集（EVT, IV, p. 32, prop. 1），因此 \(j(u(\mathrm{B}))\) 是 \(\widehat{\mathrm{F}}\) 的一个相对紧子集（TG, II, 第 2 目, p. 29），且 \(j \circ u\) 是紧线性映射。由已经证明的条件 (i) 与 (ii) 的等价性，\(j(u(\mathrm{B}))\) 是 \(\widehat{\mathrm{F}}\) 的一个紧子集。因此，\(u(\mathrm{B})\) 是 F 的一个紧子集。

\subsection{转置}

设 E 是局部凸空间，\(E^{\prime}\) 是它的对偶。回顾，我们用 \(E_{b}^{\prime}\) 和 \(E_{c}^{\prime}\) 分别表示给向量空间 \(E^{\prime}\) 赋以有界收敛拓扑和紧收敛拓扑所得到的局部凸空间（EVT, III, p. 14）。空间 \(E_{b}^{\prime}\) 也称为 E 的强对偶（同前引处, 例 4）。

命题 9. --- 设 E 是局部凸空间，F 是分离的局部凸空间，u 是从 E 到 F 的连续线性映射。

\begin{enumerate}
\def\labelenumi{\alph{enumi})}
\item
  若映射 \(u\) 是紧的，则它的转置 \(^{t}u\) 是从 \(F_c'\) 到 \(E_c'\) 的紧线性映射，且是从 \(F_c'\) 到 \(E_b'\) 的连续线性映射；
\item
  假设空间 E 是半赋范空间且空间 F 是拟完备的。若转置 \({}^{t}u\) 是从 \(F_{c}^{\prime}\) 到 \(E_{b}^{\prime}\) 的连续映射，则线性映射 u 是紧的。
\end{enumerate}

首先假设线性映射 \(u\) 是紧的。于是存在 E 中 0 的一个邻域 V，它在 \(u\) 下的像包含于 F 的一个紧子集 C。按照定义（EVT, II, p. 47, def. 2 及 p. 68, def. 1），C 的极集 \(\mathrm{C}^{\circ}\) 是 \(\mathrm{F}_c^\prime\) 中 0 的一个邻域，并且 \({}^t u(\mathrm{C}^\circ) \subset \mathrm{V}^\circ\)。而 \(\mathrm{V}^\circ\) 是 \(\mathrm{E}'\) 的一个等度连续且对弱拓扑为闭的子集，从而在 \(\mathrm{E}_c^\prime\) 中紧（EVT, III, p. 17, cor. 2 及 prop. 5）。这就证明了 \({}^t u\) 是从 \(\mathrm{F}_c^\prime\) 到 \(\mathrm{E}_c^\prime\) 的紧线性映射。

我们保留前面的假设，设 U 是 \(E_{b}^{\prime}\) 中 0 的一个邻域。它包含 E 的某个有界子集 A 的极集 \(A^{\circ}\)。由于线性映射 u 紧，集合 \(u(\mathrm{A})\) 在 F 中相对紧（III, p. 2, 注记 1），且 \((^{t}u)^{-1}(\mathrm{A}^{\circ}) = u(\mathrm{A})^{\circ}\) 是 \(F_{c}^{\prime}\) 中 0 的一个邻域。这就证明了 \(^{t}u\) 是从 \(F_{c}^{\prime}\) 到 \(E_{b}^{\prime}\) 的连续线性映射。

现在置于 \(b\) 的假设之下，并用 B 表示 E 的单位球。B 的极集 \(\mathrm{B}^{\circ}\) 是 \(\mathrm{E}_b'\) 的单位球，并且 \((^{t}u)^{-1}(\mathrm{B}^{\circ}) = u(\mathrm{B})^{\circ}\)。若 \(^{t}u\) 是从 \(\mathrm{F}_c'\) 到 \(\mathrm{E}_b'\) 的连续映射，则集合 \(u(\mathrm{B})^{\circ}\) 是 \(\mathrm{F}_c'\) 中 0 的一个邻域；于是它包含 F 的某个紧子集 C 的极集 \(\mathrm{C}^{\circ}\)。由双极定理（EVT, II, p. 49, cor. 3），集合 \(u(\mathrm{B})\) 包含于 \(\mathrm{C} \cup \{0\}\) 的闭凸包；由于空间 F 拟完备（EVT, III, p. 8），后者是紧的。这就证明了线性映射 \(u\) 是紧的。

推论 1（绍德尔）. --- 设 E 是半赋范空间，F 是巴拿赫空间，E' 和 F' 是它们各自的强对偶，u 是从 E 到 F 的连续线性映射。下列性质等价：

\begin{enumerate}
\def\labelenumi{(\roman{enumi})}
\item
  从 E 到 F 的线性映射 u 是紧的；
\item
  线性映射 \({}^{t}u\)（从 \(\mathrm{F}'\) 到 \(\mathrm{E}'\) 中）是紧的；
\item
  线性映射 \({}^{t}u\)（从 \(\mathrm{F}_{c}^{\prime}\) 到 \(\mathrm{E}'\) 中）是连续的。
\end{enumerate}

(i) 与 (iii) 的等价性由命题 9 的 a) 和 b) 得出；蕴含关系 (iii)⇒(ii) 来自从 F' 到 F'c 的恒等映射是紧的这一事实（III, p. 7, 命题 7 之推论）。

最后我们来证明条件 (ii) 蕴含 (i)。假设线性映射 \({}^{t}u\)：\(F^{\prime} \to E^{\prime}\) 是紧的。把蕴含关系 (i)⇒(ii) 应用于 \({}^{t}u\)，便知 \({}^{tt}u\) 是从 \(E''\) 到 \(F''\) 的紧线性映射。用 \(v\) 表示从 E 到 \(E''\) 的典范线性映射；它是连续的。由于 F 等同于 \(F''\) 的一个闭向量子空间，且 u 在 E 上与 \({}^{t}({}^{t}u) \circ v\) 重合，故由 III, p. 3, 注记 5 推出线性映射 \(u\) 是紧的。

推论 2. --- 设 E 是半赋范空间，F 是可数型的巴拿赫空间。设 u 是从 E 到 F 的连续线性映射。假设对 F' 中元素的每个弱趋于 0 的序列 \((y_{n})\)，序列 \((^{t}u(y_{n}))\) 都在 E' 中强趋于 0。则线性映射 u 是紧的。

用 B′ 表示 F′ 的单位球，赋有由拓扑 \(\sigma(\mathrm{F}',\mathrm{F})\) 诱导的拓扑。由于 F 是可数型的巴拿赫空间，B' 是紧的可度量化空间（EVT, III, p. 19, cor. 2）。在所作假设之下，\({}^{t}u\) 在 B' 上的限制是从 B' 到 E' 的连续映射，且 \({}^{t}u(\mathrm{B}')\) 是 E' 的一个紧子集。由推论 1，线性映射 u 是紧的。

注记. --- 设 E 和 F 是分离的局部凸空间。从 E 到 F 的紧线性映射的转置并不总是从 \(F_{b}^{\prime}\) 到 \(E_{b}^{\prime}\) 的紧线性映射（参见 III, p. 108, 习题 15）。

\subsection{勒雷--绍德尔定理}

下面的定理将在 TA（待出版）中证明。

定理 1（布劳威尔）. --- 设 B 是有限维赋范向量空间的一个非空紧凸子集。则从 B 到 B 的每个连续映射都有不动点。

我们将由它导出下面的结果。

定理 2（勒雷--绍德尔）. --- 设 X 是非空的紧拓扑空间，同胚于一个局部凸空间的凸子集。则从 X 到 X 的每个连续映射都有不动点。

我们可以假设 X 是局部凸空间 E 的一个凸子集。设 N 是 \(\{0\}\) 在 E 中的闭包，\(\pi: E \rightarrow E/N\) 是典范满射。由于 X 是分离空间，\(\pi\) 在 X 上的限制是单射；又由于 X 紧，它便定义一个从 X 到分离局部凸空间 E/N 的一个非空紧凸子集上的同胚。这就使我们能够假设空间 E 是分离的。

设 \(f: \mathrm{X} \to \mathrm{X}\) 是连续映射。映射 \(h: \mathrm{X} \to \mathrm{E}\)（由 \(h(x) = f(x) - x\) 定义）是连续的；由于 X 紧，像 \(h(\mathrm{X})\) 是闭的。因此只需证明 0 附着于 \(h(\mathrm{X})\)，也就是说，对 E 中 0 的每个开凸邻域 U，存在 X 的一个点 \(b\)，使得 \(f(b) - b \in \mathrm{U}\)。

设 U 是 E 中 0 的一个凸邻域。对每个 \(a \in X\)，用 \(V_a\) 表示 X 中这样的元素 \(x\) 所成的集：\(f(x) - f(a) \in U\)。它在 X 中是开的且包含 \(a\)。由于 X 紧且非空，存在 X 的一个非空有限子集 A，使得诸集合 \(V_a\)（\(a \in A\)）覆盖 X。设 \((\varphi_a)_{a \in A}\) 是从属于覆盖 \((V_a)_{a \in A}\) 的一个连续单位分解（TG, IX, p. 43, prop. 1 及 p. 47, th. 3）。用 F 表示由 \(f(A)\) 生成的 E 的子空间。对每个 \(x \in F \cap X\)，令

\[
g (x) = \sum_ {a \in A} \varphi_ {a} (x) f (a).
\]

这样就定义了一个连续映射 \(g: F \cap X \to E\)。它的像包含于 \(f(A)\) 的凸包，从而包含于 \(F \cap X\)。由于集合 \(F \cap X\) 是有限维向量空间的一个非空紧凸子集，由定理 1，映射 g 有一个不动点 \(b \in F \cap X\)。我们有

\[
f (b) - b = f (b) - g (b) = \sum_ {a \in \mathrm{A}} \varphi_ {a} (b) (f (b) - f (a)).
\]

对每个 \(a \in \mathrm{A}\)，若 \(\varphi_{a}(b) \neq 0\)，则我们有 \(b \in \mathrm{V}_{a}\)，从而 \(f(b) - f(a) \in \mathrm{U}\)。由于 U 是凸的，我们得到 \(f(b) - b \in \mathrm{U}\)，证明完毕。

\subsection{紧算子的不变子空间}

应当把下面的定理与舒尔引理（A, VIII, p. 43, cor. 及 V, p. 386, prop. 6）以及 I, p. 26 的推论 4 相比较。

定理 3. --- 设 E 是 C 上的分离局部凸空间，A 是 \(\mathcal{L}(\mathrm{E})\) 的一个子集。我们作下列假设：

\begin{enumerate}
\def\labelenumi{(\roman{enumi})}
\item
  不存在 E 的既异于 \(\{0\}\) 又异于 E 的、在 A 下稳定的闭向量子空间；
\item
  集合 A 含有一个非零的紧自同态。那么 A 的交换子化归于恒等映射的纯量倍数。
\end{enumerate}

用 A 生成的 \(\mathcal{L}(\mathrm{E})\) 的含单位元子代数代替 A，就化归为 A 是 \(\mathcal{L}(\mathrm{E})\) 的一个含单位元子代数的情形。于是设 \(h\) 是属于 A 的一个非零紧自同态。

引理 1. --- 存在 A 的一个元素 a，使得 \(ha - 1_{E}\) 的核不化归于 0。

存在 E 的一个元素 \(x_{0}\)，使得 \(h(x_{0}) \neq 0\)。设 V 是 \(h(x_{0})\) 的一个不包含 0 的闭邻域。由于自同态 h 紧，可以取 \(x_{0}\) 的一个开凸邻域 U，使得 \(h(\mathrm{U})\) 是 V 的一个相对紧子集。于是 \(h(\mathrm{U})\) 的闭包 C 是 E 的一个紧凸子集；由于它包含于 V，它不包含 0。

设 \(y\) 是 C 的一个点。用 \(Ay\) 表示 \(y\) 在 A 的诸元素下的像所成的集。它是 E 的一个在 A 下稳定的向量子空间；由于它包含 \(y\)，它是非零的。由定理的假设 (i)，集合 \(Ay\) 在 E 中稠密。因此存在 A 的一个元素 \(b\)，使得 \(b(y) \in U\)。由于集合 C 紧，存在 A 中元素的一个有限族 \((b_1, \ldots, b_n)\)，使得诸集合 \(b_i^{-1}(U)\) 覆盖 C。存在 C 上从属于这个覆盖的一个连续单位分解 \((\varphi_1, \ldots, \varphi_n)\)（TG, IX, p. 47, th. 3）。我们用下式定义映射 \(f: C \to E\)：

\[
f (x) = \sum_ {i = 1} ^ {n} \varphi_ {i} (x) b _ {i} (x)
\]

其中 \(x \in \mathrm{C}\)；这是一个连续映射。由于 U 是凸的，且 \(b_i(x)\) 在 \(\varphi_i(x)\) 非零时属于 U，我们有 \(f(\mathrm{C}) \subset \mathrm{U}\) 及 \(h(f(\mathrm{C})) \subset \mathrm{C}\)。由勒雷--绍德尔定理（III, p. 11, 定理 2），存在一个元素 \(x \in \mathrm{C}\)，使得 \(h(f(x)) = x\)。于是令

\[
a = \sum_ {i = 1} ^ {n} \varphi_ {i} (x) b _ {i}.
\]

我们有 \(h(a(x)) = h(f(x)) = x\)，且由于 0 不属于 C，x 非零，因此 \(ha - 1_{E}\) 的核非零。

我们来完成定理 3 的证明。设 \(a\) 是 A 的一个元素，使得 \(ha - 1_{\mathrm{E}}\) 的核 T 非零（引理 1）。用 \(i\) 表示 T 到 E 中的典范单射。线性映射 \(i\) 等于 \(h \circ a \circ i\)，因此它是紧的。于是 T 的维数有限（III, p. 2, 注记 3）。设 \(u\) 是 \(\mathcal{L}(\mathrm{E})\) 的一个与 A 交换的元素。由于 \(u\) 与 \(h \circ a\) 交换，我们有 \(u(\mathrm{T}) \subset \mathrm{T}\)。由于 T 在代数闭域 \(\mathbf{C}\) 上是有限维且非零的，\(u\) 在 T 上诱导的自同态有一个特征值 \(\lambda\)。

令 \(\mathrm{F} = \mathrm{Ker}(u - \lambda1_{\mathrm{E}})\)。这是 E 的一个非零闭向量子空间；由于 u 与 A 的诸元素交换，它在 A 下稳定。由假设 (i)，我们有 F = E，从而 \(u = \lambda1_{E}\)。

推论 1. --- 保持定理 3 的假设，并再设 A 的诸元素两两交换。那么 E 在 C 上是一维的。

由定理 3 的假设 (ii)，空间 E 非零。因此它包含一个维数为 1 的向量子空间 F。由于 E 假设为分离的，这个子空间是闭的。由定理 3，A 的每个元素都是位似变换，故使 F 稳定。于是由定理 3 的假设 (i) 推出 F 等于 E。

推论 2（洛莫诺索夫）. --- 设 E 是域 C 上维数至少为 2 的分离局部凸空间，u 是 E 的一个自同态。假设存在 E 的一个与 u 交换的紧自同态 \(h \neq 0\)。那么存在 E 的一个既异于 \(\{0\}\) 又异于 E 的闭向量子空间 F，使得 \(u(F) \subset F\)。

事实上，推论 1 表明集合 A = \{u, h\} 不可能满足定理 3 的假设 (i)。

推论 3. --- 设 E 是域 C 上的分离局部凸空间，u 是 E 的紧自同态。若 E 的维数至少为 2，则存在 E 的一个既异于 \(\{0\}\) 又异于 E 的闭向量子空间 F，使得 \(u(F) \subset F\)。

u = 0 的情形是直接的。\(u \neq 0\) 的情形由推论 2 取 h = u 得出。

\subsection{逼近空间}

给定局部凸空间 E 和 F，回顾（EVT, III, p. 14, 例 2），我们用 \(\mathcal{L}_{\mathrm{pc}}(\mathrm{E};\mathrm{F})\) 表示给 \(\mathcal{L}(\mathrm{E};\mathrm{F})\) 赋以预紧收敛拓扑所得到的局部凸空间。我们也用 \(\mathcal{L}_{\mathrm{pc}}(\mathrm{E})\) 表示局部凸空间 \(\mathcal{L}_{\mathrm{pc}}(\mathrm{E};\mathrm{E})\)。当局部凸空间 E 是豪斯多夫的且拟完备（EVT, III, p. 8, def. 6）时，\(\mathcal{L}(\mathrm{E};\mathrm{F})\) 上的预紧收敛拓扑与紧收敛拓扑重合，因为 E 的预紧子集的闭包是紧的（EVT, III, p. 8）。

定义 2. --- 局部凸空间 E 称为具有逼近性质，或又称为逼近空间，如果 E 的恒等映射 \(1_{E}\) 在空间 \(\mathcal{L}_{\mathrm{pc}}(\mathrm{E})\) 中附着于 E 的连续有限秩自同态所成的集 \(\mathcal{L}^{\mathrm{f}}(\mathrm{E})\)。

特别地，每个有限维局部凸空间都是逼近空间。

注记. --- 1) 为使局部凸空间 E 和 F 的拓扑直和是逼近空间，必须且只需 E 和 F 都是逼近空间。

\begin{enumerate}
\def\labelenumi{\arabic{enumi})}
\setcounter{enumi}{1}
\item
  设 E 是局部凸空间。设 N 是 \(\{0\}\) 在 E 中的闭包，P 是 N 在 E 中的一个代数补子空间。那么 P 是 N 在 E 中的拓扑补，且同构于局部凸空间 E/N。空间 N 是逼近空间。由注记 1，为使 E 是逼近空间，必须且只需豪斯多夫局部凸空间 E/N 是逼近空间。
\item
  设 E 是局部凸空间，A 是 \(\mathcal{L}^{\mathrm{f}}(\mathrm{E})\) 的一个等度连续子集，T 是 E 的一个完全子集。若 \(1_{\mathrm{E}}\) 对 T 上的逐点收敛拓扑附着于 A，则 \(1_{\mathrm{E}}\) 在 \(\mathcal{L}_{\mathrm{pc}}(\mathrm{E})\) 中附着于 A（EVT, III, p. 16, prop. 4 及 p. 17, prop. 5），从而 E 是逼近空间。
\item
  设 E 是 C 上的局部凸空间。用 \(E_{0}\) 表示 E 的底层 R 上的局部凸空间。为使 E 是逼近空间，必须且只需 \(E_{0}\) 本身是逼近空间。事实上，该条件是必要的；我们来证明它是充分的。于是设 \(E_{0}\) 是逼近空间。设 C 是 E 的一个预紧子集，U 是 E 中 0 的一个凸均衡邻域。令 \(C' = C \cup iC\)。存在一个从 \(E_{0}\) 到 \(E_{0}\) 的连续有限秩 R-线性映射 u，使得 \(x - u(x)\) 对一切 \(x \in C'\) 都属于 U。对 E 中的一切 x，令 \(v(x) = \frac{1}{2}(u(x) - iu(ix))\)。这样就定义了一个从 E 到 E 的连续有限秩 C-线性映射。对每个 \(x \in C\)，我们有 \(x \in C'\)，\(ix \in C'\)，并且 \(x - v(x) = \frac{1}{2}(x - u(x)) - \frac{i}{2}(ix - u(ix))\)，从而 \(x - v(x)\) 属于 U。这就证明了 E 是逼近空间。
\item
  设 E 是 R 上的局部凸空间。为使复化局部凸空间 \(\mathrm{E}_{(\mathbf{C})}\) 是逼近空间，必须且只需 E 本身是逼近空间。这由注记 1 和注记 4 得出，因为 \(\mathrm{E}_{(\mathbf{C})}\) 的底层实局部凸空间同构于 \(E \times E\)。
\end{enumerate}

命题 10. --- 设 E 是希尔伯特空间。那么 E 是逼近空间。

E 上的有限秩正交投影算子所成的集 A 是等度连续的。设 \(n \geqslant 1\) 是整数，\((x_{1}, \ldots, x_{n})\) 是 E 的元素的一个族。用 V 表示由诸 \(x_{i}\) 生成的向量子空间，用 \(p_{V}\) 表示以 V 为像的正交投影算子。我们有 \(p_{\mathrm{V}}(x_{i}) = x_{i}\)（对 \(1 \leqslant i \leqslant n\)）。这就蕴含 \(1_{E}\) 对逐点收敛拓扑附着于 A，从而 E 是逼近空间（注记 3）。

逼近空间的其他例子将在 III, p. 20, 第 7 目中给出。

引理 2. --- 设 E 是局部凸空间。那么 E 是逼近空间，当且仅当：对 E 的每个预紧子集 C 和 E 中 0 的每个邻域 U，存在一个整数 \(n \geqslant 0\)、E 的元素 \(e_{1}, \ldots, e_{n}\) 以及 E 上的连续线性形式 \(f_{1}, \ldots, f_{n}\)，使得 \(x - \sum_{i=1}^{n} f_{i}(x)e_{i}\) 对一切 \(x \in C\) 都属于 U。

预紧收敛拓扑的定义表明该条件是充分的。反之，设 E 是逼近空间。设 P 是 \(\{0\}\) 的闭包在 E 中的一个拓扑补空间（参见注记 2）。此时，\(u \in \mathcal{L}^{\mathrm{f}}(\mathrm{E})\)（满足 \(u(\mathrm{E}) \subset \mathrm{P}\)）所成的集 A 在 \(\mathcal{L}^{\mathrm{f}}(\mathrm{E})\) 中稠密。由于 P 是分离的，A 的每个元素 \(u\) 都形如 \(x \mapsto \sum_{i=1}^{n} f_i(x)e_i\)，其中 \(n \in \mathbf{N}\)，\(e_1, \ldots, e_n\) 是 E 的元素，而 \(f_1, \ldots, f_n\) 是 E 上的连续线性形式（EVT, I, p. 14, th. 2）。这就证明了该条件是必要的。

注记 6. --- 存在不具有逼近性质的巴拿赫空间，这是由 P. ENFLO 证明的（A counterexample to the approximation problem in Banach spaces, Acta Math. 130 (1973), p. 309--317; 参见 III, p. 112, 习题 25）。这回答了 S. BANACH 的一个问题（Théorie des opérations linéaires, Monografie Matematyczne I, Warszawa, 1932, 第 VI 章 § 1 中的注记）。另见 p. 21 上的注记。

命题 11. --- 设 E 和 F 是局部凸空间。假设 E 是逼近空间。

\begin{enumerate}
\def\labelenumi{\alph{enumi})}
\item
  集 \(\mathcal{L}^{\mathrm{f}}(\mathrm{E};\mathrm{F})\) 在 \(\mathcal{L}_{\mathrm{pc}}(\mathrm{E};\mathrm{F})\) 中稠密；
\item
  集 \(\mathcal{L}^{\mathrm{f}}(\mathrm{F};\mathrm{E})\) 在 \(\mathcal{L}_{\mathrm{pc}}(\mathrm{F};\mathrm{E})\) 中稠密；
\item
  设 \(\mathfrak{S}\) 是 F 的有界子集所成的一个集，\(u \in \mathcal{L}(\mathrm{F};\mathrm{E})\)。假设在 \(u\) 下，属于 \(\mathfrak{S}\) 的 F 的每个子集的像都是预紧的。那么 \(u\) 附着于 \(\mathcal{L}^{\mathrm{f}}(\mathrm{F};\mathrm{E})\)（对 \(\mathfrak{S}\)-拓扑）（EVT, III, p. 13）。
\end{enumerate}

设 \(v\) 是 \(\mathcal{L}(\mathrm{E};\mathrm{F})\) 的一个元素。映射 \(\varphi\colon w\mapsto v\circ w\)（从 \(\mathcal{L}_{\mathrm{pc}}(\mathrm{E})\) 到 \(\mathcal{L}_{\mathrm{pc}}(\mathrm{E};\mathrm{F})\) 中）是连续的（TG, X, p. 5, prop. 3）。我们有 \(\varphi(1_{\mathrm{E}})=v\) 及 \(\varphi(\mathcal{L}^{\mathrm{f}}(\mathrm{E}))\subset\mathcal{L}^{\mathrm{f}}(\mathrm{E};\mathrm{F})\)，因此 \(v\) 附着于 \(\mathcal{L}^{\mathrm{f}}(\mathrm{E};\mathrm{F})\)（在 \(\mathcal{L}_{\mathrm{pc}}(\mathrm{E};\mathrm{F})\) 中）。这就证明了 \(a\)）。

同样，在 \(c\)）的假设之下，映射 \(\psi: w \mapsto w \circ u\)（从 \(\mathcal{L}_{\mathrm{pc}}(\mathrm{E})\) 到 \(\mathcal{L}_{\mathfrak{S}}(\mathrm{F};\mathrm{E})\) 中）是连续的（同前引处）。我们有 \(\psi(1_{\mathrm{E}}) = u\) 及 \(\psi(\mathcal{L}^{\mathrm{f}}(\mathrm{E})) \subset \mathcal{L}^{\mathrm{f}}(\mathrm{F};\mathrm{E})\)，因此 \(u\) 附着于 \(\mathcal{L}^{\mathrm{f}}(\mathrm{F};\mathrm{E})\)（在 \(\mathcal{L}_{\mathfrak{S}}(\mathrm{F};\mathrm{E})\) 中）。

F 的预紧子集在从 F 到 E 的连续线性映射下的像是预紧的，故 b) 由 c) 得出。

推论. --- 设 E 是分离的逼近空间，F 是局部凸空间。从 F 到 E 的每个紧线性映射都附着于 \(\mathcal{L}^{\mathrm{f}}(\mathrm{F};\mathrm{E})\)（在赋有有界收敛拓扑的空间 \(\mathcal{L}(\mathrm{F};\mathrm{E})\) 中）。

这由命题 11 的 c) 得出，因为 F 的有界子集在从 F 到 E 的紧线性映射下的像在 E 中相对紧（III, p. 2, 注记 1），从而预紧。

命题 12. --- 设 E 是局部凸空间，I 是一个集合，且对每个 \(i \in I\)，设 \(F_{i}\) 是局部凸空间，\(v_{i}: E \to F_{i}\) 是有稠密像的连续线性映射。假设对 E 中 0 的每个邻域 U，存在 \(i \in I\) 及 \(F_{i}\) 中 0 的一个邻域 V，使得 \(v_{i}^{-1}(V) \subset U\)。若诸 \(F_{i}\) 都是逼近空间，则 E 也是逼近空间。

设 A 是 E 的一个预紧子集，U 是 E 中 0 的一个邻域。按假设，存在 \(i \in I\) 及一个连续半范数 \(p\)（定义在 \(F_i\) 上），使得 U 包含 \((p \circ v_i)^{-1}([0,1])\)。令 \(F = F_i\)，\(v = v_i\)，\(B = v(A)\)，并设 F 是逼近空间。集合 B 在 F 中预紧。因此，由引理 2，存在一个整数 \(n \geqslant 1\)、F 的元素 \(y_1, \ldots, y_n\) 以及 F 上的连续线性形式 \(f_1, \ldots, f_n\)，使得对一切 \(y \in B\)，有

\[
p \left(y - \sum_ {j = 1} ^ {n} f _ {j} (y) y _ {j}\right) \leqslant \frac {1}{2}.
\]

由于 B 有界（EVT, III, p. 3, prop. 2），存在实数 M > 0，使得 \(|f_{j}(y)| \leqslant M\) 对一切满足 \(1 \leqslant j \leqslant n\) 的 j 及每个 \(y \in B\) 都成立。此外，由于 \(v(\mathrm{E})\) 在 F 中稠密，对每个满足 \(1 \leqslant j \leqslant n\) 的整数 j，存在 E 的一个元素 \(x_{j}\)，使得 \(p(y_{j} - v(x_{j})) \leqslant (2n\mathrm{M})^{-1}\)。线性映射

\[
u \colon x \longmapsto \sum_ {j = 1} ^ {n} f _ {j} (v (x)) x _ {j}
\]

属于 \(\mathcal{L}^{\mathrm{f}}(\mathrm{E})\)。对每个 \(x\in \mathrm{A}\)，我们有

\[
v (x - u (x)) = v (x) - \sum_ {j = 1} ^ {n} f _ {j} (v (x)) y _ {j} + \sum_ {j = 1} ^ {n} f _ {j} (v (x)) \left(y _ {j} - v \left(x _ {j}\right)\right),
\]

由此 \(p(v(x - u(x))) \leqslant \frac{1}{2} + \frac{nM}{2nM} = 1\)，从而 \(x - u(x) \in U\)。命题得证。

推论 1. --- 逼近空间的一个稠密向量子空间是逼近空间。

推论 2. --- 若局部凸空间 E 的分离完备化是逼近空间，则 E 是逼近空间。

推论 3. --- 一族逼近空间的乘积是逼近空间。

事实上，设 \((\mathrm{E}_{i})_{i\in\mathrm{I}}\) 是一族逼近空间。对 I 的每个有限子集 J，令 \(E_{J}=\prod_{i\in J}E_{i}\)，并用 \(v_{J}\) 表示从 \(E=\prod_{i\in I}E_{i}\) 到 \(E_{J}\) 中的典范映射。局部凸空间 \(E_{J}\) 是逼近空间（注记 1）。于是，把命题 12 应用于局部凸空间 E、局部凸空间族 \((\mathrm{E}_{\mathrm{J}})\) 以及连续线性映射 \(v_{J}:E\to E_{J}\)，即得该推论。

推论 4. --- 设 E 是局部凸空间。若 E 上的每个连续半范数都被一个连续预希尔伯特半范数所控制，则 E 是逼近空间。

设 \(\mathcal{P}\) 是 E 上连续预希尔伯特半范数所成的集。对每个 \(p \in \mathcal{P}\)，假设蕴含：半赋范空间 \(\mathrm{E}_p\)（给 E 赋以半范数 \(p\) 而得到）是逼近空间（推论 2 与命题 10），并且从 E 到 \(\mathrm{E}_p\) 的恒等映射是连续的。因此该推论由命题 12 得出。

引理 3. --- 设 E 是可度量化的局部凸空间，\((x_{n})_{n\in\mathbf{N}}\) 是 E 中收敛于 0 的元素的一个序列。则存在 E 中收敛于 0 的元素的一个序列 \((y_{n})_{n\in\mathbf{N}}\)，以及区间 [0,1] 中收敛于 0 的元素的一个序列 \((\lambda_{n})_{n\in\mathbf{N}}\)，使得 \(x_{n}=\lambda_{n}y_{n}\) 对一切 \(n\in N\) 成立。

由于 E 可度量化，存在 E 中 0 的均衡邻域的一个基本系 \((\mathrm{V}_{m})_{m\in\mathbf{N}}\)，使得 \(V_{0}=E\)，且 \(2^{m+1}V_{m+1}\subset V_{m}\) 对一切 \(m\geqslant0\) 成立。由于 \((x_{n})\) 收敛于 0，存在整数的一个严格递增序列 \((\mathrm{N}_{m})_{m\in\mathbf{N}}\)，使得 \(N_{0}=0\)，且对一切 \(m\geqslant0\)，有 \(x_{n}\in V_{m}\)（对 \(n\geqslant N_{m}\)）。对每个整数 \(n\geqslant0\)，存在唯一的整数 \(m\geqslant0\)，使得 \(N_{m}\leqslant n<N_{m+1}\)，这时我们令 \(\lambda_{n}=2^{-m}\)，\(y_{n}=2^{m}x_{n}\)。如此定义的序列 \((\lambda_{n})\) 收敛于 0。此外，由于 \(y_{n}\in V_{m}\)（对 \(n\geqslant N_{m+1}\)），序列 \((y_{n})\) 在 E 中收敛于 0。最后，对一切 n，我们有 \(x_{n}=\lambda_{n}y_{n}\)。

回顾（EVT, II, p. 28），若 E 是向量空间而 L 是 E 的一个凸均衡子集，我们用 \(E_{L}\) 表示由 L 生成的 E 的向量子空间，赋以以 L 为单位球的半范数。当集合 L 不含任何直线时，这个半范数是一个范数。

引理 4. — 设 E 是弗雷歇空间，A 是 E 的一个紧子集。则存在 E 的一个包含 A 的紧凸均衡子集 L，使得由 E 与由 \(E_{L}\) 在 A 上诱导的拓扑重合。

集合 A 包含于 E 中收敛于 0 的元素的一个序列 \((x_{n})_{n\in\mathbf{N}}\) 的诸点的闭凸均衡包中（EVT, IV, p. 24, cor. 1）。设 \((y_{n})_{n\in\mathbf{N}}\) 与 \((\lambda_{n})_{n\in\mathbf{N}}\) 是满足引理 3 的结论的序列。序列 \((y_{n})\) 的诸点的闭凸均衡包 L 包含 A，且是 E 的一个紧子集（EVT, III, p. 8）。因此，\(E_{L}\) 是巴拿赫空间（EVT, III, p. 8, cor.）。在这个巴拿赫空间中，\(x_{n}\) 的范数被 \(\lambda_{n}\) 所控制，故序列 \((x_{n})\) 趋于 0。\(\widetilde{A}\)（序列 \((x_{n})\) 在 \(E_{L}\) 中的闭凸均衡包）是 \(E_{L}\) 的一个紧子集（EVT, III, p. 8）。由于 L 是 E 的有界子集，\(E_{L}\) 到 E 中的典范单射是连续的。在 \(\widetilde{A}\) 上，由 \(E_{L}\) 诱导的拓扑与由 E 诱导的拓扑重合，且 \(\widetilde{A}\) 是 E 的一个紧子集，从而是闭子集。由于集合 \(\widetilde{A}\) 是凸的、均衡的且包含序列 \((x_{n})\)，故它包含 A。证明完毕。

定理 4. --- 设 E 是弗雷歇空间。

\begin{enumerate}
\def\labelenumi{\alph{enumi})}
\item
  假设 E 是逼近空间。那么，对每个半赋范空间 F，空间 \(\mathcal{L}^{\mathrm{f}}(\mathrm{F};\mathrm{E})\) 对有界收敛拓扑在 \(\mathcal{L}^{\mathrm{c}}(\mathrm{F};\mathrm{E})\) 中稠密；
\item
  反之，假设对每个巴拿赫空间 F，从 F 到 E 的每个紧线性映射都对有界收敛拓扑附着于 \(\mathcal{L}^{\mathrm{f}}(\mathrm{F};\mathrm{E})\)。那么 E 是逼近空间。
\end{enumerate}

设 F 是半赋范空间。\(\mathcal{L}^{\mathrm{f}}(\mathrm{F};\mathrm{E})\) 在 \(\mathcal{L}(\mathrm{F};\mathrm{E})\) 中的闭包含于 \(\mathcal{L}^{\mathrm{c}}(\mathrm{F};\mathrm{E})\)（III, p. 4, 命题 1 与命题 2）。若 E 是逼近空间，则由命题 11 的推论，它等于 \(\mathcal{L}^{\mathrm{c}}(\mathrm{F};\mathrm{E})\)。这就证明了断言 \(a\)）。

假设 \(b\)）的假设成立。设 \(\varepsilon > 0\) 是实数。设 A 是 E 的一个紧子集，\(p\) 是 E 上的一个连续半范数。设 L 是 E 的一个紧凸均衡子集，使得 A 是赋范空间 \(\mathrm{E_L}\) 的一个紧子集（引理 4）。典范单射 \(j: \mathrm{E_L} \to \mathrm{E}\) 是紧的，且 \(\mathrm{E_L}\) 是巴拿赫空间（EVT, III, p. 8, cor.）。因此，按假设，存在一个整数 \(n \geqslant 1\)、E 的元素 \(e_1, \ldots, e_n\) 以及连续线性形式 \(\ell_1, \ldots, \ell_n\)（定义在 \(\mathrm{E_L}\) 上），使得映射 \(v\)（它从 \(\mathrm{E_L}\) 到 E，由 \(v(x) = \sum_{i=1}^{n} \ell_i(x)e_i\) 定义）满足 \(p(x - v(x)) \leqslant \frac{\varepsilon}{2}\)（对 \(x \in A\)）。映射 \({}^{t}j: E' \to (E_{L})'\) 的像对弱拓扑在 \((E_{L})'\) 中稠密（EVT, IV, p. 6, prop. 5）。在 \((E_{L})'\) 上，紧收敛拓扑与 \((E_{L})'\) 和 \(E_{L}\) 之间的对偶相容（EVT, IV, p. 3, 例）。因此 \({}^{t}j(E')\) 对紧收敛拓扑在 \((E_{L})'\) 中稠密（EVT, IV, p. 1, prop. 1），从而存在 E 上的连续线性形式 \(f_{1}, \ldots, f_{n}\)，使得

\[
| \ell_ {i} (x) - f _ {i} (x) | p (e _ {i}) \leqslant \frac {\varepsilon}{2 n}
\]

其中 \(x \in \mathrm{A}\)，\(1 \leqslant i \leqslant n\)。自同态 \(u \colon x \mapsto \sum_{i=1}^{n} f_i(x)e_i\)（\(\mathrm{E}\) 的）属于 \(\mathcal{L}^{\mathrm{f}}(\mathrm{E})\)，且对每个 \(x \in \mathrm{A}\)，我们有

\[
p (x - u (x)) \leqslant p (x - v (x)) + p (v (x) - u (x)) \leqslant \frac {\varepsilon}{2} + n \times \frac {\varepsilon}{2 n} = \varepsilon .
\]

由此推出，\(1_{E}\) 对紧收敛拓扑附着于 \(\mathcal{L}^{\mathrm{f}}(\mathrm{E})\)。而由于 E 完备，这个拓扑与预紧收敛拓扑重合。因此 E 是逼近空间。

\subsection{逼近空间的例子}

回顾，每个希尔伯特空间都是逼近空间（III, p. 15, 命题 10）。本节给出其它一些例子。

若 X 是局部紧空间，我们用 \(\mathcal{C}_{0}(\mathrm{X};\mathrm{K})\)，或简记作 \(\mathcal{C}_{0}(\mathrm{X})\)，表示在无穷远处消失的从 X 到 K 的连续函数所构成的巴拿赫空间，赋以由 \(\|f\|=\sup_{x\in\mathrm{X}}|f(x)|\)（对 \(f\in\mathcal{C}_{0}(\mathrm{X})\)）定义的范数。当 X 紧时，这个空间与从 X 到 K 的连续函数构成的空间 \(\mathcal{C}(\mathrm{X})\) 重合。

引理 5. --- 设 X 是紧拓扑空间，F 是 \(\mathcal{C}(\mathrm{X})\) 的一个有限子集，\(\varepsilon > 0\) 是实数。

存在一个有限子集 \(X_{0} \subset X\) 和一个线性映射 \(u: \mathcal{C}(X_{0}) \to \mathcal{C}(X)\)（范数 \(\leqslant 1\)），使得 \(\|u(f|X_{0}) - f\| \leqslant \varepsilon\) 对一切 \(f \in F\) 成立。

由于集合 F 是 \(\mathcal{C}(\mathrm{X})\) 的等度连续子集，存在 X 的一个有限覆盖 \((\mathrm{U}_{i})_{i\in\mathrm{I}}\)，使得对每个 \(i\in I\) 和每个 \(f\in F\)，\(f(\mathrm{U}_{i})\) 的直径 \(\leqslant\varepsilon\)。对每个 \(i\in I\)，选取一个点 \(x_{i}\)（属于 \(U_{i}\)），并用 \(X_{0}\) 表示诸 \(x_{i}\)（\(i\in I\)）所成的有限集。

设 \((\varphi_{i})_{i\in\mathrm{I}}\) 是从属于 \((\mathrm{U}_{i})_{i\in\mathrm{I}}\) 的连续单位分解。如下定义一个从 \(\mathcal{C}(\mathrm{X}_{0})\) 到 \(\mathcal{C}(\mathrm{X})\) 中的线性映射 u：令

\[
u (g) = \sum_ {i \in \mathrm{I}} g (x _ {i}) \varphi_ {i}
\]

其中 \(g \in \mathcal{C}(X_{0})\)。线性映射 u 的范数 \(\leqslant 1\)，且对 F 中的一切 f 满足 \(\|u(f|X_{0}) - f\| \leqslant \varepsilon\)，由此即得断言。

命题 13. --- 设 X 是局部紧拓扑空间。巴拿赫空间 \(\mathcal{C}_{0}(\mathrm{X})\) 是逼近空间。

首先假设 X 紧。设 \(A \subset \mathcal{L}^{\mathrm{f}}(\mathcal{C}(\mathrm{X}))\) 是形如 \(f \mapsto u(f|X_{0})\) 的映射所成的集合，其中 \(X_{0}\) 是 X 的有限子集，而 \(u : \mathcal{C}(X_{0}) \to \mathcal{C}(X)\) 是范数 \(\leqslant 1\) 的线性映射。集合 A 等度连续。由引理 5，\(\mathcal{C}(X)\) 的恒等映射对 \(\mathcal{C}(X)\) 上的逐点收敛拓扑附着于 A。于是命题由 III, p. 14 的注记 3 得出。

转到一般情形。设 Y 是 X 的亚历山德罗夫紧化，\(\omega\) 是 Y 的无穷远点（TG, I, p. 67）。把 \(\mathcal{C}_{0}(X)\) 与 \(\mathcal{C}(Y)\) 中在 \(\omega\) 处消失的元素所成的集合等同起来。那么 \(\mathcal{C}(Y)\) 是 \(\mathcal{C}_{0}(X)\) 与 Y 上常值函数的向量空间的拓扑直和。由于按前面所证 \(\mathcal{C}(Y)\) 是逼近空间，故空间 \(\mathcal{C}_{0}(X)\) 是逼近空间（III, p. 14 的注记 1）。

推论. --- 每个交换星代数都是逼近空间。

事实上，这样的代数同构于一个局部紧空间上在无穷远处消失的连续函数所构成的代数（I, p. 108, 定理 1）。

注记. --- 若 E 是无穷维希尔伯特空间，则星代数 \(\mathcal{L}(\mathrm{E})\) 不具有逼近性质（A. SZANKOWSKI, \(\mathcal{B}(\mathrm{H})\) does not have the approximation property, Acta Math. 147 (1981), p. 89--108）。

引理 6. --- 设 X 是局部紧拓扑空间。设 \(\mu\) 是 X 上的一个正测度，\(p \in [1, +\infty[\)。设 F 是 \(\mathrm{L}_{\mathrm{K}}^{p}(\mathrm{X}, \mu)\) 的一个有限子集，\(\varepsilon > 0\) 是实数。

存在 \(\mathrm{L}_{\mathrm{K}}^{p}(\mathrm{X},\mu)\) 的一个范数 \(\leqslant 1\) 的有限秩投影子 u，使得对 F 中的每个 f，\(\|u(f)-f\|\leqslant\varepsilon\)。

设 \(\mathcal{P}\) 是 X 的有限分割 \(\pi = (\mathrm{K}_1, \ldots, \mathrm{K}_n, \mathrm{H})\) 所成的集合，其中 \(n \geqslant 1\) 是整数，而 \(\mathrm{K}_1, \ldots, \mathrm{K}_n\) 是 X 的测度非零的可积子集。对每个分割 \(\pi = (\mathrm{K}_1, \ldots, \mathrm{K}_n, \mathrm{H}) \in \mathcal{P}\)，如下定义自同态 \(v_\pi\)（\(\mathcal{L}^p(\mathrm{X}, \mu)\) 的）：令

\[
v _ {\pi} (f) = \sum_ {i = 1} ^ {n} \frac {1}{\mu (\mathrm{K} _ {i})} \left(\int_ {\mathrm{K} _ {i}} f d \mu\right) \varphi_ {\mathrm{K} _ {i}},
\]

其中 \(f \in \mathcal{L}_{\mathrm{K}}^{p}(\mathrm{X}, \mu)\)，\(\varphi_{\mathrm{K}_i}\) 是 \(\mathrm{K}_i\) 的特征函数。映射 \(v_\pi\) 通过过渡到商，诱导出一个投影子 \(u_\pi\)（\(\mathrm{L}_{\mathrm{K}}^{p}(\mathrm{X}, \mu)\) 的）。容易验证，\(u_\pi\) 的像是由满足以下条件的函数 \(f \in \mathcal{L}_{\mathrm{K}}^{p}(\mathrm{X}, \mu)\) 的类所构成的空间：\(f\) 在 H 上为零，且在 \(\mathrm{K}_i\) 上取常数（对 \(1 \leqslant i \leqslant n\)）。

我们来证明 \(\| u_{\pi}\| \leqslant 1\)。对 \(f\in \mathcal{L}_{\mathrm{K}}^{p}(\mathrm{X},\mu)\)，我们有

\[
\| u _ {\pi} (f) \| _ {p} ^ {p} = \sum_ {i = 1} ^ {n} \mu (\mathrm{K} _ {i}) ^ {1 - p} \left| \int_ {\mathrm{K} _ {i}} f d \mu \right| ^ {p}.
\]

由赫尔德不等式（INT, IV, p. 208, § 6, 第 4 目, 定理 2），对每个 i 我们有不等式

\[
\left| \int_ {\mathrm{K} _ {i}} f d \mu \right| ^ {p} \leqslant \mu (\mathrm{K} _ {i}) ^ {p - 1} \int_ {\mathrm{K} _ {i}} | f | ^ {p} d \mu ,
\]

由此得 \(\| u_{\pi}(f)\| _p\leqslant \| f\| _p.\)

设 \(\mathcal{E}\) 是 X 上只取有限多个值的可积函数在 \(\mathrm{L}_{\mathrm{K}}^{p}(\mathrm{X},\mu)\) 中的类所成的集合。由于 \(\mathcal{E}\) 在 \(\mathrm{L}_{\mathrm{K}}^{p}(\mathrm{X},\mu)\) 中稠密（INT, IV, p. 162, §4, 第 10 目, 推论 1），存在一个有限集 \(\mathrm{F}'\)（含于 \(\mathcal{E}\)），使得 F 的每个元素都以至多 \(\varepsilon\) 的距离接近 \(\mathrm{F}'\) 的某个元素。考虑由诸测度非零的集合所构成的有限分割 \(\pi\)，映射 \(x\mapsto (f(x))_{f\in \mathrm{F}'}\) 在每个这样的集合上取一个给定的值，便可看出存在一个元素 \(\pi\)（属于 \(\mathcal{P}\)），使得 \(u_{\pi}(f) = f\) 对一切 \(f\)（含于 \(\mathrm{F}'\)）成立。因此投影子 \(u_{\pi}\) 具有所要的性质。

命题 14. --- 设 X 是局部紧拓扑空间，\(\mu\) 是 X 上的一个正测度，\(p \in [1, +\infty]\)。空间 \(\mathrm{L}_{\mathrm{K}}^{p}(\mathrm{X}, \mu)\) 是逼近空间。

若 \(p = +\infty\)，则空间 \(\mathrm{L}_{\mathbf{C}}^{\infty}(\mathrm{X},\mu)\) 是交换星代数（I, p. 103, 例 4），因而是逼近空间（命题 13 的推论），而由 III, p. 15 的注记 5，\(\mathrm{L}_{\mathbf{R}}^{\infty}(\mathrm{X},\mu)\) 亦然。

假设 p 有限。设 \(\mathrm{A}\subset\mathcal{L}^{\mathrm{f}}(\mathrm{L}_{\mathrm{K}}^{p}(\mathrm{X},\mu))\) 是范数 \(\leqslant1\) 的有限秩投影子所成的集合。由引理 6，\(\mathrm{L}_{\mathrm{K}}^{p}(\mathrm{X},\mu)\) 的恒等映射对逐点收敛拓扑附着于 A，于是命题由 III, p. 14 的注记 3 得出。

\section{紧线性映射的例子}

\subsection{\texorpdfstring{有限迹、希尔伯特--施密特与核型 \(p^{e}\) 幂型的自同态}{有限迹、希尔伯特--施密特与核型 p\^{}\{e\} 幂型的自同态}}

*设 E 与 F 是希尔伯特空间。E 的每个有限迹自同态（EVT, V, p. 49, 定义 7 与 p. 50, 定义 8）都是紧的（IV, p. 165, 推论 1）；从 E 到 F 的每个希尔伯特--施密特映射（EVT, V, p. 51, 定义 9）都是紧的（IV, p. 165, 推论 2）。

对每个实数 \(p \geqslant 1\)，向量空间 \(\mathcal{L}_{p}(\mathrm{E};\mathrm{F})\)（由从 E 到 F 的核型 \(p^{\mathrm{e}}\) 幂型线性映射构成）包含于 \(\mathcal{L}^{\mathrm{c}}(\mathrm{E};\mathrm{F})\)（对 \(p = 1\) 的情形参看 IV, p. 169, 注记 2）。*

\subsection{序列空间中的对角算子}

设 I 是非空集合。回顾（EVT, I, p. 4 与 p. 5），由有界族 \(x = (x_i)_{i \in \mathrm{I}}\)（\(\mathbf{C}\) 的元素）构成的向量空间，赋以范数 \(\| x \| = \sup_i |x_i|\)，是记作 \(\ell_{\mathbf{C}}^{\infty}(\mathrm{I})\) 的巴拿赫空间。

设 \(p\) 是 \([1, +\infty[\) 的一个元素。空间 \(\ell_{\mathbf{C}}^{p}(\mathrm{I})\)（由族 \(x = (x_{i})_{i \in \mathrm{I}}\)（\(\mathbf{C}\) 的元素，且使族 \((|x_{i}|^{p})\) 可和）构成）是一个向量空间，在其上映射 \(x \mapsto \| x \| = (\sum_{i} |x_{i}|^{p})^{1/p}\) 是一个范数，且使它成为巴拿赫空间。若给集合 I 赋以离散拓扑以及使 \(\mu(\{x\}) = 1\)（对一切 \(x \in \mathrm{I}\)）的测度，这个空间正是巴拿赫空间 \(\mathrm{L}_{\mathbf{C}}^{p}(\mathrm{I}, \mu)\)。当 \(p = 1\) 或 p = 2 时，这个记号于是与 EVT, I, p. 4 及 EVT, V, p. 18 中的记号一致（参看 INT, IV, p. 141, §4, 第 1 目, 例）。

对复数的族 \(x = (x_{i})_{i \in \mathrm{I}}\) 与 \(y = (y_{i})_{i \in \mathrm{I}}\)，用 xy 表示族 \((x_{i}y_{i})_{i \in \mathrm{I}}\)。

命题 1. --- 设 \(\lambda = (\lambda_i)_{i\in I}\) 是 C 的元素的一个有界族。设 \(p\) 是 \([1, +\infty]\) 的一个元素。设 E 是巴拿赫空间 \(\ell_{\mathbf{C}}^p (\mathrm{I})\)。

\begin{enumerate}
\def\labelenumi{\alph{enumi})}
\item
  对每个 \(x \in E\)，有 \(\lambda x \in E\)，且映射 u：\(x \mapsto \lambda x\) 是 E 的连续自同态，它在巴拿赫代数 \(\mathcal{L}(E)\) 中的谱等于诸 \(\lambda_{i}\) 所成集合在 C 中的闭包，其范数等于 \(\sup_{i} |\lambda_{i}|\)；
\item
  自同态 u 是紧的，当且仅当对每个实数 \(\varepsilon > 0\)，元素 \(i \in I\) 中满足 \(|\lambda_{i}| \geqslant \varepsilon\) 者所成的集合是有限的。
\end{enumerate}

采用命题中的记号，自同态 \(x \mapsto \lambda x\) 称为 \(\ell_{\mathbf{C}}^{p}(\mathrm{I})\) 的对角自同态。

我们来证明该命题。设 \(C = \sup_{i} |\lambda_{i}|\)。设 \(x \in E\)。我们有不等式

\[
\sum_ {i \in I} | \lambda_ {i} x _ {i} | ^ {p} \leqslant C ^ {p} \| x \| ^ {p}, \text {   if   } p \neq + \infty
\]

\[
\sup _ {i \in I} | \lambda_ {i} x _ {i} | \leqslant C \| x \|, \text {   if   } p = + \infty .
\]

这就证明了 \(\lambda x \in E\)。于是映射 \(x \mapsto \lambda x\) 是 E 的自同态，且这些不等式证明它具有范数 \(\leqslant C\)。

对 \(j \in \mathrm{I}\)，用 \(e_j\) 表示元素 \((x_i)_{i \in \mathrm{I}}\)（\(\ell_{\mathbf{C}}^p(\mathrm{I})\) 中的）：它满足 \(x_j = 1\)，且 \(x_i = 0\)（当 \(i \neq j\)）。于是有 \(u(e_j) = \lambda_j e_j\)（对一切 \(j \in \mathrm{I}\)），这表明 \(\lambda_j\) 属于 \(u\) 的谱。由于 \(u\) 的谱是闭的，在 \(\mathbf{C}\) 中诸 \(\lambda_i\) 所成集合的闭包包含于 \(u\) 的谱。由于 \(u\) 的谱包含在以 0 为中心、以 \(\|u\|\) 为半径的圆盘中（I, p. 24, 定理 1 与 I, p. 21 的公式 (3)），我们推出不等式 \(C \leqslant \|u\|\)，从而 \(\|u\| = C\)。

反之，若 \(\lambda \in \mathbf{C}\) 不附着于诸 \(\lambda_{i}\) 所成的集合，则族 \(((\lambda_{i} - \lambda)^{-1})_{i\in \mathrm{I}}\) 有界。于是前面的论证表明，线性映射 \(x\mapsto ((\lambda -\lambda_i)^{-1}x_i)_{i\in \mathrm{I}}\) 是 E 的一个自同态。它是 \(u - \lambda 1_{\mathrm{E}}\) 的逆，因而 \(\lambda\) 不属于 \(u\) 的谱。这就证明了 \(a\)）。

现在证明 \(b\)）。假设自同态 \(u\) 是紧的。设 \(\varepsilon > 0\)。用 J 表示 \(i \in I\) 中满足 \(|\lambda_i| \geqslant \varepsilon\) 者所成的集合。

诸元素 \(e_i\)（\(i \in \mathbf{J}\)）所成的集合在 E 中有界。它在 \(u\) 下的像由诸元素 \(\lambda_i e_i\)（\(i \in \mathbf{J}\)）组成，因此在 E 中相对紧（III, p. 2 的注记 1）。由于 \(\| \lambda_i e_i - \lambda_j e_j \| \geqslant \varepsilon\)（对 J 中每一对元素 \(i \neq j\)），这仅当集合 J 有限时才可能。

我们来证明逆命题。设 J 是 I 的一个有限子集，记 \(z_{\mathrm{J}} = (z_{\mathrm{J},i})_{i \in \mathrm{I}}\)，其中 \(z_{J,i} = \lambda_i\)（当 \(i \in J\)），而 \(z_{J,i} = 0\)（当 \(i \in I - J\)）；族 \(z_J\) 定义了 E 的一个有限秩自同态 \(u_J \colon x \mapsto z_J x\)。由前所证，线性映射 \(u - u_J\) 的范数被 \(R_+\) 中族 \((|\lambda_i|)_{i \in I-J}\) 的上确界所控制。该假设蕴含：对每个 \(\varepsilon > 0\)，存在有限子集 \(J \subset I\)，使得这个上界 \(\leqslant \varepsilon\)。由此推出映射 u 附着于 \(\mathcal{L}^{\mathrm{f}}(\mathrm{E})\)，因而它是紧的（III, p. 4, 命题 2 的推论）。

注记. --- 1) 若 I 有限，断言 b) 的假设总成立。当集合 I 无穷时，该假设意味着族 \((\lambda_{i})_{i\in\mathrm{I}}\) 沿 I 的有限子集的余集所成的滤子趋于 0；特别地，\(i\in I\) 中满足 \(\lambda_{i}\neq0\) 者所成的集合这时是可数的。

\begin{enumerate}
\def\labelenumi{\arabic{enumi})}
\setcounter{enumi}{1}
\tightlist
\item
  我们稍后将看到（参看 IV, p. 149, 定理 1），当 \(p = 2\) 时，希尔伯特型空间 \(\ell_{\mathbf{C}}^{2}(\mathrm{I})\) 的每个紧正规自同态都形如 \(u = w \circ v \circ w^{-1}\)，其中 \(v\) 是紧对角自同态，而 \(w\) 是希尔伯特空间 \(\ell_{\mathbf{C}}^{2}(\mathrm{I}).*\) 的酉自同态（EVT, V, p. 40）。
\end{enumerate}

\subsection{由核定义的取值于连续函数空间中的线性映射}

这里我们取 K = C。设 X 与 Y 是局部紧拓扑空间。设 \(\mu\) 是 X 上的一个复测度，k 是从 X × Y 到 C 中的、具有下列两条性质的映射：

\begin{enumerate}
\def\labelenumi{(\roman{enumi})}
\item
  对每个 \(y \in Y\)，从 X 到 C 中的函数 \(k_{y} \colon x \mapsto k(x, y)\) 是 \(\mu\) 可积的；
\item
  映射 \(y \mapsto k_{y}\) 从 Y 到 \(\mathcal{L}_{\mathbf{C}}^{1}(\mathrm{X}, \mu)\) 中是连续的。
\end{enumerate}

给 Y 上的复值连续函数空间 \(\mathcal{C}(\mathrm{Y})\) 赋以紧收敛拓扑。对 \(f\)（\(\mathcal{L}_{\mathbf{C}}^{\infty}(\mathrm{X},\mu)\) 中的），定义 \(\widetilde{k}(f)\colon \mathrm{Y}\to \mathbf{C}\) 如下：

\[
\widetilde {k} (f) (y) = \int_ {\mathrm{X}} k (x, y) f (x) d \mu (x) = \int_ {\mathrm{X}} k _ {y} f d \mu .\tag{1}
\]

映射 \(h \mapsto \int_{X} hfd\mu\) 是 \(\mathcal{L}_{\mathbf{C}}^{1}(X,\mu)\) 上的连续线性形式，故条件 (ii) 表明函数 \(\widetilde{k}(f)\) 是连续的。

命题 2. — 映射 \(\widetilde{k}\)（从 \(\mathcal{L}_{\mathbf{C}}^{\infty}(\mathrm{X},\mu)\) 到 \(\mathcal{C}(\mathrm{Y})\) 中）是紧的。

对 Y 中任意的 \(y\) 与 \(y'\)，我们有

\[
| \widetilde {k} (f) (y) | \leqslant \| k _ {y} \| _ {1} \| f \| _ {\infty}, \quad | \widetilde {k} (f) (y) - \widetilde {k} (f) (y ^ {\prime}) | \leqslant \| k _ {y} - k _ {y ^ {\prime}} \| _ {1} \| f \| _ {\infty}.
\]

由阿斯科利定理（TG, X, p. 18, 推论 2）可得，在 \(\widetilde{k}\) 下，\(\mathcal{L}_{\mathbf{C}}^{\infty}(\mathrm{X};\mu)\) 的单位球的像在 \(\mathcal{C}(\mathrm{Y})\) 中相对紧。

推论. --- 设 X 是紧拓扑空间，Y 是局部紧拓扑空间，\(\mu\) 是 X 上的一个复测度，k 是从 X × Y 到 C 中的连续映射。那么公式 (1) 定义了一个仍记作 \(\widetilde{k}\) 的紧线性映射，它从赋以一致收敛拓扑的 \(\mathcal{C}(X)\) 到赋以紧收敛拓扑的 \(\mathcal{C}(Y)\) 中。

我们来验证 \(k\) 的条件 (i) 与 (ii)。对每个 \(y \in \mathrm{Y}\)，映射 \(k_y\) 在紧空间 X 上连续，因而是 \(\mu\) 可积的。此外，由于 X 紧（TG, X, p. 28, 定理 3），映射 \(y \mapsto k_y\) 从 Y 到 \(\mathcal{C}(\mathrm{X})\) 中连续。由于 \(\| k_y - k_{y'} \|_1 \leqslant \| \mu \| \| k_y - k_{y'} \|\)（对一切 \((y, y') \in \mathrm{Y}^2\)），这个从 Y 到 \(\mathcal{L}_{\mathbf{C}}^1(\mathrm{X}, \mu)\) 中的映射是连续的。由命题 2，映射 \(\widetilde{k}\)（从 \(\mathcal{L}_{\mathbf{C}}^\infty(\mathrm{X}, \mu)\) 到 \(\mathcal{C}(\mathrm{Y})\) 中）是紧的。与典范映射 \(\mathcal{C}(\mathrm{X}) \to \mathcal{L}_{\mathbf{C}}^\infty(\mathrm{X}, \mu)\) 复合，便得到一个从 \(\mathcal{C}(\mathrm{X})\) 到 \(\mathcal{C}(\mathrm{Y})\) 中的紧线性映射。
\subsection{以核定义的勒贝格空间之间的线性映射}

这里我们取 K = C。设 X 是赋有正测度 \(\mu\) 的一个局部紧拓扑空间。对每个 \(r \in [1, +\infty]\)，我们记 \(\mathcal{L}^r (\mathrm{X}) = \mathcal{L}_{\mathbf{C}}^r (\mathrm{X},\mu)\)，\(\mathrm{L}^r (\mathrm{X}) = \mathrm{L}_{\mathbf{C}}^r (\mathrm{X},\mu)\)。当 \(r\in [1, + \infty [\) 时，我们把 \(\mathrm{L}^r (\mathrm{X})\) 的对偶与 \(\mathrm{L}^{r^{\prime}}(\mathrm{X})\) 等同起来，其中 \(r^\prime\) 是 \(r\) 的共轭指数，满足 \(1 / r + 1 / r' = 1\)（INT, V, p. 61, § 5, 第 8 目, th. 4）。我们用 \(\| f\| _r\) 表示 \(f\in \mathrm{L}^r (\mathrm{X})\)（或 \(f\in \mathcal{L}^r (\mathrm{X})\)）的范数（或半范数）。

回顾一下，X 的子集 A 称为 \(\mu\) -适度的（INT, V, p. 4, § 1, 第 2 目），如果它含于由一列 \(\mu\) -可积集所成的并之中；而函数 \(f\)（它定义在 X 上，取值于一个向量空间或 \(\overline{\mathbf{R}}\) 中）称为 \(\mu\) -适度的，如果它在 X 的某个 \(\mu\) -适度子集的余集上为零。若 \(1 \leqslant p < +\infty\) 且 \(f \in \mathcal{L}^p(\mathrm{X})\)，则 \(f\) 是适度的（INT, V, p. 9, § 1, 第 3 目, cor.）。若 \(f\) 是 \(\mu\) -适度的，则 \(fg\) 是 \(\mu\) -适度的（对任何函数 \(g\)）。

设 \(p\) 和 \(q\) 是 \(]1, +\infty[\) 中的元素。设 X 与 Y 是局部紧拓扑空间，分别赋有正测度 \(\mu\) 与 \(\nu\)。我们在空间 \(X \times Y\)（相应地，\(Y \times X\)）上赋以乘积测度 \(\mu \otimes \nu\)（相应地，\(\nu \otimes \mu\)）。

我们用 \(\mathcal{N}^{p,q}(\mathrm{X}\times\mathrm{Y},\mu\otimes\nu)\)，或简作 \(\mathcal{N}^{p,q}(\mathrm{X}\times\mathrm{Y})\)，表示满足下列条件的、从 X × Y 到 C 中的关于 \((\mu\otimes\nu)\) 可测的函数 k 所成的集：存在一个实数 \(C\geqslant0\)，使得

\[
\int_ {\mathrm{X} \times \mathrm{Y}} ^ {*} | k (x, y) f (x) g (y) | d (\mu \otimes \nu) (x, y) \leqslant \mathrm{C} \| f \| _ {p} \| g \| _ {q}\tag{2}
\]

对一切函数 \(f \in \mathcal{L}^p(\mathrm{X})\) 与 \(g \in \mathcal{L}^q(\mathrm{Y})\) 成立。我们随后用 \(\| k \|_{p,q}\) 表示具有该性质的实数 C 所成之集的下确界。

注记. --- 设 \(f \in \mathcal{L}^{p}(\mathrm{X})\)，\(g \in \mathcal{L}^{q}(\mathrm{Y})\)。由于 \(p < +\infty\) 且 \(q < +\infty\)，函数 f 与 g 都是适度的，而由 \((x, y) \mapsto f(x)g(y)\) 定义的函数是 \((\mu \otimes \nu)\) -适度的（参见 INT, V, p. 92, § 8, 第 3 目, cor. 2）。因此，函数 \((x, y) \mapsto k(x, y)f(x)g(y)\) 是 \((\mu \otimes \nu)\) -适度的，特别地，其绝对值的上积分与本质上积分重合（INT, V, p. 6, § 1, 第 3 目, prop. 7）。

集 \(\mathcal{N}^{p,q}(\mathrm{X}\times \mathrm{Y})\) 是从 \(\mathrm{X}\times\mathrm{Y}\) 到 \(\mathbf{C}\) 中的函数空间的一个向量子空间，而由 \(k\mapsto \| k\|_{p,q}\) 定义的映射是 \(\mathcal{N}^{p,q}(\mathrm{X}\times \mathrm{Y})\) 上的一个半范数。

引理 1. --- 设 \(k \in \mathcal{N}^{p,q}(X \times Y)\)。函数 \(k\) 是局部 \((\mu \otimes \nu)\) -可积的。

映射 \(k\) 依假设是 \((\mu \otimes \nu)\) -可测的。设 A 是 \(\mathrm{X} \times \mathrm{Y}\) 的一个紧子集。存在紧子集 \(\mathrm{B} \subset \mathrm{X}\) 与 C ⊂ Y，使得 A ⊂ B × C。把 (2) 应用于 B 的特征函数 \(\varphi_{B}\) 与 C 的特征函数 \(\varphi_{C}\)，我们得到

\[
\begin{array}{l} \int_ {X \times Y} ^ {*} | k (x, y) | \varphi_ {\mathrm{A}} (x, y) d (\mu \otimes \nu) (x, y) \\ \qquad \leqslant \int_ {X \times Y} ^ {*} | k (x, y) | \varphi_ {\mathrm{B}} (x) \varphi_ {\mathrm{C}} (y) d (\mu \otimes \nu) (x, y) <   + \infty , \end{array}
\]

由此即得结论（INT, V, p. 41, §5, 第 1 目, prop. 1 与 def. 1）。

命题 3. --- a) 设 \(k \in \mathcal{N}^{p,q}(X \times Y)\)。对 \(f \in \mathcal{L}^p(X)\) 与 \(g \in \mathcal{L}^q(Y)\)，函数

\[
(x, y) \mapsto k (x, y) f (x) g (y)
\]

是 \((\mu \otimes \nu)\) -可积的，并且存在唯一的映射 \(u_{k}\)（从 \(\mathrm{L}^{p}(\mathrm{X})\) 到 \(\mathrm{L}^{q'}(\mathrm{Y})\) 中），使得

\[
\langle g, u _ {k} (f) \rangle = \int_ {\mathrm{X} \times \mathrm{Y}} k (x, y) f (x) g (y) d (\mu \otimes \nu) (x, y)
\]

对一切 \(f \in \mathrm{L}^p(\mathrm{X})\) 与 \(g \in \mathrm{L}^q(\mathrm{Y})\) 成立。映射 \(u_k\) 是线性的且连续的，其范数 \(\leqslant \| k\|_{p,q}\)；

\begin{enumerate}
\def\labelenumi{\alph{enumi})}
\setcounter{enumi}{1}
\item
  设 \(k \in \mathcal{N}^{p,q}(X \times Y)\)。我们有 \(u_k = 0\)，当且仅当 k 是局部 \((\mu \otimes \nu)\) -零测的；
\item
  把 \(u_{k}\) 与 k 相对应的映射是从赋半范空间 \(\mathcal{N}^{p,q}(\mathrm{X}\times\mathrm{Y})\) 到巴拿赫空间 \(\mathcal{L}(\mathrm{L}^{p}(\mathrm{X});\mathrm{L}^{q^{\prime}}(\mathrm{Y}))\) 中的连续线性映射，使得 \(\|u_{k}\|\leqslant\|k\|_{p,q}\) 对一切 \(k\in\mathcal{N}^{p,q}(\mathrm{X}\times\mathrm{Y})\) 成立。
\end{enumerate}

设 \(k \in \mathcal{N}^{p,q}(\mathrm{X} \times \mathrm{X})\)。按定义，从 \(\mathrm{X} \times \mathrm{Y}\) 到 \(\mathbf{C}\) 中的、由 \((x,y) \mapsto k(x,y)f(x)g(y)\) 定义的函数是 \((\mu \otimes \nu)\) -可积的（当 \(f \in \mathcal{L}^p(\mathrm{X})\) 且 \(g \in \mathcal{L}^q(\mathrm{Y})\) 时），而映射

\[
b _ {k} \colon (f, g) \mapsto \int_ {\mathrm{X} \times \mathrm{Y}} k (x, y) f (x) g (y) d (\mu \otimes \nu) (x, y)
\]

是从 \(\mathrm{L}^{p}(\mathrm{X})\times\mathrm{L}^{q}(\mathrm{Y})\) 到 C 中的连续双线性映射。因此，对每个 \(f\in\mathrm{L}^{p}(\mathrm{X})\)，存在唯一的 \(h\in\mathrm{L}^{q^{\prime}}(\mathrm{Y})\)，使得 \(b_{k}(f,g)=\langle g,h\rangle\) 对一切 \(g\in\mathrm{L}^{q}(\mathrm{Y})\) 成立；把它记作 \(u_{k}(f)=h\)。由 (2)，我们有 \(\|u_{k}(f)\|_{q^{\prime}}\leqslant\|k\|_{p,q}\|f\|_{p}\)；映射 \(u_{k}\) 是线性的、连续的，且范数 \(\leqslant\|k\|_{p,q}\)。此外，映射 \(k\mapsto u_{k}\) 是线性的，由此即得断言 a) 与 c)。

我们来证明 \(b\) )。假设 \(k\) 是局部 \((\mu \otimes \nu)\) -零测的。设 \(f \in \mathcal{L}^p(\mathrm{X})\)，\(g \in \mathcal{L}^q(\mathrm{Y})\)；在 \(\mathrm{X} \times \mathrm{Y}\) 上由 \((x, y) \mapsto k(x, y)f(x)g(y)\) 定义的函数是 \((\mu \otimes \nu)\) -适度的。由于它是局部 \((\mu \otimes \nu)\) -零测的，故它是零测的（INT, V, p. 7, § 1, 第 2 目, cor. 1）。于是，我们有 \(\langle u_k(f), g \rangle = b_k(f, g) = 0\)。由此推出 \(u_k = 0\)。

反之，设 \(k \in \mathcal{N}^{p,q}(X \times Y)\) 满足 \(u_k = 0\)。函数 k 是局部 \((\mu \otimes \nu)\) -可积的（引理 1）。对一切函数 \(f \in \mathcal{K}(X)\) 与 \(g \in \mathcal{K}(Y)\)，我们有

\[
\langle g, u _ {k} (f) \rangle = \int_ {\mathrm{X} \times \mathrm{Y}} f (x) g (y) k (x, y) d (\mu \otimes \nu) (x, y),
\]

因此 X × Y 上的测度 \(k \cdot (\mu \otimes \nu)\) 为零（参见 INT, III, p. 82, § 4, 第 1 目, Th. 1）。这就意味着 k 是局部 \((\mu \otimes \nu)\) -零测的（INT, V, p. 46, §5, 第 3 目, Cor. 2）。

定义 1. --- 设 \(k \in \mathcal{N}^{p,q}(\mathrm{X} \times \mathrm{Y})\)。我们称映射 \(u_k\) 为以 \(k\) 为核、从 \(\mathrm{L}^p(\mathrm{X})\) 到 \(\mathrm{L}^{q'}(\mathrm{Y})\) 中的积分算子。

对 \(k \in \mathcal{N}^{p,q}(\mathrm{X} \times \mathrm{Y})\) 与 \(f \in \mathcal{L}^p(\mathrm{X})\)，我们有时也用 \(u_k(f)\) 表示经 \(u_k\) 所得的、\(f\) 在 \(\mathrm{L}^p(\mathrm{X})\) 中的类的像。

命题 4. --- 设 \(k \in \mathcal{N}^{p,q}(\mathrm{X} \times \mathrm{Y})\)，\(f \in \mathcal{L}^p(\mathrm{X})\)。在 \(\mathrm{X}\) 上由 \(x \mapsto k(x,y)f(x)\) 定义的映射是 \(\mu\) -可积的，对每个 \(y\)（它位于一个局部 \(\nu\) -零测集之外）成立，且 \(u_k(f)\) 与由

\[
y \mapsto \int_ {\mathrm{X}} k (x, y) f (x) d \mu (x)
\]

所定义的函数的类重合，该函数是局部 ν-几乎处处定义的。

用 \(Y'\) 表示由那些 \(y \in Y\) 所成的集：对这些元素，映射 \(x \mapsto k(x, y)f(x)\) 不是 \(\mu\) -可积的。设 \(y \in Y\)，并设 C 是 y 的一个紧邻域，其特征函数记作 \(\varphi\)。把条件 (2) 应用于 f 与 \(\varphi\)，则函数 \((x, y) \mapsto k(x, y)f(x)\varphi(y)\) 关于测度 \(\mu \otimes \nu\) 在 X × Y 上可积。由勒贝格--富比尼定理（INT, V, p. 96, § 8, 第 4 目, th. 1, a)），函数 \(x \mapsto k(x, y)f(x)\varphi(y)\) 是 \(\mu\) -可积的，对位于一个 \(\nu\) -零测集 \(Y_{C}\) 之外的 y 成立。由于 \(Y' \cap C \subset Y_{C}\)，由此得出 \(Y'\) 在 Y 中是局部 \(\nu\) -零测的（INT, IV, p. 172, § 5, 第 2 目, def. 3）。

用 h 表示在 Y 上局部 \(\nu\) -几乎处处由 \(y \mapsto \int k(x, y)f(x)d\mu\) 定义的映射。它是局部 \(\nu\) -可积的（INT, V, loc.

cit.）。设 \(g \in \mathcal{K}(\mathrm{Y})\)。我们有

\[
\begin{array}{l} \langle g, u _ {k} (f) \rangle = \int_ {\mathrm{X} \times \mathrm{Y}} k (x, y) f (x) g (y) d (\mu \otimes \nu) (x, y) \\ \qquad = \int_ {\mathrm{Y}} \Bigl (\int_ {\mathrm{X}} k (x, y) f (x) g (y) d \mu (x) \Bigr) d \nu (y) \\ \qquad = \int_ {\mathrm{Y}} g (y) h (y) d \nu (y) \end{array}
\]

（INT, V, loc. cit.）。因此我们有 \(u_{k}(f) = h\)，这在局部 \(\nu\) -几乎处处成立（参见 INT, V, p. 46, § 5, 第 3 目, cor. 2）。

命题 5. --- 空间 \(\mathcal{L}^{p'}(\mathrm{X}\times\mathrm{Y})\) 含于 \(\mathcal{N}^{p,p}(\mathrm{X}\times\mathrm{Y})\)。从 \(\mathcal{L}^{p'}(\mathrm{X}\times\mathrm{Y})\) 到 \(\mathcal{L}(\mathrm{L}^{p}(\mathrm{X});\mathrm{L}^{p'}(\mathrm{Y}))\) 中的、由 \(k\mapsto u_{k}\) 定义的线性映射，经过取商，诱导出从 \(\mathrm{L}^{p'}(\mathrm{X}\times\mathrm{Y})\) 到 \(\mathcal{L}(\mathrm{L}^{p}(\mathrm{X});\mathrm{L}^{p'}(\mathrm{Y}))\) 中的一个连续单射线性映射。

设 \(k \in \mathcal{L}^{p'}(\mathrm{X} \times \mathrm{Y})\)。函数 \(k\) 是可测的（INT, IV, p. 84, § 5, 第 6 目, th. 5）。设 \(f \in \mathcal{L}^p(\mathrm{X})\)，\(g \in \mathcal{L}^p(\mathrm{Y})\)。函数 \(h \colon (x,y) \mapsto f(x)g(y)\) 属于 \(\mathcal{L}^p(\mathrm{X} \times \mathrm{Y})\)，且 \(\|h\|_p = \|f\|_p \|g\|_p\)（INT, V, p. 95, § 8, 第 3 目, cor. 2）。因此函数 \(hk\) 在 \(\mathrm{X} \times \mathrm{Y}\) 上可积，且我们有

\[
(\mu \otimes \nu) ^ {*} (| h k |) \leqslant \| k \| _ {p ^ {\prime}} \| h \| _ {p} = \| k \| _ {p ^ {\prime}} \| f \| _ {p} \| g \| _ {p}
\]

这是赫尔德不等式的推论（INT, IV, p. 208, § 6, 第 4 目, th. 2），由此证明了 \(k \in \mathcal{N}^{p,p}(\mathrm{X} \times \mathrm{Y})\)，且 \(\| k \|_{p,p} \leqslant \| k \|_{p'}\)。

由于 \(p' \neq +\infty\)，每个函数 \(k \in \mathcal{L}^{p'}(\mathrm{X} \times \mathrm{Y})\) 都是适度的（INT, V, p. 9, § 1, 第 3 目, cor.），从而最后的断言由命题 3 的 b) 与 c) 得出。

对任何函数 \(k\colon X \times Y \to \mathbf{C}\) 以及每个 \(y \in Y\)，我们用 \(k_y\) 表示从 X 到 \(\mathbf{C}\) 中的、由 \(k_y(x) = k(x, y)\) 给出的映射，并用 \(k_\circ\) 表示从 Y 到 \(\mathcal{F}(X; \mathbf{C})\) 中的、由 \(y \mapsto k_y\) 定义的映射。我们同样用这一记号表示从 Y 到 X 上 \(\mu\) -几乎处处定义的函数空间中的、使 \(y\) 对应于 \(k_y\) 的类的映射。

对任何函数 \(k\)（从 \(\mathrm{X} \times \mathrm{Y}\) 到 \(\mathbf{C}\) 中），我们记

\[
\mathrm{N} _ {p ^ {\prime}, q ^ {\prime}} (k) = \left(\int_ {\mathrm{Y}} ^ {\bullet} \mathrm{N} _ {p ^ {\prime}} (k _ {y}) ^ {q ^ {\prime}} d \nu (y)\right) ^ {1 / q ^ {\prime}}.
\]

这是 \([0, +\infty]\) 中的一个元素。

引理 2. --- 对每个 \(k \in \mathcal{L}^{p'}(\mathrm{X} \times \mathrm{Y})\)，我们有 \(\mathrm{N}_{p', p'}(k) = \| k \|_{p'}\)。

设 \(k \in \mathcal{L}^{p'}(\mathrm{X} \times \mathrm{Y})\)。由 INT, V, p. 96, § 8, 第 4 目, th. 1, a)，集 \(Y'\)（由那些元素 \(y \in \mathrm{Y}\) 组成，对这些元素，函数 \(|k_y|^{p'}\) 不是 \(\mu\) -可积的）是 \(\nu\) -零测的；此外，由 \(y \mapsto \mu(|k_y|^{p'})\) 在 \(\mathrm{Y} - \mathrm{Y}'\) 上定义且在 \(\mathrm{Y}'\) 上等于零的函数是 \(\nu\) -可积的，并且满足

\[
\| k \| _ {p ^ {\prime}} ^ {p ^ {\prime}} = \int_ {\mathrm{X} \times \mathrm{Y}} | k (x, y) | ^ {p ^ {\prime}} d (\mu \otimes \nu) (x, y) = \int_ {\mathrm{Y}} \mu (| k _ {y} | ^ {p ^ {\prime}}) d \nu = \mathrm{N} _ {p ^ {\prime}, p ^ {\prime}} (k) ^ {p ^ {\prime}}.
\]

我们用 \(\mathcal{L}^{p',q'}(\mathrm{X},\mathrm{Y},\mu,\nu)\)，或简作 \(\mathcal{L}^{p',q'}(\mathrm{X},\mathrm{Y})\)，表示满足下列条件的、关于 \((\mu\otimes\nu)\) 可测的函数 \(k\)（从 \(\mathrm{X}\times\mathrm{Y}\) 到 \(\mathbf{C}\) 中）所成的复向量空间：对 \(\nu\) -几乎所有的 \(y\in\mathrm{Y}\)，映射 \(k_y\) 属于 \(\mathcal{L}^{p'}(\mathrm{X})\)，并且映射 \(k_\circ\) 属于 \(\mathcal{L}^{q'}_{\mathrm{L}^{p'}(\mathrm{X})}(\mathrm{Y},\nu)\)（INT, IV, p. 129, § 3, 第 4 目, def. 2）。我们在 \(\mathcal{L}^{p',q'}(\mathrm{X},\mathrm{Y})\) 上赋以半范数 \(k\mapsto\mathrm{N}_{p',q'}(k)\)；注意此时有 \(\mathrm{N}_{p',q'}(k)=\|k_\circ\|_{q'}\)，其中 \(k_\circ\) 被视为取值于巴拿赫空间 \(\mathrm{L}^{p'}(\mathrm{X})\) 中的映射。

命题 6. --- a) 空间 \(\mathcal{L}^{p'}(\mathrm{X} \times \mathrm{Y})\) 含于 \(\mathcal{L}^{p',p'}(\mathrm{X},\mathrm{Y})\)，且对 \(k \in \mathcal{L}^{p'}(\mathrm{X} \times \mathrm{Y})\)，我们有 \(\mathrm{N}_{p',p'}(k) = \| k\|_{p'}\)；

\begin{enumerate}
\def\labelenumi{\alph{enumi})}
\setcounter{enumi}{1}
\tightlist
\item
  空间 \(\mathcal{L}^{p',q'}(\mathrm{X},\mathrm{Y})\) 含于 \(\mathcal{N}^{p,q}(\mathrm{X}\times \mathrm{Y})\)，且对 \(k\in \mathcal{L}^{p',q'}(\mathrm{X},\mathrm{Y})\)，我们有 \(\| k\|_{p,q}\leqslant \mathrm{N}_{p',q'}(k)\)。
\end{enumerate}

我们来证明断言 \(a\) )。设 \(k \in \mathcal{L}^{p'}(\mathrm{X} \times \mathrm{Y})\)。由 INT, V, p. 96, § 8, 第 4 目, th. 1, \(a\) )，函数 \(k_y\) 属于 \(\mathcal{L}^{p'}(\mathrm{X})\)（对 \(\nu\) -几乎所有的 \(y \in \mathrm{Y}\)）。

设 \(\varepsilon > 0\)。存在 \(h \in \mathcal{K}(\mathrm{X} \times \mathrm{Y})\)，使得 \(\|k - h\|_{p'} < \varepsilon\)。函数 \(h_{\circ}\) 是从 Y 到 \(\mathrm{L}^{p'}(\mathrm{X})\) 中的、具有紧支撑的连续函数，它满足

\[
\| k _ {\circ} - h _ {\circ} \| _ {p ^ {\prime}} = \mathrm{N} _ {p ^ {\prime}, p ^ {\prime}} (k - h) = \| k - h \| _ {p ^ {\prime}} <   \varepsilon
\]

这是引理 2 的结论。由此得出 \(k_{\circ}\) 属于 \(\mathcal{L}_{\mathrm{L}^{p^{\prime}}(\mathrm{X})}^{p^{\prime}}(\mathrm{Y},\nu)\)，也就是说，\(k\) 属于 \(\mathcal{L}^{p^{\prime},p^{\prime}}(\mathrm{X},\mathrm{Y})\)。于是引理 2 表明 \(\mathrm{N}_{p^{\prime},p^{\prime}}(k)=\|k\|_{p^{\prime}}\)。

我们来证明断言 \(b\) )。设 \(k \in \mathcal{L}^{p', q'}(\mathrm{X}, \mathrm{Y})\)，并设 \(f \in \mathcal{L}^p(\mathrm{X})\)。对每个 \(y \in \mathrm{Y}\)，把赫尔德不等式（INT, IV, p. 208, § 6, 第 4 目, th. 2）应用于 \(k_y\)，便蕴含

\[
\int_ {X} ^ {*} | k (x, y) f (x) | d \mu (x) \leqslant \| k _ {y} \| _ {p ^ {\prime}} \| f \| _ {p}.
\]

设 \(g \in \mathcal{L}^{q'}(\mathrm{Y})\)。映射 \(f\) 与 \(g\) 都是适度的，因为 \(p\) 与 \(q'\) 是有限的。因此，映射 \((x,y) \mapsto f(x)g(y)\) 是 \((\mu \otimes \nu)\) -适度的（参见 INT, V, p. 92, § 8, 第 3 目, cor. 2），从而映射 \((x,y)\mapsto k(x,y)f(x)g(y)\) 是 \((\mu\otimes\nu)\) -适度的。根据 INT, V, p. 93, § 8, 第 3 目, prop. 7, a)，于是得出

\[
\begin{array}{l} \int_ {\mathrm{X} \times \mathrm{Y}} ^ {*} | k (x, y) f (x) g (y) | d (\mu \otimes \nu) \\ = \int_ {\mathrm{Y}} ^ {*} | g (y) | \Big (\int_ {\mathrm{X}} ^ {*} | k (x, y) f (x) | d \mu (x) \Big) d \nu (y) \\ \leqslant \| f \| _ {p} \int_ {\mathrm{Y}} ^ {*} \| k _ {y} \| _ {p ^ {\prime}} | g (y) | d \nu (y) \leqslant \mathrm{N} _ {p ^ {\prime}, q ^ {\prime}} (k) \| f \| _ {p} \| g \| _ {q}, \end{array}
\]

这里再次用到了赫尔德不等式。证明到此完成。

命题 7. --- 设 \(k \in \mathcal{L}^{p',q'}(\mathrm{X},\mathrm{Y})\)。线性映射 \(u_k\)（它从 \(\mathrm{L}^p (\mathrm{X})\) 到 \(\mathrm{L}^{q'}(\mathrm{Y})\) 中）是紧的。

首先假设 \(k\) 满足 \(k_{\circ} \in \mathcal{K}(\mathrm{Y}; \mathcal{L}^{p'}(\mathrm{X}))\)，并用 A 表示 \(k_{\circ}\) 的支撑。设 F 是 \(\mathcal{K}(\mathrm{Y}, \mathrm{A}; \mathcal{L}^{p'}(\mathrm{X}, \mu))\) 中由形如 \(y \mapsto f_2(y)f_1\) 的函数所生成的子空间，其中 \(f_2 \in \mathcal{K}(\mathrm{Y}, \mathrm{A})\)，\(f_1 \in \mathcal{L}^{p'}(\mathrm{X})\)。当 \(m \in \mathcal{K}(\mathrm{Y}; \mathcal{L}^{p'}(\mathrm{X}))\) 满足 \(m_{\circ} \in \mathrm{F}\) 时，线性映射 \(u_m\) 是有限秩的，因而是紧的。对 \(m_1\) 与 \(m_2\)（它们属于 \(\mathcal{K}(\mathrm{Y}, \mathrm{A}; \mathcal{L}^{p'}(\mathrm{X}))\)），我们还有

\[
\begin{array}{c} \| u _ {m _ {1}} - u _ {m _ {2}} \| \leqslant \| m _ {1} - m _ {2} \| _ {p ^ {\prime}, q ^ {\prime}} = \Big (\int_ {\mathrm{Y}} \| m _ {1, y} - m _ {2, y} \| _ {p ^ {\prime}} ^ {q ^ {\prime}} d \nu (y) \Big) ^ {1 / q ^ {\prime}} \\ \leqslant \nu (\mathrm{A}) ^ {1 / q ^ {\prime}} \sup _ {y \in \mathrm{Y}} \mathrm{N} _ {p ^ {\prime}} (m _ {1, y} - m _ {2, y}). \end{array}
\]

由于空间 F 在 \(\mathcal{K}(\mathrm{Y},\mathrm{A};\mathrm{L}^{p^{\prime}}(\mathrm{X}))\) 中对于 A 上的一致收敛拓扑是稠密的（INT, III, p. 41, § 1, 第 1 目, prop. 1 及 INT, III, p. 46, § 1, 第 2 目, prop. 5），由此得出结论：线性映射 \(u_{k}\) 是由有限秩线性映射组成的一个序列的极限。因此它是紧的（III, p. 4 的 cor. 1）。

考察一般的情形。设 \(k \in \mathcal{L}^{p', q'}(\mathrm{X}, \mathrm{Y})\)。对每个 \(\varepsilon > 0\)，存在映射 \(k_{\varepsilon, \circ} \in \mathcal{K}(\mathrm{Y}; \mathcal{L}^{p'}(\mathrm{X}))\)，使得

\[
\left\| k _ {\circ} - k _ {\varepsilon , \circ} \right\| _ {q ^ {\prime}} <   \varepsilon .
\]

相应的函数 \(k_{\varepsilon}\) : X × Y → C 于是满足

\[
\| u _ {k} - u _ {k _ {\varepsilon}} \| \leqslant \| k - k _ {\varepsilon} \| _ {p, q} \leqslant \mathrm{N} _ {p ^ {\prime}, q ^ {\prime}} (k - k _ {\varepsilon}) = \| k _ {\circ} - k _ {\varepsilon , \circ} \| _ {q ^ {\prime}} <   \varepsilon .
\]

由于 \(u_{k_{\varepsilon}}\) 依前所述是紧的，映射 \(u_{k}\) 也是紧的（III, p. 4 的 prop. 2）。

\pandocbounded{\includegraphics[keepaspectratio]{images/part001_ec4b93bc.jpg}}

取 \(p = 1\) 并设 \(q > 1\)。可能存在这样的映射 \(k \colon \mathrm{X} \times \mathrm{Y} \to \mathbf{C}\)：\(k_y \in \mathcal{L}^\infty(\mathrm{X})\) 对所有 \(y \in \mathrm{Y}\) 成立，并且积分

\[
\int_ {Y} ^ {*} \| k _ {y} \| _ {\infty} ^ {q ^ {\prime}} d \nu (y)
\]

是有限的（特别地，III, p. 27 的条件 (2) 成立），但线性映射 \(u_k\)（它从 \(L^1(X)\) 到 \(L^{q'}(X)\) 中）却不是紧的（III, p. 119 的习题 2）。

特别地，由该命题可推出下面的推论：

推论 1（希尔伯特--施密特）. --- 设 \(k \in \mathcal{L}^{2}(X \times Y)\)。线性映射 \(u_{k}\) 是从 \(L^{2}(X)\) 到 \(L^{2}(Y)\) 中的紧映射。

由命题 5，有 \(k \in \mathcal{N}^{2,2}(\mathrm{X} \times \mathrm{Y})\)，故线性映射 \(u_k\) 有定义（命题 3）。有 \(k \in \mathcal{L}^{2,2}(\mathrm{X} \times \mathrm{Y})\)（命题 6, \(a\) )），因此 \(u_k\) 是紧的（命题 7）。

注记. --- 在 p = q = 2 的情形，用纯量积来表达积分算子 \(u_{k}\) 的刻画通常更为方便：它是从 \(\mathrm{L}^{2}(\mathrm{X})\) 到 \(\mathrm{L}^{2}(\mathrm{Y})\) 中的唯一映射，使得

\[
\langle g \mid u _ {k} (f) \rangle = \int_ {\mathrm{X} \times \mathrm{Y}} k (x, y) f (x) \overline {{g (y)}} d (\mu \otimes \nu) (x, y)
\]

对一切 \(f \in \mathrm{L}^{2}(\mathrm{X})\) 与 \(g \in \mathrm{L}^{2}(\mathrm{Y})\) 成立。此外，\(u_{k}\) 的伴随是线性映射 \(u_{k^{*}}\)，其中 \(k^{*} \in \mathrm{L}^{2}(\mathrm{Y} \times \mathrm{X})\) 满足 \(k^{*}(y, x) = \overline{k(x, y)}\)（对几乎每个 \((y, x) \in \mathrm{Y} \times \mathrm{X}\)）。事实上，对一切 \(f \in \mathrm{L}^{2}(\mathrm{X})\) 与 \(g \in \mathrm{L}^{2}(\mathrm{Y})\)，我们有

\[
\langle g \mid u _ {k} (f) \rangle = \int_ {\mathrm{X} \times \mathrm{Y}} k (x, y) \overline {{g (y)}} f (x) d (\mu \otimes \nu) (x, y) = \langle u _ {k ^ {*}} (g) \mid f \rangle .
\]

推论 2. --- 设 G 是赋有哈尔测度 μ 的紧拓扑群。设 p 是实数，满足 1 \textless{} p \textless{} +∞，并设 q 是 p 的共轭指数。设 \(f \in \mathcal{L}^{q}(G)\)。对所有 (x,y) ∈ G × G，置 k(x,y) = f(x⁻¹y)。

\begin{enumerate}
\def\labelenumi{\alph{enumi})}
\item
  我们有 \(k \in \mathcal{L}^{q,p}(\mathrm{G},\mathrm{G})\)；
\item
  对 \(\varphi \in \mathrm{L}^p (\mathrm{G})\)，卷积 \(\varphi *f\) 属于 \(\mathrm{L}^p (\mathrm{G})\)，且线性映射 \(v_{f}\colon \varphi \mapsto \varphi *f\)（它从 \(\mathrm{L}^p (\mathrm{G})\) 到其自身）是连续的。它与自同态 \(u_{k}\) 重合。特别地，线性映射 \(v_{f}\) 是紧的。
\end{enumerate}

由于 G 是紧的，函数 f 属于 \(\mathcal{L}^{1}(G)\)。因此线性映射 \(v_{f}\) 有定义且连续（依据 INT, VIII, p. 167, § 4, 第 5 目, prop. 13）。

映射 \(k\) 是 \((\mu \otimes \mu)\) -可测的，且 \(\| k_y\| _q = \| f\| _q\) 对所有 \(y\in G\) 成立。由于测度 \(\mu\) 是有界的，由此得出

\[
\int_ {\mathrm{G}} \| k _ {y} \| _ {q} ^ {p} d \mu (y) = \mu (\mathrm{G}) \| f \| _ {q} ^ {p},
\]

因此 \(k \in \mathcal{L}^{q,p}(G, G)\)。于是线性映射 \(u_{k}\) 是从 \(\mathrm{L}^{p}(\mathrm{G})\) 到 \(\mathrm{L}^{p}(\mathrm{G})\) 中的紧线性映射（命题 7）。

设 \(\varphi\in\mathcal{K}(G)\)。由 INT, VIII, p. 166, § 4, 第 5 目, prop. 11 以及 III, p. 29 的 prop. 4，我们于是有

\[
\begin{array}{c} v _ {f} (\varphi) (y) = \int_ {\mathrm{G}} \varphi (x) f (x ^ {- 1} y) d \mu (x) \\ = \int_ {\mathrm{G}} \varphi (x) k (x, y) d \mu (x) = u _ {k} (\varphi) (y) \end{array}
\]

对几乎每个 \(y \in G\) 成立。这就蕴含 \(v_{f} = u_{k}\)。

\subsection{可微映射的限制}

设 n 与 r 是正整数，设 U 是 \(R^{n}\) 的开子集，并设 F 是巴拿赫空间。用 \(\mathcal{C}^{r}(\mathrm{U};\mathrm{F})\) 表示从 U 到 F 中的 \(C^{r}\) 类映射所成的向量空间，赋以 \(C^{r}\) 紧收敛拓扑。回忆它是当 K 取遍 U 的紧子集全体时，在 K 上的 \(C^{r}\) 一致收敛拓扑诸拓扑的上确界（VAR, R2, 12.3.10, p. 56）。空间 \(\mathcal{C}^{0}(\mathrm{U};\mathrm{F})\) 就是从 U 到 F 中的连续映射所成的空间 \(\mathcal{C}(\mathrm{U};\mathrm{F})\)，赋以紧收敛拓扑。

引理 3. --- 设 A 是由多重指标 \(\alpha \in \mathbf{N}^n\)（其中 \(|\alpha| \leqslant r\)）组成的集，并设 u 是线性映射 \(f \mapsto (\partial^\alpha f)_{\alpha \in \mathrm{A}}\)（它从 \(\mathcal{C}^r(\mathrm{U};\mathrm{F})\) 到 \(\mathcal{C}(\mathrm{U};\mathrm{F})^{\mathrm{A}}\) 中）。

\begin{enumerate}
\def\labelenumi{\alph{enumi})}
\item
  映射 u 是单射、连续的、严格的，且它的像是闭的。
\item
  拓扑向量空间 \(\mathcal{C}^{r}(\mathrm{U};\mathrm{F})\) 是完备的。
\end{enumerate}

映射 \(u\) 是线性的且是单射。由 \(\mathcal{C}^r (\mathrm{U};\mathrm{F})\) 的拓扑的定义，它是连续的且是严格的。

设 \((g_{\alpha})_{\alpha\in\mathrm{A}}\) 是 \(\mathcal{C}(\mathrm{U};\mathrm{F})^{\mathrm{A}}\) 中附着于 u 的像的一点。存在 \(\mathcal{C}^{r}(\mathrm{U};\mathrm{F})\) 上的一个滤子 F，使得 \(g_{\alpha}=\lim_{f,\mathcal{F}}\partial^{\alpha}f\) 在 \(\mathcal{C}(\mathrm{U};\mathrm{F})\) 中对一切 \(\alpha\in A\) 成立。设 m 是满足 \(0\leqslant m\leqslant r\) 的整数。对 m 用归纳法，由 FVR, II, p. 2 的 th. 1 可得，映射 \(g_{0}\) 是 \(C^{m}\) 类的，且 \(g_{\alpha}=\partial^{\alpha}g_{0}\) 对一切 \(\alpha\in\mathbf{N}^{n}\)（其中 \(|\alpha|\leqslant m\)）成立。于是 \(g_{0}\) 属于空间 \(\mathcal{C}^{r}(\mathrm{U};\mathrm{F})\)，且它在 u 下的像是 \((g_{\alpha})_{\alpha\in\mathrm{A}}\)。这就证明了 u 的像是闭的，从而建立了断言 a)。

断言 \(a\)）蕴含空间 \(\mathcal{C}^{r}(\mathrm{U};\mathrm{F})\) 同构于它在 \(\mathcal{C}(\mathrm{U};\mathrm{F})^{\mathrm{A}}\) 中的像；由于这个像是闭的，且空间 \(\mathcal{C}(\mathrm{U};\mathrm{F})^{\mathrm{A}}\) 是完备的（TG, X, p. 9, th. 2 的 cor. 3 以及 TG, II, p. 17, prop. 10），所以空间 \(\mathcal{C}^{r}(\mathrm{U};\mathrm{F})\) 是完备的。

命题 8. --- 假设 F 是有限维的。设 s 是满足 \(0 \leqslant s < r\) 的整数，并设 V 是 U 的相对紧开子集。线性映射 \(f \mapsto f|V\)（它从 \(\mathcal{C}^{r}(U;F)\) 到 \(\mathcal{C}^{s}(V;F)\) 中）是紧的。

设 W 是由这样的函数 \(f \in \mathcal{C}^{r}(\mathrm{U};\mathrm{F})\) 组成的集：其 \(\leqslant r\) 阶偏导数在 \(\overline{V}\) 的每一点处取范数小于 1 的值。集 W 是空间 \(\mathcal{C}^{r}(\mathrm{U};\mathrm{F})\) 中 0 的邻域。设 \(\alpha\) 是满足 \(|\alpha| \leqslant s\) 的多重指标。考虑由形如 \((\partial^{\alpha}f)|\mathrm{V}\)（其中 f 属于 W）的函数组成的集 H。由中值定理（VAR, R, 2.2.3），集 H 是 \(\mathcal{C}(\mathrm{V};\mathrm{F})\) 的等度连续子集。此外，对每个 \(x \in V\)，H 在映射 \(g \mapsto g(x)\) 下的像是 F 中的有界子集，因而是相对紧的。由阿斯科利定理（TG, X, p. 18, cor. 2），集 H 在 \(\mathcal{C}(\mathrm{V};\mathrm{F})\) 中是相对紧的。这就证明了线性映射 \(f \mapsto (\partial^{\alpha}f)|\mathrm{V}\)（它从 \(\mathcal{C}^{r}(\mathrm{U};\mathrm{F})\) 到 \(\mathcal{C}(\mathrm{V};\mathrm{F})\) 中）是紧的。命题于是由引理 3 以及 III, p. 3 的注记 5 和 6 得出。

\subsection{向量丛的可微截面的限制}

设 r 是正整数，X 是 \(C^{r}\) 类的局部有限维可微流形，而 E 是一个向量丛（实的或复的），

其基为 X 且是 \(C^{r}\) 类的。对 X 的每个开集 U，用 \(\mathcal{S}^{r}(\mathrm{U};\mathrm{E})\)（在 VAR, R1, 7.4, p. 74 中记作 \(\mathcal{S}_{\mathrm{E}}^{r}(\mathrm{U})\)）表示 E 在 U 上的 \(C^{r}\) 截面所成的向量空间，赋以 \(C^{r}\) 紧收敛拓扑。

引理 4. --- 设 \(\mathcal{U}\) 是 X 的一个开覆盖。映射 \(u: f \mapsto (f|U)_{U \in \mathcal{U}}\)（它从 \(\mathcal{S}^{r}(X;E)\) 到 \(\prod_{U \in \mathcal{U}} \mathcal{S}^{r}(U;E)\) 中）是线性的、单射、连续的且是严格的。它的像是闭的。

映射 \(u\) 是线性的、单射且连续的。依据 TG, IX, p. 43, prop. 1 与 p. 48, cor. 1，X 的每个紧子集都有有限覆盖 \((\mathrm{C}_i)_{i\in \mathrm{I}}\)，使得对每个 \(i\in \mathrm{I}\)，集 \(\mathrm{C}_i\) 是覆盖 \(\mathcal{U}\) 中某个开集的紧子集。由此得出线性映射 \(u\) 是严格的。最后，它的像是闭的，因为它由这样的族 \((f_{\mathrm{U}})_{\mathrm{U}\in \mathcal{U}}\) 组成：\(f_{\mathrm{U}}\) 与 \(f_{\mathrm{V}}\) 在 \(\mathrm{U}\cap \mathrm{V}\) 中重合（对 \(\mathcal{U}\) 中所有的 U 与 V）。

命题 9. --- 空间 \(\mathcal{S}^r(\mathrm{X};\mathrm{E})\) 是完备的。

首先假设存在整数 \(n \geqslant 0\) 与巴拿赫空间 F，使得 X 是 \(R^{n}\) 的开子集，且 E 是以 X 为基、以 F 为纤维的平凡向量丛 \(X \times F\)。在这种情形，拓扑向量空间 \(\mathcal{S}^{r}(X;E)\) 同构于 \(\mathcal{C}^{r}(X;F)\)，结论由 III, p. 34 的引理 3 得出。

在一般情形，设 \(\mathcal{U}\) 是 X 的由坐标域组成的开覆盖，使得对每个 \(\mathrm{U} \in \mathcal{U}\)，E 在 U 上的限制是可平凡化的。由上面的引理，空间 \(\mathcal{S}^r(\mathrm{X};\mathrm{E})\) 同构于线性映射 \(f \mapsto (f|\mathrm{U})_{\mathrm{U} \in \mathcal{U}}\) 的像，而这个像在诸空间 \(\mathcal{S}^r(\mathrm{U};\mathrm{E})\)（其中 U 取遍 \(\mathcal{U}\)）的乘积中是闭的。由前面的情形，每个空间 \(\mathcal{S}^r(\mathrm{U};\mathrm{E})\) 都是完备的，因此它们的乘积是完备的（TG, II, p. 17, prop. 10）。从而 \(\mathcal{S}^r(\mathrm{X};\mathrm{E})\) 是完备的。

注记. --- 设 \(\mathcal{U}\) 是 X 的由坐标域 \(c_{\mathrm{U}} = (\mathrm{U}, \varphi_{\mathrm{U}}, \mathbf{R}^{n_{\mathrm{U}}})\) 组成的开覆盖，使得对每个 \(U \in \mathcal{U}\)，E 在 U 上的限制是 \(F_{U}\) 型可平凡化的，其中 \(F_{U}\) 是巴拿赫空间。对每个 \(U \in \mathcal{U}\)，把 E 在 U 上的截面 f 与映射 \(f_{U}\)（它从 \(\varphi_{\mathrm{U}}(\mathrm{U})\) 到 \(F_{U}\) 中）等同起来。由 prop. 9 的证明可得，空间 \(\mathcal{S}^{r}(\mathrm{X};\mathrm{E})\) 的拓扑由半范数族 \(p_{U,C,\alpha}\) 定义，其中

\[
p _ {\mathrm{U,C}, \alpha} (f) = \sup _ {x \in \mathrm{C}} \| (\partial^ {\alpha} f _ {\mathrm{U}}) (x) \|,
\]

其中 U 取遍 \(\mathcal{U}\)，C 取遍 \(\varphi_{\mathrm{U}}(\mathrm{U})\) 的紧子集全体，而 \(\alpha \in \mathbf{N}^{n_{\mathrm{U}}}\) 取遍满足 \(|\alpha| \leqslant r\) 的多重指标全体。

命题 10. --- 假设向量丛 E 是局部有限秩的。设 s 是满足 \(0 \leqslant s < r\) 的整数，并设 Y 是 X 的相对紧开子集。线性映射 \(f \mapsto f|Y\)（它从 \(\mathcal{S}^{r}(X;E)\) 到 \(\mathcal{S}^{s}(Y;E)\) 中）是紧的。

首先假设存在整数 \(n \geqslant 0\) 与有限维巴拿赫空间 F，使得 X 是 \(R^{n}\) 的开子集，且 E 是以 X 为基、以 F 为纤维的平凡向量丛 \(X \times F\)。在这种情形，拓扑向量空间 \(\mathcal{S}^{r}(X;E)\) 与 \(\mathcal{S}^{s}(Y;E)\) 分别同构于 \(\mathcal{C}^{r}(X;F)\) 与 \(\mathcal{C}^{s}(Y;F)\)，而 prop. 10 是 prop. 8 的推论。

转到一般情形。设 C 是 X 的包含 Y 的紧子集。对 C 的每一点 x，选取 x 在 X 中的一个开邻域 \(U(x)\)，它是一个坐标域，且向量丛 E 在其上是可平凡化的。此外，再选取 x 的一个相对紧开邻域 \(V(x)\)，它包含在 \(U(x)\) 中。由于集 C 是紧的，它被这样一个有限族 \((\mathrm{V}(x_{1}),\ldots,\mathrm{V}(x_{m}))\) 所覆盖。对每个 i，置 \(U_{i}=U(x_{i})\) 与 \(Y_{i}=V(x_{i})\cap Y\)。那么 \(Y_{i}\) 是 \(U_{i}\) 的相对紧开子集，且我们有 \(Y=Y_{1}\cup\cdots\cup Y_{m}\)。线性映射 \(f\mapsto f|U_{i}\)（从 \(\mathcal{S}^{r}(X;E)\) 到 \(\mathcal{S}^{r}(U_{i};E)\) 中）是连续的，而线性映射 \(g\mapsto g|Y_{i}\)（从 \(\mathcal{S}^{r}(U_{i};E)\) 到 \(\mathcal{S}^{s}(Y_{i};E)\) 中）由证明的第一部分是紧的。因此线性映射 \(f\mapsto f|Y_{i}\)（从 \(\mathcal{S}^{r}(X;E)\) 到 \(\mathcal{S}^{s}(Y_{i};E)\) 中）是紧的（III, p. 5 的 prop. 3）。把引理 4（应用于 Y 被诸开集 \(Y_{i}\) 所成的覆盖）与 III, p. 3 的注记 5 和 6 一并考虑，线性映射 \(f\mapsto f|Y\)（从 \(\mathcal{S}^{r}(X;E)\) 到 \(\mathcal{S}^{s}(Y;E)\) 中）是紧的，这就完成了证明。

\subsection{向量丛的解析截面的限制}

设 X 是局部有限维的复解析流形，E 是以 X 为基的解析向量丛（VAR, p. 35 与 VAR, p. 70）。用 \(\mathcal{S}^{\omega}(\mathrm{X};\mathrm{E})\) 表示 E 在 X 上的解析截面所成的向量空间，赋以紧收敛拓扑。

引理 5. --- 设 \(X_{0}\) 是 X 的底实解析流形，而 \(E_{0}\) 是以 \(X_{0}\) 为基的、作为 E 的底的复向量丛。

\begin{enumerate}
\def\labelenumi{\alph{enumi})}
\item
  向量子空间 \(\mathcal{S}^{\omega}(\mathrm{X};\mathrm{E})\)（它包含于 \(\mathcal{S}^{0}(\mathrm{X}_{0};\mathrm{E}_{0})\) 中）是闭的，且从 \(\mathcal{S}^{\omega}(\mathrm{X};\mathrm{E})\) 到 \(\mathcal{S}^{0}(\mathrm{X}_{0};\mathrm{E}_{0})\) 中的单射是连续的且是严格的；
\item
  从 \(\mathcal{S}^{\omega}(\mathrm{X};\mathrm{E})\) 到 \(\mathcal{S}^{1}(\mathrm{X}_{0};\mathrm{E}_{0})\) 中的典范单射是连续的。
\end{enumerate}

首先假设存在整数 \(n \geqslant 0\) 与巴拿赫空间 F，使得 X 是 \(C^{n}\) 的开子集，且 E 是平凡向量丛 \(X \times F\)。在这种情形，拓扑向量空间 \(\mathcal{S}^{\omega}(X;E)\)、\(\mathcal{S}^{0}(X_{0};E_{0})\) 与 \(\mathcal{S}^{1}(X_{0};E_{0})\) 分别同构于 \(\mathcal{C}^{\omega}(X;F)\)、\(\mathcal{C}^{0}(X_{0};F)\) 与 \(\mathcal{C}^{1}(X_{0};F)\)，且空间 \(\mathcal{S}^{0}(X_{0};E_{0})\) 是可度量的（TG, X, p. 20, prop. 1 的 cor.）。引理于是由 VAR, 3.3.2, p. 28 得出。

设 \(\mathfrak{F}\) 是 \(\mathcal{S}^{\omega}(\mathrm{X};\mathrm{E})\) 上的一个滤子，它在空间 \(\mathcal{S}^{0}(\mathrm{X}_{0};\mathrm{E}_{0})\) 中收敛于极限 \(f\)。我们必须证明 \(f\) 属于 \(\mathcal{S}^{\omega}(\mathrm{X};\mathrm{E})\)，且滤子 \(\mathfrak{F}\) 收敛于 \(f\)（在空间 \(\mathcal{S}^{1}(\mathrm{X}_{0};\mathrm{E}_{0})\) 中）。这个命题是局部性质的，因而由证明的第一部分得出。

命题 11. --- 假设向量丛 E 是局部有限秩的。设 Y 是 X 的相对紧开子集。限制映射 \(f \mapsto f|Y\)（它从 \(\mathcal{S}^{\omega}(X;E)\) 到 \(\mathcal{S}^{\omega}(Y;E)\) 中）是紧的。

用引理 5 的记号，我们有交换图

\[
\begin{array}{c c c} \mathcal {S} ^ {\omega} (\mathrm{X}; \mathrm{E}) & \stackrel {{i}} {{\longrightarrow}} & \mathcal {S} ^ {1} (\mathrm{X} _ {0}; \mathrm{E} _ {0}) \\ \Big \downarrow^ {u} & & \Big \downarrow^ {v} \\ \mathcal {S} ^ {\omega} (\mathrm{Y}; \mathrm{E}) & \stackrel {{j}} {{\longrightarrow}} & \mathcal {S} ^ {0} (\mathrm{Y} _ {0}; \mathrm{E} _ {0}) \end{array}\tag{3}
\]

其中 i 与 j 是典范单射，而 u 与 v 是限制映射。映射 i 是连续的（引理 5, b)），映射 v 是紧的（Prop. 10）。由于典范单射 j 是连续的、严格的且有闭像（引理 5, a)），映射 u 是紧的（III, p. 3 的注记 5）。

推论. --- 设 X 是紧复解析流形，E 是 X 上的解析向量丛，且是局部有限秩的。向量空间 \(\mathcal{S}^{\omega}(X;E)\) 是有限维的。

解析流形 X 是紧的，因而是局部有限维的。由 prop. 11，\(\mathcal{S}^{\omega}(\mathrm{X};\mathrm{E})\) 的恒等映射是紧线性映射。这蕴含空间 \(\mathcal{S}^{\omega}(\mathrm{X};\mathrm{E})\) 是有限维的（III, p. 2 的注记 3）。

\section{弗雷德霍姆自同态与里斯自同态}

\subsection{严格态射与有限秩线性映射}

命题 1. --- 设 E 与 F 是局部凸空间，\(E_{1}\) 是 E 的有限余维闭子空间，而 u 是从 E 到 F 中的连续线性映射。要使 u 是具有闭像的严格态射，必须且只需 \(u|E_{1}\) 是这样的态射。

首先假设线性映射 \(u|E_1\) 是单射。于是我们有 \(E_1 \cap \text{Ker}(u) = \{0\}\)，从而 \(\text{Ker}(u)\) 是有限维的。设 S 是 E 的一个向量子空间，它是 \(E_1 + \text{Ker}(u)\) 的补。空间 E 是 \(E_1\)、\(\text{Ker}(u)\) 与 S 的拓扑直和（EVT, I, p. 15, cor. 4 与 prop. 3）。如果 \(u\) 是具有闭像的严格态射，那么它通过限制定义了从 \(E_1 \oplus S\) 到 F 的一个闭向量子空间上的同构，从而更加必然地定义了从 \(E_1\) 到 F 的一个闭向量子空间上的同构。反之，假设 \(u\) 通过限制定义了从 \(E_1\) 到 \(u(E_1)\) 上的同构，且 \(u(E_1)\) 在 F 中是闭的。我们有 \(u(E) = u(E_1) \oplus u(S)\)。由此得出 \(u(E)\) 是 \(u(E_1)\) 与 \(u(S)\) 的拓扑直和，且在 F 中是闭的（loc. cit.）。由于 \(u(E_1)\) 是闭的，且我们有 \(u(E_1) \cap u(S) = \{0\}\)，空间 \(u(S)\) 是分离的，且 \(u\) 通过限制定义了从 S 到 \(u(S)\) 上的同构（EVT, I, p. 15, prop. 3 的 cor.），因而也定义了从 \(E_1 \oplus S\) 到 \(u(E)\) 上的同构。这就证明了 \(u\) 是具有闭像的严格态射。

转到一般情形，置 \(N = E_{1} \cap \operatorname{Ker}(u)\) 与 G = E/N。局部凸空间 \(E_{1}/N\) 等同于 G 的一个有限余维闭向量子空间 \(G_{1}\)。用 \(v\colon G \to F\) 表示由 u 通过过渡到商而得到的连续线性映射。要使 \(u\)（相应地，\(u|E_{1}\)）是具有闭像的严格态射，必须且只需 \(v\)（相应地，\(v|G_{1}\)）是这样的态射。这就化归到已处理过的情形。

推论 1. --- 假设 F 是分离的。设 \(v \in \mathcal{L}^{\mathrm{f}}(\mathrm{E};\mathrm{F})\)。如果 u 是从 E 到 F 中的具有闭像的严格态射，那么 \(u + v\) 也是。

由于 F 是分离的，v 的核是 E 的闭向量子空间；它在 E 中有有限余维，且 \(u + v\) 与 \(u\) 在 \(\operatorname{Ker}(v)\) 上有相同的限制。推论于是由该命题得出。

推论 2. --- 设 T 是 F 的分离的有限维向量子空间。用 \(\pi\colon\mathrm{F}\to\mathrm{F}/\mathrm{T}\) 表示典范满射。要使 u 是具有闭像的严格态射，必须且只需 \(\pi\circ u\) 是这样的态射。

T 到自身的恒等映射延拓为连续映射 \(q \colon \mathrm{F} \to \mathrm{T}\)（EVT, II, p. 26, remark）。\(q\) 的核 S 是 T 的闭拓扑补。设 \(p \colon \mathrm{F} \to \mathrm{F}\) 是与分解 \(\mathrm{F} = \mathrm{T} \oplus \mathrm{S}\) 相伴的、以 S 为像的投影。要使 \(\pi \circ u\) 是具有闭像的严格态射，必须且只需 \(p \circ u\) 是这样的态射。此外，\(p \circ u\) 与 \(u\) 在 \(u^{-1}(\mathrm{S})\)（它是 E 中的有限余维闭向量子空间）上有相同的限制，断言于是由该命题得出。

\subsection{弗雷德霍姆映射}

设 E 与 F 是局部凸空间。在本节中，我们用 \(u \equiv v\) 表示模 \(\mathcal{L}^{\mathrm{f}}(\mathrm{E};\mathrm{F})\) 的同余关系，这里 u 与 v 是 \(\mathcal{L}(\mathrm{E};\mathrm{F})\) 的元素。

如果 G 是局部凸空间，且 \(u'\) 与 \(v'\) 是 \(\mathcal{L}(\mathrm{F};\mathrm{G})\) 的元素，那么关系 \(u \equiv v\) 与 \(u' \equiv v'\) 蕴含关系 \(u' \circ u \equiv v' \circ v\)。

我们称元素 \(w\)（属于 \(\mathcal{L}(\mathrm{F};\mathrm{E})\)）是元素 \(u\)（属于 \(\mathcal{L}(\mathrm{E};\mathrm{F})\)）的拟逆，如果有 \(w \circ u \equiv 1_{\mathrm{E}}\) 与 \(u \circ w \equiv 1_{\mathrm{F}}\)。

假设 w 是 u 的拟逆。如果 \(u_{1}\) 是 \(\mathcal{L}(\mathrm{E};\mathrm{F})\) 的元素，且 \(w_{1}\) 是 \(\mathcal{L}(\mathrm{F};\mathrm{E})\) 的元素，使得 \(u_{1} \equiv u\) 且 \(w_{1} \equiv w\)，那么 \(w_{1}\) 是 \(u_{1}\) 的拟逆，因为 \(w \circ u \equiv w_{1} \circ u_{1}\) 且 \(u \circ w \equiv u_{1} \circ w_{1}\)。

如果 w 与 \(w_{1}\) 都是 u 的拟逆，那么 \(w_{1} \equiv w\)，因为

\[
w _ {1} = 1 _ {\mathrm{E}} \circ w _ {1} \equiv w \circ u \circ w _ {1} \equiv w \circ 1 _ {\mathrm{F}} = w.
\]

定义 1. --- 设 E 与 F 是局部凸空间。我们称 \(\mathcal{L}(\mathrm{E};\mathrm{F})\) 中的一个元素 u 为弗雷德霍姆映射 \(^{(1)}\)，如果它拥有一个拟逆。从 E 到 E 中的弗雷德霍姆映射称为 E 的弗雷德霍姆自同态。

我们将用 \(\mathcal{F}(\mathrm{E};\mathrm{F})\) 表示从 E 到 F 中的弗雷德霍姆映射全体所成的集，并用 \(\mathcal{F}(\mathrm{E})\) 表示 E 的弗雷德霍姆自同态全体所成的集。

注记. --- 设 E、F 与 G 是局部凸空间，\(u: E \to F\) 与 \(v: F \to G\) 是连续线性映射。

\begin{enumerate}
\def\labelenumi{\arabic{enumi})}
\item
  假设 u 是弗雷德霍姆映射，并用 \(u_{1}\) 表示 u 的一个拟逆。由于 u 是 \(u_{1}\) 的拟逆，映射 \(u_{1}\) 是弗雷德霍姆映射。
\item
  假设 u 与 v 是弗雷德霍姆映射，并设 \(u_{1}\)（相应地，\(v_{1}\)）是 u（相应地，v）的拟逆。那么 \(v \circ u\) 是从 E 到 G 中的弗雷德霍姆映射，且 \(u_{1} \circ v_{1}\) 是 \(v \circ u\) 的拟逆。事实上，我们计算出
\end{enumerate}

\[
\left(u _ {1} \circ v _ {1}\right) \circ (v \circ u) = u _ {1} \circ \left(v _ {1} \circ v\right) \circ u \equiv u _ {1} \circ 1 _ {\mathrm{F}} \circ u = u _ {1} \circ u \equiv 1 _ {\mathrm{E}},
\]

\[
(v \circ u) \circ (u _ {1} \circ v _ {1}) = v \circ (u \circ u _ {1}) \circ v _ {1} \equiv v \circ 1 _ {\mathrm{F}} \circ v _ {1} = v \circ v _ {1} \equiv 1 _ {\mathrm{G}}.
\]

\begin{enumerate}
\def\labelenumi{\arabic{enumi})}
\setcounter{enumi}{2}
\tightlist
\item
  假设 u 与 \(v \circ u\) 是弗雷德霍姆映射，并设 \(w_{1}\) 是 \(v \circ u\) 的拟逆。那么 v 是弗雷德霍姆映射，且 \(u \circ w_{1}\) 是 v 的拟逆。
\end{enumerate}

事实上，由第一个注记，\(w_{1}\) 是弗雷德霍姆算子。设 \(u_{1}\) 是 u 的拟逆；由第二个注记，\(u \circ w_{1}\) 是弗雷德霍姆映射，且 \((v \circ u) \circ u_{1}\) 是一个拟逆。我们有 \(u \circ u_{1} \equiv 1_{F}\)，从而 \(v \circ u \circ u_{1} \equiv v\)；这就证明了断言。

\begin{enumerate}
\def\labelenumi{\arabic{enumi})}
\setcounter{enumi}{3}
\tightlist
\item
  假设 v 与 \(v \circ u\) 是弗雷德霍姆映射，并设 \(w_{1}\) 是 \(v \circ u\) 的拟逆。那么 u 是弗雷德霍姆算子，且 \(w_{1} \circ v\) 是 u 的拟逆。
\end{enumerate}

证明与前一个注记的证明类似。

引理 1. --- 设 E 与 F 是局部凸空间，u 是从 E 到 F 中的弗雷德霍姆映射。u 的核与余核都是有限维的。

设 \(v \colon \mathrm{F} \to \mathrm{E}\) 是 \(u\) 的拟逆。\(u\) 的核包含于有限秩线性映射 \(1_{\mathrm{E}} - v \circ u\) 的像之中，因而是有限维的。\(u\) 的像包含有限秩线性映射 \(1_{\mathrm{F}} - u \circ v\) 的核，因而在 F 中有有限余维。

命题 2. — 设 E 与 F 是分离的局部凸空间，并设 u 是 \(\mathcal{L}(\mathrm{E};\mathrm{F})\) 的一个元素。下列性质等价：

\begin{enumerate}
\def\labelenumi{(\roman{enumi})}
\item
  映射 u 是弗雷德霍姆算子；
\item
  映射 \(u\) 是严格态射，它的核是有限维的，它的像是闭的且在 F 中有有限余维；
\item
  存在 E 的闭向量子空间 \(E_{1}\) 与 F 的闭向量子空间 \(F_{1}\)，两者都有有限余维，使得 u 通过过渡到子空间诱导了从 \(E_{1}\) 到 \(F_{1}\) 上的同构；
\item
   存在拓扑直和分解 \(E = E_{1} \oplus E_{2}\) 与 \(F = F_{1} \oplus F_{2}\)，使得 \(E_{2}\) 与 \(F_{2}\) 是有限维的，u 在 \(E_{2}\) 上为零，并通过过渡到子空间定义了从 \(E_{1}\) 到 \(F_{1}\) 上的同构。
\item
  \(\Longrightarrow\) (iii)：设 v 是 u 的拟逆，\(E_{1}\) 是 \(1_{E}-v\circ u\) 的核，而 \(F_{1}\) 是 \(1_{F}-u\circ v\) 的核。由于线性映射 \(1_{E}-v\circ u\) 与 \(1_{F}-u\circ v\) 是连续的且是有限秩的，\(E_{1}\) 与 \(F_{1}\) 分别是 E 与 F 的有限余维闭向量子空间。设 \(x\in E_{1}\)。我们有
\end{enumerate}

\[
(1 _ {\mathrm{F}} - u \circ v) (u (x)) = u ((1 _ {\mathrm{E}} - v \circ u) (x)) = u (0) = 0,
\]

从而 \(u(x) \in \mathrm{F}_1\)。因此我们有 \(u(\mathrm{E}_1) \subset \mathrm{F}_1\)；同样我们有 \(v(\mathrm{F}_1) \subset \mathrm{E}_1\)。于是，连续线性映射 \(u_1: \mathrm{E}_1 \to \mathrm{F}_1\) 与 \(v_1: \mathrm{F}_1 \to \mathrm{E}_1\)（它们由 \(u\) 与 \(v\) 得到）是互逆的同构，因为 \(v \circ u\) 与 \(1_{\mathrm{E}}\)（相应地，\(u \circ v\) 与 \(1_{\mathrm{F}}\)）在 \(\mathrm{E}_1\) 上（相应地，在 \(\mathrm{F}_1\) 上）重合。

\begin{enumerate}
\def\labelenumi{(\roman{enumi})}
\setcounter{enumi}{2}
\item
  \(\Longrightarrow\) (ii)：设 \(E_{1}\) 与 \(F_{1}\) 满足 (iii) 的假设。我们有 \(\mathrm{E}_{1}\cap\mathrm{Ker}(u)=\{0\}\) 与 \(\mathrm{F}_{1}\subset\mathrm{Im}(u)\)，因而 \(\mathrm{Ker}(u)\) 是有限维的，而 \(\mathrm{Im}(u)\) 是闭的且在 F 中有有限余维。由 III, p. 39 的 prop. 1 可得映射 u 是严格态射。
\item
   \(\Longrightarrow\) (iv)：假设条件 (ii) 满足。E 的闭向量子空间 \(E_{2} = \operatorname{Ker}(u)\) 是有限维的，且存在 E 的向量子空间 \(E_{1}\)，它在拓扑上与 \(E_{2}\) 互补（EVT, II, p. 27, cor. 2）。向量子空间 \(F_{1} = \operatorname{Im}(u)\) 是闭的且有有限余维，它在 F 中有一个拓扑补 \(F_{2}\)。由 III, p. 39 的 prop. 1，映射 \(u|E_{1}\) 是严格态射，因而 u 诱导了从 \(E_{1}\) 到 \(F_{1}\) 上的同构。
\item
   \(\Longrightarrow\) (i)：在 (iv) 的假设下，从 F 到 E 中的、与 \(u_1^{-1}\) 在 \(\mathrm{F}_1\) 上重合且在 \(\mathrm{F}_2\) 上为零的线性映射是 \(u\) 的拟逆。
\end{enumerate}

注记 5. --- 设 E、F 是分离的局部凸空间，而 \(u\colon\mathrm{E}\to\mathrm{F}\) 是弗雷德霍姆算子。如果 \(u\) 是双射，那么 \(u\) 是同构（事实上，由 Prop. 2, (ii)，\(u\) 是严格态射）。

\subsection{弗雷德霍姆映射的指标}

设 E、F 与 G 是局部凸空间。

设 \(u \colon \mathrm{E} \to \mathrm{F}\) 是弗雷德霍姆算子。向量空间 \(\operatorname{Ker}(u)\) 与 \(\operatorname{Coker}(u)\) 是有限维的（III, p. 41 的引理 1）。回忆整数

\[
\dim \operatorname{Coker} (u) - \dim \operatorname{Ker} (u) = \operatorname{codim} _ {\mathrm{F}} \operatorname{Im} (u) - \dim \operatorname{Ker} (u) \tag {1}
\]

称为 u 的指标，并记作 ind(u)（A, V, p. 126）。

如果 \(u \colon \mathrm{E} \to \mathrm{F}\) 与 \(v \colon \mathrm{F} \to \mathrm{G}\) 是弗雷德霍姆映射，那么 \(v \circ u\) 也是（III, p. 41, 注记 2），并且我们有（A, V, p. 127, 引理 2）

\[
\operatorname{ind} (v \circ u) = \operatorname{ind} (v) + \operatorname{ind} (u).\tag{2}
\]

假设 E 与 F 是分离的，并设 \(u\colon\mathrm{E}\to\mathrm{F}\) 是弗雷德霍姆算子；采用 III, p. 42 的 Prop. 2 的条件 (iv) 的记号。我们于是有 \(\mathrm{ind}(u)=\mathrm{dim}(\mathrm{F}_{2})-\mathrm{dim}(\mathrm{E}_{2})\)。给这些空间中每一个的对偶赋以弱拓扑（相应地，紧开拓扑、有界收敛拓扑）。那么 \(\mathrm{E}^{\prime}\) 等同于 \(\mathrm{E}_{1}^{\prime}\) 与 \(\mathrm{E}_{2}^{\prime}\) 的拓扑直和，\(\mathrm{F}^{\prime}\) 等同于 \(\mathrm{F}_{1}^{\prime}\) 与 \(\mathrm{F}_{2}^{\prime}\) 的拓扑直和，且 \(^{t}u\) 诱导了从 \(\mathrm{F}_{1}^{\prime}\) 到 \(\mathrm{E}_{1}^{\prime}\) 上的同构并在 \(\mathrm{F}_{2}^{\prime}\) 上为零。于是转置 \(^{t}u\colon\mathrm{F}^{\prime}\to\mathrm{E}^{\prime}\) 是弗雷德霍姆映射（loc. cit.）。 \(^{t}u\) 的核是 \(\mathrm{F}_{2}^{\prime}\)，它的维数等于 \(\mathrm{F}_{2}\) 的维数，也就是 \(u\) 的余核的维数。从而

\[
\operatorname{ind} (u) = \dim \operatorname{Ker} (^ {t} u) - \dim \operatorname{Ker} (u).\tag{3}
\]

此外，\(^{t}u\) 的像是 \(E_{1}^{\prime}\)，因而 \(^{t}u\) 的余核的维数等于 \(E_{2}^{\prime}\) 的维数，也就是 u 的核 \(E_{2}\) 的维数。由此我们推出

\[
\operatorname{ind} (^ {t} u) = - \operatorname{ind} (u).\tag{4}
\]

假设 E 与 F 是分离的，并设 \(u\colon\mathrm{E}\to\mathrm{F}\) 是指标为 0 的弗雷德霍姆算子。那么由 III, p. 42 的 prop. 2，\(u\) 是严格态射。由于 \(\dim\operatorname{Ker}(u)=\operatorname{codim}_{\mathrm{F}}\operatorname{Im}(u)\)，映射 \(u\) 只要是单射或满射就是同构。

命题 3. --- 设 E 与 F 是局部凸空间，且 \(u \in \mathcal{L}(\mathrm{E};\mathrm{F})\)。设 \(\mathrm{E}_{1}\)（相应地，\(\mathrm{F}_{1}\)）是 E（相应地，F）的有限余维闭子空间。假设 u 把 \(\mathrm{E}_{1}\) 映到 \(\mathrm{F}_{1}\) 中，并用 \(u_{1} \in \mathcal{L}(\mathrm{E}_{1};\mathrm{F}_{1})\) 表示在 \(\mathrm{E}_{1}\) 上与 u 重合的映射。要使 u 是弗雷德霍姆算子，必须且只需 \(u_{1}\) 是。这时我们有

\[
\operatorname{ind} (u) - \operatorname{ind} \left(u _ {1}\right) = \operatorname{codim} _ {\mathrm{F}} \left(\mathrm{F} _ {1}\right) - \operatorname{codim} _ {\mathrm{E}} \left(\mathrm{E} _ {1}\right).\tag{5}
\]

设 \(i\colon \mathrm{E}_{1}\to \mathrm{E}\) 与 \(j\colon \mathrm{F}_{1}\to \mathrm{F}\) 是典范单射。它们是弗雷德霍姆映射，并且有

\[
\operatorname{ind} (i) = \operatorname{codim} _ {\mathrm{E}} (\mathrm{E} _ {1}), \quad \operatorname{ind} (j) = \operatorname{codim} _ {\mathrm{F}} (\mathrm{F} _ {1}).
\]

由于 \(j \circ u_{1} = u \circ i\)，可以看出 u 是弗雷德霍姆映射当且仅当 \(u_{1}\) 是（III, p. 41 的注记 3 和 4）。若是如此，我们有

\[
\operatorname{ind} (j) + \operatorname{ind} (u _ {1}) = \operatorname{ind} (j \circ u _ {1}) = \operatorname{ind} (u \circ i) = \operatorname{ind} (u) + \operatorname{ind} (i),
\]

由此得到公式 (5)。

命题 4. --- 设 E 与 F 是分离的局部凸空间，\(u: \mathrm{E} \to \mathrm{F}\) 是弗雷德霍姆算子，而 \(\widehat{u}: \widehat{\mathrm{E}} \to \widehat{\mathrm{F}}\) 是 u 到完备化上的延拓。那么 \(\widehat{u}\) 是弗雷德霍姆算子，并且有 \(\operatorname{Ker}(\widehat{u}) = \operatorname{Ker}(u)\) 与 \(\operatorname{ind}(\widehat{u}) = \operatorname{ind}(u)\)。

采用 III, p. 42 的 Prop. 2 的条件 (iv) 的记号。由于向量空间 \(E_{2}\) 与 \(F_{2}\) 是有限维的，它们是完备的。\(E_{1}\)（相应地，\(F_{1}\)）的完备化等同于 \(E_{1}\) 在 \(\widehat{E}\) 中的闭包（相应地，\(F_{1}\) 在 \(\widehat{F}\) 中的闭包），且 \(\widehat{E}\)（相应地，\(\widehat{F}\)）是 \(\widehat{E}_{1}\) 与 \(E_{2}\)（相应地，\(\widehat{F}_{1}\) 与 \(F_{2}\)）的拓扑直和。线性映射 \(\widehat{u}\) 通过限制定义了从 \(\widehat{E}_{1}\) 到 \(\widehat{F}_{1}\) 上的同构并在 \(E_{2}\) 上为零。命题于是由 loc. cit. 的蕴含关系 (iv)⇒(i) 得出。

\subsection{里斯自同态}

设 E 是向量空间，u 是 E 的自同态。

E 的向量子空间的序列 \((\operatorname{Ker}(u^{n}))_{n\in\mathbf{N}}\) 是递增的；它们的并是 E 的一个在 u 下稳定的子空间，称为 u 的幂零空间。

E 的向量子空间的序列 \((\mathrm{Im}(u^{n}))_{n\in\mathbf{N}}\) 是递减的；它们的交是 E 的一个在 u 下稳定的子空间，称为 u 的余幂零空间。

引理 2. --- 设 E 是向量空间，u 是 E 的一个有指标的自同态（A, V, p. 126）。

如果下列性质中有两条满足，那么第三条也满足：

\begin{enumerate}
\def\labelenumi{(\roman{enumi})}
\item
  自同态 u 的指标为 0；
\item
  u 的幂零空间 N 是有限维的；
\item
  u 的余幂零空间 I 有有限余维。
\end{enumerate}

对每个整数 \(n \geqslant 0\)，自同态 \(u^{n}\) 有指标，等于 \(n \operatorname{ind}(u)\)（A, V, p. 127, 引理 2）。自然整数序列 \((\dim(\operatorname{Ker}(u^{n})))_{n \in \mathbf{N}}\) 是递增的；要使它是平稳的，必须且只需 u 的幂零空间是有限维的。自然整数序列 \((\operatorname{codim}_{\operatorname{E}}(\operatorname{Im}(u^{n})))_{n \in \mathbf{N}}\) 是递增的；要使它是平稳的，必须且只需 u 的余幂零空间是有限维的。假设条件 (i) 满足；关系式

\[
\dim (\operatorname{Ker} (u ^ {n})) - \operatorname{codim} _ {\mathrm{E}} (\operatorname{Im} (u ^ {n})) = \operatorname{ind} (u ^ {n}) = n \operatorname{ind} (u) = 0
\]

于是蕴含条件 (ii) 与 (iii) 等价。反之，如果条件 (ii) 与 (iii) 满足，这个公式蕴含序列 \((n\operatorname{ind}(u))_{n\in\mathbf{N}}\) 是平稳的，因此 u 的指标为零。

引理 3. --- 设 E 是向量空间，u 是 E 的一个有指标的自同态。假设 u 的指标为零。设 N 是 u 的幂零空间，I 是它的余幂零空间。

\begin{enumerate}
\def\labelenumi{\alph{enumi})}
\item
  这时空间 E 是子空间 N 与 I 的直和（维尔--菲廷分解，参见 A, VIII, § 2, 第 2 目, p. 25）。
\item
  自同态 u 通过过渡到子空间，诱导了 N 的一个幂零自同态与 I 的一个自同构 \(u_{I}\)。
\item
  设 v 表示 E 的这样一个自同态：它在 N 上为零且在 I 上与 \(u_{I}^{-1}\) 重合。我们有 \(u \circ v = v \circ u\)，且自同态 \(1_{E} - u \circ v\) 是以 N 为像、以 I 为核的投影。
\item
  E 的每个与 u 交换的自同态都稳定 N 与 I，且与 v 交换。
\end{enumerate}

自同态 \(u\) 满足前一个引理的各条性质。

设 \(n \in \mathbf{N}\) 使得 \(\operatorname{Ker}(u^{n}) = N\)。那么 \(\operatorname{Ker}(u^{m}) = \operatorname{Ker}(u^{n})\) 对每个整数 \(m \geqslant n\) 成立；由于 \(\operatorname{ind}(u^{m}) = m \operatorname{ind}(u) = 0\)，这就蕴含 \(\operatorname{Im}(u^{m}) = \operatorname{Im}(u^{n})\) 对 \(m \geqslant n\) 成立，因而 \(\operatorname{I} = \operatorname{Im}(u^{n})\)。由 A, VIII, p. 25, prop. 2 可得断言 a) 与 b) 满足。

对每个 \(x \in N\)，我们有 \(v(x) = 0\) 与 \(u(x) \in N\)，因而 \(u(v(x)) = 0\) 与 \(v(u(x)) = 0\)。对每个 \(x \in I\)，我们有 \(u(v(x)) = v(u(x)) = x\)。这就蕴含 \(u \circ v = v \circ u\)，且 \(1_{E} - u \circ v\) 是以 N 为像、以 I 为核的投影。

设 w 是 E 的与 u 交换的自同态。设 \(n \in \mathbf{N}\) 使得 \(N = \operatorname{Ker}(u^{n})\) 且 \(I = \operatorname{Im}(u^{n})\)。设 \(x \in N\)；我们有 \(u^{n}(w(x)) = w(u^{n}(x)) = 0\)，因而 \(w(x) \in N\)。设 \(x \in I\)；存在 \(y \in E\) 使得 \(x = u^{n}(y)\)；那么 \(w(x) = w(u^{n}(y)) = u^{n}(w(y)) \in I\)。这就证明了 w 稳定 N 与 I。

最后证明 \(w\) 与 \(v\) 可交换。对每个 \(x\)（在 \(N\) 中），我们有 \(v(w(x)) = 0 = w(v(x))\)。设 \(x \in I\)；\(v(w(x))\) 与 \(w(v(x))\) 这两个 \(E\) 的元素都属于 \(I\)，且它们在 \(u\) 下的像都等于 \(w(x)\)，因而它们相等。由线性性，这就蕴含 \(v\) 与 \(w\) 交换。

注记. --- 假设存在整数 \(n \geqslant 0\) 使得 \(\operatorname{Ker}(u^{n}) = \operatorname{Ker}(u^{n+1})\)（相应地，\(\operatorname{Im}(u^{n}) = \operatorname{Im}(u^{n+1})\)）；那么 \(\operatorname{Ker}(u^{m}) = \operatorname{Ker}(u^{n})\)（相应地，\(\operatorname{Im}(u^{m}) = \operatorname{Im}(u^{n})\)）对每个整数 \(m \geqslant n\) 成立。

假设 \(u\) 的指标为零。对每个整数 \(n \geqslant 0\)，我们有

\[
n \operatorname{ind} (u) = \operatorname{ind} \left(u ^ {n}\right) = \operatorname{codim} _ {\mathrm{E}} \left(\operatorname{Im} \left(u ^ {n}\right)\right) - \dim \left(\operatorname{Ker} \left(u ^ {n}\right)\right).
\]

因此条件 \(\operatorname{Ker}(u^n) = \operatorname{Ker}(u^{n+1})\) 与 \(\operatorname{Im}(u^n) = \operatorname{Im}(u^{n+1})\) 等价。当它们满足时，向量空间 E 是 \(\operatorname{Ker}(u^n)\) 与 \(\operatorname{Im}(u^n)\) 的直和（A, VIII, p. 25, prop. 2），空间 \(\operatorname{Ker}(u^n)\) 是有限维的，且 \(u\) 通过过渡到子空间诱导了 \(\operatorname{Im}(u^n)\) 的一个自同构。

定义 2. --- 设 E 是分离的局部凸空间。E 的里斯自同态指的是 E 的这样的弗雷德霍姆自同态：其幂零空间是有限维的且其余幂零空间有有限余维。

命题 5. --- 设 E 是分离的局部凸空间。E 的每个里斯自同态的指标都为零。

这由定义和引理 2 得出。

例子. --- 1) E 的每个自同构都是 E 的里斯自同态。如果 E 是有限维的，\(\mathcal{L}(\mathrm{E})\) 的每个元素都是 E 的里斯自同态。

\begin{enumerate}
\def\labelenumi{\arabic{enumi})}
\setcounter{enumi}{1}
\item
   如果 \((\mathrm{E}_{i})_{i\in\mathrm{I}}\) 是分离局部凸空间的有限族，且对 I 中所有 i，\(u_{i}\) 是 \(\mathcal{L}(\mathrm{E}_{i})\) 的元素，那么自同态 \(u=\bigoplus_{i\in\mathrm{I}}u_{i}\)（它定义在局部凸空间 \(\bigoplus_{i\in\mathrm{I}}E_{i}\) 上）是里斯自同态当且仅当对所有 i，\(u_{i}\) 是。
\item
   设 E 是 R 上的分离局部凸空间，并设 \(\mathrm{E}_{(\mathbf{C})}\) 是它的复化。设 u 是 E 的自同态。要使 u 是里斯自同态，必须且只需它的复化 \(u_{(\mathbf{C})}\) 是 \(\mathrm{E}_{(\mathbf{C})}\) 的里斯自同态。
\end{enumerate}

命题 6. --- 设 E 是分离的局部凸空间，u 是 E 的里斯自同态。用 N 表示 u 的幂零空间，用 I 表示它的余幂零空间。

\begin{enumerate}
\def\labelenumi{\alph{enumi})}
\item
  E 的子空间 N 与 I 是闭的且在 u 下稳定，而空间 E 是它们的拓扑直和。
\item
  自同态 u 通过限制定义了 I 的一个自同构 \(u_{I}\)；
\item
   向量空间 N 是有限维的，且 u 通过限制定义了 N 的一个幂零自同态 \(u_{N}\)；
\item
  设 v 是 \(\mathcal{L}(\mathrm{E})\) 中这样的元素：它在 N 上为零且在 I 上与 \(u_{\mathrm{I}}^{-1}\) 重合。它是 E 的里斯自同态，且是与 u 交换的 u 的拟逆。E 的自同态 \(1_{\mathrm{E}} - u \circ v\) 是以 N 为像、以 I 为核的投影。\(\mathcal{L}(\mathrm{E})\) 中每个与 u 交换的元素都稳定 N 与 I，且与 v 交换。
\end{enumerate}

设 n 是自然数，使得 \(\operatorname{Ker}(u^{n}) = \operatorname{N}\) 且 \(\operatorname{Im}(u^{n}) = \operatorname{I}\)。由于 E 是分离的且 u 是连续的，\(\operatorname{N} = \operatorname{Ker}(u^{n})\) 在 E 中是闭的。由于 u 是弗雷德霍姆自同态，\(u^{n}\) 也是，且 \(\operatorname{I} = \operatorname{Im}(u^{n})\) 是闭的（III, p. 42, prop. 2）。由引理 2，空间 E 是 N 与 I 的直和，断言 a) 于是由 EVT, I, p. 15 的 prop. 3 得出。

由于 I 是闭的、在 E 中有有限余维且在 u 下稳定，由 u 诱导的映射 \(u_{I}\colon I \to I\) 是弗雷德霍姆自同态（III, p. 44, prop. 3）；它还是双射，因而是 I 的自同构（III, p. 42, prop. 2）。这就证明了 b)，其余断言由引理 2 得出。

自同态 \(v\)（它在命题的断言 \(d\)）中定义）是 \(u\) 的一个拟逆，称为 \(u\) 的典范拟逆。

里斯自同态享有弗雷德霍姆映射的某些稳定性性质。

命题 7. — 设 E 是分离的局部凸空间，u 是 E 的自同态。

\begin{enumerate}
\def\labelenumi{\alph{enumi})}
\item
   设 \(E_{1}\) 是 E 的有限余维闭子空间，且在 u 下稳定。设 \(u_{1}\) 是 \(\mathcal{L}(E_{1})\) 中在 \(E_{1}\) 上与 u 重合的元素。那么 u 是里斯自同态当且仅当 \(u_{1}\) 是；
\item
   假设 u 是指标为 0 的弗雷德霍姆自同态。设 \(\widehat{\mathbf{E}}\) 是 E 的完备化，并设 \(\widehat{u} \in \mathcal{L}(\widehat{\mathbf{E}})\) 是 u 的连续延拓。那么 u 是里斯自同态当且仅当 \(\widehat{u}\) 是；
\item
 给 E 的对偶 E' 赋以有界收敛拓扑、紧收敛拓扑或弱拓扑。如果 u 是 E 的里斯自同态，那么 \(^{t}\) u 是 E' 的里斯自同态。
\end{enumerate}
a) 要使 \(u\) 是指标为零的弗雷德霍姆自同态，必须且只需 \(u_1\) 是（III, p. 44, prop. 3）。另一方面，对每个 \(n \in \mathbf{N}\)，我们有 \(\mathrm{Ker}(u_1^n) = \mathrm{E}_1 \cap \mathrm{Ker}(u^n)\)，由此


\[
\dim \operatorname{Ker} (u _ {1} ^ {n}) \leqslant \dim \operatorname{Ker} (u ^ {n}) \leqslant \dim \operatorname{Ker} (u _ {1} ^ {n}) + \operatorname{codim} _ {\mathrm{E}} (\mathrm{E} _ {1}),
\]

因此序列 \((\dim\operatorname{Ker}(u^{n}))_{n\in\mathbf{N}}\) 有界当且仅当序列 \((\dim\operatorname{Ker}(u_{1}^{n}))_{n\in\mathbf{N}}\) 有界。于是，u 是 E 的里斯自同态当且仅当 \(u_{1}\) 是 \(E_{1}\) 的里斯自同态（引理 2）。

\begin{enumerate}
\def\labelenumi{\alph{enumi})}
\setcounter{enumi}{1}
\item
  由 III, p. 44 的 prop. 4，对每个 \(n \in \mathbf{N}\)，映射 \(\widehat{u}^n\) 是 \(\widehat{\mathrm{E}}\) 的弗雷德霍姆自同态，满足 \(\operatorname{Ker}(\widehat{u}^n) = \operatorname{Ker}(u^n)\) 且 \(\operatorname{ind}(\widehat{u}^n) = \operatorname{ind}(u^n) = n\operatorname{ind}(u)\)。于是，\(u\) 是里斯自同态当且仅当 \(\widehat{u}\) 是里斯自同态。
\item
  转置 \({}^{t}u\) 是 E' 的指标为 0 的弗雷德霍姆自同态（III, p. 43, 第 3 目）。由于 \((^{t}u)^{n}\) 的核是 \(u^{n}\) 的像的正交补（EVT, IV, p. 27, prop. 2），序列 \((\mathrm{Ker}({}^{t}u)^{n})_{n}\) 是平稳的。这就证明了 \({}^{t}u\) 是 E' 的里斯自同态。
\end{enumerate}

命题 8. --- 设 E 是分离的局部凸空间，u 是 E 的自同态。下列条件等价：

\begin{enumerate}
\def\labelenumi{(\roman{enumi})}
\item
  u 是 E 的里斯自同态；
\item
  存在 u 的拟逆 v，它与 u 交换；
\item
  存在 E 的闭向量子空间 \(E_{1}\)，它有有限余维且在 u 下稳定，使得 u 诱导了 \(E_{1}\) 的一个自同构。
\end{enumerate}

蕴含关系 (i) \(\Longrightarrow\) (ii) 由 prop. 6, d) 得出。

假设 \(u\) 有一个拟逆 \(v\)，它与 \(u\) 交换。设 \(\mathrm{E}_1\) 是 \(1_{\mathrm{E}} - u \circ v\) 的核；它是 E 的闭向量子空间，有有限余维，且在 \(u\) 与 \(v\) 下稳定。由于 \(u(v(x)) = v(u(x)) = x\) 对一切 \(x \in \mathrm{E}_1\) 成立，映射 \(u\) 与 \(v\) 诱导了 \(\mathrm{E}_1\) 的互逆的自同构。因此 (ii) 蕴含 (iii)。

最后，蕴含关系 (iii) \(\Longrightarrow\) (i) 由 Prop. 7 的断言 \(a\)）以及自同构是里斯自同态这一事实得出。

命题 9. --- 设 E 是分离的局部凸空间，u、v 是 \(\mathcal{L}(\mathrm{E})\) 中可交换的元素。下列条件等价：

\begin{enumerate}
\def\labelenumi{(\roman{enumi})}
\item
  自同态 u 与 v 都是里斯自同态。
\item
  自同态 \(u \circ v\) 是里斯自同态。
\end{enumerate}

假设 u 与 v 是里斯自同态；用 \(u'\) 与 \(v'\) 表示它们的典范拟逆。E 的自同态 u、v、\(u'\) 与 \(v'\) 两两交换（prop. 6, d)）。自同态 \(v' \circ u'\) 是 \(u \circ v\) 的拟逆且与 \(u \circ v\) 交换，因而 \(u \circ v\) 是里斯自同态（prop. 8）。

反之，假设 \(u \circ v\) 是里斯自同态，并设 w 是它的典范拟逆。由于 u 与 \(u \circ v\) 交换，由 prop. 6, d)，自同态 u 与 w 交换。同样 v 与 w 交换。因此 E 的自同态 u、v 与 w 两两交换；由此得出 E 的自同态 \(v \circ w\) 是 u 的拟逆，而 \(w \circ u\) 是 v 的拟逆。由 prop. 8，u 与 v 是里斯自同态。

\pandocbounded{\includegraphics[keepaspectratio]{images/part001_6cebf9a3.jpg}}

设 u 是 E 的里斯自同态。如果 v 是 E 的自同构，那么 E 的自同态 \(u \circ v\) 不一定是 E 的里斯自同态，即使 \(u \circ v - v \circ u\) 有有限秩也是如此。如果 h 是 E 的有限秩自同态，自同态 \(u + h\) 也不一定是里斯自同态（参见 III, p. 120, 习题 3）。

命题 10. --- 设 E 与 F 是局部凸空间，\(p \in \mathcal{L}(\mathrm{E};\mathrm{F})\) 且 \(q \in \mathcal{L}(\mathrm{F};\mathrm{E})\)。置 \(u = 1_{\mathrm{E}} - q \circ p\) 与 \(v = 1_{\mathrm{F}} - p \circ q\)。

\begin{enumerate}
\def\labelenumi{\alph{enumi})}
\item
  设 \(n \in \mathbf{N}\)。映射 \(p\) 通过限制定义了从 \(\mathrm{Ker}(u^n)\) 到 \(\mathrm{Ker}(v^n)\) 上的拓扑向量空间同构，并通过过渡到商定义了从 Coker( \(u^n\) ) 到 Coker( \(v^n\) ) 上的拓扑向量空间同构；
\item
  如果 u 的像在 E 中是闭的，那么 v 的像在 F 中是闭的；
\item
   如果 \(\operatorname{Ker}(u)\) 在 E 中有一个拓扑补，那么 \(\operatorname{Ker}(v)\) 在 F 中也有一个。如果 \(\operatorname{Im}(u)\) 在 E 中有一个拓扑补，那么 \(\operatorname{Im}(v)\) 在 F 中也有一个；
\item
  如果 u 是严格态射，那么 v 是严格态射。
\item
  如果 u 是 E 的自同构，那么 v 是 F 的自同构。
\item
  如果 \(u\) 是 E 的弗雷德霍姆自同态，那么 \(v\) 是 F 的弗雷德霍姆自同态，且我们有 \(\operatorname{ind}(u) = \operatorname{ind}(v)\)；
\item
  假设 E 与 F 是分离的。如果 u 是 E 的里斯自同态，那么 v 是 F 的里斯自同态。
\end{enumerate}

我们首先注意公式

\[
q \circ v = u \circ q, \quad v \circ p = p \circ u.\tag{6}
\]

然后我们证明两个引理。

引理 4. --- 设 \(n \in \mathbf{N}\)。置

\[
q _ {n} = \sum_ {i = 1} ^ {n} (- 1) ^ {i - 1} \binom {n} {i} q \circ (p \circ q) ^ {i - 1}.
\]

我们有 \(u^{n}=1_{E}-q_{n}\circ p\) 与 \(v^{n}=1_{F}-p\circ q_{n}\)。

\((q\circ p)^{i} = q\circ (p\circ q)^{i - 1}\circ p\) 对每个整数 \(i\geqslant 1\) 成立。然后我们计算

\[
u ^ {n} = (1 _ {\mathrm{E}} - q \circ p) ^ {n} = 1 _ {\mathrm{E}} - \sum_ {i = 1} ^ {n} (- 1) ^ {i - 1} {\binom {n} {i}} (q \circ p) ^ {i} = 1 _ {\mathrm{E}} - q _ {n} \circ p
\]

\[
v ^ {n} = (1 _ {\mathrm{F}} - p \circ q) ^ {n} = 1 _ {\mathrm{F}} - \sum_ {i = 1} ^ {n} (- 1) ^ {i - 1} \binom {n} {i} (p \circ q) ^ {i} = 1 _ {\mathrm{F}} - p \circ q _ {n},
\]

由此得到结论。

引理 5. --- 设 \(u'\) 是 \(\mathcal{L}(\mathrm{E})\) 的一个元素。置 \(v' = 1_{\mathrm{F}} + p \circ u' \circ q\)。如果 \(u'\) 是 \(u\) 的逆，那么 \(v'\) 是 \(v\) 的逆。如果 \(u'\) 是 \(u\) 的拟逆，那么 \(v'\) 是 \(v\) 的拟逆。如果 \(u'\) 与 \(u\) 交换，那么 \(v'\) 与 \(v\) 交换。

利用公式 (6)，我们计算

\[
v - v ^ {\prime} \circ v = - p \circ u ^ {\prime} \circ q \circ v = - p \circ u ^ {\prime} \circ u \circ q
\]

\[
v - v \circ v ^ {\prime} = - v \circ p \circ u ^ {\prime} \circ q = - p \circ u \circ u ^ {\prime} \circ q,
\]

由此得

\[
1 _ {\mathrm{F}} - v ^ {\prime} \circ v = p \circ (1 _ {\mathrm{E}} - u ^ {\prime} \circ u) \circ q
\]

\[
1 _ {\mathrm{F}} - v \circ v ^ {\prime} = p \circ (1 _ {\mathrm{E}} - u \circ u ^ {\prime}) \circ q,
\]

这就证明了引理。

有了引理 4，只需证明断言 \(a\)）对 \(n = 1\) 成立。依据公式 (6)，映射 \(p\) 与 \(q\) 通过过渡到子空间定义了连续线性映射 \(p'\) ： \(\operatorname{Ker}(u) \to \operatorname{Ker}(v)\) 与 \(q'\) ： \(\operatorname{Ker}(v) \to \operatorname{Ker}(u)\)，并通过过渡到商定义了连续线性映射 \(p''\) ： \(\operatorname{Coker}(u) \to \operatorname{Coker}(v)\) 与 \(q''\) ： \(\operatorname{Coker}(v) \to \operatorname{Coker}(u)\)。

由于 \(1_{\mathrm{E}} - u = q \circ p\) 且 \(1_{\mathrm{F}} - v = p \circ q\)，映射 \(p'\) 与 \(q'\) 一方面、\(p''\) 与 \(q''\) 另一方面，是互逆的同构。这就证明了 \(a\)）。

要使一个连续线性映射的像是闭的，必须且只需这个映射的余核是分离空间。于是 b) 由 a) 得出。

用 \(i\colon \operatorname{Ker}(u)\to \mathrm{E}\) 与 \(j\colon \operatorname{Ker}(v)\to \mathrm{F}\) 表示典范单射。假设 \(\operatorname{Ker}(u)\) 在 E 中有一个拓扑补，并设 \(r\colon \mathrm{E}\to \operatorname{Ker}(u)\) 是 \(i\) 的一个连续线性收缩。置 \(r' = p'\circ r\circ q\)。我们有 \(r'\in \mathcal{L}(\mathrm{F};\operatorname {Ker}(v))\) 且

\[
r ^ {\prime} \circ j = p ^ {\prime} \circ r \circ q \circ j = p ^ {\prime} \circ r \circ i \circ q ^ {\prime} = p ^ {\prime} \circ q ^ {\prime} = 1 _ {\operatorname{Ker} (v)},
\]

因而 \(r'\) 是 j 的一个连续线性收缩，且 \(\operatorname{Ker}(v)\) 在 F 中有一个拓扑补。这就证明了 c) 的第一个断言。第二个断言类似地证明，注意如果 \(s: \operatorname{Coker}(u) \to \operatorname{E}\) 是从 E 到 \(\operatorname{Coker}(u)\) 上的典范满射的一个连续线性截面，那么 \(p \circ s \circ q''\) 是从 F 到 \(\operatorname{Coker}(v)\) 上的典范满射的一个连续线性截面。

用 \(\overline{u}\colon\mathrm{E}/\mathrm{Ker}(u)\to\mathrm{Im}(u)\) 与 \(\overline{v}\colon\mathrm{F}/\mathrm{Ker}(v)\to\mathrm{Im}(v)\) 表示由 u 与 v 得到的连续线性双射。由公式 (6)，映射 p 通过过渡到商定义了一个连续线性映射 \(p_{1}\colon\mathrm{E}/\mathrm{Ker}(u)\to\mathrm{F}/\mathrm{Ker}(v)\)，且 p、q 通过过渡到子空间定义了连续线性映射 \(p_{0}\colon\mathrm{Im}(u)\to\mathrm{Im}(v)\) 与 \(q_{0}\colon\mathrm{Im}(v)\to\mathrm{Im}(u)\)。

用 \(t \colon \operatorname{Im}(v) \to \operatorname{F} / \operatorname{Ker}(v)\) 表示典范单射 \(\operatorname{Im}(v) \to \operatorname{F}\) 与典范满射 \(\operatorname{F} \to \operatorname{F} / \operatorname{Ker}(v)\) 的复合映射；它是连续的。我们有 \(\overline{v} \circ t = 1_{\operatorname{Im}(v)} - p_0 \circ q_0\) 与 \(\overline{v} \circ p_1 = p_0 \circ \overline{u}\)，由此

\[
\overline {{v}} \circ (t + p _ {1} \circ \overline {{u}} ^ {- 1} \circ q _ {0}) = \overline {{v}} \circ t + p _ {0} \circ q _ {0} = 1 _ {\mathrm{Im} (v)}.
\]

这证明了 \(t + p_{1} \circ \overline{u}^{-1} \circ q_{0}\) 是 \(\overline{v}\) 的逆双射。如果 u 是严格态射，那么映射 \(\overline{u}^{-1}\) 是连续的，\(\overline{v}^{-1}\) 亦然，这就证明 v 是严格态射。这就证明了 d)。

由引理 5，如果 \(u\) 是 E 的自同构，那么 \(v\) 是 F 的自同构。类似地，如果 \(u\) 是 E 的弗雷德霍姆自同态，那么 \(v\) 是 F 的弗雷德霍姆自同态，且这时有 \(\operatorname{ind}(u) = \operatorname{ind}(v)\)，依据 \(a\)）。这就证明了 \(e\)）与 \(f\)）。最后，如果 E 与 F 是分离的且 \(u\) 是 E 的里斯自同态，那么由引理 5 与 Prop. 8 的条件 (ii)，\(v\) 是 F 的里斯自同态；这就是断言 \(g\)）。

\subsection{弗雷歇空间之间的弗雷德霍姆映射与里斯映射}

命题 11. --- 设 E 与 F 为弗雷歇空间。为使从 E 到 F 中的连续线性映射 u 是弗雷德霍姆映射，必须且只需它的核与余核都是有限维的。

这由 III, p. 42 的 Prop. 2 的条件 (ii) 所给出的弗雷德霍姆映射的特征刻画以及下面的引理得出：

引理 6. --- 设 E 与 F 为弗雷歇空间，且 \(u \in \mathcal{L}(\mathrm{E};\mathrm{F})\)。假设 \(u\) 的像在 F 中有有限余维。那么 \(u\) 的像在 F 中是闭的，且 \(u\) 是严格态射。

设 G 是 F 的一个向量子空间，它是 \(\operatorname{Im}(u)\) 的补；子空间 G 是闭的（EVT, I, p. 14, cor. 1），因而商空间 F/G 是弗雷歇空间（TG, IX, p. 25, prop. 4）。设 \(\pi: F \to F/G\) 是典范满射。映射 \(\pi \circ u\) 是满射，因而是严格的（EVT, I, p. 19, cor. 3）。断言随即由 III, p. 40 的 cor. 2 得出。

命题 12. --- 设 E 与 F 为弗雷歇空间。给它们的对偶 E' 与 F' 赋以弱拓扑、紧收敛拓扑或有界收敛拓扑。设 \(u \in \mathcal{L}(\mathrm{E};\mathrm{F})\)。为使 u 是从 E 到 F 的弗雷德霍姆映射，必须且只需 \({}^{t}u\) 是从 F' 到 E' 的弗雷德霍姆映射。

该条件是必要的（III, p. 43, 第 3 目）。我们来证明它是充分的。如果 \(^{t}u\) 是弗雷德霍姆映射，那么它是严格态射（III, p. 42, prop. 2），因而 u 是有闭像的严格态射（EVT, IV, p. 28, th. 1 及 p. 29, cor. 3）。考虑到典范同构 Coker(\(^{t}u\)) → Ker(u)' 与 Ker(\(^{t}u\)) → Coker(u)'（EVT, IV, p. 27, prop. 2），向量空间 Ker(u) 与 Coker(u) 是有限维的，且 u 是弗雷德霍姆算子（III, p. 42 的 prop. 2）。

命题 13. --- 设 E 为弗雷歇空间，且 \(u \in \mathcal{L}(\mathrm{E})\)。

\begin{enumerate}
\def\labelenumi{\alph{enumi})}
\item
  自同态 u 是里斯自同态，当且仅当它的幂零空间 N 是有限维的且它的余幂零空间 I 有有限余维；
\item
 给 E' 赋以有界收敛拓扑、紧收敛拓扑或弱拓扑。如果 \(^{t}\) u 是 E' 的里斯自同态，那么 u 是 E 的里斯自同态。
\end{enumerate}
a) 假设幂零空间 N 是有限维的且余幂零空间 I 有有限余维。由于 \(\operatorname{Ker}(u) \subset \operatorname{N}\) 与 \(\operatorname{I} \subset \operatorname{Im}(u)\)，u 的核是有限维的，且 u 的像有有限余维。由定义与命题 11，映射 u 是里斯自同态。逆命题由定义得出。

b) 由命题 12，该假设蕴含 \(u\) 是弗雷德霍姆自同态，其指标为 \(\text{ind}(^t u) = -\text{ind}(u) = 0\)（\(n^o 3\)，公式 (4)）。设 \(n \in \mathbf{N}\)。由于 \(u^n\) 的像是闭的且以 \(^t u^n\) 的核为其正交补（EVT, IV, p. 27, prop. 2），序列 \((\text{Im}(u^n))\) 是平稳的，因而 \(u\) 是里斯自同态。


\subsection{里斯自同态的谱特征刻画}

设 E 是复巴拿赫空间，u 是 E 的自同态。回忆我们用 \(\mathrm{Sp}(u)\) 表示 u 相对于含幺巴拿赫代数 \(\mathcal{L}(\mathrm{E})\) 的谱（参见 I, p. 127 的 § 7）。

假设 0 是 \(\mathrm{Sp}(u)\) 的一个孤立点；回忆这时我们用 \(e_0(u)\) 表示与 \(u\) 以及开闭子集 \(\{0\}\)（它是 \(u\) 的谱的子集）相伴的谱投影（参见 I, p. 131 的第 3 目）。我们有 \(e_0(u) = f(u)\)，它对每个全纯函数芽 \(f\)（在 \(\mathrm{Sp}(u)\) 的邻域中定义、在 0 的一个邻域中等于 1 且在 \(\mathrm{Sp}(u) \setminus \{0\}\) 的一个邻域中为零）成立。

通过过渡到子空间，自同态 \(u\) 在 \(e_0(u)\) 的像上诱导出一个拟幂零自同态（其谱约化为 \(\{0\}\)），并在 \(e_0(u)\) 的核上诱导出一个自同构（其谱为 \(\mathrm{Sp}(u) \setminus \{0\}\)）（loc. cit.）。

命题 14. --- 设 E 是复巴拿赫空间。为使 \(\mathcal{L}(\mathrm{E})\) 的元素 u 是 E 的里斯自同态，必须且只需下列两个互斥条件之一被满足：

\begin{enumerate}
\def\labelenumi{(\roman{enumi})}
\item
  E 的自同态 u 是 E 的自同构；
\item
  点 0 是 \(\operatorname{Sp}(u)\) 的孤立点，且投影 \(e_0(u)\) 有有限秩。
\end{enumerate}

当条件 (ii) 满足时，\(e_{0}(u)\) 的像是 u 的幂零空间，而 \(e_{0}(u)\) 的核是 u 的余幂零空间。

E 的每个自同构都是 E 的里斯自同态（III, p. 47, 注记 1）。如果 u 满足条件 (ii)，\(e_{0}(u)\) 的核 F 是 E 的闭向量子空间，有有限余维，且在 u 下稳定。设 \(u_{F}\) 是由 u 导出的 F 的自同态。它的谱包含在 \(\mathbf{C}\setminus\{0\}\) 中（I, p. 131 的第 3 目），因而 \(u_{F}\) 是 F 的自同构。由此得出 u 是 E 的里斯自同态（III, p. 48, prop. 8）。

反之，设 \(u\) 是 E 的里斯自同态。用 N 表示它的幂零空间，用 I 表示它的余幂零空间。由 III, p. 47 的 prop. 6，空间 E 是 N 与 I 的拓扑直和，且映射 \(u\) 通过过渡到子空间定义了 N 的一个幂零自同态 \(u_{\mathrm{N}}\) 与 I 的一个自同构 \(u_{\mathrm{I}}\)。特别地，我们有 \(\mathrm{Sp}(u_{\mathrm{N}}) \subset \{0\}\) 与 \(0 \not\in \mathrm{Sp}(u_{\mathrm{I}})\)。如果 N 是零空间，那么 I = E 且 \(u\) 是 E 的自同构。否则，0 是 \(\mathrm{Sp}(u)\) 的孤立点，且 \(e_0(u)\) 是以 N 为像、以 I 为核的投影（I, p. 129 的 prop. 2）；特别地，它有有限秩。

注记. --- 设 E 是实巴拿赫空间，u 是 E 的自同态。假设 0 是 u 的复谱的一个孤立点，也就是说，是自同态 \(1 \otimes u\)（它作用于 E 的完备可赋范复化空间 \(\mathrm{E}_{(\mathbf{C})}\)）的谱的一个孤立点（I, p. 85, 第 13 目）。子集 \(\{0\}\)（它属于 \(u\) 的复谱）是开闭的，且在共轭下不变；用 \(e_0(u) \in \mathcal{L}(\mathrm{E})\) 表示与之相伴的谱投影（I, p. 85 的第 13 目）。命题 14 在此情形下作必要的修改后成立。

\section{巴拿赫空间中的摄动}

\subsection{直态射}

定义 1. --- 设 E 是拓扑向量空间。E 的向量子空间称为直的，如果它有一个拓扑补。

为使 E 的向量子空间 F 是直的，必须且只需在 \(\mathcal{L}(\mathrm{E})\) 中存在一个以 F 为像的投影。如果 \(p\) 是这样一个投影，E 是 F 与 \(p\) 的核的拓扑直和。\(1_{\mathrm{E}} - p\) 的核是 F；因此如果 E 是分离的，那么 F 是闭的。

设 F 是 E 的直向量子空间，\(E_{1}\) 是 E 的包含 F 的向量子空间；那么 F 是 \(E_{1}\) 的直向量子空间。

E 中每个有有限余维的闭向量子空间都是直和项（EVT, I, p. 15, prop. 3），而 E 的每个包含在 \{0\} 的闭包中的向量子空间 F 都是直的（事实上，存在一个以 F 为像的投影，且这样一个投影必然是连续的）。

命题 1. --- 设 E 是局部凸空间。E 的每个有限维向量子空间都有补。

设 F 是 E 的有限维向量子空间。用 S 表示 \(\{0\}\) 在 F 中的闭包的一个拓扑补（该闭包是 E 的直子空间）。空间 \(F_{2} = F \cap S\) 是分离的且有限维；因此存在整数 \(n \geqslant 0\) 使得 \(F_{2}\) 同构于 \(K^{n}\)（EVT, I, p. 14, th. 2）。于是存在一个连续线性映射 \(p: S \to F_{2}\)，它延拓 \(F_{2}\) 的恒等映射。p 的核是 \(F_{2}\) 在 S 中的拓扑补，也是 F 在 E 中的拓扑补。

定义 2. --- 设 E 与 F 是局部凸空间，u 是从 E 到 F 中的连续线性映射。称 u 是直态射，如果 u 是严格态射，其核是 E 的直向量子空间且其像是 F 的直向量子空间。

设 \(u \in \mathcal{L}(\mathrm{E};\mathrm{F})\)。为使 \(u\) 的核与像分别是 E 与 F 的直向量子空间，必须且只需存在分解为拓扑直和的 \(\mathrm{E} = \mathrm{E}_1 \oplus \mathrm{E}_2\) 与 \(\mathrm{F} = \mathrm{F}_1 \oplus \mathrm{F}_2\)，使得 \(u\) 由矩阵 \(\begin{pmatrix} u_1 & 0 \\ 0 & 0 \end{pmatrix}\) 表示，其中 \(u_1 \in \mathcal{L}(\mathrm{E}_1; \mathrm{F}_1)\) 是双射。由于这时 \(u\) 的核是 \(\mathrm{E}_2\)，从 \(\mathrm{E}_1\) 到 \(\mathrm{E}/\mathrm{E}_2\) 上的典范线性映射是同构，且 \(u\) 的像是 \(\mathrm{F}_1\)，可以看出 \(u\) 是严格态射当且仅当 \(u_1\) 是从 \(\mathrm{E}_1\) 到 \(\mathrm{F}_1\) 上的同构。用 \(v\) 表示 \(\mathcal{L}(\mathrm{F};\mathrm{E})\) 中由矩阵 \(\begin{pmatrix} u_1^{-1} & 0 \\ 0 & 0 \end{pmatrix}\) 表示的元素。那么线性映射 \(u \circ v\) 是 F 中以 \(\mathrm{F}_2\) 为核、以 \(\mathrm{F}_1\) 为像的投影，线性映射 \(v \circ u\) 是 E 中以 \(\mathrm{E}_2\) 为核、以 \(\mathrm{E}_1\) 为像的投影，且有 \(u \circ v \circ u = u\)。

反之，我们有下面的结果：

命题 2. --- 设 E 与 F 是局部凸空间，设 \(u \in \mathcal{L}(\mathrm{E};\mathrm{F})\) 与 \(v \in \mathcal{L}(\mathrm{F};\mathrm{E})\)。假设 \(u = u \circ v \circ u\)。那么 \(u\) 是直态射。此外，\(v \circ u\) 是 E 中以 \(\operatorname{Ker}(u)\) 为核的连续投影，而 \(u \circ v\) 是 F 中以 \(\operatorname{Im}(u)\) 为像的连续投影。

置 \(p = v \circ u\) 与 \(q = u \circ v\)。我们有

\[
p ^ {2} = v \circ (u \circ v \circ u) = v \circ u = p, q ^ {2} = (u \circ v \circ u) \circ v = u \circ v = q,
\]

因而 p 与 q 是投影。它们是连续的。设 \(E_{1}\) 与 \(E_{2}\) 是 p 的像与核，设 \(F_{1}\) 与 \(F_{2}\) 是 q 的像与核。由于 \(u \circ v \circ u = u\)，我们有

\[
\operatorname{Ker} (u) \subset \operatorname{Ker} (v \circ u) \subset \operatorname{Ker} (u \circ v \circ u) = \operatorname{Ker} (u),
\]

因而 \(\operatorname{Ker}(u) = \operatorname{Ker}(v \circ u) = \operatorname{Ker}(p) = \operatorname{E}_{2}\)。类似地，由包含关系

\[
\operatorname{Im} (u) \supset \operatorname{Im} (u \circ v) \supset \operatorname{Im} (u \circ v \circ u) = \operatorname{Im} (u)
\]

得出 \(u\) 的像是 \(\mathrm{F}_{1}\)。空间 \(\mathrm{E}_{2}\) 与 \(\mathrm{F}_{1}\) 分别是 E 与 F 的直向量子空间。

向量空间 E 是 \(E_{1}\) 与 \(E_{2}\) 的拓扑直和，且 u 通过限制定义了一个双射线性映射 \(u_{1}\)（它从 \(E_{1}\) 到 \(F_{1}\) 上）。对每个 \(x \in E_{1}\)，我们有 \(v(u(x)) = p(x) = x\)，且映射 \(u_{1}^{-1}\)（它从 \(F_{1}\) 到 \(E_{1}\) 上）在 \(F_{1}\) 上与 v 一致。因此，\(u_{1}\) 是从 \(E_{1}\) 到 \(F_{1}\) 上的同构，而 u 是从 E 到 F 中的严格态射。

设 E 与 F 是局部凸空间，且 \(u \in \mathcal{L}(\mathrm{E};\mathrm{F})\)。为使 u 是单射（相应地，满射）的直态射，必须且只需存在一个 \(v\)（它在 \(\mathcal{L}(\mathrm{F};\mathrm{E})\) 中），使得 \(v \circ u = 1_{E}\)（相应地，\(u \circ v = 1_{F}\)）（参见 TG, III, p. 47 及 48）。

命题 3. --- 设 E 与 F 是巴拿赫空间。下列集合都是巴拿赫空间 \(\mathcal{L}(\mathrm{E};\mathrm{F})\) 的开子集：

\begin{enumerate}
\def\labelenumi{\alph{enumi})}
\item
  从 E 到 F 上的同构组成的集 \(\mathcal{I}(\mathrm{E};\mathrm{F})\)；
\item
   从 E 到 F 中的单射直态射组成的集 \(\mathcal{M}\mathcal{D}(\mathrm{E};\mathrm{F})\)，且更精确地说，对 F 的每个闭向量子空间 \(F_{1}\)，集 \(\mathcal{M}_{\mathrm{F}_{1}}(\mathrm{E};\mathrm{F})\)（它由 \(\mathcal{M}\mathcal{D}(\mathrm{E};\mathrm{F})\) 中像为 \(F_{1}\) 的拓扑补的元素组成）；
\item
   从 E 到 F 中的满射直态射组成的集 \(\mathcal{E}\mathcal{D}(\mathrm{E};\mathrm{F})\)，且更精确地说，对 E 的每个闭向量子空间 \(\mathrm{E}_1\)，集 \(\mathcal{E}_{\mathrm{E}_1}(\mathrm{E};\mathrm{F})\)（它由 \(\mathcal{E}\mathcal{D}(\mathrm{E};\mathrm{F})\) 中核为 \(\mathrm{E}_1\) 的拓扑补的元素组成）。
\end{enumerate}

此外，映射 \(u \mapsto u^{-1}\)（它从 \(\mathcal{I}(\mathrm{E};\mathrm{F})\) 到 \(\mathcal{I}(\mathrm{F};\mathrm{E})\) 上）是解析的。

依定义，\(\mathcal{I}(\mathrm{E};\mathrm{E})\) 是完备赋范代数 \(\mathcal{L}(\mathrm{E})\) 的可逆元素组成的集。由 TG, IX, p. 40, prop. 14，它是 \(\mathcal{L}(\mathrm{E})\) 的开子集。映射 \(u\mapsto u^{-1}\)（它从 \(\mathcal{I}(\mathrm{E};\mathrm{E})\) 到 \(\mathcal{I}(\mathrm{E};\mathrm{E})\) 中）是解析的（LIE, III, § 1, 第 1 目, 例子 2）。

如果集 \(\mathcal{I}(\mathrm{E};\mathrm{F})\) 是空的，那么它是开的。否则，设 \(u_0\) 是从 E 到 F 上的一个同构。于是映射 \(v\mapsto u_0\circ v\) 是从 \(\mathcal{L}(\mathrm{E})\) 到 \(\mathcal{L}(\mathrm{E};\mathrm{F})\) 上的同构，它把 \(\mathcal{I}(\mathrm{E};\mathrm{E})\) 变换为 \(\mathcal{I}(\mathrm{E};\mathrm{F})\)。因此集 \(\mathcal{I}(\mathrm{E};\mathrm{F})\) 在 \(\mathcal{L}(\mathrm{E};\mathrm{F})\) 中是开的。此外，如果 \(u = u_0\circ v\) 是 \(\mathcal{I}(\mathrm{E};\mathrm{F})\) 的一个元素，我们有 \(u^{-1} = v^{-1}\circ u_0^{-1}\)，因而映射 \(u\mapsto u^{-1}\)（它从 \(\mathcal{I}(\mathrm{E};\mathrm{F})\) 到 \(\mathcal{I}(\mathrm{F};\mathrm{E})\) 上）是解析的。

设 \(F_{1}\) 是 F 的闭向量子空间。对每个 \(u \in \mathcal{L}(E;F)\)，用 \(\overline{u}\) 表示 \(\mathcal{L}(E \times F_{1};F)\) 中由 \(\overline{u}(x,y) = u(x) + y\) 定义的元素。映射 \(u \mapsto \overline{u}\)（它从 \(\mathcal{L}(E;F)\) 到 \(\mathcal{L}(E \times F_{1};F)\) 中）是连续的，而 \(\mathcal{M}_{\mathrm{F}_{1}}(E;F)\) 是由 \(\mathcal{L}(E;F)\) 中使 \(\overline{u}\) 属于 \(\mathcal{I}(E \times F_{1};F)\) 的元素 u 组成的集。由于 \(\mathcal{I}(E \times F_{1};F)\) 在 \(\mathcal{L}(E \times F_{1};F)\) 中是开的

由前所述，集 \(\mathcal{M}_{\mathrm{F}_1}(\mathrm{E};\mathrm{F})\) 在 \(\mathcal{L}(\mathrm{E};\mathrm{F})\) 中是开的。对 \(\mathcal{MD}(\mathrm{E};\mathrm{F})\) 同样如此，它是诸 \(\mathcal{M}_{\mathrm{F}_1}(\mathrm{E};\mathrm{F})\)（当 \(\mathrm{F}_1\) 取遍 F 的闭向量子空间组成的集时）的并。

设 \(E_{1}\) 是 E 的闭向量子空间，p 是从 E 到商巴拿赫空间 \(E/E_{1}\) 上的典范映射。为使一个元素 \(u\)（它在 \(\mathcal{L}(E;F)\) 中）属于 \(\mathcal{E}_{\mathrm{E}_{1}}(\mathrm{E};\mathrm{F})\)，必须且只需映射 \((u,p)\)（它从 E 到 \(F\times E/E_{1}\) 中）属于 \(\mathcal{I}(E;F\times E/E_{1})\)。与前面一样，我们推出 \(\mathcal{E}_{\mathrm{E}_{1}}(\mathrm{E};\mathrm{F})\) 在 \(\mathcal{L}(\mathrm{E};\mathrm{F})\) 中是开的；对 \(\mathcal{E}\mathcal{D}(\mathrm{E};\mathrm{F})\) 同样如此，它是诸 \(\mathcal{E}_{\mathrm{E}_{1}}(\mathrm{E};\mathrm{F})\) 的并。

\subsection{弗雷德霍姆映射的摄动}

定理 1. --- 设 E 与 F 是巴拿赫空间。从 E 到 F 的弗雷德霍姆算子组成的集 \(\mathcal{F}(\mathrm{E};\mathrm{F})\) 在巴拿赫空间 \(\mathcal{L}(\mathrm{E};\mathrm{F})\) 中是开的，且映射 \(u\mapsto\mathrm{ind}(u)\)（它从 \(\mathcal{F}(\mathrm{E};\mathrm{F})\) 到 \(\mathbf{Z}\) 中）是局部常值的。

该定理由下面更精确的命题得出：

命题 4. --- 设 E 与 F 是巴拿赫空间，\(u_{0}\colon\mathrm{E}\to\mathrm{F}\) 是弗雷德霍姆算子，\(v_{0}\) 是 \(u_{0}\) 的一个拟逆。存在 \(u_{0}\) 在 \(\mathcal{L}(\mathrm{E};\mathrm{F})\) 中的一个开邻域 U 以及一个解析映射 \(\varphi\colon\mathrm{U}\to\mathcal{L}(\mathrm{F};\mathrm{E})\)，使得

\begin{enumerate}
\def\labelenumi{(\roman{enumi})}
\item
  我们有 \(\varphi(u_{0}) = v_{0}\)；
\item
  对每个 \(u\)（在 \(U\) 中），映射 \(\varphi(u)\) 是 \(u\) 的拟逆，特别地，\(u\) 是弗雷德霍姆映射；
\item
  对 U 中的每个 \(u\)，我们有 \(\operatorname{ind}(u) = \operatorname{ind}(u_0)\)。
\end{enumerate}

由 III, p. 42 的 prop. 2 的 (ii)，存在拓扑直和分解 \(\mathrm{E} = \mathrm{E}_1 \oplus \mathrm{E}_2\) 与 \(\mathrm{F} = \mathrm{F}_1 \oplus \mathrm{F}_2\)，且存在 \(\alpha_0 \in \mathcal{I}(\mathrm{E}_1; \mathrm{F}_1)\)，使得 \(\mathrm{E}_2\) 与 \(\mathrm{F}_2\) 是有限维的，而 \(u_0\) 相对于这些分解由矩阵 \(\left( \begin{array}{cc}\alpha_0 & 0\\ 0 & 0 \end{array} \right)\) 表示。

设 U 是由这样的元素 \(u\)（它属于 \(\mathcal{L}(\mathrm{E};\mathrm{F})\)）组成的集：在矩阵表示 \(\left( \begin{array}{cc}\alpha & \beta \\ \gamma & \delta \end{array} \right)\)（即 \(u\) 相对于这些分解的矩阵表示）中，我们有 \(\alpha \in \mathcal{I}(\mathrm{E}_1;\mathrm{F}_1)\)。由于 \(\mathcal{I}(\mathrm{E}_1;\mathrm{F}_1)\) 在 \(\mathcal{L}(\mathrm{E}_1;\mathrm{F}_1)\) 中是开的

（III, p. 57 的 prop. 3），U 是 \(u_0\) 在 \(\mathcal{L}(\mathrm{E};\mathrm{F})\) 中的开邻域。对 U 中的 \(u = \begin{pmatrix} \alpha & \beta \\ \gamma & \delta \end{pmatrix}\)，置

\[
\varphi (u) = v _ {0} + \left( \begin{array}{c c} \alpha^ {- 1} - \alpha_ {0} ^ {- 1} & 0 \\ 0 & 0 \end{array} \right).\tag{1}
\]

我们有 \(\varphi(u_{0}) = v_{0}\)，且映射 \(\varphi\) 是解析的（loc. cit.）。模有限秩连续线性映射，我们有同余关系

\[
u \equiv \left( \begin{array}{c c} \alpha & 0 \\ 0 & 0 \end{array} \right) , v _ {0} \equiv \left( \begin{array}{c c} \alpha_ {0} ^ {- 1} & 0 \\ 0 & 0 \end{array} \right) , \varphi (u) \equiv \left( \begin{array}{c c} \alpha^ {- 1} & 0 \\ 0 & 0 \end{array} \right).
\]

因此，\(\varphi(u)\) 是 \(u\) 的拟逆。U 的每个元素 \(u\) 通过限制定义了一个从 \(E _1\) 到 \(F _2\) 在 F 中的一个拓扑补上的同构。由 III, p. 44 的 prop. 3，我们因此有

\[
\begin{array}{c} \operatorname{ind} (u) = \operatorname{codim} _ {\mathrm{F}} (u (\mathrm{E} _ {1})) - \operatorname{codim} _ {\mathrm{E}} (\mathrm{E} _ {1}) \\ = \dim (\mathrm{F} _ {2}) - \dim (\mathrm{E} _ {2}) = \operatorname{ind} (u _ {0}), \end{array}
\]

证毕。
\subsection{里斯自同态的扰动}

引理 1. --- 设 E 是巴拿赫空间。设 p 和 q 是 E 上的连续投影，满足 \(\|q - p\| < 1\) 。则 p 诱导从 \(\operatorname{Im}(q)\) 到 \(\operatorname{Im}(p)\) 上的一个同构，而 \(1_{E} - p\) 诱导从 \(\operatorname{Ker}(q)\) 到 \(\operatorname{Ker}(p)\) 上的一个同构。

考察连续线性映射 \(u\colon x\mapsto p(x)\) （从 \(\operatorname{Im}(q)\) 到 \(\operatorname{Im}(p)\) 中）与 \(v\colon y\mapsto q(y)\) （从 \(\operatorname{Im}(p)\) 到 \(\operatorname{Im}(q)\) 中）。对每个 \(x\in\operatorname{Im}(q)\) ，有 \(x=q(x)\) ，由此

\[
\begin{array}{c} (q - p) ^ {2} (x) = q ^ {2} (x) - q (p (x)) - p (q (x)) + p ^ {2} (x) \\ = x - q (p (x)) = x - v (u (x)). \end{array}
\]

由此可得 \(\|1_{E}-v\circ u\|\leqslant\|q-p\|^{2}<1\) 。由 I, p. 22 的推论 1，\(v\circ u\) 是 \(\operatorname{Im}(q)\) 的一个自同构。类似地可证 \(u\circ v\) 是 \(\operatorname{Im}(p)\) 的一个自同构。这蕴含 u 是从 \(\operatorname{Im}(q)\) 到 \(\operatorname{Im}(p)\) 上的一个同构，从而得出引理的第一个断言。第二个断言由此通过把 p 和 q 分别换成 \(1_{E}-p\) 和 \(1_{E}-q\) 而得到。

定理 2. --- 设 E 是巴拿赫空间。E 的里斯自同态所成的集 \(\mathcal{R}(\mathrm{E})\) 在 \(\mathcal{L}(\mathrm{E})\) 中是开的。把 \(\mathcal{R}(\mathrm{E})\) 的元素映为其幂零空间的维数的映射在 \(\mathcal{L}(\mathrm{E})\) 上是上半连续的。

定理由下面更精确的命题得出：

命题 5. --- 设 E 是巴拿赫空间，\(u_{0}\) 是 E 的一个里斯自同态。设 N（相应地，I）是 \(u_{0}\) 的幂零空间（相应地，余幂零空间）。用 d 表示 N 的维数。存在 \(u_{0}\) 在 \(\mathcal{L}(\mathrm{E})\) 中的一个开邻域 U 和一个解析映射 \(\pi\colon\mathrm{U}\to\mathcal{L}(\mathrm{E})\) ，使得

\begin{enumerate}
\def\labelenumi{(\roman{enumi})}
\item
  自同态 \(\pi(u_0)\) 是以 N 为像、以 I 为核的投影；
\item
  对每个 \(u \in U\) ，线性映射 \(\pi(u)\) 是秩为 d 的投影，它属于 u 的双交换子，特别地与 u 交换。此外，u 诱导 \(\operatorname{Ker}(\pi(u))\) 的一个自同构；
\item
  U 的每个元素都是里斯自同态，其幂零空间的维数 \(\leqslant d\) 。
\end{enumerate}

若 K = C，用 \(\mathrm{Sp}(u_{0})\) 表示 \(u_{0}\) 关于代数 \(\mathcal{L}(\mathrm{E})\) 的谱。当 \(K = R\) 时，用 \(\mathrm{Sp}(u_{0})\) 表示复谱 \(\mathrm{Sp}_{\mathcal{L}(\mathrm{E}_{(\mathbf{C})})}((u_{0})_{(\mathbf{C})})\) ，即 \(u_{0}\) 的复谱（I, p. 85, 第 13 目）。由 III, p. 54 的命题 14，点 0 在 \(\{0\} \cup \mathrm{Sp}(u_{0})\) 中是孤立的。设 r \textgreater{} 0 是实数，使得 \(\mathrm{Sp}(u_{0}) - \{0\}\) 的每个元素的模都 \textgreater{} r。用 V 表示 C 中由绝对值 ≠ r 的复数所组成的开集；用 f 表示 V 上如下定义的全纯函数：\(f(z) = 0\) （若 \(|z| > r\) ），\(f(z) = 1\) （若 \(|z| < r\) ）。若 K = C（相应地，K = R），用 U' 表示 \(\mathcal{L}(\mathrm{E})\) 中谱（相应地，复谱）包含于 V 的元素 u 所成的集。集 U' 是 \(u_{0}\) 在 \(\mathcal{L}(\mathrm{E})\) 中的一个开邻域，且映射 \(u \mapsto f(u)\) （从 U' 到 \(\mathcal{L}(\mathrm{E})\) 中）是全纯的（I, p. 76, 命题 10 及 I, p. 85, 第 13 目）。

设 \(u \in \mathrm{U}'\) 。自同态 \(f(u)\) 是谱投影 \(e_0(u)\) ；它与 \(u\) 交换，且 \(u\) 通过过渡到子空间诱导 \(\operatorname{Ker}(e_0(u))\) 的一个自同构（参见 III, p. 53, 第 6 目）。

投影 \(e_{0}(u_{0})\) 以 N 为像、以 I 为核（III, p. 54, 命题 14）；它的秩是 d。由引理 1，由 \(u \in U'\) 中使 \(e_{0}(u)\) 的秩为 d 的元素所成的集 U 是 \(u_{0}\) 的一个开邻域。集 U 与映射 \(\pi: u \mapsto e_{0}(u)\) （从 U 到 \(\mathcal{L}(\mathrm{E})\) 中）满足命题的条件 (i) 与 (ii)。

设 \(u \in U\) 。由于 u 诱导 \(\operatorname{Ker}(e_{0}(u))\) （它是 E 的一个闭的有限余维数向量子空间）的一个自同构，自同态 u 是 E 的里斯自同态（III, p. 48, 命题 8）；u 的幂零空间包含于 \(\pi(u) = e_{0}(u)\) 的像中，且其维数 \(\leqslant d\) 。

\subsection{连续线性映射的余范数}

设 E 和 F 是赋范空间，\(u: E \rightarrow F\) 是连续线性映射，N 是 u 的核，I 是它的像。用 p 表示从 E 到 E/N 上的典范满射，用 i 表示从 I 到 F 中的典范单射。用 \(\widetilde{u}\) 表示从 E/N 到 I 上的双射线性映射，它满足 \(u = i \circ \widetilde{u} \circ p\) 。向量空间 E/N 赋有商范数，即

\[
\| y \| = \inf _ {x \in p^{-1}(y)} \| x \|\tag{2}
\]

对一切 \(y \in E/N\) 成立。映射 \(\widetilde{u}\) 连续且 \(\|u\| = \|\widetilde{u}\|\) ，由此

\[
\| u\| = \sup_{\substack{y\in \mathrm{E / N}\\ y\neq 0}}\frac{\|\widetilde{u} (y)\|}{\|y\|},\tag{3}
\]

其中上确界在 \(\overline{R}_{+}\) 中取。u 的余范数是指

\[
((u)) = \inf_{\substack{y\in \mathrm{E} / \mathrm{N}\\ y\neq 0}}\frac{\| \widetilde{u}(y)\|}{\|y\|},\tag{4}
\]

其中下确界在 \(\overline{R}_{+}\) 中取。因此我们有

\[
((u)) \| y \| \leqslant \| \widetilde {u} (y) \| \leqslant \| u \| \| y \|\tag{5}
\]

对 E/N 的每个元素 \(y\) 成立。令 \(v = i \circ \widetilde{u}\) 。我们有 \(u = v \circ p\) ，并且

\[
((u)) = \inf_{\substack{y\in \mathrm{E / N}\\ y\neq 0}}\frac{\|v(y)\|}{\|y\|} = \inf_{\substack{y\in \mathrm{E / N}\\ \|y\| = 1}}\| v(y)\| .\tag{6}
\]

当 \(u\) 是零映射时，空间 E/N 缩减为 0，且有 \(((u)) = +\infty\) 与 \(\| u \| = 0\) 。当 \(u \neq 0\) 时，由 (3) 与 (4) 推出不等式

\[
0 \leqslant ((u)) \leqslant \| u \| <   + \infty .\tag{7}
\]

当 u 单射时，我们有

\[
((u)) = \inf_{\substack{x\in \mathrm{E}\\ x\neq 0}}\frac{\|u(x)\|}{\|x\|},\tag{8}
\]

且对一切 \(x \in \mathrm{E}\) ，我们有

\[
((u)) \| x \| \leqslant \| u (x) \| \leqslant \| u \| \| x \|.\tag{9}
\]

用 \(j\) 表示赋范空间 F 到其完备化 \(\widehat{\mathbf{F}}\) 中的典范单射。我们有 \(((u)) = ((j \circ u))\) 。

注记. --- 按定义，\(((u)) > 0\) 等价于双射线性映射 \(\widetilde{u}\) 是双连续的，亦即等价于 u 是严格态射（TG, III, p. 16）。此时我们有

\[
((u)) = \| \widetilde {u} ^ {- 1} \| ^ {- 1}.\tag{10}
\]

引理 2. --- 设 c 是实数。若 c \textless{} ((u))，则对每个元素 \(z \in \operatorname{Im}(u)\) ，存在 E 的元素 x，使得 \(u(x) = z\) 且 \(c \|x\| \leqslant \|z\|\) 。反之，若对每个 \(z \in \operatorname{Im}(u)\) 存在 \(x \in E\) ，使得 \(u(x) = z\) 且 \(c \|x\| \leqslant \|z\|\) ，则 \(c \leqslant ((u))\) 。

这是公式 (2)、公式 (5) 以及 \(u\) 的余范数定义的推论。

命题 6. --- 设 E 和 F 是赋范空间，\(u \in \mathcal{L}(\mathrm{E};\mathrm{F})\) 。用 B 表示 E 中范数 \(< 1\) 的元素所成的集。令

\[
\mathrm{P} = u (\mathrm{E}) - u (\mathrm{B}), \quad \mathrm{Q} = u (\mathrm{E}) - (\overline {{u (\mathrm{B})}} \cap u (\mathrm{E})).
\]

u 的余范数等于在 F 中从 0 到 P 的距离。若赋范空间 E 完备或 u 是严格态射，则 u 的余范数等于在 F 中从 0 到 Q 的距离。

设 N 是 u 的核；设 \(p: E \to E/N\) 是典范满射，并设 \(v: E/N \to F\) 是 u 在商空间上诱导的映射。集 \(p(B)\) 是 E/N 中范数 \textless{} 1 的元素所成的集。我们有

\[
\mathrm{P} = u (\mathrm{E}) - u (\mathrm{B}) = v (\mathrm{E} / \mathrm{N}) - v (p (\mathrm{B})) = v ((\mathrm{E} / \mathrm{N}) - p (\mathrm{B}))
\]

因为映射 \(v\) 是单射。换言之，P 由 F 中形如 \(v(y)\) （其中 \(y \in \mathrm{E/N}\) 且 \(\|y\| \geqslant 1\) ）的元素组成。用 \(d_{\mathrm{P}}\) 表示在 F 中从 0 到 P 的距离。我们有

\[
d_{\mathrm{P}} = \inf_{\substack{y\in \mathrm{E} / \mathrm{N}\\ \| y\| \geqslant 1}}\| v(y)\| = \inf_{\substack{y\in \mathrm{E} / \mathrm{N}\\ \| y\| = 1}}\| v(y)\| = ((u))
\]

由 (6) 即得。

设 \(u\) 是严格态射。设 \(\varepsilon > 0\) 。集 \(\varepsilon u(\mathrm{B})\) 是 \(u(\mathrm{E})\) 中 0 的一个邻域。\(u(\mathrm{B})\) 在 \(u(\mathrm{E})\) 中的闭包等于 \(\overline{u(\mathrm{B})} \cap u(\mathrm{E})\) 。它包含于集 \(u(\mathrm{B}) + \varepsilon u(\mathrm{B})\) ，后者等于 \((1 + \varepsilon)u(\mathrm{B})\) ，因为 \(u(\mathrm{B})\) 是凸的。于是有 \((1 + \varepsilon)\mathrm{P} \subset \mathrm{Q} \subset \mathrm{P}\) ，且在 F 中从 0 到 Q 的距离 \(d_{\mathrm{Q}}\) 满足不等式 \(d_{\mathrm{P}} \leqslant d_{\mathrm{Q}} \leqslant (1 + \varepsilon)d_{\mathrm{P}}\) 。由于这对每个 \(\varepsilon > 0\) 都成立，故有 \(d_{\mathrm{Q}} = d_{\mathrm{P}} = ((u))\) 。

设 \(u\) 不是严格态射，但赋范空间 E 完备。此时有 \(((u)) = 0\) 。\(u(\mathrm{B})\) 在 \(u(\mathrm{E})\) 中的闭包（它等于 \(\overline{u(\mathrm{B})} \cap u(\mathrm{E})\) ）不是 \(u(\mathrm{E})\) 中 0 的邻域（EVT, I, p. 17, 定理 1）。因此存在 Q 中任意接近 0 的点，从而 \(d_{\mathrm{Q}} = 0 = ((u))\) 。

命题 7. --- 设 E 是巴拿赫空间，F 是赋范空间。映射 \(u \mapsto ((u))\) （从 \(\mathcal{L}(\mathrm{E};\mathrm{F})\) 到 \(\overline{\mathbf{R}}\) 中）是上半连续的。

设 \(u \in \mathcal{L}(\mathrm{E};\mathrm{F})\) 。只需证明：对每个实数 \(c > ((u))\) ，由元素 \(v \in \mathcal{L}(\mathrm{E};\mathrm{F})\) （满足 \(((v)) < c\) ）所成的集是 \(u\) 的一个邻域。用 B 表示 E 中范数 \(< 1\) 的元素所成的集。由命题 6，存在 \(y \in \mathrm{E}\) ，使得 \(u(y) \notin \overline{u(\mathrm{B})}\) 且 \(\| u(y)\| < c\) 。距离 \(d\) （从 \(u(y)\) 到闭集 \(\overline{u(\mathrm{B})}\) ）是严格正的。由元素 \(v\) （\(\mathcal{L}(\mathrm{E};\mathrm{F})\) 中的、满足关系式 \(\| v(y)\| < c\) 与 \(\| u - v\|(1 + \| y\|) < d\) 者）所成的集 V 是 \(u\) 在 \(\mathcal{L}(\mathrm{E};\mathrm{F})\) 中的一个邻域。设 \(v \in \mathrm{V}\) 。对每个 \(x \in \mathrm{B}\) ，我们有

\[
\begin{array}{c} d \leqslant \| u (y) - u (x) \| \leqslant \| v (y) - v (x) \| + \| u - v \| (\| y \| + \| x \|) \\ <   \| v (y) - v (x) \| + d. \end{array}
\]

于是 \(v(y)\) 不属于 \(v(\mathrm{B})\) ，且由命题 6 有 \(((v)) \leqslant \| v(y) \| < c\) 。

命题 8. --- 设 E 是巴拿赫空间，F 是赋范空间。对每个 \(u \in \mathcal{L}(\mathrm{E};\mathrm{F})\) ，我们有 \(((u)) = ((^t u))\) 。

用 \(j\) 表示赋范空间 F 到其完备化 \(\widehat{\mathrm{F}}\) 中的典范单射。我们有 \(((u)) = ((j \circ u))\) 。由于线性映射 \(^t j: (\widehat{\mathrm{F}})' \to \mathrm{F}'\) 是双射且等距，也有 \(((^t u)) = ((^t (j \circ u)))\) 。因此只需就赋范空间 F 完备的情形证明本命题，我们在证明的其余部分即作此假设。

若 \(u\) 为零，则有 \(((u)) = ((^t u)) = +\infty\) 。若 \(u\) 不是严格态射，则 \(^t u\) 也不是（EVT, IV, p. 29, 推论 3），且有 \(((u)) = ((^t u)) = 0\) 。

现在设 u 是非零的严格态射。用 N 表示 u 的核，I 表示它的像，并考察 u 的典范分解：

\[
\mathrm{E} \xrightarrow {p} \mathrm{E/N} \xrightarrow {v} \mathrm{I} \xrightarrow {i} \mathrm{F}.
\]

\(^{t}u\) 的核是 I 在 F' 中的正交 I°（EVT, IV, p. 27, 命题 2），且 \(^{t}i\) 通过过渡到商空间诱导一个从 F'/I° 到 I' 上的等距 \(\iota\) （EVT, IV, p. 9, 命题 10）。此外，\(^{t}p\) 给出从 (E/N)' 到 N 在 E' 中的正交 N° 上的一个等距（EVT, IV, p. 8, 命题 9）。因此 \(^{t}u\) 的典范分解是

\[
\mathrm{F} ^ {\prime} \longrightarrow \mathrm{F} ^ {\prime} / \mathrm{I} ^ {\circ} \stackrel {v ^ {\prime}} {\longrightarrow} \mathrm{N} ^ {\circ} \longrightarrow \mathrm{E} ^ {\prime},
\]

其中 \(v^{\prime} = \pi \circ{}^{t}v\circ \iota\)。此时有

\[
\| (v ^ {\prime}) ^ {- 1} \| = \| (^ {t} v) ^ {- 1} \| = \| ^ {t} (v ^ {- 1}) \| = \| v ^ {- 1} \|
\]

（EVT, IV, p. 7, 命题 8），从而由公式 (10) 得 \(((u)) = ((^t u))\) 。

\subsection{赋范空间的有限维向量子空间}

下面的命题将在 TA（待出版）中作为博苏克–乌拉姆定理的推论而证明。

定理 3. — 设 n 是正整数，V 是维数为 \(n + 1\) 的实赋范向量空间，W 是 V 的维数为 n 的向量子空间。设 S 是 V 的单位球面。不存在连续映射 \(f: S \to W\) ，使得 \(\|f(x) - x\| < 1\) 对一切 \(x \in S\) 成立。

定理 4（克赖因、克拉索塞利斯基、米尔曼）. — 设 E 是赋范空间，F 和 G 是 E 的向量子空间。假设 G 的维数有限且严格小于 F 的维数。则存在 F 中范数为 1 的元素，它到 G 的距离等于 1。

只需处理域 K 等于 R 的情形。设 n 是 G 的维数。把 F 换成 F 的一个包含 G 且维数为 \(n+1\) 的向量子空间，就化归于 \(\dim(F)=n+1\) 的情形。用反证法，假设定理的结论不成立。设 S 是 F 的单位球面。对每个 \(x\in S\) ，x 到 G 的距离严格小于 \(\|x\|=1\) ，故可选取 G 的元素 \(y(x)\) ，使得 \(\|x-y(x)\|<1\) 。用 \(V_{x}\) 表示 S 中满足 \(\|z-y(x)\|<1\) 的元素 z 所成的集；它是 x 在 S 中的一个开邻域。存在 S 上的局部有限连续单位分解 \((\varphi_{x})_{x\in\mathrm{S}}\) ，它从属于 S 的覆盖 \((\mathrm{V}_{x})_{x\in\mathrm{S}}\) （TG, IX, p. 46, 命题 3 及 p. 51, 命题 6）。用 \(f\colon S\to G\) 表示由下式给出的连续映射

\[
f (z) = \sum_ {x \in \mathrm{S}} \varphi_ {x} (z) y (x).
\]

对一切 \(z \in S\)，我们有 \(\varphi_x(z) \geqslant 0\) （对一切 \(x \in S\)），且存在 \(x \in S\) 使 \(\varphi_x(z) > 0\)，因为 \(\sum_{x \in S} \varphi_x(z) = 1\)，由此

\[
\begin{array}{c} \| z - f (z) \| = \left\| \sum_ {x \in \mathrm{S}} \varphi_ {x} (z) (z - y (x)) \right\| \leqslant \sum_ {x \in \mathrm{S}} \varphi_ {x} (z) \| z - y (x) \| \\ <   \sum_ {x \in \mathrm{S}} \varphi_ {x} (z) = 1. \end{array}
\]

f 的这一性质与定理 3 矛盾。

命题 9. — 设 E 和 F 是赋范空间，\(n \in N\)，并设 u 是从 E 到 F 中的连续线性映射，其核的维数为 n。u 的余范数等于从 u 到映射 \(v \in \mathcal{L}(\mathrm{E};\mathrm{F})\) （其核的维数至少为 \(n + 1\)）所成的集的距离 d。

当 \(u = 0\) 时，我们有 \(((u)) = +\infty\) 且 \(\dim(\mathrm{E}) = n\) ，从而 \(d = +\infty\) 。以下设 \(u\) 非零，于是 \(((u)) < +\infty\) 。设 \(v\) 是 \(\mathcal{L}(\mathrm{E};\mathrm{F})\) 中满足 \(\| u - v \| < ((u))\) 的元素，我们来证明它的核的维数 \(\leqslant n\) 。设 \(x\) 是 \(\operatorname{Ker}(v)\) 的一个范数为 1 的元素，\(y\) 是它在 \(\operatorname{E}/\operatorname{Ker}(u)\) 中的像。我们有（III, p. 61 的公式 (5)）

\[
((u)) \| y \| \leqslant \| u (x) \| = \| (u - v) (x) \| \leqslant \| u - v \| <   ((u))
\]

从而 \(\|y\| < 1\) 。而 \(\|y\|\) 正是 x 到 \(\operatorname{Ker}(u)\) 的距离。于是定理 4 蕴含 \(\dim \operatorname{Ker}(v) \leqslant \dim \operatorname{Ker}(u) = n\) 。这就证明了不等式 \(((u)) \leqslant d\) 。反向不等式 \(d \leqslant ((u))\) 由下面更精确的引理得出。

引理 3. — 设 c 是满足 \(c > ((u))\) 的实数。存在连续线性映射 \(h: E \to F\) ，其秩为 1、范数 \textless{} c，使得 \(u + h\) 的核包含 u 的核且维数为 \(n + 1\) 。

用 \(p\) 表示从 E 到 \(\mathrm{E} / \mathrm{Ker}(u)\) 上的典范映射。存在 \(a \in \mathrm{E}\) ，使得 \(\| p(a) \| = 1\) 且 \(\| u(a) \| < c\) （III, p. 61 的公式 (6)）。由哈恩–巴拿赫定理（EVT, II, p. 24, 推论 2），存在连续线性泛函 \(f\) （在赋范空间 \(\mathrm{E} / \mathrm{Ker}(u)\) 上），使得 \(f(p(a)) = 1\) 且 \(\| f \| = 1\) 。线性映射 \(h: x \mapsto -(f \circ p)(x)u(a)\) （从 E 到 F 中）连续且秩为 1，并且有 \(\| h \| \leqslant \| f \| \| p \| \| u(a) \| < c\) 。映射 \(u + h\) 的核包含 \(u\) 的核和 \(a\) ；由于 \(a \notin \mathrm{Ker}(u)\) ，它的维数至少为 \(n + 1\) 。另一方面，线性映射 \(u\) 通过过渡到子空间诱导一个从 \(\mathrm{Ker}(u + h)\) 到 \(\mathrm{Im}(h)\) 中的线性映射，其核为 \(\mathrm{Ker}(u) \cap \mathrm{Ker}(u + h)\) ，于是 \(\dim(\mathrm{Ker}(u + h)) \leqslant \dim(\mathrm{Ker}(u)) + 1\) 。结论得证。

推论 1. --- 设 E 和 F 是赋范空间，u、v 是从 E 到 F 中的非零连续线性映射，其核有相同的有限维数。则我们有

\[
\left| \left((u)\right) - ((v)) \right| \leqslant \| u - v \|.
\]

用 \(n\) 表示 \(u\) 与 \(v\) 的核的共同维数。用 A 表示从 E 到 F 的、核的维数 \(\geqslant n + 1\) 的连续线性映射所成的集。由于 \(u \neq 0\) ，我们有 \(\dim(\mathrm{E}) > n\) ，且集 A 包含 0。由命题 9，\(((u))\) 与 \(((v))\) 分别是 \(u\) 与 \(v\) 到 \(\mathcal{L}(\mathrm{E};\mathrm{F})\) 中集 A 的距离。于是只需应用公式 \(|d(u,\mathrm{A}) - d(v,\mathrm{A})| \leqslant \|u - v\|\) （参见 TG, IX, p. 13）。

推论 2. --- 设 E 和 F 是巴拿赫空间，u 和 v 是从 E 到 F 中的非零连续线性映射，其像在 F 中有相同的有限余维数。则有

\[
| ((u)) - ((v)) | \leqslant \| u - v \|.
\]

态射 u 是严格的（III, p. 52, 引理 6）。连续线性映射 \({}^{t}u\) 的核是正交 \((\mathrm{Im}(u))^{\circ}\) （\(\mathrm{Im}(u)\) 的正交）（EVT, IV, p. 27, 命题 2），由此 \(\dim(\mathrm{Ker}({}^{t}u)) = \mathrm{codim}_{\mathrm{F}}(\mathrm{Im}(u))\) ，类似地 \(\dim(\mathrm{Ker}({}^{t}v)) = \mathrm{codim}_{\mathrm{F}}(\mathrm{Im}(v))\) 。因此 \({}^{t}u\) 与 \({}^{t}v\) 的核有相同的有限维数。由于

\[
\| ^ {t} u - ^ {t} v \| = \| u - v \|, \quad ((^ {t} u)) = ((u)), \quad ((^ {t} v)) = ((v))
\]

（EVT, IV, p. 7, 命题 8 及 III, p. 63, 命题 8），把推论 1 应用于 \(^{t}u\) 和 \(^{t}v\) 即得该断言。

\subsection{连续单射或满射线性映射的扰动}

在本节中，我们采用下列约定：若 E 是有限维向量空间，\(\dim(\mathrm{E})\) 表示它的维数；若 E 是无穷维向量空间，则令 \(\dim(\mathrm{E}) = +\infty \in \overline{\mathbf{R}}\) 。若 u 是核或余核为有限维的线性映射，则令 \(\operatorname{ind}(u) = \dim \operatorname{Coker}(u) - \dim \operatorname{Ker}(u)\) ，其中的运算在 \(\overline{R}\) 中进行。

设 E 和 F 是赋范空间。用 \(\mathcal{M}(\mathrm{E};\mathrm{F})\) 表示从 E 到 F 中的严格单形态射所成的集，用 \(\mathcal{Q}\mathcal{M}(\mathrm{E};\mathrm{F})\) 表示从 E 到 F 中的、核为有限维的严格态射所成的集。

命题 10. --- 设 E 和 F 是赋范空间。集 \(\mathcal{M}(\mathrm{E};\mathrm{F})\) 在 \(\mathcal{L}(\mathrm{E};\mathrm{F})\) 中是开的。它是 \(\mathcal{L}(\mathrm{E};\mathrm{F})\) 中单射映射所成的集的内部。

用 A 表示从 E 到 F 中的连续单射线性映射所成的集。为使映射 \(u \in A\) 是严格态射，必须且只需它的余范数 \(((u))\) 严格为正（III, p. 62, 注记）。而 \(((u))\) 是从 \(u\) 到 A 在 \(\mathcal{L}(\mathrm{E};\mathrm{F})\) 中的补集的距离（III, p. 65, 命题 9）。命题得证。

命题 11. --- 设 E 和 F 是巴拿赫空间。集 \(\mathcal{Q}\mathcal{M}(\mathrm{E};\mathrm{F})\) 在 \(\mathcal{L}(\mathrm{E};\mathrm{F})\) 中是开的。它是 \(\mathcal{L}(\mathrm{E};\mathrm{F})\) 中核为有限维的映射所成的集的内部。

用 A 表示从 E 到 F 中的、核为有限维的连续线性映射所成的集。

设 u 是 \(\mathcal{Q}\mathcal{M}(\mathrm{E};\mathrm{F})\) 的一个元素。此时有 \(((u)) > 0\) （III, p. 62, 注记）。每个元素 \(v \in \mathcal{L}(\mathrm{E};\mathrm{F})\) （与 u 的距离 \(<((u))\) ）都属于 A（III, p. 65, 命题 9），故 u 是 A 的一个内点。

反之，设 \(u\) 是 A 的一个不是严格态射的元素。我们有 \(((u)) = 0\) （III, p. 62, 注记）。设 \(v\) 是 \(\mathcal{L}(\mathrm{E};\mathrm{F})\) 的一个与 \(u\) 相差一个有限秩线性映射的元素；由 III, p. 40 的推论 1，态射 \(v\) 不可能是具有闭像的严格态射；由于从 E 到 F 中的严格态射在 F 中有闭像（EVT, IV, p. 28, 定理 1），这意味着 v 不是严格态射。于是有 \(((v)) = 0\) 。设 \(\varepsilon\) 为一实数，\textgreater{} 0。借助前面的考察，III, p. 66 的引理 3 使我们可以归纳地构造元素的一个序列 \((u_{m})_{m \in \mathbf{N}}\) （\(\mathcal{L}(\mathrm{E}; \mathrm{F})\) 中的），它满足下列条件：

\begin{enumerate}
\def\labelenumi{(\roman{enumi})}
\item
  有 \(u_{0}=u\) ；
\item
  对每个 \(m \geqslant 0\) ，\(u_{m+1} - u_m\) 是秩为 1 且范数 \(\leqslant 2^{-m-1}\varepsilon\) 的连续线性映射；
\item
  对每个 \(m \geqslant 0\) ，\(u_m\) 的核维数为 \(n + m\) ，且包含于 \(u_{m+1}\) 的核中。
\end{enumerate}

序列 \((u_{m})\) 是巴拿赫空间 \(\mathcal{L}(\mathrm{E};\mathrm{F})\) 中的柯西序列。设 \(u'\) 是它的极限。\(u'\) 的核包含 \(u_{m}\) 的核（对一切 \(m \geqslant 0\) ）；它是无穷维的。由于

\[
\left\| u ^ {\prime} - u \right\| \leqslant \sum_ {m = 0} ^ {\infty} \left\| u _ {m + 1} - u _ {m} \right\| \leqslant \varepsilon \sum_ {m = 0} ^ {\infty} 2 ^ {- m - 1} = \varepsilon ,
\]

u 到 A 的补集的距离小于 \(\varepsilon\) 。由于这对每个 \(\varepsilon > 0\) 都成立，由此得出 u 不是 A 的内点。

命题 12. --- 设 E 和 F 是巴拿赫空间，u 是 \(\mathcal{Q}\mathcal{M}(\mathrm{E};\mathrm{F})\) 的一个元素。每个 \(v \in \mathcal{L}(\mathrm{E};\mathrm{F})\) （满足 \(\|v - u\| < ((u))\) ）都属于 \(\mathcal{Q}\mathcal{M}(\mathrm{E};\mathrm{F})\) 且满足关系式

\[
\begin{array}{c} \dim \operatorname{Ker} (v) \leqslant \dim \operatorname{Ker} (u), \\ \dim \operatorname{Coker} (v) \leqslant \dim \operatorname{Coker} (u), \\ \operatorname{ind} (v) = \operatorname{ind} (u). \end{array}
\]

由于 \(u\) 是严格的，我们有 \(((u)) > 0\) 。用 B 表示以 \(v\) 为中心、\(((u))\) 为半径的开球（在 \(\mathcal{L}(\mathrm{E};\mathrm{F})\) 中）。对每个 \(v \in \mathrm{B}\) ，我们有 \(\dim \operatorname{Ker}(v) \leqslant \dim \operatorname{Ker}(u)\) （III, p. 65, 命题 9）且 \(v \in \mathcal{Q}\mathcal{M}(\mathrm{E};\mathrm{F})\) （命题 11）。

对 \(r \in \mathbf{Z} \cup \{+\infty\}\) ，用 \(B_r\) 表示由元素 \(v \in B\) （满足 \(\text{ind}(v) = r\) ）所成的集。若 \(r \in \mathbf{Z}\) ，则集 \(B_r\) 是从 E 到 F 的、指标为 \(r\) 且属于 B 的弗雷德霍姆映射所成的集（III, p. 52, 命题 11）；由 III, p. 58 的定理 1，它在 B 中是开的。

我们来证明诸集 \(B_{r}\) 在 B 中是闭的。设 \(v \in B\) 是附着于 \(B_{r}\) 的一点。由于诸集 \(B_{s}\) （\(s \in Z\) ）在 B 中是开的且两两不相交，我们有 \(v \in B_{r}\) 或 \(v \in B_{+\infty}\) 。

设 \(v \in B_{+\infty}\) 。用 n 表示 \(\operatorname{Ker}(u)\) 的维数。选取 F 的一个维数为 \(r + 2n + 1\) 的向量子空间 T，它与 \(\operatorname{Im}(v)\) 的交缩减为 0；再选取 \(\operatorname{Ker}(v)\) 在 E 中的一个拓扑补 S（III, p. 55, 命题 1）。考察连续线性映射 \(f: (s, t) \mapsto v(s) + t\) （从 S × T 到 F）。它是单射的。线性映射 v 是严格态射，因为 v 属于 B ⊂ \(\mathcal{Q}\mathcal{M}(\mathrm{E};\mathrm{F})\) 。于是 v 的像是闭的（EVT, IV, p. 28, 定理 1）。由 III, p. 39 的命题 1，v 在 S 上的限制是从 S 到 F 的具有闭像的严格态射，而 f 是从 S × T 到 F 的具有闭像的严格态射。

由命题 10，存在 \(v\) 在 B 中的一个邻域 U，使得对每个 \(w \in \mathrm{U}\) ，线性映射 \((s,t) \mapsto w(s) + t\) （从 \(\mathrm{S} \times \mathrm{T}\) 到 F 中）是单射的。设 \(w\) 是 U 的一个元素。此时有 \(\operatorname{Ker}(w) \cap \mathrm{S} = \{0\}\) ，从而 \(w\) 通过限制并过渡到商空间诱导一个从 \(\operatorname{Ker}(w)\) 到 E/S 中的单射线性映射，这蕴含 \(\operatorname{Ker}(w)\) 的维数至多为 \(n\) 。我们还有 \(w(\mathrm{S}) \cap \mathrm{T} = \{0\}\) ，由此

\[
\operatorname{codim} _ {\mathrm{F}} (\operatorname{Im} (w)) \geqslant \dim (\mathrm{T}) - \operatorname{codim} _ {\mathrm{E}} (\mathrm{S}) = r + n + 1,
\]

这蕴含 \(\operatorname{ind}(w) \geqslant r + 1\) ，且与 v 附着于 B \(_{r}\) 这一假设相矛盾。

若 \(r\) 是 \(\mathbf{Z}\) 中使 \(B_r\) 非空的元素，则我们有 \(B_r = B\) ，因为 \(B_r\) 在 \(B\) 中既开又闭且 \(B\) 连通。若 \(B_r\) 对每个 \(r \in \mathbf{Z}\) 都是空的，则有 \(\text{ind}(v) = +\infty\) （对一切 \(v \in B\) ）。于是映射 \(v \mapsto \text{ind}(v)\) 在 \(B\) 上是常数。最后，对每个 \(v \in B\) ，我们有

\[
\begin{array}{c} \dim \operatorname{Coker} (v) = \operatorname{ind} (v) + \dim \operatorname{Ker} (v) \\ \leqslant \operatorname{ind} (u) + \dim \operatorname{Ker} (u) = \dim \operatorname{Coker} (u). \end{array}
\]

证明到此结束。

推论 1. --- 由 \(u \mapsto \dim \operatorname{Ker}(u)\) 与 \(u \mapsto \dim \operatorname{Coker}(u)\) 在 \(\mathcal{Q}\mathcal{M}(\mathrm{E};\mathrm{F})\) 上定义的函数是上半连续的。函数 \(u \mapsto \operatorname{ind}(u)\) 在 \(\mathcal{Q}\mathcal{M}(\mathrm{E};\mathrm{F})\) 上是局部常数的。

推论 2. --- 函数 \(u \mapsto \dim \operatorname{Coker}(u)\) 在从 E 到 F 中的严格单形态射所成的集 \(\mathcal{M}(\mathrm{E};\mathrm{F})\) 上是局部常数的。

事实上，有 \(\dim \operatorname{Coker}(u) = \operatorname{ind}(u)\) （对 \(u \in \mathcal{M}(\mathrm{E};\mathrm{F})\) ）。

引理 4. --- 设 E 和 F 是巴拿赫空间。为使 \(\mathcal{L}(\mathrm{E};\mathrm{F})\) 的元素 u 是从 E 到 F 中的严格态射，必须且只需 \(^{t}u\) 是从 F' 到 E' 中的严格态射。此时有 \(\dim \operatorname{Coker}(^{t}u)=\dim\operatorname{Ker}(u)\) 与 \(\dim\operatorname{Ker}(^{t}u)=\dim\operatorname{Coker}(u)\) 。

为使 u 是严格态射，必须且只需 \(^{t}u\) 是（EVT, IV, p. 29, 推论 3）；此时 u 的像是闭的（EVT, IV, p. 28, 定理 1），且向量空间 \(\mathrm{Ker}(^{t}u)\) （相应地，\(\mathrm{Coker}(^{t}u)\) ）典范地同构于赋范空间 \(\mathrm{Coker}(u)\) （相应地，\(\mathrm{Ker}(u)\) ）的对偶（依照 EVT, IV, p. 27, 命题 2）。引理得证。

设 E 和 F 是巴拿赫空间。用 \(\mathcal{E}(\mathrm{E};\mathrm{F})\) 表示从 E 到 F 中的连续满射线性映射所成的集，用 \(\mathcal{Q}\mathcal{E}(\mathrm{E};\mathrm{F})\) 表示从 E 到 F 中的、像有有限余维数的连续线性映射所成的集。\(\mathcal{Q}\mathcal{E}(\mathrm{E};\mathrm{F})\) 的每个元素都是具有闭像的严格态射（III, p. 52, 引理 6）。由引理 4 得出，\(\mathcal{E}(\mathrm{E};\mathrm{F})\) 与 \(\mathcal{Q}\mathcal{E}(\mathrm{E};\mathrm{F})\) 分别是 \(\mathcal{M}(\mathrm{F}^{\prime};\mathrm{E}^{\prime})\) 与 \(\mathcal{Q}\mathcal{M}(\mathrm{F}^{\prime};\mathrm{E}^{\prime})\) 在连续映射 \(u\mapsto{}^{t}u\) （从 \(\mathcal{L}(\mathrm{E};\mathrm{F})\) 到 \(\mathcal{L}(\mathrm{F}^{\prime};\mathrm{E}^{\prime})\) 中）下的原像。

命题 13. --- 设 E 和 F 是巴拿赫空间。集 \(\mathcal{E}(\mathrm{E};\mathrm{F})\) 与 \(\mathcal{Q}\mathcal{E}(\mathrm{E};\mathrm{F})\) 在 \(\mathcal{L}(\mathrm{E};\mathrm{F})\) 中是开的。更精确地说，若 \(u\) 是 \(\mathcal{Q}\mathcal{E}(\mathrm{E};\mathrm{F})\) 的元素，而 \(v\) 是 \(\mathcal{L}(\mathrm{E};\mathrm{F})\) 中满足 \(\| v - u\| < ((u))\) 的元素，则我们有 \(v \in \mathcal{Q}\mathcal{E}(\mathrm{E};\mathrm{F})\) ，且

\(\dim \operatorname{Ker}(v) \leqslant \dim \operatorname{Ker}(u), \quad \dim \operatorname{Coker}(v) \leqslant \dim \operatorname{Coker}(u),\)

\[
\operatorname{ind} (v) = \operatorname{ind} (u).
\]

我们前面已经看到 \(^{t}u \in \mathcal{Q}\mathcal{M}(F';E')\) 。此外，对每个元素 \(v\) （在 \(\mathcal{L}(E;F)\) 中），有 \(\|^{t}v - ^{t}u\| = \|v - u\|\) 与 \(((^{t}u)) = ((u))\) （EVT, IV, p. 7, 命题 8 及 III, p. 63, 命题 8）。由命题 12，从这些关系式得出：若 \(\|v - u\| < ((u))\) ，则 \(^{t}v\) 属于 \(\mathcal{Q}\mathcal{M}(F';E')\) ，且有不等式

\[
\dim \operatorname{Ker} (^ {t} v) \leqslant \dim \operatorname{Ker} (^ {t} u), \quad \dim \operatorname{Coker} (^ {t} v) \leqslant \dim \operatorname{Coker} (^ {t} u)
\]

以及等式 \(\mathrm{ind}(^{t}v)=\mathrm{ind}(^{t}u)\) 。于是命题由引理 4 得出。

推论 1. --- 由 \(u \mapsto \dim \operatorname{Ker}(u)\) 与 \(u \mapsto \dim \operatorname{Coker}(u)\) 在 \(\mathcal{Q}\mathcal{E}(\mathrm{E};\mathrm{F})\) 上定义的函数是上半连续的。函数 \(u \mapsto \operatorname{ind}(u)\) 在 \(\mathcal{Q}\mathcal{E}(\mathrm{E};\mathrm{F})\) 上是局部常数的。

推论 2. --- 由 \(u \mapsto \dim \operatorname{Ker}(u)\) 定义的函数在 \(\mathcal{E}(\mathrm{E};\mathrm{F})\) 上是局部常数的。

事实上，有 \(\dim \operatorname{Ker}(u) = -\operatorname{ind}(u)\) （对一切 \(u \in \mathcal{E}(\mathrm{E};\mathrm{F})\) ）。

\section{借紧线性映射的扰动}

\subsection{严格态射与真性}

引理 1. --- 设 E 和 F 是拓扑向量空间，u 是从 E 到 F 中的连续线性映射，U 是 E 中 0 的一个邻域。假设 u 诱导一个从 U 到 F 的闭子集上的同胚。则 u 的像 I 是闭的，且 u 诱导一个从 E 到 I 上的同胚。

集 \(u(\mathrm{U})\) 是 I 中 0 的一个邻域；它在 F 中是闭的，故 F 的子群 I 在 0 处局部闭，从而它是闭的（TG, III, p. 7, 命题 4）。

由于 \(\operatorname{Ker}(u)\) 与 U 只交于 0 这一点，有 \(\operatorname{Ker}(u) = \{0\}\) ，故映射 \(u\) 是单射。设 V 是包含于 U 的 E 中 0 的闭邻域。由于 \(u(\mathring{\mathrm{V}})\) 是 \(u(\mathrm{V})\) 中 0 的一个邻域，存在 F 中 0 的均衡邻域 W，使得 \(u(\mathrm{V}) \cap \mathrm{W} \subset u(\mathring{\mathrm{V}})\) 。集 \(u^{-1}(\mathrm{W})\) 是均衡的，故是连通的。它包含 0 且包含于 \(\mathring{\mathrm{V}} \cup (\mathrm{E} - \mathrm{V})\) 。由于 \(\mathring{\mathrm{V}}\) 与 E-V 是 E 的不相交开子集，集 \(u^{-1}(\mathrm{W})\) 包含于 \(\mathring{\mathrm{V}}\) ，从而 \(\mathrm{W} \cap \mathrm{I} \subset u(\mathring{\mathrm{V}})\) 。于是 \(u(\mathring{\mathrm{V}})\) 是 I 中 0 的一个邻域。这蕴含 \(u\) 诱导一个从 E 到 I 上的同胚。

命题 1. --- 设 E 是分离局部凸空间，F 是局部凸空间，u 是从 E 到 F 中的连续线性映射。下列条件等价：

\begin{enumerate}
\def\labelenumi{(\roman{enumi})}
\item
  映射 u 是严格态射，它的核有限维且它的像在 F 中是闭的；
\item
  存在 E 中 0 的闭邻域 V，使得 u 在 V 上的限制是真映射（TG, I, p. 72）。
\item
  \(\Longrightarrow\) (ii)：假设 u 满足条件 (i)。由于 u 的核有限维，它在 E 中有一个拓扑补 \(E_{1}\) （III, p. 55, 命题 1），且在 \(\operatorname{Ker}(u)\) 中存在 0 的一个紧邻域 C。把 E 与 \(E_{1} \times \operatorname{Ker}(u)\) 等同起来；此时集 \(V = E_{1} \times C\) 是 E 中 0 的一个闭邻域。u 在 V 上的限制是 \(E_{1} \times C\) 到 \(E_{1}\) 上的投影（它是真映射，TG, I, p. 77, 推论 5）与 \(u_{1}\) （即 u 在 \(E_{1}\) 上的限制）的复合。而 \(u_{1}\) 是从 \(E_{1}\) 到 F 的闭子空间上的一个同胚，故是真映射（TG, I, p. 72, 命题 2）。两个真映射的复合是真映射（TG, I, p. 73, 命题 5, a)），所以 u 在 V 上的限制是真映射。
\item
  \(\Longrightarrow\) (i)：设 V 是 E 中 0 的一个闭邻域，使得 u 在 V 上的限制 v 是真映射。此时集 \(\mathrm{V} \cap \mathrm{Ker}(u) = \overline{v}(\{0\})\) 是紧的（TG, I, p. 75, 定理 1）；于是向量空间 \(\mathrm{Ker}(u)\) 是局部紧的，从而有限维（EVT, I, p. 15, 定理 3）。设 \(E_{1}\) 是 \(\mathrm{Ker}(u)\) 在 E 中的一个拓扑补（III, p. 55 的命题 1）；令 \(V_{1} = E_{1} \cap V\) 。集 \(V_{1}\) 在 V 中是闭的。映射 \(u|V_{1}\) 是真映射（TG, I, p. 74, 推论 1）且是单射，故是从 \(V_{1}\) 到 F 的闭子集上的一个同胚（TG, I, p. 72, 命题 2）。由引理 1，u 在 \(E_{1}\) 上的限制是从 \(E_{1}\) 到 F 的闭子空间上的一个同胚，故 u 是具有闭像的严格态射。
\end{enumerate}

\subsection{单射或满射线性映射的扰动}

定理 1. --- 设 E 和 F 是分离局部凸空间，u 是从 E 到 F 中的严格态射，其核有限维且其像是闭的。设 h 是从 E 到 F 中的紧线性映射。线性映射 \(u + h\) 是严格态射，它的核有限维，且它的像是闭的。

由 III, p. 71 的命题 1，存在 E 中 0 的一个闭邻域 V，使得 \(u\) 在 V 上的限制是真映射。由于 \(h\) 是紧的，存在包含于 V 的 0 的闭邻域 W，使得集 \(h(\mathrm{W})\) 是相对紧的。令 \(\mathrm{C} = \overline{h(\mathrm{W})}\) 。\(u\) 在 W 上的限制是真映射（TG, I, p. 74, 推论 1）。映射 \(\alpha: x \mapsto (u(x), h(x))\) （从 W 到 \(\mathrm{F} \times \mathrm{C}\) 中）是真映射，因为映射 \(\mathrm{pr}_1 \circ \alpha = u|\mathrm{W}\) （从 W 到 F 中）是真映射（TG, I, p. 73, 命题 5, \(d\) ）。映射 \(\beta: (y, z) \mapsto (y + z, z)\) （从 \(\mathrm{F} \times \mathrm{C}\) 到 \(\mathrm{F} \times \mathrm{C}\) 中）是同胚，而映射 \(\mathrm{pr}_1\) （从 \(\mathrm{F} \times \mathrm{C}\) 到 F 中）是真映射（TG, I, p. 77, 推论 5）。于是复合映射 \(\mathrm{pr}_1 \circ \beta \circ \alpha\) （从 W 到 F 中）是真映射（TG, I, p. 73, 命题 5, \(a\) ）。但这个映射正是 \(u + h\) 在 W 上的限制。由 III, p. 71 的命题 1，\(u + h\) 是严格态射，它的核有限维，且它的像是闭的。

引理 2. --- 设 E 和 F 是弗雷歇空间，且 \(u \in \mathcal{L}(\mathrm{E};\mathrm{F})\) 。下列条件等价：

\begin{enumerate}
\def\labelenumi{(\roman{enumi})}
\item
  u 的余核有限维；
\item
  \(^{t}u \colon \mathrm{F}_c' \to \mathrm{E}_c'\) 是具有闭像且核有限维的严格态射。
\item
  \(\Longrightarrow\) (ii)：假设 u 的余核有限维。此时映射 u 是严格态射（III, p. 52 的引理 6）。由 EVT, IV, p. 27 的命题 2，\(^{t}u\) 的像在赋有弱拓扑的 E' 中是闭的，从而在 \(E_{c}'\) 中更是闭的。\(^{t}u\) 的核是 u 的像的正交补（同前）；故它有限维。最后，由 EVT, IV, p. 28 的定理 1，\(^{t}u\) 是从 \(F_{c}'\) 到 \(E_{c}'\) 中的严格态射。
\item
  \(\Longrightarrow\) (i)：假设 \(^t u\colon F_c' \to E_c'\) 是严格态射。由 EVT, IV, p. 28 的定理 1，\(u\) 的像是闭的。\(^t u\) 的核是 \(\operatorname{Im}(u)\) 的正交补（依照 EVT, IV, p. 27, 命题 2）；若这个核有限维，则 \(u\) 的像在 F 中有有限余维数。
\end{enumerate}

定理 2. --- 设 E 和 F 是弗雷歇空间，u: E → F 是余核有限维的连续线性映射，h: E → F 是紧线性映射。线性映射 u + h 是严格态射，它的余核有限维，且它的像是闭的。

由 III, p. 52 的引理 6，只需证明 \(u + h\) 的余核有限维。而由 III, p. 9 的命题 9，映射 \(^{t}h\) （从 \(F_{c}^{\prime}\) 到 \(E_{c}^{\prime}\) 中）是紧的。定理 2 由定理 1 与引理 2 得出。

\subsection{弗雷德霍姆映射的扰动}

定理 3. --- 设 E 是局部凸空间，F 是分离局部凸空间，u 是从 E 到 F 中的弗雷德霍姆映射，h 是从 E 到 F 中的紧线性映射。则 \(u + h\) 是弗雷德霍姆算子，且有 \(\operatorname{ind}(u + h) = \operatorname{ind}(u)\) 。

我们将分三步证明该定理。

\begin{enumerate}
\def\labelenumi{\Alph{enumi})}
\item
  假设 E 和 F 是巴拿赫空间。III, p. 42 的命题 2 表明线性映射 u 是具有闭像且核与余核均有限维的严格态射。由于 E 和 F 是巴拿赫空间，III, p. 73 的定理 1 与定理 2 蕴含：对每个 \(t \in [0,1]\) ，线性映射 \(u_t = u + th\) 具有这些同样的性质，从而是弗雷德霍姆算子（III, p. 42, 命题 2）。映射 \(t \mapsto u_t\) （从 \([0,1]\) 到从 E 到 F 的弗雷德霍姆映射所成的集 \(\mathcal{F}(\mathrm{E};\mathrm{F})\) 中）是连续的。由 III, p. 58 的定理 1，映射 \(t \mapsto \operatorname{ind}(u_t)\) 是局部常数的。由于区间 \([0,1]\) 连通，这个映射是常数。于是我们有 \(\operatorname{ind}(u) = \operatorname{ind}(u + h)\) 。
\item
  我们假设 E = F，\(u = 1_{E}\) 。由 III, p. 5 的命题 5，存在一个巴拿赫空间 G、一个连续线性映射 \(p: E \to G\) 和一个紧线性映射 \(q: G \to E\)，使得 \(h = q \circ p\) 。G 的自同态 \(p \circ q\) 是紧的。由 A)，\(1_{G} + p \circ q\) 是 G 的指标为零的弗雷德霍姆自同态。但此时 \(1_{E} + h = 1_{E} + q \circ p\) 是 E 的指标为零的弗雷德霍姆自同态（III, p. 49, 命题 10, f)）。
\item
  一般情形。设 \(v\) 是 \(u\) 的一个拟逆。自同态 \(u \circ v\) 与 \((u + h) \circ v\) （\(F\) 的）与 \(1_{F}\) 都只差一个紧线性映射。由 B)，它们是 \(F\) 的指标为零的弗雷德霍姆自同态。由此得出 \(u + h\) 是弗雷德霍姆自同态（III, p. 40, 第 2 目）且与 \(u\) 有相同的指标（III, p. 43, 第 3 目, 公式 (2)）。
\end{enumerate}

推论 1. --- 设 E 和 F 是分离局部凸空间，且 \(u \in \mathcal{L}(\mathrm{E};\mathrm{F})\) 。为使 u 是弗雷德霍姆算子，必须且只需存在一个映射 \(v \in \mathcal{L}(\mathrm{F};\mathrm{E})\) ，使得线性映射 \(1_{\mathrm{E}} - v \circ u\) 与 \(1_{\mathrm{F}} - u \circ v\) 都是紧的。

由于有限秩的连续线性映射是紧的，所述条件是必要的。

设 \(v\) 是 \(\mathcal{L}(\mathrm{F};\mathrm{E})\) 中使得线性映射 \(1_{\mathrm{E}} - v \circ u\) 与 \(1_{\mathrm{F}} - u \circ v\) 为紧的的元素。由定理 3，\(v \circ u\) 与 \(u \circ v\) 是弗雷德霍姆映射。设 \(p\) 与 \(q\) 分别是 \(v \circ u\) 与 \(u \circ v\) 的拟逆。令 \(w = p \circ v\) ，\(w' = v \circ q\) 。采用 III, p. 40 第 2 目的记号，我们有 \(w \circ u \equiv 1_{\mathrm{E}}\) 与 \(u \circ w' \equiv 1_{\mathrm{F}}\) ，由此 \(w \equiv w \circ u \circ w' \equiv w'\) 。于是得出 \(w\) 是 \(u\) 的一个拟逆。

推论 2. --- 设 E 和 F 是分离局部凸空间，且 \(u \in \mathcal{L}(\mathrm{E};\mathrm{F})\) 。下列条件等价：

\begin{enumerate}
\def\labelenumi{(\roman{enumi})}
\item
  映射 u 是指标为零的弗雷德霍姆映射；
\item
  存在一个从 E 到 F 上的同构 v，使得 u - v 有限秩；
\item
  存在同构 \(v\) （从 \(E\) 到 \(F\) 上的），使得 \(u - v\) 是紧的。
\item
  \(\Longrightarrow\) (ii)：假设 u 是指标为零的弗雷德霍姆算子。存在拓扑直和分解 \(E = E_{1} \oplus E_{2}\) 与 \(F = F_{1} \oplus F_{2}\) （其中 \(E_{2}\) 与 \(F_{2}\) 有限维），使得 u 在 \(E_{2}\) 上为零，且通过限制定义一个同构 \(u_{1}\) （从 \(E_{1}\) 到 \(F_{1}\) 上）（III, p. 42, 命题 2）。若 \(\text{ind}(u) = 0\) ，则 \(E_{2}\) 与 \(F_{2}\) 的维数相等，且存在从 E 到 F 上的一个同构 v，它是 \(u_{1}\) 的扩张。u - v 的核包含 \(E_{1}\) ，故 u - v 有限秩。
\end{enumerate}

由于从 E 到 F 中的每个有限秩连续线性映射都是紧的，我们有 (ii) \(\Longrightarrow\) (iii)；而由定理 3 有 (iii) \(\Longrightarrow\) (i)。

\subsection{里斯自同态的扰动}

设 E 是巴拿赫空间。回忆（参见 III, p. 5, 命题 3）：E 的紧自同态所成的集 \(\mathcal{L}^c (\mathrm{E})\) 构成巴拿赫代数 \(\mathcal{L}(\mathrm{E})\) 的一个闭双边理想。商巴拿赫代数称为 E 的卡尔金代数；记作 \(\mathcal{Calk}(\mathrm{E})\) 。我们用 \(\pi\) 表示从 \(\mathcal{L}(\mathrm{E})\) 到 \(\mathcal{Calk}(\mathrm{E})\) 上的典范同态。由 III, p. 74 的推论 1，E 的自同态 \(u\) 是弗雷德霍姆自同态，当且仅当 \(\pi (u)\) 在代数 \(\mathcal{Calk}(\mathrm{E})\) 中可逆。

命题 2. --- 设 \(u \in \mathcal{L}(\mathrm{E})\) 满足 \(\|1_{\mathrm{E}} - \pi(u)\| < 1\) （在代数 \(\mathcal{Calk}(\mathrm{E})\) 中）。则 \(u\) 是 E 的里斯自同态。

设 \(r \geqslant 1\) 是使 \(\| 1_{\mathrm{E}} - \pi(u)\|^{r} < \frac{1}{2}\) 的整数。设 \(\mathrm{P} \in \mathbf{C}[\mathrm{X}]\) 是多项式 \(\frac{1 - (1 - \mathrm{X})^r}{\mathrm{X}}\) 。令 \(v = 1_{\mathrm{E}} - (1_{\mathrm{E}} - u)^r\) 。于是我们有 \(v = u\mathrm{P}(u)\) 与 \(\| 1_{\mathrm{E}} - \pi(v)\| < \frac{1}{2}\) 。由于 E 的自同态 \(u\) 与 \(\mathrm{P}(u)\) 交换，只需证明 \(v\) 是 E 的里斯自同态（III, p. 49, 命题 9）。

由于 \(\|1_{\mathrm{E}}-\pi(v)\|<1/2\) ，由空间 \(\mathcal{Calk}(\mathrm{E})\) 中商范数的定义，存在 E 的一个紧自同态 h 和 E 的一个自同态 w，使得 \(v=1_{E}+h+w\) 且 \(\|w\|<\frac{1}{2}\) 。由 I, p. 22 的推论 1，元素 \(1_{E}+w\) 是 E 的一个自同构。由于 h 是紧的，\(v=(1_{\mathrm{E}}+w)+h\) 是 E 的指标为 0 的弗雷德霍姆自同态（III, p. 74 的推论 2）。对每个整数 \(n\geqslant0\) ，用 \(N_{n}\) 表示 \(v^{n}\) 的核。为证明 v 是 E 的里斯自同态，只需证明存在一个整数 \(n \geqslant 0\) 使 \(N_{n} = N_{n+1}\) （III, p. 46, 定义 2 及 III, p. 46, 注记）。

我们用反证法，假设序列 \((\mathrm{N}_{n})\) 严格递增。对每个 \(n \in N\) ，用 \(p_{n}\) 表示从 E 到赋范空间 \(E/N_{n}\) 上的典范映射。设 c 是使 \(2\|w\| < c < 1\) 的实数。设 \(n \in N\) 。由于 \(N_{n+1}\) 异于 \(N_{n}\) ，存在一个元素 \(x_{n} \in N_{n+1}\) ，使得 \(\|p_{n}(x_{n})\| = c\) 且 \(\|x_{n}\| < 1\) （事实上，存在 \(y \in N_{n+1}/N_{n}\) 范数为 c，故对每个 \(\varepsilon > 0\) ，存在 \(x_{n} \in N_{n+1}\) 使得 \(p_{n}(x_{n}) = y\) 且 \(\|x_{n}\| \leqslant c + \varepsilon\) ）。

设 \(m\) 、\(n\) 是使 \(m > n \geqslant 0\) 的两个整数。我们有

\[
\begin{array}{r l} & {\| h (x _ {m}) - h (x _ {n}) \| = \| v (x _ {m} - x _ {n}) - (1 _ {\mathrm{E}} + w) (x _ {m} - x _ {n}) \|} \\ & {\qquad \geqslant \| (v - 1 _ {\mathrm{E}}) (x _ {m} - x _ {n}) \| - \| w (x _ {m} - x _ {n}) \|} \\ & {\qquad \geqslant \| (v - 1 _ {\mathrm{E}}) (x _ {m} - x _ {n}) \| - 2 \| w \|} \end{array}\tag{1}
\]

因为 \(\| x_m\| \leqslant 1\) 且 \(\| x_n\| \leqslant 1\) 。此外，注意到

\[
(v - 1 _ {\mathrm{E}}) (x _ {m} - x _ {n}) = v (x _ {m} - x _ {n}) + x _ {n} - x _ {m}.
\]

但由于 n \textless{} m 且 \(v(\mathrm{N}_{m+1}) \subset \mathrm{N}_{m}\) ，我们有 \(v(x_{m} - x_{n}) + x_{n} \in \mathrm{N}_{m}\) ，由此

\[
\left\| v (x _ {m} - x _ {n}) + x _ {n} - x _ {m} \right\| \geqslant \left\| p _ {m} (x _ {m}) \right\| = c.
\]

于是不等式 (1) 给出下界

\[
\left\| h \left(x _ {m}\right) - h \left(x _ {n}\right) \right\| \geqslant c - 2 \| w \| > 0.\tag{2}
\]

于是序列 \((h(x_{m}))_{m\in\mathbf{N}}\) 在 E 中没有聚点；由于 h 是紧的（III, p. 2 的注记 1），E 的单位球在 h 下的像是相对紧的，这就产生矛盾。

推论 1. --- 设 E 是分离局部凸空间，且 \(u \in \mathcal{L}(\mathrm{E})\) 。为使 u 是里斯自同态，必须且只需存在 \(\mathcal{L}(\mathrm{E})\) 的一个元素 v，它与 u 交换且使 \(1_{\mathrm{E}} - u \circ v\) 是紧的。

由于每个有限秩连续线性映射都是紧的，该条件是必要的（III, p. 47, 命题 6 d)）。我们来证明它是充分的。设 v 是 \(\mathcal{L}(\mathrm{E})\) 的一个元素，它与 u 交换，且使 E 的自同态 \(h = 1_{E} - u \circ v\) 是紧的。存在一个巴拿赫空间 G、一个连续线性映射 \(p: E \to G\) 和一个紧线性映射 \(q: G \to E\)，使得 \(h = q \circ p\) （III, p. 5, 命题 5）。G 的自同态 \(p \circ q\) 是紧的。由命题 2，线性映射 \(1_{G} - p \circ q\) 是 G 的里斯自同态。但此时 \(u \circ v = 1_{E} - q \circ p\) 是 E 的里斯自同态（III, p. 49, 命题 10, g)）。由于 u 与 v 交换，u 是 E 的里斯自同态（III, p. 49, 命题 9）。

推论 2. --- 设 E 是分离局部凸空间，u 是 E 的里斯自同态，\(h \in \mathcal{L}(\mathrm{E})\) 是紧线性映射。若 u 与 h 交换，则 \(u + h\) 是 E 的里斯自同态。

设 \(v\) 是 \(u\) 的典范拟逆。它与 \(u\) 交换，也与 \(h\) 交换（III, p. 47, 命题 6），从而与 \(u + h\) 交换。E 的自同态 \(1_{\mathrm{E}} - (u + h) \circ v\) 是有限秩连续线性映射 \(1_{\mathrm{E}} - u \circ v\) 与紧线性映射 \(-h \circ v\) 之和，故是紧的。由此得出 \(u + h\) 是 E 的里斯自同态（推论 1）。

命题 3. --- 设 E 是巴拿赫空间。设 u 是 E 的一个自同态，A 是它在 \(\mathcal{L}(\mathrm{E})\) 中的交换子。\(\pi(\mathrm{A})\) 在代数 \(\mathcal{Calk}(\mathrm{E})\) 中的闭包 B 是一个巴拿赫代数。为使 u 是 E 的里斯自同态，必须且只需 \(\pi(u)\) 在 B 中可逆。

由于 \(\pi(A)\) 是 \(\mathcal{Calk}(E)\) 的一个含单位元的子代数，其闭包 B 是一个巴拿赫代数。若 \(u\) 是 E 的里斯自同态，则元素 \(\pi(u)\) 在 \(\pi(A)\) 中有一个逆（推论 1），从而在 B 中有逆。反之，假设 \(\pi(u)\) 在 B 中有一个逆 \(s\) 。由 B 的定义，存在 A 的一个元素 \(v\) 使得

\[
\left\| s - \pi (v) \right\| <   \frac {1}{\left\| \pi (u) \right\|},
\]

由此 \(\|1_{E}-\pi(u\circ v)\|<1\) 。由命题 2，线性映射 \(u\circ v\) 是 E 的里斯自同态。由于 u 与 v 交换，u 也是（III, p. 49, 命题 9）。

\subsection{弗雷德里克·里斯理论}

设 E 是分离局部凸空间，h 是 E 的一个紧自同态。设 \(\lambda \in K\) 。自同态 \(\lambda h\) 是紧的，故 \(1_{E} - \lambda h\) 是 E 的里斯自同态（III, p. 75 的命题 2 的推论 2）；用 \(N_{\lambda}\) 与 \(I_{\lambda}\) 表示它的幂零空间与余幂零空间。由 III, p. 47 的命题 6，我们有下列性质：

\begin{enumerate}
\def\labelenumi{(\roman{enumi})}
\item
  E 的向量子空间 \(I_{\lambda}\) 与 \(N_{\lambda}\) 是闭的且在 h 下稳定；
\item
  局部凸空间 E 是 \(I_{\lambda}\) 与 \(N_{\lambda}\) 的拓扑直和；
\item
  映射 \(1_{\mathrm{E}} - \lambda h\) 通过限制诱导 \(\mathrm{I}_{\lambda}\) 的一个自同构；
\item
  向量空间 \(N_{\lambda}\) 有限维，且存在一个整数 \(n_{\lambda} \geqslant 0\) ，使得 \((1_{\mathrm{E}} - \lambda h)^{n_{\lambda}}(x) = 0\) （对一切 \(x \in N_{\lambda}\) ）。
\end{enumerate}

这些性质以唯一的方式确定空间 \(I_{\lambda}\) 与 \(N_{\lambda}\) （A, VIII, p. 26, 注记 2）。

引理 3. --- 设 X 是一个拓扑具有可数基的拓扑空间。X 的每个离散子集都是可数的。

设 \(\mathcal{B}\) 是 X 的拓扑的一个可数基，D 是 X 的一个离散子集。对 D 的每个元素 \(x\) ，存在元素 \(\mathrm{B}_x\) （属于 \(\mathcal{B}\) ），使得 \(\mathrm{B}_x \cap \mathrm{D} = \{x\}\) 。从 D 到 \(\mathcal{B}\) 的映射由 \(x \mapsto \mathrm{B}_x\) 定义，它是单射，由此引理得证。

定理 4. --- 设 E 是分离局部凸空间，h 是 E 的一个紧自同态。集 \(\Sigma\) 由那些 \(\lambda \in K\) （使 \(1_{\mathrm{E}} - \lambda h\) 不是 E 的自同构的）组成，它是域 K 的一个可数、闭且离散的子集。

集 \(\Sigma\) 也是由那些 \(\lambda\in\mathrm{K}\) （使 \(1_{\mathrm{E}}-\lambda h\) 不是单射的）所成的集，以及由那些 \(\lambda\in\mathrm{K}\) （使 \(1_{\mathrm{E}}-\lambda h\) 不是满射的）所成的集。

存在一个巴拿赫空间 G、一个连续线性映射 \(p: E \rightarrow G\) 和一个紧线性映射 \(q: G \rightarrow E\)，使得 \(h = q \circ p\) （III, p. 5, 命题 5）。G 的自同态 \(h' = p \circ q\) 是紧的，且 \(1_{E} - \lambda h\) 是 E 的自同构当且仅当 \(1_{G} - \lambda h'\) 是 G 的自同构（III, p. 49 的命题 10, e)）。因此我们只需在 E 是巴拿赫空间的情形证明该定理，以下即作此假设。

设 \(\lambda \in K\) 。由于里斯自同态的指标为零，下列说法等价：映射 \(1_{\mathrm{E}} - \lambda h\) 是 E 的自同构，或它是单射的，或它是满射的。

集 \(\Sigma\) 在 K 中是闭的，因为 E 的自同构所成的集在 \(\mathcal{L}(\mathrm{E})\) 中是开的。设 \(\lambda\) 是 \(\Sigma\) 的一个元素。我们有 \(\lambda \neq 0\) 。用 N 表示 h 的幂零空间，I 表示其余幂零空间，并用 \(h_{I}\) 与 \(h_{N}\) 表示 h 在这些子空间上诱导的 I 与 N 的自同态。则 \(1_{I} - \lambda h_{I}\) 是 I 的一个自同构，且存在 \(\lambda\) 在 K 中的一个邻域 U，使得 \(1_{I} - \mu h_{I}\) 是 I 的自同构（对一切 \(\mu \in U\) ）。N 的自同态 \(1_{N} - \lambda h_{N}\) 是幂零的；对一切 \(\mu \neq \lambda\) ，我们有

\[
1 _ {\mathrm{N}} - \mu h _ {\mathrm{N}} = \frac {\lambda - \mu}{\lambda} \Big (1 _ {\mathrm{N}} + \frac {\mu}{\lambda - \mu} (1 _ {\mathrm{N}} - \lambda h _ {\mathrm{N}}) \Big),
\]

故 \(1_{N} - \mu h_{N}\) 是 N 的一个自同构。于是集 \(U \cap \Sigma\) 缩减为单个元素 \(\lambda\) ，从而集 \(\Sigma\) 是离散的。由引理 3 它是可数的。

\subsection{弗雷德霍姆择一定理}

命题 4（弗雷德霍姆择一定理）. --- 设 E 是巴拿赫空间，h 是 E 的一个紧自同态。设 \(\lambda\) 是 K-\{0\} 的一个元素。设 F 是 \(1_{\mathrm{E}} - \lambda h\) 的核，G 是 \(1_{\mathrm{E}'} - \lambda^{t}h\) 的核。

\begin{enumerate}
\def\labelenumi{\alph{enumi})}
\item
  空间 F 与 G 有相同的有限维数 n ≥ 0；
\item
  对 \(y \in \mathrm{E}\) ，存在 \(x \in \mathrm{E}\) 使 \(x - \lambda h(x) = y\) ，当且仅当 \(\langle y, \ell \rangle = 0\) （对一切 \(\ell \in \mathrm{G}\) ）；
\item
  对 \(\ell \in \mathrm{E}'\) ，存在 \(f \in \mathrm{E}'\) 使 \(f - \lambda^t h(f) = \ell\) ，当且仅当 \(\langle x, f \rangle = 0\) （对一切 \(x \in \mathrm{F}\) ）。
\end{enumerate}

特别地，下列条件中恰有一个成立：

\begin{enumerate}
\def\labelenumi{(\roman{enumi})}
\item
  空间 F 非零；
\item
  对每个 \(y \in \mathrm{E}\) ，存在 \(x \in \mathrm{E}\) 使 \(x - \lambda h(x) = y\) 。
\end{enumerate}

通过把 h 换成 \(\lambda h\) ，我们化归到 \(\lambda = 1\) 的情形。

自同态 \(^{t}h\) 由 III, p. 9 的推论 1 是紧的。于是 E 的自同态 \(w = 1_{\mathrm{E}} - h\) 与 \(w' = 1_{\mathrm{E}'} - ^{t}h\) （\(\mathrm{E}'\) 的）都是里斯自同态（III, p. 75 的命题 2 的推论 2）。特别地，它们的核 F 与 G 是有限维的。由于 \(w' = ^{t}w\) 且 w 的指标为零，由 III, p. 43 的公式 (3)，F 的维数等于 G 的维数。

对 E 的元素 \(y\) ，方程 \(x - h(x) = y\) 有解 \(x \in \mathrm{E}\) 当且仅当 \(x\) 属于 \(w\) 的像。由于 \(w\) 的余核的对偶与 \(\operatorname{Ker}(^t w)\) 的核 G 等同（EVT, IV, p. 27, 命题 2），这件事成立当且仅当 \(\langle y, \ell \rangle = 0\) （对一切 \(\ell \in \mathrm{G}\) ）。

对 E' 的元素 \(\ell\) ，方程 \(f - {}^t h(f) = \ell\) 有解 \(f \in \mathrm{E}'\) 当且仅当 \(f\) 属于 \(w'\) 的像。由于 \({}^t w\) 的像是 \(w\) 的核 F 的正交（EVT, IV, p. 27, 命题 2），这件事成立当且仅当 \(\langle x, f \rangle = 0\) （对一切 \(x \in \mathrm{F}\) ）。

例. --- 设 X 是紧空间，\(\mu\) 是 X 上的一个测度（依 K 等于 R 或 C 而为实或复），N: \(X \times X \to K\) 是一个连续函数。

用 E 表示巴拿赫空间 \(\mathcal{C}(\mathrm{X};\mathrm{K})\) 。假设给定了函数 \(a\in\mathrm{E}\) 与 \(\lambda\in\mathrm{K}-\{0\}\) 。我们研究积分方程的解 \(f\in\mathrm{E}\) 的存在性与唯一性问题

\[
f (x) - \lambda \int_ {\mathrm{X}} \mathrm{N} (x, y) f (y) d \mu (y) = a (x) \quad (x \in \mathrm{X}).\tag{3}
\]

为此，引入集 \(F_{\lambda}\) （由相应的齐次方程的解 \(f \in E\) 组成）

\[
f (x) - \lambda \int_ {\mathrm{X}} \mathrm{N} (x, y) f (y) d \mu (y) = 0 \qquad (x \in \mathrm{X}),\tag{4}
\]

以及集 \(G_{\lambda}\) （由 (4) 的转置齐次方程的解 \(g \in E\) 组成），即

\[
g (y) - \lambda \int_ {\mathrm{X}} \mathrm{N} (x, y) g (x) d \mu (x) = 0 \qquad (y \in \mathrm{X}).\tag{5}
\]

若 (3) 的解集非空，则它是 E 的一个以 \(F_{\lambda}\) 为方向的仿射子空间。

对 \(f\in \mathrm{E}\) 与 \(x\in \mathrm{X}\) ，令

\[
h (f) (x) = \int_ {\mathrm{X}} \mathrm{N} (x, y) f (y) d \mu (y);
\]

如此定义的映射 \(f \mapsto h(f)\) 是巴拿赫空间 E 的一个紧自同态（III, p. 26 的命题 2 的推论）。E 的对偶 E' 由 X 上的测度组成（依 K 等于 R 或 C 而为实或复）（INT, III, p. 47, § 1, 第 3 目, 定义 2）。h 的转置 \(^{t}h\) 是 E' 的自同态，由关系式 \(\langle f, ^{t}h(m)\rangle = \langle h(f), m\rangle\) （对 \(f \in E\) 与 \(m \in E'\) ）刻画，即

\[
\begin{array}{r l} \int_ {\mathrm{X}} f (x) d (^ {t} h (m)) (x) & = \int_ {\mathrm{X}} h (f) (x) d m (x) \\ & = \int_ {\mathrm{X}} \Big (\int_ {\mathrm{X}} \mathrm{N} (x, y) f (y) d \mu (y) \Big) d m (x) \\ & = \int_ {\mathrm{X}} \Big (\int_ {\mathrm{X}} \mathrm{N} (x, y) d m (x) \Big) f (y) d \mu (y) \end{array}\tag{6}
\]

（对 \(m \in E'\) 与 \(f \in E\) ；INT, III, p. 84, § 4, 第 1 目, 定理 2）。

引理 4. --- 从 \(G_{\lambda}\) 到 \(E'\) 中的线性映射（由 \(g \mapsto g \cdot \mu\) 定义）是单射的，其像就是 \(1_{E'} - \lambda^{t}h\) 的核。

若 \(g \in G_{\lambda}\) 满足 \(g \cdot \mu = 0\) ，则公式 (5) 蕴含 \(g = 0\) ，故映射 \(g \mapsto g \cdot \mu\) 是单射的。

设 \(m \in \mathrm{E}'\) 是属于 \(1_{\mathrm{E}'} - \lambda^t h\) 的核的一个测度。由关系式 (6)，我们有

\[
\int_ {\mathrm{X}} f (x) d m (x) = \lambda \int_ {\mathrm{X}} \left(\int_ {\mathrm{X}} \mathrm{N} (x, y) d m (x)\right) f (y) d \mu (y)\tag{7}
\]

（对每个连续函数 \(f: \mathrm{X} \to \mathrm{K}\) ）。于是测度 \(m\) 等于测度 \(g \cdot \mu\) ，其中 \(g\) 是由下式定义的从 \(\mathrm{X}\) 到 \(\mathrm{K}\) 中的连续函数

\[
g (y) = \lambda \int_ {\mathrm{X}} \mathrm{N} (x, y) d m (x)\tag{8}
\]

（对一切 \(y \in X\) ）。于是公式 (8) 蕴含函数 g 属于 \(G_{\lambda}\) 。

反之，设 \(g: X \to K\) 是属于 \(G_{\lambda}\) 的一个连续函数，令 \(m = g \cdot \mu\) 。于是对每个连续函数 \(f: X \to K\) ，我们有

\[
\begin{array}{l} \int_ {\mathrm{X}} f (y) d m (y) = \int_ {\mathrm{X}} g (y) f (y) d \mu (y) \\ \qquad = \lambda \int_ {\mathrm{X}} \Big (\int_ {\mathrm{X}} \mathrm{N} (x, y) g (x) d \mu (x) \Big) f (y) d \mu (y) \\ \qquad = \lambda \int_ {\mathrm{X}} f (x) d m (x), \end{array}
\]

从而 \({}^{t}h(m)=\lambda m\) 。这就证明了引理 4。

把 III, p. 78 的命题 4 与定理 4 应用于自同态 h，我们便得到下列命题。

定理 5. --- 设 \(\lambda\in K\) 。

\begin{enumerate}
\def\labelenumi{\alph{enumi})}
\item
  集 \(F_{\lambda}\) 与 \(G_{\lambda}\) 是 \(\mathcal{C}(X;K)\) 的有限维向量子空间，且它们的维数相等；
\item
  为使方程 (3) 至少有一个解 \(f \in \mathcal{C}(\mathrm{X}, \mathrm{K})\) ，必须且只需有 \(\int_{\mathrm{X}} a(x) g(x) d\mu(x) = 0\) （对每个函数 \(g \in G_{\lambda}\) ）。此时 (3) 的解集是 \(\mathcal{C}(\mathrm{X}; \mathrm{K})\) 的一个以 \(F_{\lambda}\) 为方向的仿射子空间；
\item
  下列互相排斥的条件之一成立：
\end{enumerate}

\begin{enumerate}
\def\labelenumi{(\roman{enumi})}
\item
  对任何函数 \(a \in \mathcal{C}(\mathrm{X};\mathrm{K})\) ，积分方程 (3) 存在唯一的解 \(f \in \mathcal{C}(\mathrm{X};\mathrm{K})\) ；
\item
  齐次方程 (4) 有一个非零解 \(f \in \mathcal{C}(\mathrm{X};\mathrm{K})\) 。
\end{enumerate}

推论. --- 由 \(\lambda\in\mathrm{K}\) （使空间 \(\mathrm{F}_{\lambda}\) 非零的）组成的集是 K 的一个可数、闭且离散的子集。

\section{巴拿赫空间自同态的谱性质}

除另有说明外，本段所考虑的向量空间都是 C 上的向量空间。完备可赋范空间 E 的自同态的谱是相对于含单位元代数 \(\mathcal{L}(\mathrm{E})\) 的谱（参见 I, p. 127 的第 7 目）。

\subsection{谱的孤立点与敏感点}

设 E 是完备可赋范空间，u 是 E 的一个自同态。设 \(\lambda\) 是 u 的谱的一个孤立点。回忆我们用 \(e_{\lambda}(u)\) 表示与 u 及 u 的谱的既闭又开的子集 \(\{\lambda\}\) 相伴的谱投影（I, p. 131 的第 3 目）。自同态 \(u - \lambda1_{E}\) 诱导 \(\operatorname{Ker}(e_{\lambda}(u))\) 的一个自同构以及 \(\operatorname{Im}(e_{\lambda}(u))\) 的一个拟幂零自同态（同前）。

定义 1. --- 我们说 \(\lambda\) 对 \(u\) 有有限谱重数，如果谱投影 \(e_{\lambda}(u)\) 有限秩；此时，整数 \(m_{\lambda}(u) = \dim (\mathrm{Im}(e_{\lambda}(u)))\) 称为 \(\lambda\) 对 \(u\) 的谱重数。

定义 2. --- 我们把 Sp(u) 的有限谱重数的孤立点称为 u 的谱的敏感点。Sp(u) 的敏感点所成的集称为 u 的敏感谱，记作 Sp \(_{s}\) (u)。\(^{(2)}\) 我们把 u 的本质谱记作 Sp \(_{e}\) (u)，即补集 \(\mathrm{Sp}(u) - \mathrm{Sp}_{\mathrm{s}}(u)\)。

由于集 \(\mathrm{Sp}_{\mathrm{s}}(u)\) 由 \(\mathrm{Sp}(u)\) 的孤立点组成，它在 \(\mathrm{Sp}(u)\) 中是开的、离散的且可数的（III, p. 78, 引理 3）；它还是有界的。集 \(\mathrm{Sp}_{\mathrm{e}}(u)\) 在 \(\mathrm{Sp}(u)\) 中是闭的，故是紧的。

由 III, p. 54 的命题 14，\(u\) 的谱的敏感点是那些复数 \(\lambda\) （使 \(u - \lambda 1_{\mathrm{E}}\) 为 E 的里斯自同态且不是 E 的自同构的）。

对每个复数 \(\lambda\) ，用 \(\mathrm{N}_{\lambda}(u)\) 与 \(\mathrm{I}_{\lambda}(u)\) 表示 E 的自同态 \(u - \lambda 1_{\mathrm{E}}\) 的幂零空间与余幂零空间。由于 E 是弗雷歇空间，\(u - \lambda 1_{\mathrm{E}}\) 是里斯自同态当且仅当 \(\mathrm{N}_{\lambda}(u)\) 有限维且 \(\mathrm{I}_{\lambda}(u)\) 有有限余维数（III, p. 53 的命题 13）；在此情形，\(u - \lambda 1_{\mathrm{E}}\) 可逆当且仅当 \(\mathrm{N}_{\lambda}(u)\) 缩减为 0（III, p. 47, 命题 6）。于是 \(u\) 的谱的敏感点是复数 \(\lambda\) （使 \(\mathrm{N}_{\lambda}(u)\) 为非零有限维且 \(\mathrm{I}_{\lambda}(u)\) 在 E 中有有限余维数的）（III, p. 54, 命题 14）。此时谱重数 \(m_{\lambda}(u)\) （\(\lambda\) 的）就是 \(\mathrm{N}_{\lambda}(u)\) 的维数，且自同态 \(u - \lambda 1_{\mathrm{E}}\) 通过限制诱导 \(\mathrm{I}_{\lambda}(u)\) 的一个自同构与 \(\mathrm{N}_{\lambda}(u)\) 的一个幂零自同态。由于 \(\mathrm{N}_{\lambda}(u)\) 不缩减为 0，由此得出 \(\lambda\) 是 \(u\) 的一个特征值。最小的整数 \(n \geqslant 1\) （使 \(\operatorname{Ker}((u - \lambda 1_{\mathrm{E}})^n)\) 等于 \(\mathrm{N}_{\lambda}(u)\) ）是 \(u\) 的预解式在点 \(\lambda\) 处极点的阶（I, p. 131 的第 3 目），它以 \(m_{\lambda}(u)\) 为上界。

特别地，我们由此证明了：

命题 1. --- u 的谱的敏感点是有限谱重数的特征值。

\(u\) 的本质谱由那些复数 \(\lambda\) （使 \(u - \lambda 1_{\mathrm{E}}\) 不是 E 的里斯自同态的）组成。特别地，若 \(u\) 是紧的，则每个 \(\lambda \in \mathrm{Sp}(u) - \{0\}\) 都是谱的敏感点（III, p. 77 的推论 2）。

命题 2. --- 设 E 是完备可赋范空间，u 是 E 的一个自同态。设 H 是 \(\mathrm{Sp}_{\mathrm{s}}(u)\) 的一个有限子集。集 H 在 \(\mathrm{Sp}(u)\) 中既开又闭。谱投影 \(e_{\mathrm{H}}(u)\) 有限秩，其像为 \(\sum_{\lambda \in \mathrm{H}}\mathrm{N}_{\lambda}(u)\) ，核为 \(\bigcap_{\lambda \in \mathrm{H}}\mathrm{I}_{\lambda}(u)\) 。

这由 I, p. 131 的第 3 目得出，因为对每个 \(\lambda \in \mathrm{Sp}_{\mathrm{s}}(u)\) ，子空间 \(\mathrm{N}_{\lambda}(u)\) 与 \(\mathrm{I}_{\lambda}(u)\) 分别是谱投影 \(e_{\lambda}(u)\) 的像与核。

推论 1. --- 设 V 是 \(\mathrm{Sp_e}(u)\) 的一个邻域。

\begin{enumerate}
\def\labelenumi{\alph{enumi})}
\item
  存在 E 作为拓扑直和 F ⊕ G 的分解，使得 F 有限维，F 与 G 在每个与 u 交换的自同态下稳定，且 u 诱导的 G 的自同态的谱包含于 V。
\item
  假设 V 非空。存在元素 v （属于 u 在 \(\mathcal{L}(\mathrm{E})\) 中的双交换子），其谱包含于 V 且 \(v - u\) 有限秩。
\end{enumerate}

令 \(\mathrm{H} = \mathrm{Sp}(u) \cap (\mathbf{C} - \mathring{\mathrm{V}})\) 。集 H 包含于 \(\mathrm{Sp}_{\mathrm{s}}(u)\) ，故是离散的；由于它在紧空间 \(\mathrm{Sp}(u)\) 中是闭的，它是有限的。可取谱投影 \(e_{\mathrm{H}}(u)\) 的像与核作为 F 与 G（命题 2）。

假设 V 非空。设 \(\mu\in\mathrm{V}\) ，用 \(v\) 表示 E 的这样一个自同态：它在 F 上与比率为 \(\mu\) 的位似重合，在 G 上与 \(u\) 重合。其谱包含于 \(\{\mu\}\cup(\mathrm{Sp}(u)-\mathrm{H})\) ，从而包含于 V，且 \(v-u\) 有限秩，因为 F 有限维。

对 C 的每个紧子集 S，我们用 \(\mathcal{O}(S)\) 表示在 S 的一个邻域中定义、取值于 C 的全纯函数芽所成的代数（I, p. 49, §4, 第 1 目）。

推论 2. --- 设 \(f \in \mathcal{O}(\mathrm{Sp}(u))\) ，\(\mu\) 是一个复数。设 H 是 \(\lambda \in \mathrm{Sp}(u)\) （满足 \(f(\lambda) = \mu\) ）所成的集。则 \(\mu\) 是 \(f(u)\) 的谱的敏感点，当且仅当 H 有限、非空且包含于 \(u\) 的敏感谱。在这些条件下，谱投影 \(e_{\mu}(f(u))\) 等于投影 \(e_{\mathrm{H}}(u)\) （与 \(u\) 及 H 相伴的）。我们有

\[
\mathrm{N} _ {\mu} (f (u)) = \bigoplus_ {\lambda \in \mathrm{H}} \mathrm{N} _ {\lambda} (u), \quad \mathrm{I} _ {\mu} (f (u)) = \bigcap_ {\lambda \in \mathrm{H}} \mathrm{I} _ {\lambda} (u)
\]

而 \(\mu\) 对 \(f(u)\) 的谱重数是 H 的元素的谱重数之和，即

\[
m _ {\mu} (f (u)) = \sum_ {\lambda \in \mathrm{H}} m _ {\lambda} (u).
\]

复数 \(\mu\) 属于 \(f(\mathrm{Sp}(u))\) ，从而属于 \(f(u)\) 的谱（I, p. 75 的命题 8），当且仅当 H 非空。此时我们有 \(e_{\mu}(f(u)) = e_{\mathrm{H}}(u)\) （I, p. 127 的第 1 目）。其余断言由命题 2 得出。

例. --- 1) 设 E 是有限维向量空间，u 是 E 的一个自同态。u 的谱有限且与 \(\mathrm{Sp}_{\mathrm{s}}(u)\) 重合。其元素是 u 的特征多项式 \(\chi_{u}\) 的根；它们就是 u 的特征值。谱重数 \(m_{\lambda}(u)\) （一个元素 \(\lambda \in \mathrm{Sp}(u)\) 的）是 \(\lambda\) 作为多项式 \(\chi_{u}\) 的根的重数（A, VII, p. 36, 推论）。由上面的推论 2，我们有

\[
\chi_ {f (u)} (\mathrm{T}) = \prod_ {\lambda \in \operatorname{Sp} (u)} (\mathrm{T} - f (\lambda)) ^ {m _ {\lambda}}
\]

（对一切 \(f \in \mathcal{O}(\mathrm{Sp}(u))\) ）。这推广了 A, VII, p. 36 的命题 10。

\begin{enumerate}
\def\labelenumi{\arabic{enumi})}
\setcounter{enumi}{1}
\tightlist
\item
  设 X 是 C 的一个紧子集，\(\mathcal{C}(X)\) 是 X 上取复值的连续函数所成的巴拿赫空间。用 u 表示 \(\mathcal{C}(X)\) 的这样一个自同态：它把 \(f \in \mathcal{C}(X)\) 映为函数 \(z \mapsto zf(z)\) （在 \(\mathcal{C}(X)\) 中） 。u 的谱等于 X，其敏感点是 X 的孤立点，且它们的谱重数等于 1。
\end{enumerate}

\subsection{谱的一个分划}

用 \(\overline{Z}\) 表示子集 \(Z \cup \{-\infty, +\infty\}\) （\(\overline{R}\) 中的）。若 u 是一个核或余核有限维的线性映射，则称 u 的指标为由下式定义的元素 \(\text{ind}(u)\) （\(\overline{Z}\) 中的）

\[
\operatorname{ind} (u) = \dim \operatorname{Coker} (u) - \dim \operatorname{Ker} (u)
\]

（参见 III, p. 67 的第 6 目）。

定义 3. --- 设 E 是完备可赋范空间，u 是 E 的一个自同态。对每个 \(n \in \overline{Z}\) ，用 \(\mathrm{Sp}_{n}(u)\) 表示这样的复数 \(\lambda \in \mathrm{Sp}_{\mathrm{e}}(u)\) 所成的集：使 \(u - \lambda1_{E}\) 是严格态射，其核或余核有限维，且其指标为 \(n\) 。我们用 \(\mathrm{Sp}_{\omega}(u)\) 表示在 \(\mathrm{Sp}_{\mathrm{e}}(u)\) 中子集 \(\mathrm{Sp}_{n}(u)\) （\(n \in \overline{Z}\) ）之并的补集。

诸集 \(\mathrm{Sp}_{\mathrm{s}}(u)\) 、\(\mathrm{Sp}_{n}(u)\) （\(n \in \overline{Z}\) ）与 \(\mathrm{Sp}_{\omega}(u)\) 构成 u 的谱的一个分划。

E 的每个余核有限维的自同态都是严格的（III, p. 52, 引理 6）。E 的弗雷德霍姆自同态就是 E 的核与余核均有限维的自同态（III, p. 52, 命题 11）。集 \(\mathbf{C} - \mathrm{Sp}_{\mathrm{e}}(u)\) 由那些 \(\lambda \in \mathbf{C}\) （使 \(u - \lambda 1_{\mathrm{E}}\) 是 E 的里斯自同态的）组成，而这样的自同态是 E 的指标为 0 的弗雷德霍姆自同态。

于是，对 \(\lambda \in \mathbf{C}\) 及 \(n\in \mathbf{Z} - \{0\}\) ，我们有：

\(\lambda \in \mathbf{C} - \mathrm{Sp}(u) \iff u - \lambda 1_{\mathrm{E}}\) 是自同构；

\(\lambda \in \mathrm{Sp}_{\mathrm{s}}(u) \quad \Longleftrightarrow \quad u - \lambda 1_{\mathrm{E}} \text{ 是里斯自同态，但不是自同构；}\)

\(\lambda \in \mathrm{Sp}_{0}(u) \quad \Longleftrightarrow \quad u - \lambda 1_{\mathrm{E}}\) 是 E 的指标为 0 的弗雷德霍姆自同态但不是 E 的里斯自同态；

\(\lambda \in \mathrm{Sp}_n(u) \quad \Longleftrightarrow \quad u - \lambda 1_{\mathrm{E}}\) 是 E 的指标为 \(n\) 的弗雷德霍姆自同态（\(n \neq 0\)）；

\[
\lambda \in \mathrm{Sp} _ {- \infty} (u) \quad \Longleftrightarrow \quad u - \lambda 1 _ {\mathrm{E}} \text{ 是严格态射，其核无限维且其余核有限维；}
\]

\(\lambda \in \mathrm{Sp}_{+\infty}(u) \quad \Longleftrightarrow \quad u - \lambda 1_{\mathrm{E}}\) 是严格态射，其核有限维且其余核无限维；

\(\lambda \in \mathrm{Sp}_{\omega}(u) \quad \Longleftrightarrow \quad \text{要么 } u - \lambda 1_{\mathrm{E}} \text{ 不是严格态射，要么其核与余核都是无限维。}\)

注记. --- 用 \(\pi\) 表示从 \(\mathcal{L}(\mathrm{E})\) 到 E 的卡尔金代数 \(\mathcal{C}alk(\mathrm{E})\) 上的典范同态（参见 III, p. 75 的第 4 目）。\(\pi(u)\) 相对于代数 \(\mathcal{C}alk(\mathrm{E})\) 的谱有时称为 \(u\) 的稳定谱。其元素是那些复数 \(\lambda\) （使 \(u - \lambda 1_{\mathrm{E}}\) 不是 E 的弗雷德霍姆自同态的）（III, p. 73 的定理 3 的推论 1）。因此它等于 \(\mathrm{Sp}_{\omega}(u) \cup \mathrm{Sp}_{-\infty}(u) \cup \mathrm{Sp}_{+\infty}(u)\) 。

设 A 是 E 的与 u 交换的自同态所成的集。\(\pi(\mathrm{A})\) 在 \(\mathcal{C}alk(\mathrm{E})\) 中的闭包 B 是一个完备可赋范代数，且 \(\pi(u)\) 相对于 B 的谱就是 u 的本质谱 \(\mathrm{Sp}_{\mathrm{e}}(u)\) （III, p. 77, 命题 3）。把 I, p. 28 的命题 6 的推论应用于 \(\mathcal{C}alk(\mathrm{E})\) 的子代数 B，便知集 \(\mathrm{Sp}_{\mathrm{e}}(u)\) 是 \(\mathrm{Sp}_{\omega}(u) \cup \mathrm{Sp}_{-\infty}(u) \cup \mathrm{Sp}_{+\infty}(u)\) 与其补集的某些有界连通分支的并。

定理 1. --- 设 E 是完备可赋范空间，u 是 E 的一个自同态。

\begin{enumerate}
\def\labelenumi{\alph{enumi})}
\item
  集 \(\mathrm{Sp}_{\omega}(u)\) 是紧的。若 E 无限维，则它非空；
\item
  设 \(n \in \overline{Z}\) 。集 \(\mathrm{Sp}_{n}(u)\) 是 \(\mathbf{C} - \mathrm{Sp}_{\omega}(u)\) 的一族有界连通分支的并。它在 C 中是开的，且它在 C 中的边界包含于 \(\mathrm{Sp}_{\omega}(u)\) 。
\end{enumerate}

集 \(\mathbf{C} - \mathrm{Sp}_{\omega}(u)\) 由那些复数 \(\lambda \in \mathbf{C}\) （使 \(u - \lambda 1_{\mathrm{E}}\) 为核或余核有限维的严格态射的）组成。由 III, p. 67 的命题 11 与 III, p. 70 的命题 13，它是开的。于是集 \(\mathrm{Sp}_{\omega}(u)\) 是闭的。由于它有界，它是紧的。

我们来证明 \(b\) 。设 \(n \in \overline{\mathbf{Z}}\) 。集 \(\mathrm{Sp}_n(u)\) 包含于 \(\mathbf{C} - \mathrm{Sp}_{\omega}(u)\) 。设 U 是 \(\mathbf{C} - \mathrm{Sp}_{\omega}(u)\) 的一个与 \(\mathrm{Sp}_n(u)\) 相交的连通分支。由于映射 \(\lambda \mapsto \operatorname{ind}(u - \lambda 1_{\mathrm{E}})\) （从 \(\mathbf{C} - \mathrm{Sp}_{\omega}(u)\) 到 \(\overline{\mathbf{Z}}\) 中）是局部常数的（III, p. 68 的命题 12 的推论 1 及 III, p. 70 的命题 13 的推论 1），\(u - \lambda 1_{\mathrm{E}}\) 的指标等于 \(n\) （对一切 \(\lambda \in U\) ）。若 \(n \neq 0\) ，这就蕴含 U 包含于 \(\mathrm{Sp}_n(u)\) 。若 \(n = 0\) ，注意到集 U 是 \(\mathbf{C} - (\mathrm{Sp}_{\omega}(u) \cup \mathrm{Sp}_{-\infty}(u) \cup \mathrm{Sp}_{+\infty}(u))\) 的一个连通分支。由于 U 与 \(\mathrm{Sp}_0(u)\) 相交，从而与 \(\mathrm{Sp}_{\mathrm{e}}(u)\) 相交，由注记 2 得出集 U 包含于 \(\mathrm{Sp}_{\mathrm{e}}(u)\) ，从而包含于 \(\mathrm{Sp}_0(u)\) 。在所有情形，我们都得出结论：\(\mathrm{Sp}_n(u)\) 是 \(\mathbf{C} - \mathrm{Sp}_{\omega}(u)\) 的与 \(\mathrm{Sp}_n(u)\) 相交的那些连通分支的并。由于集 \(\mathrm{Sp}(u)\) 有界，这些连通分支必然有界。于是 \(\mathrm{Sp}_n(u)\) 在 \(\mathbf{C}\) 中是开的，且其边界包含于 \(\mathrm{Sp}_{\omega}(u)\) 。这就证明了 \(b\) 。

最后假设集 \(\mathrm{Sp}_{\omega}(u)\) 是空的。由 b)，诸集 \(\mathrm{Sp}_{n}(u)\) （\(n \in \overline{Z}\) ）此时也都是空的。于是我们有 \(\mathrm{Sp}(u) = \mathrm{Sp}_{\mathrm{s}}(u)\) 。因而 u 的谱是离散且紧的，故有限；又由于其所有点都有有限谱重数，向量空间 E 是有限维的（III, p. 83, 命题 2）。这就完成了 a) 的证明。

推论. --- a) 设 \(\Omega\) 是 \(\mathbf{C} - \mathrm{Sp}_{\omega}(u)\) 的无界连通分支。我们有 \(\Omega \cap \mathrm{Sp}(u) \subset \mathrm{Sp}_{\mathrm{s}}(u)\) ；

\begin{enumerate}
\def\labelenumi{\alph{enumi})}
\setcounter{enumi}{1}
\tightlist
\item
  每个附着于 \(\mathrm{Sp}_{\mathrm{s}}(u)\) 但不属于 \(\mathrm{Sp}_{\mathrm{s}}(u)\) 的点都属于 \(\mathrm{Sp}_{\omega}(u)\) 。
\end{enumerate}

断言 \(a\) 是定理 1 的断言 \(b\) 的直接推论。设 \(\lambda\) 是附着于 \(\mathrm{Sp}_{\mathrm{s}}(u)\) 但不属于 \(\mathrm{Sp}_{\mathrm{s}}(u)\) 的一点。它属于 \(u\) 的谱，因为后者是闭的。它不属于任何集 \(\mathrm{Sp}_{n}(u)\) （\(n \in \overline{\mathbf{Z}}\) ），因为这些集是开的（同前）且与 \(\mathrm{Sp}_{\mathrm{s}}(u)\) 不相交。于是我们有 \(\lambda \in \mathrm{Sp}_{\omega}(u)\) ，由此得 \(b\) 。

命题 3. --- 设 E 和 F 是完备可赋范空间，\(u\colon E\to F\) 与 \(v\colon F\to E\) 是连续线性映射。

\begin{enumerate}
\def\labelenumi{\alph{enumi})}
\item
  在 \(\mathbf{C} - \{0\}\) 上，集 \(\mathrm{Sp}(v \circ u)\) 与 \(\mathrm{Sp}(u \circ v)\) （相应地，\(\mathrm{Sp}_{\mathrm{s}}(v \circ u)\) 与 \(\mathrm{Sp}_{\mathrm{s}}(u \circ v)\) ，相应地，\(\mathrm{Sp}_{n}(v \circ u)\) 与 \(\mathrm{Sp}_{n}(u \circ v)\) （\(n \in \overline{\mathbf{Z}}\) ），相应地，\(\mathrm{Sp}_{\omega}(v \circ u)\) 与 \(\mathrm{Sp}_{\omega}(u \circ v)\) ）的迹相等；
\item
  设 \(\lambda\) 是 \(\mathrm{Sp}_{\mathrm{s}}(v \circ u)\) 中异于 0 的一个元素。\(\lambda\) 对 \(v \circ u\) 与对 \(u \circ v\) 的谱重数相等。
\end{enumerate}

设 \(\mu\) 是一个非零复数，且 \(n \in \overline{\mathbf{Z}}\) 。为使 \(\mu 1_{\mathrm{E}} - v \circ u\) 是一个自同构（相应地，一个里斯自同态，相应地，一个核或余核有限维且指标为 \(n\) 的严格态射），必须且只需 \(\mu 1_{\mathrm{F}} - u \circ v\) 也是（III, p. 49, 命题 10）。断言 \(a\) 于是由定义得出。

设 \(\lambda\) 是 \(\mathrm{Sp}_{\mathrm{s}}(v \circ u)\) 中异于 0 的一点。我们有

\[
\dim \operatorname{Ker} ((\lambda 1 _ {\mathrm{E}} - v \circ u) ^ {n}) = \dim \operatorname{Ker} ((\lambda 1 _ {\mathrm{F}} - u \circ v) ^ {n})
\]

（对一切 \(n \geqslant 0\) ；同前），故 \(\lambda\) 对 \(v \circ u\) 与 \(u \circ v\) 的谱重数相等。

\subsection{自同态的转置的谱}

命题 4. --- 设 E 是完备可赋范空间，E' 是 E 的对偶空间，u 是 E 的一个自同态。

\begin{enumerate}
\def\labelenumi{\alph{enumi})}
\item
  有 \(\mathrm{Sp}_{\mathrm{s}}(u) = \mathrm{Sp}_{\mathrm{s}}(^{t}u)\) ，\(\mathrm{Sp}_n(u) = \mathrm{Sp}_{-n}(^t u)\) （对一切 \(n \in \overline{\mathbf{Z}}\) ），且 \(\mathrm{Sp}_{\omega}(u) = \mathrm{Sp}_{\omega}(^t u)\) ；
\item
  \(\mathrm{Sp}_{\mathrm{s}}(u)\) 的每个点对 u 与 \({}^{t}u\) 有相同的谱重数。
\end{enumerate}

I, p. 131 的命题 3 的断言 \(a\) 表明 \(\mathrm{Sp}(u) = \mathrm{Sp}({}^{t}u)\) ，且同处的断言 \(c\) 蕴含 \(\mathrm{Sp}_{\mathrm{s}}(u) = \mathrm{Sp}_{\mathrm{s}}({}^{t}u)\) 以及断言 \(b\) 成立。

对每个 \(\lambda \in \mathbf{C}\) ，III, p. 69 的引理 4 蕴含：\(u - \lambda 1_{\mathrm{E}}\) 是严格态射当且仅当 \({}^{t}u - \lambda 1_{\mathrm{E}'}\) 是严格态射，且

\[
\dim \operatorname{Coker} (^ {t} (u - \lambda 1 _ {\mathrm{E}})) = \dim \operatorname{Ker} (u - \lambda 1 _ {\mathrm{E}})
\]

\[
\dim \operatorname{Ker} (^ {t} (u - \lambda 1 _ {\mathrm{E}})) = \dim \operatorname{Coker} (u - \lambda 1 _ {\mathrm{E}})
\]

（在 \(\overline{Z}\) 中）。由此并注意到谱的各部分的定义（III, p. 85 的定义 3），便得出断言 a)。

\subsection{借紧算子的扰动}

定理 2. --- 设 E 是完备可赋范空间，u 是 E 的一个自同态，h 是 E 的一个紧自同态。则 \(\mathrm{Sp}_{\omega}(u + h) = \mathrm{Sp}_{\omega}(u)\) ，且 \(\mathrm{Sp}_n(u + h) = \mathrm{Sp}_n(u)\) （对一切 \(n \in \overline{\mathbf{Z}} - \{0\}\) ）。

设 \(\lambda \in \mathbf{C}\) 。为使 \(u + h - \lambda 1_{\mathrm{E}}\) 是核（相应地，余核）有限维的严格态射，必须且只需 \(u - \lambda 1_{\mathrm{E}}\) 具有同样的性质，这由 III, p. 72 的定理 1（相应地，III, p. 73 的定理 2）得出。等式

\[
\begin{array}{c} \operatorname{Sp} _ {- \infty} (u + h) = \operatorname{Sp} _ {- \infty} (u), \quad \operatorname{Sp} _ {+ \infty} (u + h) = \operatorname{Sp} _ {+ \infty} (u), \\ \operatorname{Sp} _ {\omega} (u + h) = \operatorname{Sp} _ {\omega} (u) \end{array}
\]

由此得出。此外，我们有 \(\operatorname{Sp}_{n}(u+h)=\operatorname{Sp}_{n}(u)\) （由 III, p. 73 的定理 3，对每个 \(n\in \mathbf{Z}-\{0\}\) ）。

推论. --- 假设 \(\mathrm{Sp}_{\omega}(u)\) 的补集的无界连通分支包含 0。则 \(u + h\) 是 E 的里斯自同态。

有 \(\mathrm{Sp}_{\omega}(u+h)=\mathrm{Sp}_{\omega}(u)\) （由定理 2），故 0 属于 \(\mathbf{C}-\mathrm{Sp}_{\omega}(u+h)\) 的无界连通分支。由 III, p. 87 的定理 1 的推论，或者 0 不属于 \(u+h\) 的谱，或者它是该谱的敏感点。在两种情形下，\(u+h\) 都是 E 的里斯自同态。

\subsection{紧算子的谱}

引理 1. --- 设 E 是 R 上维数 \(\geqslant 2\) 的分离拓扑向量空间，X 是 E 的一个可数子集。补集 E - X 是连通的。

首先假设 E 的维数为 2。我们可以假设 \(E = R^{2}\) 赋有欧几里得范数（EVT, I, p. 14, 定理 2）。由于 X 可数，存在一个实数 \(r \in R_{+}^{*}\) ，使得以 0 为中心、r 为半径的圆 C 不与 X 相交。设 \(x \in E - X\) ；若 \(x \notin C\) ，存在一点 \(y \in C\) ，使得连接 x 与 y 的直线 \(L_{x}\) 不与 X 相交，因为 X 可数。集 E−X 是 C（它是连通的）与诸连通集 \(L_{x} \cup C\) （\(x \in \mathrm{E} - (\mathrm{X} \cup \mathrm{C})\) ）的并；这些集都包含 C，因此 E-X 是连通的（TG, I, p. 81, 命题 2）。

考察一般情形。对元素 \(x \in E - X\) 把 X 换成 X - x，我们化归到 \(0 \notin X\) 的情形。由于 E-X 是诸集 \(F - (X \cap F)\) （当 F 取遍 E 的 2 维子空间时所成的集）的并，且由前面的情形这些集都连通并包含 0，故集 E-X 是连通的（同前）。

引理 2. --- 设 S \(\subset \mathbf{C}\) 是 \(\mathbf{C} - \{0\}\) 中的一个无限的、离散的、有界的且闭的集。则 S 是某个序列 ( \(\lambda_{n}\) ) \(_{n\in N}\) 的值集，该序列由两两不同的非零复数组成，且序列 ( \(|\lambda_{n}|\) ) \(_{n\in N}\) 递减并收敛到 0。

对每个整数 \(i \geqslant 1\) ，集 \(\mathbf{A}_i\) （由复数 \(\lambda \in \mathrm{S}\) （满足 \(|\lambda| \geqslant \frac{1}{i}\) ）组成）在 \(\mathbf{C}\) 中是紧且离散的，故有限。用 \(a_i\) 表示它的基数。由于 S 无限，序列 \((a_i)\) 趋于 \(+\infty\) 。令 \(\mathbf{A}_0 = \varnothing\) ，\(a_0 = 0\) 。对每个 \(i \geqslant 1\) ，选取一个双射 \(n \mapsto \lambda_n\) （从区间 \([a_{i-1}, a_i[\) （\(\mathbf{N}\) 的）到 \(\mathbf{A}_i - \mathbf{A}_{i-1}\) 上），使得映射 \(n \mapsto |\lambda_n|\) 在 \([a_{i-1}, a_i[\) 上递减。序列 \((\lambda_n)_{n \in \mathbf{N}}\) 具有所要求的性质。

设 E 是无限维的完备可赋范空间。代数 \(\mathcal{L}^{\mathrm{c}}(\mathrm{E})\) 是 \(\mathcal{L}(\mathrm{E})\) 的一个不含单位元的子代数。回忆：对每个紧自同态 \(u \in \mathcal{L}^{\mathrm{c}}(\mathrm{E})\) ，谱 \(\mathrm{Sp}'_{\mathcal{L}^{\mathrm{c}}(\mathrm{E})}(u)\) 是 \(u\) 相对于含单位元子代数 \(\mathcal{L}^{\mathrm{c}}(\mathrm{E}) \oplus \mathbf{C1}_{\mathrm{E}}\) （\(\mathcal{L}(\mathrm{E})\) 的）的谱（I, p. 4, 第 4 目）。

命题 5. — 设 E 是完备可赋范空间，u 是 E 的一个紧自同态。\(\mathrm{Sp}_{\mathrm{s}}(u)\) 的每个元素都是有限谱重数的特征值。此外：

\begin{enumerate}
\def\labelenumi{\alph{enumi})}
\item
  若 E 有限维，则 \(\operatorname{Sp}(u) = \operatorname{Sp}_{\mathrm{s}}(u)\) ；
\item
  若 E 无限维，则 \(\mathrm{Sp}_{\mathrm{s}}(u) = \mathrm{Sp}(u) - \{0\}\) 且 \(\mathrm{Sp}_{\omega}(u) = \{0\}\) ；
\item
  若 \(\mathrm{Sp}_{\mathrm{s}}(u)\) 无限，则它是某个序列 \((\lambda_{n})_{n\in\mathbf{N}}\) 的值集，该序列由两两不同的非零复数组成，且序列 \((|\lambda_{n}|)_{n\in\mathbf{N}}\) 递减并收敛到 0；
\item
  若 E 无限维，则 \(\mathrm{Sp}(u) = \mathrm{Sp}'_{\mathcal{L}^c (\mathrm{E})}(u)\) 。
\end{enumerate}

\(\mathrm{Sp}_{\mathrm{s}}(u)\) 的每个元素都是有限谱重数的特征值（III, p. 83 的命题 1）。

断言 \(a\) 是初等的（III, p. 85, 例 1）。现在假设 E 无限维。

设 \(\lambda \in \mathrm{Sp}(u) - \{0\}\) 。则 \(u - \lambda 1_{\mathrm{E}}\) 是 E 的里斯自同态（III, p. 75 的命题 2 的推论 2），故 \(\lambda\) 属于 \(\mathrm{Sp}_{\mathrm{s}}(u)\) 。若 E 无限维，则 \(\mathrm{Sp}_{\omega}(u)\) 非空（III, p. 87, 定理 1, a)）。必然有 \(\mathrm{Sp}_{\omega}(u) = \{0\}\) ，由此得 \(b\) 。

集 \(\mathrm{Sp}_{\mathrm{s}}(u)\) 是离散且有界的。此外，我们有 \(\mathrm{Sp}_{\mathrm{s}}(u) = \mathrm{Sp}(u) \cap (\mathbf{C} - \{0\})\) （依照 \(b\) ），故 \(\mathrm{Sp}_{\mathrm{s}}(u)\) 在 \(\mathbf{C} - \{0\}\) 中是闭的。于是断言 \(c\) 由引理 2 得出。

\(u\) 的谱是可数的，故其在 \(\mathbf{C}\) 中的补集是连通的（引理 1）。把 I, p. 28 的命题 6 的推论应用于含单位元子代数 \(\mathcal{L}^{\mathrm{c}}(\mathrm{E})\oplus \mathbf{C}1_{\mathrm{E}}\) （\(\mathcal{L}(\mathrm{E})\) 的），我们断定 \(\mathrm{Sp}(u) = \mathrm{Sp}'_{\mathcal{L}^{\mathrm{c}}(\mathrm{E})}(u)\) 。

命题 6. --- 设 E 是 C 上的完备可赋范空间，u 是 E 的一个紧自同态。

\begin{enumerate}
\def\labelenumi{\alph{enumi})}
\item
  设 \(f \in \mathcal{O}(\mathrm{Sp}(u))\) 满足 \(f(0) = 0\) 。自同态 \(f(u)\) 是紧的；
\item
  假设 E 是复希尔伯特空间且 u 是正规的。设 f 是 \(\mathrm{Sp}(u)\) 上的一个连续函数，满足 \(f(0)=0\) 。正规自同态 \(f(u)\) 是紧的。
\end{enumerate}

此外，若 E 无限维，则逆命题成立；若 E 有限维，则条件 \(f(0) = 0\) 不是必要的。

我们可以假设 E 无限维。

我们来证明 \(a\) 及其逆。自同态 \(u\) 是巴拿赫代数 \(\mathcal{L}^{\mathrm{c}}(\mathrm{E})\) 的一个元素。由于 E 无限维，我们有等式 \(\mathrm{Sp}(u) = \mathrm{Sp}'_{\mathcal{L}^{\mathrm{c}}(\mathrm{E})}(u)\) （由命题 5 的 \(d\) ）。全纯泛函演算的元素 \(f(u)\) （该演算是把 \(\mathcal{L}^{\mathrm{c}}(\mathrm{E})\) 添加单位元后得到的含单位元巴拿赫代数的全纯泛函演算）属于 \(\mathcal{L}^{\mathrm{c}}(\mathrm{E})\) 当且仅当 \(f(0) = 0\) （I, p. 88）。而且这个元素与元素 \(f(u)\) （\(\mathcal{L}(\mathrm{E})\) 的）重合（I, p. 75 的命题 7），故 \(f(u)\) 是紧的当且仅当 \(f(0) = 0\) 。

断言 \(b\) 及其逆的证明完全相同，只需考察 C*-代数 \(\mathcal{L}^{\mathrm{c}}(\mathrm{E})\) 的连续泛函演算（I, p. 110, 定义 5）。

推论. --- 设 E 和 F 是希尔伯特空间，u 是从 E 到 F 中的一个连续线性映射。线性映射 u 是紧的，当且仅当 E 的自同态 \(|u|\) 是紧的。

设 \((j, |u|)\) 是 u 的极分解（I, p. 140 的定义 4）。由于 \(u = j|u|\) 且 \(|u| = \sqrt{u^{*}u}\) （I, p. 139 的命题 10），该等价性由 III, p. 5 的命题 3 以及前一命题的断言 b) 得出。

\subsection{希尔伯特空间的情形}

在本节中，E 表示 C 上的一个希尔伯特空间。我们用 \(\pi\) 表示从 \(\mathcal{L}(\mathrm{E})\) 到卡尔金代数 \(\mathcal{Calk}(\mathrm{E})\) 上的典范同态。

命题 7. --- 设 \(u \in \mathcal{L}(\mathrm{E})\) 。

\begin{enumerate}
\def\labelenumi{\alph{enumi})}
\item
  若 u 是正规的，则 u 是里斯自同态当且仅当 u 是指标为 0 的弗雷德霍姆自同态。
\item
  若 u 是埃尔米特的，则 u 是弗雷德霍姆自同态当且仅当 u 是里斯自同态。
\end{enumerate}

每个里斯自同态都是指标为 0 的弗雷德霍姆自同态（III, p. 46 的命题 5）。反之，假设 u 是指标为 0 的弗雷德霍姆自同态且 u 是正规的。此时它的幂零空间与它的核重合（EVT, V, p. 43, 命题 8），故有限维。于是由 III, p. 45 的引理 2 以及定义，u 是里斯自同态。

我们来证明 \(b\) 。由 \(a\) ，只需验证埃尔米特弗雷德霍姆自同态 \(u\) 的指标为零。但此时 \(u\) 的像（它是闭的）的正交补等于 \(u\) 的核（EVT, V, p. 41, 命题 4），由此得该断言。

推论. --- 设 \(u \in \mathcal{L}(\mathrm{E})\) 。若 u 是正规的，则 \(\mathrm{Sp}_{0}(u)\) 是空的；若 u 是埃尔米特的，则 \(\mathrm{Sp}_{\mathrm{e}}(u)\) 与 \(\pi(u)\) 相对于代数 \(\mathcal{Calk}(\mathrm{E})\) 的谱重合。

两个断言都由该命题以及诸集 \(\mathrm{Sp}_{n}(u)\) （\(n \in \overline{Z}\) ）与 \(\mathrm{Sp}_{\omega}(u)\) 的定义得出，它们构成 u 的本质谱的一个分划（III, p. 83 的定义 2）。

定理 3（外尔）. --- 设 \(u \in \mathcal{L}(\mathrm{E})\) 是 E 的一个正规自同态。\(u\) 的本质谱是诸集 \(\mathrm{Sp}(u + h)\) （当 \(h\) 取遍 \(\mathcal{L}^{\mathrm{c}}(\mathrm{E})\) 时）的交。

设 \(h \in \mathcal{L}^c(\mathrm{E})\) 。由于 \(\mathrm{Sp}_0(u)\) 是空的（上面的推论），III, p. 89 的定理 2 蕴含 \(\mathrm{Sp_e}(u + h) = \mathrm{Sp_e}(u)\) 。于是诸集 \(\mathrm{Sp}(u + h)\) 的交包含 \(\mathrm{Sp_e}(u)\) 。

设 \(\lambda \in \mathrm{Sp}_{\mathrm{s}}(u)\) 。用 \(\mathrm{E}_{\lambda}\) 表示 \(u\) 的相对于 \(\lambda\) 的特征空间，用 \(\mathrm{F}_{\lambda}\) 表示与 \(u\) 及 \(\mathbf{C} = \{\lambda\}\) 相伴的谱投影的像。空间 E 是 \(\mathrm{E}_{\lambda}\) 与 \(\mathrm{F}_{\lambda}\) 的拓扑直和。设 \(h\) 是 E 的这样一个有限秩自同态：它在 \(\mathrm{F}_{\lambda}\) 上为零，而在 \(\mathrm{E}_{\lambda}\) 上与恒同映射重合。自同态 \(u + h\) 可逆，故 \(\lambda \notin \mathrm{Sp}(u + h)\) 。定理由此得出。

于是，若 u 是 E 的正规自同态，且 \(h \in \mathcal{L}^{c}(\mathrm{E})\) ，则 \(u + h\) 的谱与 u 的谱只能相差一些具有有限谱重数的孤立点。
\subsection{克赖因-鲁特曼定理}

引理 3. --- 设 \((a_n)_{n \geqslant 0}\) 是一个正实数序列，使得幂级数

\[
f (z) = \sum_ {n \geqslant 0} a _ {n} z ^ {n}
\]

具有有限的收敛半径 r \textgreater{} 0。设 D ⊂ C 是以 0 为心、以 r 为半径的开圆盘。不存在定义在 C 中 r 的某个开邻域 U 上且在 U ∩ D 上与 f 一致的全纯函数 \(\widetilde{f}\)。

假设这样一个全纯函数 \(\tilde{f}\)（定义在 \(r\) 的一个开邻域 U 上）存在。存在实数 \(s < r\) 和 \(\delta > 0\)，使得 \(s + \delta > r\)，且使得以 \(s\) 为心、以 \(\delta\) 为半径的开圆盘 D' 包含于 U 中。\(\tilde{f}\) 在点 \(s\) 处的幂级数展开（VAR, R1, p. 26, 3.2.1）于是对每个 \(z\)（只要它落在以 0 为心、以 \(\delta\) 为半径的圆盘中）都收敛（VAR, R1, p. 29, 3.3.4）。此级数为

\[
\widetilde {f} _ {s} (z) = \sum_ {n = 0} ^ {+ \infty} \frac {\widetilde {f} ^ {(n)} (s)}{n !} z ^ {n} = \sum_ {n = 0} ^ {+ \infty} \frac {f ^ {(n)} (s)}{n !} z ^ {n}
\]

（VAR, R1, p. 27, 3.2.4）。由于 \(s \in D\) ，我们有

\[
f ^ {(n)} (s) = \sum_ {k = n} ^ {+ \infty} k (k - 1) \dots (k - n + 1) a _ {k} s ^ {k - n}
\]

（VAR, R1, p. 28, 3.2.11）。取 \(z\) 使得 \(0 < z < \delta\) 。由于 \(a_{k} \geqslant 0\) ，我们有

\[
\begin{array}{r l} & {\widetilde {f} _ {s} (z) = \sum_ {n = 0} ^ {+ \infty} \biggl (\sum_ {k = n} ^ {+ \infty} \frac {k (k - 1) \cdots (k - n + 1)}{n !} a _ {k} s ^ {k - n} \biggr) z ^ {n}} \\ & {\qquad = \sum_ {k = 0} ^ {+ \infty} \biggl (a _ {k} \sum_ {n = 0} ^ {k} \frac {k !}{n ! (k - n) !} s ^ {k - n} z ^ {n} \biggr) = \sum_ {k = 0} ^ {+ \infty} a _ {k} (s + z) ^ {k}} \end{array}
\]

（TG, III, p. 40）。当 \(s + z > r\) 时，这个表达式的收敛性与幂级数 \(f\) 的收敛半径等于 \(r\) 这一假设相矛盾。

设 E 是一个实巴拿赫空间。设 C 是 E 中的一个尖凸锥。由 C 生成的向量空间等于 C -- C（EVT, II, p. 12, 推论 1）。特别地，锥 C 在 E 中是全的，当且仅当 C -- C 在 E 中稠密。回顾：当 C ∩ (−C) 仅由 0 组成时，称 C 是突出的（EVT, II, p. 11）。

C 的极集 C° 是由连续线性形式 \(\ell \in E'\) 组成的集合，其中 \(\ell(x) \geqslant 0\) 对所有 \(x \in C\) 成立（EVT, II, p. 47, 命题 4）。由双极定理（EVT, II, p. 48, 定理 1），若 C 是闭尖凸锥，则 C°° = C。

引理 4. --- 设 E 是一个实巴拿赫空间，\(u \in \mathcal{L}(\mathrm{E}_{(\mathbf{C})})\) 是 \(\mathrm{E}_{(\mathbf{C})}\) 的一个非零自同态。设 C 是 E 中的一个全凸锥。则存在 \(x \in \mathrm{C}\) 使得 \(u(x) \neq 0\) 。

事实上，若 \(u\) 在 C 上为零，则 \(u\) 在 C - C 上为零，从而在 E 上为零，因此在 \(\mathrm{E}_{(\mathbf{C})}\) 上为零。

定理 4（克赖因–鲁特曼）. --- 设 E 是 R 上的一个完备可赋范空间。设 C 是 E 中一个闭的全突出凸锥，并设 \(u \in \mathcal{L}(\mathrm{E})\) 是满足 \(u(\mathrm{C}) \subset \mathrm{C}\) 的紧线性映射。若谱半径 \(\varrho(u)\) \textgreater0，则 \(\varrho(u)\) 是 u 的谱的孤立点，并且存在 u 的非零特征向量 \(x \in C\)，其特征值为 \(\varrho(u)\) 。

设 \(E_{(\mathbf{C})}\) 是 E 的复化，设 \(u_{(\mathbf{C})}\) 是 \(E_{(\mathbf{C})}\) 的由 u 经纯量扩张得到的自同态。\(u_{(\mathbf{C})}\) 的谱半径等于 \(\varrho(u)\)（I, p. 86）；以下简记为 \(\varrho\) 。由于 \(\varrho > 0\) ，\(u\) 的复谱不只由 0 组成。因此存在 \(\lambda_0 \in \mathrm{Sp}_\mathrm{s}(u_{(\mathbf{C})})\) 使得 \(|\lambda_0| = \varrho\)（III, p. 90 的命题 5）。设 \(e_0\) 是 \(u_{(\mathbf{C})}\) 的与 \(\lambda_0\) 相伴的谱投影。

预解式 \(\lambda \mapsto \mathrm{R}(u_{(\mathbf{C})},\lambda) = (\lambda -u_{(\mathbf{C})})^{-1}\) 在 \(u_{(\mathbf{C})}\) 的谱的余集上是全纯的（I, p. 24 的定理 1）。复数 \(\lambda_0\) 是预解式的一个极点，其留数是谱投影 \(e_0\)（I, p. 131）。

设 \(y\) 和 \(y'\) 是 E 中使得 \(y + iy' \in \mathrm{E}_{(\mathbf{C})}\) 是 \(u_{(\mathbf{C})}\) 的属于特征值 \(\lambda_0\) 的特征向量的元素。由于 \(e_0(y + iy') = y + iy' \neq 0\) ，存在元素 \(x_0 \in \mathrm{C}\) 使得 \(e_0(x_0) \neq 0\)（引理 4）。由于 C 是闭的且是尖的，其极集 \(\mathrm{C}^\circ\) 是全的（EVT, II, p. 48, 推论 1），因此存在线性形式 \(\ell_0 \in \mathrm{C}^\circ\) 使得 \(\langle e_0(x_0), \ell_0 \rangle \neq 0\) 。

考虑函数 \(f\)，它由 \(f(0)=0\) 以及

\[
f (z) = \langle \mathrm{R} (u _ {(\mathbf {C})}, z ^ {- 1}) x _ {0}, \ell_ {0} \rangle
\]

（其中 \(z \in \mathbf{C}\) 满足 \(z^{-1} \notin \mathrm{Sp}(u_{(\mathbf{C})})\)）定义。这个函数满足

\[
f (z) = \ell_ {0} \bigg (\Bigl (\sum_ {n = 0} ^ {+ \infty} z ^ {n + 1} u _ {(\mathbf {C})} ^ {n} \Bigr) x _ {0} \bigg) = \sum_ {n = 0} ^ {\infty} \langle u ^ {n} (x _ {0}), \ell_ {0} \rangle z ^ {n + 1}\tag{1}
\]

对 \(|z| < 1 / \varrho\) 成立（I, p. 24, 定理 1, d)），因而在以 0 为心、以 \(1 / \varrho\) 为半径的圆盘内是全纯的。

存在一个全纯函数 \(\widetilde{\mathrm{R}}\)，它定义在 \(\lambda_0\) 的一个开邻域 U 上，取值于 \(\mathcal{L}(\mathrm{E})\) 中，使得只要 \(z\) 属于 \(\mathrm{U} - \{\lambda_0\}\) ，就有

\[
\mathrm{R} (u _ {(\mathbf {C})}, z) = \widetilde {\mathrm{R}} (z) + \sum_ {n = 0} ^ {+ \infty} (z - \lambda_ {0}) ^ {- n - 1} (u _ {(\mathbf {C})} - \lambda_ {0}) ^ {n} e _ {0}
\]

（I, p. 83 的命题 17）。于是，对 \(z\) 满足 \(z^{-1} \in \mathrm{U}\) 且 \(z \neq 1 / \lambda_0\) 的情形，有

\[
f (z) = \langle \widetilde {\mathrm{R}} (z ^ {- 1}) x _ {0}, \ell_ {0} \rangle + \sum_ {n = 0} ^ {+ \infty} (z ^ {- 1} - \lambda_ {0}) ^ {- n - 1} \langle (u _ {(\mathbf {C})} - \lambda_ {0}) ^ {n} e _ {0} (x _ {0}), \ell_ {0} \rangle .
\]

该级数中对应于 \(n=0\) 的项是 \((z^{-1}-\lambda_{0})^{-1}\langle e_{0}(x_{0}),\ell_{0}\rangle\) 。由于 \(\langle e_{0}(x_{0}),\ell_{0}\rangle\neq0\) ，洛朗展开的唯一性（VAR, R1, p. 30, 3.3.9）表明 \(f\) 不能延拓为 \(1/\lambda_{0}\) 的一个邻域上的全纯函数。特别地，函数 \(f\) 在点 0 处的幂级数展开 (1) 的收敛半径等于 \(1/\varrho\) 。

幂级数 (1) 的系数是 \(\langle u^n(x_0),\ell_0\rangle\geqslant 0\)，因为 \(u(\mathrm{C})\subset\mathrm{C}\) 且 \(\ell_{0}\in\mathrm{C}^{\circ}\) 。由引理 3，函数 \(f\) 不能延拓为 \(1/\varrho\) 的一个邻域上的全纯函数。因此 \(u_{(\mathbf{C})}\) 的预解式不能延拓为 \(\varrho\) 的一个邻域上的全纯函数，即 \(\varrho\in\mathrm{Sp}(u)\) 。这蕴含 \(\varrho\) 是 \(u_{(\mathbf{C})}\) 的特征值（III, p. 90 的命题 5）。由于 \(\varrho\) 是实数，它也是 \(u\) 的特征值。

设 \(d \geqslant 1\) 是 \(u_{(\mathbf{C})}\) 的预解式在 \(\varrho\) 处极点的阶。设 \(e\) 是自同态

\[
e = \lim _ {z \to \varrho} (z - \varrho) ^ {d} \mathrm{R} (u _ {(\mathbf {C})}, z).
\]

它是非零的且与 \(u_{(\mathbf{C})}\) 交换，其像包含在 \(u_{(\mathbf{C})}\) 对应于 \(\varrho\) 的特征空间中（参见 I, p. 131, no. 3）。现在设 \(\ell\) 是 C° 的一个元素，\(x\) 是 C 的一个元素。由 I, p. 24, 定理 1, d)，我们有

\[
\langle e(x),\ell \rangle = \lim_{\substack{z\to \varrho \\ z > \varrho}}(z - \varrho)^{d}\sum_{n\geqslant 0}\langle u^{n}(x),\ell \rangle z^{-n - 1}\geqslant 0.
\]

双极定理（EVT, II, p. 48, 定理 1）表明 \(e(x) \in \mathrm{C}\) 对所有 \(x \in \mathrm{C}\) 成立。由于 \(\mathrm{C}\) 是全的，存在 \(x \in \mathrm{C}\) 使得 \(e(x) \neq 0\)（引理 4），于是 \(e(x)\) 属于 \(\mathrm{C}\) 且是 \(u\) 的属于特征值 \(\varrho\) 的特征向量，这正是所要证明的。

推论. --- 设 E 是 R 上的一个完备可赋范空间。设 C 是 E 中一个闭的、全的突出凸锥，并设 \(u \in \mathcal{L}(E)\) 是满足 \(u(\mathrm{C}) \subset \mathrm{C}\) 的紧线性映射。若 u 的谱半径 \(\varrho(u)\) \textgreater{} 0，则 \(\varrho(^{t}u) = \varrho(u)\) 是 \(^{t}u\) 的特征值，并且 \(^{t}u\) 有一个属于 \(\varrho(u)\) 的特征向量位于 \(C^{\circ}\) 中。

我们有 \(\varrho(^{t}u) = \varrho(u)\)（I, p. 131 的命题 3）。由 III, p. 9 的推论 1，E' 的自同态 \(^{t}u\) 是紧的。此外，我们有 \(^{t}u(C^{\circ}) \subset C^{\circ}\) ；于是结论由应用于 \(^{t}u\) 以及闭凸锥 \(C^{\circ}\) 的克赖因-鲁特曼定理得出，因为后者是突出的（由于 C 是全的）且是全的（由于 C 是突出的）。

注. --- 在克赖因–鲁特曼定理中，假设 \(\varrho(u) > 0\) 一般不能省略。例如，设 V 是巴拿赫空间 \(\mathcal{C}([0,1])\) 上由

\[
\mathrm{V} (f) (x) = \int_ {0} ^ {x} f (y) d y
\]

（其中 \(f \in \mathcal{C}([0,1])\)）所定义的自同态。映射 V 是紧的，且其谱半径为零（I, p. 187 的习题 1）。它保持 \(\mathcal{C}([0,1])\) 中由正函数组成的闭的、全的、突出的凸锥，并且没有特征值（出处同上）。

下面的命题刻画了保持一个锥的自同态具有严格正谱半径的一个充分条件，并进而细化了克赖因-鲁特曼定理。

命题 8. --- 设 E 是 R 上的一个非零完备可赋范空间。设 C 是 E 中内部 \(\mathring{\mathrm{C}}\) 非空的闭突出凸锥，并设 \(u \in \mathcal{L}(\mathrm{E})\) 是满足 \(u(\mathrm{C} - \{0\}) \subset \mathring{\mathrm{C}}\) 的紧线性映射。

\begin{enumerate}
\def\labelenumi{\alph{enumi})}
\item
  有 \(\varrho(u) > 0\)，并且存在一个向量 \(x_0\)，它是 \(u\) 的位于 \(\mathring{C}\) 中、属于特征值 \(\varrho(u)\) 的特征向量；
\item
  u 的对应于 \(\varrho(u)\) 的谱投影的秩为 1；
\item
  设 \(\lambda \neq \varrho(u)\) 是 \(u_{(\mathbf{C})}\) 的任意一个特征值，则 \(|\lambda| < \varrho(u)\)；
\item
  设 F 是 E 的在 u 下稳定且使 \((\mathrm{C}-\{0\})\cap\mathrm{F}\) 非空的子空间。则 \(x_{0}\in F\) 。特别地，u 在 C 中的特征向量只有 \(x_{0}\) 的倍向量。
\end{enumerate}

我们先证明两个预备引理。

引理 5. --- 设 E 是实巴拿赫空间，C 是 E 中的一个凸锥，\(u \in \mathcal{L}(\mathrm{E})\) 满足 \(u(\mathrm{C} - \{0\}) \subset \mathring{\mathrm{C}}\) 。设 \(\ell \in \mathrm{C}^{\circ}\) 是 \(^{t}u\) 的一个特征向量。则 \(\ell\) 的核在 \(u\) 下稳定且与 \(\mathrm{C} - \{0\}\) 不相交。

设 \(\lambda \in \mathbf{R}\) 满足 \(^t u(\ell) = \lambda \ell\) 。对每个 \(x \in \mathrm{Ker}(\ell)\)，我们有

\[
\langle u (x), \ell \rangle = \langle x, ^ {t} u (\ell) \rangle = \lambda \langle x, \ell \rangle = 0,
\]

因而 \(\operatorname{Ker}(\ell)\) 在 \(u\) 下稳定。

设 \(x\) 是 \(\operatorname{Ker}(\ell)\) 的一个非零元素。对 E 的每个元素 \(y\)，只要 \(\langle y, \ell \rangle < 0\)，就有 \(\langle u(x) + y, \ell \rangle < 0\)，由此推出 \(u(x) + y \notin \mathrm{C}\)，因为 \(\ell \in \mathrm{C}^\circ\) 。我们推得 \(u(x)\) 不属于 \(\mathring{\mathrm{C}}\) 。由于 \(u(\mathrm{C} - \{0\}) \subset \mathring{\mathrm{C}}\) ，这蕴含 \(x \notin \mathrm{C}\) 。

引理 6. --- 设 E 是有限维实向量空间，B 是 E 中 0 的一个紧凸邻域。设 u 是 E 的满足 \(u(B) \subset \mathring{B}\) 的自同态。则 u 的复谱包含于 C 的单位圆盘，特别地，它与 C 中以 0 为心、以 1 为半径的圆不相交。

将 B 换成 \(B \cap (-B)\) ，可化归为 B 是均衡集的情形。此时 B 的度规 p（EVT, II, p. 28, 习题 3）是 E 上的一个范数，并且它给出 E 的拓扑（EVT, I, p. 14, 定理 2）。假设蕴含 \(p(u(x)) < p(x)\) 对 E - \{0\} 的每个元素 x 成立，从而 u 的谱半径 \textless{} 1；由于它也是 \(u_{(\mathbf{C})}\) 的谱半径（参见 I, p. 86），结论得证。

现在来证明该命题。

用 \(x \preccurlyeq y\) 表示 E 上与凸锥 C 相伴的序关系（EVT, II, p. 13, n° 5），即 \(x \preccurlyeq y\) 当且仅当 \(y - x \in \mathrm{C}\) 。我们有 \(u(x) \preccurlyeq u(y)\)，只要 \(x \preccurlyeq y\)。由于 \(\mathring{\mathrm{C}}\) 非空，由 C 生成的向量空间 C -- C（EVT, II, p. 12, 推论 1）包含 0 的一个邻域，因此锥 C 是全的。

我们来证明断言 \(a\) )。设 \(\varrho = \varrho(u)\) 。设 \(y_0\) 是 \(\mathring{\mathrm{C}}\) 的一个元素。我们有 \(y_0 \neq 0\) 。取 \(r > 0\) 使得以 \(y_0\) 为心、以 \(r\) 为半径的闭球包含于 C。对每个 \(y \in \mathrm{E} - \{0\}\) ，有 \(y_0 - r \| y\|^{-1}y \in \mathrm{C}\) ，因而 \(y \preccurlyeq r^{-1}\| y\| y_0\) 对所有 \(y \in \mathrm{E}\) 成立。

由于 \(y_{0} \neq 0\)，假设蕴含 \(u(y_{0}) \in \mathring{\mathrm{C}}\) 。因此存在 t \textgreater{} 0 使得 \(tu(y_{0}) - y_{0} \in \mathrm{C}\) 。令 v = tu。这是 E 的一个紧自同态，满足 \(v(\mathrm{C}) \subset \mathrm{C}\) 且 \(v(y_{0}) \succcurlyeq y_{0}\) 。对每个整数 \(n \geqslant 1\)，我们有

\[
y _ {0} \preccurlyeq v ^ {n} (y _ {0}) \preccurlyeq r ^ {- 1} \| v ^ {n} (y _ {0}) \| y _ {0} \preccurlyeq r ^ {- 1} \| v ^ {n} \| \| y _ {0} \| y _ {0},
\]

因而 \((r^{-1}\|v^{n}\|\|y_{0}\|-1)y_{0}\in\mathrm{C}\) 。由于 C 是突出的，这蕴含 \(t^{n}\|u^{n}\|=\|v^{n}\|\geqslant r/\|y_{0}\|\) ，从而 \(\varrho\geqslant t^{-1}>0\)（I, p. 20, 命题 1）。

由克赖因–鲁特曼定理（定理 4），实数 \(\varrho\) 是 \(u\) 的特征值，且存在一个向量 \(x_0\)，它是 \(u\) 的位于 C 中、属于特征值 \(\varrho\) 的特征向量。由假设有 \(\varrho x_0 = u(x_0) \in \mathring{\mathrm{C}}\)，因而 \(x_0 \in \mathring{\mathrm{C}}\) 。这就证明了断言 \(a\) )。

设 K 是 \(u_{(\mathbf{C})}\) 的谱与 C 中以 0 为心、以 \(\varrho\) 为半径的圆的交。由于 u 是紧的且 \(\varrho > 0\)，集合 K 是有限的，并且谱投影 \(e_{K}\) 的像是 \(\mathrm{E}_{(\mathbf{C})}\) 的一个有限维子空间（III, p. 83 的命题 2 与 III, p. 90 的命题 5）。由于 K 在复共轭下稳定，\(e_{K}\) 的像是 \(\mathrm{E}_{(\mathbf{C})}\) 中在 R 上有理的有限维子空间（A, V, p.60, 命题 6）；用 F 表示 E 中使 \(\mathrm{F}_{(\mathbf{C})}\) 等于 \(e_{K}\) 的像的子空间。空间 F 在 u 下稳定且包含 u 对应于 \(\varrho\) 的特征空间，因此 F 非零。

为证明断言 \(b)\) 与 \(c),\)，只需证明 \(F\) 是一维的。

用 \(v\) 表示由 \(u\) 通过过渡到子空间得到的 F 的自同态。设 \(C_F = C \cap F\) 。这是 F 中内部非空的闭突出凸锥（因为 \(x_0 \in \mathring{C} \cap F\)），因而是全的；它满足 \(v(C_F - \{0\}) \subset \mathring{C}_F\)，特别地在 \(v\) 下稳定。由于 \(\varrho(v) \leqslant \varrho\) 且 \(x_0 \in F\)，我们有 \(\varrho(v) = \varrho\)。由定理 4 的推论（应用于 \(C_F\) 和 \(v\)），存在一个向量 \(\ell\)，它是 \(^t v\) 的位于 \(C_F^\circ\) 中、属于特征值 \(\varrho(v) = \varrho > 0\) 的特征向量。

F 的子空间 \(\mathrm{Ker}(\ell)\) 在 v 下稳定。设 w 是 \(\mathrm{Ker}(\ell)\) 的、由 \(\varrho^{-1}v\) 通过过渡到商子空间而得到的自同态；它的谱包含于单位圆。集合 B 由这样的 \(y \in \mathrm{Ker}(\ell)\) 组成：\(x_{0} + y \in C_{F}\)；它是 \(\mathrm{Ker}(\ell)\) 中 0 的一个闭凸邻域。对每个 \(y \in B\) 和每个 \(z \in \mathrm{Ker}(\ell)\)，我们有

\[
x _ {0} + (w (y) + z) = \varrho^ {- 1} (v (x _ {0} + y) + \varrho z)
\]

当 z 的范数充分小时它属于 \(C_{F}\)，因而 \(w(y) \in \mathring{\mathrm{B}}\) 。集合 B 是有界的：事实上，若存在一个序列 \((y_{n})_{n \in \mathbf{N}}\)，它含于 \(\operatorname{Ker}(\ell)\)，使得 \(\|y_{n}\| \to +\infty\) 且 \(x_{0} + y_{n} \in C_{F}\)，那么在过渡到子列后，将有 \(y_{n}/\|y_{n}\| \to y\)，其中 \(y \in \operatorname{Ker}(\ell)\) 非零，而 \(y = \lim \|y_{n}\|^{-1}(x_{0} + y_{n})\) 将属于 \(C_{F}\)，这与引理 5 矛盾。因此集合 B 是紧的。

于是，由应用于 \(\mathrm{Ker}(\ell)\)、B 和 \(w\) 的引理 6 推出，\(w\) 的复谱与以 0 为心、以 1 为半径的圆不相交；这蕴含 \(w\) 的谱是空的，即 \(\mathrm{Ker}(\ell)\) 化为 \(\{0\}\)，从而 F 的维数为 1。因此断言 \(b\) 和 \(c\) 得证。

线性映射 \({}^t u\) 是紧的（III, p. 9 的推论 1），且 \({}^t u(\mathrm{C}^\circ) \subset \mathrm{C}^\circ\)；此外 \(\varrho(^t u) = \varrho > 0\)（I, p. 131 的命题 3）。由定理 4 的推论，在 \(\mathrm{C}^\circ\) 中存在一个向量 \(\ell \neq 0\)，它是 \({}^t u\) 的属于特征值 \(\varrho\) 的特征向量。\(\ell\) 的核在 \(u\) 下稳定且与 \(\mathrm{C} - \{0\}\) 不相交（引理 5）。特别地，有 \(x_0 \notin \operatorname{Ker}(\ell)\) 。

设 F 是 E 的一个在 u 下稳定的子空间。我们有分解 \(\mathrm{F} = (\mathrm{F} \cap \mathbf{R}x_{0}) \oplus (\mathrm{F} \cap \operatorname{Ker}(\ell))\) 。若 F 不包含 \(x_{0}\)，则 \(\mathrm{F} \subset \operatorname{Ker}(\ell)\)，从而 F 与 \(C - \{0\}\) 不相交。这就证明了 d)，命题证毕。

推论. --- 设 E 是 R 上的一个非零完备可赋范空间。设 C 是 E 中内部 \(\dot{C}\) 非空的闭凸突出锥，并设 u 是 E 的满足 \(u(\mathrm{C}) \subset \mathrm{C}\) 的紧自同态。假设存在整数 \(k \geqslant 1\) 使得 \(u^{k}(\mathrm{C} - \{0\}) \subset \mathring{\mathrm{C}}\) 。

\begin{enumerate}
\def\labelenumi{\alph{enumi})}
\item
  有 \(\varrho(u) > 0\)，并且存在一个向量 \(x_0\)，它是 \(u\) 的位于 \(\dot{C}\) 中、属于特征值 \(\varrho(u)\) 的特征向量；
\item
  u 的对应于 \(\varrho(u)\) 的谱投影的秩为 1；
\item
  设 \(\lambda \neq \varrho(u)\) 是 \(u_{(\mathbf{C})}\) 的任意一个特征值，则 \(|\lambda| < \varrho(u)\)；
\item
  u 在 C 中的特征向量只有 \(x_0\) 的倍向量。
\end{enumerate}

设 \(\varrho = \varrho(u)\) 。可以把定理 4 应用于 \(u\)，把命题 8 应用于 \(u^k\) 。因此存在一个向量 \(x_0\)，它是 \(u\) 的位于 C 中、属于特征值 \(\varrho\) 的特征向量。由于 \(0 < \varrho(u^k) = \varrho^k\)（I, p. 21 的公式 (4)），特别地有 \(\varrho > 0\)；又由于 \(x_0\) 是 \(u^k\) 的属于特征值 \(\varrho(u^k)\) 的特征向量，我们有 \(x_0 \in \mathring{\mathrm{C}}\) 。

我们有 \(\mathrm{Sp}(u_{(\mathbf{C})}^{k})=\mathrm{Sp}(u_{(\mathbf{C})})^{k}\)（I, p. 2 的注 4），因此每个特征值 \(\lambda\in C\)（属于 \(u_{(\mathbf{C})}\)），只要 \(|\lambda|=\varrho\)，就满足 \(\lambda^{k}=\varrho^{k}\)；由于对应于 \(\lambda\) 的每个特征向量都是 \(u^{k}\) 的属于 \(\varrho^{k}\) 的特征向量，从而与 \(x_{0}\) 成比例，故有 \(\lambda=\varrho\) 。

若 u 的对应于特征值 \(\varrho\) 的谱投影的秩至少为 2，则 \(u^{k}\) 的对应于特征值 \(\varrho^{k}\) 的谱投影的秩也至少为 2（参见 I, p. 129, no. 2）。

最后，若 \(x \in C\) 是 u 的一个特征向量，则它也是 \(u^{k}\) 的特征向量，从而与 \(x_{0}\) 成比例。

设 n 是一个 \(\geqslant 1\) 的整数。集合 \(R_{+}^{n}\) 是 \(R^{n}\) 中的闭凸尖突出锥，其非空内部等于 \((\mathbf{R}_{+}^{*})^{n}\) 。

设 \(A = (a_{i,j})\) 是一个 \((n,n)\) 型实矩阵，\(u_A\) 是由 \(x \mapsto Ax\) 给出的 \(\mathbf{R}^n\) 的自同态。设 \(\varrho = \varrho(u_A)\) 是它的谱半径。

引理 7. --- 有 \(u_A(\mathbf{R}_+^n) \subset \mathbf{R}_+^n\)，当且仅当对所有 i 和 j 有 \(a_{i,j} \geqslant 0\)；有 \(u_A(\mathbf{R}_+^n - \{0\}) \subset (\mathbf{R}_+^*)^n\)，当且仅当对所有 i 和 j 有 \(a_{i,j} > 0\)。

设 \((e_{1},\ldots,e_{n})\) 是 \(R^{n}\) 的典范基。向量 \(e_{i}\) 属于 \(R_{+}^{n}\)，且对所有 i 有 \(u_{A}(e_{i})\in\mathbf{R}_{+}^{n}\)（相应地，对所有 i 有 \(u_{A}(e_{i})\in(\mathbf{R}_{+}^{*})^{n}\)），当且仅当对所有 i 和 j 有 \(a_{i,j}\geqslant0\)（相应地，对所有 i 和 j 有 \(a_{i,j}>0\)）。由线性性即得结论。

定理 5（佩龙–弗罗贝尼乌斯）. --- a) 若 \(a_{i,j} \geqslant 0\)，其中 \(i\) 和 \(j\) 取遍 \(\{1,\ldots,n\}\)，则实数 \(\varrho\) 是 \(A\) 的特征值；b) 若 \(a_{i,j} > 0\)，其中 \(i\) 和 \(j\) 取遍 \(\{1,\ldots,n\}\)，则 \(\varrho > 0\)，且 \(u_A\) 对应于 \(\varrho\) 的准素空间，即 \(A - \varrho 1_{\mathbf{R}^n}\) 的幂零空间，维数为 1，并由一个向量 \(x_0 \in (\mathbf{R}_+^*)^n\) 生成。\(A\) 没有其他在 \(\mathbf{R}_+^n\) 中有特征向量的特征值，并且所有复特征值 \(\lambda \neq \varrho\)（\(A\) 的）都满足 \(|\lambda| < \varrho\) 。

若 \(\varrho = 0\)，断言 \(a\) ) 是初等的，因为此时 0 是 \(A\) 的特征值。若 \(\varrho > 0\)，则由引理 7，断言 \(a\) ) 可由定理 4 得出。出于同样的理由，断言 \(b\) ) 是命题 8 的推论。

注. --- 假设存在整数 \(k \geqslant 1\) 使得 \(A^{k}\) 的所有系数都 \textgreater0（这样的矩阵有时称为本原矩阵）。那么，由命题 8 的推论，谱半径 \(\varrho\) 是 A 的特征值，\(u_{A}\) 对应于 \(\varrho\) 的准素空间维数为 1，且由一个向量 \(x_{0} \in (\mathbf{R}_{+}^{*})^{n}\) 生成。此外，A 没有其他在 \(R_{+}^{n}\) 中有特征向量的复特征值，并且 A 的所有复特征值 \(\lambda \neq \varrho\) 满足 \(|\lambda| < \varrho\) 。

\specialchapter{习题}

\exercisesection{}

仅在习题 1 至 5 中，K 表示一个完备的非阿基米德、非离散的赋值域，其赋值记为 \(x \mapsto |x|\) 。我们还用 \(\mathrm{G} \subset \mathbf{R}_{+}^{*}\) 表示 \(\mathrm{K}^{*}\) 在映射 \(x \mapsto |x|\) 下的像。所考虑的巴拿赫空间都是 \(\mathrm{K}\) 上的巴拿赫空间。

\begin{enumerate}
\def\labelenumi{\arabic{enumi})}
\tightlist
\item
  \begin{enumerate}
  \def\labelenumii{\alph{enumii})}
  \tightlist
  \item
    设 E 是 K 上的一个巴拿赫空间。E 上存在一个范数 \(x \mapsto \|x\|\)，它给出 E 的拓扑，且使得 \(\|x\| \in \overline{G}\) 对所有 \(x \in E\) 成立。这样的范数称为圆范数。若 K 的赋值不是离散的，则每个给出 E 的拓扑的范数都是圆范数。
  \end{enumerate}
\end{enumerate}

\begin{enumerate}
\def\labelenumi{\alph{enumi})}
\setcounter{enumi}{1}
\item
  设 E 和 F 是 K 上的巴拿赫空间。我们在空间 \(\mathcal{L}(\mathrm{E};\mathrm{F})\) 上赋予范数 \(\| u\| = \sup_{x\in \mathrm{E}- \{0\}}\| u(x)\| /\| x\|\) 。空间 \(\mathcal{L}(\mathrm{E};\mathrm{F})\) 是 K 上的巴拿赫空间。若 E 上的范数是圆范数，则 \(\| u\| = \sup_{\| x\| \leqslant 1}\| u(x)\|\) 。若 F 上的范数也是圆范数，则 \(\mathcal{L}(\mathrm{E};\mathrm{F})\) 上的范数是圆范数。我们记 \(\mathrm{E}' = \mathcal{L}(\mathrm{E};\mathrm{K})\)，并称 \(\mathrm{E}'\) 为 E 的对偶空间。
\item
  设 I 是一个集合，F 是 K 上的一个巴拿赫空间，\(\mathcal{C}_{0}(\mathrm{I},\mathrm{F})\) 是由函数 \(f\colon I\to F\) 组成的 K-向量空间，其中 \(\|f(i)\|\) 沿 I 的有限子集的余集滤子趋于 0。赋以范数 \(\|f\|=\sup_{i\in I}\|f(i)\|\) 后，空间 \(\mathcal{C}_{0}(\mathrm{I},\mathrm{F})\) 是 K 上的巴拿赫空间，且若 F 的范数是圆范数，则它的范数也是圆范数。我们记 \(\mathcal{C}_{0}(\mathrm{I})=\mathcal{C}_{0}(\mathrm{I},\mathrm{K})\) 。
\end{enumerate}

\(d\) ) 设 I 是一个集合。对每个 \(i\in \mathrm{I}\) ，用 \(\varphi_{i}\in \mathcal{C}_{0}(\mathrm{I})\) 表示 \(\{i\}\) 的特征函数。设 F 是 K 上的一个巴拿赫空间。映射 \(u \mapsto (u(\varphi_i))_{i \in \mathrm{I}}\) 是从巴拿赫空间 \(\mathcal{L}(\mathcal{C}_0(\mathrm{I});\mathrm{F})\) 到巴拿赫空间 \(\mathcal{B}(\mathrm{I};\mathrm{K})\) 的同构（EVT, I, p. 4, 例；参见 EVT, I, p. 23, 习题 7）。

\begin{enumerate}
\def\labelenumi{\alph{enumi})}
\setcounter{enumi}{4}
\item
  假设 K 的赋值是离散的。设 E 是 K 上的范数为圆范数的巴拿赫空间。存在一个集合 I，使得 E 等距同构于巴拿赫空间 \(\mathcal{C}_{0}(\mathrm{I})\)（参见 EVT, I, p. 23, 习题 7）。
\item
  假设 K 的赋值是离散的。设 E 是赋有圆范数的 K 上的巴拿赫空间，F 是 E 的一个闭向量子空间。存在 E 的一个连续投影，其像为 F 且范数 \(\leqslant 1\) 。
\end{enumerate}

\begin{enumerate}
\def\labelenumi{\arabic{enumi})}
\setcounter{enumi}{1}
\tightlist
\item
  设 E 和 F 是 K 上的巴拿赫空间。我们称 \(u \in \mathcal{L}(\mathrm{E};\mathrm{F})\) 是完全连续的，如果它属于 \(\mathcal{L}(\mathrm{E};\mathrm{F})\) 中由有限秩线性映射组成的空间 \(\mathcal{L}^{\mathrm{f}}(\mathrm{E};\mathrm{F})\) 的闭包。我们用 \(\mathcal{L}^{\mathrm{c}}(\mathrm{E};\mathrm{F})\) 表示从 E 到 F 的完全连续线性映射组成的空间，并令 \(\mathcal{L}^{\mathrm{c}}(\mathrm{E}) = \mathcal{L}^{\mathrm{c}}(\mathrm{E};\mathrm{E})\) 。 a) 空间 \(\mathcal{L}^{\mathrm{c}}(\mathrm{E})\) 是 \(\mathcal{L}(\mathrm{E})\) 的一个闭双边理想。设 \(u \in \mathcal{L}(\mathrm{E};\mathrm{F})\) 。映射 \(\ell \mapsto \ell \circ u\) 是从 \(F'\) 到 \(E'\) 的连续线性映射，记为 \({}^{t}u\) 。若 \(u \in \mathcal{L}^{\mathrm{c}}(\mathrm{E};\mathrm{F})\) ，则 \({}^{t}u \in \mathcal{L}^{\mathrm{c}}(\mathrm{F}'; \mathrm{E}')\) 。
\end{enumerate}

\begin{enumerate}
\def\labelenumi{\alph{enumi})}
\setcounter{enumi}{1}
\item
  设 J 是一个集合，\(\mathrm{F} = \mathcal{C}_{0}(\mathrm{J})\) 。设 \(u \in \mathcal{L}(\mathrm{E};\mathrm{F})\) 。对每个 \(j \in J\)，映射 \(\ell_{j} \colon x \mapsto u(x)(j)\) 是 \(E'\) 的一个元素。我们有 \(\|u\| = \sup_{j \in J} \| \ell_{j} \|\)，并且映射 \(u \mapsto (\ell_{j})_{j \in J}\) 通过过渡到子空间诱导一个从 \(\mathcal{L}^{\mathrm{c}}(\mathrm{E};\mathrm{F})\) 到巴拿赫空间 \(\mathcal{C}_{0}(\mathrm{J};\mathrm{E}')\) 的同构。
\item
  若 K 是局部紧的，则 \(u \in \mathcal{L}(\mathrm{E}; \mathrm{F})\) 是完全连续的，当且仅当 E 的每个有界子集在 u 下的像在 F 中是相对紧的。（可将此结果与 III, p. 16 的命题 11 相比较。）
\end{enumerate}

\begin{enumerate}
\def\labelenumi{\arabic{enumi})}
\setcounter{enumi}{2}
\tightlist
\item
  我们用 A 表示 K 的由元素 \(x \in \mathrm{K}\)（满足 \(|x| \leqslant 1\)）组成的子环；它是一个局部环，其极大理想由 \(x \in \mathrm{A}\)（满足 \(|x| < 1\)）组成。
\end{enumerate}

设 I 是一个集合，E 是巴拿赫空间 \(\mathcal{C}_{0}(\mathrm{I})\)（参见习题 1 的问题 c)）。对 \(i \in I\)，用 \(\varphi_{i} \in E\) 表示 \(\{i\}\) 的特征函数。设 \(u \in \mathcal{L}^{\mathrm{c}}(\mathrm{E})\) 满足 \(\|u\| \leqslant 1\) 。

\begin{enumerate}
\def\labelenumi{\alph{enumi})}
\item
  子空间 \(E_{0}\) 由这样的元素 \(x \in E\) 组成：\(\|x\| \leqslant 1\)；它在 u 下稳定。
\item
  对 A 的每个非零理想 a，线性映射 u 通过过渡到商定义一个自同态 \(u_{a}\)，它作用于 A/a-模 \(E_{a} = E_{0}/aE_{0}\) 。这个 A/a-模是自由的，且 \(u_{a}\) 的像包含于一个有限生成的子模中。
\item
  存在唯一的形式级数 \(f_{u} \in A[[t]]\)，使得对 A 的每个非零理想 a，\(f_{u}\) 在 \((A/\mathfrak{a})[[t]]\) 中的像等于 A, VIII, p. 455 中定义的多项式 \(\det(1 - tu_{\mathfrak{a}})\)。
\end{enumerate}

\(d)\) 设 \(v \in \mathcal{L}^{\mathrm{c}}(\mathrm{E})\)，并设 \(c \in \mathrm{K}^{*}\) 满足 \(\|cv\| \leqslant 1\) 。形式级数 \(f_{cv}\) 与满足此性质的 \(c\) 的选取无关。

我们将用 \(\det(1-tv)\) 表示形式级数 \(f_v\)，这里 \(c \in \mathbf{K}^*\) 为任意满足 \(\|cv\| \leqslant 1\) 的元素。

¶ 4) 我们沿用上一习题的记号。

\begin{enumerate}
\def\labelenumi{\alph{enumi})}
\item
  设 \(v \in \mathcal{L}^c(\mathrm{E})\) 。形式级数 \(\det(1 - tv)\) 的收敛半径是无穷大。
\item
  设 \((v_{n})_{n\in\mathbf{N}}\) 是 \(\mathcal{L}^{c}(\mathrm{E})\) 中收敛于 \(v\in\mathcal{L}^{c}(\mathrm{E})\) 的一个序列。形式级数序列 \(\det(1-tv_{n})\) 对于形式级数系数的逐点收敛拓扑收敛于 \(\det(1-tv)\)。
\item
  设 \(v \in \mathcal{L}^{c}(E)\) 。对 K 的每个扩张 L 和每个 \(x \in L\)，我们用 \(\det(1 - xv)\) 表示形式幂级数 \(\det(1 - tv)\) 在 x 处的值，它由 a) 良定义。我们有 \(\det(1 + u)\det(1 + v) = \det(1 + u + v + u \circ v)\) 。
\item
  设 J 是一个集合，\(\mathrm{F} = \mathcal{C}_{0}(\mathrm{J})\) 。对所有 \(u \in \mathcal{L}^{\mathrm{c}}(\mathrm{E};\mathrm{F})\) 和 \(v \in \mathcal{L}^{\mathrm{c}}(\mathrm{F};\mathrm{E})\)，作为形式幂级数有 \(\det(1 - tu \circ v) = \det(1 - tv \circ u)\) 。
\item
  设 \(v \in \mathcal{L}^{\mathrm{c}}(\mathrm{E})\) 。我们用 \(\operatorname{Tr}(v)\) 表示 \(\det(1 - tu)\) 中 t 的系数的相反数。若 K 的特征为零，则
\end{enumerate}

\[
\det (1 - t u) = \exp \left(- \sum_ {m \geqslant 1} \frac {\operatorname{Tr} \left(u ^ {m}\right)}{m} t ^ {m}\right)
\]

作为形式幂级数成立。

¶ 5) 我们保留上一习题的记号。设 \(v \in \mathcal{L}^{\mathrm{c}}(\mathrm{E})\) 。

\begin{enumerate}
\def\labelenumi{\alph{enumi})}
\item
  形式幂级数 \(\det(1 - tv)\sum_{k\geqslant 0}t^{k}v^{k}\)（\(\mathcal{L}(\mathrm{E})[[t]]\) 中的）的收敛半径为无穷大。
\item
  设 \(\lambda \in \mathrm{K}\) 。E 的自同态 \(1_{\mathrm{E}} - \lambda v\) 可逆，当且仅当 \(\det(1 - \lambda v) \neq 0\) 。
\item
  设 \(\lambda \in K\) 满足 \(\det(1 - \lambda v) = 0\) 。存在 E 的唯一的闭子空间 \(I_{\lambda}\) 和 \(N_{\lambda}\)，它们在 \(v\) 下稳定，使得
\end{enumerate}

\begin{enumerate}
\def\labelenumi{(\roman{enumi})}
\item
  空间 E 是 \(I_{\lambda}\) 与 \(N_{\lambda}\) 的拓扑直和；
\item
  \(I_{\lambda}\) 的由 \(1_{E} - \lambda v\) 通过过渡到子空间得到的自同态是可逆的；
\item
  \(N_{\lambda}\) 的由 \(1_{E} - \lambda v\) 通过过渡到子空间得到的自同态是幂零的。
\end{enumerate}

\(d)\) 设 \(h \geqslant 1\) 是 det(1 - tv) 的零点 \(\lambda\) 的阶。我们有 \(\dim(N_{\lambda}) = h\) 。

应将本习题的结果与 III, p. 77 的 no. 5 中关于 R 或 C 上巴拿赫空间的结果相比较。

在下文中，K 表示域 R 或 C，且所有拓扑向量空间都是 K 上的。

\begin{enumerate}
\def\labelenumi{\arabic{enumi})}
\setcounter{enumi}{5}
\tightlist
\item
  设 E 是空间 \(C^{N}\)，赋以与半范数序列 \((p_{n})_{n\in\mathbf{N}}\) 相伴的拓扑，其中 \(p_{n}(x)=|x_{n}|\) 对 E 的每个元素 \(x=(x_{n})_{n\in\mathbf{N}}\) 成立。
\end{enumerate}

\begin{enumerate}
\def\labelenumi{\alph{enumi})}
\item
  空间 E 是一个弗雷歇空间。
\item
  E 的紧自同态组成的集合 \(\mathcal{L}^{\mathrm{c}}(\mathrm{E})\) 在 \(\mathcal{L}(\mathrm{E})\) 中不是闭的。
\end{enumerate}

\begin{enumerate}
\def\labelenumi{\arabic{enumi})}
\setcounter{enumi}{6}
\item
  设 E 和 F 是局部凸拓扑向量空间。若 F 是可度量化的，则每个紧线性映射 \(E \rightarrow F\) 的像是可数型的。
\item
  构造希尔伯特空间 E 的一个非紧自同态 \(u\) 的例子，使得 \(u^2\) 是紧的。
\item
  对每一对 \((x,y)\in]0,1]^{2}\) ，存在一个整数 \(n\geqslant1\) 和一个整数 k 满足 \(0\leqslant k\leqslant2^{n-1}-1\)，使得
\end{enumerate}

\[
2 ^ {- n} <   x \leqslant 2 ^ {- n + 1}, \quad \frac {2 k}{2 ^ {n}} <   y \leqslant \frac {2 k + 2}{2 ^ {n}}.
\]

然后我们如下定义 N: ]0,1]×]0,1] → C：

\[
\mathrm{N} (x, y) = \left\{ \begin{array}{l l} 2 & \text { if } \frac {2 k}{2 ^ {n}} <   y \leqslant \frac {2 k + 1}{2 ^ {n}} \\ 0 & \text { if } \frac {2 k + 1}{2 ^ {n}} <   y \leqslant \frac {2 k + 2}{2 ^ {n}}. \end{array} \right.
\]

\begin{enumerate}
\def\labelenumi{\alph{enumi})}
\tightlist
\item
  证明公式
\end{enumerate}

\[
u (f) (x) = \int_ {[ 0, 1 ]} \mathrm{N} (x, y) f (y) d y
\]

在 \(L^{1}([0,1])\) 上定义了一个关于勒贝格测度的连续线性映射 u。

\begin{enumerate}
\def\labelenumi{\alph{enumi})}
\setcounter{enumi}{1}
\item
  证明自同态 \(u\) 不是紧的。
\item
  自同态 \(u^2\) 是紧的。
\end{enumerate}

\(d)\) 设 \((x_{n})_{n\in \mathbf{N}}\) 是 \(\mathrm{L}^1 ([0,1])\) 中的一个有界序列。序列 \((u(x_n))_{n\in \mathbf{N}}\) 有一个弱收敛的子列。

\begin{enumerate}
\def\labelenumi{\alph{enumi})}
\setcounter{enumi}{4}
\item
  设 \((x_{n})_{n\in \mathbf{N}}\) 是 \(\mathrm{L}^1 ([0,1])\) 中的一个弱收敛序列。序列 \((u(x_{n}))_{n\in \mathbf{N}}\) 在 \(\mathrm{L}^1 ([0,1])\) 中收敛。
\item
  设 E 是这样的拓扑向量空间：其元素是有界可测的函数 \(f\colon[0,1]\to\mathbf{C}\)，赋以强于 \(\mathrm{L}^{2}([0,1])\) 的拓扑以及相对于 \(\mathrm{L}^{1}([0,1])\) 的弱拓扑的最细拓扑。则由
\end{enumerate}

\[
v (f) (x) = \int_ {[ 0, 1 ]} \overline {{\mathrm{N} (y , x)}} f (y) d y
\]

定义的线性映射 v 是 E 的一个紧自同态。它的转置 \(^{t}v\) 对于 E' 上的强拓扑不是紧的。

\begin{enumerate}
\def\labelenumi{\arabic{enumi})}
\setcounter{enumi}{9}
\tightlist
\item
  设 p 和 r 是满足 \(1 \leqslant p < r\) 的实数。设 u 是从 \(\ell_{\mathrm{K}}^{r}(\mathbf{N})\) 到 \(\ell_{\mathrm{K}}^{p}(\mathbf{N})\) 的一个连续线性映射。
\end{enumerate}

我们用 \(r_{n}\)（相应地，\(p_{n}\)）表示投影

\[
(x _ {i}) _ {i \in \mathbf {N}} \mapsto (x _ {0}, \ldots , x _ {n - 1}, 0, 0, \ldots)
\]

到 \(\ell_{\mathrm{K}}^{r}(\mathbf{N})\) 上（相应地，到 \(\ell_{\mathrm{K}}^{p}(\mathbf{N})\) 上）。我们用 \(\| \cdot \| _p\)（相应地，\(\| \cdot \| _r\)）表示 \(\ell_{\mathrm{K}}^{p}(\mathbf{N})\)（相应地，\(\ell_{\mathrm{K}}^{r}(\mathbf{N})\)）上的范数。

\begin{enumerate}
\def\labelenumi{\alph{enumi})}
\item
  证明对每个 \(n \in \mathbf{N}\) ，有 \(\|p_n u(1 - r_m)\| \to 0\)（当 \(m \to +\infty\)），其中 1 表示 \(\ell_{\mathrm{K}}^{r}(\mathbf{N})\) 的恒同态射。
\item
  对 \(n\) 、 \(m \in \mathbf{N}\) ，令 \(u_{n,m} = (1 - p_n)u(1 - r_m)\) 。设 \(\alpha = (\alpha_n)_{n \in \mathbf{N}}\) 是一个严格正实数的序列，满足 \(\|\alpha\|_r = 1\) 且 \(\alpha\) 不属于 \(\ell_{\mathrm{K}}^p(\mathbf{N})\)，并设 \((\varepsilon_n)\) 是一个严格正实数的序列，使得级数 \(\sum \varepsilon_n\alpha_n\) 收敛。
\end{enumerate}

若 \(u\) 不是紧的，证明存在 \(\delta > 0\) 、整数的严格递增序列 \((n_i)_{i \in \mathbf{N}}\) 与 \((m_i)_{i \in \mathbf{N}}\)，以及两个序列：\((x_i)_{i \in \mathbf{N}}\)（由 \(\ell_{\mathrm{K}}^r(\mathbf{N})\) 的元素组成）与 \((y_i)_{i \in \mathbf{N}}\)（由 \(\ell_{\mathrm{K}}^p(\mathbf{N})\) 的元素组成），具有下列性质：

\begin{enumerate}
\def\labelenumi{(\roman{enumi})}
\item
  有 \(\|x_{i}\|_{r}=1\) 对所有 \(i\in N\) 成立；
\item
  集合 \(S_{i} = \{n \in N \mid x_{i,n} \neq 0\}\) 两两不相交，且集合 \(T_{i} = \{n \in N \mid y_{i,n} \neq 0\}\) 两两不相交；
\item
  有 \(\| y_i\| _p\geqslant \delta\)
\item
  有 \(\|u(x_{i})\|_{q} > \delta - 2\varepsilon_{i}\) 。
\end{enumerate}

（用归纳法进行证明，并利用公式

\[
u = u _ {n, m} + p _ {n} u (1 - r _ {m}) + (1 - p _ {n}) u r _ {m} + p _ {n} u q _ {m}
\]

对 N 中的 n 和 m。）

\begin{enumerate}
\def\labelenumi{\alph{enumi})}
\setcounter{enumi}{2}
\item
  用反证法推出 u 是紧的（"皮特定理"）。
\item
  证明 \(\ell_{\mathrm{K}}^{p}(\mathbf{N})\) 到 \(\ell_{\mathrm{K}}^{r}(\mathbf{N})\) 的典范单射不是紧的。
\item
  设 p 和 r 是满足 \(1 \leqslant p \leqslant r \leqslant +\infty\) 的实数。证明 \(\ell_{\mathrm{K}}^{r}(\mathbf{N})\) 同构于 \(\ell_{\mathrm{K}}^{p}(\mathbf{N})\) 当且仅当 p = r。
\end{enumerate}

\begin{enumerate}
\def\labelenumi{\arabic{enumi})}
\setcounter{enumi}{10}
\tightlist
\item
  设 E 和 F 是 K 上的巴拿赫空间。称线性映射 \(u \in \mathcal{L}(\mathrm{E};\mathrm{F})\) 为完全连续的，如果对 E 中每个弱收敛序列 \((x_{n})_{n \in \mathbb{N}}\)，序列 \((u(x_{n}))_{n \in \mathbb{N}}\) 在 F 中收敛。
\end{enumerate}

\begin{enumerate}
\def\labelenumi{\alph{enumi})}
\item
  每个紧线性映射都是完全连续的。
\item
  若 E 的对偶空间 E' 是可数生成的，则每个完全连续线性映射 E → F 都是紧的。
\item
  设 \((u_{n})_{n\in\mathbf{N}}\) 是 \(\ell_{\mathrm{K}}^{1}(\mathbf{N})\) 中元素的一个序列。证明 \((u_{n})\) 收敛当且仅当它弱收敛（"舒尔定理"；参见 EVT, IV, p. 54, 习题 15）。
\end{enumerate}

\(d)\) 由此推出 \(\ell_{\mathrm{K}}^{1}(\mathbf{N})\) 到 \(\ell_{\mathrm{K}}^{2}(\mathbf{N})\) 的单射是完全连续的，但不是紧的。

\begin{enumerate}
\def\labelenumi{\alph{enumi})}
\setcounter{enumi}{4}
\item
  对每个巴拿赫空间，每个连续线性映射 \(\mathrm{E} \to \ell_{\mathrm{K}}^{1}(\mathbf{N})\) 都是完全连续的，且每个连续线性映射 \(\ell_{\mathrm{K}}^{1}(\mathbf{N}) \to \mathrm{E}\) 都是完全连续的。
\item
  若 E 是自反的或其对偶是可数生成的，则每个连续线性映射 \(\mathrm{E} \rightarrow \ell_{\mathrm{K}}^{1}(\mathbf{N})\) 都是紧的。
\end{enumerate}

\begin{enumerate}
\def\labelenumi{\arabic{enumi})}
\setcounter{enumi}{11}
\tightlist
\item
  设 E 是一个无限维自反巴拿赫空间。设 \(u \in \mathcal{L}^{c}(E)\) 是一个紧自同态。假设 u 不是有限秩的。设 \(E_{\sigma}\) 是赋以弱拓扑的空间 E。
\end{enumerate}

证明线性映射 \(^{t}u\) 是 \((\mathrm{E}_{\sigma})'\) 的一个紧自同态。

\begin{enumerate}
\def\labelenumi{\arabic{enumi})}
\setcounter{enumi}{12}
\tightlist
\item
  设 E 是 K 上的一个可数生成的巴拿赫空间。回顾（EVT, IV, p. 70, 习题 14）：E 中元素的一个序列 \((x_{n})_{n\in \mathbf{N}}\) 称为 E 的绍德尔基，如果对每个 \(x\in \mathrm{E}\) ，存在唯一的纯量序列 \((\lambda_n(x))_{n\in \mathbf{N}}\)，使得
\end{enumerate}

\[
x = \sum_ {n = 1} ^ {+ \infty} \lambda_ {n} (x) x _ {n}
\]

其中级数在 E 中收敛。

\begin{enumerate}
\def\labelenumi{\alph{enumi})}
\item
  每个可分希尔伯特空间都具有巴拿赫基；每个空间 \(\ell^{p}(\mathrm{I})\)（其中 I 可数）都具有巴拿赫基。
\item
  由复平面中半径为 1 的开圆盘 D 上有界全纯函数组成的巴拿赫空间 E，赋以范数 \(\| f \| = \sup_{z \in \mathrm{D}} |f(z)|\)，具有绍德尔基。
\item
  每个具有巴拿赫基的巴拿赫空间 E 都是一个逼近空间。更精确地说，存在实数 \(c \geqslant 0\)，使得 \(1_{E}\) 属于 \(\mathcal{L}_{\mathrm{pc}}(\mathrm{E})\) 中由 E 中范数 \(\leqslant c\) 的有限秩自同态组成的集合的闭包。
\end{enumerate}

\begin{enumerate}
\def\labelenumi{\arabic{enumi})}
\setcounter{enumi}{13}
\tightlist
\item
  \begin{enumerate}
  \def\labelenumii{\alph{enumii})}
  \tightlist
  \item
    设 E 和 F 是希尔伯特空间。从 E 到 F 的连续线性映射 u 是紧的，当且仅当 u 的像不含无限维的闭子空间。
  \end{enumerate}
\end{enumerate}

\begin{enumerate}
\def\labelenumi{\alph{enumi})}
\setcounter{enumi}{1}
\tightlist
\item
  a) 中关于从 E 到 F 的紧自同态的刻画并非对每一对巴拿赫空间 (E, F) 都成立。（可利用 p. 121 的习题 5。）
\end{enumerate}

\begin{enumerate}
\def\labelenumi{\arabic{enumi})}
\setcounter{enumi}{14}
\tightlist
\item
  设 \(E_{0}\) 是由函数 \(f: N^{*} \times N^{*} \to C\) 组成的复向量空间。对 \(k \in N^{*}\) ，设 \(\alpha_{k} \in E_{0}\) 是满足下式的函数：
\end{enumerate}

\[
\alpha_ {k} (n, m) = \left\{ \begin{array}{l l} m ^ {k} & \text { if } 1 \leqslant n \leqslant k - 1 \\ n ^ {k} & \text { if } n \geqslant k. \end{array} \right.
\]

设 E 是 \(E_{0}\) 的由满足 \(p_{k}(f)\) 对所有 \(k \in N^{*}\) 有限的函数 f 组成的向量子空间，其中 \(p_{k}\) 由

\[
p _ {k} (f) = \sum_ {n, m} \alpha_ {k} (n, m) | f (n, m) |.
\]

\begin{enumerate}
\def\labelenumi{\alph{enumi})}
\tightlist
\item
  赋以由范数 \((p_{k})_{k\in\mathbf{N}^{*}}\) 定义的拓扑的空间 E 是一个弗雷歇空间，也是一个蒙泰尔空间。它的对偶 E' 等同于 \(E_{0}\) 中由满足下式的函数 g 组成的子空间：对至少一个 \(k\geqslant1\) ，有
\end{enumerate}

\[
\sup _ {m, n} \alpha_ {k} (m, n) ^ {- 1} | g (m, n) | <   + \infty
\]

（参见 EVT, IV, p. 63, 习题 8）。

\begin{enumerate}
\def\labelenumi{\alph{enumi})}
\setcounter{enumi}{1}
\tightlist
\item
  考虑由下式定义的线性映射 \(u\)，它把 E 映到 \(\ell^1 (\mathbf{N}^*)\)：
\end{enumerate}

\[
u (f) = \left(\sum_ {n} f (n, m)\right) _ {m \in \mathbf {N} ^ {*}}.
\]

线性映射 u 是连续的。

\begin{enumerate}
\def\labelenumi{\alph{enumi})}
\setcounter{enumi}{2}
\item
  把 \(\ell^{\infty}(\mathbf{N}^{*})\) 等同于 \(\ell^{1}(\mathbf{N}^{*})\) 的对偶。\(u\) 的转置是映射 \(v\colon \ell^{\infty}(\mathbf{N}^{*})\to \mathrm{E}^{\prime}\subset \mathrm{E}_{0}\)，由 \((x_{n})\mapsto f\) 确定，其中 \(f(n,m) = x_{m}\) 对所有 \((n,m)\in \mathbf{N}^{*}\times \mathbf{N}^{*}\) 成立。
\item
  映射 u 诱导一个同构 \(\mathrm{E}/\mathrm{Ker}(u)\to\ell^{1}(\mathbf{N}^{*})\)（参见出处同上）。
\item
  映射 u 不是紧的。
\item
  u 的转置 v 对于在 \(\ell^{\infty}(\mathbf{N}^{*})\) 的有界子集上收敛的拓扑是紧的。
\item
  映射 v 是分离局部凸空间之间一个紧线性映射的例子，其转置不是紧的。
\end{enumerate}

\begin{enumerate}
\def\labelenumi{\arabic{enumi})}
\setcounter{enumi}{15}
\item
  设 E 是局部凸空间，u 是 E 的一个紧自同态。用 \(E_{b}^{\prime}\) 表示赋有有界收敛拓扑的 E 的对偶空间。线性映射 \(^{t}u \circ ^{t}u\) 是 \(E_{b}^{\prime}\) 的一个紧自同态。
\item
  设 E 是可数型的巴拿赫空间。则 \(\mathcal{L}^{\mathrm{c}}(\mathrm{E})\) 赋以由 \(\mathcal{L}(\mathrm{E})\) 的范数诱导的拓扑后是可数型的，当且仅当对偶空间 \(\mathrm{E}'\) 是可数型的。
\item
  设 E、F、G 是巴拿赫空间，u 是 \(\mathcal{L}^{c}(\mathrm{E};\mathrm{F})\) 的一个元素。
\end{enumerate}

\begin{enumerate}
\def\labelenumi{\alph{enumi})}
\tightlist
\item
  设 v 是 \(\mathcal{L}(\mathrm{E};\mathrm{G})\) 的一个元素。假设对每个 \(\varepsilon > 0\) ，存在 C \textgreater{} 0，使得
\end{enumerate}

\[
\| v (x) \| _ {\mathrm{G}} \leqslant \varepsilon \| x \| _ {\mathrm{E}} + \mathrm{C} \| u (x) \| _ {\mathrm{F}}
\]

对所有 \(x\in \mathrm{E}\) 成立。则 \(v\in \mathcal{L}^{\mathrm{c}}(\mathrm{E};\mathrm{G})\) 。

\begin{enumerate}
\def\labelenumi{\alph{enumi})}
\setcounter{enumi}{1}
\tightlist
\item
  设 w 是 \(\mathcal{L}(\mathrm{F};\mathrm{G})\) 的一个单射元素。则对每个 \(\varepsilon > 0\) ，存在 C \textgreater{} 0，使得
\end{enumerate}

\[
\| u (x) \| _ {\mathrm{F}} \leqslant \varepsilon \| x \| _ {\mathrm{E}} + \mathrm{C} \| (w \circ u) (x) \| _ {\mathrm{G}}
\]

对所有 \(x \in E\) 成立。

\begin{enumerate}
\def\labelenumi{\arabic{enumi})}
\setcounter{enumi}{18}
\tightlist
\item
  设 X 是一个局部紧、可度量化且在无穷远处可数的拓扑空间，\(\mu\) 是 X 上的一个正测度。我们用 A 表示 \(\mu\) 的原子集。设 \(f \in \mathcal{L}^{\infty}(X, \mu)\) 且 \(p \in [1, +\infty]\) 。
\end{enumerate}

证明乘以 \(f\) 在 \(\mathrm{L}^{p}(\mathrm{X},\mu)\) 中定义一个紧自同态，当且仅当下列两个条件被满足：

\begin{enumerate}
\def\labelenumi{(\roman{enumi})}
\item
  集合 \(\{x\in X-A\mid f(x)\neq0\}\) 是 \(\mu\) -可忽略的；
\item
  对每个 \(\varepsilon > 0\) ，存在 A 的有限子集 F，使得 \(|f(a)| \leqslant \varepsilon\) 对所有 \(a \in \mathbf{A} - \mathbf{F}\) 成立。
\end{enumerate}

（考虑乘以 \(f\) 的运算在空间 \(\mathrm{L}^{p}(\mathrm{X}_{\varepsilon},\mu)\) 上的作用，其中 \(\mathrm{X}_{\varepsilon}\) 是由 \(x\in\mathrm{X}\) 且满足 \(|f(x)|\geqslant\varepsilon\) 的元素组成的集合。）

\begin{enumerate}
\def\labelenumi{\arabic{enumi})}
\setcounter{enumi}{19}
\tightlist
\item
  设 X、Y 是紧度量空间，T 是从 \(\mathcal{C}(\mathrm{Y})\) 到 \(\mathcal{C}(\mathrm{X})\) 的一个紧算子。
\end{enumerate}

证明存在一个正测度 \(\mu\)（在 \(X\) 上）以及一个可测映射 \(N\)（从 \(X \times Y\) 到 \(\mathbf{C}\)），使得

\begin{enumerate}
\def\labelenumi{\alph{enumi})}
\item
  对每个 \(y \in Y\) ，映射 \(\mathrm{N}_{y} \colon x \mapsto \mathrm{N}(x, y)\) 是从 \(X\) 到 \(\mathbf{C}\) 的 \(\mu\) -可积映射；
\item
  映射 \(y \mapsto N_{y}\) 是从 Y 到 \(\mathcal{L}^{1}(X, \mu)\) 的连续映射；
\item
  对每个 \(f \in \mathcal{C}(\mathrm{X})\) 和每个 \(y \in \mathrm{Y}\) ，我们有
\end{enumerate}

\[
\mathrm{T} f (y) = \int_ {\mathrm{X}} \mathrm{N} (x, y) f (x) d \mu (x).
\]

\begin{enumerate}
\def\labelenumi{\arabic{enumi})}
\setcounter{enumi}{20}
\tightlist
\item
  我们用 D 表示 C 中以 0 为心、以 1 为半径的开圆盘，用 H \(^{2}\) (D) 表示 D 上的哈代空间（EVT, V, p. 7, 例 5）。我们记 T = R/2πZ，并赋以规范化的哈尔测度。设 j 是由下式定义的从 H \(^{2}\) (D) 到 L \(^{2}\) (T) 的等距：
\end{enumerate}

\[
j (f) (\theta) = \sum_ {n \geqslant 0} \frac {f ^ {(n)} (0)}{n !} e ^ {i n \theta}.
\]

H \(^{2}\) (D) 在 \(j\) 下的像将记为 H \(^{2}\) (T)。

设 \(\widetilde{\mathrm{P}}\) 是 \(\mathrm{L}^2 (\mathbf{T})\) 的以 \(\mathrm{H}^2 (\mathbf{T})\) 为像的正交投影，并设 P 是线性映射 \(f\mapsto \widetilde{\mathrm{P}} (f)\)，它从 \(\mathrm{L}^2 (\mathbf{T})\) 到 \(\mathrm{H}^2 (\mathbf{T})\)。若 \(f\) 是 \(\mathrm{L}^{\infty}(\mathbf{T})\) 的一个元素，设 \(\mathrm{M}_f\) 是自同态 \(g\mapsto fg\)，它作用于 \(\mathrm{L}^2 (\mathbf{T})\)。

证明若 \(f\) 是连续函数，则换位子 \(\widetilde{\mathrm{P}}\mathrm{M}_f - \mathrm{M}_f\widetilde{\mathrm{P}}\) 是一个紧线性映射。（考虑 \(f\) 是三角多项式的特殊情形。）

\begin{enumerate}
\def\labelenumi{\arabic{enumi})}
\setcounter{enumi}{21}
\tightlist
\item
  我们保留上一习题的记号。设 R 是从 \(\mathrm{H}^{2}(\mathbf{T})\) 到 \(\mathrm{L}^{2}(\mathbf{T})\) 的由 \(\mathrm{R}f(\theta)=f(-\theta)\) 定义的连续线性映射。对 \(g\in\mathrm{L}^{\infty}(\mathbf{T})\)，我们令 \(\Gamma_{g}=\mathrm{PM}_{g}\mathrm{R}\) 。
\end{enumerate}

\begin{enumerate}
\def\labelenumi{\alph{enumi})}
\item
  证明线性映射 \(\Gamma\colon g\mapsto\Gamma_{g}\) 是从 \(\mathrm{L}^{\infty}(\mathbf{T})\) 到 \(\mathcal{L}(\mathrm{H}^{2}(\mathbf{T}))\) 的连续映射，范数 \(\leqslant 1\)，且其核 F 等于 \(\mathrm{H}^{2}(\mathbf{T})^{\circ}\cap\mathrm{L}^{\infty}(\mathbf{T})\) 。
\item
  若 \(g\) 是连续函数，则 \(\Gamma_g\) 是 \(\mathrm{H}^2 (\mathbf{T})\) 的一个紧自同态。
\item
  设 \(g \in \mathrm{L}^{\infty}(\mathbf{T})\) 。我们用 \(\alpha_{g}\) 表示 \(\mathrm{H}^{2}(\mathbf{T})\) 上由下式定义的线性形式
\end{enumerate}

\[
\alpha_ {g} (f) = \int_ {\mathbf {T}} \Gamma_ {g} (f) (\theta) d \theta .
\]

设 \(\mathcal{A}\) 是在 \(j\) 下 \(\mathrm{H}^2 (\mathrm{D})\) 中由可解析延拓到 \(\overline{\mathrm{D}}\) 的一个邻域的全纯函数组成的子空间的像。证明若 \(f_{1}\) 和 \(f_{2}\) 属于 \(\mathcal{A}\) ，则

\[
\alpha_ {g} (f _ {1} f _ {2}) = \int_ {\mathbf {T}} \Gamma_ {g} (f _ {1}) (\theta) f _ {2} (- \theta) d \theta .
\]

\(d)\) 设 \(\varphi\) 是 \(\overline{\mathrm{D}}\) 的一个邻域上的全纯函数，在 \(\partial\mathrm{D}\) 上不为零，并设 \(f = j(\varphi|\mathrm{D})\) 。则存在 \(f_{1}\) 和 \(f_{2}\)（都在 \(\mathscr{A}\) 中），使得 \(f = f_{1}f_{2}\)，且

\[
\| f \| _ {\mathrm{L} ^ {1} (\mathbf {T})} = \| f _ {1} \| _ {\mathrm{L} ^ {2} (\mathbf {T})} \| f _ {2} \| _ {\mathrm{L} ^ {2} (\mathbf {T})}.
\]

（设 \(z_{1},\ldots,z_{n}\) 是 \(\varphi\) 在 D 中的零点；则存在一个在 D 的邻域内全纯的函数 \(\psi\)，使得

\[
\varphi (z) = \prod_ {k = 1} ^ {n} \frac {z - z _ {k}}{1 - \overline {{z}} _ {k} z} \psi (z) ^ {2}
\]

对所有 z 成立。）

\begin{enumerate}
\def\labelenumi{\alph{enumi})}
\setcounter{enumi}{4}
\tightlist
\item
  证明对所有 \(f \in \mathrm{H}^{2}(\mathbf{T})\) ，我们有不等式
\end{enumerate}

\[
\left| \alpha_ {g} (f) \right| \leqslant \left\| \Gamma_ {g} \right\| _ {\mathscr {L} \left(\mathrm{H} ^ {2} (\mathbf {T})\right)} \| f \| _ {\mathrm{L} ^ {1} (\mathbf {T})}.
\]

（若 \(\varphi \in \mathrm{H}^2 (\mathrm{D})\) ，则存在 ]0,1[ 中的一个序列 \((r_n)_{n\in \mathbf{N}}\)，使得对每个 \(n\in \mathbf{N}\) ，函数 \(\varphi_{n}\)（由 \(\varphi_{n}(z) = \varphi (r_{n}z)\) 定义）在 \(\partial \mathrm{D}\) 上没有零点，且 \(\lim \varphi_n = \varphi\) 在 \(\mathrm{H}^2 (\mathrm{D}).\) 中

由此推出存在 \(h \in \mathrm{L}^{\infty}(\mathbf{T})\)，使得

\[
\alpha_ {g} (f) = \int_ {\mathbf {T}} f (\theta) h (\theta) d \theta
\]

对所有 \(f \in \mathrm{H}^2(\mathbf{T})\) 成立，且 \(\|h\|_{\mathrm{L}^\infty(\mathbf{T})} \leqslant \| \Gamma_g \|_{\mathcal{L}(\mathrm{H}^2(\mathbf{T}))}\) 。

\(f)\) 证明存在 \(k \in \mathrm{L}^{\infty}(\mathbf{T})\)，使得 \(\Gamma_g = \Gamma_k\) 且 \(\| k \|_{\mathrm{L}^{\infty}(\mathbf{T})} = \| \Gamma_g \|_{\mathcal{L}(\mathrm{H}^2(\mathbf{T}))}\) 。

于是映射 \(\widehat{\Gamma}\colon\mathrm{L}^{\infty}(\mathbf{T})/\mathrm{F}\to\mathcal{L}(\mathrm{H}^{2}(\mathbf{T}))\) 是一个等距（“内哈里定理”）。

\(g)\) 设 \(g\) 是 \(\mathbf{T}\) 上的连续函数。则我们有 \(d(g, \mathrm{F}) = d(g, \mathrm{F} \cap \mathcal{C}(\mathbf{T}))\) 。（对 \(\mathrm{R} > 1\) ，用下式定义线性映射 \(\mathrm{P}_{\mathrm{R}}\)，即从 \(\mathrm{L}^{\infty}(\mathbf{T})\) 到 \(\mathcal{C}(\mathbf{T})\) 的映射：

\[
\mathrm{P} _ {\mathrm{R}} f (\theta) = \int_ {\mathbf {T}} \frac {\mathrm{R} ^ {2} - 1}{| \mathrm{R} e ^ {i (\theta - \varphi)} - 1 | ^ {2}} f (\varphi) d \varphi .
\]

则有 \(\| \mathrm{P}_{\mathrm{R}}f\|_{\mathcal{C}(\mathbf{T})}\leqslant \| f\|_{\mathrm{L}^{\infty}(\mathbf{T})}\)，且若 \(f\in \mathcal{C}(\mathbf{T})\) ，则 \(\mathrm{P_R}f\to f\) 在 \(\mathcal{C}(\mathbf{T})\) 中成立（当 \(\mathrm{R}\rightarrow 1.\)）。

由此推出 \(\Gamma(\mathcal{C}(\mathbf{T}))\) 是 \(\mathcal{L}(\mathrm{H}^{2}(\mathbf{T}))\) 的一个闭子空间（“萨拉森定理”）。

\begin{enumerate}
\def\labelenumi{\alph{enumi})}
\setcounter{enumi}{7}
\tightlist
\item
  设 \(g \in \mathrm{L}^{\infty}(\mathbf{T})\) 满足 \(\Gamma_{g}\) 是一个紧自同态。对 \(m \in Z\)，我们令 \(e_{m}(\theta) = e^{im\theta}\) 。证明 \(\|\Gamma_{ge_{-m}}\|_{\mathcal{L}(\mathrm{H}^{2}(\mathbf{T}))}\) 当 m 趋于 \(+\infty\) 时趋于零，然后证明存在 F 中的一个序列 \((h_{m})_{m \in \mathbf{N}}\)，使得序列 \((e_{m}h_{m})_{m \in \mathbf{N}}\) 在 \(\mathrm{L}^{\infty}(\mathbf{T})\) 中收敛于 g。由此推出在 g 模 F 的类中存在一个连续函数（“哈特曼定理”）。
\end{enumerate}

\begin{enumerate}
\def\labelenumi{\arabic{enumi})}
\setcounter{enumi}{22}
\item
  设 B 是一个无限维巴拿赫空间。我们用 E 表示 B 的强对偶，用 F 表示赋有拓扑 \(\sigma(\mathrm{B}^{\prime},\mathrm{B})\) 的 B 的对偶。证明 B 的对偶的恒同映射从 E 到 F 是紧的，但其转置从 \(\mathrm{F}_{b}^{\prime}\) 到 \(\mathrm{E}_{b}^{\prime}\) 不是紧的。
\item
  设 E 是一个实或复的巴拿赫空间，\(u \in \mathcal{L}(\mathrm{E})\) 。假设序列 \((u^{n})_{n \in \mathbf{N}}\) 有界，且 u 从 E 到赋有弱拓扑 \(\sigma(\mathrm{E}, \mathrm{E}')\) 的 E 是紧的。对每个整数 \(N \geqslant 1\)，我们令
\end{enumerate}

\[
v _ {\mathrm{N}} = \frac {1}{\mathrm{N}} (1 _ {\mathrm{E}} + u + \dots + u ^ {\mathrm{N-1}}) \in \mathcal {L} (\mathrm{E}).
\]

\begin{enumerate}
\def\labelenumi{\alph{enumi})}
\item
  设 F 是 \(1_{\mathrm{E}} - u\) 的像的闭包。对每个 \(x \in \mathrm{F}\) ，我们有 \(v_{\mathrm{N}}(x) \to 0\)（在 E 中，当 \(\mathrm{N} \to +\infty\) 时）。
\item
  设 \(x \in E\) 。序列 \((v_{\mathrm{N}}(x))\) 收敛到一个元素 \(y \in E\)，使得 \(u(y) = y\) 。（存在 \((v_{\mathrm{N}}(x))\) 的一个子序列，它弱收敛到 \(y \in E\) ；利用哈恩–巴拿赫定理证明 \(x - y \in F\) 。）
\item
  映射 \(w \colon \mathrm{E} \to \mathrm{E}\) 满足 \(w(x)\) 是上一问题给出的元素 \(y\)，它是一个连续投影，其像是 \(u\) 的对应于特征值 1 的特征空间。
\end{enumerate}

\(d)\) 设 \(z \in \mathbf{C}\) 满足 \(|z| = 1\) 。由下式定义的 \(\mathcal{L}(\mathrm{E})\) 中的序列

\[
v _ {z, \mathrm{N}} = \frac {1}{\mathrm{N}} (1 _ {\mathrm{E}} + z ^ {- 1} u \dots + z ^ {- (\mathrm{N-1})} u ^ {\mathrm{N-1}})
\]

在赋有逐点收敛拓扑的 \(\mathcal{L}(\mathrm{E})\) 中收敛。用 \(w_{z}\) 表示它的极限。则 \(w_{z}\) 是一个投影，我们有 \(w_{z}w = ww_{z} = zw_{z}\)，以及 \(w_{z_1}w_{z_2} = 0\)（若 \(z_{1} \neq z_{2}\)）。

\begin{enumerate}
\def\labelenumi{\arabic{enumi})}
\setcounter{enumi}{24}
\tightlist
\item
  \begin{enumerate}
  \def\labelenumii{\alph{enumii})}
  \tightlist
  \item
    设 \(N \geqslant 1\) 是一个整数。我们用 \(G_{N}\)（相应地，\(H_{N}\)）表示集合 \(\{-1,1\}^{N}\)（相应地，集合 \(\{-1,2\}^{N}\)）。我们用 \(\mu_{N}\)（相应地，\(\nu_{N}\)）表示 \(G_{N}\) 上（相应地，\(H_{N}\) 上）总质量为 1 的正测度，使得对每个 \((t_{i})_{1 \leqslant i \leqslant N}\)（属于 \(G_{N}\)，相应地，属于 \(H_{N}\)），有 \(\mu_{\mathrm{N}}((t_{i})_{1 \leqslant i \leqslant N}) = \frac{1}{2^{\mathrm{N}}}\) ，相应地，
  \end{enumerate}
\end{enumerate}

\[
\nu_ {\mathrm{N}} ((t _ {i}) _ {1 \leqslant i \leqslant \mathrm{N}}) = \left(\frac {1}{3}\right) ^ {j} \left(\frac {2}{3}\right) ^ {\mathrm{N} - j},
\]

其中 j 是使得 \(t_{i}=2\) 的指标 i 的个数。

存在一个实数 \(c > 0\)，使得对每个由复数组成的族 \((a_{i})_{1\leqslant i\leqslant N}\)，我们有

\[
\mu_ {\mathrm{N}} \Big (\Big \{(t _ {i}) _ {1 \leqslant i \leqslant \mathrm{N}} \in \mathrm{G} _ {\mathrm{N}} | \Big | \sum_ {i = 1} ^ {\mathrm{N}} a _ {i} t _ {i} \Big | > c (\log \mathrm{N}) ^ {1 / 2} \Big (\sum_ {i = 1} ^ {\mathrm{N}} | a _ {i} | ^ {2} \Big) ^ {1 / 2} \Big \} \Big) <   \frac {c}{\mathrm{N} ^ {3}},
\]

以及

\[
\nu_ {\mathrm{N}} \Big (\Big \{(t _ {i}) _ {1 \leqslant i \leqslant \mathrm{N}} \in \mathrm{H} _ {\mathrm{N}} | \left| \sum_ {i = 1} ^ {\mathrm{N}} a _ {i} t _ {i} \right| > c (\log \mathrm{N}) ^ {1 / 2} \Big (\sum_ {i = 1} ^ {\mathrm{N}} | a _ {i} | ^ {2} \Big) ^ {1 / 2} \Big \} \Big) <   \frac {c}{\mathrm{N} ^ {3}}.
\]

（我们可以假设 \(a_{i} \in \mathbf{R}\) 对所有 \(i\) 成立，且 \(\sum |a_{i}|^{2} = 1\) 。对第一个不等式，证明并利用上界估计

\[
\int_ {G _ {N}} \exp \left(\lambda \left| \sum_ {i = 1} ^ {N} a _ {i} t _ {i} \right|\right) d \mu_ {N} \leqslant 2 \exp (\lambda^ {2})
\]

对每个实数 \(\lambda > 0\) 成立；对第二个上界估计，利用不等式 \(\frac{1}{3}(e^{2x} + 2e^{-x}) \leqslant e^{2x^2}\)（对所有 \(x \in \mathbf{R}\)）以类似的方式推理。）

\begin{enumerate}
\def\labelenumi{\alph{enumi})}
\setcounter{enumi}{1}
\tightlist
\item
  对 \(k \in N\)，令 \(\mathrm{A}_{k} = \mathbf{Z}/(3 \cdot 2^{k})\mathbf{Z}\) 。存在一个不依赖于 k 的实数 \(c' > 0\)，以及一个函数 \(f_{k} \colon \widehat{\mathrm{A}}_{k} \to \{-1, 2\}\)，使得
\end{enumerate}

\[
\sum_ {\chi \in \widehat {\mathrm{A}} _ {k}} f _ {k} (\chi) = 0
\]

且对所有 \(x \in A_{k}\) ，有

\[
\Big | \sum_ {\chi \in \widehat {\mathrm{A}} _ {k}} \chi (x) f _ {k} (\chi) \Big | \leqslant c ^ {\prime} (k + 1) ^ {1 / 2} 2 ^ {k / 2}.
\]

（利用 \(a\)），确定满足第二个条件的函数 \(f_{k}\) 的存在性，然后修改 \(f_{k}\) 以得到第一个条件。）

我们用 \(B_{k}\)（相应地，\(C_{k}\)）表示特征标 \(\chi\)（\(A_{k}\) 上的）中满足 \(f_{k}(\chi)=2\)（相应地，\(f_{k}(\chi)=-1\)）者所成的集合。我们有 \(\mathrm{Card(B_{k})}=2^{k}\) 和 \(\mathrm{Card(C_{k})}=2^{k+1}\) 。

对每个 \(k \geqslant 1\) ，我们固定一个双射 \(\iota_k \colon B_k \to C_{k-1}\) 。

我们用 A 表示诸 \(A_{k}\)（\(k \in N\)）的不交并；我们赋给 A 离散拓扑，并记 \(\mathrm{E} = \mathcal{C}_{b}(\mathrm{A})\) 。

我们固定一个函数族 \(\varepsilon = (\varepsilon_k)_{k\geqslant \mathbf{N}}\)，其中的函数 \(\varepsilon_{k}\colon \mathrm{B}_{k}\to \{-1,1\}\)。然后我们定义函数 \(e_{k,\chi}\in \mathrm{E}\)（对 \(k\in \mathbf{N}\) 和 \(\chi \in \mathrm{B}_k\)）如下：

\[
e _ {k, \chi} (x) = \left\{ \begin{array}{l l} \iota_ {k} (\chi) (x) & \text { if } k \geqslant 1 \text {   and   } x \in \mathrm{A} _ {k - 1}, \\ \varepsilon_ {k} (\chi) \chi (x) & \text { if } x \in \mathrm{A} _ {k} \\ 0 & \text { otherwise. } \end{array} \right.
\]

我们用 \(E_{\varepsilon}\) 表示由函数 \((e_{k,\chi})_{k\in\mathbf{N},\chi\in\mathbf{B}_{k}}\) 生成的 E 的闭子空间。这是一个可数型的复巴拿赫空间。习题的其余部分旨在证明：对 \(\varepsilon\) 的适当选取，空间 \(\mathrm{E}_{\varepsilon}\) 不具有逼近性质；它不拥有巴拿赫基。

\begin{enumerate}
\def\labelenumi{\alph{enumi})}
\setcounter{enumi}{2}
\tightlist
\item
  对 \(k \in N\) 和 \(\chi \in B_{k}\)，存在一个连续线性形式 \(\lambda_{k,\chi} \in E_{\varepsilon}'\)，使得对每个 \(f \in E_{\varepsilon}\)，
\end{enumerate}

\[
\lambda_ {k, \chi} (f) = \frac {1}{\operatorname{Card} (\mathrm{A} _ {k})} \sum_ {x \in \mathrm{A} _ {k}} \varepsilon_ {k} (\chi) \chi (x ^ {- 1}) f (x).
\]

对 \(k \geqslant 1\) 和 \(f \in \mathrm{E}_{\varepsilon}\) ，我们有

\[
\lambda_ {k, \chi} (f) = \frac {2}{\operatorname{Card} (\mathrm{A} _ {k})} \sum_ {x \in \mathrm{A} _ {k}} \iota_ {k} (\chi) (x ^ {- 1}) f (x).
\]

\begin{enumerate}
\def\labelenumi{\alph{enumi})}
\setcounter{enumi}{3}
\tightlist
\item
  对每个 \(k \in N\) ，用 \(\beta_{k}\) 表示 \(\mathcal{L}(E_{\varepsilon})\) 上由下式定义的连续线性形式
\end{enumerate}

\[
\beta_ {k} (u) = \frac {1}{2 ^ {k}} \sum_ {\chi \in \mathrm{B} _ {k}} \lambda_ {k, \chi} (u (e _ {k, \chi})).
\]

我们有 \(\beta_{k}(1_{\mathrm{E}_{\varepsilon}}) = 1\)

\begin{enumerate}
\def\labelenumi{\alph{enumi})}
\setcounter{enumi}{4}
\tightlist
\item
  设 \(u \in \mathcal{L}(\mathrm{E}_{\varepsilon})\) 且 \(k \in \mathbf{N}\) 。我们有
\end{enumerate}

\[
\beta_ {k + 1} (u) - \beta_ {k} (u) = \frac {1}{\operatorname{Card} \left(\mathrm{A} _ {k}\right)} \sum_ {x \in \mathrm{A} _ {k}} u \left(\varphi_ {k, x}\right) (x)
\]

其中，对 \(x \in \mathrm{A}_k\) ，我们令

\[
\varphi_ {k, x} = \frac {1}{2 ^ {k + 1}} \sum_ {\chi \in B _ {k + 1}} \iota_ {k + 1} (\chi) (x ^ {- 1}) e _ {k + 1, \chi} - \frac {1}{2 ^ {k}} \sum_ {\chi \in B _ {k}} \varepsilon_ {k} (\chi) \chi (x ^ {- 1}) e _ {k, \chi}.
\]

我们有 \(|\varphi_{k,x}(y)| \leqslant c'(k + 1)^{1/2}2^{-k/2}\)，对 \((x,y) \in \mathrm{A}_k^2\) 成立。

\begin{enumerate}
\def\labelenumi{\alph{enumi})}
\setcounter{enumi}{5}
\item
  存在一个实数 \(c'' \geqslant c'\)，使得对每个 \(k \geqslant 1\) ，存在一个函数 \(\varepsilon_k \colon B_k \to \{-1,1\}\)，满足 \(|\varphi_{k,x}(y)| \leqslant c''3(k+1)^{1/2}2^{-k/2}\) 对 \((x,y) \in A_k \times A_{k-1}\) 成立。（应用 a) 的第一个不等式。）我们现在假设 \(\varepsilon = (\varepsilon_k)\) 满足这些性质。
\item
  对每个 \(k \in \mathbf{N}\) 和 \(x \in \mathrm{A}_k\) ，我们有 \(\| \varphi_{k,x} \| \leqslant c''(k + 1)^{1/2} 2^{-k/2}\) 。
\item
  设 K 是 \(E_{\varepsilon}\) 的子集，其元素为 \(e_{0,\chi_{0}}\)（其中 \(\chi_{0}\) 是 \(B_{0}\) 的唯一元素）以及函数 \((k+1)^{2}\varphi_{k,x}\)（\(k\in N\)，\(x\in A_{k}\)）。集合 K 在 \(E_{\varepsilon}\) 中是相对紧的，且我们有
\end{enumerate}

\[
| \beta_ {k + 1} (u) - \beta_ {k} (u) | \leqslant \frac {1}{(k + 1) ^ {2}} \sup _ {x \in K} \| u (x) \|
\]

对所有 \(k \in \mathbf{N}\) 和每个 \(u \in \mathcal{L}(\mathrm{E}_{\varepsilon})\) 成立。

\begin{enumerate}
\def\labelenumi{\roman{enumi})}
\tightlist
\item
  序列 \((\beta_{k})_{k\in \mathbf{N}}\) 在赋有逐点收敛拓扑的空间 \(\mathcal{L}(\mathrm{E}_{\varepsilon})\) 的对偶中收敛。用 \(\beta\) 表示它的极限。对每个 \(u\in \mathcal{L}(\mathrm{E}_{\varepsilon})\) ，我们有 \(|\beta (u)|\leqslant 3\sup_{x\in \mathrm{K}}\| u(x)\|\) 。
\end{enumerate}

\begin{enumerate}
\def\labelenumi{\alph{enumi})}
\setcounter{enumi}{9}
\tightlist
\item
  我们有 \(\beta(1_{\mathrm{E}_{\varepsilon}})=1\)，且 \(\beta(u)=0\)（若 \(u\) 是有限秩的）。
\end{enumerate}

\(k)\) 空间 \(\mathrm{E}_{\varepsilon}\) 不具有逼近性质。

\begin{enumerate}
\def\labelenumi{\alph{enumi})}
\setcounter{enumi}{11}
\tightlist
\item
  设 \(p \in]2, +\infty[\) 。我们赋给 A 计数测度，并用 \(E_{\varepsilon}^{p}\) 表示 \(L^{p}(A)\) 的由 \((e_{k,\chi})_{k \in \mathbf{N}, \chi \in \mathrm{B}_{k}}\) 生成的闭子空间。空间 \(E_{\varepsilon}^{p}\) 不具有逼近性质。（验证 \(\varphi_{k,x}\) 的支撑的基数 \(\leqslant 21 \cdot 2^{k}\)，并由此推出 \(\|\varphi_{k,x}\|_{p} = \mathrm{O}((k+1)^{1/2}2^{-k(p-2)/2p})\) 对 \(k \in N\) 和 \(x \in A_{k}\) 成立。）
\end{enumerate}

\begin{enumerate}
\def\labelenumi{\arabic{enumi})}
\setcounter{enumi}{25}
\tightlist
\item
  设 E 和 F 是巴拿赫空间，u 是从 E 到 F 的一个连续线性映射。
\end{enumerate}

\begin{enumerate}
\def\labelenumi{\alph{enumi})}
\tightlist
\item
  证明下列条件等价：
\end{enumerate}

\begin{enumerate}
\def\labelenumi{(\roman{enumi})}
\item
  映射 u 是紧的；
\item
  存在一个序列 \((x_n')_n\)（在 E 的强对偶 \(\mathrm{E}'\) 中趋于 0），满足 \(\| u(x)\| \leqslant \sup |\langle x,x_n' \rangle|\) 对所有 \(x \in \mathrm{E}\) 成立；
\item
  存在巴拿赫空间 \(c_{0}\) 的一个闭子空间 L 以及紧线性映射 \(f: E \to L\) 和 \(g: L \to F\)，使得 \(u = g \circ f\)（若 u 是紧的，对 F 的单位球的极集在 \(^{t}u\) 下的像应用 EVT, IV, p. 24 的推论 1。在假设 (ii) 下，设 \((\lambda_{n})\) 是 \(c_{0}\) 的一个元素，使得序列 \((\lambda_{n}^{-1}x_{n}^{\prime})\) 趋于 0；由 \(f(x) = (\langle x, \lambda_{n}^{-1}x_{n}^{\prime} \rangle)_{n}\) 定义 f，并取 L 为 f 的像的闭包。）
\end{enumerate}

\begin{enumerate}
\def\labelenumi{\alph{enumi})}
\setcounter{enumi}{1}
\tightlist
\item
  若 \(u\) 是紧的，我们有 \(\|u\| = \inf(\sup_n \|x_n'\|)\)，其中下确界取遍满足 (ii) 的序列 \((x_n')\) 所成的集合，且 \(\|u\| = \inf(\|f\|\|g\|)\) ，下确界取遍满足 (iii) 的条件的三元组 (L, \(f,g\) ) 所成的集合。
\end{enumerate}

\begin{enumerate}
\def\labelenumi{\arabic{enumi})}
\setcounter{enumi}{26}
\tightlist
\item
  设 E 和 F 是局部凸空间。映射 \(u \in \mathcal{L}(\mathrm{E};\mathrm{F})\) 称为弱紧的，如果它是从 E 到赋有拓扑 \(\sigma(\mathrm{F},\mathrm{F}')\) 的 F 的紧映射。
\end{enumerate}

\begin{enumerate}
\def\labelenumi{\alph{enumi})}
\item
  弱紧映射构成一个子空间 \(\mathcal{L}^{\mathrm{fc}}(\mathrm{E};\mathrm{F})\)，它是 \(\mathcal{L}(\mathrm{E};\mathrm{F})\) 的子空间。设 \(\mathrm{E}_1\) 、\(\mathrm{F}_1\) 是局部凸空间，\(f\in \mathcal{L}(\mathrm{E}_1;\mathrm{E})\)，\(g\in \mathcal{L}(\mathrm{F};\mathrm{F}_1)\) 且 \(u\in \mathcal{L}(\mathrm{E};\mathrm{F})\)；若 \(u\) 是弱紧的，则 \(g\circ u\circ f\) 亦然。
\item
  我们有 \(\mathcal{L}^{c}(\mathrm{E};\mathrm{F})\subset\mathcal{L}^{\mathrm{fc}}(\mathrm{E};\mathrm{F})\) 。若 E 是无限维自反巴拿赫空间，则 E 的恒同映射是弱紧的但不是紧的。
\end{enumerate}

\begin{enumerate}
\def\labelenumi{\arabic{enumi})}
\setcounter{enumi}{27}
\tightlist
\item
  设 E 和 F 是局部凸空间，且 \(u \in \mathcal{L}(\mathrm{E};\mathrm{F})\) 。
\end{enumerate}

\begin{enumerate}
\def\labelenumi{\alph{enumi})}
\tightlist
\item
  证明下列性质等价：
\end{enumerate}

\begin{enumerate}
\def\labelenumi{(\roman{enumi})}
\item
  E 的每个有界子集在 \(u\) 下的像对于弱拓扑是相对紧的；
\item
  映射 \(^{tt}u\) 把 F'' 映到 E 中。
\end{enumerate}

\begin{enumerate}
\def\labelenumi{\alph{enumi})}
\setcounter{enumi}{1}
\tightlist
\item
  若 F 是拟完备的，证明这些性质等价于：
\end{enumerate}

\begin{enumerate}
\def\labelenumi{(\roman{enumi})}
\setcounter{enumi}{2}
\tightlist
\item
  映射 \(^{t}u\) 把 F' 的每个等度连续子集映到 E' 的一个对于 \(\sigma(\mathrm{E}',\mathrm{E}'')\) 相对紧的子集。
\end{enumerate}

\begin{enumerate}
\def\labelenumi{\alph{enumi})}
\setcounter{enumi}{2}
\tightlist
\item
  若 E 和 F 是巴拿赫空间，则 (i) 至 (iii) 也等价于
\end{enumerate}

(i') 映射 u 是弱紧的；

(iii') 映射 \(^{t}u\colon F_{b}^{\prime}\to E_{b}^{\prime}\) 是弱紧的。

\begin{enumerate}
\def\labelenumi{\arabic{enumi})}
\setcounter{enumi}{28}
\tightlist
\item
  \begin{enumerate}
  \def\labelenumii{\alph{enumii})}
  \tightlist
  \item
    证明巴拿赫空间 \(\ell^{1}(\mathbf{N})\) 的每个弱紧子集是紧的（参见 EVT, IV, p. 69, 习题 11）。特别地，从局部凸空间 E 到 \(\ell^{1}(\mathbf{N})\) 的每个弱紧线性映射是紧的。
  \end{enumerate}
\end{enumerate}

\begin{enumerate}
\def\labelenumi{\alph{enumi})}
\setcounter{enumi}{1}
\tightlist
\item
  设 E 是巴拿赫空间。由 a) 推出，从 \(c_{0}\) 到 E 的每个弱紧线性映射是紧的。
\end{enumerate}

\begin{enumerate}
\def\labelenumi{\arabic{enumi})}
\setcounter{enumi}{29}
\tightlist
\item
  \begin{enumerate}
  \def\labelenumii{\alph{enumii})}
  \tightlist
  \item
    设 E 和 F 是巴拿赫空间，\(u \in \mathcal{L}(\mathrm{E};\mathrm{F})\) 。证明下列性质等价：
  \end{enumerate}
\end{enumerate}

\begin{enumerate}
\def\labelenumi{(\roman{enumi})}
\item
  E 的每个弱紧子集在 \(u\) 下的像是（强）紧的；
\item
  E 的每个弱相对紧子集在 \(u\) 下的像在 F 中相对紧；
\item
  E 中每个弱收敛序列在 \(u\) 下的像在 F 中强收敛。
\end{enumerate}

（利用施穆良定理，EVT, IV, p. 36, 定理 2。）

如果 u 满足这些性质，就说 u 是半紧的 \(^{(3)}\) 。

\begin{enumerate}
\def\labelenumi{\alph{enumi})}
\setcounter{enumi}{1}
\item
  设 E 和 F 是巴拿赫空间。证明由半紧映射组成的集合 \(\mathcal{L}^{\mathrm{sc}}(\mathrm{E};\mathrm{F})\) 是 \(\mathcal{L}(\mathrm{E};\mathrm{F})\) 的一个闭子空间。
\item
  设 E 和 F 是巴拿赫空间，\(u \in \mathcal{L}(\mathrm{E};\mathrm{F})\) 。设 \(E_{1}\) 、\(F_{1}\) 是巴拿赫空间，\(f \in \mathcal{L}(\mathrm{E}_{1};\mathrm{E})\) 且 \(g \in \mathcal{L}(\mathrm{F};\mathrm{F}_{1})\) 。若 \(u: E \to F\) 是半紧的，则 \(g \circ u \circ f\) 亦然。
\item
  设 E 和 F 是巴拿赫空间，\(u \in \mathcal{L}(\mathrm{E};\mathrm{F})\) 。若 u 是紧的，则它是半紧的。反之，若 E 是自反的且 u 是半紧的，则 u 是紧的。
\item
  设 E 和 F 是巴拿赫空间。若 E 的对偶空间 E′ 是可数型的，则每个完全连续线性映射 E → F 都是紧的。
\item
  设 I 是无限集；\(\ell^{1}(\mathbf{I})\) 的恒同映射是半紧的，但不是弱紧的（参见 EVT, IV, p. 54, 习题 15）。g) \(\ell_{\mathrm{K}}^{1}(\mathbf{N})\) 到 \(\ell_{\mathrm{K}}^{2}(\mathbf{N})\) 的单射是半紧的但不是紧的（参见出处同上）。
\item
  设 E 是巴拿赫空间。每个连续线性映射 \(E \rightarrow \ell_{K}^{1}(\mathbf{N})\) 都是半紧的，且每个连续线性映射 \(\ell_{\mathrm{K}}^{1}(\mathbf{N}) \rightarrow \mathrm{E}\) 都是半紧的。
(3) 一些作者称它们为完全连续映射。
\item
  设 E 是巴拿赫空间。若 E 是自反的或其对偶是可数型的，则每个连续线性映射 \(\mathrm{E} \to \ell_{\mathrm{K}}^{1}(\mathbf{N})\) 都是紧的。
\end{enumerate}

\begin{enumerate}
\def\labelenumi{\arabic{enumi})}
\setcounter{enumi}{30}
\tightlist
\item
  设 E 是巴拿赫空间。
\end{enumerate}

\begin{enumerate}
\def\labelenumi{\alph{enumi})}
\tightlist
\item
  证明下列条件等价：
\end{enumerate}

\begin{enumerate}
\def\labelenumi{(\roman{enumi})}
\item
  对任何巴拿赫空间 F，从 E 到 F 的每个弱紧映射都是半紧的。
\item
  若 \((x_{n})\) 是 E 中对于拓扑 \(\sigma(\mathrm{E},\mathrm{E}^{\prime})\) 趋于 0 的序列，而 \((x_{n}^{\prime})\) 是 \(E^{\prime}\) 中对于拓扑 \(\sigma(\mathrm{E}^{\prime},\mathrm{E}^{\prime\prime})\) 趋于 0 的序列，则序列 \(\langle x_{n},x_{n}^{\prime}\rangle\) 趋于 0。
\end{enumerate}

（为证明 (i) 蕴含 (ii)，取 \(F = c_0\) 。若 (i) 不被满足，则存在从 E 到某个巴拿赫空间 F 的一个弱紧映射 \(u\) 以及 E 中弱趋于 0 的一个序列 \((x_n)\)，使得 \(\inf(\|u(x_n)\|) > 0\) ；构造一个序列 \((y_n)\)（在 \(E'\) 的单位球中），它对于 \(\sigma(E', E'')\) 收敛且满足 \(\langle u(x_n), y_n \rangle = \|u(x_n)\|\) ，并取 \(x_n' = y_n - \lim y_n\) 。）

这些条件特别在下列情形成立：\(E = \mathcal{C}(T)\)（其中 T 是紧空间），以及 \(E = L^{1}(X, \mu)\) 或 \(E = L^{\infty}(X, \mu)\)（其中 X 是局部紧空间，\(\mu\) 是 X 上的一个测度）（INT, VI, § 2, 习题 23）。b) 证明具有上述性质的巴拿赫空间不拥有无限维的自反直因子。

\begin{enumerate}
\def\labelenumi{\alph{enumi})}
\setcounter{enumi}{2}
\tightlist
\item
  设 H 是自反巴拿赫空间，B 是 H' 的赋有弱拓扑的单位球。证明 H 可以等同于 \(\mathcal{C}(\mathrm{B})\) 的一个闭的非直接子空间。
\end{enumerate}

\begin{enumerate}
\def\labelenumi{\arabic{enumi})}
\setcounter{enumi}{31}
\item
  一个拥有绍德尔基的可数型巴拿赫空间（EVT, IV, p. 70, 习题 14）具有逼近性质。
\item
  \begin{enumerate}
  \def\labelenumii{\alph{enumii})}
  \tightlist
  \item
    证明逼近空间序列的严格归纳极限是逼近空间。
  \end{enumerate}
\end{enumerate}

\begin{enumerate}
\def\labelenumi{\alph{enumi})}
\setcounter{enumi}{1}
\tightlist
\item
  证明一族逼近空间的局部凸直和是逼近空间。
\end{enumerate}

\begin{enumerate}
\def\labelenumi{\arabic{enumi})}
\setcounter{enumi}{33}
\tightlist
\item
  设 E 是巴拿赫空间。证明下列性质等价：
\end{enumerate}

\begin{enumerate}
\def\labelenumi{(\roman{enumi})}
\item
  空间 E' 是逼近空间；
\item
  对每个巴拿赫空间 F，\(\mathcal{L}^{\mathrm{f}}(\mathrm{E};\mathrm{F})\) 在 \(\mathcal{L}(\mathrm{E};\mathrm{F})\) 中的闭包是 \(\mathcal{L}^{\mathrm{c}}(\mathrm{E};\mathrm{F})\) ；
\item
  对 \(c_{0}\) 的每个闭子空间 F，\(\mathcal{L}^{\mathrm{f}}(\mathrm{E};\mathrm{F})\) 在 \(\mathcal{L}(\mathrm{E};\mathrm{F})\) 中的闭包是 \(\mathcal{L}^{\mathrm{c}}(\mathrm{E};\mathrm{F})\) 。
\end{enumerate}

\begin{enumerate}
\def\labelenumi{\arabic{enumi})}
\setcounter{enumi}{34}
\tightlist
\item
  设 E 是巴拿赫空间。我们称 E 具有度量逼近性质，如果 E 的恒同映射属于范数 \(\leqslant 1\) 的元素（\(\mathcal{L}^{\mathrm{f}}(\mathrm{E})\) 中的）组成的集合在赋有紧收敛拓扑的空间 \(\mathcal{L}(\mathrm{E})\) 中的闭包。
\end{enumerate}

证明下列性质等价：

\begin{enumerate}
\def\labelenumi{(\roman{enumi})}
\item
  空间 E 具有度量逼近性质；
\item
  \(\mathcal{L}(\mathrm{E})\) 的单位球是范数 \(\leqslant 1\) 的元素（\(\mathcal{L}^{\mathrm{f}}(\mathrm{E})\) 中的）组成的集合对于紧收敛拓扑的闭包；
\item
  对每个巴拿赫空间 F，\(\mathcal{L}(\mathrm{E};\mathrm{F})\) 的单位球是范数 \(\leqslant 1\) 的元素（\(\mathcal{L}^{\mathrm{f}}(\mathrm{E};\mathrm{F})\) 中的）组成的集合对于紧收敛拓扑的闭包；
\item
  对每个巴拿赫空间 F，\(\mathcal{L}(\mathrm{F};\mathrm{E})\) 的单位球是范数 \(\leqslant 1\) 的元素（\(\mathcal{L}^{\mathrm{f}}(\mathrm{F};\mathrm{E})\) 中的）组成的集合对于紧收敛拓扑的闭包。
\end{enumerate}

\begin{enumerate}
\def\labelenumi{\arabic{enumi})}
\setcounter{enumi}{35}
\tightlist
\item
  \begin{enumerate}
  \def\labelenumii{\alph{enumii})}
  \tightlist
  \item
    设 X 是局部紧拓扑空间。证明空间 \(\mathcal{C}_{0}(\mathrm{X})\) 具有度量逼近性质（改造 III, p. 21 命题 13 的证明）。
  \end{enumerate}
\end{enumerate}

\begin{enumerate}
\def\labelenumi{\alph{enumi})}
\setcounter{enumi}{1}
\tightlist
\item
  设 X 是赋有正测度的局部紧拓扑空间，且 \(p \in [1, +\infty]\) 。证明空间 \(L^{p}(X)\) 具有度量逼近性质（同样的方法）。
\end{enumerate}

\begin{enumerate}
\def\labelenumi{\arabic{enumi})}
\setcounter{enumi}{36}
\tightlist
\item
  设 X 是 \(C^{r}\) 类（\(1 \leqslant r \leqslant \infty\)）微分流形，仿紧且局部有限维，并设 E 是 X 上的 \(C^{r}\) 类向量丛，局部有限维。证明空间 \(\mathcal{S}^{r}(\mathrm{X};\mathrm{E})\) 具有逼近性质。
\end{enumerate}

\exercisesection{}

\begin{enumerate}
\def\labelenumi{\arabic{enumi})}
\tightlist
\item
  设 \(I = [a, b]\) 是 R 的一个紧区间。对每个连续函数 \(f: I \to C\) 和每个 \(y \in I\)，令
\end{enumerate}

\[
\mathrm{V} (f) (y) = \int_ {a} ^ {y} f (x) d x.
\]

证明从 \(\mathcal{C}(\mathrm{I})\) 到 \(\mathcal{C}(\mathrm{I})\) 的线性映射 V 是紧的（“沃尔泰拉算子”）。

\begin{enumerate}
\def\labelenumi{\arabic{enumi})}
\setcounter{enumi}{1}
\tightlist
\item
  设 X = {]}0, 1{[}，μ 是 X 上的勒贝格测度。用下式定义一个函数序列 \((f_{n})_{n \geqslant 1}\)（X → R）：
\end{enumerate}

\[
f _ {2 n} (x) = \left\{ \begin{array}{l l} \sqrt {2} \sin (2 n \pi x) & \text { if } 0 <   x <   1 / 2 \\ - 2 ^ {n} & \text { if } 1 - 2 ^ {- 2 n} <   x <   1 - 2 ^ {- 2 n - 1} \\ 0 & \text { otherwise }, \end{array} \right.
\]

以及

\[
f _ {2 n - 1} (x) = \left\{ \begin{array}{l l} \sqrt {2} \sin (2 n \pi x) & \text { if } 0 <   x <   1 / 2 \\ 2 ^ {n} & \text { if } 1 - 2 ^ {- 2 n} <   x <   1 - 2 ^ {- 2 n - 1} \\ 0 & \text { otherwise }, \end{array} \right.
\]

对所有 \(n \geqslant 1\) 和每个 \(x \in X\) 成立。

\begin{enumerate}
\def\labelenumi{\alph{enumi})}
\item
  族 \((f_n)_{n\geqslant 1}\) 在 \(\mathcal{L}^2 (\mathrm{X},\mu)\) 中是规范正交的。
\item
  对每个 \((x,y)\in \mathrm{X}\times \mathrm{X}\)，我们令
\end{enumerate}

\[
\mathrm{N} (x, y) = \sum_ {n = 1} ^ {+ \infty} 2 ^ {- n} (f _ {2 n - 1} (x) f _ {2 n - 1} (y) - f _ {2 n} (x) f _ {2 n} (y)).
\]

证明 \(|\mathrm{N}(x,y)| \leqslant 2^{3/2}\) 对所有 \((x,y) \in \mathrm{X} \times \mathrm{X}\) 成立。由此推出 \(\mathrm{N}_y \in \mathcal{L}^\infty(\mathrm{X})\) 对所有 \(y \in \mathrm{X}\) 成立，且

\[
\int_ {X} ^ {*} \| N _ {y} \| _ {\infty} ^ {2} d y <   + \infty .
\]

\begin{enumerate}
\def\labelenumi{\alph{enumi})}
\setcounter{enumi}{2}
\tightlist
\item
  证明通过令
\end{enumerate}

\[
u _ {\mathrm{N}} (f) (y) = \int_ {\mathrm{X}} \mathrm{N} (x, y) f (x) d \mu
\]

对 \(f \in \mathcal{L}^{1}(X, \mu)\) 和 \(y \in X\) ，通过过渡到商空间，可以定义一个从 \(\mathrm{L}^{1}(X, \mu)\) 到 \(\mathrm{L}^{2}(X, \mu)\) 的连续线性映射。

\begin{enumerate}
\def\labelenumi{\alph{enumi})}
\setcounter{enumi}{3}
\tightlist
\item
  线性映射 \(u_{N}\) 不是紧的。（考虑序列 \(g_{n}=2^{n}(f_{2n-1}-f_{2n})\) ；证明 \(g_{n}\) 可积，且序列 \((u_{\mathrm{N}}(g_{n}))\)（在 \(\mathrm{L}^{2}(\mathrm{X},\mu)\) 中）没有收敛子列。）
\end{enumerate}

\begin{enumerate}
\def\labelenumi{\arabic{enumi})}
\setcounter{enumi}{2}
\tightlist
\item
  设 X 是局部紧拓扑空间，\(\mu\) 是 X 上的一个正测度，\(Y \subset X\) 是局部 \(\mu\) -可忽略的子集。设 N 是 \(X \times Y\) 的特征函数。
\end{enumerate}

\begin{enumerate}
\def\labelenumi{\alph{enumi})}
\tightlist
\item
  我们有
\end{enumerate}

\[
\int_ {X \times X} ^ {*} | N (x, y) f (x) g (y) | d (\mu \otimes \mu) (x, y) = 0
\]

对 \(f, g \in \mathrm{L}^2(\mathrm{X}, \mu)\) 成立。

\begin{enumerate}
\def\labelenumi{\alph{enumi})}
\setcounter{enumi}{1}
\tightlist
\item
  若 \(f \in \mathrm{L}^{2}(\mathrm{X}, \mu) - \mathrm{L}^{1}(\mathrm{X}, \mu)\)，则由 \(y \in X\) 中使 \(x \mapsto \mathrm{N}(x, y)f(x)\) 为 \(\mu\) -可积的那些 y 组成的集合与 Y 重合。
\end{enumerate}

\begin{enumerate}
\def\labelenumi{\arabic{enumi})}
\setcounter{enumi}{3}
\tightlist
\item
  设 \(\alpha \in ]0,1]\) 。设 I 是 \(\mathbf{R}\) 的一个紧区间。我们用 \(\mathcal{C}^{\alpha}(\mathrm{I})\) 表示 \(\mathcal{C}(\mathrm{I})\) 的由满足下列条件的函数 \(f\) 组成的向量子空间：存在 \(\mathrm{C} \geqslant 0\) 满足
\end{enumerate}

\[
| f (x) - f (y) | \leqslant \mathrm{C} | x - y | ^ {\alpha}
\]

对所有 \((x,y)\in\mathrm{I}\times\mathrm{I}\) 成立。

我们赋给 \(\mathcal{C}^{\alpha}(\mathrm{I})\) 范数

\[
\| f\|_{\alpha} = \sup_{x\in \mathrm{I}}|f(x)| + \sup_{\substack{(x,y)\in \mathrm{I}\times \mathrm{I}\\ x\neq y}}\frac{|f(x) - f(y)|}{|x - y|^{\alpha}}.
\]

\begin{enumerate}
\def\labelenumi{\alph{enumi})}
\item
  空间 \(\mathcal{C}^{\alpha}(\mathrm{I})\) 是巴拿赫空间。
\item
  典范单射 \(\mathcal{C}^{\alpha}(\mathrm{I})\to \mathcal{C}(\mathrm{I})\) 是紧的。
\item
  若 \(0 < \alpha < \beta \leqslant 1\)，则典范单射 \(\mathcal{C}^{\beta}(\mathrm{I}) \to \mathcal{C}^{\alpha}(\mathrm{I})\) 是紧的。\\
\item
  索伯列夫空间 \(\mathrm{H}^{1}(\mathring{\mathrm{I}})\)（参见 IV, p. 221, 第 14 目）等同于 \(\mathcal{C}^{1/2}(\mathrm{I})\) 的一个子空间，但典范单射 \(\mathrm{H}^{1}(\mathring{\mathrm{I}}) \to \mathcal{C}^{1/2}(\mathrm{I})\) 不是紧的。
\end{enumerate}

\exercisesection{}

\begin{enumerate}
\def\labelenumi{\arabic{enumi})}
\item
  设 E 和 F 是希尔伯特空间，u 是从 E 到 F 的一个连续线性映射。证明 u 是弗雷德霍姆映射，当且仅当 u 的像是闭的且 u 和 \(u^{*}\) 的核是有限维的。
\item
  设 \(p\) 是满足 \(1 \leqslant p < +\infty\) 的实数。设 \(u\) 是 \(\ell_{\mathrm{K}}^{p}(\mathbf{N})\) 的一个对角自同态。证明 \(u\) 是弗雷德霍姆自同态，当且仅当它是里斯自同态。特别地，\(\ell_{\mathrm{K}}^{p}(\mathbf{N})\) 的每个对角弗雷德霍姆自同态的指标为零。
\item
  \begin{enumerate}
  \def\labelenumii{\alph{enumii})}
  \tightlist
  \item
    存在一个巴拿赫空间 E、E 的一个连续自同构 u 以及 E 的一个里斯自同态 v，使得 \(u \circ v\) 不是里斯自同态。甚至存在这样的例子，其中 \(u \circ v - v \circ u\) 是有限秩的。
  \end{enumerate}
\end{enumerate}

\begin{enumerate}
\def\labelenumi{\alph{enumi})}
\setcounter{enumi}{1}
\tightlist
\item
  存在一个巴拿赫空间 E、E 的一个自同构 u 以及一个有限秩自同态 h，使得 \(u + h\) 不是里斯自同态。
\end{enumerate}

\begin{enumerate}
\def\labelenumi{\arabic{enumi})}
\setcounter{enumi}{3}
\tightlist
\item
  设 E 是无限维复希尔伯特空间。
\end{enumerate}

\begin{enumerate}
\def\labelenumi{\alph{enumi})}
\item
  E 上的酉算子群 \(\mathbf{U}(\mathrm{E})\)，赋以由 \(\mathcal{L}(\mathrm{E})\) 的范数诱导的拓扑，是道路连通的。（证明若 u 是酉的，则存在一个埃尔米特算子 v，使得 \(u = \exp(iv)\) 。）
\item
  E 的线性同构群是道路连通的。（利用极分解，参见 I, p. 139, 定义 3。）
\item
  对每个 \(n \in Z\)，E 的指标为 n 的弗雷德霍姆自同态组成的集合是开且连通的。它是非空的。
\end{enumerate}

\begin{enumerate}
\def\labelenumi{\arabic{enumi})}
\setcounter{enumi}{4}
\tightlist
\item
  设 E 和 F 是无限维巴拿赫空间。连续线性映射 \(u: \mathrm{E} \to \mathrm{F}\) 称为严格奇异的，如果不存在无限维的闭子空间 \(\mathrm{F}_1 \subset \mathrm{E}\)，使得 \(u\) 通过过渡到子空间诱导一个从 \(\mathrm{F}_1\) 到 \(u(\mathrm{F}_1)\) 上的同构。
\end{enumerate}

\begin{enumerate}
\def\labelenumi{\alph{enumi})}
\tightlist
\item
  设 u 是从 E 到 F 的一个连续线性映射。下列条件等价：
\end{enumerate}

\begin{enumerate}
\def\labelenumi{(\roman{enumi})}
\item
  映射 \(u\) 是严格奇异的；
\item
  对 E 的每个无限维闭子空间 \(F_{1}\)，不存在实数 \(\delta > 0\)，使得 \(\|u(x)\| \geqslant \delta\|x\|\) 对所有 \(x \in F_{1}\) 成立。
\end{enumerate}

\begin{enumerate}
\def\labelenumi{\alph{enumi})}
\setcounter{enumi}{1}
\item
  设 \(u \colon \mathrm{E} \to \mathrm{F}\) 是一个连续线性映射，使得 \(u\) 的像是闭的且 \(u\) 的核是有限维的。若 \(v \colon \mathrm{E} \to \mathrm{F}\) 是严格奇异的，则 \(u + v\) 的像是闭的且其核是有限维的。
\item
  E 的严格奇异自同态组成的集合是 \(\mathcal{L}(\mathrm{E})\) 的一个闭双边理想。
\item
  每个紧线性映射都是严格奇异的。若 E 和 F 是希尔伯特空间，则每个严格奇异映射 \(E \rightarrow F\) 都是紧的。
\item
  存在巴拿赫空间 E 和 F 以及一个严格奇异但非紧的映射 \(E \rightarrow F\) 。
\end{enumerate}

\(f)\) 设 \(p\) 和 \(r\) 是满足 \(1 \leqslant p < r\) 的实数。每个连续线性映射 \(\ell_{\mathrm{K}}^{r}(\mathbf{N}) \to \ell_{\mathrm{K}}^{p}(\mathbf{N})\) 都是严格奇异的。（利用 III, p. 106 习题 10。）

\begin{enumerate}
\def\labelenumi{\arabic{enumi})}
\setcounter{enumi}{5}
\item
  设 E 是复希尔伯特空间。设 u 是 E 的一个弗雷德霍姆自同态。若 u 是正规的，则 \(\text{ind}(u) = 0\) 。
\item
  设 E 和 F 是巴拿赫空间，u 是 \(\mathcal{L}(\mathrm{E};\mathrm{F})\) 的一个元素。证明 u 是弗雷德霍姆映射，当且仅当下列条件被满足：
\end{enumerate}

\begin{enumerate}
\def\labelenumi{(\roman{enumi})}
\tightlist
\item
  存在 E 上的一个连续半范数 p 和 \(C \geqslant 0\)，使得从 E 到 \((E, p)\) 的分离完备化的典范映射是紧的，且
\end{enumerate}

\[
\| x \| \leqslant \mathrm{C} \| u (x) \| + p (x)
\]

对所有 \(x \in E\) 成立；

\begin{enumerate}
\def\labelenumi{(\roman{enumi})}
\setcounter{enumi}{1}
\tightlist
\item
  存在一个连续半范数 \(q\)（\(\mathrm{F}'\) 上的）和 \(\mathrm{C}' \geqslant 0\)，使得从 \(\mathrm{F}'\) 到 \((\mathrm{F}', q)\) 的分离完备化的典范映射是紧的，且
\end{enumerate}

\[
\| y ^ {\prime} \| \leqslant \mathrm{C} ^ {\prime} \| ^ {t} u (y ^ {\prime}) \| + q (y ^ {\prime})
\]

对所有 \(y' \in F'\) 成立。

\begin{enumerate}
\def\labelenumi{\arabic{enumi})}
\setcounter{enumi}{7}
\tightlist
\item
  设 U 是 C 的一个连通开子集。我们用 \(\mathcal{H}(\mathrm{U})\) 表示 U 上的赋有紧收敛拓扑的全纯函数空间（EVT III, p. 10, 例 \(d\)）。
\end{enumerate}

\begin{enumerate}
\def\labelenumi{\alph{enumi})}
\item
  设 \(a \in \mathcal{H}(U)\) 。乘以 a 是 \(\mathcal{H}(U)\) 的一个弗雷德霍姆自同态，当且仅当 a 的零点集是有限的。然后计算它的指标。
\item
  设 \(\mathcal{U} = (\mathrm{U}_{i})_{i \in \mathrm{I}}\) 是 U 的一个由非空连通开集组成的覆盖。我们令
\end{enumerate}

\[
\begin{array}{l} \mathrm{I} _ {2} = \{(i, j) \in \mathrm{I} \times \mathrm{I} | \mathrm{U} _ {i} \cap \mathrm{U} _ {j} \neq \varnothing \}, \\ \mathrm{I} _ {3} = \{(i, j, k) \in \mathrm{I} \times \mathrm{I} \times \mathrm{I} | \mathrm{U} _ {i} \cap \mathrm{U} _ {j} \cap \mathrm{U} _ {k} \neq \varnothing \}. \end{array}
\]

我们考虑向量空间

\[
\begin{array}{l} \mathrm{Z} ^ {1} (\mathcal {U}; \mathbf {C}) = \{(c _ {i, j}) \in \mathbf {C} ^ {\mathrm{I} _ {2}} | c _ {i j} + c _ {j k} + c _ {k i} = 0 \text { for all } (i, j, k) \in \mathrm{I} _ {3} \} \\ \mathrm{B} ^ {1} (\mathcal {U}; \mathbf {C}) = \{(c _ {i, j}) \in \mathbf {C} ^ {\mathrm{I} _ {2}} | \text { there exists } c ^ {\prime} \in \mathbf {C} ^ {\mathrm{I}} \text { such that } \end{array}
\]

\[
c _ {i j} = c _ {i} ^ {\prime} - c _ {j} ^ {\prime} \text {   for all   } (i, j) \in \mathrm{I} _ {2} \}.
\]

我们有 \(\mathrm{B}^1 (\mathcal{U};\mathbf{C})\subset \mathrm{Z}^1 (\mathcal{U};\mathbf{C})\) 。我们令

\[
\mathrm{H} ^ {1} (\mathcal {U}; \mathbf {C}) = \mathrm{Z} ^ {1} (\mathcal {U}; \mathbf {C}) / \mathrm{B} ^ {1} (\mathcal {U}; \mathbf {C}).
\]

设 \(\mathcal{V}=(\mathrm{V}_{\alpha})_{\alpha\in\mathrm{A}}\) 是 U 的一个由非空连通开集组成且比 U 更细的覆盖，并设 \(\varrho\colon A\to I\) 是一个映射，使得 \(V_{\alpha}\subset U_{\varrho(\alpha)}\) 对所有 \(\alpha\in A\) 成立。

映射 \((c_{ij})_{i,j} \mapsto (c_{\varrho(\alpha)\varrho(\beta)})_{\alpha,\beta}\) 诱导一个不依赖于 \(\varrho\) 的选取的线性映射，从 \(\mathrm{H}^1(\mathcal{U};\mathbf{C})\) 到 \(\mathrm{H}^1(\mathcal{V};\mathbf{C})\) 。

我们用 \(H^{1}(U; \mathbf{C})\) 表示当 U 取遍 U 的由非空连通开集组成的覆盖时所得到的赋有这些线性映射的向量空间 \(H^{1}(\mathcal{U}; \mathbf{C})\) 的归纳极限。

\begin{enumerate}
\def\labelenumi{\alph{enumi})}
\setcounter{enumi}{2}
\tightlist
\item
  设 D 是 \(\mathcal{H}(\mathrm{U})\) 的由 \(Df = f'\) 定义的自同态。证明存在一个从 D 的余核到 \(\mathrm{H}^1 (\mathrm{U};\mathbf{C})\) 上的向量空间同构。由此推出 D 是弗雷德霍姆自同态，当且仅当空间 \(\mathrm{H}^1 (\mathrm{U};\mathbf{C})\) 是有限维的。然后计算 D 的指标。
\end{enumerate}

\exercisesection{}

\begin{enumerate}
\def\labelenumi{\arabic{enumi})}
\item
  设 E 和 F 是巴拿赫空间，u 是一个连续线性映射 \(E \rightarrow F\) 。我们有 \(((u)) > 0\)，当且仅当 u 的像在 F 中是闭的。
\item
  设 \(p \in [1, +\infty]\) ，且 \(\mathrm{E} = \ell_{\mathbf{C}}^{p}(\mathbf{N})\) 。设 u 是 E 的一个对角自同态。对每个 \(\lambda \in C\) ，用 u 的谱计算 \(((u - \lambda 1_{\mathrm{E}}))\) 。
\item
  设 \(p \in [1, +\infty]\) 且 \(\mathrm{E} = \ell_{\mathrm{K}}^{p}(\mathbf{N})\)，并设 \(u\) 是移位自同态
\end{enumerate}

\[
(x _ {0}, x _ {1}, \dots ,) \mapsto (0, x _ {0}, x _ {1}, \dots)
\]

（E 上的）。证明 u 是指标为 1 的弗雷德霍姆自同态，且 \(((u)) = 1\) 。

\begin{enumerate}
\def\labelenumi{\arabic{enumi})}
\setcounter{enumi}{3}
\item
  设 E 是巴拿赫空间。映射 \(u \mapsto ((u))\)（从 \(\mathcal{L}(\mathrm{E})\) 到 R）是连续的，当且仅当 E 的维数 \(\leqslant 1\) 。
\item
  设 E 是希尔伯特空间，\(u \in \mathcal{L}(\mathrm{E})\) 。我们有 \(((u)) = \inf(\mathrm{Sp}(|u|) - \{0\})\)（\(|u|\) 的定义参见 I, p. 139, 定义 3）。
\item
  设 \(E = \ell^{2}(\mathbf{N})\)，\((e_{n})_{n \in \mathbf{N}}\) 是 E 的规范正交基。设 u 是 E 的对角自同态，满足 \(u(e_{0}) = e_{0}\) 且 \(u(e_{n}) = 2e_{n}\)（若 \(n \geqslant 1\)）。证明 \(((u)) = 1\)，且 \(u + h\) 对每个满足 \(\|h\| < 2\) 的 h 是弗雷德霍姆自同态。
\item
  设 E 和 F 是巴拿赫空间，u 是 \(\mathcal{Q}\mathcal{M}(\mathrm{E};\mathrm{F})\)（相应地，\(\mathcal{Q}\mathcal{E}(\mathrm{E};\mathrm{F})\)）的一个元素。设 h 是 \(\mathcal{L}(\mathrm{E};\mathrm{F})\) 的一个元素。假设存在正实数 a 和 b，使得 \(a + b((u))^{-1} < 1\)，从而
\end{enumerate}

\[
\| h (x) \| \leqslant a \| u (x) \| + b \| x \|\tag{2}
\]

对所有 \(x \in E\) 成立。

证明 \(u + h \in \mathcal{Q}\mathcal{M}(\mathrm{E};\mathrm{F})\)（相应地，\(u + h \in \mathcal{Q}\mathcal{E}(\mathrm{E};\mathrm{F})\)）且

\[
\begin{array}{c} \dim \operatorname{Ker} (u + h) \leqslant \dim \operatorname{Ker} (u), \\ \dim \operatorname{Coker} (u + h) \leqslant \dim \operatorname{Coker} (u), \\ \operatorname{ind} (u + h) = \operatorname{ind} (u). \end{array}
\]

（化归到 \(a\) 和 \(b\) 都 \(>0\) 的情形，并设 \(c > 0\) 满足 \(a + b/c < 1\) ；考虑 E 上由下式定义的范数 \(p\)

\[
p (x) = a \| u (x) \| + b \| x \| \quad \text {   for   } x \in \mathrm{E},
\]

并设 \(E_{1}\) 是赋有范数 p 的巴拿赫空间 E；我们有 \(\mathcal{Q}\mathcal{M}(E_{1};F)=\mathcal{Q}\mathcal{M}(E;F)\)，\(\mathcal{Q}\mathcal{E}(E_{1};F)=\mathcal{Q}\mathcal{E}(E;F)\) 且 \(\mathcal{L}(E_{1};F)=\mathcal{L}(E;F)\) 。然后应用 III, p. 62 的引理 2。）

\begin{enumerate}
\def\labelenumi{\arabic{enumi})}
\setcounter{enumi}{7}
\item
  设 E 是希尔伯特空间。不用博苏克–乌拉姆定理，直接证明 III, p. 64 的定理 4。
\item
  设 \(r \geqslant 1\) 是整数。若 \(f: \mathrm{M}_1 \to \mathrm{M}_2\) 是 \(\mathrm{C}^r\) 类流形之间的一个 \(\mathrm{C}^r\) 类态射，则称 \(x \in \mathrm{M}_1\) 是 \(f\) 的临界点，如果 \(f\) 在 \(x\) 处不是淹没（VAR, R1, 5.9.1）。设 X 是 \(f\) 的临界点集。集合 \(f(\mathrm{X})\)（在 \(\mathrm{M}_2\) 中）称为 \(f\) 的临界值集，而 \(f(\mathrm{M}_1) = f(\mathrm{X})\) 称为 \(f\) 的正则值集。
\end{enumerate}

若 \(M_{1}\) 和 \(M_{2}\) 是 \(C^{r}\) 类实流形（VAR, R1, 5.1.5, p. 35），流形的一个态射 \(f: M_{1} \to M_{2}\)（\(C^{r}\) 类）称为弗雷德霍姆的，如果对每个 \(x \in M_{1}\) ，切线性映射 \(\mathrm{T}_{x}(f)\) 是从巴拿赫空间 \(\mathrm{T}_{x}(\mathrm{M}_{1})\) 到巴拿赫空间 \(\mathrm{T}_{f(x)}(\mathrm{M}_{2})\) 的弗雷德霍姆自同态。a) 设 \(f: M_{1} \to M_{2}\) 是一个弗雷德霍姆态射。若 \(M_{1}\) 是连通的，则指标 \(\operatorname{ind}(\mathrm{T}_{x}(f))\) 不依赖于 \(x \in M_{1}\) 。它称为弗雷德霍姆态射 f 的指标，记作 \(\operatorname{ind}(f)\) 。

\begin{enumerate}
\def\labelenumi{\alph{enumi})}
\setcounter{enumi}{1}
\tightlist
\item
  设 U 是巴拿赫空间 E 的一个开子集，满足 \(0 \in U\)，并设 F 是巴拿赫空间。设 \(f: U \to F\) 是一个 \(C^{r}\) 类弗雷德霍姆态射，并设 \(u = \mathrm{T}_{0}(f)\) 。存在一个巴拿赫空间 \(E_{1}\)、U 的一个卡 \((\mathrm{U}_{1}, \varphi, \mathrm{E}_{1} \times \mathrm{Ker}(u))\)（满足 \(0 \in U_{1}\)）、F 的一个卡 \((\mathrm{U}_{2}, \psi, \mathrm{E}_{1} \times \mathrm{Coker}(u))\)（满足 \(f(0) \in U_{2}\)），以及一个弗雷德霍姆态射 \(g: \varphi(\mathrm{U}_{1}) \to \mathrm{E}_{1} \times \mathrm{Coker}(u)\)（\(C^{r}\) 类），使得
\end{enumerate}

\[
(\psi \circ f \circ \varphi^ {- 1}) (a, b) = (a, g (a, b))
\]

对所有 \((a,b)\in \varphi (\mathrm{U}_1)\subset \mathrm{E}_1\times \mathrm{Ker}(u).\) 成立。

\begin{enumerate}
\def\labelenumi{\alph{enumi})}
\setcounter{enumi}{2}
\item
  设 \(f \colon \mathrm{M}_1 \to \mathrm{M}_2\) 是 \(\mathbf{C}^r\) 类流形之间的一个 \(\mathbf{C}^r\) 类弗雷德霍姆态射。对每个 \(y \in \mathbf{M}_2\)，存在一个开邻域 \(\mathrm{V}\)（\(y\) 的），使得 \(f\) 通过过渡到子空间定义一个逆紧映射 \(f^{-1}(\mathrm{V}) \to \mathrm{V}\) 。
\item
  设 U 是巴拿赫空间 E 的一个包含 0 的开子集，并设 F 是巴拿赫空间。设 \(f \colon \mathrm{U} \to \mathrm{F}\) 是一个 \(\mathrm{C}^{\infty}\) 类弗雷德霍姆态射。对每个 \(x \in \mathrm{U}\) 和 F 中 \(f(x)\) 的每个邻域 V，存在 \(f\) 在 V 中的一个正则值。（利用第二萨德定理，VAR, R2, 10.1.3, p. 37。）e) 设 \(f \colon \mathrm{M}_1 \to \mathrm{M}_2\) 是 \(\mathrm{C}^{\infty}\) 类流形的一个 \(\mathrm{C}^{\infty}\) 类弗雷德霍姆态射。若 \(\mathrm{M}_1\) 和 \(\mathrm{M}_2\) 具有可数基，则 \(f\) 的临界值集是贫集（“斯梅尔定理”）。
\end{enumerate}

\(f)\) 设 \(f\colon \mathrm{M}_1\to \mathrm{M}_2\) 是具有可数基的 \(\mathrm{C}^{\infty}\) 类流形的一个 \(\mathrm{C}^{\infty}\) 类弗雷德霍姆态射。若 \(\operatorname {ind}(f) < 0\)，则 \(f\) 的像的内部是空的。

\(g)\) 设 \(f\colon \mathrm{M}_1\to \mathrm{M}_2\) 是具有可数基的 \(\mathrm{C}^{\infty}\) 类流形的一个 \(\mathrm{C}^{\infty}\) 类弗雷德霍姆态射。对每个位于一个贫集之外的 \(y\in \mathrm{M}_2\)，集合 \(f^{-1}(y)\) 是维数为 \(\operatorname {ind}(f)\) 的纯子流形。
\exercisesection{}

在下面的习题中，当 E 是巴拿赫空间时，我们用 \(\mathcal{C}alk(E)\) 表示卡尔金代数 \(\mathcal{L}(E)/\mathcal{L}^{c}(E)\) 。用 \(\pi\) 表示从 \(\mathcal{L}(E)\) 到 \(\mathcal{C}alk(E)\) 上的典范投影 。

若 E 是复希尔伯特空间，我们给 \(\mathcal{C}alk(\mathrm{E})\) 赋以由 \(\mathcal{L}(\mathrm{E})\) 的对合经取商而诱导的对合。

\begin{enumerate}
\def\labelenumi{\arabic{enumi})}
\tightlist
\item
  设 E 是局部凸空间。
\end{enumerate}

\begin{enumerate}
\def\labelenumi{\alph{enumi})}
\item
  假设 E 中 0 的每个邻域都包含 E 的一个有限余维数的向量子空间。证明：E 的弗雷德霍姆自同态所成的空间 \(\mathcal{F}(\mathrm{E})\) 在赋以有界收敛拓扑的 \(\mathcal{L}(\mathrm{E})\) 中不是闭的。
\item
  给出一个无穷维弗雷歇空间的例子，使得其中 0 的每个邻域都包含 E 的一个有限余维数的向量子空间。c) 设 \(E_{\sigma}\) 是赋以弱拓扑的无穷维巴拿赫空间。集合 \(\mathcal{F}(E_{\sigma})\) 在 \(\mathcal{L}(E_{\sigma})\) 中不是开的。
\end{enumerate}

\begin{enumerate}
\def\labelenumi{\arabic{enumi})}
\setcounter{enumi}{1}
\tightlist
\item
  设 E 是无穷维巴拿赫空间，且 \(u \in \mathcal{L}(\mathrm{E})\) 。我们用 \(\mathrm{S}(u)\) 表示如下的复数 \(\lambda\) 所成的集：存在实数 \(\delta > 0\) 和 E 的一个有限余维数子空间 F，使得 \(\| u(x) - \lambda x \| \geqslant \delta \| x \|\) 对所有 \(x \in \mathrm{F}\) 成立。
\end{enumerate}

\begin{enumerate}
\def\labelenumi{\alph{enumi})}
\item
  集合 \(S(u)\) 在 C 中是开的，且包含 u 的谱在 C 中的补集。
\item
  若 \(\lambda \in \mathrm{S}(u)\) ，则自同态 \(u - \lambda 1_{\mathrm{E}}\) 具有闭的像和有限维的核。
\item
  补集 \(\mathrm{F}(u) = \mathbf{C} - \mathrm{S}(u)\) 非空。（证明 \(\mathrm{Sp}_{\mathrm{e}}(u)\) 中一个模最大的元素属于 \(\mathrm{F}(u)\) 。）
\end{enumerate}

\begin{enumerate}
\def\labelenumi{\arabic{enumi})}
\setcounter{enumi}{2}
\tightlist
\item
  设 E 是无穷维复希尔伯特空间。
\end{enumerate}

\begin{enumerate}
\def\labelenumi{\alph{enumi})}
\item
  对合代数 \(\mathcal{C}alk(E)\) 是星代数。
\item
  设 u 是 \(\mathcal{L}(\mathrm{E})\) 的使 \(u^{*}u\) 为紧算子的元素。证明 u 是紧算子。
\item
  设 G 是 \(\mathcal{C}alk(\mathrm{E})\) 中可逆元素所成的群，\(G_{0}\) 是它的单位分量。设 \(p: G \to G/G_{0}\) 是典范投影。存在唯一的群同构 \(\alpha: Z \to G/G_{0}\) ，使得对 E 的每个弗雷德霍姆自同态 u 都有 \(\alpha(\operatorname{ind}(u)) = p(\pi(u))\) 。
\end{enumerate}

\begin{enumerate}
\def\labelenumi{\arabic{enumi})}
\setcounter{enumi}{3}
\tightlist
\item
  设 E 是复希尔伯特空间。设 \(\xi\) 是 \(\mathcal{C}alk(E)\) 的一个可逆元素，我们称 \(\xi\) 的指标为 E 的任一弗雷德霍姆自同态 \(u\) （满足 \(\pi(u) = \xi\) ）的指标。
\end{enumerate}

\begin{enumerate}
\def\labelenumi{\alph{enumi})}
\item
  证明：对每个埃尔米特的 \(\xi \in \mathcal{C}alk(\mathrm{E})\) ，存在埃尔米特的 \(u \in \mathcal{L}(\mathrm{E})\) 使得 \(\pi(u) = \xi\) 。
\item
  证明：对每个酉的且指标为零的 \(\xi \in \mathcal{C}alk(\mathrm{E})\) ，存在酉的 \(u \in \mathcal{L}(\mathrm{E})\) 使得 \(\pi(u) = \xi\) 。（设 \(v\) 可逆且满足 \(\pi(v) = \xi\) ；考察 \(v = j|v|\) 这一 \(v\) 的极分解，验证 \(|v| - 1_{\mathrm{E}}\) 是紧的，再推出 \(v - j\) 是紧的。）
\item
  证明：对每个酉的 \(\xi \in \mathcal{C}alk(\mathrm{E})\) （满足 \(\xi^2 = 1\) ），存在酉的 \(u \in \mathcal{L}(\mathrm{E})\) 使得 \(u^2 = 1\) 且 \(\pi(u) = \xi\) 。
\end{enumerate}

\(d)\) 对每个元素 \(p \in \mathcal{C}alk(\mathrm{E})\) ，只要 \(p^2 = p\) ，就存在投影 \(\mathrm{P} \in \mathcal{L}(\mathrm{E})\) 使得 \(\pi(\mathrm{P}) = p\) 。

\begin{enumerate}
\def\labelenumi{\arabic{enumi})}
\setcounter{enumi}{4}
\tightlist
\item
  设 E 是复希尔伯特空间。
\end{enumerate}

\begin{enumerate}
\def\labelenumi{\alph{enumi})}
\tightlist
\item
  设 \(n\) 是满足 \(n \geqslant 2\) 的整数。我们置 \(\mathrm{I} = \mathbf{N} \times \{0, \ldots, n\}\) 。设 \((\sigma_i)_{1 \leqslant i \leqslant n}\) 是 I 的由下式定义的置换族
\end{enumerate}

\[
\begin{array}{c} \sigma_ {i} (m, i) = (m - 1, i) \text {for} m \geqslant 1 \\ \sigma_ {i} (0, i) = (0, i - 1) \\ \sigma_ {i} (m, i - 1) = (m + 1, i - 1) \text {for} m \geqslant 0 \\ \sigma_ {i} (m, j) = (m, j) \text {for} m \geqslant 0 \text {and} j \notin \{i, i - 1 \}. \end{array}
\]

\begin{enumerate}
\def\labelenumi{\alph{enumi})}
\setcounter{enumi}{1}
\item
  用 \((e_{m,i})\) 表示 \(\ell^{2}(\mathrm{I})\) 的典范标准正交基。设 \(u_{i}\) 是 \(\ell^{2}(\mathrm{I})\) 的自同态，使得 \(u_{i}(e_{m,j}) = e_{\sigma_{i}(m,j)}\) 对所有 \((m,j) \in \mathrm{I}\) 成立。自同态 \(u_{i}\) 是酉的。
\item
  元素 \(\pi(u_i)\) 两两交换。
\end{enumerate}

\(d)\) 不存在 E 的由两两交换的酉自同态组成的族 \((v_{1},\ldots,v_{n})\) ，使得 \(\pi(v_{i})=\pi(u_{i})\) 对所有 \(i\) 成立。

\begin{enumerate}
\def\labelenumi{\arabic{enumi})}
\setcounter{enumi}{5}
\tightlist
\item
  设 E 是复希尔伯特空间。回顾 \(\mathcal{C}alk(E)\) 是星代数（习题 3）。
\end{enumerate}

  设 \(\mathfrak{F}\) 是 \(\mathbf{N}\) 上的一个非主超滤子。设 F 是 E 中沿 \(\mathfrak{F}\) 弱收敛于 0 的有界序列所成的向量空间，\(F_0\) 是 F 中由沿 \(\mathfrak{F}\) 依范数收敛于 0 的序列组成的子空间。我们置 \(\widetilde{\mathrm{H}} = \mathrm{F} / \mathrm{F}_0\) 。a) 对 F 的所有元素 \(x = (x_n)\) 和 \(y = (y_n)\) ，证明序列 \((\langle x_n | y_n \rangle)_{n \in \mathbf{N}}\) 沿 \(\mathfrak{F}\) 收敛。用 \(\langle x | y \rangle_{\mathrm{F}}\) 表示其极限。

\begin{enumerate}
\def\labelenumi{\alph{enumi})}
\setcounter{enumi}{1}
\item
  证明 \((x,y)\mapsto\langle x\mid y\rangle_{\mathrm{F}}\) 是 F 上的正埃尔米特形式，且 \(\langle x\mid x\rangle_{F}=0\) 当且仅当 \(x\in F_{0}\) ；因而这个埃尔米特形式在 \(\widetilde{H}\) 上定义了一个分离的预希尔伯特空间结构。
\item
  设 H 是由分离的预希尔伯特空间 \(\widetilde{H}\) 完备化得到的希尔伯特空间。证明：对 \(\mathcal{L}(E)\) 中的每个 u，都得到一个元素 \(\varrho(u)\) （属于 \(\mathcal{L}(H)\) ），这只需把 F 到 F 的线性映射 \((x_{n})\mapsto(u(x_{n}))\) 过渡到商。d) 映射 \(u\mapsto\varrho(u)\) 诱导 C*-代数的一个等距态射 \(\mathcal{C}alk(E)\to\mathcal{L}(H)\) 。
\end{enumerate}

（我们以后将看到，每个 C*-代数 A 都有一个等距表示 \(A \to \mathcal{L}(E)\) ，其中 E 是某个希尔伯特空间，参见 V, p. 442 的定理 2。）

\begin{enumerate}
\def\labelenumi{\arabic{enumi})}
\setcounter{enumi}{6}
\tightlist
\item
  我们给单位圆 \(S_{1} \subset \mathbf{C}\) 赋以规范化哈尔测度。设 E 是复希尔伯特空间 \(\mathrm{L}^{2}(\mathbf{S}_{1})\) 。对 \(n \in \mathbf{Z}\) ，设 \(e_{n}\) 是 E 的元素 \(z \mapsto z^{n}\) 。设 F ⊂ E 是由诸 \(e_{n}\) （其中 \(n \geqslant 0\) ）生成的闭子空间，并设 \(p: E \to F\) 是正交投影。对 \(f \in L^{\infty}(\mathbf{S}_{1})\) ，我们在 E（相应地，F）上定义线性映射 \(M_{f}\)（相应地，\(T_{f}\)）如下：\(\mathrm{M}_{f}(g) = fg\) （\(g \in E\) ），且 \(T_{f} = p \circ M_{f}\) 。我们称 \(T_{f}\) 为以 f 为符号的特普利茨算子。(a) 线性映射 \(M_{f}\) 与 \(T_{f}\) 连续。\(M_{f}\) 的谱是 f 的傅里叶系数 \(c_{n} = \langle f | e_{n} \rangle\) 所成之集的闭包。自同态 \(M_{f}\) 是弗雷德霍姆自同态，当且仅当它是里斯自同态。
\end{enumerate}

\begin{enumerate}
\def\labelenumi{\alph{enumi})}
\setcounter{enumi}{1}
\item
  自同态 \(T_{e_{n}}\) （其中 \(n \in Z\) ）都是弗雷德霍姆自同态，且 \(\text{ind}(\mathrm{T}_{e_{n}}) = -n\) 对所有 \(n \in Z\) 成立。
\item
  \(\mathrm{T}_{\overline{f}} = \mathrm{T}_f^*\) 对每个函数 \(f\in \mathrm{L}^{\infty}(\mathbf{S}_1)\) 都成立。
\item
  若 \(T_{f}\) 是同构，则 \(M_{f}\) 是同构，且 f 在 \(\mathrm{L}^{\infty}(\mathbf{S}_{1})\) 中可逆。由此推出 \(M_{f}\) 的谱包含于 \(T_{f}\) 的谱，且 \(\|T_{f}\|=\|f\|\) 。
\item
  设 \(f \in \mathcal{C}(\mathbf{S}_1)\) 是连续函数，\(g \in \mathrm{L}^{\infty}(\mathbf{S}_1)\) 。自同态 \(\mathrm{T}_f\mathrm{T}_g - \mathrm{T}_{fg}\) 与 \(\mathrm{T}_g\mathrm{T}_f - \mathrm{T}_{gf}\) 是紧算子。（先考虑 \(f = e_n\) 的情形。）
\item
  设 \(\mathcal{C}alk(\mathrm{F})\) 是 F 的卡尔金代数，\(\pi\colon\mathcal{L}(\mathrm{F})\to\mathcal{C}alk(\mathrm{F})\) 是典范投影。证明 \(f\mapsto\pi(\mathrm{T}_{f})\) 是从 \(\mathcal{C}(\mathbf{S}_{1})\) 到 \(\mathcal{C}alk(\mathrm{F})\) 中的对合代数的单态射。（为证单性，可利用 I, p. 31 的命题 1。）
\end{enumerate}

\(g)\) 设 \(f \in \mathcal{C}(\mathbf{S}_1)\) 。则 \(\mathrm{T}_f\) 是弗雷德霍姆自同态，当且仅当 \(f\) 在 \(\mathbf{S}_1\) 上不取零值。（利用 I, p. 106 的命题 5。）

\(h)\) 设 \(f \in \mathcal{C}(\mathbf{S}_1)\) 且 \(f\) 在 \(\mathbf{S}_1\) 上不取零值。\(\mathrm{T}_f\) 的指标等于 \(n \in \mathbf{Z}\) ，当且仅当道路 \(t \mapsto f(e^{2i\pi t})\) （在 \(\mathbf{C} - \{0\}\) 中）同伦于道路 \(t \mapsto e_n(t)\) （参见 TA, III, p. 229, 定义 1）。

\begin{enumerate}
\def\labelenumi{\arabic{enumi})}
\setcounter{enumi}{7}
\tightlist
\item
  我们沿用上一习题的记号。
\end{enumerate}

\begin{enumerate}
\def\labelenumi{\alph{enumi})}
\item
  用 \(\mathrm{T}(\mathbf{S}_{1})\) 表示 \(\mathcal{L}(\mathrm{F})\) 中由特普利茨算子 \(T_{e_{1}}\) 与 \(T_{e_{-1}}\) 生成的闭子代数。证明 \(\mathrm{T}(\mathbf{S}_{1})\) 是 \(\mathcal{L}(\mathrm{F})\) 的星子代数。
\item
  设 \(f \in \mathcal{C}(\mathbf{S}_1)\) 。证明 \(T_f\) 属于 \(T(\mathbf{S}_1)\) 。
\item
  \(\mathrm{T}(\mathbf{S}_{1})\) 在 F 上的表示是不可约的；星代数 \(\mathrm{T}(\mathbf{S}_{1})\) 包含 \(\mathcal{L}^{\mathrm{c}}(\mathrm{F})\) 。
\end{enumerate}

\(d)\) \(u \in \mathrm{T}(\mathbf{S}_1)\) 当且仅当存在 \(f \in \mathcal{C}(\mathbf{S}_1)\) 与 \(h \in \mathcal{L}^c(\mathrm{F})\) 使得 \(u = \mathrm{T}_f + h\) 。（证明形如 \(\mathrm{T}_f + h\) 的自同态所成之集在 \(\mathcal{L}(\mathrm{F})\) 中构成一个闭星代数。）

\begin{enumerate}
\def\labelenumi{\alph{enumi})}
\setcounter{enumi}{4}
\tightlist
\item
  存在对合代数态射的一个正合列
\end{enumerate}

\[
0 \longrightarrow \mathcal {L} ^ {\mathrm{c}} (\mathrm{F}) \longrightarrow \mathrm{T} (\mathbf {S} _ {1}) \stackrel {\varrho} {\longrightarrow} \mathcal {C} (\mathbf {S} _ {1}) \longrightarrow 0
\]

使得 \(\varrho(\mathrm{T}_f) = f\) 对每个函数 \(f \in \mathcal{C}(\mathbf{S}_1)\) 成立。

\begin{enumerate}
\def\labelenumi{\arabic{enumi})}
\setcounter{enumi}{8}
\item
  设 B 是无穷维巴拿赫空间。用 E 表示 B 的强对偶，用 F 表示赋以拓扑 \(\sigma(\mathrm{B}^{\prime},\mathrm{B})\) 的 B 的对偶。证明：存在从 E 到 F 中的连续双射线性映射 u，以及从 E 到 F 中的紧线性映射 h，使得 \(u + h\) 有无穷维的余核。
\item
  设 E 是分离的局部凸空间，u 与 h 是 \(\mathcal{L}(\mathrm{E})\) 中交换的元素。假设存在整数 n 使得 \(h^{n}\) 是紧自同态。
\end{enumerate}

\begin{enumerate}
\def\labelenumi{\alph{enumi})}
\item
  若 u 是弗雷德霍姆自同态，则 u-h 是弗雷德霍姆自同态。
\item
  若 \(u\) 是里斯自同态，则 \(u - h\) 是里斯自同态。
\end{enumerate}

\begin{enumerate}
\def\labelenumi{\arabic{enumi})}
\setcounter{enumi}{10}
\tightlist
\item
  设 \(\mathrm{E} = \ell^2 (\mathbf{N})\) ，\((e_n)\) 是 E 的典范基。我们定义 \(u\in \mathcal{L}(\mathrm{E})\) 为满足 \(u(e_i) = e_{i + 1}\) （\(i\in \mathbf{N}\) ）的自同态。设 \(\mathrm{F} = \mathrm{E}\oplus \mathrm{E}\) ，\(v = u\oplus (u^{*} + 2\cdot 1_{\mathrm{E}})\in \mathcal{L}(\mathrm{F})\) 。设 \(\overline{v} = \pi (v)\) 是 \(v\) 在卡尔金代数 \(\mathcal{Calk}(\mathrm{F})\) 中的像。
\end{enumerate}

\begin{enumerate}
\def\labelenumi{\alph{enumi})}
\item
  自同态 \(v\) 是指标为 -1 的弗雷德霍姆自同态，而 \(v - 2 \cdot 1_{\mathrm{E}}\) 是指标为 1 的弗雷德霍姆自同态。
\item
  设 \(p = \mathrm{X}(\mathrm{X} - 2) \in \mathbf{C}[\mathrm{X}]\) ，S 是 \(\overline{v}\) 的指数谱（I, p. 173 的习题 5）。我们有 \(0 \in p(\mathrm{S})\) ，但 0 不属于 \(p(\overline{v})\) 的指数谱。（利用习题 3。）
\end{enumerate}

\exercisesection{}

\begin{enumerate}
\def\labelenumi{\arabic{enumi})}
\tightlist
\item
  设 A 是 \(\mathbf{C}\) 的紧子集，它等于其内部 V 的闭包。我们置 \(\mathcal{C}(\mathrm{A}) = \mathcal{C}(\mathrm{A};\mathbf{C})\) 。设 \(B(A)\) 是 \(\mathcal{C}(\mathrm{A})\) 中由在 V 上的限制为全纯的函数 \(f\in \mathcal{C}(\mathrm{A})\) 所组成的向量子空间；
\end{enumerate}

它是 \(\mathcal{C}(\mathrm{A})\) 的一个巴拿赫子代数。设 \(u \in \mathcal{L}(\mathrm{B}(\mathrm{A}))\) 是满足 \(u(f)(z) = zf(z)\) （对所有 \(z \in \mathrm{A}\) 及每个 \(f \in \mathrm{B}(\mathrm{A})\) ）的线性映射。

\begin{enumerate}
\def\labelenumi{\alph{enumi})}
\item
  \(u\) 的谱等于 A。
\item
  对每个 \(z \in V\) ，B(A) 的自同态 u - z 是指标为 1 的弗雷德霍姆自同态。
\item
  对每个 \(z \in A - V\) ，B(A) 的自同态 u - z 不是弗雷德霍姆自同态。
\item
  u 的稳定谱是 A 的边界，u 的敏感谱是空集，而 u 的本质谱是 A。
\item
  我们假设 A 在 \(\mathbf{C}\) 中的余集是连通的。对每个 \(h \in \mathcal{L}^{c}(\mathrm{B}(\mathrm{A}))\) ，\(u + h\) 的本质谱等于 A。
\end{enumerate}

\begin{enumerate}
\def\labelenumi{\arabic{enumi})}
\setcounter{enumi}{1}
\tightlist
\item
  我们保留上一习题的记号。假设 A 包含单位圆 \(U\) （\(\mathbf{C}\) 中的），且 \(0 \notin A\) 。用 \(\varphi_{A}\) 表示 A 的特征函数。
\end{enumerate}

\begin{enumerate}
\def\labelenumi{\alph{enumi})}
\tightlist
\item
  映射
\end{enumerate}

\[
\ell \colon f \mapsto \frac {1}{2 i \pi} \int_ {\mathbf {U}} f (z) d z
\]

是 B(A) 上的一个连续线性形式。

\begin{enumerate}
\def\labelenumi{\alph{enumi})}
\setcounter{enumi}{1}
\tightlist
\item
  设 \(h\) 是 B(A) 的有限秩自同态，满足 \(h(f) = \ell(f)\varphi_{\mathrm{A}}\) 。\(u - h\) 的本质谱异于 \(u\) 的本质谱。（证明它包含 0。）
\end{enumerate}

\begin{enumerate}
\def\labelenumi{\arabic{enumi})}
\setcounter{enumi}{2}
\tightlist
\item
  设 E 是复希尔伯特空间。设 U 是 \(\mathbf{C}\) 的连通开子集。设 f 是从 U 到空间 \(\mathcal{L}^{c}(\mathrm{E})\) 中的全纯函数。假设存在 \(z \in U\) 使得 1 属于 \(f(z)\) 的预解集。
\end{enumerate}

\begin{enumerate}
\def\labelenumi{\alph{enumi})}
\item
  由 \(z \in U\) （满足 \(1 \in \operatorname{Sp}(f(z))\) ）组成的集 D 是 U 的一个离散子集。
\item
  从 U-D 到 \(\mathcal{L}(\mathrm{E})\) 中的由 \(z\mapsto \mathrm{R}(f(z),1)\) 定义的映射是全纯的。
\end{enumerate}

我们固定一点 \(z_0 \in \mathrm{D}\) 。

\begin{enumerate}
\def\labelenumi{\alph{enumi})}
\setcounter{enumi}{2}
\item
  存在整数 \(k \geqslant 1\) ，使得从 U-D 到 \(\mathcal{L}(\mathrm{E})\) 中的由 \(z \mapsto (z - z_{0})^{k}\mathrm{R}(f(z), 1)\) 定义的映射可延拓为 \((\mathrm{U}-\mathrm{D}) \cup \{z_{0}\}\) 上的一个全纯映射。
\item
  存在 E 的自同态的唯一序列 \((u_{n})_{n\in\mathbf{N}}\) （当 n 充分大时 \(u_{n}=0\) ），以及从 \((\mathrm{U}-\mathrm{D})\cup\{z_{0}\}\) 到 \(\mathcal{L}(\mathrm{E})\) 中的唯一全纯函数 S，使得对所有 \(z\in U-D\) 有
\end{enumerate}

\[
\mathrm{R} (f (z), 1) = \sum_ {j = - k} ^ {- 1} (z - z _ {0}) ^ {j} u _ {j} + \mathrm{S} (z).
\]

\begin{enumerate}
\def\labelenumi{\alph{enumi})}
\setcounter{enumi}{4}
\tightlist
\item
  自同态 \(u_{-1}\) 是 \(f(z_0)\) 相对于 \(\{1\}\) （\(f(z_0)\) 的谱的开且闭子集）的谱投影。
\end{enumerate}

\begin{enumerate}
\def\labelenumi{\arabic{enumi})}
\setcounter{enumi}{3}
\tightlist
\item
  设 E 是可数型的巴拿赫空间。设 \(u \in \mathcal{L}(\mathrm{E})\) 。假设序列 \((\| u^k \|)_{k \in \mathbf{N}}\) 有界。
\end{enumerate}

对每个 \(\lambda\in\mathrm{Sp}(u)\cap\mathbf{U}\) ，设 \(x_{\lambda}\in\mathrm{E}\) 是 \(u\) 的属于 \(\lambda\) 的范数为 1 的特征向量 。

\begin{enumerate}
\def\labelenumi{\alph{enumi})}
\tightlist
\item
  对所有 \(\lambda_1\) 、\(\lambda_2\) （属于 \(\mathrm{Sp}(u) \cap \mathbf{U}\) ）以及每个整数 \(k \in \mathbf{N}\) ，我们有
\end{enumerate}

\[
\left| \lambda_ {1} ^ {k} - \lambda_ {2} ^ {k} \right| \leqslant \left(1 + \| u ^ {k} \|\right) \| x _ {\lambda_ {1}} - x _ {\lambda_ {2}} \|.
\]

\begin{enumerate}
\def\labelenumi{\alph{enumi})}
\setcounter{enumi}{1}
\item
  存在实数 c \textgreater{} 0 ，使得 \(\|x_{\lambda_{1}} - x_{\lambda_{2}}\| \geqslant c\) 对所有 \(\lambda_{1} \neq \lambda_{2}\) （\(\operatorname{Sp}(u) \cap \mathbf{U}\) 中的）成立。
\item
  集合 \(\operatorname{Sp}(u) \cap \mathbf{U}\) 是可数的（“贾米森定理”）。
\item
  若 E 不是可数型的，贾米森定理一般不成立。
\end{enumerate}

\begin{enumerate}
\def\labelenumi{\arabic{enumi})}
\setcounter{enumi}{4}
\tightlist
\item
  本习题的目标是在空间 E 为巴拿赫空间的情形给出洛莫诺索夫定理（III, p. 13 的推论 2）的另一个证明。
\end{enumerate}

设 E 是维数至少为 2 的复巴拿赫空间。设 h 是 E 的一个非零紧自同态，A 是 h 在 \(\mathcal{L}(\mathrm{E})\) 中的交换子。a) 下列条件之一成立：

\begin{enumerate}
\def\labelenumi{(\roman{enumi})}
\item
  存在 \(y \in E\) ，使得 Ay 是 E 中的一个非零且不稠密的子空间；
\item
  对 A 的每个内部非空的闭有界子集 C，存在 A 的有限子集 I，使得 \(h(\mathrm{C})\) 包含于诸集合 \(u^{-1}(\mathrm{C})\) （\(u \in I\) ）的并之中。
\end{enumerate}

\begin{enumerate}
\def\labelenumi{\alph{enumi})}
\setcounter{enumi}{1}
\tightlist
\item
  假设上一问题的条件 (ii) 成立。设 C 是 E 的一个内部非空的闭有界子集。对每个整数 \(m \geqslant 1\) ，存在 A 中的有限序列 \((u_{1}, \ldots, u_{m})\) 使得
\end{enumerate}

\[
u _ {1} \circ \dots \circ u _ {m} \circ h ^ {m} (x _ {0}) \in \mathrm{C}.
\]

（打乒乓球。）

\begin{enumerate}
\def\labelenumi{\alph{enumi})}
\setcounter{enumi}{2}
\item
  在同一假设下，证明 h 的谱半径非零（在相反的情形，考虑一个元素 \(x_{0} \in E\) ，满足 \(\|h(x_{0})\| > \|h\|\) ，以及以 \(x_{0}\) 为心、\(\|h\|\) 为半径的闭球 C；注意到 \(\|( \lambda h)^{n}\|\) 对每个 \(\lambda \in R\) 趋于 0，证明上一问题的结论将蕴含 0 属于 C 的闭包）。
\item
  对 E 的每个与 E 的某个非零紧自同态交换的自同态 \(u\) ，存在异于 E 的非零闭子空间 \(F \subset E\) 使得 \(u(F) \subset F\) 。
\end{enumerate}

\begin{enumerate}
\def\labelenumi{\arabic{enumi})}
\setcounter{enumi}{5}
\tightlist
\item
  设 A 是一个幺 Z-代数。假设 A 是自由 Z-模。设 B 是 A 的一个子集，它作为 Z-模是 A 的基。我们称对子 \((A, B)\) 是一个自然代数，如果 A 关于 B 的结构常数（A, III, p. 10）属于 N，且 \(1_{A}\) 在基 B 中的系数属于 N。此外，我们称 (A, B) 是幺的，如果 \(1_{A}\) 是 B 的元素。
\end{enumerate}

设 (A, B) 是一个自然代数。设 \(B_{0}\) 是由 \(b \in B\) （使 b 在 \(1_{A}\) 中的系数非零）组成的集。用 \(\tau_{A}\) 表示从 A 到 Z 的唯一的 Z-模同态，使得 \(\tau_{\mathrm{A}}(b) = 1\) （若 \(b \in B_{0}\) ），否则 \(\tau_{\mathrm{A}}(b) = 0\) 。

\begin{enumerate}
\def\labelenumi{\alph{enumi})}
\item
  \(\tau_{\mathrm{A}}(xy) = \tau_{\mathrm{A}}(yx)\) 对所有 \((x,y)\in \mathrm{A}^2\) 成立。
\item
  对所有 \(b \neq b'\) （在 \(B_{0}\) 中），有 \(b^2 = b\) 且 \(bb' = 0\) 。
\end{enumerate}

\begin{enumerate}
\def\labelenumi{\arabic{enumi})}
\setcounter{enumi}{6}
\tightlist
\item
  设 (A, B) 是一个自然代数。
\end{enumerate}

\begin{enumerate}
\def\labelenumi{\alph{enumi})}
\tightlist
\item
  至多存在一个对合 \(\iota\) （\(\mathcal{B}\) 的），使得从 A 到 A 的唯一 \(\mathbf{Z}\)-模同态（把 \(b\) 映为 \(\iota(b)\) ）是 A 的一个反自同构，且满足条件 \(\tau_{\mathrm{A}}(bb') = 1\) （若 \(b' = \iota(b)\) ），否则 \(\tau_{\mathrm{A}}(bb') = 0\) 。
\end{enumerate}

当这样一个对合 \(\iota\) 存在时，我们称 (A, \(\mathcal{B}\) ) 是一个自然基础代数。

我们现在假设 (A, B) 是一个自然基础代数。

\begin{enumerate}
\def\labelenumi{\alph{enumi})}
\setcounter{enumi}{1}
\tightlist
\item
  于是 \(\iota(b) = b\) 对所有 \(b \in B_{0}\) 成立，并有公式
\end{enumerate}

\[
1 _ {\mathrm{A}} = \sum_ {b \in \mathcal {B} _ {0}} b.
\]

\begin{enumerate}
\def\labelenumi{\alph{enumi})}
\setcounter{enumi}{2}
\tightlist
\item
  A 的结构常数 \(\gamma_{bb'}^{b''}\) 满足 \(\gamma_{bb'}^{\iota(b'')} = \tau_{\mathrm{A}}(bb'b'')\) （对 \((b, b', b'')\) ，\(\mathcal{B}^{3}\) 中的）。
\end{enumerate}

\begin{enumerate}
\def\labelenumi{\arabic{enumi})}
\setcounter{enumi}{7}
\tightlist
\item
  设 (A, B) 是一个自然基础代数，且 A 的秩有限。
\end{enumerate}

\begin{enumerate}
\def\labelenumi{\alph{enumi})}
\tightlist
\item
  对每个 \(x \in A\) ，元素
\end{enumerate}

\[
\sum_ {b \in \mathcal {B}} b x \iota (b)
\]

是中心元。

在本习题的其余部分，我们进一步假设 (A, \(\mathcal{B}\) ) 是幺的。

\begin{enumerate}
\def\labelenumi{\alph{enumi})}
\setcounter{enumi}{1}
\item
  对所有 \(b\) 与 \(b'\) （\(\mathcal{B}\) 中的），存在 \(b_1\) 与 \(b_2\) （\(\mathcal{B}\) 中的），使得 \(b'\) 在 \(bb_1\) 中以及在 \(b_2b\) 中的系数非零。
\item
  对每个 \(x \in B\) ，在基 B 中表示从 A 到 A 的映射 \(y \mapsto xy\) 的矩阵是一个实矩阵，其所有系数都是正的。用 pf 表示从 A 到 R 的唯一同态，使得 \(\operatorname{pf}(b)\) （当 \(b \in B\) 时）是这个矩阵的谱半径。
\end{enumerate}

\(d)\) 对每个 \(b \in \mathcal{B}\) ，实数 \(\alpha = \mathrm{pf}(b)\) 是一个代数整数，满足 \(\alpha \geqslant 1\) ；对每个共轭数 \(\alpha' \in \mathbf{C}\) （\(\alpha\) 的），有 \(|\alpha'| \leqslant \alpha\) 。

\begin{enumerate}
\def\labelenumi{\alph{enumi})}
\setcounter{enumi}{4}
\tightlist
\item
  设 \(a_0 \in \mathrm{A}\) 是 \(\mathcal{B}\) 的诸元素之和。表示映射 \(f: y \mapsto ya_0\) 的矩阵在基 \(\mathcal{B}\) 中的所有系数都严格为正。存在一个元素 \(r \in \mathbf{C} \otimes \mathrm{A}\) ，它在基 \((1 \otimes b)_{b \in \mathcal{B}}\) 中的系数都严格为正，使得 \(f_{\mathbf{C}}(r) = \varrho(f_{\mathbf{C}})r\) 。
\end{enumerate}

\(f)\) 对每个 \(x \in \mathrm{A}\) ，存在唯一的复数 \(\delta(x)\) ，使得 \(xr = \delta(x)r\) 在 \(\mathbf{C} \otimes \mathrm{A}\) 中成立。映射 \(\delta: \mathrm{A} \to \mathbf{C}\) 是 \(\mathrm{A}\) 的一个特征标。

\begin{enumerate}
\def\labelenumi{\alph{enumi})}
\setcounter{enumi}{6}
\tightlist
\item
  \(\delta(x) = \mathrm{pf}(x)\) 对所有 \(x \in \mathrm{A}\) 成立。特别地，映射 \(x \mapsto \mathrm{pf}(x)\) 是 A 的一个特征标。
\end{enumerate}

\(h)\) 映射 pf 是 A 的唯一特征标，使得 \(\mathrm{pf}(b) \geqslant 0\) 对所有 \(b \in \mathcal{B}\) 成立；此外，\(\mathrm{pf}(b) > 0\) 对所有 \(b \in \mathcal{B}\) 成立。

\begin{enumerate}
\def\labelenumi{\roman{enumi})}
\tightlist
\item
  对每个 \(y \in \mathrm{A}\) ，有 \(ry = \delta(y)r\) （在 \(\mathbf{C} \otimes \mathbf{A}\) 中）。
\end{enumerate}

\begin{enumerate}
\def\labelenumi{\alph{enumi})}
\setcounter{enumi}{9}
\tightlist
\item
  元素
\end{enumerate}

\[
r _ {0} = \sum_ {b \in \mathcal {B}} \mathrm{pf} (b) b
\]

具有问题 i) 所要求的性质。

\begin{enumerate}
\def\labelenumi{\arabic{enumi})}
\setcounter{enumi}{8}
\tightlist
\item
  \begin{enumerate}
  \def\labelenumii{\alph{enumii})}
  \tightlist
  \item
    设 \(n\) 是整数，\(M\) 是系数属于 \(\mathbf{N}\) 的一个方阵。假设 \(\varrho(M^t M) = \varrho(M)^2\) 且 \(|\varrho(M)| < 2\) 。则存在整数 \(m \geqslant 2\) 使得 \(\varrho(M) = 2\cos(\pi/m)\) 。（利用克罗内克定理，参见 A, V, p. 155, 习题 17。）
  \end{enumerate}
\end{enumerate}

\begin{enumerate}
\def\labelenumi{\alph{enumi})}
\setcounter{enumi}{1}
\item
  对每个整数 \(m \geqslant 2\) ，存在一个系数属于 N 的方阵 M，使得 \(\varrho(M) = 2 \cos(\pi/m)\) 。
\item
  设 (A, B) 是一个秩有限的幺自然基础代数。设 \(b \in B\) 。若 \(\operatorname{pf}(b) < 2\) ，则存在整数 \(m \geqslant 3\) 使得 \(\operatorname{pf}(b) = 2 \cos(\pi/m)\) 。
\end{enumerate}

习题 10 至 24 给出 III, p. 100 的推论 5 的另一个证明以及若干强化。其中将用到下列概念。

设 I 是有限集，n 是它的基数。

我们用 \(1_{\mathrm{I}}\) 表示 (I, I) 型的实单位矩阵。

- 一个 (I, I) 型的实系数方阵称为正的，如果它的所有系数都 \(\geqslant 0\) ；称为严格正的，如果它的所有系数都 \(>0\) 。

- (I,I) 型的实系数方阵 \((a_{ij})_{i,j\in \mathrm{I}}\) 称为可约的，如果存在把 I 分成两个非空子集 I' 与 I'\,' 的划分，使得 \(a_{ij} = 0\) 对所有 \((i,j)\in \mathrm{I}'\times \mathrm{I}''\) 成立；若它不是可约的，则称为不可约的。

- 我们把 \(\mathbf{R}^{\mathrm{I}}\) 的由典范基向量的某个子集生成的任一向量子空间称为 \(\mathbf{R}^{\mathrm{I}}\) 的坐标子空间。坐标子空间共有 \(2^{n}\) 个；对每个整数 \(k\in\{0,\ldots,n\}\) ，其中维数为 \(k\) 的坐标子空间有 \(\binom{n}{k}\) 个。

因此，(I, I) 型的实矩阵是不可约的，当且仅当它所稳定的坐标子空间只有 \(\{0\}\) 和 \(R^{I}\) 。

对所有 \(x\) 与 \(y\) （\(\mathbf{R}^{\mathrm{I}}\) 中的），我们记 \(x \preccurlyeq y\) （及 \(y \succcurlyeq x\) ），如果 \(y - x\) 属于 \(\mathbf{R}^{\mathrm{I}}\) 中坐标为正的元素所成的锥，亦即 \(x_i \leqslant y_i\) 对所有 \(i \in \mathrm{I}\) 成立。我们记 \(x \prec y\) （及 \(y \succ x\) ），如果 \(y - x\) 的所有坐标都严格为正。

\begin{enumerate}
\def\labelenumi{\arabic{enumi})}
\setcounter{enumi}{9}
\item
  设 \(A\) 是 (I, I) 型的正且不可约的实矩阵。对每个整数 \(n\) （满足 \(n \geqslant \operatorname{Card}(\mathrm{I})\) ），矩阵 \((1_{\mathrm{I}} + A)^{n-1}\) 是严格正的。（设 \(y \in \mathbf{R}_{+}^{n}\) ，\(z = (1_{\mathrm{I}} + A)y\) ；若 \(y \neq 0\)，证明由 \(i \in \mathrm{I}\) （使 \(z_{i} = 0\) ）所成的集严格包含于 \(\{i \in \mathrm{I} \mid y_{i} = 0\}\) 。）
\item
  设 \(A\) 是 (I, I) 型的正且不可约的实矩阵。用 \(\psi \in \mathbf{R}[\mathrm{X}]\) 表示 \(A\) 的极小多项式，并置 \(m = \deg(\psi)\) 。
\end{enumerate}

对每个整数 \(q\)，置 \(A^q = (a_{ij}^{(q)})_{i,j \in \mathrm{I}}\) 。

对每个 \((i,j)\in\mathrm{I}^{2}\) ，存在整数 q 使得 \(1\leqslant q\leqslant m\) 且 \(a_{ij}^{(q)}>0\) 。（设 r 是 \((1+\mathrm{X})^{n-1}\) 除以 \(\psi\) 的欧几里得除法的余式；考察矩阵 \(r(A)\) 并注意它是严格正的。）

\begin{enumerate}
\def\labelenumi{\arabic{enumi})}
\setcounter{enumi}{11}
\tightlist
\item
  设 A 是 (I, I) 型的不可约且正的实矩阵。对每个非零的 \(x \in R^{I}\) ，我们置
\end{enumerate}

\[
r_{x} = \inf_{\substack{i\in \mathrm{I}\\ x_{i}\neq 0}}\frac{(Ax)_{i}}{x_{i}},
\]

其中 \((Ax)_i = \sum_{k\in I}a_{ik}x_k\) 是 \(Ax\) 的以 \(i\) 为指标的坐标。

\begin{enumerate}
\def\labelenumi{\alph{enumi})}
\item
  若 \(x \succcurlyeq 0\) 且 \(x \neq 0\) ，则 \(r_x\) 是最大的实数 \(\varrho\) ，使得 \(\varrho x \preccurlyeq Ax\) 。
\item
  假设 \(x \succcurlyeq 0\) ，并置 \(y = (1_{\mathrm{I}} + A)x\) 。则有 \(r_x \leqslant r_y\) 。
\item
  限制在 \((\mathbf{R}_{+}^{*})^{\mathrm{I}}\) 上时，函数 \(x \mapsto r_{x}\) 是连续的。
\end{enumerate}

\(d)\) 由此推出存在 \(z \in \mathbf{R}^{\mathrm{I}}\) 满足 \(z \succcurlyeq 0\) 且 \(r_z = \sup_{x \succcurlyeq 0} r_x\) （利用紧集 \(\mathrm{M} = \{x \in \mathbf{R}_{+}^{\mathrm{I}} \mid \sum_{i \in \mathrm{I}} x_i = 1\}\) 以及 \((1_{\mathrm{I}} + A)^{n-1}(\mathrm{M})\) ，再用习题 10。）

我们置 \(r = r_{z}\) 。\(R^{I}\) 中满足 \(z \succcurlyeq 0\) 且 \(r_{z} = r\) 的元素 z 称为极端元素。

\begin{enumerate}
\def\labelenumi{\alph{enumi})}
\setcounter{enumi}{4}
\item
  我们有 r \textgreater{} 0 。
\item
  若 z 是极端向量，则 Az = rz。（若 \(u = Az - rz \succcurlyeq 0\) 非零，则由习题 10，\((1_{\mathrm{I}} + A)^{n-1}u\) 将为 \(\succ 0\) ，且对 \(x = (1_{\mathrm{I}} + A)^{n-1}z\) 将有 \(Ax \succ rx\) 。）
\item
  每个极端元素都 \(\succ 0\) 。
\end{enumerate}

对每个向量 \(y\) （在 \(\mathbf{C}^{\mathrm{I}}\) 中），用 \(y^{+}\) 表示向量 \((|y_{i}|)_{i\in \mathrm{I}}\) ；于是有 \(y^{+}\in \mathbf{R}^{\mathrm{I}}\) 且 \(y^{+}\succcurlyeq 0\) 。

\begin{enumerate}
\def\labelenumi{\alph{enumi})}
\setcounter{enumi}{7}
\item
  设 \(\alpha \in \mathbf{C}\) 是把 \(A\) 视为复矩阵时的一个特征值，\(y\) 是 \(\mathbf{C}^{\mathrm{I}}\) 中相应的一个特征向量。我们有 \(|\alpha|y^{+} \preccurlyeq Ay^{+}\) 且 \(|\alpha| \leqslant r\) 。
\item
  推出：若 \(\alpha = r\) ，则 y 的所有坐标都非零，进而 A 的属于特征值 r 的特征空间（在 \(C^{I}\) 中）的维数为 1。
\end{enumerate}

用 \(\Delta(\mathrm{X})\) 表示 \(\mathrm{X} \cdot 1_{\mathrm{I}} - A\) 的行列式，用 \(B(\mathrm{X})\) 表示 \(\mathrm{X} \cdot 1_{\mathrm{I}} - A\) 的余子式矩阵的转置。我们记 \(B(\mathrm{X}) = (B_{ij}(\mathrm{X}))_{i,j \in \mathrm{I}}\) 。于是有 \(B(\mathrm{X})(\mathrm{X} \cdot 1_{\mathrm{I}} - A) = (\mathrm{X} \cdot 1_{\mathrm{I}} - A)B(\mathrm{X}) = \Delta(\mathrm{X}) \cdot 1_{\mathrm{E}}\) （A, III, p. 99, 命题 12）。

\begin{enumerate}
\def\labelenumi{\alph{enumi})}
\setcounter{enumi}{9}
\tightlist
\item
  对每个 \(i \in I\) ，\(B(r)\) 的以 i 为指标的列非零。（若不然，构造 A 的属于特征值 r 的一个特征向量，其第 i 个系数为零，这与 f) 矛盾。）
\end{enumerate}

\(k\) ) \(B(r)\) 的所有系数都有相同的符号。（首先利用关系式 \((r1_{\mathrm{I}} - A)B(r) = 0\) 证明：对所有 \(i\in \mathrm{I}\) ，以 \(i\) 为指标的 \(B(r)\) 的列的所有系数非零且同号，然后再对转置进行论证。）

\(l)\) \(\Delta'(r) > 0\) 且 \(B_{ij}(r) > 0\) 对所有 \((i,j) \in \mathrm{I}^2\) 成立（利用公式 \(\Delta'(X) = \sum_i B_{ii}(X)\) ，并注意 \(\Delta'\) 在区间 \([r, +\infty]\) 上不变号。）

\begin{enumerate}
\def\labelenumi{\alph{enumi})}
\setcounter{enumi}{12}
\tightlist
\item
  由此得出 r 是 A 的特征多项式的单根。
\end{enumerate}

\begin{enumerate}
\def\labelenumi{\arabic{enumi})}
\setcounter{enumi}{12}
\tightlist
\item
  设 A（相应地，C）是 (I, I) 型的实（相应地，复）方阵。假设 A 是正的且不可约，并且 \(C^{+} \preccurlyeq A\) ，这里与习题 12 一样，用 \(C^{+}\) 表示其系数为 C 的系数的绝对值的矩阵，而记号 \(C^{+} \preccurlyeq A\) 表示 \(A - C^{+}\) 的系数都是正的。
\end{enumerate}

我们沿用上一习题的记号。

\begin{enumerate}
\def\labelenumi{\alph{enumi})}
\item
  设 \(\gamma \in \mathbf{C}\) 是 \(C\) 的一个特征值，\(y\) 是 \(\mathbf{C}^{\mathrm{I}}\) 中相应的一个特征向量。我们有 \(|\gamma|y^{+} \preccurlyeq C^{+}y^{+} \preccurlyeq Ay^{+}\) 且 \(|\gamma| \leqslant r_{y^{+}} \leqslant r\) 。
\item
  我们进一步假设 \(|\gamma| = r\) 。则 \(y^{+}\) 是 A 的一个极端向量（参见习题 12 的问题 f)），且有关系式 \(Ay^{+} = C^{+}y^{+} = ry^{+}\) 与 \(A = C^{+}\) 。
\end{enumerate}

设 \(u\)（相应地，\(u_j\) ，\(j \in I\) ）是使 \(\gamma = |\gamma|u\) （相应地，\(y_j = |y_j|u_j\) ，\(j \in I\) ）成立的唯一复数。用 \(D\) 表示 (I, I) 型的复方阵，其对角系数为 \((u_j)_{j \in I}\) ，其余系数为零，于是 \(y = Dy^+\)。设 \(F = u^{-1}D^{-1}CD\) 。

\begin{enumerate}
\def\labelenumi{\alph{enumi})}
\setcounter{enumi}{2}
\item
  有 \(Fy^{+} = ry^{+}\) 且 \(F^{+}y^{+} = Fy^{+}\) （注意 \(F^{+} = C^{+}\) 。）
\item
  有 \(F = F^{+}\) 且 \(C = uDAD^{-1}\) 。
\end{enumerate}

\begin{enumerate}
\def\labelenumi{\arabic{enumi})}
\setcounter{enumi}{13}
\tightlist
\item
  设 A 是 (I, I) 型的正且不可约的实方阵。我们沿用习题 12 的记号。
\end{enumerate}

设 \((\lambda_{h})_{h\in\mathrm{H}}\) 是 A 的模为 r 的复特征值按重数重复而得的族；置 \(v_{h}=\lambda_{h}/r\) （对所有 \(h\in H\) ）。

设 \(h \in \mathrm{H}\) ；把习题 13 应用于 \(C = A\) 与 \(\gamma = \lambda_h\) ，存在矩阵 \(D_{(h)}\) 使得 \(A = v_h D_{(h)} AD_{(h)}^{-1}\) 且 \(D_{(h)}^+ = 1_{\mathrm{I}}\) 。

\begin{enumerate}
\def\labelenumi{\alph{enumi})}
\item
  设 \(z \succ 0\) 是 A 的一个极端向量。对 \(h \in H\) ，向量 \(y_{(h)} = D_{(h)}z\) 是 A 的属于特征值 \(\lambda_{h}\) 的特征向量。
\item
  证明：对 \(h \in H\) ，A 的特征值 \(\lambda_{h}\) 是单重的。由此推出矩阵 \(D_{(h)}\) 与向量 \(y_{(h)}\) 在相差一个纯量因子的意义下唯一确定。
\item
  证明：对 \(H^{2}\) 中所有 h 与 k，有 \(A = u_{h}u_{k}^{-1}D_{(h)}D_{(k)}^{-1}AD_{(k)}D_{(h)}^{-1}\) ，从而向量 \(D_{(h)}D_{(k)}^{-1}z\) 是 A 的属于特征值 \(u_{h}u_{k}^{-1}r\) 的特征向量。
\item
  由此推出 \(\{u_h\}_{h\in \mathrm{H}}\) 是 \(\mathbf{C}^*\) 的一个子群，且存在非零复数的一个族 \((\mu_h)_{h\in \mathrm{H}}\) ，使得 \(\{\mu_h^{-1}D_{(h)}\}_{h\in \mathrm{H}}\) 是 (I,I) 型复系数可逆矩阵所成的群的一个子群，并且这两个子群都是阶为 Card(H) 的循环群。（选取一个指标 \(i\in \mathrm{I}\) ，并把 \(\mathrm{D}_{(h)}\) 规范化，使其行指标与列指标均为 \(i\) 的系数等于 1。）
\end{enumerate}

在本习题的其余部分，我们假设 H = Z/dZ，且存在 \(\omega \in C^{*}\) 使得 \(\omega^{d} = 1\) 且 \(u_{h} = \omega^{h}\) （\(h \in H\) ）。置 \(D = D_{(1)}\) ，则 \(D^{d} = 1_{I}\) 且 \(D_{(h)} = D^{h}\) 对所有 \(h \in H\) 成立。

\begin{enumerate}
\def\labelenumi{\alph{enumi})}
\setcounter{enumi}{4}
\tightlist
\item
  由此推出 \(A = \omega DAD^{-1}\) 。
\end{enumerate}

我们考虑 I 的划分 \((\mathrm{I}_{h})_{h\in\mathrm{H}}\) ，使得对 \(i\in I\) ，D 的以 \((i,i)\) 为指标的系数等于 \(\omega^{h}\) 当且仅当 \(i\in I_{h}\) 。对每个 h，置 \(n_{h}=\mathrm{Card}(I_{h})\) 。矩阵 D 写成“分块”形式

\[
D = \left( \begin{array}{c c c c} 1 _ {n _ {0}} & & & \\ & \omega 1 _ {n _ {1}} & & \\ & & \ddots & \\ & & & \omega^ {d - 1} 1 _ {n _ {d - 1}} \end{array} \right)
\]

我们考虑 \(A\) 的相应的分解：

\[
A = \left( \begin{array}{c c c c} A _ {0, 0} & A _ {0, 1} & \dots & A _ {0, d - 1} \\ A _ {1, 0} & A _ {1, 1} & \dots & A _ {1, d - 1} \\ \vdots & \vdots & \ddots & \vdots \\ A _ {d - 1, 0} & A _ {d - 1, 1} & \dots & A _ {d - 1, d - 1} \end{array} \right),
\]

使得对 \(h, k \in \mathrm{H}\) ，\(A_{h,k}\) 是 \((\mathrm{I}_h, \mathrm{I}_k)\) 型的矩阵。

\(f)\) 对所有 \((h,k)\in\mathrm{H}\) ，有 \(\omega A_{h,k}=\omega^{k-h}A_{h,k}\) 。由此推出 \(\mathrm{A}_{h,k}=0\) （若 \(k-h\neq1\) ）。

\begin{enumerate}
\def\labelenumi{\arabic{enumi})}
\setcounter{enumi}{14}
\tightlist
\item
  我们沿用习题 12 的假设与记号。设 \(\Delta_1(\mathrm{X})\) 是 \(B(\mathrm{X})\) 的诸系数的最大公因子；假设多项式 \(\Delta_1(\mathrm{X})\) 是首一的。
\end{enumerate}

\begin{enumerate}
\def\labelenumi{\alph{enumi})}
\item
  多项式 \(\Delta_1(X)\) 整除 \(\Delta(X)\) ，且 \(\Delta_1(r) \neq 0\) 。
\item
  由此推出 \(\Delta_{1}(X)\) 的实根都严格小于 r。从而得出 \(\Delta_{1}(r)>0\) 。
\item
  置 \(C(\mathrm{X}) = \frac{1}{\Delta_1(\mathrm{X})} B(\mathrm{X})\) 。这是一个系数属于 \(\mathbf{R}[\mathrm{X}]\) 的矩阵。证明 \(C(r) \succ 0\) 。
\end{enumerate}

\begin{enumerate}
\def\labelenumi{\arabic{enumi})}
\setcounter{enumi}{15}
\tightlist
\item
  设 A 是 (I, I) 型的正且不可约的实方阵。对所有 \(j \in I\) 置 \(s_{j} = \sum_{k \in I} a_{jk}\) ，
\end{enumerate}

\[
s = \inf _ {j \in \mathrm{I}} s _ {j}, \quad \text {and} \quad \mathrm{S} = \sup _ {j \in \mathrm{I}} s _ {j}.
\]

我们也采用习题 12 的记号 \(r\) 与 \(B(\mathbf{X})\) 。

\begin{enumerate}
\def\labelenumi{\alph{enumi})}
\item
  有 \(\sum_{j\in \mathrm{I}}(r - s_j)B_{kj}(r) = 0\)
\item
  由此推出 \(s \leqslant r \leqslant S\) ，且这两个不等式中任何一个取等号都蕴含 s = S。
\end{enumerate}

\begin{enumerate}
\def\labelenumi{\arabic{enumi})}
\setcounter{enumi}{16}
\item
  设 \(A\) 是 (I, I) 型的正且不可约的实方阵。设 \(y \in \mathbf{R}^{\mathrm{I}}\) 是 \(A\) 的属于特征值 \(\alpha\) 的特征向量，且满足 \(y \succcurlyeq 0\) 。则 \(\alpha = r\) （沿用习题 12 的记号）且 \(y \succ 0\) 。（考虑 \(u\) （\(^{t}A\) 的一个极端向量），并计算 \(\langle y|^{t}Au\rangle = \langle Ay|u\rangle\) 。）
\item
  设 A 是 (I, I) 型的正且不可约的实方阵。设 \(\varrho\) 是 A 的谱半径。对每个 \(x \in R^{I}\) （满足 \(y \succcurlyeq 0\) 且 \(x \neq 0\) ），我们置
\end{enumerate}

\[
\varrho^ {x} = \sup _ {i \in \mathrm{I}} \frac {(A x) _ {i}}{x _ {i}},
\]

（约定 \(\varrho^x = +\infty\) ，如果存在指标 \(i\) 使 \(x_i = 0\) 且 \((Ax)_i \neq 0\) ）。设

\[
\varrho^{\prime} = \inf_{\substack{x\succ 0\\ x\neq 0}}\varrho^{x}.
\]

沿用与习题 12 类似的途径，证明 \(\varrho'\) 是 \(A\) 的一个特征值且 \(\varrho' = \varrho\) 。

\begin{enumerate}
\def\labelenumi{\arabic{enumi})}
\setcounter{enumi}{18}
\tightlist
\item
  设 A 是 (I, I) 型的正且不可约的实方阵；\(\varrho\) 是它的谱半径。设 \(z \in R^{I}\) 满足 \(z \succcurlyeq 0\) 且 \(z \neq 0\) 。
\end{enumerate}

若 \(\varrho z \preccurlyeq Az\) （相应地，\(\varrho z \succcurlyeq Az\) ），则 \(\varrho z = Az\) 且 \(z \succ 0\) 。

\begin{enumerate}
\def\labelenumi{\arabic{enumi})}
\setcounter{enumi}{19}
\tightlist
\item
  设 A 是 (I, I) 型的正但未必不可约的实方阵；\(\varrho\) 是 A 的谱半径。我们保留前面诸习题的记号 \(B(\mathrm{X})\) 与 \(C(\mathrm{X})\) 。
\end{enumerate}

\begin{enumerate}
\def\labelenumi{\alph{enumi})}
\item
  实数 \(\varrho\) 是 \(A\) 的特征值，且存在至少一个特征向量 \(\succcurlyeq 0\) 相伴于 \(\varrho\) 。
\item
  对每个实数 \(\lambda \geqslant r\) ，矩阵 \(B(\lambda)\) 与 \(B'(\lambda)\) 都是正的。（对第二个矩阵，可对 I 的基数用归纳法。）
\item
  矩阵 \(C(\lambda)\) 对所有 \(\lambda \geqslant r\) 都是正的。
\item
  若 A 是不可约的，则对每个实数 \(\lambda > r\) ，矩阵 \(B(\lambda)\) 、\(C(\lambda)\) 与 \((\lambda1_{\mathrm{I}} - A)^{-1}\) 都是严格正的。
\item
  设 J 是 I 的异于 I 的子集。设 \(\varrho_{\mathrm{J}}\) 是子矩阵 \((a_{i,j})_{i,j\in \mathrm{J}}\) 的谱半径。我们有 \(\varrho_{\mathrm{J}}\leqslant \varrho\) ，且当 \(A\) 不可约时该不等式是严格的。反之，若 \(A\) 是可约的，则存在异于 I 的子集 \(\mathrm{J}\subset \mathrm{I}\) 使得 \(\varrho_{\mathrm{J}} = \varrho\) 。
\item
  设 J 是 I 的异于 I 的子集，且子矩阵 \((a_{i,j})_{i,j\in\mathrm{J}}\) 的谱半径等于 \(\varrho\) 。对每个满足 \(J\subset K\subset I\) 的子集 K，矩阵 \((a_{i,j})_{i,j\in\mathrm{K}}\) 的谱半径等于 \(\varrho\) 。由此推出 \(A\succ0\) 当且仅当 \(B(r)\) 的所有对角系数都严格为正。
\end{enumerate}

\(g)\) 对每个实数 \(\lambda >\varrho\) 以及 I 的每个满足 \(\mathrm{J}\neq \mathrm{I}\) 的子集 J，矩阵 \((\lambda \delta_{i,j} - a_{i,j})_{i,j\in \mathrm{J}}\) 的行列式严格为正。

\begin{enumerate}
\def\labelenumi{\alph{enumi})}
\setcounter{enumi}{7}
\tightlist
\item
  设 \(\lambda\) 是实数；假设 \((\lambda \delta_{i,j} - a_{i,j})_{i,j \in \mathbf{J}}\) 的行列式对每个子集 \(\mathbf{J} \subset \mathbf{I}\) （异于 \(\mathbf{I}\) ）都严格为正。则有 \(\lambda > \varrho\) 。
\end{enumerate}

在下面的问题中，我们假设 \(\mathrm{I}\) 是区间 \([1,n]\) （在 \(\mathbf{N}\) 中）。i) 设 \((g_{i,j})_{1\leqslant i,j\leqslant n}\) 是一个实矩阵，使得 \(g_{i,j}\leqslant0\) 对所有 \(i\neq j\) 成立，且矩阵 \((g_{i,j})_{1\leqslant i,j\leqslant m}\) 的行列式对每个整数 \(m\) （满足 \(1\leqslant m\leqslant n\) ）都严格为正。则对 I 的每个子集 J 有 \(\det(g_{i,j})_{i,j\in\mathrm{J}}>0\) 。

\begin{enumerate}
\def\labelenumi{\alph{enumi})}
\setcounter{enumi}{9}
\item
  设 \(\lambda\) 是实数；假设对每个整数 \(m\) （满足 \(1 \leqslant m \leqslant n\) ），行列式 \(\det(\lambda \delta_{i,j} - a_{i,j})_{1 \leqslant i,j \leqslant m}\) 严格为正。则 \(\lambda > \varrho\) 。
\item
  设 \(C = (c_{i,j})_{1 \leqslant i,j \leqslant n}\) 是一个实矩阵，使得 \(c_{i,j} \geqslant 0\) 对所有 \(i \neq j\) 成立。若 \((-1)^{k} \det(c_{i,j})_{1 \leqslant i,j \leqslant m}\) 对每个满足 \(1 \leqslant m \leqslant n\) 的整数 m 都严格为正，则 C（视为复矩阵）的特征值的实部都是负的。
\end{enumerate}

\begin{enumerate}
\def\labelenumi{\arabic{enumi})}
\setcounter{enumi}{20}
\tightlist
\item
  设 A 是 (I, I) 型的正的实方阵；\(\varrho\) 是它的谱半径。
\end{enumerate}

\begin{enumerate}
\def\labelenumi{\alph{enumi})}
\tightlist
\item
  证明：存在整数 p 与 q 满足 \(0 \leqslant q \leqslant p\) ，以及 I 的一个分成非空部分的划分 \((\mathrm{I}_{s})_{1 \leqslant s \leqslant p}\) ，使得相应的分块分解 \(A = (A_{s,t})_{1 \leqslant s, t \leqslant p}\) 具有下列性质：
\end{enumerate}

- 矩阵 \(A_{s,s}\) 都是不可约的；

- 我们有 \(A_{s,t} = 0\) （若 \(t > s\) ）；

- 我们有 \(A_{s,t} = 0\) （若 \(1 \leqslant t < s \leqslant q\) ）；

- 对每个整数 \(s\) （满足 \(q \leqslant s \leqslant p\) ），存在 \(t\) 使得 \(A_{s,t} \neq 0\) 。

\begin{enumerate}
\def\labelenumi{\alph{enumi})}
\setcounter{enumi}{1}
\item
  确定这种划分的所有可能情形。（特别地，证明整数 q 与 p 都是唯一确定的，集组 \(\{I_{1},\ldots,I_{g}\}\) 与集组 \(\{I_{q+1},\ldots,I_{p}\}\) 亦然。）
\item
  存在 A 的一个特征向量 \(z \succ 0\) 属于特征值 \(\varrho\) ，当且仅当对每个整数 \(s \in [1, p]\) ，\(A_{s,s}\) 的谱半径等于 \(\varrho\) （若 \(s \leqslant q\) ），否则严格小于 \(\varrho\) 。
\item
  为使 \(A\) 与 \({}^t A\) 二者都承认一个严格正的极端向量，必须且只需存在如上的一个划分，使得 \(A_{s,t} = 0\) （对 \(1 \leqslant t < s\) 且 \(q < s \leqslant p\) ），且矩阵 \(A_{s,s}\) （\(s \leqslant q\) ）的谱半径等于 \(\varrho\) 。
\item
  若 A 与 \(^{t}A\) 都承认一个严格正的极端向量，且 \(\varrho\) 的谱重数等于 1，则 A 是不可约的。
\end{enumerate}

\begin{enumerate}
\def\labelenumi{\arabic{enumi})}
\setcounter{enumi}{21}
\tightlist
\item
  设 A 是 (I, I) 型的正且不可约的实方阵；\(\varrho\) 是它的谱半径。设 \(d \geqslant 1\) 是由习题 14 确定的整数。当 d = 1 时，我们称 A 是本原的。设 \(\Delta(\mathrm{X})\) 是多项式 \(\det(\mathrm{X} \cdot 1_{\mathrm{I}} - A)\) 。
\end{enumerate}

\begin{enumerate}
\def\labelenumi{\alph{enumi})}
\item
  我们写 \(\Delta(\mathrm{X}) = \mathrm{X}^{n} + a_{1}\mathrm{X}^{n_{1}} + \cdots + a_{t}\mathrm{X}^{n_{t}}\) ，其中 \(a_{1}, \ldots, a_{t}\) 是非零实数，\(n_{1}, \ldots, n_{t}\) 是满足 \(n > n_{1} > \cdots > n_{t}\) 的自然数。证明 d 是 \(n - n_{1}, n_{1} - n_{2}, \ldots, n_{t-1} - n_{t}\) 的最大公因子。
\item
  存在非零整数 m 与多项式 g，使得 \(\Delta(\mathrm{X}) = g(\mathrm{X}^{d})\mathrm{X}^{m}\) 。
\item
  若存在整数 p \textgreater{} 0 使得 \(A^{p}\) 严格正，则 A 是本原的。
\item
  假设 \(A\) 是本原的。引入习题 15 的矩阵 \(C(\mathbf{X})\) ，并按该习题的记号置 \(\mathsf{X}(\mathbf{X}) = \Delta (\mathbf{X}) / \Delta_1(\mathbf{X})\) 。序列 \((A^k /\varrho^k)_{k\in \mathbf{N}}\) 的极限等于 \(C(\varrho) / \mathsf{X}'(\varrho)\) 。（可利用 \(A\) 的若尔当块分解。）利用习题 15 的结果，推出存在整数 \(p > 0\) 使得 \(A^p\) 严格正。
\item
  假设存在整数 q \textgreater{} 0 使得 \(A^{q}\) 是可约的。设 p 是 q 与 d 的最大公因子。证明：存在 I 的一个分成非空部分的划分 \((\mathrm{I}_{s})_{1 \leqslant s \leqslant p}\) ，使 \(A^{q}\) 关于此划分是分块对角的，其对角块不可约且谱半径为 \(\varrho^{q}\) 。
\item
  若 A 是本原的，则 \(A^{p}\) 对每个 p \textgreater{} 0 都是不可约的。
\end{enumerate}

\(g)\) 存在 I 的一个分成非空部分的划分 \((\mathrm{I}_s)_{1 \leqslant s \leqslant d}\) ，使 \(A^d\) 关于此划分是分块对角的，其对角块是本原的且谱半径等于 \(\varrho^d\) 。

\begin{enumerate}
\def\labelenumi{\arabic{enumi})}
\setcounter{enumi}{22}
\tightlist
\item
  设 P 是 (I, I) 型的实方阵；若 P 是非负的，且向量 \((1, \ldots, 1) \in \mathbf{R}^{\mathrm{I}}\) 是 P 的属于特征值 1 的特征向量，则称 P 是随机的。
\end{enumerate}

设 A 是 (I, I) 型的实方阵。证明：A 是正的且承认一个严格正的极端向量，当且仅当存在一个随机矩阵 P、一个实数 r \textgreater{} 0 以及一个对角系数严格为正的对角矩阵 Z，使得 \(A = rZPZ^{-1}\) 。在这种情形，\(A\) 的谱半径等于 \(r\) ，而 \(A\) 的一个极端向量是 \(Z(1, \ldots, 1)\) 。

\begin{enumerate}
\def\labelenumi{\arabic{enumi})}
\setcounter{enumi}{23}
\tightlist
\item
  \begin{enumerate}
  \def\labelenumii{\alph{enumii})}
  \tightlist
  \item
    设 P 是 n 阶的实随机方阵。实数 1 是 P 的极小多项式的单根。（可利用习题 20。）此外，P 的极小多项式的每个模为 1 的根都是单根。
  \end{enumerate}
\end{enumerate}

\begin{enumerate}
\def\labelenumi{\alph{enumi})}
\setcounter{enumi}{1}
\tightlist
\item
  设 A 是一个正实矩阵，它承认一个极端向量 \(z \succ 0\) ，谱半径为 \(\varrho\) 。A 的极小多项式的所有模为 \(\varrho\) 的根都是单根。
\end{enumerate}

在习题 25 至 30 中，我们考虑一个实弗雷歇空间 E 和 E 中的一个闭凸锥 C，并假设 C 的内部非空。由 C 诱导的 E 上的序关系记作 \(\preccurlyeq_{\mathrm{C}}\) 。于是，对 E 中的 \(x\) 、\(y\) ，\(x \succcurlyeq_{\mathrm{C}} y\) 当且仅当 \(x - y \in \mathrm{C}\) 。我们将用 \(x \succ_{\mathrm{C}} y\) 表示关系 \(x - y \in \mathring{\mathrm{C}}\) 。

我们考虑满足下列条件的连续映射 u: E → E：

\begin{enumerate}
\def\labelenumi{(\roman{enumi})}
\item
  映射 u 是正齐次的：\(u(\lambda x) = \lambda u(x)\) 对每个实数 \(\lambda \geqslant 0\) 成立；
\item
  映射 u 关于 C 是递增的，即 \(x \succsim_{C} y\) 蕴含 \(u(x) \succsim_{C} u(y)\) ；
\item
  对每个 \(x \in \mathrm{C} - \{0\}\) ，有 \(u(x) \in \mathring{\mathrm{C}}\) 。
\end{enumerate}

这些习题取自 P.-L. 利翁斯 2020/2021 学年在法兰西公学院的课程，其目标是在某些条件下证明：存在元素 \(x \in \mathrm{C} - \{0\}\) 与 \(\lambda > 0\) 使得 \(x = \lambda u(x)\) ；特别地，这给出了克赖因–鲁特曼定理（III, p. 94, 定理 4 与 III, p. 97, 命题 8）的另一个证明。

对 \(f \in \mathring{\mathrm{C}}\) 与 \(\lambda \in \mathbf{R}_{+}\) ，我们把方程 \(y = \lambda u(y) + f\) 的上解（相应地，下解）称为一个元素 \(y \in \mathrm{C}\) ，它满足 \(y \succcurlyeq_{\mathrm{C}} \lambda u(y) + f\) （相应地，满足 \(y \preccurlyeq_{\mathrm{C}} \lambda u(y) + f\) ）。我们用 \(U_{f}\) 表示由 \(\lambda \in \mathbf{R}_{+}\) （使方程 \(y = \lambda u(y) + f\) 的上解集合非空）所成的集。

\begin{enumerate}
\def\labelenumi{\arabic{enumi})}
\setcounter{enumi}{24}
\tightlist
\item
  \begin{enumerate}
  \def\labelenumii{\alph{enumii})}
  \tightlist
  \item
    证明：对每个 \(x \in C\) 与每个 \(y \in \mathring{C}\) ，存在实数 \(\lambda \geqslant 0\) 使得 \(x \preccurlyeq_{C} \lambda y\) 。
  \end{enumerate}
\end{enumerate}

\begin{enumerate}
\def\labelenumi{\alph{enumi})}
\setcounter{enumi}{1}
\tightlist
\item
  集合 \(U_{f}\) 是包含 0 的一个区间。
\end{enumerate}

我们对映射 u 引入下列条件。

\begin{enumerate}
\def\labelenumi{(\roman{enumi})}
\setcounter{enumi}{3}
\item
  对每个 \(\lambda \in \mathbf{R}_{+}\) 、每个 \(f \in \mathring{\mathrm{C}}\) ，以及每个序列 \((x_{n})\) （由 C 中元素组成，满足 \(x_{n+1} = \lambda u(x_{n}) + f\) （对所有 \(n\) ），且递增并上有界或递减），它都是收敛的；
\item
  E 的每个有界子集在 u 下的像是相对紧的。
\item
  对每个映射 \(\lambda \mapsto \varepsilon_{\lambda}\) （从 \(U_f\) 到 C 中），若它满足 \(\lim_{\lambda \to \lambda_m} \varepsilon_{\lambda} = 0\) ，则集合
\end{enumerate}

\[
\{z \in \mathrm{C} \mid \exists \lambda \in \mathrm{U}_{f}\text{ 使得 } z = \lambda u(z) + \varepsilon_{\lambda} \}
\]

在 E 中是相对紧的。

\begin{enumerate}
\def\labelenumi{\arabic{enumi})}
\setcounter{enumi}{25}
\tightlist
\item
  我们沿用前面的记号与假设，并另外假设条件 (iv) 成立。
\end{enumerate}

\begin{enumerate}
\def\labelenumi{\alph{enumi})}
\item
  设 \(\lambda \in \mathrm{U}_f\) 。我们定义一个序列 \((x_n)_{n \in \mathbf{N}}\) ：\(x_0 = 0\) 且 \(x_{n+1} = \lambda u(x_n) + f\) （对所有 \(n \in \mathbf{N}\) ）。它递增且上有界；我们有 \(x_n \in \mathring{\mathrm{C}}\) （对 \(n \geqslant 1\) ）。
\item
  由条件 (iv)，序列 \((x_{n})\) 收敛；用 \(x(\lambda)\) 表示它的极限。证明 \(x(\lambda) \in \mathring{\mathrm{C}}\) 与 \(x(\lambda) = \lambda u(x(\lambda)) + f\) 。
\item
  设 \(\lambda \in \mathrm{U}_f\) ，并设 \(y \in \mathrm{C}\) 满足 \(y \succcurlyeq_{\mathrm{C}} \lambda u(y) + f\) ；由 \(x_0 = y\) 、\(x_{n+1} = \lambda u(x_n) + f\) 所定义的序列是递减的，取值于 C 中，其极限 \(x\) 满足方程 \(x = \lambda u(x) + f\) 。
\end{enumerate}

\(d)\) \(\mathrm{U}_{f}\) 在 \(\mathbf{R}_{+}\) 中的上确界不属于 \(\mathrm{U}_{f}\) 。

\begin{enumerate}
\def\labelenumi{\alph{enumi})}
\setcounter{enumi}{4}
\item
  \(U_{g} = U_{f}\) 对所有 \(g \in \mathring{C}\) 成立（对每个 \(\lambda \in U_{f}\) ，设 \(x \in C\) 满足 \(x \succsim_{C} \lambda u(x) + f\) ；则存在 M \textgreater{} 0 使得 y = Mx 满足 \(y \succsim_{C} \lambda u(y) + g\) 。） f) 设 \(\lambda \leqslant \mu\) 是实数。设 \(x\) 是方程 \(y = \lambda u(y) + f\) 的一个下解，设 \(x'\) 是方程 \(y = \mu u(y) + f\) 的一个上解。则 \(x \preccurlyeq_{C} x'\) 。（证明由 \(\varepsilon \in [0,1]\) （使 \(\varepsilon x \preccurlyeq_{C} x'\) ）组成的集与 {[}0,1{]} 重合。）
\item
  若 C 是尖锥，则对每个 \(\lambda \in \mathrm{U}_f\) ，存在唯一的元素 \(x \in \mathring{\mathrm{C}}\) 使得 \(x = \lambda u(x) + f\) 。
\end{enumerate}

\(h\) ) 设 \(\mu\in\mathbf{R}_{+}\) 满足 \(f\preccurlyeq_{\mathrm{C}}\mu u(f)\) 。则 \(\mu\notin\mathrm{U}_{f}\) 。（序列 \((x_{n})_{n\in\mathbf{N}}\) （由 \(x_{0}=0\) 与 \(x_{n+1}=\mu u(x_{n})+f\) 定义）满足 \(x_{n+1}-x_{n}\succcurlyeq_{\mathrm{C}}f\) （对所有 \(n\in\mathbf{N}\) ）。）特别地，区间 \(\mathrm{U}_{f}\) 有界。此后我们用 \(\lambda_{\mathrm{m}}\) 表示 \(\mathrm{U}_{f}\) 的上确界。

\begin{enumerate}
\def\labelenumi{\roman{enumi})}
\tightlist
\item
  设 \(z \in \mathrm{C} - \{0\}\) 与 \(\mu \in \mathbf{R}_+\) 满足 \(\mu u(z) \succsim_{\mathrm{C}} z\) 。则 \(\lambda_{\mathrm{m}} \leqslant \mu\) ，从而
\end{enumerate}

\[
\lambda_ {\mathrm{m}} \leqslant \inf _ {z \in \mathrm{C} - \{0 \}} \inf \{\mu \mid z \leqslant \mu u (z) \}.
\]

\begin{enumerate}
\def\labelenumi{\alph{enumi})}
\setcounter{enumi}{9}
\tightlist
\item
  映射 \(\lambda \mapsto x(\lambda)\) （其中 \(x(\lambda)\) 在问题 \(a\) 中定义）是递增的。
\end{enumerate}

\(k)\) 映射 \(\lambda \mapsto x(\lambda)\) 当 \(\lambda\) 趋于 \(\lambda_{\mathrm{m}}\) 时没有极限。更确切地说，若 \((\lambda_n)_n\) 是 \(\mathrm{U}_f\) 中趋于 \(\lambda_{\mathrm{m}}\) 的一个序列，则序列 \((x(\lambda_n))_n\) 不收敛。

\begin{enumerate}
\def\labelenumi{\arabic{enumi})}
\setcounter{enumi}{26}
\tightlist
\item
  我们假设 (i) 至 (v) 各条假设成立。
\end{enumerate}

\begin{enumerate}
\def\labelenumi{\alph{enumi})}
\item
  证明：\(x(\lambda)\) （对 \(\lambda \in \mathrm{U}_f\) ）所成的集不是有界的。
\item
  设 d 是 E 上定义其拓扑的一个距离。对 \(\lambda \in R_{+}\) ，定义 \(\hat{x}(\lambda) = d(0, x(\lambda))^{-1}x(\lambda)\) 。证明映射 \(\hat{x}\) 有界，且 \(\hat{x}\) 的值集有一个聚点 \(x_{m}\) 满足
\end{enumerate}

\begin{enumerate}
\def\labelenumi{(\roman{enumi})}
\item
  \(d(0, x_{\mathrm{m}}) = 1,\)
\item
  \(x_{\mathrm{m}} = \lambda_{\mathrm{m}} u(x_{\mathrm{m}}),\)
\item
  \(x_{m} \in \overset{\circ}{C}.\)
\end{enumerate}

\begin{enumerate}
\def\labelenumi{\arabic{enumi})}
\setcounter{enumi}{27}
\tightlist
\item
  在本习题中我们假设条件 (i) 至 (iv) 成立。 a) 我们固定 \(x_0 \in \mathring{\mathrm{C}}\) 。对 \(\lambda \in \mathrm{U}_f\) ，令 \(\mathrm{A}(\lambda) = \inf \{\mathrm{M} \in \mathbf{R}_+ \mid x(\lambda) \preccurlyeq_{\mathrm{C}} \mathrm{M}x_0\}\) 。设 \(\mathrm{C}_0 \in \mathbf{R}_+\) 满足 \(f \preccurlyeq_{\mathrm{C}} \mathrm{C}_0 u(x_0)\) 。对每个 \(\lambda\) （\(\mathrm{U}_f\) 中的），区间 \([\lambda, \lambda + \mathrm{C}_0 / \mathrm{A}(\lambda)[\) 包含于 \(\mathrm{U}_f\) （参见习题 26 的问题 \(d\) )）。
\end{enumerate}

\begin{enumerate}
\def\labelenumi{\alph{enumi})}
\setcounter{enumi}{1}
\tightlist
\item
  由此推出：函数 \(\lambda \mapsto \mathrm{A}(\lambda)\) 趋于 \(+\infty\) （当 \(\lambda \in \mathrm{U}_f\) 趋于 \(\lambda_{\mathrm{m}}\) 时）。
\end{enumerate}

对每个 \(\lambda\) （\(\mathrm{U}_f\) 中的），定义 \(\hat{x} (\lambda) = x(\lambda) / \mathrm{A}(\lambda)\) 。

我们现在假设条件 (vi) 成立。

\begin{enumerate}
\def\labelenumi{\alph{enumi})}
\setcounter{enumi}{2}
\tightlist
\item
  证明映射 \(\hat{x}\) 的值集有一个聚点 \(x_{\mathrm{m}}\) 满足
\end{enumerate}

\[
0 \preccurlyeq_ {\mathrm{C}} x _ {\mathrm{m}} \preccurlyeq_ {\mathrm{C}} x _ {0} \quad \text {and} \quad x _ {\mathrm{m}} = \lambda_ {\mathrm{m}} u (x _ {\mathrm{m}}).
\]

\begin{enumerate}
\def\labelenumi{\alph{enumi})}
\setcounter{enumi}{3}
\item
  证明 \(x_{m} \neq 0\) 。（若 \(x_{m} = 0\) ，设 U 是 E 中 0 的一个邻域，使得 \(x_{0} + U \subset C\) ；则对每个 \(\mu \geqslant 0\) ，将存在一个序列 \((\lambda_{n})\) （\(U_{f}\) 中的），收敛于 \(\lambda_{m}\) ，使 \(x_{0} - \mu \hat{x}(\lambda_{n}) \succsim_{C} 0\) ，这与 \(A(\lambda)\) 的定义矛盾。）
\item
  设 \(y \in \mathrm{C} - \{0\}\) 满足 \(y = \lambda_{\mathrm{m}} u(y)\) 。置 \(\theta = \sup \{\gamma \mid \gamma y \preccurlyeq_{\mathrm{C}} x_{\mathrm{m}}\}\) 。我们有 \(\theta \in \mathbf{R}_{+}^{*}\) 且 \(x_{\mathrm{m}} = \theta y\) 。（我们有 \(\theta y \preccurlyeq_{\mathrm{C}} x_{\mathrm{m}}\) 与 \(\theta y = \lambda_{\mathrm{m}} u(\theta y)\) ；若 \(x_{\mathrm{m}} \neq \theta y\) ，则将有 \(x_{\mathrm{m}} - \theta y \succ_{\mathrm{C}} 0\) ，且存在 \(\varepsilon > 0\) 使 \(x_{\mathrm{m}} - \theta y \succcurlyeq_{\mathrm{C}} \varepsilon y\) ，这与 \(\theta\) 的定义矛盾。）
\item
  设 \(y \in C - \{0\}\) 与 \(\mu \in R\) 满足 \(y = \mu u(y)\) 。则 \(\mu = \lambda_{m}\) 。（注意 \(0 < \mu\) 且 \(\mu \geqslant \lambda_{m}\) ；若有 \(\mu > \lambda_{m}\) ，令 \(\theta = \sup\{\gamma \mid \gamma x_{m} \preccurlyeq_{C} y\}\) ；则将有 \(\theta \in R_{+}^{*}\) ，进而 \(\theta x_{m} \preccurlyeq_{C} y\) ，于是 \(\theta x_{m} = \lambda_{m} u(\theta x_{m}) \prec_{C} \mu u(\theta x_{m}) \preccurlyeq_{C} \mu u(y) = y\) ，这与 \(\theta\) 的定义矛盾。）
\end{enumerate}

\begin{enumerate}
\def\labelenumi{\arabic{enumi})}
\setcounter{enumi}{28}
\tightlist
\item
  我们沿用前一习题的记号；特别地，假设 C 的内部非空。我们假设 u 是线性映射，且条件 (i)、(ii)、(iii) 成立。
\end{enumerate}

\begin{enumerate}
\def\labelenumi{\alph{enumi})}
\tightlist
\item
  若 u 是 E 的紧自同态，则条件 (iv) 与 (v) 成立。
\end{enumerate}

我们假设条件 (iv) 与 (vi) 成立。设 \((\lambda_{\mathrm{m}}, x_{\mathrm{m}})\) 是由前一习题给出的。

\begin{enumerate}
\def\labelenumi{\alph{enumi})}
\setcounter{enumi}{1}
\tightlist
\item
  对每个 \(\lambda \in \mathbf{R}_{+}\) （满足 \(\lambda < \lambda_{\mathrm{m}}\) ），自同态 \(1_{\mathrm{E}} - \lambda u\) 是可逆的。（首先证明 \(1_{\mathrm{E}} - \lambda u\) 通过过渡到子空间在 \(\mathring{\mathrm{C}}\) 上诱导一个双射，参见前一习题的问题 \(a\) ) 与 \(f\) )；然后推出
\end{enumerate}

\(1_{E} - \lambda u\) 是满射，然后通过把 \(x_{1} - x_{2}\) 写成 \(1_{E} - \lambda u\) 的核中元素（其中 \(x_{i} \in \mathring{C}\) ）并验证存在 \(z \in E\) 使 \(x_{1} + z - \lambda u(x_{1} + z) \in \mathring{C}\) ，推出它是单射。

\begin{enumerate}
\def\labelenumi{\alph{enumi})}
\setcounter{enumi}{2}
\tightlist
\item
  设 E 中的 y 满足 \(y = \lambda_{\mathrm{m}} u(y)\) 。则 \(y \in C\) 。（存在 c 使 \(cx_{m} + y \succ_{C} 0\) ；利用前一习题的 (e) 部分得出结论。）
\end{enumerate}

\(d)\) 设 \(\mathrm{C}^{\circ}\subset\mathrm{E}^{\prime}\) 是 C 的极集。对每个 \(x\in\mathring{\mathrm{C}}\) 与每个 \(\xi\in\mathrm{C}^{\circ}-\{0\}\) ，有 \(\langle x,\xi\rangle>0\) 。

\begin{enumerate}
\def\labelenumi{\alph{enumi})}
\setcounter{enumi}{4}
\tightlist
\item
  C 的极集 C° 在 \(^{t}u\colon E'\to E'\) 下不变；对每个 \(\lambda\in R_{+}\) （满足 \(\lambda<\lambda_{m}\) ），自同态 \(1_{E'}-\lambda^{t}u\) （\(E'\) 的）是一个同构。
\end{enumerate}

\begin{enumerate}
\def\labelenumi{\arabic{enumi})}
\setcounter{enumi}{29}
\tightlist
\item
  我们沿用前一习题的假设，即 u 是线性、紧的，且假设条件 (i) 至 (v) 成立。自同态 \(^{t}u\) 也是紧的。
\end{enumerate}

\begin{enumerate}
\def\labelenumi{\alph{enumi})}
\item
  设 \(\eta \in \mathrm{C}^{\circ} - \{0\}\) 。对每个 \(\lambda \in \mathbf{R}_{+}\) （满足 \(\lambda < \lambda_{\mathrm{m}}\) ），唯一解 \(\xi(\lambda)\) （方程 \(\xi - \lambda^{t}u(\xi) = \eta\) 的）属于 \(\mathrm{C}^{\circ}\) 。（若 \(x \in \mathrm{C}\) ，设 \(y \in \mathrm{C}\) 满足 \(y - \lambda u(y) = x\) ；则 \(\langle y, \eta \rangle = \langle x, \xi \rangle\) 。）
\item
  设 \(x_{0} \in \mathring{C}\) 。映射 \(\lambda \mapsto \langle x_{0}, \xi(\lambda) \rangle\) 当 \(\lambda\) 趋于 \(\lambda_{m}\) 时趋于无穷。证明值 \(\xi(\lambda)/\langle x_{0}, \xi(\lambda) \rangle\) 所成的集有一个聚点 \(\xi_{m}\) 满足：
\end{enumerate}

\begin{enumerate}
\def\labelenumi{(\roman{enumi})}
\item
  \(\xi_{\mathrm{m}}\in \mathrm{C}^{\circ}\) ;
\item
  \(\langle x_{0},\xi_{m}\rangle=1;\)
\item
  \(\xi_{\mathrm{m}} = \lambda_{\mathrm{m}}^{t}u(\xi_{\mathrm{m}})\) .
\end{enumerate}

\begin{enumerate}
\def\labelenumi{\alph{enumi})}
\setcounter{enumi}{2}
\item
  设 \(f \in E\) 满足 \(\langle f, \xi_{m} \rangle = 0\) 。对 \(\lambda \in R_{+}\) （满足 \(\lambda < \lambda_{m}\) ），用 \(x(\lambda)\) 表示唯一解 \(x\) （方程 \(x - \lambda u(x) = f\) 在 E 中的）。证明 \(x(\lambda)\) 所成的集在 E 中有界（设 d 是定义 E 的拓扑的距离；若存在序列 \(\lambda_{n} \to \lambda_{m}\) 使 \(x(\lambda_{n})\) 趋于无穷，置 \(\hat{x}_{n} = x(\lambda_{n}) / d(0, x(\lambda_{n}))\) ，并证明序列 \((\hat{x}_{n})\) 收敛于一个元素 \(\hat{x}\) ，它满足 \(d(0, \hat{x}) = 1\) 、\(\hat{x} = \lambda u(\hat{x})\) 与 \(\langle \hat{x}, \xi_{m} \rangle = 0\) ，这与必有 \(\hat{x} \in \mathring{C}\) 的事实矛盾）。
\item
  存在唯一的元素 \(x \in E\) 满足 \(f = x - \lambda_{\mathrm{m}} u(x)\) 与 \(\langle x, \xi_{\mathrm{m}} \rangle = 0\) 。（关于存在性，利用前一问题证明存在一个序列 \((\lambda_{n})\) 收敛于 \(\lambda_{m}\) ，使 \((x(\lambda_{n}))\) 收敛于一个满足所需条件的元素 x。）
\item
  自同态 \(1_{\mathrm{E}'} - \lambda_{\mathrm{m}}{}^{t}u\) （\(\mathrm{E}'\) 的）通过过渡到子空间诱导极集 \(x_{\mathrm{m}}^{\circ}\) （\(x_{\mathrm{m}}\) 的）上的一个双射。
\item
  特征空间 \(\operatorname{Ker}(1_{\mathrm{E}} - \lambda_{\mathrm{m}}u)\) 的维数为 1，且等于 \(1_{E} - \lambda_{m}u\) 的像的极集。
\end{enumerate}

第四章

\chapter{希尔伯特谱理论}
\setcounter{exercisepar}{0}

对每个向量空间 E，我们用 \(1_{E}\) 表示 E 的恒等映射。若 E、F 与 G 是向量空间，且 v: \(E \rightarrow F\) 与 u: \(F \rightarrow G\) 是线性映射，则从 E 到 G 的线性映射 \(u \circ v\) 有时记作 uv。

设 E 是一个预希尔伯特空间。我们用 \(\langle x|y\rangle_{E}\) ，或简作 \(\langle x|y\rangle\) ，表示 E 上的内积。E 的子集 A 的正交记作 A°。它是 E 的一个闭子空间；它化为 0 当且仅当 A 在 E 中是全的（EVT, V, p. 16, 推论 1）。若 F ⊂ E 是 E 的向量子空间，则 F°° = \(\overline{F}\) （参见 EVT, V, p. 13）。E 的自同态的谱总是相对于巴拿赫代数 \(\mathcal{L}(E)\) 而言的。

设 F 是希尔伯特空间。我们用 \(\mathcal{L}_1(\mathrm{E})\) 表示 E 的有限迹自同态所成的空间（EVT, V, p. 50, 定义 8），用 \(\mathcal{L}_2(\mathrm{E};\mathrm{F})\) 表示从 E 到 F 的希尔伯特–施密特映射所成的空间（EVT, V, p. 51, 定义 9）。映射 \(u\mapsto \| u\| _2 = (\mathrm{Tr}(u^{*}u))^{1 / 2}\) 是 \(\mathcal{L}_2(\mathrm{E};\mathrm{F})\) 上的一个希尔伯特范数（EVT, V, p. 52, 定理 1, (i)）。

有限秩的连续线性映射所成的空间 \(\mathcal{L}^{\mathrm{f}}(\mathrm{E};\mathrm{F})\) 在 \(\mathcal{L}^{\mathrm{c}}(\mathrm{E};\mathrm{F})\) 中稠密（III, p. 15, 命题 10 与 III, p. 16, 命题 11 的推论）。特别地，空间 \(\mathcal{L}_2(\mathrm{E};\mathrm{F})\) 在 \(\mathcal{L}^{\mathrm{c}}(\mathrm{E};\mathrm{F})\) 中稠密。

E 与 F 的希尔伯特张量积（EVT, V, p. 28, 定义 1）记作 E \(\widehat{\otimes}_2\) F。

若 X 是局部紧拓扑空间，\(\mu\) 是 X 上的一个测度，我们记 \(\mathcal{L}^{p}(\mathrm{X},\mu)=\mathcal{L}_{\mathbf{C}}^{p}(\mathrm{X},\mu)\) 与 \(\mathrm{L}^{p}(\mathrm{X},\mu)=\mathrm{L}_{\mathbf{C}}^{p}(\mathrm{X},\mu)\) （对 \(p\in[1,+\infty]\) ）。

\section{希尔伯特空间上的紧算子}

在本节中，K 是等于 R 或 C 的域，E 表示域 K 上的一个希尔伯特空间。

\subsection{对角自同态}

定义 1. --- 设 \(\mathrm{B} = (e_i)_{i \in \mathrm{I}}\) 是 E 的一个标准正交基。称一个自同态 \(u\) 在基 B 中是对角的或相对于 B 是对角的，如果存在由 K 的元素组成的族 \(\lambda = (\lambda_i)_{i \in \mathrm{I}}\) ，使得 \(u(e_i) = \lambda_i e_i\) （对所有 \(i \in \mathrm{I}\) ）。

设 \(\mathrm{B} = (e_i)_{i \in \mathrm{I}}\) 是 E 的一个标准正交基。我们用 \(\mathcal{D}_{\mathrm{B}}(\mathrm{E})\) 表示 E 的相对于 B 为对角的自同态所成的集。这是 \(\mathcal{L}(\mathrm{E})\) 的一个闭的、幺的交换子代数。设 \(u \in \mathcal{D}_{\mathrm{B}}(\mathrm{E})\) 。族 \(\lambda = (\lambda_i)_{i \in \mathrm{I}}\) （满足 \(u(e_i) = \lambda_i e_i\) ）由 u 和基 B 唯一确定，称为 u 相对于 B 的特征值族。

假设 K = R，并把 E 与 \(\mathrm{E}_{(\mathbf{C})}\) 的一个子空间等同。设 \(\mathrm{B} = (e_i)_{i \in \mathrm{I}}\) 是 E 的一个标准正交基。用 \(\mathrm{B}_{(\mathbf{C})}\) 表示标准正交基 \((1 \otimes e_i)_{i \in \mathrm{I}}\) （\(\mathrm{E}_{(\mathbf{C})}\) 的）（参见 EVT, V, p. 29, 推论 1）。映射 \(u \mapsto u_{(\mathbf{C})}\) 诱导从 \(\mathcal{D}_{\mathrm{B}}(\mathrm{E})\) 到 \(\mathcal{D}_{\mathrm{B}_{(\mathbf{C})}}(\mathrm{E}_{(\mathbf{C})})\) 中的一个单的代数态射；它的像是 \(u \in \mathcal{D}_{\mathrm{B}_{(\mathbf{C})}}(\mathrm{E}_{(\mathbf{C})})\) （满足 \(u(\mathrm{E}) \subset \mathrm{E}\) ）所成的集；换言之，即基 \(\mathrm{B}_{(\mathbf{C})}\) 中特征值为实的对角自同态 u 所成的集。

对每个 \(i \in \mathrm{I}\) ，用 \(p_i\) 表示 E 的以 \(\mathrm{Ke}_i\) 为像的正交投影。我们有 \(\| p_i \| = 1\) （对所有 \(i \in \mathrm{I}\) ）。自同态 \(p_i\) 在基 B 中是对角的，其特征值族是族 \((\delta_{ij})_{j\in\mathrm{I}}\) （克罗内克符号，参见 A, II, p. 24）。

设 \(\mathrm{B}^{\prime}=(f_{i})_{i\in\mathrm{I}}\) 是 E 的一个标准正交基，\(u:E\to E\) 是满足 \(u(f_{i})=e_{i}\) 的等距同构。则 \(\mathcal{D}_{\mathrm{B}^{\prime}}(\mathrm{E})=u^{-1}\mathcal{D}_{\mathrm{B}}(\mathrm{E})u\) 。

我们给 I 赋以离散拓扑，并用 \(\mathcal{B}(\mathrm{I}) = \mathcal{C}_{b}(\mathrm{I}; \mathrm{K})\) 表示 I 上取值于 K 的有界函数所成的幺巴拿赫代数（I, p. 17, 例 2）；若 K = C，它是一个星代数（I, p. 102, 例 2）。

引理 1. --- 设 \(u \in \mathcal{D}_{\mathrm{B}}(\mathrm{E})\) ，\(\lambda = (\lambda_i)_{i \in \mathrm{I}}\) 是它的特征值族。族 \(\lambda\) 有界，且有 \(\sup_i |\lambda_i| = \| u\|\) 。

有 \(|\lambda_i| \leqslant \|u\|\) （对所有 \(i\) ），从而 \(\sup_{i \in I} |\lambda_i| \leqslant \|u\|\) 。此外，由于

\[
\| u (x) \| ^ {2} = \sum_ {i \in \mathrm{I}} | \lambda_ {i} | ^ {2} | \langle e _ {i} | x \rangle | ^ {2} \leqslant \left(\sup _ {i \in \mathrm{I}} | \lambda_ {i} | ^ {2}\right) \| x \| ^ {2}
\]

对所有 \(x \in \mathrm{E}\) 成立（EVT, V, p. 22, 命题 5），由此得出 \(\| u\| = \sup |\lambda_i|\) 。

命题 1. --- a) 映射 \(\alpha\) （从 \(\mathcal{D}_{\mathrm{B}}(\mathrm{E})\) 到 \(\mathcal{B}(\mathrm{I})\) 中的）把对角自同态 \(u\) 对应到 \(u\) 的特征值族，它是 K 上的巴拿赫代数的一个等距同构。若 K = C，它是对合代数的一个态射。

\begin{enumerate}
\def\labelenumi{\alph{enumi})}
\setcounter{enumi}{1}
\item
  对任何有界族 \(\lambda = (\lambda_i)_{i\in \mathrm{I}}\) ，族 \((\lambda_ip_i)_{i\in \mathrm{I}}\) 在赋以逐点收敛拓扑的空间 \(\mathcal{L}(\mathrm{E})\) 中是可和的，且该族之和是唯一的映射 \(u\in \mathcal{D}_{\mathrm{B}}(\mathrm{E})\) ，它满足 \(\alpha (u) = \lambda\) ；
\item
  设 \(u \in \mathcal{D}_{\mathrm{B}}(\mathrm{E})\) ，\(\lambda\) 是它相对于 B 的特征值族。\(u\) 的谱是 \(\lambda\) 的值集在 K 中的闭包；
\item
  若 K = C，则从 \(\mathcal{C}(\mathrm{Sp}(u))\) 到 \(\mathcal{L}(\mathrm{E})\) 中的连续函数演算映射把连续函数 \(f \in \mathcal{C}(\mathrm{Sp}(u))\) 对应到自同态 \(\alpha(f \circ \lambda)\) （\(\mathcal{D}_{\mathrm{B}}(\mathrm{E})\) 中的）。
\end{enumerate}

由引理 1，在 B 中为对角的自同态的特征值族是有界的，且如此定义的从 \(\mathcal{D}_{\mathrm{B}}(\mathrm{E})\) 到 \(\mathcal{B}(\mathrm{I})\) 中的映射是幺巴拿赫代数的一个连续等距态射。

设 \(i \in I\) 。对每个 \(j \in I\) ，我们有 \(\langle u^{*}(e_{i}) | e_{j} \rangle = \lambda_{j} \langle e_{i} | e_{j} \rangle = \langle \overline{\lambda_{j}} e_{i} | e_{j} \rangle\) 。由此得出 \(u^{*}(e_{i}) = \overline{\lambda}_{i} e_{i}\) 。于是 u 的伴随是基 B 中的对角自同态，其特征值族是 \((\overline{\lambda}_{i})_{i \in I}\) 。断言 a) 由此得出。

设 \(\lambda = (\lambda_i)_{i\in \mathrm{I}}\in \mathcal{B}(\mathrm{I})\) 。对每个 \(x\in \mathrm{E}\) ，族 \((|\langle e_i|x\rangle |^2)_{i\in \mathrm{I}}\) 可和，其和为 \(\| x\|^2\) （EVT, V, p. 22, 命题 5），从而，对 I 的每个有限子集 J：

\[
\begin{array}{r l r} & & {\left\| \sum_ {i \in \mathrm{J}} \lambda_ {i} p _ {i} (x) \right\| ^ {2} = \sum_ {i \in \mathrm{J}} | \lambda_ {i} | ^ {2} | \langle e _ {i} | x \rangle | ^ {2} \leqslant \Bigl (\underset {i \in \mathrm{I}} {\sup} | \lambda_ {i} | ^ {2} \Bigr) \sum_ {i \in \mathrm{J}} | \langle e _ {i} | x \rangle | ^ {2}} \\ & & {\leqslant \Bigl (\underset {i \in \mathrm{I}} {\sup} | \lambda_ {i} | \Bigr) ^ {2} \| x \| ^ {2}.} \end{array}
\]

因此，族 \((\lambda_{i}p_{i})\) 在赋以简单收敛拓扑的 \(\mathcal{L}(\mathrm{E})\) 中可和。它的和 \(u_{\lambda}\) 满足 \(u_{\lambda}(e_{i}) = \lambda_{i}e_{i}\) ；于是它是基 B 中的对角自同态，以 \(\lambda\) 为特征值族。断言 b) 由此得出。

最后的断言由 \(a\) )、I, p. 17 的例 3 与 I, p. 111 的例 4 得出。

注记. --- 巴拿赫代数 \(\mathcal{D}_{\mathrm{B}}(\mathrm{E})\) 是 \(\mathcal{L}(\mathrm{E})\) 的一个极大交换闭子代数。事实上，设 u 是与 \(\mathcal{D}_{\mathrm{B}}(\mathrm{E})\) 交换的 E 的自同态。设 \(i \in I\) 。由于正交投影 \(p_{i}\) 在基 B 中是对角的，我们有 \(p_{i}(u(e_{i})) = u(p_{i}(e_{i})) = u(e_{i})\) ，这蕴含 \(u(e_{i})\) 与 \(e_{i}\) 成比例。于是 u 在基 B 中是对角的。

若 E 是无穷维的，则存在 \(\mathcal{L}(\mathrm{E})\) 的极大交换对合子代数，它们不同构于 \(\mathcal{D}_{\mathrm{B}}(\mathrm{E})\) （IV, p. 314, 习题 5）。

命题 2. --- 设 \(u \in \mathcal{D}_{\mathrm{B}}(\mathrm{E})\) ，\(\lambda = (\lambda_i)_{i \in \mathrm{I}}\) 是它的特征值族。

\begin{enumerate}
\def\labelenumi{\alph{enumi})}
\tightlist
\item
  下列条件等价：
\end{enumerate}

\begin{enumerate}
\def\labelenumi{(\roman{enumi})}
\item
  有 \(\lambda \in \mathcal{C}_{0}(\mathrm{I};\mathrm{K})\) ；
\item
  自同态 u 是紧的；
\item
  族 \((\lambda_{i}p_{i})_{i\in\mathrm{I}}\) 在巴拿赫空间 \(\mathcal{L}(\mathrm{E})\) 中可和。其和这时等于 u。
\end{enumerate}

\begin{enumerate}
\def\labelenumi{\alph{enumi})}
\setcounter{enumi}{1}
\tightlist
\item
  假设 \(u\) 是紧的，用 \(\Lambda\) 表示 \(\lambda\) 的值集。由 \(i \in I\) （使 \(\lambda_i \neq 0\) ）组成的集是可数的。若 E 是无穷维的，我们有 \(\mathrm{Sp}_{\mathrm{s}}(u) = \Lambda - \{0\}\) 与 \(\mathrm{Sp}(u) = \Lambda \cup \{0\}\) 。若 E 是有限维的，则 \(\mathrm{Sp}_{\mathrm{s}}(u) = \mathrm{Sp}(u) = \Lambda\) 。
\end{enumerate}

由引理 1，对 I 的每个有限子集 J，我们有 \(\| \sum_{j\in J}\lambda_jp_j\| = \sup_{j\in J}|\lambda_j|\) 。由此，按柯西判别准则，族 \((\lambda_ip_i)\)

在 \(\mathcal{L}(\mathrm{E})\) 中可和当且仅当族 \(\lambda\) 在无穷远处趋于 0，这蕴含条件 (i) 与 (iii) 等价。

条件 (iii) 蕴含 \(u\) 是紧的（III, p. 4, 命题 2 的推论）。反之，假设自同态 \(u\) 是紧的。由于族 B 在 E 中有界，它在 \(u\) 下的像在 E 中相对紧（III, p. 2），从而是准紧的。设 \(\varepsilon > 0\) ，设 \(J\) 是 I 的一个有限子集，使得 B 在 \(u\) 下的像包含于半径为 \(\varepsilon\) 、心为 \(u(e_j)\) （\(j \in \mathrm{J}\) ）的球之并。设 \(i \in \mathrm{I} - \mathrm{J}\) 。存在 \(j \in \mathrm{J}\) 使 \(\| u(e_i) - u(e_j) \| \leqslant \varepsilon\) ，从而

\[
\left| \lambda_ {i} \right| ^ {2} \leqslant \left| \lambda_ {i} \right| ^ {2} + \left| \lambda_ {j} \right| ^ {2} = \left\| u \left(e _ {i}\right) - u \left(e _ {j}\right) \right\| ^ {2} \leqslant \varepsilon^ {2}.
\]

因此，我们有 \(\lambda \in \mathcal{C}_0(\mathrm{I};\mathrm{K})\) ，即条件 (i)。

最后，u 的谱及其敏感谱利用 \(\lambda\) 由命题 1, c) 与 III, p. 90 命题 5 计算。

注记. --- 当 I 无限时，条件 \(\lambda\in\mathcal{C}_{0}(\mathrm{I};\mathrm{K})\) 也可表述为“族 \(\lambda\) 沿 I 的有限子集的余集的滤子趋于 0”。

\subsection{紧自同态的对角化}

定理 1. --- 假设 K = C。设 u 是 E 的紧且正规的自同态。存在 E 的一个标准正交基 B，使 u 在基 B 中是对角的。

集 \(\mathrm{Sp}_{\mathrm{s}}(u)\) 是可数的，且当 E 无限时它不含 0（III, p. 90, 命题 5, \(b\) ）。对每个元素 \(\lambda \in \mathrm{Sp}_{\mathrm{s}}(u)\) ，用 \(\mathrm{N}_{\lambda}\) 表示 \(u - \lambda 1_{\mathrm{E}}\) 的幂零空间。它是有限维的（同处），且由于 \(u\) 正规，它与 \(u\) 的相对于 \(\lambda\) 的特征空间重合（EVT, V, p. 43, 命题 8 的推论）。诸空间 \(\mathrm{N}_{\lambda}\) 两两正交（I, p. 132, 第 5 目）。

设 F 是 E 的这样的子空间：它是诸空间 \(N_{\lambda}\) （对 \(\lambda \in \mathrm{Sp}_{\mathrm{s}}(u) - \{0\}\) ）的希尔伯特和。它是可数型的空间，在 u 下稳定，因为每个子空间 \(N_{\lambda}\) 在 u 下稳定。由于 \(N_{\lambda}\) 也是 \(u^{*}\) 的相对于 \(\overline{\lambda}\) 的特征空间（EVT, V, 同处），自同态 u 通过过渡到子空间诱导一个自同态 \(\widetilde{u}\) （\(F^{\circ}\) 的）。自同态 \(\widetilde{u}\) 是紧的（III, p. 5, 命题 3）且正规的（I, p. 135, 引理 4）。按构造，\(\widetilde{u}\) 的敏感谱包含于 \(\{0\}\) （事实上，\(\widetilde{u}\) 的每个特征向量也会是 u 的特征向量，从而若对应的特征值非零，它就属于某个空间 \(N_{\lambda}\) ）。于是 \(\widetilde{u}\) 的谱半径为零，从而 \(\widetilde{u}=0\) （因为 \(\widetilde{u}\) 正规；I, p. 108, 推论 1）。因此我们有 \(F^{\circ}\subset\operatorname{Ker}(u)\) 。

对每个 \(\lambda \in \mathrm{Sp}_{\mathrm{s}}(u)\) ，设 \(B_{\lambda}\) 是 \(N_{\lambda}\) 的一个标准正交基，设 \((e_j)_{j \in J}\) 是由诸 \(B_{\lambda}\) 的并组成的族；它是 F 的一个标准正交基。设 B 是 \((e_j)_{j \in J}\) 与 \(F^{\circ}\) 的一个标准正交基之并。这是 E 的一个标准正交基，且 \(u\) 在基 B 中是对角的。

推论 1. --- 设 u 是 E 的紧埃尔米特自同态。存在 E 的一个标准正交基 B，使 u 在基 B 中是对角的且其特征值为实数。

若 K = C，这立即由定理得出。假设 K = R。空间 E 是 \(\mathrm{E}_{(\mathbf{C})}\) 上的一个 R-结构（参见 A, II, p. 119）。自同态 \(u_{(\mathbf{C})}\) （\(\mathrm{E}_{(\mathbf{C})}\) 的）是紧的（III, p. 2, 注记 4）且埃尔米特的。设 \(B = (e_j)_{j \in J}\) 是 \(\mathrm{E}_{(\mathbf{C})}\) 的一个标准正交基，使得 \(u_{(\mathbf{C})} \in \mathcal{D}_{\mathrm{B}}(\mathrm{E}_{(\mathbf{C})})\) ，设 \(\lambda\) 是 \(u_{(\mathbf{C})}\) 的特征值族（定理 1）。我们有 \(\lambda \in R^{J}\) （I, p. 106, 命题 4）；由于线性映射 \(u_{(\mathbf{C})}\) 是 R-有理的，\(u_{(\mathbf{C})}\) 的相对于 \(\lambda_j\) 的特征空间对每个 \(j \in J\) 都是 R-有理的（参见 A, V, p. 60, 命题 6）。因此存在它的一个属于 E 的基，且更有理由存在 E 中的一个标准正交基。这些基之并是 E 的一个标准正交基 \(B_R\) ，使得 \(u \in \mathcal{D}_{\mathrm{B}_R}(\mathrm{E})\) 。

推论 2. --- 设 F 是希尔伯特空间，u 是从 E 到 F 的紧连续线性映射。存在 u 的始空间的一个标准正交基 \((e_{i})_{i\in\mathrm{I}}\) （\(\operatorname{Ker}(u)^{\circ}\) 的）、F 的一个标准正交族 \((f_{i})_{i\in\mathrm{I}}\) 以及一个族 \((\alpha_{i})_{i\in\mathrm{I}}\in(\mathbf{R}_{+}^{*})^{\mathrm{I}}\) ，使得 \(u(e_{i})=\alpha_{i}f_{i}\) （对所有 \(i\in I\) ）。

设 \(v = u^{*} \circ u\) 。这是 E 的一个紧的、正的（从而埃尔米特的）自同态。由推论 1，存在 E 的一个标准正交基 \((e_{j})_{j \in \mathrm{J}}\) ，使得 \(v\) 在此基中是对角的。它的特征值族 \((\lambda_{j})_{j \in \mathrm{J}}\) 包含于 \(\mathbf{R}_{+}^{\mathrm{J}}\) 。置 \(\alpha_{j} = \sqrt{\lambda_{j}}\) （对所有 \(j \in \mathrm{J}\) ）。设 I 是由 \(j \in \mathrm{J}\) （使 \(\alpha_{j} \neq 0\) ）组成的集。由于 \(v\) 紧，这是一个可数集。族 \((e_{i})_{i \in \mathrm{I}}\) 是 \(v\) 的始空间（即始空间 \(\operatorname{Ker}(u)^{\circ}\) ，\(u\) 的）的一个标准正交基（EVT, V, p. 43, 命题 8）。置 \(f_{i} = \frac{1}{\alpha_{i}} u(e_{i})\) （对 \(i\in \mathrm{I}\) ）。对 I 中任何 \(i\) 与 \(j\) ，我们有

\[
\langle f _ {i} | f _ {j} \rangle = \frac {1}{\alpha_ {i} \alpha_ {j}} \langle u (e _ {i}) | u (e _ {j}) \rangle = \frac {1}{\alpha_ {i} \alpha_ {j}} \langle v (e _ {i}) | e _ {j} \rangle = \frac {\lambda_ {i}}{\alpha_ {i} \alpha_ {j}} \langle e _ {i} | e _ {j} \rangle ,
\]

由此得出族 \((f_{i})_{i\in\mathrm{I}}\) 在 F 中是标准正交的。由于 \(u(e_{i})=\alpha_{i}f_{i}\) （对所有 \(i\in I\) ），推论得证。

定义 2. --- 沿用推论的记号，族 \((\alpha_{i})_{i\in I}\) 是 \(u\) 的相对于始空间的标准正交基 \((e_{i})_{i\in I}\) （该始空间为 \(u\) 的始空间）的奇异值族。

注记. --- 1) 该推论推广了 EVT, V, p. 54 的定理 2，后者对应于希尔伯特–施密特映射。

\begin{enumerate}
\def\labelenumi{\arabic{enumi})}
\setcounter{enumi}{1}
\tightlist
\item
  沿用推论的记号，我们有公式
\end{enumerate}

\[
u (x) = \sum_ {i \in I} \alpha_ {i} \langle e _ {i} | x \rangle f _ {i}\tag{1}
\]

对所有 \(x \in E\) 成立。

\begin{enumerate}
\def\labelenumi{\arabic{enumi})}
\setcounter{enumi}{2}
\tightlist
\item
  设 \(u\) 是 E 的紧正自同态。设 \(B = (f_i)_{i \in I}\) 是 E 的一个标准正交基，使得 \(u\) 在基 B 中是对角的（定理 1），设 \((\lambda_i)_{i \in I}\) 是 \(u\) 在此基中的特征值族。设 J 是由 \(i \in I\) （使 \(\lambda_i > 0\) ）组成的集；族 \((e_i)_{i \in J}\) 是空间 \(\text{Ker}(u)^{\circ}\) 的一个标准正交基。对每个 \(i \in J\) ，置 \(e_i = f_i\) 与 \(\alpha_i = \lambda_i\) 。我们得到 \(u(e_i) = \alpha_if_i\) ：族 \((\alpha_i)_{i \in J}\) 是 \(u\) 的相对于基 \((e_i)_{i \in J}\) 的奇异值族。
\end{enumerate}

\subsection{特征值的递减序列}

在本节中，我们假设 K = C。

我们置 \(\overline{\mathbf{N}} = \mathbf{N} \cup \{+\infty\} \subset \overline{\mathbf{R}}\) 。在本节中，若向量空间 E 不是有限维的，就说它有维数 \(+\infty \in \overline{\mathbf{N}}\) 。设 \(\mathrm{I_E} \subset \mathbf{N}\) 是 E 的满足 \(\mathrm{F} \neq \mathrm{E}\) 的有限维子空间 F 的维数所成的集。若 E 无穷维，我们有 \(\mathrm{I_E} = \mathbf{N}\) ，否则 \(\mathrm{I_E} = \{0, \ldots, \dim(\mathrm{E}) - 1\}\) 。当不致引起混淆时，我们置 \(\mathrm{I} = \mathrm{I_E}\) 。

设 \(u\) 是 E 的紧正（特别地，埃尔米特）自同态。\(u\) 的敏感谱是正实数的一个严格递减序列 \((\nu_{k})_{0\leqslant k<\mathrm{Card}(\mathrm{Sp}_{\mathrm{s}}(u))}\) 的值集（参见

III, p. 90, 命题 5）。对每个整数 \(k\) （满足 \(0 \leqslant k < \text{Card}(\text{Sp}_s(u))\) ），用 \(n_k\) 表示 \(\nu_k\) 的谱重数，其中 \(n_k \geqslant 1\) 。设 \(\text{M} \in \overline{\mathbf{N}}\) 是谱重数 \(n_k\) 之和；它是 \(u\) 的像的维数。我们有 \(\text{M} \leqslant \text{Card}(\text{I})\) 。

对 \(0 \leqslant n < M\) ，定义 \(\lambda_n(u) = \nu_k\) ，其中 \(k \geqslant 0\) 是唯一使下式成立的整数：

\[
n _ {0} + \dots + n _ {k - 1} \leqslant n <   n _ {0} + \dots + n _ {k}.
\]

置 \(\lambda_{n}(u)=0\) （若 \(n\in I\) 满足 \(n\geqslant M\) ）。这种情形只能在 I=N 且 \(\mathrm{Sp}_{\mathrm{s}}(u)\) 有限时发生（或者等价地，E 无穷维且 u 有限秩）。

序列 \((\lambda_{n}(u))_{n\in\mathrm{I}}\) 是递减的；对每个 \(\lambda\in\mathrm{Sp}_{\mathrm{s}}(u)\) ，使 \(\lambda_{n}(u)=\lambda\) 的整数 n 的个数等于 u 的特征值 \(\lambda\) 的谱重数。

我们约定俗成地说，\((\lambda_{n}(u))_{n\in\mathrm{J}}\) 是 u 的按重数重复的特征值递减序列。

命题 3. --- 设 u 是 E 的紧正自同态。存在 E 中的一个标准正交族 \((e_{n})_{n\in\mathrm{I}}\) ，使得对每个 \(x\in E\) ，有

\[
u (x) = \sum_ {n \in \mathrm{I}} \lambda_ {n} (u) \langle e _ {n} | x \rangle e _ {n}, \quad \langle x | u (x) \rangle = \sum_ {n \in \mathrm{I}} \lambda_ {n} (u) | \langle e _ {n} | x \rangle | ^ {2}.
\]

设 \(\mathrm{B} = (f_{j})_{j \in \mathrm{J}}\) 是 E 的使 u 为对角的标准正交基（IV, p. 150, 推论 1），\((\lambda_{j})_{j \in \mathrm{J}}\) 是 u 在基 B 中的特征值族。设 J' 是由 \(j \in J\) （使 \(\lambda_{j}\) 属于 u 的敏感谱）组成的集。

对每个 \(\lambda \in \mathrm{Sp}_{\mathrm{s}}(u)\) ，在整数 \(n\) （满足 \(0 \leqslant n < \mathrm{M}\) 与 \(\lambda_{n}(u) = \lambda\) ）与 \(j \in \mathrm{J}'\) （满足 \(\lambda_{j} = \lambda\) ）之间存在一个双射，因为这两个集有相同的基数，即 \(\lambda\) 的谱重数。对每个 \(\lambda\) 选取这样的双射，就定义了双射 \(\iota\) （从满足 \(0 \leqslant n < \mathrm{M}\) 的整数所成的集到集 \(\mathrm{J}'\) 的）。序列 \((e_n)_{0 \leqslant n < \mathrm{M}}\) 由置 \(e_n = f_{\iota(n)}\) （对 \(0 \leqslant n < \mathrm{M}\) ）在 E 中定义。它是 E 中的一个标准正交族。

在 M \textless{} Card(I) = +∞ 的情形，由 \(\{e_{0},\ldots,e_{M-1}\}\) 生成的空间 F 是有限维的，而它的正交 F° 是无穷维的；对 \((e_{n})_{n\geqslant M}\) 选取 F° 中的一个标准正交族。

对每个 \(x \in \mathrm{E}\) ，我们有

\[
u (x) = \sum_ {j \in J ^ {\prime}} \lambda_ {j} \langle f _ {j} \mid x \rangle f _ {j} = \sum_ {0 \leqslant n <   M} \lambda_ {n} (u) \langle e _ {n} \mid x \rangle e _ {n}
\]

因为 \(u(f_{j}) = 0\) （当 \(j \in J - J'\) ）。若 \(n \in I\) 满足 \(n \geqslant M\) ，我们有 \(\lambda_{n}(u) = 0\) ，由此得到命题的第一个公式。第二个公式由它得出。

命题 4. --- 设 u 是 E 的紧正自同态。设 \(f \in \mathcal{C}(\mathbf{R}_{+})\) 是连续递增函数，满足 \(f(0) = 0\) 。

自同态 \(f(u)\) 是紧且正的，且对每个 \(n \in \mathrm{I}_{\mathrm{E}}\) ，我们有 \(\lambda_n(f(u)) = f(\lambda_n(u))\) 。

由 III, p. 91, 命题 6, b) 与 I, p. 117, 命题 15, a)，自同态 \(f(u)\) 是紧且正的。\(f(u)\) 的谱是在 \(f\) 下 \(u\) 的谱的像（I, p. 111, 推论 2）。若 \(\lambda \in \mathrm{Sp}_{\mathrm{s}}(f(u))\) ，则 \(\lambda\) 的谱重数是 \(\mu \in \mathrm{Sp}(u)\) （满足 \(f(\mu) = \lambda\) ）的谱重数之和（III, p. 84, 推论 2）。由于 \(f\) 递增，序列 \((f(\lambda_n(u)))_{n \in \mathrm{I_E}}\) 是递减的。断言于是由序列 \((\lambda_n(f(u)))_{n \in \mathrm{I_E}}\) 的定义得出。

\subsection{特征值的变分刻画}

在本节中，我们假设 K = C。

设 \(u\) 是 E 的紧正自同态，\((\lambda_n(u))_{n\in \mathrm{I_E}}\) 是 \(u\) 的特征值递减序列。我们置 \(\mathrm{I} = \mathrm{I_E}\) 。

对 E 的每个闭子空间 F，我们记

\[
r _ {\mathrm{F}} (u) = \inf _ {x \in \mathrm{F} - \{0 \}} \frac {\langle x \mid u (x) \rangle}{\| x \| ^ {2}}, \quad \mathrm{R} _ {\mathrm{F}} (u) = \sup _ {x \in \mathrm{F} ^ {\circ} - \{0 \}} \frac {\langle x \mid u (x) \rangle}{\| x \| ^ {2}},
\]

其中下确界（相应地，上确界）是在 \([0, +\infty]\) 中取的。

对每个 \(n \in \mathbf{N}\) ，用 \(\mathcal{F}_n\) 表示 E 的向量子空间 \(\mathrm{F} \subset \mathrm{E}\) （维数为 \(n\) ）所成的集。我们称子空间 \(\mathrm{F} \in \mathcal{F}_n\) 适应于 \(u\) ，如果它有一个标准正交基 \((f_i)_{0 \leqslant i \leqslant n-1}\) ，使得 \(u(f_i) = \lambda_i(u)f_i\) （对 \(0 \leqslant i \leqslant n-1\) ）。

命题 5. --- 设 \(n \in I\) 。

\begin{enumerate}
\def\labelenumi{\alph{enumi})}
\item
  对每个子空间 \(\mathrm{F} \in \mathcal{F}_{n+1}\) （适应于 \(u\) 的），我们有 \(\lambda_n(u) = r_{\mathrm{F}}(u)\) ；
\item
  对每个子空间 \(F \in F_{n}\) （适应于 u 的），我们有 \(\lambda_{n}(u) = R_{\mathrm{F}}(u)\) 。
\end{enumerate}

设 \(F \in F_{n+1}\) 是适应于 u 的子空间，\((f_i)_{0 \leqslant i \leqslant n}\) 是 F 的一个标准正交基，使得 \(u(f_i) = \lambda_i(u)f_i\) （对所有 i）。对 F 中所有 x，我们有

\[
\langle x \mid u (x) \rangle = \sum_ {0 \leqslant i \leqslant n} \lambda_ {i} (u) | \langle f _ {i} \mid x \rangle | ^ {2} \geqslant \lambda_ {n} (u) \sum_ {0 \leqslant i \leqslant n} | \langle f _ {i} \mid x \rangle | ^ {2} = \lambda_ {n} (u) \| x \| ^ {2},
\]

其中等式当 \(x = f_{n}\) 时成立。这蕴含 \(r_{\mathrm{F}}(u) = \lambda_{n}(u)\) ，从而得到断言 (a)。

设 \(F \in F_{n}\) 是适应于 u 的子空间。闭子空间 \(F^{\circ}\) 非零（因为 \(n = \dim(F) < \dim(E)\) ）且在 u 下稳定；由 u 通过过渡到子空间得到的自同态 \(\widetilde{u}\) （\(F^{\circ}\) 的）是紧且正的。

有 \(\mathrm{Sp}(\widetilde{u})\subset \mathrm{Sp}(u)\) 。我们还有 \(\mathrm{Sp}(\widetilde{u})\subset [0,\lambda_n(u)]\) 。事实上，只需验证 \(\lambda \leqslant \lambda_{n}(u)\) （对所有 \(\lambda \in \mathrm{Sp}_{\mathrm{s}}(\widetilde{u})\) ）。数 \(\lambda\) 这时是 \(u\) 的特征值，故存在 \(j\in I\) 使 \(\lambda = \lambda_j(u)\) 。于是 \(u\) 的相对于 \(\lambda_{j}(u)\) 的特征空间不包含于 F，这蕴含 \(\lambda_{j}(u)\leqslant \lambda_{n}(u)\) 。

假设 \(\lambda_n(u) > 0\) 。\(u\) 的相对于 \(\lambda_n(u)\) 的特征空间于是不包含于 F，从而 \(\lambda_n(u)\) 属于 \(\widetilde{u}\) 的敏感谱。我们推出 \(\sup (\mathrm{Sp}(\widetilde{u})) = \lambda_n(u)\) 。若 \(\lambda_n(u) = 0\) ，由于 \(\widetilde{u}\) 的谱这时化为 \(\{0\}\) ，我们得到同样的结果。

此外，按定义我们有 \(\mathrm{R}_{\mathrm{F}}(u)=\mathrm{R}_{\{0\}}(\widetilde{u})\) ，最后按 I, p. 139, 命题 9, a)，有 \(\mathrm{R}_{\{0\}}(\widetilde{u})=\sup(\mathrm{Sp}(\widetilde{u}))\) 。断言 b) 得证。

命题 6. --- 对所有 \(n \in I\) ，我们有

\[
\lambda_ {n} (u) = \sup _ {\mathrm{F} \in \mathscr {F} _ {n + 1}} r _ {\mathrm{F}} (u) = \inf _ {\mathrm{F} \in \mathscr {F} _ {n}} \mathrm{R} _ {\mathrm{F}} (u).
\]

设 \((e_{n})_{n\in\mathrm{I}}\) 是具有 IV, p. 152 命题 3 的性质的标准正交族。对每个整数 \(n\) （满足 \(1\leqslant n<M+1\) ），设 \(F_{n}\in\mathcal{F}_{n}\) 是由 \((e_{0},\ldots,e_{n-1})\) 生成的 E 的 n 维子空间；按构造，空间 \(F_{n}\) 适应于 u。于是按命题 5，我们有

\[
\lambda_ {n} (u) = r _ {\mathrm{F} _ {n + 1}} (u) = \mathrm{R} _ {\mathrm{F} _ {n}} (u).\tag{2}
\]

设 \(n \in I\) 。设 \(F \in \mathcal{F}_{n+1}\) 。在 \(F\) 上的限制（到 \(F_n\) 上的正交投影的）不是单射，故存在 \(x \neq 0\) （\(F\) 中的，正交于 \(F_n\) 的）。

由于 \(x \in \mathrm{F}_{n}^{\circ}\) ，我们这时有（IV, p. 152, 命题 3）

\[
\begin{array}{l}\langle x\mid u(x)\rangle = \sum_{\substack{m\in \mathrm{I}\\ m\geqslant n}}\lambda_{m}(u)|\langle e_{m}\mid x\rangle |^{2}\\ \leqslant \lambda_{n}(u)\sum_{\substack{m\in \mathrm{I}\\ m\geqslant n}}|\langle e_{m}\mid x\rangle |^{2} = \lambda_{n}(u)\| x\|^{2}. \end{array}
\]

这证明 \(r_{\mathrm{F}}(u)\leqslant\lambda_{n}(u)\) ，从而特别地有不等式

\[
\sup _ {\mathrm{F} \in \mathscr {F} _ {n + 1}} r _ {\mathrm{F}} (u) \leqslant \lambda_ {n} (u).\tag{3}
\]

设 \(F \in F_{n}\) 。在 \(F_{n+1}\) 上的限制（到 F 上的正交投影的）不是单射，故存在向量 \(x \neq 0\) （\(F_{n+1}\) 中的，正交于 F 的）。由于 \(x \in F_{n+1}\) ，我们有（同处）

\[
\begin{array}{c} \langle x \mid u (x) \rangle = \sum_ {0 \leqslant m \leqslant n} \lambda_ {m} (u) | \langle e _ {m} \mid x \rangle | ^ {2} \\ \geqslant \lambda_ {n} (u) \sum_ {0 \leqslant m \leqslant n} | \langle e _ {m} \mid x \rangle | ^ {2} = \lambda_ {n} (u) \| x \| ^ {2} \end{array}
\]

因而 \(\mathrm{R_F}(u) \geqslant \lambda_n(u)\) 。特别地，由此得出

\[
\inf _ {\mathrm{F} \in \mathcal {F} _ {n}} \mathrm{R} _ {\mathrm{F}} (u) \geqslant \lambda_ {n} (u).\tag{4}
\]

鉴于公式 (2)、(3) 与 (4)，命题得证。

\subsection{特征值变分刻画的应用}

在本节中，我们考虑 C 上的希尔伯特空间。

命题 7. --- 设 u 与 v 是 E 的正紧自同态。

\begin{enumerate}
\def\labelenumi{\alph{enumi})}
\item
  有 \(|\lambda_n(u) - \lambda_n(v)|\leqslant \| u - v\|\) （对所有 \(n\in \mathbf{I}\) ）；
\item
  若 \(u \leqslant v\) ，则 \(\lambda_n(u) \leqslant \lambda_n(v)\) （对所有 \(n \in I\) ）。
\end{enumerate}

设 \(n \in I\) ，设 F 是 E 的 n 维向量子空间。对所有 \(x \in F\) ，我们有

\[
| \langle x | v (x) \rangle - \langle x | u (x) \rangle | \leqslant \| u - v \| \| x \| ^ {2},
\]

因此有不等式

\[
\mathrm{R} _ {\mathrm{F}} (v) - \| u - v \| \leqslant \mathrm{R} _ {\mathrm{F}} (u) \leqslant \mathrm{R} _ {\mathrm{F}} (v) + \| u - v \|.
\]

断言 \(a\) ) 由此按 IV, p. 154 命题 6 得出。

若 \(u \leqslant v\) ，我们有 \(\langle x | u(x) \rangle \leqslant \langle x | v(x) \rangle\) （对所有 \(x \in E\) ）。因此对每个 \(n \in I\) 与每个子空间 \(F\) （维数为 \(n\) 的），我们有 \(R_F(u) \leqslant R_F(v)\) ，从而 \(\lambda_n(u) \leqslant \lambda_n(v)\) （同处）。

命题 8. --- 设 u 是 E 的正紧自同态。设 H 是 E 的一个闭子空间，\(i_{\mathrm{H}}: \mathrm{H} \to \mathrm{E}\) 是典范单射。用 \(u_{\mathrm{H}}\) 表示 H 的自同态 \(i_{\mathrm{H}}^{*}ui_{\mathrm{H}}\) 。它是紧且正的。

\begin{enumerate}
\def\labelenumi{\alph{enumi})}
\item
  有 \(\mathrm{I_H}\subset \mathrm{I_E}\) 与 \(\lambda_n(u_{\mathrm{H}})\leqslant \lambda_n(u)\) （对所有 \(n\in \mathrm{I_H}\) ）；
\item
  若 H 在 E 中有有限余维数 \(k \in N\) ，则我们有 \(I_{H} + k \subset I_{E}\) 与 \(\lambda_{n+k}(u) \leqslant \lambda_{n}(u_{\mathrm{H}})\) （对所有 \(n \in I_{H}\) ）。
\end{enumerate}

自同态 \(u_{\mathrm{H}}\) 是紧的（III, p. 5, 命题 3）。它是正的，因为 \(\langle x \mid u_{\mathrm{H}}(x) \rangle = \langle i_{\mathrm{H}}(x) \mid u(i_{\mathrm{H}}(x)) \rangle \geqslant 0\) （对所有 \(x \in \mathrm{H}\) ）。

设 \(n \in I_{H} \subset I_{E}\) 。设 F 是 H 的一个维数为 \(n + 1\) 的适应于 \(u_{H}\) 的子空间。于是我们有 \(\lambda_{n}(u_{\mathrm{H}}) = r_{\mathrm{F}}(u_{\mathrm{H}})\) （IV, p. 153, 命题 5, a)），又由于 \(r_{\mathrm{F}}(u_{\mathrm{H}}) = r_{\mathrm{F}}(u) \leqslant \lambda_{n}(u)\) （IV, p. 154, 命题 6），我们得到断言 a)。

假设 H 在 E 中有余维数 \(k \in N\) 且 \(n \in I_{H}\) 。设 F 是 H 的一个维数为 n 的适应于 \(u_{H}\) 的子空间。它在 H 中的正交等于 \(H \cap F^{\circ}\) ，且它是 E 中子空间 \(F + H^{\circ}\) （维数为 \(n + k\) ）的正交。因此 \(n + k \in I_{E}\) ，且（IV, p. 153, 命题 5, b)）

\[
\lambda_{n}(u_{\mathrm{H}}) = \sup_{\substack{x\in \mathrm{H}\cap \mathrm{F}^{\circ}\\ x\neq 0}}\frac{\langle x\mid u_{\mathrm{H}}(x)\rangle}{\|x\|^{2}} = \mathrm{R}_{\mathrm{F} + \mathrm{H}^{\circ}}(u)
\]

从而 \(\lambda_{n}(u_{\mathrm{H}})\leqslant \lambda_{n + k}(u)\) （IV, p. 154, 命题 6）。

定义 3. --- 设 F 是希尔伯特空间，u 是从 E 到 F 中的一个紧线性映射。

对每个整数 \(n \in \mathrm{I}_{\mathrm{E}}\) ，置 \(\alpha_{n}(u) = \lambda_{n}(u^{*} \circ u)\) 。设 J 是由 \(n \in \mathrm{I}_{\mathrm{E}}\) （使 \(\alpha_{n}(u) > 0\) ）组成的集。族 \((\alpha_{n}(u))_{n \in \mathrm{J}}\) 称为 \(u\) 的按重数重复的奇异值序列。

我们称序列 \((\alpha_{n}(u))_{n\in\mathrm{I}_{\mathrm{E}}}\) 为 u 的增广奇异值序列。

序列 \((\alpha_{n}(u))_{n\in\mathrm{I}_{\mathrm{E}}}\) 是完全定义的，因为 E 的自同态 \(u^{*}\circ u\) 是紧的（III, p. 5, 命题 3）；由于 \(u^{*}\circ u\) 是正的，它是一个正实数的递减族。

命题 9. --- 设 F 是希尔伯特空间，u 是从 E 到 F 中的一个紧线性映射。

\begin{enumerate}
\def\labelenumi{\alph{enumi})}
\tightlist
\item
  对 \(n \in \mathrm{I}_{\mathrm{E}}\) ，我们有
\end{enumerate}

\[
\alpha_ {n} (u) = \sup _ {\mathrm{F} \in \mathscr {F} _ {n + 1}} \inf _ {x \in \mathrm{F} - \{0 \}} \frac {\| u (x) \|}{\| x \|} = \inf _ {\mathrm{F} \in \mathscr {F} _ {n}} \sup _ {x \in \mathrm{F} ^ {\circ} - \{0 \}} \frac {\| u (x) \|}{\| x \|}.
\]

\begin{enumerate}
\def\labelenumi{\alph{enumi})}
\setcounter{enumi}{1}
\tightlist
\item
  设 J 是由 \(n \in I_{E}\) （使 \(\alpha_{n}(u) \neq 0\) ）组成的集。存在 E 中的标准正交族 \((e_{n})_{n \in \mathrm{J}}\) 与 F 中的标准正交族 \((f_{n})_{n \in \mathrm{J}}\) ，使得对每个 \(x \in E\) ，我们有
\end{enumerate}

\[
u (x) = \sum_ {n \in J} \alpha_ {n} (u) \langle e _ {n} | x \rangle f _ {n}.
\]

由于 \(\langle x \mid u^{*}u(x)\rangle = \|u(x)\|^{2}\) （对所有 \(x \in E\) ），诸定义与 IV, p. 154, 命题 6, a) 蕴含第一个断言中的等式。设 \((e_{n})_{n\in\mathrm{I}_{\mathrm{E}}}\) 是 E 的一个标准正交族，满足把 IV, p. 152 命题 3 应用于 E 的紧正自同态 \(u^{*} \circ u\) 所得的结论。置 \(f_{n} = \alpha_{n}^{-1}u(e_{n})\) （对 \(n \in J\) ）。按 IV, p. 150 推论 2 的证明中的推理，我们得到断言 c)。

推论. --- 设 F 是希尔伯特空间，u 是从 E 到 F 中的一个紧线性映射。设 \(w \in \mathcal{L}(\mathrm{E})\) 与 \(v \in \mathcal{L}(\mathrm{F})\) 。对每个 \(n \in I_{E}\) ，我们有 \(\alpha_{n}(w \circ u \circ v) \leqslant \|v\|\) \(\|w\|\) \(\alpha_{n}(u)\) 。

这是前一命题断言 \(a\) ) 的推论。

注记. --- 1) 序列 \((\alpha_{n}(u))_{n\in\mathrm{I}_{\mathrm{E}}}\) 是递减的，故集 J 或等于 \(\mathrm{I}_{\mathrm{E}}\) ，或等于线段 \(\{0,\ldots,m\}\) （\(\mathrm{I}_{\mathrm{E}}\) 中的）（其中 \(m\in\mathbf{N}\) ）。后一情形成立当且仅当 \(u\) 有限秩。

\begin{enumerate}
\def\labelenumi{\arabic{enumi})}
\setcounter{enumi}{1}
\item
  由于 \(\| u(x)\| = \| |u|(x)\|\) （对所有 \(x\in \mathrm{E}\) ；I, p. 139, 命题 10），我们有 \(\alpha_{n}(u) = \alpha_{n}(|u|)\) （对所有 \(n\in \mathrm{I}_{\mathrm{E}}\) ）。
\item
  若 \(u\) 是正的，则 \(\alpha_{n}(u)=\lambda_{n}(u)\) （对所有 \(n\in\mathrm{I}_{\mathrm{E}}\) ；事实上，按 IV, p. 153 命题 4，我们这时有 \(\alpha_{n}(u)^{2}=\lambda_{n}(u^{*}u)=\lambda_{n}(u^{2})=\lambda_{n}(u)^{2}\) ）。
\item
  有可能有 \(\alpha_{n}(v\circ u\circ w)=0\) ，即使 \(\alpha_{n}(u)\) 非零；在这种情形，\(\alpha_{n}(v\circ u\circ w)\) 不是 \(v\circ u\circ w\) 的奇异值。
\end{enumerate}
\subsection{外尔不等式}

在本节中，我们考虑 K = C 上的希尔伯特空间。

设 E 是希尔伯特空间。对每个 \(n \in N\) ，我们回顾希尔伯特外幂 \(\hat{\Lambda}^{n}E\) 已在 EVT, V, p. 34 中定义。对每个希尔伯特空间 F 以及每个 \(u \in \mathcal{L}(E;F)\) ，我们还定义了线性映射 \(\hat{\Lambda}^{n}u \in \mathcal{L}(\hat{\Lambda}^{n}E;\hat{\Lambda}^{n}F)\) （同前）。这些构造是函子性的：对每个希尔伯特空间 G 以及每个线性映射 \(v \in \mathcal{L}(F;G)\) ，公式

\[
\widehat {\Lambda} ^ {n} v \circ \widehat {\Lambda} ^ {n} u = \widehat {\Lambda} ^ {n} (v \circ u)
\]

成立（同前，公式 (28)）。

设 H 是 E 的闭子空间，并设 \(i_{H}\) 是从 H 到 E 中的典范单射。对 E 的每个自同态 u，我们用 \(u_{H}\) 表示 H 的自同态 \(i_{H}^{*}ui_{H}\) 。

引理 2. --- 设 \(n \in N\) 。映射 \(\hat{\Lambda}^{n}i_{H}\) 是从 \(\hat{\Lambda}^{n}H\) 到 \(\hat{\Lambda}^{n}E\) 中的线性等距映射。

设 \((e_{j})_{j\in\mathrm{J}}\) 是 H 的标准正交基，\((e_{j})_{j\in\mathrm{J}^{\prime}}\) 是 E 的标准正交基，其中 \(J\subset J^{\prime}\) 。给 \(J^{\prime}\) 赋予一个全序。元素 \(e_{j_{1}}\wedge\cdots\wedge e_{j_{n}}\) （其中 \(j_{1}<\cdots<j_{n}\) 在 \(J^{\prime}\) 中（相应地，在 J 中）变化）构成 \(\widehat{\Lambda}^{n}E\) （相应地，\(\widehat{\Lambda}^{n}H\) ）的标准正交基（EVT, V, p. 34 的命题 5）。引理得证。

在下面，我们将把 \(\hat{\Lambda}^n\mathrm{H}\) 等同于 \(\hat{\Lambda}^n\mathrm{E}\) 的一个闭子空间，借助的是映射 \(\hat{\Lambda}^n i_{\mathrm{H}}\) 。

引理 3. --- 设 F 是希尔伯特空间，且 \(u \in \mathcal{L}(\mathrm{E};\mathrm{F})\)。设 \(n \in \mathbf{N}\) 。a) 我们有 ( \(\widehat{\Lambda}^n u\) ) \(^*\) = \(\widehat{\Lambda}^n (u^*)\) ；

\begin{enumerate}
\def\labelenumi{\alph{enumi})}
\setcounter{enumi}{1}
\item
  若 F = E 且 u 是埃尔米特的（相应地，正的、正规的、酉的），则 \(\hat{\Lambda}^{n}u\) 是埃尔米特的（相应地，正的、正规的、酉的）；
\item
  我们有 \(|\widehat{\Lambda}^n u| = \widehat{\Lambda}^n |u|\) ；
\item
  设 H 是 E 的闭子空间。限制 \((\widehat{\Lambda}^{n}u)|\widehat{\Lambda}^{n}\mathrm{H}\) （即 \(\widehat{\Lambda}^{n}u\) 在 \(\widehat{\Lambda}^{n}\mathrm{H}\) 上的限制）等于 \(\widehat{\Lambda}^{n}(u|\mathrm{H})\) （等式在 \(\mathcal{L}(\widehat{\Lambda}^{n}\mathrm{H};\widehat{\Lambda}^{n}\mathrm{F})\) 中成立）；
\item
  假设 F = E。设 H 是 E 的闭子空间。我们有 \((\hat{\Lambda}^{n}u)_{\hat{\Lambda}^{n}\mathrm{H}}=\hat{\Lambda}^{n}(u_{\mathrm{H}})\) （在 \(\mathcal{L}(\hat{\Lambda}^{n}\mathrm{H})\) 中）。
\end{enumerate}

断言 \(a\)）由 EVT, V, p. 39 的公式 (13) 得出。它立即蕴含：\(\widehat{\Lambda}^{n}u\) 在 \(u\) 是埃尔米特的（相应地，酉的）时是埃尔米特的（相应地，酉的）。若 \(u\) 是正的，则 \(u^{1/2}\) 是埃尔米特的，因而 \(\widehat{\Lambda}^{n}(u^{1/2})\) 也是；关系式 \(\widehat{\Lambda}^{n}u = (\widehat{\Lambda}^{n}u^{1/2})^{2}\) 蕴含 \(\widehat{\Lambda}^{n}u\) 是正的（I, p. 138，命题 8）。这就证明了 \(b\)）。

利用前述结果，可以算得

\[
\begin{array}{r l} & {\left(\widehat {\Lambda} ^ {n} | u |\right) ^ {2} = \widehat {\Lambda} ^ {n} (| u | ^ {2}) = \widehat {\Lambda} ^ {n} (u ^ {*} u)} \\ & {\qquad = \widehat {\Lambda} ^ {n} u ^ {*} \widehat {\Lambda} ^ {n} u = \left(\widehat {\Lambda} ^ {n} u\right) ^ {*} (\widehat {\Lambda} ^ {n} u) = | \widehat {\Lambda} ^ {n} u | ^ {2}.} \end{array}
\]

由于 \(\widehat{\Lambda}^n |u|\) 依 \(b\)）是正的，故由 I, p. 118 的命题 16 得 \(\widehat{\Lambda}^n |u| = |\widehat{\Lambda}^n u|\) ，由此得 \(c\)）。

设 \(i_{\mathrm{H}}\) 是从 H 到 E 中的典范单射。空间 \(\widehat{\Lambda}^{n}\mathrm{H}\) 等同于 \(\widehat{\Lambda}^{n}\mathrm{E}\) 的一个闭子空间，借助的是映射 \(\widehat{\Lambda}^{n}i_{\mathrm{H}}\) ，由此得

\[
(\widehat {\Lambda} ^ {n} u) | \widehat {\Lambda} ^ {n} \mathrm{H} = \widehat {\Lambda} ^ {n} u \circ \widehat {\Lambda} ^ {n} i _ {\mathrm{H}} = \widehat {\Lambda} ^ {n} (u \circ i _ {\mathrm{H}}) = \widehat {\Lambda} ^ {n} (u | \mathrm{H}),
\]

以及

\[
(\widehat {\Lambda} ^ {n} u) _ {\widehat {\Lambda} ^ {n} \mathrm{H}} = (\widehat {\Lambda} ^ {n} i _ {\mathrm{H}}) ^ {*} \circ \widehat {\Lambda} ^ {n} u \circ \widehat {\Lambda} ^ {n} i _ {\mathrm{H}} = \widehat {\Lambda} ^ {n} (i _ {\mathrm{H}} ^ {*} u i _ {\mathrm{H}}) = \widehat {\Lambda} ^ {n} (u _ {\mathrm{H}}),
\]

这就证明了 \(d\)）与 \(e\)）。

设 F 是希尔伯特空间，且 \(u \in \mathcal{L}^{c}(\mathrm{E};\mathrm{F})\) 。如同 IV, p. 151 的 no. 3 中那样，我们用 \(I_{E}\) 表示 E 的满足 \(F \neq E\) 的有限维子空间 F 的维数所成的集。我们回顾，\((\alpha_{n}(u))_{n \in I_{E}}\) 表示 u 的奇异值的扩张序列（IV, p. 156 的定义 3）。

命题 10. --- 我们有等式

\[
\prod_ {i = 0} ^ {n} \alpha_ {i} (u) = \left\| \widehat {\Lambda} ^ {n + 1} u \right\|
\]

对所有 \(n \in I_{E}\) 。

引理 3 的 c) 表明 \(\|\widehat{\Lambda}^{n+1}u\|=\|\widehat{\Lambda}^{n+1}|u\|\) ；此外，由于 \(\alpha_{i}(u)=\alpha_{i}(|u|)\) 对所有 \(i\in I_{E}\) 成立（IV, p. 157 的注 5, (2)），只需对 \(|u|\) 证明断言即可。于是可设 F=E 且 u 是正的。

于是设 \(\mathrm{B}=(e_{j})_{j\in\mathrm{J}}\) 是 E 的标准正交基，使得 u 在基 B 中可对角化（IV, p. 149 的定理 1），并设 \((\lambda_{j})_{j\in\mathrm{J}}\) 是 u 在基 B 中的特征值族。给 J 赋予一个全序，并用 \(J_{n}\) 表示 J 的元素的严格递增族 \((j_{0},\ldots,j_{n})\in\mathbf{J}^{n+1}\) 所成的集。向量

\[
e _ {\iota} = e _ {j _ {0}} \wedge \dots \wedge e _ {j _ {n}}
\]

（其中 \(\iota=(j_{0},\ldots,j_{n})\in\mathrm{J}_{n}\)）构成标准正交基 \(B_{n}\) ，它是 \(\widehat{\Lambda}^{n+1}\mathrm{E}\) 的标准正交基（EVT, V, p. 34，命题 5）。对每个 \(\iota=(j_{0},\ldots,j_{n})\in\mathrm{J}_{n}\) ，记

\[
\lambda_ {\iota} = \prod_ {j = 0} ^ {n} \lambda_ {i _ {j}}.
\]

由此可得 \((\hat{\Lambda}^{n + 1}u)e_{\iota} = \lambda_{\iota}e_{\iota}\) ，因而 \(\hat{\Lambda}^{n + 1}u\) 在基 \(\mathrm{B}_n\) 中是对角的。于是，我们有

\[
\| \widehat {\Lambda} ^ {n + 1} u \| = \sup _ {\iota \in \mathrm{I} _ {n}} \lambda_ {\iota}
\]

（IV, p. 147 的引理 1），它等于 \(\lambda_{0}(u)\cdots\lambda_{n}(u)\) ，即最大的 \(n+1\) 个特征值（\(u\) 的）之积。由于 \(\lambda_{i}(u)=\alpha_{i}(u)\) 对所有 \(i\in\mathrm{I}_{\mathrm{E}}\) 在 \(u\) 为正时成立（IV, p. 157 的注 5, (3)），即得所欲证的公式。

特别地，若 \(\alpha_{1}(u) < \alpha_{0}(u) = \| u\|\) ，则可看出，不等式 \(\| \widehat{\Lambda}^{n}(u)\| \leqslant \| u\|^{n}\) （EVT, V, p. 34，公式 (29)）在 \(n \geqslant 2\) 时一般不取等号。

推论. --- 设 G 是希尔伯特空间，且 \(v \in \mathcal{L}^{\mathrm{c}}(\mathrm{F};\mathrm{G})\) 。我们有

\[
\prod_ {i = 0} ^ {n} \alpha_ {i} (v \circ u) \leqslant \prod_ {i = 0} ^ {n} \alpha_ {i} (u) \alpha_ {i} (v)
\]

对所有 \(n \in \mathrm{I}_{\mathrm{E}}\) 。

只需注意

\[
\| \widehat {\Lambda} ^ {n + 1} (v u) \| = \| \widehat {\Lambda} ^ {n + 1} v \circ \widehat {\Lambda} ^ {n + 1} u \| \leqslant \| \widehat {\Lambda} ^ {n + 1} v \| \| \widehat {\Lambda} ^ {n + 1} u \|
\]

并应用命题 10。

引理 4. --- 设 A 是环。设 \(n \in \mathbf{N}\) ，并设 \((a_i)_{0 \leqslant i \leqslant n}\) 与 \((b_i)_{0 \leqslant i \leqslant n}\) 是 A 的元素族。对 \(0 \leqslant j \leqslant n\) ，置

\[
\mathrm{A} _ {j} = \sum_ {i = 0} ^ {j} a _ {i}.
\]

我们有

\[
\sum_ {i = 0} ^ {n} a _ {i} b _ {i} = \mathrm{A} _ {n} b _ {n} - \sum_ {i = 0} ^ {n - 1} \mathrm{A} _ {i} (b _ {i + 1} - b _ {i}).
\]

置 \(A_{-1}=0\) 。我们有 \(a_{i}=A_{i}-A_{i-1}\) （\(0\leqslant i\leqslant n\)），因而

\[
\sum_ {i = 0} ^ {n} a _ {i} b _ {i} = \sum_ {i = 0} ^ {n} (\mathrm{A} _ {i} - \mathrm{A} _ {i - 1}) b _ {i} = \sum_ {i = 0} ^ {n - 1} \mathrm{A} _ {i} (b _ {i} - b _ {i + 1}) + \mathrm{A} _ {n} b _ {n},
\]

即得所欲证。

设 I 是 \(\overline{R}\) 的一个不含 \(+\infty\) 的区间。回顾（FVR, I, p. 38，注），称连续函数 \(f: I \to [-\infty, +\infty[\) 是凸的，如果它在 I 的内部上的限制是取实值的凸函数；此时它在 inf I 处有（有限或无穷的）右极限。在下面的陈述中，按约定，0 与 \(\{-\infty, +\infty\}\) 中元素之积等于 0。

引理 5. --- 设 I ⊂ R 是一个不含 +∞ 的区间，f 是从 I 到 \([-\infty, +\infty[\) 中的一个递增凸函数。

设 n 是自然数，且

\[
a _ {0} \geqslant a _ {1} \geqslant \dots \geqslant a _ {n}, \quad b _ {0} \geqslant b _ {1} \geqslant \dots \geqslant b _ {n}
\]

是 I 的元素。

设 \(\varrho_{0} \geqslant \varrho_{1} \geqslant \cdots \geqslant \varrho_{n}\) 是正实数。假设

\[
\sum_ {i = 0} ^ {j} a _ {i} \leqslant \sum_ {i = 0} ^ {j} b _ {i}.\tag{5}
\]

对所有满足 \(0 \leqslant j \leqslant n\) 的整数 j 成立。则我们有

\[
\sum_ {i = 0} ^ {n} \varrho_ {i} f (a _ {i}) \leqslant \sum_ {i = 0} ^ {n} \varrho_ {i} f (b _ {i}).\tag{6}
\]

首先假设 \(a_i \in \mathring{\mathrm{I}}\) （从而 \(f(a_i) \in \mathbf{R}\) ）对所有 \(i\) 成立。设 \(i\) 满足 \(0 \leqslant i \leqslant n\) ，并设 \((\alpha_i, \beta_i)\) 是实数，使得方程为 \(y = \alpha_i x + \beta_i\) 的直线是 \(f\) 的图形在点 \((a_i, f(a_i))\) 处的支撑直线（FVR, I, p. 37）。于是有 \(f(a_i) = \alpha_i a_i + \beta_i\) 与 \(f(b_i) \geqslant \alpha_i b_i + \beta_i\) ，因为 \(f\) 的图形位于支撑直线的上方，由此得

\[
f (b _ {i}) - f (a _ {i}) \geqslant \alpha_ {i} (b _ {i} - a _ {i}).\tag{7}
\]

此外，若 \(i < n\) ，我们有

\[
\alpha_ {i} \geqslant f _ {g} ^ {\prime} (a _ {i}) \geqslant f _ {d} ^ {\prime} (a _ {i + 1}) \geqslant \alpha_ {i + 1}
\]

（同前及 FVR, I, p. 36，推论 1）；此外 \(\alpha_{i} \geqslant 0\) ，因为 \(f\) 递增（FVR, I, p. 22，推论），从而 \(\varrho_{i}\alpha_{i} \geqslant \varrho_{i+1}\alpha_{i+1} \geqslant 0\) （\(0 \leqslant i < n\)）。

设 \(j\) 是满足 \(0 \leqslant j \leqslant n\) 的整数。置

\[
\mathrm{A} _ {j} = \sum_ {i = 0} ^ {j} (b _ {i} - a _ {i}),
\]

于是由假设 (5) 知 \(A_{j} \geqslant 0\) 。依次应用不等式 (7) 与引理 4，我们导出

\[
\begin{array}{l} \sum_ {i = 0} ^ {n} \varrho_ {i} (f (b _ {i}) - f (a _ {i})) \geqslant \sum_ {i = 0} ^ {n} \varrho_ {i} \alpha_ {i} (b _ {i} - a _ {i}) \\ \qquad = \varrho_ {n} \alpha_ {n} A _ {n} + \sum_ {j = 0} ^ {n - 1} (\varrho_ {j} \alpha_ {j} - \varrho_ {j + 1} \alpha_ {j + 1}) A _ {j} \geqslant 0. \end{array}
\]

考虑一般情形，并对 \(n\) 作归纳推理。若某个 \(a_{i}\) 不属于 I 的内部，则必有 \(a_{0} = \sup I\) 或 \(a_{n} = \inf I\) 。

先设 \(a_{n}=\inf I\) 。此时有 \(a_{n}\leqslant b_{n}\) ，从而 \(\varrho_{n}f(a_{n})\leqslant\varrho_{n}f(b_{n})\) 。把归纳假设应用于族 \((a_{0},\ldots,a_{n-1})\) 、\((b_{0},\ldots,b_{n-1})\) 与 \((\varrho_{0},\ldots,\varrho_{n-1})\) ，便得

\[
\sum_ {i = 0} ^ {n - 1} \varrho_ {i} f (a _ {i}) \leqslant \sum_ {i = 0} ^ {n - 1} \varrho_ {i} f (b _ {i}),
\]

再加上 \(\varrho_{n}f(a_{n})\)即得所欲证的不等式。\(a_{0} = \sup I\) 的情形可类似处理。

命题 11（外尔不等式）. --- 设 G 是希尔伯特空间，且 \(v \in \mathcal{L}^{\mathrm{c}}(\mathrm{F};\mathrm{G})\)。设 \(g: \mathbf{R}_{+} \to [-\infty, +\infty[\) 是不减函数，使得函数 \(g \circ \exp\) 是凸的。我们有

\[
\sum_ {i = 0} ^ {n} g \left(\alpha_ {i} (v \circ u)\right) \leqslant \sum_ {i = 0} ^ {n} g \left(\alpha_ {i} (v) \alpha_ {i} (u)\right)
\]

对所有 \(n \in I_{E} \cap I_{F}\) 。

置 \(I = [-\infty, +\infty[\) 且 \(f = g \circ \exp\)。可以对下列取值应用引理 5

\[
a _ {i} = \log (\alpha_ {i} (v \circ u)) \in \mathrm{I}, \quad b _ {i} = \log (\alpha_ {i} (v) \alpha_ {i} (u)) \in \mathrm{I}
\]

并取 \(\varrho_{i} = 1\) （\(0\leqslant i\leqslant n\)），因为

\[
\sum_ {i = 0} ^ {j} a _ {i} = \log \Bigl (\prod_ {i = 0} ^ {j} \alpha_ {i} (v \circ u) \Bigr) \leqslant \log \Bigl (\prod_ {i = 0} ^ {j} \alpha_ {i} (v) \alpha_ {i} (u) \Bigr) = \sum_ {i = 0} ^ {j} b _ {i}
\]

对所有 \(0 \leqslant j \leqslant n\) 成立，依据是命题 10 的推论。于是不等式 (6) 就是所要的结论。

引理 6. --- 设 g 与 h 分别是定义在 R 的区间 I 与 J 上的凸函数。若 g 递增且定义在 h 的像上，则函数 g ∘ h 在 J 上是凸的。

事实上，对 \(t \in [0,1]\) 与 \((x,y) \in \mathrm{J} \times \mathrm{J}\) ，我们有

\[
\begin{array}{c} g (h (t x + (1 - t) y)) \leqslant g (t h (x) + (1 - t) h (y)) \\ \leqslant t g (h (x)) + (1 - t) g (h (y)). \end{array}
\]

推论. --- 设 G 是希尔伯特空间，且 \(v \in \mathcal{L}^{\mathrm{c}}(\mathrm{F};\mathrm{G})\)。设 \(n \in I_{E} \cap I_{F}\) 。

\begin{enumerate}
\def\labelenumi{\alph{enumi})}
\tightlist
\item
  设 \(r \in \mathbf{R}_{+}^{*}\) 。我们有
\end{enumerate}

\[
\sum_ {i = 0} ^ {n} \alpha_ {i} (v \circ u) ^ {r} \leqslant \sum_ {i = 0} ^ {n} \alpha_ {i} (v) ^ {r} \alpha_ {i} (u) ^ {r}.
\]

\begin{enumerate}
\def\labelenumi{\alph{enumi})}
\setcounter{enumi}{1}
\tightlist
\item
  设 \(p, q, r \in \mathbf{R}_{+}^{*}\) 满足 \(\frac{1}{p} + \frac{1}{q} = \frac{1}{r}\) 。则
\end{enumerate}

\[
\left(\sum_ {i = 0} ^ {n} \alpha_ {i} (v \circ u) ^ {r}\right) ^ {1 / r} \leqslant \left(\sum_ {i = 0} ^ {n} \alpha_ {i} (v) ^ {p}\right) ^ {1 / p} \left(\sum_ {i = 0} ^ {n} \alpha_ {i} (u) ^ {q}\right) ^ {1 / q}.
\]

\begin{enumerate}
\def\labelenumi{\alph{enumi})}
\setcounter{enumi}{2}
\tightlist
\item
  假设 F = E。对每个整数 \(m \geqslant 2\) ，我们有
\end{enumerate}

\[
\sum_ {i = 0} ^ {n} \alpha_ {i} (u ^ {m}) \leqslant \left(\sum_ {i = 0} ^ {n} \alpha_ {i} (u) ^ {2}\right) ^ {m / 2}.
\]

设 \(r \in R_{+}^{*}\) ，并设 g 是在 \(\mathbf{R}_{+}\) 上由 \(g(x) = x^{r}\) 定义的函数。函数 \(g \circ \exp\) 是凸的（引理 6），故把命题 11 应用于函数 g 即得断言 a)。

断言 \(b\)）由 \(a\)）及赫尔德不等式（INT, I，命题 4）得出。

设 \(m \geqslant 2\) 是整数。我们把 \(b)\) 应用于 \(r = 1\) 、\(p = q = 2\) 与 \(v = u^{m-1}\) ，得到

\[
\sum_ {i = 0} ^ {n} \alpha_ {i} (u ^ {m}) \leqslant \left(\sum_ {i = 0} ^ {n} \alpha_ {i} (u ^ {m - 1}) ^ {2}\right) ^ {1 / 2} \left(\sum_ {i = 0} ^ {n} \alpha_ {i} (u) ^ {2}\right) ^ {1 / 2}.
\]

我们对 \(m \geqslant 2\) 作归纳来证明 c)。前一不等式给出了 m = 2 时的断言。设 \(m \geqslant 3\) 且断言对 m - 1 成立；由于我们有

\[
\sum_ {i = 0} ^ {n} \alpha_ {i} (u ^ {m - 1}) ^ {2} \leqslant \left(\sum_ {i = 0} ^ {m} \alpha_ {i} (u ^ {m - 1})\right) ^ {2},
\]

上面的不等式蕴含

\[
\sum_ {i = 0} ^ {n} \alpha_ {i} (u ^ {m}) \leqslant \left(\sum_ {i = 0} ^ {n} \alpha_ {i} (u ^ {m - 1})\right) ^ {1 / 2} \left(\sum_ {i = 0} ^ {n} \alpha_ {i} (u) ^ {2}\right) ^ {1 / 2} \leqslant \left(\sum_ {i = 0} ^ {n} \alpha_ {i} (u) ^ {2}\right) ^ {m / 2}
\]

其中应用了归纳假设。

\subsection{有限迹自同态}

我们回顾，E 的正自同态 \(u\) 是有限迹的，当且仅当存在 E 的标准正交基 \((e_i)_{i \in I}\) ，使得

\[
\sum_ {i \in \mathrm{I}} \langle e _ {i} | u (e _ {i}) \rangle <   + \infty
\]

（EVT, V, p. 48，引理 3 及 p. 49，定义 7）。若 K = C，E 的有限迹自同态所成的空间 \(\mathcal{L}_1(\mathrm{E})\) 是由正的有限迹自同态组成的集生成的向量空间（EVT, V, p. 50，定义 8）；若 K = R，空间 \(\mathcal{L}_1(\mathrm{E})\) 定义为交 \(\mathcal{L}(\mathrm{E}) \cap \mathcal{L}_1(\mathrm{E}_{(\mathbf{C})})\) （EVT, V, p. 50）。

若 \(u\in \mathcal{L}_1(\mathrm{E})\) ，则级数

\[
\sum_ {i \in \mathrm{I}} \langle e _ {i} | u (e _ {i}) \rangle
\]

对 E 的每个标准正交基 \((e_{i})_{i\in\mathrm{I}}\) 收敛，且其和与标准正交基无关；这个和称为 u 的迹 \(\operatorname{Tr}(u)\) （EVT, V, p. 50）。若 K = R，我们有 \(\operatorname{Tr}(u) = \operatorname{Tr}(u_{(\mathbf{C})})\) 。

设 \(u \in \mathcal{L}_1(\mathrm{E})\) 。我们有 \(u^* \in \mathcal{L}_1(\mathrm{E})\) ，且 \(\operatorname{Tr}(u^*) = \overline{\operatorname{Tr}(u)}\) （同前）。

命题 12. --- 设 \(B = (e_i)_{i \in I}\) 是 E 的标准正交基。设 \(u \in \mathcal{D}_B(E)\) ，并用 \(\lambda = (\lambda_i)_{i \in I}\) 表示它的特征值族。自同态 \(u\) 是有限迹的，当且仅当族 \(\lambda\) 在 K 中可和。此时我们有 \(\operatorname{Tr}(u) = \sum \lambda_i\) 。

如有必要，把 \(u\) 换成 \(u_{(\mathbf{C})}\) ，我们可设 \(K = \mathbf{C}\) 。

先设 \(u\) 是有限迹的。根据 EVT, V, p. 50 及 p. 49 的公式 (25)，族 \((\langle e_i | u(e_i) \rangle)_{i \in I} = (\lambda_i)_{i \in I}\) 是可和的。

反之，设族 \(\lambda\) 是可和的。

于是族 \((\mathcal{R}(\lambda_{i})^{+})\) 、\((\mathcal{R}(\lambda_{i})^{-})\) 、\((\mathcal{I}(\lambda_{i})^{+})\) 、\((\mathcal{I}(\lambda_{i})^{-})\) 都是可和的。自同态 u 是 \(\mathcal{D}_{\mathrm{B}}(\mathrm{E})\) 中以这些族为特征值的元素的线性组合。\(\mathcal{D}_{\mathrm{B}}(\mathrm{E})\) 的这些元素都是正的。由于按定义，有限迹自同态空间由正的有限迹自同态生成，因此可设对所有 i 有 \(\lambda_{i} \geqslant 0\) 。由于 \(\langle e_{i} | u(e_{i}) \rangle = \lambda_{i}\)，族 \((\langle e_{i} | u(e_{i}) \rangle)_{i \in \mathrm{I}}\) 是可和的，故 u 是有限迹的（EVT, V, p. 48，引理 3）。

最后，若 u 有有限迹，则由 EVT, V, p. 50，我们有

\[
\operatorname{Tr} (u) = \sum_ {i \in \mathrm{I}} \left\langle e _ {i} \mid u \left(e _ {i}\right) \right\rangle = \sum_ {i \in \mathrm{I}} \lambda_ {i}.
\]

推论 1. --- 设 u 是 E 的有限迹自同态。

\begin{enumerate}
\def\labelenumi{\alph{enumi})}
\item
  自同态 u 是紧的；
\item
  设 \(\mathrm{B} = (e_i)_{i \in \mathrm{I}}\) 是 u 的始空间 \(\operatorname{Ker}(u)^{\circ}\) 的标准正交基，\(\mathrm{C} = (f_i)_{i \in \mathrm{I}}\) 是 E 的标准正交族，且 \((\alpha_i)_{i \in \mathrm{I}}\) 是 \((\mathbf{R}_+^*)^{\mathrm{I}}\) 中的族，使得 \(u(e_i) = \alpha_i f_i\) 对所有 \(i \in I\) 成立（IV, p. 150 的推论 2）。我们有
\end{enumerate}

\[
\operatorname{Tr} (u) = \sum_ {i \in I} \alpha_ {i} \langle e _ {i} | f _ {i} \rangle .
\]

为证明 \(a\)），可设 \(\mathbf{K} = \mathbf{C}\) （EVT, V, p. 50 及 III, p. 2 的注 4），且 \(u\) 是正的（EVT, V, p. 50，定义 8）。

由 EVT, V, p. 56 的推论 1，存在 E 的标准正交基 \(B = (e_i)_{i \in I}\) ，使得 \(u\) 在基 B 中是对角的，且特征值族 \(\lambda\) （\(u\) 的特征值族）属于 \(\mathbf{R}_+^I\) 并可和。因此族 \(\lambda\) 属于 \(\mathcal{C}_0(\mathrm{I};\mathrm{K})\) （TG, III, p. 38，命题 1），从而自同态 \(u\) 是紧的（IV, p. 148 的命题 2）。

我们来证明 \(b\)）。设 \((e_i)_{i \in \mathbf{J}}\) 是 E 的标准正交基，它扩张族 B，使得 \(u(e_i) = 0\) （当 \(i \in \mathbf{J} - \mathbf{I}\) 时）。我们得到

\[
\operatorname{Tr} (u) = \sum_ {i \in \mathrm{J}} \langle e _ {i} \mid u (e _ {i}) \rangle = \sum_ {i \in \mathrm{I}} \langle e _ {i} \mid u (e _ {i}) \rangle = \sum_ {i \in \mathrm{I}} \alpha_ {i} \langle e _ {i} \mid f _ {i} \rangle .
\]

推论 2. --- 设 F 是希尔伯特空间。设 \(u \in \mathcal{L}(\mathrm{E};\mathrm{F})\) 是希尔伯特--施密特算子。则 u 是紧线性算子。

设 \((j,|u|)\) 是 u 的极分解（I, p. 140，定义 4）。按定义（EVT, V, p. 50，定义 9），自同态 \(u^{*}u\) 有有限迹，因而是紧的（推论 1, a)）；对 \(|u|=\sqrt{u^{*}u}\) （III, p. 91 的命题 6, b)）与 \(u=j|u|\) （III, p. 5 的命题 3）结论同样成立。

推论 3. --- 假设 K = C。设 u 是 E 的正紧自同态。自同态 u 是有限迹的，当且仅当它的特征值的递减族 \((\lambda_{n}(u))_{n\in\mathrm{I}_{\mathrm{E}}}\) 可和；此时 u 的迹就是这个族的和。

由 IV, p. 149 的定理 1，这可由前一命题以及序列 \((\lambda_n(u))_{n\in \mathrm{I_E}}\) 的定义（IV, p. 151 的 no. 3）得出，并用到 EVT, V, p. 49 的公式 (28)。

推论 4. --- 设 F 是希尔伯特空间。设 \(u \in \mathcal{L}(\mathrm{E};\mathrm{F})\) 是紧线性映射。设 \(\mathrm{B} = (e_i)_{i \in \mathrm{I}}\) 是始空间 \(\operatorname{Ker}(u)^{\circ}\) （\(u\) 的始空间）的标准正交基，\(\mathrm{C} = (f_i)_{i \in \mathrm{I}}\) 是 F 的标准正交族，且 \((\alpha_{i})_{i\in\mathrm{I}}\) 是 \((\mathbf{R}_{+}^{*})^{\mathrm{I}}\) 中的族，使得 \(u(e_{i})=\alpha_{i}f_{i}\) 对所有 \(i\in I\) 成立（IV, p. 150 的推论 2）。

E 的自同态 \(|u|\) 有有限迹，当且仅当族 \((\alpha_{i})_{i\in \mathrm{I}}\) 可和。此时我们有 \(u\in \mathcal{L}_2(\mathrm{E};\mathrm{F})\) ，并且

\[
\operatorname{Tr} (| u |) = \sum_ {i \in \mathrm{I}} \alpha_ {i}, \quad \| u \| _ {2} ^ {2} = \operatorname{Tr} (u ^ {*} u) = \sum_ {i \in \mathrm{I}} \alpha_ {i} ^ {2},\tag{8}
\]

特别地有 \(\|u\|_{2}\leqslant\mathrm{Tr}(|u|)\) 。

\(|u|\) 的非零特征值族是 \((\alpha_{i})_{i\in \mathrm{I}}\) ，故 \(|u|\) 有有限迹当且仅当族 \((\alpha_{i})_{i\in \mathrm{I}}\) 可和（命题 12）。若确实如此，非零特征值族 \((\alpha_{i}^{2})\) （\(u^{*}u = |u|^{2}\) 的非零特征值族）可和，且公式 (8) 由同前得出。

引理 7. --- 设 \(\lambda = (\lambda_i)_{i \in \mathrm{I}}\) 是复数族。对每个 \(t \in \mathbf{C}^*\) ，用 \(n_t\) 表示 \(i \in \mathrm{I}\) 中使 \(\lambda_i = t\) 者所成之集的基数。族 \((\lambda_i)_{i \in \mathrm{I}}\) 可和，当且仅当 \(n_t\) 对所有 \(t \in \mathbf{C}^*\) 有限，且族 \((n_t t)_{t \in \mathbf{C}^*}\) 可和。此时这两个族的和相等。

设族 \((\lambda_{i})_{i\in\mathrm{I}}\) 可和。对每个 \(t\in\mathbf{C}\) ，用 \(I_{t}\) 表示 \(i\in I\) 中使 \(\lambda_{i}=t\) 者所成之集。把 TG, III, p. 39 的定理 2 应用于 I 被诸集 \(I_{t}\) 所作的分划，可知集 \(I_{t}\) 对每个 \(t\in\mathbf{C}^{*}\) 有限，且族 \((n_{t}t)_{t\in\mathbf{C}^{*}}\) 可和，其和等于族 \((\lambda_{i})_{i\in\mathrm{I}}\) 的和。

反之，设 \(n_{t}\) 对所有 \(t \in \mathbf{C}^{*}\) 有限，且族 \((n_{t}t)_{t \in \mathbf{C}^{*}}\) 可和。设 J 是 I 的有限子集，\(\Lambda\) 是诸 \(\lambda_{i}\) （\(i \in J\)）所成之集。我们有

\[
\left| \sum_ {i \in J} \lambda_ {i} \right| \leqslant \sum_ {t \in \Lambda - \{0 \}} n _ {t} | t | \leqslant \sum_ {t \in \mathbf {C} ^ {*}} n _ {t} | t |,
\]

因此族 \((\lambda_{i})_{i\in\mathrm{I}}\) 可和（TG, VII, p. 17，推论）。

命题 13. --- 假设 K = C。设 u 是 E 的紧正规自同态。对 \(t \in \mathrm{Sp}_{\mathrm{s}}(u)\) ，用 \(n_{t} \geqslant 1\) 表示 u 的特征值 t 的谱重数。u 有有限迹的必要充分条件是族 \((n_{t}t)_{t \in \mathrm{Sp}_{\mathrm{s}}(u)}\) 可和。此时 u 的迹就是这个族的和。

设 \(\mathrm{B} = (e_i)_{i \in \mathrm{I}}\) 是 E 的标准正交基，使得 u 在基 B 中是对角的（IV, p. 149 的定理 1），并设 \(\lambda = (\lambda_i)_{i \in \mathrm{I}}\) 是它的特征值族。由于对 \(t \in \mathrm{Sp}_\mathrm{s}(u)\) ，重数 \(n_t\) 等于元素 \(i \in I\) 中使 \(\lambda_i = t\) 者的个数，且 \(\lambda\) 的非零元素属于 \(\mathrm{Sp}_{\mathrm{s}}(u)\) ，断言即由命题 12 与前一引理得出。

\subsection{核型映射}

在本节中，F 表示 K 上的希尔伯特空间。当 K = C 时，我们回顾（I, p. 139，定义 3），对 \(u \in \mathcal{L}(\mathrm{E}; \mathrm{F})\) ，用 \(|u|\) 表示 E 的正自同态 \(\sqrt{u^{*} \circ u}\) 。当 K = R 时，元素 \(|u_{(\mathbf{C})}|\) （\(\mathcal{L}(\mathrm{E}_{(\mathbf{C})})\) 中的元素）具有形式 \(v_{(\mathbf{C})}\) ，其中自同态 \(v \in \mathcal{L}(\mathrm{E})\) 是唯一的且仍记作 \(|u|\) （参见 I, p. 87）。

我们用 \((u,v)\mapsto\langle u|v\rangle=\mathrm{Tr}(u^{*}v)\) 表示希尔伯特空间 \(\mathcal{L}_{2}(\mathrm{E};\mathrm{F})\) 中的内积（EVT, V, p. 53，注 1 与注 2）。

对每个 \(u \in \mathcal{L}(\mathrm{E};\mathrm{F})\) ，用 \(\| u\| _1\) 表示量 \(\mathrm{Tr}(|u|)\) （若 \(|u|\) 有有限迹）；否则置 \(\| u\| _1 = +\infty\) 。于是我们有 \(\| u\| _1 \in \mathbf{R}_+ \cup \{+\infty\}\) 。由于 \(\| u\| = \| |u|\|\) （I, p. 139 的命题 10, a)），且当 \(v\) 是正的并有限迹时有 \(\mathrm{Tr}(v) \geqslant \| v\|\) （EVT, V, p. 49，公式 (24bis) 及 p. 44，命题 9），由此可得

\[
\| u \| \leqslant \| u \| _ {1}.\tag{9}
\]

若 \(\| u\| _1\) 有限，则 \(u\) 是希尔伯特--施密特算子，且 \(\| u\| _2 \leqslant \| u\| _1\) （IV, p. 165 的推论 4）。

命题 14. --- 设 \(u \in \mathcal{L}_2(\mathrm{E};\mathrm{F})\) 。我们有

\[
\| u\|_{1} = \sup_{\substack{v\in \mathcal{L}_{2}(\mathrm{E};\mathrm{F})\\ \| v\| \leqslant 1}}|\langle v  |  u\rangle |,
\]

其中的上确界在 \(\mathbf{R}_{+} \cup \{+\infty\}\) 中计算。

设 \(\mathrm{B} = (e_i)_{i \in \mathrm{I}}\) 是 u 的始空间 \(\operatorname{Ker}(u)^{\circ}\) 的标准正交基，\(\mathrm{C} = (f_i)_{i \in \mathrm{I}}\) 是 F 中的标准正交族，\((\alpha_i)_{i \in \mathrm{I}}\) 是 \((\mathbf{R}_+^*)^{\mathrm{I}}\) 中的族，使得 \(u(e_i) = \alpha_i f_i\) 对所有 \(i \in I\) 成立（IV, p. 150 的推论 2）。族 \((\alpha_i)_{i \in \mathrm{I}}\) 平方可积，因为 u 属于 \(\mathcal{L}_2(\mathrm{E};\mathrm{F})\) （IV, p. 165 的推论 4）。设 \((e_j)_{j \in \mathrm{J}}\) 是 \((e_i)_{i \in \mathrm{I}}\) 的扩张，即 E 的标准正交基。

设 L 是 I 的有限子集。设 \(v_{L}\) 是由

\[
v _ {\mathrm{L}} (x) = \sum_ {i \in \mathrm{L}} \langle e _ {i} | x \rangle f _ {i}
\]

对所有 \(x\) 属于 E 定义的从 E 到 F 中的连续有限秩线性映射。我们有 \(\| v_{\mathrm{L}}\| \leqslant 1\) 且 \(v_{\mathrm{L}}\in \mathcal{L}_{2}(\mathrm{E};\mathrm{F})\) 。此外，对每个 \(j\in \mathrm{J}\) ，有 \(v_{\mathrm{L}}(e_j) = 0\) （若 \(j\notin \mathrm{L}\) ），且有 \(v_{\mathrm{L}}(e_j) = f_j\) （若 \(j\in \mathrm{L}\) ）。于是

\[
| \langle v _ {\mathrm{L}} | u \rangle | = | \operatorname{Tr} (v _ {\mathrm{L}} ^ {*} u) | = \left| \sum_ {j \in \mathrm{J}} \langle v _ {\mathrm{L}} (e _ {j}) | u (e _ {j}) \rangle \right| = \sum_ {j \in \mathrm{L}} \alpha_ {j}.
\]

由同前，我们推出

\[
\| u\|_{1} = \sup_{\text{L}}\sum_{j\in \text{L}}\alpha_{j}\leqslant \sup_{\substack{v\in \mathcal{L}_{2}(\text{E};\text{F})\\ \| v\| \leqslant 1}}|\langle v\mid u\rangle |\tag{10}
\]

在 \(\mathbf{R}_{+}\cup\{+\infty\}\) 中成立。

当 \(\|u\|_{1}=+\infty\) 时，这即给出命题中的等式。

设 \(\|u\|_{1}\) 有限。对每个 \(i \in I\) ，设 \(p_{i}\) 是从 E 到 F 的线性映射，满足 \(p_{i}(e_{i}) = f_{i}\) 且 \(p_{i}(x) = 0\) 对所有 \(x \in e_{i}^{\circ}\) 成立。对 I 中所有的 j 和 k，我们有

\[
\langle p _ {j} \mid p _ {k} \rangle = \mathrm{Tr} (p _ {j} ^ {*} p _ {k}) = \sum_ {i \in \mathrm{J}} \langle p _ {j} (e _ {i}) \mid p _ {k} (e _ {i}) \rangle ,
\]

除非 \(j = k\) ，它为零；在此情形这个量等于 \(\| f_k\|^2 = 1\) 。因此族 \((p_i)_{i \in \mathrm{I}}\) 在 \(\mathcal{L}_2(\mathrm{E};\mathrm{F})\) 中标准正交。从而族 \((\alpha_i p_i)_{i \in \mathrm{I}}\) 在 \(\mathcal{L}_2(\mathrm{E};\mathrm{F})\) 中可和；其和等于 \(u\) ，因为这两个连续线性映射在元素 \(e_i\) 上（对所有 \(i \in \mathrm{J}\)）相合。

设 \(v \in \mathcal{L}_2(\mathrm{E};\mathrm{F})\) 。对每个 \(i \in \mathrm{I}\) ，由此可得

\[
\langle v \mid p _ {i} \rangle = \operatorname{Tr} (v ^ {*} p _ {i}) = \sum_ {j \in \mathrm{J}} \langle v (e _ {j}) \mid p _ {i} (e _ {j}) \rangle = \langle v (e _ {i}) \mid f _ {i} \rangle .
\]

若 \(\|v\| \leqslant 1\) ，我们有

\[
\begin{array}{c} | \langle v | u \rangle | = \left| \langle v | \sum_ {i \in I} \alpha_ {i} p _ {i} \rangle \right| \leqslant \sum_ {i \in I} \alpha_ {i} | \langle v | p _ {i} \rangle | \\ = \sum_ {i \in I} \alpha_ {i} | \langle v (e _ {i}) | f _ {i} \rangle | \leqslant \sum_ {i \in I} \alpha_ {i} = \mathrm{Tr} (| u |), \end{array}
\]

由此得到不等式

\[
\sup_{\substack{v\in \mathcal{L}_{2}(\mathrm{E};\mathrm{F})\\ \| v\| \leqslant 1}}|\langle v\mid u\rangle |\leqslant \| u\|_{1},
\]

它与公式 (10) 相结合，便完成了命题的证明。

由这个命题可得，集 \(\mathcal{L}_1(\mathrm{E};\mathrm{F})\) （从 E 到 F 的连续线性映射 \(u\) 中使 \(\| u\| _1\) 有限者所成之集）是 \(\mathcal{L}_{2}(\mathrm{E};\mathrm{F})\) 的向量子空间，且映射 \(u\mapsto \| u\| _1\) 是 \(\mathcal{L}_1(\mathrm{E};\mathrm{F})\) 上的半范数；不等式 (9) 表明它是一个范数。

定义 4. --- 我们称向量空间 \(\mathcal{L}_{1}(\mathrm{E};\mathrm{F})\) （赋以范数 \(u\mapsto \| u\|_{1}\)）为从 E 到 F 的核型映射所成的空间。若 \(u\in \mathcal{L}_1(\mathrm{E};\mathrm{F})\) ，则称 \(u\) 是核型的。

注记. --- 1) 假设 K = R。设 \(u \in \mathcal{L}(E; F)\) 。我们有 \(u \in \mathcal{L}_{1}(E; F)\) 当且仅当 \(u_{(\mathbf{C})} \in \mathcal{L}_{1}(E_{(\mathbf{C})}; F_{(\mathbf{C})})\) ；此时有 \(\|u\|_{1} = \|u_{(\mathbf{C})}\|_{1}\) 。

\begin{enumerate}
\def\labelenumi{\arabic{enumi})}
\setcounter{enumi}{1}
\tightlist
\item
  设 \(u\) 是从 E 到 F 的核型映射。由于映射 \(u\) 是希尔伯特--施密特的，它是紧的（IV, p. 165 的推论 2）。此外，命题 14 蕴含对每个 \(v \in \mathcal{L}_2(\mathrm{E};\mathrm{F})\) ，我们有
\end{enumerate}

\[
| \langle v | u \rangle | \leqslant \| v \| \| u \| _ {1}.\tag{11}
\]

\begin{enumerate}
\def\labelenumi{\arabic{enumi})}
\setcounter{enumi}{2}
\item
  假设 \(\mathrm{K} = \mathbf{C}\)。设 \(u \in \mathcal{L}_1(\mathrm{E};\mathrm{E})\) 且 \(u\) 是正的。范数 \(\| u\| _1\) （\(u\) 的范数）是序列 \((\lambda_n(u))_{n\in \mathrm{I_E}}\) 的和（IV, p. 165 的推论 3）。
\item
  \(\mathcal{L}_1(\mathrm{E};\mathrm{F})\) 到 \(\mathcal{L}(\mathrm{E};\mathrm{F})\) 中的典范单射的范数 \(\leqslant 1\) （p. 167 的不等式 (9)）。
\end{enumerate}

命题 15. --- 设 G 是 K 上的希尔伯特空间。映射 \((u,v)\mapsto v\circ u\) （从 \(\mathcal{L}(\mathrm{E};\mathrm{F})\times\mathcal{L}(\mathrm{F};\mathrm{G})\) 到 \(\mathcal{L}(\mathrm{E};\mathrm{G})\) 中的映射）过渡到商，定义了范数 \(\leqslant1\) 的从 \(\mathcal{L}_{2}(\mathrm{E};\mathrm{F})\times\mathcal{L}_{2}(\mathrm{F};\mathrm{G})\) 到 \(\mathcal{L}_{1}(\mathrm{E};\mathrm{G})\) 中的连续双线性映射。

设 \(u \in \mathcal{L}_2(\mathrm{E};\mathrm{F})\) 且 \(v \in \mathcal{L}_2(\mathrm{F};\mathrm{G})\)。设 \(w \in \mathcal{L}_2(\mathrm{E};\mathrm{G})\) 。我们得到 \(\langle uv | w \rangle = \operatorname{Tr}(v^* u^* w) = \langle v | u^* w \rangle\) ，由此

\[
| \langle u v | w \rangle | \leqslant \| v \| _ {2} \| u ^ {*} w \| _ {2} \leqslant \| v \| _ {2} \| u \| _ {2} \| w \|
\]

此处用到柯西--施瓦茨不等式及 EVT, V, p. 52 的公式 (37)。因此结论由命题 14 得出。

引理 8. --- 映射 \(u \mapsto u^{*}\) （从 \(\mathcal{L}(\mathrm{E};\mathrm{F})\) 到 \(\mathcal{L}(\mathrm{F};\mathrm{E})\) 中的映射）过渡到商，定义了从 \(\mathcal{L}_{1}(\mathrm{E};\mathrm{F})\) 到 \(\mathcal{L}_{1}(\mathrm{F};\mathrm{E})\) 中的等距线性映射。

设 \(u \in \mathcal{L}_1(\mathrm{E};\mathrm{F})\) 。由于 \(u\) 是希尔伯特--施密特映射，\(u^*\) 也是（EVT, V, p. 54）。映射 \(v \mapsto v^*\) 是从 \(v \in \mathcal{L}_2(\mathrm{E};\mathrm{F})\) 中满足 \(\|v\| \leqslant 1\) 者所成之集到 \(w \in \mathcal{L}_2(\mathrm{F};\mathrm{E})\) 中满足 \(\|w\| \leqslant 1\) 者所成之集上的双射；对每个 \(v \in \mathcal{L}_2(\mathrm{E};\mathrm{F})\)

满足 \(\|v\|\leqslant1\) 者，我们有 \(\langle v|u\rangle=\langle u^{*}|v^{*}\rangle\) （EVT, V, p. 54，公式 (42)），于是结论由命题 14 得出。

命题 16. --- 设 \(E_{1}\) 与 \(F_{1}\) 是希尔伯特空间。设 \(u\) 属于 \(\mathcal{L}_{1}(\mathrm{E};\mathrm{F})\) ，\(v\) 属于 \(\mathcal{L}(\mathrm{E}_{1};\mathrm{E})\) ，\(v_{2}\) 属于 \(\mathcal{L}(\mathrm{F};\mathrm{F}_{1})\)。则有 \(v_{2}uv_{1} \in \mathcal{L}_{1}(\mathrm{E}_{1};\mathrm{F}_{1})\) 且 \(\|v_{2}uv_{1}\|_{1} \leqslant \|v_{2}\|\|v_{1}\|\|u\|_{1}\) 。

设 \(w \in \mathcal{L}_2(\mathrm{E};\mathrm{F}_1)\) 满足 \(\| w\| \leqslant 1\) 。我们有 \(v_{2}u \in \mathcal{L}_{2}(\mathrm{E};\mathrm{F}_{1})\) （EVT, V, p. 52，公式 (36)）。由于 \(v_{2}^{*}w \in \mathcal{L}_{2}(\mathrm{E};\mathrm{F})\) （同前），由此可得

\[
| \langle w | v _ {2} u \rangle | = | \langle v _ {2} ^ {*} w | u \rangle | \leqslant \| v _ {2} \| \| u \| _ {1}
\]

（公式 (11)），于是 \(v_2u \in \mathcal{L}_1(\mathrm{E};\mathrm{F}_1)\) 且 \(\| v_2u\| _1 \leqslant \| v_2\| \| u\| _1\) （命题 14）。设 \(v_{1} \in \mathcal{L}(\mathrm{E}_{1};\mathrm{E})\) ；由于 \(uv_{1} = (v_{1}^{*}u^{*})^{*}\) ，我们有 \(uv_{1} \in \mathcal{L}_{1}(\mathrm{E}_{1};\mathrm{F})\) 且 \(\| uv_{1}\|_{1} \leqslant \| v_{1}^{*}\| \| u^{*}\|_{1} = \| v_{1}\| \| u\|_{1}\) （引理 8）。

命题由这些不等式立即得出。

命题 17. — 空间 \(\mathcal{L}_{1}(\mathrm{E};\mathrm{E})\) 与 E 的有限迹自同态所成的空间 \(\mathcal{L}_{1}(\mathrm{E})\) 相重合，且我们有 \(|\operatorname{Tr}(u)|\leqslant \| u\|_{1}\) ，对 E 的每个有限迹自同态 \(u\) 成立。

我们可以假设 K = C。按定义，空间 \(\mathcal{L}_{1}(E;E)\) 包含有限迹正自同态所成之集，因此 \(\mathcal{L}_{1}(E)\subset\mathcal{L}_{1}(E;E)\) （EVT, p. 50，定义 8）。反过来，我们要证明 \(\mathcal{L}_{1}(E;E)\) 含于 \(\mathcal{L}_{1}(E)\) 。

设 \(u \in \mathcal{L}_{1}(\mathrm{E};\mathrm{E})\) 。我们有 \(u^{*} \in \mathcal{L}_{1}(\mathrm{E};\mathrm{E})\) （引理 8）；因此只需证明 \(\mathcal{L}_{1}(\mathrm{E};\mathrm{E})\) 中的埃尔米特元素有有限迹（I, p. 96 的引理 2）。

设 u 是这样一个自同态。它是紧的（p. 169 的注 2）。设 B 是 E 的标准正交基，使得 u 在基 B 中是对角的（IV, p. 149 的定理 1），并设 \(\lambda\) 是 u 关于 B 的特征值族。自同态 \(|u|\) 在基 B 中可对角化，其特征值族是 \(|\lambda|\) （IV, p. 147 的命题 1, d)）。由于 \(|u|\) 有有限迹，后一族可和（IV, p. 164 的命题 12），因此族 \(\lambda\) 可和。于是 u 有有限迹（同前），且我们有 \(|\mathrm{Tr}(u)| \leqslant \mathrm{Tr}(|u|) = \|u\|_{1}\) 。

从而，空间 \(\mathcal{L}_{1}(\mathrm{E})\) 是 C*-代数 \(\mathcal{L}(\mathrm{E})\) 的自伴双边理想。若 E 是无限维的，这个理想在 \(\mathcal{L}(\mathrm{E})\) 中不闭。

命题 17 与命题 16 表明，映射 \(b\) ——从 \(\mathcal{L}_1(\mathrm{E};\mathrm{F})\times \mathcal{L}(\mathrm{F};\mathrm{E})\) 到 \(\mathbf{C}\) 中的映射——由 \(b(u,v) = \operatorname {Tr}(vu)\) 定义，是连续双线性型，它使空间 \(\mathcal{L}_1(\mathrm{E};\mathrm{F})\) 与 \(\mathcal{L}(\mathrm{F};\mathrm{E})\) 处于对偶。

引理 9. --- 设 \(u \in \mathcal{L}_1(\mathrm{E};\mathrm{F})\) 。我们有

\[
\| u\|_{1} = \sup_{\substack{v\in \mathcal{L}^{c}(\mathrm{F};\mathrm{E})\\ \| v\| \leqslant 1}}|b(u,v)|.
\]

由于每个希尔伯特--施密特映射都是紧的（IV, p. 165 的推论 2），我们有

\[
\| u\|_{1}\leqslant \sup_{\substack{v\in \mathcal{L}^{\mathrm{c}}(\mathrm{E};\mathrm{F})\\ \| v\| \leqslant 1}}|\operatorname{Tr}(vu)| = \sup_{\substack{v\in \mathcal{L}^{\mathrm{c}}(\mathrm{F};\mathrm{E})\\ \| v\| \leqslant 1}}|b(u,v)|
\]

（命题 14）。

设 \(v \in \mathcal{L}^{c}(\mathrm{E};\mathrm{F})\) 的范数 \(\leqslant 1\) 。由于 \(\mathcal{L}_{2}(\mathrm{E};\mathrm{F})\) 在 \(\mathcal{L}^{c}(\mathrm{E};\mathrm{F})\) 中稠密，存在序列 \((v_{n})_{n\in\mathbf{N}}\) （\(\mathcal{L}_{2}(\mathrm{E};\mathrm{F})\) 中的序列）在 \(\mathcal{L}(\mathrm{E};\mathrm{F})\) 中收敛到 v。序列 \((b(u,v_{n}))_{n\in\mathbf{N}}\) 收敛到 \(b(u,v)\) ；由于 \(|b(u,v_{n})| = |\operatorname{Tr}(v_{n}u)| = |\langle v_{n}^{*}|u\rangle| \leqslant \|u\|_{1}\) （命题 14）对所有 \(n \in N\) 成立，由此可得 \(|b(u,v)| \leqslant \|u\|_{1}\) 。

用 \(\theta\) 表示连续线性映射

\[
\theta \colon \mathcal {L} _ {1} (\mathrm{E}; \mathrm{F}) \to \mathcal {L} ^ {\mathrm{c}} (\mathrm{F}; \mathrm{E}) ^ {\prime}
\]

使得 \(\theta(u)(v) = b(u,v)\) 。

命题 18. --- 映射 θ 是等距同构。

由引理 9，映射 \(\theta\) 是等距的，只需证明 \(\theta\) 是满射。

设 \(\lambda \in \mathcal{L}^{\mathrm{c}}(\mathrm{F};\mathrm{E})'\) 。由于 \(\| v\| \leqslant \| v\|_{2}\) 对所有 \(v\in \mathcal{L}_2(\mathrm{F};\mathrm{E})\) 成立（EVT, V, p. 52，公式 (33)），\(\lambda\) 在子空间 \(\mathcal{L}_2(\mathrm{F};\mathrm{E})\) 上的限制是希尔伯特空间 \(\mathcal{L}_2(\mathrm{F};\mathrm{E})\) 上的连续线性型。因此存在元素 \(u\) （\(\mathcal{L}_2(\mathrm{F};\mathrm{E})\) 的元素），使得 \(\lambda (v) = \langle u|v\rangle\) 对所有 \(v\in \mathcal{L}_2(\mathrm{F};\mathrm{E})\) 成立（EVT, V, p. 15，定理 3）。

对每个 \(w \in \mathcal{L}_2(\mathrm{E};\mathrm{F})\) （范数 \(\leqslant 1\) 者），我们有 \(|\langle w|u\rangle| = |\lambda(w)| \leqslant \| \lambda \|\) 。于是 \(u\) 是从 \(\mathrm{E}\) 到 \(\mathrm{F}\) 中的核型映射（命题 14）。

连续线性型 \(\lambda\) 与 \(\theta(u)\) 在稠密子空间 \(\mathcal{L}_{2}(\mathrm{F},\mathrm{E})\) （\(\mathcal{L}^{\mathrm{c}}(\mathrm{F};\mathrm{E})\) 的稠密子空间）上相等。由此可得 \(\lambda=\theta(u)\) ，证明完成。

推论. --- 赋范空间 \(\mathcal{L}_{1}(\mathrm{E};\mathrm{F})\) 是巴拿赫空间。

由于空间 \(\mathcal{L}^{\mathrm{c}}(\mathrm{F};\mathrm{E})\) 是赋范空间，断言由本命题及 EVT, III, p. 24 的推论 2 得出。

\subsection{希尔伯特--施密特积分算子}

在本节中，X 与 Y 是局部紧拓扑空间。设 \(\mu\) 是 X 上的正测度，\(\nu\) 是 Y 上的正测度。我们分别用 \(\mathrm{L}^{2}(\mathrm{X})\) 、\(\mathrm{L}^{2}(\mathrm{Y})\) 与 \(\mathrm{L}^{2}(\mathrm{X} \times \mathrm{Y})\) 表示空间 \(\mathrm{L}^{2}(\mathrm{X}, \mu)\) 、\(\mathrm{L}^{2}(\mathrm{Y}, \nu)\) 与 \(\mathrm{L}^{2}(\mathrm{X} \times \mathrm{Y}, \mu \otimes \nu)\) ，对 \(\mathcal{L}^{2}(\mathrm{X})\) 、\(\mathcal{L}^{2}(\mathrm{Y})\) 与 \(\mathcal{L}^{2}(\mathrm{X} \times \mathrm{Y})\) 亦同样处理。

我们回顾（参见 III, p. 26 的 no. 4），对每个 \(N \in \mathcal{L}^{2}(X \times Y)\) ，存在唯一的连续线性映射 \(u_{\mathrm{N}}\) （从 \(L^{2}(X)\) 到 \(L^{2}(Y)\) 中的连续线性映射），使得

\[
\langle g \mid u _ {\mathrm{N}} (f) \rangle = \int_ {\mathrm{X} \times \mathrm{Y}} \overline {{g (y)}} \mathrm{N} (x, y) f (x) d (\mu \otimes \nu) (x, y)\tag{12}
\]

对所有 \(f \in \mathrm{L}^{2}(\mathrm{X})\) 与每个 \(g \in \mathrm{L}^{2}(\mathrm{Y})\) 成立。映射 \(u_{N}\) 是紧的（III, p. 33 的推论 1）。映射 \(N \mapsto u_{N}\) 过渡到商，诱导出从 \(\mathrm{L}^{2}(\mathrm{X} \times \mathrm{Y})\) 到 \(\mathcal{L}(\mathrm{L}^{2}(\mathrm{X});\mathrm{L}^{2}(\mathrm{Y}))\) 中的连续单射线性映射（III, p. 30 的命题 5）。

我们用 \(\theta\) 表示从 \(\mathcal{L}^2 (\mathrm{X})\otimes \mathcal{L}^2 (\mathrm{Y})\) 到 \(\mathcal{L}^2 (\mathrm{X}\times \mathrm{Y})\) 中的唯一线性映射，它对所有 \(f\in \mathcal{L}^2 (\mathrm{X})\) 与 \(g\in \mathcal{L}^2 (\mathrm{Y})\) 把 \(f\otimes g\) 映为由 \((x,y)\mapsto f(x)g(y)\) 定义的函数。

引理 10. --- 映射 \(\theta\) 过渡到商，定义了线性映射 \(\widetilde{\theta}\) （从 \(\mathrm{L}^2 (\mathrm{X})\otimes \mathrm{L}^2 (\mathrm{Y})\) 到 \(\mathrm{L}^2 (\mathrm{X}\times \mathrm{Y})\) 中的线性映射），且存在唯一的从 \(\mathrm{L}^2 (\mathrm{X})\widehat{\otimes}_2\mathrm{L}^2 (\mathrm{Y})\) 到 \(\mathrm{L}^2 (\mathrm{X}\times \mathrm{Y})\) 上的等距同构，它在 \(\mathrm{L}^2 (\mathrm{X})\otimes \mathrm{L}^2 (\mathrm{Y})\) 上与之相合。

第一个断言是初等的。我们来证明第二个。

设 \((f_i)_{i\in \mathrm{I}}\) 与 \((g_j)_{j\in \mathrm{J}}\) 分别是 \(\mathrm{L}^2 (\mathrm{X})\) 与 \(\mathrm{L}^2 (\mathrm{Y})\) 的标准正交基。族 \((\overline{f}_i\otimes g_j)_{(i,j)\in \mathrm{I}\times \mathrm{J}}\) 是 \(\mathrm{L}^2 (\mathrm{X})\widehat{\otimes}_2\mathrm{L}^2 (\mathrm{Y})\) 的标准正交基（EVT, V, p. 29，推论 1），而族 \((\widetilde{\theta} (\overline{f}_i\otimes g_j))_{(i,j)\in \mathrm{I}\times \mathrm{J}}\) 在 \(\mathrm{L}^2 (\mathrm{X}\times \mathrm{Y})\) 中标准正交。因此映射 \(\widetilde{\theta}\) 由连续性可扩张为从 \(\mathrm{L}^2 (\mathrm{X})\widehat{\otimes}_2\mathrm{L}^2 (\mathrm{Y})\) 到 \(\mathrm{L}^2 (\mathrm{X}\times \mathrm{Y})\) 中的等距线性映射。剩下要证明这个扩张是满射。

为此，只需证明 \(\widetilde{\theta}\) 的像在 \(\mathrm{L}^{2}(\mathrm{X}\times\mathrm{Y})\) 中稠密。设 N 是 \(\mathrm{L}^{2}(\mathrm{X}\times\mathrm{Y})\) 中与 \(\widetilde{\theta}\) 的像正交的元素。对所有 \(i\in I\) 与 \(j\in J\) ，我们有 \(\langle g_{j}\mid u_{\mathrm{N}}(f_{i})\rangle=\langle\widetilde{\theta}(\overline{f}_{i}\otimes g_{j})\mid\mathrm{N}\rangle=0\)。由公式 (12) 可得 \(u_{N}=0\) ，从而 N=0。

前一引理证实了 EVT, V, p. 29 例 2 中所述的断言。

在下文中，我们将借助引理 10 的等距同构把 \(\mathrm{L}^{2}(\mathrm{X})\widehat{\otimes}_{2}\mathrm{L}^{2}(\mathrm{Y})\) 与 \(\mathrm{L}^{2}(\mathrm{X}\times\mathrm{Y})\) 等同起来。

命题 19. --- 映射 \(N\mapsto u_{\mathrm{N}}\) 是从 \(\mathrm{L}^{2}(\mathrm{X}\times\mathrm{Y})\) 到空间 \(\mathcal{L}_{2}(\mathrm{L}^{2}(\mathrm{X});\mathrm{L}^{2}(\mathrm{Y}))\) （从 \(\mathrm{L}^{2}(\mathrm{X})\) 到 \(\mathrm{L}^{2}(\mathrm{Y})\) 中的希尔伯特--施密特映射所成的空间）上的等距同构。

从 \(\mathrm{L}^{2}(\mathrm{Y})\otimes\overline{\mathrm{L}^{2}(\mathrm{X})}\) 到 \(\mathcal{L}_{2}(\mathrm{L}^{2}(\mathrm{X});\mathrm{L}^{2}(\mathrm{Y}))\) 中把 \(g\otimes f\) 映为希尔伯特--施密特映射 \(h\mapsto\langle f|h\rangle g\) 的线性映射，可扩张为一个等距同构 \(\theta_{1}\) ，即从 \(\mathrm{L}^{2}(\mathrm{Y})\widehat{\otimes}_{2}\overline{\mathrm{L}^{2}(\mathrm{X})}\) 到 \(\mathcal{L}_{2}(\mathrm{L}^{2}(\mathrm{X});\mathrm{L}^{2}(\mathrm{Y}))\) 上的等距同构（EVT, V, p. 52，定理 1）。

另外，用 \(\theta_{2}\) 表示 \(\mathrm{L}^{2}(\mathrm{X})\widehat{\otimes}_{2}\mathrm{L}^{2}(\mathrm{Y})\) 到 \(\mathrm{L}^{2}(\mathrm{Y})\widehat{\otimes}_{2}\overline{\mathrm{L}^{2}(\mathrm{X})}\) 上的等距同构，它把 \(f\otimes g\) 映为 \(g\otimes \overline{f}\) ，对所有 \(f\in \mathrm{L}^2 (\mathrm{X})\) 与每个 \(g\in \mathrm{L}^2 (\mathrm{Y})\) 成立。

线性映射 \(\theta_{3} = \theta_{1} \circ \theta_{2}\) 等同于 \(\mathrm{L}^{2}(\mathrm{X} \times \mathrm{Y})\) 到 \(\mathcal{L}_{2}(\mathrm{L}^{2}(\mathrm{X}); \mathrm{L}^{2}(\mathrm{Y}))\) 上的一个等距同构。

设 \(f \in \mathrm{L}^2(\mathrm{X})\) 且 \(g \in \mathrm{L}^2(\mathrm{Y})\) ，并设 \(\mathrm{N}\) 是 \(\mathrm{L}^2(\mathrm{X} \times \mathrm{Y})\) 的元素（与 \(f \otimes g\) 等同）。由 INT, V, p. 95, § 8, no. 3 的推论 2 及公式 (12)，我们有

\[
\langle g _ {1} \mid u _ {\mathrm{N}} (f _ {1}) \rangle = \langle \overline {{f}} \mid f _ {1} \rangle \langle g _ {1} \mid g \rangle
\]

对所有 \(f_1 \in \mathrm{L}^2(\mathrm{X})\) 与 \(g_1 \in \mathrm{L}^2(\mathrm{Y})\) 成立。映射 \(u = \theta_3(\mathrm{N})\) 就是线性映射 \(\theta_1(g \otimes \overline{f})\) ；因此它满足 \(u(h) = \langle \overline{f} | h \rangle g\) ，对所有 \(h \in \mathrm{L}^2(\mathrm{X})\) 成立。于是，对所有 \(f_1 \in \mathrm{L}^2(\mathrm{X})\) 与 \(g_1 \in \mathrm{L}^2(\mathrm{Y})\) ，可得

\[
\langle g _ {1} \mid u (h _ {1}) \rangle = \langle \overline {{f}} \mid h _ {1} \rangle \langle g _ {1} \mid g \rangle = \langle g _ {1} \mid u _ {\mathrm{N}} (h _ {1}) \rangle ,
\]

由此得 \(\theta_3(N) = u_N\) 。由于 \(\theta_3\) 与 \(N \mapsto u_N\) 都连续，我们断定 \(\theta_3(N) = u_N\) 对所有 \(N \in L^2(X \times Y)\) 成立，证明完成。

推论. --- 对每个 \(N \in L^{2}(X \times Y)\) ，线性映射 \(u_{N}\) 是希尔伯特--施密特算子，且我们有 \(\mathrm{Tr}(u_{\mathrm{N}}^{*}u_{\mathrm{N}}) = \|N\|^{2}\) 。

事实上，\(\| u\| _2 = \sqrt{\mathrm{Tr}(u^*u)}\) 对所有 \(u\in \mathcal{L}_2(\mathrm{L}^2 (\mathrm{X});\mathrm{L}^2 (\mathrm{Y}))\) 成立。

注. --- 设 \(N \in L^{2}(X \times Y)\) 。由 IV, p. 150 的推论 2，存在可数集 I、标准正交族 \((f_{i})_{i \in I}\) （\(L^{2}(X)\) 中的）与标准正交族 \((g_{i})_{i \in I}\) （\(L^{2}(Y)\) 中的），以及族 \((\alpha_{i})_{i \in I}\) （\(\mathbf{R}_{+}^{*}\) 中的族），使得

\[
u _ {\mathrm{N}} (f) = \sum_ {i \in \mathrm{I}} \alpha_ {i} \langle f _ {i} | f \rangle g _ {i}
\]

对所有 \(f \in \mathrm{L}^2(\mathrm{X})\) 成立，其中级数在 \(\mathrm{L}^2(\mathrm{Y})\) 中收敛。由 IV, p. 165 的推论 4 及上述推论，我们有

\[
\sum_ {i \in \mathrm{I}} \alpha_ {i} ^ {2} = \mathrm{Tr} (u _ {\mathrm{N}} ^ {*} u _ {\mathrm{N}}) = \| u _ {\mathrm{N}} \| _ {2} ^ {2} = \int_ {\mathrm{X} \times \mathrm{Y}} | \mathrm{N} (x, y) | ^ {2} d (\mu \otimes \nu) (x, y).
\]

另外，注意 \(h_{i,j} = \overline{f}_{i} \otimes g_{j} \in \mathrm{L}^{2}(\mathrm{X} \times \mathrm{Y})\) 对所有 \((i, j) \in \mathrm{I} \times \mathrm{I}\) 成立。于是我们有

\[
\mathrm{N} = \sum_ {i \in \mathrm{I}} \alpha_ {i} h _ {i, i}
\]

in \(\mathrm{L}^2 (\mathrm{X}\times \mathrm{Y})\) .

事实上，设 \((f_{j})_{j\in\mathrm{J}}\) 与 \((g_{k})_{k\in\mathrm{K}}\) 分别是 \(\mathrm{L}^{2}(\mathrm{X})\) 与 \(\mathrm{L}^{2}(\mathrm{Y})\) 的标准正交基，且分别是族 \((f_{i})_{i\in\mathrm{I}}\) 与 \((g_{i})_{i\in\mathrm{I}}\) 的扩张。置 \(h_{j,k}=\overline{f}_{j}\otimes g_{k}\) ，对所有 \((j,k)\in\mathrm{J}\times\mathrm{K}\) 。由引理 10，族 \((h_{j,k})_{(j,k)\in\mathrm{J}\times\mathrm{K}}\) 是 \(\mathrm{L}^{2}(\mathrm{X}\times\mathrm{Y})\) 的标准正交基。我们有 \(\langle h_{j,k}\mid\mathrm{N}\rangle=\langle g_{k}\mid u_{\mathrm{N}}(f_{j})\rangle\) ，对所有 \((j,k)\in\mathrm{J}\times\mathrm{K}\) 。若 \(j\notin I\) ，这个量为零。若 \(j\in I\) ，它等于 \(\alpha_{j}\langle g_{k}\mid g_{j}\rangle\) ，故除 k=j 外为零，而当 k=j 时 \(\langle h_{j,j}\mid N\rangle=\alpha_{j}\) 。因此

\[
\mathrm{N} = \sum_ {(j, k) \in \mathrm{J} \times \mathrm{K}} \left\langle h _ {j, k} \mid \mathrm{N} \right\rangle h _ {j, k} = \sum_ {i \in \mathrm{I}} \alpha_ {i} h _ {i, i}.
\]

\subsection{连续核积分算子的迹}

在本节中，我们保持前一节的约定，取 Y = X 且 \(\nu = \mu\)。我们假设 X 是在无穷远处可数的局部紧拓扑空间（TG, I, p. 68，定义 5）。

特别地，我们把空间 \(\mathrm{L}^{2}(\mathrm{X}\times\mathrm{X})\) 与 \(\mathrm{L}^{2}(\mathrm{X})\widehat{\otimes}_{2}\mathrm{L}^{2}(\mathrm{X})\) 等同起来（IV, p. 172 的引理 10）。我们用 \(\widetilde{f}\) 表示 \(\mathrm{L}^{2}(\mathrm{X})\) （相应地，\(\mathrm{L}^{2}(\mathrm{X}\times\mathrm{X})\) ）中函数 \(f\in\mathcal{L}^{2}(\mathrm{X})\) （相应地，\(\mathcal{L}^{2}(\mathrm{X}\times\mathrm{X})\) 中的函数）的类。

命题 20. --- 设 \((f_n)_{n \in \mathbf{N}}\) 是从 X 到可度量化空间 F 中的可测映射序列。假设 \(f_n(x)\) 的极限在 X 的一个 \(\mu\)-零测子集之外存在。则存在 X 的紧子集序列 \((\mathrm{C}_m)_{m \in \mathbf{N}}\) ，其并 \(\widetilde{\mathrm{X}}\) 满足 \(|\mu|(\mathrm{X} - \widetilde{\mathrm{X}}) = 0\) ，使得函数 \(f_n\) 在 \(\mathrm{C}_m\) 上连续（对所有 \(n \in \mathbf{N}\) 与每个 \(m \in \mathbf{N}\)），且序列 \((f_n)_{n \in \mathbf{N}}\) 一致收敛到 \(f\) （在 \(\mathrm{C}_m\) 上，对所有 \(m \in \mathbf{N}\)）。

由 INT, IV, p. 175, § 5, no. 4 的定理 2 及 INT, IV, p. 188, § 5, no. 8 的命题 12, b)，X 中使得函数 \(f_n\) 对每个 \(n\) 在 C 上连续且 \((f_n)\) 在 C 上一致收敛到 \(f\) 的紧子集 C 所成之集在 X 中 \(\mu\)-稠密（INT, IV, p. 189, § 5, no. 8，定义 6）。断言随即由 INT, IV, p. 189, § 5, no. 8 的注得出。

设 \(N\) 是属于 \(\mathcal{L}^2 (\mathrm{X}\times \mathrm{X})\) 的函数，\(u_{\mathrm{N}}\) 是从 \(\mathrm{L}^2 (\mathrm{X})\) 到 \(\mathrm{L}^2 (\mathrm{X})\) 中的以 \(N\) 为核的希尔伯特--施密特映射（IV, p. 173 的命题 19）。由于映射 \(u_{\mathrm{N}}\) 是紧的，由 IV, p. 156 的命题 9，存在元素 \(M\) （\(\overline{\mathbf{N}}\) 中的元素），并且用 \(I\) 表示整数 \(n\in \mathbf{N}\) 中满足 \(n\leqslant \mathrm{M}\) 者所成之集，存在标准正交族 \((f_i)_{i\in \mathrm{I}}\) 与 \((g_{i})_{i\in \mathrm{I}}\) （\(\mathcal{L}^2 (\mathrm{X})\) 中的标准正交族）以及族 \((\alpha_{i})_{i\in \mathrm{I}}\) （\(\mathbf{R}_+^*\) 中的族），使得

\[
u _ {\mathrm{N}} (f) = \sum_ {i \in \mathrm{I}} \alpha_ {i} \langle \widetilde {f} _ {i} | f \rangle \widetilde {g} _ {i}\tag{13}
\]

对所有 \(f \in \mathrm{L}^2(\mathrm{X})\) 成立，其中级数在 \(\mathrm{L}^2(\mathrm{X})\) 中收敛。

对每个 \(i \in I\) ，用 \(h_i \in \mathcal{L}^2(\mathrm{X} \times \mathrm{X})\) 表示由 \(h_i(x, y) = \overline{f_i(x)} g_i(y)\) 定义的函数，从而 \(\widetilde{h}_i\) 是 \(\overline{f}_i \otimes g_i\) 的类。由 IV, p. 173 的注 9，通项为 \(\alpha_i h_i\) 的级数在 \(\mathrm{L}^2(\mathrm{X} \times \mathrm{X})\) 中收敛到 N。

在本节其余部分，我们假设 N 是连续函数。

命题 21. --- 假设 \(u_{\mathrm{N}}\) 有有限迹。则存在一个集 \(\widetilde{\mathrm{X}} \subset \mathrm{X}\) ，其余集 \(\mathrm{X} - \widetilde{\mathrm{X}}\) 是 \(\mu\)-零测的，还存在函数 \(\mathrm{H} \in \mathcal{L}^2(\mathrm{X} \times \mathrm{X})\) ，满足以下条件：

\begin{enumerate}
\def\labelenumi{(\roman{enumi})}
\item
  对每个 \((x,y)\in\widetilde{\mathrm{X}}\times\widetilde{\mathrm{X}}\) ，族 \((\alpha_{i}\overline{f_{i}(x)}g_{i}(y))_{i\in\mathrm{I}}\) 在 \(\mathbf{C}\) 中可和，其和为 \(\mathrm{N}(x,y)\) ；
\item
  对 I 的每个有限子集 J 与每个 \((x,y)\in\widetilde{\mathbf{X}}\times\widetilde{\mathbf{X}}\) ，我们有
\end{enumerate}

\[
\left| \sum_ {i \in \mathrm{J}} \alpha_ {i} h _ {i} (x, y) \right| \leqslant \mathrm{H} (x, y);\tag{14}
\]

\begin{enumerate}
\def\labelenumi{(\roman{enumi})}
\setcounter{enumi}{2}
\tightlist
\item
  函数 \(x \mapsto \mathrm{H}(x, x)\) 属于 \(\mathcal{L}^1(\mathrm{X})\) 。
\end{enumerate}

由于 \(u_{N}\) 有有限迹，族 \((\alpha_{i})\) 可和（IV, p. 165 的推论 4）。级数

\[
\sum_ {i \in \mathrm{I}} \alpha_ {i} | f _ {i} | ^ {2}, \quad \sum_ {i \in \mathrm{I}} \alpha_ {i} | g _ {i} | ^ {2}
\]

因此在 \(\mathcal{L}^{1}(X)\) 中收敛，从而 \(\mu\)-几乎处处收敛。用 F 与 G 分别表示由这些级数的和在 \(\mu\)-几乎处处意义下定义的函数，其中置 \(F(x) = 0\) 与 \(G(x) = 0\) ，对相应级数不收敛的每个 \(x\) 。由于 F 与 G 属于 \(\mathcal{L}^{1}(X)\)，由 \(H(x,y) = \sqrt{F(x)G(y)}\) 定义的函数 H 属于 \(\mathcal{L}^{2}(X \times X)\) （INT, V, p. 95, § 8, no. 3，推论 2）。

把命题 20 应用于空间 X，取 \(F = \mathbf{C}^{2}\) ，并取可测映射

\[
s_{n}\colon x\mapsto \Bigl (\sum_{\substack{i\in \mathrm{I}\\ i\leqslant n}}\alpha_{i}|f_{i}(x)|^{2},\sum_{\substack{i\in \mathrm{I}\\ i\leqslant n}}\alpha_{i}|g_{i}(x)|^{2}\Bigr)
\]

它们定义在 X 上。因此存在 X 的紧子集序列 \((\mathrm{C}_{m})_{m\in\mathbf{N}}\) ，满足以下条件：

\begin{enumerate}
\def\labelenumi{(\arabic{enumi})}
\item
  并 \(\widetilde{X}\) 满足 \(\mu(\mathrm{X}-\widetilde{\mathrm{X}})=0\) ；
\item
  对所有 \(m \in \mathbf{N}\) 与每个 \(n \in \mathbf{N}\) ，函数 \(s_n\) 在 \(\mathrm{C}_m\) 上连续；
\item
  对所有 \(m \in \mathbf{N}\) ，级数 \(\sum \alpha_{i}|f_{i}|^{2}\) 与 \(\sum \alpha_{i}|g_{i}|^{2}\) 在 \(C_{m}\) 上分别一致收敛到 F 与 G；
\item
  由 \(\mu\) 在 \(C_{m}\) 上诱导的测度的支撑等于 \(C_{m}\) （必要时把 \(C_{m}\) 换成这个测度的支撑）。
\end{enumerate}

我们来证明集 \(\widetilde{\mathbf{X}}\) 与函数 H 满足条件 (i)、(ii) 与 (iii)。

设 \((x,y)\in \widetilde{\mathbf{X}}\times \widetilde{\mathbf{X}}\)。对每个有限子集 \(\mathrm{J}\subset \mathrm{I}\) ，我们有

\[
\left| \sum_ {i \in \mathrm{J}} \alpha_ {i} \overline {{f _ {i} (x)}} g _ {i} (y) \right| \leqslant \left(\sum_ {i \in \mathrm{J}} \alpha_ {i} | f _ {i} (x) | ^ {2}\right) ^ {1 / 2} \left(\sum_ {i \in \mathrm{J}} \alpha_ {i} | g _ {i} (y) | ^ {2}\right) ^ {1 / 2}\tag{15}
\]

进而有

\[
\left| \sum_ {i \in \mathrm{J}} \alpha_ {i} \overline {{f _ {i} (x)}} g _ {i} (y) \right| \leqslant \mathrm{H} (x, y),\tag{16}
\]

这已给出性质 (ii)。

由不等式 (15)，级数 \(\sum_{i} \alpha_{i} h_{i}\) 在 \(\mathrm{K}_{m} \times \mathrm{K}_{n}\) 上一致收敛（对所有 \((m, n) \in \mathbf{N}^{2}\)）。设 \(\widetilde{\mathrm{N}}\) 是 \(\mathrm{X} \times \mathrm{X}\) 上由

\[
\widetilde {\mathrm{N}} (x, y) = \sum_ {i \in \mathrm{I}} \alpha_ {i} h _ {i} (x, y) = \sum_ {i \in \mathrm{I}} \alpha_ {i} \overline {{f _ {i} (x)}} g _ {i} (y)
\]

对所有 \((x,y)\in\widetilde{\mathrm{X}}\times\widetilde{\mathrm{X}}\) 定义、否则置 \(\widetilde{\mathrm{N}}(x,y)=0\) 的函数。函数 \(\widetilde{N}\) 可测（INT, IV, p. 175, § 5, n°4，定理 2），且在 \(K_{m}\times K_{n}\) 上连续（对所有 \((m,n)\in\mathbf{N}^{2}\)）。

由 (16)，我们有 \(|\widetilde{\mathrm{N}}(x,y)| \leqslant \mathrm{H}(x,y)\) ，对所有 \((x,y) \in \mathrm{X} \times \mathrm{X}\) 成立，特别地 \(\widetilde{\mathrm{N}}\) 属于 \(\mathcal{L}^2(\mathrm{X} \times \mathrm{X})\) （INT, IV, p. 84, § 5, no. 6，定理 5）。

我们来证明 \(N = \widetilde{N}\) 在 \(\widetilde{X} \times \widetilde{X}\) 上成立，这将给出性质 (i)。设 f 与 g 是 \(\mathcal{L}^{2}(X)\) 的元素。我们有

\[
\langle g \mid u _ {\widetilde {N}} (f) \rangle = \int_ {X \times X} \widetilde {N} (x, y) f (x) \overline {{g (y)}} d (\mu \otimes \mu) (x, y)
\]

（IV, p. 172 的公式 (12)）。对每个 \((x,y)\in \widetilde{\mathrm{X}}\times \widetilde{\mathrm{X}}\) 与 I 的每个有限子集 J，我们有

\[
\left| \sum_ {i \in J} \alpha_ {i} \overline {{f _ {i} (x)}} g _ {i} (y) f (x) \overline {{g (y)}} \right| \leqslant | f (x) g (y) | H (x, y)
\]

（公式 (16)）。由于这个不等式的右端在 X × X 上可积（INT, V, p. 95, § 8, n° 3，推论 2），可以应用勒贝格定理（INT, IV, p. 137, § 3, n° 7，定理 6）及公式 (13)，推出

\[
\begin{array}{c} \langle g | u _ {\widetilde {N}} (f) \rangle = \sum_ {i \in I} \alpha_ {i} \int_ {X \times X} \overline {{f _ {i} (x)}} g _ {i} (y) f (x) \overline {{g (y)}} d (\mu \otimes \mu) (x, y) \\ = \sum_ {i \in I} \alpha_ {i} \langle f _ {i} | f \rangle \langle g | g _ {i} \rangle = \langle g | u _ {N} (f) \rangle . \end{array}
\]

因此可得 \(u_{\widetilde{\mathrm{N}}} = u_{\mathrm{N}}\) ，从而 \(\mathrm{N} = \widetilde{\mathrm{N}}\) 在 \(\mathrm{L}^2(\mathrm{X} \times \mathrm{X})\) 中成立 （III, p. 28 的命题 3, b)））。

对每个 \((m,n)\in\mathbf{N}^{2}\)，函数 N 与 \(\widetilde{N}\) 在 \(K_{m}\times K_{n}\) 上连续，故在 \(K_{m}\times K_{n}\) 上相等（INT, III, p. 69, § 2, n°2 的命题 9），因为 \(\mu\otimes\mu\) 在 \(K_{m}\times K_{n}\) 上诱导的测度的支撑等于 \(K_{m}\times K_{n}\) 。因此 N 与 \(\widetilde{N}\) 在 \(\widetilde{X}\times\widetilde{X}\) 上相合。

最后，上界 \(\sqrt{\mathrm{FG}} \leqslant (\mathrm{F} + \mathrm{G}) / 2\) 蕴含函数 \(x \mapsto \sqrt{\mathrm{F}(x)\mathrm{G}(x)} = \mathrm{H}(x,x)\) 在 X 上可积，由此得到性质 (iii)。

在本节其余部分，我们保持该命题的记号。

定理 2. --- 假设 \(u_{\mathrm{N}}\) 有有限迹。则函数 \(x \mapsto \mathrm{N}(x, x)\) 属于 \(\mathcal{L}^1(\mathrm{X})\) ，且我们有

\[
\mathrm{Tr} (u _ {\mathrm{N}}) = \int_ {\mathrm{X}} \mathrm{N} (x, x) d \mu (x).
\]

由条件 (i) 与 (ii)，我们有 \(|\mathrm{N}(x,x)| \leqslant \mathrm{H}(x,x)\) （对 \(x \in \widetilde{X}\)），故由条件 (iii)，函数 \(x \mapsto \mathrm{N}(x,x)\) 属于 \(\mathcal{L}^{1}(\mathrm{X})\) 。

条件 (i) 给出等式

\[
\mathrm{N} (x, x) = \sum_ {i \in \mathrm{I}} \alpha_ {i} \overline {{f _ {i} (x)}} g _ {i} (x)
\]

对所有 \(x \in \widetilde{\mathrm{X}}\) 成立。由条件 (ii) 与 (iii)，可以应用勒贝格定理（INT, IV, p. 137, § 3, no. 7，定理 6），由此可得

\[
\begin{array}{r} \int_ {\mathrm{X}} \mathrm{N} (x, x) d \mu (x) = \sum_ {i \in \mathrm{I}} \alpha_ {i} \int_ {\mathrm{X}} \overline {{f _ {i} (x)}} g _ {i} (x) d \mu (x) \\ = \sum_ {i \in \mathrm{I}} \alpha_ {i} \langle f _ {i} | g _ {i} \rangle = \mathrm{Tr} (u _ {\mathrm{N}}), \end{array}
\]

最后一个等式来自 IV, p. 170 命题 17 的推论。

引理 11. --- 若自同态 \(u_{\mathrm{N}}\) （\(\mathrm{L}^2(\mathrm{X})\) 的自同态）是正的，则 \(\mathrm{N}(x,x) \geqslant 0\) 对每个 \(x\) （\(\mu\) 的支撑中的元素）成立。

由于 \(u_{N}\) 是正的，可以选取族 \((f_{i})\) 与 \((g_{i})\) ，使得 \(f_{i} = g_{i}\) 对所有 \(i \in I\) 成立（参见 IV, p. 151 的注 3）。于是对每个 \(x \in \widetilde{X}\) ，我们有

\[
\mathrm{N} (x, x) = \sum_ {i \in \mathrm{I}} \alpha_ {i} | f _ {i} (x) | ^ {2} \geqslant 0.
\]

由于 \(N\) 连续且 \(\mu(\mathbf{X}-\widetilde{\mathbf{X}})=0\) ，函数 \(x \mapsto N(x,x)\) 在 \(\mu\) 的支撑上非负。

注. --- 该引理的逆命题不成立（IV, p. 316 的习题 14）。

命题 22. --- 假设 \(u_{\mathrm{N}}\) 是 \(\mathrm{L}^2(\mathrm{X})\) 的正自同态。则自同态 \(u_{\mathrm{N}}\) 有有限迹当且仅当函数 \(x \mapsto \mathrm{N}(x,x)\) 是 \(\mu\)-可积的。此时我们有

\[
\mathrm{Tr} (u _ {\mathrm{N}}) = \int_ {\mathrm{X}} \mathrm{N} (x, x) d \mu (x).
\]

若 \(u_{N}\) 有有限迹，定理 2 蕴含 \(x \mapsto \mathrm{N}(x, x)\) 在 X 上可积，且其积分是 \(u_{N}\) 的迹。

反之，设函数 \(x \mapsto \mathrm{N}(x, x)\) 可积。由于 \(u_{N}\) 是正的，可以选取标准正交族 \((f_i)_{i \in I}\) 与 \((g_i)_{i \in I}\) ，使得 \(f_i = g_i\) 对所有 \(i \in I\) 成立（IV, p. 151 的注 3）。对每个 \(x \in \widetilde{X}\) ，我们有

\[
\mathrm{N} (x, x) = \sum_ {i \in \mathrm{I}} \alpha_ {i} | f _ {i} (x) | ^ {2},
\]

由此，对 I 的每个有限子集 J，我们推出

\[
\begin{array}{c} \sum_ {i \in \mathrm{J}} \alpha_ {i} = \sum_ {i \in \mathrm{J}} \alpha_ {i} \int_ {\mathrm{X}} | f _ {i} (x) | ^ {2} d \mu (x) \\ = \int_ {\mathrm{X}} \sum_ {i \in \mathrm{J}} \alpha_ {i} | f _ {i} (x) | ^ {2} d \mu (x) \leqslant \int_ {\mathrm{X}} \mathrm{N} (x, x) d \mu (x). \end{array}
\]

因此族 \((\alpha_{i})_{i\in\mathrm{I}}\) 可和，这蕴含自同态 \(u_{N}\) 有有限迹（IV, p. 164 的命题 12）。

注. --- 即使当 X 紧且 N 连续时，自同态 \(u_{\mathrm{N}}\) （\(\mathrm{L}^{2}(\mathrm{X})\) 的自同态）也不总有有限迹（参见 IV, p. 314 的习题 8）。

\section{正规自同态}

在本章中，除另有说明外，所考虑的希尔伯特空间都是复的。

\subsection{\texorpdfstring{关于空间 \(L^{p}(X,\mu)\) 的补充}{关于空间 L\^{}\{p\}(X,\textbackslash mu) 的补充}}

命题 1. --- 设 X 是局部紧拓扑空间。设 \(\mu\) 是 X 上的正测度，\(p \in [1, +\infty[\)。设 \((\mathrm{X}_i)_{i \in \mathrm{I}}\) 是 X 中两两不相交的局部紧子集的局部可数族（INT, IV, p. 190, § 5, no. 9，定义 7），使得诸 \(\mathrm{X}_i\) 之并的余集是局部 \(\mu\)-零测的。设 \(\varphi_i\) 是 \(\mathrm{X}_i\) 的特征函数，\(\mu_i\) 是 \(\mu\) 在 \(\mathrm{X}_i\) 上诱导的测度（INT, IV, p. 186, §5, n° 7，定义 4）。

\begin{enumerate}
\def\labelenumi{\alph{enumi})}
\item
  映射 \(p_{i}\) ：\(f \mapsto f\varphi_{i}\) 是 \(\mathcal{L}^{p}(\mathrm{X}, \mu)\) 的投影子，且过渡到商，定义了投影子 \(\widetilde{p}_{i}\) （\(\mathrm{L}^{p}(\mathrm{X}, \mu)\) 的投影子）。当 p = 2 时，后者是 \(\mathrm{L}^{2}(\mathrm{X}, \mu)\) 的正交投影子；
\item
  映射 \(f \mapsto f|X_i\) 诱导出从 \(p_i\) 的像到空间 \(\mathcal{L}^p(X_i, \mu_i)\) 上的等距同构，且过渡到商，定义了从 \(\widetilde{p}_i\) 的像到 \(\mathrm{L}^p(X_i, \mu_i)\) 上的等距同构，借助它这两个空间互相等同；
\item
  诸空间 \(\mathrm{L}^{p}(\mathrm{X}_{i},\mu_{i})\) 之和在空间 \(\mathrm{L}^{p}(\mathrm{X},\mu)\) 中是完全的。当 p = 2 时，空间 \(\mathrm{L}^{2}(\mathrm{X},\mu)\) 是诸空间 \(\mathrm{L}^{2}(\mathrm{X}_{i},\mu_{i})\) 的希尔伯特和。
\end{enumerate}

断言 \(a\)）是初等的。断言 \(b\)）由 INT, V, p. 84, § 7, n° 1 的评注得出。

设 \(q \in ]1, +\infty]\) 是 \(p\) 的共轭指数。我们把 \(L^p(X, \mu)\) 的对偶与 \(L^q(X, \mu)\) 等同起来（INT, V, p. 61, § 5, n° 8，定理 4）。设 \(f\) 是 \(\mathcal{L}^q(X, \mu)\) 中的函数，其类 \(\widetilde{f}\) 在 \(L^q(X, \mu)\) 中满足 \(\langle \widetilde{f}, \widetilde{p}_i(\varphi) \rangle = 0\) （对所有 \(i \in I\) 与每个 \(\varphi \in L^p(X, \mu)\)）。由于 \(p_i\) 的像包含 \(\mathcal{K}(X_i)\) ，可得测度 \((f|X_i) \cdot \mu_i\) 在 \(X_i\) 上为零，即 \(f\) 在 \(X_i\) 上的限制关于 \(\mu_i\) 在 \(X_i\) 上局部零测（INT, V, p. 46, § 5, n° 3，推论 2）。由于诸 \(X_i\) 之并在 X 中的余集是局部 \(\mu\)-零测的，可以断定函数 \(f\) 在 X 上局部 \(\mu\)-零测（INT, IV, p. 190, § 5, n° 9 及 p. 172, n° 2，命题 5）。于是 \(f\) 的类在 \(L^q(X, \mu)\) 中为零（当 \(p \neq 1\) 时用到 INT, V, p. 8, § 1, n° 3 的引理 1 及命题 9 的推论）。由哈恩--巴拿赫定理（EVT, II, p. 46，推论 1），断言 \(c\)）的第一部分得证。

若 \(p = 2\) ，\(p_i\) 的像与 \(p_j\) 的像正交（对所有 \(i \neq j\) ），因为此时 \(X_i\) 与 \(X_j\) 不相交，由此得到最后一个断言。

命题 2. --- 设 X 是在无穷远处可数的局部紧空间。设 \(\mu\) 是 X 上的正测度。对每个 \(p \in [1, +\infty[\)，空间 \(\mathrm{L}^{p}(\mathrm{X},\mu)\) 是可数型的。

设 \((\mathrm{U}_{n})_{n\in\mathbf{N}}\) 是 X 中相对紧开集的序列，其并等于 X 且满足 \(\overline{U}_{n}\subset U_{n+1}\) 对所有 \(n\in\mathbf{N}\) 成立（TG, I, p. 68，命题 15）。对每个 \(n\in\mathbf{N}\)，空间 \(\mathcal{K}(\mathrm{X},\overline{\mathrm{U}}_{n})\) 等同于巴拿赫空间 \(\mathcal{C}(\overline{\mathrm{U}}_{n})\) 的闭子空间（INT, III, p. 40, § 1, n° 1）；由于后一空间是可数型的（TG, X, p. 25，推论），\(\mathcal{K}(\mathrm{X},\overline{\mathrm{U}}_{n})\) 亦然（TG, IX, p. 19, cor., (i)）。设 \(F_{n}\) 是 \(\mathcal{K}(\mathrm{X},\overline{\mathrm{U}}_{n})\) 的可数稠密子集。

设 \(f \in \mathcal{L}^p(\mathrm{X},\mu)\) 且 \(\varepsilon > 0\) 。存在整数 \(n \in \mathbf{N}\) 使得 \(\int_{\mathrm{X}-\mathrm{U}_n}|f|^p < \varepsilon / 2\) ，且存在 \(g \in \mathcal{F}_n\) 使得 \(\int_{\overline{\mathrm{U}}_n}|f - g|^p < \varepsilon / 2\) 。因此，\(\mathrm{L}^p(\mathrm{X},\mu)\) 中诸集合 \(\mathcal{F}_n\) 的元素之类所成的并在 \(\mathrm{L}^p(\mathrm{X},\mu)\) 中稠密，证明完成（TG, IX, p. 18，命题 12）。

\subsection{可测函数的本质像}

在本节中，我们用 X 表示局部紧拓扑空间，用 \(\mu\) 表示 X 上的正测度。

定义 1. --- 设 Y 是拓扑空间。对每个从 X 到 Y 中的可测函数 g，\(\mu\)-本质像（\(g\) 的）是满足如下条件的 \(y \in Y\) 所成之子集：对每个开邻域 \(U\) （\(y\) 的），集 \(g^{-1}(\mathrm{U})\) 在 X 中不是局部 \(\mu\)-零测的（INT, IV, p. 172, § 5, no. 2，定义 3）。

引理 1. --- 设 \(g\) 是从 X 到拓扑空间 Y 中的 \(\mu\)-可测函数，并设 \(S\) 是它的 \(\mu\)-本质像。元素 \(x \in X\) 中使得 \(g(x)\) 不属于 S 者所成之集在 X 中局部 \(\mu\)-零测。

设 \(Z = g^{-1}(Y - S)\) 是所考虑的集。

首先假设 X 紧且 g 连续。在此情形，前推测度 \(g(\mu)\) 有定义（INT, V, p. 69, § 6, no. 1，定义 1），因为 \(\mu\) 是有界测度。由定义可得，\(\mu\)-本质像（\(g\) 的）是测度 \(g(\mu)\) 的支撑（参见 INT, V, p. 70, § 6, no. 2，推论 2），于是 \(\mu(Z) = g(\mu)(Y - S) = 0\) （INT, IV, p. 118, § 2, no. 2，命题 5）。

现在考虑一般情形。设 C 是 X 的紧子集，\(\varepsilon > 0\) 是实数。由于 \(g\) 可测，存在紧子集 \(C_1 \subset C\) ，使得 \(\mu(C - C_1) \leqslant \varepsilon\) 且 \(g|C_1\) 连续（INT, IV, p. 169, § 5, no. 1，命题 1）。于是我们有

\[
\mu (\mathrm{Z} \cap \mathrm{C}) \leqslant \mu (\mathrm{C} - \mathrm{C} _ {1}) + \mu (\mathrm{Z} \cap \mathrm{C} _ {1}) \leqslant \varepsilon + \mu (\mathrm{Z} \cap \mathrm{C} _ {1}).
\]

设 \(\mu_{1}\) 是 \(\mu\) 在 \(C_{1}\) 上诱导的测度（INT, IV, p. 186, § 5, no. 7，定义 4）。设 \(S_{1}\) 是 \(\mu_{1}\)-本质像（\(g|C_{1}\) 的）。我们有 \(Z\cap C_{1}\subset(g|C_{1})^{-1}(Y-S_{1})\) 。由第一种情形，集 \(Z\cap C_{1}\) 因此是 \(\mu_{1}\)-零测的，从而是 \(\mu\)-零测的（INT, IV, p. 187, § 5, no. 7，引理 2, (i)）。上面的不等式变成 \(\mu(Z\cap C)\leqslant\varepsilon\) ；由于 \(\varepsilon\) 与 C 是任意的，集 Z 局部 \(\mu\)-零测（INT, IV, p. 172, § 5, no. 2，命题 5）。

引理 2. --- 设 \(g\) 是从 X 到 \(\mathbf{C}\) 中的连续函数。\(\mu\)-本质像（\(g\) 的）是 \(\mu\) 的支撑在 \(g\) 下的像的闭包。

设 \(Y\) 是 \(\mu\) 的支撑。若 \(z \in \mathbf{C}\) 不附着于 \(g(Y)\)，则存在开邻域 \(U\) （\(z\) 的），使得开集 \(g^{-1}(U)\) 不与 Y 相交，从而是局部 \(\mu\)-零测的；这意味着 \(z\) 不属于 \(\mu\)-本质像（\(g\) 的）。

反之，若 \(z \in \mathbf{C}\) 附着于 \(g(Y)\)，则对 z 的每个开邻域 \(U\) ，集 \(g^{-1}(U)\) 是 X 中与 Y 相交的开集。因此它不是 \(\mu\)-零测的，又由于它是开的，它不是局部 \(\mu\)-零测的（INT, IV, p. 172, § 5, n°2，推论 2）。于是 z 属于 g 的 \(\mu\)-本质像。

引理 3. --- 设 \(g\) 是从 X 到 \(\mathbf{C}\) 中的连续函数，S 是它的 \(\mu\)-本质像。假设 S 非空且 \(0 \notin S\) ，并用 \(\delta\) 表示 0 到 S 的距离。则 \(\delta > 0\)。

置 \(h(x) = 1 / g(x)\) （若 \(g(x) \neq 0\)），否则置 \(h(x) = 0\) 。函数 \(h\) 属于 \(\mathcal{L}^{\infty}(\mathrm{X},\mu)\)。设 \(\widetilde{h}\) 是 \(h\) 在 \(\mathrm{L}^{\infty}(\mathrm{X},\mu)\) 中的类。我们有公式 \(\delta^{-1} = \| \widetilde{h}\|_{\infty}\) 。

设 U 是 0 的开邻域，使得开集 \(Z = g^{-1}(U)\) 局部 \(\mu\)-零测，从而是零测的（INT, IV, p. 172, § 5, no. 2，推论 2）。设 Y 是 \(\mu\) 的支撑；我们有 \(Y \subset X - Z\) 。函数 \(h\) 在 \(X - Z\) 上的限制连续且有界，故 \(h \in \mathcal{L}^{\infty}(X, \mu)\) ，且 \(\widetilde{h}\) 在 \(\mathcal{L}^{\infty}(X, \mu)\) 中的范数等于它在 \(X - Z\) 上的限制的范数。此外，对每个 \(\alpha \in R_{+}\) ，\(x \in X - Z\) 中满足 \(|h(x)| > \alpha\) 者所成之集在 \(X - Z\) 中是开的；因此它局部 \(\mu\)-零测当且仅当它不与 Y 相交（INT, IV，同前及 INT, III, p. 66, § 3, no. 2，定义 1）。于是，我们有

\[
\| \widetilde {h} \| _ {\infty} = \sup _ {x \in \mathrm{Y}} \frac {1}{| g (x) |},
\]

由此得

\[
\frac {1}{\| \widetilde {h} \| _ {\infty}} = \inf _ {x \in \mathrm{Y}} | g (x) | = \inf _ {\lambda \in g (\mathrm{Y})} | \lambda | = \inf _ {\lambda \in \overline {{g (\mathrm{Y})}}} | \lambda |
\]

由前一引理，它等于 δ。

\subsection{普遍可测函数}

在本节中，X 是局部紧拓扑空间。我们回顾（INT, V, p. 28, § 3, n° 4，定义 2），从 X 到拓扑空间 Y 中的映射 f 称为普遍可测的，如果它对 X 上的每个正测度 \(\mu\) 都是 \(\mu\)-可测的。为此只需 f 对 X 上每个具有紧支撑的正测度 \(\mu\) 是 \(\mu\)-可测的（INT, V, p. 28, § 4, no. 3，命题 6）。

引理 4. --- 设 Y 与 Z 是局部紧拓扑空间，\(f\colon X\to Y\) 与 \(g\colon Y\to Z\) 是普遍可测映射。映射 \(g\circ f\colon X\to Z\) 是普遍可测的。

设 \(\mu\) 是 X 上具有紧支撑的正测度，C 是它的支撑。\(f\) 在 C 上的限制是 \((\mu|\mathrm{C})\)-正常的。映射 \(g\) 关于像测度 \((f|\mathrm{C})(\mu|\mathrm{C})\) 可测，故映射 \((g \circ f)|\mathrm{C} = g \circ (f|\mathrm{C})\) 是 \((\mu|\mathrm{C})\)-可测的。因此映射 \(g \circ f\) 是 \(\mu\)-可测的。引理得证。

我们用 \(\mathcal{L}_{\mathrm{u}}(\mathrm{X})\) 表示在 X 上普遍可测的复值函数所成的向量空间，用 \(\mathcal{L}_{\mathrm{u}}^{\infty}(\mathrm{X})\) 表示在 X 上有界的函数 \(f \in \mathcal{L}_{\mathrm{u}}(\mathrm{X})\) 所成的子空间。对 \(f \in \mathcal{L}_{\mathrm{u}}^{\infty}(\mathrm{X})\)，我们记 \(\|f\|_{\infty} = \sup_{x \in X} |f(x)|\) 。

命题 3. --- a) 集 \(\mathcal{L}_{\mathrm{u}}(\mathrm{X})\) 是从 X 到 C 中的函数所成的对合代数的对合子代数；

\begin{enumerate}
\def\labelenumi{\alph{enumi})}
\setcounter{enumi}{1}
\item
  集 \(\mathcal{L}_{\mathrm{u}}^{\infty}(\mathrm{X})\) 是 \(\mathcal{L}_{\mathrm{u}}(\mathrm{X})\) 的子代数，且由 \(f\mapsto \| f\|_{\infty}\) 定义的映射是 \(\mathcal{L}_{\mathrm{u}}^{\infty}(\mathrm{X})\) 上的一个范数，对此范数 \(\mathcal{L}_{\mathrm{u}}^{\infty}(\mathrm{X})\) 是有单位元的星代数；
\item
  代数 \(\mathcal{L}_{\mathrm{u}}(\mathrm{X})\) 包含 \(\mathcal{C}(\mathrm{X})\) ，且 \(\mathcal{L}_{\mathrm{u}}^{\infty}(\mathrm{X})\) 包含 \(\mathcal{C}_b(\mathrm{X})\) ；
\item
  从 X 到 \(\mathbf{C}\) 中的每个博雷尔函数（TG, IX, p. 61，定义 5）都属于 \(\mathcal{L}_{\mathrm{u}}(\mathrm{X})\) ；
\item
  对任意序列 \((f_n)_{n\in \mathbf{N}}\) ——\(\mathcal{L}_{\mathrm{u}}(\mathrm{X})\) 中逐点收敛到函数 \(f\) （从 X 到 \(\mathbf{C}\) 中的函数）的序列——我们有 \(f\in \mathcal{L}_{\mathrm{u}}(\mathrm{X})\) 。
\end{enumerate}

断言 \(a\)）与 \(c\)）由定义得出（参见 INT, IV, p. 175, § 5, no. 3，推论 3）。断言 \(d\)）是 INT, IV, p. 179, § 5, no. 5 定理 4 的推论，因为 \(\mathbf{C}\) 的每个博雷尔子集在博雷尔函数下的原像是 X 的博雷尔子集，从而是普遍可测的（INT, V, p. 28, § 3, no. 4）。最后，断言 \(e\)）由叶戈罗夫定理得出（INT, IV, p. 175, § 5, no. 4，定理 2）。

为证明断言 \(b\)），只需注意断言 \(e\)）蕴含，\(a\) 更不用说，普遍可测函数的一致极限是普遍可测的。

X 上可能存在不是博雷尔的普遍可测函数。事实上，若 X 可度量化，X 的苏斯林子集的特征函数是普遍可测的（参见 INT, IV, p. 171, § 5, no. 1，推论 2）；然而在每个不可数的波兰空间中，都存在不是博雷尔的苏斯林子集（“苏斯林定理”；这由 TG, IX, p. 120 习题 8 以及对每个不可数波兰空间 X 存在逆映射为博雷尔的双射 \(\mathrm{X}\to\mathbf{N}^{\mathbf{N}}\) 这一事实得出）。

引理 5. --- 设 \(\mu\) 是 X 上的正测度。设 \(g\) 是 \(\mu\)-可测函数（从 X 到 \(\mathbf{C}\) 中的），并设 \(S\) 是它的 \(\mu\)-本质像。对每个函数 \(f \in \mathcal{L}_{\mathrm{u}}(\mathrm{S})\)，函数 \(h\) （从 X 到 \(\mathbf{C}\) 中的函数），若满足 \(h(x) = 0\) （当 \(g(x) \notin \mathrm{S}\)）与 \(h(x) = f(g(x))\) （当 \(g(x) \in \mathrm{S}\)），则是 \(\mu\)-可测的。

设 \(x \in X\) ，\(U\) 是 x 的相对紧开邻域。集 \(\mathrm{N} = \mathrm{U} - (\mathrm{U} \cap g^{-1}(\mathrm{S}))\) 是 \((\mu|\mathrm{U})\)-零测的。对每个整数 \(n \geqslant 1\)，设 \(V_{n}\) 是满足 \(N \subset V_{n} \subset U\) 且 \(\mu(V_{n} - N) < 1/n\) 的开集（INT, IV, p. 116, § 1, no. 4，命题 19）。

限制 \(\widetilde{g}\) （\(g\) 在 \(\mathrm{U}-\mathrm{V}_n\) 上的限制）是 \(\mu|(\mathrm{U}-\mathrm{V}_n)\)-正常的，且其像含于 S。由于 \(f\) 普遍可测，映射 \(x\mapsto f(g(x))\) （从 \(\mathrm{U}-\mathrm{V}_n\) 到 \(\mathbf{C}\) 中的映射）关于测度 \(\mu|(\mathrm{U}-\mathrm{V}_n)\) 可测（INT, V, p. 71, § 6, no. 2，命题 3）。因此，映射 \(h_n\) （从 X 到 \(\mathbf{C}\) 中的映射）——若 \(h_n(x) = 0\) （当 \(x\notin \mathrm{U}-\mathrm{V}_n\)）且 \(h_n(x) = f(g(x))\) （当 \(x\in \mathrm{U}-\mathrm{V}_n\)）——是 \(\mu\)-可测的（INT, IV, p. 193, § 5, no. 10，命题 16）。

由于 \(h_{n}(x) \to h(x)\) 对 \(\mu\)-几乎所有的 \(x \in U\) 成立，\(h\) 在 U 上的限制是 \(\mu\)-可测的（INT, IV, p. 175, § 5, no. 4，定理 2）。于是可断定 \(h\) 是 \(\mu\)-可测的，依据是局部化原理（INT, IV, p. 171, § 4, no. 2，命题 4）。

注. --- 在该引理的假设下，我们约定滥用记号，用 \(f \circ g\) 表示引理中定义的函数 \(h\) 。

若 \(g = g_{1}\) 在 X 上局部 \(\mu\)-几乎处处成立，则函数 \(g\) 与 \(g_{1}\) 有相同的 \(\mu\)-本质像，且有 \(f \circ g = f \circ g_{1}\) 。

若 \(g\) 是在 X 上 \(\mu\)-几乎处处定义的 \(\mu\)-可测函数，我们也用 \(f \circ g\) 表示函数 \(f \circ g_{1}\) ，其中 \(g_{1}\) 是定义在 X 上且局部 \(\mu\)-几乎处处等于 g 的 \(\mu\)-可测函数。

\subsection{\texorpdfstring{星代数 \(\mathrm{L}^{\infty}(\mathrm{X},\mu)\)}{星代数 \textbackslash mathrm\{L\}\^{}\{\textbackslash infty\}(\textbackslash mathrm\{X\},\textbackslash mu)}}

设 X 是局部紧拓扑空间，\(\mu\) 是 X 上的正测度。我们考虑交换星代数 \(\mathrm{L}^{\infty}(\mathrm{X},\mu)\) （I, p. 103 的例 4）。

引理 6. --- 设 \(g \in \mathcal{L}^{\infty}(\mathrm{X}, \mu)\)。\(g\) 在 \(\mathrm{L}^{\infty}(\mathrm{X}, \mu)\) 中的类的谱等于 \(\mu\)-本质像（\(g\) 的）。

用 \(\widetilde{g}\) 表示 \(g\) 在 \(\mathrm{L}^{\infty}(\mathrm{X},\mu)\) 中的类，用 \(S\) 表示 \(\mu\)-本质像（\(g\) 的）。设 \(z\in \mathbf{C} - \mathbf{S}\)。按定义，存在开邻域 \(U\) （\(z\) 的），使得 \(\mathrm{Y} = g^{-1}(\mathrm{U})\) 局部 \(\mu\)-零测。函数 \(h\) ——由 \(h(x) = (g(x) - z)^{-1}\) （当 \(x\notin \mathrm{Y}\)）与 \(h(x) = 0\) （当 \(x\in \mathrm{Y}\)）定义——属于 \(\mathcal{L}^{\infty}(\mathrm{X},\mu)\)。它的类 \(\widetilde{h}\) 在 \(\mathrm{L}^{\infty}(\mathrm{X},\mu)\) 中满足 \((\widetilde{g} -z)\widetilde{h} = 1\) ，因为 \((g(x) - z)h(x) = 1\) 对每个 \(x\) （只要它不属于局部 \(\mu\)-零测集 Y）成立。因此 \(z\in \mathbf{C} - \mathrm{Sp}(\widetilde{g})\) 。

反之，设 \(z \in \mathbf{C} - \mathrm{Sp}(\widetilde{g})\)。设 \(h \in \mathcal{L}^{\infty}(\mathrm{X},\mu)\) 是这样的函数：它的类是 \(\widetilde{g} - z\) 在 \(\mathrm{L}^{\infty}(\mathrm{X},\mu)\) 中的逆。存在实数 \(\mathrm{M} > 0\) ，使得 \(|h(x)| \leqslant \mathrm{M}\) 局部 \(\mu\)-几乎处处成立，此外我们有 \((g(x) - z)h(x) = 1\) 局部 \(\mu\)-几乎处处成立。设 U 是以 \(z\) 为心、\(\mathrm{M}^{-1}\) 为半径的开球（在 \(\mathbf{C}\) 中）；则 \(g^{-1}(\mathrm{U})\) 含于局部 \(\mu\)-零测集

\[
h ^ { - 1 } ( ] \mathrm{M} , + \infty [ ) \cup \{ x \in \mathrm{X} \mid ( g ( x ) - z ) h ( x ) \neq 1 \} ,
\]

从而 \(z \in \mathbf{C} - \mathrm{S}\) 。引理得证。

命题 4. --- 设 \(g \in \mathcal{L}^{\infty}(\mathrm{X},\mu)\) 。用 \(\widetilde{g}\) 表示 \(g\) 在 \(\mathrm{L}^{\infty}(\mathrm{X},\mu)\) 中的类，用 \(\mathrm{S}\) 表示 \(\widetilde{g}\) 的谱。映射 \(f \mapsto f \circ g\) （p. 184 的注）是从 \(\mathcal{C}(\mathrm{S})\) 到 \(\mathrm{L}^{\infty}(\mathrm{X},\mu)\) 中的与 \(\widetilde{g}\) 相伴的泛函演算映射。

设 \(f \in \mathcal{C}(\mathrm{S})\) 。函数 \(h = f \circ g\) 是 \(\mu\)-可测的（IV, p. 184 的引理 5）且有界。用 \(\widetilde{f \circ g}\) 表示它在 \(\mathrm{L}^{\infty}(\mathrm{X},\mu)\) 中的类。映射 \(f \mapsto \widetilde{f \circ g}\) 是从 \(\mathcal{C}(\mathrm{S})\) 到 \(\mathrm{L}^{\infty}(\mathrm{X},\mu)\) 中的连续有单位元的对合代数同态。由于 \(g(x)\) 属于 \(\mathrm{S}\) （局部 \(\mu\)-几乎处处；IV, p. 181 的引理 1 及引理 6），这个同态把 \(\mathrm{S}\) 上的恒等映射映到类 \(\widetilde{g}\) （函数 \(g\) 的类）。因此它与 \(\widetilde{g}\) 的泛函演算映射相合（I, p. 111 的命题 7）。
\subsection{乘法自同态}

设 X 是局部紧拓扑空间，\(\mu\) 是 X 上的正测度。设 \(p \in [1, +\infty[\) 。

设 \(g \in \mathcal{L}^{\infty}(\mathrm{X},\mu)\) 。用 \(m_{g}\) 表示由 \(f \mapsto g \cdot f\) 所给出的从 \(\mathcal{L}^{p}(\mathrm{X},\mu)\) 到自身的映射。映射 \(m_{g}\) 是线性的且连续的，因为

\[
\mathrm{N} _ {p} (m _ {g} (f)) \leqslant \mathrm{N} _ {\infty} (g) \mathrm{N} _ {p} (f)\tag{1}
\]

对每个函数 \(f \in \mathcal{L}^{p}(\mathrm{X}, \mu)\) 成立。特别地，函数 \(m_{g}(f)\) 是 \(\mu\)-零测的，如果 f 是 \(\mu\)-零测的。于是，映射 \(m_{g}\) 通过过渡到商诱导出 \(\mathrm{L}^{p}(\mathrm{X}, \mu)\) 的一个自同态，我们将它记作 \(\widetilde{m}_{g}\) 。

引理 7. --- 设 \(g \in \mathcal{L}^{\infty}(\mathrm{X},\mu)\) 。自同态 \(\tilde{m}_g\) （\(\mathrm{L}^{p}(\mathrm{X},\mu)\) 的）是单射的，当且仅当由 \(x \in \mathrm{X}\) （满足 \(g(x) = 0\) ）所成的集是局部 \(\mu\)-零测的。

用 A 表示由 \(x \in X\) （满足 \(g(x) = 0\) ）所成的集。

假设 A 是局部 \(\mu\)-零测的。设 \(f \in \mathcal{L}^p(\mathrm{X}, \mu)\) ，\(\widetilde{f}\) 是它在 \(\mathrm{L}^p(\mathrm{X}, \mu)\) 中的类。假设 \(\widetilde{m}_g(\widetilde{f}) = 0\) 。按定义，这等价于函数 \(fg\) 是 \(\mu\)-零测的，于是由 \(x \in \mathrm{X}\) （满足 \(f(x)g(x) \neq 0\) ）所成的集 B 是 \(\mu\)-零测的。函数 \(f\) 在 \(\mathrm{A} \cup \mathrm{B}\) 之外为零，故是 \(\mu\)-零测的。这就蕴含自同态 \(\widetilde{m}_g\) 是单射的。

反之，假设 A 不是局部 \(\mu\)-零测的。设 C 是 X 的紧子集，使得 A∩C 不是 \(\mu\)-零测的，并设 \(\varphi\) 是 A∩C 的特征函数。函数 \(\varphi\) 在 L \(^{p}\) (X, \(\mu\) ) 中的类不为零，但 \(m_{g}(\varphi)\) 的类为零，故 \(m_{g}\) 不是单射的。

由于局部 \(\mu\)-零测函数与 \(\mu\)-零测函数之积是 \(\mu\)-零测的，自同态 \(\widetilde{m}_g\) 只依赖于 \(g\) 在 \(\mathrm{L}^{\infty}(\mathrm{X},\mu)\) 中的类。对 \(\widetilde{g}\in \mathrm{L}^{\infty}(\mathrm{X},\mu)\) ，我们也用 \(\widetilde{m}_{\widetilde{g}}\) 记自同态 \(\widetilde{m}_g\) （\(\mathrm{L}^p (\mathrm{X},\mu)\) 的），这里 \(g\in \mathcal{L}^{\infty}(\mathrm{X},\mu)\) 是以 \(\widetilde{g}\) 为类的任一函数。我们称 \(\widetilde{m}_{\widetilde{g}}\) 为由 \(\widetilde{g}\) 所作的、\(\mathrm{L}^p (\mathrm{X},\mu)\) 中的乘法自同态。

命题 5. --- 设 \(p \in [1, +\infty[\) 。映射 \(m: g \mapsto \widetilde{m}_g\) （从 \(\mathrm{L}^{\infty}(\mathrm{X}, \mu)\) 到 \(\mathcal{L}(\mathrm{L}^{p}(\mathrm{X}, \mu))\) 中的）是含幺巴拿赫代数的一个等距态射。

对 \(a\) ， \(b\) （属于 \(\mathbf{C}\) ）以及 \(g_{1}\) ， \(g_{2}\) （属于 \(\mathrm{L}^{\infty}(\mathrm{X},\mu)\) ），可直接验证 \(\widetilde{m}_{ag_1+bg_2}=a\widetilde{m}_{g_1}+b\widetilde{m}_{g_2}\) ，故映射 \(\widetilde{m}\) 是线性的。此外，我们有 \(\widetilde{m}_{1}=1_{\mathrm{L}^{p}(\mathrm{X},\mu)}\) 与 \(\widetilde{m}_{g_{1}g_{2}}=\widetilde{m}_{g_{1}}\widetilde{m}_{g_{2}}\) ，故映射 \(\widetilde{m}\) 是从 \(\mathrm{L}^{\infty}(\mathrm{X},\mu)\) 到 \(\mathcal{L}(\mathrm{L}^{p}(\mathrm{X},\mu))\) 中的含幺代数态射。

上面的公式 (1) 表明 \(\widetilde{m}\) 的范数 \(\leqslant 1\) 。

设 \(g \in \mathcal{L}^{\infty}(\mathrm{X},\mu)\) ，\(\widetilde{g}\) 是它在 \(\mathrm{L}^{\infty}(\mathrm{X},\mu)\) 中的类。设 \(\varepsilon > 0\) 。由 \(x \in X\) （满足 \(|g(x)| > \| \widetilde{g}\|_{\infty} - \varepsilon\) ）所成的集 Y 不是局部 \(\mu\)-零测的。因此存在 X 的紧子集 C 使得 \(\mu (\mathrm{Y} \cap \mathrm{C}) > 0\) 。设 \(\varphi\) 是 \(\mathrm{Y} \cap \mathrm{C}\) 的特征函数。由于

\[
\begin{array}{c} \| \widetilde {m} _ {\widetilde {g}} (\varphi) \| _ {p} = \Bigl (\int_ {\mathrm{X}} | \varphi g | ^ {p}   d \mu \Bigr) ^ {1 / p} \geqslant \Bigl (\int_ {\mathrm{Y}} | \varphi g | ^ {p}   d \mu \Bigr) ^ {1 / p} \\ \geqslant (\| \widetilde {g} \| _ {\infty} - \varepsilon) \mu (\mathrm{Y} \cap \mathrm{C}) ^ {1 / p} = (\| \widetilde {g} \| _ {\infty} - \varepsilon) \| \varphi \| _ {p}, \end{array}
\]

便得到 \(\|\widetilde{m}_{\widetilde{g}}\|\geqslant\|\widetilde{g}\|_{\infty}-\varepsilon\) 。由于 \(\varepsilon\) 是任意的，可推出 \(\|\widetilde{m}_{\widetilde{g}}\|\geqslant\|\widetilde{g}\|_{\infty}\) ，这就证明了 \(\widetilde{m}\) 是等距的。

引理 8. --- 设 \((g_n)\) 是 \(\mathcal{L}^\infty(\mathrm{X},\mu)\) 中 \(\mu\)-几乎处处逐点收敛的有界序列。设 \(\widetilde{m} \in \mathrm{L}^\infty(\mathrm{X},\mu)\) 是其极限的类。那么 \(\widetilde{m}_{g_n}\) 收敛于 \(\widetilde{m}_{\widetilde{g}}\) ，这里是在赋有逐点收敛拓扑的空间 \(\mathcal{L}(\mathrm{L}^p(\mathrm{X},\mu))\) 中。

设 \(f \in \mathcal{L}^p(\mathrm{X},\mu)\) 。序列 \((g_{n}f)\) 在 \(\mathcal{L}^p(\mathrm{X},\mu)\) 中有界，且 \(\mu\)-几乎处处逐点收敛于 \(gf\) 。由勒贝格定理（INT, IV, p. 137, § 3, no. 7，定理 6），序列 \((g_{n}f)\) 收敛于 \(gf\) （在 \(\mathcal{L}^p(\mathrm{X},\mu)\) 中），断言得证。

现在考虑 p = 2 的情形。

命题 6. --- 映射 m: \(g \mapsto \widetilde{m}_{g}\) 是从星代数 \(\mathrm{L}^{\infty}(\mathrm{X}, \mu)\) 到星代数 \(\mathcal{L}(\mathrm{L}^{2}(\mathrm{X}, \mu))\) 中的含幺等距态射。

特别地，对每个 \(g \in \mathrm{L}^{\infty}(\mathrm{X}, \mu)\) ，乘法自同态 \(\widetilde{m}_{g}\) 是正规的；它是埃尔米特的，当且仅当 g \(\mu\)-几乎处处局部取实值。

由命题 5，映射 \(\tilde{m}\) 是从 \(\mathrm{L}^{\infty}(\mathrm{X},\mu)\) 到 \(\mathcal{L}(\mathrm{L}^{2}(\mathrm{X},\mu))\) 中的含幺巴拿赫代数的单射等距态射。设 \(g\in \mathcal{L}^{\infty}(\mathrm{X},\mu)\) 。对 \(f_{1}\) 与 \(f_{2}\in \mathcal{L}^{2}(\mathrm{X},\mu)\) ，我们有

\[
\langle f _ {1} \mid \widetilde {m} _ {g} (f _ {2}) \rangle = \int_ {\mathrm{X}} \overline {{f _ {1} (x)}} g (x) f _ {2} (x) d \mu (x) = \langle \widetilde {m} _ {\overline {{g}}} (f _ {1}) \mid f _ {2} \rangle ,
\]

由此得出 \(\widetilde{m}_g^* = \widetilde{m}_{\overline{g}}\) ，这就证明了 \(m\) 是对合态射。最后的断言由此得出（参见 I, p. 106，命题 5）。

推论. --- 设 \(g \in \mathrm{L}^{\infty}(\mathrm{X},\mu)\) 。

\begin{enumerate}
\def\labelenumi{\alph{enumi})}
\item
  \(\widetilde{m}_g\) 在 \(\mathcal{L}(\mathrm{L}^2 (\mathrm{X},\mu))\) 中的谱是 \(\mu\)-本质意义下 \(g\) 的像；
\item
  自同态 \(\widetilde{m}_g\) 是正的，当且仅当 \(g\) \(\mu\)-几乎处处局部为正；
\item
  对任意函数 \(f \in \mathcal{C}(\mathrm{Sp}(\widetilde{m}_g))\) ，我们有 \(f(\widetilde{m}_g) = \widetilde{m}_{f \circ g}\) 。
\end{enumerate}

由于 \(\widetilde{m}\) 是单射的，由 I, p. 106 的命题 5，我们有 \(\mathrm{Sp}(\widetilde{m}_{g}) = \mathrm{Sp}(g)\) ；结论由 IV, p. 185 的引理 6、IV, p. 185 的命题 4 以及 I, p. 115 的定义 6 得出。

命题 7. --- 态射 \(\tilde{m}\) （从 \(\mathrm{L}^{\infty}(\mathrm{X},\mu)\) 到 \(\mathcal{L}(\mathrm{L}^{2}(\mathrm{X},\mu))\) 中的）的像是代数 \(\mathcal{L}(\mathrm{L}^{2}(\mathrm{X},\mu))\) 的一个极大交换子代数。

只须证明：若 \(u \in \mathcal{L}(\mathrm{L}^{2}(\mathrm{X}, \mu))\) 是与 \(\widetilde{m}_{g}\) （对每个函数 \(g \in \mathcal{L}^{\infty}(\mathrm{X}, \mu)\) ）都交换的自同态，则 u 属于 \(\widetilde{m}\) 的像。设 u 是这样一个自同态。

我们用 \(\tilde{f}\) 表示 \(\mathrm{L}^2 (\mathrm{X},\mu)\) 中函数 \(f\in \mathcal{L}^2 (\mathrm{X},\mu)\) 的类。用 \(v\) 表示从 \(\mathcal{L}^2 (\mathrm{X},\mu)\) 到 \(\mathrm{L}^2 (\mathrm{X},\mu)\) 中的、由 \(f\mapsto u(\widetilde{f})\) 所定义的线性映射。我们有 \(v\circ m_g = \widetilde{m}_g\circ v\) （对所有 \(g\in \mathcal{L}^{\infty}(\mathrm{X},\mu)\) ）。

对 X 的任意 \(\mu\)-可测子集 Y，我们用 \(\varphi_{\mathrm{Y}}\) 表示它的特征函数；自同态 \(\widetilde{m}_{\varphi_{\mathrm{Y}}}\) 是 \(\mathrm{L}^2 (\mathrm{X},\mu)\) 的一个正交投影算子。

首先假设 X 是紧的，从而常值函数 1 的类属于 \(\mathrm{L}^{2}(\mathrm{X},\mu)\) 。设 g 是 \(\mathcal{L}^{2}(\mathrm{X},\mu)\) 中的函数，它在 \(\mathrm{L}^{2}(\mathrm{X},\mu)\) 中的类为 \(v(1)\) 。

设 \(c\) 是正实数，Y 是由 \(x \in \mathrm{X}\) （满足 \(|g(x)| \geqslant c\) ）所成的集；它是一个 \(\mu\)-可测集。我们在 \(\mathrm{L}^2(\mathrm{X}, \mu)\) 中有等式 \(\widetilde{\varphi_{\mathrm{Y}}} g = \widetilde{m}_{\varphi_{\mathrm{Y}}}(v(1)) = v(m_{\varphi_{\mathrm{Y}}}(1)) = v(\varphi_{\mathrm{Y}})\) 。

又由于 \(|c\varphi_{\mathrm{Y}}| \leqslant |g\varphi_{\mathrm{Y}}|\) ，我们得到

\[
c ^ {2} \mu (\mathrm{Y}) = \int_ {\mathrm{X}} (c \varphi_ {\mathrm{Y}}) ^ {2} d \mu \leqslant \int_ {\mathrm{X}} | \varphi_ {\mathrm{Y}} g | ^ {2} d \mu = \| \widetilde {u} (\varphi_ {\mathrm{Y}}) \| ^ {2} \leqslant \| u \| ^ {2} \mu (\mathrm{Y}).
\]

这个不等式表明 \(\mu(\mathrm{Y}) = 0\) （若 \(c > \|u\|\) ），从而函数 \(g\) \(\mu\)-几乎处处以 \(\|u\|\) 为界。最后，对每个函数 \(f \in \mathcal{C}(\mathrm{X})\) ，我们有

\[
u (\widetilde {f}) = v (m _ {f} (1)) = \widetilde {m} _ {f} (v (1)) = \widetilde {m} _ {f} (\widetilde {g}) = \widetilde {m} _ {g} (\widetilde {f}),
\]

由此得到等式 \(u = \tilde{m}_{g}\) ，特别地 \(\mathrm{N}_{\infty}(g) = \| u\|\) 。

现在考虑一般情形。存在 X 的两两不相交紧子集所成的局部可数族 \((\mathrm{X}_{i})_{i\in\mathrm{I}}\) ，使得 \(Z = X - \bigcup_{i \in I} X_i\) 是局部 \(\mu\)-零测的（INT, IV, p.190, § 5, no. 9，命题 14）。设 \(\mu_i\) 是 \(\mu\) 在 \(X_i\) 上诱导的测度（INT, IV, p. 186, § 5, no. 7，定义 4）。

由 IV, p. 179 的命题 1，对每个 \(i \in I\) ，像 \(\mathrm{E}_i\) （\(\widetilde{m}_{\varphi_{\mathrm{X}_i}}\) 的）与 \(\mathrm{L}^2(\mathrm{X}_i, \mu_i)\) 等同——通过 \(f \mapsto f|\mathrm{X}_i\) 。对每个函数 \(g \in \mathcal{L}^\infty(\mathrm{X}, \mu)\) ，乘法自同态 \(\widetilde{m}_g\) 与 \(\widetilde{m}_{\varphi_{\mathrm{X}_i}}\) 交换，因而使 \(\mathrm{E}_i\) 保持稳定，且 \(\mathrm{E}_i\) 的、由 \(\widetilde{m}_g\) 过渡到子空间而得的自同态与乘以 \(g|\mathrm{X}_i\) 的乘法自同态（\(\mathrm{L}^2(\mathrm{X}_i, \mu_i)\) 上的）重合。

由于 \(u\) 与 \(m_{\varphi_{X_i}}\) 交换，它也使子空间 \(\mathrm{E}_i\) 保持稳定。用 \(u_i\) 表示 \(\mathrm{L}^2 (\mathrm{X}_i,\mu_i)\) 的、由 \(u\) 过渡到子空间而得的自同态。

由于从 \(\mathcal{L}^{\infty}(\mathrm{X},\mu)\) 到 \(\mathcal{L}^{\infty}(\mathrm{X}_{i},\mu_{i})\) 中的、由 \(g\mapsto g|\mathrm{X}_{i}\) 所定义的映射是满射的，假设蕴含 \(u_{i}\) 与 \(\widetilde{m}_{g}\) 交换（对每个函数 \(g\in \mathcal{L}^{\infty}(\mathrm{X}_{i},\mu_{i})\) ）。由证明的第一部分，存在函数 \(g_{i}\in \mathcal{L}^{\infty}(\mathrm{X}_{i},\mu_{i})\) ，使得 \(u_{i} = \widetilde{m}_{g_{i}}\) 且 \(\mathrm{N}_{\infty}(g_{i})\leqslant \| u_{i}\| \leqslant \| u\|\) 。

X 上与 \(g_{i}\) 在 \(X_{i}\) 上重合、在 Z 上为零的函数 g 是有界且 \(\mu\)-可测的（INT, IV, p. 193, § 5, no. 10，命题 16），故 \(g \in \mathcal{L}^{\infty}(X, \mu)\) 。对每个 \(i \in I\) ，u 在 \(E_{i}\) 上的限制与 \(\tilde{m}_{g}\) 的限制重合；因此，由 IV, p. 179 的命题 1, c)，我们有 \(u = m_{g}\) 。

推论. --- 设 u 是希尔伯特空间 \(\mathrm{L}^{2}(\mathrm{X},\mu)\) 的自同态，它与 \(\widetilde{m}_{g}\) （对每个函数 \(g\in\mathcal{K}(\mathrm{X})\) ）交换。那么存在唯一的元素 \(f\in\mathrm{L}^{\infty}(\mathrm{X},\mu)\) ，使得 \(u=\widetilde{m}_{f}\) 。

由命题 7，只须证明 \(u\) 与 \(\widetilde{m}_{g}\) 交换（对所有 \(g\in\mathcal{L}^{\infty}(\mathrm{X},\mu)\) ）。设 \(h_{1}\) 与 \(h_{2}\) 是 \(\mathcal{L}^{2}(\mathrm{X},\mu)\) 中的元素，它们的类 \(\widetilde{h}_{1}\) 与 \(\widetilde{h}_{2}\) 属于 \(\mathrm{L}^{2}(\mathrm{X},\mu)\) 。设 \(k_{1}\) （相应地， \(k_{2}\) ）是 \(\mathcal{L}^{2}(\mathrm{X},\mu)\) 中的函数，其类为 \(u(\widetilde{h}_{1})\) （相应地， \(u^{*}(\widetilde{h}_{2})\) ）。

设 \(h = h_1 \overline{k}_2 - k_1 \overline{h}_2\) ；我们有 \(h \in \mathcal{L}^1(\mathrm{X}, \mu)\) 。定义测度 \(\nu = h \cdot \mu\) （在 \(\mathrm{X}\) 上）；它是有界的。对每个 \(g \in \mathcal{L}^\infty(\mathrm{X}, \mu)\) ，我们有

\[
\begin{array}{c} \langle \widetilde {h} _ {2} | u (\widetilde {m} _ {g} (\widetilde {h} _ {1})) - \widetilde {m} _ {g} (u (\widetilde {h} _ {1})) \rangle = \langle u ^ {*} (\widetilde {h} _ {2}) | \widetilde {m} _ {g} (\widetilde {h} _ {1}) \rangle - \langle \widetilde {h} _ {2} | \widetilde {m} _ {g} (u (\widetilde {h} _ {1})) \rangle \\ = \int_ {\mathrm{X}} g \cdot h _ {1} \cdot \overline {{k}} _ {2} d \mu - \int_ {\mathrm{X}} g \cdot k _ {1} \cdot \overline {{h}} _ {2} d \mu = \nu (g). \end{array}
\]

于是，按假设有 \(\nu(g) = 0\) （对每个函数 \(g \in \mathcal{K}(\mathrm{X})\) ），即 \(\nu = 0\) 。由此可得

\[
\langle \widetilde {h} _ {2} | u (\widetilde {m} _ {g} (\widetilde {h} _ {1})) - \widetilde {m} _ {g} (u (\widetilde {h} _ {1})) \rangle = \nu (g) = 0
\]

对所有 \(g \in \mathcal{L}^{\infty}(\mathrm{X},\mu)\) 成立。由于这对所有 \(h_1\) 与 \(h_2\) （\(\mathcal{L}^2 (\mathrm{X},\mu)\) 中的）都成立，我们有 \(u\circ \widetilde{m}_g = \widetilde{m}_g\circ u\) 。

注. --- 在下文中，我们常径直用 \(m_{g}\) 表示乘以 g 的乘法算子（\(\mathrm{L}^{p}(\mathrm{X},\mu)\) 上的）。

\subsection{谱测度}

在本节中，我们固定一个复希尔伯特空间 E。

回顾：若 A 是交换含幺星代数，我们用 X(A) 表示其幺特征标构成的紧拓扑空间（I, p. 29 的推论 1），并用 \(\mathcal{G}_{\mathrm{A}}\) 表示 A 的盖尔范德变换（I, p. 7 的定义 5），它是从 A 到 \(\mathcal{C}(\mathrm{X}(\mathrm{A}))\) 上的含幺对合代数的等距同构（I, p. 108 的定理 1）。我们将用 \(\mathcal{H}_{\mathrm{A}}\) 表示其逆同构。

引理 9. --- 设 A 是 \(\mathcal{L}(\mathrm{E})\) 的交换含幺星子代数。对 E 中所有 \(x\) 与 \(y\) ，映射 \(\mu\) ——从 \(\mathsf{X}(\mathrm{A})\) 到 \(\mathbf{C}\) 中，由 \(\mu(f) = \langle x | \mathcal{H}_{\mathrm{A}}(f)y \rangle\) （对每个 \(f \in \mathcal{C}(\mathsf{X}(\mathrm{A}))\) ）所定义——是紧空间 \(\mathsf{X}(\mathrm{A})\) 上的有界测度。当 \(x = y\) 时它是正的。其范数 \(\leqslant \|x\|\ \|y\|\) ，且当 \(x = y\) 时取等号。

映射 \(\mu\) 是线性的。对每个函数 \(f\in\mathcal{C}(\mathsf{X}(\mathsf{A}))\) ，我们有

\[
| \mu (f) | \leqslant \| x \| \| y \| \| \mathcal {H} _ {\mathrm{A}} (f) \| = \| x \| \| y \| \| f \|,
\]

因而 \(\mu\) 是 \(\mathsf{X}(\mathsf{A})\) 上范数 \(\leqslant \| x\|\| y\|\) 的有界测度。

假设 \(x = y\) 。对每个正函数 \(f \in \mathcal{C}(\mathsf{X}(\mathsf{A}))\) ，元素 \(\mathcal{H}_{\mathsf{A}}(f)\) 是 A 的正元素（I, p. 115 的例 1），因而是 \(\mathcal{L}(\mathsf{E})\) 的正元素，于是 \(\mu(f) = \langle x | \mathcal{H}_{\mathsf{A}}(f)(x) \rangle \geqslant 0\) （I, p. 138 的命题 8）。这表明 \(\mu\) 是正测度。由于 \(\mu(1) = \|x\|^2\) ，\(\mu\) 的总质量是 \(\|x\|^2\) （INT, III, p. 58, § 1, no. 8，推论 2）。

定义 2. --- 设 A 是 \(\mathcal{L}(\mathrm{E})\) 的交换含幺星子代数。对 E 中所有 \(x\) 与 \(y\) ，线性形式 \(f \mapsto \langle x | \mathcal{H}_{\mathrm{A}}(f)y \rangle\) （\(\mathcal{C}(\mathrm{X}(\mathrm{A}))\) 上的）称为 \((x,y)\) 关于 A 的谱测度。若 \(x = y\) ，我们称它为 \(x\) 关于 A 的谱测度。

注. --- 设 A 是 \(\mathcal{L}(\mathrm{E})\) 的交换含幺星子代数。对 E 中所有 \(x\) 与 \(y\) ，用 \(\mu_{x,y}\) 表示 \((x,y)\) 关于 A 的谱测度。由 \((x,y) \mapsto \mu_{x,y}\) 定义的从 \(\mathrm{E} \times \mathrm{E}\) 到 \(\mathcal{M}^1(\mathrm{X}(\mathrm{A}))\) 中的映射是半双线性的。

设 \(u\) 是 E 的正规自同态，A 是 \(\mathcal{L}(\mathrm{E})\) 的、由 \(u\) 所生成的含幺星子代数。它是交换的。空间 \(\mathsf{X}(\mathrm{A})\) 与 \(\mathrm{Sp}(u)\) 等同——通过映射 \(\chi \mapsto \chi(u)\) （I, p. 109 的引理 10）。于是，谱测度 \(\mu_{x,y}\) ——\((x,y)\) 关于 A 的——与 \(\mathrm{Sp}(u)\) 上的一个测度等同，该测度称为 \((x,y)\) 关于 \(u\) 的谱测度。于是，对每个函数 \(f \in \mathcal{C}(\mathrm{Sp}(u))\) ，我们有

\[
\int_ {\mathrm{Sp} (u)} f d \mu_ {x, y} = \langle x | f (u) y \rangle
\]

依据 I, p. 110 的定义 5。

\subsection{希尔伯特空间上自同态的交换星代数}

定义 3. --- 设 A 是 \(\mathbf{C}\) 上的代数，E 是复拓扑向量空间。设 π 是 A 在 E 中的一个表示。若对 a ∈ A 的元素 π(a)x 所成的集在 E 中是全的，则称 E 的元素 x 是 π 的循环向量。

设 \(u \in \mathcal{L}(\mathrm{E})\) 是 E 的自同态。我们称 \(x \in \mathrm{E}\) 是 \(u\) 的循环向量，如果它是由 \(u\) 在 \(\mathcal{L}(\mathrm{E})\) 中所生成的星子代数的恒等表示的循环向量。

命题 8. --- 设 E 是复希尔伯特空间。设 A 是 \(\mathcal{L}(\mathrm{E})\) 的交换含幺星子代数。设 \(x\) 是 E 的元素。令 \(\mathrm{E}_x = \overline{\mathrm{A} \cdot x}\) ，并用 \(\mu_x\) 表示 \(x\) 关于 A 的谱测度。

存在唯一的等距同构 \(\theta_x\) ——从 \(\mathrm{L}^2 (\mathsf{X}(\mathrm{A}),\mu_x)\) 到 \(\mathrm{E}_x\) 上——使得 \(\theta_{x}(f) = \mathcal{H}_{\mathrm{A}}(f)(x)\) 对每个函数 \(f\in \mathcal{C}(\mathsf{X}(\mathrm{A}))\) 成立。对每个元素 \(a\in \mathrm{A}\) ，空间 \(\mathrm{E}_x\) 在 \(a\) 下稳定，且若用 \(a_{x}\) 表示由 a 过渡到子空间而得的 \(\mathrm{E}_x\) 的自同态，则我们有 \(a_{x}\circ \theta_{x} = \theta_{x}\circ m_{\mathcal{G}_{\mathrm{A}}(a)}\) 。

设 \(\widetilde{\theta}_x\) 是从 \(\mathcal{C}(\mathrm{X(A)})\) 到 \(\mathrm{E}_x\) 中的线性映射，定义为

\[
\widetilde {\theta} _ {x} (f) = \mathcal {H} _ {\mathrm{A}} (f) (x).
\]

对每个函数 \(f \in \mathcal{C}(\mathsf{X}(\mathsf{A}))\) ，我们有

\[
\| \widetilde {\theta} _ {x} (f) \| ^ {2} = \langle \mathcal {H} _ {\mathrm{A}} (f) x | \mathcal {H} _ {\mathrm{A}} (f) x \rangle = \langle x | \mathcal {H} _ {\mathrm{A}} (| f | ^ {2}) x \rangle = \int_ {\mathsf {X} (\mathrm{A})} | f | ^ {2} d \mu_ {x}.
\]

由于 \(\mathcal{C}(\mathrm{X(A)})\) 在 \(\mathcal{L}^2 (\mathrm{X(A)},\mu_x)\) 中稠密，存在唯一的从 \(\mathrm{L}^2 (\mathrm{X(A)},\mu_x)\) 到 \(\mathrm{E}_x\) 中的等距线性映射，它延拓 \(\widetilde{\theta}_x\) ；我们把它记作 \(\theta_{x}\) 。映射 \(\theta_{x}\) 是满射的，因为它的像是闭的（I, p. 107 的引理 8）且包含 \(\mathrm{A}\cdot x\) 。因而它是一个等距同构。

设 \(a \in \mathrm{A}\) 。令 \(g = \mathcal{G}_{\mathrm{A}}(a) \in \mathcal{C}(\mathsf{X}(\mathrm{A}))\) 。空间 \(\mathrm{E}_x\) 在 \(a\) 下稳定。对 \(f \in \mathcal{C}(\mathsf{X}(\mathrm{A}))\) ，我们有

\[
\widetilde {\theta} _ {x} (m _ {g} (f)) = \mathcal {H} _ {\mathrm{A}} (\mathcal {G} _ {\mathrm{A}} (a) f) x = (a \circ \mathcal {H} _ {\mathrm{A}} (f)) x = a (\widetilde {\theta} _ {x} (f)),
\]

由此得出 \(\theta_x \circ m_g = a_x \circ \theta_x\) 。

推论. — 设 E 是复希尔伯特空间。设 A 是 \(\mathcal{L}(\mathrm{E})\) 的交换含幺星子代数，且具有循环向量 \(x\) 。设 \(\mu_x\) 是 \(x\) 关于 A 的谱测度。存在唯一的等距同构

\[
\theta_ {x} \colon \mathrm{L} ^ {2} (\mathsf {X} (\mathrm{A}), \mu_ {x}) \to \mathrm{E}
\]

使得 \(\theta_x(f) = \mathcal{H}_{\mathrm{A}}(f)\) 对每个函数 \(f \in \mathcal{C}(\mathsf{X}(\mathrm{A}))\) 成立。对 A 中每个 a，我们有 \(a \circ \theta_x = \theta_x \circ m_{\mathcal{G}_{\mathrm{A}}(a)}\) 。

事实上，这时我们有 \(E_{x}=E\) ，且 \(a_{x}=a\) （对所有 \(a\in A\) ）。

设 E 是复希尔伯特空间。设 \(x \in E\) ，A 是 \(\mathcal{L}(\mathrm{E})\) 的交换含幺星子代数。用 \(E_{x}\) 表示 E 的闭子空间 \(\overline{A \cdot x}\) 。对每个 \(a \in A\) ，我们有 \(a(\mathrm{E}_{x}) \subset \mathrm{E}_{x}\) 。用 \(A_{x}\) 表示 \(\mathcal{L}(\mathrm{E}_{x})\) 的交换含幺星子代数，它由 A 的元素经过渡到子空间而得的 \(E_{x}\) 的自同态组成。向量 x 是 \(A_{x} \subset \mathcal{L}(\mathrm{E}_{x})\) 的循环向量。

命题 9. --- 存在 E 的子集 C，使得 E 是诸空间 \(E_{x}\) （\(x \in C\)）的希尔伯特直和。若 E 是可数生成的，则集 C 是可数的。

设 \(\mathcal{O}\) 是由 E 的这样的子集 C 组成的集：诸子空间 \(\mathrm{E}_x\) （\(x\in \mathrm{C}\)）两两正交。集 \(\mathcal{O}\) 按包含关系有序，它是有限特征的（E, III, p. 34，定义 2），因为 C 属于 \(\mathcal{O}\) 当且仅当由 C 的两个元素组成的集属于 \(\mathcal{O}\) 。由 E, III, p. 35 的定理 1，存在 \(\mathcal{O}\) 的一个极大元 C。

设 F 是由诸子空间 \(E_{x}\) （\(x \in C\)）生成的 E 的闭子空间。只须证明 \(F^{\circ}\) 归结为 0 即可完成命题的证明。设 y 是 \(F^{\circ}\) 的元素。对每个 \(x \in C\) 及 A 中所有 a 与 b，我们有 \(\langle a(y) | b(x) \rangle = \langle y | a^{*}b(x) \rangle = 0\) ，因为 \(a^{*}b(x)\) 属于 \(E_{x}\) ，从而属于 F。由于元素 \(a(y)\) （相应地， \(b(x)\) ）生成 \(E_{y}\) （相应地， \(E_{x}\) ）的稠密子空间，我们有 \(E_{y} \subset E_{x}^{\circ}\) 。由于 C 在 O 中是极大的，这就意味着 \(E_{y}\) 归结为 0，从而 y = 0。

假设 E 是可数生成的。由于 \(E_{x}\) 对所有非零 \(x \in C\) 非零，任何使 E 成为诸空间 \(E_{x}\) （\(x \in C\)）的希尔伯特直和的集 C 都是可数的。

定理 1. --- 设 A 是 \(\mathcal{L}(\mathrm{E})\) 的交换含幺星子代数。存在局部紧拓扑空间 X、X 上的正测度 \(\mu\) 、等距同构 \(\theta\) ——从 \(\mathrm{L}^2 (\mathrm{X},\mu)\) 到 E 上——以及星代数等距态射 \(\pi\) ——从 A 到 \(\mathcal{C}_b(\mathrm{X})\) 中——使得对每个 \(a\in \mathrm{A}\) ，有

\[
a \circ \theta = \theta \circ m _ {\pi (a)}.
\]

由命题 9，存在 E 的子集 C，使得 E 是诸子空间 \(E_{x}\) （\(x \in C\)）的希尔伯特直和。设 X 为局部紧拓扑空间 \(\mathsf{X}(\mathsf{A}) \times \mathsf{C}\) ，其中 C 赋有离散拓扑。存在 X 上唯一的测度 \(\mu\) ，使得 \(\mu|(\mathsf{X}(\mathsf{A}) \times \{x\})\) 与 x 的谱测度 \(\mu_{x}\) 等同（对每个 \(x \in C\) ；INT, III, p. 65, § 2, n°1，命题 1）；这个测度是正的（参见同前）。于是有希尔伯特直和分解

\[
\mathrm{L} ^ {2} (\mathrm{X}, \mu) = \bigoplus_ {x \in \mathrm{C}} \mathrm{L} ^ {2} (\mathrm{X} (\mathrm{A}), \mu_ {x})
\]

（IV, p. 179 的命题 1, c) 应用于诸集 \(\mathsf{X}(\mathsf{A}) \times \{x\}\) （\(x \in \mathrm{C}\)））。因此存在唯一的等距 \(\theta: \mathrm{L}^2(\mathrm{X}, \mu) \to \mathrm{E}\) ，它与 \(\theta_x: \mathrm{L}^2(\mathsf{X}(\mathsf{A}), \mu_x) \to \mathrm{E}_x\) 在 \(\mathrm{L}^2(\mathsf{X}(\mathsf{A}), \mu_x)\) 上重合（对所有 \(x \in \mathrm{C}\) ）。

设 \(p \colon \mathrm{X} \to \mathrm{X}(\mathrm{A})\) 是典范投影；它是满射的，因而线性映射 \(p^* \colon \mathcal{C}_b(\mathrm{X}(\mathrm{A})) \to \mathcal{C}(\mathrm{X})\) ——由 \(f \mapsto f \circ p\) 所定义——是单射的。令 \(\pi = p^* \circ \mathcal{G}_{\mathrm{A}}\) ；这是从 A 到 \(\mathcal{C}_b(\mathrm{X})\) 中的星代数的单射态射；因而它是等距的（I, p. 112 的命题 9）。

设 \(a \in \mathrm{A}\) 。对每个 \(x \in \mathrm{C}\) ，我们有 \(a_x \circ \theta_x = \theta_x \circ m_{\mathcal{G}_A(a)}\) 。于是连续线性映射 \(a \circ \theta\) 与 \(\theta \circ m_{\pi(a)}\) 在诸子空间 \(\mathrm{L}^2(\mathrm{X}(\mathrm{A}), \mu_x)\) 上重合，因而相等。

推论（谱定理）. --- 设 E 是复希尔伯特空间。设 \(u \in \mathcal{L}(\mathrm{E})\) 是 E 的正规自同态。存在局部紧拓扑空间 X、X 上的正测度 \(\mu\) 、等距同构 \(\theta\) ——从 \(\mathrm{L}^{2}(\mathrm{X},\mu)\) 到 E 上——以及 X 上的连续有界函数 g，使得 \(u = \theta \circ m_{g} \circ \theta^{-1}\) 。

若 u 有循环向量 x，则可取 \(X = \mathrm{Sp}(u)\) ，并取 g 为 X 的恒等函数。

只须对 \(\mathcal{L}(\mathrm{E})\) 的、由 \(u\) 所生成的星子代数 A 应用定理 1（相应地，命题 8 的推论），并令 \(g = \pi(u)\) 。

这个命题把关于希尔伯特空间单个正规自同态的每个问题化归为关于乘法算子的类似问题，这往往使其研究得到简化（例如参见 IV, p. 325 的习题 19）。

\subsection{函数演算的连续性}

在本节中，我们考虑函数演算关于两个变量的连续性质。

对星代数 A 及 A 的正规元素 \(a\) ，以及对函数 \(f\) ——在 \(\mathbf{C}\) 的包含 \(\mathrm{Sp}_{\mathrm{A}}(a)\) 的子集上连续——我们用 \(f(a)\) 表示把 \(a\) 的连续函数演算应用于 \(f\) 在 \(a\) 的谱上的限制所得到的元素。

命题 10. --- 设 A 是含幺星代数。设 U 是 \(\mathbf{C}\) 中的相对紧开集。用 \(\Omega_{\mathrm{n}}\) 表示 A 中满足 \(\mathrm{Sp}_{\mathrm{A}}(a)\) 含于 U 的正规元素所成的集。映射 \((f,a)\mapsto f(a)\) ——从 \(\mathcal{C}(\overline{\mathrm{U}})\times \Omega_{\mathrm{n}}\) 到 A 中——是连续的。

用 \(q\) 表示映射 \((f, a) \mapsto f(a)\) ——从 \(\mathcal{C}(\overline{\mathrm{U}}) \times \Omega_{\mathrm{n}}\) 到 A 中。函数演算映射 \(f \mapsto f(a)\) （\(a \in \Omega_{\mathrm{n}}\)）所成的族在 \(\mathcal{L}(\mathcal{C}(\overline{\mathrm{U}}); \mathrm{A})\) 中是等度连续的，因为其中每个都是范数 \(\leqslant 1\) 的连续线性映射。因此，为证明断言，只须验证：对每个 \(f \in \mathcal{C}(\overline{\mathrm{U}})\) ，从 \(\Omega_{\mathrm{n}}\) 到 A 中的、由 \(a \mapsto f(a)\) 所定义的映射是连续的（TG, X, p. 13，推论 3）。

设 \(\mathcal{A}\) 是由 \(f\in\mathcal{C}(\overline{\mathrm{U}})\) （使映射 \(a\mapsto f(a)\) ——从 \(\Omega_{\mathrm{n}}\) 到 A 中——连续）组成的集；要证明的是 \(\mathcal{A}=\mathcal{C}(\overline{\mathrm{U}})\) 。

集 \(\mathcal{A}\) 是 \(\mathcal{C}(\overline{\mathrm{U}})\) 的含幺对合子代数。它包含 \(\overline{\mathrm{U}}\) 上的恒等函数，因而分离点。于是，它在 \(\mathcal{C}(\overline{\mathrm{U}})\) 中稠密（TG, X, p. 39，命题 7）。我们来证明它是闭的。

设 \((f_{\mathrm{n}})\) 是 \(\mathcal{A}\) 中收敛于 \(f \in \mathcal{C}(\overline{\mathrm{U}})\) 的序列。设 \(\varepsilon > 0\) ，并选取 \(n \in \mathbf{N}\) 使得 \(\|f - f_{n}\|_{\infty} \leqslant \varepsilon/4\) 。对每个 \((a_{1}, a_{2}) \in \Omega_{\mathrm{n}}^{2}\) ，我们有

\[
\begin{array}{c} \| f (a _ {1}) - f (a _ {2}) \| \leqslant 2 \| f - f _ {n} \| _ {\infty} + \| f _ {n} (a _ {1}) - f _ {n} (a _ {2}) \| \\ \leqslant \frac {\varepsilon}{2} + \| f _ {n} (a _ {1}) - f _ {n} (a _ {2}) \|. \end{array}
\]

由于 \(f_n\) 属于 \(\mathcal{A}\) ，存在 \(a_1\) 在 \(\Omega_n\) 中的邻域 V，使得 \(\| f_n(a_1) - f_n(a_2) \| \leqslant \varepsilon / 2\) 对所有 \(a_2 \in V\) 成立。于是，我们有 \(\| f(a_1) - f(a_2) \| \leqslant \varepsilon\) （对所有 \(a_2 \in V\) ）。因而映射 \(a \mapsto f(a)\) 在 \(a_1\) 处连续；断言得证。

推论 1. --- 设 A 是含幺星代数，\(A_{n}\) 是 A 的正规元素所成的集。给空间 \(\mathcal{C}(\mathbf{C})\) 赋以紧收敛拓扑。映射 \((f,a)\mapsto f(a)\) ——从 \(\mathcal{C}(\mathbf{C})\times\mathrm{A}_{\mathrm{n}}\) 到 A 中——是连续的。

设 \((f_{0}, a_{0}) \in \mathcal{C}(\mathbf{C}) \times \mathrm{A}_{\mathrm{n}}\) 。设 U 是 \(a_{0}\) 的谱的相对紧邻域。设 V 是 A 中满足 \(\mathrm{Sp}_{\mathrm{A}}(a) \subset \mathrm{U}\) 的正规元素 a 所成的集；它是 \(a_{0}\) 在 \(A_{n}\) 中的开邻域（I, p. 76，命题 10）。对每个 \(a \in V\) 及每个函数 \(f \in \mathcal{C}(\mathbf{C})\) ，我们有 \(f(a) = (f|\overline{\mathrm{U}})(a)\) 。由于映射 \(f \mapsto \|f|\overline{U}\|_{\infty}\) 是赋以紧收敛拓扑的空间 \(\mathcal{C}(\mathbf{C})\) 上的连续半范数，映射 \((f, a) \mapsto f(a)\) 在 \((f_{0}, a_{0})\) 处的连续性由命题 10 得出。

推论 2. --- 设 E 是复希尔伯特空间，\(\mathcal{L}(\mathrm{E})_{\mathrm{n}}\) 是 E 的正规自同态所成的集。给空间 \(\mathcal{C}(\mathbf{C})\) 赋以紧收敛拓扑。从 \(\mathcal{C}(\mathbf{C}) \times \mathcal{L}(\mathrm{E})_{\mathrm{n}}\) 到 \(\mathcal{L}(\mathrm{E})\) 中的、由 \((f, u) \mapsto f(u)\) 所定义的映射是连续的。

\section{广义函数与缓增广义函数}

在本节中，n 表示一个自然数。我们用 \(\mu\) 表示 \(\mathbf{R}^{n}\) 上的勒贝格测度，以及它在 \(\mathbf{R}^{n}\) 的每个局部紧子集上的限制。

我们用

\[
x \cdot y = \sum_ {i = 1} ^ {n} x _ {i} y _ {i}
\]

表示欧几里得空间 \(\mathbf{R}^{n}\) 上的数量积；欧几里得范数记作 \(x \mapsto \|x\|\) （TG, VI, p. 7）。我们回顾，群 \(\mathbf{R}^{n}\) 通过映射 \((x, y) \mapsto \exp(2i\pi x \cdot y)\) 与自身对偶，且勒贝格测度的对偶测度这时与勒贝格测度等同（II, p. 236 的推论 3）。我们用 F （相应地， \(\overline{F}\) ）表示 \(\mathbf{R}^{n}\) 上的傅里叶变换（相应地，傅里叶逆变换）（参见 II, p. 237 的 no. 9）。

对每个 \(\alpha = (\alpha_{i})_{1\leqslant i\leqslant n}\in \mathbf{N}^{n}\) ，我们用 \(X^{\alpha}\) 表示从 \(\mathbf{R}^n\) 到 \(\mathbf{R}\) 中的、由 \(x = (x_{i})_{1\leqslant i\leqslant n}\mapsto x^{\alpha} = \prod_{i = 1}^{n}x_{i}^{\alpha_{i}}\) 所定义的函数。

设 \(\alpha\) 与 \(\beta\) 是 \(\mathbf{N}^n\) 的元素。我们记

\[
| \alpha | = \sum_ {i = 1} ^ {n} \alpha_ {i}, \qquad \binom{\alpha}{\beta} = \prod_ {i = 1} ^ {n} \binom{\alpha_ {i}}{\beta_ {i}}.
\]

引理 1. --- 设 U 是 \(\mathbf{R}^{n}\) 的开子集。

\begin{enumerate}
\def\labelenumi{\alph{enumi})}
\item
  设 K 是 U 的紧子集，V 是 K 在 U 中的开邻域。存在 \(\varphi \in \mathcal{C}^{\infty}(\mathrm{U})\) ，其紧支撑含于 V，使得 \(0 \leqslant \varphi \leqslant 1\) ，且 \(\varphi(x) = 1\) 对所有 \(x \in \mathrm{K}\) 成立；
\item
  对 U 的每个局部有限开覆盖，存在一个由 \(C^{\infty}\) 类函数组成、从属于它的单位分解。
\end{enumerate}

这由 VAR, R1, p. 40, 5.3.6 得出。

命题 1. --- 设 U 是 \(\mathbf{R}^{n}\) 的开子集，\(k\in\mathbf{N}\) 。设 E、F、G 是拓扑向量空间，b: E×F→G 是连续双线性映射。设 f 与 g 分别是 U 到 E 与 F 中的 \(\mathbf{C}^{k}\) 类函数。那么映射 h: x↦b(f(x),g(x)) 在 U 上是 \(\mathbf{C}^{k}\) 类的；此外，对每个 \(\alpha\in\mathbf{N}^{n}\) （满足 \(|\alpha|\leqslant k\) ）以及

对每个 \(x \in \mathrm{U}\) ，我们有

\[
\partial^{\alpha}h(x) = \sum_{\substack{\beta \in \mathbf{N}^{n}\\ \beta \leqslant \alpha}}\binom {\alpha}{\beta}b\big(\partial^{\beta}f(x),\partial^{\alpha -\beta}g(x)\big).
\]

映射 h 是映射 \((f,g)\) : \(U \rightarrow E \times F\) （它是 \(C^{k}\) 类的）与映射 b: \(E \times F \rightarrow G\) （它是 \(C^{\infty}\) 类的）的复合。因而它是 \(C^{k}\) 类的。其偏导数的表达式由莱布尼茨公式得出（FVR, I, p. 28，命题 2，参见 A, III, p. 122，推论）。

我们回顾，分离的局部凸拓扑向量空间若为桶型空间且每个有界子集都相对紧，则称为蒙泰尔空间（EVT, IV, p. 18，定义 4）。

我们还回顾，若 E 与 F 是局部凸空间，且 E 是有界型的（EVT, III, p. 12，定义 1），则线性映射 \(u\colon\mathrm{E}\to\mathrm{F}\) 连续当且仅当 E 的每个有界子集在 \(u\) 下的像在 F 中有界（EVT, III, p. 11，命题 1, (iiibis)；当 F 是可赋范时，一般情形可形式地由此推出，参见 EVT, II, p. 7，命题 5, b)）。

\subsection{积分号下求导}

我们固定一个局部紧拓扑空间 X 和 X 上的一个测度 \(\nu\) 。设 E 是巴拿赫空间。

命题 2. --- 设 I 是 \(\mathbf{R}\) 的区间，f 是从 X × I 到 E 中的映射，满足

\begin{enumerate}
\def\labelenumi{(\roman{enumi})}
\item
  对每个 \(t \in I\) ，从 X 到 E 中的映射 \(x \mapsto f(x, t)\) 是 \(\nu\)-可积的；
\item
  对每个 \(x \in X\) ，从 I 到 E 中的映射 \(t \mapsto f(x, t)\) 有导数，记作 \(t \mapsto f'(x, t)\) ；
\item
  存在 X 上的正的 \(\nu\)-可积函数 g，使得 \(\|f'(x,t)\| \leqslant g(x)\) 对所有 \((x,t) \in \mathrm{X} \times \mathrm{I}\) 成立。
\end{enumerate}

那么由

\[
\mathrm{F} (t) = \int_ {\mathrm{X}} f (x, t) d \nu (x)
\]

定义的从 I 到 E 中的映射 F 在 I 中可微，且对每个 \(t \in \mathrm{I}\) ，我们有

\[
\mathrm{F} ^ {\prime} (t) = \int_ {\mathrm{X}} f ^ {\prime} (x, t) d \nu (x).
\]

设 \(t_0 \in \mathrm{I}\) ，\(\mathrm{J}\) 是 \(\mathbf{R}\) 的区间，含于 \(\mathrm{I}\) ，且是 \(t_0\) 在 \(\mathrm{I}\) 中的邻域。设 \(h: \mathrm{X} \times \mathrm{J} \to \mathrm{E}\) 是由 \((x, t) \mapsto (f(x, t) - f(x, t_0)) / (t - t_0)\) （对 \((x, t) \in \mathrm{X} \times (\mathrm{J} - \{t_0\})\) ）以及 \((x, t_0) \mapsto f'(x, t_0)\) （对 \(x \in \mathrm{X}\) ）所定义的映射。设 \(x \in \mathrm{X}\) 。我们有 \(\|h(x, t_0)\| \leqslant g(x)\) ，且对所有 \(t \neq t_0\) ，可得 \(\|h(x, t)\| \leqslant g(x)\) 。依据 FVR, I, p. 23 的定理 2。由导数的定义，命题由应用于映射 \(h\) 的 INT, IV, p. 144, § 4, no. 3 的推论 1 得出。

我们采用 VAR, R1, 2.4, p. 19 中关于偏导数的记号。

推论 1. --- 设 U 是 \(\mathbf{R}^{n}\) 的开子集。设 \(k \in \mathbf{N}\) ，\(f\) 是从 \(\mathrm{X} \times \mathrm{U}\) 到 E 中的映射，满足下列条件：

\begin{enumerate}
\def\labelenumi{(\roman{enumi})}
\item
  对每个 \(t \in U\) ，从 X 到 E 中的映射 \(x \mapsto f(x, t)\) 是 \(\nu\)-可积的；
\item
  对每个 \(x \in X\) ，从 U 到 E 中的映射 \(t \mapsto f(x, t)\) 是 \(C^{k}\) 类的，其偏导数（记作 \(\partial^{\alpha}f(x, t)\) ，对每个 \(\alpha \in \mathbf{N}^{n}\) 满足 \(|\alpha| \leqslant k\) ）；
\item
  存在 \(\nu\)-可积的函数 \(g\) （X 上的），使得对每个 \(\alpha \in \mathbf{N}^n\) （满足 \(|\alpha| \leqslant k\) ）以及每个 \((x,t) \in \mathrm{X} \times \mathrm{U}\) ，有 \(\|\partial^\alpha f(x,t)\| \leqslant g(x)\) 。
\end{enumerate}

那么由

\[
\mathrm{F} (t) = \int_ {\mathrm{X}} f (x, t) d \nu (x)
\]

定义的从 U 到 E 中的映射 F 在 U 中是 \(C^{k}\) 类的，且对每个 \(\alpha\in \mathbf{N}^{n}\) （满足 \(|\alpha|\leqslant k\) ）及每个 \(t\in U\) ，我们有

\[
\partial^ {\alpha} \mathrm{F} (t) = \int_ {\mathrm{X}} \partial^ {\alpha} f (x, t) d \nu (x).
\]

这由命题对 k 作归纳得出。

推论 2. --- 设 k 是自然数。设 \(f \in \mathcal{L}^1(\mathbf{R}^n, \mu)\) ，\(g \in \mathcal{C}^k(\mathbf{R}^n)\) 。假设对每个 \(\alpha \in \mathbf{N}^n\) （满足 \(|\alpha| \leqslant k\) ），函数 \(\partial^\alpha g\) 都有界。那么卷积 \(f * g\) 属于 \(\mathcal{C}^k(\mathbf{R}^n)\) ，且对每个 \(\alpha\) （满足 \(|\alpha| \leqslant k\) ），有 \(\partial^\alpha(f * g) = f * \partial^\alpha g\) 。

我们可以对赋勒贝格测度的空间 \(X = \mathbf{R}^{n}\) 以及由 \((x, t) \mapsto f(x)g(t - x)\) 定义的从

\(\mathbf{R}^{n} \times \mathbf{R}^{n}\) 到 \(\mathbf{C}\) 中的映射 h 应用推论 1；事实上，对每个 \(\alpha \in \mathbf{N}^{n}\) （满足 \(|\alpha| \leqslant k\) ），我们有不等式

\[
| \partial^ {\alpha} h (x, t) | \leqslant \Bigl (\sup _ {| \beta | \leqslant k} \sup _ {y \in \mathbf {R} ^ {n}} | \partial^ {\beta} g (y) | \Bigr) | f (x) |
\]

其右端是 \(R^{n}\) 上的可积函数。

\subsection{\texorpdfstring{\(\mathbf{R}^{n}\) 与 \(\mathbf{Z}^{n}\) 中的可积性判据}{\textbackslash mathbf\{R\}\^{}\{n\} 与 \textbackslash mathbf\{Z\}\^{}\{n\} 中的可积性判据}}

命题 3. --- 设 N 是 \(\mathbf{R}^{n}\) 上的一个范数，\(r\) 是实数，\(p \in [1, +\infty[\) 。

\begin{enumerate}
\def\labelenumi{\alph{enumi})}
\tightlist
\item
  函数 \((1 + \mathrm{N})^{r}\) 属于 \(\mathcal{L}^{p}(\mathbf{R}^{n}, \mu)\) 当且仅当 \(rp < -n\)；
\end{enumerate}

a') 在 \(\mathbf{Z}^{n}\) 上，函数 \((1+N)^{r}\) 的限制属于 \(\ell^{p}(\mathbf{Z}^{n})\) 当且仅当 \(rp < -n\)；

\begin{enumerate}
\def\labelenumi{\alph{enumi})}
\setcounter{enumi}{1}
\tightlist
\item
  设 V 是 \(\mathbf{R}^{n}\) 中 0 的有界可测邻域。函数 \(\mathbf{N}^{r}\) 属于 \(\mathcal{L}^{p}(\mathbf{R}^{n}-\mathbf{V},\mu)\) 当且仅当 \(rp<-n\) 。
\end{enumerate}

由 \(\mathbf{R}^{n}\) 上范数的等价性（EVT I, p. 14，定理 2 及 TG, IX, p. 32，命题 8），可以假设范数 N 由 \(\mathrm{N}(x)=\sup|x_{i}|\) 给出，其中 \(x=(x_{i})\in\mathbf{R}^{n}\) 。用 B 表示 \(\mathbf{R}^{n}\) 对此范数的单位球。

设 V 是 \(\mathbf{R}^n\) 中 0 的有界可测邻域。函数 \(N^r\) 在 \((\mathrm{B - V})\cup (\mathrm{V - B})\) 上连续且有界，这表明断言 \(b)\) 成立当且仅当 \(V = B\) 时它成立，以下我们就这样假设。函数 \((1 + N)^r\) 在 B 上连续且有界，且满足不等式 \(N^r\leqslant (1 + N)^r\leqslant 2^r N^r\) （在 \(\mathbf{R}^n -\mathbf{B}\) 上），由此可看出断言 \(a)\) 与 \(b)\) 等价。

我们来证明 V = B 时的 b)。可以假设 p = 1。\(r > 0\) 的情形是初等的，故设 \(r \leqslant 0\) 。对每个整数 \(j \geqslant 1\) ，可积集

\[
\mathrm{C} _ {j} = \{x \in \mathbf {R} ^ {n} \mid 2 ^ {j - 1} \leqslant \mathrm{N} (x) <   2 ^ {j} \}
\]

的测度为 \(2^{nj}(2^n - 1)\) 。诸集合 \((\mathrm{C}_j)_{j \geqslant 1}\) 构成 \(\mathbf{R}^n - \mathbf{B}\) 的一个划分。依据 INT, V, p. 4, § 1, no. 1，推论，我们因此有

\[
\begin{array}{c} \int_ {\mathbf {R} ^ {n} - \mathrm{B}} ^ {*} \mathrm{N} ^ {r} d \mu = \sum_ {j \geqslant 1} \int_ {\mathrm{C} _ {j}} ^ {*} \mathrm{N} ^ {r} d \mu \\ \leqslant \sum_ {j \geqslant 1} 2 ^ {n + n j} 2 ^ {r (j - 1)} = 2 ^ {n - r} \sum_ {j \geqslant 1} 2 ^ {(n + r) j} \end{array}
\]

它在 \(r < -n\) 时有限。另一方面，我们有（同前）

\[
\int_ {\mathbf {R} ^ {n} - \mathrm{B}} ^ {*} \mathrm{N} ^ {r} d \mu \geqslant \sum_ {j \geqslant 1} 2 ^ {r j} 2 ^ {n j} (2 ^ {n} - 1) = (2 ^ {n} - 1) \sum_ {j \geqslant 1} 2 ^ {(r + n) j}
\]

它在 \(r \geqslant -n\) 时无穷。

我们类似地证明 \(a'\) )，所考虑的集合 \(\mathrm{C}_{j} \cap \mathbf{Z}^{n}\) 覆盖 \(\mathbf{Z}^{n} - \{0\}\) 并满足

\[
\operatorname{Card} \left(\mathrm{C} _ {j} \cap \mathbf {Z} ^ {n}\right) = (2 ^ {j + 1} - 1) ^ {n} - (2 ^ {j} - 1) ^ {n},
\]

该数属于区间 \([(1-2^{-n})(2^{j+1}-1)^{n},2^{n(j+1)}]\) 。

\subsection{测试函数}

在本章中，U 表示 \(\mathbf{R}^{n}\) 的开子集。对 \(p \in [1, +\infty]\) ，我们将写 \(\mathcal{L}^{p}(U)\) 与 \(L^{p}(U)\) ，而不写 \(\mathcal{L}^{p}(U, \mu)\) 与 \(L^{p}(U, \mu)\) 。我们把属于 \(\mathcal{L}^{p}(U)\) 的连续函数与其在 \(L^{p}(U)\) 中的像等同起来。

给 U 上的复值无穷次可微函数空间 \(\mathcal{C}^{\infty}(\mathrm{U})\) 赋以由半范数族 \(p_{\alpha,K}\) ——对 \(\alpha \in \mathbf{N}^{n}\) 和 \(K \subset U\) 定义为

\[
p _ {\alpha , \mathrm{K}} (\varphi) = \sup _ {x \in \mathrm{K}} | \partial^ {\alpha} \varphi (x) |.
\]

该空间是完备的（参见 EVT, III, p. 9，例 b））。

对 U 的每个紧子集 K，我们用 \(\mathcal{C}_{\mathrm{K}}^{\infty}(\mathrm{U})\) 表示 \(\mathcal{C}^{\infty}(\mathrm{U})\) 中由支集含于 K 的函数构成的子空间。空间 \(\mathcal{C}_{\mathrm{K}}^{\infty}(\mathrm{U})\) 赋以由 \(\mathcal{C}^{\infty}(\mathrm{U})\) 的拓扑诱导的拓扑。它是一个弗雷歇空间，定义其拓扑的可数半范数族即族 \((p_{\alpha,\mathrm{K}})_{\alpha\in\mathbf{N}^{n}}\) 。特别地，它是有界型空间（EVT, III, p. 12，命题 2）。

我们用 \(\mathcal{D}(\mathrm{U})\) 表示向量空间 \(\mathcal{K}(\mathrm{U})\cap \mathcal{C}^{\infty}(\mathrm{U})\) ——由 U 中具有紧支集的 \(\mathrm{C}^{\infty}\) 类函数所构成。我们称 \(\mathcal{D}(\mathrm{U})\) 为 U 上的测试函数空间。空间 \(\mathcal{D}(\mathrm{U})\) 是诸空间 \(\mathcal{C}_{\mathrm{K}}^{\infty}(\mathrm{U})\) 的归纳极限，并赋以相应的局部凸归纳极限拓扑（EVT, II, p. 31）。

该空间在 EVT, III, p. 9 中记作 \(\mathcal{C}_{\circ}^{\infty}(U)\) 。

设 \((\mathrm{K}_{m})\) 是 U 的紧子集的递增序列，其内部构成 U 的覆盖。那么空间 \(\mathcal{D}(\mathrm{U})\) 是诸空间 \(\mathcal{C}_{\mathrm{K}_{m}}^{\infty}(\mathrm{U})\) 的严格归纳极限（EVT, III, p. 9）。因此它是完备空间（EVT, II, p. 35，命题 9）。\(\mathcal{D}(\mathrm{U})\) 的每个有界子集都含于某个子空间 \(\mathcal{C}_{\mathrm{K}_{m}}^{\infty}(\mathrm{U})\) 之中（EVT, III, p. 5，命题 6）。这就是说，\(\mathrm{B}\subset\mathcal{D}(\mathrm{U})\) 有界当且仅当存在 U 的紧子集 K 以及族 \((\mathrm{M}_{\alpha})_{\alpha\in\mathbf{N}^{n}}\) （\(R_{+}\) 中的），使得 B 含于诸函数 \(\varphi\in\mathcal{C}_{\mathrm{K}}^{\infty}(\mathrm{U})\) ——满足

\[
p _ {\alpha , \mathrm{K}} (\varphi) \leqslant \mathrm{M} _ {\alpha}
\]

——所成之集（对所有 \(\alpha\in\mathbf{N}^{n}\) ）。

空间 \(\mathcal{D}(\mathrm{U})\) 是有界型空间（EVT, III, p. 12，例 3）和蒙泰尔空间（EVT, IV, p. 18，例 4）。

设 V 是含于 U 的 \(\mathbf{R}^{n}\) 的开子集。若 \(K \subset V\) 是紧的，则函数到 V 上的限制定义了从 \(\mathcal{C}_{\mathrm{K}}^{\infty}(\mathrm{U})\) 到 \(\mathcal{C}_{\mathrm{K}}^{\infty}(\mathrm{V})\) 上的连续满线性映射。定义在 V 上的函数用零延拓，诱导一个从 \(\mathcal{D}(\mathrm{V})\) 到 \(\mathcal{D}(\mathrm{U})\) 中的单连续线性映射，称为典范映射。

注. --- 我们类似地定义 U 上的实值测试函数空间 \(\mathcal{D}_{\mathbf{R}}(\mathrm{U})\) 。使得 \(z \otimes \varphi \mapsto z\varphi\) （对所有 \(z \in \mathbf{C}\) 及每个 \(\varphi \in \mathcal{D}_{\mathbf{R}}(\mathrm{U})\) ）的线性映射是 \(\mathbf{C} \otimes \mathcal{D}_{\mathbf{R}}(\mathrm{U})\) 到 \(\mathcal{D}(\mathrm{U})\) 上的同构。

引理 2. --- 设 \(\mathcal{U}\) 是 U 的开覆盖。存在比 \(\mathcal{U}\) 更细的、由相对紧集合组成的 U 的局部有限开覆盖。

由于 U 是局部紧的，存在 \(\mathcal{V}\) ——U 的覆盖，比 \(\mathcal{U}\) 更细，由在 U 中相对紧的开集组成。由于 U 是仿紧的（TG, IX, p. 51，定理 4），存在局部有限开覆盖 \(\mathcal{W}\) ，比 \(\mathcal{V}\) 更细。于是 \(\mathcal{W}\) 比 \(\mathcal{U}\) 更细，且由相对紧的开集组成。

引理 3. --- 设 \(f \in \mathcal{K}(\mathbf{R}^n)\) ，V 是 f 的支集 K 的开邻域。存在整数 \(m_0 \geqslant 1\) ，使得对每个 \(x \in K\) ，集合 V 包含以 x 为中心、半径为 \(m_0^{-1}\) 的球。对每个整数 \(m > m_0\) ，设 \(V_m\) 是以 0 为中心、半径为 \(m^{-1}\) 的闭球（在 \(\mathbf{R}^{n}\) 中），\(\varphi_{m}\) 是支集含于 \(V_{m}\) 且积分为 1 的正测试函数。令 \(f_{m} = \varphi_{m} * f\) 。

\begin{enumerate}
\def\labelenumi{\alph{enumi})}
\item
  有 \(f_{m}\in \mathcal{D}(\mathbf{R}^{n})\) ，且 \(f_{m}\) 的支集含于 V；
\item
  设 \(p \in [1, +\infty]\) 。序列 \((f_m)_{m > m_0}\) 收敛于 \(f\) （在 \(L^p(\mathbf{R}^n)\) 中）。
\end{enumerate}

依据 TG, IX, p. 14，注，我们有 \(d(\mathrm{K}, \mathbf{R}^{n}-\mathrm{V}) > 0\) 。设 \(m_{0} \geqslant 1\) 是满足 \(m_{0}^{-1} < d(\mathrm{K}, \mathbf{R}^{n}-\mathrm{V})\) 的整数；它满足所需的条件。对所有 \(m > m_{0}\) ，函数 \(f_{m}\) 连续（INT, VIII, p. 166, § 4, n° 5，命题 11），且其支集含于 \(\mathrm{K} + \mathrm{V}_{m}\) （INT, VIII, p. 126, § 1, n° 4，命题 5），从而含于 V。依据 IV, p. 198 的推论 2，我们有 \(f_{m} \in \mathcal{D}(\mathbf{R}^{n})\) 。若 \(p \neq +\infty\) ，则序列 ( \(f_{m}\) ) 收敛于 \(f\) （在 \(\mathrm{L}^{p}(\mathbf{R}^{n})\) 中；INT, VIII, p. 172, § 4, no. 7，命题 20）。\(^{(1)}\)

设 \(p = +\infty\) 。设 \(\varepsilon > 0\) 。函数 \(f\) 在 \(\mathbf{R}^n\) 上一致连续；因此存在整数 \(m_1 > m_0\) ，使得对每个 \(m \geqslant m_1\) 以及所有 \(x \in \mathbf{R}^n\) 和 \(y \in \mathrm{V}_m\) ，有 \(|f(x) - f(x - y)| < \varepsilon\) 。

对每个 \(m \geqslant m_{1}\) ，函数 \(\varphi_{m}\) 在 \(\mathrm{V}_{m}\) 外为零、是正的且积分为 1。因此我们推出不等式

\[
| f (x) - f _ {m} (x) | = \left| \int_ {\mathbf {R} ^ {n}} (f (x) - f (x - y)) \varphi_ {m} (y) d \mu (y) \right| \leqslant \varepsilon
\]

（对所有 \(x \in \mathbf{R}^{n}\) ），由此得结论。

命题 4. --- a) \(\mathcal{D}(\mathrm{U})\) 到 \(\mathcal{K}(\mathrm{U})\) 中的典范单射是连续的，且 \(\mathcal{D}(\mathrm{U})\) 在 \(\mathcal{K}(\mathrm{U})\) 中稠密；

\begin{enumerate}
\def\labelenumi{\alph{enumi})}
\setcounter{enumi}{1}
\tightlist
\item
  对每个 \(p \in [1, +\infty[\) ，\(\mathcal{D}(\mathrm{U})\) 到 \(\mathcal{L}^p(\mathrm{U})\) 中的典范单射是连续的，且 \(\mathcal{D}(\mathrm{U})\) 在 \(\mathrm{L}^p(\mathrm{U})\) 中稠密。
\end{enumerate}

\(\mathcal{K}(\mathrm{U})\) 上的每个连续半范数在 \(\mathcal{D}(\mathrm{U})\) 上连续，因而 \(\mathcal{D}(\mathrm{U})\) 到 \(\mathcal{K}(\mathrm{U})\) 中的典范单射是连续的。

设 \(f \in \mathcal{K}(\mathrm{U})\) 。存在序列 \((f_{m})_{m \in \mathbf{N}}\) （在 \(\mathcal{D}(\mathrm{U})\) 中），它在 U 上一致收敛于 f，且诸 \(f_{m}\) 的支集都含于 f 的支集的一个固定的相对紧邻域（引理 3）。因此序列 \((f_{m})\) 收敛于 f （在 \(\mathcal{K}(\mathrm{U})\) 中）。这就完成了 a) 的证明。

设 \(p \in [1, +\infty[\) 。设 \(f \in \mathcal{D}(U)\) ，K 是 f 的支集。我们于是有不等式 \(\mathrm{N}_{p}(f) \leqslant \mu(\mathrm{K})^{1/p} \sup_{x \in \mathrm{K}} |f(x)|\) ，因而 \(\mathcal{D}(U)\) 到 \(\mathcal{L}^{p}(U)\) 中的典范单射是连续的。断言 b) 的后半部分由 a) 得出。

两个测试函数的乘积仍是测试函数，故莱布尼茨公式（FVR, I, p. 28，命题 2）表明空间 \(\mathcal{D}(\mathrm{U})\) 是一个拓扑代数。

更一般地，设 \(f \in \mathcal{C}^{\infty}(\mathrm{U})\) 。那么线性映射 \(\varphi \mapsto f\varphi\) ——从 \(\mathcal{D}(\mathrm{U})\) 到 \(\mathcal{D}(\mathrm{U})\) 中——是连续的。事实上，对每个紧子集 \(\mathbf{K}\) （\(\mathrm{U}\) 的）及每个 \(\alpha \in \mathbf{N}^{n}\) ，我们有

\[
p _ {\alpha , \mathrm{K}} (f \varphi) \leqslant \sum_ {0 \leqslant \beta \leqslant \alpha} \binom {\alpha} {\beta} p _ {\beta , \mathrm{K}} (f) p _ {\alpha - \beta , \mathrm{K}} (\varphi)
\]

（参见同前）。

\subsection{广义函数}

我们沿用上一节的记号；于是 U 表示 \(R^{n}\) 的开子集。

定义 1. --- \(\mathcal{D}(\mathrm{U})\) 的对偶空间，赋以有界收敛拓扑，称为 U 上的广义函数空间。记作 \(\mathcal{D}'(\mathrm{U})\) 。

于是，U 上的广义函数就是 \(\mathcal{D}(\mathrm{U})\) 上的连续线性形式。

若 \(f\) 是 \(\mathcal{D}(\mathrm{U})\) 上的线性形式（特别地，若 \(f\) 是广义函数），且 \(\varphi\) 属于 \(\mathcal{D}(\mathrm{U})\) ，我们用 \(\langle f, \varphi \rangle\) 表示 \(f\) 在 \(\varphi\) 处的取值。

由于 \(\mathcal{D}(\mathrm{U})\) 是有界型的，空间 \(\mathcal{D}'(\mathrm{U})\) 是完备的（TVS, III, p. 24，推论 1）。由于 \(\mathcal{D}(\mathrm{U})\) 是蒙泰尔空间，\(\mathcal{D}'(\mathrm{U})\) 亦然（EVT, IV, p. 19，命题 9）。特别地，空间 \(\mathcal{D}'(\mathrm{U})\) 是自反的（EVT, IV, p. 16，定理 2）。

引理 4. --- 设 f 是从 \(\mathcal{D}(\mathrm{U})\) 到 \(\mathbf{C}\) 中的线性映射。那么 \(f\) 是广义函数当且仅当：对每个紧子集 \(\mathrm{K}\) （\(\mathrm{U}\) 的）以及每个族 \((\mathrm{M}_{\alpha})_{\alpha \in \mathbf{N}^{n}}\) （\(\mathbf{R}_{+}\) 中的），线性形式 \(f\) 在由这样的函数 \(\varphi \in \mathcal{C}_{\mathrm{K}}^{\infty}(\mathrm{U})\) ——对每个 \(\alpha \in \mathbf{N}^{n}\) 有

\[
\sup _ {x \in \mathrm{K}} | \partial^ {\alpha} \varphi (x) | \leqslant \mathrm{M} _ {\alpha}.
\]

——所成之集上有界。事实上，空间 \(\mathcal{D}(\mathrm{U})\) 是有界型的，且 \(\mathcal{D}(\mathrm{U})\) 的每个有界子集都含于陈述中所描述的有界集之一。

引理 5. --- 设 F 是 \(\mathcal{D}'(\mathrm{U})\) 上具有可数基或包含一个单有界集合的滤子。那么 F 收敛于某个广义函数当且仅当 \(\langle f, \varphi \rangle\) 对 U 上的每个测试函数 \(\varphi\) 依 F 收敛。

特别地，广义函数序列 \((f_{m})_{m\in\mathbf{N}}\) 收敛于广义函数 f 当且仅当：对每个 \(\varphi\in\mathcal{D}(\mathrm{U})\) ，序列 \((\langle f_{m},\varphi\rangle)_{m\in\mathbf{N}}\) 收敛于 \(\langle f,\varphi\rangle\) 。

由于 \(\mathcal{D}'(\mathrm{U})\) 是蒙泰尔空间，故为桶型空间；这由巴拿赫–斯坦豪斯定理得出（EVT, III, p. 26，定理 1 的推论 3）。

设 V 是含于 U 的开集。从 \(\mathcal{D}(\mathrm{V})\) 到 \(\mathcal{D}(\mathrm{U})\) 中的典范线性映射的转置是从 \(\mathcal{D}'(\mathrm{U})\) 到 \(\mathcal{D}'(\mathrm{V})\) 中的连续线性映射，称为广义函数从 U 到 V 的限制，记作 \(r_{\mathrm{VU}}\) ，有时也记作 \(r_{\mathrm{V,U}}\) 。

有 \(r_{\mathrm{VV}} = 1_{\mathcal{D}'(\mathrm{V})}\) 。若 \(W \subset V \subset U\) 是 U 的开子集，则我们有 \(r_{\mathrm{WV}} \circ r_{\mathrm{VU}} = r_{\mathrm{WU}}\) 。换言之，若用 \(\mathcal{T}\) 表示 U 的拓扑，则投射系 \(\underline{\mathcal{D}'}(\mathrm{U}) = ((\mathcal{D}'(\mathrm{V}))_{\mathrm{V} \in \mathcal{T}}, (r_{\mathrm{WV}}))\) 是 U 上取值于局部凸拓扑向量空间的结构种的预层（TA, I, p. 42，定义 1 及 p. 66, §10）。

命题 5. --- 预层 \(\underline{\mathcal{D}}^{\prime}(\mathrm{U})\) 是一个层。

回顾（TA, I, p. 43，定义 2），这意味着对 U 的每个开子集 V 及 V 的每个开覆盖 \((\mathrm{V}_{i})_{i\in\mathrm{I}}\) ，下列条件满足：

\begin{enumerate}
\def\labelenumi{(\roman{enumi})}
\item
  映射 \((r_{\mathrm{V}_i\mathrm{V}})_{i\in \mathrm{I}}\colon \mathcal{D}'(\mathrm{V})\to \prod_{i\in \mathrm{I}}\mathcal{D}'(\mathrm{V}_i)\) 是单射；
\item
  对每个族 \((f_i)\in \prod_{i\in \mathrm{I}}\mathcal{D}'(\mathrm{V}_i)\) ——满足
\end{enumerate}

\[
r _ {(\mathrm{V} _ {i} \cap \mathrm{V} _ {i ^ {\prime}}) \mathrm{V} _ {i}} (f _ {i}) = r _ {(\mathrm{V} _ {i} \cap \mathrm{V} _ {i ^ {\prime}}) \mathrm{V} _ {i ^ {\prime}}} (f _ {i ^ {\prime}})
\]

——对每对 \((i, i') \in \mathrm{I} \times \mathrm{I}\) 都成立，则存在广义函数 \(f \in \mathcal{D}'(\mathrm{V})\) ，使得对每个 \(i \in I\) ，有 \(r_{\mathrm{V}_{i}\mathrm{V}}(f) = f_{i}\) 。

设 V 是 U 的开子集，\(\mathcal{V} = (\mathrm{V}_{i})_{i \in \mathrm{I}}\) 是 V 的开覆盖。设 \((\mathrm{W}_{j})_{j \in \mathrm{J}}\) 是 V 的开覆盖，它局部有限、比 \(\mathcal{V}\) 更细，且由相对紧的部分组成（IV, p. 201，引理 2）。取定单位分解 \((\varphi_{j})_{j \in \mathrm{J}}\) ——从属于覆盖 \((\mathrm{W}_{j})_{j \in \mathrm{J}}\) ——使得 \(\varphi_{j}\) 的支集含于 \(W_{j}\) （对所有 \(j \in J\) ；IV, p. 196，引理 1）。

设 \(\varphi\in\mathcal{D}(V)\) 。由于这样的 \(j\in J\) ——使 \(W_{j}\) 与 \(\varphi\) 的支集相交——组成的集合是有限的（TG, I, §9, no. 1, p. 59），我们有

\[
\varphi = \sum_ {j \in \mathrm{J}} \varphi \varphi_ {j}
\]

其中和式只有有限多个非零项。

我们来证明 (i)。设 \(f \in \mathcal{D}'(\mathrm{V})\) 满足 \(r_{\mathrm{V}_i\mathrm{V}}(f) = 0\) （对所有 \(i \in \mathrm{I}\) ）。这意味着 \(\langle f, \varphi \rangle = 0\) 对每个测试函数 \(\varphi\) （其支集含于某个开集 \(\mathrm{V}_i\) ）成立。若 \(\varphi\) 的支集含于某个开集 \(\mathrm{W}_j\) ，则更是如此。但这时，对每个 \(\varphi \in \mathcal{D}(\mathrm{V})\) ，我们有

\[
\langle f, \varphi \rangle = \langle f, \sum_ {j \in \mathrm{J}} \varphi \varphi_ {j} \rangle = \sum_ {j \in \mathrm{J}} \langle f, \varphi \varphi_ {j} \rangle = 0,
\]

因而 f = 0。

我们来证明 (ii)。设 \((f_{i})_{i\in\mathrm{I}}\) 是这样的族：\(f_{i}\in\mathcal{D}^{\prime}(\mathrm{V}_{i})\) （对所有 \(i\in I\) ），且 \(r_{\mathrm{V}_{i}\cap\mathrm{V}_{i^{\prime}},\mathrm{V}_{i}}(f_{i})=r_{\mathrm{V}_{i}\cap\mathrm{V}_{i^{\prime}},\mathrm{V}_{i^{\prime}}}(f_{i^{\prime}})\) 对 I 中所有 i 与 \(i^{\prime}\) 成立。设 \(\iota:J\to I\) 是满足 \(\mathrm{W}_{j}\subset\mathrm{V}_{\iota(j)}\) （对所有 \(j\in J\) ）的映射。对每个 \(j\in J\) ，令 \(\widetilde{f}_{j}=r_{\mathrm{W}_{j},\mathrm{V}_{\iota(j)}}(f_{\iota(j)})\) 。我们于是有

\[
\begin{array}{r l} & r _ {\mathrm{W} _ {j} \cap \mathrm{W} _ {j ^ {\prime}}, \mathrm{W} _ {j}} (\widetilde {f} _ {j}) = r _ {\mathrm{W} _ {j} \cap \mathrm{W} _ {j ^ {\prime}}, \mathrm{W} _ {j}} (r _ {\mathrm{W} _ {j}, \mathrm{V} _ {\iota (j)}} (f _ {\iota (j)})) \\ & \qquad = r _ {\mathrm{W} _ {j} \cap \mathrm{W} _ {j ^ {\prime}}, \mathrm{V} _ {\iota (j)}} (f _ {\iota (j)}) \\ & \qquad = r _ {\mathrm{W} _ {j} \cap \mathrm{W} _ {j ^ {\prime}}, \mathrm{V} _ {\iota (j)} \cap \mathrm{V} _ {\iota (j ^ {\prime})}} \circ r _ {\mathrm{V} _ {\iota (j)} \cap \mathrm{V} _ {\iota (j ^ {\prime})}, \mathrm{V} _ {\iota (j)}} (f _ {\iota (j)}). \end{array}
\]

交换 j 与 j' 的角色，并注意到

\[
r _ {\mathrm{V} _ {\iota (j)} \cap \mathrm{V} _ {\iota (j ^ {\prime})}, \mathrm{V} _ {\iota (j)}} (f _ {\iota (j)}) = r _ {\mathrm{V} _ {\iota (j)} \cap \mathrm{V} _ {\iota (j ^ {\prime})}, \mathrm{V} _ {\iota (j ^ {\prime})}} (f _ {\iota (j ^ {\prime})}),
\]

由假设成立，我们推出

\[
r _ {\mathrm{W} _ {j} \cap \mathrm{W} _ {j ^ {\prime}}, \mathrm{W} _ {j}} (\widetilde {f} _ {j}) = r _ {\mathrm{W} _ {j} \cap \mathrm{W} _ {j ^ {\prime}}, \mathrm{W} _ {j}} (\widetilde {f} _ {j ^ {\prime}})\tag{1}
\]

对所有 \((j, j') \in \mathrm{J}^{2}\) 成立。

对 \(\varphi \in \mathcal{D}(\mathrm{V})\) ，令

\[
\lambda (\varphi) = \sum_ {j \in \mathrm{J}} \langle \widetilde {f} _ {j}, \varphi \varphi_ {j} | \mathrm{W} _ {j} \rangle ,
\]

其中和式是有限的，因为只有有限多项可能非零。

映射 \(\lambda\) 是 \(\mathcal{D}(\mathrm{V})\) 上的线性形式。我们来证明它是广义函数。设 B 是 \(\mathcal{D}(\mathrm{V})\) 的有界子集，K 是 V 的紧子集且 \(B \subset \mathcal{C}_{K}^{\infty}(V)\) 。依据 TG, I, §9, no. 1, p. 59，存在 J 的有限子集 J' 使得

\[
\lambda (\varphi) = \sum_ {j \in J ^ {\prime}} \left\langle \widetilde {f} _ {j}, \varphi \varphi_ {j} \mid W _ {j} \right\rangle
\]

对所有 \(\varphi\in\mathrm{B}\) 成立。由于 \(\widetilde{f}_{j}\) 对每个 \(j\in\mathrm{J}^{\prime}\) 是广义函数，且 \(\mathcal{D}(\mathrm{V})\) 是拓扑代数，可知 \(\lambda\) 在 B 上有界，由此得结论（引理 4）。

设 \(j \in J\) ，\(\varphi \in \mathcal{D}(V)\) 的支集含于 \(W_{j}\) 。我们于是有

\[
\langle \lambda , \varphi \rangle = \sum_ {j ^ {\prime} \in \mathrm{J}} \left\langle \widetilde {f} _ {j ^ {\prime}}, \varphi \varphi_ {j ^ {\prime}} \right| \mathrm{W} _ {j} \rangle = \sum_ {j ^ {\prime} \in \mathrm{J}} \left\langle \widetilde {f} _ {j}, \varphi \varphi_ {j ^ {\prime}} \right| \mathrm{W} _ {j} \rangle
\]

这依据 (1)，因为 \(\varphi\varphi_{j'}\) 的支集含于 \(W_j \cap W_{j'}\) 。从而

\[
\langle \lambda , \varphi \rangle = \langle \widetilde {f} _ {j}, \sum_ {j ^ {\prime} \in \mathrm{J}} \varphi \varphi_ {j ^ {\prime}} | \mathrm{W} _ {j} \rangle = \langle \widetilde {f} _ {j}, \varphi | \mathrm{W} _ {j} \rangle ,
\]

因而 \(r_{\mathrm{W}_j\mathrm{V}}(\lambda) = \widetilde{f}_j\) 对所有 \(j\in \mathbf{J}\) 成立。

设 \(i \in I\) 。最后我们来证明 \(\lambda\) 在 \(V_i\) 上的限制与 \(f_i\) 一致。将条件 (i) 应用于 \(V_i\) 的由诸开集 \(W_j\) 组成的覆盖，只需验证对每个 \(j \in J\) ，\(\lambda\) 在 \(V_i \cap W_j\) 上的限制与 \(f_i\) 的限制一致。由前所述，只需验证 \(f_i\) 在 \(V_i \cap W_j\) 上的限制与 \(\widetilde{f}_j\) 的限制一致。而我们有

\[
\begin{array}{c} r _ {\mathrm{V} _ {i} \cap \mathrm{W} _ {j}, \mathrm{V} _ {i}} (f _ {i}) = r _ {\mathrm{V} _ {i} \cap \mathrm{W} _ {j}, \mathrm{V} _ {i} \cap \mathrm{V} _ {\iota (j)}} (r _ {\mathrm{V} _ {i} \cap \mathrm{V} _ {\iota (j)}, \mathrm{V} _ {i}} (f _ {i})) \\ = r _ {\mathrm{V} _ {i} \cap \mathrm{W} _ {j}, \mathrm{V} _ {i} \cap \mathrm{V} _ {\iota (j)}} (r _ {\mathrm{V} _ {i} \cap \mathrm{V} _ {\iota (j)}, \mathrm{V} _ {\iota (j)}} (f _ {\iota (j)})) = r _ {\mathrm{V} _ {i} \cap \mathrm{W} _ {j}, \mathrm{V} _ {\iota (j)}} (f _ {\iota (j)}), \end{array}
\]

其中我们用到了关于族 \((f_{i})\) 的假设。于是

\[
r _ {\mathrm{V} _ {i} \cap \mathrm{W} _ {j}, \mathrm{V} _ {\iota (j)}} (f _ {\iota (j)}) = r _ {\mathrm{V} _ {i} \cap \mathrm{W} _ {j}, \mathrm{W} _ {j}} (r _ {\mathrm{W} _ {j}, \mathrm{V} _ {\iota (j)}} (f _ {\iota (j)})) = r _ {\mathrm{V} _ {i} \cap \mathrm{W} _ {j}, \mathrm{W} _ {j}} (\widetilde {f} _ {j}),
\]

由此即可得出结论。

\subsection{函数解释为广义函数}

命题 6. --- 设 ν 是 U 上的测度。ν 在 D(U) 上的限制是一个广义函数，它为零当且仅当测度 ν 为零。

设 K 是 U 的紧子集。对每个函数 \(\varphi\in\mathcal{C}_{\mathrm{K}}^{\infty}(\mathrm{U})\) ，我们有 \(|\langle\nu,\varphi\rangle|\leqslant p_{0,\mathrm{K}}(\varphi)|\nu|(\mathrm{K})\) ，因而 \(\nu\) 在 \(\mathcal{D}(\mathrm{U})\) 上的限制是连续的。由于 \(\mathcal{D}(\mathrm{U})\) 在 \(\mathcal{K}(\mathrm{U})\) 中稠密，依据 IV, p. 202 的命题 4(a)，\(\nu\) 在 \(\mathcal{D}(\mathrm{U})\) 上的限制为零当且仅当 \(\nu\) 为零。

我们将把 U 上的复测度空间 \(\mathcal{M}(\mathrm{U};\mathbf{C})\) 与 \(\mathcal{D}'(\mathrm{U})\) 的一个子空间等同起来。我们还将把空间 \(\mathrm{L}_{\mathrm{loc}}^{1}(\mathrm{U})\) 与 \(\mathcal{D}'(\mathrm{U})\) 的一个子空间等同起来，所用的映射是由把 \(f\in \mathcal{L}_{\mathrm{loc}}^{1}(\mathrm{U})\) 映为测度 \(f\cdot \mu\) 的映射过渡到商而诱导的映射（INT, V, p. 44, § 5, no. 2，定义 2）。换言之，对 \(f\in \mathrm{L}_{\mathrm{loc}}^{1}(\mathrm{U})\) 与 \(\varphi \in \mathcal{D}(\mathrm{U})\) ，我们有

\[
\langle f, \varphi \rangle = \int_ {\mathrm{U}} f \varphi d \mu .
\]

特别地，对 \(p \in [1, +\infty]\) ，这使我们可以把空间 \(\mathrm{L}^p(\mathrm{U})\) 与 \(\mathcal{D}'(\mathrm{U})\) 的一个子空间等同起来（参见 INT, V, p. 43, § 5, no. 1）。

命题 7. --- 设 \(p \in [1, +\infty]\) 。\(\mathrm{L}^p(\mathrm{U})\) 到 \(\mathcal{D}'(\mathrm{U})\) 中的单射是连续的。

设 \(f \in \mathrm{L}^{p}(\mathrm{U})\) 。设 K 是 U 的紧子集，用 \(\varphi_{K}\) 表示它的特征函数。对每个支集含于 K 的测试函数 \(\varphi\) ，赫尔德不等式给出

\[
| \langle f, \varphi \rangle | = \left| \int_ {\mathrm{U}} f \varphi d \mu \right| \leqslant \mathrm{N} _ {p} (f) \mathrm{N} _ {q} (\varphi) \leqslant \mathrm{N} _ {p} (f) \mathrm{N} _ {q} (\varphi_ {\mathrm{K}}) p _ {0, \mathrm{K}} (\varphi)
\]

其中 q 是 p 的共轭指数。

特别地，可以把 \(\mathcal{D}(\mathrm{U})\) 与 \(\mathcal{D}'(\mathrm{U})\) 的一个子空间等同起来。

命题 8. --- 空间 \(\mathcal{D}(\mathrm{U})\) 在 \(\mathcal{D}'(\mathrm{U})\) 中稠密。

设 \(\lambda\) 是 \(\mathcal{D}'(\mathrm{U})\) 上在 \(\mathcal{D}(\mathrm{U})\) 上为零的线性形式。由于 \(\mathcal{D}(\mathrm{U})\) 是自反的，存在测试函数 \(\varphi \in \mathcal{D}(\mathrm{U})\) ，使得 \(\lambda(f) = \langle f, \varphi \rangle\) 对所有 \(f \in \mathcal{D}'(\mathrm{U})\) 成立。我们得到

\[
0 = \lambda (\overline {{\varphi}}) = \langle \overline {{\varphi}}, \varphi \rangle = \int_ {\mathrm{U}} | \varphi | ^ {2} d \mu ,
\]

因而 \(\varphi = 0\) 。命题由哈恩--巴拿赫定理得出（EVT, II, p. 49，推论 3 (ii)）。

对 \(h \in \mathcal{C}^{\infty}(\mathrm{U})\) ，连续线性映射 \(\varphi \mapsto h\varphi\) （\(\mathcal{D}(\mathrm{U})\) 到自身的）的转置是 \(\mathcal{D}'(\mathrm{U})\) 到自身的连续线性映射，记作 \(f \mapsto hf\) 。这个定义的合理性在于：若 \(f\) 是与测度 \(\nu\) （\(\mathrm{U}\) 上的）相伴的广义函数，则 \(hf\) 与测度 \(h \cdot \nu\) 相伴。事实上，对每个测试函数 \(\varphi \in \mathcal{D}(\mathrm{U})\) ，可以算出

\[
\langle h f, \varphi \rangle = \langle f, h \varphi \rangle = \int_ {\mathrm{U}} h \varphi d \nu = \int_ {\mathrm{U}} \varphi d (h \cdot \nu).
\]

\subsection{广义函数的求导}

设 \(\alpha \in \mathbf{N}^n\) 。线性映射 \(\varphi \mapsto \partial^\alpha \varphi\) 从 \(\mathcal{D}(\mathrm{U})\) 到 \(\mathcal{D}(\mathrm{U})\) 中是连续的。它的转置 \(^t\partial^\alpha\) 是从 \(\mathcal{D}'(\mathrm{U})\) 到 \(\mathcal{D}'(\mathrm{U})\) 中的连续线性映射（EVT, IV, p. 6，推论，\(b\) ））。

我们用 \(\partial^{\alpha}\) 表示连续线性映射 \((-1)^{|\alpha|}{}^{t}\partial^{\alpha}\) ——从 \(\mathcal{D}'(\mathrm{U})\) 到自身。

定义 2. --- 若 \(f \in \mathcal{D}'(\mathrm{U})\) 是广义函数，则称 \(\partial^{\alpha} f\) 为 \(\alpha\) 阶迭代偏导数（\(f\) 的）。

于是，按定义

\[
\langle \partial^ {\alpha} f, \varphi \rangle = (- 1) ^ {| \alpha |} \langle f, \partial^ {\alpha} \varphi \rangle
\]

（对每个函数 \(\varphi\in\mathcal{D}(\mathrm{U})\) ）。等式 \(\partial^{\alpha+\beta}=\partial^{\alpha}\circ\partial^{\beta}\) 对所有 \(\alpha\) 与 \(\beta\) （\(\mathbf{N}^{n}\) 中的）成立。

若 \(n=1\) ，也用 \(f'\) 记广义函数 \(f\in\mathcal{D}'(U)\) 的导数。该定义的合理性由下述引理给出。

引理 6. --- 设 k 是自然数。设 \(f \in \mathcal{C}^k(\mathrm{U})\) ，\(\lambda\) 是与 \(f\) 相伴的广义函数。对每个 \(\beta \in \mathbf{N}^n\) （满足 \(|\beta| \leqslant k\) ），广义函数 \(\partial^\beta\lambda\) 是与函数 \(\partial^\beta f\) 相伴的广义函数。

对 \(k\) 作归纳，只需在 \(\beta\) 满足 \(|\beta| = 1\) 时证明此性质，甚至可以假设 \(\beta = (0, \ldots, 0, 1)\) 。由于广义函数构成一个层（IV, p. 204，命题 5），只需在存在开集 \(V \subset \mathbf{R}^{n-1}\) 与开区间 \(I \subset \mathbf{R}\) 使得 \(U = V \times I\) 时验证断言。

对每个测试函数 \(\varphi\in\mathcal{D}(\mathrm{U})\) ，按定义有

\[
\langle \partial^ {\beta} \lambda , \varphi \rangle = - \int_ {\mathrm{U}} f (x) \partial^ {\beta} \varphi (x) d x = - \int_ {\mathrm{V}} \Bigl (\int_ {\mathrm{I}} f (y, t) \partial^ {\beta} \varphi (y, t) d t \Bigr) d y
\]

这里用到勒贝格--富比尼定理（INT, V, p. 96, § 8, n° 4，定理 1）。由分部积分（FVR, II, p. 10），有

\[
- \int_ {\mathrm{I}} f (y, t) \partial^ {\beta} \varphi (y, t) d t = \int_ {\mathrm{I}} \partial^ {\beta} f (y, t) \varphi (y, t) d t,
\]

因为 \(t \mapsto \varphi(y, t)\) 在 I 中有紧支集。于是得到

\[
\langle \partial^ {\beta} \lambda , \varphi \rangle = \int_ {V} \left(\int_ {I} \partial^ {\beta} f (y, t) \varphi (y, t) d t\right) d y = \langle \partial^ {\beta} f, \varphi \rangle .
\]

命题 9（莱布尼茨公式）. --- 设 f 是 U 上的广义函数，g 是 U 上的无穷次可微函数。设 \(\alpha \in \mathbf{N}^n\) 。我们有关系式

\[
\partial^ {\alpha} (f g) = \sum_ {\beta \leqslant \alpha} \binom {\alpha} {\beta} \partial^ {\beta} f   \partial^ {\alpha - \beta} g.
\]

如同 FVR, I, p. 28，命题 2 的证明那样对 \(|\alpha|\) 作归纳，只需考虑 \(|\alpha|=1\) 的情形。结论由计算

\[
\begin{array}{c} \langle \partial^ {\alpha} (f g), \varphi \rangle = \langle f g, - \partial^ {\alpha} \varphi \rangle = \langle f, - g \partial^ {\alpha} \varphi \rangle \\ = \langle f, - \partial^ {\alpha} (g \varphi) + \varphi \partial^ {\alpha} g \rangle = \langle g \partial^ {\alpha} f + f \partial^ {\alpha} g, \varphi \rangle \end{array}
\]

（对 \(\varphi\in\mathcal{D}(\mathrm{U})\) 成立）得出。

\subsection{施瓦茨函数}

我们用 \(\mathcal{S}(\mathbf{R}^{n})\) 表示无穷次可微函数 \(\varphi\) （在 \(\mathbf{R}^{n}\) 上、取复值）所成的空间，其中对所有 \(\alpha\) 与 \(\beta\) （属于 \(\mathbf{N}^{n}\) ），函数 \(\mathrm{X}^{\beta}\partial^{\alpha}\varphi\) 在 \(\mathbf{R}^{n}\) 上有界。我们给 \(\mathcal{S}(\mathbf{R}^{n})\) 赋以由可数半范数族 \((q_{\alpha,\beta})_{(\alpha,\beta)\in\mathbf{N}^{n}\times\mathbf{N}^{n}}\) 定义的局部凸拓扑，其中 \(q_{\alpha,\beta}\) 定义为

\[
q _ {\alpha , \beta} (\varphi) = \sup _ {x \in \mathbf {R} ^ {n}} | x ^ {\beta} \partial^ {\alpha} \varphi (x) | = \| \mathrm{X} ^ {\beta} \partial^ {\alpha} \varphi \| _ {\infty}
\]

（对 \(\varphi\in\mathcal{S}(\mathbf{R}^{n})\) ）。该拓扑是分离的。它也由半范数 \(\widetilde{q}_{\alpha,k}\) 定义，其中

\[
\widetilde {q} _ {\alpha , k} (\varphi) = \sup _ {x \in \mathbf {R} ^ {n}} \| x \| ^ {k} | (\partial^ {\alpha} \varphi) (x) |.
\]

对所有 \(\varphi\in\mathcal{S}(\mathbf{R}^{n})\) ，其中 \(k\in\mathbf{N}\) ，\(\alpha\in\mathbf{N}^{n}\) 。

我们称 \(\mathcal{S}(\mathbf{R}^{n})\) 为 \(\mathbf{R}^{n}\) 上的施瓦茨空间或施瓦茨函数空间。

注. --- 设 \(\varphi\in\mathcal{S}(\mathbf{R}^{n})\) 。对每个 \(k\in\mathbf{N}\) ，我们有

\[
\lim _ {\| x \| \to + \infty} \| x \| ^ {k} \varphi (x) = 0,
\]

因为函数 \(x \mapsto \|x\|^{k+1} \varphi(x)\) 有界。

例. --- 函数 \(\gamma_{n}\) ——在 \(\mathbf{R}^{n}\) 上由 \(\gamma_{n}(x)=\exp(-\|x\|^{2})\) 定义——属于 \(\mathcal{S}(\mathbf{R}^{n})\) 。事实上，对 k 作归纳可以证明：对每个整数 \(k\in \mathbf{N}\) ，存在多项式 \(P_{k}\in\mathbf{R}[X]\) 使得 \(\partial_{k}\gamma_{1}=P_{k}\gamma_{1}\) 。

于是，对所有 \(\alpha = (\alpha_{i})\in \mathbf{N}^{n}\) 与 \(\beta = (\beta_{i})\in \mathbf{N}^{n}\) ，以及每个 \(x = (x_{i})\in \mathbf{R}^{n}\) ，可得

\[
| (\mathrm{X} ^ {\beta} \partial^ {\alpha} \gamma_ {n}) (x) | = \prod_ {i = 1} ^ {n} | x _ {i} | ^ {\beta_ {i}} | \mathrm{P} _ {\alpha_ {i}} (x _ {i}) | \gamma_ {1} (x _ {i}),
\]

它在 x 于 \(\mathbf{R}^{n}\) 中变化时是有界量。

设 \(\alpha \in \mathbf{N}^n\) 与 \(\beta \in \mathbf{N}^n\) 。若 \(\varphi \in \mathcal{S}(\mathbf{R}^n)\) ，则 \(\mathrm{X}^{\beta}\partial^{\alpha}(\varphi)\) 仍是施瓦茨函数；如此定义的映射 \(\varphi \mapsto \mathrm{X}^{\beta}\partial^{\alpha}\varphi\) ——从 \(\mathcal{S}(\mathbf{R}^n)\) 到自身——是连续的。

空间 \(\mathcal{S}(\mathbf{R}^{n})\) 是一个拓扑代数。更确切地说，若 \(\varphi_{1}\) 与 \(\varphi_{2}\) 属于 \(\mathcal{S}(\mathbf{R}^{n})\) ，则 \(\varphi_{1}\varphi_{2}\) 是满足

\[
q _ {\alpha , \beta} (\varphi_ {1} \varphi_ {2}) \leqslant \sum_ {0 \leqslant \gamma \leqslant \alpha} \binom{\alpha}{\gamma} q _ {\gamma , \beta} (\varphi_ {1}) q _ {\alpha - \gamma , 0} (\varphi_ {2})\tag{2}
\]

（对所有 \(\alpha\) 与 \(\beta \in \mathbf{N}^n\) ；IV, p. 196，命题 1）的施瓦茨函数。

\(\mathcal{S}(\mathbf{R}^{n})\) 到空间 \(\mathcal{C}^{\infty}(\mathbf{R}^{n})\) （赋以 IV, p. 200 的 §3 中所述的拓扑）中的典范包含是连续的，因为

\[
\sup _ {x \in K} | \partial^ {\alpha} \varphi (x) | \leqslant q _ {\alpha , 0} (\varphi)
\]

对 \(\mathbf{R}^{n}\) 的每个紧子集 K、所有 \(\alpha \in \mathbf{N}^{n}\) 及每个施瓦茨函数 \(\varphi\) 都成立。

引理 7. --- 设 \(k \in \mathbf{N}\) ，\(\alpha \in \mathbf{N}^n\) 。对每个函数 \(\varphi \in \mathcal{S}(\mathbf{R}^n)\) 及每个实数 \(T > 0\)，我们有

\[
\widetilde {q} _ {\alpha , k} (\varphi) \leqslant \mathrm{T} ^ {k} \sup _ {\| x \| \leqslant \mathrm{T}} | \partial^ {\alpha} \varphi (x) | + \frac {1}{\mathrm{T}} \widetilde {q} _ {\alpha , k + 1} (\varphi).\tag{3}
\]

事实上，有

\[
\begin{array}{c} \widetilde {q} _ {\alpha , k} (\varphi) \leqslant \sup _ {\| x \| \leqslant T} \| x \| ^ {k} | \partial^ {\alpha} \varphi (x) | + \sup _ {\| x \| > T} \| x \| ^ {k} | \partial^ {\alpha} \varphi (x) | \\ \leqslant T ^ {k} \sup _ {\| x \| \leqslant T} | \partial^ {\alpha} \varphi (x) | + \frac {1}{T} \sup _ {\| x \| > T} \| x \| ^ {k + 1} | \partial^ {\alpha} \varphi (x) |. \end{array}
\]

命题 10. --- 设 B 是 \(\mathcal{S}(\mathbf{R}^{n})\) 的有界子集。\(\mathcal{S}(\mathbf{R}^{n})\) 在 B 上诱导的拓扑与 \(\mathcal{C}^{\infty}(\mathbf{R}^{n})\) 诱导的拓扑一致。

由于 \(\mathcal{S}(\mathbf{R}^{n})\) 到 \(\mathcal{C}^{\infty}(\mathbf{R}^{n})\) 中的包含是连续的，只需证明：对 \(\mathcal{S}(\mathbf{R}^{n})\) 的每个开子集 V，交集 \(V \cap B\) 对 \(\mathcal{C}^{\infty}(\mathbf{R}^{n})\) 诱导的拓扑而言在 B 中是开的。

设 V 是 \(\mathcal{S}(\mathbf{R}^{n})\) 的开子集。设 \(\varphi_{0} \in \mathrm{V} \cap \mathrm{B}\) 。存在有限集 I、族 \((\alpha_{i}, k_{i})_{i \in \mathrm{I}} \in (\mathbf{N}^{n} \times \mathbf{N})^{\mathrm{I}}\) 及实数 \(\varepsilon > 0\) ，使得 V 包含由 \(\varphi \in \mathcal{S}(\mathbf{R}^{n})\) ——满足

\[
\sup _ {i \in I} \widetilde {q} _ {\alpha_ {i}, k _ {i}} (\varphi - \varphi_ {0}) \leqslant \varepsilon .
\]

——所成之集。由于 B 在 \(\mathcal{S}(\mathbf{R}^{n})\) 中有界，存在 \(M > 0\) ，使得半范数 \(\widetilde{q}_{\alpha_{i},k_{i} + 1}\) （对 \(i\in I\) ）以 \(M\) 为界（在 \(B\) 上）。设 \(\delta >0\) 与 \(T > 0\) 为实数。依据不等式 (3)，只要 \(\varphi \in B\) 满足界

\[
\sup _ {i \in I} \sup _ {\| x \| \leqslant T} | \partial^ {\alpha_ {i}} (\varphi - \varphi_ {0}) | \leqslant \delta ,\tag{4}
\]

我们就有

\[
\sup _ {i \in I} \widetilde {q} _ {\alpha_ {i}, k _ {i}} (\varphi - \varphi_ {0}) \leqslant \delta T ^ {k} + \frac {2 M}{T}.
\]

取 \(T = \frac{4M}{\varepsilon}\) ，再取 \(\delta = \frac{\varepsilon}{2T^{k}}\) 。可以看出，\(V \cap B\) 包含 \(\varphi_{0}\) 在 B 中的、对 \(\mathcal{C}^{\infty}(\mathbf{R}^{n})\) 诱导拓扑而言由 (4) 定义的邻域。证明至此完成。

推论. --- 设 \((\varphi_{m})_{m\in\mathbf{N}}\) 是 \(\mathcal{S}(\mathbf{R}^{n})\) 中的有界序列。对每个函数 \(\varphi\in\mathcal{S}(\mathbf{R}^{n})\) ，下列断言等价：

\begin{enumerate}
\def\labelenumi{\alph{enumi})}
\item
  序列 \((\varphi_{m})\) 收敛于 \(\varphi\) （在 \(\mathcal{S}(\mathbf{R}^n)\) 中）；
\item
  序列 \((\varphi_{m})\) 收敛于 \(\varphi\) （在 \(\mathcal{C}^{\infty}(\mathbf{R}^{n})\) 中）。
\end{enumerate}

注. --- 序列 \((\varphi_{m})\) （在 \(\mathcal{C}^{\infty}(\mathbf{R}^{n})\) 中）收敛当且仅当：对每个 \(\alpha\in \mathbf{N}^{n}\) ，序列 \((\partial^{\alpha}\varphi_{m})\) 收敛于某个函数 \(\varphi^{(\alpha)}\) （在赋以紧收敛拓扑的 \(\mathcal{C}(\mathbf{R}^{n})\) 中）。这时我们有 \(\varphi^{(\alpha)}=\partial^{\alpha}\varphi\) ，且 \((\varphi_{m})\) 收敛于 \(\varphi^{(0)}\) 。

事实上，该条件是必要的。反之，若诸序列 \((\partial^{\alpha}\varphi_{m})\) 收敛于函数 \(\varphi^{(\alpha)}\) ——对所有 \(\alpha\in \mathbf{N}^{n}\) ——，则由 FVR, II, p. 2，定理 1 可得 \(\varphi^{(\alpha)}=\partial^{\alpha}\varphi^{(0)}\) ，这意味着序列 \((\varphi_{m})\) 收敛于 \(\varphi^{(0)}\) （在 \(\mathcal{C}^{\infty}(\mathbf{R}^{n})\) 中）。

命题 11. --- 空间 \(\mathcal{S}(\mathbf{R}^{n})\) 是弗雷歇空间和蒙泰尔空间。

由于空间 \(\mathcal{C}^{\infty}(\mathbf{R}^{n})\) 是完备的（EVT, III, p. 9，例 \(b\) ），命题 10 的推论表明 \(\mathcal{S}(\mathbf{R}^{n})\) 中的每个柯西序列都在 \(\mathcal{S}(\mathbf{R}^{n})\) 中收敛，因为它有界且在 \(\mathcal{C}^{\infty}(\mathbf{R}^{n})\) 中收敛。

于是空间 \(\mathcal{S}(\mathbf{R}^n)\) 是弗雷歇空间；特别地，它是桶型的（EVT, III, p. 25，命题 2 的推论）。设 B 是 \(\mathcal{S}(\mathbf{R}^n)\) 的有界子集，\((\varphi_m)_{m \in \mathbf{N}}\) 是取值于 B 的序列。由于 \(\mathcal{C}^\infty(\mathbf{R}^n)\) 是蒙泰尔空间（EVT, IV, p. 18，例 (4)），存在 \((\varphi_m)_{m \in \mathbf{N}}\) 的子序列在 \(\mathcal{C}^\infty(\mathbf{R}^n)\) 中收敛，从而在 \(\mathcal{S}(\mathbf{R}^n)\) 中收敛（命题 10）。因而 B 在 \(\mathcal{S}(\mathbf{R}^n)\) 中相对紧。由此得出 \(\mathcal{S}(\mathbf{R}^n)\) 是蒙泰尔空间。

\subsection{函数空间在施瓦茨函数空间中的包含}

命题 12. — 空间 \(\mathcal{D}(\mathbf{R}^{n})\) 含于 \(\mathcal{S}(\mathbf{R}^{n})\) ，且 \(\mathcal{D}(\mathbf{R}^{n})\) 到 \(\mathcal{S}(\mathbf{R}^{n})\) 中的包含是连续的且像稠密。

设 \(B \subset \mathcal{D}(\mathbf{R}^{n})\) 是有界子集，K 是 \(\mathbf{R}^{n}\) 的紧子集且 \(\mathrm{B} \subset \mathcal{C}_{\mathrm{K}}^{\infty}(\mathbf{R}^{n})\) 。设 \(\alpha \in \mathbf{N}^{n}\) 与 \(k \in \mathbf{N}\) 。对每个函数 \(\varphi \in B\) ，可得

\[
\widetilde {q} _ {\alpha , k} (\varphi) \leqslant \Bigl (\sup _ {x \in \mathrm{K}} \| x \| ^ {k} \Bigr) p _ {\alpha , \mathrm{K}} (\varphi),
\]

因而 B 在 \(\mathcal{S}(\mathbf{R}^{n})\) 中有界。包含的连续性则由空间 \(\mathcal{S}(\mathbf{R}^{n})\) 与 \(\mathcal{D}(\mathbf{R}^{n})\) 都是有界型的这一事实得出。

我们来证明 \(\mathcal{D}(\mathbf{R}^n)\) 在 \(\mathcal{S}(\mathbf{R}^n)\) 中稠密。设 B 是 \(\mathbf{R}^n\) 的单位球，\(\eta \in \mathcal{D}(\mathbf{R}^n)\) 是支集含于 2B 且满足 \(0 \leqslant \eta \leqslant 1\) 及 \(\eta(x) = 1\) （对所有 \(x \in \mathrm{B}\) ）的测试函数（IV, p. 196，引理 1, \(a\) ）。

设 \(\varphi \in \mathcal{S}(\mathbf{R}^n)\) 。对每个整数 \(m \geqslant 1\) 及每个 \(x \in \mathbf{R}^n\) ，令 \(\eta_m(x) = \eta(x/m)\) 。最后定义 \(\varphi_m = \eta_m\varphi\) ；我们有 \(\varphi_m \in \mathcal{D}(\mathbf{R}^n)\) 。

由于 \(\partial^{\alpha}\eta_{m} = m^{-|\alpha|}(\partial^{\alpha}\eta)(x / m)\) （对所有 \(\alpha \in \mathbf{N}^{n}\) 与 \(x\in \mathbf{R}^n\) ），由 IV, p. 210 的公式 (2) 推出序列 \((\varphi_{m})\) 在 \(\mathcal{S}(\mathbf{R}^n)\) 中有界。

设 C 是 \(R^{n}\) 的紧子集。序列 \((\varphi_{m})_{m\geqslant1}\) 收敛于 \(\varphi\) （在 \(\mathcal{C}_{\mathrm{K}}^{\infty}(\mathbf{R}^{n})\) 中），因为对所有充分大的 m，\(\varphi_{m}\) 在 C 上与 \(\varphi\) 一致。于是序列 \((\varphi_{m})\) 收敛于 \(\varphi\) （在 \(\mathcal{C}^{\infty}(\mathbf{R}^{n})\) 中），而 IV, p. 211 的命题 10 的推论使我们得出序列 \((\varphi_{m})\) 收敛于 \(\varphi\) （在 \(\mathcal{S}(\mathbf{R}^{n})\) 中）。

引理 8. --- 设 B 是 \(\mathcal{D}(\mathbf{R}^{n})\) 的有界子集。\(\mathcal{S}(\mathbf{R}^{n})\) 的拓扑在 B 上诱导的拓扑与 \(\mathcal{D}(\mathbf{R}^{n})\) 诱导的拓扑一致。

由于 \(\mathcal{D}(\mathbf{R}^{n})\) 到 \(\mathcal{S}(\mathbf{R}^{n})\) 中的包含是连续的，\(\mathcal{D}(\mathbf{R}^{n})\) 在 B 上诱导的拓扑比 \(\mathcal{S}(\mathbf{R}^{n})\) 诱导的拓扑更细。此外，存在 \(\mathbf{R}^{n}\) 的紧子集 K 使得 \(\mathrm{B} \subset \mathcal{C}_{\mathrm{K}}^{\infty}(\mathbf{R}^{n})\) 。于是对每个 \(\alpha \in \mathbf{N}^{n}\) ，有 \(p_{\alpha,K}(\varphi) \leqslant \widetilde{q}_{\alpha,0}(\varphi)\) ，这蕴含 \(\mathcal{S}(\mathbf{R}^{n})\) 诱导的拓扑比 \(\mathcal{D}(\mathbf{R}^{n})\) 诱导的拓扑更细。

命题 13. --- 设 \(p \in [1, +\infty]\) 。空间 \(\mathcal{S}(\mathbf{R}^{n})\) 含于 \(\mathcal{L}^{p}(\mathbf{R}^{n})\) ，且 \(\mathcal{S}(\mathbf{R}^{n})\) 到 \(\mathcal{L}^{p}(\mathbf{R}^{n})\) 中的典范单射是连续的。则 \(\mathcal{S}(\mathbf{R}^{n})\) 在 \(\mathrm{L}^{p}(\mathbf{R}^{n})\) 中的像（若 \(p \neq +\infty\) ）稠密。

对 \(p = +\infty\) ，第一个断言是直接的。现设 \(p \in [1, +\infty[\) 。设 m 是满足 \(n + 1 < mp\) 的整数。对每个 \(\varphi \in \mathcal{S}(\mathbf{R}^{n})\) 与 \(x \in \mathbf{R}^{n}\) ，我们有

\[
\| x \| ^ {n + 1} | \varphi (x) | ^ {p} \leqslant \widetilde {q} _ {0, m} (\varphi) ^ {p}
\]

因而依据 IV, p. 199 的命题 3，\(\varphi \in \mathcal{L}^p (\mathbf{R}^n)\) 。此外，我们得到

\[
\mathrm{N} _ {p} (\varphi) \leqslant a _ {n} ^ {1 / p} \widetilde {q} _ {0, 0} (\varphi) + b _ {n} \widetilde {q} _ {0, m} (\varphi),
\]

其中

\[
a _ {n} = \int_ {\| x \| \leqslant 1} d \mu (x), \quad b _ {n} = \int_ {\| x \| \geqslant 1} \frac {1}{\| x \| ^ {n + 1}} d \mu (x),
\]

因此 \(\mathcal{S}(\mathbf{R}^n)\) 到 \(\mathcal{L}^p (\mathbf{R}^n)\) 中的单射是连续的。

由于空间 \(\mathcal{D}(\mathbf{R}^n)\) 含于 \(\mathcal{S}(\mathbf{R}^n)\) ，IV, p. 202 的命题 4 蕴含 \(\mathcal{S}(\mathbf{R}^n)\) 在 \(\mathrm{L}^p (\mathbf{R}^n)\) 中稠密（若 \(p\neq +\infty\) ）。

\subsection{多项式增长的函数}

定义 3. --- 函数 \(f\colon\mathbf{R}^{n}\to\mathbf{C}\) 称为具有多项式增长的，如果存在整数 \(k\geqslant1\) ，使得由 \(x\mapsto(1+\|x\|)^{-k}f(x)\) 定义的映射在 \(\mathbf{R}^{n}\) 上有界。

每个具有多项式增长的函数都是局部有界的。\(\mathbf{R}^{n}\) 上的每个多项式函数都具有多项式增长。

命题 14. --- 设 \(f \in \mathcal{C}^{\infty}(\mathbf{R}^{n})\) 。假设对每个 \(\alpha\) （\(\mathbf{N}^{n}\) 中的），函数 \(\partial^{\alpha}f\) 都具有多项式增长。那么从空间 \(\mathcal{S}(\mathbf{R}^{n})\) 到自身的、由 \(\varphi \mapsto f\varphi\) 所定义的线性映射是连续的。

对每个 \(\varphi\) （在 \(\mathcal{S}(\mathbf{R}^{n})\) 中），函数 \(f\varphi\) 属于 \(\mathcal{C}^{\infty}(\mathbf{R}^{n})\) 。由假设，对每个 \(\alpha \in \mathbf{N}^{n}\) ，存在整数 \(k_{\alpha} \geqslant 0\) 与实数 \(\mathrm{C}_{\alpha}\) ，使得 \(|\partial^{\alpha}f(x)| \leqslant \mathrm{C}_{\alpha}(1 + \|x\|)^{k_{\alpha}}\) 对所有 \(x\) （在 \(\mathbf{R}^{n}\) 中）成立。设 \(\varphi \in \mathcal{S}(\mathbf{R}^{n})\) 。设 \(\alpha \in \mathbf{N}^{n}\) 与 \(k \in \mathbf{N}\) 。依据 IV, p. 196 的命题 1，可得

\[
\begin{array}{l} \widetilde {q} _ {\alpha , k} (f \varphi) \leqslant \sum_ {0 \leqslant \beta \leqslant \alpha} \binom {\alpha} {\beta} \sup _ {x \in \mathbf {R} ^ {n}} \| x \| ^ {k} | \partial^ {\beta} f (x)   \partial^ {\alpha - \beta} \varphi (x) | \\ \leqslant \sum_ {0 \leqslant \beta \leqslant \alpha} \binom {\alpha} {\beta} C _ {\beta} \sup _ {x \in \mathbf {R} ^ {n}} \| x \| ^ {k} (1 + \| x \|) ^ {k _ {\beta}} | \partial^ {\alpha - \beta} \varphi (x) |. \end{array}
\]

设 \(\beta \in \mathbf{N}^n\) 满足 \(0 \leqslant \beta \leqslant \alpha\) 。对每个 \(x \in \mathbf{R}^n\) ，我们有

\[
\begin{array}{c} \| x \| ^ {k} (1 + \| x \|) ^ {k _ {\beta}} | \partial^ {\alpha - \beta} \varphi (x) | \leqslant \sup _ {\| x \| \leqslant 1} 2 ^ {k _ {\beta}} | \partial^ {\alpha - \beta} \varphi (x) | \\ + \sup _ {x \in \mathbf {R} ^ {n}} 2 ^ {k _ {\beta}} \| x \| ^ {k + k _ {\beta}} | \partial^ {\alpha - \beta} \varphi (x) | \end{array}
\]

于是最终得到

\[
\widetilde {q} _ {\alpha , k} (f \varphi) \leqslant \sum_ {0 \leqslant \beta \leqslant \alpha} \binom {\alpha} {\beta} 2 ^ {k _ {\beta}} \mathrm{C} _ {\beta} \Bigl (\widetilde {q} _ {\alpha - \beta , 0} (\varphi) + \widetilde {q} _ {\alpha - \beta , k + k _ {\beta}} (\varphi) \Bigr),
\]

由此即得命题。

\subsection{缓增广义函数}

定义 4. --- 我们把 \(\mathbf{R}^{n}\) 上的缓增广义函数空间称为赋以有界收敛拓扑的 \(\mathcal{S}(\mathbf{R})\) 的对偶空间。记作 \(\mathcal{S}'(\mathbf{R}^{n})\) 。

由于 \(\mathcal{S}(\mathbf{R}^n)\) 是有界型的，空间 \(\mathcal{S}'(\mathbf{R}^n)\) 是完备的（TVS, III, p. 24，推论 1）。由于 \(\mathcal{S}(\mathbf{R}^n)\) 是蒙泰尔空间，\(\mathcal{S}'(\mathbf{R}^n)\) 亦然（EVT, IV, p. 19，命题 9）。于是空间 \(\mathcal{S}'(\mathbf{R}^n)\) 是自反的（EVT, IV, p. 16，定理 2）。

引理 9. --- 从 \(\mathcal{S}(\mathbf{R}^{n})\) 到 \(\mathbf{C}\) 中的线性映射 f 是缓增广义函数当且仅当：对每个族 \((\mathrm{M}_{\alpha,k})_{(\alpha,k)\in\mathbf{N}^{n}\times\mathbf{N}}\) （\(\mathbf{R}_{+}\) 中的），线性形式 f 在由这样的函数 \(\varphi\in\mathcal{S}(\mathbf{R}^{n})\) ——满足 \(\widetilde{q}_{\alpha,k}(\varphi)\leqslant\mathrm{M}_{\alpha,k}\) （对所有 \((\alpha,k)\in\mathbf{N}^{n}\times\mathbf{N}\) ）——所成之集上有界。

事实上，空间 \(\mathcal{S}(\mathbf{R}^{n})\) 是有界型的（EVT, III, p. 12，命题 2），且 \(\mathcal{S}(\mathbf{R}^{n})\) 的每个有界子集都含于陈述中所描述的有界集之一。

引理 10. --- 设 F 是 \(\mathcal{S}'(\mathbf{R}^n)\) 上具有可数基或包含一个单有界集合的滤子。那么 F 收敛于某个缓增广义函数当且仅当 \(\langle f, \varphi \rangle\) 对每个施瓦茨函数 \(\varphi\) 依 F 收敛。

特别地，缓增广义函数序列 \((f_{m})_{m\in\mathbf{N}}\) 收敛于缓增广义函数 f 当且仅当 \(\langle f_{m},\varphi\rangle\to\langle f,\varphi\rangle\) 对所有 \(\varphi\in\mathcal{S}(\mathbf{R}^{n})\) 成立。

由于 \(\mathcal{S}'(\mathbf{R}^n)\) 是弗雷歇空间，故为桶型空间（EVT, III, p. 25，命题 2 的推论），引理由巴拿赫–斯坦豪斯定理得出（EVT, III, p. 26，定理 1 的推论 3）。

典范单射 \(j\) ——\(\mathcal{D}(\mathbf{R}^{n})\) 到 \(\mathcal{S}(\mathbf{R}^{n})\) 中的——是连续的且其像稠密（IV, p. 212，引理 12），因而 \(j\) 的转置——即缓增广义函数到子空间 \(\mathcal{D}(\mathbf{R}^{n})\) 的限制映射——是从 \(\mathcal{S}'(\mathbf{R}^{n})\) 到 \(\mathcal{D}'(\mathbf{R}^{n})\) 中的单连续线性映射。我们将通过此映射把 \(\mathcal{S}'(\mathbf{R}^{n})\) 与 \(\mathcal{D}'(\mathbf{R}^{n})\) 的一个子空间等同起来。

设 \(\alpha \in \mathbf{N}^n\) 。线性映射 \(\varphi \mapsto \partial^\alpha \varphi\) ——从 \(\mathcal{S}(\mathbf{R}^n)\) 到 \(\mathcal{S}(\mathbf{R}^n)\) 中——是连续的。因此它的转置是从 \(\mathcal{S}'(\mathbf{R}^n)\) 到 \(\mathcal{S}'(\mathbf{R}^n)\) 中的连续线性映射（EVT, IV, p. 6，推论，\(b\) ））。我们用 \(\partial^\alpha\) 表示连续线性映射 \((-1)^{|\alpha|}{}^{t}\partial^\alpha\) ——从 \(\mathcal{S}'(\mathbf{R}^n)\) 到 \(\mathcal{S}'(\mathbf{R}^n)\) 中。该定义与 IV, p. 208 中关于广义函数的定义 2 相容。

设 \(h \in \mathcal{C}^{\infty}(\mathbf{R}^{n})\) 是 \(\partial^{\alpha}h\) （对每个 \(\alpha \in \mathbf{N}^{n}\) ）都具有多项式增长的函数。连续线性映射 \(\varphi \mapsto h\varphi\) （IV, p. 214，命题 14）的转置是 \(\mathcal{S}'(\mathbf{R}^{n})\) 上的连续线性映射，记作 \(f \mapsto hf\) 。
\subsection{将函数解释为缓增广义函数}

定义 5. --- 设 \(\nu\) 是 \(\mathbf{R}^{n}\) 上的测度。如果存在整数 \(r \in \mathbf{N}\) ，使得连续映射 \(x \mapsto (1 + \|x\|)^{-r}\) 关于 \(\nu\) 在 \(\mathbf{R}^{n}\) 上可积，就称它是缓增的。

换言之，测度 \(\nu\) 是缓增的，如果存在 \(r \in \mathbf{N}\) ，使得由 \(x \mapsto \| x\|^{-r}\) 所定义的函数关于 \(\nu\) 在 \(\mathbf{R}^n\) 中单位球的余集上可积。特别地， \(\mathbf{R}^n\) 上的每个有界测度都是缓增的。更一般地，若 \(f\) 是多项式增长的 \(\mu\)-可测函数，且 \(\nu\) 是缓增的，则测度 \(f \cdot \nu\) 是缓增的。

\(\mathcal{M}^{\mathrm{t}}(\mathbf{R}^{n})\) （\(\mathbf{R}^{n}\) 上缓增测度所成的集）是 \(\mathcal{M}(\mathbf{R}^{n};\mathbf{C})\) （\(\mathbf{R}^{n}\) 上复测度所成的空间）的一个向量子空间。

命题 15. --- 设 \(\nu\) 是 \(\mathbf{R}^{n}\) 上的缓增测度。 \(\nu\) 在 \(\mathcal{S}(\mathbf{R}^{n})\) 上的限制是一个缓增广义函数。它为零当且仅当测度 \(\nu\) 为零。

由于 \(\nu\) 是缓增的，存在正整数 \(k\) ，使得映射 \(x \mapsto \| x\|^{-k}\) 关于 \(\nu\) 在 \(\mathbf{R}^n\) 中单位球的余集上可积。对每个施瓦茨函数 \(\varphi \in \mathcal{S}(\mathbf{R}^n)\) ，我们有

\[
| \langle \nu , \varphi \rangle | \leqslant \Bigl (\int_ {\| x \| \leqslant 1} d \nu \Bigr) \widetilde {q} _ {0, 0} (\varphi) + \Bigl (\int_ {\| x \| > 1} \| x \| ^ {- k} d \nu \Bigr) \widetilde {q} _ {0, k} (\varphi),
\]

因此 \(\nu\) 在 \(\mathcal{S}(\mathbf{R}^n)\) 上的限制是一个缓增广义函数。

最后的断言由 IV, p. 206 的命题 6 得出，因为 \(\mathcal{D}(\mathbf{R}^n)\) 包含在 \(\mathcal{S}(\mathbf{R}^n)\) 中。

我们将把空间 \(\mathcal{M}^{\mathrm{t}}(\mathbf{R}^{n})\) 等同于 \(\mathcal{S}'(\mathbf{R}^{n})\) 的一个子空间。

命题 16. --- 设 \(p \in [1, +\infty]\) ， \(f \in \mathcal{L}^{p}(\mathbf{R}^{n})\) 。那么测度 \(f \cdot \mu\) （以 \(f\) 为密度，关于勒贝格测度）是缓增的。映射 \(f \mapsto f \cdot \mu\) （从 \(\mathrm{L}^{p}(\mathbf{R}^{n})\) 到 \(\mathcal{S}'(\mathbf{R}^{n})\) 中的）是连续的。

设 \(q\) 是 \(p\) 的共轭指数，并设 \(r \geqslant 0\) 满足 \(rq > n\) 。对每个 \(x \in \mathbf{R}^n\) ，令 \(g(x) = (1 + \|x\|)^{-r}\) 。函数 \(g\) 属于 \(\mathcal{L}^q(\mathbf{R}^n)\) （由 IV, p. 199 的命题 3）。由赫尔德不等式，我们有

\[
\int_ {\mathbf {R} ^ {n}} ^ {*} (1 + \| x \|) ^ {- r} | f (x) | d \mu (x) \leqslant \mathrm{N} _ {q} (g) \mathrm{N} _ {p} (f) <   + \infty ,
\]

因此测度 \(f \cdot \mu\) 是缓增的。

设 \(f \in \mathrm{L}^p(\mathbf{R}^n)\) ， \(\varphi \in \mathcal{S}(\mathbf{R}^n)\) 。赫尔德不等式蕴含

\[
| \langle f \cdot \mu , \varphi \rangle | = \left| \int_ {\mathbf {R} ^ {n}} f (x) \varphi (x) d \mu (x) \right| \leqslant \| f \| _ {p} \| \varphi \| _ {q}
\]

而映射 \(f \mapsto f \cdot \mu\) 的连续性则由 IV, p. 213 的命题 13 得出。

对每个 \(p \in [1, +\infty]\) ，我们将把 \(\mathrm{L}^p(\mathbf{R}^n)\) 等同于 \(\mathcal{S}'(\mathbf{R}^n)\) 的一个子空间（借助线性映射 \(f \mapsto f \cdot \mu\) ）。

命题 17. --- 空间 \(\mathcal{D}(\mathbf{R}^{n})\) 与 \(\mathcal{S}(\mathbf{R}^{n})\) 在 \(\mathcal{S}'(\mathbf{R}^{n})\) 中稠密。

只须证明 \(\mathcal{D}(\mathbf{R}^{n})\) 在 \(\mathcal{S}'(\mathbf{R}^{n})\) 中稠密。设 \(\lambda\) 是 \(\mathcal{S}'(\mathbf{R}^{n})\) 上的连续线性形式，它在 \(\mathcal{D}(\mathbf{R}^{n})\) 上为零。由于空间 \(\mathcal{S}(\mathbf{R}^{n})\) 是自反的，存在函数 \(\varphi \in \mathcal{S}(\mathbf{R}^{n})\) ，使得对所有 \(\lambda(f) = \langle f, \varphi \rangle\) 有 \(f \in \mathcal{S}'(\mathbf{R}^{n})\) 。于是我们有

\[
0 = \lambda (\psi) = \langle \psi , \varphi \rangle = \int_ {\mathbf {R} ^ {n}} \psi \varphi d \mu
\]

对所有 \(\psi\in\mathcal{D}(\mathbf{R}^{n})\) 成立。因此 \(\varphi\cdot\mu\) 上的测度 \(\mathbf{R}^{n}\) 为零（IV, p. 206 的命题 6），从而 \(\varphi=0\) 。命题于是由哈恩–巴拿赫定理（EVT, II, p. 46，推论 1）得出。

\subsection{缓增广义函数的傅里叶变换}

由于每个施瓦茨函数 \(\varphi\) 在 \(\mathbf{R}^n\) 上可积（IV, p. 213 的命题 13），它有傅里叶变换 \(\mathcal{F}(\varphi)\) （相应地，余傅里叶变换 \(\overline{\mathcal{F}}(\varphi)\) ），后者等同于 \(\mathbf{R}^n\) 上由下式定义的连续有界函数

\[
y \mapsto \int_ {\mathbf {R} ^ {n}} \varphi (x) \exp (- 2 i \pi x \cdot y) d \mu (x)
\]

（相应地，函数 \(y \mapsto \int_{\mathbf{R}^n} \varphi(x) \exp(2i\pi x \cdot y) d\mu(x)\) ）。

引理 11. --- 设 \(\varphi\in\mathcal{S}(\mathbf{R}^{n})\) 。函数 \(\mathcal{F}(\varphi)\) 在 \(\mathbf{R}^{n}\) 上是无穷次可微的，且我们有

\[
\mathcal {F} (\mathrm{X} ^ {\alpha} \varphi) = (- 2 i \pi) ^ {- | \alpha |} \partial^ {\alpha} (\mathcal {F} (\varphi)),
\]

\[
\mathcal {F} (\partial^ {\alpha} \varphi) = (2 i \pi) ^ {| \alpha |} X ^ {\alpha} \mathcal {F} (\varphi)
\]

对所有 \(\alpha\) （属于 \(\mathbf{N}^n\) ）成立。

不妨设 \(n \geqslant 1\) 。设 \(\varphi \in \mathcal{S}(\mathbf{R}^{n})\) 。由 \((x, y) \mapsto \varphi(x) \exp(2i\pi x \cdot y)\) （从 \(\mathbf{R}^{n} \times \mathbf{R}^{n}\) 到 \(\mathbf{C}\) 中的函数）对每个整数 k 满足 IV, p. 198 的推论 1 的假设。因此 \(\varphi\) 的傅里叶变换是无穷次可微的，且满足

\[
\partial^ {\alpha} (\mathcal {F} \varphi) (y) = (- 2 i \pi) ^ {| \alpha |} \int_ {\mathbf {R} ^ {n}} x ^ {\alpha} \varphi (x) \exp (- 2 i \pi x \cdot y) d \mu (x)
\]

对所有 \(y \in \mathbf{R}^{n}\) 成立，由此即得第一个公式。

我们来证明第二个公式。对 \(|\alpha|\) 作归纳，只须处理 \(|\alpha| = 1\) 的情形，而易化为 \(\alpha = (1,0,\ldots,0)\) 的情形。把每个 \(x \in \mathbf{R}^n\) 写成 \(x = (x_1,x')\) 的形式，其中 \(x' \in \mathbf{R}^{n-1}\) ，并用 \(\mu_1\) （相应地， \(\mu'\) ）表示 \(\mathbf{R}^{n-1}\) （相应地， \(\mathbf{R}\) ）上的勒贝格测度。由勒贝格–富比尼定理（INT, V, p. 96, § 8, no. 4，定理 1），对每个 \(y = (y_1,y') \in \mathbf{R} \times \mathbf{R}^{n-1}\) ，得到

\[
\begin{array}{l} \mathcal {F} (\partial_ {1} \varphi) (y) = \int_ {\mathbf {R} ^ {n - 1}} \exp (- 2 i \pi   x ^ {\prime} \cdot y ^ {\prime}) \\ \qquad \times \Big (\int_ {\mathbf {R}} (\partial_ {1} \varphi) (x _ {1}, x ^ {\prime}) \exp (- 2 i \pi   x _ {1} y _ {1}) d \mu_ {1} (x _ {1}) \Big) d \mu^ {\prime} (x ^ {\prime}). \end{array}
\]

对每个紧区间 \([a,b]\) （R 中的）以及每个 \(x'\in R^{n-1}\) ，我们有

\[
\begin{array}{r l} & {\int_ {a} ^ {b} (\partial_ {1} \varphi) (x _ {1}, x ^ {\prime}) \exp (- 2 i \pi x _ {1} y _ {1}) d \mu_ {1} (x _ {1}) =} \\ & {\qquad \left[ \varphi (x _ {1}, x ^ {\prime}) \exp (- 2 i \pi x _ {1} y _ {1}) \right] _ {a} ^ {b}} \\ & {\qquad + 2 i \pi y _ {1} \int_ {a} ^ {b} \varphi (x _ {1}, x ^ {\prime}) \exp (- 2 i \pi x _ {1} y _ {1}) d \mu_ {1} (x _ {1})} \end{array}
\]

由分部积分（FVR, II, p. 10，公式 (10)）。当 \(a\) 趋于 \(-\infty\) 且 \(b\) 趋于 \(+\infty\) 时，右端第一项趋于 0，因为 \(\varphi \in \mathcal{S}(\mathbf{R}^n)\) 。第二项由勒贝格定理（INT, IV, p. 137, § 3, no. 7，定理 6）收敛于

\[
2 i \pi y _ {1} \int_ {\mathbf {R}} \varphi (x _ {1}, x ^ {\prime}) \exp (- 2 i \pi x _ {1} y _ {1}) d \mu_ {1} (x _ {1}),
\]

因为映射 \(x_{1} \mapsto \varphi(x_{1}, x')\) 在 R 上可积。由于 \(x_{1} \mapsto \partial_{1}\varphi(x_{1}, x')\) 也在 R 上可积，由此得出

\[
\begin{array}{r l r} & & {\int_ {\mathbf {R}} \partial_ {1} \varphi (x _ {1}, x ^ {\prime}) \exp (- 2 i \pi x _ {1} y _ {1}) d \mu_ {1} (x _ {1})} \\ & & {= 2 i \pi y _ {1} \int_ {\mathbf {R}} \varphi (x _ {1}, x ^ {\prime}) \exp (- 2 i \pi x _ {1} y _ {1}) d \mu_ {1} (x _ {1})} \end{array}
\]

最后，再次应用勒贝格–富比尼定理，我们得出

\[
\mathcal {F} (\partial_ {1} \varphi) (y) = 2 i \pi y _ {1} \int_ {\mathbf {R} ^ {n}} \varphi (x) \exp (- 2 i \pi x \cdot y) d \mu (x),
\]

这正是所要证明的。

命题 18. --- 傅里叶变换在 \(\mathcal{S}(\mathbf{R}^{n})\) 上的限制是拓扑向量空间的一个自同构，其逆是余傅里叶变换的限制。

设 \(\varphi\in\mathcal{S}(\mathbf{R}^{n})\) 。由前一引理， \(\varphi\) 的傅里叶变换属于 \(\mathcal{C}^{\infty}(\mathbf{R}^{n})\) 。此外，对 \(\alpha\in\mathbf{N}^{n}\) 与 \(\beta\in\mathbf{N}^{n}\) ，我们有

\[
\begin{array}{r} \mathrm{X} ^ {\beta} \partial^ {\alpha} (\mathcal {F} (\varphi)) = (- 2 i \pi) ^ {| \alpha |} \mathrm{X} ^ {\beta} \mathcal {F} (\mathrm{X} ^ {\alpha} \varphi) \\ = (- 1) ^ {| \alpha |} (2 i \pi) ^ {| \alpha | - | \beta |} \mathcal {F} (\partial^ {\beta} (\mathrm{X} ^ {\alpha} \varphi)). \end{array}
\]

特别地，函数 \(\mathrm{X}^{\beta}\partial^{\alpha}(\mathcal{F}(\varphi))\) 有界。由于 \(\alpha\) 与 \(\beta\) 在 \(\mathbf{N}^{n}\) 中任意，这表明 \(\mathcal{F}(\varphi)\) 属于 \(\mathcal{S}(\mathbf{R}^{n})\) 。此外，这一计算蕴含

\[
q _ {\alpha , \beta} (\mathcal {F} (\varphi)) \leqslant (2 \pi) ^ {| \alpha | - | \beta |} \| \partial^ {\beta} (X ^ {\alpha} \varphi) \| _ {1}
\]

对 \((\alpha, \beta) \in \mathbf{N}^n \times \mathbf{N}^n\) 与 \(\varphi \in \mathcal{S}(\mathbf{R}^n)\) 成立。

由于空间 \(\mathcal{S}(\mathbf{R}^{n})\) 在 \(\mathrm{L}^{1}(\mathbf{R}^{n})\) 中的包含映射是连续的（IV, p. 213 的命题 13），映射 \(q_{\alpha,\beta} \circ \mathcal{F}\) （从 \(\mathcal{S}(\mathbf{R}^{n})\) 到 \(\mathbf{R}\) 中的）是连续的。由此得出，傅里叶变换是从空间 \(\mathcal{S}(\mathbf{R}^{n})\) 到其自身的连续映射（参见 EVT, II, p. 7，命题 5, c)）。类似地可以验证，余傅里叶变换是从空间 \(\mathcal{S}(\mathbf{R}^{n})\) 到其自身的连续映射。由傅里叶反演公式（II, p. 222 的定理 3），傅里叶变换与余傅里叶变换是互逆的同构。

定义 6. --- \(\mathcal{S}'(\mathbf{R}^n)\) 上的傅里叶变换（相应地，余傅里叶变换）称为 \(\mathcal{S}(\mathbf{R}^n)\) 上的傅里叶变换（相应地，余傅里叶变换）的转置。

我们仍用 \(\mathcal{F}\) （相应地， \(\overline{\mathcal{F}}\) ）表示 \(\mathcal{S}'(\mathbf{R}^n)\) 上的傅里叶变换（相应地，余傅里叶变换）。 \(\mathcal{S}'(\mathbf{R}^n)\) 上的傅里叶变换于是是一个拓扑向量空间自同构，其逆是余傅里叶变换。对每个 \(f \in \mathcal{S}'(\mathbf{R}^n)\) ，缓增广义函数 \(\mathcal{F}(f)\) （相应地， \(\overline{\mathcal{F}}(f) )\) ）由下述公式定义

\[
\langle \mathcal {F} (f), \varphi \rangle = \langle f, \mathcal {F} (\varphi) \rangle \quad (\text { resp. } \langle \overline {{\mathcal {F}}} (f), \varphi \rangle = \langle f, \overline {{\mathcal {F}}} (\varphi) \rangle)
\]

对所有 \(\varphi\in\mathcal{S}(\mathbf{R}^{n})\) 成立。

命题 19. --- 设 f 是与有界测度 \(\nu \in \mathcal{M}^{1}(\mathbf{R}^{n})\) （相应地，与 \(g \in \mathrm{L}^{2}(\mathbf{R}^{n})\) ）相伴的缓增广义函数。 f 在 \(\mathcal{S}'(\mathbf{R}^{n})\) 中的傅里叶变换是与测度 \(\nu\) 的傅里叶变换（相应地，与 g 的傅里叶变换）相伴的缓增广义函数。

设 \(\nu\) 是 \(\mathbf{R}^{n}\) 上的有界测度， f 是与 \(\nu\) 相伴的缓增广义函数。傅里叶变换 \(\mathcal{F}(\nu)\) 是 \(\mathbf{R}^{n}\) 上的连续有界函数（II, p. 207 的命题 3）。与这一函数相伴的缓增广义函数满足

\[
\langle \mathcal {F} (\nu), \varphi \rangle = \int_ {\mathbf {R} ^ {n}} \mathcal {F} (\nu) \varphi d \mu = \int_ {\mathbf {R} ^ {n}} \mathcal {F} (\varphi) d \nu = \langle f, \mathcal {F} (\varphi) \rangle = \langle \mathcal {F} (f), \varphi \rangle
\]

对所有 \(\varphi\in\mathcal{S}(\mathbf{R}^{n})\) 成立，其中第二个等式就是 II, p. 221 的命题 13，由于测度 \(\varphi\cdot\mu\) 有界，它在此适用。因此与 \(\mathcal{F}(\nu)\) 相伴的缓增广义函数是 \(\mathcal{F}(f)\) 。

当 \(f\) 是与 \(g \in \mathrm{L}^2(\mathbf{R}^n)\) 相伴的缓增广义函数时，我们用完全类似的方法，借助 II, p. 221 的公式 (29) 处理。

对余傅里叶变换也有类似的命题。

注. --- 关于测度的傅里叶变换的初等公式对缓增广义函数的傅里叶变换仍然有效。例如，对 \(f \in \mathcal{S}'(\mathbf{R}^{n})\) 与 \(\alpha \in \mathbf{N}^{n}\) ，我们有

\[
\begin{array}{l} \mathcal {F} (\partial^ {\alpha} f) = (2 i \pi) ^ {| \alpha |} X ^ {\alpha} \mathcal {F} (f) \\ \mathcal {F} (X ^ {\alpha} f) = (- 2 i \pi) ^ {- | \alpha |} \partial^ {\alpha} (\mathcal {F} (f)) \end{array}
\]

这由引理 11 得出，

\subsection{向量空间上的广义函数与缓增广义函数}

设 \(u\) 是从 \(\mathbf{R}^{n}\) 到 \(\mathbf{R}^{n}\) 中的双射线性映射。映射 \(\varphi \mapsto \varphi \circ u\) 是 \(\mathcal{S}(\mathbf{R}^{n})\) （相应地， \(\mathcal{D}(\mathbf{R}^{n})\) ）的自同构；它的转置是 \(\mathcal{S}'(\mathbf{R}^{n})\) （相应地， \(\mathcal{D}'(\mathbf{R}^{n})\) ）的自同构。

设 E 是 n 维实向量空间。设 \(v: \mathbf{R}^{n} \to E\) 是向量空间同构。我们用 \(\mathcal{S}(E)\) （相应地， \(\mathcal{D}(E)\) ）表示映射 \(\varphi: E \to \mathbf{C}\) 中满足 \(\varphi \circ v \in \mathcal{S}(\mathbf{R}^{n})\) （相应地，满足 \(\varphi \circ v \in \mathcal{D}(\mathbf{R}^{n})\) ）者所成的集。由前面的注，这些空间不依赖于 v 的选取；它们同构于 \(\mathcal{S}(\mathbf{R}^{n})\) （相应地， \(\mathcal{D}(\mathbf{R}^{n})\) ）。我们用 \(\mathcal{S}'(E)\) （相应地， \(\mathcal{D}'(E)\) ）表示 \(\mathcal{S}(E)\) （相应地， \(\mathcal{D}(E)\) ）的对偶，赋有有界收敛拓扑。它是同构于 \(\mathcal{S}'(\mathbf{R}^{n})\) （相应地，同构于 \(\mathcal{D}'(\mathbf{R}^{n})\) ）的拓扑向量空间。

设 E 与 F 是 n 维实向量空间，它们关于双线性形式 \(b: E \times F \to \mathbf{R}\) 处于对偶。局部紧交换群 E 关于映射 

\[
(x, y) \mapsto \exp (2 i \pi b (x, y))
\]

（从 E × F 到 U 中；参见 II, p. 235 的推论 1）与 F 处于对偶。我们给 E 与 F 赋以关于这一映射互为对偶的哈尔测度。

空间 \(\mathcal{S}(\mathrm{E})\) 包含在 \(\mathrm{L}^1 (\mathrm{E})\) 中； E 的傅里叶变换通过过渡到子空间与对偶，诱导出从 \(\mathcal{S}(\mathrm{E})\) 到 \(\mathcal{S}(\mathrm{F})\) 中的拓扑向量空间同构，其转置是从 \(\mathcal{S}'(\mathrm{F})\) 到 \(\mathcal{S}'(\mathrm{E})\) 中的拓扑向量空间同构。

\subsection{索伯列夫空间}

设 U 是 \(\mathbf{R}^{n} \) 的开子集。设 p 是实数 \(\geqslant 1\) ， k 是自然数。我们用 \(W^{k,p}(U)\) 表示满足如下条件的广义函数 \(f \in \mathcal{D}'(U)\) 所成的空间：对每个 \(\alpha \in N^{n}\) （满足 \(|\alpha| \leqslant k\) ），广义函数 \(\partial^{\alpha}f\) 与 \(\mathrm{L}^{p}(U) \) 的一个元素相伴。

特别地，当 \(U = \mathbf{R}^{n}\) 时， \(W^{k,p}(U)\) 中的元素是缓增广义函数。

从 \(\mathrm{W}^{k,p}(\mathrm{U})\) 到 \(\mathbf{R}_{+}\) 中的、把 \(f\in \mathrm{W}^{k,p}(\mathrm{U})\)

\[
\| f \| _ {k, p} = \Bigl (\sum_ {| \alpha | \leqslant k} \| \partial^ {\alpha} f \| _ {p} ^ {p} \Bigr) ^ {1 / p}
\]

是 \(W^{k,p}(U)\) 上的一个范数。空间 \(W^{k,p}(U)\) 将始终赋有这一范数；这一赋范空间称为指标为 k 、指数为 p 的索伯列夫空间。

空间 \(\mathcal{D}(\mathrm{U})\) 包含在 \(\mathrm{W}^{k,p}(\mathrm{U})\) 中。我们用 \(\mathrm{W}_0^{k,p}(\mathrm{U})\) 表示 \(\mathcal{D}(\mathrm{U})\) 在 \(\mathrm{W}^{k,p}(\mathrm{U})\) 中的闭包。它是 \(\mathrm{W}^{k,p}(\mathrm{U})\) 的一个闭子空间。

有 \(\mathrm{W}_{0}^{k,p}(\mathbf{R}^{n})=\mathrm{W}^{k,p}(\mathbf{R}^{n})\) ，但空间 \(\mathrm{W}^{k,p}(\mathrm{U})\) 与 \(\mathrm{W}_{0}^{k,p}(\mathrm{U})\) 一般是不同的（参见 IV, p. 334 的习题 12 与 IV, p. 334 的习题 14）。

我们也记 \(\mathrm{H}^k (\mathrm{U}) = \mathrm{W}^{k,2}(\mathrm{U})\) ， \(\mathrm{H}_0^k (\mathrm{U}) = \mathrm{W}_0^{k,2}(\mathrm{U})\) 。

 \(\mathrm{H}^{k}(\mathrm{U})\) 的范数是准希尔伯特范数，它由 \(\mathrm{H}^{k}(\mathrm{U})\) 上的、由下式定义的正埃尔米特形式所导出。

\[
\left(f _ {1}, f _ {2}\right) \mapsto \sum_ {| \alpha | \leqslant k} \int_ {\mathrm{U}} \overline {{\partial^ {\alpha} f _ {1}}} \partial^ {\alpha} f _ {2} d \mu .
\]

希尔伯特空间 \(\mathrm{H}^k (\mathrm{U})\) 与 EVT, V, p. 6 例 (3) 中记作 \(\mathcal{H}^k\) 的空间重合。

按定义有 \(\mathrm{W}^{0,p}(\mathrm{U})=\mathrm{L}^{p}(\mathrm{U})\) ， \(\mathrm{H}^{0}(\mathrm{U})=\mathrm{L}^{2}(\mathrm{U})\) ；此外，由 IV, p. 202 的命题 4, b)， \(\mathrm{W}_{0}^{0,p}(\mathrm{U})=\mathrm{L}^{p}(\mathrm{U})\) 。

命题 20. — 索伯列夫空间 \(W^{k,p}(U)\) 与 \(W_{0}^{k,p}(U)\) 是可数型的巴拿赫空间。特别地，空间 \(H^{k}(U)\) 与 \(H_{0}^{k}(U)\) 是可数型的希尔伯特空间。

只须证明关于 \(\mathrm{W}^{k,p}(\mathrm{U})\) 的断言。

设 I 是满足 \(\alpha \in \mathbf{N}^n\) 的 \(|\alpha| \leqslant k\) 所成的集。线性映射 \(u\) （从 \(\mathrm{W}^{k,p}(\mathrm{U})\) 到 \(\mathrm{L}^p (\mathrm{U})^{\mathrm{I}}\) 中的）把 \(f\) 对应到族 \((\partial^{\alpha}f)_{\alpha \in \mathrm{I}}\) ，它是单射；按 \(\mathrm{W}^{k,p}(\mathrm{U})\) 上范数的定义，它是连续且严格的。为证 \(\mathrm{W}^{k,p}(\mathrm{U})\) 完备，只须证它在 \(u\) 下的像是闭的。设 \((f_n)_{n \in \mathbf{N}}\) 是 \(\mathrm{W}^{k,p}(\mathrm{U})\) 中的序列，使得 \((u(f_n))_{n \in \mathbf{N}}\) 收敛。设 \((g_{\alpha})_{\alpha \in \mathrm{I}} \in \mathrm{L}^p (\mathrm{U})^{\mathrm{I}}\) 是它的极限。对 \(\alpha \in \mathrm{I}\) ，序列 \((\partial^{\alpha}f_n)_{n \in \mathbf{N}}\) 在 \(\mathrm{L}^p (\mathrm{U})\) 中收敛于 \(g_{\alpha}\) 。当然，收敛也在 \(\mathcal{D}'(\mathrm{U})\) 中发生。令 \(f = g_0\) 。对每个 \(\varphi \in \mathcal{D}(\mathrm{U})\) ，得到

\[
\begin{array}{r l} & {\langle \partial^ {\alpha} f, \varphi \rangle = (- 1) ^ {| \alpha |} \langle f, \partial^ {\alpha} \varphi \rangle} \\ & {\qquad = \lim _ {n \to + \infty} (- 1) ^ {| \alpha |} \langle f _ {n}, \partial^ {\alpha} \varphi \rangle = \lim _ {n \to + \infty} \langle \partial^ {\alpha} f _ {n}, \varphi \rangle = \langle g _ {\alpha}, \varphi \rangle .} \end{array}
\]

这就证明了 \(g_{\alpha} = \partial^{\alpha} f\) 对所有 \(\alpha \in I\) 成立，从而 \((g_{\alpha})_{\alpha \in I} = u(f)\) 属于 u 的像。

空间 \(\mathrm{W}^{k,p}(\mathrm{U})\) 由 u 等同于空间 \(\mathrm{L}^{p}(\mathrm{U})^{\mathrm{I}}\) 的一个子空间；后者是可数型的（IV, p. 180 的命题 2 与 TG, IX, p. 19，cor., (ii))），因而 \(\mathrm{W}^{k,p}(\mathrm{U})\) 也是可数型的（loc. cit., (i)）。

命题 21. --- 设 N 是 \(\mathbf{R}^{n}\) 上的欧几里得范数。设 \(k\) 是整数 \(\geqslant 0\) 。索伯列夫空间 \(\mathrm{H}^{k}(\mathbf{R}^{n})\) 是这样的 \(f \in \mathcal{S}'(\mathbf{R}^{n})\) 所成的空间： \((1 + \mathrm{N}^{k})\mathcal{F}(f)\) 属于 \(\mathrm{L}^{2}(\mathbf{R}^{n})\) 。

我们对 \(k\) 作归纳。当 \(k = 0\) 时，结论由普朗谢雷尔定理（II, p. 215，定理 1）得出。设 \(k = 1\) 。对 \(f \in \mathcal{S}'(\mathbf{R}^n)\) ，我们有 \((1 + \mathrm{N})\mathcal{F}f \in \mathrm{L}^2(\mathbf{R}^n)\) 当且仅当 \(f \in \mathrm{L}^2(\mathbf{R}^n)\) 且 \(\mathrm{N}\mathcal{F}f \in \mathrm{L}^2(\mathbf{R}^n)\) 。此外， \(\mathrm{N}\mathcal{F} \in \mathrm{L}^2(\mathbf{R}^n)\) 当且仅当：对每个 \(\alpha \in \mathbf{N}^n\) （满足 \(|\alpha| = 1\) ），有 \(\mathrm{X}^\alpha\mathcal{F}f \in \mathrm{L}^2(\mathbf{R}^n)\) 。由于 \(\mathcal{F}(\partial^\alpha f) = 2i\pi\mathrm{X}^\alpha\mathcal{F}f\) ，这一条件意味着对每个 \(\mathcal{F}(\partial^\alpha f) \in \mathrm{L}^2(\mathbf{R}^n)\) （满足 \(\alpha\) ）有 \(|\alpha| = 1\) ，亦即对每个 \(\partial^\alpha f \in \mathrm{L}^2(\mathbf{R}^n)\) （满足 \(\alpha\) ）有 \(|\alpha| = 1\) 。由此得出断言对 \(k = 1\) 为真。

现在设 \(k \geqslant 2\) ，并且关于 \(\mathrm{H}^{\ell}(\mathbf{R}^{n})\) 的断言对每个正整数 \(\ell \leqslant k - 1\) 成立。设 \(f \in \mathcal{S}'(\mathbf{R}^{n})\) 。按定义， \(f \in \mathrm{H}^{k}(\mathbf{R}^{n})\) 当且仅当 \(f \in \mathrm{L}^{2}(\mathbf{R}^{n})\) ，且对每个 \(\beta \in \mathbf{N}^{n}\) （满足 \(|\beta| \leqslant 1\) ），广义函数 \(\partial^{\beta} f\) 属于 \(\mathrm{H}^{k-1}(\mathbf{R}^{n})\) 。由归纳假设，这等价于 \(f \in \mathrm{L}^{2}(\mathbf{R}^{n})\) ，且对每个 \((1 + \mathrm{N}^{k-1})\mathcal{F}(\partial^{\beta} f) \in \mathrm{L}^{2}(\mathbf{R}^{n})\) （满足 \(\beta \in \mathbf{N}^{n}\) ）有 \(|\beta| \leqslant 1\) 。由于 \(\mathcal{F}(\partial^{\beta} f) = (2i\pi X)^{|\beta|}\mathcal{F}(f)\) ，条件 \(f \in \mathrm{H}^{k}(\mathbf{R}^{n})\) 等价于说： \(\mathcal{F}f \in \mathrm{L}^{2}(\mathbf{R}^{n})\) ，且对每个 \((1 + \mathrm{N}^{k-1})\mathrm{X}^{\beta}\mathcal{F}f \in \mathrm{L}^{2}(\mathbf{R}^{n})\) （满足 \(\beta \in \mathbf{N}^{n}\) ）有 \(|\beta| \leqslant 1\) 。

不等式

\[
1 + \mathrm{N}^{k}\leqslant 1 + \mathrm{N}^{k - 1}\sum_{\substack{\beta \in \mathbf{N}^{n}\\ |\beta |\leqslant 1}}|\mathrm{X}^{\beta}|\leqslant 1 + n^{1 / 2}\mathrm{N}^{k}\leqslant (1 + n^{1 / 2})(1 + \mathrm{N}^{k})
\]

于是蕴含 \(f \in \mathrm{H}^k(\mathbf{R}^n)\) 当且仅当 \((1 + \mathrm{N}^k)\mathcal{F} \in \mathrm{L}^2(\mathbf{R}^n)\) 。

\section{部分算子}

\subsection{部分算子}

在本节中， K 是交换域。

我们回顾（E, II, §3, p. 9–10）：\(\text{graph}^{(2)}\) 是这样一个集合，其元素都是有序偶。设 A 与 B 是集合， A 与 B 之间的一个 \(\text{correspondence}\) 是一个三元组 \((\Gamma, A, B)\) ，其中 \(\Gamma\) 是包含在 A × B 中的一个图像；它的定义集（也称为 \(\text{domain}\) ）是 \(\text{pr}_1(\Gamma)\) ，而它的值集是 \(\text{pr}_2(\Gamma)\) 。一个对应是一个 \(\text{function}\) （E, II, p. 13，定义 9），如果它的图像是泛函的，且它的出发集与它的定义域重合。泛函图像的每个子集都是泛函图像。

定义 1. --- 设 E 与 F 是 K 上的向量空间。从 E 到 F 的部分算子 u 是 E 与 F 之间的一个对应 (Γ, E, F) ，它满足下述条件：

\begin{enumerate}
\def\labelenumi{(\roman{enumi})}
\item
  图像 \(\Gamma\) 是 \(\mathrm{E} \times \mathrm{F}\) 的一个向量子空间；
\item
  图像 \(\Gamma\) 是泛函的。
\end{enumerate}

若 E = F ，就说 u 是 E 上的部分算子。

设 u 是从 E 到 F 的部分算子。对应 u 的图像 \(\Gamma\) 称为部分算子 u 的图像，也记作 \(\Gamma_{u}\) 。我们用 \(\mathcal{P}(\mathrm{E};\mathrm{F})\) 表示从 E 到 F 的部分算子所成的集；我们简记 \(\mathcal{P}(\mathrm{E}) = \mathcal{P}(\mathrm{E};\mathrm{E})\) 。

给出一个从 E 到 F 的部分算子，相当于给出 E 的一个向量子空间 D 与一个从 D 到 F 的线性映射 u ；相应的部分算子是对应 \((\Gamma, \mathrm{E}, \mathrm{F})\) ，其中 \(\Gamma \subset D \times F\) 是 u 的图像。

部分算子 \(u\) 的定义域简称为 \(u\) 的定义域，并记作 \(\operatorname{dom}(u)\) 。

每个从 E 到 F 的线性映射 u 都是从 E 到 F 的部分算子。

若 D ⊂ E 是向量子空间，我们将用 1 \(_{D}\) 表示以 D 为定义域、在 D 上为恒等映射的部分算子，即对应 (Δ \(_{D}\) , E, E) ，其中 Δ \(_{D}\) 是 D × D 的对角线（E, II, p. 13，定义 8）。我们将用 \(0_{\mathrm{D}}\) 表示以 D 为定义域、在 D 上为零的部分算子，即对应 \((\mathrm{D} \times \{0\}, \mathrm{E}, \mathrm{F})\) 。

部分算子 \(u = (\Gamma, \mathrm{E}, \mathrm{F})\) 与 \(u' = (\Gamma', \mathrm{E}, \mathrm{F})\) （都是从 E 到 F 的）相等，当且仅当 \(\operatorname{dom}(u) = \operatorname{dom}(u')\) ，且从 \(u\) 到 F 中的线性映射 \(u'\) 与 \(\operatorname{dom}(u)\) 重合。

依照 E, II, §3 ，定义下列概念：

\begin{enumerate}
\def\labelenumi{(\roman{enumi})}
\tightlist
\item
  u 下 E 的子集 A 的像是 F 的子集 \(u(A \cap D)\) ；简记为 \(u(A)\) 。 u 下 F 的子集 B 的原像是 D 的子集 \(\overline{u}(B)\) 。
\end{enumerate}

若 A （相应地， B ）是 E （相应地， F ）的向量子空间，则它在 u 下的像（相应地，原像）是 F （相应地， E ）的向量子空间。

 \(u\) 的像是 \(u(\mathrm{D})\) 的向量子空间 \(\mathrm{F}\) ，也记作 \(\operatorname{Im}(u)\) 。称 \(u\) 为满射部分算子，如果 \(\operatorname{Im}(u) = \mathrm{F}\) 。 \(u\) 的核是 \(u^{-1}(\{0\})\) 的向量子空间 \(\mathrm{E}\) ，也记作 \(\operatorname{Ker}(u)\) 。 \(u\) 的核化为 0 ，当且仅当从 \(u\) （ \(\operatorname{dom}(u)\) 到 \(\mathrm{F}\) 中的）线性映射是单射的。这时称 \(u\) 是单射的。若 \(u\) 既单射又满射，就称它是双射的。

\begin{enumerate}
\def\labelenumi{(\roman{enumi})}
\setcounter{enumi}{1}
\item
  若 E 、 F 与 G 是 K 上的向量空间，且 \(u = (\Gamma, \mathrm{E}, \mathrm{F})\) ， \(v = (\Gamma', \mathrm{F}, \mathrm{G})\) 分别是从 E 到 F 与从 F 到 G 的部分算子，则复合对应 \(v \circ u = (\Gamma' \circ \Gamma, \mathrm{E}, \mathrm{G})\) 是从 E 到 G 中的部分算子。它的定义域是 \(u^{-1}(\mathrm{dom}(v))\) 。若 H 是 K 上的向量空间，且 \(w = (\Gamma'', \mathrm{G}, \mathrm{H})\) 是从 G 到 H 中的部分算子，则有 \(w \circ (v \circ u) = (w \circ v) \circ u\) 。有时把 \(v \circ u\) 简写作 vu 。
\item
  特别地，对每个从 E 到 F 中的部分算子 \(u\) 以及每个 \(a \in \mathrm{K}\) ，部分算子 \(au = (a1_{\mathrm{F}}) \circ u\) 与 \(ua = u \circ (a1_{\mathrm{E}})\) 有定义。若 \(a \neq 0\) 或者 \(u\) 的定义域等于 E ，则二者相等；我们有 \(u0 = 0_{\mathrm{E}}\) 与 \(0u = 0_{\mathrm{dom}(u)}\) 。
\end{enumerate}

设 E 是向量空间。由前所述，集合 \(\mathcal{P}(\mathrm{E})\) （赋有由 \((u,v)\mapsto u\circ v\) 所定义的复合律）是一个结合的含幺群胚（A, I, p. 4，定义 5 与 A, I, p. 12，定义 2），其单位元为 \(1_{\mathrm{E}}\) 。对每个 \(n\in \mathbf{N}\) ，我们用 \(u^n\) 表示复合 \({}^n u\) （A, I, p. 13）。

此外，我们定义下列概念：

\begin{enumerate}
\def\labelenumi{(\roman{enumi})}
\item
  若 \(u = (\Gamma, \mathrm{E}, \mathrm{F})\) 是从 E 到 F 的部分算子， G 是 E 的向量子空间，则 \(u\) 到 G 上的约化是从 E 到 F 的部分算子（ \(\Gamma \cap (\mathrm{G} \times \mathrm{F}), \mathrm{E}, \mathrm{F}\) ）。它的定义域是 \(\operatorname{dom}(u) \cap \mathrm{G}\) ；当不致与 \(u|_{\mathrm{G}}\) 到子空间 G 上的限制相混淆时，有时把它记作 \(u\) 。
\item
  设 \(v\) 是从 F 到 E 的单射部分算子。这时 \(v^{-1} = (\Gamma^{-1}, \mathrm{E}, \mathrm{F})\) 的逆对应 \(v\) 是一个部分算子，满足 \(\operatorname{dom}(v^{-1}) = \operatorname{Im}(v)\) 。称 \(v^{-1}\) 为 \(v\) 的逆部分算子。我们有等式 \(v \circ v^{-1} = 1_{\operatorname{dom}(v^{-1})}\) （在 \(\mathcal{P}(\mathrm{E})\) 中）与 \(v^{-1} \circ v = 1_{\operatorname{dom}(v)}\) （在 \(\mathcal{P}(\mathrm{F})\) 中）。部分算子 \(v^{-1}\) 是单射的，且有 \((v^{-1})^{-1} = v\) 。
\item
  设 E 、 F 与 G 是向量空间。设 u （相应地， v ）是从 E 到 F （相应地，从 F 到 G ）的单射部分算子。那么部分算子 \(v \circ u\) 是单射的，且 \((v \circ u)^{-1} = u^{-1} \circ v^{-1}\) 。
\item
  若 u 与 v 是从 E 到 F 的部分算子，我们说 v 是 u 的扩张，并记 \(u \subset v\) ，如果 u 的图像包含在 v 的图像中。这蕴含 \(\text{dom}(u) \subset \text{dom}(v)\) ，且 u 是 v 到 \(\text{dom}(u)\) 上的约化。关系 “ \(u \subset v\) ” 是 \(\mathcal{P}(\text{E};\text{F})\) 上的一个序关系。例如，对所有 \(au \subset ua\) 与每个 \(a \in K\) ，有 \(u \in \mathcal{P}(\text{E};\text{F})\) 。
\item
  设 E 是 K 上的向量空间， \((\mathrm{F}_{i})_{i\in\mathrm{I}}\) 是 K 上向量空间的一个族。对 \(i\in I\) ，设 \(u_{i}\) 是从 E 到 \(F_{i}\) 的一个部分算子。诸 \(u_{i}\) 的乘积部分算子是从 E 到诸空间 \(F_{i}\) 的乘积向量空间中的部分算子，其定义域是诸 \(\mathrm{dom}(u_{i})\) 的交 D ，它把 \(x\in D\) 对应到族 \((u_{i}(x))_{i\in\mathrm{I}}\) 。把它记作 \((u_{i})_{i\in\mathrm{I}}\) 。
\item
  设 A: F×F → F 是线性映射 \((x,y)\mapsto x+y\) 。设 u 与 v 是从 E 到 F 的部分算子。和 \(u+v\) 是从 E 到 F 的部分算子 A ∘ \((u,v)\) 。它的定义域是 dom \((u)\cap\text{dom}(v)\) 。对 \(\mathcal{P}(\text{E};\text{F})\) 中的 u 、 v 、 w ，有 \((u+v)+w=u+(v+w)\) 。
\end{enumerate}

设 G 是 K 上的向量空间。对 \(\mathcal{P}(\mathrm{E};\mathrm{F})\) 中所有 u 与 v 以及每个 \(w \in \mathcal{P}(\mathrm{F};\mathrm{G})\) ，有 \(w \circ u + w \circ v \subset w \circ (u + v)\) 。一般而言，这一公式中没有等号（IV, p. 344 的习题 1），但当 w 的定义域等于 F 时就是这种情形。对 \(w \in \mathcal{P}(\mathrm{G};\mathrm{E})\) ，有 \(u \circ w + v \circ w = (u + v) \circ w\) 。

\begin{enumerate}
\def\labelenumi{(\roman{enumi})}
\setcounter{enumi}{6}
\tightlist
\item
  设 L 是域 K 的一个扩张。设 E 与 F 是 K 上的向量空间， \(\mathrm{E}_{(\mathrm{L})} = \mathrm{L} \otimes_{\mathrm{K}} \mathrm{E}\) ， \(\mathrm{F}_{(\mathrm{L})} = \mathrm{L} \otimes_{\mathrm{K}} \mathrm{F}\) 是由 K 到 L 的纯量扩张得到的 L-向量空间（A, II, p. 82）。对每个从 E 到 F 的部分算子 u ，我们用 \(u_{(L)}\) 表示从 \(\mathrm{E}_{(\mathrm{L})}\) 到 \(\mathrm{F}_{(\mathrm{L})}\) 中的、其图像为 \(L \otimes_{K} \Gamma_{u}\) （ \(\mathrm{E}_{(\mathrm{L})} \times \mathrm{F}_{(\mathrm{L})}\) 的向量子空间）的部分算子；它的定义域是 \(L \otimes_{K} \operatorname{dom}(u)\) ，且它在这一定义域上与把 \(1 \otimes x\) 映到 \(1 \otimes u(x)\) （对所有 \(x \in \operatorname{dom}(u)\) ）的唯一线性映射重合。
\end{enumerate}

设 v 是从 E 到 F 的部分算子。我们有 \(u \subset v\) 当且仅当 \(u_{(L)} \subset v_{(L)}\) 。

\begin{enumerate}
\def\labelenumi{(\roman{enumi})}
\setcounter{enumi}{7}
\tightlist
\item
  设 \(E_{1}\) ， \(F_{1}\) ， \(E_{2}\) ， \(F_{2}\) 都是 K-向量空间。设 u （相应地， v ）是从 \(E_{1}\) 到 \(F_{1}\) （相应地，从 \(E_{2}\) 到 \(F_{2}\) ）的部分算子。我们用 \(u \otimes v\) 表示从 \(E_{1} \otimes F_{1}\) 到 \(E_{2} \otimes F_{2}\) 中的、定义域为 dom(u) \(\otimes\) dom(v) 且满足 \((u \otimes v)(x \otimes y) = u(x) \otimes v(y)\) （对所有 \((x, y) \in E_{1} \times E_{2}\) ）的部分算子。
\end{enumerate}

\subsection{闭算子、可闭算子与稠定算子}

在本节中， K 表示一个交换拓扑域（TG, III, p. 54）。

定义 2. --- 设 E 与 F 是 K 上的拓扑向量空间（EVT, I, p. 1，定义 1）。设 \(u \in \mathcal{P}(\mathrm{E};\mathrm{F})\) 是从 E 到 F 的部分算子。

若 u 的定义域在 E 中稠密，就称 u 是稠定算子。

若 u 的图像在拓扑向量空间 E × F 中是闭的，就称 u 是闭算子。若 u 有一个闭扩张，就称 u 是可闭的。

设 E 、 F 与 G 是 K 上的拓扑向量空间。稠定算子 \(u \in \mathcal{P}(\mathrm{E};\mathrm{F})\) 的每个扩张都是稠定的。此外，若 \(v \in \mathcal{L}(\mathrm{E};\mathrm{F})\) ，则 \(u + v\) 是稠定算子。若 \(v: F \to G\) （相应地， \(w: G \to E\) ）是拓扑向量空间的同构，则 \(v \circ u\) （相应地， \(u \circ w\) ）是稠定算子。

例. --- 设 E 与 F 是 K 上的拓扑向量空间。

\begin{enumerate}
\def\labelenumi{\arabic{enumi})}
\item
  设空间 F 是豪斯多夫的。设 u 是从 E 到 F 中的线性映射。若 u 连续，则部分算子 u 是闭的（TG, I, p. 53，推论 2）。再假设 K 是非离散赋值域，且 E 与 F 是 K 上的可度量化的完备拓扑向量空间。由闭图像定理（EVT, I, p. 19，推论 5），这时由 u 所定义的部分算子是闭的，当且仅当 u 连续。
\item
  设空间 E 是豪斯多夫的。设 v 是从 F 到 E 中的单射连续线性映射。这时部分算子 \(v^{-1} \in \mathcal{P}(\mathrm{E};\mathrm{F})\) （参见 IV, p. 226）是闭的，因为它的图像是 v 的图像在从 \(F \times E\) 到 \(E \times F\) 的、由 \((y,x) \mapsto (x,y)\) 所定义的拓扑向量空间同构下的像，而 v 的图像是闭的。
\end{enumerate}

命题 1. --- 设 E 与 F 是 K 上的拓扑向量空间。从 E 到 F 中的部分算子 u 是可闭的，当且仅当 u 的图像 \(\Gamma_{u}\) 在 E × F 中的闭包是一个泛函图像。这时存在唯一的从 E 到 F 的部分算子 v ，其图像为 \(\overline{\Gamma}_{u}\) ，且它是 u 的最小闭扩张。

若 \(u\) 的图像在 \(\mathrm{E} \times \mathrm{F}\) 中的闭包是泛函图像，则它就是某个部分算子的图像，而后者是 \(u\) 的一个闭扩张，因而它是可闭的。反之，设 \(u \subset w\) ，其中 \(w\) 是闭的。闭包 \(\overline{\Gamma}_u\) （即 \(u\) 的图像在 \(\mathrm{E} \times \mathrm{F}\) 中的闭包）包含在 \(\Gamma_w\) 中，因而 \(\overline{\Gamma}_u\) 是泛函图像。

最后的断言由下述事实得出：若 w 是 u 的一个闭扩张，则 w 的图像包含 \(\overline{\Gamma}_{u}\) 。

定义 3. --- 设 E 与 F 是 K 上的拓扑向量空间。设 u 是从 E 到 F 中的可闭算子。图像为 \(\overline{\Gamma}_{u}\) 的闭算子称为 \(u\) 的闭包，记作 \(\overline{u}\) 。

注. --- 设 E 与 F 是 K 上的拓扑向量空间。设 u 是从 E 到 F 的可闭算子。u 的闭包的定义域包含在 u 的定义域在 E 中的闭包内。一般而言，二者并不相等（IV, p. 344 的习题 1, b)）。

若 \(u \in \mathcal{P}(\mathrm{E};\mathrm{F})\) 是可闭的且 \(\operatorname{dom}(u) = \operatorname{E}\) ，则 \(u = \overline{u}\) 是闭的，因为这时 \(\operatorname{dom}(\overline{u}) = \operatorname{dom}(u)\) 。

命题 2. --- 设 K = R ，并设 E 与 F 是 R 上的拓扑向量空间。设 u 是从 E 到 F 中的部分算子。则 u 是稠定的（相应地，是闭的、是可闭的），当且仅当 \(u_{(\mathbf{C})}\) 是稠定的（相应地，是闭的、是可闭的）。

设 \(u\) 的定义域在 E 中稠密。 \(\mathrm{E}_{(\mathbf{C})}\) 中 0 的每个邻域都包含形如 \(\mathrm{V} + i\mathrm{V}\) 的一个邻域（EVT, II, p. 65），其中 V 是 E 中 0 的一个邻域，因而它包含 \(u_{(\mathbf{C})}\) 的定义域中的一个元素；所以这一部分算子是稠定的。反过来也成立，因为从 \(\mathrm{E}_{(\mathbf{C})}\) 到 E 中的、把 \(x + iy\) 映到 \(x\) （对所有 \((x,y) \in \mathrm{E} \times \mathrm{E}\) ）的映射是连续且满射的。

我们将 \(u_{(\mathbf{C})}\) 的图像等同于 \(u\) 的图像的复化拓扑向量空间。于是，它在 \(\mathrm{E}_{(\mathbf{C})} \times \mathrm{E}_{(\mathbf{C})}\) 中是闭的，如果 \(u\) 的图像在 \(\mathrm{E} \times \mathrm{E}\) 中是闭的。反之，我们有 \(\Gamma_u = \Gamma_{u_{(\mathbf{C})}} \cap (\mathrm{E} \times \mathrm{E})\) （在 \(\mathrm{E}_{(\mathbf{C})} \times \mathrm{E}_{(\mathbf{C})}\) 中）；由于 \(\mathrm{E} \times \mathrm{E}\) 在 \(\mathrm{E}_{(\mathbf{C})} \times \mathrm{E}_{(\mathbf{C})}\) 中是闭的，故部分算子 \(u\) 在 \(u_{(\mathbf{C})}\) 闭时是闭的。

从 E 到 F 中的部分算子 \(v\) 是 \(u\) 的扩张，当且仅当 \(v_{(\mathbf{C})}\) 是 \(u_{(\mathbf{C})}\) 的扩张；因此，部分算子 \(u_{(\mathbf{C})}\) 在 \(u\) 可闭时是可闭的。反之，若 \(u_{(\mathbf{C})}\) 可闭，则部分算子 \(u\) 也可闭，因为 \(\overline{\Gamma_u} = \overline{\Gamma_{u_{(\mathbf{C})}}} \cap (\mathrm{E} \times \mathrm{E})\) 这时是泛函图像（命题 1）。

引理 1. — 设 E 、 F 与 G 是 K 上的拓扑向量空间。设 u 是从 E 到 F 的闭算子。

\begin{enumerate}
\def\labelenumi{\alph{enumi})}
\item
  对每个 \(v \in \mathcal{L}(\mathrm{E};\mathrm{F})\) ，部分算子 \(u + v\) 是闭的。
\item
  对每个 \(v \in \mathcal{L}(\mathrm{G};\mathrm{E})\) ，部分算子 \(u \circ v\) 是闭的。
\end{enumerate}

我们来证明 \(a\) )。设 \(\gamma\) 是映射 \((x,y)\mapsto (x,y - v(x))\) ，它从 \(\mathrm{E}\times \mathrm{F}\) 到其自身，并且是连续的。对每个 \((x,y)\in \mathrm{E}\times \mathrm{F}\) ，我们有 \(\gamma (x,y)\in \Gamma_u\) 当且仅当 \(x\in \operatorname {dom}(u)\) 且 \(y = u(x) + v(x)\) ，亦即 \(\overline{\gamma}^{1}(\Gamma_{u}) = \Gamma_{u + v}\) 。断言得证。

我们来证明 \(b\) )。映射 \(\eta = (v, 1_{\mathrm{F}})\) 从 \(\mathbf{G} \times \mathbf{F}\) 到 \(\mathbf{E} \times \mathbf{F}\) 中是连续的；对每个 \((z, y) \in \mathbf{G} \times \mathbf{F}\) ，我们有 \(\eta(z, y) = (v(z), y)\) ，从而 \(\overline{\eta}(\Gamma_u) = \Gamma_{u \circ v}\) 。断言得证。

设 E 与 F 是 K 上的拓扑向量空间，且 F 是豪斯多夫的。设 \(a \in \mathrm{K}\) 。若 \(u \in \mathcal{P}(\mathrm{E};\mathrm{F})\) 是可闭的，则 \(au\) 也是可闭的。若 \(a \neq 0\) ，我们有 \(\overline{au} = a\overline{u}\) ，且 \(u\) 是闭的当且仅当 \(au\) 是闭的。若 \(a = 0\) ，则 \(au\) 的闭包是 \(0_{\overline{\operatorname{dom}(u)}}\) ，而 \(a\overline{u}\) 等于 \(0_{\operatorname{dom}(\overline{u})}\) ；因此可能出现这样的情形： \(u\) 是闭的，而 \(au\) 不是闭的。

命题 3. --- 设 E 与 F 是 K 上的拓扑向量空间，其中 F 是豪斯多夫的。设 u 是从 E 到 F 的闭部分算子。u 的核是 E 的一个闭向量子空间。

事实上，u 的核是 E × F 的闭子空间 \(\Gamma_{u} \cap (\mathrm{E} \times \{0\})\) 在从 E 到 E × F 中的连续线性映射 \(x \mapsto (x, 0)\) 下的原像。

在本节余下部分，我们假设 K 是非离散赋值域。

定义 4. --- 设 E 与 F 是 K 上的赋范空间，u 是从 E 到 F 的部分算子。对 dom(u) 中的 x ，我们记

\[
\| x \| _ {u} = (\| x \| _ {\mathrm{E}} ^ {2} + \| u (x) \| _ {\mathrm{F}} ^ {2}) ^ {1 / 2}.
\]

如此定义的映射 \(x \mapsto \|x\|_{u}\) 是 \(\text{dom}(u)\) 上的一个范数。我们把这样得到的赋范空间记作 \(E_{u}\) 。

注. --- 设 E 与 F 是 K 上的赋范空间，u 是从 E 到 F 中的部分算子。

\begin{enumerate}
\def\labelenumi{\arabic{enumi})}
\item
  \(E_{u}\) 到 E 中的典则单射是连续的，因为我们有 \(\|x\|\leqslant\|x\|_{u}\) （对所有 \(x\in E\) ）。特别地，\(\operatorname{dom}(u)\) 的在 E 中闭的每个子空间在 \(E_{u}\) 中也是闭的。
\item
  若 E 与 F 是希尔伯特空间，则空间 \(E_{u}\) 是准希尔伯特空间，因为 \(E_{u}\) 上的范数来自下述正埃尔米特形式
\end{enumerate}

\[
(x, y) \mapsto (x \mid y) _ {u} = \langle x \mid y \rangle + \langle u (x) \mid u (y) \rangle
\]

（它定义在 dom(u) 上）。

命题 4. --- 设 E 与 F 是巴拿赫空间（相应地，希尔伯特空间），u 是从 E 到 F 中的部分算子。则 u 是闭的，当且仅当赋范空间 \(E_{u}\) 是巴拿赫空间（相应地，希尔伯特空间）。

只须处理巴拿赫空间的情形。 \(E_{u}\) 的范数，是经由双射线性映射 \((x,y)\mapsto x\) （它把 \(\Gamma_{u}\) 映到 dom(u) 中）作结构迁移，从下述范数得到的：即限制到子空间 \(\Gamma_{u}\) （赋以范数 \((x,y)\mapsto(\|x\|_{\mathrm{E}}^{2}+\|y\|_{\mathrm{F}}^{2})^{1/2}\) ，它由巴拿赫空间 \(E\oplus F\) 上的范数限制而得）上所得到的范数。因此，空间 \(E_{u}\) 是巴拿赫空间，当且仅当子空间 \(\Gamma_{u}\) （在 \(E\oplus F\) 中）是闭的。

若 \(u\) 与 \(v\) 分别是从 E 到 F 与从 F 到 G 的部分算子，且若 \(u\) 是闭的，则部分算子 \(v \circ u\) 一般说来不是闭的，即使 \(v\) 是连续的（IV, p. 344 的习题 1, c)）。尽管如此，我们有如下的充分条件：

引理 2. --- 设 E 、 F 与 G 是 K 上的赋范空间，其中 F 是巴拿赫空间。设 u 是从 E 到 F 的闭部分算子，并设 \(v \in \mathcal{L}(\mathrm{F};\mathrm{G})\) 。如果存在 \(\mathrm{C} \in \mathbf{R}_{+}\) 使得

\[
\| u (x) \| \leqslant C (\| x \| + \| (v \circ u) (x) \|)
\]

对所有 \(x \in \text{dom}(v \circ u) = \text{dom}(u)\) 成立；也就是说，若从 \(x \mapsto u(x)\) 到 F 中的线性映射 \(E_{v \circ u}\) 是连续的，则 \(v \circ u\) 是闭的。

记 \(w = v \circ u\) 。设 \((x_n, w(x_n))_{n \in \mathbf{N}}\) 是 \(w\) 的图像中的一个序列，它在 \(E \times G\) 中收敛。设 \(x\) 是序列 \((x_n)\) 的极限。假设表明，序列 \((u(x_n))_{n \in \mathbf{N}}\) 这时是 \(F\) 中的一个柯西序列；它收敛到元素 \(y\) ，该元素属于 \(F\) 。由于 \(u\) 是闭的，我们有 \(x \in \text{dom}(u)\) 与 \(y = u(x)\) 。于是 \(w(x_n) = v(u(x_n))\) 趋于 \(v(y) = v(u(x))\) ，因为 \(v\) 连续；因此 \(w\) 的图像是闭的。

设 u 是巴拿赫空间 E 上的闭部分算子， F 是 \(\text{dom}(u)\) 的一个子空间。若 F 在巴拿赫空间 \(E_{u}\) 中稠密，则 u 到 F 上的约化是可闭的，且其闭包等于 u 。这时称 F 为 u 的一个核。

\subsection{部分算子的例子}

在本节中， K = R 或 C 。

例. --- 1) 设 X 是局部紧拓扑空间， \(\mu\) 是 X 上的一个正测度。我们固定元素 \(p_1\) 与 \(p_2\) ，它们属于 \([1, +\infty[\) 。

设 \(g\) 是 \(\mu\)-可测的函数，定义在 \(X\) 上，取值于 \(K\) 中。用 \(D\) 表示 \(\mathcal{L}_{K}^{p_1}(X, \mu)\) 中满足如下条件的函数 \(f\) （属于 \(\mathcal{L}_{K}^{p_1}(X, \mu)\) ）所成的子空间： \(gf \in \mathcal{L}_{K}^{p_2}(X, \mu)\) 。由 \(D\) 到 \(\mathcal{L}_{K}^{p_2}(X, \mu)\) 中的、由 \(f \mapsto gf\) 所定义的线性映射确定了 \(\mathcal{L}_{K}^{p_1}(X, \mu)\) 到 \(\mathcal{L}_{K}^{p_2}(X, \mu)\) 中的一个部分算子，我们把它记作 \(m_g\) 。

\(\mu\)-可忽略函数在 \(\mathcal{L}_{\mathrm{K}}^{p_1}(\mathrm{X},\mu)\) 中所成的向量子空间包含在 D 中；并且， \(m_g\) 把 \(\mu\)-可忽略函数映为仍为 \(\mu\)-可忽略的函数。我们将用 \(\tilde{m}_g\) 表示从 \(\mathrm{L}_{\mathrm{K}}^{p_1}(\mathrm{X},\mu)\) 到 \(\mathrm{L}_{\mathrm{K}}^{p_2}(\mathrm{X},\mu)\) 中的、由 \(m_g\) 通过过渡到商诱导出的部分算子。它称为以 \(g\) 为乘子的乘法算子（从 \(\mathrm{L}_{\mathrm{K}}^{p_1}(\mathrm{X},\mu)\) 到 \(\mathrm{L}_{\mathrm{K}}^{p_2}(\mathrm{X},\mu)\) 中）。凡是 \(g_1\) 与 \(g_2\) 局部 \(\mu\)-几乎处处相等的函数，都定义同一个乘法算子。

命题 5. — 乘法算子 \(\widetilde{m}_{g}\) （从 \(\mathrm{L}_{\mathrm{K}}^{p_{1}}(\mathrm{X},\mu)\) 到 \(\mathrm{L}_{\mathrm{K}}^{p_{2}}(\mathrm{X},\mu)\) ）是稠定的闭算子。

我们先证明部分算子 \(\widetilde{m}_g\) 是闭的。设 \((f_n, h_n)_{n \in \mathbf{N}}\) 是 \(\mathcal{L}_{\mathrm{K}}^{p_1}(\mathrm{X}, \mu) \times \mathcal{L}_{\mathrm{K}}^{p_2}(\mathrm{X}, \mu)\) 中的一个序列，使得 \((\widetilde{f}_n, \widetilde{h}_n)\) 的类所作成的序列（即 \(f_n\) 与 \(h_n\) 的类）属于 \(\widetilde{m}_g\) 的图像，且当 \(\mathrm{L}_{\mathrm{K}}^{p_1}(\mathrm{X}, \mu) \times \mathrm{L}_{\mathrm{K}}^{p_2}(\mathrm{X}, \mu)\) 趋于无穷时在 \(n\) 中收敛。设 \((f, h)\) 是 \(\mathcal{L}_{\mathrm{K}}^{p_1}(\mathrm{X}, \mu) \times \mathcal{L}_{\mathrm{K}}^{p_2}(\mathrm{X}, \mu)\) 中的一对元素，使得它们的类所成的对 \((\widetilde{f}, \widetilde{h})\) 是 \((\widetilde{f}_n, \widetilde{h}_n)\) 的极限。

存在序列 \((f_{n_{k}})_{k\in\mathbf{N}}\) （从序列 \((f_{n})_{n}\) 中抽取的）使得 \(f_{n_{k}}(x)\) 收敛到 \(f(x)\) ，对 \(\mu\)-几乎所有的 x 成立（INT, IV, p. 131, § 3, no. 4，定理 3）。由此推出， \(h_{n_{k}}(x)=g(x)f_{n_{k}}(x)\) \(\mu\)-几乎处处收敛于 \(g(x)f(x)\) 。此外，序列 \((h_{n_{k}})\) 在空间 \(\mathcal{L}_{\mathrm{K}}^{p_{2}}(\mathrm{X},\mu)\) 中收敛于 h 。因此，函数 h 与 gf \(\mu\)-几乎处处相等（loc. cit.）。于是 \(\widetilde{f}\) 属于 \(\widetilde{m}_{g}\) 的定义域，且 \(\widetilde{h}=\widetilde{m}_{g}(\widetilde{f})\) 。这就证明了 \(\widetilde{m}_{g}\) 是闭的。

我们来证明 \(\widetilde{m}_g\) 的定义域在 \(\mathrm{L}_{\mathrm{K}}^{p_1}(\mathrm{X},\mu)\) 中稠密。只须验证：函数 \(f\in \mathcal{K}(\mathrm{X};\mathrm{K})\) 的类属于 \(\widetilde{m}_g\) 的定义域在 \(\mathrm{L}_{\mathrm{K}}^{p_1}(\mathrm{X},\mu)\) 中的闭包。设 \(f\in \mathcal{K}(\mathrm{X};\mathrm{K})\) ，并用 \(\widetilde{f}\) 表示它在 \(\mathrm{L}_{\mathrm{K}}^{p_1}(\mathrm{X},\mu)\) 中的类。对每个整数 \(n\in \mathbf{N}\) ，用 \(\varphi_{n}\) 表示由元素 \(x\in \mathrm{X}\) （满足 \(|g(x)|\leqslant n\) ）所成集合的特征函数，并令 \(f_{n} = f\varphi_{n}\) 。这时我们有 \(|gf_n|\leqslant n|f|\) ，而它属于 \(\mathcal{L}_{\mathrm{K}}^{p_2}(\mathrm{X},\mu)\) ，因此 \(f_{n}\) 属于 \(\widetilde{m}_g\) 的定义域。对 X 的每个元素 \(x\) ，序列 \((f_n(x))_{n\in \mathbf{N}}\) 收敛于 \(f(x)\) （当 \(n\to +\infty\) 时）；此外，我们有 \(|f_n|\leqslant |f|\) ，它属于 \(\mathcal{L}_{\mathrm{K}}^{p_1}(\mathrm{X},\mu)\) 。由勒贝格定理（INT, IV, p. 137, § 3, n°7，定理 6）， \(f_{n}\) 的类的序列在 \(\widetilde{f}\) 中收敛于 \(\mathrm{L}_{\mathrm{K}}^{p_1}(\mathrm{X},\mu)\) 。因此， \(f\) 的类属于 \(\widetilde{m}_g\) 的定义域的闭包。

在下述命题中，我们假设 \(p_1 = p_2 = 2\) 。

命题 6. --- a) 设 \(g'\) 是 X 上的 \(\mu\)-可测函数，满足 \(|g|\leqslant |g'|\) 。我们有 \(\mathrm{dom}(\widetilde{m}_{g'})\subset\mathrm{dom}(\widetilde{m}_{g})\) ，且 \(\mathrm{dom}(\widetilde{m}_{g'})\) 是部分算子 \(\widetilde{m}_{g}\) 的一个核；

\begin{enumerate}
\def\labelenumi{\alph{enumi})}
\setcounter{enumi}{1}
\tightlist
\item
  设 F 是 \(\mathcal{L}_{\mathrm{K}}^{2}(\mathrm{X},\mu)\) 的一个子空间，它与 \(\mathcal{K}(\mathrm{X};\mathrm{K})\) 的交在 \(\mathcal{K}(\mathrm{X};\mathrm{K})\) 中稠密，且它在 \(\mathrm{L}_{\mathrm{K}}^{2}(\mathrm{X},\mu)\) 中的像 G 包含在 \(\mathrm{dom}(\widetilde{m}_g)\) 中。若 \(|g|^2\) 局部 \(\mu\)-可积，则 \(\mathcal{K}(\mathrm{X};\mathrm{K})\) 包含在 \(\widetilde{m}_g\) 的定义域中，且 G 是 \(m_g\) 的一个核。
\end{enumerate}

我们来证明 \(a\) )。若 \(f \in \mathcal{L}_{\mathrm{K}}^{2}(\mathrm{X}, \mu)\) 属于 \(m_{g'}\) 的定义域（从而 \(fg' \in \mathcal{L}_{\mathrm{K}}^{2}(\mathrm{X}, \mu)\) ），则假设蕴含 \(fg \in \mathcal{L}_{\mathrm{K}}^{2}(\mathrm{X}, \mu)\) ，由此即得结论。

我们来证明 \(\mathrm{dom}(\widetilde{m}_{g'})\) 是 \(\widetilde{m}_g\) 的核，也就是说， \(\widetilde{m}_{g'}\) 的定义域在希尔伯特空间 \(\mathrm{E}_{\widetilde{m}_g}\) 中稠密。设 \(h \in \mathcal{L}_{\mathrm{K}}^2(\mathrm{X}, \mu)\) ，其类 \(\widetilde{h}\) 属于 \(\mathrm{E}_{\widetilde{m}_g}\) 且与 \(\mathrm{dom}(\widetilde{m}_{g'})\) 正交。这就意味着

\[
(\widetilde {h} \mid \widetilde {h} ^ {\prime}) _ {\widetilde {m} _ {g}} = \int_ {\mathrm{X}} \overline {{h}} h ^ {\prime} (1 + | g | ^ {2}) d \mu = 0
\]

它对每个这样的函数 \(h' \in \mathcal{L}_{\mathrm{K}}^{2}(\mathrm{X}, \mu)\) （其类 \(\widetilde{h}'\) 属于 \(\operatorname{dom}(\widetilde{m}_{g'})\) ）成立。

设 C 是 X 的紧子集， \(\varphi\) 是它的特征函数。设 \(n \in \mathbf{N}\) 。设 \(\varphi_n\) 是 \(\mu\)-可积集合 \(\mathrm{C}_n\) （由 C 中满足 \(x \in \mathrm{C}\) 的元素 \(|h(x)| \leqslant n\) 所成）的特征函数，并令 \(h_n' = \varphi_n h\) 。 \(h_n'\) 的类属于 \(\tilde{m}_{g'}\) 的定义域，因为 \(|g'h_n'| \leqslant n\varphi\) ；由此得出

\[
0 = \int_ {\mathrm{X}} \overline {{h}} h _ {n} ^ {\prime} (1 + | g | ^ {2}) d \mu = \int_ {\mathrm{X}} | h | ^ {2} \varphi_ {n} (1 + | g | ^ {2}) d \mu .
\]

这蕴含：对 \(\mu\)-几乎所有 \(x \in C_{n}\) ， h 为零；由于 n 任意，对 \(\mu\)-几乎所有 \(x \in C\) ， h 为零。最终推出 h \(\mu\) \(\mu\)-几乎处处为零，因为 C 任意且 h 是缓增的（INT, V, p. 9, § 1, n°3，推论）。

现在考虑断言 \(b\) )。由于 \(|g|^2\) 局部 \(\mu\)-可积，函数 \(fg\) 属于 \(\mathcal{L}_{\mathrm{K}}^{2}(\mathrm{X},\mu)\) （对 \(f\in\mathcal{K}(\mathrm{X};\mathrm{K})\) ），因此 \(\mathcal{K}(\mathrm{X};\mathrm{K})\) 包含在 \(\widetilde{m}_{g}\) 的定义域中。

设 \(h \in \mathcal{L}_{\mathrm{K}}^{2}(\mathrm{X}, \mu)\) ，其类 \(\widetilde{h}\) 属于 \(\mathrm{E}_{\widetilde{m}_{g}}\) 且与 G 正交。这时我们有

\[
0 = (\widetilde {h} \mid \widetilde {h} ^ {\prime}) _ {\widetilde {m} _ {g}} = \int_ {\mathrm{X}} h \overline {{h}} ^ {\prime} (1 + | g | ^ {2}) d \mu
\]

对所有 \(h' \in F\) ，其类为 \(h'\) 。鉴于关于 F 的假设，这意味着测度 \(h(1 + |g|^{2}) \cdot \mu\) 为零，因而 h \(\mu\)-几乎处处为零，因为 h 是缓增的。

设 \(p\) 是实数 \(\geqslant 1\) 。设 \(h\) 是 \(\mathcal{L}_{\mathrm{K}}^{\infty}(\mathrm{X},\mu)\) 中的元素。以 \(\widetilde{m}_h\) 为乘子的乘法算子 \(h\) 是 \(\mathrm{L}_{\mathrm{K}}^{p}(\mathrm{X},\mu)\) 的一个自同态。假设由 \(x\in X\) 中满足 \(h(x)=0\) 的 x 所成的集 Y 是局部 \(\mu\)-可忽略的。这时自同态 \(\widetilde{m}_h\) 是单射的（IV, p. 186 的引理 7）。用 \(h^{-1}\) 表示 X 上在 Y 上等于 0 、在 \(x \mapsto 1/h(x)\) 上等于 \(\mathrm{X} - \mathrm{Y}\) 的函数。逆部分算子 \(\widetilde{m}_{h}^{-1}\) 就是以 \(h^{-1}\) 为乘子的、从 \(\mathrm{L}_{\mathrm{K}}^{p}(\mathrm{X},\mu)\) 到 \(\mathrm{L}_{\mathrm{K}}^{p}(\mathrm{X},\mu)\) 中的乘法算子，即 \(\widetilde{m}_{h}^{-1} = \widetilde{m}_{h^{-1}}\) 。事实上， \(\widetilde{m}_{h}\) 的像是由形如 \(g \in \mathcal{L}_{\mathrm{K}}^{p}(\mathrm{X},\mu)\) （其中 \(g = hf\) ）的函数 \(f \in \mathcal{L}_{\mathrm{K}}^{p}(\mathrm{X},\mu)\) 的类所成的空间。这一条件等价于：对所有 \(g(x)/h(x) = f(x)\) 有 \(x \in \mathrm{X} - \mathrm{Y}\) ，且若 \(g(x) = 0\) 则 \(x \in \mathrm{Y}\) 。这蕴含： \(\widetilde{m}_{h}^{-1}\) 在 \(\mathrm{L}_{\mathrm{K}}^{p}(\mathrm{X},\mu)\) 中的定义域就是 \(\widetilde{m}_{h^{-1}}\) 的定义域，且等式 \(\widetilde{m}_{h}^{-1} = \widetilde{m}_{h^{-1}}\) 成立。

在下文中，我们有时把部分乘法算子简记为 \(m_{h}\) ，它以 h 为乘子，是从 \(\mathrm{L}_{\mathrm{K}}^{p_{1}}(\mathrm{X},\mu)\) 到 \(\mathrm{L}_{\mathrm{K}}^{p_{2}}(\mathrm{X},\mu)\) 中的。

\begin{enumerate}
\def\labelenumi{\arabic{enumi})}
\setcounter{enumi}{1}
\tightlist
\item
  设 E 是 K 上的希尔伯特空间， \(\mathrm{B} = (e_i)_{i \in \mathrm{I}}\) 是 E 的一个标准正交基。设 \((\lambda_i)_{i \in \mathrm{I}}\) 是 K 中元素的一个族。设 D 是 E 中由元素 \(x \in E\) （使得族 \((\lambda_i \langle e_i | x \rangle)_{i \in \mathrm{I}}\) 在 K 中平方可和）所成的向量子空间。空间 D 在 E 中稠密，因为它包含向量 \(e_i\) （对所有 \(i \in I\) ）。以 D 为定义域的部分算子 u 由下式给出
\end{enumerate}

\[
x \mapsto \sum_ {i \in I} \lambda_ {i} \langle e _ {i} | x \rangle e _ {i}
\]

称为基 B 中的部分对角算子，而 \((\lambda_{i})_{i\in I}\) 称为 u 的特征值族。

算子 \(u\) 是闭的。事实上，设 \((x_{n}, u(x_{n}))_{n \in \mathbf{N}}\) 是 \(u\) 的图像中在 \(\mathrm{E} \times \mathrm{E}\) 内收敛的一个元素序列，并设 \((x, y)\) 为其极限。这时我们有 \(\langle e_{i} | x_{n} \rangle \to \langle e_{i} | x \rangle\) （对所有 \(i \in \mathrm{I}\) ），并且

\[
\langle e _ {i} | u (x _ {n}) \rangle = \lambda_ {i} \langle e _ {i} | x _ {n} \rangle \rightarrow \langle e _ {i} | y \rangle
\]

对所有 \(i \in I\) 。于是， \(\lambda_{i} \langle e_{i} \mid x \rangle = \langle e_{i} \mid y \rangle\) 对所有 \(i \in I\) 成立；这就证明了 \(x \in D\) 且 \(u(x) = y\) ，即 \(u\) 是闭的。

这个例子实际上是前一例子应用于拓扑空间 X = I （赋以离散拓扑与计数测度 \(\mu\) ）的特殊情形，因为 E 通过映射 \(\ell^{2}(\mathrm{I}) = \mathrm{L}^{2}(\mathrm{I}, \mu)\) 等同于空间 \(x \mapsto (\langle e_{i} \mid x \rangle)_{i \in \mathrm{I}}\) （EVT, V, p. 23，推论 2），而 u 这时等同于乘法算子 \(m_{\lambda}\) ，其中 \(\lambda\) 是函数 \(i \mapsto \lambda_{i}\) 。

\begin{enumerate}
\def\labelenumi{\arabic{enumi})}
\setcounter{enumi}{2}
\tightlist
\item
  集合 \(N_{R} = N \cup \{\infty, \omega\}\) 赋有 VAR, R2, p. 10 中所述的全序，使得 \(n < \infty < \omega\) 对所有 \(n \in N\) 成立。设 \(r \in N_{R}\) 。设 \(n \in N\) ，并设 U 是 \(R^{n}\) 的一个开子集。设 \(k \in N\) 满足 \(k \leqslant r\) 。设 \((n_{\alpha})_{|\alpha| \leqslant k}\) 是 \(\mathcal{C}^{r}(U)\) 中元素的一个族，其中所考虑的重指标属于 \(N^{n}\) 。族 \((n_{\alpha})\) 在 U 上定义一个阶 \(\leqslant k\) 的纯量微分算子 D （参见 VAR, R2, 14.1.6, 14.1.4）。
\end{enumerate}

对每个满足 \(k \leqslant m \leqslant r\) 的整数 m ，微分算子 D 定义一个从 \(\mathrm{C}^{m}(\mathrm{U})\) 到 \(\mathrm{C}^{m-k}(\mathrm{U})\) 中的线性映射，它把 \(f \in \mathrm{C}^{m}(\mathrm{U})\) 映为

\[
\mathrm{D} (f) = \sum_ {| \alpha | \leqslant k} n _ {\alpha} \partial^ {\alpha} (f).
\]

同一公式定义了一个从 \(\mathcal{D}(\mathrm{U})\) 到 \(\mathcal{D}'(\mathrm{U})\) 中的连续线性映射（IV, p. 208 的定义 2）。

定义 5. --- 设 E 是 \(\mathcal{D}'(\mathrm{U})\) 的一个包含 \(\mathcal{D}(\mathrm{U})\) 的向量子空间。 E 上与 D 相伴的微分算子，指的是 E 上这样一个部分算子：它是以下述部分算子的扩张，该算子以 \(\mathcal{D}(\mathrm{U})\) 为定义域、由 \(\varphi \mapsto \mathrm{D}(\varphi)\) 所定义。

例如，设系数 \(n_{\alpha}\) 是 U 上的有界函数。设 \(\mu\) 是 U 上的勒贝格测度。若 p 是 \([1,+\infty[\) 中的一个元素，则可在 \(\mathrm{L}^{p}(\mathrm{U})\) 上定义一个与 D 相伴的微分算子，其定义域是索伯列夫空间 \(\mathrm{W}^{k,p}(\mathrm{U})\) （IV, p. 221 的第 14 节），因为这时我们有 \(n_{\alpha}\partial^{\alpha}(f)\in\mathrm{L}^{p}(\mathrm{U})\) （对所有 \(f\in\mathrm{W}^{k,p}(\mathrm{U})\) 与每个 \(|\alpha|\leqslant k\) ）。

\subsection{伴随}

在本节中， K 是域 R 或 C 之一， E 与 F 表示 K 上的希尔伯特空间。

设 \(u\) 是从 E 到 F 的稠定算子。用 \(D\) 表示 \(u\) 的定义域。对 \(y \in \mathrm{F}\) ，令 \(\lambda_y\) 为 D 上由 \(\lambda_y(x) = \langle y \mid u(x) \rangle\) （对所有 \(x \in \mathrm{D}\) ）所给的线性形式。我们用 \(\mathrm{D}^*\) 表示满足下述条件的向量 \(y \in \mathrm{F}\) 所成的集： \(\lambda_y\) 在 D 上连续。它是 F 的一个向量子空间。设 \(y \in \mathrm{D}^*\) ；由于 D 在 E 中稠密，线性形式 \(\lambda_y\) 唯一地扩张为 E 上的连续线性形式，我们仍把它记作 \(\lambda_y\) 。依照 EVT, V, p. 15，定理 3，存在唯一的元素 \(u^*(y)\) 属于 E ，使得 \(\lambda_y(x) = \langle u^*(y) \mid x \rangle\) 对所有 \(x \in \mathrm{E}\) 成立。映射 \(y \mapsto u^*(y)\) 是从 \(\mathrm{D}^*\) 到 E 中的线性映射。

定义 6. --- 以 \(D^*\) 为定义域、由 \(y \mapsto u^*(y)\) 所定义的从 F 到 E 的部分算子称为 \(u\) 的伴随，记作 \(u^*\) 。

注. --- 于是， \(y \in D^{*}\) 当且仅当存在 \(c \in R_{+}\) 使得 \(|\langle y | u(x) \rangle| \leqslant c \|x\|\) 对所有 \(x \in D\) 成立。这时元素 \(u^{*}(y)\) 由下述关系式刻画：

\[
\langle y \mid u (x) \rangle = \langle u ^ {*} (y) \mid x \rangle\tag{1}
\]

对所有 \(x \in \mathrm{D}\) 成立。这时我们有 \(|\langle y | u(x) \rangle| \leqslant \| u^{*}(y) \| \| x \|\) 对所有 \(x \in \mathrm{D}\) 成立。

在 \(u\) 是从 E 到 F 中的连续线性映射的情形，按前一定义意义的伴随与 EVT, V, p. 38，定义 1 中定义的伴随重合，因为这时 D* 等于 F 。

设 \(v \in \mathcal{P}(\mathrm{E};\mathrm{F})\) 满足 \(u \subset v\) 。这时我们有 \(v^{*} \subset u^{*}\) 。

设 \(v \in \mathcal{P}(\mathrm{E};\mathrm{F})\) 满足 \(u + v\) 稠定。这时部分算子 \(v\) 便是稠定的，且有 \(u^{*} + v^{*} \subset (u + v)^{*}\) 。一般而言等式不成立（IV, p. 347 的习题 9）。若 \(v \in \mathcal{L}(\mathrm{E};\mathrm{F})\) ，则 \(u + v\) 稠定，且 \((u + v)^{*} = u^{*} + v^{*}\) 。例如，若 \(\mathrm{F} = \mathrm{E}\) 且 \(v = \lambda 1_{\mathrm{E}}\) （其中 \(\lambda \in \mathrm{K}\) ），就是这种情形。

设 G 是 K 上的希尔伯特空间， \(v \in \mathcal{P}(\mathrm{F};\mathrm{G})\) 是稠定算子。若 \(v \circ u\) 稠定，则 \(u^{*} \circ v^{*} \subset (v \circ u)^{*}\) 。一般而言等式不成立（loc. cit.）。若 u （相应地， v ）是同构，则 \(v \circ u\) 稠定且有 \((v \circ u)^{*} = u^{*} \circ v^{*}\) 。例如，若 E = F （相应地， F = G ）且 \(u = \lambda 1_{E}\) （相应地， \(v = \lambda 1_{F}\) ），其中 \(\lambda \in K^{*}\) ，就是这种情形。

我们用 \(s\) 表示从 \(\mathrm{E} \oplus \mathrm{F}\) 到 \(\mathrm{F} \oplus \mathrm{E}\) 的希尔伯特空间等距同构，它由 \(s(x, y) = (-y, x)\) （对所有 \((x, y) \in \mathrm{E} \oplus \mathrm{F}\) ）定义。

命题 7. --- 设 u 是从 E 到 F 中的稠定部分算子。

\begin{enumerate}
\def\labelenumi{\alph{enumi})}
\item
  \(u^{*}\) 的图像等于 \(s(\Gamma_{u}^{\circ}) = s(\Gamma_{u})^{\circ}\) ；
\item
  部分算子 \(u^{*}\) 是闭的；
\item
  \(u^{*}\) 的核是 \(u\) 的像的正交补。
\end{enumerate}

记 \(\mathrm{W} = s(\Gamma_u)^\circ\) 。由于线性映射 \(s\) 是酉的，我们有 \(\mathrm{W} = s(\Gamma_u^\circ)\) 。

我们有 \((y,x)\in \mathrm{W}\) 当且仅当

\[
\langle (y, x) \mid (- u (x ^ {\prime}), x ^ {\prime}) \rangle = 0
\]

对所有 \(x' \in \mathrm{dom}(u)\) 成立，亦即若

\[
\langle y \mid u (x ^ {\prime}) \rangle = \langle x \mid x ^ {\prime} \rangle
\]

（对所有 \(x' \in \text{dom}(u)\) ）。当 \(y \in \text{dom}(u^{*})\) 且 \(x = u^{*}(y)\) 时，这一性质成立（参见 p. 236 的公式 (1)）。反之，若这一条件成立，我们推出 \(|\langle y \mid u(x') \rangle| \leqslant \|x\| \|x'\|\) 对所有 \(x' \in \text{dom}(u)\) 成立，这蕴含 \(y\) 属于 \(\text{dom}(u^*)\) ；这时我们有 \(u^*(y) = x\) 。因此 \(W = \Gamma_{u^*}\) 。

算子 \(u^{*}\) 是闭的，因为空间 \(s(\Gamma_{u})^{\circ}\) 在 \(\mathrm{F} \oplus \mathrm{E}\) 中是闭的。

我们来证明断言 \(c\) )。若 \(y\) 与 \(u\) 的像正交，则 D 上的线性形式 \(\lambda_y: x \mapsto \langle y | u(x) \rangle\) 为零，于是 \(y \in \text{dom}(u^*)\) 且 \(u^*(y) = 0\) 。反之，设 \(y \in \text{dom}(u^*)\) 。这时 \(u^*(y) = 0\) 当且仅当 \(y\) 与 \(u(x)\) （对所有 \(x \in D\) ）正交（p. 236 的公式 (1)）。

命题 8. --- 设 u 是从 E 到 F 的稠定算子。则 \(u^{*}\) 稠定当且仅当 u 可闭。在这种情形， u 的闭包 \(\overline{u}\) 等于 \(u^{**}\) ，且 \(\overline{u}\) 的伴随等于 \(u^{*}\) 。

由命题 7，部分算子 \(u^*\) 是闭的。假设 \(u^*\) 的定义域 D* 在 F 中稠密。设 \(u^{**}\) 是 \(u^*\) 的伴随；它是从 E 到 F 的闭部分算子。我们来证明 \(u \subset u^{**}\) ，这将蕴含 \(u\) 可闭。设 \(x \in \text{dom}(u)\) 。按 \(u^*\) 的定义， D* 上由 \(y \mapsto \langle x | u^*(y) \rangle\) 与 \(y \mapsto \langle u(x) | y \rangle\) 所给的线性形式相等；因此我们有 \(x \in \text{dom}(u^{**})\) 且 \(u^{**}(x) = u(x)\) ，由此即得断言。

反之，设 \(u\) 可闭；我们有 \(\Gamma_{\overline{u}} = \overline{\Gamma}_u\) （IV, p. 228 的命题 1）。设 \(y \in F\) 是与 \(\mathrm{dom}(u^*)\) 正交的一个向量。 \((y,0)\) （ \(F \oplus E\) 中的元素）这时属于 \(u^*\) 的图像的正交补。而由命题 7, a) ，我们有

\[
\Gamma_ {u ^ {*}} ^ {\circ} = (s (\Gamma_ {u}) ^ {\circ}) ^ {\circ} = \overline {{s (\Gamma_ {u})}} = s (\overline {{\Gamma}} _ {u})
\]

由此得出 \((0,y)\in\Gamma_{\overline{u}}\) ，从而 \(y=\overline{u}(0)=0\) 。由于 dom \((u^{*})\) 的正交补约化为 0 ，空间 dom \((u^{*})\) 在 F 中稠密。

最后，把命题 7 应用于 \(u^{*}\) ，即得

\[
\Gamma_ {u ^ {* *}} = s ^ {- 1} (\Gamma_ {u ^ {*}} ^ {\circ}) = s ^ {- 1} (s (\Gamma_ {u} ^ {\circ \circ})) = \overline {{\Gamma}} _ {u},
\]

因而 \(u^{**} = \overline{u}\) ，进而 \(\overline{u}^{*} = (u^{*})^{**} = \overline{u^{*}} = u^{*}\) ，因为 \(u^{*}\) 是闭的。

推论. --- 若 u 是从 E 到 F 中的闭稠定部分算子，则 \(u^{*}\) 稠定，且有 \(u^{**} = u\) 。

定义 7. --- 设 u 是 E 上的部分算子。若 u 稠定且 \(u^{*}\) 是 u 的扩张，就称 u 是对称的。若 u 稠定且 \(u^{*} = u\) ，就称 u 是自伴的。

称 u 是本质自伴的，如果它可闭且其闭包 \(\overline{u}\) 是自伴的。

称 u 是下有界部分算子，如果 u 是对称的，且存在实数 c 使得 \(\langle x | u(x) \rangle \geqslant c \|x\|^{2}\) 对属于 u 的定义域的所有 x 成立。这时称 c 是 u 的一个下界。若 c = 0 ，也称 u 是正部分算子。

我们用 \(\mathcal{A}(\mathrm{E})\) 表示 E 上自伴部分算子所成的集。

注. --- 1) 算子 \(u\) 在 E 上稠定且对称，当且仅当下式成立：

\[
\langle x \mid u (y) \rangle = \langle u (x) \mid y \rangle\tag{2}
\]

（对所有 \((x,y)\in\mathrm{dom}(u)^{2}\) ）。这一公式表明 u 的定义域包含在 \(u^{*}\) 的定义域中，进而 \(u^{*}\) 与 u 在 u 的定义域上重合。特别地，由此得出 \(\langle x|u(x)\rangle\in\mathbf{R}\) 对所有 \(x\in\mathrm{dom}(u)\) 成立。

正如在多个例子中将看到的那样，公式 (2) 常常可以通过直接计算加以验证。另一方面，伴随的定义域的精确确定——唯有它才能判别一个对称算子是否自伴——可能是非常精细的。

\begin{enumerate}
\def\labelenumi{\arabic{enumi})}
\setcounter{enumi}{1}
\item
  自伴部分算子 u 是本质自伴的（参见命题 8）。
\item
  设 \(u\) 是 E 上的对称部分算子。算子 \(u\) 可闭（命题 7, b)）。它满足 \(\text{dom}(u) \subset \text{dom}(u^*)\) ，且 \(u\) 自伴当且仅当 \(\text{dom}(u) = \text{dom}(u^*)\) 。此外， \(\overline{u}\) 的闭包 \(u\) 是对称的，因为 \(\overline{u} \subset u^* = \overline{u}^*\) （命题 8）。
\item
  设 K = C 。设 \(u \in \mathcal{P}(\mathrm{E};\mathrm{E})\) 是稠定部分算子。条件 \(\langle x | u(x) \rangle \in \mathbf{R}\) （对所有 \(x \in \operatorname{dom}(u)\) ）蕴含 u 是对称的（EVT, V, p. 2，注）；特别地，若 \(\langle x | u(x) \rangle \in \mathbf{R}_{+}\) 对所有 \(x \in \operatorname{dom}(u)\) 成立，则 u 是正的。
\item
  设 u 与 v 是 E 上的对称部分算子。若 u 自伴且 \(u \subset v\) ，则 \(v \subset v^{*} \subset u^{*} = u\) ，因此 u = v 。
\item
  本质自伴部分算子 \(u\) 是对称的，因为 \(u \subset \overline{u}\) 蕴含 \(\overline{u} = \overline{u}^{*} \subset u^{*}\) ，从而 \(u \subset u^{*}\) 。
\item
  设 u 与 v 是 E 上的对称部分算子。若 \(u+v\) 稠定（例如 u 或 v 属于 \(\mathcal{L}(\mathrm{E})\) ），则 \(u+v\) 是对称的。一般而言，部分算子 \(u + v\) 不是自伴的，即使 u 与 v 都自伴（IV, p. 347 的习题 9）。
\item
  设 \(u\) 是 E 上的对称部分算子。实数 \(c\) 是 \(u\) 的下界，当且仅当算子 \(u - c \cdot 1_{\mathrm{E}}\) 是正的。
\end{enumerate}

引理 3. --- 假设 K = R 。设 u 是从 E 到 F 中的稠定部分算子。

\begin{enumerate}
\def\labelenumi{\alph{enumi})}
\item
\(u_{(\mathbf C)}\) 的伴随是 \((u^*)_{(\mathbf C)}\) ；
\item
  设 E = F ；部分算子 u 是对称的（相应地，自伴的），当且仅当部分算子 \(u_{(\mathbf{C})}\) 是对称的（相应地，自伴的）。
\end{enumerate}

我们来证明 \(a\) )。设 \(y \in \mathrm{F}_{(\mathbf{C})}\) ，并把它写成 \(y = y_1 + iy_2\) ，其中 \(y_1, y_2 \in \mathrm{F}\) 。对所有 \((x_1, x_2) \in \mathrm{E} \times \mathrm{E}\) ，我们有

\[
\begin{array}{c} \langle u _ {(\mathbf {C})} (x _ {1} + i x _ {2})   |   y \rangle = \langle u (x _ {1})   |   y _ {1} \rangle + i \langle u (x _ {1})   |   y _ {2} \rangle \\ - i \langle u (x _ {2})   |   y _ {1} \rangle + \langle u (x _ {2})   |   y _ {2} \rangle . \end{array}
\]

若 \(y \in \text{dom}(u^*)(\mathbf{C})\) ，我们推出 \(y \in \text{dom}((u_{(\mathbf{C})})^*)\) 且 \(u_{(\mathbf{C})}^*(y) = (u_{(\mathbf{C})})^*(y)\) ，从而 \(u_{(\mathbf{C})}^* \subset (u_{(\mathbf{C})})^*\) 。

反之，设 \(y\) 属于 \(\text{dom}((u_{(\mathbf{C})})^*)\) 。在上式中取 \(x_2 = 0\) （相应地， \(x_1 = 0\) ），即可验证 \(y_1 \in \text{dom}(u^*)\) （相应地， \(y_2 \in \text{dom}(u^*)\) ），从而 \(y \in \text{dom}(u^*)(\mathbf{C})\) 。

断言 \(a\) ) 蕴含：若 \(u_{(\mathbf{C})}\) 是对称的（相应地，自伴的），则 \(u\) 也是对称的（相应地，自伴的）。

反之，设 \(u_{(\mathbf{C})}\) 是对称的。关系式 \(\langle u(x)|y\rangle = \langle x|u(y)\rangle\) （对所有 \((x,y) \in \mathrm{dom}(u_{(\mathbf{C})}) \times \mathrm{dom}(u_{(\mathbf{C})})\) ）蕴含：在取 x 与 y 属于 \(\mathrm{dom}(u)\) （ \(\mathrm{dom}(u_{(\mathbf{C})})\) 的子空间）时， u 是对称的。若 \(u_{(\mathbf{C})}\) 自伴，则前述论证表明 u 是对称的；由于 \(\mathrm{dom}(u^{*}) = \mathrm{dom}(u_{(\mathbf{C})}^{*}) \cap \mathrm{F}\) ，断言 a) 蕴含 \(\mathrm{dom}(u^{*}) = \mathrm{dom}(u_{(\mathbf{C})}) \cap \mathrm{F} = \mathrm{dom}(u)\) ，因此 u 是自伴的。

\subsection{自伴算子的初等判别法}

命题 9. --- 设 \(v \in \mathcal{L}(\mathrm{F};\mathrm{E})\) 是从 F 到 E 中的连续单射线性映射，其像在 E 中稠密。 \(v\) 的伴随是从 E 到 F 中的连续单射线性映射，且有 \((v^{*})^{-1} = (v^{-1})^{*}\) 。特别地，若 E = F ，则自同态 \(v\) 是埃尔米特的，当且仅当部分算子 \(v^{-1}\) 是自伴的。

部分算子 \(v^{-1}\) 是从 E 到 F 中的闭稠定算子（IV, p. 228 的例 2），而 v 的伴随 \(v^{*}\) 是从 E 到 F 中的连续线性映射；它是单射的，因为 v 的像在 E 中稠密（EVT, V, p. 41，命题 4）。设 s （相应地， \(s'\) ）是从 E ⊕ F 到 F ⊕ E 上的等距同构 \((x,y)\mapsto(-y,x)\) （相应地，从 F ⊕ E 到 E ⊕ F 上的等距同构 \((y,x)\mapsto(-x,y)\) ），并设 \(\iota\) （相应地， \(\iota'\) ）是从 F ⊕ E 到 E ⊕ F 上的等距同构 \((y,x)\mapsto(x,y)\) （相应地，从 E ⊕ F 到 F ⊕ E 上的等距同构 \((x,y)\mapsto(y,x)\) ）。这时有 \(s\circ\iota=-\iota'\circ s'\) ，由此

\[
\Gamma_ {(v ^ {- 1}) ^ {*}} = s \left(\Gamma_ {v ^ {- 1}}\right) ^ {\circ} = s \left(\iota \left(\Gamma_ {v}\right)\right) ^ {\circ} = - \iota^ {\prime} \left(s ^ {\prime} \left(\Gamma_ {v}\right)\right) ^ {\circ} = - \iota^ {\prime} \left(\Gamma_ {v ^ {*}}\right) = \Gamma_ {(v ^ {*}) ^ {- 1}}
\]

（依据 IV, p. 236 的命题 7）。命题得证。

命题 10（海林格–特普利茨）. --- 设 u 是希尔伯特空间 E 上的对称部分算子。若 u 的定义域等于 E ，则 \(u \in \mathcal{L}(E)\) ，且 u 是埃尔米特的。

事实上，部分算子 u 是可闭的（IV, p. 237 的命题 8），而由于它的定义域是 E ，它是闭的（IV, p. 228，注）。于是引用 EVT, I, p. 19，推论 5 即得结论。

推论. --- 设 u 是 E 上的对称部分算子。若部分算子 u 诱导一个从 dom(u) 到 E 上的双射线性映射，则 u 是自伴的。

事实上， u 的逆部分算子 \(u^{-1}\) 是对称的且以 E 为定义域（命题 9），于是 \(u^{-1}\) 是 \(\mathcal{L}(\mathrm{E})\) 中的自伴元素（命题 10），从而 u 是埃尔米特的（命题 9）。

命题 11. --- 设 u 是 E 上的对称部分算子，并设 \(\lambda \in \mathbf{C}\) 。若 \(u + \lambda 1_{\mathrm{E}}\) 与 \(u + \overline{\lambda} 1_{\mathrm{E}}\) 都是满射的，则 u 是自伴的。

只须证明 \(\operatorname{dom}(u^{*}) \subset \operatorname{dom}(u)\) 。设 \(x \in \operatorname{dom}(u^{*})\) 。由假设，存在 \(y \in \operatorname{dom}(u)\) 使得 \(u(y) + \overline{\lambda}y = u^{*}(x) + \overline{\lambda}x\) 。我们来证明 y = x 。对所有 \(z \in \operatorname{dom}(u)\) ，得到

\[
\begin{array}{r l} & {\langle (u + \lambda 1 _ {\mathrm{E}}) (z) | x \rangle = \langle z | (u ^ {*} + \overline {{\lambda}} 1 _ {\mathrm{E}}) (x) \rangle} \\ & {\qquad = \langle z | (u + \overline {{\lambda}} 1 _ {\mathrm{E}}) (y) \rangle = \langle (u + \lambda 1 _ {\mathrm{E}}) (z) | y \rangle ,} \end{array}
\]

（因为 u 是对称的）。由于算子 \(u + \lambda1_{E}\) 是满射的，确有 y = x ，从而 \(x \in \text{dom}(u)\) 。

我们稍后将看到（参见 IV, p. 248 的命题 17）：若 \(\lambda \in \mathbf{C} - \mathbf{R}\) ，则逆命题成立。

命题 12. --- 设 u 是从 E 到 F 中的闭稠定算子。 E 上的部分算子 \(u^{*}\circ u\) 是自伴且正的。它的定义域是 u 的一个核。

用 \(v\) 表示部分算子 \(1_{\mathrm{E}} + u^{*} \circ u\) 。它的定义域是 \(\operatorname{dom}(u^{*} \circ u)\) ，包含在 \(\operatorname{dom}(u)\) 中。对所有 \(x \in \operatorname{dom}(u)\) 与 \(y \in \operatorname{dom}(v)\) ，我们有 \(u(y) \in \operatorname{dom}(u^{*})\) 以及

\begin{enumerate}
\def\labelenumi{(\arabic{enumi})}
\setcounter{enumi}{2}
\tightlist
\item
\end{enumerate}

\[
\langle x \mid v (y) \rangle = \langle x \mid y \rangle + \langle x \mid (u ^ {*} \circ u) (y) \rangle = \langle x \mid y \rangle + \langle u (x) \mid u (y) \rangle ,\tag{4}
\]

\[
\langle y \mid v (y) \rangle = \| y \| ^ {2} + \| u (y) \| ^ {2}.
\]

公式 (4) 蕴含部分算子 \(v\) 是单射的。此外，由 IV, p. 236 的命题 7 \(a\) )，有 \(\mathrm{F} \oplus \mathrm{E} = \Gamma_{u^*} \oplus s(\Gamma_u)\) 。设 \(x \in \mathrm{E}\) 。存在 \(y' \in \mathrm{dom}(u^*)\) 与 \(x' \in \mathrm{dom}(u)\) 使得

\[
(0, x) = (y ^ {\prime}, u ^ {*} (y ^ {\prime})) + (- u (x ^ {\prime}), x ^ {\prime}) = (y ^ {\prime} - u (x ^ {\prime}), x ^ {\prime} + u ^ {*} (y ^ {\prime})).
\]

因此我们有 \(y' = u(x')\) ，从而 \(x' \in \text{dom}(u^* \circ u) = \text{dom}(v)\) 以及

\[
x = x ^ {\prime} + u ^ {*} (y ^ {\prime}) = x ^ {\prime} + (u ^ {*} \circ u) (x ^ {\prime}) = v (x ^ {\prime}).
\]

因此， E 上的部分算子 v 是满射的，并诱导一个从 dom(v) 到 E 上的双射线性映射。

设 \(x \in \mathrm{E}\) 与 \(\operatorname{dom}(v)\) 正交。把 x 写成 \(x = v(x')\) ，其中 \(x' \in \operatorname{dom}(v)\) 。由公式 (4) ，我们得到

\[
0 = \langle x ^ {\prime} | x \rangle = \langle x ^ {\prime} | v (x ^ {\prime}) \rangle = \| x ^ {\prime} \| ^ {2} + \| u (x ^ {\prime}) \| ^ {2}
\]

由此 \(x' = 0\) ，进而 \(x = 0\) 。因此 \(v\) 的定义域在 E 中稠密。

部分算子 \(v\) 稠定；它是双射的，且公式 (3) 表明它是对称的。对 \(v\) 应用命题 10 的推论，我们断定它是自伴的。于是 \(u^{*} \circ u = v - 1_{\mathrm{E}}\) 是自伴的。此外，公式

\[
\langle x \mid (u ^ {*} \circ u) (x) \rangle = \| u (x) \| ^ {2}
\]

（对所有 \(x \in \text{dom}(u^{*} \circ u)\) ）蕴含 \(u^{*} \circ u\) 是正的。

最后，设 \(y \in \mathrm{E}_u\) 与 \(\operatorname{dom}(u^* \circ u)\) 正交。存在 \(x\) 的定义域中的元素 \(u^* \circ u\) 使得 \(y = v(x) = x + (u^* \circ u)(x)\) 。这时我们有 \(0 = (x \mid y)_u = \langle y \mid y \rangle\) ，由此 \(y = 0\) 。

\subsection{微分算子}

设 \(n \in N\) ， U 是 \(R^{n}\) 的开子集。我们给 \(R^{n}\) 与 U 赋以勒贝格测度，记作 \(\mu\) 。

设 \(k \in N\) 与 \(h \in N\) 满足 \(h \geqslant k\) 。设 D 是 U 上阶 \(\leqslant k\) 的纯量微分算子，其系数 \((n_{\alpha})_{|\alpha| \leqslant k}\) 在 U 上属于类 \(C^{h}\) 。假设对每个满足 \(\alpha\) 的 \(|\alpha| \leqslant k\) 以及每个满足 \(\beta\) 的 \(0 \leqslant \beta \leqslant \alpha\) ，函数 \(\partial^{\beta} n_{\alpha}\) 在 U 上有界。

设 \({}^{t}D\) 是 D 的转置纯量微分算子（VAR, R2, 14.3.2）；它的阶 \(\leqslant k\) ，属于类 \(C^{h-k}\) ；对 \(\varphi\in\mathcal{D}(U)\) ，我们有

\[
{ } ^ { t } \mathrm{D} ( \varphi ) = \sum _ { | \alpha | \leqslant k } ( - 1 ) ^ { | \alpha | } \partial ^ { \alpha } ( \overline { { n } } _ { \alpha } \varphi )
\]

（loc. cit.）；特别地， \({}^{t}D\) 的系数在 U 上有界。

我们用 D\_ 表示 \(L^{2}(U)\) 上以 \(\mathcal{D}(U)\) 为定义域、由下式定义的部分算子：

\[
\varphi \mapsto \mathrm{D} (\varphi) = \sum_ {| \alpha | \leqslant k} n _ {\alpha} \partial^ {\alpha} \varphi .\tag{5}
\]

设 \(H_{D}\) 是这样的 \(f \in L^{2}(U)\) 所成的空间：下列广义函数

\[
\mathrm{D} (f) = \sum_ {| \alpha | \leqslant k} n _ {\alpha} \partial^ {\alpha} f
\]

属于 \(L^{2}(U)\) ；我们用 \(D_{+}\) 表示以 \(H_{D}\) 为定义域、由 \(f \mapsto \mathrm{D}(f)\) 所定义的部分算子。

由于我们有 \(\partial^{\alpha}f\in \mathrm{L}^{2}(\mathrm{U})\) （若 \(f\in \mathrm{H}^k (\mathrm{U})\) 且 \(|\alpha |\leqslant k\) ），索伯列夫空间 \(\mathrm{H}^k (\mathrm{U})\) 包含在 \(\mathrm{H}_{\mathrm{D}}\) 中；一般而言，这两个空间是不同的。

我们有 \(D_{-} \subset D_{+}\) ，它们都是 D 在 \(L^{2}(U)\) 上的相伴微分算子（IV, p. 235 的定义 5）。

命题 13. --- 设 u 是 L²(U) 上的部分算子。若 D₋ ⊂ u ⊂ D₊ ，则 u 可闭，且 (ᵗD)₋ ⊂ u* ⊂ (ᵗD)₊ 。

设 \(\varphi\) 与 \(\psi\) 属于 \(\mathcal{D}(\mathrm{U})\) 。这时我们有

\[
\begin{array}{c} \langle \varphi | \mathrm{D} (\psi) \rangle = \sum_ {| \alpha | \leqslant k} \int_ {\mathrm{U}} n _ {\alpha} \overline {{\varphi}} \partial^ {\alpha} \psi d \mu \\ = \sum_ {| \alpha | \leqslant k} (- 1) ^ {| \alpha |} \langle \partial^ {\alpha} (\overline {{n}} _ {\alpha} \varphi) | \psi \rangle = \langle {} ^ {t} \mathrm{D} (\varphi) | \psi \rangle \end{array}
\]

（参见 VAR, R2, 14.3.8）。由于 \(\mathcal{D}(\mathrm{U})\) 在 \(\mathrm{L}^2 (\mathrm{U})\) 中稠密，这蕴含 \(\varphi \in \mathrm{dom}(u^{*})\) 且 \(u^{*}(\varphi) = {}^{t}\mathrm{D}(\varphi)\) 。因此我们有 \(({}^{t}\mathrm{D})_{-}\subset u^{*}\) 。特别地， \(u^{*}\) 稠定而 \(u\) 可闭（IV, p. 237 的命题 8）。

设 \(f \in \mathrm{dom}(u^{*})\) 与 \(\varphi \in \mathcal{D}(\mathrm{U})\) 。由于 \(\mathcal{D}(\mathrm{U}) \subset \mathrm{dom}(u)\) ，与 \(u^{*}(f)\) 相伴的广义函数满足

\[
\langle u ^ {*} (f), \varphi \rangle = \langle \overline {{\varphi}} | u ^ {*} (f) \rangle = \langle u (\overline {{\varphi}}) | f \rangle .
\]

由于 \(D_{-} \subset u\) ，可由公式 (5) 计算 \(u(\overline{\varphi})\) ，由此

\[
\begin{array}{l} \langle u ^ {*} (f), \varphi \rangle = \sum_ {\alpha} \langle n _ {\alpha} \partial^ {\alpha} \overline {{\varphi}} | f \rangle = \sum_ {\alpha} \langle \partial^ {\alpha} \overline {{\varphi}} | \overline {{n}} _ {\alpha} f \rangle \\ \qquad = \sum_ {\alpha} \langle \overline {{n}} _ {\alpha} f, \partial^ {\alpha} \varphi \rangle = \sum_ {\alpha} (- 1) ^ {| \alpha |} \langle \partial^ {\alpha} (\overline {{n}} _ {\alpha} f), \varphi \rangle = \langle {} ^ {t} D (f), \varphi \rangle . \end{array}
\]

于是广义函数 \(u^{*}(f)\) 与 \({}^{t}\mathrm{D}(f)\) 相等；因而广义函数 \(f\) 属于 \(H_{tD}\) ，且 \(u^{*}(f)={}^{t}\mathrm{D}(f)\) ，由此 \(u^{*}\subset({}^{t}\mathrm{D})_{+}\) 。

注. --- 这一命题表明，部分算子 \(u\) （它满足 \(\mathrm{D}_{-} \subset u \subset \mathrm{D}_{+}\) ）的伴随总可以在广义函数的意义下计算：元素 \(f\) （它们属于 \(u^{*}\) 的定义域）是 \(\mathrm{L}^{2}(\mathrm{U})\) 中这样的元素：广义函数 \(^{t}\mathrm{D}(f)\) 属于 \(\mathrm{L}^{2}(\mathrm{U})\) ，且我们有 \(u^{*}(f) = ^{t}\mathrm{D}(f)\) 。

若作为纯量微分算子有 \({}^{t}D = D\) ，就称 D 是形式对称的。在这种情形，部分算子 D\_ 是对称的。

考虑由下式定义的二阶纯量微分算子的特殊情形：

\[
\Delta = - \sum_ {i = 1} ^ {n} \partial_ {i} ^ {2}.
\]

拉普拉斯算子（ U 上的），指的是 \(L^{2}(U)\) 上自伴的部分算子 u ，满足 \(\Delta_{-} \subset u\) （参见 VAR, R2, 14.4.3, p. 83）。我们稍后将看到（IV, p. 261，例）， \(L^{2}(U)\) 上总存在至少一个拉普拉斯算子；也可能存在不止一个（IV, p. 358 的习题 17）。
\subsection{谱与预解式}

引理 4. --- 设 E 为复巴拿赫空间，u 为 E 上的闭单射部分算子，满足 \(u^{-1} \in \mathcal{L}(\mathrm{E})\)。设 \(v \in \mathcal{L}(\mathrm{E})\)

满足 \(\|v\| < \|u^{-1}\|^{-1}\)。则部分算子 \(u + v\) 是单射的，我们有 \((u + v)^{-1} \in \mathcal{L}(\mathrm{E})\)，并且

\[
(u + v) ^ {- 1} = u ^ {- 1} \circ \sum_ {k = 0} ^ {+ \infty} (- v u ^ {- 1}) ^ {k},\tag{6}
\]

其中该级数在 \(\mathcal{L}(\mathrm{E})\) 中绝对收敛。此外，我们有

\[
\| (u + v) ^ {- 1} \| \leqslant \frac {\| u ^ {- 1} \|}{1 - \| v \| \| u ^ {- 1} \|}.
\]

由于 \(\|v\|\|u^{-1}\|<1\)，通项为 \((-vu^{-1})^{k}\) 的级数在 \(\mathcal{L}(\mathrm{E})\) 中绝对收敛，且其和是自同态 \(1_{E}+vu^{-1}\) 的逆（I，第 22 页，命题 2）。于是，部分算子 \((1_{\mathrm{E}}+vu^{-1})\circ u=u+v\)（IV，第 226 页，注 vi）是单射的，且逆部分算子 \((u+v)^{-1}=u^{-1}\circ(1_{\mathrm{E}}+vu^{-1})^{-1}\)（IV，第 226 页，注 iii）属于 \(\mathcal{L}(\mathrm{E})\)。由于

\[
\left\| \left(1 _ {\mathrm{E}} + v u ^ {- 1}\right) ^ {- 1} \right\| \leqslant \sum_ {k = 0} ^ {+ \infty} (\left\| v \right\| \left\| u ^ {- 1} \right\|) ^ {k} = \frac {1}{1 - \left\| v \right\| \left\| u ^ {- 1} \right\|},
\]

引理得证。

定义 8. --- 设 u 是复巴拿赫空间 E 上的稠定闭算子。u 的预解集是由这样的复数 λ 组成的集合：部分算子 \(\lambda 1_{\mathrm{E}} - u\) 是单射的，且其逆 \((\lambda 1_{\mathrm{E}} - u)^{-1}\) 属于 \(\mathcal{L}(\mathrm{E})\)。

u 的谱，记作 \(\mathrm{Sp}(u)\)，是预解集在 C 中的余集。

若 \(\lambda \in \mathbf{C} - \mathrm{Sp}(u)\)，我们用 \(\mathrm{R}(u,\lambda) \in \mathcal{L}(\mathrm{E})\) 表示 \(\lambda 1_{\mathrm{E}} - u\) 的逆。从 \(\mathbf{C} - \mathrm{Sp}(u)\) 到 \(\mathcal{L}(\mathrm{E})\) 中把 \(\lambda\) 映为 \(\mathrm{R}(u,\lambda)\) 的映射称为 \(u\) 的预解式。

注记。--- 设 E 为复巴拿赫空间，u 为 E 上稠定的闭算子。

\begin{enumerate}
\def\labelenumi{\arabic{enumi})}
\item
  若 \(u \in \mathcal{L}(\mathrm{E})\)，则它的谱与元素 \(u\) 作为代数 \(\mathcal{L}(\mathrm{E})\) 中元素的谱一致（I，第 2 页，定义 1）。
\item
  设 D 为 \(u\) 的定义域。对每个 \(\lambda \in \mathbf{C} - \mathrm{Sp}(u)\)，我们有
\end{enumerate}

\[
1 _ {\mathrm{E}} = (\lambda 1 _ {\mathrm{E}} - u) \circ \mathrm{R} (u, \lambda), \quad 1 _ {\mathrm{D}} = \mathrm{R} (u, \lambda) \circ (\lambda 1 _ {\mathrm{E}} - u).\tag{7}
\]

此外，对于 u 的预解集中的 \(\lambda_{1}\) 与 \(\lambda_{2}\)，我们有

\[
\mathrm{R} (u, \lambda_ {1}) - \mathrm{R} (u, \lambda_ {2}) = (\lambda_ {2} - \lambda_ {1}) \mathrm{R} (u, \lambda_ {2}) \circ \mathrm{R} (u, \lambda_ {1}),\tag{8}
\]

\[
\mathrm{R} (u, \lambda_ {1}) \circ \mathrm{R} (u, \lambda_ {2}) = \mathrm{R} (u, \lambda_ {2}) \circ \mathrm{R} (u, \lambda_ {1}),\tag{9}
\]

事实上，有

\[
\mathrm{R} (u, \lambda_ {1}) - \mathrm{R} (u, \lambda_ {2}) = \mathrm{R} (u, \lambda_ {1}) \circ 1 _ {\mathrm{E}} - 1 _ {\mathrm{D}} \circ \mathrm{R} (u, \lambda_ {2}).
\]

由于 \(1_{\mathrm{E}} = (\lambda_21_{\mathrm{E}} - u)\circ \mathrm{R}(u,\lambda_2)\) 与 \(1_{\mathrm{D}} = \mathrm{R}(u,\lambda_1)\circ (\lambda_11_{\mathrm{E}} - u)\)，我们得到

\[
\begin{array}{c} \mathrm{R} (u, \lambda_ {1}) - \mathrm{R} (u, \lambda_ {2}) = \lambda_ {2} \mathrm{R} (u, \lambda_ {1}) \circ \mathrm{R} (u, \lambda_ {2}) \\ - \mathrm{R} (u, \lambda_ {1}) \circ u \circ \mathrm{R} (u, \lambda_ {2}) - \lambda_ {1} \mathrm{R} (u, \lambda_ {1}) \circ \mathrm{R} (u, \lambda_ {2}) \\ + \mathrm{R} (u, \lambda_ {1}) \circ u \circ \mathrm{R} (u, \lambda_ {2}), \end{array}
\]

由此得到第一个公式。交换 \(\lambda_{1}\) 与 \(\lambda_{2}\) 的角色，便推出第二个公式。

\begin{enumerate}
\def\labelenumi{\arabic{enumi})}
\setcounter{enumi}{2}
\item
  设 \(\lambda \in \mathbf{C}\)。由闭图像定理（EVT，I，第 19 页，推论 5），若从 \(\mathrm{dom}(u)\) 到 E 中的由 \(x \mapsto (\lambda 1_{\mathrm{E}} - u)(x)\) 定义的线性映射是双射，则它的逆（其图像是闭的）从 E 到 E 连续。因此，\(\lambda\) 属于 \(u\) 的预解集，当且仅当由 \(x \mapsto (\lambda 1_{\mathrm{E}} - u)(x)\) 给出的从 \(\mathrm{dom}(u)\) 到 E 中的映射是双射。
\item
  若 \(\lambda\) 属于 u 的谱，则下列性质之一成立：
\end{enumerate}

\begin{enumerate}
\def\labelenumi{(\roman{enumi})}
\item
  \(\lambda 1_{\mathrm{E}} - u\) 的核不化为 0；这时称 \(\lambda\) 是 \(u\) 的特征值，并称 \(\operatorname{Ker}(\lambda 1_{\mathrm{E}} - u)\) 的维数为其重数；
\item
  部分算子 \(\lambda 1_{\mathrm{E}} - u\) 是单射的，而其像在 E 中不稠密；这时称 \(\lambda\) 属于 \(u\) 的剩余谱；
\item
  部分算子 \(\lambda 1_{\mathrm{E}} - u\) 是单射的，其像在 E 中稠密，但 \(\lambda 1_{\mathrm{E}} - u\) 不是满射；这时称 \(\lambda\) 属于 \(u\) 的连续谱。
\end{enumerate}

\begin{enumerate}
\def\labelenumi{\arabic{enumi})}
\setcounter{enumi}{4}
\tightlist
\item
  设 \(\lambda\) 是属于 \(u\) 的预解集的复数。预解式 \(\mathrm{R}(u,\lambda)\) 是从 E 到 E 中的单射线性映射；其像是 \(u\) 的定义域，且 \(u = \lambda 1_{\mathrm{E}} - \mathrm{R}(u,\lambda)^{-1}\)（参见 IV，第 226 页）。反之，这一性质刻画了预解集与预解式。准确地说，设 \(\lambda \in \mathbf{C}\)；若存在从 E 到 E 中的连续单射线性映射 w，使得 \(u = \lambda 1_{E} - w^{-1}\)，则 \(\lambda\) 属于 u 的预解集，且 \(w = \mathrm{R}(u, \lambda)\)。
\end{enumerate}

特别地，若 v 是 E 上的闭稠定算子，且 \(\lambda \in C\) 是不属于 \(\operatorname{Sp}(u) \cup \operatorname{Sp}(v)\) 的复数，则等式 \(\operatorname{R}(u, \lambda) = \operatorname{R}(v, \lambda)\) 蕴含 u = v。

\begin{enumerate}
\def\labelenumi{\arabic{enumi})}
\setcounter{enumi}{5}
\tightlist
\item
  我们把可闭部分算子 \(u\) 的谱定义为其闭包的谱。
\end{enumerate}

存在谱为空集、或谱等于 C 的闭算子（IV，第 347 页，习题 12）。

设 E 为复巴拿赫空间，u 为 E 上的闭稠定算子。若 F 是复巴拿赫空间，且 \(v: E \to F\) 是同构，则我们有 \(\mathrm{Sp}(v \circ u \circ v^{-1}) = \mathrm{Sp}(u)\) 与 \(\mathrm{R}(v \circ u \circ v^{-1}, \lambda) = v \circ \mathrm{R}(u, \lambda) \circ v^{-1}\)，对一切 \(\lambda \notin \mathrm{Sp}(u)\) 成立。

命题 14. --- 设 E 为复巴拿赫空间。设 u 为 E 上稠定的闭算子，U 为其预解集。a) 对每个 \(\lambda\in\mathrm{U}\)，C 中以 \(\lambda\) 为心、以 \(\|\mathrm{R}(u,\lambda)\|^{-1}\) 为半径的开圆盘都包含于 U；

\begin{enumerate}
\def\labelenumi{\alph{enumi})}
\setcounter{enumi}{1}
\item
  集合 U 是 C 中的开集；
\item
  设 \(\operatorname{Sp}(u)\) 非空。设 \(\lambda \in \mathrm{U}\)，\(\delta\) 是 \(\mathbf{C}\) 中从 \(\lambda\) 到 \(u\) 的谱的距离。我们有 \(\delta > 0\) 及 \(\| \mathrm{R}(u, \lambda) \| \geqslant 1 / \delta\)；
\item
  映射 \(\lambda \mapsto \mathrm{R}(u,\lambda)\) 是从 U 到 \(\mathcal{L}(\mathrm{E})\) 中的全纯映射。对每个整数 \(k\in \mathbf{N}\) 与每个 \(\lambda \in \mathrm{U}\)，我们有
\end{enumerate}

\[
\frac {\partial^ {k}}{\partial \lambda^ {k}} \mathrm{R} (u, \lambda) = (- 1) ^ {k} k! \mathrm{R} (u, \lambda) ^ {k + 1}.
\]

设 \(\lambda \in \mathrm{U}\)。对每个 \(\mu \in \mathbf{C}\)，只要 \(\| (\mu -\lambda)1_{\mathrm{E}}\| < \| \mathrm{R}(u,\lambda)\|^{-1}\)，将引理 4 应用于单射部分算子 \(\lambda 1_{\mathrm{E}} - u\) 并取 \(v = (\mu -\lambda)1_{\mathrm{E}}\)，便知部分算子 \(\mu 1_{\mathrm{E}} - u = \lambda 1_{\mathrm{E}} - u + v\) 是单射的且具有连续的逆。这就蕴含 \(a\)）。由同上所引，我们还有公式

\[
\mathrm{R} (u, \mu) = \mathrm{R} (u, \lambda) \circ \sum_ {k \in \mathbf {N}} (\lambda - \mu) ^ {k} \mathrm{R} (u, \lambda) ^ {k},
\]

由此推出 u 的预解式在 U 中全纯。

断言 \(b\)）由 \(a\)）立得。若 \(\mathrm{Sp}(u)\) 非空，则 \(\lambda\) 到 \(\mathrm{Sp}(u)\) 的距离严格为正（TG，IX，第 13 页，命题 2），由此得 \(c\)）。

断言 \(d\)）的最后一部分对 \(k\) 用归纳法证明，其中 \(k = 1\) 的情形是第 245 页公式 (8) 的推论。

命题 15. --- 设 u 是复巴拿赫空间 E 上稠定的闭算子，λ 是属于 u 的预解集的复数。

\begin{enumerate}
\def\labelenumi{\alph{enumi})}
\item
  子集 \(\mathrm{Sp}(\mathrm{R}(u,\lambda)) \setminus \{0\}\)（\(\mathbf{C}\) 中的）是 \(u\) 的谱经映射 \(\mu \mapsto (\lambda -\mu)^{-1}\)（从 \(\mathbf{C} \setminus \{\lambda \}\) 到 \(\mathbf{C}^*\) 中）所成的像；
\item
  对每个 \(\mu \neq \lambda\)（在 \(\mathbf{C}\) 中），有
\end{enumerate}

\[
\operatorname{Ker} \left(\mu 1 _ {\mathrm{E}} - u\right) = \operatorname{Ker} \left(\left(\lambda - \mu\right) ^ {- 1} 1 _ {\mathrm{E}} - \mathrm{R} (u, \lambda)\right).
\]

我们来证明断言 \(a\)）。对每个 \(\mu \neq \lambda\)，计算

\[
\mu 1 _ {\mathrm{E}} - u = (\lambda - \mu) ((\lambda - \mu) ^ {- 1} 1 _ {\mathrm{E}} - \mathrm{R} (u, \lambda)) (\lambda 1 _ {\mathrm{E}} - u).
\]

由于 \(\lambda \notin \mathrm{Sp}(u)\) 且 \(\mu \neq \lambda\)，线性映射 \((\lambda - \mu)(\lambda 1_{\mathrm{E}} - u)\) 是从 \(\mathrm{dom}(u)\) 到 E 上的双射。于是，该公式表明：\(\mu 1_{\mathrm{E}} - u\) 是从 \(\mathrm{dom}(u)\) 到 E 上的双射，当且仅当 \((\lambda - \mu)^{-1} 1_{\mathrm{E}} - \mathrm{R}(u, \lambda)\) 是从 E 到 E 上的双射，由此即得断言。

我们来证明 \(b\)）。若 \(\mu \neq \lambda\) 且 \(x \in \mathrm{Ker}((\lambda - \mu)^{-1}1_{\mathrm{E}} - \mathrm{R}(u, \lambda))\)，则 \(x \in \mathrm{dom}(u)\)，且公式 \(1_{\mathrm{E}} = (\lambda 1_{\mathrm{E}} - u) \circ \mathrm{R}(u, \lambda)\) 蕴含 \(x \in \mathrm{Ker}(\mu 1_{\mathrm{E}} - u)\)。反之，若 \(x \in \mathrm{Ker}(\mu 1_{\mathrm{E}} - u)\) 且 \(\mu \neq \lambda\)，则公式 \(1_{\mathrm{dom}(u)} = \mathrm{R}(u, \lambda) \circ (\lambda 1_{\mathrm{E}} - u)\) 蕴含 \(\mathrm{R}(u, \lambda)(x) = (\lambda - \mu)^{-1}x\)。

命题 16. --- 设 u 是复希尔伯特空间 E 上稠定的闭算子。\(u^{*}\) 的谱是 u 的谱在复共轭下的像，且对 u 的预解集的每个元素 \(\lambda\)，有 \(\mathrm{R}(u,\lambda)^{*}=\mathrm{R}(u^{*},\overline{\lambda})\)。特别地，若 u 自伴，则自同态 \(\mathrm{R}(u,\lambda)\) 对每个 \(\lambda\notin\mathrm{Sp}(u)\) 都是正规的。

设 \(\lambda\in\mathbf{C} - \mathrm{Sp}(u)\) 是 \(u\) 的预解集中的一个元素。我们有 \(u=\lambda1_{\mathrm{E}}-\mathrm{R}(u,\lambda)^{-1}\)，于是

\[
u ^ {*} = \overline {{\lambda}} 1 _ {\mathrm{E}} - (\mathrm{R} (u, \lambda) ^ {- 1}) ^ {*} = \overline {{\lambda}} 1 _ {\mathrm{E}} - (\mathrm{R} (u, \lambda) ^ {*}) ^ {- 1}
\]

（IV，第 236 页及 IV，第 239 页命题 9）。由注 5 推出 \(\overline{\lambda} \in \mathbf{C} - \mathrm{Sp}(u^{*})\) 且 \(\mathrm{R}(u, \lambda)^{*} = \mathrm{R}(u^{*}, \overline{\lambda})\)。因此，\(u^{*}\) 的谱包含于 \(\mathrm{Sp}(u)\) 经复共轭所成的像。把该性质应用于 \(u^{*}\) 即得等式，因为 \(u^{**} = u\)（IV，第 237 页命题 8 的推论）。最后的断言则由第 245 页公式 (9) 得出。

推论。--- 设 u 是复希尔伯特空间 E 上的自伴部分算子。若 E 非零，则 u 的谱非空。

设 \(\operatorname{Sp}(u)\) 为空集。则 u 是单射的，且 \(u^{-1} = -\operatorname{R}(u,0)\) 是 E 的单射自同态，满足 \(\operatorname{Sp}(u^{-1}) \subset \{0\}\)（命题 15，a）），从而 \(\operatorname{Sp}(u^{-1}) = \{0\}\)（I，第 26 页，推论 1）。由于 \(u^{-1}\) 正规（命题 16），这就蕴含 \(u^{-1}\) 为零（I，第 110 页，例 1），此为矛盾。

引理 5. --- 设 E 为复希尔伯特空间，u 为 E 上的闭稠定部分算子。设 \(\lambda \in \mathbf{C}\)。假设存在实数 \(c > 0\)，使得

\[
\| u (x) - \lambda x \| \geqslant c \| x \| \quad \text{对一切 } x \in \operatorname{dom} (u),\tag{10}
\]

\[
\| u ^ {*} (x) - \overline {{\lambda}} x \| \geqslant c \| x \| \quad \text{对一切 } x \in \operatorname{dom} (u ^ {*}),\tag{11}
\]

则 \(\lambda\) 属于 \(u\) 的预解集，且 \(\| \mathrm{R}(u,\lambda)\| \leqslant c^{-1}\)。

假设蕴含 \(u - \lambda 1_{\mathrm{E}}\) 与 \(u^{*} - \overline{\lambda} 1_{\mathrm{E}}\) 都是单射的。设 F 为 \(u - \lambda 1_{\mathrm{E}}\) 的像。空间 F 在 E 中稠密，因为其正交补等于 \(\operatorname{Ker}(u^{*} - \overline{\lambda} 1_{\mathrm{E}})\)（IV，第 236 页，命题 7，c）），而后者为零。

我们来证明空间 F 是闭的。设 \((x_{n})_{n\in\mathbf{N}}\) 是 dom(u) 中的序列，使序列 \((u(x_{n})-\lambda x_{n})_{n\in\mathbf{N}}\) 收敛于 \(y\in F\)。不等式 (10) 蕴含序列 \((x_{n})_{n\in\mathbf{N}}\) 是 E 中的柯西序列。设 \(x\in E\) 为其极限。序列 \((x_{n},u(x_{n}))\) 在 E×E 中收敛于 \((x,y+\lambda x)\)；由于 u 的图像是闭的，我们因此有 \(x\in\text{dom}(u)\) 与 \(u(x)=y+\lambda x\)，这就证明了 \(y\in F\)。

由此断定 F = E。于是部分算子 \(u - \lambda1_{E}\) 是双射，从而 \(\lambda \notin \operatorname{Sp}(u)\)（注 3）。不等式 (10) 进而又蕴含 \(\|R(u, \lambda)\| \leqslant c^{-1}\)。

命题 17. — 设 E 为复希尔伯特空间，u 为 E 上的自伴部分算子。

\begin{enumerate}
\def\labelenumi{\alph{enumi})}
\item
  u 的谱包含于 \(\mathbf{R}\)；
\item
  若 u 是正的，则 u 的谱包含于 \(\mathbf{R}_{+}\)；
\item
  设 E 非零。设 \(\lambda \notin \operatorname{Sp}(u)\)，且 \(\delta > 0\) 为 \(\lambda\) 到 \(u\) 的谱的距离。我们有 \(\| \mathrm{R}(u, \lambda) \| = \delta^{-1}\)。
\end{enumerate}

设 \((a,b)\in\mathbf{R}\times\mathbf{R}\)，\(\lambda=a+ib\)。设 \(x\in\mathrm{dom}(u)\)。由于 u 自伴，我们有 \(\langle x|u(x)\rangle\in\mathbf{R}\)，由此

\[
\begin{array}{r} \| u (x) - \lambda x \| ^ {2} = \| u (x) \| ^ {2} - 2 a \langle x | u (x) \rangle + (a ^ {2} + b ^ {2}) \| x \| ^ {2} \\ = \| u (x) - \overline {{\lambda}} x \| ^ {2}. \end{array}
\]

设 \(b \neq 0\)。于是得到

\[
\begin{array}{c} \| u (x) - \lambda x \| ^ {2} = \| u (x) - \overline {{\lambda}} x \| ^ {2} \geqslant (\| u (x) \| - a \| x \|) ^ {2} + b ^ {2} \| x \| ^ {2} \\ \geqslant b ^ {2} \| x \| ^ {2}. \end{array}
\]

于是由引理 5 得 \(\lambda \notin \operatorname{Sp}(u)\)，即断言 \(a\)）。

设 u 还是正的。若 b = 0 且 a \textless{} 0，则对 \(x \in \text{dom}(u)\) 类似地得到不等式

\[
\left\| u (x) - \lambda x \right\| ^ {2} = \left\| u (x) - \overline {{\lambda}} x \right\| ^ {2} \geqslant \left(\left\| u (x) \right\| - a \| x \|\right) ^ {2} \geqslant a ^ {2} \| x \| ^ {2},
\]

从而 \(\lambda\notin\mathrm{Sp}(u)\)（同上所引），这就证明了 \(b\)）。

最后证明 \(c\)）。由命题 16，预解式 \(\mathrm{R}(u,\lambda)\) 是 E 的正规自同态。因此其范数等于其谱半径（I，第 108 页，推论 1），由此

\[
\| \mathrm{R} (u, \lambda) \| = \sup _ {\mu \in \operatorname{Sp} (\mathrm{R} (u, \lambda))} | \mu |
\]

（I，第 24 页，定理 1）。R(u,λ) 的谱不可能化为 \{0\}，因为否则将有 \(\| \mathrm{R}(u,\lambda)\| = 0\)，于是作为 R(u,λ) 之像的 \(\operatorname{dom}(u)\) 将为零，E 亦然。因此命题 15 蕴含

\[
\| \mathrm{R} (u, \lambda) \| = \sup _ {\mu \in \mathrm{Sp} (u)} \frac {1}{| \lambda - \mu |} = \frac {1}{\delta}.
\]

推论。--- 设 u 是 E 上的自伴算子。

\begin{enumerate}
\def\labelenumi{\alph{enumi})}
\item
  u 的剩余谱是空集；
\item
  对所有 \(\lambda \neq \mu\)（在 \(\mathbf{C}\) 中），\(u\) 的属于 \(\lambda\) 与 \(\mu\) 的特征空间相互正交。
\end{enumerate}

我们来证明 a)。设 \(\lambda\) 属于 \(u\) 的谱；它是实数（命题 17）。我们有 \(\mathrm{Ker}(\lambda 1_{\mathrm{E}} - u) = \mathrm{Im}(\lambda 1_{\mathrm{E}} - u)^{\circ}\)（IV，第 236 页，命题 7），因此 \(\lambda\) 不是 \(u\) 的特征值，只要部分算子 \(\lambda 1_{\mathrm{E}} - u\) 的像稠密。按定义，这就蕴含 \(u\) 的剩余谱是空集。

我们来证明 \(b\)）。由该命题，可以假设 \(\lambda\) 与 \(\mu\) 是实数。设 \(x \in \operatorname{dom}(u)\) 满足 \(u(x) = \lambda x\)，且 \(y \in \operatorname{dom}(u)\) 满足 \(u(y) = \mu y\)。于是有

\[
\lambda \langle x \mid y \rangle = \langle u (x) \mid y \rangle = \langle x \mid u (y) \rangle = \mu \langle x \mid y \rangle ,
\]

由此 \(\langle x \mid y \rangle = 0.\)

注。--- 若 u 是非自伴的闭对称算子，则其谱不包含于 R（参见下文 IV，第 257 页推论 10），而且 u 的对应于 λ 与 \(\overline{\lambda}\) 的特征空间有可能不正交（IV，第 347 页，习题 11）。

\subsection{伪谱}

定义 9. --- 设 E 为复巴拿赫空间，u 为 E 上的闭部分算子。设 \(\varepsilon\) 为严格正的实数。我们把 u 的 \(\varepsilon\)-伪谱称为 u 的谱与如下的复数 \(\lambda\) 的集合之并：这些复数属于 u 的预解集且满足 \(\| \mathrm{R}(u,\lambda)\| >\varepsilon^{-1}\)。我们把这个集合记作 \(\mathrm{PSp}_{\varepsilon}(u)\)。

\pandocbounded{\includegraphics[keepaspectratio]{images/part002_14758b8b.jpg}}

有些作者把 \(\varepsilon\)-伪谱（对 \(u\)）定义为集合 \(\mathrm{T}_{\varepsilon}(u)\)，即 \(\mathrm{Sp}(u)\) 与如下的 \(\lambda \in \mathbf{C} - \mathrm{Sp}(u)\) 所成集合之并：它们满足 \(\| \mathrm{R}(u,\lambda)\| \geqslant \varepsilon^{-1}\)。\(\mathrm{PSp}_{\varepsilon}(u)\) 的闭包包含于 \(\mathrm{T}_{\varepsilon}(u)\)，但这一包含可以是严格的，即使 E 是希尔伯特空间亦然（参见 IV，第 348 页习题 18 及第 349 页习题 19）。

命题 18. --- 设 E 为复巴拿赫空间，u 为 E 上的闭部分算子。设 \(\varepsilon \in \mathbf{R}_{+}^{*}\)。集合 \(\mathrm{PSp}_{\varepsilon}(u)\) 是 C 的开子集。当 E 非零时它非空。

若 E 为零，则 \(\mathrm{PSp}_{\varepsilon}(u)\) 是空集。设 E 非零。由 IV，第 246 页命题 14，集合 \(\mathrm{PSp}_{\varepsilon}(u)\) 是 C 的开子集。

集合 \(\mathrm{PSp}_{\varepsilon}(u)\) 非空，如果 \(u\) 的谱非空。若 \(\mathrm{Sp}(u)\) 为空集，则刘维尔定理（VAR，R1，第 29 页，3.3.6）蕴含 \(\mathbf{C}\) 上由 \(\lambda \mapsto \mathrm{R}(u,\lambda)\) 定义的全纯函数不有界，故存在 \(\lambda\) 在 \(\mathbf{C}\) 中，使得 \(\| \mathrm{R}(u,\lambda)\| >\varepsilon^{-1}\)。

命题 19. --- 设 E 为复巴拿赫空间，u 为 E 上的闭部分算子。则我们有 \(\mathrm{PSp}_{\varepsilon}(u) \subset \mathrm{PSp}_{\delta}(u)\)（当 \(0 < \varepsilon < \delta\) 时），并且

\[
\bigcap_ {\varepsilon \in \mathbf {R} _ {+} ^ {*}} \mathrm{PSp} _ {\varepsilon} (u) = \mathrm{Sp} (u).
\]

第一个断言由定义立得。至于第二个断言，按定义，u 的谱包含于 \(\mathrm{PSp}_{\varepsilon}(u)\)，对一切 \(\varepsilon > 0\) 成立；而若 \(\lambda \notin \mathrm{Sp}(u)\)，则 \(\lambda \notin \mathrm{PSp}_{\varepsilon}(u)\)，只要 \(\varepsilon < \|R(u, \lambda)\|^{-1}\)。

命题 20. --- 设 E 为复巴拿赫空间，u 为 E 上的闭部分算子。设 \(\varepsilon\in\mathbf{R}_{+}^{*}\)。对每个 \(\lambda\in\mathbf{C}\)，下列条件等价：

\begin{enumerate}
\def\labelenumi{(\roman{enumi})}
\item
  有 \(\lambda \in \mathrm{PSp}_{\varepsilon}(u)\)；
\item
  或者 \(\lambda \in \mathrm{Sp}(u)\)，或者存在 \(x \in \mathrm{dom}(u)\) 满足 \(\| x \| = 1\) 且 \(\| (\lambda 1_{\mathrm{E}} - u)(x) \| < \varepsilon\)；
\item
  存在 \(v \in \mathcal{L}(\mathrm{E})\) 满足 \(\| v \| < \varepsilon\) 且 \(\lambda \in \mathrm{Sp}(u + v)\)。
\end{enumerate}

可以假设 E 非零。条件 (ii) 是定义的重新表述，因而也就是条件 (i)。

设条件 (i) 成立，我们来证明 (iii)。若 \(\lambda\) 属于 \(u\) 的谱，则在 (iii) 中可取 \(v = 0\)。

于是设 \(\lambda \notin \mathrm{Sp}(u)\)。由 \(\mathrm{PSp}_{\varepsilon}(u)\) 的定义，存在 \(y \in \mathrm{E}\) 满足 \(\| y \| = 1\) 且 \(\| \mathrm{R}(u, \lambda)y \| > \varepsilon^{-1}\)。令 \(x = \mathrm{R}(u, \lambda)y\)。我们有 \(x \neq 0\)。由哈恩–巴拿赫定理（EVT，II，第 67 页，推论 2），线性形式 \(\ell\)（在 \(\mathbf{C}x\) 上、满足 \(\ell(x) = 1\)）容许连续延拓 \(\ell_1 \in \mathrm{E}'\) 满足 \(\| \ell_1 \| = \| \ell \|\)；于是 \(\| \ell_1 \| = \| x \|^{-1} < \varepsilon\)。对每个 \(e \in \mathrm{E}\)，令 \(v(e) = \ell_1(e)y\)。我们有 \(v \in \mathcal{L}(\mathrm{E})\) 且 \(v(x) = y\)。我们得到 \((u + v)x = u(x) + y = \lambda x\)，故 \(\lambda \in \mathrm{Sp}(u + v)\)。又因 \(\| v \| = \| \ell_1 \| < \varepsilon\)，条件 (iii) 得到满足。

反之，设 \(v \in \mathcal{L}(\mathrm{E})\) 满足 \(\|v\| < \varepsilon\) 且 \(\lambda \in \mathrm{Sp}(u + v)\)。于是部分算子 \(\lambda1_{\mathrm{E}} - (u + v)\) 或者不是单射的，或者没有连续的逆；把 IV，第 243 页的引理 4 应用于单射部分算子 \(\lambda1_{E} - u\) 与 \(-v\)，便得 \(\|v\| \geqslant \|R(u, \lambda)\|^{-1}\)。由此推出条件 (iii) 蕴含 (i)。

推论。— 设 E 为复巴拿赫空间，u 为 E 上的闭部分算子。设 \(\varepsilon > 0\)。

\begin{enumerate}
\def\labelenumi{\alph{enumi})}
\item
  对每个 \(v \in \mathcal{L}(\mathrm{E})\)，我们有 \(\mathrm{PSp}_{\varepsilon}(u) \subset \mathrm{PSp}_{\varepsilon + \| v\|}(u + v)\)；
\item
  设 \(\delta > 0\)，\(D_{\delta}\) 为以 0 为心、以 \(\delta\) 为半径的 \(\mathbf{C}\) 中的开圆盘。我们有 \(\mathrm{PSp}_{\varepsilon}(u) + D_{\delta} \subset \mathrm{PSp}_{\varepsilon + \delta}(u)\)。
\end{enumerate}

设 \(\lambda \in \mathrm{PSp}_{\varepsilon}(u)\)。存在自同态 \(w \in \mathcal{L}(\mathrm{E})\) 满足 \(\|w\| < \varepsilon\) 且 \(\lambda \in \mathrm{Sp}(u+w)\)（命题 20）。由于 \(u+w = (u+v)+(w-v)\) 且 \(\|w-v\| < \varepsilon + \|v\|\)，我们有 \(\lambda \in \mathrm{PSp}_{\varepsilon+ \|v\|}(u+v)\)（同上所引）。

设 \(\mu \in \mathrm{D}_{\delta}\)；我们有 \(\lambda + \mu \in \mathrm{Sp}(u + (w + \mu 1_{\mathrm{E}}))\) 且 \(\| w + \mu 1_{\mathrm{E}} \| < \varepsilon + \delta\)，故 \(\lambda + \mu \in \mathrm{PSp}_{\varepsilon + \delta}(u)\)（同上所引）。

命题 21. --- 设 E 为复巴拿赫空间，u 为 E 上的闭部分算子。设 \(\varepsilon\in\mathbf{R}_{+}^{*}\)。\(\mathrm{PSp}_{\varepsilon}(u)\) 的每个有界连通分支都与 \(u\) 的谱相交。

设 U 是 \(\mathrm{PSp}_{\varepsilon}(u)\) 的一个不与 \(\mathrm{Sp}(u)\) 相交的连通分支。集合 U 在 \(\mathrm{PSp}_{\varepsilon}(u)\) 中既开又闭，因此它在 C 中的闭包 \(\overline{U}\) 满足 \(\overline{U} \cap \mathrm{PSp}_{\varepsilon}(u) = U\)。

由于 \(\mathrm{Sp}(u)\) 包含于 \(\mathrm{PSp}_{\varepsilon}(u)\) 而 U 与 \(\mathrm{Sp}(u)\) 不相交，这一等式表明集合 \(\overline{U}\) 与 \(\mathrm{Sp}(u)\) 不相交，从而包含于 u 的预解集。此外，它还蕴含集合 \(\overline{U}-U\) 与 \(\mathrm{PSp}_{\varepsilon}(u)\) 不相交。

于是我们有 \(\|\mathrm{R}(u,\lambda)\| \leqslant \varepsilon^{-1}\)，对一切 \(\lambda\)（属于 \(\overline{U} - U\)）；而 \(\|\mathrm{R}(u,\lambda)\| > \varepsilon^{-1}\)，对 \(\lambda \in U\)。若集合 U 有界，则其闭包 \(\overline{U}\) 是紧的，且存在 \(\lambda_{0} \in \overline{U}\) 使得 \(\|\mathrm{R}(u,\lambda)\| \leqslant \|\mathrm{R}(u,\lambda_{0})\|\) 对 \(\lambda \in \overline{U}\) 成立。前述结果表明 \(\lambda_{0} \in U\)，这与极大值原理（VAR，R1，第 29 页，3.3.7）矛盾，因为 u 的预解式在 u 的预解集中是全纯的（IV，第 246 页，命题 14）。

\subsection{乘法算子}

设 X 是局部紧空间，\(\mu\) 是 X 上的正测度。我们考察 \(\mathrm{L}^{2}(\mathrm{X},\mu)\) 上的乘法算子；它们是稠定的闭算子（IV，第 232 页，命题 5）。对 X 上任何 \(\mu\)-可测函数 g，我们都用 \(m_{g}\) 表示 \(\mathrm{L}^{2}(\mathrm{X},\mu)\) 上以 g 为乘子的乘法部分算子。

命题 22. — 设 \(g\) 是 X 上的 \(\mu\)-可测函数。

\begin{enumerate}
\def\labelenumi{\alph{enumi})}
\item
  \(m_{g}\) 的谱是相对于 \(\mu\) 的 \(g\) 的本质像 S；
\item
  设 \(\lambda \in \mathbf{C} - \mathrm{Sp}(m_g)\)。预解式 \(\mathrm{R}(m_g, \lambda)\) 是乘法算子 \(m_h\)，其中 \(h\) 是 X 上如下定义的函数：\(h(x) = 0\)（若 \(g(x) = \lambda\)），否则 \(h(x) = (\lambda - g(x))^{-1}\)。
\end{enumerate}

我们来证明 \(\mathbf{C} - \mathbf{S}\) 包含于 \(m_g\) 的预解集。设 \(\lambda \in \mathbf{C} - \mathbf{S}\)。存在 \(\lambda\) 的开邻域 U，使 X 的子集 \(\mathrm{Y} = g^{-1}(\mathrm{U})\) 是局部 \(\mu\)-零测的。X 上如下定义的函数 \(k\)：\(k(x) = (\lambda - g(x))^{-1}\)（若 \(x \notin \mathrm{Y}\)），而 \(k(x) = 0\)（若 \(x \in \mathrm{Y}\)），它属于 \(\mathcal{L}^{\infty}(\mathrm{X}, \mu)\)（IV，第 184 页，引理 5）；因此以 \(k\) 为乘子的乘子算子是 \(\mathrm{L}^{2}(\mathrm{X}, \mu)\) 的自同态。

由于 \(|gk| \leqslant 1 + |\lambda k|\)，我们有

\[
| g k f | \leqslant | f | + | \lambda k f |
\]

对 \(f \in \mathcal{L}^2(\mathrm{X},\mu)\) 成立，这蕴含 \(m_k\) 的像包含于 \(m_g\) 的定义域。反之，设 \(f \in \mathcal{L}^2(\mathrm{X},\mu)\)，其类 \(\widetilde{f}\) 属于 \(m_g\) 的定义域。在局部 \(\mu\)-零测集 Y 之外，我们有 \(f(x) = k(x)(\lambda - g(x))f(x)\)，故 \(\widetilde{f}\) 在 \(m_k\) 的像中。同一公式证明 \(\lambda\) 属于 \(m_g\) 的预解集且 \(\mathrm{R}(m_g,\lambda) = m_k\)。由于集合 Y 是局部 \(\mu\)-零测的，乘法算子 \(m_k\) 与算子 \(m_h\)（如断言 \(b\)）中所描述的）一致。

反之，我们来证明 \(\mathbf{C} - \mathrm{Sp}(m_g)\) 包含于 \(\mathbf{C} - \mathrm{S}\)。设 \(\lambda \in \mathbf{C} - \mathrm{Sp}(m_g)\)。取实数 \(\mathrm{M} > \| \mathrm{R}(m_g, \lambda) \|\)。用 Y 表示由这样的 \(x \in \mathrm{X}\) 组成的集合：它们满足 \(|\lambda - g(x)| < \mathrm{M}^{-1}\)。我们来证明 Y 是局部 \(\mu\)-零测的，这将蕴含 \(\lambda\) 不属于 S，并完成证明。

设 \(K\) 是 \(X\) 的紧子集。设 \(\varphi\) 是 \(Y \cap K\) 的特征函数；它是 \(\mathcal{L}^{2}(X,\mu)\) 中的元素，我们把它在 \(\mathrm{L}^{2}(X,\mu)\) 中的类记作 \(\widetilde{\varphi}\)。设 \(\psi\) 是 \(\mathcal{L}^{2}(X,\mu)\) 中的函数，它在 \(\mathrm{L}^{2}(X,\mu)\) 中的类是 \(\mathrm{R}(m_g,\lambda)(\widetilde{\varphi})\)。我们有 \(\mathrm{R}(m_g,\lambda)(\widetilde{\varphi}) \in \operatorname{dom}(m_g)\) 且 \((\lambda - m_g)(\mathrm{R}(m_g,\lambda)(\widetilde{\varphi})) = \widetilde{\varphi}\)，故 \((\lambda - g(x))\psi(x) = 1\) 对 \(\mu\)-几乎所有的 \(x \in Y \cap K\) 成立。由此可得

\[
\| \mathrm{R} (m _ {g}, \lambda) \| ^ {2} \| \widetilde {\varphi} \| ^ {2} \geqslant \| \mathrm{R} (m _ {g}, \lambda) (\widetilde {\varphi}) \| ^ {2} \geqslant \mathrm{M} ^ {2} \mu (\mathrm{Y} \cap \mathrm{K}) = \mathrm{M} ^ {2} \| \widetilde {\varphi} \| ^ {2}.
\]

鉴于 M 的取法，这表明 \(\varphi\) \(\mu\)-几乎处处为零。于是集合 Y \(\cap\) K 是 \(\mu\)-零测的。因此集合 Y 是局部 \(\mu\)-零测的（INT，IV，第 172 页，§ 5，第 2 目，命题 5）。

命题 23. --- 设 g 是 X 上的 μ-可测函数。乘法算子 \(m_{g}\) 的伴随算子是 \(m_{\overline{g}}\)。

对每个整数 \(n \geqslant 1\)，设 \(\varphi_{n}\) 是由这样的元素 \(x \in X\) 组成的集合的特征函数：它们满足 \(|g(x)| \leqslant n\)，并设 \(\widetilde{\varphi}_{n}\) 是它在 \(\mathrm{L}^{2}(\mathrm{X}, \mu)\) 中的类。设 \(f \in \mathcal{L}^{2}(\mathrm{X}, \mu)\)，其类 \(\widetilde{f}\) 属于 \(\operatorname{dom}(m_{g}^{*})\)，又设 \(\psi\) 是以 \(m_{g}^{*}(\widetilde{f})\) 为类的函数。

对每个 \(h \in \mathcal{L}^2(\mathrm{X}, \mu)\)，若其类 \(\widetilde{h}\) 属于 \(\operatorname{dom}(m_g)\)，则也有 \(\widetilde{h}\widetilde{\varphi}_n \in \operatorname{dom}(m_g)\)，从而 \(\langle \widetilde{f} | m_g(\widetilde{h}\widetilde{\varphi}_n) \rangle = \langle m_g^*(\widetilde{f}) | \widetilde{h}\widetilde{\varphi}_n \rangle\)。由此得到等式

\[
\int_ {\mathrm{X}} \overline {{(f \overline {{g}} - \psi)}} \varphi_ {n} h d \mu = 0.
\]

由于 \(m_{g}\) 的定义域在 \(\mathrm{L}^{2}(\mathrm{X},\mu)\) 中稠密，我们推出 \((f\overline{g}-\psi)\varphi_{n}\) \(\mu\)-几乎处处为零。由于 n 是任意的，这就意味着 \(m_{g}^{*}(\widetilde{f})\) 在 \(\mathrm{L}^{2}(\mathrm{X},\mu)\) 中等于 \(f\overline{g}\) 的类。特别地，由 \(m_{g}^{*}(\widetilde{f}) \in \mathrm{L}^{2}(\mathrm{X},\mu)\)，我们断定 f 属于 \(m_{\overline{g}}\) 的定义域，且 \(m_{g}^{*}(\widetilde{f}) = m_{\overline{g}}(\widetilde{f})\)。

因此 \(m_{g}\) 的伴随算子是 \(m_{\overline{g}}\) 的扩张。此外，我们有

\[
\langle f \mid m _ {g} (h) \rangle = \int_ {\mathrm{X}} \overline {{f}} \cdot (g h) d \mu = \int_ {\mathrm{X}} f \overline {{g}} h d \mu
\]

对一切 \(f \in \mathrm{L}^{2}(\mathrm{X}, \mu)\) 与 \(h \in \mathrm{dom}(m_{g})\) 成立，这就证明线性形式 \(h \mapsto \langle f | m_{g}(h) \rangle\) 当 \(m_{\overline{g}}(f) g\) 属于 \(\mathrm{L}^{2}(\mathrm{X}, \mu)\) 时是连续的。于是，\(m_{\overline{g}}\) 的定义域包含于 \(m_{g}^{*}\) 的定义域，证明至此完成。

推论。--- 设 g 是 X 上的 μ-可测函数。乘法算子 \(m_{g}\)（\(\mathrm{L}^{2}(\mathrm{X},\mu)\) 上的）是自伴的（相应地，正的），当且仅当函数 g 局部 μ-几乎处处取实值（相应地，局部 μ-几乎处处为正）。

第一个断言由该命题得出。若 \(g\) 局部 \(\mu\)-几乎处处为正，则有 \(\langle f|m_{g}(f)\rangle=\int_{\mathrm{X}}g|f|^{2}d\mu\geqslant0\)，对一切 \(f\in\mathrm{L}^{2}(\mathrm{X},\mu)\) 成立，故部分算子 \(m_{g}\) 是正的。

反之，若 \(m_{g}\) 是正的，则其谱包含于 \(\mathbf{R}_{+}\)（IV，第 248 页，命题 17）；由于该谱是 g 的 \(\mu\)-本质像（命题 22，a）），这就意味着 g 局部 \(\mu\)-几乎处处为正。

引理 6. --- 设 \(g_{1}\) 与 \(g_{2}\) 是 X 上的 \(\mu\)-可测函数。

\begin{enumerate}
\def\labelenumi{\alph{enumi})}
\item
  部分算子 \(m_{g_{1}} + m_{g_{2}}\) 是可闭的，其闭包是乘法算子 \(m_{g_{1}+g_{2}}\)；
\item
  有 \(m_{g_{1}} \circ m_{g_{2}} \subset m_{g_{1}g_{2}}\)；
\item
  设 \(g_{1}\) 有界。这时我们有 \(m_{g_{2}} \circ m_{g_{1}} = m_{g_{1}g_{2}}\)。此外，\(m_{g_{2}}\) 的定义域包含于 \(m_{g_{1}g_{2}}\) 的定义域，且 \(m_{g_{1}} \circ m_{g_{2}}\) 是 \(m_{g_{1}g_{2}}\) 到 \(\text{dom}(m_{g_{2}})\) 上的限制。
\end{enumerate}

\(m_{g_{1}+g_{2}}\) 是 \(m_{g_{1}} + m_{g_{2}}\) 的扩张，这是初等的；于是后一算子是可闭的，且 \(\overline{m_{g_{1}} + m_{g_{2}}} \subset m_{g_{1}+g_{2}}\)。

设 \(f \in \mathcal{L}^2(\mathrm{X},\mu)\)，使函数 \(h = (g_1 + g_2)f\) 属于 \(\mathcal{L}^2(\mathrm{X},\mu)\)。对每个整数 \(n \geqslant 1\)，用 \(\mathrm{X}_n\) 表示由这样的 \(x \in \mathrm{X}\) 组成的集合：它们满足 \(|g_1(x)| + |g_2(x)| \leqslant n\)，并设 \(\varphi_n\) 是 \(\mathrm{X}_n\) 的特征函数。我们有 \(\varphi_n f \in \mathrm{dom}(m_{g_1} + m_{g_2})\)。由于 \((\varphi_n f)(x)\) 收敛于 \(f(x)\)，对一切 \(x \in \mathrm{X}\)，且对每个 \(n \in \mathbf{N}\) 有 \(|\varphi_n f| \leqslant |f|\) 与 \(|(g_1 + g_2)\varphi_n f| \leqslant |(g_1 + g_2)f| = |h|\)（由假设 \(h \in \mathcal{L}^2(\mathrm{X},\mu)\)），勒贝格定理（INT，IV，第 137 页，§ 3，第 7 目，定理 6）蕴含：\(\mathrm{L}^{2}(\mathrm{X},\mu)\times\mathrm{L}^{2}(\mathrm{X},\mu)\) 中 \((\varphi_{n}f,(g_{1}+g_{2})\varphi_{n}f)\) 的类对（它们属于 \(m_{g_{1}}+m_{g_{2}}\) 的图像）所成的序列，收敛于 \((f,h)\) 在 \(\mathrm{L}^{2}(\mathrm{X},\mu)\) 中的类对。因此 \(m_{g_{1}}+m_{g_{2}}\) 的闭包确实等于 \(m_{g_{1}+g_{2}}\)。

\(m_{g_{1}} \circ m_{g_{2}} \subset m_{g_{1}g_{2}}\) 是初等的。设 \(g_{1}\) 有界，于是 \(m_{g_{1}}\) 是 \(\mathrm{L}^{2}(\mathrm{X},\mu)\) 上的连续线性映射。此时 \(m_{g_{2}} \circ m_{g_{1}}\) 的定义域是这样的函数 \(f \in \mathcal{L}^{2}(\mathrm{X},\mu)\) 的等价类所成的集合：这些函数使 \(g_{2}(g_{1}f) \in \mathcal{L}^{2}(\mathrm{X},\mu)\)，亦即 \(m_{g_{1}g_{2}}\) 的定义域。因此我们有 \(m_{g_{2}} \circ m_{g_{1}} = m_{g_{1}g_{2}}\)。

同样初等的是：\(\operatorname{dom}(m_{g_{1}} \circ m_{g_{2}}) = \operatorname{dom}(m_{g_{2}})\) 包含于 \(\operatorname{dom}(m_{g_{1}g_{2}})\)，且 \(m_{g_{1}g_{2}}\) 到这个空间的限制等于 \(m_{g_{1}} \circ m_{g_{2}}\)。

应当警惕，不要以为一般地有 \(m_{g_{1}} \circ m_{g_{2}} = m_{g_{1}g_{2}}\)（IV，第 347 页，习题 10）。尽管如此，我们有如下的部分结果：

命题 24. --- 设 g 是 X 上的 μ-可测函数。

\begin{enumerate}
\def\labelenumi{\alph{enumi})}
\item
  有 \(m_{\overline{g}} \circ m_g = m_{|g|^2}\)；
\item
  对一切整数 \(k, \ell \in \mathbf{N}\)，有 \(m_{g^k\overline{g}^\ell} = m_g^k m_{\overline{g}}^\ell\)。
\end{enumerate}

由引理 6 有 \(m_{\overline{g}} \circ m_g \subset m_{|g|^2}\)。反之，由不等式 \(|g| \leqslant 1 + |g|^2\) 推出 \(\text{dom}(m_{|g|^2}) = \text{dom}(m_{\overline{g}} \circ m_g)\)，由此得第一个断言。

设 \(k, \ell \in \mathbf{N}\)。我们有 \(m_g^k m_{\overline{g}}^\ell \subset m_{g^k \overline{g}^\ell}\)（同上所引）。\(m_g^k m_{\overline{g}}^\ell\) 的定义域是在 \(\mathrm{L}^2(\mathrm{X}, \mu)\) 中这样的函数 \(h \in \mathcal{L}^2(\mathrm{X}, \mu)\) 的类所成的集合：\(|g|^j h\) 属于 \(\mathcal{L}^2(\mathrm{X}, \mu)\)，对每个整数 \(j\)（满足 \(0 \leqslant j \leqslant k + \ell\)）。不等式

\[
\left| g ^ {j} h \right| \leqslant | h | + \left| g ^ {k + \ell} h \right|,
\]

对 \(0 \leqslant j \leqslant k + \ell\) 成立，由它可以看出 \(\mathrm{dom}(m_g^k m_{\overline{g}}^\ell)\) 等于 \(\mathrm{dom}(m_{g^k\overline{g}^\ell})\)，由此得断言 \(b\)）。

\subsection{对称算子的自伴扩张}

在本节中，我们打算对复希尔伯特空间上对称算子的自伴扩张进行分类，特别是找出保证这类算子的自伴扩张存在的条件。

设 E 是复希尔伯特空间。对 E 上每个部分算子 u，本节中我们将把 \(u + i\) 与 \(u - i\) 分别写作 \(u + i1_{E}\) 与 \(u - i1_{E}\) 的简写。

设 u 是 E 上稠定的闭算子。我们回顾（参见 IV，第 230 页定义 4 及 IV，第 230 页命题 4）：我们用 \(E_{u}\) 表示赋有埃尔米特形式的希尔伯特空间 dom(u)

\[
(x \mid y) _ {u} = \langle x \mid y \rangle + \langle u (x) \mid u (y) \rangle .
\]

我们用 \(\|x\|_{u}\) 表示希尔伯特空间 \(E_{u}\) 中元素 x 的范数。dom(u) 的每个在 E 中闭的子空间都在 \(E_{u}\) 中闭（IV，第 230 页，注 1）。

定义 10. --- 设 u 是复希尔伯特空间 E 上的闭对称算子。设 \(\mathrm{E}_{+} = \mathrm{Ker}(u^{*} - i)\)，\(\mathrm{E}_{-} = \mathrm{Ker}(u^{*} + i)\)。子空间偶 \((\mathrm{E}_{+}, \mathrm{E}_{-})\)（\(\mathrm{dom}(u^{*})\) 的）称为 u 的亏偶。子空间 \(E_{+}\) 与 \(E_{-}\) 是 E 的闭子空间（IV，第 229 页，命题 3）。这些子空间的希尔伯特维数所成的偶 \((\dim(\mathrm{E}_{+}), \dim(\mathrm{E}_{-}))\) 称为 u 的亏指标。

命题 25. --- 设 u 是复希尔伯特空间 E 上的闭对称算子。我们有

\[
\begin{array}{l} \operatorname{Ker} (u ^ {*} - i) ^ {\circ} = \operatorname{Im} (u + i), \\ \operatorname{Ker} (u ^ {*} + i) ^ {\circ} = \operatorname{Im} (u - i). \end{array}
\]

由 IV，第 236 页命题 7，c)，只需证明 \(u + i\)（相应地，\(u - i\)）的像在 E 中闭。由于 \(u\) 对称，我们有 \(\langle (u + i)(x) \mid (u + i)(y) \rangle = (x \mid y)_u\)，对一切 \(x\) 与 \(y\)（\(\text{dom}(u)\) 中的）。于是映射 \(x \mapsto u(x) + ix\)（从 \(E_u\) 到 E 中的）是等距的。它的像在 E 中闭（I，第 107 页，引理 8）。由于这也是 \(u + i\) 的像，故后者在 E 中闭。同理，\(u - i\) 的像在 E 中闭。

引理 7. --- 设 u 是复希尔伯特空间 E 上的闭对称算子，v 是 u 的闭对称扩张。u 的定义域是希尔伯特空间 \(E_{v}\) 的闭子空间。

由于 v 是 u 的扩张，我们有 \((x \mid y)_{v} = (x \mid y)_{u}\)，对 \(\text{dom}(u)\) 中的 x 与 y。于是 \(E_{u}\) 到 \(E_{v}\) 中的典范单射是等距的，结论由此得出（I，第 107 页，引理 8）。

命题 26. --- 设 u 是复希尔伯特空间 E 上的闭对称算子。设 \((\mathrm{E}_{+}, \mathrm{E}_{-})\) 是 u 的亏偶。空间 \(E_{+}\)、\(E_{-}\) 与 dom(u) 是 \(E_{u^{*}}\) 中两两正交的闭子空间，其希尔伯特直和等于 \(E_{u^{*}}\)。

由于 u 对称，我们有 \(u \subset u^{*}\)，故空间 \(\text{dom}(u)\) 在 \(E_{u^{*}}\) 中闭（引理 7）。子空间 \(E_{+}\) 与 \(E_{-}\) 在 E 中闭且包含于 \(\text{dom}(u^{*})\)，故在 \(E_{u^{*}}\) 中闭。

设 \(x \in \mathrm{E}_{+}\)。我们有 \(u^{*}(x) = ix\)，故 \(\langle x | u(y) \rangle = -i\langle x | y \rangle\) 对一切 \(y \in \operatorname{dom}(u)\) 成立。于是，对每个 \(y \in \operatorname{dom}(u)\)，有

\[
(x \mid y) _ {u ^ {*}} = \langle x \mid y \rangle + \langle u ^ {*} (x) \mid u ^ {*} (y) \rangle = i (\langle x \mid u (y) \rangle - \langle x \mid u ^ {*} (y) \rangle) = 0
\]

因为 u 对称。因此空间 \(E_{+}\) 与 \(\operatorname{dom}(u)\) 在 \(E_{u^{*}}\) 中正交。同理，\(E_{-}\) 与 \(\operatorname{dom}(u)\) 在 \(E_{u^{*}}\) 中正交。

设 \(x \in \mathrm{E}_{+}\)，\(y \in \mathrm{E}_{-}\)。则

\[
(x \mid y) _ {u ^ {*}} = \langle x \mid y \rangle + \langle u ^ {*} (x) \mid u ^ {*} (y) \rangle = \langle x \mid y \rangle - \langle x \mid y \rangle = 0,
\]

故 \(\mathrm{E}_{+}\) 与 \(\mathrm{E}_{-}\) 在 \(\mathrm{E}_{u^{*}}\) 中正交。

设 \(x \in \mathrm{E}_{u^*}\) 与闭子空间 \(\operatorname{dom}(u) \oplus \mathrm{E}_+ \oplus \mathrm{E}_-\) 正交。对每个 \(y \in \operatorname{dom}(u)\)，有 \(u^*(y) = u(y)\)（因为 \(u \subset u^*\)），于是

\[
0 = (x \mid y) _ {u ^ {*}} = \langle x \mid y \rangle + \langle u ^ {*} (x) \mid u ^ {*} (y) \rangle = \langle x \mid y \rangle + \langle u ^ {*} (x) \mid u (y) \rangle .
\]

这蕴含 \(z = u^{*}(x)\) 属于 \(u^{*}\) 的定义域且满足 \(u^{*}(z) = -x\)。令 \(x_{-} = z - ix\)。它是 \(\text{dom}(u^{*})\) 中满足 \(u^{*}(x_{-}) = -x - iz = -ix_{-}\) 的元素。故 \(x_{-}\) 属于 \(E_{-}\)。但是，对每个 \(w \in E_{-}\)，有

\[
\begin{array}{c} - i \langle x _ {-} | w \rangle = - i \langle z | w \rangle + \langle x | w \rangle \\ = \langle u ^ {*} (x) | u ^ {*} (w) \rangle + \langle x | w \rangle = (x | w) _ {u ^ {*}} = 0. \end{array}
\]

取 \(w = x_{-}\)，便得 \(x_{-} = 0\)，即 \(z = ix\)；由于 \(z = u^{*}(x)\)，我们有 \(x \in \mathrm{E}_{+}\)，从而 \(x = 0\)，因为 \(x\) 与 \(\mathrm{E}_{+}\) 在 \(\mathrm{E}_{u^{*}}\) 中正交。这就完成了直和 \(\operatorname{dom}(u) \oplus \mathrm{E}_{+} \oplus \mathrm{E}_{-}\) 等于 \(\mathrm{E}_{u^{*}}\) 的证明。

推论。--- 设 u 是 E 上的闭对称算子。则 u 自伴，当且仅当 u 的亏指标等于 (0,0)。

事实上，对称算子 \(u\) 自伴当且仅当 \(\text{dom}(u^*) = \text{dom}(u)\)；而命题 26 表明 \(\text{dom}(u^*)\) 是 \(\text{dom}(u)\) 与 \(u\) 的亏子空间的希尔伯特直和。

\pandocbounded{\includegraphics[keepaspectratio]{images/part002_f39b57bf.jpg}}

设 \(u\) 是闭对称部分算子。\(u\) 的谱包含于 \(\mathbf{R}\)，当且仅当 \(u\) 自伴。事实上，我们知道 \(\mathrm{Sp}(u) \subset \mathbf{R}\) 当 \(u\) 自伴时成立（IV，第 248 页，命题 17）；若 \(u\) 不自伴，则子空间 \(\mathrm{Ker}(u^{*} + i)\) 与 \(\mathrm{Ker}(u^{*} - i)\) 之一非零，于是 \(u + i\) 或 \(u - i\) 的像是 E 的真子空间（IV，第 256 页，命题 25），从而 \(i\) 或 \(-i\) 属于 \(\mathrm{Sp}(u)\)。

定义 11. --- 设 u 是复希尔伯特空间 E 上的闭对称算子。u 的边界条件是从 \(\operatorname{Ker}(u^{*} - i)\) 到 \(\operatorname{Ker}(u^{*} + i)\) 中的一个部分等距线性映射（EVT，V，第 41 页，定义 3）。

设 \(u\) 是复希尔伯特空间 E 上的闭对称算子，\(b\) 是 \(u\) 的一个边界条件。设 \(I = \text{Ker}(b)^{\circ}\) 为 \(b\) 的初始子空间。我们用 \(u_b\) 表示 \(u^*\) 到 \(\text{dom}(u^*)\) 的如下子空间上的缩减：该子空间是 \(\text{dom}(u)\) 与一个图像的直和，该图像位于 \(E_+ \oplus E_-\) 中，是 \(b\) 在 I 上的限制的图像。由于 \(E_{u^*} = \text{dom}(u) \oplus E_+ \oplus E_-\)（命题 26），部分算子 \(u_b\) 是完全确定的。

换句话说，\(u_{b}\) 的定义域是这样的元素 \(x \in E\) 所成的空间：它们形如 \(x = x_{0} + y + b(y)\)，其中 \(x_{0} \in \text{dom}(u)\)，y 属于 b 的初始子空间。此时我们有 \(u_{b}(x_{0} + y + b(y)) = u(x_{0}) + iy - ib(y)\)。

定理 1（冯·诺伊曼）。--- 设 u 是复希尔伯特空间 E 上的闭对称算子，\((\mathrm{E}_{+}, \mathrm{E}_{-})\) 是它的亏偶。

\begin{enumerate}
\def\labelenumi{\alph{enumi})}
\item
  对 u 的每个边界条件 b，部分算子 \(u_{b}\) 是 u 的闭对称扩张；
\item
  映射 \(b \mapsto u_{b}\) 是从 u 的边界条件集合到 u 的闭对称扩张集合上的双射；
\item
  对 u 的每个边界条件 b，我们有
\end{enumerate}

\[
\begin{array}{l} \operatorname{Im} (u _ {b} + i) = \operatorname{Im} (u + i) \oplus \operatorname{Im} (b) \\ \operatorname{Im} (u _ {b} - i) = \operatorname{Im} (u - i) \oplus \operatorname{Ker} (b) ^ {\circ}; \end{array}
\]

\begin{enumerate}
\def\labelenumi{\alph{enumi})}
\setcounter{enumi}{3}
\tightlist
\item
  对 u 的每个边界条件 b，\(u_{b}\) 的亏偶是 \((\operatorname{Ker}(b), \operatorname{Ker}(b^{*}))\)。
\end{enumerate}

设 \(b\) 是 \(u\) 的一个边界条件，\(\mathrm{I} \subset \mathrm{E}_{+}\) 是它的初始子空间。\(u^{*}\) 到 \(\Gamma_{b|I} \subset \mathrm{E}_{+} \oplus \mathrm{E}_{-}\) 的限制是由 \(x + b(x) \mapsto ix - ib(x)\)（\(x \in I\)）定义的连续线性映射。\(u_{b}\) 的图像是 u 的图像与该线性映射的图像的直和；因此它是闭的。闭算子 \(u_{b}\) 是 u 的扩张且满足 \(u_{b} \subset u^{*}\)。

设 \(\gamma_{1}\) 与 \(\eta_{1}\) 是 I 的元素。考虑元素

\[
\gamma = \gamma_ {1} + b (\gamma_ {1}), \quad \eta = \eta_ {1} + b (\eta_ {1})
\]

它们是图像 \(\Gamma_{b|I}\) 中的元素。我们计算

\[
\begin{array}{c} \langle \gamma | u ^ {*} (\eta) \rangle = i \big (\langle b (\gamma_ {1}) | \eta_ {1} \rangle - \langle \gamma_ {1} | b (\eta_ {1}) \rangle \big) + i \big (\langle \gamma_ {1} | \eta_ {1} \rangle - \langle b (\gamma_ {1}) | b (\eta_ {1}) \rangle \big) \\ = i \big (\langle b (\gamma_ {1}) | \eta_ {1} \rangle - \langle \gamma_ {1} | b (\eta_ {1}) \rangle \big) \end{array}
\]

因为 b 在 I 上等距。由此我们推出

\[
\langle u ^ {*} (\gamma) \mid \eta \rangle = \langle \gamma \mid u ^ {*} (\eta) \rangle\tag{12}
\]

对一切 \(\gamma\) 与 \(\eta\)（在 \(\Gamma_{b|I}\) 中）成立。

设 x 与 y 是 \(\operatorname{dom}(u)\) 的元素，\(\gamma\) 与 \(\eta\) 是 \(\Gamma_{b|I}\) 的元素。我们得到

\[
\begin{array}{c} \langle x + \gamma | u _ {b} (y + \eta) \rangle = \langle x | u (y) \rangle + \langle x | u ^ {*} (\eta) \rangle + \langle \gamma | u (y) \rangle + \langle \gamma | u ^ {*} (\eta) \rangle \\ = \langle u (x) | y \rangle + \langle u (x) | \eta \rangle + \langle u ^ {*} (\gamma) | y \rangle + \langle u ^ {*} (\gamma) | \eta \rangle \\ = \langle u _ {b} (x + \gamma) | y + \eta \rangle , \end{array}
\]

其中用到了 u 对称这一事实、\(\gamma\) 与 \(\eta\) 属于 \(u^{*}\) 的定义域这一事实，以及公式 (12)。于是算子 \(u_{b}\) 是 u 的闭对称扩张。断言 a) 得证。

映射 \(b \mapsto u_b\) 是单射。事实上，从 \(E_+\) 到 \(E_-\) 中的部分等距映射由其初始子空间及其在该子空间上的限制唯一确定。而 \(u_b\) 的定义域确定了 \(b\) 的初始子空间 I 以及 \(b\) 在 I 上限制的图像。

我们来证明映射 \(b \mapsto u_b\) 是满射。设 \(w\) 是 \(u\) 的闭对称扩张。我们有 \(w \subset w^* \subset u^*\)。于是 \(w\) 的定义域是 \(\text{dom}(u^*)\) 中包含 \(\text{dom}(u)\) 的子空间，且 \(w\) 是 \(u^*\) 到该子空间上的缩减。设 G 是 \(w\) 的定义域与 \(E_+ \oplus E_-\) 的交。它是 \(E_{u^*}\) 的闭子空间（引理 7），且由命题 26 有 \(\text{dom}(w) = \text{dom}(u) \oplus G\)。

设 \(x \in \mathrm{E}_{+}\)，\(y \in \mathrm{E}_{-}\) 满足 \(x + y \in \mathrm{G}\)。由于 \(w\) 对称且是 \(u^{*}\) 到 \(\operatorname{dom}(w)\) 上的缩减，我们有

\[
\langle x + y \mid u ^ {*} (x + y) \rangle = \langle u ^ {*} (x + y) \mid x + y \rangle .
\]

此等式等价于

\[
\begin{array}{c} \langle x | i x \rangle + \langle x | - i y \rangle + \langle y | i x \rangle + \langle y | - i y \rangle \\ = \langle i x | x \rangle + \langle i x | y \rangle + \langle - i y | x \rangle + \langle - i y | y \rangle , \end{array}
\]

亦即等价于 \(\|x\|^{2} = \|y\|^{2}\)。特别地，典范投影 \(p_{+}: G \to E_{+}\) 是单射，且若用 I 表示它的像，则空间 G 是等距映射 \(b_{0}\)（从 I 到 \(E_{-}\) 中的）的图像。对 I 中的 x 与 \(E_{-}\) 中满足 \(x + y \in G\) 的 y，我们有

\[
\begin{array}{r} \| x + y \| _ {u ^ {*}} ^ {2} = \| x \| _ {u ^ {*}} ^ {2} + \| y \| _ {u ^ {*}} ^ {2} = 2 (\| x \| ^ {2} + \| y \| ^ {2}) \\ = 4 \| x \| ^ {2} = 4 \| p _ {+} (x + y) \| ^ {2}. \end{array}
\]

于是，把 G 看作希尔伯特空间 \(E_{u^{*}}\) 的闭子空间时，映射 \(x \mapsto \frac{1}{2}p_{+}(x)\) 是从 G 到 I 上的等距映射。特别地，I 在 \(E_{+}\) 中闭（I，第 107 页，引理 8）。于是存在 u 的唯一的边界条件 b，以 I 为初始空间且在 I 上与 \(b_{0}\) 相合。我们有 \(\text{dom}(w) = \text{dom}(u_{b})\)，从而 \(w = u_{b}\)。这就建立了断言 b)。

我们来证明断言 \(c\)）与 \(d\)）。设 \(b\) 是 \(u\) 的一个边界条件，I 是它的初始子空间。对 \(x_0 \in \mathrm{dom}(u)\) 与 \(y \in \mathrm{I}\)，有

\[
u _ {b} (x _ {0} + y + b (y)) - i (x _ {0} + y + b (y)) = (u - i) (x _ {0}) - 2 i b (y).
\]

我们有 \(\operatorname{Im}(u - i) = \operatorname{E}_{-}^{\circ}\)（命题 25）。由于 \(b\) 的像包含在 \(\operatorname{E}_{-}\) 中，该公式表明 \(\operatorname{Im}(u_b - i) = \operatorname{Im}(u - i) \oplus \operatorname{Im}(b)\)。由同上所引，我们因此有 \(\operatorname{Ker}(u_b^* + i) = \operatorname{E}_{-} \cap \operatorname{Im}(b)^{\circ} = \operatorname{Ker}(b^*)\)。

类似地，可以验证 \(\operatorname{Im}(u_b + i) = \operatorname{Im}(u + i) \oplus \operatorname{I}\)，从而（同上所引）有 \(\operatorname{Ker}(u_b^* - i) = \operatorname{I}^\circ \cap \operatorname{E}_+ = \operatorname{Ker}(b)\)。

推论 1. --- 设 u 是复希尔伯特空间 E 上的闭对称算子。映射 \(b \mapsto u_b\) 诱导出从 \(\text{Ker}(u^* - i)\) 到 \(\text{Ker}(u^* + i)\) 上的等距同构集合与 u 的自伴扩张集合之间的双射。特别地，存在 u 的自伴扩张，当且仅当 u 的亏指标的两个分量相等。

这由命题 26 的推论与该定理以及 EVT，V，第 25 页，推论 2 得出。

推论 2. --- 设 E 是实希尔伯特空间。设 u 是 E 上的闭对称算子。部分算子 \(u_{(\mathbf{C})}\)（\(\mathrm{E}_{(\mathbf{C})}\) 上的）对称且闭，并且至少有一个自伴扩张。

部分算子 \(u_{(\mathbf{C})}\) 闭且对称（IV，第 228 页命题 2 与 IV，第 239 页引理 3）。同构 \(j\)（从 \(\mathrm{E}_{(\mathbf{C})}\) 到 \(\mathrm{E}_{(\mathbf{C})}\) 上的 R-线性同构）满足 \(j(z \otimes x) = \overline{z} \otimes x\)，对一切 \(z \in \mathbf{C}\) 与每个 \(x \in \mathrm{E}\)；它诱导出从 \(\mathrm{Ker}(u_{(\mathbf{C})}^* - i)\) 到 \(\mathrm{Ker}(u_{(\mathbf{C})}^* + i)\) 上的希尔伯特空间同构，于是可以应用推论 1。

推论 3. --- 设 u 是 E 上的对称算子。下列条件等价：

\begin{enumerate}
\def\labelenumi{(\roman{enumi})}
\item
  算子 u 本质自伴；
\item
  空间 \(\operatorname{Ker}(u^{*} + i)\) 与 \(\operatorname{Ker}(u^{*} - i)\) 为零；
\item
  空间 \(\operatorname{Im}(u + i)\) 与 \(\operatorname{Im}(u - i)\) 在 E 中稠密；
\item
  对称部分算子 u 有唯一的自伴扩张。
\end{enumerate}

由 IV，第 236 页命题 7，c)，条件 (ii) 与 (iii) 等价。

部分算子 \(u\) 本质自伴，当且仅当 \(\overline{u}\) 自伴，亦即当且仅当 \(\text{Ker}(\overline{u}^* - i) = \text{Ker}(\overline{u}^* + i) = \{0\}\)（命题 26 的推论）。由于 \(\overline{u}^* = u^*\)（IV，第 237 页，命题 8），条件 (i) 因此等价于条件 (ii)。此外，前一推论表明 (ii) 蕴含 \(u\) 有唯一的自伴扩张，此即条件 (iv)。

最后，设 \(u\) 有唯一的自伴扩张 \(v\)。那么 \(\overline{u}\) 亦然，因为 \(\overline{u}\) 的每个自伴扩张都是 \(u\) 的自伴扩张，从而等于 \(v\)。由前一推论，空间 \(\mathrm{Ker}(\overline{u}^{*} + i)\) 与 \(\mathrm{Ker}(\overline{u}^{*} - i)\) 必为零，再利用 \(\overline{u}^{*} = u^{*}\) 即得条件 (ii)。

例。--- 设 U 是 \(\mathbf{R}^{n}\) 中赋有勒贝格测度 \(\mu\) 的开子集。设 \(\Delta\) 是数量微分算子

\[
\Delta = - \sum_ {i = 1} ^ {n} \partial_ {i} ^ {2}
\]

此算子定义在 U 上。用 \(u\) 表示实希尔伯特空间 \(\mathrm{L}_{\mathbf{R}}^{2}(\mathrm{U},\mu)\) 上以 \(\mathcal{D}_{\mathbf{R}}(\mathrm{U})\)（IV，第 201 页，注）为定义域、由 \(\varphi \mapsto -\sum_{i=1}^{n}\partial_{i}^{2}\varphi\) 定义的部分算子。我们有 \(u_{(\mathbf{C})} = \Delta_{-}\)，它是可闭的部分算子（IV，第 242 页，命题 13）且对称（因为 \(\Delta\) 形式上对称），故 \(u\) 可闭且对称（IV，第 228 页命题 2 与 IV，第 239 页引理 3）。于是由上述推论，存在 \(u_{(\mathbf{C})}\) 的自伴扩张。这就是 U 上的一个拉普拉斯算子（IV，第 243 页）。

\section{正规部分算子与谱定理}

\subsection{有界化}

设 E 是复希尔伯特空间。我们用 s 表示酉映射 \((x,y)\mapsto(-y,x)\)（\(E\oplus E\) 上的）。设 u 是 E 上具有稠密定义域的闭部分算子。部分算子 \(u^{*}u\) 自伴且正（IV，第 241 页，命题 12），故 \(-1\notin\mathrm{Sp}(u^{*}u)\)（IV，第 248 页，命题 17）。我们记 \(\mathrm{W}(u)=(1_{\mathrm{E}}+u^{*}u)^{-1}=-\mathrm{R}(u^{*}u,-1)\)；它是 E 的正的单射自同态。

用 \(p_{1}\) 与 \(p_{2}\) 表示 \(\Gamma_{u}\) 到 E 中的两个典范投影。它们是 \(\mathcal{L}(\Gamma_{u};\mathrm{E})\) 的元素，且有对应关系的等式 \(p_{2}=u\circ p_{1}\)。设 \((j,|p_{1}^{*}|)\) 是自同态 \(p_{1}^{*}\in\mathcal{L}(\mathrm{E};\Gamma_{u})\) 的极分解（I，第 140 页，定义 4），即 \(p_{1}^{*}=|p_{1}^{*}|\circ j\)。映射 j 是从 E 到 \(\Gamma_{u}\) 中的部分等距。

定义 1. --- E 的自同态 \(p_2 \circ j\) 称为 \(u\) 的有界化；我们用 \(b(u)\) 表示这个自同态。

我们有 \(\|b(u)\|\leqslant1\)。

命题 1. --- 我们有下列公式

\[
| p _ {1} ^ {*} | = \mathrm{W} (u) ^ {1 / 2},\tag{1}
\]

\[
b (u) = u \circ \mathrm{W} (u) ^ {1 / 2} = u \circ | p _ {1} ^ {*} |,\tag{2}
\]

\[
1 _ {\mathrm{E}} - b (u) ^ {*} b (u) = \mathrm{W} (u),\tag{3}
\]

\[
b (u) \mathrm{W} (u) = \mathrm{W} (u ^ {*}) b (u), \quad b (u) \mathrm{W} (u) ^ {1 / 2} = \mathrm{W} (u ^ {*}) ^ {1 / 2} b (u),\tag{4}
\]

设 \(x \in E\)。令 \(y = \mathrm{W}(u)(x) \in \mathrm{dom}(u^{*}u)\)。我们有 \(y \in \mathrm{dom}(u)\) 且 \(u(y) \in \mathrm{dom}(u^{*})\)。对形如 \((y_{1}, u(y_{1}))\) 的、\(\Gamma_{u}\) 中的每个元素（其中 \(y_{1} \in \mathrm{dom}(u)\)），计算得

\[
\begin{array}{c} \langle (y, u (y)) | (y _ {1}, u (y _ {1})) \rangle = \langle y | y _ {1} \rangle + \langle u (y) | u (y _ {1}) \rangle \\ = \langle y + u ^ {*} u (y) | y _ {1} \rangle = \langle x | p _ {1} (y _ {1}, u (y _ {1})) \rangle . \end{array}
\]

这就表明 \(p_{1}^{*}(x) = (y, u(y))\)，从而 \(p_{1} \circ p_{1}^{*}(x) = y = \mathrm{W}(u)(x)\)。于是有 \(p_{1} \circ p_{1}^{*} = \mathrm{W}(u)\)，故 \(|p_{1}^{*}| = \mathrm{W}(u)^{1/2}\)，即公式 (1)。特别地，由 I，第 140 页命题 11 的推论有 \(\operatorname{Im}(p_{1}) = \operatorname{Im}(|p_{1}^{*}|)\)，由此得出 \(\operatorname{dom}(u) = \operatorname{Im}(p_{1}) = \operatorname{Im}(\mathrm{W}(u)^{1/2})\)。

于是部分算子 \(u \circ \mathrm{W}(u)^{1/2}\) 以 E 为定义域。关系式 \(p_1 \circ j = (p_1^*)^* \circ j = |p_1^*| = \mathrm{W}(u)^{1/2}\)（I，第 140 页，命题 11，a)）进而蕴含

\[
b (u) = p _ {2} \circ j = u \circ p _ {1} \circ j = u \circ \mathrm{W} (u) ^ {1 / 2},
\]

此即公式 (2)。

由 \(\operatorname{Ker}(j) = \operatorname{Ker}(|p_1^*|)\)（I，第 139 页，命题 10，\(b\)）与公式 (1) 便推出部分等距 \(j: \mathrm{E} \to \Gamma_u\) 是单射，从而是等距的，于是 \(j^* \circ j = 1_{\mathrm{E}}\)。把 j 写成 \(j = (p_1 \circ j, b(u))\)，我们得到

\[
1 _ {\mathrm{E}} = j ^ {*} \circ j = (p _ {1} \circ j) ^ {*} (p _ {1} \circ j) + b (u) ^ {*} b (u) = \mathrm{W} (u) + b (u) ^ {*} b (u),
\]

由此得公式 (3)。

设 \(x \in \mathrm{dom}(u)\)，令 \(y = \mathrm{W}(u)(x)\)。我们有 \(y \in \mathrm{dom}(u)\)，且公式 \(y + u^{*}u(y) = x\) 蕴含 \(u^{*}u(y) \in \mathrm{dom}(u)\)。于是我们有

\[
\begin{array}{r l} & {\mathrm{W} (u ^ {*}) ^ {- 1} (u (y)) = (1 _ {\mathrm{E}} + u u ^ {*}) (u (y))} \\ & {\qquad = u (y) + u u ^ {*} u (y) = u (y + u ^ {*} u (y)) = u (x),} \end{array}
\]

这表明 \(u \circ \mathrm{W}(u)\) 到 u 的定义域上的缩减等于 \(\mathrm{W}(u^{*}) \circ u\)。由于 \(|p_{1}^{*}|\) 的像等于 \(p_{1}\) 的像，从而等于 u 的定义域（I，第 140 页命题 11 的推论），由此得

\[
u \circ \mathrm{W} (u) \circ | p _ {1} ^ {*} | = \mathrm{W} (u ^ {*}) \circ u \circ | p _ {1} ^ {*} | = \mathrm{W} (u ^ {*}) \circ b (u)
\]

（公式 (2)）。此外，\(|p_{1}^{*}| = \mathrm{W}(u)^{1/2}\) 与 \(\mathrm{W}(u)\) 交换，于是我们得到

\[
b (u) \circ \mathrm{W} (u) = u \circ | p _ {1} ^ {*} | \circ \mathrm{W} (u) = u \circ \mathrm{W} (u) \circ | p _ {1} ^ {*} | = \mathrm{W} (u ^ {*}) \circ b (u).
\]

此关系式蕴含 \(b(u) \circ f(\mathrm{W}(u)) = f(\mathrm{W}(u^{*})) \circ b(u)\)，对每个函数 \(f \in \mathcal{C}(\mathbf{R}_{+})\)（I，第 113 页，命题 11），特别地有 \(b(u) \circ \mathrm{W}(u)^{1/2} = \mathrm{W}(u^{*})^{1/2} \circ b(u)\)。这就证明了公式 (4) 并完结了证明。

推论。--- 我们有 \(b(u)^{*} = b(u^{*})\)。

用 \(q_{1}\) 与 \(q_{2}\) 表示投影 \(\Gamma_{u^{*}}\to\mathrm{E}\)，并设 \((k,|q_{1}^{*}|)\) 是 \(q_{1}^{*}\) 的极分解，于是 \(b(u^{*})=q_{2}\circ k\)。对 E 中所有 \(x\) 与 \(y\)，我们有 \(j(x)\in\Gamma_{u}\) 与 \(k(y)\in\Gamma_{u^{*}}\)，且由 IV，第 236 页命题 7，\(\Gamma_{u}\) 与 \(s(\Gamma_{u^{*}})\) 正交，由此得

\[
0 = \langle j (x) | s (k (y)) \rangle = \langle p _ {1} \circ j (x) | - q _ {2} \circ k (y) \rangle + \langle p _ {2} \circ j (x) | q _ {1} \circ k (y) \rangle .
\]

由于 \(p_1 \circ j = |p_1^*|\) 且 \(q_1 \circ k = |q_1^*|\)（I，第 140 页，命题 11，b)），我们得到

\[
\langle b (u) x \mid | q _ {1} ^ {*} | y \rangle = \langle | p _ {1} ^ {*} | x \mid b (u ^ {*}) y \rangle ,
\]

故 \(|q_1^*|b(u) = b(u^*)^* |p_1^*|\)。利用公式 (1) 与 (4)，我们由此得出 \(b(u)|p_1^*| = |q_1^*|b(u) = b(u^*)^* |p_1^*|\)，又因 \(|p_1^*|\) 的像是 \(u\) 的定义域，它在 E 中稠密，故得 \(b(u) = b(u^*)^*\)。

我们用 \(\Omega(\mathrm{E})\) 表示由 \(v \in \mathcal{L}(\mathrm{E})\)（满足 \(1_{\mathrm{E}} - v^{*}v\) 正且单射的）所成的集合。对 \(v \in \Omega(\mathrm{E})\)，埃尔米特自同态 \((1_{\mathrm{E}} - v^{*}v)^{1/2}\) 是单射的，因为它的平方是单射的；我们用 \(\mathrm{B}(v)\) 表示部分算子 \(v \circ ((1_{\mathrm{E}} - v^{*}v)^{1/2})^{-1}\)。

注意 \(1_{\mathrm{E}} - v^{*}v\) 正当且仅当 \(\| v\| \leqslant 1\)，因为 \(\langle x|(1_{\mathrm{E}} - v^{*}v)(x)\rangle = \| x\| ^2 -\| v(x)\| ^2\) 对一切 \(x\in \mathrm{E}\) 成立。

引理 1. --- 子集 \(\Omega(\mathrm{E})\) 在 \(\mathcal{L}(\mathrm{E})\) 中是自伴的。

设 \(v \in \Omega(\mathrm{E})\)。我们有 \(\|v^{*}\| = \|v\| \leqslant 1\)。于是自同态 \(1_{\mathrm{E}} - vv^{*}\) 是正的。它是单射的：若 \(x \in \operatorname{Ker}(1_{\mathrm{E}} - vv^{*})\)，我们有 \(vv^{*}(x) = x\)，故 \(v^{*}(v(v^{*}(x))) = v^{*}(x)\)，于是 \(v^{*}(x) = 0\)（因为 \(1_{\mathrm{E}} - v^{*}v\) 单射），最后 \(x = v(v^{*}(x)) = 0\)。引理由此得证。

命题 2. — 映射 \(u \mapsto b(u)\) 定义了从 E 上闭的稠定部分算子集合到集合 \(\Omega(E)\) 上的一个双射。逆双射由 \(v \mapsto B(v)\) 给出。

关系式 \(1_{\mathrm{E}} - b(u)^{*}b(u) = \mathrm{W}(u)\)（公式 (3)）蕴含 \(b(u)\) 属于 \(\Omega (\mathrm{E})\)，因为 \(\mathrm{W}(u)\) 正且单射。此外，由 \(b(u) = u\circ \mathrm{W}(u)^{1 / 2}\)（公式 (2)），我们得到

\[
u = b (u) \circ (\mathrm{W} (u) ^ {1 / 2}) ^ {- 1} = b (u) \circ ((1 - b (u) ^ {*} b (u)) ^ {1 / 2}) ^ {- 1} = \mathrm{B} (b (u)).
\]

反之，设 \(v \in \Omega(\mathrm{E})\)。令 \(w = (1_{\mathrm{E}} - v^{*}v)^{1/2}\)，并令 \(u = \mathrm{B}(v) = v \circ w^{-1}\)。u 的定义域是 w 的像，由于 w 埃尔米特且单射，它在 E 中稠密（EVT，V，第 41 页，命题 4 (i)）。于是对每个 \(x \in \mathrm{dom}(u)\)，有

\[
\begin{array}{r l} & {\| u (x) \| ^ {2} = \langle (v ^ {*} v) (w ^ {- 1} (x)) | w ^ {- 1} (x) \rangle} \\ & {\qquad = - \langle (1 _ {\mathrm{E}} - v ^ {*} v) (w ^ {- 1} (x)) | w ^ {- 1} (x) \rangle + \| w ^ {- 1} (x) \| ^ {2}} \\ & {\qquad = - \langle w (x) | w ^ {- 1} (x) \rangle + \| w ^ {- 1} (x) \| ^ {2}} \end{array}
\]

因为 \((1_{\mathrm{E}} - v^{*}v) \circ w^{-1}\) 是 w 到 \(w^{-1}\) 的定义域上的缩减。由于 w 自伴，由此得

\[
\| w ^ {- 1} (x) \| ^ {2} = \| x \| ^ {2} + \| v \circ w ^ {- 1} (x) \| ^ {2} = \| x \| ^ {2} + \| u (x) \| ^ {2},
\]

于是由 IV，第 231 页引理 2，部分算子 \(\mathrm{B}(v) = u = v \circ w^{-1}\) 是闭的。

最后我们来证明 \(b(\mathrm{B}(v)) = v\)。由于 \(u = v \circ w^{-1}\)，我们有 \(v = u \circ w\)，因此由公式 (2)，只需验证 \(w = \mathrm{W}(u)^{1/2}\)，甚至只需验证 \(1_{\mathrm{E}} - v^{*}v = \mathrm{W}(u)\)。

自同态 \(v^*\) 属于 \(\Omega(\mathrm{E})\)。由引理 1，令 \(w' = (1_{\mathrm{E}} - vv^{*})^{1/2}\)。它是 E 的单射正自同态，且我们有 \(v \circ w = w' \circ v\)（I，第 113 页，命题 11）。\(u\) 的图像是形如 \((w(x), v(x))\)（\(x \in \mathrm{E}\)）的元素所成的集合。由 IV，第 236 页命题 7，\(u^*\) 的图像是 \(s(\Gamma_u)^{\circ}\)。由于对一切 \(x\) 与 \(y \in \mathrm{E}\)，有

\[
\langle (w ^ {\prime} (x), v ^ {*} (x)) | (- v (y), w (y)) \rangle = - \langle x | w ^ {\prime} v (y) \rangle + \langle x | v w (y) \rangle = 0,
\]

故 \(u^*\) 的图像包含形如 \((w'(x), v^*(x))\)（\(x \in \mathrm{E}\)）的元素，并且此时有 \(u^*(w'(x)) = v^*(x)\)。特别地，\(u^*\) 的定义域包含 \(w'\) 的像。

设 \(x \in E\)，令 \(y = w^{2}(x) = (1_{\mathrm{E}} - v^{*}v)(x)\)，于是 \(x = y + v^{*}v(x)\)。我们有 \(y \in \operatorname{dom}(u)\) 且 \(u(y) = v(w^{-1}(y)) = v(w(x)) = w'(v(x))\)。特别地，\(u(y) \in \operatorname{Im}(w) \subset \operatorname{dom}(u^{*})\)，且 \(u^{*}(u(y)) = v^{*}(v(x))\)。由此得 \(x = y + v^{*}v(x) = y + u^{*}u(y)\)，从而 \(y = \operatorname{W}(u)(x)\)，即

\[
(1 _ {\mathrm{E}} - v ^ {*} v) (x) = \mathrm{W} (u) (x).
\]

我们于是证明了 \(1_{\mathrm{E}} - v^{*}v = \mathrm{W}(u)\)，这正是所要的。

推论。--- 自同态 \(b(u)\) 是埃尔米特的，当且仅当 u 自伴。

这由映射 \(u \mapsto b(u)\) 的单射性（命题 2）以及公式 \(b(u^{*}) = b(u)^{*}\)（命题 1 的推论）得出。

\subsection{正规部分算子与谱定理}

在本节中，E 是复希尔伯特空间。

定义 2. --- 设 u 是 E 上的部分算子。称 u 为正规的，如果 u 闭且稠定，并且它的有界变换 b(u) 是 E 的正规自同态。

若 \(u \in \mathcal{L}(\mathrm{E})\)，则由公式 \(1_{\mathrm{E}} - b(u)^{*}b(u) = \mathrm{W}(u) = (1_{\mathrm{E}} + u^{*}u)^{-1}\)（IV，第 262 页，命题 1）与 \(b(u^{*}) = b(u)^{*}\)（同上所引之推论），此定义与 EVT，V，第 42 页，定义 4 相一致。

若 \(u\) 是 \(\mathrm{E}\) 上的自伴部分算子，则 \(b(u)\) 是埃尔米特的（IV，第 264 页命题 2 的推论），故 \(u\) 正规。

设 D 是 C 中的开单位圆盘。我们用 \(\beta\) 表示由 \(\beta(z)=z/\sqrt{1+|z|^{2}}\) 定义的从 C 到 D 中的函数。它是一个同胚，其逆满足 \(\beta^{-1}(z)=z/\sqrt{1-|z|^{2}}\)，对 \(z\in D\)。

设 \(u \in \mathcal{L}(\mathrm{E})\)。由 IV，第 262 页命题 1 的公式 (2) 与 (3) 得 \(u = \beta^{-1}(b(u))\)，从而

\[
b (u) = \beta (u).\tag{5}
\]

引理 2. --- 设 X 是局部紧拓扑空间，μ 是 X 上的正测度。

\begin{enumerate}
\def\labelenumi{\alph{enumi})}
\item
  设 \(g\) 是 X 上 \(\mu\)-可测的函数。乘法算子 \(m_g\)（\(\mathrm{L}^2 (\mathrm{X},\mu)\) 上的）是正规的，且 \(b(m_g) = m_{\beta \circ g}\)；
\item
  设 \(h\colon X\to D\) 是 \(\mu\)-可测函数。自同态 \(m_h\) 属于 \(\Omega(L^2(X,\mu))\)，且 \(B(m_h)=m_{\beta^{-1}\circ h}\)。
\end{enumerate}

部分算子 \(m_g\) 闭且稠定。由于 \(m_g^* m_g = m_{\overline{g}} m_g = m_{|g|^2}\)（IV，第 253 页命题 23 与 IV，第 255 页命题 24），我们有

\[
\mathrm{W} (m _ {g}) = - \mathrm{R} (m _ {| g | ^ {2}}, - 1) = m _ {(1 + | g | ^ {2}) ^ {- 1}}
\]

（IV，第 252 页，命题 22）。由此得

\[
b (m _ {g}) = m _ {g} \circ \mathrm{W} (m _ {g}) ^ {1 / 2} = m _ {g} \circ m _ {(1 + | g | ^ {2}) ^ {- 1 / 2}} = m _ {\beta \circ g}
\]

这里依次用到 IV，第 262 页的公式 (2)、IV，第 187 页命题 6 的推论以及 IV，第 254 页引理 6。最后应用 IV，第 187 页命题 6，即得断言 a)。

我们来证明 \(b\)）。由于 \(1 - |h(x)|^2 > 0\) 对一切 \(x \in \mathrm{X}\) 成立，我们有 \(m_h \in \Omega(\mathrm{L}^2(\mathrm{X}, \mu))\)（IV，第 187 页命题 6 的推论）。由于映射 \(u \mapsto b(u)\) 是单射（IV，第 264 页，命题 2），且自同态 \(b(\mathrm{B}(m_h))\) 与 \(b(m_{\beta^{-1}\circ h})\) 都等于 \(m_h\)（由断言 \(a\)），故确有 \(\mathrm{B}(m_h) = m_{\beta^{-1}\circ h}\)。

定理 1（“谱定理”）。--- 设 u 是希尔伯特空间 E 上的正规部分算子。存在局部紧拓扑空间 X、X 上的正测度 μ、从 L²(X,μ) 到 E 上的等距同构 θ 以及 X 上的连续函数 g，使得 \(u = \theta \circ m_g \circ \theta^{-1}\)。

按定义，E 的自同态 \(b(u)\) 是正规的。由 E 的正规自同态的谱定理（IV，第 193 页，推论 1），存在局部紧拓扑空间 \(\widetilde{X}\)、\(\mu\)（\(\widetilde{X}\) 上的正测度）、\(\widetilde{\theta}\)（从 \(\mathrm{L}^2 (\widetilde{\mathrm{X}},\mu)\) 到 E 上的等距同构），以及 \(h\)（\(\widetilde{X}\) 上的有界连续函数），使 \(b(u)\) 与 \(\widetilde{\theta}\circ m_h\circ \widetilde{\theta}^{-1}\) 相合。

设 N 是由 \(x \in \widetilde{X}\)（满足 \(|h(x)| \geqslant 1\) 的）所成的闭子集。由于自同态 \(1 - b(u)^{*}b(u) = \widetilde{\theta} \circ m_{1-|h|^2} \circ \widetilde{\theta}^{-1}\) 正且单射，集合 N 是局部 \(\mu\)-零测的（IV，第 186 页引理 7 与 IV，第 187 页命题 6 的推论）。于是令 X = \(\widetilde{X} - N\)。它是局部紧空间，且函数到 X 上的限制诱导出从 \(\mathrm{L}^{2}(\widetilde{\mathrm{X}}, \mu)\) 到 \(\mathrm{L}^{2}(\mathrm{X}, \mu|\mathrm{X})\) 上的等距同构（IV，第 179 页，命题 1）。于是 \(\widetilde{\theta}\) 与函数到 X 上的限制的复合诱导出等距同构 \(\theta\)（从 \(\mathrm{L}^{2}(\mathrm{X}, \mu|\mathrm{X})\) 到 E 上的）。由于 \(b(u) = \theta \circ m_{h|\mathrm{X}} \circ \theta^{-1}\)，得 \(u = \theta \circ m_{\beta^{-1} \circ (h|\mathrm{X})} \circ \theta^{-1}\)（引理 2）。定理由此公式得出。

引理 3. --- 设 u 是 E 上的正规部分算子。我们有

\[
\operatorname{Sp} (b (u)) \cap \mathrm{D} = \beta (\operatorname{Sp} (u)).
\]

我们有 \(b(u) = u \circ \mathrm{W}(u)^{1/2}\)（IV，第 262 页，命题 1）。自同态 \(\mathrm{W}(u)^{1/2}\) 是单射的，且它的像是 u 的定义域（IV，第 262 页，公式 (1)）。

设 \(\lambda \in \mathbf{C}\)。数 \(\lambda\) 属于 \(u\) 的预解集，当且仅当 \((u - \lambda 1_{\mathrm{E}}) \circ \mathrm{W}(u)^{1/2}\) 是从 E 到 E 上的双射线性映射（IV，第 245 页，注 3）。而由 IV，第 262 页的公式 (3)，有

\[
(u - \lambda 1 _ {\mathrm{E}}) \circ \mathrm{W} (u) ^ {1 / 2} = b (u) - \lambda (1 _ {\mathrm{E}} - b (u) ^ {*} b (u)) ^ {1 / 2} = f _ {\lambda} (b (u))
\]

其中 \(f_{\lambda}\) 是在 \(\overline{\mathrm{D}}\) 上由 \(z \mapsto z - \lambda (1 - |z|^2)^{1/2}\) 定义的连续函数。自同态 \(f_{\lambda}(b(u))\) 是双射的，当且仅当它的谱不含 0。由于 \(\mathrm{Sp}(f_{\lambda}(b(u))) = f_{\lambda}(\mathrm{Sp}(b(u)))\)（I，第 111 页，命题 7，推论 2），此事成立当且仅当 0 不属于集合 \(f_{\lambda}(\mathrm{Sp}(b(u)))\)。于是，\(\lambda \in \mathrm{Sp}(u)\) 当且仅当存在 \(z \in \mathrm{Sp}(b(u)) \cap \overline{\mathrm{D}}\) 使得

\[
z - \lambda (1 - | z | ^ {2}) ^ {1 / 2} = 0.
\]

此等式若成立，便蕴含 \(z \in D\)，且意味着 \(\lambda = \beta^{-1}(z)\)。我们得出结论 \(\operatorname{Sp}(u) = \beta^{-1}(\operatorname{Sp}(b(u)) \cap D)\)，正如所述。

设 u 是 E 上的正规部分算子。设 \(f \in \mathcal{K}(\mathrm{Sp}(u))\)。映射 \(z \mapsto f(\beta^{-1}(z))\)（从 \(\mathrm{Sp}(b(u)) \cap D\) 到 C 中的）连续且具有紧支集；于是唯一的映射 \(f_{\beta}\)（从 \(\mathrm{Sp}(b(u))\) 到 C 中、用零延拓它的）是连续的。

定义 3. --- 对每个函数 \(f \in \mathcal{K}(\mathrm{Sp}(u))\)，我们定义 E 的自同态 \(f(u)\) 如下：\(f(u) = f_{\beta}(b(u))\)。

映射 \(f \mapsto f_{\beta}\) 是从 \(\mathcal{K}(\mathrm{Sp}(u))\) 到 \(\mathcal{C}(\mathrm{Sp}(b(u)))\) 中的复代数同态，故映射 \(f \mapsto f(u)\) 是从 \(\mathcal{K}(\mathrm{Sp}(u))\) 到 \(\mathcal{L}(\mathrm{E})\) 中的复代数同态。我们有 \(\overline{f}(u) = f(u)^*\)，因为 \(\overline{f}_{\beta} = \overline{f_{\beta}}\)。若 \(f \geqslant 0\)，则 \(f_{\beta} \geqslant 0\)，故 \(f(u)\) 是 E 的正自同态。

注记。--- 设 u 是 E 的正规自同态，于是 \(\mathrm{Sp}(u)\) 是 C 的紧子集。由于有界变换 \(b(u)\) 与 \(\beta(u)\) 相合（公式 (5)），且 \(\mathrm{Sp}(b(u)) = \beta(\mathrm{Sp}(u))\)（I，第 111 页，命题 7，推论 2），\(b(u)\) 的谱是 D 的紧子集。对每个连续函数 \(f \in \mathcal{C}(\mathrm{Sp}(u))\)，函数 \(f_{\beta}\) 与 \(f \circ \beta^{-1}\) 到 \(\mathrm{Sp}(u)\) 上的限制相合。于是，\(f_{\beta}(\beta(u))\) 与由 u 的连续泛函演算所定义的自同态 \(f(u)\) 相合。因此上述定义与星代数 \(\mathcal{L}(\mathrm{E})\) 的连续泛函演算的定义相容。

设 \(f\in \mathcal{K}(\mathrm{Sp}(u))\)。我们有

\[
\| f (u) \| \leqslant \| f _ {\beta} \| _ {\infty} = \| f \| _ {\infty},
\]

于是，对所有 \(x\) 与 \(y\)（E 中的），我们得到估计 \(|\langle x | f(u)y\rangle| \leqslant \|x\|\|y\|\|f\|_{\infty}\)。故映射 \(f \mapsto \langle x | f(u)y\rangle\) 是 \(\mathrm{Sp}(u)\) 上总质量 \(\leqslant \|x\|\|y\|\) 的有界测度（INT，IV，第 154 页，§ 4，第 7 目）。若 \(x = y\)，它是正测度，因为 \(f(u)\) 当 \(f \geqslant 0\) 时是正的。

定义 4. --- 设 u 是复希尔伯特空间 E 上的正规部分算子。设 x 与 y 在 E 中。由 \(f \mapsto \langle x | f(u)y \rangle\)（\(f \in \mathcal{K}(\mathrm{Sp}(u))\)）定义的 Sp(u) 上的有界测度称为 \((x,y)\) 关于 \(u\) 的谱测度。当 \(x = y\) 时，我们称它为 \(x\) 关于 \(u\) 的谱测度。

注记。--- 1) 当 u 连续时，此定义与 IV，第 190 页定义 2 相一致。

\begin{enumerate}
\def\labelenumi{\arabic{enumi})}
\setcounter{enumi}{1}
\tightlist
\item
  设 \(j\) 是 \(\mathrm{Sp}(u)\) 到 \(\mathbf{C}\) 中的典范单射。由于 \(\mathrm{Sp}(u)\) 是闭的，映射 \(j\) 是 \(\mu\)-逆紧的，对每个有界测度 \(\mu\)（\(\mathrm{Sp}(u)\) 上的）（INT，V，第 68 页，§ 6，第 1 目，命题 1）。人们常把 \(u\) 的谱测度与 \(\mathbf{C}\) 上作为它们在 \(j\) 下的像的测度等同起来（参见 INT，V，第 84 页，§ 7，第 2 目）。
\end{enumerate}

从 E×E 到巴拿赫空间 \(\mathcal{M}^{1}(\mathrm{Sp}(u))\) 中把 \((x,y)\) 映为 \((x,y)\) 关于 u 的谱测度的映射是半双线性的，并且是范数 \(\leqslant 1\) 的连续映射。

引理 4. --- 设 X 是局部紧拓扑空间，\(\mu\) 是 X 上的正测度，\(g\) 是 X 上 \(\mu\)-可测的函数。设 \(m_g\) 是乘以 \(g\) 的乘法算子（\(\mathrm{L}^2 (\mathrm{X},\mu)\) 上的）。

\begin{enumerate}
\def\labelenumi{\alph{enumi})}
\item
  映射 \(f\mapsto f(m_{g})\)（从 \(\mathcal{K}(\mathrm{Sp}(m_g))\) 到 \(\mathcal{L}(\mathrm{L}^{2}(\mathrm{X},\mu))\) 中的）由 \(f\mapsto m_{f\circ g}\) 给出；
\item
  对 \(f_1\) 与 \(f_2\)（\(\mathcal{L}^2(\mathrm{X},\mu)\) 中的，其类 \(\widetilde{f}_1\) 与 \(\widetilde{f}_2\) 在 \(\mathrm{L}^2(\mathrm{X},\mu)\) 中），\((\widetilde{f}_1,\widetilde{f}_2)\) 关于 \(m_g\) 的谱测度是限制到 \(\mathrm{Sp}(m_g)\) 上的像测度 \(g(\overline{f}_1f_2\cdot \mu)\)。
\end{enumerate}

部分算子 \(m_g\) 是正规的，且 \(b(m_g) = m_{\beta \circ g} = m_{g(1+|g|^2)^{-1/2}}\)（引理 2）。此外，\(\beta(g(x))\) 属于 \(\mathrm{Sp}(m_{\beta \circ g})\)，对每个 \(x\)（位于局部 \(\mu\)-零测集 \(Y \subset X\) 之外）（IV，第 252 页命题 22 与 IV，第 181 页引理 1）。由于 \(f_\beta(\beta(g(x))) = f(g(x))\) 对一切 \(x \in X - Y\) 成立，由定义 3 与 IV，第 187 页命题 6 的推论，我们推出 \(f(m_g) = m_{f \circ g}\)。

设 \(f_1\) 与 \(f_2\) 是 \(\mathcal{L}^2(\mathrm{X},\mu)\) 中的元素，其类 \(\widetilde{f}_1\) 与 \(\widetilde{f}_2\) 在 \(\mathrm{L}^2(\mathrm{X},\mu)\) 中。由于测度 \(\nu = \overline{f}_1f_2\cdot \mu\) 有界，像测度 \(g(\nu)\) 有定义（INT，V，第 69 页，§ 6，第 1 目，注 1）。设 \(f\in \mathcal{K}(\mathrm{Sp}(m_g))\)。我们有 \(f(m_{g}) = m_{f\circ g}\)（由 \(a\)），于是

\[
\langle \widetilde {f} _ {1} \mid f (m _ {g}) \widetilde {f} _ {2} \rangle = \int_ {\mathrm{X}} \overline {{f}} _ {1} (f \circ g) f _ {2} d \mu = \int_ {\mathrm{X}} (f \circ g) (\overline {{f}} _ {1} f _ {2} d \mu) = \int_ {\mathbf {C}} f d \nu
\]

（INT，V，第 69 页，§ 6，第 1 目，公式 (1)），这就建立了断言 b)。

例。--- 设 \(n \in \mathbf{N}\)。我们给 \(\mathbf{R}^{n}\) 赋以勒贝格测度，并按 II，第 236 页推论 3 把 \(\mathbf{R}^{n}\) 与它的对偶群等同起来。用 N 表示 \(\mathbf{R}^{n}\) 上的欧几里得范数，用 \(\mathcal{F}\) 表示 \(\mathcal{S}(\mathbf{R}^{n})\) 上与 \(\mathcal{S}'(\mathbf{R}^{n})\) 上的傅里叶变换（IV，第 217 页，第 12 目）。

设 \(\Delta_{\mathcal{S}}\) 是 \(\mathrm{L}^2 (\mathbf{R}^n)\) 上以空间 \(\mathcal{S}(\mathbf{R}^n)\) 为定义域且满足

\[
\Delta_ {\mathcal {S}} (\varphi) = - \sum_ {i = 1} ^ {n} \partial_ {i} ^ {2} \varphi
\]

的部分算子，对一切 \(\varphi\in\mathcal{S}(\mathbf{R}^{n})\) 成立。于是 \(\mathcal{F}(\Delta_{\mathcal{S}}(\varphi))=4\pi^{2}\mathrm{N}^{2}\mathcal{F}(\varphi)\)，对每个函数 \(\varphi\in\mathcal{S}(\mathbf{R}^{n})\)（IV，第 220 页，注 12）。由于傅里叶变换是 \(\mathcal{S}(\mathbf{R}^{n})\) 的自同构（IV，第 219 页，命题 18），这表明部分算子 \(u=\mathcal{F}\circ\Delta_{\mathcal{S}}\circ\mathcal{F}^{-1}\) 是乘法算子 \(m_{4\pi^{2}\mathrm{N}^{2}}\) 到空间 \(\mathcal{S}(\mathbf{R}^{n})\) 上的缩减。部分算子 \(u\) 可闭且对称；它的闭包是正的自伴部分算子 \(m_{4\pi^{2}\mathrm{N}^{2}}\)（IV，第 232 页命题 6，对空间 \(\mathcal{D}(\mathbf{R}^{n})\subset\mathcal{S}(\mathbf{R}^{n})\) 应用，并用到 IV，第 202 页命题 4）。因此，部分算子 \(\Delta_{\mathcal{S}}\) 本质自伴。我们用 \(\Delta\) 表示它的闭包；它是正的自伴算子，并且是 \(\mathbf{R}^{n}\) 上唯一的拉普拉斯算子（\(cf.\) IV，第 243 页第 6 目的例以及 IV，第 257 页命题 26 的推论）。

由 IV，第 223 页命题 21，索伯列夫空间 \(\mathrm{H}^2 (\mathbf{R}^n)\) 是由 \(f\in \mathrm{L}^2 (\mathbf{R}^n)\)（使 \((1 + \mathrm{N}^2)\mathcal{F}(f)\) 属于 \(\mathrm{L}^2 (\mathbf{R}^n)\) 的）所成的集合，亦即使 \(\mathcal{F}(f)\) 属于 \(m_{4\pi^2\mathrm{N}^2}\) 的定义域。由于傅里叶变换是空间 \(\mathrm{L}^2 (\mathbf{R}^n)\) 到其自身上的等距同构（II，第 215 页，定理 1），拉普拉斯算子 \(\Delta\) 的定义域是 \(\mathrm{H}^2 (\mathbf{R}^n)\)。我们有 \(\Delta = \mathcal{F}^{-1}\circ m_{4\pi^2\mathrm{N}^2}\circ \mathcal{F}\)。特别地，\(\Delta\) 的谱等于 \(\mathbf{R}_{+}\)。
\subsection{普遍可测函数演算}

在本节中，我们所考虑的希尔伯特空间都是复的。

设 u 是希尔伯特空间 E 上的正规部分算子。我们用 \(\mu_{x,y}\) （相应地，\(\mu_{y}\) ）表示 \((x,y)\in\mathrm{E}\times\mathrm{E}\) （相应地，y）关于 u 的谱测度。

设 \(y \in \mathrm{E}\) 。映射 \(\mathcal{K}(\mathrm{Sp}(u)) \to \mathrm{E}\) 由 \(f \mapsto f(u)(y)\) 所定义，它满足

\[
\| f (u) (y) \| ^ {2} = \int_ {\mathrm{Sp} (u)} | f | ^ {2} \mu_ {y} = \| f \| _ {\mathcal {L} ^ {2} (\mathrm{Sp} (u), \mu_ {y})} ^ {2}.
\]

因此，存在唯一的等距线性映射 \(ev_{y}\) ，它从空间 \(\mathrm{L}^{2}(\mathrm{Sp}(u),\mu_{y})\) 映入 E，使得 \(\mathrm{ev}_{y}(\widetilde{f})=f(u)(y)\) ，这里 \(\widetilde{f}\) 是函数 \(f\in\mathcal{K}(\mathrm{Sp}(u))\) 的类。

设 \(f\) 是定义在 \(u\) 的谱上的普遍可测函数。我们用 \(\mathrm{D}_{f}\) 表示这样的元素 \(y \in \mathrm{E}\) 所成的集：\(f\) 属于 \(\mathcal{L}^{2}(\mathrm{Sp}(u), \mu_{y})\) 。

命题 3. --- 设 f 是定义在 u 的谱上的普遍可测函数。集合 D \(_{f}\) 是 E 的稠密向量子空间。映射 y ↦ ev \(_{y}\) (f) 是 E 上的正规部分算子，其定义域为 D \(_{f}\) ，记作 f(u)。

对每个 \(x \in \mathrm{E}\) 及每个 \(y \in \mathrm{D}_f\) ，我们有 \(f \in \mathcal{L}^1(\mathrm{Sp}(u), \mu_{x,y})\) ，并且

\[
\langle x \mid f (u) y \rangle = \int_ {\operatorname{Sp} (u)} f \mu_ {x, y}.\tag{6}
\]

我们将证明下列更精确的陈述：

命题 4. --- 设 f 是 u 的谱上的普遍可测函数。

\begin{enumerate}
\def\labelenumi{\alph{enumi})}
\item
  集合 \(D_{f}\) 是 E 的稠密向量子空间。映射 \(f(u): y \mapsto \operatorname{ev}_{y}(f)\) 从 \(D_{f}\) 到 E 中是线性的，且当 f = 1 时与 \(1_{E}\) 重合；
\item
  设 \(g \in \mathcal{L}_{\mathrm{u}}(\mathrm{Sp}(u))\) 。对每个 \(y \in \mathrm{D}_f\) 及每个 \(x \in \mathrm{D}_g\) ，我们有 \(fg \in \mathcal{L}^1(\mathrm{Sp}(u), \mu_{x,y})\) ，并且
\end{enumerate}

\[
\langle g (u) x \mid f (u) y \rangle = \int_ {\mathrm{Sp} (u)} \overline {{g}} f \mu_ {x, y}.\tag{7}
\]

\begin{enumerate}
\def\labelenumi{\alph{enumi})}
\setcounter{enumi}{2}
\tightlist
\item
  假设 \(E = L^{2}(X, \mu)\) ，其中 X 是局部紧拓扑空间，\(\mu\) 是 X 上的正测度；又假设 \(u = m_{h}\) ，其中 \(h: X \to \mathbf C\) 是 \(\mu\)-可测的。我们有 \(f(m_{h}) = m_{f \circ h}\) 。
\end{enumerate}

由 IV, p. 266 的定理 1，我们可以假设处于断言 \(c\) ) 的情形，即 \(\mathrm{E} = \mathrm{L}^2 (\mathrm{X},\mu)\) 且 \(u = m_h\) ，其中 X 是局部紧拓扑空间，\(\mu\) 是 X 上的正测度，而 \(h\colon \mathrm{X}\to \mathbf{C}\) 是 \(\mu\)-可测的。我们将用 \(\widetilde{\varphi}\) 表示 \(\mathrm{L}^2 (\mathrm{X},\mu)\) 中函数 \(\varphi \in \mathcal{L}^2 (\mathrm{X},\mu)\) 的类。

设 \(\mathrm{S} = \mathrm{Sp}(m_h)\) 。我们用 \(\mu_{\widetilde{\varphi}_1, \widetilde{\varphi}_2}\) （相应地，\(\mu_{\widetilde{\varphi}}\) ）表示 \((\widetilde{\varphi}_1, \widetilde{\varphi}_2)\) （相应地，\(\widetilde{\varphi}\) ）关于 \(m_h\) 的谱测度，对所有 \((\varphi_1, \varphi_2) \in \mathcal{L}^2(\mathrm{X}, \mu)^2\) （相应地，每个 \(\varphi \in \mathcal{L}^2(\mathrm{X}, \mu)\) ）。

设 \(\varphi\in\mathcal{L}^{2}(X,\mu)\) 。谱测度 \(\mu_{\widetilde{\varphi}}\) 等于 S 上的像测度 \(h(|\varphi|^{2}\cdot\mu)\) （IV, p. 269 的引理 4）。由于函数 \(f\) 是 \(\mu_{\varphi}\)-可测的，我们有 \(\varphi\in D_{f}\) 当且仅当积分

\[
\int_ {\mathrm{S}} ^ {*} | f | ^ {2} d \mu_ {\widetilde {\varphi}} = \int_ {\mathrm{S}} ^ {*} | f | ^ {2} h (| \varphi | ^ {2} d \mu) = \int_ {\mathrm{X}} ^ {*} | f \circ h | ^ {2} | \varphi | ^ {2} d \mu
\]

都是有限的（INT, V, p. 70, § 6, 第 2 目, 命题 2）。这表明 \(D_{f}\) 是乘法算子 \(m_{f\circ h}\) 的定义域，而该定义域是 E 的稠密子空间（IV, p. 232 的命题 5）。

\(\mathrm{ev}_{\widetilde{\varphi}}\) 在函数类 \(g \in \mathcal{K}(\mathrm{S})\) 上的限制是把类 \(\widetilde{g}\) 对应到元素 \(g(m_h)(\widetilde{\varphi}) = m_{g \circ h}(\widetilde{\varphi})\) 的映射（IV, p. 269 的引理 4）。因此，映射 \(\mathrm{ev}_{\widetilde{\varphi}}\) 与从 \(\mathrm{L}^2(\mathrm{S}, \mu_{\widetilde{\varphi}})\) 到 \(\mathrm{L}^2(\mathrm{X}, \mu)\) 中的等距映射重合，后者是由映射 \(g \mapsto (g \circ h) \cdot \varphi\) （从 \(\mathcal{L}^2(\mathrm{S}, \mu_{\widetilde{\varphi}})\) 到 \(\mathcal{L}^2(\mathrm{X}, \mu)\) 中）取商得到的。

特别地，我们因此有 \(\mathrm{ev}_{\widetilde{\varphi}}(f)=m_{f\circ h}(\widetilde{\varphi})\) ，于是 \(f(m_{h})\) 作为从 \(D_{f}\) 到 E 中的映射与部分算子 \(m_{f\circ h}\) 重合。这就证明了断言 c)；若 f=1，则得到 \(\mathrm{ev}_{\widetilde{\varphi}}(1)=\widetilde{\varphi}\) ，这也就完成了 a) 的证明。

我们来证明断言 \(b\) )。设 \(\varphi_{1}\) 与 \(\varphi_{2}\) 是使 \(\widetilde{\varphi}_{1} \in \mathrm{D}_{f}\) 且 \(\widetilde{\varphi}_{2} \in \mathrm{D}_{g}\) 的函数。由（INT, V, loc. cit.）可得

\[
\begin{array}{l} \int_ {\mathrm{S}} ^ {*} | f g |   | \mu_ {\widetilde {\varphi_ {1}}, \widetilde {\varphi_ {2}}} | = \int_ {\mathrm{X}} ^ {*} | (f \circ h) (g \circ h) \varphi_ {1} \varphi_ {2} | d \mu \\ \qquad \leqslant \| (f \circ h) \varphi_ {1} \| \| (g \circ h) \varphi_ {2} \| \end{array}
\]

它是有限的，因此 \(fg \in \mathcal{L}^1(\mathrm{S}, \mu_{\widetilde{\varphi}_1, \widetilde{\varphi}_2})\) 。此外，这时我们有

\[
\begin{array}{c} \langle g (m _ {h}) (\widetilde {\varphi} _ {2})   |   f (m _ {h}) (\widetilde {\varphi} _ {1}) \rangle = \int_ {\mathrm{X}} \overline {{(g \circ h)   \varphi_ {2}}}    (f \circ h)   \varphi_ {1}    d \mu \\ = \int_ {\mathrm{S}} \overline {{g}}   f    d \mu_ {\widetilde {\varphi} _ {1}, \widetilde {\varphi} _ {2}}, \end{array}
\]

证明完毕。

定义 5. --- 映射 \(f \mapsto f(u)\) 由 IV, p. 271 的命题 3 所定义，它从 \(\mathcal{L}_{\mathrm{u}}(\mathrm{Sp}(u))\) 映到 E 上正规部分算子所成的集，称为与 \(u\) 相伴的普遍可测函数演算映射。

更一般地，设 T 是一个集合，并设 \(g\colon\operatorname{Sp}(u)\times\operatorname{T}\to\mathbf{C}\) 是一个映射；设 \(t\in T\) 使得函数 \(g_{t}\colon z\mapsto g(z,t)\) 在 \(\operatorname{Sp}(u)\) 上普遍可测。这时我们记 \(g(u,t)=g_{t}(u)\) 。

公式 (7) 特别地蕴含

\[
\| f (u) (y) \| ^ {2} = \int_ {\operatorname{Sp} (u)} | f | ^ {2} \mu_ {y}\tag{8}
\]

对所有 \(f \in \mathcal{L}_{\mathrm{u}}(\mathrm{Sp}(u))\) 及每个 \(y \in \mathrm{D}_f\) 成立。由此，取 \(f = 1\) ，我们推出 \(\mu_y\) 是总质量为 \(\| y\|^2\) 的正测度。

若 \(u\) 是希尔伯特空间 E 上的正规部分算子且 \(f \in \mathcal{K}(\mathrm{Sp}(u))\) ，则 \(\mathrm{D}_f = \mathrm{E}\) ，且部分算子 \(f(u)\) 这时是连续的，因为

\[
\| f (u) y \| \leqslant \| f \| _ {\infty} \| y \|
\]

对所有 \(y \in E\) 成立（依据 (8)）。自同态 \(f(u)\) 与 IV, p. 268 的定义 3 中的那个自同态重合。

推论 1. --- a) 对每个 \(f \in \mathcal{L}_{\mathrm{u}}(\mathrm{Sp}(u))\) ，部分算子 \(f(u)\) 是正规的，且 \(f(u)^{*} = \overline{f}(u)\) 。此外，\(f(u)\) 在 \(f \geqslant 0\) 时是正的，在 \(f\) 取实值时是自伴的；

\begin{enumerate}
\def\labelenumi{\alph{enumi})}
\setcounter{enumi}{1}
\item
  设 \(k \in \mathbf{N}\) ，令 \(f(z) = z^{k}\) ，\(z \in \operatorname{Sp}(u)\) 。我们有 \(f(u) = u^{k}\) ；
\item
  设 \(\lambda \in \mathbf{C} - \mathrm{Sp}(u)\) ，令 \(f(z) = (\lambda - z)^{-1}\) ，\(z \in \mathrm{Sp}(u)\) 。我们有 \(f(u) = \mathrm{R}(u, \lambda)\) ；
\item
  我们有 \(\beta(u) = b(u)\) ；
\item
  设 \(f \in \mathcal{L}_{\mathrm{u}}(\mathrm{Sp}(u))\) 与 \(g \in \mathcal{L}_{\mathrm{u}}(\mathrm{Sp}(u))\) 满足 \(|g| \leqslant 1 + |f|\) 。\(g(u)\) 的定义域是 \(f(u)\) 的核。
\end{enumerate}

鉴于谱定理（IV, p. 266 的定理 1），这由前面的命题分别结合以下各项而得：

\begin{enumerate}
\def\labelenumi{\alph{enumi})}
\item
  IV, p. 266 的引理 2, a)、IV, p. 253 的命题 23 及其推论；
\item
  IV, p. 255 的命题 24, b)；
\item
  IV, p. 252 的命题 22, b)；
\item
  IV, p. 266 的引理 2, a)；
\item
  IV, p. 232 的命题 6, b)。
\end{enumerate}

注. --- 对 \(f = \operatorname{Id}_{\operatorname{Sp}(u)}\) ，我们有 \(f(u) = u\) （断言 b) 取 \(k = 1\) ）。因此，u 的定义域与这样的 \(x \in E\) 所成的集重合：\(\operatorname{Sp}(u)\) 的恒等映射属于 \(\mathcal{L}^{2}(\operatorname{Sp}(u), \mu_{x})\) ；它特别包含这样的元素 \(x \in E\) ：测度 \(\mu_{x}\) 具有紧支撑。

下面的推论推广了 IV, p. 247 的命题 16 的推论以及 IV, p. 248 的命题 17。

推论 2. --- 设 u 是希尔伯特空间 E 上的正规部分算子。假设 E 非零。

\begin{enumerate}
\def\labelenumi{\alph{enumi})}
\item
  u 的谱非空；
\item
  对每个 \(\lambda \in \mathbf{C} - \mathrm{Sp}(u)\) ，预解式 \(\mathrm{R}(u, \lambda)\) 的范数等于 \(1/\delta\) ，其中 \(\delta > 0\) 是在 \(\mathbf{C}\) 中从 \(\lambda\) 到 \(u\) 的谱的距离。
\item
  对每个 \(\varepsilon > 0\) ，\(\varepsilon\)-伪谱 \(\mathrm{PSp}_{\varepsilon}(u)\) 是由 \(\lambda \in \mathbf{C}\) （距离 \(< \varepsilon\) ，从 \(\mathrm{Sp}(u)\) 量起）所成的集。
\end{enumerate}

假设 \(u\) 的谱为空。那么 \(u\) 是单射的，且自同态 \(u^{-1} = -\mathrm{R}(u,0)\) 是正规的（推论 1, c) 与 a)）。由于 \(\mathrm{Sp}(u^{-1}) \subset \{0\}\) （命题 15, a)），我们有 \(\mathrm{Sp}(u^{-1}) = \{0\}\) （I, p. 26, 推论 1）。由于 \(u^{-1}\) 是正规的，这就蕴含 \(u^{-1} = 0\) （I, p. 110, 例 1），而这是矛盾，因为 E 不为零。

为证明 \(b\) )，可以假设 \(u\) 是乘法算子 \(m_g\) ，作用于 \(\mathrm{L}^2 (\mathrm{X},\mu)\) ，其中 \(g\) 是赋以正测度 \(\mu\) 的局部紧拓扑空间 X 上的连续函数（IV, p. 266 的定理 1）。设 \(\lambda \in \mathbf{C} - \mathrm{Sp}(u)\) ，并设 \(\delta >0\) 是 \(\lambda\) 到 \(u\) 的谱的距离。为证明 \(\| \mathrm{R}(u,\lambda)\| = \delta^{-1}\) ，我们归结到 \(\lambda = 0\) 的情形，办法是把 \(u\) 换成 \(u - \lambda 1_{\mathrm{E}}\) 。于是实数 \(\delta\) 是 0 到 \(m_g\) 的谱的距离。由于后者与以 \(\mu\) 为测度的 \(g\) 的本质像重合，结果是 IV, p. 182 的引理 3 的推论。

最后，断言 \(c\) ) 由 \(b\) ) 及 \(\mathrm{PSp}_{\varepsilon}(u)\) 的定义（IV, p. 250, 定义 9）得出。

推论 3. --- 设 u 是希尔伯特空间 E 上的正规部分算子。设 \(f \in \mathcal{L}_{\mathrm{u}}(\mathrm{Sp}(u))\) 。对 E 中所有 x 与 y，\((x, y)\) 关于 \(f(u)\) 的谱测度是像测度 \(f(\mu)\) ，其中 \(\mu\) 是 \((x, y)\) 关于 u 的谱测度。

鉴于谱定理（IV, p. 266 的定理 1），可以假设 \(u\) 是乘法算子 \(m_g\) （作用于 \(\mathrm{L}^2 (\mathrm{X},\mu)\) ），其中 X 是局部紧拓扑空间，\(\mu\) 是 X 上的正测度，且 \(g\in \mathcal{C}(\mathrm{X})\) 。设 \(f_{1}\) 与 \(f_{2}\) 属于 \(\mathcal{L}^2 (\mathrm{X},\mu)\) ，其类 \(\widetilde{f}_1\) 与 \(\widetilde{f}_2\) 在 \(\mathrm{L}^2 (\mathrm{X},\mu)\) 中。由于 \(f(m_g) = m_{f\circ g}\) （命题 4, c)），\((\widetilde{f}_1,\widetilde{f}_2)\) 关于 \(f(m_g)\) 的谱测度是像测度 \((f\circ g)(\overline{f}_1f_2\cdot \mu)\) （IV, p. 269 的引理 4, b)）。这个测度等于像测度 \(f(g(\overline{f}_1f_2\cdot \mu))\) （INT, V, p. 72, § 6, 第 4 目, 命题 4, a)），由此即得断言（IV, p. 269 的引理 4, b)）。

推论 4. --- 设 \(g \in \mathcal{L}_{\mathrm{u}}(\mathrm{Sp}(f(u)))\) 。我们有 \(g(f(u)) = (g \circ f)(u)\) 。

我们有 \(g \circ f \in \mathcal{L}_{\mathrm{u}}(\mathrm{Sp}(u))\) （IV, p. 184 的引理 5）。对 E 中所有 \(x\) 与 \(y\) ，用 \(\mu_{x,y}^{\prime}\) 表示 \((x,y)\) 关于 \(f(u)\) 的谱测度。由前面的推论以及 INT, V, p. 71, § 6, 第 2 目, 定理 1，我们有 \(g \in \mathcal{L}^{2}(\mathrm{Sp}(f(u)), \mu_{x,y}^{\prime})\) 当且仅当 \(g \circ f \in \mathcal{L}^{2}(\mathrm{Sp}(u), \mu_{x,y})\) ，因此 \(g(f(u))\) 的定义域等于 \((g \circ f)(u)\) 的定义域。于是，对所有 \(x \in \mathrm{E}\) 及 \(y \in \operatorname{dom}(g(f(u)))\) ，有公式

\[
\begin{array}{c} \langle x \mid g (f (u)) y \rangle = \int_ {\mathrm{Sp} (f (u))} g \mu_ {x, y} ^ {\prime} \\ = \int_ {\mathrm{Sp} (f (u))} g f (\mu_ {x, y}) = \langle x \mid (g \circ f) y \rangle \end{array}
\]

（loc. cit.），由此 \(g(f(u)) = (g \circ f)(u)\) 。

命题 5. --- 设 u 是希尔伯特空间 E 上的正规部分算子。

\begin{enumerate}
\def\labelenumi{\alph{enumi})}
\item
  \(If f \in \mathcal{L}_{\mathrm{u}}^{\infty}(\mathrm{Sp}(u)),\) 则 \(f(u) \in \mathcal{L}(\mathrm{E})\) ；
\item
  \(If f \in \mathcal{L}_{\mathrm{u}}^{\infty}(\mathrm{Sp}(u))\) 且 \(g \in \mathcal{L}_{\mathrm{u}}(\mathrm{Sp}(u))\) ，则 \(f(u) \circ g(u) \subset (fg)(u)\) ；
\item
  映射 \(f \mapsto f(u)\) 从 \(\mathcal{L}_{\mathrm{u}}^{\infty}(\mathrm{Sp}(u))\) 到 \(\mathcal{L}(\mathrm{E})\) 中是星代数的连续单位态射。特别地，有 \(\|f(u)\|\leqslant\|f\|_{\infty}\) ，其中 \(f\in\mathcal{L}_{\mathrm{u}}^{\infty}(\mathrm{Sp}(u))\) ；
\item
  若 \(u \in \mathcal{L}(\mathrm{E})\) ，则对所有 \(f \in \mathcal{L}_{\mathrm{u}}^{\infty}(\mathrm{Sp}(u))\) ，自同态 \(f(u)\) 属于 \(u\) 在 \(\mathcal{L}(\mathrm{E})\) 中的双交换子。
\end{enumerate}

设 \(f \in \mathcal{L}_{\mathrm{u}}^{\infty}(\mathrm{Sp}(u))\) 。由 p. 272 的公式 (8)，我们有 \(\mathrm{D}_f = \mathrm{E}\) ，且 \(f(u)\) 是 \(\mathrm{E}\) 的连续自同态，由此得断言 \(a\) )。

设 \(g \in \mathcal{L}_{\mathrm{u}}(\mathrm{Sp}(u))\) 。设 \(y \in \mathrm{D}_g\) 。我们有 \(y \in \mathrm{D}_{fg}\) ，并且对所有 \(x \in \mathrm{E}\) ，由 p. 271 的公式 (7) 可得 \(\langle \overline{f}(u)x | g(u)y \rangle = \langle x | (fg)(u)y \rangle\) ，由此 \(f(u)(g(u)y) = (fg)(u)(y)\) ，这就证明了 \(b\) 。

由 \(a\) ) 与 \(b\) )，映射 \(f \mapsto f(u)\) 从 \(\mathcal{L}_{\mathrm{u}}^{\infty}(\mathrm{Sp}(u))\) 到 \(\mathcal{L}(\mathrm{E})\) 中，是对合代数的单位态射；因而它是范数 \(\leqslant 1\) 的连续态射（I, p. 104, 命题 2），由此得断言 \(c\) )。

假设 \(u\) 是 E 的自同态。设 \(v \in \mathcal{L}(\mathrm{E})\) 与 \(u\) 交换。于是，由连续函数演算的性质（I, p. 110, 注），我们有 \(v \circ f(u) = f(u) \circ v\) ，这里 \(f \in \mathcal{C}(\mathrm{Sp}(u))\) 。设 \(x\) 与 \(y\) 属于 E。公式

\[
\langle x \mid (v \circ f (u)) y \rangle = \langle x \mid (f (u) \circ v) y \rangle ,\tag{9}
\]

对每个函数 \(f \in \mathcal{C}(\mathrm{Sp}(u))\) 成立，它表示 \((v^{*}(x), y)\) 与 \((x, v(y))\) 关于 u 的谱测度相等。由 p. 271 的公式 (6)，这个等式蕴含公式 (9) 对所有 \(f \in \mathcal{L}_{\mathrm{u}}^{\infty}(\mathrm{Sp}(u))\) 成立。因此我们有 \(v \circ f(u) = f(u) \circ v\) 。

推论. --- 设 \(f\) 与 \(g\) 属于 \(\mathcal{L}_{\mathrm{u}}(\mathrm{Sp}(u))\) ，并设 \((x,y)\in\mathrm{D}_{f}\times\mathrm{D}_{g}\) 。\((f(u)x,g(u)y)\) 关于 \(u\) 的谱测度是测度 \(\overline{f}g\cdot\mu_{x,y}\) ，其中 \(\mu_{x,y}\) 是 \((x,y)\) 关于 \(u\) 的谱测度。

用 \(\nu\) 表示 \((f(u)x,g(u)y)\) 关于 \(u\) 的谱测度。对每个函数 \(\varphi \in \mathcal{H}(\mathrm{Sp}(u))\) ，我们有

\[
\begin{array}{r l} & {\int_ {\mathrm{Sp} (u)} \varphi \nu = \langle f (u) x | \varphi (u) (g (u) y) \rangle} \\ & {\qquad = \langle f (u) x | (\varphi g) (u) y \rangle = \int_ {\mathrm{Sp} (u)} \overline {{f}} \varphi g \mu_ {x, y},} \end{array}
\]

（命题 5, b)）由此 \(\nu = \overline{f} g \cdot \mu_{x,y}\) 。

命题 6. --- 设 u 是希尔伯特空间 E 上的正规部分算子。设 \((f_n)_{n \in \mathbf{N}}\) 是 \(\mathcal{L}_{\mathrm{u}}(\mathrm{Sp}(u))\) 中的序列，它逐点收敛于 \(f \in \mathcal{L}_{\mathrm{u}}(\mathrm{Sp}(u))\) ，且存在 \(g \in \mathcal{L}_{\mathrm{u}}(\mathrm{Sp}(u))\) 满足 \(|f_n| \leqslant g\) ，对所有 \(n \in \mathbf{N}\) 。这时我们有 \(\operatorname{dom}(g(u)) \subset \operatorname{dom}(f_n(u))\) ，对所有 \(n \in \mathbf{N}\) ，以及 \(\operatorname{dom}(g(u)) \subset \operatorname{dom}(f(u))\) 。此外，对 \(g(u)\) 的定义域的每个元素 y，我们有

\[
f (u) (y) = \lim _ {n \to + \infty} f _ {n} (u) (y).
\]

特别地，若 \(f_{n} \in \mathcal{L}_{\mathrm{u}}^{\infty}(\mathrm{Sp}(u))\) 对所有 \(n \in \mathbf{N}\) 成立，且函数 \(f_{n}\) 一致有界，则 \(f_{n}(u)\) 收敛于 \(f(u)\) ，这里的空间 \(\mathcal{L}(\mathrm{E})\) 赋以逐点收敛拓扑。

用 \(\mu_{x,y}\) （相应地，\(\mu_{x}\) ）表示 \((x,y)\) （相应地，x）关于 u 的谱测度。

设 \(y \in \mathrm{dom}(g(u))\) ，于是 \(g \in \mathcal{L}^2(\mathrm{Sp}(u), \mu_y)\) 。条件 \(|f_n| \leqslant g\) 蕴含 \(f_n \in \mathcal{L}^2(\mathrm{Sp}(u), \mu_y)\) ，因此 \(y \in \mathrm{dom}(f_n(u))\) 。

由于 \((f_{n})\) 逐点收敛于 f 且 \(|f_{n}| \leqslant g\) ，由勒贝格定理（INT, IV, p. 137, § 3, 第 7 目, 定理 6），我们有 \(f \in \mathcal{L}^{2}(\mathrm{Sp}(u), \mu_{y})\) ，因此 \(y \in \mathrm{dom}(f(u))\) 。此外，序列 \((f_{n})\) 在 \(\mathcal{L}^{2}(\mathrm{Sp}(u), \mu_{y})\) 中收敛于 f，于是 \(f_{n}(u)y\) 的范数收敛于 \(f(u)y\) 的范数。

设 \(x \in \mathrm{E}\) 。函数 \(f\) 与 \(g\) ，以及函数 \(f_n\) （对所有 \(n \in \mathbf{N}\) ），都属于 \(\mathcal{L}^1(\mathrm{Sp}(u), \mu_{x,y})\) （命题 3）。由勒贝格定理（INT, IV, loc. cit.），序列 \((f_n)\) 收敛于 \(f\) （在 \(\mathcal{L}^1(\mathrm{Sp}(u), \mu_{x,y})\) 中），由此

\[
\langle x \mid f _ {n} (u) y \rangle = \int_ {\operatorname{Sp} (u)} f _ {n} \mu_ {x, y} \rightarrow \int_ {\operatorname{Sp} (u)} f \mu_ {x, y} = \langle x \mid f (u) y \rangle .
\]

由此得出 \(f_{n}(u)(y)\) 收敛于 \(f(u)(y)\) （EVT, V, p. 17, 命题 10）。

命题 7. --- 设 X 是局部紧拓扑空间，ν 是 X 上的测度。设 g: C × X → C 是具有紧支撑的连续函数。对 z ∈ C，置

\[
h (z) = \int_ {\mathrm{X}} g (z, x) d \nu (x).
\]

\begin{enumerate}
\def\labelenumi{\alph{enumi})}
\item
  从 C 到 C 的映射 h 是连续且有界的；
\item
  从 X 到 \(\mathcal{L}(\mathrm{E})\) 的、由 \(x\mapsto g(u,x)\) 定义的映射是 \(\nu\)-可积的，并且我们有
\end{enumerate}

\[
h (u) = \int_ {\mathrm{X}} g (u, x) d \nu (x).
\]

函数 h 是有界的，因为 g 是具有紧支撑的连续函数；且由 INT, IV, p. 144, § 4, 第 3 目, 推论 1，它是连续的。由于 g 连续且具有紧支撑，映射 \(x \mapsto g(u, x)\) 从 X 到 \(\mathcal{L}(\mathrm{E})\) 中连续。这个映射具有紧支撑，故在 X 上关于 \(\nu\) 有界且可积。我们记

\[
v = \int_ {\mathrm{X}} g (u, x) d \nu (x) \in \mathscr {L} (\mathrm{E}).
\]

设 y 与 z 是 E 的元素，并设 \(\mu\) 是 \((y, z)\) 关于 u 的谱测度。我们有

\[
\langle y \mid v (z) \rangle = \int_ {\mathrm{X}} \langle y \mid g (u, x) z \rangle d \nu (x) = \int_ {\mathrm{X}} \left(\int_ {\operatorname{Sp} (u)} g (\lambda , x) d \mu (\lambda)\right) d \nu (x)
\]

（p. 271 的公式 (6)）。由于 \(g \in \mathcal{K}(\mathbf{C} \times \mathbf{X})\) ，由此可得

\[
\langle y \mid v (z) \rangle = \int_ {\operatorname{Sp} (u)} \Bigl (\int_ {\mathrm{X}} g (\lambda , x) d \nu (x) \Bigr) d \mu (\lambda) = \langle y \mid h (u) z \rangle
\]

依据 INT, III, p. 84, § 4, 第 1 目, 定理 2 以及 p. 271 的公式 (6)。这就证明了 \(v = h(u)\) ，此即所要证者。

\subsection{谱投影}

设 \(u\) 是复希尔伯特空间 E 上的正规部分算子。设 A 是 \(\mathrm{Sp}(u)\) 的普遍可测子集，\(\varphi_{\mathrm{A}}\) 是它的特征函数。由于 \(\varphi_{\mathrm{A}}\) 有界且满足 \(\varphi_{\mathrm{A}}^{2} = \varphi_{\mathrm{A}}\) ，E 的自同态 \(\varphi_{\mathrm{A}}(u)\) 是 E 的正交投影。它称为由 A 定义的 \(u\) 的谱投影。我们记 \(p_{\mathrm{A}} = \varphi_{\mathrm{A}}(u)\) 。我们有 \(p_{\emptyset} = 0\) 及 \(p_{\mathrm{Sp}(u)} = 1_{\mathrm{E}}\) 。

设 A 是 C 的普遍可测子集。由 A 定义的 u 的谱投影是谱投影 \(p_{\mathrm{Sp}(u)\cap\mathrm{A}}\) 。它也简单地记作 \(p_{A}\) 。对每个函数 \(f \in \mathcal{L}_{\mathrm{u}}(\mathrm{Sp}(u))\) ，对所有 \(x \in E\) 及每个 \(y \in \mathrm{dom}(f(u))\) ，我们有公式

\[
\langle p _ {\mathrm{A}} (x) \mid f (u) y \rangle = \int_ {\operatorname{Sp} (u) \cap \mathrm{A}} f d \mu ,\tag{10}
\]

其中 \(\mu\) 是 \((x,y)\) 关于 \(u\) 的谱测度（p. 271 的公式 (7)）。设 A 与 B 是 \(\mathrm{Sp}(u)\) 的普遍可测子集。由于 \(\varphi_{\mathrm{A}}\varphi_{\mathrm{B}} = \varphi_{\mathrm{A}\cap \mathrm{B}}\) ，我们有 \(p_{\mathrm{A}} \circ p_{\mathrm{B}} = p_{\mathrm{A}\cap \mathrm{B}}\) （IV, p. 275 的命题 5, \(c\) )）。特别地，若 A 与 B 不相交，则 \(p_{\mathrm{A}}\) 的像与 \(p_{\mathrm{B}}\) 的像正交。

命题 8. --- 设 \((\mathrm{A}_{i})_{i\in\mathrm{I}}\) 是 \(\mathrm{Sp}(u)\) 的两两不相交的普遍可测子集组成的可数族，并设 \(p_{i}\) 是由 \(A_{i}\) 定义的 u 的谱投影。诸集合 \(A_{i}\) 的并 A 是 \(\mathrm{Sp}(u)\) 的普遍可测子集，且级数 \(\sum_{i}p_{i}\) 收敛于 \(p_{A}\) ，这里的 \(\mathcal{L}(\mathrm{E})\) 赋以逐点收敛拓扑。

由 INT, IV, p. 177, § 5, 第 4 目, 推论 2，集合 A 是普遍可测的。级数 \(\sum_{i} \varphi_{\mathrm{A}_{i}}\) 逐点收敛于 \(\varphi_{\mathrm{A}}\) ，且其部分和以 1 为界。因此断言由 IV, p. 276 的命题 6 得出。

命题 9. --- 设 A 是 Sp(u) 的闭子集。设 \(\varphi_{\mathrm{A}}\) 是 A 的特征函数，\(p_{\mathrm{A}} = \varphi_{\mathrm{A}}(u)\) 是由 A 定义的 u 的谱投影。用 \(\mathrm{E}_{\mathrm{A}}\) 表示 \(p_{\mathrm{A}}\) 的像。它是 E 的闭子空间。

\begin{enumerate}
\def\labelenumi{\alph{enumi})}
\item
  子空间 \(E_{A}\) 是使 x 关于 u 的谱测度的支撑包含在 A 中的那些 \(x \in E\) 所成的集；
\item
  若 A 在 \(\mathbf C\) 中有界，则 \(E_{A}\) 包含在 u 的定义域中；
\item
  对属于 u 的定义域的所有 x，我们有 \(p_{\mathrm{A}}(x) \in \operatorname{dom}(u)\) 且 \(u(p_{\mathrm{A}}(x)) \in \mathrm{E}_{\mathrm{A}}\) ；特别地，\(u(x) \in \mathrm{E}_{\mathrm{A}}\) ，只要 \(x \in \operatorname{dom}(u) \cap \mathrm{E}_{\mathrm{A}}\) 。
\end{enumerate}

对 E 中的 \(x\) ，我们用 \(\mu_x\) 表示 \(x\) 关于 \(u\) 的谱测度。

我们来证明 \(a\) )。设 \(x\) 是 \(\mathrm{E}_{\mathrm{A}}\) 的元素。设 \(z \in \mathbf{C} - \mathbf{A}\) ，并设 U 是 \(z\) 的与 A 不相交的相对紧开邻域。对每个支撑包含在 U 中的函数 \(f \in \mathcal{K}(\mathbf{C})\) ，我们有

\[
\int_ {\mathrm{Sp} (u)} f \mu_ {x} = \langle x | f (u) x \rangle = \langle p _ {\mathrm{A}} (x) | f (u) x \rangle = \int_ {\mathrm{Sp} (u) \cap \mathrm{A}} f \mu_ {x} = 0
\]

（依据公式 (10)）。这表明 \(z\) 不属于 \(\mu_x\) 的支撑。因此，\(\mu_x\) 的支撑包含在 A 中。

反之，设 \(x \in E\) 使 \(\mu_{x}\) 的支撑包含在 A 中。由此可得

\[
\langle x \mid p _ {\mathrm{A}} (x) \rangle = \int_ {\operatorname{Sp} (u) \cap \mathrm{A}} \mu_ {x} = \int_ {\operatorname{Sp} (u)} \mu_ {x} = \langle x \mid x \rangle ,
\]

\((loc. \ cit.) \ \text{因此} \ \|p_{\mathrm{A}}(x)-x\|^{2} = \|p_{\mathrm{A}}(x)\|^{2} - \|x\|^{2} \leqslant 0 \ \text{从而 } p_{\mathrm{A}}(x) = x, \ \text{即 } x \in \mathrm{E}_{\mathrm{A}}.\)

断言 b) 由 a) 及 IV, p. 273 的注得出。

我们来证明 \(c\) )。\(u\) 的定义域是这样的 \(x \in \mathrm{E}\) 所成的集：\(\operatorname{Sp}(u)\) 的恒等函数属于 \(\mathcal{L}^2(\operatorname{Sp}(u), \mu_x)\) （loc. cit.）。由于 \(\mu_{p_{\mathrm{A}}(x)} = \varphi_{\mathrm{A}} \cdot \mu_x\) 对 \(x \in \mathrm{E}\) 成立（IV, p. 275 的命题 5 的推论），我们有 \(p_{\mathrm{A}}(x) \in \operatorname{dom}(u)\) ，只要 \(x \in \operatorname{dom}(u)\) 。

设 \(x \in \operatorname{dom}(u)\) 且 \(y = p_{\operatorname{A}}(x)\) ；我们有 \(y \in \operatorname{dom}(u) \cap \operatorname{E}_{\operatorname{A}}\) ，依据 \(c\) )。\(u(y)\) 关于 \(u\) 的谱测度是 \(|\operatorname{Id}_{\operatorname{Sp}(u)}|^2 \cdot \mu_y\) （loc. cit.）。由于 \(\mu_y\) 的支撑在 A 中，\(\mu_{u(y)}\) 亦然，故 \(u(y) \in \operatorname{E}_{\operatorname{A}}\) ，依据 \(a\) )。

推论. --- 设 \(\lambda \in \mathrm{Sp}(u)\) 。\(p_{\{\lambda\}}\) 的像是 \(u\) 关于 \(\lambda\) 的特征子空间。

设 \(x \in \operatorname{dom}(u)\) 。用 \(\mu\) （相应地，\(\nu\) ）表示元素 \(x\) 关于 \(u\) 的谱测度（相应地，\((x, u(x))\) 关于 \(u\) 的谱测度）。我们有 \(\nu = \operatorname{Id}_{\operatorname{Sp}(u)} \cdot \mu\) （IV, p. 275 的命题 5 的推论）。若 \(u(x) = \lambda x\) ，我们还有 \(\nu = \lambda \mu\) ，由此得等式 \(\operatorname{Id}_{\operatorname{Sp}(u)} \cdot \mu = \lambda \mu\) 。这蕴含 \(\mu\) 的支撑包含在 \(\{\lambda\}\) 中（参见 INT, V, p. 46, § 5, 第 3 目, 推论 2），因此 \(x\) 属于 \(p_{\{\lambda\}}\) 的像（命题 9, a)）。

反之，假设 \(x\) 属于 \(\mathrm{E}_{\lambda}\) ，即 \(p_{\{\lambda\}}\) 的像。那么 \(x\) 属于 \(\operatorname{dom}(u)\) ，且 \(u(x)\) 也属于 \(\mathrm{E}_{\lambda}\) （loc. cit., \(b\) )）。由于 \(\varphi_{\{\lambda\}} \cdot (\operatorname{Id}_{\operatorname{Sp}(u)} - \lambda) = 0\) ，我们有关系式 \(p_{\{\lambda\}} \circ (u - \lambda 1_{\mathrm{E}}) \subset 0\) （IV, p. 275 的命题 5, \(c\) )），由此 \(0 = p_{\{\lambda\}}(u(x)) - \lambda p_{\{\lambda\}}(x) = u(x) - \lambda x\) 。

命题 10. — 设 \(\lambda \in \mathbf{C}\) 。我们有 \(\lambda \in \mathrm{Sp}(u)\) 当且仅当：对 \(\lambda\) 在 \(\mathbf{C}\) 中的每个开邻域 V，u 关于 V 的谱投影 \(p_{\mathrm{V}}\) 非零。

若 \(\lambda \notin \mathrm{Sp}(u)\) ，则存在开邻域 \(V\) ，它包含 \(\lambda\) ，位于 \(\mathbf{C}\) 中，且与 \(\mathrm{Sp}(u)\) 不相交，于是 \(p_{\mathrm{V}} = p_{\varnothing} = 0\) 。

反之，假设存在 \(\lambda\) 在 \(\mathbf{C}\) 中的开邻域 V 使得 \(p_{V} = 0\) 。设 c > 0 使得以 \(\lambda\) 为中心、以 c 为半径的圆盘包含在 V 中。

设 \(x \in \mathrm{dom}(u)\) ，并设 \(\mu_x\) 是 \(x\) 关于 \(u\) 的谱测度。由于 \(\mu_x(\mathrm{V}) = \langle x | p_{\mathrm{V}}(x) \rangle = 0\) ，我们计算出

\[
\begin{array}{r l} & {\int_ {\mathbf {C}} | z - \lambda | ^ {2} d \mu_ {x} (z) = \int_ {\mathbf {C - V}} | z - \lambda | ^ {2} d \mu_ {x} (z)} \\ & {\qquad \geqslant c ^ {2} \int_ {\mathbf {C - V}} d \mu_ {x} (z) = c ^ {2} \int_ {\mathbf {C}} d \mu_ {x} (z) = c ^ {2} \| x \| ^ {2}.} \end{array}
\]

但是，另一方面，我们有

\[
\| u (x) - \lambda x \| ^ {2} = \int_ {\mathbf {C}} | z - \lambda | ^ {2} d \mu_ {x} (z) = \| u ^ {*} (x) - \overline {{\lambda}} x \| ^ {2}
\]

（p. 272 的公式 (8)），由此 \(\|u(x)-\lambda x\|\geqslant c\|x\|\) 且 \(\|u^{*}(x)-\overline{\lambda}x\|\geqslant c\|x\|\) 。于是 \(\lambda\) 属于 u 的预解集（IV, p. 248 的引理 5）。证明完毕。

推论. — 设 A 是 \(\mathbf{C}\) 中的开集，且 \(n \in \mathbf N\) 。若 A 含有 \(\mathrm{Sp}(u)\) 的 n 个元素，则 u 关于 A 的谱投影 \(p_{A}\) 的像的维数至少等于 n。

设 \(\lambda_1, \ldots, \lambda_n\) 是 \(u\) 的谱中属于 A 的互不相同的元素。存在两两不相交的开集族 \((\mathrm{V}_i)_{1 \leqslant i \leqslant n}\) （\(\mathbf{C}\) 中的开集），使得 \(\lambda_i \in \mathrm{V}_i\) 对 \(1 \leqslant i \leqslant n\) 成立。设 B 是诸集合 \(\mathrm{V}_i\) 的并。\(p_{\mathrm{A}}\) 的像包含谱投影 \(p_{\mathrm{B}}\) 的像；此外，由于 \(p_{\mathrm{B}}\) 是诸投影 \(p_{\mathrm{V}_i}\) 之和，且 \(p_{\mathrm{V}_i}\) 的像与 \(p_{\mathrm{V}_j}\) 的像正交（对所有 \(i \neq j\) ），结果由命题 10 得出。

\subsection{埃尔弗--舍斯特兰德公式}

在本节中，E 表示一个复希尔伯特空间。我们将对自伴部分算子的函数演算的某些情形得到一个公式，它直接用所论算子的预解式来表达。

我们给 R （相应地，C）赋以勒贝格测度，记作 \(\mu\) ，并用映射 \((x,y)\mapsto\exp(2i\pi xy)\) 把群 R 与它的对偶群等同起来（参见 II, p. 236 的推论 3）。

对函数 \(f\) ——它在 \(\mathbf{R}^{2}\) 的开集 U 上定义且可微，该开集等同于 \(\mathbf{C}\) ，赋以实坐标 \(x\) 与 \(y\) ——我们记

\[
{\frac {\partial f}{\partial \overline {{{z}}}}} = {\frac {1}{2}} {\Bigl (} {\frac {\partial f}{\partial x}} + i {\frac {\partial f}{\partial y}} {\Bigr)}
\]

（参见 VAR, R2, 8.8.10, p. 24）。

引理 5. --- 设 \(f \in \mathcal{D}(\mathbf{R})\) 。存在函数 \(\tilde{f}\) 属于 \(\mathcal{D}(\mathbf{C})\) ，它与 \(f\) 在 \(\mathbf{R}\) 上重合，且满足

\[
\frac {\partial \widetilde {f}}{\partial \overline {{z}}} (x, 0) = 0\tag{11}
\]

对所有 \(x \in \mathbf{R}\) 成立。这时我们有

\[
\frac {\partial \widetilde {f}}{\partial y} (x, 0) = i f ^ {\prime} (x)\tag{12}
\]

对所有 \(x \in R\) 成立，且存在实数 C ≥ 0 使得

\[
\left| \frac {\partial \widetilde {f}}{\partial \overline {{z}}} (x, y) \right| \leqslant C | y |\tag{13}
\]

对所有 \((x,y)\in \mathbf{R}^2\)

存在 \(\varphi\in\mathcal{D}(\mathbf{R})\) ，其支撑包含在 \([-2,2]\) 中，且在 \([-1,1]\) 上等于 1（IV, p. 196 的引理 1）。置

\[
\widetilde {f} (x, y) = \bigl (f (x) + i y f ^ {\prime} (x) \bigr) \varphi (y)
\]

对 \((x,y)\in\mathbf{R}^{2}\) 。我们有 \(\widetilde{f}\in\mathcal{D}(\mathbf{C})\) ，且 \(\widetilde{f}\) 在 R 上与 f 重合。此外，对任意 \((x,y)\in\mathbf{R}^{2}\) ，可得

\[
\frac {\partial \widetilde {f}}{\partial \overline {{z}}} (x, y) = \frac {1}{2} \Big ((i f (x) - y f ^ {\prime} (x)) \varphi^ {\prime} (y) + i y f ^ {\prime \prime} (x) \varphi (y) \Big).
\]

由于函数 \(\varphi\) 在 0 的一个邻域中等于 1，我们有 \(\varphi'(0) = 0\) ，由此得 (11)。

设 \(\tilde{f}\) 是 \(\mathcal{D}(\mathbf{C})\) 中满足 (11) 的函数。由此即得公式 (12)，而估计式 (13) 则借助中值定理（FVR, I, p. 23, 定理 2）得到。

我们称 \(\tilde{f}\) 是 f 的一个几乎解析延拓。

引理 6. --- 设 \(\varepsilon > 0\) 。函数 \(\sigma_{\varepsilon}\) 定义在 \(\mathbf{R}\) 上，由下式给出

\[
\sigma_ {\varepsilon} (x) = \frac {2 i \varepsilon x}{x ^ {2} + \varepsilon^ {2}}
\]

属于 \(L^{2}(\mathbf{R})\) 。它的傅里叶变换是 \(L^{2}(\mathbf{R})\) 中函数 \(\eta_{\varepsilon}\) 的类，该函数在 0 处取值 0 且满足

\[
\eta_ {\varepsilon} (y) = \frac {2 \pi \varepsilon y}{| y |} e ^ {- 2 \pi \varepsilon | y |}
\]

对所有 \(y \neq 0\) 成立。

由 IV, p. 199 的命题 3，有 \(\sigma_{\varepsilon} \in \mathrm{L}^{2}(\mathbf{R})\) ，且函数 \(\eta_{\varepsilon}\) 的类属于 \(\mathrm{L}^{2}(\mathbf{R}) \cap \mathrm{L}^{1}(\mathbf{R})\) 。对每个 \(x \in \mathbf{R}\) ，我们有

\[
\begin{array}{r} \overline {{\mathcal {F}}} (\eta_ {\varepsilon}) (x) = 2 \pi \varepsilon \int_ {\mathbf {R} _ {+}} e ^ {2 \pi (i x - \varepsilon) y} d y - 2 \pi \varepsilon \int_ {\mathbf {R} _ {-}} e ^ {2 \pi (i x + \varepsilon) y} d y \\ = \frac {\varepsilon}{\varepsilon - i x} - \frac {\varepsilon}{\varepsilon + i x} = \frac {2 i \varepsilon x}{x ^ {2} + \varepsilon^ {2}}, \end{array}
\]

由此 \(\overline{\mathcal{F}}(\eta_{\varepsilon}) = \sigma_{\varepsilon}\) 。于是结果由 \(\mathrm{L}^2(\mathbf{R})\) 中的傅里叶反演公式（II, p. 220 的定理 2 的推论）得出。

引理 7. --- 设 \(f \in \mathcal{D}(\mathbf{R})\) ，并设 \(\widetilde{f} \in \mathcal{D}(\mathbf{C})\) 是 \(f\) 的一个几乎解析延拓。对 \(\varepsilon > 0\) ，定义 \(f_{\varepsilon}\) 于 \(\mathbf{R}\) 上如下

\[
f _ {\varepsilon} (x) = - \frac {1}{2 i \pi} \int_ {\mathbf {R}} \biggl (\frac {\widetilde {f} (y + i \varepsilon)}{y - x + i \varepsilon} - \frac {\widetilde {f} (y - i \varepsilon)}{y - x - i \varepsilon} \biggr) d y.
\]

那么 \(f_{\varepsilon}\) 是连续且有界的，且 \(f_{\varepsilon}\) 收敛于 \(f\) （在 \(\mathcal{C}_b(\mathbf{R})\) 中，当 \(\varepsilon\) 趋于 0 时）。

\(f_{\varepsilon}\) 的连续性是 INT, IV, p. 144, § 4, 第 3 目, 推论 1 的推论，因为 \(\widetilde{f}\) 具有紧支撑。此外，若 \(r\) 使得 \(\widetilde{f}\) 的支撑包含在 \([-r, r] \times \mathbf{R}\) 中，我们有

\[
| f _ {\varepsilon} (x) | \leqslant 2 \| \widetilde {f} \| _ {\infty} \int_ {- r} ^ {r} \frac {1}{\sqrt {(x - y) ^ {2} + \varepsilon^ {2}}} d y \leqslant \frac {4 r}{\varepsilon} \| \widetilde {f} \| _ {\infty},
\]

因此 \(f_{\varepsilon}\) 是有界的。

一阶泰勒展开（FVR, I, p. 30, 命题 3）与公式 (13) 表明，存在 M ≥ 0 及一个函数 \(\varrho_{1}\) ，它取值于 \(R^{2}\) 中，使得

\[
\widetilde {f} (y + i \gamma) = f (y) + i \gamma f ^ {\prime} (y) + \gamma^ {2} \varrho_ {1} (y; \gamma), \text {   and   } | \varrho_ {1} (y; \gamma) | \leqslant M,
\]

对所有 \(y \in \mathbf{R}\) 及 \(\gamma \in \mathbf{R}\) 成立。由于映射 \(y \mapsto \widetilde{f}(y + i\gamma)\) 的支撑包含在 \([-r, r]\) 中（对所有 \(\gamma \in \mathbf{R}\) ），映射 \(y \mapsto \varrho_{1}(y; \gamma)\) 的支撑也包含在 \([-r, r]\) 中，对所有 \(\gamma \in \mathbf{R}\) 。

设 \(g_{\varepsilon}\) 是定义在 \(\mathbf{R}^2\) 上的函数

\[
g _ {\varepsilon} (x, y) = \frac {\widetilde {f} (y + i \varepsilon)}{y - x + i \varepsilon} - \frac {\widetilde {f} (y - i \varepsilon)}{y - x - i \varepsilon}.
\]

对每个 \((x,y)\in \mathbf{R}^2\) 及 \(\varepsilon >0\) ，我们得到

\[
\begin{array}{r l} g _ {\varepsilon} (x, y) = - \frac {2 i \varepsilon}{(x - y) ^ {2} + \varepsilon^ {2}} f (y) + \frac {2 i \varepsilon (y - x)}{(x - y) ^ {2} + \varepsilon^ {2}} f ^ {\prime} (y) \\ & \quad + \varepsilon^ {2} \bigg (\frac {\varrho_ {1} (y ; \varepsilon)}{y - x + i \varepsilon} - \frac {\varrho_ {1} (y ; - \varepsilon)}{y - x - i \varepsilon} \bigg). \end{array}
\]

设 \(x \in \mathbf{R}\) 且 \(\varepsilon > 0\) 。我们有

\[
\begin{array}{c} \left| \varepsilon^ {2} \int_ {\mathbf {R}} \left(\frac {\varrho_ {1} (y ; \varepsilon)}{y - x + i \varepsilon} - \frac {\varrho_ {1} (y ; - \varepsilon)}{y - x - i \varepsilon}\right) d y \right| \leqslant 2 \mathrm{M} \varepsilon^ {2} \int_ {- r} ^ {r} \frac {d y}{\sqrt {(x - y) ^ {2} + \varepsilon^ {2}}} \\ \leqslant 4 \mathrm{Mr} \varepsilon . \end{array}
\]

因此，对每个 \(x \in \mathbf{R}\) ，可得

\[
f _ {\varepsilon} (x) = - \frac {1}{2 i \pi} \int_ {\mathbf {R}} g _ {\varepsilon} (x, y) d y = (f * \delta_ {\varepsilon}) (x) - \frac {1}{2 i \pi} (f ^ {\prime} * \sigma_ {\varepsilon}) (x) + k _ {\varepsilon} (x)
\]

其中 \(\delta_{\varepsilon}\) 与 \(\sigma_{\varepsilon}\) 是 \(\mathbf{R}\) 上由下式定义的函数

\[
\delta_ {\varepsilon} (x) = \frac {1}{\pi} \frac {\varepsilon}{x ^ {2} + \varepsilon^ {2}}, \qquad \sigma_ {\varepsilon} (x) = \frac {2 i \varepsilon x}{x ^ {2} + \varepsilon^ {2}},
\]

且 \(\| k_{\varepsilon}\|_{\infty}\leqslant 2\mathrm{Mr}\varepsilon .\)

函数 \(\mathcal{F}(f') \in \mathcal{S}(\mathbf{R})\) 是可积的（IV, p. 219 的命题 18 及 IV, p. 213 的命题 13）。由引理 6，可得

\[
\begin{array}{r l r} & & {\int_ {\mathbf {R}} | \mathcal {F} (f ^ {\prime}) (y) \mathcal {F} (\sigma_ {\varepsilon}) (y) | d y = 2 \pi \varepsilon \int_ {\mathbf {R}} | \mathcal {F} (f ^ {\prime}) (y) | e ^ {- 2 \pi \varepsilon | y |} d y} \\ & & {\leqslant 2 \pi \varepsilon \int_ {\mathbf {R}} | \mathcal {F} (f ^ {\prime}) (y) | d y,} \end{array}
\]

因此 \(\mathcal{F}(f')\mathcal{F}(\sigma_{\varepsilon})\) 在 \(\mathrm{L}^{1}(\mathbf{R})\) 中收敛于 0，当 \(\varepsilon\) 趋于 0 时。由于 \(f'\) 与 \(\sigma_{\varepsilon}\) 属于 \(\mathrm{L}^{2}(\mathbf{R})\) ，II, p. 223 的命题 14 蕴含 \(f'*\sigma_{\varepsilon}=\overline{\mathcal{F}}(\mathcal{F}(f')\mathcal{F}(\sigma_{\varepsilon}))\) 在 \(\mathcal{C}_{b}(\mathbf{R})\) 中收敛于 0。

对每个 \(\varepsilon > 0\) ，正测度 \(\delta_{\varepsilon} \cdot dx\) 在 \(\mathbf{R}\) 上的总质量为 1（参见 FVR, III, p. 7）。测度 \(\delta_{\varepsilon} \cdot dx\) （对 \(\varepsilon > 0\) ）所成的集，与 \(\mathbf{R}_{+}^{*}\) 中 0 的邻域滤子在该集上诱导的滤子，满足 INT, VIII, p. 137, § 2, 第 7 目的引理 4 的假设。因此，测度 \(\delta_{\varepsilon} \cdot dx\) 在 \(\mathcal{M}^{1}(\mathbf{R})\) 中收敛于点测度 \(\varepsilon_{0}\) ，当 \(\varepsilon\) 趋于 0 时。由此可得 \(f * \delta_{\varepsilon} \to f\) 在 \(\mathcal{C}_{b}(\mathbf{R})\) 中成立（INT, VIII, p. 163, § 4, 第 4 目）。引理得证。

我们用 \(\mu\) 表示 \(\mathbf{C}\) 上的勒贝格测度。

定理 2（埃尔弗--舍斯特兰德）. --- 设 u 是 E 上的自伴部分算子。设 \(f \in \mathcal{D}(\mathbf{R})\) ，且 \(\widetilde{f} \in \mathcal{D}(\mathbf{C})\) 是 \(f\) 的一个几乎解析延拓。设 h 是由下式定义的从 \(\mathbf{C}\) 到 \(\mathcal{L}(\mathrm{E})\) 中的映射

\[
h (\lambda) = \frac {\partial \widetilde {f}}{\partial \overline {{z}}} (\lambda) \mathrm{R} (u, \lambda)
\]

若 \(\lambda\in\mathbf{C} - \mathbf{R}\) ；而 \(h(\lambda)=0\) ，若 \(\lambda\in\mathbf{R}\) 。那么 h 关于 \(\mu\) 在 \(\mathbf{C}\) 上可积，且

\[
f (u) = - \frac {1}{\pi} \int_ {\mathbf {C}} h (\lambda) d \mu (\lambda).
\]

映射 h 是可测的，且其支撑是紧的。由于 u 是自伴的，我们有 \(\|\mathrm{R}(u,\lambda)\| \leqslant |\mathcal{I}(\lambda)|^{-1}\) ，对所有 \(\lambda \in \mathbf{C} - \mathbf{R}\) （IV, p. 248 的命题 17）。因此，依据公式 (13)，映射 h 是有界的，从而它在 \(\mathbf{C}\) 上可积。

设 \(\varepsilon \in \mathbf{R}_{+}^{*}\) 。我们用 \(\mathrm{F}_{\varepsilon}^{+}\) （相应地，\(\mathrm{F}_{\varepsilon}^{-}\) ）表示 \(\mathbf{C}\) 中由这样的 \(\lambda \in \mathbf{C}\) 所成的闭集：\(\mathcal{I}(\lambda) \geqslant \varepsilon\) （相应地，\(\mathcal{I}(\lambda) \leqslant -\varepsilon\) ）。我们有

\[
\int_ {\mathbf {C}} h (\lambda) d \mu (\lambda) = \lim _ {\varepsilon \to 0} \Bigl (\int_ {\mathrm{F} _ {\varepsilon} ^ {+}} h (\lambda) d \mu (\lambda) + \int_ {\mathrm{F} _ {\varepsilon} ^ {-}} h (\lambda) d \mu (\lambda) \Bigr).
\]

设 r > 0 使得 h 的支撑包含在 \(\mathbf{C} = [-r, r]^{2}\) 中。用 \(R_{\varepsilon}^{+}\) 表示 \([-2r, 2r] \times [\varepsilon, 2r]\) 这一 \(\mathbf{C}\) 中的长方体。它是 \(\mathbf{C}\) 的局部多面体子集（VAR, R2, 11.3, p. 48）。它满足下列条件：

\begin{enumerate}
\def\labelenumi{(\roman{enumi})}
\item
  我们有 \(R_{\varepsilon}^{+} \subset 2C \cap F_{\varepsilon}^{+}\) ；
\item
  集合 \(R_{\varepsilon}^{+}\) 包含 \(F_{\varepsilon}^{+}\) 与 h 的支撑的交；
\item
  正则边界 \(\partial R_{\varepsilon}^{+}\) （VAR, R2, 11.3.2, p. 49）包含线段 \(S_{\varepsilon} = [-r, r] + i\varepsilon \subset C\) ；
\item
  有 \(h(\lambda) = 0\) ，只要 \(\lambda \in \partial \mathrm{R}_{\varepsilon}^{+} - \mathrm{S}_{\varepsilon}\) 。
\end{enumerate}

用 \(d\lambda\) （相应地，\(d\overline{\lambda}\) ）表示 \(\mathbf{C}\) 上的、\(\mathbf{C}\) 的恒等映射（相应地，复共轭）的微分 1-形式。用 \(g\) 表示 \(\mathbf{C} - \mathbf{R}\) 上取值于 \(\mathcal{L}(\mathrm{E})\) 的、使得 \(g(\lambda) = \widetilde{f}(\lambda)\mathrm{R}(u,\lambda)\) （对 \(\lambda \in \mathbf{C} - \mathbf{R}\) ）的函数。设 \(\omega = g d\lambda\) ；它是 \(\mathbf{C} - \mathbf{R}\) 上具有紧支撑且取值于 \(\mathcal{L}(\mathrm{E})\) 的 1 次微分形式。由于 \(u\) 的预解式是全纯的（IV, p. 246 的命题 14），我们有

\[
d \omega = \left(\frac {\partial \widetilde {f}}{\partial \overline {{z}}} (\lambda) \mathrm{R} (u, \lambda) + \widetilde {f} (\lambda) \frac {\partial}{\partial \overline {{z}}} \mathrm{R} (u, \lambda)\right) d \overline {{\lambda}} \wedge d \lambda = - h (\lambda) d \lambda \wedge d \overline {{\lambda}}.
\]

与 \(d\omega\) 相伴的向量测度（VAR, R2, 10.4.3, p. 43）是关于勒贝格测度以 -2ih 为密度的测度。对局部多面体集 \(R_{\varepsilon}^{+}\) 及微分形式 \(\omega\) 应用斯托克斯公式（VAR, R2, 11.3.4, p. 49），我们便得到

\[
\frac {i}{2} \int_ {\mathrm{R} _ {\varepsilon} ^ {+}} d \omega = \frac {i}{2} \int_ {\partial \mathrm{R} _ {\varepsilon} ^ {+}} \omega = \frac {i}{2} \int_ {\mathrm{S} _ {\varepsilon}} \omega = \frac {i}{2} \int_ {- r} ^ {r} \widetilde {f} (y + i \varepsilon) \mathrm{R} (u, y + i \varepsilon) d y,
\]

由此

\[
\int_ {\mathrm{F} _ {\varepsilon} ^ {+}} h (\lambda) d \mu (\lambda) = - \frac {i}{2} \int_ {\mathbf {R}} \widetilde {f} (y + i \varepsilon) \mathrm{R} (u, y + i \varepsilon) d y.
\]

对 \(F_{\varepsilon}^{-}\) 作同样的推理，我们得到

\[
\int_ {\mathrm{F} _ {\varepsilon} ^ {-}} h (\lambda) d \mu (\lambda) = - \frac {i}{2} \int_ {\mathbf {R}} \widetilde {f} (y - i \varepsilon) \mathrm{R} (u, y - i \varepsilon) d y,
\]

并且我们得出，\(h\) 在 \(\mathbf{C}\) 上的积分是当 \(\varepsilon\to0\) 时下列量的极限

\[
v _ {\varepsilon} = \frac {i}{2} \int_ {\mathbf {R}} \Bigl (\widetilde {f} (y + i \varepsilon) \mathrm{R} (u, y + i \varepsilon) - \widetilde {f} (y - i \varepsilon) \mathrm{R} (u, y - i \varepsilon) \Bigr) d y.
\]

由命题 7，我们有 \(v_{\varepsilon} = \pi f_{\varepsilon}(u)\) ，其中 \(f_{\varepsilon}\) 是定义在 \(\mathbf{R}\) 上的函数

\[
f _ {\varepsilon} (x) = - \frac {1}{2 i \pi} \int_ {\mathbf {R}} \biggl (\frac {\widetilde {f} (y + i \varepsilon)}{y - x + i \varepsilon} - \frac {\widetilde {f} (y - i \varepsilon)}{y - x - i \varepsilon} \biggr) d y.
\]

由于 \(f_{\varepsilon} \to f\) 当 \(\varepsilon \to 0\) 时在 \(\mathcal{C}_{b}(\mathbf{R})\) 中成立（引理 7），自同态 \(v_{\varepsilon} = \pi f_{\varepsilon}(u)\) 收敛于 \(\pi f(u)\) ，在 \(\mathcal{L}(\mathrm{E})\) 中（IV, p. 275 的命题 5）。定理得证。

\subsection{预解拓扑与函数演算的连续性}

在本节中，E 表示一个复希尔伯特空间。回顾我们用 \(\mathcal{A}(\mathrm{E})\) 表示 E 上自伴部分算子所成的集。我们将把 IV, p. 194 第 8 目的连续性性质推广到 \(\mathcal{A}(\mathrm{E})\) 上。

定义 6. --- 设 T 是 \(\mathcal{L}(\mathrm{E})\) 上的一个局部凸拓扑。\(\mathcal{A}(\mathrm{E})\) 上的 T-预解拓扑是使映射 \(u \mapsto \mathrm{R}(u, \lambda)\) （从 \(\mathcal{A}(\mathrm{E})\) 到赋以拓扑 T 的 \(\mathcal{L}(\mathrm{E})\) 中）对每个 \(\lambda \in \mathbf{C} - \mathbf{R}\) 都连续的最粗拓扑。

命题 11. --- 设 T 是 \(\mathcal{L}(\mathrm{E})\) 上的局部凸拓扑，它比 \(\mathcal{L}(\mathrm{E})\) 的巴拿赫空间拓扑粗。

\begin{enumerate}
\def\labelenumi{\alph{enumi})}
\item
  设 \(f \in \mathcal{C}_0(\mathbf{R})\) 。对每个序列 \((u_n)_{n \in \mathbf{N}}\) （属于 \(\mathcal{A}(\mathrm{E})\) ，按 \(\mathcal{T}\)-预解拓扑收敛于 u），序列 \((f(u_n))_{n \in \mathbf{N}}\) 收敛于 \(f(u)\) ，在空间 \(\mathcal{L}(\mathrm{E})\) 中，该空间赋以拓扑 \(\mathcal{T}\) ；
\item
  假设每个 \(u \in \mathcal{A}(\mathrm{E})\) 对 \(\mathcal{T}\)-预解拓扑都具有可数的基本邻域系。从 \(\mathcal{C}_0(\mathbf{R}) \times \mathcal{A}(\mathrm{E})\) 到空间 \(\mathcal{L}(\mathrm{E})\) （赋以拓扑 \(\mathcal{T}\) ）中的、由 \((f, u) \mapsto f(u)\) 所定义的映射是连续的。
\end{enumerate}

我们来证明 \(a\) )。拓扑 \(\mathcal{T}\) 在 \(\mathcal{L}(\mathrm{E})\) 上是由对拓扑 \(\mathcal{T}\) 连续的半范数所定义的诸拓扑的上确界（EVT, II, p. 26, 推论及其后的注）。因此，只须证明：对每个半范数 \(p\) ，只要它在 \(\mathcal{L}(\mathrm{E})\) 上对拓扑 \(\mathcal{T}\) 连续，序列 \(p(f(u_n) - f(u))\) 就收敛于 0（TG, I, p. 51, 命题 10）。

设 p 是这样一个半范数。用 \(\mathcal{L}(\mathrm{E})_{p}\) 表示赋以半范数 p 的空间 \(\mathcal{L}(\mathrm{E})\) 的分离巴拿赫空间完备化，并用 \(\varpi\) 表示从 \(\mathcal{L}(\mathrm{E})\) 到 \(\mathcal{L}(\mathrm{E})_{p}\) 中的典范映射；它是连续的，且有 \(p(u) = \|\varpi(u)\|_{p}\) ，对所有 \(u \in \mathcal{L}(\mathrm{E})\) ，其中范数是空间 \(\mathcal{L}(\mathrm{E})_{p}\) 的范数。

由于 T 比 \(\mathcal{L}(\mathrm{E})\) 的巴拿赫空间拓扑粗，存在 \(c \geqslant 0\) 使得 \(p(u) \leqslant c \|u\|\) 对所有 \(u \in \mathcal{L}(\mathrm{E})\) 成立。

先假设 \(f \in \mathcal{D}(\mathbf{R})\) 。设 \(\tilde{f}\) 是 \(f\) 的一个几乎解析延拓。用 \(\mu\) 表示 \(\mathbf{C}\) 上的勒贝格测度。由埃尔弗--舍斯特兰德公式（IV, p. 284 的定理 2），有

\[
\varpi (f (u _ {n})) = - \frac {1}{\pi} \int_ {\mathbf {C}} \frac {\partial \widetilde {f}}{\partial \bar {z}} (\lambda) \varpi (\mathrm{R} (u _ {n}, \lambda)) d \mu (\lambda)\tag{14}
\]

对所有 \(n \in \mathbf{N}\) 成立。

设 \(\lambda \in \mathbf{C} - \mathbf{R}\) 。依据假设，预解式序列 \(\mathrm{R}(u_n,\lambda)\) 收敛于 \(\mathrm{R}(u,\lambda)\) ，在空间 \(\mathcal{L}(\mathrm{E})\) 中，该空间赋以拓扑 \(\mathcal{T}\) ，因此序列 \((\varpi (\mathrm{R}(u_n,\lambda)))_n\) 收敛于 \(\varpi (\mathrm{R}(u,\lambda))\) ，在巴拿赫空间 \(\mathcal{L}(\mathrm{E})_p\) 中。

对每个 \(n \in \mathbf{N}\) ，我们有

\[
\left\| \frac {\partial \widetilde {f}}{\partial \bar {z}} (\lambda) \mathrm{R} (u _ {n}, \lambda) \right\| \leqslant \left| \frac {1}{\mathcal{I} (\lambda)} \frac {\partial \widetilde {f}}{\partial \bar {z}} (\lambda) \right|
\]

（IV, p. 248 的命题 17），因此

\[
\left\| \frac {\partial \widetilde {f}}{\partial \overline {{z}}} (\lambda) \varpi (\mathrm{R} (u _ {n}, \lambda)) \right\| _ {p} \leqslant c \left| \frac {1}{\mathcal {I} (\lambda)} \frac {\partial \widetilde {f}}{\partial \overline {{z}}} (\lambda) \right|.
\]

IV, p. 281 的估计式 (13) 表明，这个不等式的右端对 \(\lambda \in \mathbf{C} - \mathbf{R}\) 是有界函数；它在 \(\mathbf{C}\) 上可积，因为 \(\widetilde{f}\) 具有紧支撑。由勒贝格定理（INT, IV, p. 137, § 3, 第 7 目, 定理 6）以及应用于 \(f\) 的埃尔弗--舍斯特兰德公式，我们推出 \(\varpi(f(u_n))\) 收敛于 \(\varpi(f(u))\) ，从而 \(p(f(u_n) - f(u))\) 趋于 0。

现在假设 \(f \in \mathcal{C}_0(\mathbf{R})\) 。设 \(\varepsilon > 0\) 。存在函数 \(f_{\varepsilon}\) ，它属于 \(\mathcal{D}(\mathbf{R})\) ，使得 \(\| f_{\varepsilon} - f \|_{\infty} \leqslant \varepsilon\) （参见 IV, p. 202 的命题 4, a) 及 INT, III, p. 45, § 1, 第 2 目, 命题 3）。于是，依据 IV, p. 275 的命题 5, c)，有 \(\| f(u) - f_{\varepsilon}(u) \| \leqslant \varepsilon\) 及 \(\| f(u_n) - f_{\varepsilon}(u_n) \| \leqslant \varepsilon\) ，对所有 \(n \in \mathbf{N}\) 。因此，可得

\[
p (f (u _ {n}) - f (u)) \leqslant 2 c \varepsilon + p (f _ {\varepsilon} (u _ {n}) - f _ {\varepsilon} (u))
\]

对所有 \(n \in \mathbf{N}\) 成立。由于 \(f_{\varepsilon} \in \mathcal{D}(\mathbf{R})\) ，由前一种情形，序列 \((p(f_{\varepsilon}(u_n) - f_{\varepsilon}(u)))_{n \in \mathbf{N}}\) 收敛于 0，因此 \(p(f(u_n) - f(u))\) 收敛于 0。这就完成了 a) 的证明。

我们来证明断言 \(b\) )。在关于 \(\mathcal{T}\) 的假设下，由 TG, IX, p. 17 的命题 10 及其后的注，以及断言 \(a\) )，可得映射 \(u \mapsto f(u)\) 从 \(\mathcal{A}(\mathrm{E})\) 到 \(\mathcal{L}(\mathrm{E})\) 中是连续的，当 \(f \in \mathcal{C}_0(\mathbf{R})\) 时。由于映射 \(f \mapsto f(u)\) 从 \(\mathcal{C}_0(\mathbf{R})\) 到 \(\mathcal{L}(\mathrm{E})\) 中连续且范数 \(\leqslant 1\) ，对所有 \(u \in \mathcal{A}(\mathrm{E})\) （IV, p. 275 的命题 5, \(c\) )），断言 \(b\) ) 由 TG, X, p. 13, 推论 3 得出。

例. --- 1) 设 \(\mathscr{S}\) 是 E 的有界子集所成的一个集。\(\mathscr{S}\)-拓扑在 \(\mathcal{L}(\mathrm{E})\) 上（EVT, III, p. 13）是比 \(\mathcal{L}(\mathrm{E})\) 的巴拿赫空间拓扑更粗的局部凸拓扑。

\begin{enumerate}
\def\labelenumi{\arabic{enumi})}
\setcounter{enumi}{1}
\tightlist
\item
  设 \(T_{b}\) 是 \(\mathcal{L}(\mathrm{E})\) 的巴拿赫空间拓扑。对 \(T_{b}\)-预解拓扑，每个 \(u \in \mathcal{A}(\mathrm{E})\) 都具有可数的基本邻域系。
\end{enumerate}

事实上，只须证明：对每个 \(\lambda \in \mathbf{C} - \mathbf{R}\) 及每个 \(\varepsilon > 0\) ，存在 \(\lambda' \in \mathbf{Q} + i\mathbf{Q}^*\) 及整数 \(n \geqslant 1\) ，使得每个部分算子 \(v \in \mathcal{A}(\mathrm{E})\) 若满足 \(\| \mathrm{R}(v, \lambda') - \mathrm{R}(u, \lambda')\| < 1/n\) ，则也满足条件 \(\|R(v,\lambda)-R(u,\lambda)\|<\varepsilon\) ；写出

\[
\begin{array}{c} \| \mathrm{R} (v, \lambda) - \mathrm{R} (u, \lambda) \| \leqslant \| \mathrm{R} (v, \lambda) - \mathrm{R} (v, \lambda^ {\prime}) \| \\ + \| \mathrm{R} (v, \lambda^ {\prime}) - \mathrm{R} (u, \lambda^ {\prime}) \| + \| \mathrm{R} (u, \lambda^ {\prime}) - \mathrm{R} (u, \lambda) \|, \end{array}
\]

这个性质由 IV, p. 245 的公式 (8)、IV, p. 248 的命题 17 的诸估计式，以及 \(\mathbf{Q} + i\mathbf{Q}^*\) 在 \(\mathbf{C} - \mathbf{R}\) 中处处稠密这一事实得出。

若把 \(\mathcal{C}_{0}(\mathbf{R})\) 换成 \(\mathcal{C}_{b}(\mathbf{R})\) ，命题的结论一般不成立（IV, p. 366 习题 29, e)）。尽管如此，我们有下列结果。

推论. --- 设 \(\mathcal{T}_{s}\) 是 \(\mathcal{L}(\mathrm{E})\) 上的逐点收敛拓扑。设 \(f \in \mathcal{C}_{b}(\mathbf{R})\) 。对每个序列 \((u_{n})_{n \in \mathbf{N}}\) （属于 \(\mathcal{A}(\mathrm{E})\) ，收敛于 \(u\) ，按 \(\mathcal{T}_{s}\)-预解拓扑），序列 \((f(u_{n}))_{n \in \mathbf{N}}\) 收敛于 \(f(u)\) ，在 \(\mathcal{L}(\mathrm{E})\) 中，该空间赋以拓扑 \(\mathcal{T}_{s}\) 。

设 \((u_{n})\) 是 \(\mathcal{A}(\mathrm{E})\) 中按 \(T_{s}\)-预解拓扑收敛于 u 的序列。设 \(x \in E\) ，并设 \(\varepsilon > 0\) 。

对每个整数 \(N \in \mathbf{N}\) ，设 \(\varphi_{N} \in \mathcal{K}(\mathbf{R})\) 是支撑包含在 \([-(N+1), N+1]\) 中的函数，使得 \(0 \leqslant \varphi_{N} \leqslant 1\) ，且 \(\varphi_{N}(t) = 1\) 对所有 \(t \in [-N, N]\) 成立。当 N 趋于无穷时，函数 \(\varphi_{N}\) 逐点收敛于 1，且满足 \(|\varphi_{N}| \leqslant 1\) ；因此 IV, p. 276 的命题 6 蕴含存在 \(N \in \mathbf{N}\) 使得 \(\|\varphi_{N}(u)x - x\| \leqslant \varepsilon\) 。置 \(f_{N} = f\varphi_{N}\) 。

对每个 \(n\in \mathbf{N}\) ，我们有

\[
\begin{array}{r l} \| f (u _ {n}) x - f (u) x \| \leqslant & \| f (u _ {n}) x - f _ {\mathrm{N}} (u _ {n}) x \| + \\ & \| f _ {\mathrm{N}} (u _ {n}) x - f _ {\mathrm{N}} (u) x \| + \| f _ {\mathrm{N}} (u) x - f (u) x \|. \end{array} \tag {15}
\]

对每个 \(v\in \mathcal{A}(\mathrm{E})\) ，我们有

\[
\begin{array}{r l} & {\| f (v) x - f _ {\mathrm{N}} (v) x \| = \| (f (1 - \varphi_ {\mathrm{N}})) (v) x \|} \\ & {\qquad \leqslant \| f (v) \|   \| x - \varphi_ {\mathrm{N}} (v) x \| \leqslant \| f \| _ {\infty}   \| x - \varphi_ {\mathrm{N}} (v) x \|} \end{array}
\]

（IV, p. 275 的命题 5, c)）。由于 \(\varphi_{\mathrm{N}} \in \mathcal{C}_0(\mathbf{R})\) ，命题 11 的断言 \(a\) ) 应用于 \(\mathcal{T}_s\)-预解拓扑，蕴含 \(\varphi_{\mathrm{N}}(u_n)x\) 收敛于 \(\varphi_{\mathrm{N}}(u)x\) ，当 \(n \to +\infty\) 时。因此，对所有充分大的 \(n\) ，有

\[
\left\| x - \varphi_ {\mathrm{N}} (u _ {n}) x \right\| \leqslant \left\| x - \varphi_ {\mathrm{N}} (u) x \right\| + \varepsilon \| x \| \leqslant (1 + \| x \|) \varepsilon .
\]

于是不等式 (15) 变成

\[
\| f (u _ {n}) x - f (u) x \| \leqslant \| f \| _ {\infty} (2 + \| x \|) \varepsilon + \| f _ {\mathrm{N}} (u _ {n}) x - f _ {\mathrm{N}} (u) x \|
\]

对所有充分大的 n 成立。对函数 \(f_{\mathrm{N}} \in \mathcal{C}_{0}(\mathbf{R})\) 及 \(T_{s}\)-预解拓扑应用 loc. cit. 的断言 a)，即完成证明。

\subsection{极分解}

引理 8. --- 设 E 是分离的拓扑向量空间。设 \((\mathrm{E}_{1},\|\cdot\|_{1})\) 与 \((\mathrm{E}_{2},\|\cdot\|_{2})\) 是 E 的稠密子空间，其上赋以希尔伯特结构，使得 \(E_{1}\) 与 \(E_{2}\) 到 E 中的典范单射都是连续的。设 F 是 \(E_{1}\cap E_{2}\) 的子空间，它对于这些希尔伯特结构在 \(E_{1}\) 与 \(E_{2}\) 中稠密。若 \(\|x\|_{1}=\|x\|_{2}\) 对所有 \(x\in F\) 成立，则 \(E_{1}=E_{2}\) 且 \(\|\cdot\|_{1}=\|\cdot\|_{2}\) 。

设 \(x \in E_{1}\) 。存在 F 中的序列 \((x_{n})_{n \in \mathbf{N}}\) 使得 \(x_{n}\) 在希尔伯特空间 \(E_{1}\) 中收敛于 x。由于对所有整数 n 与 m 有 \(\|x_{n} - x_{m}\|_{2} = \|x_{n} - x_{m}\|_{1}\) ，序列 \((x_{n})\) 是 \(E_{2}\) 中的柯西序列。用 y 表示它的极限。我们有 \(\|y\|_{2} = \lim\|x_{n}\|_{2} = \|x\|_{1}\) 。由于 \(E_{1}\) 与 \(E_{2}\) 到 E 中的典范单射都是连续的，我们有 \(x_{n} \to x\) 在 E 中成立，同样 \(x_{n} \to y\) 也在 E 中成立。于是由此得出 x = y 且 \(\|x\|_{1} = \|x\|_{2}\) 。特别地，\(E_{1} \subset E_{2}\) 。由对称性即得结论。

设 E 是一个复希尔伯特空间。设 u 是 E 上的正自伴部分算子。对每个 \(\alpha \in \mathbf R_{+}\) ，我们置 \(u^{\alpha} = f(u)\) ，其中 f 是连续映射 \(x \mapsto x^{\alpha}\) ，从 \(\mathbf R_{+}\) 到 \(\mathbf R\) 中。这是一个正自伴部分算子（IV, p. 273 的推论 1, a)）。当 \(u \in \mathcal{L}(E)\) 时，这个记号与关于 C*-代数 \(\mathcal{L}(E)\) 的记号相容（参见 I, p. 118 的命题 16）。若 \(\alpha\) 是正整数，则部分算子 \(u^{\alpha}\) 与由复合 \(u \circ \cdots \circ u\) 所定义的部分算子重合（IV, p. 273 的推论 1, b)）。

设 \(\beta \in \mathbf{R}_{+}\) 。我们有 \(u^{\alpha \beta} = (u^{\alpha})^{\beta}\) （IV, p. 274 的推论 4）。特别地，若 \(\alpha > 0\) ，则部分算子 \(u^{1 / \alpha}\) 是 E 上唯一的正自伴部分算子 \(v\) ，使得 \(v^{\alpha} = u\) 。

假设 \(0 \leqslant \alpha \leqslant \beta\) 。这时我们有 \(\operatorname{dom}(u^{\beta}) \subset \operatorname{dom}(u^{\alpha})\) ：事实上，对每个实数 \(x \geqslant 0\) ，有 \(x^{\alpha} \leqslant 1 + x^{\beta}\) ，于是结论由 \(u^{\alpha}\) 的定义域的定义得出（参见 IV, p. 271 的命题 3）。

设 u 是从 E 到复希尔伯特空间 F 中的闭稠定算子。部分算子 \(u^{*}u\) 是自伴且正的（IV, p. 241 的命题 12）。我们记 \(|u| = (u^{*}u)^{1/2}\) 。

命题 12. --- 设 u 是从 E 到复希尔伯特空间 F 中的闭稠定算子。

\begin{enumerate}
\def\labelenumi{\alph{enumi})}
\item
  正自伴部分算子 \(|u|\) 的定义域与 \(u\) 的定义域重合；
\item
  存在唯一的从 E 到 F 中的部分等距线性映射 \(j\) ，使得 \(u = j|u|\) 且 \(\operatorname{Ker}(j) = \operatorname{Ker}(u)\) ；
\item
  设 \(u_{1}\) 是 E 上的正自伴算子，\(j_{1}\) 是从 E 到 F 中的部分等距线性映射，使得 \(u = j_{1}u_{1}\) 且 \(\operatorname{Ker}(j_{1}) = \operatorname{Ker}(u_{1})\) 。这时我们有 \(u_{1} = |u|\) 且 \(j_{1} = j\) 。
\end{enumerate}

\(u^{*}u\) 的定义域包含在 \(\operatorname{dom}(u)\) 与 \(\operatorname{dom}(|u|)\) 中。它在希尔伯特空间 \(E_{u}\) （loc. cit.）与 \(E_{|u|}\) （IV, p. 273 的推论 1, e)）中稠密，且对每个 \(x \in \operatorname{dom}(u^{*}u)\) ，我们有

\[
\begin{array}{c} \| x \| _ {u} ^ {2} = \| x \| ^ {2} + \langle u (x) | u (x) \rangle = \| x \| ^ {2} + \langle x | (u ^ {*} u) (x) \rangle \\ = \| x \| ^ {2} + \langle | u | (x) | | u | (x) \rangle = \| x \| _ {| u |} ^ {2}. \end{array}\tag{16}
\]

因此，希尔伯特空间 \(E_{u}\) 与 \(E_{|u|}\) 相等（引理 8），从而 \(\text{dom}(u) = \text{dom}(|u|)\) ，且 \(\|u(x)\| = \||u|(x)\|\) 对所有 \(x \in \text{dom}(u)\) 成立。

公式 (16) 蕴含 \(\operatorname{Ker}(u) = \operatorname{Ker}(|u|)\) ，并且存在唯一的等距线性映射 \(v\) ，它从 \(\overline{\operatorname{Im}(|u|)}\) 到 \(\overline{\operatorname{Im}(u)}\) 上，且满足 \(v(|u|(x)) = u(x)\) 对所有 \(x \in \operatorname{dom}(|u|)\) 成立。由于 \(|u|\) 是自伴的，我们有 \(\operatorname{Im}(|u|)^{\circ} = \operatorname{Ker}(|u|)\) （IV, p. 236 的命题 7, c)）。因此存在唯一的从 E 到 F 中的部分等距 \(j\) ，它延拓 \(v\) 且在 \(\operatorname{Ker}(|u|) = \operatorname{Ker}(u)\) 上为零。由于 \(\operatorname{E} = \operatorname{Ker}(u) \oplus \overline{\operatorname{Im}(|u|)}\) ，这个映射是使 \(u = j|u|\) 且 \(\operatorname{Ker}(j) = \operatorname{Ker}(u)\) 的唯一的部分等距映射。

我们来证明 c)。我们有 \(u_{1}j_{1}^{*}j_{1}u_{1}\subset(j_{1}u_{1})^{*}j_{1}u_{1}=u^{*}u\) 。线性映射 \(j_{1}^{*}j_{1}\) 是以 \(\mathrm{Ker}(j_{1})=\mathrm{Ker}(u_{1})\) 为核（EVT, V, p. 41, 命题 5 (ii)），从而以 \(\mathrm{Ker}(u_{1})^{\circ}=\overline{\mathrm{Im}(u_{1})}\) 为像（IV, p. 236 的命题 7, c)）的正交投影算子。因此，\(u_{1}^{2}\subset u^{*}u\) ，又由于这两个算子都是自伴的，故 \(u_{1}^{2}=u^{*}u\) 。于是由此得出 \(u_{1}=(u^{*}u)^{1/2}\) ，而 b) 中的唯一性断言最终表明 \(j_{1}=j\) 。

定义 7. --- 设 u 是从 E 到复希尔伯特空间 F 中的稠定闭算子。由命题 12 所确定的偶 \((j, |u|)\) 称为 u 的极分解。

假设 \(u \in \mathcal{L}(\mathrm{E};\mathrm{F})\) 。它在这个定义意义下的极分解与 I, p. 140 的定义 4 中的极分解重合。

\subsection{由正部分埃尔米特形式定义的自伴算子}

在本节中，E 表示一个复希尔伯特空间。

定义 8. --- 设 D 是 E 的向量子空间。E 上以 D 为定义域的部分埃尔米特形式是指一个对应 \(q = (\Gamma, \mathrm{E} \times \mathrm{E}, \mathbf{C})\) ，其定义域为 \(\mathrm{D} \times \mathrm{D}\) ，其图像 \(\Gamma\) 是函数性的，且从 \(\mathrm{D} \times \mathrm{D}\) 到 \(\mathbf{C}\) 中的、由 \(\Gamma\) 所定义的映射是埃尔米特形式。我们记 \(\operatorname{dom}(q) = \mathrm{D}\) 。

若 \(q\) 所定义的埃尔米特形式是正的，则称它为正部分形式。

设 u 是 E 上的自伴部分算子。对 E 的所有元素 x 与 y，我们用 \(\mu_{x,y}\) （相应地，\(\mu_{x}\) ）表示 \((x,y)\) （相应地，x）关于 u 的谱测度。

设 \((j, |u|)\) 是 \(u\) 的极分解（IV, p. 290 的定义 7）。用 D' 表示 \(|u|^{1/2}\) 的定义域。按定义，它是这样的 \(x \in E\) 所成的空间：函数 \(z \mapsto |z|\) 在 \(\mathrm{Sp}(u)\) 上关于测度 \(\mu_x\) 可积。它包含 \(\mathrm{dom}(u)\) 。

设 \((x,y)\in\mathrm{D}^{\prime}\times\mathrm{D}^{\prime}\) 。\(\mathrm{Sp}(u)\) 的恒等映射关于测度 \(\mu_{x,y}\) 可积（IV, p. 271 的命题 4, b)）；我们置

\[
q _ {u} (x, y) = \int_ {\operatorname{Sp} (u)} t d \mu_ {x, y} (t).
\]

若 \(y \in \operatorname{dom}(u)\) ，则由 IV, p. 271 的公式 (6)，我们有 \(q_u(x, y) = \langle x \mid u(y) \rangle\) 。映射 \(q_u\) 是 D' 上的埃尔米特形式。它定义了 E 上的一个部分埃尔米特形式，称为与 u 相伴的部分埃尔米特形式。若算子 u 是正的，则形式 \(q_u\) 是正部分埃尔米特形式。

定义 9. --- 设 \(q\) 是 E 上的正部分埃尔米特形式。我们用 \(\mathrm{E}_q\) 表示赋以下述纯量积的分离预希尔伯特空间 \(\operatorname{dom}(q)\)

\[
(x \mid y) _ {q} = \langle x \mid y \rangle + q (x, y).
\]

我们用 \(\|x\|_{q}\) 表示 \(x \in E_{q}\) 的范数。若 \(E_{q}\) 是希尔伯特空间，则称形式 q 是闭的。

命题 13. --- 设 u 是 E 上的正自伴部分算子。设 \(q_{u}\) 是与 u 相伴的正部分形式。

\begin{enumerate}
\def\labelenumi{\alph{enumi})}
\item
  \(q_{u}\) 的定义域是 \(u^{1/2}\) 的定义域，且对所有 \(\operatorname{dom}(u^{1/2})\) 中的 x 与 y，我们有 \(q_{u}(x,y)=\langle u^{1/2}(x)\mid u^{1/2}(y)\rangle\) ；
\item
  正部分形式 \(q_{u}\) 是闭的。u 的定义域是 \(q_{u}\) 的核。
\end{enumerate}

由于 \(u\) 是正的，我们有 \(|u| = u\) 。因此 \(\mathrm{dom}(|u|^{1/2})\) （即 \(q_u\) 的定义域）与 \(\mathrm{dom}(u^{1/2})\) 重合，并且我们有

\[
\begin{array}{c} q _ {u} (x, y) = \int_ {\operatorname{Sp} (u)} t d \mu_ {x, y} (t) = \int_ {\operatorname{Sp} (u)} t ^ {1 / 2} \cdot t ^ {1 / 2} d \mu_ {x, y} (t) \\ = \langle u ^ {1 / 2} (x) | u ^ {1 / 2} (y) \rangle \end{array}
\]

对 \(x\) 与 \(y\) （属于 \(\mathrm{dom}(q_u)\) ）成立（IV, p. 271 的命题 4, b)）。这就证明了断言 \(a\) )。

此外，预希尔伯特空间 \(E_{q_{u}}\) 这时与预希尔伯特空间 \(E_{u^{1/2}}\) 重合，后者相伴于 \(u^{1/2}\) （IV, p. 230 的定义 4）。由于 \(u^{1/2}\) 是闭部分算子，这个空间是希尔伯特空间（IV, p. 230 的命题 4），且 \(\text{dom}(u)\) 在 \(E_{u^{1/2}}\) 中稠密（依据 IV, p. 273 的推论 1 的断言 \(e\) )）。

设 \(q\) 是 \(E\) 上的部分埃尔米特形式，其定义域 \(D\) 在 \(E\) 中稠密。对 \(y \in D\) ，用 \(\lambda_y\) 表示线性形式 \(x \mapsto q(y, x)\) （定义在 \(D\) 上）。我们用 \(\widetilde{D}\) 表示这样的 \(y \in D\) 所成的集：线性形式 \(\lambda_y\) 在 \(D\) 上连续。

设 \(y \in \widetilde{D}\) 。由于 D 在 E 中稠密，存在 E 上唯一的连续线性形式，它延拓 \(\lambda_{y}\) 。我们仍用 \(\lambda_{y}\) 表示它。存在 E 中唯一的元素 \(u(y)\) ，使得 \(\lambda_{y}(x) = \langle u(y) | x \rangle\) 对所有 \(x \in E\) 成立（EVT, V, p. 15, 定理 3）。映射 \(y \mapsto u(y)\) 从 \(\widetilde{D}\) 到 E 中是线性的；我们称 E 上以 \(\widetilde{D}\) 为定义域的部分算子 u 为表示 q 的部分算子。因此我们有

\[
q (x, y) = \langle x \mid u (y) \rangle
\]

对 \(y \in \operatorname{dom}(u)\) 及 \(x \in D\) 成立。

注记. --- 设 \(q\) 是闭正部分形式，\(q'\) 是 E 上的连续正埃尔米特形式。在 \(\text{dom}(q)\) 上由 \((x, y) \mapsto q(x, y) + q'(x, y)\) 定义的正埃尔米特形式是以与 \(q\) 相同的定义域为定义域的闭正部分埃尔米特形式。我们把它记作 \(q + q'\) 。

依据 EVT, V, p. 16, 推论 2，存在唯一的线性映射 \(u' \in \mathcal{L}(\mathrm{E})\) ，使得 \(q'(x,y) = \langle x \mid u'(y) \rangle\) 对所有 \((x,y) \in \mathrm{E} \times \mathrm{E}\) 成立。自同态 \(u'\) 是正的，而表示闭正部分埃尔米特形式 \(q + q'\) 的部分算子是 \(u + u'\) 。

命题 14. --- 设 q 是 E 上稠定闭正部分形式，u 是表示 q 的部分算子。

\begin{enumerate}
\def\labelenumi{\alph{enumi})}
\item
  u 的定义域在希尔伯特空间 \(E_{q}\) 中稠密；
\item
  部分算子 u 是自伴且正的。
\end{enumerate}

由于 q 是闭的，定义 9 中的预希尔伯特空间 \(E_{q}\) 是希尔伯特空间。

我们来证明部分算子 \(u + 1_{E}\) （以 \(\operatorname{dom}(u)\) 为定义域）是双射的。由于

\[
\langle x \mid (u + 1 _ {\mathrm{E}}) (x) \rangle = q (x, x) + \| x \| ^ {2} \geqslant \| x \| ^ {2}
\]

对所有 \(x \in \operatorname{dom}(u)\) 成立，这个部分算子是单射的。

设 \(y \in E\) 。\(E_{q}\) 上由 \(x \mapsto \langle y | x \rangle\) 所定义的线性形式是连续的，因为 \(\|x\| \leqslant \|x\|_{q}\) 。因此存在 \(y' \in E_{q}\) 使得

\[
\langle y \mid x \rangle = (y ^ {\prime} \mid x) _ {q} = \langle y ^ {\prime} \mid x \rangle + q (y ^ {\prime}, x)
\]

对所有 \(x \in E_{q}\) 成立（EVT, V, p. 15, 定理 3）。按定义，这意味着 \(y'\) 属于部分算子 \(\widetilde{u}\) 的定义域，该算子表示以 dom(q) 为定义域的部分形式 \((x, y) \mapsto (x \mid y)_{q}\) ，且 \(\widetilde{u}(y') = y\) 。由于由上面的注记有 \(\widetilde{u} = u + 1_{E}\) ，我们由此推出部分算子 \(u + 1_{E}\) 也是满射的，从而是双射的。

我们来证明 \(u\) 的定义域在 \(\mathrm{E}_q\) 中稠密。设 \(y \in \mathrm{E}_q\) 与 \(\operatorname{dom}(u)\) 正交（在 \(\mathrm{E}_q\) 中）。存在 \(y' \in \operatorname{dom}(u)\) 使得 \(u(y') + y' = y\) 。这时我们有

\[
0 = (y \mid y ^ {\prime}) _ {q} = \langle y \mid y ^ {\prime} \rangle + q (y, y ^ {\prime}) = \langle y \mid y ^ {\prime} + u (y ^ {\prime}) \rangle = \| y \| ^ {2}
\]

因此 \(y = 0\) 。于是断言 \(a\) ) 得证。

由于 \(\mathrm{dom}(q)\) 在 E 中稠密，且 \(\| x\| \leqslant \| x\| _q\) 对所有 \(x\in \mathrm{dom}(q)\) 成立，断言 \(a\) ) 蕴含 \(\mathrm{dom}(u)\) 在 E 中稠密。

由于形式 \(q\) 是埃尔米特的（相应地，正的），对所有 \(x\) 与 \(y\) （属于 \(\operatorname{dom}(u)\) ），我们有

\[
\langle y \mid u (x) \rangle = q (y, x) = \overline {{q (x , y)}} = \overline {{\langle x \mid u (y) \rangle}} = \langle u (y) \mid x \rangle ,
\]

（相应地，\(\langle x \mid u(x) \rangle = q(x, x) \geqslant 0\) ），从而 u 是对称的（相应地，正的）。最后，由 IV, p. 240 的命题 10 的推论，部分算子 \(u + 1_{E}\) 是自伴的，u 亦然。

定理 3. --- 使 E 上的正自伴算子 u 对应于正部分形式 \(q_u\) 的映射，是 E 上正自伴部分算子所成的集与 E 上稠定闭正部分形式所成的集之间的一个双射。逆双射使正部分形式 q 对应于表示 q 的部分算子。

由命题 13, b) 与命题 14, b)，陈述中所描述的映射是妥善定义的。我们来证明它们互为逆双射。

设 \(u\) 是 E 上的正自伴部分算子，\(q\) 是与 \(u\) 相伴的正部分形式。用 \(v\) 表示表示 \(q\) 的正自伴算子。设 \(y \in \mathrm{dom}(u) \subset \mathrm{dom}(u^{1/2}) = \mathrm{dom}(q)\) 。对每个 \(x \in \mathrm{dom}(q) = \mathrm{dom}(u^{1/2})\) ，我们有

\[
q (y, x) = \langle u ^ {1 / 2} (y) | u ^ {1 / 2} (x) \rangle = \langle u (y) | x \rangle ,
\]

因此 y 属于 v 的定义域且满足 \(v(y) = u(y)\) 。于是部分算子 v 是 u 的一个延拓；由于 u 与 v 都是自伴的，它们相等。

反之，设 q 是稠定闭正部分形式，u 是表示 q 的正自伴部分算子。我们有 \(\operatorname{dom}(u) \subset \operatorname{dom}(u^{1/2})\) 。对 \(\operatorname{dom}(u)\) 中的 x 与 y，我们有

\[
q _ {u} (x, y) = \langle u ^ {1 / 2} (x) | u ^ {1 / 2} (y) \rangle = \langle x | u (y) \rangle = q (x, y)
\]

（依据命题 13, a)）。于是，希尔伯特空间 \(\mathrm{E}_q\) 与 \(\mathrm{E}_{q_u}\) 都包含 \(\operatorname{dom}(u)\) 作为稠密子空间（分别为命题 14, a) 与命题 13, b)），且 \(\|x\|_q = \|x\|_{q_u}\) 对所有 \(x \in \operatorname{dom}(u)\) 成立。由此得出 \(\mathrm{E}_q = \mathrm{E}_{q_u}\) 且 \(q = q_u\) （引理 8）。

推论. --- 设 \(q\) 是 \(E\) 上的闭正部分形式，而 \(u\) 是表示 \(q\) 的正自伴算子。\(q\) 的定义域等于 \((1_{E} + u)^{1/2}\) 的定义域，且我们有

\[
\| x \| _ {q} = \| (1 _ {\mathrm{E}} + u) ^ {1 / 2} x \|
\]

对所有 \(x \in \operatorname{dom}(q)\) 成立。

\(q\) 的定义域等于 \(u^{1/2}\) 的定义域（命题 13, \(a\) )），后者与 \((1_{\mathrm{E}} + u)^{1/2}\) 的定义域重合（参见 IV, p. 271 的命题 3）。对每个 \(x \in \operatorname{dom}(u) \subset \operatorname{dom}(u^{1/2})\) ，我们有

\[
\| (1 _ {\mathrm{E}} + u) ^ {1 / 2} x \| ^ {2} = \langle x | (1 _ {\mathrm{E}} + u) (x) \rangle = \| x \| ^ {2} + \langle x | u (x) \rangle = \| x \| _ {q} ^ {2}.
\]

由于 \(u\) 的定义域在希尔伯特空间 \(\mathrm{E}_q\) 中稠密（命题 14, a)），这个公式由连续性延拓到所有 \(x \in \mathrm{dom}(u^{1/2})\) 。

例. --- 设 u 是 E 上的正部分算子，它未必是闭的。用下式定义以 dom(u) 为定义域的正部分形式 q

\[
q (x, y) = \langle x \mid u (y) \rangle
\]

对 \(x\) 与 \(y\) （属于 \(\mathrm{dom}(u)\) ）成立。设 \(\mathrm{E}_q\) 是 IV, p. 292 的定义 9 中的分离预希尔伯特空间。设 \(\widetilde{\mathrm{E}}_q\) 是 \(\mathrm{E}_q\) 的完备化希尔伯特空间，其纯量积仍记作 \((x,y) \mapsto (x \mid y)_q\) 。由于 \(\iota\) （\(\mathrm{E}_q\) 到 E 中的典范包含）是连续的，它具有唯一的连续延拓，记作 \(\widetilde{\iota}\) （从 \(\widetilde{\mathrm{E}}_q\) 到 E 中）。埃尔米特形式 \(q\) 在 \(\mathrm{E}_q\) 上是连续的，因而也延拓为唯一的连续正埃尔米特形式 \(\widetilde{q}\) ，定义在 \(\widetilde{\mathrm{E}}_q\) 上。

我们来证明线性映射 \(\widetilde{\iota}\) 是单射的。设 \(x\in\mathrm{Ker}(\widetilde{\iota})\) 。考虑序列 \((x_{n})_{n\in\mathbf{N}}\) （属于 \(\mathrm{E}_{q}\) ），它收敛于 \(x\) ，在 \(\widetilde{\mathrm{E}}_{q}\) 中。这时我们有

\[
\lim _ {n \to + \infty} \iota (x _ {n}) = \widetilde {\iota} (x) = 0,
\]

因此序列 \((x_{n})_{n\in\mathbf{N}}\) 在 E 中收敛于 0。设 \(y\in E_{q}\) 。由于 \(\widetilde{q}\) 在 \(\widetilde{E}_{q}\) 上连续，由此可得

\[
\begin{array}{r l} (x \mid y) _ {q} = \lim _ {n \to + \infty} (x _ {n} \mid y) _ {q} = & \lim _ {n \to + \infty} \Big (\langle x _ {n} \mid y \rangle + q (x _ {n}, y) \Big) \\ & = \lim _ {n \to + \infty} \Big (\langle x _ {n} \mid y \rangle + \langle x _ {n} \mid u (y) \rangle \Big) = 0. \end{array}
\]

由于空间 \(E_{q}\) 在 \(\widetilde{E}_{q}\) 中稠密，我们推出 x = 0，此即所要证者。

把 \(\widetilde{E}_{q}\) 与它在 \(\widetilde{\iota}\) 下的像等同于 E 的子空间，我们便把 \(\widetilde{q}\) 解释为延拓 q 的稠定闭正部分形式。与 \(\widetilde{q}\) 相伴的正自伴算子是 u 的一个自伴延拓。它称为 u 的弗里德里希斯扩张。

注. --- 设 \(c \in \mathbf R_{+}\) 。设 q 是稠定闭部分埃尔米特形式，使得对所有属于 q 的定义域的 x 有 \(q(x, x) \geqslant -c \|x\|^2\) 。这意味着由 \(\widetilde{q}(x, y) = q(x, y) + c \langle x \mid y \rangle\) 定义的以 dom(q) 为定义域的闭部分埃尔米特形式是正部分形式。于是由定理 3，这样的形式对应于 E 上使 \(u + c1_{E}\) 为正的自伴部分算子 u，亦即使得 u 以 -c 为下界（IV, p. 237 的定义 7）者。

反之，设 u 是 E 上的对称部分算子，使得

\[
\langle x \mid u (x) \rangle \geqslant - c \| x \| ^ {2}\tag{17}
\]

对所有 \(x \in \operatorname{dom}(u)\) 成立。这时 \(v = u + c1_{E}\) 是 E 上的正部分算子。u 的弗里德里希斯扩张是自伴算子 \(\widetilde{v} - c1_{E}\) ，其中 \(\widetilde{v}\) 是 v 的弗里德里希斯扩张；它是 u 的一个自伴延拓，不依赖于满足 (17) 的实数 c 的选取。
\subsection{正算子谱的变分原理}

本条中，E 是一个非零复希尔伯特空间。

命题 15. --- 设 u 是 E 上的自伴部分算子。假设 u 以 \(c \in \mathbf{R}\) 为下界（IV，第 237 页定义 7）。有

\[
\inf (\mathrm{Sp} (u)) = \inf _ {x \in \mathrm{dom} (u) - \{0 \}} \frac {\langle x \mid u (x) \rangle}{\| x \| ^ {2}} \in [ c, + \infty [.
\]

令 m 为待证等式的右端。设 \(x \in E\)，并令 \(\mu_{x}\) 为 x 关于 u 的谱测度。则有

\[
\begin{array}{r l} \langle x \mid u (x) \rangle & = \int_ {\mathrm{Sp} (u)} t d \mu_ {x} (t) \\ & \geqslant \inf (\mathrm{Sp} (u)) \int_ {\mathrm{Sp} (u)} d \mu_ {x} (t) = \inf (\mathrm{Sp} (u)) \| x \| ^ {2}, \end{array}
\]

（公式 (6)，第 271 页），故 \(\inf(\mathrm{Sp}(u)) \leqslant m\)。反之，部分算子 \(u - m\) 是正的，因此 \(\inf(\mathrm{Sp}(u)) - m = \inf(\mathrm{Sp}(u - m \cdot 1_{\mathrm{E}})) \geqslant 0\)（IV，第 248 页命题 17）。

I 第 139 页的命题 9 对应于本命题在 u 是 \(\mathcal{L}(\mathrm{E})\) 的埃尔米特元素时的特殊情形，此时 u 必然下有界。在此情形下，我们还有

\[
\sup (\operatorname{Sp} (u)) = \sup _ {x \in \operatorname{dom} (u) - \{0 \}} \frac {\langle x \mid u (x) \rangle}{\| x \| ^ {2}}
\]

（同上）；若 u 不属于 \(\mathcal{L}(\mathrm{E})\)，则此上确界为 \(+\infty\)。

定义 10. --- 设 u 是 E 上的正规部分算子。若复数 \(\lambda \in \mathrm{Sp}(u)\) 属于 \(u\) 的敏感谱，当且仅当 \(\lambda\) 是 \(\mathrm{Sp}(u)\) 的孤立点且 \(\lambda\) 是有限重数的特征值。

记 \(\mathrm{Sp}_{\mathrm{s}}(u)\) 为 u 的敏感谱。它在 \(\mathrm{Sp}(u)\) 中的补集称为 u 的本质谱，记作 \(\mathrm{Sp}_{\mathrm{e}}(u)\)。

E 上正规部分算子 u 的本质谱在 C 中是闭的，因为 \(\mathrm{Sp}(u)\) 在 C 中是闭的，而本质谱的补集在 \(\mathrm{Sp}(u)\) 中是开的。u 的敏感谱不一定在 C 中闭（IV，第 369 页习题 36）。

引理 9. --- 设 E 是复希尔伯特空间，u 是 E 上的正规部分算子。设 \(\lambda \in \mathrm{Sp}(u)\)。有 \(\lambda \in \mathrm{Sp}_{\mathrm{s}}(u)\) 当且仅当存在 C 中 \(\lambda\) 的一个开邻域 V，使得由 V 定义的 u 的谱投影 \(p_{\mathrm{V}}\) 具有有限秩。

设 \(\lambda \in \mathrm{Sp}_{\mathrm{s}}(u)\)。存在一个开邻域 \(V\)，它是 \(\lambda\) 在 \(\mathbf{C}\) 中的开邻域，使得 \(\mathrm{Sp}(u) \cap V = \{\lambda\}\)；于是谱投影 \(p_V = p_{\{\lambda\}}\) 具有有限秩（IV，第 278 页命题 9 的推论）。

反之，假设存在 \(\lambda\) 的一个开邻域 V，使得谱投影 \(p_{V}\) 具有有限秩 \(n \in N\)。根据 IV 第 279 页命题 10 的推论，交集 \(V \cap \mathrm{Sp}(u)\) 至多含有 n 个元素，因而 \(\lambda\) 是 \(\mathrm{Sp}(u)\) 的孤立点。因此，根据 IV 第 278 页命题 9 的推论，它是 u 的一个谱重数至多为 n 的特征值。

在本条的其余部分，我们假设 E 是无限维的；当 E 有限维时，下面结果的类似结论由 IV 第 153 页第 4 号推出。

设 u 是 E 上下有界的自伴算子。令 c 为 u 的一个下界；u 的谱包含于 \([c, +\infty[\)（命题 15）。假设 u 的本质谱非空。记 \(\varrho_{e}\) 为 u 的本质谱的下界，于是 \(\varrho_{e} \geqslant c\)，且 \(\varrho_{e}\) 是算子 \(u\) 的谱中的一个元素。记 \(\mathrm{Sp}_{\mathrm{h}}(u) = \mathrm{Sp}(u) \cap [\varrho_{\mathrm{e}}, +\infty[\)；它是 \(\mathrm{Sp}(u)\) 的闭子集，因而也是 \(\mathbf{C}\) 的闭子集，并满足 \(\inf (\mathrm{Sp}_{\mathrm{h}}(u)) = \varrho_{\mathrm{e}}\)；称其为算子 \(u\) 的谱的上端。若算子 \(u\) 的本质谱为空，则记 \(\mathrm{Sp}_{\mathrm{h}}(u) = \varnothing\) 且 \(\varrho_{\mathrm{e}} = +\infty\)。

交集 \(\mathrm{Sp}_{\mathrm{b}}(u) = \mathrm{Sp}(u) \cap [0, \varrho_{\mathrm{e}}[\) 包含于算子 \(u\) 的敏感谱中，因而其元素都是有限重数的孤立特征值；称其为算子 \(u\) 的谱的下端。令 \(\mathrm{E}_{\mathrm{b}}\) 为由对应于特征值 \(\lambda \in \mathrm{Sp}_{\mathrm{b}}(u)\) 的特征空间生成的 E 的闭子空间。根据 IV 第 278 页命题 9 的推论和第 278 页命题 8，空间 \(\mathrm{E}_{\mathrm{b}}\) 是谱投影 \(p_{\mathrm{Sp}_{\mathrm{b}}(u)}\)（由 \(\mathrm{Sp}_{\mathrm{b}}(u)\) 定义）的像。由于 \(\mathrm{Sp}(u)\) 是 \(\mathrm{Sp}_{\mathrm{b}}(u)\) 与 \(\mathrm{Sp}_{\mathrm{h}}(u)\) 的不交并，\(\mathrm{E}_{\mathrm{h}}\) 是 \(\mathrm{E}_{\mathrm{b}}\) 的正交补，也是谱投影 \(p_{\mathrm{Sp}_{\mathrm{h}}(u)}\)（由 \(\mathrm{Sp}_{\mathrm{h}}(u)\) 定义）的像。若 \(\mathrm{Sp}_{\mathrm{b}}(u)\) 有限，则空间 \(\mathrm{E}_{\mathrm{b}}\) 是有限维的；由于假设 E 是无限维的，此时空间 \(\mathrm{E}_{\mathrm{h}}\) 非零。

引理 10. --- 有

\[
\langle x \mid u (x) \rangle \geqslant \varrho_ {\mathrm{e}} \| x \| ^ {2}.\tag{18}
\]

对所有 \(x \in \mathrm{E}_{\mathrm{h}} \cap \operatorname{dom}(u)\)。

若 \(x \in \mathrm{E}_{\mathrm{h}} \cap \mathrm{dom}(u)\)，则 x 的谱测度 \(\mu_{x}\)（即 \(x\) 关于 \(u\) 的谱测度）的支撑包含于区间 \([\varrho_{\mathrm{e}}, +\infty[\)（IV 第 278 页命题 9 \(a\)），故

\[
\langle x \mid u (x) \rangle = \int_ {\mathbf {R}} t d \mu_ {x} (t) \geqslant \varrho_ {\mathrm{e}} \| x \| ^ {2}.
\]

引理 11. --- 假设 \(E_{b}\) 是有限维的。则对每个实数 \(\varepsilon > 0\)，由区间 \([\varrho_{e}, \varrho_{e} + \varepsilon]\) 定义的 u 的谱投影的像都是无限维的。

若 \(E_{b}\) 是有限维的，则 u 的本质谱非空，从而 \(\varrho_{e}\) 有限且属于 \(\mathrm{Sp}_{\mathrm{e}}(u)\)。此外，下谱 \(\mathrm{Sp}_{\mathrm{b}}(u)\) 有限，故存在 \(\delta > 0\)，使得 \([\varrho_{e} - \delta, \varrho_{e} + \delta] \cap \mathrm{Sp}(u) \subset [\varrho_{e}, \varrho_{e} + \delta]\)。断言于是由引理 9 得出。

命题 16. --- 集合 \(\mathrm{Sp}_{\mathrm{b}}(u) \subset [c, \varrho_{\mathrm{e}}[\) 是可数的、离散的，并且在 \([0, \varrho_{\mathrm{e}}[\) 中闭。它是正实数严格递增序列 \((\nu_{k})_{0 \leqslant k < \mathrm{Card}(\mathrm{Sp}_{\mathrm{b}}(u))}\) 的值集；若 \(\mathrm{Sp}_{\mathrm{b}}(u)\) 无限，则序列 \((\nu_{k})\) 收敛到 \(\varrho_{\mathrm{e}}\)。

令 \(\mathrm{T} = \mathrm{Sp}_{\mathrm{b}}(u) \cap [0, \varrho_{\mathrm{e}}[\)。按定义，它是 \([c, \varrho_{\mathrm{e}}[\) 的闭离散子集。对每个整数 \(i \geqslant 1\)，集合 \(\mathrm{Sp}_{\mathrm{b}}(u) \cap [c, \varrho_{\mathrm{e}} - 1/i]\) 是紧的且离散的，因而是有限的。由于 T 是这些集合（\(i \geqslant 1\)）的并，集合 T 是可数的。

若 \(\mathrm{Sp}_{\mathrm{b}}(u)\) 有限，证明即告完成。若 \(\mathrm{Sp}_{\mathrm{b}}(u)\) 无限，则将 III 第 90 页引理 2 应用于 T 在映射 \(\lambda \mapsto \varrho_{e} - \lambda\) 下的像 S，即可得出结论。

对每个整数 \(k\)，满足 \(0 \leqslant k < \text{Card}(\text{Sp}_b(u))\)，记 \(n_k \geqslant 1\) 为 u 的特征值 \(\nu_k\)（算子 \(u\)）的重数。记 \(\overline{\mathbf{N}} = \mathbf{N} \cup \{+\infty\} \subset \overline{\mathbf{R}}\)。令 \(\mathbf{M} \in \overline{\mathbf{N}}\) 为重数 \(n_k\) 之和。这是 \(\mathbf{E}_b\) 的希尔伯特维数；若约定无限希尔伯特维数的空间的维数为 \(+\infty \in \overline{\mathbf{N}}\)，则上述说法成立。

对满足 \(0 \leqslant n < M\) 的每个整数 n，定义 \(\lambda_{n}(u) = \nu_{k}\)，其中 \(k \geqslant 0\) 是满足下式的唯一整数：

\[
n _ {0} + \dots + n _ {k - 1} \leqslant n <   n _ {0} + \dots + n _ {k}.
\]

若 \(n \in \mathbf{N}\) 满足 \(n \geqslant \mathrm{M}\)，则置 \(\lambda_n(u) = \varrho_{\mathrm{e}}\)。此情形只会在 \(\mathrm{Sp}_{\mathrm{b}}(u)\) 有限时发生；由于假设 E 是无限维的，此时 \(\varrho_{\mathrm{e}}\) 有限。

按构造，序列 \((\lambda_{n}(u))_{n\in\mathbf{N}}\) 是递增的；对每个元素 \(\lambda\)（属于 \(\mathrm{Sp}_{\mathrm{b}}(u)\)），满足 \(\lambda_{n}(u)=\lambda\) 的整数 n 的个数等于 \(\lambda\) 作为 u 的特征值的重数。当且仅当 u 的本质谱为空时，序列 \((\lambda_{n}(u))_{n\in\mathbf{N}}\) 趋于 \(+\infty\)。

对每个 \(n \in N\)，记 \(F_{n}\)（分别为 \(F_{n}^{u}\)）为所有维数为 n 的向量子空间 \(F \subset E\) 的集合（分别为所有维数为 n 的向量子空间 \(F \subset \text{dom}(u)\) 的集合）。

设 \(n \in \mathbf{N}\) 满足 \(n < \mathrm{M}\)，且 \(\mathrm{F} \in \mathcal{F}_n^u\)。称 \(\mathrm{F}\) 适应于 \(u\)，如果 \(\mathrm{F}\) 有一个标准正交基 \((f_i)_{0 \leqslant i \leqslant n-1}\)，满足 \(u(f_i) = \lambda_i(u)f_i\)（\(0 \leqslant i \leqslant n-1\)）。

设 \(F \in F_{n}^{u}\) 适应于 u。空间 F 包含于 \(E_{b}\)；它由对应于满足 \(\lambda \leqslant \lambda_{n-1}(u)\) 的特征值的 u 的特征向量生成，并包含对应于满足 \(\lambda < \lambda_{n-1}(u)\) 的特征值的特征空间。此外，对 \(\mathrm{Sp}(u)\) 的每个普遍可测子集 A，空间 F 在由 A 定义的谱投影 \(p_{A}\) 下稳定（IV，第 278 页命题 9，c）。

引理 12. --- 设 \(F \in F_{n}^{u}\) 是适应于 u 的子空间。属于空间 \(F^{\circ} \cap E_{b}\) 的 u 的每个特征向量，其特征值都 \(\geqslant \lambda_{n}(u)\)。

令 \(\ell\) 满足 \(\lambda_{n-1}(u) = \nu_{\ell}\)。设 \(x \in \mathrm{F}^{\circ} \cap \mathrm{E}_{\mathrm{b}}\) 是 \(u\) 关于特征值 \(\lambda\) 的一个特征向量，并令 \(k < \operatorname{Card}(\operatorname{Sp}_{\mathrm{b}}(u))\) 满足 \(\lambda = \nu_{k}\)。

不可能有 \(k < \ell\)，因为 \(\nu_k < \nu_\ell\)，此时空间 F 将包含特征值 \(\nu_k\) 对应的特征空间，这与 \(x\) 正交于 F 矛盾。若 \(k = \ell\)，则 \(x\) 是特征值 \(\lambda_{n-1}(u)\) 的特征向量；由于 \(x \in \mathrm{F}^\circ\)，重数 \(n_k\) 严格大于满足 \(i < n\) 且 \(\lambda_i(u) = \nu_k\) 的整数 i 的个数，从而 \(\lambda_n(u) = \lambda_{n-1}(u)\)。最后，若 \(k > \ell\)，则 \(\lambda = \nu_k > \nu_\ell = \lambda_{n-1}(u)\)，故 \(\lambda \geqslant \lambda_n(u)\)。

对 E 的每个子空间 F，记

\[
\widetilde{r}_{\mathrm{F}}(u) = \inf_{\substack{x\in \operatorname{dom}(u)\cap \mathrm{F}^{\circ}\\ x\neq 0}}\frac{\langle x\mid u(x)\rangle}{\|x\|^{2}},
\]

\[
\widetilde{\mathrm{R}}_{\mathrm{F}}(u) = \sup_{\substack{x\in \operatorname{dom}(u)\cap \mathrm{F}\\ x\neq 0}}\frac{\langle x\mid u(x)\rangle}{\|x\|^ {2}}.
\]

命题 17. --- a) 对每个整数 \(n \in \mathbf{N}\)，有

\[
\lambda_ {n} (u) = \sup _ {\mathrm{F} \in \mathcal {F} _ {n}} \widetilde {r} _ {\mathrm{F}} (u) = \inf _ {\mathrm{F} \in \mathcal {F} _ {n + 1} ^ {u}} \widetilde {\mathrm{R}} _ {\mathrm{F}} (u);
\]

\begin{enumerate}
\def\labelenumi{\alph{enumi})}
\setcounter{enumi}{1}
\item
  对每个整数 \(n < \mathrm{M}\)，以及每个子空间 \(\mathrm{F} \in \mathcal{F}_n^u\)，若其适应于 \(u\)，则有 \(\lambda_n(u) = \widetilde{r}_{\mathrm{F}}(u)\)；
\item
  对每个整数 \(n < \mathrm{M}\)，以及每个子空间 \(\mathrm{F} \in \mathcal{F}_{n+1}^{u}\)，若其适应于 \(u\)，则有 \(\lambda_{n}(u) = \widetilde{\mathrm{R}}_{\mathrm{F}}(u)\)。
\end{enumerate}

存在一个标准正交基 \((e_{n})_{0\leqslant n<M}\)，它是空间 \(E_{b}\) 的基，使得 \(e_{n}\in\operatorname{dom}(u)\)，且 \(u(e_{n})=\lambda_{n}(u)e_{n}\) 对每个满足 \(0\leqslant n<M\) 的 n 都成立。因此，对每个 \(x\in E_{b}\cap\operatorname{dom}(u)\)，有

\[
\langle x \mid u (x) \rangle = \sum_ {0 \leqslant n <   M} \lambda_ {n} (u) | \langle e _ {n} \mid x \rangle | ^ {2}.
\]

 对每个整数 \(n\)，满足 \(1 \leqslant n < M + 1\)，令 \(F_n\) 为维数为 \(n\) 的 \(E_b\) 中由 \((e_0, \ldots, e_{n-1})\) 生成的子空间。我们有 \(F_n \subset \text{dom}(u)\)，且 \(F_n\) 适应于 \(u\)。

设 \(n \in \mathbf{N}\) 且 \(\mathrm{F} \in \mathcal{F}_n\)。我们来证明 \(\widetilde{r}_{\mathrm{F}}(u) \leqslant \lambda_n(u)\)，从而有

\[
\sup _ {\mathrm{F} \in \mathcal {F} _ {n}} \widetilde {r} _ {\mathrm{F}} (u) \leqslant \lambda_ {n} (u).\tag{19}
\]

若 \(0 \leqslant n < M\)，特别是当 M 无限时，E 到 F 的正交投影在 \(\mathrm{F}_{n+1}\) 上的限制不是单射，故存在向量 \(x \neq 0\)，属于 \(F_{n+1}\) 且与 F 正交。由于

\[
\langle x \mid u (x) \rangle = \sum_ {0 \leqslant i \leqslant n} \lambda_ {i} (u) | \langle e _ {i} \mid x \rangle | ^ {2} = \lambda_ {n} (u) | \langle e _ {n} \mid x \rangle | ^ {2} \leqslant \lambda_ {n} (u) \| x \| ^ {2},
\]

我们有 \(\widetilde{r}_{\mathrm{F}}(u) \leqslant \lambda_n(u)\)。

若 \(M \in N\) 且 \(n \geqslant M\)，则对每个实数 \(\varepsilon > 0\)，存在 \(\text{dom}(u)\) 中范数为 1、与 F 正交的 x，使其谱测度 \(\mu_{x}\) 的支撑包含于 \([\varrho_{e}, \varrho_{e} + \varepsilon]\)（IV 第 278 页引理 11 和命题 9，a），于是

\[
\widetilde {r} _ {\mathrm{F}} (u) \leqslant \langle x \mid u (x) \rangle = \int_ {\operatorname{Sp} (u)} t d \mu_ {x} (t) \leqslant \varrho_ {\mathrm{e}} + \varepsilon = \lambda_ {n} (u) + \varepsilon
\]

按定义。由于 \(\varepsilon > 0\) 任意，因而 \(\widetilde{r}_{\mathrm{F}}(u) \leqslant \lambda_n(u)\)。不等式 (19) 得证。

设 \(n \in \mathbf{N}\) 且 \(\mathrm{F} \in \mathcal{F}_{n+1}^{u}\)。我们来证明 \(\widetilde{\mathrm{R}}_{\mathrm{F}}(u) \geqslant \lambda_{n}(u)\)，从而有

\[
\inf _ {\mathrm{F} \in \mathcal {F} _ {n + 1} ^ {u}} \widetilde {\mathrm{R}} _ {\mathrm{F}} (u) \geqslant \lambda_ {n} (u).\tag{20}
\]

若 \(0 \leqslant n < M\)，注意到正交投影到 \(\mathrm{F}_n\) 在 F 上的限制不是单射，故存在向量 \(x \neq 0\)，属于 F 且与 \(\mathrm{F}_n\) 正交。置 \(x_{\mathrm{b}} = p_{\mathrm{Sp}_{\mathrm{b}}(u)}(x)\) 且 \(x_{\mathrm{h}} = p_{\mathrm{Sp}_{\mathrm{h}}(u)}(x)\)。于是 \(x = x_{\mathrm{b}} + x_{\mathrm{h}}\)。元素 \(x_{\mathrm{b}}\) 和 \(x_{\mathrm{h}}\) 属于 \(u\) 的定义域（IV 第 278 页命题 9，c），并且都与 \(\mathrm{F}_n\) 正交。我们有下界

\[
\langle x _ {\mathrm{b}} \mid u (x _ {\mathrm{b}}) \rangle = \sum_ {i \geqslant n} \lambda_ {i} (u) | \langle e _ {i} \mid x _ {\mathrm{b}} \rangle | ^ {2} \geqslant \lambda_ {n} (u) \| x _ {\mathrm{b}} \| ^ {2},
\]

且 \(\langle x_{\mathrm{h}}|u(x_{\mathrm{h}})\rangle \geqslant \varrho_{\mathrm{e}}\| x_{\mathrm{h}}\|^2\)（公式 (18)），于是

\[
\begin{array}{r l} & {\langle x \mid u (x) \rangle = \langle x _ {\mathrm{b}} \mid u (x _ {\mathrm{b}}) \rangle + \langle x _ {\mathrm{h}} \mid u (x _ {\mathrm{h}}) \rangle} \\ & {\qquad \geqslant \lambda_ {n} (u) \| x _ {\mathrm{b}} \| ^ {2} + \varrho_ {\mathrm{e}} \| x _ {\mathrm{h}} \| ^ {2} \geqslant \lambda_ {n} (u) \| x \| ^ {2}.} \end{array}
\]

若 M 有限且 \(n \geqslant M\)，则存在 F 中的 \(x \neq 0\)，与 \(E_{b}\) 正交，从而 \(x \in E_{h}\)，并且

\[
\langle x \mid u (x) \rangle \geqslant \varrho_ {\mathrm{e}} \| x \| ^ {2} = \lambda_ {n} (u) \| x \| ^ {2}
\]

由 (18) 可知。不等式 (20) 得证。

现在证明断言 \(b\)）和 \(c\)），它们说明当 \(n < M\) 时不等式 (19) 和 (20) 都是等式。

证明断言 \(b\)。设 \(\mathrm{F} \in \mathcal{F}_n^u\) 是适应于 \(u\) 的空间。我们有希尔伯特直和

\[
\mathrm{F}^{\circ} = \mathrm{E}_{\mathrm{h}}\oplus \bigoplus_{\substack{\lambda \in \operatorname{Sp}_{\mathrm{b}}(u)\\ \lambda >\lambda_{n - 1}(u)}}\operatorname{Ker}(u - \lambda \cdot 1_{\mathrm{E}})\oplus (\mathrm{F}^{\circ}\cap \operatorname{Ker}(u - \lambda_{n - 1}(u)\cdot 1_{\mathrm{E}})).
\]

设 \(x \in \mathrm{F}^{\circ} - \{0\}\)。写成 \(x = x_{\mathrm{h}} + y + z\)，其中 \(x_{\mathrm{h}} \in \mathrm{E}_{\mathrm{h}}\)，而 \(y\)（分别为 \(z\)）属于上述分解中的第二（分别为第三）个空间。再次利用 (18) 以及每个满足 \(\lambda > \lambda_{n-1}(u)\) 的特征值（关于算子 \(u\)）都 \(\geqslant \lambda_n(u)\) 这一事实，得到

\[
\begin{array}{r l} & {\langle x \mid u (x) \rangle = \langle x _ {\mathrm{h}} \mid u (x _ {\mathrm{h}}) \rangle + \langle y \mid u (y) \rangle + \langle z \mid u (z) \rangle} \\ & {\qquad \geqslant \varrho_ {\mathrm{e}} \| x _ {\mathrm{h}} \| ^ {2} + \lambda_ {n} (u) \| y \| ^ {2} + \lambda_ {n - 1} (u) \| z \| ^ {2}.} \end{array}
\]

若 \(z \neq 0\)，则由引理 12 有 \(\lambda_n(u) = \lambda_{n-1}(u)\)。因此得到 \(\langle x | u(x) \rangle \geqslant \lambda_n(u) \| x \|^2\)，从而 \(\widetilde{r}_{\mathrm{F}}(u) \geqslant \lambda_n(u)\)。结合 (19)，这就证明了断言 \(b\)）。

证明断言 \(c\)。设 \(\mathrm{F} \in \mathcal{F}_{n+1}^{u}\) 是适应于 \(u\) 的子空间。设 \((f_{i})_{0 \leqslant i \leqslant n}\) 是 \(\mathrm{F}\) 的一个标准正交基，满足 \(u(f_{i}) = \lambda_{i}(u)f_{i}\)（\(0 \leqslant i \leqslant n\)）。对每个 \(x \in \mathrm{F}\)，有

\[
\langle x \mid u (x) \rangle = \sum_ {0 \leqslant i \leqslant n} \lambda_ {i} (u) | \langle f _ {i} \mid x \rangle | ^ {2} \leqslant \lambda_ {n} (u) \| x \| ^ {2},
\]

当 \(x = f_{n}\) 时取等号，故 \(\widetilde{\mathrm{R}}_{\mathrm{F}}(u) = \lambda_{n}(u)\)。

最后证明当 \(n \geqslant M\) 时 (19) 和 (20) 是等式。在此情形下，M 有限，故 \(E_{b}\) 包含于 u 的定义域；此外，按定义有 \(\lambda_{n}(u) = \varrho_{\mathrm{e}}\)。

 存在 \(F \in F_{n}\)，包含 \(E_{b}\)。因此，每个元素 \(x \neq 0\)（属于 \(\text{dom}(u)\) 且与 F 正交）都与 \(E_{b}\) 正交，并满足

\[
\frac {\langle x \mid u (x) \rangle}{\| x \| ^ {2}} \geqslant \varrho_ {\mathrm{e}}
\]

（公式 (18)），故 \(\widetilde{r}_{\mathrm{F}}(u) \geqslant \varrho_{\mathrm{e}} = \lambda_n(u)\)，从而

\[
\sup _ {\mathrm{F} \in \mathcal {F} _ {n}} \widetilde {r} _ {\mathrm{F}} (u) \geqslant \lambda_ {n} (u).
\]

令 \(\varepsilon > 0\)。由于 \(\mathrm{E_b}\) 有限维，根据引理 11，存在 E 中的标准正交族 \((x_i)_{i \in \mathrm{I}}\)，使得由 \(\mathrm{E_b}\) 和 \((x_i)_{i \in \mathrm{I}}\) 生成的子空间 F 具有维数 \(n + 1\)，包含于 \(\operatorname{dom}(u)\)，并且满足 \(\langle x_{i} \mid u(x_{i})\rangle \leqslant \varrho_{\mathrm{e}} + \varepsilon\)（对所有 \(i \in I\)）。于是 \(\widetilde{\mathrm{R}}_{\mathrm{F}}(u) \leqslant \varrho_{\mathrm{e}} + \varepsilon\)。由于 \(\varepsilon > 0\) 任意，我们得到

\[
\inf _ {\mathrm{F} \in \mathcal {F} _ {n + 1} ^ {u}} \widetilde {\mathrm{R}} _ {\mathrm{F}} (u) \leqslant \varrho_ {\mathrm{e}} = \lambda_ {n} (u).
\]

命题得证。

\subsection{紧扰动与本质谱}

本条中，E 是一个无限维复希尔伯特空间。

引理 13. --- 设 I 是 E 中的标准正交族。沿 I 的有限子集补集滤子，族 I 弱收敛于 0。

设 \(x \in E\)。根据 Bessel 不等式（EVT，第 V 卷，第 21 页，命题 4）和 TG，第 IV 卷，第 37 页定理 1，族 \((|\langle e_{i} | x\rangle|^{2})_{i\in I}\) 可求和，故 \(\langle e_{i} | x\rangle\) 沿 I 的有限子集补集滤子趋于 0（TG，第 III 卷，第 38 页，命题 1）。

命题 18. --- 设 \(u\) 是 E 上的自伴部分算子，\(\lambda\) 是实数。有 \(\lambda \in \mathrm{Sp}_{\mathrm{e}}(u)\) 当且仅当 E 中存在一个标准正交序列 \((x_{n})_{n \in \mathbf{N}}\)，使得 \(u(x_{n}) - \lambda x_{n}\) 在 E 中趋于 0。

首先假设 \(\lambda \in \mathrm{Sp}_{\mathrm{e}}(u)\)。若 \(u\) 关于 \(\{\lambda\}\) 的谱投影具有无限秩，则其像中的每个标准正交序列 \((x_{n})\) 都满足要求（参见 IV 第 278 页命题 9 的推论）。以下假设情形并非如此。

对每个 \(k \in \mathbf{N}\)，令 \(\mathrm{J}_k\) 为满足下式的 \(t \in [\lambda - 1, \lambda + 1]\) 的集合：

\[
{\frac {1}{k + 2}} <   | t - \lambda | \leqslant {\frac {1}{k + 1}}.
\]

集合 \(J_{k}\) 两两不交。此外，对每个整数 \(K \in N\)，集合族 \((\mathrm{J}_{k})_{k \geqslant K}\) 构成集合

\[
\mathrm{I} _ {\mathrm{K}} = [ \lambda - 1 / (\mathrm{K} + 1), \lambda + 1 / (\mathrm{K} + 1) ] - \{0 \}.
\]

 由于算子 \(u\) 关于 \(\mathrm{I_K} \cup \{0\}\) 的谱投影具有无限秩（IV 第 297 页引理 9），而 u 关于 \(\{\lambda\}\) 的谱投影按假设具有有限秩，根据 IV 第 278 页命题 8，存在严格递增的序列 \((k_n)_{n \in \mathbf{N}}\)（在 \(\mathbf{N}\) 中），使得谱投影 \(p_{n}\)（u 关于 \(J_{k_{n}}\) 的）非零。令 \(x_{n}\) 为 \(p_{n}\) 的像中范数为 1 的向量。序列 \((x_{n})\) 是标准正交的，因为 \(p_{n}\) 的像与 \(p_{m}\) 的像对于 N 中所有 \(n \neq m\) 都正交。

设 \(n \in N\)。记 \(\mu_{n}\) 为 \(x_{n}\) 关于 u 的谱测度；其支撑包含于 \(J_{k_{n}}\)（IV 第 278 页命题 9）。由此

\[
\| u (x _ {n}) - \lambda x _ {n} \| ^ {2} = \int_ {\mathbf {C}} | t - \lambda | ^ {2} d \mu_ {n} (t) \leqslant \frac {1}{k _ {n} ^ {2}},
\]

因此序列 \((x_{n})_{n\in\mathbf{N}}\) 具有所要求的性质。

反之，假设 E 中存在标准正交序列 \((x_{n})_{n\in \mathbf{N}}\)，使得 \(u(x_{n}) - \lambda x_{n}\to 0\)。记 \(\mu_{n}\) 为 \(x_{n}\) 关于 \(u\) 的谱测度。

令 \(\varepsilon > 0\)。记 \(p_{\varepsilon}\) 为 \(u\) 关于开区间 \(\mathrm{I}_{\varepsilon} = ]\lambda - \varepsilon, \lambda + \varepsilon[\) 的谱投影。对每个 \(n \in \mathbf{N}\)，有

\[
\begin{array}{r l} & 1 = \| x _ {n} \| ^ {2} = \mu_ {n} (\mathrm{I} _ {\varepsilon}) + \mu_ {n} (\mathbf {R} - \mathrm{I} _ {\varepsilon}) \\ & \qquad \leqslant \mu_ {n} (\mathrm{I} _ {\varepsilon}) + \frac {1}{\varepsilon^ {2}} \int_ {\mathbf {R} - \mathrm{I} _ {\varepsilon}} | t - \lambda | ^ {2} d \mu_ {n} (t) \\ & \qquad = \| p _ {\varepsilon} (x _ {n}) \| ^ {2} + \frac {1}{\varepsilon^ {2}} \| u (x _ {n}) - \lambda x _ {n} \| ^ {2}. \end{array}
\]

因此，关于序列 \((x_{n})\) 的假设说明 \(p_{\varepsilon}(x_{n})\) 的范数不可能趋于 0。由于标准正交序列 \((x_{n})\) 在 E 中弱收敛于 0（引理 13），投影 \(p_{\varepsilon}\) 不可能是紧的（III 第 6 页命题 6 的推论），因而具有无限秩。由于这对每个 \(\varepsilon > 0\) 都成立，IV 第 297 页引理 9 使我们可以得出 \(\lambda \in \mathrm{Sp}_{\mathrm{e}}(u)\)。

下面的定理是 III 第 93 页定理 3 对自伴部分算子的对应结果。

定理 4. --- 设 u 是 E 上的自伴部分算子。u 的本质谱是集合 Sp(u + v) 的交，其中 v 遍历 E 的紧埃尔米特自同态的集合。

对 E 的每个埃尔米特自同态 v，部分算子 \(u + v\) 都是自伴的，因为 \((u + v)^{*} = u^{*} + v^{*}\)（参见 IV，第 236 页）。

证明若 \(v\) 是紧的，则 \(\mathrm{Sp_e}(u) \subset \mathrm{Sp_e}(u + v)\)。设 \(\lambda \in \mathrm{Sp_e}(u)\)，并令 \((x_n)_{n \in \mathbf{N}}\) 为 E 中满足 \(u(x_n) - \lambda x_n\) 趋于 0 的标准正交序列（命题 18）。序列 \((x_n)\) 在 E 中弱收敛于 0（引理 13）；由于自同态 \(v\) 是紧的，序列 \((v(x_n))_{n \in \mathbf{N}}\) 在 E 中趋于 0（III 第 6 页命题 6 的推论）。因此，序列 \((u+v)(x_{n})-\lambda x_{n}\) 在 E 中趋于 0，命题 18 蕴含 \(\lambda\in\mathrm{Sp}_{\mathrm{e}}(u+v)\)。

因此，算子 \(u\) 的本质谱包含于 \(\mathrm{Sp}(u + v)\) 的交，其中 \(v \in \mathcal{L}^{\mathrm{c}}(\mathrm{E})\) 遍历埃尔米特元素。

反之，设 \(\lambda \in \mathrm{Sp}_{\mathrm{s}}(u)\)。记 \(\mathrm{E}_{\lambda}\) 为算子 \(u\) 关于 \(\lambda\) 的特征空间，记 \(p_{\lambda}\) 为像为 \(\mathrm{E}_{\lambda}\) 的正交投影；它是算子 \(u\) 关于 \(\{\lambda\}\) 的谱投影（IV 第 278 页命题 9 的推论）。根据敏感谱的定义，投影 \(p_{\lambda}\) 具有有限秩，因而是紧的。置 \(w = u + p_{\lambda}\)；它是自伴部分算子。为完成证明，我们验证 \(\lambda\) 属于 \(w\) 的预解集。

有 \(\mathrm{E}_{\lambda} \subset \operatorname{dom}(u)\)，且空间 \(\mathrm{E}_{\lambda}\) 和 \(\operatorname{dom}(u) \cap \mathrm{E}_{\lambda}^{\circ}\) 在 \(u\) 下稳定（IV 第 278 页命题 9）。

设 \(x \in \text{dom}(u)\)。写成 \(x = p_{\lambda}(x) + y\)，其中 \(y \in \text{dom}(u) \cap \text{E}_{\lambda}^{\circ}\)。由上文可得

\[
\left\| w (x) - \lambda x \right\| ^ {2} = \left\| p _ {\lambda} (x) \right\| ^ {2} + \left\| w (y) - \lambda y \right\| ^ {2}.\tag{21}
\]

令 V 为不与 u 的谱相交的 \(\lambda\) 的一个开邻域，并令 c > 0，使得以 \(\lambda\) 为中心、半径为 c 的圆盘包含于 V。令 \(\mu_{y}\) 为 y 关于 u 的谱测度。由于 y 与 \(E_{\lambda}\) 正交，\(\mu_{y}\) 的支撑不与 V 相交（同上）。于是有

\[
\| w (y) - \lambda y \| ^ {2} = \| u (y) - \lambda y \| ^ {2} = \int_ {\mathbf {C}} | t - \lambda | ^ {2} d \mu_ {y} (t) \geqslant c ^ {2} \| y \| ^ {2}.\tag{22}
\]

由 (21) 和 (22)，我们得到 \(\|w(x)-\lambda x\|^{2}\geqslant\inf(c^{2},1)\|x\|^{2}\)，因为 w 是自伴的且 \(\lambda\in R\)。结论于是由 IV 第 248 页引理 5 得出。

推论。 --- 设 \(u\) 是 E 上的自伴部分算子，\(v\) 是 E 的紧埃尔米特自同态。有 \(\mathrm{Sp}_{\mathrm{e}}(u + v) = \mathrm{Sp}_{\mathrm{e}}(u)\)。

这立即由定理得出。

\subsection{扰动}

本条中，E 是一个复希尔伯特空间。

若 \(u\) 是 E 上的自伴部分算子且 \(v \in \mathcal{L}(\mathrm{E})\) 是埃尔米特自同态，则 \(u + v\) 是自伴的（参见 IV，第 236 页）。当 \(v\) 是对称部分算子时，一般并非如此（IV，第 347 页习题 9）。然而，我们将在本条中得到这类的正面结果。

定义 11. --- 设 u 是 E 上的部分算子。如果 dom(u) ⊂ dom(v)，且 v 通过限制到该子空间定义了从 \(E_{u}\) 到 E 的连续线性映射，则称 E 上的部分算子 v 相对于 u 有界。

设 u 是 E 上的部分算子，v 是相对于 u 有界的部分算子。存在实数 m，使得

\[
\| v (x) \| \leqslant m (\| x \| + \| u (x) \|)
\]

对所有 \(x \in \text{dom}(u)\) 都成立。v 关于 u 的相对范数记为 \(\|v\|_{u}\)，它是满足以下条件的实数 \(a \geqslant 0\) 的下界：存在实数 b，使得

\[
\| v (x) \| \leqslant a \| u (x) \| + b \| x \|
\]

对所有 \(x \in \text{dom}(u)\) 都成立。因此，v 的相对范数小于 v 限制到空间 \(E_{u}\) 上所得映射的范数。

注。 --- 设 u 是 E 上的部分算子。每个自同态 \(v \in \mathcal{L}(\mathrm{E})\) 都相对于 u 有界，且其相对范数为零，因为有 \(\|v(x)\| \leqslant \|v\|\|x\|\) 对所有 \(x \in \operatorname{dom}(u)\) 都成立，这使我们可以在上面的不等式中取 a = 0。

定理 5（加藤–雷利希）。— 设 u 是 E 上的自伴部分算子，v 是相对于 u 有界的对称部分算子。若相对范数 \(\|v\|_{u}\) 为 \(< 1\)，则定义域为 dom(u) 的部分算子 \(u + v\) 是自伴的。

由于 \(\| v \|_u < 1\)，按定义存在正实数 \(a\) 和 \(b\)，使得 \(a < 1\)，且 \(\| v(x) \| \leqslant a \| u(x) \| + b \| x \|\) 对所有 \(x \in \text{dom}(u)\) 都成立。

令 \(t \in R^{*}\)。置 \(w_{t} = v \circ \mathrm{R}(u, it)\)。有 \(\operatorname{dom}(w_{t}) = E\)，因为 \(\mathrm{R}(u, it)\) 的像包含于 u 的定义域，而按假设 u 的定义域包含于 v 的定义域。设 \(x \in \operatorname{dom}(u)\)。由于 u 是自伴的，有

\[
\begin{array}{c} \| (i t - u) x \| ^ {2} = \| i t x \| ^ {2} + \| u (x) \| ^ {2} - i t \langle x | u (x) \rangle + i t \langle u (x) | x \rangle \\ = | t | ^ {2} \| x \| ^ {2} + \| u (x) \| ^ {2}, \end{array}
\]

从而有不等式 \(\|u(x)\|\leqslant\|(it-u)x\|\) 和 \(\|x\|\leqslant|t|^{-1}\|(it-u)x\|\)。于是得到

\[
\| v (x) \| \leqslant (a + b | t | ^ {- 1}) \| (i t - u) x \|.
\]

特别地，取 \(x = \mathrm{R}(u, it)y\)，其中 \(y \in \mathrm{E}\)；得到

\[
\| w _ {t} (y) \| \leqslant (a + b | t | ^ {- 1}) \| y \|.
\]

由此得到 \(w_{t}\in\mathcal{L}(\mathrm{E})\)，且 \(\|w_{t}\|\leqslant a+|t|^{-1}b\)。

由于 \(a < 1\)，存在 \(t \in \mathbf{R}_{+}^{*}\)，使得 \(a + b|t|^{-1} < 1\)。于是 \(\|w_{t}\| < 1\) 且 \(\|w_{-t}\| < 1\)；因此 E 的自同态 \(1_{\mathrm{E}} - w_{t}\) 和 \(1_{\mathrm{E}} - w_{-t}\) 都可逆（I，第 22 页命题 2）。由于有公式 \((1_{\mathrm{E}} - w_{t}) \circ (it - u) = it - (u + v)\)，这说明 \(u + v - it\) 是满射；同样，部分算子 \(u + v + it\) 是满射。因此 \(u + v\) 是自伴的（IV，第 240 页命题 11）。

\subsection{具有紧预解式的算子}

本条中，E 表示一个无限维复希尔伯特空间。

命题 19. --- 设 u 是 E 上的自伴部分算子。以下条件等价：

\begin{enumerate}
\def\labelenumi{(\roman{enumi})}
\item
  存在 E 的一个标准正交基 \(\mathrm{B} = (e_{j})_{j \in \mathrm{J}}\)，使得 u 在基 B 中是对角的，并且 u 的特征值族的绝对值沿 J 的有限子集补集滤子趋于无穷；
\item
  对每个 \(\lambda\)（属于 \(u\) 的预解集），预解式 \(\mathrm{R}(u,\lambda)\) 都是紧的；
\item
  存在复数 \(\lambda\) 属于 \(u\) 的预解集，使得预解式 \(\mathrm{R}(u,\lambda)\) 是紧的；
\item
  u 的谱与 u 的敏感谱重合。
\end{enumerate}

假设 (i) 成立，并令 \((\lambda_{j})_{j\in\mathrm{J}}\) 为 u 的特征值族。令 \(\mu\) 为 J 上的计数测度。u 的谱是该族的 \(\mu\)-本质像（IV，第 252 页命题 22），即该族 \((\lambda_{j})\) 的值集。对每个 \(\lambda\notin\mathrm{Sp}(u)\)，预解式 \(\mathrm{R}(u,\lambda)\) 在基 B 中是对角的，其特征值族为 \(((\lambda-\lambda_{j})^{-1})_{j\in\mathrm{J}}\)（同上）。由于该族趋于 0，自同态 \(\mathrm{R}(u,\lambda)\) 是紧的（IV，第 148 页命题 2(iii)）。这证明了 (i) 蕴含 (ii)。

由于 \(u\) 的预解集非空（参见 IV，第 248 页命题 17），根据 IV 第 245 页公式 (8) 和 III 第 5 页命题 3，性质 (ii) 与 (iii) 等价。

由于 \(\mathrm{R}(u,\lambda)\) 在 \(\lambda\notin\mathrm{Sp}(u)\) 时是正规的（IV，第 247 页命题 16），根据 IV 第 247 页命题 15 和 III 第 90 页命题 5，条件 (iii) 蕴含 (iv)。

最后证明 (iv) 蕴含 (i)。令 \(\mathcal{O}\) 为由 \(u\) 的特征向量组成的 E 的标准正交子集的集合。按包含关系排序的集合 \(\mathcal{O}\) 具有有限特征（E，III，第 34 页，定义 2），因为 O 属于 \(\mathcal{O}\) 当且仅当由 O 中至多两个元素组成的集合属于 \(\mathcal{O}\)。根据 E，III，第 35 页定理 1，存在极大元 O，且 O 属于 \(\mathcal{O}\)。令 F 为由 O 生成的 E 的闭子空间。对 \(e \in O\)，存在唯一的 \(\lambda(e) \in \mathbf{R}\)，使得 \(u(e) = \lambda(e)e\)（IV，第 248 页命题 17）。

 映射 \(\lambda\) 从 O 到 \(\mathbf{R}\) 的值集与算子 \(u\) 的谱重合。事实上，一方面该集合包含于算子 \(u\) 的谱；另一方面，若存在 \(\lambda_0 \in \mathrm{Sp}(u)\) 不是 \(\lambda\) 的值，则根据假设 \(\lambda_0\) 是算子 \(u\) 的特征值。存在一个范数为 1 的向量 \(e \in \mathrm{dom}(u)\)，使得 \(u(e) = \lambda_0 e\) 且 \(e \notin \mathrm{O}\)；子集 \(\mathrm{O} \cup \{e\}\) 是标准正交的（因为算子 \(u\) 的对应于不同特征值的特征空间彼此正交，这是 IV 第 248 页命题 17 的推论中的断言 \(b\)），故它属于 \(\mathcal{O}\)，这与 O 的极大性矛盾。

假设 \(F \neq E\)。则 \(F^{\circ}\) 非零。E 的自同态 \(\mathrm{R}(u,i)\) 是正规的（IV，第 247 页命题 16）。它使 E 的子空间 \(F^{\circ}\) 稳定（I，第 135 页引理 4）。令 v 为 \(F^{\circ}\) 的自同态，由 \(\mathrm{R}(u,i)\) 限制到子空间得到。它是 \(F^{\circ}\) 的连续正规自同态（同上），且其谱包含于 \(\mathrm{R}(u,i)\) 的谱中。由于 \(F^{\circ}\) 非零，v 的谱非空（I，第 26 页推论 1）。此外，v 的谱不只含 0，因为正规自同态 v 非零（I，第 110 页例 1）。取 \(s \in \mathrm{Sp}(v) - \{0\}\)。由于 s 属于 \(\mathrm{R}(u,i)\) 的谱，存在 \(e \in O\)，使得 \(s = (i - \lambda(e))^{-1}\)，且 s 是 \(\mathrm{R}(u,i)\) 的特征值（IV，第 247 页命题 15，a）。由于 s 非零，根据假设它是 \(\mathrm{Sp}(\mathrm{R}(u,i))\) 的孤立点，因而也是 \(\mathrm{Sp}(v)\) 的孤立点。因此，s 是 v 的特征值（I，第 134 页命题 5，c）。令 \(e \in F^{\circ}\) 为 v 的范数为 1 的特征向量；它也是 u 的特征向量（IV，第 247 页命题 15，b），而集合 \(O \cup \{e\}\) 与 O 在 O 中极大这一事实矛盾。因此 F = E。

因此，O 中元素组成的族是 E 的一个由 u 的特征向量组成的标准正交基。u 的谱在 R 中是闭的，敏感谱是离散的；由于这两个集合重合，对每个 R > 0，满足 \(\lambda \in \mathrm{Sp}(u)\) 且 \(|\lambda| \leqslant R\) 的元素集合是紧的，因而是有限的。因此，u 的特征值绝对值族沿 O 的有限子集补集滤子趋于无穷。故 (iv) 蕴含 (i)。

定义 12. --- 设 u 是 E 上的自伴部分算子。如果命题 19 的等价条件成立，则称 u 具有紧预解式。

命题 20. --- 设 u 是 E 上的自伴部分算子。u 具有紧预解式，当且仅当希尔伯特空间 \(E_{u}\) 到 E 的典范嵌入 j 是紧的。

假设 u 具有紧预解式。存在 \(\lambda \notin \operatorname{Sp}(u)\)，使得 \(\mathrm{R}(u, \lambda)\) 是紧的（命题 19(iii)）。令 B 为希尔伯特空间 \(E_{u}\) 的单位球。由于 u 是从 \(E_{u}\) 到 E 的连续线性映射，E 的子集 \(\mathrm{B}' = (\lambda 1_{\mathrm{E}} - u)(\mathrm{B})\) 有界。由于 \(\mathrm{R}(u, \lambda)\) 是紧的，子集 \(\mathrm{B} = \mathrm{R}(u, \lambda)(\mathrm{B}')\) 在 E 中相对紧（III，第 2 页注 1）。这证明了 j 是紧的。

 反之，假设线性映射 \(j\) 是紧的。复数 \(i\) 属于 \(u\) 的预解集（IV，第 248 页命题 17）。有 \(u \circ \mathrm{R}(u, i) = -1_{\mathrm{E}} + i\mathrm{R}(u, i)\)。令 B 为 E 的单位球。对每个 \(x \in \mathrm{B}\)，有

\[
\| u \circ \mathrm{R} (u, i) (x) \| = \| - x + i \mathrm{R} (u, i) (x) \| \leqslant 1 + \| \mathrm{R} (u, i) \|,
\]

从而

\[
\begin{array}{r l} & {\| \mathrm{R} (u, i) x \| _ {u} ^ {2} = \| \mathrm{R} (u, i) x \| ^ {2} + \| u \circ \mathrm{R} (u, i) x \| ^ {2}} \\ & {\qquad \leqslant \| \mathrm{R} (u, i) \| ^ {2} + (1 + \| \mathrm{R} (u, i) \|) ^ {2}.} \end{array}
\]

因此，B 在 \(R(u,i)\) 下的像 C 在 \(E_{u}\) 中有界。由于按假设 j 是紧的，集合 \(C = j(C)\) 在 E 中相对紧（III，同上）。所以预解式 \(R(u,i)\) 是紧的，且 u 具有紧预解式（命题 19(iii)）。

推论。 --- 设 q 是 E 上的闭正部分形式，u 是表示 q 的正自伴算子。以下条件等价：

\begin{enumerate}
\def\labelenumi{(\roman{enumi})}
\item
  部分算子 u 具有紧预解式；
\item
  E 的正自同态 \(\sqrt{(1_{\mathrm{E}} + u)^{-1}} = (1_{\mathrm{E}} + u)^{-1 / 2}\) 是紧的；
\item
  希尔伯特空间 \(E_{q}\) 到 E 的典范嵌入 j（IV，第 292 页定义 9）是紧的。
\end{enumerate}

由于 u 是正的，实数 -1 属于 u 的预解集（IV，第 248 页命题 17），故正自同态 \(v = \sqrt{(1_{\mathrm{E}} + u)^{-1}}\) 有定义，并且根据函数演算有 \(v = (1_{\mathrm{E}} + u)^{-1/2}\)。

自同态 \(v\) 是紧的，当且仅当 \(v^2 = (1_{\mathrm{E}} + u)^{-1}\) 是紧的（III，第 91 页命题 6），也就是当且仅当算子 \(u\) 具有紧预解式（命题 19(iii)）。这证明了条件 (i) 和 (ii) 等价。

假设自同态 \(v\) 是紧的。令 B 为希尔伯特空间 \(\mathrm{E}_q\) 的单位球。由于 \((1_{\mathrm{E}} + u)^{1/2}\) 是从 \(\mathrm{E}_q\) 到 E 的连续线性映射（IV，第 294 页定理 3 的推论），E 的子集 \(\mathrm{B}' = (1_{\mathrm{E}} + u)^{1/2}(\mathrm{B})\) 有界，故子集 \(j(\mathrm{B}) = \mathrm{B} = v(\mathrm{B}')\) 在 E 中相对紧（III，第 2 页注 1）。这证明了 \(j\) 是紧的。因此 (ii) 蕴含 (iii)。

线性映射 \(\widetilde{v}\colon x\mapsto(1_{\mathrm{E}}+u)^{-1/2}(x)\) 从 E 到 \(\mathrm{E}_{q}\) 定义良好且为等距映射（IV 第 294 页定理 3 的推论）。由于 \(v=j\circ\widetilde{v}\)，条件 (iii) 蕴含 (ii)（III，第 5 页命题 3）。

例。 --- 设 \(n \in N\)。令 U 为 \(R^{n}\) 的一个开子集。赋予 U 以勒贝格测度，记为 \(\mu\)。令 \(\Delta\) 为 U 上的数量微分算子 \(-\sum_{i=1}^{n}\partial_{i}^{2}\)。部分算子 \(\Delta_{-}\) 的定义域为 \(\mathcal{D}(U)\)，由 \(\varphi \mapsto \Delta(\varphi)\) 定义；它是可闭的（IV，第 242 页命题 13）且是对称的（IV，第 243 页）。它是正的，因为对每个 \(\varphi \in \mathcal{D}(U)\)，有

\[
\begin{array}{r l} & {\langle \varphi | \Delta_ {-} (\varphi) \rangle = \int_ {\mathrm{U}} \overline {{\varphi}} \Delta (\varphi) d \mu = - \sum_ {i = 1} ^ {n} \int_ {\mathrm{U}} \overline {{\varphi}} \partial_ {i} ^ {2} \varphi d \mu} \\ & {\qquad = \sum_ {i = 1} ^ {n} \int_ {\mathrm{U}} \overline {{\partial_ {i} \varphi}} \partial_ {i} \varphi d \mu = \int_ {\mathrm{U}} \sum_ {i = 1} ^ {n} | \partial_ {i} \varphi | ^ {2} d \mu \geqslant 0.} \end{array}
\]

记 \(\Delta_{D}\) 为正对称部分算子 \(\Delta_{-}\) 的弗里德里希斯扩张（IV，第 295 页，例）；它是 U 上的拉普拉斯算子，称为 U 上的狄利克雷拉普拉斯算子。

令 q 为与 \(\Delta_{D}\) 相关的正部分形式。q 的定义域是 \(\mathcal{D}(U)\) 关于正定埃尔米特形式的希尔伯特空间完备化

该形式定义为

\[
(\varphi_ {1}, \varphi_ {2}) \mapsto \int_ {\mathrm{U}} \overline {{\varphi}} _ {1} \varphi_ {2} + \sum_ {i = 1} ^ {n} \int_ {\mathrm{U}} \overline {{\partial_ {i} \varphi_ {1}}} \partial_ {i} \varphi_ {2}
\]

对所有 \((\varphi_{1},\varphi_{2})\in\mathcal{D}(\mathrm{U})\times\mathcal{D}(\mathrm{U})\) 都成立。换言之，\(q\) 的定义域是索伯列夫空间 \(\mathrm{H}_{0}^{1}(\mathrm{U})\)（IV 第 221 页第 14 号）。

*假设 U 有界。典范嵌入 \(\mathrm{H}_0^1 (\mathrm{U})\) 到 \(\mathrm{L}^2 (\mathrm{U})\) 是紧的；因此 U 上的狄利克雷拉普拉斯算子是具有紧预解式的算子（命题 20 的推论）。由于希尔伯特空间 \(\mathrm{H}_0^1 (\mathrm{U})\) 是可数型的（IV，第 222 页命题 20），且 \(\Delta_{\mathrm{D}}\) 的像是无限维的，存在一个递增的实数序列 \((\lambda_n)_{n\geqslant 0}\)，其趋于 \(+\infty\)，以及一个标准正交基 \((f_n)_{n\in \mathbf{N}}\)，它是 \(\mathrm{L}^2 (\mathrm{U})\) 的基，其元素属于 \(\Delta_{\mathrm{D}}\) 的定义域，并满足 \(\Delta_{\mathrm{D}}(f_n) = \lambda_n f_n\) 对所有 \(n\in \mathbf{N}\) 都成立。可以证明（“外尔定律”），当 T 趋于 \(+\infty\) 时，有

\[
\sum_{\substack{n\in \mathbf{N}\\ \lambda_{n}\leqslant \mathrm{T}}}1\sim \frac{c_{n}}{(2\pi)^{n}} m\mathrm{T}^{n / 2}
\]

其中 \(c_{n} = \pi^{n/2}/\Gamma(1 + n/2)\) 是 \(R^{n}\) 中单位球的体积（INT，第 V 卷，第 101 页，第 8 节第 7 号），而 \(m > 0\) 是 U 的勒贝格测度。*

\specialchapter{习题}

\exercisesection{}

\begin{enumerate}
\def\labelenumi{\arabic{enumi})}
\tightlist
\item
  设 \(E = \ell_{\mathbf{C}}^{2}(\mathbf{N})\)，令 \(u \in \mathcal{L}(\mathrm{E})\) 为满足 \(u(x_{0}, \ldots, x_{n}, \ldots) = (0, x_{0}, \ldots, x_{n}, \ldots)\) 的自同态，其中该关系对 E 中所有 \((x_{n})_{n \in \mathbf{N}}\) 都成立。令 \((e_{n})\) 为 E 的标准正交基。
\end{enumerate}

\begin{enumerate}
\def\labelenumi{\alph{enumi})}
\item
  对每个 \(v \in \mathcal{L}(\mathrm{E})\) 满足 \(\|v\| < 1/2\)，向量 \(e_{0}\) 不在 \(u + v\) 的像的闭包中。
\item
  \(\mathcal{L}(\mathrm{E})\) 中以 u 为中心、半径为 1/2 的开球不包含 E 的任何可对角化自同态。
\end{enumerate}

\begin{enumerate}
\def\labelenumi{\arabic{enumi})}
\setcounter{enumi}{1}
\item
  设 E 是复希尔伯特空间。直接证明 E 的自同态 \(u\) 是弗雷德霍姆自同态，当且仅当 \(\pi(u)\) 在 \(\mathcal{C}alk(E)\) 中可逆（阿特金森定理）。（使用紧算子的谱性质，参见 IV 第 146 页第 1 节。）
\item
  设 E 是复希尔伯特空间，\(\mathrm{B}=(e_{i})_{i\in\mathrm{I}}\) 是 E 的标准正交基。对每个 \(i\in I\)，记 \(p_{i}\) 为像为 \(Ce_{i}\) 的正交投影。
\end{enumerate}

\begin{enumerate}
\def\labelenumi{\alph{enumi})}
\item
  设 K 是 E 的相对紧子集。对每个 \(\varepsilon > 0\)，存在 I 的有限子集 J，使得对每个 \(x \in \mathbf{K}\)，都有 \(\|x - \sum_{i \in \mathbf{J}} p_j(x)\| \leqslant \varepsilon\)。
\item
  设 \(u \in \mathcal{D}_{\mathrm{B}}(\mathrm{E})\)，并记其特征值族为 \(\lambda = (\lambda_{i})_{i \in \mathrm{I}}\)。族 \((\lambda_{i} p_{i})\) 在赋予相对紧收敛拓扑的 \(\mathcal{L}(\mathrm{E})\) 中可求和，且其和等于 u（参见 IV 第 147 页命题 1）。
\end{enumerate}

\begin{enumerate}
\def\labelenumi{\arabic{enumi})}
\setcounter{enumi}{3}
\item
  设 E 是无限维复希尔伯特空间。存在 E 的不是有限迹的紧自同态。
\item
  \begin{enumerate}
  \def\labelenumii{\alph{enumii})}
  \tightlist
  \item
  设 \(E = L^{2}([0,1], \mu)\)，其中 \(\mu\) 是勒贝格测度。代数 \(L^{\infty}([0,1], \mu)\) 在如下映射下的像 A 是 \(\mathcal{L}(E)\) 的极大交换子代数：该映射把 f 对应到乘以 f 的自同态；A 不具有某个 E 的标准正交基 B 所对应的 \(\mathcal{D}_{\mathrm{B}}(\mathrm{E})\) 形式。（验证代数 A 不含极小投影。）
  \end{enumerate}
\end{enumerate}

\begin{enumerate}
\def\labelenumi{\alph{enumi})}
\setcounter{enumi}{1}
\tightlist
\item
  设 E 是无限维希尔伯特空间。存在 \(\mathcal{L}(\mathrm{E})\) 的极大交换子代数，它不是对于 E 的某个标准正交基 B 而具有 \(\mathcal{D}_{\mathrm{B}}(\mathrm{E})\) 的形式。
\end{enumerate}

\begin{enumerate}
\def\labelenumi{\arabic{enumi})}
\setcounter{enumi}{5}
\tightlist
\item
  设 E 是复希尔伯特空间。
\end{enumerate}

\begin{enumerate}
\def\labelenumi{\alph{enumi})}
\item
  设 u 是 E 的紧正规自同态。E 的每个在 u 下不变的闭子空间也在 \(u^{*}\) 下不变。
\item
  若 E 无限维，则存在 E 的正规自同态 u 和闭子空间 \(F \subset E\)，使得：
\end{enumerate}

\begin{enumerate}
\def\labelenumi{(\roman{enumi})}
\item
  有 \(u(\mathrm{F}) \subset \mathrm{F}\)；
\item
  空间 F 不是 \(u^{*}\)-不变的；
\item
  空间 F° 不是在 u 下不变的；
\item
  由 \(u\) 通过限制到子空间在 F 上诱导的自同态不是正规的。
\end{enumerate}

\begin{enumerate}
\def\labelenumi{\arabic{enumi})}
\setcounter{enumi}{6}
\tightlist
\item
  本习题的目标是给出 IV 第 149 页定理 1 的一个更直接的证明。设 E 是无限维复希尔伯特空间，u 是 E 的紧正规自同态。
\end{enumerate}

\begin{enumerate}
\def\labelenumi{\alph{enumi})}
\item
  证明，对每个非零复数 \(\lambda\)，特征空间 \(\operatorname{Ker}(u - \lambda 1_{\mathrm{E}})\) 是有限维的。
\item
  假设 \(u\) 是埃尔米特的。证明 E 中存在 \(x_0 \neq 0\)，使得
\end{enumerate}

\[
\frac {\| u (x _ {0}) \| ^ {2}}{\| x _ {0} \| ^ {2}} = \sup _ {x \in \mathrm{E} - \{0 \}} \frac {\| u (x) \| ^ {2}}{\| x \| ^ {2}}.
\]

\begin{enumerate}
\def\labelenumi{\alph{enumi})}
\setcounter{enumi}{2}
\item
  证明 \(x_{0}\) 是 \(u\) 关于某个特征值 \(\lambda\) 的特征向量，且 \(|\lambda| = \|u\|\)。
\item
  证明存在一个可数集 I、E 中的标准正交族 \((e_{i})_{i\in\mathrm{I}}\) 以及实数族 \((\lambda_{i})_{i\in\mathrm{I}}\)，使得
\end{enumerate}

\[
u (x) = \sum_ {i \in I} \lambda_ {i} \langle e _ {i} | x \rangle e _ {i}
\]

对所有 \(x \in E\) 都成立。

\begin{enumerate}
\def\labelenumi{\alph{enumi})}
\setcounter{enumi}{4}
\tightlist
\item
  将结论推广到 u 紧且正规的情形。（写成 \(u = u_{1} + iu_{2}\)，其中 \(u_{1}\) 和 \(u_{2}\) 是紧的埃尔米特自同态且彼此可交换。）
\end{enumerate}

\begin{enumerate}
\def\labelenumi{\arabic{enumi})}
\setcounter{enumi}{7}
\tightlist
\item
  我们赋予群 R/Z 以其归一化的哈尔测度 \(\mu\)。
\end{enumerate}

\begin{enumerate}
\def\labelenumi{\alph{enumi})}
\tightlist
\item
  设 \(h \in \mathcal{L}^{2}(\mathbf{R}/\mathbf{Z})\)。对 \(f \in \mathcal{L}^{2}(\mathbf{R}/\mathbf{Z})\)，按下式在 R/Z 上定义函数 \(\mathrm{C}_{h}(f)\)：
\end{enumerate}

\[
\mathrm{C} _ {h} (f) (x) = \int_ {\mathbf {R} / \mathbf {Z}} f (t) h (x - t) d \mu (t).
\]

映射 \(f \mapsto \mathrm{C}_h(f)\) 是 \(\mathcal{L}^2(\mathbf{R}/\mathbf{Z})\) 的连续自同态，并通过取商定义 \(\mathrm{L}^2(\mathbf{R}/\mathbf{Z})\) 的紧正规自同态。

\begin{enumerate}
\def\labelenumi{\alph{enumi})}
\setcounter{enumi}{1}
\item
  确定 \(C_{h}\) 的谱以及该谱中各元素的谱重数。（将 h 展开为傅里叶级数。）
\item
  存在 \(h \in \mathcal{C}(\mathbf{R}/\mathbf{Z})\)，使得 \(C_{h}\) 不是有限迹的。
\end{enumerate}

\begin{enumerate}
\def\labelenumi{\arabic{enumi})}
\setcounter{enumi}{8}
\tightlist
\item
  记 E 为复希尔伯特空间 \(L^{2}([0,1])\)，其元素是在 \([0,1]\) 上关于勒贝格测度 \(\mu\) 平方可积的函数。
\end{enumerate}

令 \(u_{0} \in \mathcal{L}(E)\) 为 E 上满足下述条件的沃尔泰拉自同态：

\[
u _ {0} (f) (x) = \int_ {[ 0, x ]} f (t) d \mu (t)
\]

对所有 \(f \in \mathrm{L}^2([0,1])\) 以及几乎所有 \(x \in [0,1]\) 都成立。

\begin{enumerate}
\def\labelenumi{\alph{enumi})}
\item
  自同态 \(u_{0}\) 是紧的且 \(\operatorname{Sp}(u_{0})=\{0\}\)。
\item
  自同态 \(u_0^* u_0\) 具有有限迹；计算其迹并由此求出级数
\end{enumerate}

\[
\sum_ {n = 1} ^ {+ \infty} \frac {1}{n ^ {2}}.
\]

\begin{enumerate}
\def\labelenumi{\alph{enumi})}
\setcounter{enumi}{2}
\tightlist
\item
  令 \(v_{0}\) 为 E 上由下式定义的秩 1 自同态
\end{enumerate}

\[
v _ {0} (f) = \left(\int_ {0} ^ {1} f (t) d \mu (t)\right) \cdot 1,
\]

其中 1 表示常值函数 1 的等价类。令 \(u = u_0 - v_0\)。证明 \(u\) 是希尔伯特–施密特算子，并计算满足下式的核 N：

\[
u (f) (x) = \int_ {0} ^ {1} \mathrm{N} (x, y) f (y) d \mu (y)
\]

对 \(f \in \mathrm{E}\) 和几乎所有 \(x \in [0,1]\) 都成立。

\begin{enumerate}
\def\labelenumi{\alph{enumi})}
\setcounter{enumi}{3}
\item
  确定 u 的特征值及其谱重数。（计算 \(u(e_{n})\)，其中 \(e_{n}(t) = e^{2i\pi nt}\) 对所有 \(n \in Z\) 都成立。）
\item
  计算 \(u^{*}u\) 的迹，并再次求出通项为 \(1/n^{2}\)（\(n \geqslant 1\)）的级数的值。
\item
  将这些论证推广，证明若 \(k \geqslant 1\) 是严格正整数，则
\end{enumerate}

\[
\sum_ {n = 1} ^ {+ \infty} \frac {1}{n ^ {2 k}} \in \pi^ {2 k} {\bf Q} ^ {*}.
\]

\begin{enumerate}
\def\labelenumi{\arabic{enumi})}
\setcounter{enumi}{9}
\tightlist
\item
  设 E 是复希尔伯特空间。令 A 为 \(\mathcal{L}(\mathrm{E})\) 的星子代数。假设 A 在 E 中的表示是不可约的，且 A 含有非零紧自同态。
\end{enumerate}

\begin{enumerate}
\def\labelenumi{\alph{enumi})}
\item
  星子代数 A 含有非零有限秩投影。（证明 A 含有非零紧埃尔米特自同态，并应用函数演算。）
\item
  星子代数 A 含有秩 1 投影。（若 \(p \in \mathrm{A}\) 是最小秩的非零投影，证明代数 \(p\mathrm{A}p\) 包含于 \(\mathbf{C}p\)。）
\item
  从而 A 含有 \(\mathcal{L}^{\mathrm{c}}(\mathrm{E})\)。
\item
  反之，\(\mathcal{L}(\mathrm{E})\) 的每个含有 \(\mathcal{L}^{\mathrm{c}}(\mathrm{E})\) 的星子代数都是不可约的。
\end{enumerate}

\begin{enumerate}
\def\labelenumi{\arabic{enumi})}
\setcounter{enumi}{10}
\tightlist
\item
  设 E 是复希尔伯特空间。令 \(A \subset \mathcal{L}(\mathrm{E})\) 为紧自同态的集合，使其满足 \(u \circ v = v \circ u\)：A 中每一对元素 \((u, v)\) 都满足此关系，且 \(u^{*} \in A\) 对所有 \(u \in A\) 成立。
\end{enumerate}

存在一个可数标准正交族 \((e_i)_{i\in \mathrm{I}}\) 以及一个映射族 \((\Lambda_i)_{i\in \mathrm{I}}\)，其元素都是 \(\mathrm{A}\to \mathbf{C}\) 的映射，使得

\[
u (x) = \sum_ {i \in \mathrm{I}} \Lambda_ {i} (u) \langle e _ {i} | x \rangle e _ {i}
\]

对 A 中的每个元素 u 和 E 中的每个 x 都有。若 A 是 \(\mathcal{L}(\mathrm{E})\) 的含单位元子代数，则 \(\Lambda_{i}\) 对每个 \(i \in I\) 都是 A 的一个特征。

\begin{enumerate}
\def\labelenumi{\arabic{enumi})}
\setcounter{enumi}{11}
\tightlist
\item
  若 E 是复希尔伯特空间且 u 是 E 的紧正规自同态，证明 u 有循环向量（参见 IV 第 191 页定义 3）当且仅当满足以下条件：
\end{enumerate}

\begin{enumerate}
\def\labelenumi{(\roman{enumi})}
\item
  u 的核为 0；
\item
  u 的每个非零特征值的谱重数都等于 1。
\end{enumerate}

\begin{enumerate}
\def\labelenumi{\arabic{enumi})}
\setcounter{enumi}{12}
\tightlist
\item
  设 E 是可数型复希尔伯特空间，\((e_{n})_{n\in\mathbf{N}}\) 是 E 的标准正交基，u 是 E 的自同态。证明 u 紧，当且仅当
\end{enumerate}

\[
\lim _ {n \to + \infty} \sup _ {x \in \{e _ {0}, \dots , e _ {n - 1} \} ^ {\circ} - \{0 \}} \frac {\| u (x) \| ^ {2}}{\| x \| ^ {2}} = 0.
\]

\begin{enumerate}
\def\labelenumi{\arabic{enumi})}
\setcounter{enumi}{13}
\item
  令 \(\mathrm{X} = [0,1]\) 赋以勒贝格测度，并令 \(\mathrm{N}\colon \mathrm{X}\times \mathrm{X}\to \mathbf{R}_{+}\) 满足 \(\mathrm{N}(x,y) = |x - y|\)。自同态 \(u_{\mathrm{N}}\)（作用于 \(\mathrm{L}^2 (\mathrm{X})\)）不是正的。
\item
  设 X 是局部紧拓扑空间，\(\mu\) 是 X 上的正有界测度且支撑为 X，并令 \(N \in \mathcal{C}_{b}(X \times X)\)。
\end{enumerate}

\begin{enumerate}
\def\labelenumi{\alph{enumi})}
\tightlist
\item
  存在一个自同态 \(u\)，其作用于 \(\mathcal{C}_b(\mathrm{X})\)，使得
\end{enumerate}

\[
u (f) (y) = \int_ {\mathrm{X}} \mathrm{N} (x, y) f (x) d \mu (x)
\]

对每个函数 \(f \in \mathcal{C}_b(\mathrm{X})\) 和每个 \(y \in \mathrm{X}\) 都成立。下文假设 \(u\) 是紧的；例如当空间 \(\mathrm{X}\) 紧时就是如此。

\begin{enumerate}
\def\labelenumi{\alph{enumi})}
\setcounter{enumi}{1}
\item
  \(u_{\mathrm{N}}\) 的像包含于 \(\mathcal{C}_b(\mathrm{X})\) 中，且 \(u_{\mathrm{N}}\) 定义了从 \(\mathrm{L}_{\mathbf{C}}^{2}(\mathrm{X})\) 到 \(\mathcal{C}_b(\mathrm{X})\) 的连续线性映射。
\item
  设 \(f \in \mathrm{L}_{\mathbf{C}}^{2}(\mathrm{X})\) 且 \(\lambda \in \mathbf{C}^*\)，满足 \(u_{\mathrm{N}}(f) = \lambda f\)。则 \(f \in \mathcal{C}_b(\mathrm{X})\)。
\end{enumerate}

\(d)\) 存在一个可数族 \((f_i)_{i\in I}\)，其元素属于 \(\mathcal{C}_b(\mathrm{X})\)，以及一个非零实数族 \((\lambda_i)_{i\in I}\)，使得对所有 \(f\in \mathrm{L}_{\mathbf{C}}^2 (\mathrm{X})\)，有

\[
u (f) = \sum_ {i \in I} \lambda_ {i} \langle f _ {i} | f \rangle f _ {i}
\]

在 \(\mathcal{C}_b(\mathrm{X})\) 中。

\begin{enumerate}
\def\labelenumi{\alph{enumi})}
\setcounter{enumi}{4}
\tightlist
\item
  假设 \(u_{\mathrm{N}}\) 是 \(\mathrm{L}_{\mathbf{C}}^{2}(\mathrm{X})\) 的正自同态。对每个 \(i \in \mathrm{I}\)，记 \(h_i \in \mathcal{C}_b(\mathrm{X} \times \mathrm{X})\) 为由 \(h_i(x,y) = \overline{f_i(x)} f_i(y)\) 定义的函数。则
\end{enumerate}

\[
N = \sum_ {i \in I} \lambda_ {i} h _ {i}
\]

在赋予紧收敛拓扑的 \(\mathcal{C}(\mathrm{X}\times\mathrm{X})\) 中。（使用迪尼定理（TG，第 X 卷，第 34 页，定理 1）证明该级数收敛，然后证明等式。）

\begin{enumerate}
\def\labelenumi{\arabic{enumi})}
\setcounter{enumi}{15}
\tightlist
\item
  令 \(\mathrm{X} = [0,1]\) 赋以勒贝格测度。
\end{enumerate}

\begin{enumerate}
\def\labelenumi{\alph{enumi})}
\item
  对序列 \((f_n)_{n\in \mathbf{N}}\)（其元素是 X 上的连续函数，且在 \(\mathrm{L}_{\mathbf{C}}^2 (\mathrm{X})\) 中标准正交），存在核 \(k\in \mathcal{C}(\mathrm{X}\times \mathrm{X})\)，使得
\item
  自同态 \(u_{k}\)（作用于 \(\mathrm{L}_{\mathbf{C}}^{2}(\mathrm{X})\)）是正的；
\item
  对每个 \(n \in \mathbf{N}\)，函数 \(f_n\) 是 \(u_k\) 的特征函数，其特征值非零；
\item
  族 \((f_n)\) 是自同态 \(u_k\) 的像的一个基。
\item
  存在核 \(k \in \mathcal{C}(\mathrm{X} \times \mathrm{X})\) 以及特征函数标准正交族 \((f_n)\)，这些是自同态 \(u_k\)（作用于 \(\mathrm{L}_{\mathbf{C}}^2(\mathrm{X})\)）的有界特征函数，使得序列 \((\|f_n\|_\infty)\) 无界。
\end{enumerate}

\begin{enumerate}
\def\labelenumi{\arabic{enumi})}
\setcounter{enumi}{16}
\item
  设 E、F 是希尔伯特空间。双线性映射 \((u,v)\mapsto\) \(\mathrm{Tr}(uv)\) 从巴拿赫空间 \(\mathcal{L}_1(\mathrm{E};\mathrm{F})'\) 到 \(\mathcal{L}(\mathrm{F};\mathrm{E})\) 定义了一个等距同构。
\item
  设 E 是复希尔伯特空间，令 \(I_{E}\) 为 IV 第 151 页第 3 号中定义的集合。
\end{enumerate}

令 \(n \in I_{E}\)。令 F 为复希尔伯特空间。映射 \(u \mapsto \lambda_{n}(|u|)\) 从 \(\mathcal{L}^{c}(E;F)\) 到 \(R_{+}\) 是连续的。

\begin{enumerate}
\def\labelenumi{\arabic{enumi})}
\setcounter{enumi}{18}
\item
  保持 IV 第 161 页引理 5 的记号。证明若 \(f\) 的像包含于 \(\mathbf{R}\)，且 \(f\) 严格凸、\(\varrho_n > 0\)，则不等式 (6) 中取等号，当且仅当 \(a_i = b_i\) 对 \(0 \leqslant i \leqslant n\) 成立。
\item
  设 E 是预希尔伯特空间，F 是巴拿赫空间，且 \(u: E \to F\) 是线性映射。
\end{enumerate}

\begin{enumerate}
\def\labelenumi{\alph{enumi})}
\item
  证明 u 是紧线性映射，当且仅当 \(\|u(e_{n})\|\to0\) 对 E 中的每个标准正交序列 \((e_{n})\) 都成立。（为证明充分性，使用 TG，第 IX 卷，第 20 页命题 14。）
\item
  假设 E 是可数型希尔伯特空间且 F=E。存在 E 的一个标准正交基 \((e_{n})_{n\in\mathbf{N}}\)，使 \(\|u(e_{n})\|\to0\)，其充要条件是 u 不是弗雷德霍姆自同态。
\item
  假设 E 是希尔伯特空间且 F=E。若 \(\langle u(e_{n}) | e_{n} \rangle \to 0\) 对 E 的每个标准正交序列 \((e_{n})\) 都成立，则 u 是紧线性映射。
\end{enumerate}

\begin{enumerate}
\def\labelenumi{\arabic{enumi})}
\setcounter{enumi}{20}
\tightlist
\item
  设 E 是可数型希尔伯特空间。
\end{enumerate}

\begin{enumerate}
\def\labelenumi{\alph{enumi})}
\item
  设 A 是 E 中一个沿有限集补集滤子弱收敛于 0 的无限单位向量集。证明存在一个取值于 A 的序列 \((x_{n})_{n\geqslant1}\) 和 E 中的标准正交序列 \((e_{n})_{n\geqslant1}\)，使得 \(\|x_{n}-e_{n}\|\leqslant1/n\) 对所有 \(n\geqslant1\) 都成立。
\item
  设 u 是从 E 到 E 的线性映射，不必连续。假设对 E 中每个标准正交序列 \((e_{n})\)，序列 \((u(e_{n}))\) 都弱收敛于 0。证明 u 连续且紧。
\item
  设 u 是 E 的自同态。存在 E 的一个标准正交基 \((e_{n})_{n\in\mathbf{N}}\)，使序列 \((u(e_{n}))_{n\in\mathbf{N}}\) 弱收敛于 0，当且仅当不存在 E 的自同态 v，使得 \(v\circ u-1_{E}\) 是紧的。
\end{enumerate}

\begin{enumerate}
\def\labelenumi{\arabic{enumi})}
\setcounter{enumi}{21}
\tightlist
\item
  设 E 是复希尔伯特空间。令 u 和 v 为 E 的正自同态。
\end{enumerate}

\begin{enumerate}
\def\labelenumi{\alph{enumi})}
\tightlist
\item
  对每个满足 \(0 < p < 1\) 的实数 p 以及每个 \(x \in R_{+}\)，有
\end{enumerate}

\[
x ^ {p} = \frac {\sin (p \pi)}{\pi} \int_ {0} ^ {+ \infty} t ^ {p} \left(\frac {1}{t} - \frac {1}{t + x}\right) d t.
\]

\begin{enumerate}
\def\labelenumi{\alph{enumi})}
\setcounter{enumi}{1}
\item
  假设 \(u \leqslant v\)。对每个实数 \(p\)，满足 \(0 < p < 1\)，有 \(u^p \leqslant v^p\)。（使用 I 第 119 页引理 14 的 b) 和上一问。）\\
\item
  令 \(n \in \mathbf{N}\)。置 \(a = v^{n/2} \circ u^n \circ v^{n/2}\)，且 \(b = (v^{1/2} \circ u \circ v^{1/2})^n\)。若 \(a \leqslant 1_{\mathrm{E}}\)，则 \(b \leqslant 1_{\mathrm{E}}\)。
\item
  假设 u 和 v 具有有限迹。对每个整数 \(n \in N\)，有
\end{enumerate}

\[
\operatorname{Tr} \left(\left(v ^ {1 / 2} \circ u \circ v ^ {1 / 2}\right) ^ {n}\right) \leqslant \operatorname{Tr} \left(v ^ {n / 2} \circ u ^ {n} \circ v ^ {n / 2}\right).
\]

（令 \(a\) 和 \(b\) 如问题 \(c\) 所定义；利用 IV 第 161 页引理 5，把问题归结为证明 \(\widehat{\Lambda}^{n}b\) 的最大特征值小于 \(\widehat{\Lambda}^{n}a\) 的最大特征值；以 \(\widehat{\Lambda}^{n}\mathrm{E}\) 替代 E，并以 \(u\)、\(v\) 分别替换为 \(\widehat{\Lambda}^{n}u\)、\(\widehat{\Lambda}^{n}v\)，然后应用上一问。）
\exercisesection{}

除另有说明外，在本节的各习题中，E 表示一个复希尔伯特空间。

\begin{enumerate}
\def\labelenumi{\arabic{enumi})}
\tightlist
\item
  设 \(u\) 是 E 的自同态，其像在 E 中是闭的。设 \(q\) 是 E 的正交投影算子，并设 \(v\) 是 E 的像的自同态 \(x \mapsto q(u(x))\) 。
\end{enumerate}

\begin{enumerate}
\def\labelenumi{\alph{enumi})}
\item
  \(v\) 的伴随是 E 的像的自同态 \(x \mapsto q(u^{*}(x))\) 。
\item
  若 \(u\) 是埃尔米特的，则 \(v\) 是埃尔米特的。
\item
  有可能 \(u\) 是正规的而 \(v\) 不是正规的。
\end{enumerate}

\begin{enumerate}
\def\labelenumi{\arabic{enumi})}
\setcounter{enumi}{1}
\item
  直接利用谱定理（IV, p. 193 的定理 1 的推论）证明 IV, p. 149 的定理 1。
\item
  设 \(S \subset \mathbf{C}\) 是 \(\mathbf{C}\) 的一个紧子集。设 \(\varphi \colon \mathcal{C}(S) \to \mathcal{L}(E)\) 是对合代数的一个连续的含单位元态射。存在唯一的正规自同态 \(u\) ，它作用在 \(E\) 上，谱为 \(S\)，并且使得 \(\varphi(f) = f(u)\) 对所有 \(f \in \mathcal{L}^{\infty}(S)\) 成立。
\item
  设 \(u \in \mathcal{L}(\mathrm{E})\) 是正规的。\(u\) 的谱的凸包是 \(\mathbf{C}\) 中数值值域 \(\iota(u)\) （相应于 \(u\) ；参见 I, p. 135 的定义 2）的闭包。
\item
  设 \(u \in \mathcal{L}(\mathrm{E})\) 满足 \(\| u \| \leqslant 1\) 。对每个 \(\lambda \in \mathbf{C}\) ，用 \(p_{\lambda}\) 表示 E 到 \(u\) 的对应于 \(\lambda\) 的特征空间上的正交投影算子。对 E 中所有的 \(x\) 与 \(y\)，我们有
\end{enumerate}

\[
\lim_{N\to +\infty}\frac{1}{N}\sum_{n = 1}^{N}|\langle x \mid u^{n}(y)\rangle |^{2} = \sum_{\substack{\lambda \in \mathbf{C}\\ |\lambda | = 1}}|\langle p_{\lambda}(y) \mid p_{\lambda}(x)\rangle^{2}|\\ = \sum_{\substack{\lambda \in \mathbf{C}\\ |\lambda | = 1}}|\langle y \mid p_{\lambda}(x)\rangle |^{2}
\]

（可以利用 II, p. 300 的习题 64）。

\begin{enumerate}
\def\labelenumi{\arabic{enumi})}
\setcounter{enumi}{5}
\tightlist
\item
  设 E 是预希尔伯特空间，F 是巴拿赫空间，\(u: E \rightarrow F\) 是线性映射。
\end{enumerate}

\begin{enumerate}
\def\labelenumi{\alph{enumi})}
\item
  证明：u 是紧线性映射当且仅当 \(\|u(e_{n})\|\to0\) 对 E 的每个标准正交序列 \((e_{n})\) 成立。（为证明该条件是充分的，利用 TG, IX, p. 20 的命题 14。）
\item
  假设 E 是可分希尔伯特空间且 F = E。为使存在 E 的标准正交基 \((e_{n})\) 使得 \(\|u(e_{n})\| \to 0\)，必须且只需 u 不是弗雷德霍姆算子。
\item
  假设 E 是希尔伯特空间，并且 \(\langle u(e_{n}), e_{n} \rangle \to 0\) 对 H 的每个标准正交序列 \((e_{n})\) 成立。证明 u 是紧线性映射。
\end{enumerate}

\begin{enumerate}
\def\labelenumi{\arabic{enumi})}
\setcounter{enumi}{6}
\tightlist
\item
  设 \(u\) 与 \(v\) 是 E 的连续自同态。
\end{enumerate}

\begin{enumerate}
\def\labelenumi{\alph{enumi})}
\item
  证明 \(\operatorname{Im}(u) + \operatorname{Im}(v) = \operatorname{Im}((u \circ u^* + v \circ v^*)^{1/2})\) 。
\item
  若 \(u\) 是正的且具有闭像，证明 \(\operatorname{Im}(u) = \operatorname{Im}(u^{1/2})\) 。
\item
  若 \(u, v\) 是正的且具有闭像，并且 \(u + v\) 也具有闭像，则我们有 \(\operatorname{Im}(u + v) = \operatorname{Im}(u) + \operatorname{Im}(v)\) 。
\end{enumerate}

\begin{enumerate}
\def\labelenumi{\arabic{enumi})}
\setcounter{enumi}{7}
\tightlist
\item
  设 F 是复希尔伯特空间，\(u \colon \mathrm{E} \to \mathrm{F}\) 是具有闭像的连续线性映射。
\end{enumerate}

\begin{enumerate}
\def\labelenumi{\alph{enumi})}
\tightlist
\item
  证明存在唯一的连续线性映射 \(v \in \mathcal{L}(\mathrm{F};\mathrm{E})\) 满足下列条件：
\end{enumerate}

\begin{enumerate}
\def\labelenumi{(\roman{enumi})}
\item
  \(u \circ v \circ u = u\) 且 \(v \circ u \circ v = v\) ；
\item
  \(u \circ v\) 与 \(v \circ u\) 是自伴的。
\end{enumerate}

我们把这一映射记作 \(u^{\dagger}\) （穆尔–彭罗斯伪逆）。

\begin{enumerate}
\def\labelenumi{\alph{enumi})}
\setcounter{enumi}{1}
\tightlist
\item
  证明 \(\operatorname{Im}(u) = \operatorname{Im}(u \circ u^{\dagger})\)。证明：对 \(x \in E\) 与 \(y \in \operatorname{Im}(u)\) ，条件 \(u(x) = y\) 与 \(x - u^{\dagger}(y) \in \operatorname{Im}(1_{\operatorname{E}} - u^{\dagger} \circ u)\) 等价。
\end{enumerate}

\begin{enumerate}
\def\labelenumi{\arabic{enumi})}
\setcounter{enumi}{8}
\tightlist
\item
  设 \(u\) 与 \(v\) 是 E 的具有闭像的自同态。若 \(u + v\) 也具有闭像，我们令 \(u: v = u \circ (u + v)^{\dagger} \circ v\) ，其中 \((u + v)^{\dagger}\) 是 \(u + v\) 的穆尔–彭罗斯伪逆。
\end{enumerate}

\begin{enumerate}
\def\labelenumi{\alph{enumi})}
\item
  证明 \(u:v\) 是自伴的。
\item
  我们有 \(u:v=v:u\) 且 \(\operatorname{Im}(u:v)=\operatorname{Im}(u)\cap\operatorname{Im}(v)\) 。
\item
  假设 u 与 v 是正交投影算子。证明 2u : v 是以 \(\operatorname{Im}(u) \cap \operatorname{Im}(v)\) 为像的正交投影算子。
\end{enumerate}

\begin{enumerate}
\def\labelenumi{\arabic{enumi})}
\setcounter{enumi}{9}
\tightlist
\item
  E 的连续自同态 \(u\) 称为亚正规的，如果 \(u^{*} \circ u \geqslant u \circ u^{*}\)；称为仿正规的，如果 \(\|u^{2}(x)\|\|x\| \geqslant \|u(x)\|^{2}\) 对所有 \(x \in E\) 成立。
\end{enumerate}

\begin{enumerate}
\def\labelenumi{\alph{enumi})}
\item
  若 \(u\) 是仿正规的，则 \(u^n\) 对每个 \(n \in \mathbf{N}\) 都是仿正规的。
\item
  若 u 是仿正规的且可逆，则 \(u^{-1}\) 是仿正规的。
\item
  亚正规自同态是仿正规的。
\item
  设 u 是谱包含在单位圆周中的仿正规自同态；则 u 是酉的。
\item
  若 \(u\) 是仿正规自同态，则它的范数 \(\| u\|\) 等于它的谱半径 \(\varrho(u)\) 。
\item
  我们假设 u 是仿正规的。
\end{enumerate}

\begin{enumerate}
\def\labelenumi{\arabic{enumi})}
\setcounter{enumi}{10}
\tightlist
\item
  设 u 是 E 的连续自同态。
\end{enumerate}

\begin{enumerate}
\def\labelenumi{\alph{enumi})}
\tightlist
\item
  为使 \(u\) 是仿正规的，必须且只需
\end{enumerate}

\[
(u ^ {*}) ^ {2} \circ u ^ {2} - 2 t (u ^ {*} \circ u) + t ^ {2} 1 _ {\mathrm{E}} \geqslant 0
\]

对每个实数 t \textgreater{} 0 成立。

\begin{enumerate}
\def\labelenumi{\alph{enumi})}
\setcounter{enumi}{1}
\tightlist
\item
  为使 \(u\) 是仿正规的，必须且只需
\end{enumerate}

\[
p \circ q ^ {2} \circ p - 2 t p ^ {2} + t 1 _ {\mathrm{E}} \geqslant 0
\]

对所有 \(t > 0\) 成立，其中 \(p = (u \circ u^{*})^{1/2}\) ，\(q = (u^{*} \circ u)^{1/2}\) 。

\begin{enumerate}
\def\labelenumi{\arabic{enumi})}
\setcounter{enumi}{11}
\tightlist
\item
  \begin{enumerate}
  \def\labelenumii{\alph{enumii})}
  \tightlist
  \item
    设 u 与 v 是 E 的连续正自同态。假设对每个 t \textgreater{} 0，我们有
  \end{enumerate}
\end{enumerate}

\[
2 t u ^ {2} \circ (u ^ {2} + t ^ {2} 1 _ {\mathrm{E}}) ^ {- 1} \leqslant v \leqslant \frac {1}{2 t} (u ^ {2} + t ^ {2} 1 _ {\mathrm{E}}).
\]

证明 u = v。

\begin{enumerate}
\def\labelenumi{\alph{enumi})}
\setcounter{enumi}{1}
\tightlist
\item
  为使 \(u\) 是正规自同态，必须且只需 \(u\) 与 \(u^*\) 都是仿正规的并且有相同的核。
\end{enumerate}

\begin{enumerate}
\def\labelenumi{\arabic{enumi})}
\setcounter{enumi}{12}
\tightlist
\item
  \begin{enumerate}
  \def\labelenumii{\alph{enumii})}
  \tightlist
  \item
    设 A 是 E 的单位向量所成的无限集，它沿有限子集的余集滤子弱收敛于 0。证明存在取值于 A 的序列 \((x_{n})_{n\geqslant 1}\) 以及 E 中的标准正交序列 \((e_n)_{n\geqslant 1}\) ，使得 \(\| x_n - e_n\| \leqslant 1 / n\) 对所有 \(n\geqslant 1\) 成立 。
  \end{enumerate}
\end{enumerate}

\begin{enumerate}
\def\labelenumi{\alph{enumi})}
\setcounter{enumi}{1}
\item
  设 u 是 E 的自同态，未必连续。假设对 E 中每个标准正交序列 \((e_{n})\) ，序列 \((u(e_{n}))\) 都弱收敛于 0。证明 u 是连续且紧的。
\item
  设 u 是 E 的连续自同态。为使存在 E 的标准正交基 \((e_{n})\) 使序列 \((u(e_{n}))\) 弱收敛于 0，必须且只需：对 E 的每个自同态 v，\(v \circ u - 1_{E}\) 都不是紧的。
\end{enumerate}

\begin{enumerate}
\def\labelenumi{\arabic{enumi})}
\setcounter{enumi}{13}
\tightlist
\item
  设 X 是局部紧空间，\(\mu\) 是 X 上的正有界测度且 \(\mu(\mathrm{X}) = 1\) ，并设 \(f: X \to X\) 是 \(\mu\)-可测映射且满足 \(f(\mu) = \mu\) 。我们用 \(u_{f}\) 表示 \(\mathrm{L}_{\mathbf{C}}^{2}(\mathrm{X}, \mu)\) 的由 \(\varphi \mapsto f \circ \varphi\) 定义的酉自同态（参见 I, p. 190 的习题 16）。对每个 \(n \in N\)，我们记 \(f^{n} = f^{\circ n}\) （\(n\) 次复合），并且对 X 的每个子集 A，记 \(f^{-n}(\mathrm{A}) = g^{-1}(\mathrm{A})\) ，其中 \(g = f^{n}\) 。
\end{enumerate}

我们称 \(f\) 关于 \(\mu\) 是弱混合的，如果对 X 的所有 \(\mu\)-可测子集 A 与 B，我们有

\[
\lim _ {N \rightarrow + \infty} \frac {1}{N} \sum_ {n = 0} ^ {N} | \mu (A \cap f ^ {- n} (B) - \mu (A) \mu (B) | = 0.
\]

\begin{enumerate}
\def\labelenumi{\alph{enumi})}
\tightlist
\item
  下列条件等价：
\end{enumerate}

\begin{enumerate}
\def\labelenumi{(\roman{enumi})}
\item
  映射 f 关于 \(\mu\) 是弱混合的；
\item
  对所有 \(\varphi_{1},\varphi_{2} \in \mathrm{L}_{\mathbf{C}}^{2}(\mathrm{X}, \mu)\) ，我们有
\end{enumerate}

\[
\lim _ {N \to + \infty} \frac {1}{N} \sum_ {n = 0} ^ {N} | \langle \varphi_ {1} | u _ {f} ^ {n} (\varphi_ {2}) \rangle - \langle 1 | \varphi_ {2} \rangle \langle \varphi_ {1} | 1 \rangle | = 0;
\]

\begin{enumerate}
\def\labelenumi{(\roman{enumi})}
\setcounter{enumi}{2}
\item
  映射 \(f \times f: X \times X \to X \times X\) 是 \((\mu \otimes \mu)\) -遍历的（参见 I, p. 190 的习题 16）；
\item
  相应于特征值 1 的特征空间的维数为 1，且对每个 \(\lambda \in \mathbf{C} - \{1\}\) ，\(u_{f}\) 相应于 \(\lambda\) 的特征空间化为 0。
\end{enumerate}

（为证明 (iii) 蕴含 (i)，利用 II, p. 300 的习题 64。）

\begin{enumerate}
\def\labelenumi{\alph{enumi})}
\setcounter{enumi}{1}
\item
  考虑 I, p. 190 习题 16 中 e) 部分的例子。映射 f 关于 \(\mu\) 是弱混合的。
\item
  考虑 I, p. 190 习题 16 中问题 f) 的例子。映射 f 关于 \(\mu\) 不是弱混合的。
\end{enumerate}

\begin{enumerate}
\def\labelenumi{\arabic{enumi})}
\setcounter{enumi}{14}
\tightlist
\item
  设 \(G = (S, F, o, t, \iota)\) 是一个图（TA, II, p. 155 的定义 1）。我们假设 G 是连通的，S 是可数的，并且 S 至少有两个元素。对每个 \(x \in S\) ，我们用 \(v(x)\) 表示 G 的以 x 为端点的边的数目；我们有 \(v(x) \geqslant 1\) 对所有 \(x \in S\) 成立，并且假设 \(v(x)\) 对所有 \(x \in S\) 是有限的。对 S 中所有 x 与 y，我们用 \(a(x, y)\) 表示 G 中以 x 与 y 为端点的箭的数目。我们假设 \(a(x, y)\) 对所有 \((x, y) \in S^{2}\) 是有限的。
\end{enumerate}

我们定义 S 上的测度 \(\mu\) 如下：\(\mu(\{x\}) = v(x)\) 对所有 \(x \in S\) 成立。测度 \(\mu\) 是正的。我们用 \(\mathcal{F}(\mathrm{G})\) 表示 S 上取复值的函数所成的向量空间。我们用 \(\widetilde{\mathrm{M}}_{\mathrm{G}}\) 表示 \(\mathcal{F}(\mathrm{G})\) 的满足下式的自同态

\[
\widetilde {\mathrm{M}} _ {\mathrm{G}} (\varphi) (x) = \frac {1}{v (x)} \sum_ {y} a (x, y) \varphi (y),
\]

对所有 \(\varphi\in\mathcal{F}(\mathrm{G})\) 及每个 \(x\in\mathrm{S}\) 成立（“马尔可夫算子”）。

\begin{enumerate}
\def\labelenumi{\alph{enumi})}
\item
  自同态 \(\widetilde{\mathrm{M}}_{\mathrm{G}}\) 通过向子空间的过渡诱导出 \(\mathrm{M_G}\) ，即 \(\mathrm{L}_{\mathbf{C}}^{2}(\mathrm{S},\mu)\) 的一个自同态。
\item
  自同态 \(M_{G}\) 是自伴的。它满足
\end{enumerate}

\[
\langle \varphi | 1 - \mathrm{M} _ {\mathrm{G}} (\varphi) \rangle = \frac {1}{2} \sum_ {(x, y) \in \mathrm{S} ^ {2}} a (x, y) | \varphi (x) - \varphi (y) | ^ {2}
\]

\[
\langle \varphi | 1 + \mathrm{M} _ {\mathrm{G}} (\varphi) \rangle = \frac {1}{2} \sum_ {(x, y) \in \mathrm{S} ^ {2}} a (x, y) | \varphi (x) + \varphi (y) | ^ {2}
\]

对所有 \(\varphi\in\mathrm{L}^{2}(\mathrm{S},\mu)\) 成立。此外，我们有 \(-1\leqslant\mathrm{M}_{\mathrm{G}}\leqslant1\) 。

\begin{enumerate}
\def\labelenumi{\alph{enumi})}
\setcounter{enumi}{2}
\tightlist
\item
  设 \(d \geqslant 1\) 是整数，并假设 G 是 \(Z^{d}\) 相对于集合 \(\{e_{i}\}\) 的凯莱图，其中 \((e_{i})\) 是 \(Z^{d}\) 的典范基（TA, II, p. 221 的习题 7）。\(\varrho(\mathrm{M}_{\mathrm{G}})\) （\(M_{G}\) 的谱半径）等于 1。
\end{enumerate}

\(d\) ) 设 \(d \geqslant 3\) 是整数。假设 G 是一棵 \(d\) -正则的无限树（TA, II, p. 157），即其中每个顶点都有 \(d\) 个邻点的无限树。设 \(x_0\) 是 G 的一个固定顶点。对每个 \(\lambda \in \mathbf{C}\) ，存在唯一的函数 \(\varphi_{\lambda} \colon S \to \mathbf{C}\) ，使得 \(\varphi_{\lambda}(x_0) = 1\) 且 \(\widetilde{\mathrm{M}}_{\mathrm{G}}(\varphi_{\lambda}) = \lambda \varphi_{\lambda}\) 。此外，对每个 \(x \in S\) ，\(\varphi_{\lambda}\) 在 \(x\) 处的值只依赖于 \(x\) 到 \(x_0\) 的距离（参见 TA, II, p. 222 的习题 10）。

\begin{enumerate}
\def\labelenumi{\alph{enumi})}
\setcounter{enumi}{4}
\tightlist
\item
  \(M_{G}\) 的谱是区间 \([-2\sqrt{d-1}/d, 2\sqrt{d-1}/d]\) （“凯斯滕定理”）。特别地，我们有 \(\varrho(\mathrm{M}_{\mathrm{G}})=2\sqrt{d-1}/d\) 。
\end{enumerate}

\begin{enumerate}
\def\labelenumi{\arabic{enumi})}
\setcounter{enumi}{15}
\tightlist
\item
  我们沿用上一个习题的记号，并假设 S 与 F 是有限的。我们用 \(h_{G}\) 表示图 G 的切格常数，它定义为
\end{enumerate}

\[
h _ {\mathrm{G}} = \inf _ {\varnothing \neq \mathrm{V} \subset \mathrm{S}} \frac {\operatorname{Card} (\mathrm{E} (\mathrm{V} , \mathrm{S} - \mathrm{V}))}{\inf (\operatorname{Card} (\mathrm{V}) , \operatorname{Card} (\mathrm{S} - \mathrm{V}))},
\]

其中 E(V, W) 表示 G 的一个端点在 V 中、另一个端点在 W 中的箭所成的集。

\begin{enumerate}
\def\labelenumi{\alph{enumi})}
\tightlist
\item
  我们有 \(\varrho(\mathrm{M}_{\mathrm{G}})=1\) ，并且 \(\lambda=1\) 是 \(\mathrm{M}_{\mathrm{G}}\) 的重数为 1 的特征值。我们有 \(-1\in\mathrm{Sp}(\mathrm{M}_{\mathrm{G}})\) 当且仅当 G 是二部的（参见 TA, II, p. 219 的习题 2）；这时，\(\mathrm{M}_{\mathrm{G}}\) 的谱在 \(\lambda\mapsto-\lambda\) 下不变，并且 \(\lambda\) 的谱重数等于 \(-\lambda\) 的谱重数对所有 \(\lambda\in\mathbf{C}\) 成立。
\end{enumerate}

我们用 \(\mathrm{L}_{0}^{2}(\mathrm{S},\nu)\) 表示在 \(\mathrm{L}_{\mathbf{C}}^{2}(\mathrm{S},\nu)\) 中、\(M_{G}\) 的相应于特征值 1 与 -1 的诸特征空间之并的正交补。通过向子空间的过渡诱导的自同态 \(M_{G}\) 是 \(\mathrm{L}_{0}^{2}(\mathrm{S},\nu)\) 的自同态，记作 \(M_{G,0}\)。我们有 \(\varrho(\mathrm{M}_{\mathrm{G},0}) < 1\)。我们用 \(\lambda_{1}(\mathrm{G})\) 表示 \(1 - M_{G}\) 的最小非零特征值。

\begin{enumerate}
\def\labelenumi{\alph{enumi})}
\setcounter{enumi}{1}
\tightlist
\item
  设 \(v_{+} = \sup v(x)\) ，\(v_{-} = \inf v(x)\) 。我们有
\end{enumerate}

\[
\lambda_ {1} (\mathrm{G}) \leqslant \frac {2 v _ {+}}{v _ {-} ^ {2}} h _ {\mathrm{G}}.
\]

\begin{enumerate}
\def\labelenumi{\alph{enumi})}
\setcounter{enumi}{2}
\tightlist
\item
  设 \(\varphi\colon S\to\mathbf{R}\) 是 \(S\) 上的非常数实值函数；令 \(m=\inf\varphi(x)\) ，\(M=\sup\varphi(x)\) 。设 \(t_{0}\in\mathbf{R}\) 使得：对每个 \(t\in\mathbf{R}\) ，集合 \(W_{t}=\overline{\varphi}^{-1}(]-\infty,t_{0}[)\) 满足 \(\mathrm{Card}(W_{t})\leqslant\mathrm{Card}(S)/2\) 当且仅当 \(t<t_{0}\) 。设 \(\nu\) 是支撑含于 \([m,M]\) 的无原子测度（INT, V, p. 67 的定义 5）。令
\end{enumerate}

\[
\mathrm{A} = \frac {1}{2} \sum_ {(x, y) \in \mathrm{S} ^ {2}} a (x, y) \nu ([ \varphi (x), \varphi (y) ])
\]

\[
\mathrm{B} = \sum_ {x \in \mathrm{S}} \nu ([ t _ {0}, \varphi (x) ]),
\]

其中，当 \(a < b\) 是实数时，我们记 \(\nu([b, a]) = \nu([a, b])\) 。存在 \(t \in \mathbf{R}\) 使得

\[
\frac {\operatorname{Card} \left(\mathrm{E} \left(\mathrm{W} _ {t} , \mathrm{S} - \mathrm{W} _ {t}\right) \right.}{\inf \left(\operatorname{Card} \left(\mathrm{W} _ {t}\right) , \operatorname{Card} \left(\mathrm{S} - \mathrm{W} _ {t}\right)\right)} \leqslant \frac {\mathrm{A}}{\mathrm{B}}.
\]

\begin{enumerate}
\def\labelenumi{\alph{enumi})}
\setcounter{enumi}{3}
\tightlist
\item
  我们有
\end{enumerate}

\[
h _ {\mathrm{G}} \leqslant v _ {+} \sqrt {2 \lambda_ {1} (\mathrm{G})}.
\]

（把上一问题应用于 \(\varphi\) 取为 \(\mathrm{M_G}\) 的相应于特征值 \(\lambda_1(\mathrm{G})\) 的特征函数并且 \(\nu\) 取适当选取的测度的情形。）

\begin{enumerate}
\def\labelenumi{\alph{enumi})}
\setcounter{enumi}{4}
\tightlist
\item
  我们称一个族 \((\mathrm{G}_{j})_{j\in\mathrm{J}}\) （由有限图 \(\mathrm{G}_{j}=(\mathrm{S}_{j},\mathrm{F}_{j},o_{j},t_{j},i_{j})\) 组成）为扩张图族，如果下列条件成立：
\end{enumerate}

\begin{enumerate}
\def\labelenumi{(\roman{enumi})}
\item
  沿着 J 的有限子集的余集滤子，我们有 \(\mathrm{Card}(\mathrm{S}_{j}) \to +\infty\) ；
\item
  存在 \(C \geqslant 0\) ，使得对每个 j 及每个 \(x \in S_{j}\) ，\(G_{j}\) 在 x 处的度 \(\leqslant C\) ；
\item
  存在 \(\delta > 0\) ，使得 \(h_{\mathrm{G}_j} \geqslant \delta\) 对所有 \(j \in \mathrm{J}\) 成立。
\end{enumerate}

证明 \((\mathrm{G}_{j})\) 是扩张图族当且仅当 (i)、(ii) 以及条件

(iii') 存在 \(\delta > 0\) ，使得 \(\lambda_1(G_j) \geqslant \delta\) 对所有 \(j \in J\) 成立。

\begin{enumerate}
\def\labelenumi{\arabic{enumi})}
\setcounter{enumi}{16}
\tightlist
\item
  \begin{enumerate}
  \def\labelenumii{\alph{enumii})}
  \tightlist
  \item
    设 \(n \in \mathbf{N}\)。设 E 是维数为 \(n\) 的希尔伯特空间，\(u \in \mathcal{L}(\mathrm{E})\) ，并设 \(x \in \mathrm{E}\) 是范数为 1 的向量且满足 \(u(x) = 0\)。对每个范数为 1 的向量 \(y \in \mathrm{E}\) ，我们用 \(\Delta_u(y)\) 表示在 \(y^\circ\) 上由 \(z \mapsto \langle z | u(z) \rangle\) 定义的二次形式的判别式。
  \end{enumerate}
\end{enumerate}

对每个范数为 1 的 \(y \in E\) ，我们有 \(|\langle x | y \rangle|^{2} \Delta_{u}(x) = \Delta_{u}(y)\) 。（注意，\(\Delta_{u}(y) = \langle a | \bigwedge u^{n-1}(a) \rangle / \langle a | a \rangle\) 对空间 \(\bigwedge^{n-1} y^{\circ} \subset \bigwedge^{n-1} E\) 的每个非零元素 a 成立，其中空间 \(\bigwedge^{n-1} E\) 赋有其典范希尔伯特结构，参见 EVT, V, p.33；然后分别在 x 与 y 正交以及 x = y 的情形验证该公式。）

\begin{enumerate}
\def\labelenumi{\alph{enumi})}
\setcounter{enumi}{1}
\tightlist
\item
  设 A 是复系数的 \((n,n)\) 型矩阵，满足 \({}^{t}\overline{A}=A\) 。对每个满足 \(1\leqslant i\leqslant n\) 的整数 i，用 \(A_{i}\) 表示 \((n-1,n-1)\) 型的、通过删除 A 的第 \(i\) 行与第 \(i\) 列所得的矩阵。设 \((e_{1},\ldots,e_{n})\) 是 \(C^{n}\) 的、赋有其典范希尔伯特结构且由 A 的特征向量组成的标准正交基，并用 \(\lambda_{i}\) 表示 u 相应于 \(e_{i}\) 的特征值。记 \(e_{i}=(e_{i,j})_{1\leqslant j\leqslant n}\) 。
\end{enumerate}

对所有 \((i,j)\) ，我们有

\[
\left| e _ {i, j} \right| ^ {2} \prod_ {k \neq i} \left(\lambda_ {i} - \lambda_ {k}\right) = \det \left(\lambda_ {i} - A _ {j}\right).
\]

（化归到 \(\lambda_{i}=0\) 的情形，然后把上一问题应用于空间 \(C^{n}\) 、自同态 \(u\colon y\mapsto Ay\) 以及向量 \(x=e_{i}\) 。）

\begin{enumerate}
\def\labelenumi{\arabic{enumi})}
\setcounter{enumi}{17}
\tightlist
\item
  设 X 是紧可度量化空间，\(\mu\) 是 X 上的正测度。记 \(\mathrm{E} = \mathrm{L}_{\mathbf{C}}^{2}(\mathrm{X}, \mu)\)。设 \((\mathcal{X}_{n})_{n \in \mathbf{N}}\) 是把 X 分成可测集的有限分划的序列，使得分划 \(X_{n+1}\) 比 \(X_{n}\) 更细对所有 \(n \in N\) 成立，并且使得 \(X_{n}\) 的各部分的直径中的最大者当 \(n \to +\infty\) 时趋于 0 。
\end{enumerate}

设 \(E_{n} \subset E\) 是由满足如下条件的 \(f \in \mathrm{L}^{2}(\mathrm{X}, \mu)\) 组成的空间：f \(\mu\)-几乎处处在 \(X_{n}\) 的每个部分 Y 上为常数。我们令 \(E_{-1} = \{0\}\)。对每个 \(n \geqslant -1\)，用 \(p_{n}\) 表示 E 的以 \(E_{n}\) 为像的正交投影算子。

\begin{enumerate}
\def\labelenumi{\alph{enumi})}
\item
  序列 \((p_{n})_{n\in\mathbf{N}}\) 在赋有逐点收敛拓扑的 E 的自同态所成的空间 \(\mathcal{L}_{s}(\mathrm{E})\) 中收敛于恒等映射。（首先考虑 \(p_{n}(f)\) 当 \(f\in\mathcal{C}(\mathrm{X})\) 时的极限。）
\item
  对每个 \(n \in N\) ，自同态 \(q_{n} = p_{n} - p_{n-1}\) 是到 \(F_{n} = E_{n} \cap E_{n-1}^{\circ}\) 上的正交投影算子，并且诸空间 \(F_{n}\) 两两正交。
\item
  我们有 \(\sum_{n} q_{n} = 1_{\mathrm{E}}\) 在 \(\mathcal{L}_s(\mathrm{E})\) 中成立，并且空间 E 是诸空间 \(\mathrm{F}_n\) 的希尔伯特直和。
\end{enumerate}

\begin{enumerate}
\def\labelenumi{\arabic{enumi})}
\setcounter{enumi}{18}
\tightlist
\item
  设 \(\mathrm{X} \subset \mathbf{R}\) 是 \(\mathbf{R}\) 的紧子集，\(\mu\) 是 \(\mathrm{X}\) 上的正测度，而 \(u\) 是由 \(\mathrm{X}\) 的恒等函数给出的乘法自同态。
\end{enumerate}

我们记 \(\mathrm{E}=\mathrm{L}^{2}(\mathrm{X},\mu)\) ，并用 \(\mathcal{L}_{s}(\mathrm{E})\) 表示赋有逐点收敛拓扑的 E 的自同态所成的空间。

设 \(\varepsilon > 0\)。取 \(r > 0\) 使得 \(\mathrm{X} \subset ] - r, r[\) ，并取 \(\ell \geqslant 1\) 使得 \(r2^{-\ell} < \varepsilon\) 。对每个整数 \(n \in \mathbf{N}\) 以及每个整数 \(j\) ，其中 \(|j| \leqslant 2^{n + \ell}\) ，令

\[
\mathrm{X} _ {n, j} = \mathrm{X} \cap \Big ] \frac {r j}{2 ^ {n + \ell}}, \frac {r (j + 1)}{2 ^ {n + \ell}} \Big ]
\]

\begin{enumerate}
\def\labelenumi{\alph{enumi})}
\item
  对每个整数 \(n \in N\)，族 \(\mathcal{X}_{n} = (\mathrm{X}_{n,j})_{j}\) 是 X 的有限分划；诸集合 \(X_{n,j}\) 是可测的且直径 \(\leqslant r2^{-n-\ell}\) 。对每个 \(n \in N\)，分划 \(X_{n+1}\) 比分划 \(X_{n}\) 更细。因此我们可以对 \((\mathcal{X}_{n})\) 应用上一个习题的结果并沿用其记号，特别是空间 \(E_{n}\) 与 \(F_{n}\) 以及正交投影算子 \(p_{n}\) 与 \(q_{n}\) 。诸空间 \(E_{n}\) 是有限维的。
\item
  级数
\end{enumerate}

\[
v = \sum_ {n \in \mathbf {N}} q _ {n} u q _ {n}
\]

在空间 \(\mathcal{L}_{s}(\mathrm{E})\) 中收敛；存在 E 的标准正交基 B 使得 \(v \in \mathcal{D}_{\mathrm{B}}(\mathrm{E})\) 。

\begin{enumerate}
\def\labelenumi{\alph{enumi})}
\setcounter{enumi}{2}
\tightlist
\item
  自同态 \(w = v - u\) 是紧的。（考虑乘法自同态 \(u_n\) ，它由函数 \(g_n \in \mathrm{E}_n\) 给出，使得对每个整数 \(j\) （\(|j| \leqslant 2^{n + \ell}\) ）以及所有 \(x \in \mathrm{X}_{n,j}\) ，有 \(g_n(x) = rj2^{-n - \ell}\) ；然后验证 \(\| u - u_n \| \leqslant r2^{-n - \ell}\) 以及 \(u_n p_n = p_n u_n\) 。）
\end{enumerate}

\begin{enumerate}
\def\labelenumi{\arabic{enumi})}
\setcounter{enumi}{19}
\tightlist
\item
  设 E 是可数型复希尔伯特空间，\(u \in \mathcal{L}(\mathrm{E})\) 是埃尔米特自同态。设 \(\varepsilon\) 是严格正的实数。存在 E 的标准正交基 B、埃尔米特对角自同态 \(v \in \mathcal{D}_{\mathrm{B}}(\mathrm{E})\) 以及紧自同态 \(w \in \mathcal{L}^{\mathrm{c}}(\mathrm{E})\) ，使得 \(u = v + w\) 且 \(\|w\| < \varepsilon\) （首先考虑 u 有循环向量的情形，然后应用上一个习题。）
\end{enumerate}

21) 设 E 是可数型复希尔伯特空间，\(u \in \mathcal{L}(\mathrm{E})\) 是正规自同态。设 \(\varepsilon\) 是严格正的实数。存在 E 的标准正交基 B、对角自同态 \(v \in \mathcal{D}_{\mathrm{B}}(\mathrm{E})\) 以及紧自同态 \(w \in \mathcal{L}^{\mathrm{c}}(\mathrm{E})\) ，使得 \(u = v + w\) 且 \(\|w\| < \varepsilon\) 。

\begin{enumerate}
\def\labelenumi{\arabic{enumi})}
\setcounter{enumi}{21}
\tightlist
\item
  设 G 是局部紧交换拓扑群，\(\mu\) 是 G 上的哈尔测度。
\end{enumerate}

设 \(u \in \mathcal{L}(\mathrm{L}^2(\mathrm{G},\mu))\) 满足：\(u\) 与乘法算子 \(m_{\chi}\) 对每个特征标 \(\chi \in \widehat{\mathrm{G}}\) 交换。证明存在 \(g \in \mathrm{L}^{\infty}(\mathrm{G},\mu)\) 使得 \(u = m_g\) 。

\begin{enumerate}
\def\labelenumi{\arabic{enumi})}
\setcounter{enumi}{22}
\tightlist
\item
  设 \(u\) 是 E 的自同态，并设 A 是 \(\mathcal{L}(\mathrm{E})\) 的由 \(u\) 生成的闭含单位元子代数。
\end{enumerate}

若 v 是 E 的自同态且 \(u \circ v - v \circ u\) 属于 A，则 \(u \circ v - v \circ u\) 是拟幂零的。特别地，若 E 非零，我们有 \(u \circ v - v \circ u \neq 1_{E}\) 。（利用 I, p. 167 的习题 10；注意该结果对任何巴拿赫空间 E 都成立。）

\begin{enumerate}
\def\labelenumi{\arabic{enumi})}
\setcounter{enumi}{23}
\tightlist
\item
  设 \(\mathrm{A} \subset \mathcal{L}(\mathrm{E})\) 是 \(\mathcal{L}(\mathrm{E})\) 的含单位元的自伴子代数。我们用 \(\mathrm{A}''\) 表示 A 的双交换子。我们用 \(\mathcal{L}(\mathrm{E})_s\) 表示赋有逐点收敛拓扑的空间 \(\mathcal{L}(\mathrm{E})\) 。
\end{enumerate}

\begin{enumerate}
\def\labelenumi{\alph{enumi})}
\item
  空间 \(A''\) 是 \(\mathcal{L}(\mathrm{E})\) 的含单位元的自伴子代数，并且它在 \(\mathcal{L}(\mathrm{E})_s\) 中是闭的。
\item
  设 \(v \in A''\) 且 \(x \in E\)。向量 \(v(x)\) 属于诸 \(u(x)\) （\(u \in A\)）所成之集在 E 中的闭包。（设 F 是该闭包；注意 E 的以 F 为像的正交投影算子属于 \(A'\)。）
\item
  代数 \(A''\) 与 A 在 \(\mathcal{L}(\mathrm{E})_s\) 中的闭包重合。（对 E 的元素的每个有限族 \((x_i)_{i \in \mathrm{I}}\) ，把上一问题应用于空间 \(\mathrm{E}^{\mathrm{I}}\) 、像 \(\widetilde{\mathrm{A}}\) （即 A 在含单位元对合态射 \(f: \mathcal{L}(\mathrm{E}) \to \mathcal{L}(\mathrm{E}^{\mathrm{I}})\) （其中 \(f(u)(y_i) = (u(y_i))_{i \in \mathrm{I}}\) ）下的像），以及元素 \((x_i)\) （它属于 \(\mathrm{E}^{\mathrm{I}}\) ）。）
\end{enumerate}

\(A''\) 与 A 在 \(\mathcal{L}(\mathrm{E})_s\) 中闭包的相等就是冯·诺伊曼双交换子定理。

\begin{enumerate}
\def\labelenumi{\arabic{enumi})}
\setcounter{enumi}{24}
\tightlist
\item
  设 U 是 R 的开子集。普遍可测函数 \(f: U \to R\) 称为勒夫纳函数，如果对每个复希尔伯特空间 E 以及 E 的满足 \(\operatorname{Sp}(u) \subset \operatorname{U}\) 与 \(\operatorname{Sp}(v) \subset \operatorname{U}\) 的埃尔米特自同态 u 与 v，条件 \(u \leqslant v\) 在 \(\mathcal{L}(\operatorname{E})\) 中蕴含 \(f(u) \leqslant f(v)\) 。我们用 \(\operatorname{Loe}(\operatorname{U})\) 表示 U 上勒夫纳函数所成的集。
\end{enumerate}

\begin{enumerate}
\def\labelenumi{\alph{enumi})}
\item
  设 \(f \in \text{Loe}(U)\)。函数 f 在 U 上递增。常值函数 1 与 U 的恒等函数都属于 \(\text{Loe}(U)\) 。
\item
  设 \(f \in \text{Loe}(\mathbf{R}_+^*)\)。设 E 是复希尔伯特空间，u 是 E 的正自同态。对 E 的每个正交投影算子 p 以及每个 \(\varepsilon > 0\) ，我们有 \(f(u)_p \leqslant f((1 + \varepsilon)u_p)\) ，其中用 \(v_p\) 表示 \(\text{Im}(p)\) 的由 \(x \mapsto p(v(x))\) 定义的自同态（对所有 \(v \in \mathcal{L}(\text{E})\) ）。（把 E 等同于希尔伯特直和 \(\text{Im}(p) \oplus \text{Im}(1_{\text{E}} - p)\) ，并将 u 表示为矩阵 \(\begin{pmatrix} a & b \\ c & d \end{pmatrix}\)。注意 \(u \leqslant \begin{pmatrix} (1 + \varepsilon)a & 0 \\ 0 & (1 + \varepsilon^{-1})d \end{pmatrix}\) 对所有 \(\varepsilon > 0\) 成立。）
\item
  设 \(f \in \text{Loe}(\mathbf{R}_{+}^{*})\)。函数 f 在 \(R_{+}^{*}\) 上是凹的（设 v 是 \(C^{2}\) 的在典范基下为对角的自同态，并设 \(g \in \text{SO}(\mathbf{R}^{2})\) 视为 \(C^{2}\) 的自同态。把上一问题应用于以 \(Ce_{1}\) 为像的正交投影算子且取 \(u = g^{*}vg\) 。）
\item
  集合 \(\operatorname{Loe}(\mathbf{R})\) 是形如 \(f(x)=ax+b\) （其中 \((a,b)\in\mathbf{R}_{+}\times\mathbf{R}\) ）的函数所成的集。（把上一问题既应用于 f，也应用于函数 \(x\mapsto-f(-x)\) 。）
\item
  设 \(t \in \mathbf{R}_{+}^{*}\) 。函数 \(f: x \mapsto -(t + x)^{-1}\) 属于 \(\operatorname{Loe}(\mathbf{R}_{+}^{*})\)。
\end{enumerate}

\(f)\) 对任何有界正测度 \(\nu\) （在 \(\mathbf{R}_{+}^{*}\) 上且使函数 \(x\mapsto x^{-1}\) 为 \(\nu\) -可积），函数 \(f\colon\mathbf{R}_{+}^{*}\to\mathbf{R}\) 由

\[
f (x) = \int_ {\mathbf {R} _ {+} ^ {*}} \frac {x}{x + t} d \nu (t)
\]

所给出的函数属于 \(\text{Loe}(\mathbf{R}_{+}^{*})\)；它在 \(R_{+}^{*}\) 上可微，并且还满足下列条件：

\begin{enumerate}
\def\labelenumi{(\roman{enumi})}
\item
  \(\lim_{t\to0}f(t)=0;\)
\item
  \(\lim_{t\to0}f'(t)\) 在 R 中存在；
\item
  f 在 \(+\infty\) 处有有限极限。
\end{enumerate}

\begin{enumerate}
\def\labelenumi{\arabic{enumi})}
\setcounter{enumi}{25}
\tightlist
\item
  设 E 是复希尔伯特空间，a 是 E 的正自同态。我们记 \(\|x\|_{a} = \langle x | a(x) \rangle\) ；它是 E 上的半范数。我们同样用 \(\|u\|_{a}\) 表示 E 赋有该半范数时其自同态 u 的半范数。
\end{enumerate}

我们称函数 \(f \colon \mathbf{R}_+^* \to \mathbf{R}_+^*\) 为插值函数，如果对每个复希尔伯特空间 E、E 的每个正自同态 \(a\) 以及每个 \(u \in \mathcal{L}(\mathrm{E})\) ，条件 \(\| u \| \leqslant 1\) 与 \(\| u \|_a \leqslant 1\) 蕴含 \(\| f(u) \|_a \leqslant 1\) 。

\begin{enumerate}
\def\labelenumi{\alph{enumi})}
\item
  函数 \(f \colon \mathbf{R}_+^* \to \mathbf{R}_+^*\) 是插值函数，当且仅当对每个复希尔伯特空间 E、E 的每个正自同态 \(a\) 以及每个 \(u \in \mathcal{L}(\mathrm{E})\) ，条件 \(u^*u \leqslant 1_{\mathrm{E}}\) 与 \(u^*au \leqslant a\) 蕴含 \(u^*f(a)u \leqslant f(a)\) 。
\item
  \(R_{+}^{*}\) 上的任何勒夫纳函数都是插值函数。（给定 E 以及 \(u \in \mathcal{L}(E)\) （范数 \(\leqslant 1\) ），令 \(w = \sqrt{1_{E} - u^{*}u}\) ，\(w' = \sqrt{1_{E}-uu^{*}}\) ；考虑 \(F = E \oplus E\) ，并注意由 \(v(x,y) = (u(x)+w'(y), -w(x)+u(y))\) 定义的 F 的自同态 v 是等距的；然后考虑 \(a' \colon (x,y) \mapsto (a(x),0)\) ，并把上一个习题的 b) 应用于以 \(E \oplus \{0\}\) 为像的正交投影算子。）
\item
  每个插值函数都是 \(\mathbf{R}_{+}^{*}\) 上的勒夫纳函数。（设 \(0 \leqslant u \leqslant v\) 是复希尔伯特空间 E 的自同态且 \(v\) 可逆；令 \(\mathrm{F} = \mathrm{E} \oplus \mathrm{E}\) ，考虑 F 的自同态 \(a = u \oplus v\) 以及自同态 \(w\) ，它由矩阵 \(\left( \begin{array}{cc}0 & 1_{\mathrm{E}}\\ v^{-1/2}u^{1/2} & 0 \end{array} \right)\) 表示；证明 \(w^{*}w \leqslant 1_{\mathrm{F}}\) 且 \(w^{*}aw \leqslant a\) 。）
\item
  设 f 是插值函数。设 E 是复希尔伯特空间；我们记 \(\mathcal{R}(u)=\frac{1}{2}(u+u^{*})\) 对所有 \(u\in\mathcal{L}(\mathrm{E})\) 。对 \(\mathcal{L}(\mathrm{E})\) 中的 u 与 v，证明：\(e^{-tv^{*}}\circ u\circ e^{-tv}\leqslant u\) 对所有 \(t\in R_{+}\) 成立当且仅当 \(\mathcal{R}(u\circ v)\geqslant0\) 。（对给定的 \(x\in E\)，计算函数 \(t\mapsto\langle e^{-tv}x|(ue^{-tv})x\rangle\) 的导数。）
\end{enumerate}

由此推出：对 E 的所有自同态 \(u\) 与 \(v\) ，当条件

\[
u \geqslant 0, \quad \mathcal {R} (v) \geqslant 0, \quad \mathcal {R} (u v) \geqslant 0
\]

得到满足时，我们有 \(\mathcal{R}(f(u)v) \geqslant 0\) 。

27) 我们打算证明习题 25 问题 \(b\) 中论断的逆命题。为此，将利用下面 V, p. 503 习题 23 的结果与记号。特别地，G 表示群 \(\mathbf{R}_{+}^{*}\) ，它赋有哈尔测度 \(\mu = x^{-1}dx\) ，并与群 \(i\mathbf{R}\) 通过映射 \((x, it) \mapsto x^{it}\) 对偶。我们用 X 表示 G 的恒等映射。

我们分别用 co、si 与 ta 表示 iR 上的函数 \(it \mapsto \cos(i\pi t)\)、\(it \mapsto \sin(i\pi t)\) 与 \(it \mapsto \tan(i\pi t)\) 。

设 \(f \in \text{Loe}(\mathbf{G})\) 满足习题 25 问题 \(f\) 的条件 (i)、(ii) 与 (iii)。

\begin{enumerate}
\def\labelenumi{\alph{enumi})}
\item
  设 \(\alpha\) 是由 \(x \mapsto x / (1 + x)\) 给出的 \(G\) 上的函数。我们有 \(\mathrm{I}_{\alpha} = ] - 1,0[\) ，并且 \(\widehat{\alpha}(s) = -\pi / \sin(\pi s)\) 对 \(s \in \mathrm{B}_f\) 成立。
\item
  函数 ta 在 \(i\mathbf{R}\) 上有界。我们令 \(v = F \circ m_{ta} \circ \overline{F}\) ；它是 \(L^{2}(G)\) 的自同态，满足 \(\|v\| \leqslant 1\) 且 \(\mathcal{R}(v) = 0\) 。
\item
  设 \(\varepsilon > 0\) 是实数。我们用 \(u_{\varepsilon}\) 表示由 \(h_{\varepsilon} = \mathrm{X}/(1 + \varepsilon\mathrm{X})\) 给出的 \(\mathrm{L}^2(\mathrm{G})\) 上的乘法自同态。我们有 \(u_{\varepsilon} \geqslant 0\) ，并且对每个 \(\varphi \in \mathrm{L}^2(\mathrm{G})\) 以及每个有界连续函数 \(f\) （在 \(\mathrm{G}\) 上），有
\end{enumerate}

\[
\langle \varphi | \mathcal {R} (f (u _ {\varepsilon}) v) \varphi \rangle = \frac {1}{2} \langle f \circ h _ {\varepsilon} | \mathrm{Q} (\varphi) \rangle
\]

\(\mathrm{Q}(\varphi) = \overline{\varphi} v(\varphi) + \overline{v(\varphi)}\varphi .\)

\begin{enumerate}
\def\labelenumi{\alph{enumi})}
\setcounter{enumi}{3}
\item
  设 F 是 \(\mathrm{L}^{2}(i\mathbf{R})\) 的由诸函数 \(t \mapsto e^{\alpha(it - s_{0})^{2}}\) （其中 \(\alpha \in R_{+}^{*}\) ，\(s_{0} \in C\) ）生成的向量子空间；它在 \(\mathrm{L}^{2}(i\mathbf{R})\) 中稠密。
\item
  设 \(\psi \in \mathrm{F}\) 。存在 \(\varphi \in \mathrm{L}^2 (\mathbf{G})\) 使得 \(\widehat{\varphi} = \mathrm{co}\cdot \psi\)。令 \(g = \mathrm{Q}(\varphi)\)。我们有 \(g\in \mathrm{L}^{1}(\mathrm{G})\) ，并且存在 \(\mathbf{C}\) 上的全纯函数，它到 \(i\mathbf{R}\) 的限制与 \(\widehat{g}\) 一致。此外
\end{enumerate}

\[
\widehat {g} = \frac {1}{2 i \pi} (\psi^ {*} * \psi) \mathrm{si}.
\]

\(f)\) 对每个 \(\varepsilon > 0\) ，我们有 \(\mathcal{R}(u_{\varepsilon} \circ v) \geqslant 0\) 。（计算 \(\langle \varphi | \mathcal{R}(u_{\varepsilon} u) \varphi \rangle\) ，当 \(\varphi\) 是上一问题中那样的函数时；并像 V, p. 503 习题 23 的问题 \(d\) 中那样使用 G 上的卷积。）

\(g)\) 设 \(f \in \mathrm{Loe}(\mathbf{R}_{+}^{*})\)。设 \(g\) 是与 \(\psi \in F\) 如问题 \(e)\) 中那样相伴的函数。我们则有

\[
\int_ {\mathrm{G}} f g d \mu \geqslant 0
\]

以及

\[
\int_ {\mathbf {R}} \widehat {f} (- \sigma - i t) \sin (\pi (\sigma + i t)) (\psi^ {*} * \psi) (\sigma + i t) d t \geqslant 0
\]

对所有 \(\sigma\in]0,1[\) 成立。

\begin{enumerate}
\def\labelenumi{\alph{enumi})}
\setcounter{enumi}{7}
\tightlist
\item
  设 \(\sigma \in ]0,1[\) 。函数 \(q\) 由 \(q(s) = \widehat{f}(-s)\sin(\pi s)\) （对 \(s\) 满足 \(\mathcal{R}(s) = \sigma\) ）定义，它满足
\end{enumerate}

\[
\sum_ {i = 0} ^ {n} \sum_ {j = 0} ^ {n} \bar {t} _ {i} t _ {j} q \left(\frac {1}{2} (\bar {z} _ {i} + z _ {j})\right) \geqslant 0
\]

对所有 \(n \in N\) 以及所有复数族 \((t_{i})_{0 \leqslant i \leqslant n}\) 与 \((z_{i})_{0 \leqslant i \leqslant n}\) （它们满足对每个 i 有 \(\mathcal{R}(z_{i}) = \sigma\) ）成立。

\begin{enumerate}
\def\labelenumi{\roman{enumi})}
\tightlist
\item
  存在 G 上的正测度 \(\nu_0\) ，使得 \(\widehat{\nu}_0(it) = \widehat{f}(-it)\sin(i\pi t)/\pi\) 对所有 \(t\in\mathbf{R}\) 成立；我们有 \(\mathrm{X}^{\sigma}\in\mathrm{L}^{1}(\mathrm{G},\nu_{0})\) ，若 \(0<\sigma<1\) 。
\end{enumerate}

\begin{enumerate}
\def\labelenumi{\alph{enumi})}
\setcounter{enumi}{9}
\tightlist
\item
  设 \(\nu = \check{\nu}_0\)。测度 \(\nu\) 是有界的。（先验证 \(]-1,0[ \subset I_{\nu}\) ，再验证 \(\widehat{\nu}(-t)\) 当 \(t \to 0\) 时有界，以此证明 \(\nu([1,+\infty[)\) 有限。）对每个 \(x\) ，我们有
\end{enumerate}

\[
f (x) = \int_ {\mathrm{G}} \frac {x}{x + t} d \nu (t).
\]

\begin{enumerate}
\def\labelenumi{\arabic{enumi})}
\setcounter{enumi}{27}
\tightlist
\item
  \begin{enumerate}
  \def\labelenumii{\alph{enumii})}
  \tightlist
  \item
    设 \(D \subset C\) 是开单位圆盘。设 f 是 \(\overline{D}\) 的一个邻域中的全纯函数的芽。对每个 \(z \in D\) ，我们有
  \end{enumerate}
\end{enumerate}

\[
f (z) = i \mathcal {I} (f (0)) + \frac {1}{2 \pi} \int_ {0} ^ {2 \pi} \mathcal {R} (f (e ^ {i \theta})) k (e ^ {i \theta}, z) d \theta
\]

其中 \(k(z, w) = (z + w) / (z - w)\) 对 \((z, w) \in \mathrm{D} \times \mathrm{D}\) 。（注意

\[
k (e ^ {i \theta}, z) = 1 + 2 \sum_ {n = 1} ^ {+ \infty} z ^ {n} e ^ {- i n \theta}
\]

并使用 f 在 0 附近的泰勒展开。）

\begin{enumerate}
\def\labelenumi{\alph{enumi})}
\setcounter{enumi}{1}
\tightlist
\item
  设 \(f\) 是 D 上的全纯函数，使得 \(\mathcal{R}(f(z)) > 0\) 对所有 \(z \in \mathrm{D}\) 成立。测度 \(\mu_r = \frac{1}{2\pi} \mathcal{R}(f(re^{i\theta})) d\theta\) （\([0, 2\pi]\) 上的）当 \(r \to 1\) 时淡收敛于一个有界正测度 \(\nu\) ，并且我们有
\end{enumerate}

\[
f (z) = i \mathcal {I} (f (0)) + \frac {1}{2 \pi} \int_ {0} ^ {2 \pi} \mathcal {R} (f (e ^ {i \theta})) k (e ^ {i \theta}, z) d \nu (\theta)
\]

对所有 \(z \in D\) 成立。（化归到 \(f(0) = 1\) 的情形，并注意这时诸测度 \(\mu_{r}\) 是总质量为 1 的正测度。）

\begin{enumerate}
\def\labelenumi{\alph{enumi})}
\setcounter{enumi}{2}
\tightlist
\item
  设 H 是由这样的 \(z \in C\) 组成的 C 的开子集：\(\mathcal{I}(z) > 0\)。设 \(f: H \to C\) 是全纯函数。f 的像含于 H 当且仅当存在 R 上的有界正测度 \(\nu\) 以及实数 \(a \geqslant 0\) 使得
\end{enumerate}

\[
f (z) = \mathcal {R} (f (i)) + a z + \int_ {\mathbf {R}} \frac {1 + x z}{x - z} d \nu (x)
\]

对所有 \(z \in \mathbf{H}\) 成立。

\begin{enumerate}
\def\labelenumi{\alph{enumi})}
\setcounter{enumi}{3}
\tightlist
\item
  设 I 是 \(\mathbf{R}\) 中的非空开区间，并设 \(f\colon \mathrm{I}\to \mathbf{R}\) 。记 \(\mathrm{J} = \mathbf{R} - \mathrm{I}\) 。下列条件等价（“勒夫纳定理”）：
\end{enumerate}

\begin{enumerate}
\def\labelenumi{(\roman{enumi})}
\item
  函数 \(f\) 属于 \(\operatorname{Loe}(\mathbf{I})\) ；
\item
  存在解析函数 \(\widetilde{f}\) ，它定义在 \((\mathbf{C} - \mathbf{R}) \cup I\) 上、到 I 的限制与 \(f\) 一致，并且满足 \(\mathcal{I}(f(z)) > 0\) 若 \(\mathcal{I}(z) > 0\) ；
\item
  存在 J 上的有界测度 \(\nu\) 以及偶 \((a,b)\in\mathbf{R}_{+}\times\mathbf{R}\) 使得
\end{enumerate}

\[
f (x) = a x + b + \int_ {\mathrm{J}} \frac {1 + x y}{y - x} d \nu (y)
\]

对所有 \(x \in I\) 成立。

（首先直接验证 (iii) 蕴含 (i) 与 (ii)；为证明 (i) 蕴含 (iii)，单独处理 I = R 的情形，并把其他情形化归为 I = R \(_{+}^{*}\) 的情形，对此应用上一个习题；最后，为证明 (ii) 蕴含 (iii)，应用 c) 部分并验证由此给出的 R 上的测度 \(\nu\) 的支撑含于 J。）

\exercisesection{}

\begin{enumerate}
\def\labelenumi{\arabic{enumi})}
\item
  函数 \(\varphi_0\) 在 \(\mathbf{R}\) 上由 \(\varphi_0(x) = 0\) （若 \(x \leqslant 0\) ）与 \(\varphi_0(x) = e^{-1/x}\) （若 \(x > 0\) ）定义，它是 \(\mathrm{C}^\infty\) 类的，并且满足 \(\varphi_0^{(k)}(0) = 0\) 对每个整数 \(k \in \mathbf{N}\) 。
\item
  设 \(x_0 \in \mathbf{R}^n\) ，\(r > 0\) 且 \(t > 1\) 。存在 \(\varphi \in \mathcal{D}(\mathbf{R}^n)\) 使得
\end{enumerate}

\begin{enumerate}
\def\labelenumi{(\roman{enumi})}
\item
  我们有 \(0 \leqslant \varphi \leqslant 1\) ；
\item
  \(\varphi(x)=1\) 对 \(\|x-x_{0}\|\leqslant r\) 成立，且 \(\varphi(x)=0\) 对 \(\|x-x_{0}\|\geqslant tr\) 成立；
\item
  对每个 \(\alpha \in \mathbf{N}^n\) ，我们有
\end{enumerate}

\[
\sup _ {x \in \mathbf {R} ^ {n}} | \partial^ {\alpha} \varphi (x) | \leqslant \alpha ! ((t - 1) r) ^ {| \alpha |}.
\]

\begin{enumerate}
\def\labelenumi{\arabic{enumi})}
\setcounter{enumi}{2}
\tightlist
\item
  设 \(n \in \mathbf{N}\) 。我们用 \(\mu_n\) 表示 \(\mathbf{R}^n\) 上的勒贝格测度。
\end{enumerate}

\begin{enumerate}
\def\labelenumi{\alph{enumi})}
\tightlist
\item
  设 \(\mathcal{H}_1\) 是映射 \(\partial_1\colon \mathcal{D}(\mathbf{R}^n)\to \mathcal{D}(\mathbf{R}^n)\) 的像；它是满足如下条件的 \(\varphi \in \mathcal{D}(\mathbf{R}^n)\) 所成的空间：
\end{enumerate}

\[
\int_ {\mathbf {R}} \varphi (x, y) d \mu_ {1} (x) = 0
\]

对所有 \(y \in \mathbf{R}^{n-1}\) 成立。

\begin{enumerate}
\def\labelenumi{\alph{enumi})}
\setcounter{enumi}{1}
\tightlist
\item
  设 \(\varphi_0\in \mathcal{D}(\mathbf{R})\) 使其在 \(\mathbf{R}\) 上的积分等于 1。对每个 \(\varphi \in \mathcal{D}(\mathbf{R}^n)\) ，存在 \(\varphi_{1}\in \mathcal{H}_{1}\) 使得
\end{enumerate}

\[
\varphi (x, y) = \varphi_ {1} (x, y) + \varphi_ {0} (x) \int_ {\mathbf {R}} \varphi (t, y) d \mu_ {1} (t)
\]

对所有 \((x,y)\in \mathbf{R}\times \mathbf{R}^{n - 1}\) 成立。

\begin{enumerate}
\def\labelenumi{\alph{enumi})}
\setcounter{enumi}{2}
\tightlist
\item
  对每个广义函数 \(f \in \mathcal{D}'(\mathbf{R}^{n})\) ，存在 \(g \in \mathcal{D}'(\mathbf{R}^{n})\) 使得 \(\partial_{1}g = f\) 。若 \(g_{1}\) 与 \(g_{2}\) 属于 \(\mathcal{D}'(\mathbf{R}^{n})\) 且满足 \(\partial_{1}g_{1} = \partial_{1}g_{2} = f\) ，则存在广义函数 \(h \in \mathcal{D}'(\mathbf{R}^{n-1})\) 使得
\end{enumerate}

\[
\langle g _ {1} - g _ {2}, \varphi \rangle = \langle h, \psi \rangle
\]

对所有 \(\varphi\in\mathcal{D}(\mathbf{R}^{n})\) 成立，其中 \(\psi\in\mathcal{D}(\mathbf{R}^{n-1})\) 由下式定义：

\[
\psi (y) = \int_ {\mathbf {R}} \varphi (t, y) d \mu_ {1} (t).
\]

\(d)\) 设 U 是 \(\mathbf{R}^{n}\) 的连通开子集。设 \(f \in \mathcal{D}'(\mathrm{U})\) 是满足 \(\partial_{i}f = 0\) （对 \(1 \leqslant i \leqslant n\) ）的广义函数。广义函数 \(f\) 是与一个常值函数相伴的广义函数。

\begin{enumerate}
\def\labelenumi{\arabic{enumi})}
\setcounter{enumi}{3}
\tightlist
\item
  我们用 \(\mu\) 表示 \(\mathbf{R}\) 上的勒贝格测度。
\end{enumerate}

\begin{enumerate}
\def\labelenumi{\alph{enumi})}
\tightlist
\item
  设 \(\nu\) 是 \(\mathbf{R}\) 上的测度，\(a \in \mathbf{R}\) 且 \(\nu(\{a\}) = 0\) 。对 \(x \in \mathbf{R}\) 我们令
\end{enumerate}

\[
f (x) = \left\{ \begin{array}{l l} \nu ([ a, x ]) & \text { if } x \geqslant a, \\ \nu ([ x, a ]) & \text { if } x \leqslant a. \end{array} \right.
\]

函数 f 是连续的，并且与之相伴的广义函数满足 \(f' = \nu\) 。

\begin{enumerate}
\def\labelenumi{\alph{enumi})}
\setcounter{enumi}{1}
\item
  设 f 是 R 上的广义函数，使得 \(f'\) 是一个测度 \(\nu\) 。广义函数 f 与 R 上按上一问题方式定义的连续函数相伴。
\item
  设 \(f\) 是 \(\mathbf{R}\) 上的广义函数，使得 \(f^{(k)}\) 对每个整数 \(k \geqslant 0\) 都是测度。那么 \(f\) 是与 \(\mathcal{C}^{\infty}(\mathbf{R})\) 中的函数相伴的广义函数。
\end{enumerate}

\begin{enumerate}
\def\labelenumi{\arabic{enumi})}
\setcounter{enumi}{4}
\item
  与点测度 \(\varepsilon_{n}\) （\(n\in\mathbf{Z}\) ）相伴的广义函数族在 \(\mathcal{S}^{\prime}(\mathbf{R})\) 中可和。其和 \(\omega\) 满足 \(\mathcal{F}(\omega)=\omega\) 。（该公式等价于关于赋有勒贝格测度的局部紧群 \(\mathrm{G}=\mathbf{R}^{n}\) 与子群 \(\mathrm{H}=\mathbf{Z}^{n}\) 的泊松公式，参见 II, § 1, 第 7 目。）
\item
  设 A 是 C 上的交换含单位元代数，满足下列条件：
\end{enumerate}

\begin{enumerate}
\def\labelenumi{(\roman{enumi})}
\item
  代数 \(\mathcal{C}(\mathbf{R})\) 是 A 的含单位元子代数；
\item
  存在导子 \(d\colon\mathrm{A}\to\mathrm{A}\) ，它到 \(\mathcal{C}^{1}(\mathbf{R})\) 的限制与函数的通常导子一致。
\end{enumerate}

\begin{enumerate}
\def\labelenumi{\alph{enumi})}
\item
  R 的恒等映射 X 在 A 中可逆。（证明映射 \(g: \mathbf{R} \to \mathbf{R}\) （它满足 \(g(0) = 0\) ，并且 \(g(x) = x(\log(|x|) - 1)\) 对 \(x \neq 0\) ）是连续的且满足 \(gX = 1\) 。）
\item
  我们有 \(d \circ d(|\mathrm{X}|) = 0\) 。
\item
  不存在双线性映射 \(m\colon\mathcal{D}^{\prime}(\mathbf{R})\times\mathcal{D}^{\prime}(\mathbf{R})\to\mathcal{D}^{\prime}(\mathbf{R})\) 使得 \(m(f,g)=fg\) 对每个偶 \((f,g)\in\mathcal{D}(\mathbf{R})^{2}\) 成立，并且使得
\end{enumerate}

\[
m (f, g) ^ {\prime} = m \left(f ^ {\prime}, g\right) + m \left(f, g ^ {\prime}\right)
\]

对所有 \((f,g)\in\mathcal{D}^{\prime}(\mathbf{R})^{2}\) 成立（“广义函数乘法的不可能性”）。

\begin{enumerate}
\def\labelenumi{\alph{enumi})}
\setcounter{enumi}{3}
\tightlist
\item
  给出 R 上试验函数的序列 \((f_{n})\) 与 \((g_{n})\) 的例子，使得 \(f_{n} \to \varepsilon_{0}\) 且 \(g_{n} \to \varepsilon_{0}\) 在 \(\mathcal{D}'(\mathbf{R})\) 中，并且 \(f_{n}^{2}\) 与 \(g_{n}^{2}\) 在 \(\mathcal{D}'(\mathbf{R})\) 中收敛于不同的极限。
\end{enumerate}

\begin{enumerate}
\def\labelenumi{\arabic{enumi})}
\setcounter{enumi}{6}
\tightlist
\item
  设 \(n \in N\) 。我们用 \(\widetilde{\gamma}\) 表示 \(R^{n}\) 在从 \(R^{n}\) 到 C 的函数所成空间上的线性表示，它由下式定义：
\end{enumerate}

\[
\widetilde {\gamma} (h) f (x) = f (x + h)
\]

对所有 \(h \in \mathbf{R}^n\) 以及每个函数 \(f: \mathbf{R}^n \to \mathbf{C}\) 成立。

\begin{enumerate}
\def\labelenumi{\alph{enumi})}
\item
  对每个 \(h \in R^{n}\) ，映射 \(\widetilde{\gamma}(h)\) 通过向子空间的过渡定义了局部凸空间 \(\mathcal{D}(\mathbf{R}^{n})\) 的一个自同构，它仍记作 \(\widetilde{\gamma}(h)\) 。
\item
  我们用 \(\pmb{\gamma}\) 表示 \(\mathcal{D}'(\mathbf{R}^n)\) 的与 \(\widetilde{\pmb{\gamma}}(-h)\) 转置相应的自同构。若 \(f \in \mathcal{D}(\mathbf{R}^n)\) ，则 \(\pmb{\gamma}(h)f = \widetilde{\pmb{\gamma}}(h)f\) 对所有 \(h \in \mathbf{R}^n\) 成立。
\item
  映射 \(h \mapsto \gamma(h)\) 是 \(\mathbf{R}^{n}\) 在 \(\mathcal{D}'(\mathbf{R}^{n})\) 中的线性表示。
\end{enumerate}

\(d)\) 对每个 \(h \in \mathbf{R}^n\) 以及每个 \(\alpha \in \mathbf{N}^n\) ，我们有 \(\partial^\alpha \circ \boldsymbol{\gamma}(h) = \boldsymbol{\gamma}(h) \circ \partial^\alpha\) 。

\begin{enumerate}
\def\labelenumi{\alph{enumi})}
\setcounter{enumi}{4}
\tightlist
\item
  设 i 是满足 \(1 \leqslant i \leqslant n\) 的整数，\(e_{i}\) 是第 \(i\) 个向量，属于 \(R^{n}\) 的典范基。对每个广义函数 \(f \in \mathcal{D}'(\mathbf{R}^{n})\) ，我们有
\end{enumerate}

\[
\lim_{\substack{h\to 0\\ h\neq 0}}\frac{1}{h} (\boldsymbol {\gamma}(he_{i})f - f) = \partial_{i}f
\]

在赋有弱拓扑的 \(\mathcal{D}'(\mathbf{R}^n)\) 中成立。

\begin{enumerate}
\def\labelenumi{\arabic{enumi})}
\setcounter{enumi}{7}
\item
  设 \(n \in \mathbf{N}\) ，U 是 \(R^{n}\) 的开子集。广义函数 \(f \in \mathcal{D}'(U)\) 称为正的，如果 \(\langle f, \varphi \rangle \geqslant 0\) 对每个 \(\varphi \in \mathcal{D}(U)\) （它满足 \(\varphi \geqslant 0\) ）成立。证明 f 是正的当且仅当 f 是与 U 上一个正测度相伴的广义函数。
\item
  设 \(n \in \mathbf{N}\) ，\(U\) 是 \(\mathbf{R}^n\) 的开子集。设 \(k \in \mathbf{N}\) 且 \(p \in ]1, +\infty[\) 。空间 \(W^{k,p}(U)\) 是自反的。
\item
  设 \(n \in \mathbf{N}\) ，\(U\) 是 \(\mathbf{R}^n\) 的开子集。设 \(k \in \mathbf{N}\) 且 \(p \in [1, +\infty[\)。设 \(V \subset U\) 是相对紧的开子集。
\end{enumerate}

设 \(\varphi \in \mathcal{D}(\mathbf{R}^n)\) 是正函数，其支撑含于 \(\mathbf{R}^n\) 的单位球且积分为 1。对每个 \(\varepsilon > 0\) 以及 \(x \in \mathbf{R}^n\) ，我们令 \(\varphi_{\varepsilon}(x) = \varepsilon^{-n}\varphi(x/\varepsilon)\)。设 \(f \in \mathrm{W}^{k,p}(\mathrm{U})\) 。

\begin{enumerate}
\def\labelenumi{\alph{enumi})}
\item
  \(f_{\mathrm{V}}\) （即 \(f\) 到 V 的限制）属于 \(\mathrm{W}^{k,p}(\mathrm{V})\) 。
\item
  设 \(\widetilde{f}\) 是 \(\mathbf{R}^n\) 上在 U 上与 \(f\) 一致且在 U 外为零的函数。当 \(\varepsilon \to 0\) 时，\(\widetilde{f} * \varphi_{\varepsilon}\) 到 V 的限制属于 \(\mathrm{W}^{k,p}(\mathrm{V})\) 且收敛于 \(f_{\mathrm{V}}\) 在 \(\mathrm{W}^{k,p}(\mathrm{V})\) 中。
\end{enumerate}

\begin{enumerate}
\def\labelenumi{\arabic{enumi})}
\setcounter{enumi}{10}
\tightlist
\item
  设 \(n \in \mathbf{N}\) ，U 是 \(\mathbf{R}^n\) 的开子集。设 \(k \in \mathbf{N}\) 且 \(p \in [1, +\infty[\) 。我们用 \(\mathrm{H}^{k,p}(\mathrm{U})\) 表示 \(\mathcal{C}^{\infty}(\mathrm{U})\) 在 \(\mathrm{W}^{k,p}(\mathrm{U})\) 中的闭包。
\end{enumerate}

对 \(x \in \mathrm{U}\) ，用 \(\delta(x)\) 表示从 \(x\) 到 \(\mathrm{U}\) 在 \(\mathbf{R}^n\) 中的边界的距离；令 \(\mathrm{U}_{-1} = \mathrm{U}_0 = \varnothing\) 并且

\[
\mathrm{U} _ {k} = \{x \in \mathrm{U} | \| x \| <   k, \delta (x) > k ^ {- 1} \}
\]

对每个整数 \(k \geqslant 1\) 。令 \(\mathrm{V}_k = \mathrm{U}_{k+1} \cap (\mathbf{R}^n - \overline{\mathrm{U}}_{k-1})\) 对 \(k \geqslant 1\)。诸开集 \(\mathrm{V}_k\) 覆盖 U；设 \((\psi_k)_{k \geqslant 1}\) 是从属于覆盖 \((\mathrm{V}_k)_{k \geqslant 1}\) 且由 \(\mathcal{D}(\mathrm{U})\) 中的函数组成的单位分解。

\begin{enumerate}
\def\labelenumi{\alph{enumi})}
\tightlist
\item
  对 \(k \geqslant 1\) ，函数
\end{enumerate}

\[
\varphi_{k} = \sum_{\substack{j\geqslant 1\\ \operatorname{Supp}(\psi_{j})\subset \mathrm{V}_{k}}}\psi_{j}
\]

属于 \(\mathcal{D}(\mathrm{U})\) 且支撑含于 \(\mathrm{V}_k\) ；我们有

\[
\sum_ {k \geqslant 1} \varphi_ {k} = 1
\]

在 U 上成立。

\begin{enumerate}
\def\labelenumi{\alph{enumi})}
\setcounter{enumi}{1}
\tightlist
\item
  设 \(f \in \mathrm{W}^{k,p}(\mathrm{U})\) 且 \(\varepsilon > 0\)。设 \(\varphi\) 以及 \(\varphi_{\varepsilon}\) （对 \(\varepsilon > 0\) ）是如习题 10 中那样的函数。对每个 \(k \geqslant 1\)，存在 \(\varepsilon_{k} > 0\) 使得
\end{enumerate}

\[
\left\| \varphi_ {\varepsilon_ {k}} * (\psi_ {k} f) - \psi_ {k} f \right\| _ {k, p} <   \frac {\varepsilon}{2 ^ {k}}.
\]

\begin{enumerate}
\def\labelenumi{\alph{enumi})}
\setcounter{enumi}{2}
\tightlist
\item
  函数
\end{enumerate}

\[
f _ {\varepsilon} = \sum_ {k \geqslant 1} \varphi_ {\varepsilon_ {k}} * (\psi_ {k} f)
\]

属于 \(\mathcal{C}^{\infty}(\mathrm{U})\) 且 \(\| f - f_{\varepsilon}\|_{k,p} < \varepsilon\) 。

\(d)\). 空间 \(\mathrm{H}^{k,p}(\mathrm{U})\) 等于 \(\mathrm{W}^{k,p}(\mathrm{U})\) 。

\begin{enumerate}
\def\labelenumi{\alph{enumi})}
\setcounter{enumi}{4}
\item
  空间 \(\mathrm{H}^{k,p}(\mathrm{U})\) 可等同于 \(\mathcal{C}^{k}(\mathrm{U})\) 关于范数 \(f \mapsto \|f\|_{k,p}\) 的完备化。
\item
  设 V 是 \(R^{n}\) 的开子集，\(\varphi\colon U\to V\) 是 \(C^{k}\) 类微分同胚。假设 \(\leqslant k\) 阶偏导数——\(\varphi\) 的（相应地 \(\varphi^{-1}\) 的）——在 U 上（相应地在 V 上）有界且一致连续。映射 \(f\mapsto f\circ\varphi\) 是从 \(W^{k,p}(V)\) 到 \(W^{k,p}(U)\) 的拓扑向量空间同构。（利用上一问题的结果。）
\end{enumerate}

\begin{enumerate}
\def\labelenumi{\arabic{enumi})}
\setcounter{enumi}{11}
\tightlist
\item
  \begin{enumerate}
  \def\labelenumii{\alph{enumii})}
  \tightlist
  \item
    对所有整数 \(n \in \mathbf{N}\) 与 \(k \in \mathbf{N}\) ，以及每个实数 \(p \geqslant 1\) ，我们有 \(\mathrm{W}^{k,p}(\mathbf{R}^n) = \mathrm{W}_0^{k,p}(\mathbf{R}^n)\) 。
  \end{enumerate}
\end{enumerate}

\begin{enumerate}
\def\labelenumi{\alph{enumi})}
\setcounter{enumi}{1}
\tightlist
\item
  设 \(p \geqslant 1\) 是实数。设 \(U \subset R^{2}\) 是由 \((x, y) \in \mathbf{R}^{2}\) （满足 0 \textless{} \textbar x\textbar{} \textless{} 1 与 0 \textless{} y \textless{} 1）组成的集。特征函数 \(\varphi\) （\(U \cap (\mathbf{R}_{+}^{*} \times \mathbf{R})\) 的）属于 \(W^{1,p}(U)\) ；若 \(\varepsilon > 0\) 充分小，则不存在 \(C^{1}\) 类的函数 f （在 \(\overline{U}\) 的邻域中）使得 \(\|f - \varphi\|_{1,p} < \varepsilon\) 。
\end{enumerate}

\begin{enumerate}
\def\labelenumi{\arabic{enumi})}
\setcounter{enumi}{12}
\tightlist
\item
  设 U 是 \(\mathbf{R}^n\) 的开子集，\(j\colon \mathrm{U}\to \mathbf{R}^n\) 是典范单射。对每个函数 \(f\colon \mathrm{U}\to \mathbf{C}\) ，我们用 \(j_{!}f\) 表示 \(f\) 到 \(\mathbf{R}^n\) 的零延拓，即从 \(\mathbf{R}^n\) 到 \(\mathbf{C}\) 的、在 U 上与 \(f\) 一致且在 U 外为零的函数。
\end{enumerate}

设 \(k \geqslant 0\) 是整数，\(p \geqslant 1\) 是实数。

\begin{enumerate}
\def\labelenumi{\alph{enumi})}
\item
  设 \(f \in \mathrm{W}^{k,p}(\mathrm{U})\) 。对每个 \(\alpha \in \mathbf{N}^n\) （满足 \(|\alpha| \leqslant k\) ），我们有 \(\partial^\alpha(j_{!}f) = j_{!}(\partial^\alpha(f))\) 在 \(\mathcal{D}'(\mathbf{R}^n)\) 中。
\item
  由向子空间的过渡诱导的映射 \(f \mapsto j_{!}f\) 是从 \(\mathrm{W}_{0}^{k,p}(\mathrm{U})\) 到 \(\mathrm{W}^{k,p}(\mathbf{R}^{n})\) 中的等距线性映射。
\end{enumerate}

\begin{enumerate}
\def\labelenumi{\arabic{enumi})}
\setcounter{enumi}{13}
\tightlist
\item
  设 U 是 \(R^{n}\) 的开子集，\(k \geqslant 1\) 是整数。设 \(p \geqslant 1\) 是实数。假设余集 \(R^{n} - U\) 关于勒贝格测度不是可忽略的。用 j 表示 U 到 \(R^{n}\) 中的典范单射。
\end{enumerate}

\begin{enumerate}
\def\labelenumi{\alph{enumi})}
\item
  存在 \(R^{n}\) 中相对紧的开长方体 B，使得 \(B \cap U\) 与 \(B \cap (\mathbf{R}^{n} - \mathbf{U})\) 都不是可忽略的。
\item
  设 \(\widetilde{f} \in \mathcal{D}(\mathbf{R}^n)\) 在 \(B \cap U\) 上等于 1。设 \(f\) 是 \(\widetilde{f}\) 到 U 的限制。那么 \(f\) 不属于 \(\mathrm{W}_0^{k,p}(\mathrm{U})\) （否则，零延拓 \(j_{!}f\) 将属于 \(\mathrm{W}^{k,p}(\mathbf{R}^n)\) ；\(j_{!}f\) 到 B 的限制将是一个满足 \(\partial_i(j_{!}f) = 0\) （对 \(1 \leqslant i \leqslant n\) ）的广义函数，从而将是与一个常数相伴的广义函数。）
\end{enumerate}

\begin{enumerate}
\def\labelenumi{\arabic{enumi})}
\setcounter{enumi}{14}
\tightlist
\item
  设 \(n \geqslant 1\) 是整数，p 是满足 \(1 \leqslant p < n\) 的实数。
\end{enumerate}

\begin{enumerate}
\def\labelenumi{\alph{enumi})}
\item
  设 \(\alpha\) 是满足 \(0 < \alpha < n/p - 1\) 的实数。设 \(g\) 是 \(\mathbf{R}_{+}^{*}\) 上的无穷可微正函数，使得 \(g(x) = x^{-\alpha}\) （当 \(x\) 充分小）且 \(g(x) = 0\) （当 \(x\) 充分大）；从 \(\mathbf{R}^{n}\) 到 \(\mathbf{C}\) 的、由 \(x \mapsto g(\|x\|)\) 定义的函数属于 \(\mathrm{W}^{1,p}(\mathbf{R}^{n})\) 。
\item
  设 \(i \mapsto q_{i}\) 是从 N 到 \(Q^{n}\) 的双射，f 是由下式定义的从 \(R^{n}\) 到 C 的映射：
\end{enumerate}

\[
f (x) = \sum_ {i \in \mathbf {N}} 2 ^ {- i} g (\| x - q _ {i} \|)
\]

对 \(x \in \mathbf{R}^n\) 成立。定义 \(f\) 的级数在 \(\mathrm{W}^{1,p}(\mathbf{R}^n)\) 中收敛；对 \(\mathbf{R}^n\) 的每个开集 U，\(f\) 到 U 的限制都不是几乎处处可微的。\(c)\) 存在在任一点都不连续的函数 \(f \in \mathrm{W}^{1,n}(\mathbf{R}^n)\) 。（对函数 \(x \mapsto \log(|\log(|x|)|)\) 用类似的方法。）

\begin{enumerate}
\def\labelenumi{\arabic{enumi})}
\setcounter{enumi}{15}
\tightlist
\item
  设 \(n \geqslant 1\) 与 \(k \geqslant 1\) 是整数。设 \(p \geqslant 1\) 是实数。设 U 是 \(\mathbf{R}^n\) 的有界开子集。
\end{enumerate}

\begin{enumerate}
\def\labelenumi{\alph{enumi})}
\item
  假设 \(n \neq kp\) 。设 q 是满足 \(1 \leqslant q \leqslant np/(n-kp)\) 的实数。\(\mathcal{D}(U)\) 到 \(L^{q}(U)\) 中的单射容许一个连续延拓 \(i\) （从 \(W_{0}^{k,p}(U)\) 到 \(L^{q}(U)\) 中）。
\item
  假设 \(1 \leqslant q \leqslant np/(n-kp)\) 。线性映射 \(i\) 是紧的（“雷利希–孔德拉绍夫定理”）。
\item
  假设 \(k \geqslant 1\) 且 \(q = np/(n - p)\) 。线性映射 i 不是紧的。
\end{enumerate}

\begin{enumerate}
\def\labelenumi{\arabic{enumi})}
\setcounter{enumi}{16}
\tightlist
\item
  设 \(n \geqslant 1\) 是整数。我们用 \(\mu\) 表示 \(\mathbf{R}^n\) 上的勒贝格测度。
\end{enumerate}

\begin{enumerate}
\def\labelenumi{\alph{enumi})}
\tightlist
\item
  \(\mathrm{L}^{2}(\mathbf{R}^{n},\mu)\) 的有界子集 B 是紧的，当且仅当对每个 \(\varepsilon>0\) ，存在 \(R^{n}\) 的紧子集 K 使得
\end{enumerate}

\[
\int_ {\mathbf {R} ^ {n} - K} | f | ^ {2} + \int_ {\mathbf {R} ^ {n} - K} | \widehat {f} | ^ {2} <   \varepsilon
\]

对所有 \(f\in \mathrm{B}\) 成立。

\begin{enumerate}
\def\labelenumi{\alph{enumi})}
\setcounter{enumi}{1}
\tightlist
\item
  设 \(k \in N\) 。设 K 是 \(R^{n}\) 的紧子集。我们用 \(\mathrm{H}_{k}^{2}(\mathrm{K})\) 表示索博列夫空间 \(\mathrm{H}_{k}^{2}(\mathbf{R}^{n})\) 的由在紧集 K 外为零的函数组成的闭子空间。典范单射 \(\mathrm{H}_{k}^{2}(\mathrm{K}) \to \mathrm{H}_{k}^{2}(\mathbf{R}^{n})\) 是紧的。
\end{enumerate}

\begin{enumerate}
\def\labelenumi{\arabic{enumi})}
\setcounter{enumi}{17}
\tightlist
\item
  设 \(n \in \mathbf{N}\) ，\(U\) 是 \(\mathbf{R}^n\) 的开子集。本习题的目标是证明 \(\mathcal{D}'(U)\) 是有界型空间。
\end{enumerate}

若 K 是 U 的紧子集，我们给 \(\mathcal{C}_{\mathrm{K}}^{\infty}(\mathrm{U})'\) 赋予有界收敛拓扑。我们用 \(\pi_{\mathrm{K}}\colon\mathcal{D}'(\mathrm{U})\to\mathcal{C}_{\mathrm{K}}^{\infty}(\mathrm{U})'\) 表示把广义函数映为其到 \(\mathcal{C}_{\mathrm{K}}^{\infty}(\mathrm{U})\) 的限制的线性映射；它是连续的。

设 \(S \subset \mathcal{D}'(U)\) 是凸的、均衡的且吸收有界集的子集。

\begin{enumerate}
\def\labelenumi{\alph{enumi})}
\item
  存在 U 的紧子集 K 使得 \(\mathrm{Ker}(\pi_{\mathrm{K}}) \subset \mathrm{S}\) 。（否则，考虑 U 的紧子集的递增序列 \((\mathrm{K}_{m})_{m \in \mathbf{N}}\) ，其诸内部覆盖 U，并选取 \(f_{m}\) 于 \(\mathrm{Ker}(\pi_{\mathrm{K}_{m}}) - \mathrm{S}\) 中；那么诸元素 \(mf_{m}\) （\(m \in N\) ）所成之集是一个不被 S 吸收的有界子集。）
\item
  由广义函数 \(f \in \mathcal{D}'(\mathrm{U})\) （满足 \(f + \operatorname{Ker}(\pi_{\mathrm{K}}) \subset \mathrm{S}\) ）组成的集 T 是凸的、均衡的且吸收有界集的集。
\item
  集 \(\pi_{\mathrm{K}}(\mathrm{T})\) 是 \(\mathcal{C}_{\mathrm{K}}^{\infty}(\mathrm{U})'\) 的凸的、均衡的且吸收有界集的子集。（设 A 是 \(\mathcal{C}_{\mathrm{K}}^{\infty}(\mathrm{U})'\) 的有界子集；证明存在 \(\mathcal{D}(\mathrm{U})\) 中 0 的邻域 W 使得 \(\mathrm{A} \subset (\mathrm{W} \cap \mathcal{C}_{\mathrm{K}}^{\infty}(\mathrm{U}))^{\circ}\) ，并由此推出 A 含于 \(\pi_{\mathrm{K}}(\mathrm{W}^{\circ})\) ，从而被 \(\pi_{\mathrm{K}}(\mathrm{T})\) 吸收。）
\end{enumerate}

\(d)\) 由此推出 S 是 0 的邻域，并断言空间 \(\mathcal{D}'(U)\) 是有界型的。（注意 \(\pi_{\mathrm{K}}(\mathrm{T})\) 是 0 的邻域，并注意 S 包含 \(\pi_{\mathrm{K}}^{-1}(\pi_{\mathrm{K}}(\mathrm{T}))\) 。）

\begin{enumerate}
\def\labelenumi{\arabic{enumi})}
\setcounter{enumi}{18}
\tightlist
\item
  设 \(n \in \mathbf{N}\) ，\(\mathrm{U} \subset \mathbf{R}^n\) 是开集。
\end{enumerate}

\begin{enumerate}
\def\labelenumi{\alph{enumi})}
\tightlist
\item
  设 \(f \in \mathcal{D}'(\mathrm{U})\) 。开集 \(\mathrm{W} \subset \mathrm{U}\) （使得 \(f\) 到 W 的限制为零）的并 V 在 U 中是开的。
\end{enumerate}

f 的支撑称为 V 的余集；它是 U 的闭子集。

\begin{enumerate}
\def\labelenumi{\alph{enumi})}
\setcounter{enumi}{1}
\item
  设 \(f \in \mathcal{D}'(\mathrm{U})\) ，\(\varphi \in \mathcal{D}(\mathrm{U})\) 的支撑与 f 的支撑不相交。我们有 \(\langle f, \varphi \rangle = 0\) 。
\item
  设 \(f \in \mathcal{D}'(\mathrm{U})\) 且 \(\alpha \in N^{n}\) 。\(\partial^{\alpha}f\) 的支撑含于 f 的支撑。
\item
  若 \(f \in \mathcal{D}'(\mathrm{U})\) 是与测度 \(\nu\) 相伴的广义函数，则 f 的支撑与 \(\nu\) 的支撑重合。
\item
  设 \(f\) 是 U 中具有紧支撑的广义函数。线性形式 \(f\) 可延拓为空间 \(\mathcal{C}^{\infty}(\mathrm{U})\) 上的连续线性形式。
\item
  U 中具有紧支撑的广义函数所成的空间可等同于空间 \(\mathcal{C}^{\infty}(\mathrm{U})\) 的对偶。
\end{enumerate}

\(g)\) 设 \(f \in \mathcal{D}'(\mathbf{R}^n)\) 是具有紧支撑的广义函数。我们有 \(f \in \mathcal{S}'(\mathbf{R}^n)\) 。\(h)\) 设 \(f \in \mathcal{D}'(\mathbf{R}^n)\) 是具有紧支撑的广义函数。\(f\) 的傅里叶变换属于 \(\mathcal{C}^\infty(\mathbf{R}^n)\) ；它是连同其所有导数都具有多项式增长的函数。

\begin{enumerate}
\def\labelenumi{\arabic{enumi})}
\setcounter{enumi}{19}
\tightlist
\item
  设 \(n \in \mathbf{N}\) 。我们用 \(V_n\) 表示 \(\mathbf{C}\) 的由复数 \(\lambda\) （它们满足 \(\mathcal{R}(\lambda) > -n - 1\) 且 \(\lambda \notin \{-1, -2, \ldots, -n\}\) ）组成的开集。
\end{enumerate}

设 \(\lambda\in V_{0}\) 。我们用 \(x_{+}^{\lambda}\) 表示 \(\mathbf{R}\) 上的由局部可积函数

\[
x \mapsto \left\{ \begin{array}{l l} 0 & \text { if } x \leqslant 0, \\ x ^ {\lambda} & \text { if } x > 0. \end{array} \right.
\]

\begin{enumerate}
\def\labelenumi{\alph{enumi})}
\item
  映射 \(\lambda \mapsto x_{+}^{\lambda}\) （从 \(\mathrm{V}_0\) 到 \(\mathcal{D}'(\mathbf{R})\) 中的）是解析的。
\item
  设 \(n \in \mathbf{N}\) 。对每个 \(\lambda \in V_n\) ，在 \(\mathcal{D}(\mathbf{R})\) 上由
\end{enumerate}

\[
\begin{array}{l} \varphi \mapsto \int_ {0} ^ {1} x ^ {\lambda} \Big (\varphi (x) - \varphi (0) - x \varphi^ {\prime} (0) - \dots - \frac {x ^ {n - 1}}{(n - 1) !} \varphi^ {(n - 1)} (0) \Big) d x \\ \qquad + \int_ {1} ^ {+ \infty} x ^ {\lambda} \varphi (x) d x + \sum_ {k = 1} ^ {n} \frac {\varphi^ {(k - 1)} (0)}{(k - 1) ! (\lambda + k)} \end{array}
\]

定义的映射是缓增广义函数；它不依赖于使 \(\lambda \in V_{n}\) 的整数 n，并且与 \(x_{+}^{\lambda}\) 一致当 \(\lambda \in V_{0}\) 时。

我们将用 \(x_{+}^{\lambda}\) 表示这样定义的广义函数，对每个 \(\lambda\) 属于诸集 \(V_{n}\) （\(n \in N\) ）之并。

\begin{enumerate}
\def\labelenumi{\alph{enumi})}
\setcounter{enumi}{2}
\tightlist
\item
  设 \(n \in N\) ，\(\lambda \in C\) 满足 \(-n - 1 < \mathcal{R}(\lambda) < -n\)。对 R 上每个施瓦茨函数 \(\varphi\) ，我们有
\end{enumerate}

\[
\langle x _ {+} ^ {\lambda}, \varphi \rangle = \int_ {0} ^ {+ \infty} x ^ {\lambda} \left(\varphi (x) - \varphi (0) - x \varphi^ {\prime} (0) - \dots - \frac {x ^ {n - 1}}{(n - 1) !} \varphi^ {(n - 1)} (0)\right) d x.
\]

  从 \(V_{n}\) 到 \(\mathcal{D}'(\mathrm{U})\) 中的由 \(\lambda\mapsto x_{+}^{\lambda}\) 定义的映射是解析的。

\(d)\) 对每个 \(\lambda\) （属于诸集 \(V_{m}\) 之并的），我们有公式

\[
(x _ {+} ^ {\lambda}) ^ {\prime} = \lambda x _ {+} ^ {\lambda - 1}.
\]

\begin{enumerate}
\def\labelenumi{\arabic{enumi})}
\setcounter{enumi}{20}
\tightlist
\item
  \begin{enumerate}
  \def\labelenumii{\alph{enumii})}
  \tightlist
  \item
    设 \(f \in \mathcal{S}(\mathbf{R})\) 。R 上满足 \(g(0) = 2f'(0)\) 且 \(g(x) = (f(x) - f(-x))/x\) （对所有 \(x \neq 0\) ）的函数 g 在 R 上关于勒贝格测度可积。把 f 映为 g 在 R 上的积分的映射是一个缓增广义函数，记作 vp，称为 1/x 的柯西主值。
  \end{enumerate}
\end{enumerate}

\begin{enumerate}
\def\labelenumi{\alph{enumi})}
\setcounter{enumi}{1}
\tightlist
\item
  对每个 \(f \in \mathcal{S}(\mathbf{R})\) ，我们有
\end{enumerate}

\[
\langle \mathrm{vp}, f \rangle = \lim _ {\varepsilon \rightarrow 0} \int_ {\mathbf {R} - [ - \varepsilon , \varepsilon ]} \frac {f (x)}{x} d x.
\]

\begin{enumerate}
\def\labelenumi{\alph{enumi})}
\setcounter{enumi}{2}
\item
  vp 到 \(\mathbf{R} - \{0\}\) 的限制是与局部可积函数 \(x \mapsto 1/x\) （\(\mathbf{R} - \{0\}\) 上的）相伴的广义函数。
\item
  广义函数 vp 与下述广义函数的导数重合：该广义函数与 R 上在 0 处取值 0 且对 \(x \neq 0\) 由 \(x \mapsto \log(|x|)\) 定义的局部可积函数相伴。
\item
  设 \(\varepsilon > 0\) 。我们用 \(p_{\varepsilon}\) 表示函数 \(x \mapsto 1 / (x + i\varepsilon)\) （\(\mathbf{R}\) 上的）。我们有
\end{enumerate}

\[
\lim _ {\varepsilon \to 0} p _ {\varepsilon} = \mathrm{vp} - i \pi \varepsilon_ {0}
\]

在赋有弱拓扑的空间 \(\mathcal{S}'(\mathbf{R})\) 中成立，其中 \(\varepsilon_0\) 是与 0 处的点质量相伴的缓增广义函数。

\begin{enumerate}
\def\labelenumi{\arabic{enumi})}
\setcounter{enumi}{21}
\item
  R 上与指数函数相伴的广义函数不是缓增的。
\item
  设 \(n \in \mathbf{N}\) ，U 是 \(\mathbf{R}^n\) 的开集。我们用 \(\mu\) 表示 U 上的勒贝格测度。对 \(m \in \mathbf{N}\) ，我们记 \(\boldsymbol{m} = (m, \ldots, m) \in \mathbf{N}^n\) 。
\end{enumerate}

\begin{enumerate}
\def\labelenumi{\alph{enumi})}
\tightlist
\item
  设 \(f \in \mathcal{D}(\mathrm{U})\) 。假设存在整数 \(m \geqslant 1\) 以及实数 \(c \geqslant 0\) 使得
\end{enumerate}

\[
| \langle f, \varphi \rangle | \leqslant c \int_ {\mathrm{U}} | \partial^ {\boldsymbol {m}} \varphi | d \mu
\]

对所有 \(\varphi\in\mathcal{D}(\mathrm{U})\) 成立。于是存在函数 \(g\in\mathrm{L}^{\infty}(\mathrm{U})\) 使得 \(f=\partial^{m}g\) 。（线性映射 \(\varphi\mapsto\partial^{m}\varphi\) 在 \(\mathcal{D}(\mathrm{U})\) 上是单射；设 E 是它的像；考虑 E 上把 \(\partial^{m}\varphi\) 映为 \(\langle f,\varphi\rangle\) 的到 C 中的线性形式，并应用哈恩–巴拿赫定理。）

\begin{enumerate}
\def\labelenumi{\alph{enumi})}
\setcounter{enumi}{1}
\tightlist
\item
  设 V 是含于 U 的相对紧开集。存在整数 \(m \geqslant 1\) 以及 \(c \geqslant 0\) 使得
\end{enumerate}

\[
| \langle f, \varphi \rangle | \leqslant c \int_ {\mathrm{U}} | \partial^ {\boldsymbol {m}} \varphi | d \mu
\]

对所有支撑含于 V 的 \(\varphi\in\mathcal{D}(\mathrm{U})\) 成立。

\begin{enumerate}
\def\labelenumi{\alph{enumi})}
\setcounter{enumi}{2}
\tightlist
\item
  对每个相对紧开集 \(V \subset U\) ，存在连续函数 \(g \in \mathcal{C}(V)\) 以及 \(\alpha \in N^{n}\) ，使得 f 到 V 的限制等于 \(\partial^{\alpha}g\)。\(d)\) 给出一个广义函数 \(f \in \mathcal{D}'(\mathbf{R}^{n})\) 的例子，使得 f 不是 \(f = \partial^{\alpha}g\) 的形式，其中 \(g \in \mathcal{C}(\mathbf{R}^{n})\) 是连续函数且 \(\alpha \in N^{n}\) 。
\end{enumerate}

\begin{enumerate}
\def\labelenumi{\arabic{enumi})}
\setcounter{enumi}{23}
\tightlist
\item
  设 \(n \in N\) ，U 是 \(R^{n}\) 的开集。设 \(f \in \mathcal{D}'(U)\) 是支撑 C 在 U 中为紧的广义函数。
\end{enumerate}

\begin{enumerate}
\def\labelenumi{\alph{enumi})}
\tightlist
\item
  存在整数 \(m \in N\) ，使得广义函数 f 延拓为空间 \(\mathcal{C}^{m}(U)\) 上的连续线性形式，该空间赋有函数及其阶 \(\leqslant m\) 的偏导数在紧子集上一致收敛的拓扑。
\end{enumerate}

满足该性质的最小整数 m 称为 f 的阶。

\begin{enumerate}
\def\labelenumi{\alph{enumi})}
\setcounter{enumi}{1}
\item
  设 \(\varphi \in \mathcal{D}(\mathrm{U})\) 使得 \(\partial^{\alpha}\varphi\) 对每个 \(\alpha \in \mathbf{N}^{n}\) （满足 \(|\alpha| \leqslant m\) ）在 C 上为零。我们则有 \(\langle f, \varphi \rangle = 0\) 。（对充分小的 \(\delta > 0\) ，考虑 C 的邻域 \(\mathrm{U}_{\delta}\) ，它由 \(x \in \mathrm{U}\) （满足 \(d(x, \mathrm{C}) \leqslant \delta\) ）组成，并证明存在 \(\eta \geqslant 0\) （当 \(\delta\) 趋于 0 时趋于 0），使得对某个 \(\mathrm{C}_{0} \geqslant 0\) ，有 \(|\partial^{\alpha}\varphi(x)| \leqslant \mathrm{C}_{0}\delta^{m-|\alpha|}\eta\) 对 \(x \in \mathrm{U}_{\delta}\) 与 \(|\alpha| \leqslant m\) ；证明存在函数 \(\beta_{\delta} \in \mathcal{D}(\mathrm{U})\) ，它在 \(\mathrm{U}_{\delta/4}\) 上等于 1 且支撑含于 \(\mathrm{U}_{\delta}\) ，并满足 \(|\partial^{\alpha}\beta_{\delta}| \leqslant \mathrm{C}_{1}\delta^{-|\alpha|}\) 对某个 \(\mathrm{C}_{1} \geqslant 0\) ；最后注意 \(\langle f, \varphi \rangle = \langle f, \varphi\beta_{\delta} \rangle\) 且 \(\varphi\beta_{\delta}\) 在 \(\mathcal{C}^{m}(\mathrm{U})\) 中趋于 0。）
\item
  设 \(x \in U\) 。支撑化为单点 x 的广义函数 f 所成的空间以诸广义函数 \(\partial^{\alpha}\varepsilon_{x}\) （其中 \(\alpha \in N^{n}\) ）为代数基。
\item
  设 \(n = 1\) 。公式
\end{enumerate}

\[
\langle f, \varphi \rangle = \lim _ {m \rightarrow + \infty} \left(\sum_ {1 \leqslant k \leqslant m} \varphi (k ^ {- 1}) - m \varphi (0) - \log (m) \varphi^ {\prime} (0)\right)
\]

在 \(\mathbf{R}\) 上定义了一个阶为 1 的广义函数，其支撑是紧的，等于集 \(C = \{0\} \cup \{1/(n+1) \mid n \in \mathbf{N}\}\) 。存在序列 \((\varphi_n)\) （在 \(\mathcal{D}(\mathbf{R})\) 中），使得 \(\partial^{\alpha}\varphi_{n}\) 对每个 \(\alpha\in\mathbf{N}^{n}\) （满足 \(|\alpha|\leqslant1\) ）在 C 上一致收敛于 0，但 \(\langle f,\varphi_{n}\rangle\to+\infty\) 。

\begin{enumerate}
\def\labelenumi{\arabic{enumi})}
\setcounter{enumi}{24}
\tightlist
\item
  设 \(n \in \mathbf{N}\) 。我们用 \(N\) 表示 \(\mathbf{R}^n\) 上的欧几里得范数，并用 \(\mu\) 表示勒贝格测度。
\end{enumerate}

我们用 \(\mathcal{D}_{\mathrm{L}^{1}}(\mathbf{R}^{n})\) 表示 \(\mathcal{C}^{\infty}(\mathbf{R}^{n})\) 的向量子空间，其元素是函数 \(\varphi\) ，它满足 \(\partial^{\alpha}\varphi\in\mathrm{L}^{1}(\mathbf{R}^{n})\) 且 \(\partial^{\alpha}\varphi\in\mathcal{C}_{0}(\mathbf{R}^{n})\) 对所有 \(\alpha\in\mathbf{N}^{n}\) 。\(a)\) 空间 \(\mathcal{D}_{\mathrm{L}^{1}}(\mathbf{R}^{n})\) 赋有由半范数

\[
q _ {\alpha} (\varphi) = \int_ {\mathbf {R} ^ {n}} | \partial^ {\alpha} \varphi | d \mu , \quad \alpha \in \mathbf {N} ^ {n},
\]

定义的拓扑后是弗雷歇空间。

\begin{enumerate}
\def\labelenumi{\alph{enumi})}
\setcounter{enumi}{1}
\item
  空间 \(\mathcal{D}(\mathbf{R}^n)\) 含于 \(\mathcal{D}_{\mathrm{L}^1}(\mathbf{R}^n)\) 且在其中稠密。可以把对偶空间 \(\mathcal{D}_{\mathrm{L}^1}'(\mathbf{R}^n)\) （\(\mathcal{D}_{\mathrm{L}^1}(\mathbf{R}^n)\) 的）等同于 \(\mathcal{D}'(\mathbf{R}^n)\) 的一个子空间。
\item
  设 \(f \in \mathcal{D}'(\mathbf{R}^n)\) 。我们有 \(f \in \mathcal{D}_{\mathrm{L}^1}'(\mathbf{R}^n)\) 当且仅当存在有限集 I、族 \((\alpha_i)_{i \in \mathrm{I}}\) （在 \(\mathbf{N}^n\) 中）以及族 \((g_i)_{i \in \mathrm{I}}\) （在 \(\mathrm{L}^\infty(\mathbf{R}^n)\) 中）使得
\end{enumerate}

\[
f = \sum_ {i \in I} \partial^ {\alpha_ {i}} g _ {i}.
\]

（确定 I 与族 \((\alpha_{i})\) ，使得 \(|\langle f,\varphi\rangle |\) 有界，当 \(\varphi\) 满足 \(q_{\alpha_i}(\varphi)\leqslant 1\) 对所有 \(i\in \mathrm{I}\) ；把 \(\mathcal{D}_{\mathrm{L}^1}(\mathbf{R}^n)\) 等同于 \(\mathrm{L}^1 (\mathbf{R}^n)^m\) 的一个向量子空间，借助映射 \(\varphi \mapsto (\partial^{\alpha_i}\varphi)_{i\in \mathrm{I}}\) ，并应用哈恩–巴拿赫定理。）

\begin{enumerate}
\def\labelenumi{\alph{enumi})}
\setcounter{enumi}{3}
\tightlist
\item
  设 \(f \in \mathcal{D}'(\mathbf{R}^n)\) 。我们有 \(f \in \mathcal{D}_{\mathrm{L}^1}'(\mathbf{R}^n)\) 当且仅当存在具有缓和增长的连续函数 \(g\) （在 \(\mathbf{R}^n\) 上）以及 \(\alpha \in \mathbf{N}^n\) 使得 \(f = \partial^\alpha g\)。\(e)\) 设 \(f \in \mathcal{D}'(\mathbf{R}^n)\) 。存在整数 \(m \in \mathbf{N}\) 使得 \((1 + \mathrm{N}^2)^{-m/2}f\) 属于 \(\mathcal{D}_{\mathrm{L}^1}'(\mathbf{R}^n)\) 。
\end{enumerate}

\(f)\) 设 \(f \in \mathcal{D}'(\mathbf{R}^n)\) 。我们有 \(f \in \mathcal{S}'(\mathbf{R}^n)\) 当且仅当存在有界连续函数 \(g\colon \mathbf{R}^n \to \mathbf{C}\) 、整数 \(m \in \mathbf{N}\) 以及 \(\alpha \in \mathbf{N}^n\) 使得 \(f = \partial^\alpha((1 + \mathrm{N}^2)^{m/2}g)\) 。

\begin{enumerate}
\def\labelenumi{\arabic{enumi})}
\setcounter{enumi}{25}
\tightlist
\item
  设 \(n \in N\) 。我们给 \(R^{n}\) 赋予勒贝格测度 \(\mu\) ，并把 \(R^{n}\) 与其对偶通过映射 \((x, y) \mapsto \exp(2i\pi x \cdot y)\) 等同。
\end{enumerate}

\begin{enumerate}
\def\labelenumi{\alph{enumi})}
\item
  对任何函数 \(f \in \mathcal{S}(\mathbf{R}^n)\) ，映射 \(g \mapsto f * g\) 是局部凸空间 \(\mathcal{S}(\mathbf{R}^n)\) 的自同构。
\item
  设 \(f \in \mathcal{S}'(\mathbf{R}^{n})\) 且 \(g \in \mathcal{S}(\mathbf{R}^{n})\) 。映射 \(\varphi \mapsto \langle f, \check{g} * \varphi \rangle\) 是一个缓增广义函数，称为 f 与 g 的卷积。若 \(f \in \mathcal{S}(\mathbf{R}^{n})\) ，它与同 \(f * g\) 相伴的广义函数一致。我们仍用 \(f * g\) 表示这样定义的缓增广义函数。
\end{enumerate}

在下面的问题中，我们固定 \(f \in \mathcal{S}'(\mathbf{R}^n)\) 与 \(g \in \mathcal{S}(\mathbf{R}^n)\) 。\(c)\) 我们有 \(f * g \in \mathcal{C}^\infty(\mathbf{R})\) 并且

\[
\partial^ {\alpha} (f * g) = \partial^ {\alpha} (f) * g = f * \partial^ {\alpha} (g)
\]

对所有 \(\alpha\in\mathbf{N}^{n}\) 成立。

\(d)\) 对每个 \(y \in \mathbf{R}^n\)，用 \(g_y\) 表示函数 \(x \mapsto g(y - x)\) （\(\mathbf{R}^n\) 上的）。函数 \(f * g\) 满足

\[
(f * g) (y) = \langle f, g _ {y} \rangle
\]

对所有 \(y \in R^{n}\) 成立。

\begin{enumerate}
\def\labelenumi{\alph{enumi})}
\setcounter{enumi}{4}
\item
  设 \(h \in \mathcal{S}(\mathbf{R}^n)\) 。我们有 \((f*g)*h = f*(g*h)\) 。
\item
  函数 \(f*g\) 具有多项式增长，其所有导数亦然。（利用习题 25。）
\item
  我们有 \(\mathcal{F}(f*g) = \mathcal{F}(f)\mathcal{F}(g)\) 。
\item
  fg 的傅里叶变换是 \(R^{n}\) 上的无穷可微函数，并且连同其所有导数都具有多项式增长。此外，对每个 \(y \in R^{n}\) ，我们有
\end{enumerate}

\[
\mathcal {F} (f g) (y) = \langle f, g e _ {- y} \rangle
\]

其中 \(e_y\) 表示特征标 \(x \mapsto \exp(2i\pi x \cdot y)\) （\(\mathbf{R}^n\) 的）。

\begin{enumerate}
\def\labelenumi{\arabic{enumi})}
\setcounter{enumi}{26}
\tightlist
\item
  设 \(n \in \mathbf{N}\)。我们用 N 表示 \(\mathbf{R}^{n}\) 上的欧几里得范数，并用 \(\mu\) 表示 \(\mathbf{R}^{n}\) 的每个开子集上的勒贝格测度。设 \(\Delta\) 是微分算子
\end{enumerate}

\[
\Delta = - \sum_ {i = 1} ^ {n} \partial_ {i} ^ {2}
\]

将它视为 \(\mathcal{D}'(\mathbf{R}^{n})\) 的自同态。

对 \(k \in \mathbf{R}\) ，我们用 \(\mathrm{S}^k(\mathbf{R}^n)\) 表示这样的缓增广义函数 \(f\) （\(\mathcal{S}'(\mathbf{R}^n)\) 中的）所成的空间：\((1 + \mathrm{N}^k)\mathcal{F}(f)\) 属于 \(\mathrm{L}^2(\mathbf{R}^n)\) 。

设 \(U \subset R^{n}\) 是 \(R^{n}\) 的开子集。我们用 \(\mathrm{S}_{\mathrm{U}}^{k}(\mathbf{R}^{n})\) 表示满足如下条件的广义函数 \(f \in \mathcal{D}'(\mathbf{R}^{n})\) 所成的空间：对每个支撑含于 U 的试验函数 \(\varphi \in \mathcal{D}(\mathbf{R}^{n})\) ，我们有 \(\varphi f \in \mathrm{S}^{k}(\mathbf{R}^{n})\) 。

\begin{enumerate}
\def\labelenumi{\alph{enumi})}
\tightlist
\item
  空间 \(\mathrm{S}^{k}(\mathbf{R}^{n})\) 赋有内积
\end{enumerate}

\[
\langle f _ {1} \mid f _ {2} \rangle = \int_ {\mathbf {R} ^ {n}} (1 + \mathrm{N} ^ {k}) \overline {{\mathcal {F} (f _ {1})}} \mathcal {F} (f _ {2}) d \mu .
\]

当 \(k \in \mathbf{N}\) 时，我们有 \(W^{k}(\mathbf{R}^{n}) = S^{k}(\mathbf{R}^{n})\) ，并且这两个巴拿赫空间的范数等价；此外，这时我们有 \(W^{k}(\mathbf{R}^{n}) \subset S_{\mathrm{U}}^{k}(\mathbf{R}^{n})\) 。

\begin{enumerate}
\def\labelenumi{\alph{enumi})}
\setcounter{enumi}{1}
\tightlist
\item
  \(\mathrm{S}^{k}(\mathbf{R}^{n})\) 的对偶通过双线性形式等同于 \(\mathrm{S}^{-k}(\mathbf{R}^{n})\)
\end{enumerate}

\[
(f, g) \mapsto \int_ {\mathbf {R} ^ {n}} \mathcal {F} (f) (- x) \mathcal {F} (g) (x) d \mu (x)
\]

（对 \((f,g)\in \mathrm{S}^{-k}(\mathbf{R}^n)\times \mathrm{S}^k (\mathbf{R}^n).\)）

\begin{enumerate}
\def\labelenumi{\alph{enumi})}
\setcounter{enumi}{2}
\tightlist
\item
  假设 U 是相对紧的。我们有
\end{enumerate}

\[
\bigcup_ {k \in \mathbf {R}} \mathrm{S} _ {\mathrm{U}} ^ {k} (\mathbf {R} ^ {n}) = \mathcal {D} ^ {\prime} (\mathbf {R} ^ {n}).
\]

（设 \(f \in \mathcal{D}'(\mathbf{R}^n)\) ；把习题 19 的结果应用于 \(\varphi f\) ，其中 \(\varphi \in \mathcal{D}(\mathbf{R}^n)\) 在 U 上等于 1。）

\(d)\) 设 \(k \in \mathbf{R}\) 且 \(f \in \mathrm{S}^k(\mathbf{R}^n)\) 。对每个整数 \(i\) （满足 \(1 \leqslant i \leqslant n\) ），我们有 \(\partial_i f \in \mathrm{S}^{k-1}(\mathbf{R}^n)\) 。此外，若 \(\Delta(f) \in \mathrm{S}^k(\mathbf{R}^n)\) ，则 \(f \in \mathrm{S}^{k+2}(\mathbf{R}^n)\) 。

\begin{enumerate}
\def\labelenumi{\alph{enumi})}
\setcounter{enumi}{4}
\tightlist
\item
  设 \(k \in \mathbf{R}\) 且 \(f \in \mathrm{S}_{\mathrm{U}}^{k}(\mathbf{R}^{n})\) 。若 \(\Delta(f) \in \mathrm{S}_{\mathrm{U}}^{k}(\mathbf{R}^{n})\) ，则 \(f \in \mathrm{S}_{\mathrm{U}}^{k+2}(\mathbf{R}^{n})\) 。
\end{enumerate}

\(f)\) 设 \(k > n/2\) 是实数，\(m \in \mathbf{N}\) 是满足 \(m < k - n/2\) 的整数。空间 \(\mathrm{S}^k(\mathbf{R}^n)\) 含于 \(\mathcal{C}^m(\mathbf{R}^n)\) （利用柯西–施瓦茨不等式，验证 \(x \mapsto x^\alpha \mathscr{F} f(x)\) 属于 \(\mathrm{L}^1(\mathbf{R}^n)\) 对所有 \(\alpha \in \mathbf{N}^n\) （满足 \(|\alpha| \leqslant m\) ）。然后利用傅里叶反演公式证明 \(\mathscr{F}(f)\) 在 \(\mathbf{R}\) 上可微。）

\(g)\) 设 \(k > n/2\) 是实数，\(m \in \mathbf{N}\) 是满足 \(m < k - n/2\) 的整数。\(\mathrm{S}_{\mathrm{U}}^{k}(\mathbf{R}^{n})\) 的元素到 U 的限制属于 \(\mathcal{C}^{m}(\mathrm{U})\) （对每个 \(x \in \mathrm{U}\) ，验证得到一个良定义的函数 \(g \in \mathcal{C}^{m}(\mathrm{U})\) ，它由 \(g(x) = (\varphi_{x}f)(x)\) 定义，其中 \(\varphi_{x} \in \mathcal{D}(\mathbf{R}^{n})\) 是支撑在 U 中且在 \(x\) 的一个邻域中等于 1 的函数；最后验证 \(f\) 到 U 的限制等于 \(g\) 。）

\begin{enumerate}
\def\labelenumi{\alph{enumi})}
\setcounter{enumi}{7}
\tightlist
\item
  设 \(g \in \mathcal{D}'(\mathbf{R}^n)\) 。若 \(f \in \mathcal{D}'(\mathbf{R}^n)\) 满足 \(\Delta(f) = g\) 且 \(g\) 到 U 的限制属于 \(\mathcal{C}^\infty(\mathrm{U})\) ，则 \(f\) 到 U 的限制属于 \(\mathcal{C}^\infty(\mathrm{U})\) （“椭圆正则性定理”）。
\end{enumerate}

\begin{enumerate}
\def\labelenumi{\arabic{enumi})}
\setcounter{enumi}{27}
\tightlist
\item
  设 \(n\) 与 \(m\) 是整数，\(\mathrm{U} \subset \mathbf{R}^n\) 与 \(\mathrm{V} \subset \mathbf{R}^m\) 是开集。\(a)\) 存在唯一的单射线性映射 \(u\) ，它从 \(\mathcal{D}(\mathrm{U}) \otimes \mathcal{D}(\mathrm{V})\) 到 \(\mathcal{D}(\mathrm{U} \times \mathrm{V})\) 中，使得 \(u(\varphi \otimes \psi)\) 是函数 \((x, y) \mapsto \varphi(x)\psi(y)\) 对所有 \((\varphi, \psi) \in \mathcal{D}(\mathrm{U}) \times \mathcal{D}(\mathrm{V})\) 。此外，\(u\) 的像在 \(\mathcal{D}(\mathrm{U} \times \mathrm{V})\) 中稠密。我们把空间 \(\mathcal{D}(\mathrm{U}) \otimes \mathcal{D}(\mathrm{V})\) 等同于 \(\mathcal{D}(\mathrm{U} \times \mathrm{V})\) 的一个子空间，借助映射 \(u\) 。
\end{enumerate}

\begin{enumerate}
\def\labelenumi{\alph{enumi})}
\setcounter{enumi}{1}
\item
  对 U 上的所有广义函数 \(f\) 与 V 上的所有广义函数 \(g\) ，存在 U × V 上唯一的广义函数 \(f \otimes g\) ，使得 \(\langle f \otimes g, \varphi \otimes \psi \rangle = \langle f, \varphi \rangle \langle g, \psi \rangle\) 对所有 \((\varphi, \psi) \in \mathcal{D}(\mathrm{U}) \times \mathcal{D}(\mathrm{V})\) 成立。
\item
  设 \(f \in \mathcal{D}(\mathrm{U})\) 且 \(g \in \mathcal{D}(\mathrm{V})\) 。对每个 \(\eta \in \mathcal{D}(\mathrm{U} \times \mathrm{V})\) 以及每个 \(x \in U\)，用 \(\eta_x\) 表示映射 \(y \mapsto \eta(x, y)\)。我们有 \(\eta_x \in \mathcal{D}(\mathrm{V})\) ，并且映射 \(\varphi\) （它把每个 \(x \in U\) 映为 \(\langle f, \eta_x \rangle\) ）属于 \(\mathcal{D}(\mathrm{U})\) 且满足
\end{enumerate}

\[
\langle f \otimes g, \eta \rangle = \langle f, \varphi \rangle .
\]

\begin{enumerate}
\def\labelenumi{\arabic{enumi})}
\setcounter{enumi}{28}
\tightlist
\item
  设 \(n\) 与 \(m\) 是整数，\(\mathrm{U} \subset \mathbf{R}^n\) 与 \(\mathrm{V} \subset \mathbf{R}^m\) 是开集。我们沿用上一个习题的记号。设 \(u\) 是从 \(\mathcal{D}(\mathrm{U})\) 到 \(\mathcal{D}'(\mathrm{V})\) 中的连续线性映射。
\end{enumerate}

\begin{enumerate}
\def\labelenumi{\alph{enumi})}
\item
  存在唯一的连续线性形式 \(\widetilde{\kappa}\) ，它定义在子空间 \(\mathcal{D}(\mathrm{U})\otimes \mathcal{D}(\mathrm{V})\) （\(\mathcal{D}(\mathrm{U}\times \mathrm{V})\) 的）上，使得 \(\widetilde{\kappa} (\varphi \otimes \psi) = \langle u(\varphi),\psi \rangle\) 对所有 \((\varphi ,\psi)\) （\(\mathcal{D}(\mathrm{U})\times \mathcal{D}(\mathrm{V})\) 中的）成立。
\item
  存在广义函数 \(\kappa(u)\) （在 \(U \times V\) 上）延拓 \(\widetilde{\kappa}\) ，并且映射 \(u \mapsto \kappa(u)\) （从 \(\mathcal{L}(\mathcal{D}(U), \mathcal{D}'(V))\) 到 \(\mathcal{D}'(U \times V)\) 中的）是连续的。
\item
  映射 \(\kappa\) 是双射（“施瓦茨核定理”；为证明 \(\kappa\) 的满射性，给定 U × V 上的广义函数 \(\eta\) ，把 \(u(\varphi)\) 定义为 V 上的广义函数 \(\psi \mapsto \langle \eta, \varphi \otimes \psi \rangle\) 。）
\end{enumerate}

\begin{enumerate}
\def\labelenumi{\arabic{enumi})}
\setcounter{enumi}{29}
\tightlist
\item
  设 \(n \in \mathbf{N}\) 。我们用 \((x_1, \ldots, x_n)\) 表示 \(\mathbf{R}^n\) 的典范基的对偶基。
\end{enumerate}

\begin{enumerate}
\def\labelenumi{\alph{enumi})}
\tightlist
\item
  设 \(c \geqslant 0\) 与 r \textgreater{} 0 是实数。存在在 \(R^{n}\) 中 0 的开邻域 U 上定义且解析的函数 u 使得
\end{enumerate}

\[
\left(r - \sum_ {j = 1} ^ {n - 1} x _ {j} - u\right) \partial_ {n} u - c r \sum_ {j = 1} ^ {n - 1} \partial_ {j} u = c r
\]

在 U 上成立，并且 \(u(x,0)=0\) 对 \(x\in\mathbf{R}^{n-1}\) （满足 \((x,0)\in\mathrm{U}\) ）。（先考虑 n=2 的情形。）

\begin{enumerate}
\def\labelenumi{\alph{enumi})}
\setcounter{enumi}{1}
\tightlist
\item
  设 U 是 \(\mathbf{R}^{n}\) 中 0 的开邻域。设 \((a_{1},\ldots,a_{n-1},b)\) 是从 \(\mathrm{U}\times\mathbf{R}\) 到 \(\mathbf{C}\) 的一组函数，它们在 \((0,0)\in\mathrm{U}\times\mathbf{R}\) 的一个邻域中解析。我们用 Q 表示从 U 上解析函数所成的向量空间到其自身的映射，使得
\end{enumerate}

\[
\mathrm{Q} (u) = \partial_ {n} u - \sum_ {j = 1} ^ {n - 1} a _ {j} (x _ {1}, \dots , x _ {n - 1}, u) \partial_ {j} u + b (x _ {1}, \dots , x _ {n - 1}, u).
\]

至多存在一个函数 \(v\) ，它在 0 的开邻域 \(V\) （含于 \(U\) ）中定义且解析，使得 \(Q(v) = 0\) 并且 \(v(x,0) = 0\) 对所有 \(x \in \mathbf{R}^{n-1}\) （满足 \((x,0) \in V\) ）。（验证这样的解 \(v\) 在 0 处的泰勒展开的系数被唯一确定。）

\begin{enumerate}
\def\labelenumi{\alph{enumi})}
\setcounter{enumi}{2}
\tightlist
\item
  设 \(f\) 与 \(g\) 是 \(\mathbf{C}[[z_1, \ldots, z_n]]\) 中的形式幂级数。称 \(f\) 优控 \(g\) ，并记作 \(f \ll g\) ，如果
\end{enumerate}

\[
f = \sum_ {\alpha \in \mathbf {N} ^ {n}} a _ {\alpha} z ^ {\alpha}, \qquad g = \sum_ {\alpha \in \mathbf {N} ^ {n}} b _ {\alpha} z ^ {\alpha}
\]

其中 \(|a_{\alpha}| \leqslant b_{\alpha}\) 对所有 \(\alpha \in \mathbf{N}^{n}\) 。若形式级数 g 在 \(R^{n}\) 中 0 的一个邻域中收敛，则 f 亦然。

\(d)\) 若存在在 0 的邻域中解析的一组函数 \((\widetilde{a}_1, \ldots, \widetilde{a}_{n-1}, \widetilde{b})\) ，使得 \(a_j \ll \widetilde{a}_j\) 对 \(1 \leqslant j \leqslant n-1\) 且 \(b \ll \widetilde{b}\) （即在 0 处的相应泰勒展开满足这些条件的意义下），并且若方程

\[
\partial_ {n} v - \sum_ {j = 1} ^ {n - 1} \widetilde {a} _ {j} (x _ {1}, \dots , x _ {n - 1}, v) \partial_ {j} v + \widetilde {b} (x _ {1}, \dots , x _ {n - 1}, v) = 0
\]

关于未知函数 \(v\) 在 0 的一个邻域中有解析解，则方程 \(\mathrm{Q}(v) = 0\) 在 0 的一个邻域中有唯一的解析解。

\begin{enumerate}
\def\labelenumi{\alph{enumi})}
\setcounter{enumi}{4}
\tightlist
\item
  方程 \(Q(v) = 0\) 存在唯一的在 0 的邻域中解析的解（“柯西–柯瓦列夫斯卡娅定理”）。
\end{enumerate}

\begin{enumerate}
\def\labelenumi{\arabic{enumi})}
\setcounter{enumi}{30}
\tightlist
\item
  我们用 L 表示微分算子
\end{enumerate}

\[
\mathrm{L} = - i \partial_ {1} + \partial_ {2} - 2 (x _ {1} + i x _ {2}) \partial_ {3}
\]

它是 \(\mathbf{R}^{3}\) 上阶为 1 的算子，其中 \((x_{1}, x_{2}, x_{3})\) 是 \(\mathbf{R}^{3}\) 的典范基的对偶基，诸 \(x_{i}\) 被视为 \(\mathbf{R}^{3}\) 上的多项式（“卢伊算子”）。

\begin{enumerate}
\def\labelenumi{\alph{enumi})}
\item
  设 U 是 \(\mathbf{R}^3\) 中 0 的开邻域。存在函数 \(\varphi \in \mathcal{D}(\mathrm{U})\) 使得 \(\varphi(0) = 1\) 并且 \(\mathrm{L}(\varphi)\) 在含于 U 的 0 的邻域 V 中为零。（利用柯西–柯瓦列夫斯卡娅定理。）
\item
  设 \(u \in \mathcal{C}^{\infty}(\mathrm{U})\) 是满足 \(\mathrm{L}(u) = 0\) 且 \(u(0) = 0\) 的函数。假设存在 \(\delta > 0\) 使得 \(\mathcal{R}(u(x)) < -2\delta\) 对 \(x \in U - V\) 成立。我们令 \(v_{t} = \varphi \exp(t(u + \delta))\)。我们有 \(v_{t} \in \mathcal{C}^{\infty}(\mathrm{U})\) ，并且 \(\mathrm{L}(v_{t})\) 在 \(\mathcal{D}(\mathrm{U})\) 中趋于 0 当 \(t \to +\infty\) 时。
\item
  设 F 是由这样的函数 \(f \in \mathcal{C}^{\infty}(U)\) 组成的空间：\(\partial^{\alpha}f\) 对每个 \(\alpha \in N^{3}\) 在 U 上有界。空间 F 赋有半范数 \(q_{\alpha}(f) = \sup_{x \in U} |\partial^{\alpha}f|\) 后是弗雷歇空间。
\end{enumerate}

\(d)\) 设 \(\psi\in\mathcal{D}(\mathbf{R}^{3})\) 。用 \(f_{t}\) 表示函数 \(x\mapsto t^{3}\psi(tx)\) 到 U 的限制。我们有 \(f_{t}\in\mathrm{F}\) 对 \(t\geqslant1\) ，并且存在实数 \(c\geqslant0\) ，使得 \(q_{\alpha}(f_{t})\leqslant ct^{|\alpha|+3}\) 对所有 \(\alpha\in\mathbf{N}^{3}\) 以及每个 \(t\geqslant1\) 。

\begin{enumerate}
\def\labelenumi{\alph{enumi})}
\setcounter{enumi}{4}
\tightlist
\item
  对每个 \(t \geqslant 1\) ，我们用 \(\lambda_{t}\) 表示 F 上的线性形式 \(f \mapsto \int_{U} fv_{t}\)。它是连续的，并且对充分大的 t，我们有
\end{enumerate}

\[
\lambda_ {t} (f _ {t}) = e ^ {\delta t} \left(\int_ {\mathrm{U}} \psi (x) \exp \left(\sum_ {j = 1} ^ {3} \xi_ {j} x _ {j}\right) d x _ {1} d x _ {2} d x _ {3} + \varepsilon (t)\right)
\]

其中 \(\xi_{j} = \partial_{j}u(0)\) ，并且 \(\varepsilon (t)\to 0\) 当 \(t\rightarrow +\infty\)。

\(f)\) 假设 \((\xi_j) \neq 0\) ；存在 \(f \in \mathrm{F}\) 使得积分 \(\int_{\mathrm{U}} f v_t\) 当 \(t \to +\infty\) 时不是有界的。（若该论断不真，则利用上一问题与巴拿赫–斯坦豪斯定理得出矛盾。）

\(g)\) 诸函数 \(v_{t}\) 在 \(\mathcal{D}'(\mathrm{U})\) 中不收敛于 0 当 \(t\to+\infty\) 时。\(h)\) 设 \(\xi=(0,0,\xi_{3})\in\mathbf{R}^{3}\) 满足 \(\xi_{3}<0\)。存在 \(\mathbf{R}^{3}\) 中 0 的邻域 U 以及函数 \(u\in\mathcal{C}^{\infty}(\mathrm{U})\) ，使得 \(\mathrm{L}(u)=0\) 且 \(u(0)=0\) ，并满足问题 \(b)\) 与 \(f)\) 的条件。（利用柯西–柯瓦列夫斯卡娅定理。）

\begin{enumerate}
\def\labelenumi{\roman{enumi})}
\tightlist
\item
  设 \((v_{t})\) 是 \(\mathcal{D}(\mathrm{U})\) 中满足 \(\mathrm{L}(v_{t})\to0\) 的函数，并设 \(f\in\mathcal{D}(\mathrm{U})\) 使得 \(\int_{\mathrm{U}}fv_{t}\) 不是有界的当 \(t\to+\infty\) 时（问题 g)）；方程 Lv=f 没有解 \(v\in\mathcal{D}'(\mathrm{U})\) （“卢伊定理”）。（注意 L 的形式伴随与 L 一致。）
\end{enumerate}

\begin{enumerate}
\def\labelenumi{\arabic{enumi})}
\setcounter{enumi}{31}
\tightlist
\item
  设 \(n \geqslant 1\) 是整数。我们用 \(\Delta\) 表示 \(\mathcal{D}'(\mathbf{R}^n)\) 上的拉普拉斯算子，它定义为
\end{enumerate}

\[
\Delta (f) = - \sum_ {i = 1} ^ {n} \partial_ {i} ^ {2} f.
\]

设 u 是 \(R^{n}\) 上的局部可积函数，使得 \(\Delta(u)\) 是局部可积函数。

设 \(v\in \mathcal{L}^{\infty}(\mathbf{R}^{n})\) 是由下式定义的函数：

\[
v (x) = \left\{ \begin{array}{l l} 0 & \text { if } u (x) = 0, \\ \overline {{u (x)}} / | u (x) | & \text { if } u (x) \neq 0. \end{array} \right.
\]

\begin{enumerate}
\def\labelenumi{\alph{enumi})}
\item
  对 \(\varepsilon > 0\)，令 \(u_{\varepsilon} = (|u|^2 + \varepsilon^2)^{1/2}\) 。证明广义函数 \(-\Delta(u_{\varepsilon}) - \mathcal{R}(\overline{u}/u_{\varepsilon} \Delta(u))\) 是正的。（先考虑 \(u \in \mathcal{C}^{\infty}(\mathbf{R}^{n})\) 的情形，再用卷积处理一般情形。）
\item
  广义函数 \(-\Delta(|u|) - \mathcal{R}(v\Delta(u))\) 是正的（“加藤不等式”）。
\end{enumerate}

\begin{enumerate}
\def\labelenumi{\arabic{enumi})}
\setcounter{enumi}{32}
\tightlist
\item
  设 \(m \geqslant 1\) 是整数。
\end{enumerate}

\begin{enumerate}
\def\labelenumi{\alph{enumi})}
\tightlist
\item
  存在唯一的线性映射 \(A^{+}\) （相应地 \(A^{-}\) ），它从 \(\mathrm{H}^{m}(\mathbf{R})\) 到 \(\mathrm{H}^{m-1}(\mathbf{R})\) ，使得
\end{enumerate}

\[
\mathrm{A} ^ {+} f (x) = f ^ {\prime} (x) - x f (x) \quad (\text { resp.   A } ^ {-} f (x) = f ^ {\prime} (x) + x f (x))
\]

对所有 \(f \in \mathcal{D}(\mathbf{R})\) 以及每个 \(x \in \mathbf{R}\) 成立。

\begin{enumerate}
\def\labelenumi{\alph{enumi})}
\setcounter{enumi}{1}
\tightlist
\item
  映射 \(A^{+}\) 是指标为 1 的弗雷德霍姆算子，而 \(A^{-}\) 是指标为 -1 的弗雷德霍姆算子。（考虑埃尔米特函数在 \(A^{+}\) 与 \(A^{-}\) 下的像，参见 II, p. 263 的习题 3。）
\end{enumerate}
\exercisesection{}

除非另有说明，以下巴拿赫空间与希尔伯特空间均假定为复空间。

\begin{enumerate}
\def\labelenumi{\arabic{enumi})}
\tightlist
\item
  \begin{enumerate}
  \def\labelenumii{\alph{enumii})}
  \tightlist
  \item
    存在向量空间 E、F、G，从 E 到 F 的部分算子 u、v，以及从 F 到 G 的部分算子 w，使得 \(w \circ (u + v) \neq w \circ u + w \circ v\) 。
  \end{enumerate}
\end{enumerate}

\begin{enumerate}
\def\labelenumi{\alph{enumi})}
\setcounter{enumi}{1}
\item
  给出拓扑向量空间上一个可闭算子 u 的例子，使得 u 的闭包的定义域严格包含于 u 的定义域的闭包之中。
\item
  给出希尔伯特空间上闭算子 u 与 v 的一个例子，使得 u 连续但 \(u \circ v\) 不是闭的。
\end{enumerate}

\begin{enumerate}
\def\labelenumi{\arabic{enumi})}
\setcounter{enumi}{1}
\tightlist
\item
  设 E 与 F 为巴拿赫空间，G 为 E 的一个子空间。设 u 为从 G 到 F 的有限秩线性映射。假设 u 在 G 上不连续。
\end{enumerate}

\begin{enumerate}
\def\labelenumi{\alph{enumi})}
\item
  由 u 定义的、以 G 为定义域的从 E 到 F 的部分算子是不可闭的。
\item
  假设 E 与 G 为希尔伯特空间。计算由 u 定义的部分算子的伴随。
\end{enumerate}

\begin{enumerate}
\def\labelenumi{\arabic{enumi})}
\setcounter{enumi}{2}
\tightlist
\item
  设 K = R 或 C。设 E、F 为 K 上的拓扑向量空间，u 为从 E 到 F 的一个稠定算子。设 D' 为满足下列条件的连续线性形式 \(\lambda \in F'\) 组成的集合：存在连续线性形式 \(\mu \in E'\) 在 u 的定义域上与 \(\lambda \circ u\) 重合。
\end{enumerate}

\begin{enumerate}
\def\labelenumi{\alph{enumi})}
\tightlist
\item
  集合 D' 是 F' 的向量子空间，线性形式 \(\mu\) 由 \(\lambda\) 唯一确定，且从 D' 到 E' 的映射 \(\lambda \mapsto \mu\) 是线性的。
\end{enumerate}

我们把 u 的转置记作 \(^{t}u\) ，它是从 F' 到 E' 、以 D' 为定义域、由 \(\lambda\mapsto\mu\)（\(\lambda\in D'\) ）给出的部分算子。

\begin{enumerate}
\def\labelenumi{\alph{enumi})}
\setcounter{enumi}{1}
\item
  假设 E 与 F 是局部凸的。算子 u 可闭当且仅当对偶 \(\mathrm{F} \times \mathrm{dom}(^{t}u) \to \mathrm{K}\) 是分离的。
\item
  假设 E 与 F 是局部凸的且 F 是半自反的。算子 u 可闭当且仅当 \(^{t}u\) 稠定。
\end{enumerate}

\begin{enumerate}
\def\labelenumi{\arabic{enumi})}
\setcounter{enumi}{3}
\item
  设 E 为复希尔伯特空间，u 为 E 上的一个单射自伴部分算子。部分算子 \(u^{-1}\) 是自伴的。
\item
  设 E 为复巴拿赫空间。E 上的连续半群是指一个映射 \(s\colon\mathbf{R}_{+}\to\mathcal{L}(\mathrm{E})\) ，它满足
\end{enumerate}

\begin{enumerate}
\def\labelenumi{(\roman{enumi})}
\item
  \(s(0) = 1_{\mathrm{E}}\) ；
\item
  从 \(R_{+} \times E \to E\) 、由 \((t, x) \mapsto s(t)x\) 定义的映射是连续的；
\item
  对 \(t_1\) 与 \(t_2\) （均属于 \(\mathbf{R}_{+}\) ） ，有 \(s(t_1 + t_2) = s(t_1)s(t_2)\) 。
\end{enumerate}

我们考虑 E 上的一个连续半群 s。

\begin{enumerate}
\def\labelenumi{\alph{enumi})}
\tightlist
\item
  存在实数 C ≥ 0 与 M ≥ 0，使得
\end{enumerate}

\[
\| s (t) \| \leqslant \mathrm{Ce} ^ {\mathrm{Mt}}
\]

对所有 \(t \in R_{+}\) 成立。（设 a \textgreater{} 0；证明 \(\|s(t)\|\) 在 \(t \in [0, a]\) 上有界。）在本习题的其余部分，我们固定这样的数。

\begin{enumerate}
\def\labelenumi{\alph{enumi})}
\setcounter{enumi}{1}
\item
  设 D 为 \(x \in E\) 中使得映射 \(\psi_{x}\) （从 \(R_{+}\) 到 E 、由 \(t \mapsto s(t)x\) 定义）在 0 处可微的那些 x 组成的集合。空间 D 在 E 中稠密；由 \(u(x) = -\psi_{x}'(0)\) 定义的从 D 到 E 的映射 u 是 E 上的一个闭部分算子，称为半群 s 的无穷小生成元。
\item
  设 \(\lambda\) 为使 \(\mathcal{R}(\lambda) < -\mathrm{M}\) 的复数。公式
\end{enumerate}

\[
v _ {\lambda} = - \int_ {\mathbf {R} _ {+}} e ^ {\lambda t} s (t) d t
\]

定义 E 的一个自同态。

\(d)\) \(v_{\lambda}\) 的像包含于 \(\operatorname{dom}(u)\) ，并且我们有 \(u \circ v_{\lambda} = -1_{\mathrm{E}} - \lambda v_{\lambda}\) 。

\begin{enumerate}
\def\labelenumi{\alph{enumi})}
\setcounter{enumi}{4}
\tightlist
\item
  我们有 \(u \circ v_{\lambda} = v_{\lambda} \circ u\) 。
\end{enumerate}

\(f)\) 复数 \(\lambda\) 属于 \(u\) 的预解集，并且我们有 \(v_{\lambda} = -\mathrm{R}(u,\lambda)\) 。

\begin{enumerate}
\def\labelenumi{\alph{enumi})}
\setcounter{enumi}{6}
\tightlist
\item
  我们有
\end{enumerate}

\[
\| v _ {\lambda} ^ {n} \| \leqslant \frac {\mathrm{C}}{| \mathcal {R} (\lambda) + \mathrm{M} | ^ {n}}
\]

对所有 \(n \in N\) 成立。

\begin{enumerate}
\def\labelenumi{\arabic{enumi})}
\setcounter{enumi}{5}
\tightlist
\item
  设 E 为复巴拿赫空间，u 为 E 上的一个稠定闭部分算子。假设存在实数 \(C \geqslant 0\) 与 \(M \geqslant 0\) ，使得
\end{enumerate}

\begin{enumerate}
\def\labelenumi{(\roman{enumi})}
\item
  集合 {]}−∞, −M{[} 包含于 u 的预解集；
\item
  对每个整数 \(n \in \mathbf{N}\) 与每个实数 \(\lambda < -\mathrm{M}\) ，我们有
\end{enumerate}

\[
\left\| \mathrm{R} (u, \lambda) ^ {n} \right\| \leqslant \mathrm{C} | \lambda + \mathrm{M} | ^ {- n}.
\]

\begin{enumerate}
\def\labelenumi{\alph{enumi})}
\tightlist
\item
  对每个实数 \(y < -\mathrm{M}\) ，我们记
\end{enumerate}

\[
u _ {y} = - y 1 _ {\mathrm{E}} + y ^ {2} \mathrm{R} (u, y) \in \mathcal {L} (\mathrm{E}).
\]

诸自同态 \(u_{y}\) 两两可交换；对每个 \(x \in \text{dom}(u)\) ，我们有 \(u_{y}(x) \to u(x)\) （\(y \to -\infty\) ）。

\begin{enumerate}
\def\labelenumi{\alph{enumi})}
\setcounter{enumi}{1}
\item
  从 \(R_{+}\) 到 \(\mathcal{L}(E)\) 、由 \(s_{y}(t)=e^{-tu_{y}}\) 定义的映射是 E 上的连续半群，使得 \(\|s_{y}(t)\|\leqslant Ce^{Mt}\) 对所有 \(t\in R_{+}\) 成立。
\item
  对所有 \(y_{1}, y_{2}\) 与 \(t\) （属于 \(\mathbf{R}_{+}^{*}\) ），以及每个 \(x\) （属于 \(\mathrm{dom}(u)\) ），我们有
\end{enumerate}

\[
\left\| s _ {y _ {1}} (t) x - s _ {y _ {2}} (t) x \right\| \leqslant t \mathrm{C} ^ {2} \left\| u _ {y _ {1}} (x) - u _ {y _ {2}} (x) \right\|.
\]

\(d)\) 设 \(t \in \mathbf{R}_{+}\) 。对 E 中每个 \(x\) ，定义在 \(\mathbf{R}_{+}^{*}\) 上的函数 \(y \mapsto s_{y}(t)x\) 当 \(y \to -\infty\) 时收敛于 E 的一个元素 \(s(t)x\) 。映射 \(s(t)\) 线性且连续，范数 \(\leqslant \mathrm{Ce}^{\mathrm{M}t}\) 。

\begin{enumerate}
\def\labelenumi{\alph{enumi})}
\setcounter{enumi}{4}
\item
  映射 \(t \mapsto s(t)\) 是 E 上的连续半群。
\item
  部分算子 u 是连续半群 s 的无穷小生成元（“希尔–吉田定理”）。（首先验证 s 的无穷小生成元是 u 的一个扩张。）
\end{enumerate}

\(g)\) 设 \(\widetilde{s}\) 为以 \(u\) 为无穷小生成元的一个连续半群。我们有 \(\widetilde{s} = s\) 。

\begin{enumerate}
\def\labelenumi{\arabic{enumi})}
\setcounter{enumi}{6}
\tightlist
\item
  设 E 为复巴拿赫空间，u 为 E 上的一个部分算子。我们称 u 是增生的，如果对每个 \(x \in \text{dom}(u)\) 存在连续线性形式 \(\ell \in E'\) 使得
\end{enumerate}

\begin{enumerate}
\def\labelenumi{(\roman{enumi})}
\item
  \(\mathcal{R}(\ell(u(x))) \geqslant 0;\)
\item
  我们有 \(\|\ell\|=\|x\|\) 且 \(\ell(x)=\|x\|^{2}\) 。
\end{enumerate}

我们称 E 上的连续半群 \(s\) 是压缩的，如果 \(\|s(t)\|\leqslant1\) 对所有 \(t\in\mathbf{R}_{+}\) 成立。

\begin{enumerate}
\def\labelenumi{\alph{enumi})}
\item
  部分算子 \(u\) 是 E 上一个连续压缩半群的无穷小生成元，当且仅当 \(u\) 是增生的并且存在 \(\lambda > 0\) 使部分算子 \(\lambda 1_{\mathrm{E}} + u\) 是满射的。（为证明条件必要，设 \(s\) 为以 \(u\) 为无穷小生成元的压缩半群；对 \(x \in \operatorname{dom}(u)\) ，考虑由 \(t \mapsto \ell(s(t)x)\) 定义的函数；它在 0 处可微，且其在 0 处的导数为 \(-\mathcal{R}(\ell(u(x)))\) 。）
\item
  假设 u 是闭的，且 \(u\) 与 \(^{t}u\) 都是增生的。则 u 是 E 上一个连续压缩半群的无穷小生成元。
\end{enumerate}

\begin{enumerate}
\def\labelenumi{\arabic{enumi})}
\setcounter{enumi}{7}
\tightlist
\item
  设 E 为复巴拿赫空间。设 s 为 E 上的一个连续压缩半群，u 为其无穷小生成元。我们有
\end{enumerate}

\[
\lim _ {n \to + \infty} \left(1 + \frac {t}{n} u\right) ^ {- n} x = s (t) x
\]

对所有 \(x \in E\) 成立。

\begin{enumerate}
\def\labelenumi{\arabic{enumi})}
\setcounter{enumi}{8}
\tightlist
\item
  \begin{enumerate}
  \def\labelenumii{\alph{enumii})}
  \tightlist
  \item
    求希尔伯特空间 E 与 F，以及从 E 到 F 的部分算子 u 与 v，使得 u 与 \(u + v\) 都稠定，且 \(u^{*} + v^{*}\) 不等于 \((u + v)^{*}\) 。
  \end{enumerate}
\end{enumerate}

\begin{enumerate}
\def\labelenumi{\alph{enumi})}
\setcounter{enumi}{1}
\item
  求希尔伯特空间 E、F、G，以及分别从 E 到 F 与从 F 到 G 的稠定部分算子 u 与 v，使得 \(v \circ u\) 稠定且 \(u^{*} \circ v^{*}\) 不等于 \((v \circ u)^{*}\) 。
\item
  求一个希尔伯特空间 E 与 E 上的自伴部分算子 u 与 v，使得 \(u + v\) 稠定但不是自伴的。
\end{enumerate}

\begin{enumerate}
\def\labelenumi{\arabic{enumi})}
\setcounter{enumi}{9}
\item
  给出乘法算子 \(m_{g_{1}}\) 与 \(m_{g_{2}}\) 的一个例子，使得 \(m_{g_{1}} \circ m_{g_{2}} \neq m_{g_{1}g_{2}}\) 。
\item
  我们把空间 \(\mathcal{D}(]0,1[)\) 与 \(\mathrm{L}^2 ([0,1])\) 的一个子空间等同起来。设 \(u\) 为 \(\mathrm{L}^2 ([0,1])\) 上以 \(\mathcal{D}(]0,1[)\) 为定义域、满足 \(u(f) = i^{-1}f'\) （对每个函数 \(f \in \mathcal{D}(]0,1[)\) ）的部分算子。
\end{enumerate}

\begin{enumerate}
\def\labelenumi{\alph{enumi})}
\item
  部分算子 \(u\) 是可闭的且 \(\overline{u}\) 是对称的。
\item
  \(\overline{u}\) 的属于 i 与属于 -i 的特征空间不是正交的。
\end{enumerate}

\begin{enumerate}
\def\labelenumi{\arabic{enumi})}
\setcounter{enumi}{11}
\item
  设 \(S \subset \mathbf{C}\) 为闭子集。给出一个稠定闭算子 \(u\) （定义在希尔伯特空间 \(E\) 上）的例子，使得 \(\operatorname{Sp}(u) = S\) 。
\item
  设 \(u\) 为希尔伯特空间 E 上的一个不可闭部分算子。不存在复数 \(\lambda\) 使得从 \(\text{dom}(u)\) 到 E 的、由 \(x \mapsto \lambda x - u(x)\) 定义的线性映射是双射且具有连续的逆。
\item
  设 \(u\) 为复希尔伯特空间 E 上的一个对称部分算子。映射 \(\lambda \mapsto \dim \operatorname{Ker}(\lambda 1_{\mathrm{E}} - u^{*})\) （从 \(\mathbf{C}\) 到 \(\overline{\mathbf{N}}\) ）在 \(u\) 的预解集上是局部常数的。
\item
  设 E 为复希尔伯特空间 \(\mathbf{C}^2\) ，\(u \in \mathcal{L}(\mathrm{E})\) 为满足 \(u(ae_1 + be_2) = be_1\) 的自同态。
\end{enumerate}

\begin{enumerate}
\def\labelenumi{\alph{enumi})}
\tightlist
\item
  证明 \(u\) 的谱退缩为 0，并且对所有 \(\lambda \in \mathbf{C}^*\) ，我们有公式
\end{enumerate}

\[
\| \mathrm{R} (u, \lambda) \| ^ {2} = \frac {2}{1 + 2 | \lambda | ^ {2} - \sqrt {1 + 4 | \lambda | ^ {2}}}.
\]

由此导出伪谱 \(\mathrm{PSp}_{\varepsilon}(u)\) （对所有 \(\varepsilon > 0\) ）的形状。

\begin{enumerate}
\def\labelenumi{\arabic{enumi})}
\setcounter{enumi}{15}
\item
  给出一个例子，说明 IV 第 251 页命题 21 的结论对谱的有界分支一般不成立。
\item
  \begin{enumerate}
  \def\labelenumii{\alph{enumii})}
  \tightlist
  \item
    设 U 为 C 的连通开子集，V ⊂ U 为开集。设 E 为复巴拿赫空间。设 \(f\colon\) U → E 为全纯映射，\(z_0 \in\) U 且 M \textgreater{} 0 满足：\(\|f(z)\| \leqslant\) M 对所有 \(z \in\) V 成立，且 \(\|f(z_0)\| <\) M 。则 \(\|f(z)\| <\) M 对所有 \(z \in\) U 成立。
  \end{enumerate}
\end{enumerate}

\begin{enumerate}
\def\labelenumi{\alph{enumi})}
\setcounter{enumi}{1}
\tightlist
\item
  设 E 为复巴拿赫空间，\(u \in \mathcal{L}(\mathrm{E})\) 。设 \(U \subset C\) 为 \(\mathbf{C}-\mathrm{Sp}(u)\) 的无界连通分量的一个开子集。若 M \textgreater{} 0 使得 \(\|R(u, \lambda)\| \leqslant M\) 对所有 \(\lambda \in U\) 成立，则 \(\|R(u, \lambda)\| < M\) 对所有 \(\lambda \in U\) 成立。特别地，若 \(\mathbf{C}-\mathrm{Sp}(u)\) 连通，则 u 的预解式的范数不可能在 \(\mathbf{C}-\mathrm{Sp}(u)\) 的非空开子集上取常值。
\end{enumerate}

\begin{enumerate}
\def\labelenumi{\arabic{enumi})}
\setcounter{enumi}{17}
\tightlist
\item
  设 E 为由复数组成的有界序列 \((x_{k})_{k\in\mathbf{Z}}\) 构成的复向量空间。对 E 中每个 \(x=(x_{k})\) ，我们记
\end{enumerate}

\[
\| x \| _ {\mathrm{E}} = | x _ {0} | + \sup _ {k \in \mathbf {Z} - \{0 \}} | x _ {k} |.
\]

\begin{enumerate}
\def\labelenumi{\alph{enumi})}
\item
  映射 \(x \mapsto \|x\|_{\mathrm{E}}\) 是 E 上的一个范数；赋以此范数的空间 E 是巴拿赫空间。
\item
  设 M \textgreater{} 2 为实数。设 \(\lambda = (\lambda_k)_{k \in \mathbf{Z}}\) 为由 \(\lambda_0 = M^{-1}\) 与 \(\lambda_k = 1\) （\(k \neq 0\) ）定义的序列。我们用 u 表示由 \(u((x_k)_{k \in \mathbf{Z}}) = (\lambda_k x_{k+1})_{k \in \mathbf{Z}}\) 定义的从 E 到 E 的线性映射。它是 E 的一个连续自同态。
\item
  以 0 为心、1/M 为半径的开圆盘包含于 u 的预解集。
\end{enumerate}

\(d)\) 对每个 \(\lambda\in\mathbf{C}\) （满足 \(|\lambda|<\inf(M^{-1},\frac{1}{2}-M^{-1})\) ），我们有 \(\|\mathrm{R}(u,\lambda)\|\leqslant\mathrm{M}\) 。（利用 IV 第 244 页公式 (6)。）

\begin{enumerate}
\def\labelenumi{\alph{enumi})}
\setcounter{enumi}{4}
\tightlist
\item
  设 \(e_{0} \in E\) 为序列 \((x_{k})_{k \in \mathbf{N}}\) ，它由 \(x_{0} = 1\) 与 \(x_{k} = 0\) （\(k \neq 0\) ）定义。设 \(\lambda \in C\) 满足 \(|\lambda| < \inf(M^{-1}, \frac{1}{2} - M^{-1})\) 。我们有 \(\|R(u, \lambda)e_{0}\|_{E} = M\) ，从而 \(\|R(u, \lambda)\| = M\) 。
\end{enumerate}

（本例与下一习题中的例子都归于 E. Shargorodsky，“On the level sets of the resolvent norm of a linear operator”，Bull. London Math. Soc. 40 (2008)，493–504。）

\begin{enumerate}
\def\labelenumi{\arabic{enumi})}
\setcounter{enumi}{18}
\tightlist
\item
  设 E 为希尔伯特空间 \(\ell_2(\mathbf{N}^*)\) 。设 \((\alpha_n)_{n \in \mathbf{N}^*}\) 为一列 \(\geqslant 2\) 的实数，使得 \(\alpha_n\) 趋于 \(+\infty\) （当 \(n \to +\infty\) ）。令 \(\beta_n = 1 + \alpha_n^{-1}\) （对所有 \(n \in \mathbf{N}^*\) ）。
\end{enumerate}

\begin{enumerate}
\def\labelenumi{\alph{enumi})}
\tightlist
\item
  设 \(\mathrm{F}\subset \mathrm{E}\) 为由满足下列条件的序列 \((x_{k})_{k\in \mathbf{N}^{*}}\in \mathrm{E}\) 组成的空间：级数
\end{enumerate}

\[
\sum_ {n \in \mathbf {N} ^ {*}} \alpha_ {n} ^ {2} | x _ {2 n} | ^ {2}
\]

收敛。它是 E 的一个稠密子空间。

\begin{enumerate}
\def\labelenumi{\alph{enumi})}
\setcounter{enumi}{1}
\tightlist
\item
  存在 E 上以 F 为定义域的一个闭部分算子 \(u\) ，使得对 F 中每个 \(x = (x_{k})_{k \in \mathbf{N}^{*}}\) ，我们有 \(u(x) = (y_{k})\) ，其中
\end{enumerate}

\[
\binom{y _ {2 k - 1}}{y _ {2 k}} = \left( \begin{array}{c c} 0 & \alpha_ {k} \\ \beta_ {k} & 0 \end{array} \right) \binom{x _ {2 k - 1}}{x _ {2 k}}
\]

对所有 \(k \in \mathbf{N}^*\) 成立。

\begin{enumerate}
\def\labelenumi{\alph{enumi})}
\setcounter{enumi}{2}
\tightlist
\item
  设 \(u_{n}\) 为自同态
\end{enumerate}

\[
x \mapsto \left( \begin{array}{c c} 0 & \alpha_ {n} \\ \beta_ {n} & 0 \end{array} \right) x
\]

（配备典范内积的希尔伯特空间 \(C^{2}\) 上的）。对每个 \(\lambda \in C\) （满足 \(|\lambda| < 1\) ），\(u_{n}\) 在 \(\lambda\) 处的预解式有定义，并且我们有

\[
\lim _ {n \to + \infty} \| \mathrm{R} (u _ {n}, \lambda) \| = 1.
\]

\begin{enumerate}
\def\labelenumi{\alph{enumi})}
\setcounter{enumi}{3}
\tightlist
\item
  对每个 \(n \in N^{*}\) 与每个 \(\lambda \in C\) （满足 \(|\lambda| < \frac{1}{2}\) ），我们有 \(\|R(u_{n}, \lambda)\| < 1\) 。
\item
  以 0 为心、半径 \(\frac{1}{2}\) 的开圆盘包含于 u 的预解集，且我们有 \(\|R(u,\lambda)\|=1\) 对所有 \(\lambda\in C\) （满足 \(|\lambda|<\frac{1}{2}\) ）成立。
\end{enumerate}

\begin{enumerate}
\def\labelenumi{\arabic{enumi})}
\setcounter{enumi}{19}
\tightlist
\item
  若 \(a\) 与 \(b\) 为满足 \(a < b\) 的实数，我们用 \(\mathcal{C}^2([a,b])\) 表示 \([a,b]\) 上连续、且其限制到 \(]a,b[\) 上为 \(\mathrm{C}^2\) 类的函数所成的向量空间。我们用 \(\mu\) 表示勒贝格测度。
\end{enumerate}

\begin{enumerate}
\def\labelenumi{\alph{enumi})}
\tightlist
\item
  设 \(q_{1}\) 与 \(q_{2}\) 为 \(\mathcal{C}^{2}([a,b])\) 中满足 \(q_{1} \geqslant q_{2}\) 的实值函数。对 \(i \in \{1,2\}\) ，设 \(f_{i} \in \mathcal{C}^{2}([a,b])\) 为方程
\end{enumerate}

\[
f _ {i} ^ {\prime \prime} + q _ {i} f _ {i} = 0.
\]

假设 \(f_{2}(a) = f_{2}(b) = 0\) 且 \(f_{2} \neq 0\) 。若函数 \(f_{1}\) 在 \(]a, b[\) 中不取零值，则 \(q_{1} = q_{2}\) ，且存在 \(c \in \mathbf{C}\) 使 \(f_{1} = cf_{2}\) 。（化归到 \(f_1\) 与 \(f_2\) 在 \(]a, b[;\) 上严格为正的情形；考虑朗斯基行列式 \(w = f_1f_2' - f_1'f_2\) ，先证明 \(w(a) = w(b)\) ，再证明 \(w = 0.\) ）

\begin{enumerate}
\def\labelenumi{\alph{enumi})}
\setcounter{enumi}{1}
\tightlist
\item
  设 \(q \in \mathcal{C}^2([0,1])\) 满足 \(q \geqslant 0\) 。我们用 \(\widetilde{\mathrm{SL}}_q\) 表示 \(\mathrm{L}^2([0,1],\mu)\) 上以
\end{enumerate}

\[
\mathrm{D} = \{f \in \mathcal {C} ^ {2} ([ 0, 1 ]) \mid f (0) = f (1) = 0 \}
\]

为定义域、满足

\[
\widetilde {\mathrm{SL}} _ {q} (f) = - f ^ {\prime \prime} + q f
\]

的部分微分算子（斯图姆–刘维尔算子）。算子 \(\widetilde{SL}_{q}\) 是对称的且本质自伴。我们用 \(SL_{q}\) 表示 \(\widetilde{SL}_{q}\) 的闭包。

\begin{enumerate}
\def\labelenumi{\alph{enumi})}
\setcounter{enumi}{2}
\tightlist
\item
  设 \(\lambda \in \mathrm{Sp}(\mathrm{SL}_q)\) 为 \(\mathrm{SL}_q\) 的一个特征值。与 \(\lambda\) 相伴的特征空间维数为 1，且由一个属于 D 的实值函数生成。
\end{enumerate}

\(d\) ) \(\lambda\) 作为 \(\mathrm{SL}_q\) 的特征值，满足 \(\lambda \geqslant \pi^2\) 。（先证明 \(\lambda\) 是正的：利用问题 \(a\) ）取 \(q_{1} = q\) 与 \(q_{2} = 0\) 。）

\begin{enumerate}
\def\labelenumi{\alph{enumi})}
\setcounter{enumi}{4}
\tightlist
\item
  对 s \textgreater{} 0，存在唯一的函数 \(f_{s} \in \mathcal{C}^{2}([0,1])\) 是线性微分方程
\end{enumerate}

\[
- f _ {s} ^ {\prime \prime} + q f - s ^ {2} f = 0
\]

的解，且满足 \(f_{s}(0)=0\) 与 \(f_{s}^{\prime}(0)=s\) 。令 \(f(x,s)=f_{s}(x)\) ；映射 f 是 \(C^{2}\) 类函数，定义在 {]}0,1{[}× \(R_{+}^{*}\) 上。\(SL_{q}\) 的特征值形如 \(\lambda=s^{2}\) ，其中 s 满足 \(f(1,s)=0\) 。

\(f)\) 对 \(t \in [0,1]\) 与 \(s > 0\) ，设 \(\mathrm{N}(t,s) \in \overline{\mathbf{N}}\) 为 \(x \in [0,t]\) 中满足 \(f(x,s) = 0\) 者的个数。我们有 \(\mathrm{N}(t,s) \in \mathbf{N}\) ，且函数 \(s \mapsto \mathrm{N}(t,s)\) 对每个固定的 \(t\) 是递增的。（利用问题 \(a\) ）。）

\(g)\) 设 \(t \in [0,1]\) 且 \(s_0 > 0\) 。设 \(m = \mathrm{N}(t,s_0) \in \mathbf{N}\) 。用

\[
0 = x _ {1} (s _ {0}) <   x _ {2} (s _ {0}) <   \dots <   x _ {m} (s _ {0})
\]

表示 \(f_{s_0}\) 在 \([0,t]\) 中的零点。对 \(s \geqslant 0\) ，用

\[
0 = x _ {1} (s) <   \dots <   x _ {m} (s)
\]

表示前 \(m\) 个零点（\(f_{s}\) 在 \([0,t]\) 中的）。函数 \(s\mapsto x_{j}(s)\) 是 \(\mathbf{C}^1\) 类的，并且我们有

\[
\partial_ {1} (x _ {j} (s), s) x _ {j} ^ {\prime} (s) = - \partial_ {2} f (x _ {j} (s), s).
\]

\begin{enumerate}
\def\labelenumi{\alph{enumi})}
\setcounter{enumi}{7}
\tightlist
\item
  \(SL_{q}\) 的特征值之集是某个严格正实数递增序列的取值集
\end{enumerate}

\[
0 <   \lambda_ {0} <   \lambda_ {1} <   \dots <   \lambda_ {n} <   \dots
\]

该序列趋于 \(+\infty\) 。此外，与 \(\lambda_{n}\) 相应的特征空间维数为 1；若 \(\varphi_{n}\) 是它的一个基，则 \(\varphi_{n} \in \mathcal{C}^{2}([0,1])\) ，且函数 \(\varphi_{n}\) 恰在 \(n\) 个不同点处（\([0,1]\) 中）取零。（为证明 \(\mathrm{SL}_{q}\) 存在无穷多个特征值，利用问题 \(a\) ）取 \(q_{1} = M - s^{2}\) ，其中 \(M \in R\) 选得适当，验证 \(N(1, s)\) 趋于 \(+\infty\) （当 \(s \to +\infty\) ）。）

\begin{enumerate}
\def\labelenumi{\arabic{enumi})}
\setcounter{enumi}{20}
\tightlist
\item
  我们沿用上一习题的记号。我们用 \(\varphi_{n}\) 表示 \(\mathrm{SL}_{q}\) 的属于特征值 \(\lambda_{n}\) 且满足 \(\varphi_{n}^{\prime}(0)=\lambda_{n}\) 的唯一特征函数。
\end{enumerate}

\begin{enumerate}
\def\labelenumi{\alph{enumi})}
\tightlist
\item
  设 \(\mathrm{M} = \sup \{q(x)\}\) 。对每个 \(n \in \mathbf{N}\) ，我们有
\end{enumerate}

\[
\pi^ {2} (n + 1) ^ {2} \leqslant \lambda_ {n} \leqslant \pi^ {2} (n + 1) ^ {2} + \mathrm{M},
\]

并且此外

\[
\lambda_ {n} = \pi^ {2} (n + 1) ^ {2} \Bigl (1 + \mathrm{O} \Bigl (\frac {1}{n ^ {2}} \Bigr) \Bigr)
\]

当 \(n \rightarrow +\infty\) 时成立。

\begin{enumerate}
\def\labelenumi{\alph{enumi})}
\setcounter{enumi}{1}
\tightlist
\item
  设 \(s > 0\) 。设 \(f \in \mathcal{C}^2([0,1])\) 满足
\end{enumerate}

\[
- f ^ {\prime \prime} + s ^ {2} f = - q f, \qquad f (0) = 0, \quad f ^ {\prime} (0) = s.
\]

于是我们有 \((1-s^{-1}u_{s})f=g_{s}\) ，其中 \(g_{s}(x)=\sin(sx)\) ，而 \(u_{s}\) 是 \(\mathcal{C}([0,1])\) 上由下式给出的积分算子：

\[
u _ {s} (f) (x) = \int_ {0} ^ {x} \sin (s (x - y)) q (y) f (y) d \mu (y).
\]

\begin{enumerate}
\def\labelenumi{\alph{enumi})}
\setcounter{enumi}{2}
\item
  当 \(s\) 充分大时，我们有 \(s \notin \operatorname{Sp}(u_s)\) 。
\item
  对每个固定的 x，我们有
\end{enumerate}

\[
\varphi_ {n} (x) = \sin ((n + 1) x) + \mathrm{O} (n ^ {- 1})
\]

当 \(n \rightarrow +\infty\) 时成立。

\begin{enumerate}
\def\labelenumi{\arabic{enumi})}
\setcounter{enumi}{21}
\tightlist
\item
  我们沿用上一习题的记号。
\end{enumerate}

\begin{enumerate}
\def\labelenumi{\alph{enumi})}
\tightlist
\item
  存在函数 \(f_{1}\) 与 \(f_{2}\) 属于 \(\mathcal{C}^{2}([a,b])\) ，满足 \(-f_{i}^{\prime\prime} + qf_{i} = 0\) 且
\end{enumerate}

\[
\begin{array}{c c} f _ {1} (0) = 0, & f _ {1} ^ {\prime} (0) = 1 \\ f _ {2} (1) = 0, & f _ {2} ^ {\prime} (1) = - 1. \end{array}
\]

这些函数是这些方程的唯一解，且函数 \(w = f_{1}f_{2}' - f_{1}'f_{2}\) 是常数。

\begin{enumerate}
\def\labelenumi{\alph{enumi})}
\setcounter{enumi}{1}
\tightlist
\item
  设 \(g \in \mathcal{C}([0,1])\) 。存在唯一的函数 \(\varphi \in \mathcal{C}^2([0,1])\) 满足
\end{enumerate}

\[
- \varphi^ {\prime \prime} + q \varphi = g,
\]

并且我们有 \(\varphi = c_{1}f_{1} + c_{2}f_{2}\) ，其中

\[
c _ {1} (x) = - \frac {1}{w} \int_ {x} ^ {1} g (y) f _ {2} (y) d \mu (y), \qquad c _ {2} (x) = - \frac {1}{w} \int_ {0} ^ {x} g (y) f _ {1} (y) d \mu (y)
\]

对所有 \(x \in [0,1]\) 成立。

\begin{enumerate}
\def\labelenumi{\alph{enumi})}
\setcounter{enumi}{2}
\tightlist
\item
  部分算子 \(\mathrm{SL}_q\) 是满射的；其逆是映射 \(g\mapsto \varphi\) ，且存在函数 \(\mathrm{G}\in \mathcal{C}([0,1]^2)\) 使得
\end{enumerate}

\[
\mathrm{SL} _ {q} ^ {- 1} (g) (x) = \int_ {0} ^ {1} \mathrm{G} (x, y) g (y) d \mu (y)
\]

对所有 \(g \in \mathcal{C}([0,1])\) 与几乎所有 \(x \in [0,1]\) 成立。我们有 \(G(x,y) = G(y,x)\) 对所有 \((x,y) \in [0,1]^2\) 成立。

\(d)\) 算子 \(\mathrm{SL}_q^{-1}\) 在 \(\mathrm{L}^2 ([0,1],\mu)\) 上是紧的且单射的。

\begin{enumerate}
\def\labelenumi{\alph{enumi})}
\setcounter{enumi}{4}
\tightlist
\item
  诸函数 \((\varphi_{n})\) 构成 \(\mathrm{L}^{2}([0,1])\) 的一个标准正交基，且 \(SL_{q}\) 的谱与其特征值之集重合。
\end{enumerate}

\begin{enumerate}
\def\labelenumi{\arabic{enumi})}
\setcounter{enumi}{22}
\tightlist
\item
  设 \(E = \ell_{\mathbf{C}}^{2}(\mathbf{N})\) ，u 为由
\end{enumerate}

\[
u ((x _ {i}) _ {i \in \mathbf {N}}) = (0, x _ {1}, \dots , x _ {n}, \dots)
\]

对所有 \((x_{i})\in \mathrm{E}\) 定义的自同态。

\begin{enumerate}
\def\labelenumi{\alph{enumi})}
\item
  \(u\) 的谱是 \(\mathbf{C}\) 中的闭单位圆盘。
\item
  u 的剩余谱是 C 中的开单位圆盘。
\item
  \(u\) 的本质谱是 \(\mathbf{C}\) 中的单位圆周。
\end{enumerate}

\begin{enumerate}
\def\labelenumi{\arabic{enumi})}
\setcounter{enumi}{23}
\tightlist
\item
  设 E 为复希尔伯特空间，u 为 E 上的一个对称部分算子。设 \(F = E \oplus E\) ，v 为 F 上以 \(\text{dom}(u^{*}) \oplus \text{dom}(\overline{u})\) 为定义域、满足 \(v(x, y) = (\overline{u}(y), u^{*}(x))\) 的部分算子。
\end{enumerate}

\begin{enumerate}
\def\labelenumi{\alph{enumi})}
\item
  部分算子 v 是自伴的。
\item
  设 \(\widetilde{\mathrm{E}}\) 为线性映射 \(\delta\colon x\mapsto(x,x)\) 的像。\(v\) 到 \(\widetilde{\mathrm{E}}\) 的约化是以 \(\delta(\mathrm{dom}(u))\) 为定义域、由 \((x,x)\mapsto(u(x),u(x))\) 给出的部分算子。
\end{enumerate}

\exercisesection{}

除非另有说明，以下的希尔伯特空间均假定为复空间。

\begin{enumerate}
\def\labelenumi{\arabic{enumi})}
\tightlist
\item
  设 E 为希尔伯特空间，\(\mathcal{P}(\mathrm{E})\subset \mathcal{L}(\mathrm{E})\) 为 E 中正交投影组成的集合。设 S 为 C 的一个普遍可测子集。我们用 \(\mathcal{T}(\mathrm{S})\) 表示由满足 \(\mathrm{T}\subset \mathrm{S}\) 的 C 的普遍可测子集 T 组成的集合。S 上一个取值于 E 的投影的测度是指一个映射 \(\varpi\) （从 \(\mathcal{T}(\mathrm{S})\) 到 \(\mathcal{P}(\mathrm{E})\) ），它满足
\end{enumerate}

\begin{enumerate}
\def\labelenumi{(\roman{enumi})}
\item
  我们有 \(\varpi(\varnothing)=0\) 且 \(\varpi(S)=1_{E}\) 。
\item
  若 \((\mathrm{T}_{n})\) 是 \(\mathcal{T}(\mathrm{S})\) 中两两不相交的集合的序列，且 T 是诸集合 \(T_{n}\) 的并，则
\end{enumerate}

\[
\varpi (\mathrm{T}) = \sum_ {n \in \mathbf {N}} \varpi (\mathrm{T} _ {n})
\]

在赋有逐点收敛拓扑的空间 \(\mathcal{L}(\mathrm{E})\) 中成立。

对 E 上每个谱等于 S 的正规部分算子 u，以及对 \(T \in \mathcal{T}(S)\) ，我们用 \(p_{T,u}\) 表示由 T 定义的 u 的谱投影。

把 E 上谱为 S 的正规部分算子 \(u\) 对应于映射 \(\mathrm{T} \mapsto p_{\mathrm{T}, u}\) 的对应，是谱为 S 的正规部分算子之集与 S 上取值于 E 的投影的测度之集之间的一个双射。

\begin{enumerate}
\def\labelenumi{\arabic{enumi})}
\setcounter{enumi}{1}
\tightlist
\item
  设 E 为复希尔伯特空间，u 为 E 上的一个正规部分算子。我们用 \(p_{A}\) 表示 u 在 A 上的谱投影。设 a \textless{} b 为实数。我们有
\end{enumerate}

\[
\frac {1}{2} (p _ {[ a, b ]} + p _ {] a, b [}) = \lim _ {\varepsilon \rightarrow 0} \frac {1}{2 i \pi} \int_ {a} ^ {b} \big (\mathrm{R} (u, t - i \varepsilon) - \mathrm{R} (u, t + i \varepsilon) \big) d t,
\]

在赋有逐点收敛拓扑的空间 \(\mathcal{L}(\mathrm{E})\) 中成立（“斯通公式”）。

\begin{enumerate}
\def\labelenumi{\arabic{enumi})}
\setcounter{enumi}{2}
\item
  设 E 为可数型的希尔伯特空间。\(\mathcal{L}(\mathrm{E})\) 仅有的闭双边理想是 \(\{0\}\) 、\(\mathcal{L}^{\mathrm{c}}(\mathrm{E})\) 与 \(\mathcal{L}(\mathrm{E})\) 。（利用 IV 第 316 页习题 10；若 I 是这样一个理想，假设 I 中存在一个非紧的非零自伴元 \(u\) ，并证明：对 R 中每个使 \(0 \notin J\) 的开区间 J ，\(u\) 在 J 上的谱投影属于 I。）
\item
  设 E 为希尔伯特空间，u 为 E 上的一个自伴部分算子。映射 \(\sigma\colon f\mapsto f(u)\) （从 \(\mathcal{L}_{\mathrm{u}}^{\infty}(\mathrm{Sp}(u))\) 到 \(\mathcal{L}(\mathrm{E})\) ）是满足下列三个条件的唯一的对合代数态射：
\end{enumerate}

\begin{enumerate}
\def\labelenumi{(\roman{enumi})}
\item
  \(\| \sigma (f)\| \leqslant \| f\|_{\infty}\) 对所有 \(f\in \mathcal{L}_{\mathrm{u}}^{\infty}(\mathrm{Sp}(u))\) 成立
\item
  若 \((f_n)_{n\in \mathbf{N}}\) 是 \(\mathcal{L}_{\mathrm{u}}^{\infty}(\mathrm{Sp}(u))\) 中的有界序列，且 \(f_{n}(t)\to f(t)\) 对所有 \(t\in \mathrm{Sp}(u)\) 成立，则 \(\sigma (f_n)\rightarrow \sigma (f)\) 在赋有逐点收敛拓扑的 \(\mathcal{L}(\mathrm{E})\) 中成立。
\item
  若 \((f_n)_{n\in \mathbf{N}}\) 是 \(\mathcal{L}_{\mathrm{u}}^{\infty}(\mathrm{Sp}(u))\) 中的序列，满足 \(|f_n(t)|\leqslant |t|\) 且 \(f_{n}(t)\to t\) 对所有 \(t\in \mathrm{Sp}(u)\) 成立，则 \(u(x) = \lim \sigma (f_n)(x)\) 对所有 \(x\in \mathrm{dom}(u)\) 成立。
\end{enumerate}

\begin{enumerate}
\def\labelenumi{\arabic{enumi})}
\setcounter{enumi}{4}
\item
  设 u 为希尔伯特空间 E 上的一个正规部分算子。u 的剩余谱是空的。
\item
  设 E 与 F 为希尔伯特空间。设 u 为从 E 到 F 的一个稠定闭部分算子。设 \(\lambda \in C^{*}\) 属于 E 上自伴算子 \(u^{*} \circ u\) 的预解集（参见 IV 第 241 页命题 12）。
\end{enumerate}

\begin{enumerate}
\def\labelenumi{\alph{enumi})}
\item
  从 F 到 F 的部分算子 \(u \circ u^{*}\) 是自伴的。
\item
  对每个向量 \(x \in \mathrm{dom}(u^{*})\) ，向量
\end{enumerate}

\[
y = (1 _ {\mathrm{F}} + u \circ \mathrm{R} (u ^ {*} \circ u, \lambda) \circ u ^ {*}) x
\]

属于 dom \((u\circ u^{*})\) ，且公式

\[
(u \circ u ^ {*} - \lambda 1 _ {\mathrm{F}}) y = x
\]

成立。

\begin{enumerate}
\def\labelenumi{\alph{enumi})}
\setcounter{enumi}{2}
\tightlist
\item
  部分算子 \(u \circ \mathrm{R}(u^* \circ u, \lambda) \circ u^*\) 稠定，且在其定义域上是连续的。
\end{enumerate}

\(d)\) 部分算子 \(u \circ u^{*} - \lambda 1_{\mathrm{F}}\) 是满射的。

\begin{enumerate}
\def\labelenumi{\alph{enumi})}
\setcounter{enumi}{4}
\item
  我们有 \(\lambda \notin \operatorname{Sp}(u \circ u^{*})\) 。
\item
  我们有恒等式
\end{enumerate}

\[
\lambda \mathrm{R} (u \circ u ^ {*}, \lambda) - u \circ \mathrm{R} (u ^ {*} \circ u, \lambda) \circ u ^ {*} = 1 _ {\mathrm{F}}
\]

在 \(\mathcal{L}(\mathrm{F})\) 中成立（其中 \(u\circ \mathrm{R}(u^{*}\circ u,\lambda)\circ u^{*}\) 按记号滥用表示 F 的与此部分算子在其定义域上相重合的那个自同态，参见问题 \(c\) ））。

\begin{enumerate}
\def\labelenumi{\alph{enumi})}
\setcounter{enumi}{6}
\tightlist
\item
  我们有关系式
\end{enumerate}

\[
u ^ {*} \circ \mathrm{R} (u \circ u ^ {*}, \lambda) \subset \mathrm{R} (u ^ {*} \circ u, \lambda) \circ u ^ {*}.
\]

\begin{enumerate}
\def\labelenumi{\alph{enumi})}
\setcounter{enumi}{7}
\tightlist
\item
  部分算子 u 通过过渡到子空间，诱导出从 \(\operatorname{Ker}(u^{*} \circ u - \lambda1_{\mathrm{E}})\) 到 \(\operatorname{Ker}(u \circ u^{*} - \lambda1_{\mathrm{F}})\) 的一个同构。
\end{enumerate}

\begin{enumerate}
\def\labelenumi{\arabic{enumi})}
\setcounter{enumi}{6}
\tightlist
\item
  设 A 为希尔伯特空间 E 上一族稠定闭算子。设 \(\lambda \in \mathbf{C}\) 。我们称 A 是 \(\lambda\) -交换的，如果 \(\lambda\) 属于 \(u\) 的预解集（对所有 \(u \in \mathrm{A}\) ），且 E 的诸自同态 \(\mathrm{R}(u, \lambda)\) （\(u \in \mathrm{A}\) ）两两可交换。
\end{enumerate}

我们称 A 是预解意义上交换的，如果存在 \(\lambda \in \mathbf{C}\) 使 A 是 \(\lambda\) -交换的。

由单个稠定闭算子 u 组成的集是交换的，当且仅当 u 的预解集非空。

\begin{enumerate}
\def\labelenumi{\alph{enumi})}
\item
  若 A ⊂ L(E)，则 A 是预解意义上交换的，当且仅当 A 中各元素的谱之并不等于 C，且 A 是交换的。
\item
  设 A 为 E 上一族 λ-交换的稠定闭算子。则对属于每个 \(u \in A\) 的预解集的每个复数 μ，A 是 μ-交换的。
\item
  设 A 为 E 上一族预解意义上交换的稠定闭算子，且 \(u^{*} \in A\) 对所有 \(u \in A\) 成立。存在局部紧拓扑空间 X、X 上的正测度 \(\mu\) 、一个等距同构 \(v\) （从 \(\mathrm{L}^2 (\mathrm{X},\mu)\) 到 E ），以及映射 \(u\mapsto g_u\) （从 A 到 \(\mu\) -可测函数之集） ，使得 \(u = v\circ m_{g_u}\circ v^{-1}\) 对所有 \(u\in \mathrm{A}\) 成立。
\end{enumerate}

\begin{enumerate}
\def\labelenumi{\arabic{enumi})}
\setcounter{enumi}{7}
\tightlist
\item
  设 u 与 v 为希尔伯特空间 E 上的正规部分算子。下列断言等价：
\end{enumerate}

\begin{enumerate}
\def\labelenumi{\alph{enumi})}
\item
  自同态 \(b(u)\) 与 \(b(v)\) 可交换；
\item
  对 C 的每个普遍可测子集 A，u 与 v 在 A 上的谱投影可交换；
\end{enumerate}

若 \(u\) 与 \(v\) 的预解集非空，这些断言等价于

\begin{enumerate}
\def\labelenumi{\alph{enumi})}
\setcounter{enumi}{2}
\tightlist
\item
  对每个 \(\lambda \notin \operatorname{Sp}(u)\) 与 \(\lambda' \notin \operatorname{Sp}(v)\) ，预解式 \(\operatorname{R}(u, \lambda)\) 与 \(\operatorname{R}(v, \lambda')\) 可交换。
\end{enumerate}

若 \(u\) 与 \(v\) 自伴，这些断言等价于

\begin{enumerate}
\def\labelenumi{\alph{enumi})}
\setcounter{enumi}{3}
\tightlist
\item
  对每个 \(t \in R\) ，酉自同态 \(e^{itu}\) 与 \(e^{itv}\) 可交换。
\end{enumerate}

\begin{enumerate}
\def\labelenumi{\arabic{enumi})}
\setcounter{enumi}{8}
\tightlist
\item
  设 u 为希尔伯特空间 E 上的一个稠定闭部分算子。下列断言等价：
\end{enumerate}

\begin{enumerate}
\def\labelenumi{\alph{enumi})}
\item
  算子 u 是正规的；
\item
  存在 E 上预解意义上可交换的自伴算子 \(u_{1}\) 与 \(u_{2}\) ，使得 \(u = u_{1} + iu_{2}\) ；
\item
  我们有 \(\operatorname{dom}(u) = \operatorname{dom}(u^*)\) ，且 \(\| u(x)\| = \| u^{*}(x)\|\) 对所有 \(x\in \operatorname{dom}(u)\) 成立；
\item
  我们有 \(uu^{*} = u^{*}u\) 。
\end{enumerate}

若 \(u\) 的预解集非空，这些断言等价于

\begin{enumerate}
\def\labelenumi{\alph{enumi})}
\setcounter{enumi}{4}
\tightlist
\item
  对每个 \(\lambda \notin \operatorname{Sp}(u)\) ，预解式 \(\operatorname{R}(u, \lambda)\) 是正规的。
\end{enumerate}

\begin{enumerate}
\def\labelenumi{\arabic{enumi})}
\setcounter{enumi}{9}
\tightlist
\item
  设 E 为复希尔伯特空间，A 为 \(\mathcal{L}(\mathrm{E})\) 的一个含单位的自伴子代数，它在逐点收敛拓扑下是闭的。（这样的代数称为冯·诺伊曼代数。）
\end{enumerate}

\begin{enumerate}
\def\labelenumi{\alph{enumi})}
\item
  存在 \(\mathcal{L}(\mathrm{E})\) 的一个子集 X，使 A 与 X 在 \(\mathcal{L}(\mathrm{E})\) 中的交换子相重合。（利用 IV 第 326 页习题 24。）反之，对 X 的每个自伴子集，X 在 \(\mathcal{L}(\mathrm{E})\) 中的交换子是一个含单位的自伴子代数，且在逐点收敛拓扑下是闭的。
\item
  设 \(u \in \mathrm{A}\) ，\((j, |u|)\) 为 \(u\) 的极分解（I 第 140 页定义 4）。元素 \(j\) 与 \(u\) 属于 A。
\item
  设 \(u \in A\) 且 u 正规。对每个函数 \(f \in \mathcal{L}_{\mathrm{u}}^{\infty}(\mathrm{Sp}(u))\) ，自同态 \(f(u)\) 属于 A。
\item
  设 \(u \in A\) 。元素 u 是 A 中的右（相应地，左）零因子（参见 A，I，第 93 页）当且仅当 u 的像在 E 中不稠密（相应地，当且仅当 u 不是单射的）。
\item
  设 F 为 E 上部分算子组成的集，其元素是满足下列条件的稠定闭部分算子 u：对 X 中每个 x，部分算子 \(x \circ u\) 是部分算子 \(u \circ x\) 的一个约化。对 F 中每个 u 与 A 中每个 v，部分算子 \(u^{*}\) 、\(u + v\) 、\(v \circ u\) 与 \(u \circ v\) 都属于 F。
\end{enumerate}

\(f\) ) 设 \(u\in \mathrm{F}\) ，\((j,|u|)\) 为 \(u\) 的极分解（IV 第 290 页定义 7）。部分算子 \(|u|\) 属于 F，而 \(j\) 属于 A。

\(g)\) 设 \(u \in \mathrm{F}\) 且 \(u\) 正规。对每个函数 \(f \in \mathcal{L}_{\mathrm{u}}(\mathrm{Sp}(u))\) ，部分算子 \(f(u)\) 属于 F；若 \(f\) 有界，则 \(f(u)\) 属于 A。

\begin{enumerate}
\def\labelenumi{\alph{enumi})}
\setcounter{enumi}{7}
\item
  设 \(u \in A\) 单射且其像在 E 中稠密。部分算子 \(u^{-1}\) 属于 F。
\item
  对 F 中每个 u，存在 A 的元素 \(u_{1}\) 与 \(u_{2}\) ，使 \(u_{2}\) 单射、像稠密，且 \(u = u_{1}u_{2}^{-1}\) 。（先化归到 u 自伴且正的情形，再对函数 \(f_{1}: t \mapsto t\varphi(t) + \psi(t)\) 与 \(f_{2}: t \mapsto \varphi(t) + t^{-1}\psi(t)\) 应用普遍可测函数演算，其中 \(\varphi\) 与 \(\psi\) 分别是区间 {[}0, 1{[} 与 {[}1, +∞{[} 的特征函数。{]})
\item
  假设 A 中每个像稠密的单射元素都是满射的。则代数 A 满足 Ore 条件：对 A 中任意 u 与 s，若 s 不是零因子（既非右零因子也非左零因子），则存在 A 中 \(u'\) 与 \(s'\) ，使 \(s'\) 不是零因子且 \(us' = su'\) 。
\end{enumerate}

\(k\) ) 上一问题的结论并不总是成立。（考虑 \(\mathrm{A} = \mathcal{L}(\mathrm{E}).\) ）

\begin{enumerate}
\def\labelenumi{\arabic{enumi})}
\setcounter{enumi}{10}
\tightlist
\item
  \begin{enumerate}
  \def\labelenumii{\alph{enumii})}
  \tightlist
  \item
    设 \(q\) 为复希尔伯特空间 E 上的一个正部分形式。则 \(q\) 是闭的，当且仅当：对每个序列 \((x_{n})\) （由 \(\text{dom}(q)\) 中元素组成），若它收敛于 E 中 \(x\) 且满足
  \end{enumerate}
\end{enumerate}

\[
\lim _ {n, m \to + \infty} q (x _ {n} - x _ {m}, x _ {n} - x _ {m}) = 0,
\]

，我们有 \(x \in \text{dom}(q)\) 且 \(q(x_{n} - x, x_{n} - x) \to 0\) 。

\begin{enumerate}
\def\labelenumi{\alph{enumi})}
\setcounter{enumi}{1}
\tightlist
\item
  设 \(E = L^{2}(\mathbf{R}, \mu)\) ，其中 \(\mu\) 为勒贝格测度。以 \(\mathcal{D}(\mathbf{R})\) 为定义域、由 \(q(\varphi_{1}, \varphi_{2}) = \overline{\varphi_{1}(0)}\varphi_{2}(0)\) 定义的正埃尔米特形式 q 不是闭的，且不存在扩张 q 的闭的正部分形式。
\end{enumerate}

\begin{enumerate}
\def\labelenumi{\arabic{enumi})}
\setcounter{enumi}{11}
\tightlist
\item
  我们在 \(\mathbf{R}\) 上配备勒贝格测度。设 E 为 \(\mathrm{L}^2 (\mathbf{R})\) 中由那些 \(f\in \mathrm{L}^2 (\mathbf{R})\) （满足 \(f(-x) = f(x)\) ，对几乎所有 \(x\in \mathbf{R}\) ）组成的闭子空间。设 \(\mathcal{D}_{+}(\mathbf{R})\) 为 \(\mathbf{R}\) 上满足 \(f(x) = f(-x)\) 对所有 \(x\in \mathbf{R}\) 成立的测试函数组成的空间。我们把 \(\mathcal{D}_{+}(\mathbf{R})\) 与 E 的一个子空间等同起来。我们用 \(\mathcal{F}\in \mathcal{L}(\mathrm{E})\) 表示由 \(\mathrm{L}^2 (\mathbf{R})\) 的傅里叶变换通过过渡到子空间而导出的 E 的自同态。
\end{enumerate}

\begin{enumerate}
\def\labelenumi{\alph{enumi})}
\item
  对每个 \(f \in \mathcal{D}_{+}(\mathbf{R})\) ，我们有 \(\mathcal{F}(f) + \mathcal{F}^{*}(f) \in \mathcal{D}_{+}(\mathbf{R})\) 。
\item
  设 u 为 E 上以 \(\mathcal{D}_{+}(\mathbf{R})\) 为定义域、由 \(u(f) = \mathcal{F}(f) + \mathcal{F}^{*}(f)\) （对 \(f \in \mathcal{D}_{+}(\mathbf{R})\) ）定义的部分算子。算子 u 本质自伴。
\item
  我们有 \(\operatorname{dom}(u^{2})=\{0\}\) （利用 II 第 273 页习题 28。）
\end{enumerate}

\begin{enumerate}
\def\labelenumi{\arabic{enumi})}
\setcounter{enumi}{12}
\tightlist
\item
  设 E 为复希尔伯特空间，u 为 E 上的一个对称部分算子。假设存在实数 c \textgreater{} 0 ，使 \(\langle x \mid u(x) \rangle \geqslant c \|x\|^{2}\) 对所有 \(x \in \text{dom}(u)\) 成立。
\end{enumerate}

证明下列条件等价：

\begin{enumerate}
\def\labelenumi{(\roman{enumi})}
\item
  部分算子 \(u\) 本质自伴；
\item
  u 的像在 E 中稠密；
\item
  \(u^{*}\) 的核退缩为 0。
\end{enumerate}

\begin{enumerate}
\def\labelenumi{\arabic{enumi})}
\setcounter{enumi}{13}
\tightlist
\item
  设 \(n \in \{1, 2, 3\}\) 。我们用 \(\Delta\) 表示 \(\mathrm{L}^2(\mathbf{R}^n)\) 上的拉普拉斯算子，用 D 表示其定义域。它是一个正的部分算子。
\end{enumerate}

\begin{enumerate}
\def\labelenumi{\alph{enumi})}
\tightlist
\item
  存在实数 \(c \geqslant 0\) ，使对每个 \(f \in D\) ，我们有
\end{enumerate}

\[
\| f \| _ {\infty} \leqslant c (\| \Delta f \| + \| f \|).
\]

\begin{enumerate}
\def\labelenumi{\alph{enumi})}
\setcounter{enumi}{1}
\tightlist
\item
  对每个实数 \(a > 0\) ，存在实数 \(b \geqslant 0\) ，使对每个 \(f \in \mathrm{D}\) ，我们有
\end{enumerate}

\[
\| f \| _ {\infty} \leqslant a \| \Delta f \| + b \| f \|.
\]

（对 \(x \mapsto f(\alpha x)\) 应用上一问题，其中 \(\alpha \in \mathbf{R}_{+}^{*}\) 。）

\begin{enumerate}
\def\labelenumi{\alph{enumi})}
\setcounter{enumi}{2}
\item
  设 \(g \in \mathcal{L}^{\infty}(\mathbf{R}^{n}) + \mathcal{L}^{2}(\mathbf{R}^{n})\) 。若 \(g\) 是实值的，则部分算子 \(\Delta + m_{g}\) 自伴。
\item
  设 \(\alpha\) 为满足 \(0 \leqslant \alpha < n/2\) 的实数。设 N 为 \(\mathbf{R}^n\) 上的欧几里得范数，\(g\) 为 \(\mathbf{R}^n\) 上由 \(g(x) = 1/\mathrm{N}(x)^\alpha\) （\(x \neq 0\) 时）与 \(g(0) = 0\) 定义的函数。部分算子 \(\Delta + m_g\) 自伴。
\end{enumerate}

\begin{enumerate}
\def\labelenumi{\arabic{enumi})}
\setcounter{enumi}{14}
\tightlist
\item
  设 \(n \in \mathbf{N}\) ，\(\Delta\) 为 \(\mathrm{L}^2(\mathbf{R}^n)\) 上唯一的拉普拉斯算子。用 B 表示 \(\mathbf{R}^n\) 中的单位球。
\end{enumerate}

若 \(g \in \mathcal{L}_{\mathrm{u}}(\mathbf{R}^{n})\) 使 \(\Delta\) 的定义域包含于 \(m_{g}\) 的定义域，则

\[
\sup _ {y \in \mathbf {R} ^ {n}} \int_ {\mathrm{B}} | g (x + y) | ^ {2} d x
\]

是有限的。（首先验证 \(m_{g}\) 是 \(\Delta\) -有界的。）

\begin{enumerate}
\def\labelenumi{\arabic{enumi})}
\setcounter{enumi}{15}
\tightlist
\item
  \begin{enumerate}
  \def\labelenumii{\alph{enumii})}
  \tightlist
  \item
    设 I = {]}0, 1{[}。空间 H \(^{1}\) (I) 包含于 C(Ī)。
  \end{enumerate}
\end{enumerate}

\begin{enumerate}
\def\labelenumi{\alph{enumi})}
\setcounter{enumi}{1}
\item
  微分算子 \(u\colon f\mapsto -if'\) （\(\mathrm{L}^2 (\mathrm{I})\) 上以 \(\mathcal{D}(\mathrm{I})\) 为定义域）是对称且可闭的。对所有 \(f\in \mathrm{dom}(u^{*})\) ，导出广义函数 \(f'\) （由 \(f\) 导出）属于 \(\mathrm{L}^2 (\mathrm{I})\) ，且我们有 \(u^{*}(f) = if'\) 。
\item
  \(\overline{u}\) 的亏空间都是 1 维的；由此给出 \(\overline{u}\) 的所有自伴扩张的一个明显描述。
\end{enumerate}

17) 对 I ={]}0, 1{[} 上的拉普拉斯算子进行分类。

¶ 18) 设 n 为整数 ≥ 1。我们用 D 表示 L²(T\^{}n) 上以 C∞(T\^{}n) 为定义域、满足

\[
\mathrm{D} (\varphi) = - \sum_ {i = 1} ^ {n} \partial_ {i} ^ {2} \varphi
\]

（对所有 \(\varphi\in\mathcal{C}^{\infty}(\mathbf{T}^{n})\) ）的部分算子，

\begin{enumerate}
\def\labelenumi{\alph{enumi})}
\item
  部分算子 D 本质自伴。我们用 \(\Delta\) 表示它的闭包。
\item
  部分算子 \(\Delta\) 自伴且正。
\item
  部分算子 \(\Delta\) 有紧预解式。
\end{enumerate}

\(d)\) 设 \(\lambda\in\mathbf{R}\) 。我们有 \(\lambda\in\mathrm{Sp}(\Delta)\) 当且仅当存在自然数 \(k_{1},\ldots,k_{n}\) 使 \(\lambda=4\pi^{2}(k_{1}^{2}+\cdots+k_{n}^{2})\) 。

\begin{enumerate}
\def\labelenumi{\alph{enumi})}
\setcounter{enumi}{4}
\tightlist
\item
  设 \(\lambda\) 为 \(\Delta\) 的一个特征值；用 \(N(\lambda)\) 表示相应特征空间的维数。我们有
\end{enumerate}

\[
\sum_ {\lambda \leqslant \mathrm{T}} \mathrm{N} (\lambda) = c _ {n} \mathrm{T} ^ {n / 2} + \mathrm{O} (\mathrm{T} ^ {(n - 1) / 2})
\]

当 T → +∞ 时成立。

\(f)\) 若 \(n \geqslant 2\) ，则 \(\limsup_{\lambda} N(\lambda) = +\infty\) 。（在 \(n = 2\) 的情形，考虑 \(\lambda = 4\pi^2 m\) ，其中 \(m\) 是形如 \(4k + 1\) 的素数的乘积。） \(g)\) 若 \(n = 1\) 或 \(n = 2\) ，我们有

\[
\operatorname{Card} (\operatorname{Sp} (u) \cap [ 0, \mathrm{T} ]) = o (\mathrm{T}),
\]

并且若 \(n\geqslant 3\) ，我们有

\[
\liminf _ {T \to + \infty} \frac {1}{T} \operatorname{Card} (\operatorname{Sp} (u) \cap [ 0, T ]) > 0.
\]

\begin{enumerate}
\def\labelenumi{\arabic{enumi})}
\setcounter{enumi}{18}
\tightlist
\item
  设 u 为希尔伯特空间 E 上的一个自伴部分算子。
\end{enumerate}

\begin{enumerate}
\def\labelenumi{\alph{enumi})}
\tightlist
\item
  设 \(t \in R\) 且 \(x \in E\) 。若 x 关于 u 的谱测度具有紧支撑，则我们有
\end{enumerate}

\[
e ^ {i u} (x) = \sum_ {n \in \mathbf {N}} \frac {u ^ {n} (x)}{n !},
\]

其中级数在 E 中收敛。

\begin{enumerate}
\def\labelenumi{\alph{enumi})}
\setcounter{enumi}{1}
\tightlist
\item
  由使得 x 关于 u 的谱测度具有紧支撑的那些向量 \(x \in E\) 组成之集是 E 的一个稠密子空间。
\end{enumerate}

\begin{enumerate}
\def\labelenumi{\arabic{enumi})}
\setcounter{enumi}{19}
\tightlist
\item
  设 \(u\) 为希尔伯特空间 E 上的一个部分算子。我们用 \(\mathrm{dom}(u)^{\mathrm{an}}\) 表示由满足下列条件的向量 \(x \in \mathrm{E}\) 组成的空间：\(x \in \mathrm{dom}(u^n)\) 对每个整数 \(n \in \mathbf{N}\) 成立，且幂级数
\end{enumerate}

\[
\sum_ {n \in \mathbf {N}} \frac {\| u ^ {n} (x) \|}{n !} z ^ {n}\tag{23}
\]

的收敛半径 \textgreater{} 0。\(\text{dom}(u)^{\text{an}}\) 的元素称为 u 的解析向量。

设 \(x\) 为 E 中的一个元素，\(x\) 属于 \(\mathrm{dom}(u^n)\) （对所有 \(n \in \mathbf{N}\) ）。设 \(\mathrm{E}_x\) 为由诸向量 \(u^n(x)\) （\(n \in \mathbf{N}\) ）生成的 E 的子空间的闭包，\(u_x\) 为 \(\overline{\mathrm{E}}_x\) 上以 \(\mathrm{E}_x\) 为定义域、由 \(u\) 通过过渡到子空间导出的部分算子。我们称 \(x\) 是 \(u\) 的一个唯一性向量，如果 \(u_x\) 本质自伴。

\begin{enumerate}
\def\labelenumi{\alph{enumi})}
\tightlist
\item
  假设 u 对称，且 E 含有 u 的唯一性向量的一个完全子集。则 u 本质自伴。（证明 \(u + i1_{E}\) 的像与 \(u - i1_{E}\) 的像都在 E 中稠密。）
\end{enumerate}

以下设 u 对称，且 E 含有解析向量的一个完全子集。

\begin{enumerate}
\def\labelenumi{\alph{enumi})}
\setcounter{enumi}{1}
\item
  对 u 的每个解析向量 x，部分算子 \(u_{x}\) （在 \(\overline{E}_{x}\) 上）有一个自伴扩张。
\item
  设 \(x\) 为 \(u\) 的一个解析向量，\(r > 0\) 为幂级数 (23) 的收敛半径。设 \(v\) 为 \(u_x\) 的一个自伴扩张。对每个 \(t \in \mathbf{R}\) （满足 \(|t| < r\) ），我们有
\end{enumerate}

\[
\langle e ^ {i t v} (x) \mid x \rangle = \sum_ {n \in \mathbf {N}} \frac {(i t) ^ {n}}{n !} \langle u ^ {n} (x) \mid x \rangle .
\]

（利用 \(x\) 关于 \(v\) 的谱测度，证明左端是定义在由 \(t \in \mathbf{C}\) （满足 \(|t| < r\) ）组成之集上的一个解析函数的限制，并计算该函数在 0 处的幂级数展开。）

\begin{enumerate}
\def\labelenumi{\alph{enumi})}
\setcounter{enumi}{3}
\item
  对每个 \(t \in R\) ，酉算子 \(e^{itv}\) 只依赖于 \(u_x\) ，且算子 \(u_x\) 本质自伴。（利用 V 第 428 页定理。） e) 由此推出 u 本质自伴。
\item
  E 上闭的对称部分算子 u 自伴，当且仅当 \(\text{dom}(u)^{\text{an}}\) 在 E 中稠密。
\end{enumerate}

\begin{enumerate}
\def\labelenumi{\arabic{enumi})}
\setcounter{enumi}{20}
\tightlist
\item
  我们用 \(x\) 表示 \(\mathbf{R}\) 的恒等函数。我们用 \(\mathcal{M}_{\mathrm{m}}(\mathbf{R})\) 表示由正测度 \(\nu\) （在 \(\mathbf{R}\) 上）组成的、满足 \(x^n \in \mathcal{L}^1(\mathbf{R}, \nu)\) （对每个整数 \(n \in \mathbf{N}\) ）的集。 \(a)\) 设 \(\nu \in \mathcal{M}_{\mathrm{m}}(\mathbf{R})\) ，并令 \(\mu_n = \nu(x^n)\) （对所有 \(n \in \mathbf{N}\) ）（“测度 \(\nu\) 的矩”）。对每个整数 \(n \geqslant 0\) 与所有复数 \(a_0, \ldots, a_n\) ，我们有
\end{enumerate}

\[
\sum_ {i = 0} ^ {n} \sum_ {j = 0} ^ {n} \overline {{a}} _ {i} a _ {j} \mu_ {i + j} \geqslant 0.\tag{24}
\]

在下文中，我们打算研究问题 \(a\) ) 的逆命题（“汉堡矩问题”）。因此，在本习题的其余部分，我们给定一个满足条件 (24) 的实数序列 \(\boldsymbol{\mu} = (\mu_n)_{n \geqslant 0}\) 。

\begin{enumerate}
\def\labelenumi{\alph{enumi})}
\setcounter{enumi}{1}
\tightlist
\item
  存在唯一的正半双线性形式 \(b_{\mu}\) （在 C{[}X{]} 上），使对所有整数 i 与 j 有 \(b_{\mu}(X^{i}, X^{j}) = \mu_{i+j}\) 。
\end{enumerate}

我们用 \(E_{\mu}\) 表示 C{[}X{]} 关于形式 \(b_{\mu}\) 的分离完备希尔伯特空间。 c) 线性映射 \(u_{\mu}\) （从 C{[}X{]} 到 C{[}X{]} ，由 \(u_{\mu}(P) = XP\) 定义）通过过渡到商，诱导出一个对称部分算子 \(u_{\mu}\) ，位于希尔伯特空间 \(E_{\mu}\) 上。

\begin{enumerate}
\def\labelenumi{\alph{enumi})}
\setcounter{enumi}{3}
\item
  部分算子 \(u_{\mu}\) 有一个自伴扩张。
\item
  存在测度 \(\nu\in\mathcal{M}_{\mathrm{m}}(\mathbf{R})\) 使 \(\mu_{n}=\nu(x^{n})\) （对所有 \(n\in N\) ）。 f) 半双线性形式 \(b_{\mu}\) 是分离的，当且仅当对每个测度 \(\nu\in\mathcal{M}_{\mathrm{m}}(\mathbf{R})\) （满足 \(\mu_{n}=\nu(x^{n})\) ，对所有 \(n\in N\) ），\(\nu\) 的支撑都是无穷的。（注意不等式 (24) 的左端等于 \(b_{\mu}(p,p)\) ，其中 \(p=a_{0}+\cdots+a_{n}X^{n}\) 。）
\end{enumerate}

\begin{enumerate}
\def\labelenumi{\arabic{enumi})}
\setcounter{enumi}{21}
\tightlist
\item
  我们沿用上一习题的记号。
\end{enumerate}

设 \(\boldsymbol{\mu} = (\mu_n)_{n\in \mathbf{N}}\) 为满足条件 (24) 的实数序列。我们考虑测度 \(\nu \in \mathcal{M}_{\mathrm{m}}(\mathbf{R})\) （使 \(\mu_{n} = \nu (x^{n})\) ，对所有 \(n\in \mathbf{N}\) ）的唯一性问题。我们假设半双线性形式 \(b_{\boldsymbol{\mu}}\) 是分离的（参见上一习题的问题 \(f\) ））。

\begin{enumerate}
\def\labelenumi{\alph{enumi})}
\item
  若存在具有紧支撑的测度 \(\nu \in \mathcal{M}_{\mathrm{m}}(\mathbf{R})\) 使 \(\nu(x^{n}) = \mu_{n}\) （对所有 \(n \in \mathbf{N}\) ），则 \(\nu\) 是 \(\mathcal{M}_{\mathrm{m}}(\mathbf{R})\) 中使 \(\nu(x^{n}) = \mu_{n}\) （对所有 \(n \in \mathbf{N}\) ）成立的唯一测度。
\item
  设 \(\nu \in \mathcal{M}_{\mathrm{m}}(\mathbf{R})\) 满足 \(\nu (x^n) = \mu_n\) （对所有 \(n\in \mathbf{N}\) ）。对每个 \(p\in \mathbf{C}[\mathrm{X}]\) ，我们有 \(p\in \mathcal{L}^2 (\mathbf{R},\nu)\) 且 \(\| p\| = b_{\mu}(p,p)\) 。
\item
  从 C{[}X{]} 到 \(\mathcal{L}^{2}(\mathbf{R},\nu)\) 的典范单射通过连续性与过渡到商，定义了从 \(E_{\mu}\) 到 \(\mathrm{L}^{2}(\mathbf{R},\nu)\) 的一个等距。在下文中，我们把 \(E_{\mu}\) 与它的像等同起来，该像是 \(\mathrm{L}^{2}(\mathbf{R},\nu)\) 的一个闭子空间。
\end{enumerate}

\(d)\) 我们有 \(u_{\mu} \subset m_x\) ，其中 \(m_x\) 表示乘以 \(x\) 的算子（在 \(\mathrm{L}^2(\mathbf{R}, \nu)\) 中）。

\begin{enumerate}
\def\labelenumi{\alph{enumi})}
\setcounter{enumi}{4}
\item
  测度 \(\nu\) 是常函数 1 关于 \(m_x\) 的谱测度。
\item
  设 \(\nu \in \mathcal{M}_{\mathrm{m}}(\mathbf{R})\) 。我们有 \(\nu(x^n) = \mu_n\) （对所有 \(n \in \mathbf{N}\) ），当且仅当存在希尔伯特空间 F 与 F 上满足下列条件的自伴部分算子 \(v\) ：
\end{enumerate}

\begin{enumerate}
\def\labelenumi{(\roman{enumi})}
\item
  空间 \(E_{\mu}\) 是 F 的闭子空间，且从 \(E_{\mu}\) 到 F 的典范单射是等距的；
\item
  我们有 \(u_{\mu} \subset v\) ；
\item
  测度 \(\nu\) 是常函数 \(1 \in E_{\mu}\) 关于 v 的谱测度。
\end{enumerate}

\begin{enumerate}
\def\labelenumi{\alph{enumi})}
\setcounter{enumi}{6}
\item
  存在唯一的测度 \(\nu \in \mathcal{M}_{\mathrm{m}}(\mathbf{R})\) 使 \(\nu(x^{m}) = \mu_{n}\) （对所有 \(n \in \mathbf{N}\) ），当且仅当部分算子 \(u_{\boldsymbol{\mu}}\) 本质自伴。
\item
  设 \(\nu_{lg}\) 为 R 上的正测度，其支撑含于 \(R_{+}\) ，且它在 \(R_{+}^{*}\) 上的限制相对于勒贝格测度具有密度
\end{enumerate}

\[
t \mapsto \frac {1}{t} \exp (- \log (t) ^ {2} / 2)
\]

设 \(\widetilde{\nu}_{lg}\) 为 R 上的正测度，其支撑含于 \(R_{+}\) ，它在 \(R_{+}^{*}\) 上的限制以密度 \(t \mapsto 1 + \sin(2\pi \log(t))\) 相对于 \(\nu_{lg}\) 定义。

我们有 \(\nu_{\mathrm{lg}} \in \mathcal{M}_{\mathrm{m}}(\mathbf{R})\) 且 \(\widetilde{\nu}_{\mathrm{lg}} \in \mathcal{M}_{\mathrm{m}}(\mathbf{R})\) 。

\begin{enumerate}
\def\labelenumi{\roman{enumi})}
\tightlist
\item
  我们有 \(\nu_{\mathrm{lg}}(x^n) = \widetilde{\nu}_{\mathrm{lg}}(x^n)\) （对所有 \(n\in \mathbf{N}\) ）。
\end{enumerate}

\begin{enumerate}
\def\labelenumi{\arabic{enumi})}
\setcounter{enumi}{22}
\tightlist
\item
  我们沿用上一习题的记号；特别地，我们考虑实数序列 \(\boldsymbol{\mu} = (\mu_n)_{n\in \mathbf{N}}\) ，它满足条件 (24) 且使半双线性形式 \(b_{\mu}\) 分离。我们把 \(\mathbf{C}[\mathrm{X}]\) 与 \(\mathrm{E}_{\mu}\) 的一个稠密子空间等同起来。
\end{enumerate}

\begin{enumerate}
\def\labelenumi{\alph{enumi})}
\tightlist
\item
  存在唯一的多项式序列 \((p_{n})_{n\in\mathbf{N}}\) （属于 C{[}X{]}），使下列条件得到满足（“与 \(\mu\) 相关联的第一类正交多项式”）：
\end{enumerate}

\begin{enumerate}
\def\labelenumi{(\roman{enumi})}
\item
  族 \((p_n)\) 是 \(\mathrm{E}_{\mu}\) 的标准正交基；
\item
  多项式 \(p_n\) 的次数为 \(n\) （对所有 \(n \in \mathbf{N}\) ）；
\item
  \(\mathrm{X}^n\) 在 \(p_n\) 中的系数属于 \(\mathbf{R}_+^*\) 。
\end{enumerate}

此外，我们有 \(p_n \in \mathbf{R}[\mathbf{X}]\) （对所有 \(n \in \mathbf{N}\) ），且 \(p_0 = 1 / \sqrt{\mu_0}\) 。我们还记 \(p_{-1} = 0\) 。

\begin{enumerate}
\def\labelenumi{\alph{enumi})}
\setcounter{enumi}{1}
\tightlist
\item
  存在序列 \((a_{n})_{n\geqslant -1}\) （在 \(\mathbf{R}_{+}^{*}\) 中）与序列 \((b_{n})_{n\in \mathbf{N}}\) （在 \(\mathbf{R}\) 中），使 \(a_{-1} = 1\) 且
\end{enumerate}

\[
u _ {\pmb {\mu}} (p _ {n}) = a _ {n} p _ {n + 1} + b _ {n} p _ {n} + a _ {n - 1} p _ {n - 1}
\]

对所有 \(n \in \mathbf{N}\) 成立。这些序列是唯一的。

\begin{enumerate}
\def\labelenumi{\alph{enumi})}
\setcounter{enumi}{2}
\item
  对每个 \(n \geqslant 1\) ，\(X^{n}\) 在正交多项式 \(p_{n}\) 中的系数等于 \((a_{0} \cdots a_{n-1})^{-1}\) 。
\item
  设 \(\mathcal{F}\) 为由复数序列 \((z_{n})_{n\geqslant -1}\) （满足 \(z_{-1} = 0\) ）组成的复向量空间。我们用 \(\widetilde{v}_{\mu}\) 表示从 \(\mathcal{F}\) 到其自身的线性映射，它满足 \(\widetilde{v}_{\mu}(z_n) = (w_n)_{n\in \mathbf{N}}\) ，其中 \(w_{-1} = 0\) 且
\end{enumerate}

\[
w _ {n} = a _ {n} z _ {n + 1} + b _ {n} z _ {n} + a _ {n - 1} z _ {n - 1}
\]

对所有 \(n \in N\) 成立。

对每个 \(s \in \mathbf{C}\) ，序列 \(z_s = (p_n(s))_{n \geqslant -1}\) 是 \(\mathcal{F}\) 中唯一的满足 \((z_s)_0 = p_0(z)\) 且 \(\widetilde{v}_{\boldsymbol{\mu}}(z_s) = sz_s\) 的序列。

\begin{enumerate}
\def\labelenumi{\alph{enumi})}
\setcounter{enumi}{4}
\tightlist
\item
  设 \(z\) 与 \(w\) 为 \(\mathcal{F}\) 的元素。我们通过令下式成立来定义序列 \(\mathrm{W}(z,w) = (\mathrm{W}_n(z,w))_{n\in \mathbf{N}}\) ：
\end{enumerate}

\[
\mathrm{W} _ {n} (z, w) = a _ {n} (z _ {n + 1} w _ {n} - z _ {n} w _ {n + 1})
\]

对 \(n\in \mathbf{N}\) 成立。我们有

\[
\mathrm{W} _ {n} (z, w) = \sum_ {j = 0} ^ {n} (\widetilde {v} _ {\boldsymbol {\mu}} (z) _ {j} w _ {j} - z _ {j} \widetilde {v} _ {\boldsymbol {\mu}} (w) _ {j})
\]

对所有 \(n \in N\) 成立。

\(f)\) 设 \(\mathrm{P} \subset \ell^2(\mathbf{N})\) 为具有有限支撑的序列组成的空间。对每个 \(z \in \mathrm{P}\) ，序列 \(\widetilde{v}_{\boldsymbol{\mu}}(z)\) 属于 \(\ell^2(\mathbf{N})\) 。

我们用 \(v_{\mu}\) 表示 \(\ell^2(\mathbf{N})\) 上以 P 为定义域、由 \(z \mapsto \widetilde{v}_{\mu}(z)\) 给出的部分算子。它是对称的。

\begin{enumerate}
\def\labelenumi{\alph{enumi})}
\setcounter{enumi}{6}
\tightlist
\item
  \(v_{\mu}\) 的伴随是部分算子，其定义域为
\end{enumerate}

\[
\mathrm{dom} (v _ {\boldsymbol {\mu}} ^ {*}) = \{z \in \ell^ {2} (\mathbf {N}) | \widetilde {v} _ {\boldsymbol {\mu}} (z) \in \ell^ {2} (\mathbf {N}) \}
\]

且满足 \(v_{\mu}^{*}(z) = \widetilde{v}_{\mu}(z)\) （对 \(z \in \mathrm{dom}(v_{\mu}^{*})\) ）。

\(h)\) 对所有 \(z\) 与 \(w\) （属于 \(\mathrm{dom}(v_{\mu}^{*})\) ），\(\mathrm{W}_{n}(\overline{z}, w)\) 当 \(n \to +\infty\) 时的极限存在，且等于 \(\langle w | v_{\mu}^{*}(z) \rangle - \langle v_{\mu}^{*}(w) | z \rangle\) 。

\begin{enumerate}
\def\labelenumi{\roman{enumi})}
\tightlist
\item
  设 \(s \in C\) 。\(v_{\boldsymbol{\mu}}^{*} - s \cdot 1_{\ell^{2}(\mathbf{N})}\) 的核维数 \(\leqslant 1\) ；其维数为 1 当且仅当 \(z_{s} \in \ell^{2}(\mathbf{N})\) ，且此时它由 \(z_{s}\) 生成。
\end{enumerate}

\begin{enumerate}
\def\labelenumi{\alph{enumi})}
\setcounter{enumi}{9}
\tightlist
\item
  部分算子 \(v_{\mu}\) （相应地 \(u_{\mu}\) ）的闭包的亏指标为 (0, 0) 或 (1, 1)；第一种情形成立当且仅当存在 \(s \in \mathbf{C} - \mathbf{R}\) 使 \(z_s \notin \ell^2(\mathbf{N})\) 。（对 \(u_{\mu}\) 的情形，注意映射 \(X^n \mapsto e_n\) （其中 \((e_n)\) 是 \(\ell^2(\mathbf{N})\) 的典范基）诱导出 \(E_{\mu}\) 到 \(\ell^2(\mathbf{N})\) 的一个等距同构，使我们可以把 \(u_{\mu}\) 与 \(v_{\mu}\) 等同起来。）
\end{enumerate}

\(k)\) 设 \(s \in \mathbf{C}\) ，\(f \in \mathrm{Ker}(u_{\mu}^{*} - s \cdot 1_{\mathrm{E}_{\mu}})\) 。若 \(\langle f | 1 \rangle = 0\) ，则 \(f = 0\) 。

\begin{enumerate}
\def\labelenumi{\alph{enumi})}
\setcounter{enumi}{11}
\item
  设 \(u_1\) 与 \(u_2\) 为 \(u_\mu\) 在 \(\mathrm{E}_\mu\) 上的两个自伴扩张。我们有 \(u_1 = u_2\) ，当且仅当存在 \(s \in \mathbf{C} - \mathbf{R}\) 使 \(\langle \mathrm{R}(u_1, s)1|1\rangle = \langle \mathrm{R}(u_2, s)1|1\rangle\) 。（假设这样一个等式成立；令 \(f_1 = \mathrm{R}(u_1, s)1\) ，并证明 \(f_1\) 不属于 \(\overline{u}_\mu\) 的定义域；由此推出 \(u_1\) 与 \(u_1\) 到 \(\mathrm{dom}(\overline{u}_\mu) \oplus \mathbf{C}f_1\) 的限制的闭包重合。对 \(u_2\) 与 \(f_2 = \mathrm{R}(u_2, s)1\) 作同样的推理，然后证明 \(f_1 = f_2\) ，从而得出 \(u_1 = u_2\) 。） m) 假设 \(u_\mu\) 不是本质自伴的。则至少存在两个测度 \(\nu \in \mathcal{M}_{\mathrm{m}}(\mathbf{R})\) （满足 \(\nu(x^n) = \mu_n\) ，对所有 \(n \in \mathbf{N}\) ）。 n) 以下设 \(u_\mu\) 本质自伴。对每个 \(s \in \mathbf{C} - \mathbf{R}\) ，存在多项式序列 \((r_{s,m})_{m \in \mathbf{N}}\) ，使 \((x - s)r_{s,m}\) 在 \(\mathrm{E}_\mu\) 中当 \(m \to +\infty\) 时收敛于常函数 1。
\item
  设 \(\nu\in\mathcal{M}_{\mathrm{m}}(\mathbf{R})\) 为满足 \(\nu(x^{n})=\mu_{n}\) （对所有 \(n\in\mathbf{N}\) ）的测度。对每个 \(m\in\mathbf{N}\) ，不等式
\end{enumerate}

\[
\left| \int_ {\mathbf {R}} \frac {d \nu (t)}{t - s} - \int_ {\mathbf {R}} r _ {s, m} (t) d \nu (t) \right| \leqslant \frac {\mu_ {0} ^ {2}}{| \mathcal {I} (s) ^ {2} |} \| 1 - (x - s) r _ {s, m} \| _ {\mathrm{E} _ {\mu}}
\]

成立。

\begin{enumerate}
\def\labelenumi{\alph{enumi})}
\setcounter{enumi}{15}
\tightlist
\item
  测度 \(\nu\) 是唯一的。
\end{enumerate}

\begin{enumerate}
\def\labelenumi{\arabic{enumi})}
\setcounter{enumi}{23}
\tightlist
\item
  我们沿用前面诸习题的记号；特别地，我们考虑实数序列 \(\boldsymbol{\mu} = (\mu_n)_{n\in \mathbf{N}}\) ，它满足条件 (24) 且使半双线性形式 \(b_{\boldsymbol{\mu}}\) 分离。
\end{enumerate}

\begin{enumerate}
\def\labelenumi{\alph{enumi})}
\tightlist
\item
  存在唯一的具有实系数的多项式序列 \((q_{n})_{n\in\mathbf{N}}\) ，使对每个 \(s\in C\) ，公式
\end{enumerate}

\[
q _ {n} (s) = b _ {\pmb {\mu}} \Big (1, \frac {p _ {n} - p _ {n} (s)}{\mathrm{X} - s} \Big)
\]

成立（“第二类正交多项式”）。

\begin{enumerate}
\def\labelenumi{\alph{enumi})}
\setcounter{enumi}{1}
\tightlist
\item
  设 \(\nu \in \mathcal{M}_{\mathrm{m}}(\mathbf{R})\) 为 \(\mathbf{R}\) 上满足 \(\nu(x^n) = \mu_n\) （对所有 \(n \in \mathbf{N}\) ）的测度。对每个 \(s \in \mathbf{C} - \mathbf{R}\) ，我们有
\end{enumerate}

\[
\int_ {\mathbf {R}} \frac {p _ {n} (t)}{t - s} d \nu (t) = q _ {n} (s) + p _ {n} (s) \int_ {\mathbf {R}} \frac {d \nu (t)}{t - s}.
\]

\begin{enumerate}
\def\labelenumi{\alph{enumi})}
\setcounter{enumi}{2}
\item
  设 \(s \in C-R\) 。序列 \((p_{n}(s))_{n \in \mathbf{N}}\) 属于 \(\ell^{2}(\mathbf{N})\) ，当且仅当序列 \((q_{n}(s))_{n \in \mathbf{N}}\) 属于 \(\ell^{2}(\mathbf{N})\) 。（令 \(\sigma(s) = \int(t-s)^{-1} d\nu(t)\) ；利用上一问题，证明序列 \((q_{n}(s) + \sigma(s)p_{n}(s))_{n \in \mathbf{N}}\) 属于 \(\ell^{2}(\mathbf{N})\) 。）
\item
  对每个 \(n \in \mathbf{N}\) ，公式
\end{enumerate}

\[
a _ {n} (p _ {n + 1} (\mathrm{X}) p _ {n} (\mathrm{Y}) - p _ {n} (\mathrm{X}) p _ {n + 1} (\mathrm{Y})) = (\mathrm{X} - \mathrm{Y}) \sum_ {j = 0} ^ {n} p _ {j} (\mathrm{X}) p _ {j} (\mathrm{Y})
\]

在 \(\mathbf{C}[\mathrm{X},\mathrm{Y}]\) 中成立。

\begin{enumerate}
\def\labelenumi{\alph{enumi})}
\setcounter{enumi}{4}
\tightlist
\item
  对每个 \(n \in \mathbf{N}\) ，我们有关系式
\end{enumerate}

\[
p _ {n} q _ {n + 1} - p _ {n + 1} q _ {n} = \frac {1}{a _ {n}}.
\]

\begin{enumerate}
\def\labelenumi{\alph{enumi})}
\setcounter{enumi}{5}
\tightlist
\item
  若通项为 \(a_{n}^{-1}\) 的级数发散，则测度 \(\nu\) 是满足 \(\nu(x^{n}) = \mu_{n}\) （对所有 \(n \in N\) ）的唯一测度。
\end{enumerate}

\begin{enumerate}
\def\labelenumi{\arabic{enumi})}
\setcounter{enumi}{24}
\tightlist
\item
  我们沿用前面诸习题的记号；特别地，我们考虑实数序列 \(\boldsymbol{\mu} = (\mu_{n})_{n \in \mathbf{N}}\) ，它满足条件 (24) 且使半双线性形式 \(b_{\mu}\) 分离。
\end{enumerate}

假设 \(\mu_{0}=1\) ，且卡尔曼条件成立，即通项为 \(\mu_{2n}^{-1/(2n)}\) （\(n\geqslant1\) ）的级数发散。 a) 我们有

\[
\mu_ {2 n} ^ {- 1 / (2 n)} \leqslant (a _ {0} \dots a _ {n - 1}) ^ {- 1 / n} \leqslant \frac {e}{n ^ {2}} \sum_ {j = 1} ^ {n} \frac {j}{a _ {j - 1}}
\]

对每个整数 \(n \geqslant 1\) 成立。（为建立第二个不等式，可以证明并利用不等式 \((n/e)^{n} \leqslant n!\) （\(n \geqslant 1\) ）。）

\begin{enumerate}
\def\labelenumi{\alph{enumi})}
\setcounter{enumi}{1}
\item
  通项为 \(a_{n}^{-1}\) （\(n \geqslant 1\) ）的级数发散。
\item
  存在唯一的测度 \(\nu \in \mathcal{M}_{\mathrm{m}}(\mathbf{R})\) ，满足 \(\nu(x^n) = \mu_n\) （对所有 \(n \in \mathbf{N}\) ）。
\end{enumerate}

\begin{enumerate}
\def\labelenumi{\arabic{enumi})}
\setcounter{enumi}{25}
\tightlist
\item
  我们用 \(x\) 表示 \(\mathbf{R}_{+}\) 的恒等函数。
\end{enumerate}

\begin{enumerate}
\def\labelenumi{\alph{enumi})}
\tightlist
\item
  设 \(\nu\) 为 \(R_{+}\) 上的正测度，使 \(x^{n}\) 是 \(\nu\) -可积的（对每个整数 \(n \geqslant 0\) ）。记 \(\mu_{n} = \nu(x^{n})\) 。对每个整数 \(n \geqslant 0\) 与所有复数 \(a_{0}, \ldots, a_{n}\) ，我们有
\end{enumerate}

\[
\sum_ {i = 0} ^ {n} \sum_ {j = 0} ^ {n} \overline {{a}} _ {i} a _ {j} \mu_ {i + j} \geqslant 0
\]

\[
\sum_ {i = 0} ^ {n} \sum_ {j = 0} ^ {n} \overline {{a}} _ {i} a _ {j} \mu_ {i + j + 1} \geqslant 0.
\]

\begin{enumerate}
\def\labelenumi{\alph{enumi})}
\setcounter{enumi}{1}
\tightlist
\item
  若 \((\mu_{n})_{n\geqslant0}\) 是满足上一问题条件的实数序列，则存在正测度 \(\nu\) （在 \(R_{+}\) 上），使 \(\mu_{n}=\nu(x^{n})\) （对所有 \(n\in N\) ）。（“斯蒂尔切斯矩问题”；利用习题 21 的方法，并借助弗里德里希斯扩张。）
\end{enumerate}

\begin{enumerate}
\def\labelenumi{\arabic{enumi})}
\setcounter{enumi}{26}
\tightlist
\item
  我们记 \(M = X^{2}Y^{2}(X^{2} + Y^{2} - 3) \in \mathbf{R}[X, Y]\) （“莫茨金多项式”）。
\end{enumerate}

\begin{enumerate}
\def\labelenumi{\alph{enumi})}
\item
  我们有 M(x,y) ≥ 0 （对所有 x, y ∈ R）。
\item
  不存在有限族 \((p_{i})_{i\in I}\) （在 R{[}X,Y{]} 中）使 \(M=\sum p_{i}^{2}\) 。（若这样一个族存在，证明每个 \(p_{i}\) 必形如 \(a+bXY\) ，其中 \(a\in R\) 且 \(b\in R[X,Y]\) 的次数至多为 1；然后计算 \((\mathrm{XY})^{2}\) 的系数。）
\item
  设 \(\tau\colon\mathbf{N}^{2}\to\mathbf{N}-\{0\}\) 为由下式定义的映射：
\end{enumerate}

\[
\tau (0, 0) = 1, \quad \tau (1, 2) = 2, \quad \tau (2, 1) = 3, \quad \tau (1, 1) = 4, \quad \tau (1, 0) = 5,
\]

\[
\tau (0, 1) = 6, \quad \tau (2, 0) = 7, \quad \tau (0, 2) = 8, \quad \tau (3, 0) = 9, \quad \tau (0, 3) = 1 0
\]

以及

\[
\tau (i, j) = j + 1 + (i + j) (i + j + 1) / 2 \text {   if   } i + j \geqslant 4.
\]

映射 \(\tau\) 是一个双射。

\begin{enumerate}
\def\labelenumi{\alph{enumi})}
\setcounter{enumi}{3}
\tightlist
\item
  设 \((\nu_{n})_{n\geqslant 1}\) 为由下式定义的序列：
\end{enumerate}

\[
\nu_ {i} = 1 \text {if} i \in \{1, 2, 3 \}, \quad \nu_ {4} = 4, \quad \nu_ {n} = (n!) ^ {(n + 1)!} \text {if} n \geqslant 5.
\]

对 \(n\) 与 \(m\) （属于 \(\mathbf{N}\) ），我们令 \(\mu_{n,m} = 0\) （若 \(n\) 或 \(m\) 为奇数），否则令 \(\mu_{n,m} = \nu_{\tau(n/2,m/2)}\) 。线性映射 \(\alpha\) （从 \(\mathbf{C}[\mathrm{X},\mathrm{Y}]\) 到 \(\mathbf{C}\) ，满足 \(\alpha(\mathrm{X}^n\mathrm{Y}^m) = \mu_{n,m}\) ）满足 \(\alpha(p^2) \geqslant 0\) （对所有 \(p \in \mathbf{R}[\mathrm{X},\mathrm{Y}]\) ）。（只需验证：对每个整数 \(k \in \mathbf{N}\) ，矩阵 \((a_{m,m})_{1 \leqslant n,m \leqslant k}\) 的行列式为正，其中 \(a_{n,m}\) 由 \(a_{\tau(i,j),\tau(i',j')} = \alpha(\mathrm{X}^{i+i'}\mathrm{Y}^{j+j'})\) （对所有 \((i,j,i',j')\) ）定义；然后对 \(k\) 归纳。）

\begin{enumerate}
\def\labelenumi{\alph{enumi})}
\setcounter{enumi}{4}
\tightlist
\item
  不存在正测度 \(\nu\) （在 \(R^{2}\) 上），使得 \((x,y)\mapsto x^{n}y^{m}\) 关于 \(\mu\) 可积（对所有 \((n,m)\in\mathbf{N}^{2}\) ），且使得
\end{enumerate}

\[
\int_ {\mathbf {R} ^ {2}} x ^ {n} y ^ {m} d \nu (x, y) = \mu_ {n, m}
\]

对所有 \((n,m)\in\mathbf{N}^{2}\) 成立。

\begin{enumerate}
\def\labelenumi{\alph{enumi})}
\setcounter{enumi}{5}
\tightlist
\item
  存在希尔伯特空间 E 与 E 上一个稠定部分算子 u，满足下列条件：
\end{enumerate}

\begin{enumerate}
\def\labelenumi{(\roman{enumi})}
\item
  我们有 \(\operatorname{dom}(u) \subset \operatorname{dom}(u^*)\) ；
\item
  我们有 \(\|u(x)\|=\|u^{*}(x)\|\) （对所有 \(x\in\mathrm{dom}(u)\) ）；
\item
  不存在包含 E 为子空间的希尔伯特空间 F（使 E 到 F 的典范单射是等距的）以及 F 上的正规部分算子 v 使 \(u \subset v\) 。
\end{enumerate}

把这一结果与 IV 第 352 页习题 24 相比较。

\begin{enumerate}
\def\labelenumi{\arabic{enumi})}
\setcounter{enumi}{27}
\tightlist
\item
  对每个整数 \(k \geqslant 0\) ，我们用 \(b_k\) 表示集合 \(\{1, \ldots, k\}\) 上等价关系的个数。
\end{enumerate}

对每个 \(\lambda > 0\) ，我们用 \(\varpi_{\lambda}\) 表示参数为 \(\lambda\) 的泊松测度（在 \(\mathbf{R}\) 上）（参见 II 第 282 页习题 41），即总质量为 1 、支撑在 \(\mathbf{N}\) 上且满足 \(\varpi_{\lambda}(j) = e^{-\lambda} \lambda^{j} / j!\) （对所有 \(j \in \mathbf{N}\) ）的测度。

\begin{enumerate}
\def\labelenumi{\alph{enumi})}
\tightlist
\item
  关系式
\end{enumerate}

\[
b _ {k + 1} = \sum_ {j = 0} ^ {k} {\binom {k} {j}} b _ {j}
\]

对每个整数 \(k \in N\) 成立。

\begin{enumerate}
\def\labelenumi{\alph{enumi})}
\setcounter{enumi}{1}
\tightlist
\item
  对任意函数 \(f: N \to C\) （使函数 \(x \mapsto xf(x)\) 是 \(\varpi\) -可积的），我们有
\end{enumerate}

\[
\int_ {\mathbf {R}} x f (x) d \varpi_ {\lambda} (x) = \lambda \int_ {\mathbf {R}} f (x + 1) d \varpi_ {\lambda} (x)
\]

（“陈氏公式”）。

此外，若 \(\mu\) 是 R 上满足这一性质的有界测度，则存在 \(c \in R\) 使 \(\mu = c\varpi_{\lambda}\) 。

\begin{enumerate}
\def\labelenumi{\alph{enumi})}
\setcounter{enumi}{2}
\tightlist
\item
  对每个 \(k \in \mathbf{N}\) ，函数 \(x \mapsto x^k\) 是 \(\varpi_1\) -可积的，且我们有
\end{enumerate}

\[
\varpi_ {1} (x \mapsto x ^ {k}) = b _ {k}.
\]

此外，测度 \(\varpi_{1}\) 是 R 上满足这些等式的唯一正测度。

\begin{enumerate}
\def\labelenumi{\alph{enumi})}
\setcounter{enumi}{3}
\tightlist
\item
  对每个整数 \(k \geqslant 0\) ，我们有
\end{enumerate}

\[
b _ {k} = \frac {1}{e} \sum_ {j \geqslant 0} \frac {j ^ {n}}{j !}
\]

（“多宾斯基公式”）以及

\[
b _ {k} = \frac {2 k !}{e \pi} \mathcal {I} \left(\int_ {0} ^ {\pi} e ^ {e ^ {e ^ {i t}}} \sin (k t) d t\right)
\]

（“切萨罗公式”）。

\begin{enumerate}
\def\labelenumi{\alph{enumi})}
\setcounter{enumi}{4}
\tightlist
\item
  对每个整数 \(n \geqslant 1\) ，我们用 \(\mu_{n}\) 表示 R 上的正测度，它是对称群 \(S_{n}\) 的规范化哈尔测度在映射 \(f_{n}\) （从 \(S_{n}\) 到 R）下的像，该映射把 \(\sigma \in S_{n}\) 映为 \(\sigma\) 的不动点个数。设 \(k \in N\) 。对每个整数 \(n \geqslant k\) ，我们有
\end{enumerate}

\[
\mu_ {n} (x \mapsto x ^ {k}) = b _ {k}.
\]

\begin{enumerate}
\def\labelenumi{\alph{enumi})}
\setcounter{enumi}{5}
\tightlist
\item
  序列 \((\mu_{n})\) 当 \(n \to +\infty\) 时弱收敛于泊松测度 \(\varpi_{1}\) 。
\end{enumerate}

\begin{enumerate}
\def\labelenumi{\arabic{enumi})}
\setcounter{enumi}{28}
\tightlist
\item
  设 E 为复希尔伯特空间。我们用 \(\mathcal{T}_b\) （相应地 \(\mathcal{T}_s\) 、 \(\mathcal{T}_f\) ）表示 \(\mathcal{L}(\mathrm{E})\) 的巴拿赫空间拓扑（相应地逐点收敛拓扑、由半范数 \(u \mapsto \langle x \mid u(y) \rangle\) （对 \((x, y) \in \mathrm{E} \times \mathrm{E}\) ）定义的局部凸拓扑）。
\end{enumerate}

\begin{enumerate}
\def\labelenumi{\alph{enumi})}
\item
  设 \(\lambda_0\in \mathbf{C} - \mathbf{R}\) 。 \(\mathcal{T}_b\) -预解拓扑（相应地 \(\mathcal{T}_s\) -预解拓扑）是 \(\mathcal{A}(\mathrm{E})\) 上使映射 \(u\mapsto \mathrm{R}(u,\lambda_0)\) （从 \(\mathcal{A}(\mathrm{E})\) 到 \(\mathcal{L}(\mathrm{E})\) （赋有拓扑 \(\mathcal{T}_b\) （相应地 \(\mathcal{T}_s\) ）））连续的最粗拓扑。
\item
  \(T_{f}\) -预解拓扑与 \(T_{s}\) -预解拓扑重合。
\item
  设 \(X \subset \mathcal{L}(E)\) 为 E 的埃尔米特自同态组成的有界集。由 \(T_{b}\) -预解拓扑（相应地 \(T_{s}\) -预解拓扑）在 X 上诱导的拓扑，与由 \(T_{b}\) （相应地 \(T_{s}\) ）诱导的拓扑重合。
\item
  a) 中的性质对 \(T_{f}\) -预解拓扑不成立。
\item
  给出一个序列 \((u_{n})\) （由部分自伴算子组成，且对 \(T_{b}\) -预解拓扑收敛）以及函数 \(f \in \mathcal{C}_{b}(\mathbf{R})\) 的例子，使 \(f(u_{n})\) 不收敛于 \(f(u)\) （在 \(\mathcal{L}(\mathrm{E})\) 中）。
\end{enumerate}

\(f\) 设 \((u_n)\) 为 E 上部分自伴算子的序列，\(u\) 为 E 上的一个自伴部分算子。若存在 \(\mathrm{E}_u\) 的稠密子空间 F（包含于 \(\mathrm{dom}(u_n)\) ，对所有 \(n\) ），使 \(u_n(x)\) 收敛于 \(u(x)\) （对所有 \(x \in \mathrm{F}\) ），则 \((u_n)\) 收敛于 \(u\) （对 \(\mathcal{T}_s\) -预解拓扑）。

\begin{enumerate}
\def\labelenumi{\arabic{enumi})}
\setcounter{enumi}{29}
\tightlist
\item
  设 E 为复希尔伯特空间。用 \(T_{b}\) 表示 \(\mathcal{L}(\mathrm{E})\) 上有界收敛拓扑。
\end{enumerate}

\begin{enumerate}
\def\labelenumi{\alph{enumi})}
\tightlist
\item
  设 \(u\) 为 E 上的自伴部分算子。对每个 \(\lambda \in \mathbf{C} - \mathbf{R}\) （使 \(\lambda\) 的虚部为负），我们有
\end{enumerate}

\[
\mathrm{R} (u, \lambda) = i \int_ {0} ^ {\infty} e ^ {- i t \lambda} e ^ {i t u} d t.
\]

\begin{enumerate}
\def\labelenumi{\alph{enumi})}
\setcounter{enumi}{1}
\tightlist
\item
  \(T_{b}\) -预解拓扑是 \(\mathcal{A}(E)\) 上使映射 \(u \mapsto e^{itu}\) （从 \(\mathcal{A}(E)\) 到 \(\mathcal{L}(E)\) ）对每个 \(t \in R\) 都连续的最粗拓扑。
\end{enumerate}

\begin{enumerate}
\def\labelenumi{\arabic{enumi})}
\setcounter{enumi}{30}
\tightlist
\item
  设 E 为复希尔伯特空间。用 \(T_{s}\) 表示 \(\mathcal{L}(\mathrm{E})\) 的逐点收敛拓扑。用 H 表示虚部严格为正的复数组成的集。
\end{enumerate}

设 \((u_{n})\) 为 \(\mathcal{A}(\mathrm{E})\) 中的序列。我们假设

\begin{enumerate}
\def\labelenumi{(\roman{enumi})}
\item
  存在 \(\lambda_{+} \in \mathbf{H}\) 使 \(\mathrm{R}(u_n, \lambda_{+})\) 在 \(\mathcal{L}(\mathrm{E})\) （赋有拓扑 \(\mathcal{T}_s\) ）中收敛；我们用 \(v_{+}\) 表示其极限；
\item
  存在 \(\lambda_{-} \in \overline{\mathbf{H}}\) 使 \(\mathrm{R}(u_n, \lambda_{-})\) 在 \(\mathcal{L}(\mathrm{E})\) （赋有拓扑 \(\mathcal{T}_s\) ）中收敛；
\item
  自同态 \(v_{+}\) 的像 F 在 E 中稠密。
\end{enumerate}

设 D 为 C 中由满足下式的 \(\lambda \in H\) 组成的开圆盘

\[
| \lambda - \lambda_ {+} | <   \frac {1}{| \mathcal {I} (\lambda_ {+}) |}.
\]

\begin{enumerate}
\def\labelenumi{\alph{enumi})}
\tightlist
\item
  对每个 \(\lambda \in D\) ，级数
\end{enumerate}

\[
r _ {\lambda} = \sum_ {n = 0} ^ {\infty} (\lambda_ {+} - \lambda) v _ {+} ^ {n + 1}
\]

在 \(\mathcal{L}(\mathrm{E})\) 中收敛。此外，对每个 \(\lambda \in \mathrm{D}\) ，序列 \((\mathrm{R}(u_n,\lambda))_n\) 收敛于 \(r_{\lambda}\) （在 \(\mathcal{L}(\mathrm{E})\) 中，赋有拓扑 \(\mathcal{T}_s\) ）。

\begin{enumerate}
\def\labelenumi{\alph{enumi})}
\setcounter{enumi}{1}
\item
  存在唯一的从 H 到 \(\mathcal{L}(\mathrm{E})\) 的全纯映射，它在 D 上与映射 \(\lambda\mapsto r_{\lambda}\) 重合。我们仍把它记作 \(\lambda\mapsto r_{\lambda}\) 。
\item
  存在唯一的从 \(\mathbf{C} - \mathbf{R}\) 到 \(\mathcal{L}(\mathrm{E})\) 的全纯映射，它在 \(\mathbf{H}\) 上与映射 \(\lambda \mapsto r_{\lambda}\) 重合，且满足 \(r_{\lambda}^{*} = r_{\overline{\lambda}}\) （对所有 \(\lambda \in \mathbf{C} - \mathbf{R}\) ）。它满足 \(r_{\lambda_{-}} = v_{-}\) ，且对所有 \(\lambda \in \mathbf{C} - \mathbf{R}\) ，序列 \((\mathrm{R}(u_n,\lambda))\) 收敛于 \(r_{\lambda}\) （在 \(\mathcal{L}(\mathrm{E})\) 中，赋有拓扑 \(\mathcal{T}_s\) ）。
\item
  对所有 \(\lambda\) 与 \(\mu\) （属于 \(\mathbf{C} - \mathbf{R}\) ），我们有 \(r_{\lambda} - r_{\mu} = (\mu - \lambda) = r_{\mu}r_{\lambda}\) 以及 \(r_{\lambda}r_{\mu} = r_{\mu}r_{\lambda}\) 。\\
\item
  \(r_{\lambda}\) 的像对所有 \(\lambda \in \mathbf{C} - \mathbf{R}\) 都等于 F。
\end{enumerate}

\(f)\) 对每个 \(\lambda\in\mathbf{C}\text{—}\mathbf{R}\) ，自同态 \(r_{\lambda}\) 是单射；以 F 为定义域的部分算子 \(\lambda1_{\mathrm{E}}-r_{\lambda}^{-1}\) 不依赖于 \(\lambda\in\mathbf{C}\text{—}\mathbf{R}\) 的选取；我们把它记作 \(u\) 。 \(g)\) 部分算子 \(u\) 是自伴的；对所有 \(\lambda\in\mathbf{C}\text{—}\mathbf{R}\) ，我们有 R( \(u,\lambda\) ) = \(r_{\lambda}\) 。

\begin{enumerate}
\def\labelenumi{\alph{enumi})}
\setcounter{enumi}{7}
\tightlist
\item
  序列 \((u_{n})\) 对 \(T_{s}\) -预解拓扑收敛于 u。（“加藤–特罗特定理”）
\end{enumerate}

\begin{enumerate}
\def\labelenumi{\arabic{enumi})}
\setcounter{enumi}{31}
\tightlist
\item
  设 E 为复希尔伯特空间。设 u 与 v 为 E 上的部分自伴算子，使 \(u + v\) 自伴。
\end{enumerate}

\begin{enumerate}
\def\labelenumi{\alph{enumi})}
\item
  对 \(t \in \mathbf{R}^*\) ，令 \(v(t) = t^{-1}(e^{itu}e^{itv} - e^{it(u + v)}) \in \mathcal{L}(\mathrm{E})\) ，并用 \(w(t)\) 表示 \(v(t)\) 到希尔伯特空间 \(\mathrm{E}_{u + v}\) 的限制。用 B 表示连续线性映射 \(w(t)\) 组成的集。集 B 在巴拿赫空间 \(\mathcal{L}(\mathrm{E}_{u + v},\mathrm{E})\) 中有界。（利用巴拿赫–斯坦豪斯定理。）
\item
  我们有 \(w(t) \to 0\) （当 \(t \to 0\) ，在 \(R^{*}\) 中），收敛是在赋有紧收敛拓扑的空间 \(\mathcal{L}(\mathrm{E}_{u+v},\mathrm{E})\) 中取的。
\item
  设 \(x\in \mathrm{E}_{u + v}\) 。我们有
\end{enumerate}

\[
\lim _ {n \to + \infty} (e ^ {i t u / n} e ^ {i t v / n}) ^ {n} x = e ^ {i t (u + v)} x.
\]

对 \(n\in \mathbf{N}\) ，记

\[
\begin{array}{l} (e ^ {i t u / n} e ^ {i t v / n}) ^ {n} - e ^ {i t (u + v)} = \\ \sum_ {i = 0} ^ {n - 1} (e ^ {i t u / n} e ^ {i t v / n}) ^ {k} \Big (e ^ {i t u / n} e ^ {i t v / n} - e ^ {i t (u + v) / n} \Big) (e ^ {i t (u + v) / n}) ^ {n - 1 - i} \end{array}
\]

并把上一问题的结果应用于映射 \(t \mapsto e^{it(u+v)}x\) （从 \([-1,1]\) 到 \(E_{u+v}\) ）的紧像。

\begin{enumerate}
\def\labelenumi{\arabic{enumi})}
\setcounter{enumi}{32}
\tightlist
\item
  设 \(n \in \mathbf{N}\) 。我们用 \(\Delta\) 表示 \(\mathrm{L}^2(\mathbf{R}^n)\) 上的拉普拉斯算子。
\end{enumerate}

\begin{enumerate}
\def\labelenumi{\alph{enumi})}
\tightlist
\item
  对每个 \(t \in R^{*}\) ，自同态 \(e^{it\Delta}\) （属于 \(\mathrm{L}^{2}(\mathbf{R}^{n})\) ）由下式给出
\end{enumerate}

\[
e ^ {i t \Delta} f (x) = \frac {1}{(2 i t) ^ {n / 2}} \int_ {\mathbf {R} ^ {n}} f (y) e ^ {i | x - y | ^ {2} / (4 t)} d y
\]

若 \(f\in \mathcal{L}^1 (\mathbf{R}^n)\cap \mathcal{L}^2 (\mathbf{R}^n)\) （利用傅里叶变换。）

\begin{enumerate}
\def\labelenumi{\alph{enumi})}
\setcounter{enumi}{1}
\tightlist
\item
  设 C 为 \(\mathrm{L}^{2}(\mathbf{R}^{n})\) 上的酉算子，它由映射 \(f \mapsto \mathrm{C}(f)\) （对 \(f \in \mathcal{L}^{2}(\mathbf{R}^{n})\) 定义）通过过渡到商得到
\end{enumerate}

\[
\mathrm{C} (f) (x) = \frac {1}{(2 i t) ^ {n / 2}} e ^ {\frac {i x ^ {2}}{4 t}} \mathcal {F} f \Big (\frac {x}{2 t} \Big)
\]

对 \(x\in \mathbf{R}^n\)

我们有 \(e^{it\Delta} \to C\) （当 \(t \to +\infty\) ），收敛是在赋有单收敛拓扑的空间 \(\mathcal{L}(\mathrm{L}^{2}(\mathbf{R}^{n}))\) 中取的。

\begin{enumerate}
\def\labelenumi{\alph{enumi})}
\setcounter{enumi}{2}
\tightlist
\item
  假设 \(1 \leqslant n \leqslant 3\) 。设 \(g \in \mathcal{L}^{\infty}(\mathbf{R}^{n}) + \mathcal{L}^{2}(\mathbf{R}^{n})\) 为实值函数。用 u 表示部分算子 \(u = \Delta + m_{g}\) ；它在 \(\mathrm{L}^{2}(\mathbf{R}^{n})\) 上是自伴的（习题 14）。对 \(f \in \mathcal{L}^{1}(\mathbf{R}^{n}) \cap \mathcal{L}^{2}(\mathbf{R}^{n})\) 与 \(t \in R^{*}\) ，我们有
\end{enumerate}

\[
e ^ {i t u} f (x) = \lim _ {k \rightarrow + \infty} \frac {1}{(2 i t / k) ^ {k n / 2}} \int_ {\mathbf {R} ^ {n k}} f (y _ {k}) \exp \left(\frac {i t}{k} \sum_ {i = 1} ^ {k} \left(\frac {| y _ {i} - y _ {i - 1} | ^ {2}}{4 (t / k) ^ {2}} - g (y _ {i})\right) d y \right.
\]

其中我们记 \(y = (y_{1}, \ldots, y_{k})\) ，\(y_{i} \in R^{n}\) 。

\begin{enumerate}
\def\labelenumi{\alph{enumi})}
\setcounter{enumi}{3}
\tightlist
\item
  对上一公式的解释。
\end{enumerate}

\begin{enumerate}
\def\labelenumi{\arabic{enumi})}
\setcounter{enumi}{33}
\tightlist
\item
  设 E 为复希尔伯特空间。对 E 的所有自同态 \(u\) 与 \(v\) ，我们记 \([u,v] = uv - vu\) ，并定义 \(\mathrm{E}_{u,v} = \overline{\mathrm{Im}(u)} \cap \mathrm{Ker}(v)\) 。
\end{enumerate}

设 \(p\) 与 \(q\) 为 E 的正交投影算子。我们令 \(a = p - q\) 且 \(b = 1_{\mathrm{E}} - p - q\) 。 a) 若存在 E 的酉自同态 \(u\) 使 \(upu^{-1} = q\) 且 \(uqu^{-1} = p\) ，则 \(\mathrm{E}_{p,q}\) 与 \(\mathrm{E}_{1 - p,1 - q}\) 的希尔伯特维数相等。

\begin{enumerate}
\def\labelenumi{\alph{enumi})}
\setcounter{enumi}{1}
\tightlist
\item
  我们有关系式
\end{enumerate}

\[
a ^ {2} + b ^ {2} = 1 _ {\mathrm{E}}, \qquad a b + b a = 0, \qquad [ p, a ] = [ q, a ] = [ p, b ] = [ q, b ] = 0.
\]

\begin{enumerate}
\def\labelenumi{\alph{enumi})}
\setcounter{enumi}{2}
\tightlist
\item
  直到并包括问题 \(f\) ) 为止，我们都假设 \(\mathrm{E}_{p,q} = \mathrm{E}_{1 - p,1 - q} = \{0\}\) 。自同态 \(b\) 是单射；设 \(s\) 为从 \(\mathbf{R}\) 到 \(\mathbf{R}\) 的、满足下式的函数
\end{enumerate}

\[
s (x) = \left\{ \begin{array}{l l} 1 & \text { if } x > 0 \\ 0 & \text { if } x = 0 \\ - 1 & \text { if } x <   0. \end{array} \right.
\]

自同态 \(u = s(b)\) 是酉的。

\(d)\) 对每个 \(\varepsilon > 0\) ，我们有 \((|b| + \varepsilon)a = a(|b| + \varepsilon)\) 。

\begin{enumerate}
\def\labelenumi{\alph{enumi})}
\setcounter{enumi}{4}
\item
  我们有 \(uau^{-1} = -a\) （验证 \(s(b) = \lim_{\varepsilon \to 0} b(|b| + \varepsilon)^{-1}\) ）以及 \(ubu^{-1} = b\) 。
\item
  我们有 \(upu^{-1} = q\) 且 \(uqu^{-1} = p\) 。
\item
  我们假设 \(E_{p,q}\) 的希尔伯特维数等于 \(E_{1-p,1-q}\) 的希尔伯特维数。则存在 E 的酉自同态 u 使 \(upu^{-1} = q\) 且 \(uqu^{-1} = p\) 。（空间 \(E_{p,q}\) 与 \(E_{1-p,1-q}\) 正交；考虑希尔伯特空间 \(F = E_{p,q} \oplus E_{1-p,1-q}\) ，并对由 p 与 q 通过过渡到子空间导出的 \(F^{\circ}\) 的正交投影算子应用前面的情形。）
\item
  假设 \(\|p-q\|<1\) 。存在 E 的酉自同态 u 使 \(upu^{-1}=q\) 且 \(uqu^{-1}=p\) 。
\item
  假设 \(a = p - q\) 是紧的。自同态 \(k = qp\) 通过过渡到子空间诱导出弗雷德霍姆线性映射 \(\widetilde{k}\) （从 \(p\) 的像到 \(q\) 的像）。存在 E 的酉自同态 \(u\) 使 \(upu^{-1} = q\) 且 \(uqu^{-1} = p\) ，当且仅当 \(\widetilde{k}\) 的指标为零。
\end{enumerate}

\begin{enumerate}
\def\labelenumi{\arabic{enumi})}
\setcounter{enumi}{34}
\tightlist
\item
  设 E 为希尔伯特空间，u 为 E 上的正自伴部分算子。
\end{enumerate}

\begin{enumerate}
\def\labelenumi{\alph{enumi})}
\item
  对每个 \(t \in \mathbf{R}_{+}^{*}\) ，我们有 \(u \circ \mathrm{R}(u, -t) \in \mathcal{L}(\mathrm{E})\) 。
\item
  映射 \(f\) （从 \(\mathbf{R}_{+}^{*}\) 到 \(\mathcal{L}(\mathrm{E})\) ，由 \(f(t) = -t^{-1/2}u \circ \mathrm{R}(u, -t)\) 定义）关于 \(\mathbf{R}_{+}^{*}\) 上的勒贝格测度可积。
\item
  我们有
\end{enumerate}

\[
\sqrt {u} = \frac {1}{\pi} \int_ {\mathbf {R} _ {+} ^ {*}} f (t) d t.
\]

\begin{enumerate}
\def\labelenumi{\arabic{enumi})}
\setcounter{enumi}{35}
\item
  给出一个复希尔伯特空间上的正规部分算子的例子，其本质谱在 C 中不闭。
\item
  设 E 与 F 为复希尔伯特空间，u 为从 E 到 F 的稠定闭部分算子。若 u 是满射的，则存在 \(v \in \mathcal{L}(\mathrm{F}; \mathrm{E})\) 使 \(u \circ v = 1_{F}\) 。
\item
  设 u 为复希尔伯特空间 E 中的部分算子。
\end{enumerate}

\begin{enumerate}
\def\labelenumi{\alph{enumi})}
\item
  假设 u 是正的对称算子。设 \(F(u)\) 为 u 的弗里德里希斯扩张，q 为与 \(F(u)\) 相伴的正部分形式。证明 \(\mathrm{F}(u)\) 是 \(u\) 的唯一一个定义域包含于 \(\operatorname{dom}(q)\) 的自伴扩张。
\item
  假设 \(u\) 是具有稠密定义域的闭部分算子。设 \(q\) 为以 \(\operatorname{dom}(u)\) 为定义域、由 \(q(x,y) = \langle u(x) \mid u(y) \rangle\) 定义的正部分形式。证明表示 \(q\) 的自伴部分算子是 \(v = u^{*}u\) 。
\item
  我们假设 u 对称且 \(u^{2}\) 有稠密定义域。部分算子 \(u^{2}\) 对称且正；\(u^{2}\) 的弗里德里希斯扩张等于 \(u^{*}u\) 。
\end{enumerate}

\begin{enumerate}
\def\labelenumi{\arabic{enumi})}
\setcounter{enumi}{38}
\tightlist
\item
  设 \(n \geqslant 1\) 为整数。我们在 \(\mathbf{R}^n\) 上配备勒贝格测度。
\end{enumerate}

设 v 为 \(R^{n}\) 上局部平方可积的正函数。我们用 \(\Delta\) 表示 \(\mathcal{D}(\mathbf{R}^{n})\) 上或 \(\mathcal{D}'(\mathbf{R}^{n})\) 上的拉普拉斯算子，它由下式定义

\[
\Delta (f) = - \sum_ {i = 1} ^ {n} \partial_ {i} ^ {2} f.
\]

设 D 为 \(L^{2}(\mathbf{R}^{n})\) 上以 \(\mathcal{D}(\mathbf{R}^{n})\) 为定义域的部分算子，它满足

\[
\mathrm{D} (\varphi) = \Delta (\varphi + v \varphi).
\]

\begin{enumerate}
\def\labelenumi{\alph{enumi})}
\item
  部分算子 D 对称且正。
\item
  设 \(u \in \text{dom}(D^{*})\) 满足 \(D^{*}u + u = 0\) 。我们有 \(\Delta(u) + u = 0\) ，其中 \(u \in \text{L}^{2}(\mathbf{R}^{n})\) 的拉普拉斯算子是在广义函数意义下计算的。
\item
  广义函数 \(\Delta(u)\) 与 \((v+1)u\) 是局部可积函数，且广义函数 \(-\Delta(u)\) 是正的。（利用 IV 第 344 页习题 32。）
\item
  我们有 u = 0。（把 \(|u|\) 用无穷次可微函数 \(w_{\varepsilon}\) 通过卷积正则化，使得 \(\Delta(w_{\varepsilon}) \geqslant 0\) ，并应用上一问题。） e) 部分算子 D 是本质自伴的。
\end{enumerate}

\begin{enumerate}
\def\labelenumi{\arabic{enumi})}
\setcounter{enumi}{39}
\tightlist
\item
  设 E 为复希尔伯特空间。
\end{enumerate}

\begin{enumerate}
\def\labelenumi{\alph{enumi})}
\item
  设 u 为 E 上的正规部分算子，\((j, |u|)\) 为它的极分解。我们有 \(|u^{*}| = |u|\) ，且 \(u^{*}\) 的极分解是 \((j^{*}, |u|)\) 。我们有 \(j|u| = |u|j\) 。
\item
  设 u 为 E 上的自伴部分算子，\((j, |u|)\) 为它的极分解。我们有 \(j = \mathrm{s}(u)\) ，其中函数 s 由下式定义
\end{enumerate}

\[
\mathrm{s} (x) = \left\{ \begin{array}{l l} 1 & \text { if } x > 0 \\ 0 & \text { if } x = 0 \\ - 1 & \text { if } x <   0. \end{array} \right.
\]

\begin{enumerate}
\def\labelenumi{\alph{enumi})}
\setcounter{enumi}{2}
\tightlist
\item
  设 u 为 E 上的自伴部分算子，\((j, |u|)\) 为它的极分解。我们用 D 表示 \(|u|^{1/2}\) 的定义域，并通过令下式成立来定义以 D 为定义域的部分埃尔米特形式 q
\end{enumerate}

\[
q (x, y) = \langle | u | ^ {1 / 2} (x) \mid j | u | ^ {1 / 2} (y) \rangle .
\]

表示 \(q\) 的部分算子与 \(u\) 重合；若 \(u\) 是正的，则 \(q\) 是与 \(u\) 相伴的正部分形式。

\begin{enumerate}
\def\labelenumi{\alph{enumi})}
\setcounter{enumi}{3}
\tightlist
\item
  设 u 与 v 为 E 上的自伴部分算子。用 \(q_{u}\) 与 \(q_{v}\) 表示按问题 c) 的方式定义的部分埃尔米特形式。假设 \(\text{dom}(q_{u}) \cap \text{dom}(q_{v})\) 在 E 中稠密。u 与 v 在埃尔米特形式意义下的和是 E 上表示部分形式的部分算子 \(u + v\)
\end{enumerate}

\[
q (x, y) = q _ {u} (x, y) + q _ {v} (x, y)
\]

其定义域为 \(\mathrm{dom}(q_{u})\cap\mathrm{dom}(q_{v})\) 。证明部分算子 \(u+v\) 是 \(u+v\) 的扩张；若 u 与 v 是正的，则 \(u+v\) 是正的自伴部分算子。

\begin{enumerate}
\def\labelenumi{\alph{enumi})}
\setcounter{enumi}{4}
\item
  给出一个例子，其中 \(u + v\) 的定义域不稠密，但 \(u + v\) 有定义。
\item
  设 \(n \geqslant 1\) 为整数，\(S \subset R^{n}\) 为勒贝格测度为零的闭集。设 v 为 \(R^{n}\) 上关于勒贝格测度可测的正函数，且 v 在 \(R^{n}-S\) 上局部可积。自伴部分算子 \(\Delta + m_{v}\) 有定义。
\end{enumerate}

\begin{enumerate}
\def\labelenumi{\arabic{enumi})}
\setcounter{enumi}{40}
\tightlist
\item
  设 \(n \geqslant 1\) 为整数。我们用 \(\Delta\) 表示 \(\mathbf{R}^n\) 上的拉普拉斯算子，用 \(u\) 表示 \(\Delta\) 到子空间 \(\mathcal{D}(\mathbf{R}^n - \{0\})\) 的约化。
\end{enumerate}

\begin{enumerate}
\def\labelenumi{\alph{enumi})}
\item
  u 的定义域在希尔伯特空间 \(E_{q}\) （它同与 \(\Delta\) 相伴的部分形式 q 相伴）中稠密，当且仅当 \(n \geqslant 2\) 。（把 \(E_{q}\) 的对偶与一个缓增广义函数空间等同起来；若 \(g \in E_{q}'\) 在 dom(u) 上为零，证明 g 是支撑包含于 \(\{0\}\) 的广义函数；然后利用 IV 第 338 页习题 24 得出结论。）
\item
  \(u\) 的定义域在希尔伯特空间 \(\mathrm{E}_{\Delta}\) 中稠密，当且仅当 \(n \geqslant 4\) 。
\end{enumerate}

\begin{enumerate}
\def\labelenumi{\arabic{enumi})}
\setcounter{enumi}{41}
\tightlist
\item
  设 E 为复希尔伯特空间，u 为 E 上具有稠密定义域的闭部分算子。我们称 E 上的部分算子 v 关于 u 是相对紧的，如果 v 的定义域包含 u 的定义域，且 v 到 \(\text{dom}(u)\) 的限制是 \(E_{u}\) 到 E 的紧线性映射。
\end{enumerate}

\begin{enumerate}
\def\labelenumi{\alph{enumi})}
\item
  每个部分算子 \(v\) （关于 \(u\) 相对紧）都关于 \(u\) 有界，且其相对范数 \(\| v \|_u\) 为零。（若相对范数不为零，构造收敛于 0 的序列 \((x_n)\) （在 \(\text{dom}(u)\) 中），使 \(\| v(x_n) \| = 1\) （对所有 \(n\) ），且使 \((u(x_n))\) 弱收敛于 E 中的元素 \(y\) ；证明 \(y = 0\) ，由此推出 \(v\) 不是相对紧的。）
\item
  设 \(\lambda \in \mathbf{C} - \mathrm{Sp}(u)\) 。设 \(v\) 为 E 上的部分算子，其定义域包含 \(\mathrm{dom}(u)\) 。则 \(v\) 关于 \(u\) 相对紧，当且仅当 \(v \circ \mathrm{R}(u, \lambda)\) 是 E 的紧自同态。
\end{enumerate}

\begin{enumerate}
\def\labelenumi{\arabic{enumi})}
\setcounter{enumi}{42}
\tightlist
\item
  设 E 为复希尔伯特空间。
\end{enumerate}

\begin{enumerate}
\def\labelenumi{\alph{enumi})}
\tightlist
\item
  设 u 为 E 上的自伴部分算子。对每个 \(\lambda \in R\) ，我们有 \(\lambda \in \mathrm{Sp}_{\mathrm{e}}(u)\) ，当且仅当存在 E 中的序列 \((x_{n})\) ，它在 E 中弱收敛于 0 且满足条件
\end{enumerate}

\[
\liminf _ {n \to + \infty} \| x _ {n} \| > 0, \quad \lim _ {n \to + \infty} (u (x _ {n}) - \lambda x _ {n}) = 0.
\]

\begin{enumerate}
\def\labelenumi{\alph{enumi})}
\setcounter{enumi}{1}
\item
  设 \(u\) 与 \(v\) 为 E 上的自伴部分算子。若存在 \(\lambda \in \mathbf{C} - (\mathrm{Sp}(u) \cup \mathrm{Sp}(v))\) 使 \(\mathrm{R}(u, \lambda) - \mathrm{R}(v, \lambda)\) 紧，则 \(\mathrm{Sp}_{\mathrm{e}}(u) = \mathrm{Sp}_{\mathrm{e}}(v)\) 。
\item
  设 u 为 E 上的闭对称部分算子。设 \(u_{1}\) 与 \(u_{2}\) 为 u 的自伴扩张。若空间 \(\operatorname{Ker}(u^{*}-i)\) 是有限维的，则 \(\operatorname{Sp}_{\mathrm{e}}(u_{1})=\operatorname{Sp}_{\mathrm{e}}(u_{2})\) 。
\item
  设 u 为 E 上的自伴部分算子，v 为关于 u 相对紧的部分算子（见上一习题）。部分算子 \(u + v\) 自伴，且 \(\mathrm{Sp}_{\mathrm{e}}(u + v) = \mathrm{Sp}_{\mathrm{e}}(u)\) 。
\end{enumerate}

\chapter{酉表示}
\setcounter{exercisepar}{0}

对每个向量空间 E，我们用 \(1_{E}\) 表示 E 的恒等映射。若 E、F 与 G 为向量空间，且 u: \(F \rightarrow E\) 与 v: \(G \rightarrow F\) 为线性映射，我们有时用 uv 表示从 G 到 E 的线性映射 \(u \circ v\) 。

对希尔伯特空间 E，我们用 \(\langle x | y \rangle_{E}\) ，或简单地 \(\langle x | y \rangle\) ，表示 E 上的标量积。

若 X 为局部紧拓扑空间，\(\mu\) 为 X 上的测度，我们记 \(\mathcal{L}^{p}(\mathrm{X},\mu)=\mathcal{L}_{\mathbf{C}}^{p}(\mathrm{X},\mu)\) 与 \(\mathrm{L}^{p}(\mathrm{X},\mu)=\mathrm{L}_{\mathbf{C}}^{p}(\mathrm{X},\mu)\) （对 \(p\in[1,+\infty]\) ）。

我们用 e 表示群 G 的单位元。

设 A 为结合代数，E 为拓扑向量空间，\(\pi\colon\mathrm{A}\to\mathcal{L}(\mathrm{E})\) 为 A 在 E 中的一个表示。我们称 \(\pi\) 是非退化的，如果由向量 \(\pi(a)x\) （\(a\in\mathrm{A}\) ，\(x\in\mathrm{E}\) ）组成的集在 E 中是完全的。

\section{酉表示}

在本节中，向量空间都取在 K = R 或 C 上。

回忆一下，群 G 在 K-向量空间 E 中的线性表示是一个同态 \(\varrho\) （从 G 到 E 的自同构群 \(\mathbf{GL}(\mathrm{E})\) 中）（A, VIII, 第 387 页, def. 1）。我们称 E 为表示 \(\varrho\) 的空间，称 E 的维数为 \(\varrho\) 的维数，也记作 \(\dim(\varrho)\) 。

我们可以把 G 在 K-向量空间中的一个表示与一个 K{[}G{]}-模等同起来（loc. cit.），并将使用相应的术语，例如关于表示的直和、表示的态射，或由 R-向量空间中的表示通过纯量扩张得到的 C 上的表示。

我们称表示 \(\varrho\colon\mathbf{G}\to\mathbf{GL}(\mathbf{E})\) 是忠实的，如果它是单射同态。

\pandocbounded{\includegraphics[keepaspectratio]{images/part003_3d1522a8.jpg}}

注意这并不意味着 K{[}G{]}-模块（与 \(\varrho\) 相伴）是忠实 K-模（A, II, 第 28 页）。

G 在有限维 K-向量空间中的表示的特征标是从 G 到 K 的映射 \(\chi_{\varrho}\) ，它满足 \(\chi_{\varrho}(g) = \operatorname{Tr}(\varrho(g))\) （对所有 \(g \in G\) ），参见 A, VIII, 第 388 页。

\subsection{关于连续线性表示的回顾}

设 G 为拓扑群。我们回忆，G 在拓扑 K-向量空间 E 中的连续线性表示是 G 在 E 中的线性表示 \(\varrho\) ，使得由 \((g, x) \mapsto \varrho(g)x\) 定义的从 G × E 到 E 的映射连续（INT, VIII, 第 128 页, § 2, n° 1, def. 1）。这个映射定义了 G 在 E 上的一个作用，且 \(\varrho\) 的像包含于 \(\mathcal{L}(\mathrm{E})\) 。我们称连续表示 \(\varrho\) 是有界的，如果它的像在赋有有界收敛拓扑的空间 \(\mathcal{L}(\mathrm{E})\) 中有界。

注. —— 若 K = R 且 \(\varrho\) 是 G 在 R-向量空间 E 中的连续线性表示，则映射 \(\varrho_{(\mathbf{C})}: g \mapsto 1 \otimes \varrho(g)\) 是 G 在 \(\mathrm{E}_{(\mathbf{C})}\) 中的连续线性表示；若 \(\varrho\) 有界，则 \(\varrho_{(\mathbf{C})}\) 有界。

对每个拓扑 K-向量空间 E，把每个 \(g \in G\) 映为 \(1_{E}\) 的同态是 G 在 E 中的连续线性表示，称为 G 在 E 中的平凡表示。

设 H 为 G 的子群。\(\varrho\) 到 H 的限制是 H 在 E 中的连续线性表示，记作 \(\mathrm{Res}_{\mathrm{H}}^{\mathrm{G}}(\varrho)\) 。

设 \(\varrho_{1}\) 与 \(\varrho_{2}\) 为 G 在拓扑 K-向量空间 E \(_{1}\) 与 E \(_{2}\) 中的连续线性表示。G-态射 \(u\) （从 \(\varrho_{1}\) 到 \(\varrho_{2}\) ）是连续的线性表示态射，即连续线性映射 \(u\colon\) E \(_{1} \to\) E \(_{2}\) ，满足 \(u \circ \varrho_{1}(g) = \varrho_{2}(g) \circ u\) （对所有 \(g \in\) G）。我们用 Hom \(_{G}\) ( \(\varrho_{1}, \varrho_{2}\) ) 表示从 \(\varrho_{1}\) 到 \(\varrho_{2}\) 的 G-态射组成的向量空间。它是赋有逐点收敛拓扑的空间 \(\mathcal{L}\) (E \(_{1}\) ; E \(_{2}\) ) 的闭子空间。

设 \(\varrho\) 为群 G 在拓扑 K-向量空间 E 中的连续线性表示。E 的恒等映射 \(1_{\mathrm{E}}\) 是 \(\varrho\) 到 \(\varrho\) 的 G-态射，记作 \(1_{\varrho}\) 。若 \(\varrho_1\) 、 \(\varrho_2\) 与 \(\varrho_3\) 分别是 G 在拓扑 K-向量空间 \(\mathrm{E}_1\) 、 \(\mathrm{E}_2\) 、 \(\mathrm{E}_3\) 中的连续线性表示，且 \(u: \mathrm{E}_1 \to \mathrm{E}_2\) 与 \(v: \mathrm{E}_2 \to \mathrm{E}_3\) 是 G-态射，则 \(v \circ u\) 是 G-态射。

G-态射 \(u\) （从 \(\varrho_{1}\) 到 \(\varrho_{2}\) ）是 G-同构，如果存在 G-态射 \(v\) （从 \(\varrho_{2}\) 到 \(\varrho_{1}\) ），使 \(v\circ u=1_{\varrho_{1}}\) 且 \(u\circ v=1_{\varrho_{2}}\) 。为此必要且充分的是：\(u\) 是 G-态射，且是空间 \(\varrho_{1}\) 到空间 \(\varrho_{2}\) 的拓扑向量空间同构；其逆 \(v=u^{-1}\) 那时确实是 G-态射。若存在从 \(\varrho_{1}\) 到 \(\varrho_{2}\) 的 G-同构，我们就说这些表示是同构的。

定义 1.—— 设 \(\varrho\) 为拓扑群 G 在拓扑向量空间 E 中的连续线性表示。我们称 \(\pi\) 是 \(\varrho\) 的一个子表示，如果 F 是 E 的闭子空间，且对每个 \(g \in G\) ，空间 F 在 \(\varrho(g)\) 下稳定，而 \(\pi(g)\) 是 \(\varrho(g)\) 通过过渡到子空间诱导的 F 的自同态。

设 \(\varrho\) 为拓扑群 G 在拓扑向量空间 E 中的连续线性表示。\(\{0\}\) 在 E 中的闭包是 \(\varrho\) 的子表示。更一般地，在诸自同态 \(\varrho(g)\) 下稳定的子空间的闭包是 \(\varrho\) 的子表示。

子表示 \(\pi\) （属于 \(\varrho\) ）由一个在所有自同态 \(\varrho(g)\) 下稳定的闭子空间 F 唯一确定。后一条件常表述为“F 是 E 的 G-不变子空间”。这时我们也说 F 定义了 \(\varrho\) 的一个子表示，或者有时滥用语言，说 F 是 \(\varrho\) 或 E 的子表示。

设 F 为 E 的一个子空间，它定义了子表示 \(\pi\) （\(\varrho\) 的）。对每个 \(g \in G\) ，线性映射 \(\varrho(g)\) 通过过渡到商定义了 E/F 的一个自同态 \(\widetilde{\varrho}(g)\) 。映射 \(g \mapsto \widetilde{\varrho}(g)\) 是 G 在 E/F 中的连续线性表示，且 E 到 E/F 的典范投影是 G-态射。我们称 \(\widetilde{\varrho}\) 为 G 在 E/F 上的商表示；它也记作 \(\varrho/\pi\) 。

E 中在 G 的作用下不变的元素组成的子空间是 \(\varrho\) 的平凡子表示。我们把它记作 \(\varrho^{G}\) 或 \(E^{G}\) 。

设 A 为 E 的子集。由元素 \(\varrho(g)x\) （其中 \(g \in G\) ，\(x \in A\) ）生成的闭子空间 F 是 \(\varrho\) 的子表示，称为由 A 生成的 \(\varrho\) 的子表示；若 F = E，就说 A 生成 \(\varrho\) 。若子表示 F 是有限维的，就说子集 A 是 G-有限的。

假设 A 缩减为单个元素 \(x \in E\) 。若 A 生成 \(\varrho\) ，就称 x 是 \(\varrho\) 的循环向量，并称 \(\varrho\) 是循环表示；若 A 是 G-有限的，就称 x 是 G-有限向量。

\(\varrho\) 的 G-有限向量组成的集是 E 的在 \(\varrho\) 下稳定的向量子空间；它不一定是闭的。

设 \((\varrho_{i})_{i\in\mathrm{I}}\) 为 G 在局部凸拓扑 K-向量空间 \((\mathrm{E}_{i})_{i\in\mathrm{I}}\) 中的一族连续线性表示。设 E 为诸空间 \(E_{i}\) 的直和空间。映射 \(g\mapsto(\varrho_{i}(g))\) 是 G 在 E 中的线性表示，称为诸表示 \(\varrho_{i}\) 的和或直和。若 I 有限，则它是连续的。

若表示 \(\varrho_{i}\) 都等于某个表示 \(\varrho\) （对所有 \(i \in I\) ），我们就说诸表示 \(\varrho_{i}\) 的直和是表示 \(\varrho\) 的 Card(I) 个拷贝之和，它也记作 \(\varrho^{\mathrm{Card(I)}}\) ，或 Card(I) \(\varrho\) 。

设 \((\varrho_{i})_{i\in\mathrm{I}}\) 与 \((\pi_{j})_{j\in\mathrm{J}}\) 分别为 G 在拓扑 K-向量空间 \((\mathrm{E}_{i})_{i\in\mathrm{I}}\) 与 \((\mathrm{F}_{j})_{j\in\mathrm{J}}\) 中的有限族连续线性表示。K-向量空间的典范同构

\[
\operatorname{Hom} _ {\mathrm{K}} \left(\bigoplus_ {i \in \mathrm{I}} \mathrm{E} _ {i}, \bigoplus_ {j \in \mathrm{J}} \mathrm{F} _ {j}\right)\rightarrow \bigoplus_ {(i, j) \in \mathrm{I} \times \mathrm{J}} \operatorname{Hom} _ {\mathrm{K}} (\mathrm{E} _ {i}, \mathrm{F} _ {j})
\]

（A, II, 第 13 页, cor. 1）通过过渡到子空间诱导出同构

\[
\mathrm{Hom} _ {\mathrm{G}} \Bigl (\bigoplus_ {i \in \mathrm{I}} \varrho_ {i}, \bigoplus_ {j \in \mathrm{J}} \pi_ {j} \Bigr) \to \bigoplus_ {(i, j) \in \mathrm{I} \times \mathrm{J}} \mathrm{Hom} _ {\mathrm{G}} (\varrho_ {i}, \pi_ {j})\tag{1}
\]

同样称它是典范的。

设 \(\varrho_{1}\) 与 \(\varrho_{2}\) 为 G 在 K-向量空间 E \(_{1}\) 与 E \(_{2}\) 中的线性表示。对 \(u \in \mathcal{L}(\mathrm{E}_{1};\mathrm{E}_{2})\) 与 \(g \in \mathrm{G}\) ，我们令 \(\varrho(g)u = \varrho_{2}(g) \circ u \circ \varrho_{1}(g^{-1})\) 。映射 \(g \mapsto \varrho(g)\) 是 G 在 \(\mathcal{L}(\mathrm{E}_{1};\mathrm{E}_{2})\) 中的线性表示。这个表示的不变元素组成的空间与 \(\mathrm{Hom}_{\mathrm{G}}(\varrho_{1},\varrho_{2})\) 重合。

设 \(\varrho\) 为 G 在局部凸空间 E 中的连续线性表示。回忆（INT, VIII, 第 131 页, § 2, no. 2）逆步表示 \(\breve{\varrho}\) （\(\varrho\) 的）是 G 在 E 的对偶 E' 中的线性表示，由 \(\breve{\varrho}(g) = {}^{t} \varrho(g^{-1})\) 定义。

\subsection{连续性的一个判别法}

下述结果使人能够验证拓扑群的直积的某些线性表示是连续的。

引理 1.—— 设 G 与 H 为拓扑群，E 为巴拿赫空间。设 \(\varrho\) 为 G 在 E 中的有界表示，\(\pi\) 为 H 在 E 中的有界表示。假设 \(\varrho(g)\) 与 \(\pi(h)\) 交换（对所有 \((g,h)\in\mathrm{G}\times\mathrm{H}\) ）。映射 \(\varpi\) （从 G \(\times\) H 到 \(\mathbf{GL}(\mathrm{E})\) ，由 \((g,h)\mapsto\varrho(g)\circ\pi(h)\) 定义）是 G \(\times\) H 在 E 中的连续且有界的线性表示。

映射 \(\varpi\) 是 G × H 在 E 中的线性表示；我们来验证它连续。由于 \(\| \varpi(g, h) \| \leqslant \| \varrho(g) \| \| \pi(h) \|\) （对所有 \((g, h) \in \mathrm{G} \times \mathrm{H}\) ），表示 \(\varpi\) 有界，故其像在 \(\mathcal{L}(\mathrm{E})\) 中等度连续。于是只需证明：对每个 \(x \in \mathrm{E}\) ，映射 \((g, h) \mapsto \varpi(g, h)x\) 连续（INT, VIII, 第 129 页, § 2, no. 1 的注 2）。设 \((g_0, h_0) \in \mathrm{G} \times \mathrm{H}\) 。令 \(y = \pi(h_0)x\) 。对每个 \((g, h) \in \mathrm{G} \times \mathrm{H}\) ，我们得到

\[
\begin{array}{r l} & {\| \varpi (g, h) x - \varpi (g _ {0}, h _ {0}) x \| \leqslant \| \varrho (g) (\pi (h) x - y) \| + \| \varrho (g) y - \varrho (g _ {0}) y \|} \\ & {\qquad \leqslant \| \varrho (g) \| \| \pi (h) x - y \| + \| \varrho (g) y - \varrho (g _ {0}) y \|.} \end{array}
\]

由于 \(\varrho\) 有界且 \(\varrho\) 与 \(\pi\) 连续，这就得到结论。
\subsection{连续有限维表示}

设 \(\varrho\) 是拓扑群 G 在有限维分离拓扑 K-向量空间 E 中的连续线性表示。我们赋予 \(\mathcal{L}(\mathrm{E})\) 其在 K 上唯一的分离拓扑向量空间结构；于是，态射 \(\varrho\colon\mathrm{G}\to\mathrm{GL}(\mathrm{E})\) 连续，因为 \(\mathrm{GL}(\mathrm{E})\) 的拓扑由 \(\mathcal{L}(\mathrm{E})\) 的拓扑诱导，而该拓扑与逐点收敛拓扑重合（EVT, I, p. 14, th. 2）。因此，表示 \(\varrho\) 是连续的。如果 G 是实李群，那么 \(\varrho\) 的特征标是 G 上的解析函数（LIE, III, p. 225, § 8, no. 1, th. 1）。当 E' 赋予其作为 K 上分离拓扑向量空间的唯一拓扑时，逆步表示也是连续的。此外，对每个整数 \(n\geqslant0\)，当相应空间赋予其作为 K 上分离拓扑向量空间的拓扑时，表示 \(\mathrm{T}^{n}(\varrho)\) 、\(\mathrm{S}^{n}(\varrho)\) 和 \(\wedge^{n}(\varrho)\) （同上）都是连续的。

设 \(\varrho_{1}\) 和 \(\varrho_{2}\) 是拓扑群 G 在有限维分离拓扑向量空间 E \(_{1}\) 和 E \(_{2}\) 中的连续表示。由 \(\varrho_{1} \otimes \varrho_{2}\) 给出的 G 在 E \(_{1} \otimes\) E \(_{2}\) 中的线性表示（LIE, III, p. 256, 附录）是连续的；空间 E \(_{1} \otimes\) E \(_{2}\) 赋予其作为 K 上分离拓扑向量空间的拓扑。

\subsection{不可约表示}

本节中，G 是拓扑群。

定义 2. --- 拓扑 K-向量空间 E 中 G 的表示 \(\varrho\) 称为不可约的，是指 \(\{0\}\) 在 E 中不稠密，并且 \(\varrho\) 的唯一子表示是 \(\varrho\) 本身和在 \(\{0\}\) 的闭包中诱导的表示。

如果 \(\varrho\) 是分离拓扑向量空间 E 中 G 的不可约表示，那么 E 的每个非零元素都是 \(\varrho\) 的循环向量。

引理 2. --- 设 \(\pi\) 和 \(\varrho\) 是分离拓扑 K-向量空间中 G 的连续线性表示。假设 \(\pi\) 不可约。从 \(\pi\) 到 \(\varrho\) 的每个非零 G-态射都是单射，而从 \(\varrho\) 到 \(\pi\) 的每个非零 G-态射都有稠密像。

特别地，如果 \(\pi\) 和 \(\varrho\) 都不可约且有限维，那么从 \(\pi\) 到 \(\varrho\) 的每个非零 G-态射都是同构。

由于 \(\varrho\) 的表示空间是分离的，G-态射 \(u\)（从 \(\pi\) 到 \(\varrho\)）的核是闭的，因而定义 \(\pi\) 的一个子表示。如果态射 \(u\) 非零，其核必为 \(\{0\}\)，因为 \(\pi\) 是分离拓扑向量空间中的不可约表示。同样，从 \(\varrho\) 到 \(\pi\) 的非零 G-态射的像的闭包是 \(\pi\) 的一个非零子表示，因而等于 \(\pi\) 的表示空间。

最后一个断言由上述内容推出。

引理 3. --- 设 \(\varrho\) 是分离拓扑 K-向量空间 E 中的有限维连续线性表示。如果 E 非零，那么 \(\varrho\) 存在一个不可约子表示。

由于 E 有限维且非零，存在维数最小的非零 G-不变子空间 F。该子空间在 E 中闭，并定义 E 的一个子表示；F 的每个子表示也是 E 的子表示，因此由极小性，F 中的表示不可约。

注. --- G 的一个非零表示 E 并不总是包含不可约子表示（参见 V, p. 426, 注）。不过将会看到，当 G 是紧群、K = C，且 E 是 K 上非零的分离拟完备局部凸空间时，情况确实如此（V, p. 464 的命题 7）。

\subsection{酉表示}

定义 3. --- 设 G 是拓扑群，E 是 K 上的希尔伯特空间。G 在 E 中的酉表示，是 G 在 E 中的连续线性表示 \(\varrho\)，使得对 G 中每个 g，E 的自同态 \(\varrho(g)\) 都是 E 的酉自同态（EVT, V, p. 40）。

换言之，酉表示是希尔伯特空间中的等距表示。特别地，酉表示是有界的。

拓扑群在希尔伯特空间中的平凡表示是酉的。酉表示的每个子表示都是酉的。酉表示对子群的限制是酉的。

定义 4. --- 设 G 是拓扑群，\(\varrho\) 和 \(\pi\) 分别是 G 在希尔伯特空间 E 和 F 中的酉表示。E × F 上的 G-不变半双线性形式，是 E × F 上的连续半双线性形式 q，使得

\[
q (x, y) = q (\varrho (g) x, \pi (g) y)
\]

对所有 \(g \in \mathrm{G}\) 及每个 \((x, y) \in \mathrm{E} \times \mathrm{F}\) 成立。

E × F 上的 G-不变半双线性形式组成的向量空间记为 \(\mathrm{Sesq}_{\mathrm{G}}(\varrho,\pi)\) 或 \(\mathrm{Sesq}_{\mathrm{G}}(\mathrm{E},\mathrm{F})\)。

例. --- 设 G 是拓扑群，\(\varrho\) 是 G 在 E 中的酉表示。E 上的内积 \(q\) 是 \(\mathrm{E} \times \mathrm{E}\) 上的 G-不变半双线性形式。

引理 4. --- 设 G 是拓扑群，\(\varrho\) 是 G 到希尔伯特空间 E 的酉群 U(E) 中的同态。那么，\(\varrho\) 是 G 的酉表示，当且仅当对于 U(E) 上的逐点收敛拓扑，\(\varrho\) 在 G 的元素 e 处连续。只需该性质对 E 的一个总子集中的每个 \(x\) 成立即可。

该条件显然是必要的。根据 INT, VIII, p. 129, § 2, no. 1 的注 2，它也是充分的，因为 \(\varrho\) 的像在 \(\mathcal{L}(\mathrm{E})\) 中等度连续，并且映射 \(g \mapsto \varrho(g)x\) 在 \(e\) 处的连续性蕴含其在 G 上的连续性。

设 G 是拓扑群。如果 \(\varrho_{1}\) 和 \(\varrho_{2}\) 是 G 的酉表示，且 u 属于 \(\operatorname{Hom}_{\mathrm{G}}(\varrho_{1},\varrho_{2})\)，那么 \(u^{*}\) 属于 \(\operatorname{Hom}_{\mathrm{G}}(\varrho_{2},\varrho_{1})\)。事实上，由于 \(\varrho_{1}\) 和 \(\varrho_{2}\) 是酉表示，对每个 \(g\in G\) 有

\[
\begin{array}{r} u ^ {*} \circ \varrho_ {2} (g) = u ^ {*} \circ \varrho_ {2} (g ^ {- 1}) ^ {*} = (\varrho_ {2} (g ^ {- 1}) \circ u) ^ {*} \\ = (u \circ \varrho_ {1} (g ^ {- 1})) ^ {*} = \varrho_ {1} (g) \circ u ^ {*}. \end{array}
\]

引理 5. --- 设 G 是拓扑群，\(\varrho\) 是 G 在复希尔伯特空间 E 中的酉表示。空间 \(\mathrm{Hom}_{\mathrm{G}}(\varrho, \varrho)\) 是 \(\mathcal{L}(\mathrm{E})\) 的含单位星子代数。

前述内容表明，\(\operatorname{Hom}_{\mathrm{G}}(\varrho,\varrho)\) 是 \(\mathcal{L}(\mathrm{E})\) 的含单位自伴子代数。由于它在 \(\mathcal{L}(\mathrm{E})\) 中是闭的，所以它是 \(\mathcal{L}(\mathrm{E})\) 的星子代数。

引理 6. --- 设 \(\pi\) 和 \(\varrho\) 分别是拓扑群 G 在希尔伯特空间 E 和 F 中的酉表示。设 D 是 E 的稠密子空间，并且在 \(\pi\) 下稳定。设 u 是从 E 到 F、定义域为 D 的闭部分算子，并满足关系式 \(u \circ \pi(g) = \varrho(g) \circ u\) 对所有 \(g \in G\) 成立。则 \(u^{*}\) 的定义域在 \(\varrho\) 下稳定，并且有关系式 \(u^{*} \circ \varrho(g) = \pi(g) \circ u^{*}\) 对所有 \(g \in G\) 成立。

设 \(g \in \mathrm{G}\)，且 \(x \in \mathrm{dom}(u^*)\)。对于每个 \(y \in \mathrm{dom}(u)\)，有

\[
\begin{array}{r l} & {\langle \varrho (g) x \mid u (y) \rangle = \langle x \mid \varrho (g ^ {- 1}) u (y) \rangle = \langle x \mid u (\pi (g ^ {- 1}) y) \rangle} \\ & {\qquad = \langle u ^ {*} (x) \mid \pi (g ^ {- 1}) y \rangle = \langle \pi (g) (u ^ {*} (x)) \mid y \rangle ,} \end{array}
\]

因为 \(\varrho(g)^{*} = \varrho(g)^{-1} = \varrho(g^{-1})\)。这证明了 \(\varrho(g)x \in \mathrm{dom}(u^{*})\)，且 \(u^{*}(\varrho(g)x) = \pi(g)(u^{*}(x))\)。特别地，\(u^{*} \circ \varrho(g)\) 的定义域包含 \(u^{*}\) 的定义域，并且 \(\pi(g) \circ u^{*} \subset u^{*} \circ \varrho(g)\)。

此外，若 \(x \in \text{dom}(u^{*} \circ \varrho(g))\)，则 \(x = \varrho(g^{-1})(\varrho(g)x)\) 属于 \(\text{dom}(u^{*})\)，这是将前述结论应用于 \(g^{-1}\) 的结果。因此 \(u^{*} \circ \varrho(g) = \pi(g) \circ u^{*}\)。

命题 1. --- 设 \(\varrho_{1}\) 和 \(\varrho_{2}\) 分别是拓扑群 G 在希尔伯特空间 E \(_{1}\) 和 E \(_{2}\) 中的酉表示。从 Hom \(_{G}(\varrho_{1},\varrho_{2})\) 到 Sesq \(_{G}(\varrho_{2},\varrho_{1})\) 的映射将 u 关联到半双线性形式 \(q_{u}\)，该形式由 \(q_{u}(x,y)=\langle x|u(y)\rangle\) 定义；它是向量空间的同构。

若 \(u\) 是 G-态射，则有

\[
\begin{array}{c} q _ {u} (\varrho_ {2} (g) x, \varrho_ {1} (g) y) = \langle \varrho_ {2} (g) x \mid u (\varrho_ {1} (g)y) \rangle \\ = \langle \varrho_ {2} (g) x \mid \varrho_ {2} (g) u (y) \rangle = \langle x \mid u (y) \rangle = q _ {u} (x, y) \end{array}
\]

对所有 \(g \in G\) 及每个 \((x, y) \in \mathrm{E}_{2} \times \mathrm{E}_{1}\) 成立，因此所指映射是从 \(\operatorname{Hom}_{\mathrm{G}}(\varrho_{1}, \varrho_{2})\) 到 \(\operatorname{Sesq}_{\mathrm{G}}(\varrho_{2}, \varrho_{1})\) 的线性映射。由 EVT, V, p. 16, cor. 2，它是单射。

反之，设 \(q\) 是 G-不变半双线性形式，且 \(u\) 是从 \(\mathrm{E}_1\) 到 \(\mathrm{E}_2\) 的唯一线性映射，使得 \(q(x,y) = \langle x \mid u(y) \rangle\) 对所有 \((x,y) \in \mathrm{E}_2 \times \mathrm{E}_1\) 成立（同上）。对于每个 \(g\)（\(G\) 中）及每个 \((x,y)\)

在 \(\mathrm{E}_2\times \mathrm{E}_1\) 中，我们有

\[
\begin{array}{c} \langle x | (\varrho_ {2} (g) \circ u \circ \varrho_ {1} (g ^ {- 1})) (y) \rangle = \langle \varrho_ {2} (g ^ {- 1}) x | u (\varrho_ {1} (g ^ {- 1}) y) \rangle \\ = q (\varrho_ {2} (g ^ {- 1}) x, \varrho_ {1} (g ^ {- 1}) y) = q (x, y) = \langle x | u (y) \rangle , \end{array}
\]

因此 \(u = \varrho_{2}(g) \circ u \circ \varrho_{1}(g^{-1})\)。于是 \(u \in \operatorname{Hom}_{\mathrm{G}}(\varrho_{1}, \varrho_{2})\)。命题得证。

设 \(\varrho\) 是拓扑群 G 在希尔伯特空间 E 中的酉表示。若 K = C，记 \(\overline{\mathrm{E}}\) 为 E 的共轭空间（EVT, V, p. 6）。若 K = R，令 \(\overline{\mathrm{E}} = \mathrm{E}\)。共轭表示 \(\overline{\varrho}\) 是 G 在 \(\overline{\mathrm{E}}\) 中的表示，由 \(\overline{\varrho}(g) = \varrho(g)\)（对所有 \(g \in \mathrm{G}\)）定义。它是 G 的酉表示；按定义，\(\overline{\mathrm{E}}\) 的一个子空间是 \(\overline{\mathrm{E}}\) 的子表示，当且仅当它是 E 的子表示。

命题 2. --- 设 \(\varrho\) 是拓扑群 G 在希尔伯特空间 E 中的酉表示。记 u 为从 \(\overline{\mathrm{E}}\) 到 \(\mathrm{E}'\) 的等距同构，它将 \(x\) 关联到线性形式 \(y\mapsto \langle x|y\rangle\)。

\begin{enumerate}
\def\labelenumi{\alph{enumi})}
\item
  映射 u 是从共轭表示 \(\overline{\varrho}\) 到逆步表示 \(\breve{\varrho}\) 的同构；
\item
  给 E' 赋予通过 u 传递结构得到的希尔伯特空间结构；逆步表示 \(\breve{\varrho}\) 是 G 在 E' 中的酉表示。
\end{enumerate}

由 EVT, V, p. 15, th. 3 及其后注记，映射 \(u\) 是等距同构。

对 G 中每个 g、\(\overline{E}\) 中每个 x 以及 E 中每个 y，有

\[
\begin{array}{c} \langle y, (\breve {\varrho} (g) \circ u) (x) \rangle = \langle \varrho (g ^ {- 1}) y, u (x) \rangle \\ = \langle x | \varrho (g ^ {- 1}) y \rangle = \langle \varrho (g) x | y \rangle = \langle y, u (\overline {{\varrho}} (g) x) \rangle , \end{array}
\]

因此 \(\breve{\varrho}(g) \circ u = u \circ \overline{\varrho}(g)\)，这证明了 \(a\)；断言 \(b\) 立即得到。

推论. --- 设 \(\varrho\) 是拓扑群 G 在希尔伯特空间 E 中的酉表示。以下条件等价：

\begin{enumerate}
\def\labelenumi{(\roman{enumi})}
\item
  表示 \(\varrho\) 不可约；
\item
  逆步表示 \(\breve{\varrho}\) 不可约；
\item
  共轭表示 \(\overline{\varrho}\) 不可约。
\end{enumerate}

这由命题 2 以及命题前关于 \(\overline{E}\) 的子表示的注记推出。

命题 3. --- 设 \(\varrho\) 是拓扑群 G 在希尔伯特空间 E 中的酉表示。设 \(\pi\) 是 \(\varrho\) 的子表示，其表示空间为 F。F 在 E 中的正交补 F° 是 E 的子表示，且 E = F ⊕ F°，F° 与 E/F 同构。

空间 F° 是闭的。对于每个 \(x \in F^{\circ}\)、所有 \(g \in G\) 以及每个 \(y \in F\)，有 \(\langle \varrho(g)x \mid y \rangle = \langle x \mid \varrho(g^{-1})y \rangle = 0\)，因为 \(\varrho\) 是酉表示且 F 是子表示。因此 \(\varrho(g)x \in F^{\circ}\)。

E 到 E/F 的规范投影是 G-态射，所以它在 F° 上的限制是从 F° 到 E/F 的 G-态射；由 EVT, V, p. 13，它是等距同构。

我们也把 G 在 F° 中的表示记作 \(\pi^{\circ}\)。

命题 4. --- 设 \(\varrho\) 是拓扑群 G 在希尔伯特空间 E 中的酉表示。E 的闭子空间 F 是 \(\varrho\) 的子表示，当且仅当 E 上像为 F 的正交投影 p 是从 \(\varrho\) 到 \(\varrho\) 的 G-态射。

设 F 是 \(\varrho\) 的子表示。取 \(x \in E\)，写成 \(x = p(x) + y\)，其中 \(y \in F^{\circ}\)。对于 G 中每个 \(g\)，有

\[
\varrho (g) x = \varrho (g) (p (x)) + \varrho (g) y,
\]

并且由于 \(\varrho(g)(p(x))\) 属于 F 而 \(\varrho(g)y\) 属于 F°（命题 3），有 \(p(\varrho(g)x) = \varrho(g)(p(x))\)。因此 \(p\) 属于 \(\operatorname{Hom}_{\mathrm{G}}(\varrho, \varrho)\)。

反之，若 \(p \in \operatorname{Hom}_{\mathrm{G}}(\varrho, \varrho)\)，则 \(1_{E} - p\) 是 G-态射，因而 \(\mathrm{F} = \operatorname{Ker}(1_{\mathrm{E}} - p)\) 是 \(\varrho\) 的子表示。

\subsection{酉表示的希尔伯特直和与张量积}

设 G 是拓扑群。设 \((\varrho_{i})_{i\in\mathrm{I}}\) 是 G 在希尔伯特空间 \(E_{i}\) 中的一族酉表示。设 E 是各 \(E_{i}\) 的外部希尔伯特直和空间。对于 G 中每个 g 以及 E 中每个 \(x=(x_{i})_{i\in\mathrm{I}}\)，有

\[
\sum_ {i} \| \varrho_ {i} (g) x _ {i} \| ^ {2} = \sum_ {i} \| x _ {i} \| ^ {2} = \| x \| ^ {2},
\]

这证明了 \((\varrho_i(g)x_i)_{i\in \mathrm{I}}\) 属于 E 且与 \(x\) 有相同的范数。因此，元素 \(\varrho (g)\colon (x_{i})_{i\in \mathrm{I}}\mapsto (\varrho_{i}(g)x_{i})_{i\in \mathrm{I}}\) 是 \(\mathcal{L}(\mathrm{E})\) 中的酉元素。

引理 7. --- 映射 \(g \mapsto \varrho(g)\) 是 G 在 E 中的酉表示。

由 V, p. 380 的引理 4，只需证明对于 I 中每个 i 及 \(E_{i}\) 中每个 x，映射 \(g \mapsto \varrho_{i}(g)x\) 在 G 的单位元 e 处连续。该映射是连续映射 \(g \mapsto \varrho_{i}(g)x\) 与 \(E_{i}\) 到 E 的规范嵌入的复合，因而连续。

表示 \(\varrho\) 称为酉表示 \((\varrho_{i})_{i\in I}\) 的希尔伯特直和，记为 \(\varrho = \bigoplus \varrho_{i}\)。

设 G 和 H 是拓扑群。设 \(\varrho_{1}\)（分别为 \(\varrho_{2}\)）是 G（分别为 H）在希尔伯特空间 E \(_{1}\)（分别为 E \(_{2}\)）中的酉表示。设 E = E \(_{1}\) \(\widehat{\otimes}_{2}\) E \(_{2}\) 是 E \(_{1}\) 与 E \(_{2}\) 的希尔伯特张量积空间（EVT, V, p. 28, 定义 1）。对于 \((g,h)\in\mathrm{G}\times\mathrm{H}\)，令 \(\varrho(g,h)\) 为 E 的连续自同态 \(\varrho_{1}(g)\widehat{\otimes}_{2}\varrho_{2}(h)\)（TVS, V, p. 28）；无歧义时仍简记为 \(\varrho_{1}(g)\otimes\varrho_{2}(h)\)。

引理 8. --- a) 映射 \(\varrho\colon(g,h)\mapsto\varrho(g,h)\) 是 G × H 在 E 中的酉表示；

\begin{enumerate}
\def\labelenumi{\alph{enumi})}
\setcounter{enumi}{1}
\item
  对于任意标准正交基 \((e_i)_{i\in \mathrm{I}}\)，它是 \(\mathrm{E}_1\) 的基，从希尔伯特直和 \(\bigoplus_{i\in \mathrm{I}}\mathrm{E}_2\) 到 E 的映射 \((y_i)_{i\in \mathrm{I}}\mapsto \sum_{i\in \mathrm{I}}e_i\otimes y_i\) 是希尔伯特直和 \(\bigoplus_{i\in \mathrm{I}}\varrho_2\) 在 \(\operatorname {Res}_{\{e\} \times \mathrm{H}}^{\mathrm{G}\times \mathrm{H}}(\varrho)\) 中的等距 H-同构；
\item
  对于任意标准正交基 \((f_j)_{j\in \mathrm{J}}\)，它是 \(\mathrm{E}_2\) 的基，从希尔伯特直和 \(\bigoplus_{j\in \mathrm{J}}\mathrm{E}_1\) 到 E 的映射 \((x_{j})_{j\in \mathrm{J}}\mapsto \sum_{j\in \mathrm{J}}x_{j}\otimes f_{j}\) 是希尔伯特直和 \(\bigoplus_{j\in \mathrm{J}}\varrho_1\) 在 \(\operatorname {Res}_{\mathbf{G}\times \{e\}}^{\mathbf{G}\times \mathbf{H}}(\varrho)\) 中的等距 G-同构。
\end{enumerate}

映射 \((g,h)\mapsto\varrho(g,h)\) 是从 \(\mathrm{G}\times\mathrm{H}\) 到 \(\mathbf{GL}(\mathrm{E})\) 的同态（参见 EVT, V, p. 28, no. 2）。设 \((e_{i})_{i\in\mathrm{I}}\) 是 \(\mathrm{E}_{1}\) 的标准正交基。由 EVT, V, p. 29, prop. 3 和 cor. 2，映射

\[
u \colon (y _ {i}) \mapsto \sum_ {i \in I} e _ {i} \otimes y _ {i}
\]

是从希尔伯特直和 \(F_{2} = \bigoplus_{i \in I} E_{2}\) 到 E 上的等距映射。借助此等距映射，H 在 E 中的表示 \(\varrho_{\mathrm{H}} \colon h \mapsto \varrho(e, h)\) 被认同为表示族 \((\varrho_{2})_{i\in I}\) 在 \(F_{2}\) 中的直和表示。特别地，它是 H 在 E 中的酉表示（引理 7）。同样，同态 \(\varrho_{G}: g \mapsto \varrho(g,e)\) 是 G 在 E 中的酉表示。

取 \((g,h)\in\mathrm{G}\times\mathrm{H}\)。有 \(\varrho(g,h)=\varrho_{\mathrm{G}}(g)\circ\varrho_{\mathrm{H}}(h)\)，故 \(\varrho(g,h)\) 是酉的；此外，\(\varrho_{\mathrm{G}}(g)\) 与 \(\varrho_{\mathrm{H}}(h)\) 交换，所以由 V, p. 377 的引理 1，\(\varrho\) 是 \(G\times H\) 的酉表示。

最后，关于 \(\varrho\) 限制到子群 \(\{e\} \times \mathrm{H}\) 和 \(\mathrm{G} \times \{e\}\) 的断言已在上述论证过程中得到。

G×H 的表示 \((g,h)\mapsto\varrho_{1}(g)\otimes\varrho_{2}(h)\) 称为酉表示 \(\varrho_{1}\) 与 \(\varrho_{2}\) 的外张量积，记为 \(\varrho_{1}\boxtimes\varrho_{2}\)。

设 \(n \in \mathbf{N}\)，且 \((\mathrm{G}_i)_{1 \leqslant i \leqslant n}\) 是一族有限个拓扑群。设 \(\varrho_i\) 是 \(\mathrm{G}_i\) 在希尔伯特空间 \(\mathrm{E}_i\) 中的酉表示，对 \(1 \leqslant i \leqslant n\) 成立。类似地定义表示

\[
\varrho_ {1} \boxtimes \dots \boxtimes \varrho_ {n}
\]

\(G_{1} \times \cdots \times G_{n}\) 在希尔伯特空间 \(E = E_{1} \widehat{\otimes}_{2} \cdots \widehat{\otimes}_{2} E_{n}\) 中的表示（EVT, V, p. 27）。

若 \(G_{i} = G\)（\(1 \leqslant i \leqslant n\)），记 \(\Delta_{n}: G \to G^{n}\) 为由 \(g \mapsto (g, \ldots, g)\)（对所有 \(g \in G\)）定义的同态。我们用 \(\varrho_{1} \otimes \cdots \otimes \varrho_{n}\) 表示 G 的酉表示 \((\varrho_{1} \boxtimes \cdots \boxtimes \varrho_{n}) \circ \Delta_{n}\)。这称为酉表示 \(\varrho_{i}\) 的张量积。

对于每个置换 \(\sigma\)，即 \(\{1,\ldots,n\}\) 上的置换，规范同构

\[
\mathrm{E} _ {1} \widehat {\otimes} _ {2} \dots \widehat {\otimes} _ {2} \mathrm{E} _ {n} \rightarrow \mathrm{E} _ {\sigma (1)} \widehat {\otimes} _ {2} \dots \widehat {\otimes} _ {2} \mathrm{E} _ {\sigma (n)}
\]

（EVT, V, p. 28）是等距同构

\[
\varrho_ {1} \otimes \dots \otimes \varrho_ {n} \rightarrow \varrho_ {\sigma (1)} \otimes \dots \otimes \varrho_ {\sigma (n)}
\]

且是 G 的表示之间的同构。

若 \(\varrho_{i}=\varrho\) 对 \(1\leqslant i\leqslant n\) 成立，其中 \(\varrho\) 是 G 的酉表示，我们也用 \(\varrho^{\otimes n}\) 表示各 \(\varrho_{i}\) 的张量积表示，并称其为 \(n\) 次张量幂，即 \(\varrho\) 的张量幂。

\subsection{矩阵系数}

定义 5. --- 设 G 是拓扑群，\(\varrho\) 是 G 在希尔伯特空间 E 中的酉表示。设 \(x\) 和 \(y\) 是 E 中的元素。由 \(g \mapsto \langle x | \varrho(g)y \rangle\) 给出的从 G 到 K 的函数称为 \(\varrho\) 的矩阵系数，也称代表函数。若 \(x = y\)，则称其为对角矩阵系数。若 \(\varrho\) 是有限维的，则称其为有限维矩阵系数。

\(\varrho\) 的矩阵系数是 G 上的连续有界函数。我们用 \(\Upsilon(G)\)（分别用 \(\Theta(G)\)）表示 G 的复酉表示（分别为复有限维表示）的矩阵系数集合。

命题 5. — 设 G 是拓扑群。集合 \(\Theta(\mathrm{G})\) 和 \(\Upsilon(\mathrm{G})\) 是 \(\mathcal{C}_{\mathrm{b}}(\mathrm{G})\) 的含单位对合子代数。

常值函数 1 是 G 在 C 上的平凡表示的矩阵系数。设 \(\varrho\) 是 G 的酉表示，且 \(x\) 和 \(y\) 是 \(\varrho\) 的表示空间中的向量。对每个 \(\lambda \in \mathbf{C}\) 及每个 \(g \in \mathrm{G}\)，有 \(\lambda \langle x \mid \varrho(g)y \rangle = \langle x \mid \varrho(g)(\lambda y) \rangle\)。此外，有

\[
\langle x \mid \varrho (g) y \rangle = \langle \varrho (g ^ {- 1}) x \mid y \rangle = \overline {{\langle y \mid \varrho (g ^ {- 1}) x \rangle}}
\]

对所有 \(g \in G\) 成立。因此，集合 \(\Theta(G)\) 和 \(\Upsilon(G)\) 在标量乘法和共轭下稳定。

设 \(\varrho_{1}\) 和 \(\varrho_{2}\) 是 G 的酉表示；设 \((x_{1}, y_{1})\) 是 \(\varrho_{1}\) 的表示空间中的向量，\((x_{2}, y_{2})\) 是 \(\varrho_{2}\) 的表示空间中的向量。于是对每个 \(g \in G\)，有

\[
\begin{array}{c} \langle x _ {1}   |   \varrho_ {1} (g) y _ {1} \rangle + \langle x _ {2}   |   \varrho_ {2} (g) y _ {2} \rangle = \langle (x _ {1}, x _ {2})   |   (\varrho_ {1} \oplus \varrho_ {2}) (g) (y _ {1}, y _ {2}) \rangle , \\ \langle x _ {1}   |   \varrho_ {1} (g) y _ {1} \rangle \langle x _ {2}   |   \varrho_ {2} (g) y _ {2} \rangle = \langle x _ {1} \otimes x _ {2}   |   (\varrho_ {1} \otimes \varrho_ {2}) (g) (y _ {1} \otimes y _ {2}) \rangle , \end{array}
\]

这证明了 \(\Theta(G)\) 和 \(\Upsilon(G)\) 在加法和乘法下稳定。命题得证。

\subsection{舒尔引理}

本节中，希尔伯特空间均为复的。

命题 6（舒尔引理）。--- 设 \(\varrho\) 是拓扑群 G 在非零希尔伯特空间 E 中的酉表示。那么，\(\varrho\) 不可约，当且仅当 \(\mathrm{Hom}_{\mathrm{G}}(\varrho, \varrho)\) 等于 \(\mathbf{C} \cdot 1_{\mathrm{E}}\)。

\(\mathrm{Hom}_{\mathrm{G}}(\varrho, \varrho)\) 是 \(\mathcal{L}(\mathrm{E})\) 的含单位星子代数（V, p. 380 的引理 5）。对于 E 的每个子表示 F，像为 F 的正交投影 \(p\) 是星代数 \(\mathrm{Hom}_{\mathrm{G}}(\varrho, \varrho)\) 中的幂等元（V, p. 383 的命题 4）。若该代数等于 \(\mathbf{C} \cdot 1_{\mathrm{E}}\)，这意味着 \(p = 0\) 或 \(p = 1_{\mathrm{E}}\)，也即 F = 0 或 F = E。因此 \(\varrho\) 不可约。

反之，假设 \(\varrho\) 不可约。设 \(u\) 和 \(v\) 是星代数 \(\mathrm{Hom}_{\mathrm{G}}(\varrho, \varrho)\) 中满足 \(uv = 0\) 的交换元素。设 \(u\) 非零。于是 \(u\) 的核 F 定义 \(\varrho\) 的一个不同于 E 的子表示，故 F 必为 \(\{0\}\)；由于 F 包含 \(v\) 的像，推出 \(v = 0\)。因此由 I, p. 113 的命题 10，有 \(\mathrm{Hom}_{\mathrm{G}}(\varrho, \varrho) = \mathbf{C} \cdot 1_{\mathrm{E}}\)。

推论 1. --- 设 π 是拓扑群 G 在希尔伯特空间 E 中的不可约酉表示。设 u 是 E 上的闭部分算子。假设 u 的定义域稠密、dom(u) 在 π 下稳定，并且对所有 g ∈ G 有 u ∘ π(g) = π(g) ∘ u。那么 dom(u) = E，且 u 是恒等映射的标量倍数。

部分算子 \(u^{*} \circ u\) 是 E 上的正自伴部分算子（IV, p. 241 的命题 12）；对于 \(v = 1_{\mathrm{E}} + u^{*} \circ u\) 也同样如此，并且后者是单射的，因为 \(-1 \notin \operatorname{Sp}(u^{*} \circ u)\)（IV, p. 248 的命题 17）。有关系式 \(u^{*} \circ \pi(g) = \pi(g) \circ u^{*}\) 对所有 \(g \in G\) 成立（V, p. 381 的引理 6）；从而关系式 \(v \circ \pi(g) = \pi(g) \circ v\) 对所有 \(g \in G\) 成立。由于 \(v\) 是单射，故关系式 \(v^{-1} \circ \pi(g) = \pi(g) \circ v^{-1}\) 对所有 \(g \in G\) 成立。由于 \(v^{-1}\) 属于 \(\mathcal{L}(\mathrm{E})\)，所以由命题 6，存在 \(\lambda \in \mathbf{C}\) 使得 \(v^{-1} = \lambda 1_{\mathrm{E}}\)。必有 \(\lambda \neq 0\)，这意味着 \(\mathrm{E} = \mathrm{Im}(v^{-1}) = \mathrm{dom}(v) \subset \mathrm{dom}(u)\)，因而 \(\mathrm{dom}(u) = \mathrm{E}\)。由于 \(u\) 是闭的，有 \(u \in \mathcal{L}(\mathrm{E})\)，于是 \(u \in \mathrm{Hom}_{\mathrm{G}}(\pi, \pi)\)，并且由命题 6，\(u\) 是一个伸缩变换。

推论 2. --- 设 \(\varrho\) 和 \(\pi\) 是拓扑群 G 在希尔伯特空间 E 和 F 中的不可约酉表示。空间 \(\mathrm{Hom}_{\mathrm{G}}(\varrho, \pi)\) 的维数在 \(\varrho\) 与 \(\pi\) 同构时为 1，否则为零。特别地，若 \(\varrho\) 与 \(\pi\) 同构，则从 \(\varrho\) 到 \(\pi\) 的每个非零 G-态射都是同构。

假设存在非零 G-态射 \(u\)，它从 \(\varrho\) 到 \(\pi\)。线性映射 \(u^{*} \circ u\) 是 \(\mathrm{Hom}_{\mathrm{G}}(\varrho, \varrho)\) 的元素；因此存在复数 \(\lambda\)，使得 \(u^{*} \circ u = \lambda \cdot 1_{\mathrm{E}}\)（命题 6）。

于是有

\[
\langle u (x) \mid u (y) \rangle = \langle x \mid u ^ {*} u (y) \rangle = \lambda \langle x \mid y \rangle
\]

对 E 中所有 \(x\) 和 \(y\) 成立。特别地，\(\lambda \neq 0\)，因为 \(u\) 非零。由于 \(\| u(x)\| = |\lambda |^{1 / 2}\| x\|\) 对所有 \(x\in \mathrm{E}\) 成立，线性映射 \(u\) 是单射，且 \(u\) 的像在 F 中闭（I, p. 107 的引理 8）；于是它是不可约表示 \(\pi\) 的一个非零子表示，从而 \(u\) 是满射。因此，\(u\) 是从 \(\varrho\) 到 \(\pi\) 的同构。于是映射 \(v\mapsto u^{*}\circ v\) 是从空间 \(\mathrm{Hom}_{\mathrm{G}}(\varrho ,\pi)\) 到空间 \(\mathrm{Hom}_{\mathrm{G}}(\varrho ,\varrho)\) 的同构，而后一个空间的维数为 1（命题 6）。

推论 3. --- 设 \(\varrho\) 是拓扑群 G 在希尔伯特空间 E 中的酉表示。表示 \(\varrho\) 不可约，当且仅当空间 \(\mathrm{Sesq}_{\mathrm{G}}(\varrho, \varrho)\) 的维数为 1。此时，该空间由 E 上的内积生成。

根据 V, p. 381 的命题 1，这由命题 6 得出。

推论 4. --- 设 \(\varrho_{1}\) 和 \(\varrho_{2}\) 是拓扑群 G 在希尔伯特空间 E \(_{1}\) 和 E \(_{2}\) 中的非同构不可约酉表示。空间 Sesq \(_{G}\)（\(\varrho_{1}, \varrho_{2}\)）为零。

由 V, p. 381 的命题 1，这由推论 2 得出。

推论 5. --- 设 \(\varrho\) 是拓扑群 G 在希尔伯特空间 E 中的酉表示。设 \(\pi\) 是 G 在希尔伯特空间 F 中的不可约酉表示。

\begin{enumerate}
\def\labelenumi{\alph{enumi})}
\item
  从 \(\pi\) 到 \(\varrho\) 的每个非零 G-态射 u，都存在 \(\lambda \in \mathbf{R}_{+}^{*}\)，使得 \(\lambda u\) 是等距的；
\item
  从 π 到 ρ 的每个 G-态射 u 都有闭像；
\item
  从 \(\varrho\) 到 \(\pi\) 的每个非零 G-态射 v 都是满射。
\end{enumerate}

证明断言 \(a\)，它蕴含断言 \(b\)（I, p. 107 的引理 8）。由于态射 \(u\) 非零，它是单射（V, p. 378 的引理 2）。由公式 \(q(x,y) = \langle u(x)|u(y)\rangle\)，其中 \(x\) 和 \(y\) 属于 F，定义了 F 上的连续半双线性形式。由于 \(u\) 是 G-态射，有 \(q(\pi(g)x,\pi(g)y) = q(x,y)\) 对所有 \(g\)（属于 G）及 \((x,y)\in\mathrm{F}\times\mathrm{F}\) 成立，故 \(q\) 是 G-不变的。由推论 3，存在 \(\alpha\in\mathbf{R}_{+}^{*}\)，使得

\[
q (x, y) = \langle u (x) | u (y) \rangle = \alpha \langle x | y \rangle
\]

对所有 \((x,y)\in\mathrm{F}\times\mathrm{F}\) 成立，因而 \(\alpha^{-1/2}u\) 是等距的。

最后证明断言 \(c\)。由于 \(\pi\) 不可约，\(v\) 的像在 F 中稠密（V, p. 378 的引理 2），故伴随 \(v^*\) 是单射（TVS, V, p. 41 的命题 4）。由 \(b\)，\(v^*\) 的像 H 是 E 的闭子空间；因此它是 \(\varrho\) 的子表示，而由限制到子空间得到的 \(v^*\) 是从 F 到 H 的 G-同构。映射 \(w\) 是 \(v\) 在 H 上的限制，它定义了从 H 到 F 的 G-态射；由于其核是 H 与 \(\text{Ker}(v) = \text{H}^\circ\) 的交，它是单射（同上）。因此 \(w\) 是同构（推论 2）；特别地，\(v\) 是满射。

推论 6. --- 设 \(G_{1}\) 和 \(G_{2}\) 是拓扑群。设 \(\varrho_{1}\) 是 \(G_{1}\) 在希尔伯特空间 \(E_{1}\) 中的不可约酉表示，\(\varrho_{2}\) 是 \(G_{2}\) 在希尔伯特空间 \(E_{2}\) 中的不可约酉表示。外张量积 \(\varrho_{1} \boxtimes \varrho_{2}\) 是 \(G_{1} \times G_{2}\) 在 \(E_{1} \widehat{\otimes}_{2} E_{2}\) 中的不可约酉表示。

由 V, p. 384 的引理 8，外张量积 \(\varrho = \varrho_{1} \otimes \varrho_{2}\) 是 \(G = G_{1} \times G_{2}\) 在空间 \(E = E_{1} \widehat{\otimes}_{2} E_{2}\) 中的酉表示。

设 \(q\) 是 E 上的 \((\mathrm{G}_1 \times \mathrm{G}_2)\)-不变半双线性形式。对于每个 \((x_1, y_1) \in \mathrm{E}_1^2\)，映射 \((x_2, y_2) \mapsto q(x_1 \otimes x_2, y_1 \otimes y_2)\) 属于 \(\operatorname{Sesq}_{\mathrm{G}_2}(\varrho_2, \varrho_2)\)。记 \(b(x_1, y_1)\) 为满足下式的唯一复数：

\[
q (x _ {1} \otimes x _ {2}, y _ {1} \otimes y _ {2}) = b (x _ {1}, y _ {1}) \left\langle x _ {2} \mid y _ {2} \right\rangle
\]

对所有 \((x_{2},y_{2})\in\mathrm{E}_{2}^{2}\) 成立（推论 3）。取 \(\varepsilon\in E_{2}\) 使其范数为 1，于是有 \(b(x_{1},y_{1})=q(x_{1}\otimes\varepsilon,y_{1}\otimes\varepsilon)\) 对所有 \((x_{1},y_{1})\in\mathrm{E}_{1}^{2}\) 成立。该公式说明映射 b 是 \(E_{1}\) 上的半双线性映射。此外

\[
\| b (x _ {1}, y _ {1}) \| \leqslant \| q \| \| x _ {1} \otimes \varepsilon \| \| y _ {1} \otimes \varepsilon \| = \| q \| \| x _ {1} \| \| y _ {1} \|
\]

对所有 \((x_{1},y_{1})\in\mathrm{E}_{1}^{2}\) 成立（EVT, V, p. 26, 公式（5）），因此 b 是连续的。

设 \(g \in G\)，且 \((x_{1}, y_{1}) \in \mathrm{E}_{1}^{2}\)。由于 \(\varrho_{2}\) 是酉的且 q 不变，得到

\[
b \left(\varrho_ {1} (g) x _ {1}, \varrho_ {1} (g) y _ {1}\right) = q \left(x _ {1} \otimes \varrho_ {2} \left(g ^ {- 1}\right) \varepsilon , y _ {1} \otimes \varrho_ {2} \left(g ^ {- 1}\right) \varepsilon\right) = b \left(x _ {1}, y _ {1}\right),
\]

因此 \(b \in \mathrm{Sesq}_{\mathrm{G}_1}(\varrho_1, \varrho_1)\)。于是存在唯一的 \(\lambda \in \mathbf{C}\)，使得 \(b(x_1, y_1) = \lambda \langle x_1 | y_1 \rangle\) 对所有 \((x_1, y_1) \in \mathrm{E}_1^2\) 成立（推论 3），即

\[
q (x _ {1} \otimes x _ {2}, y _ {1} \otimes y _ {2}) = \lambda \langle x _ {1} | y _ {1} \rangle \langle x _ {2} | y _ {2} \rangle
\]

对所有 \((x_{1},y_{1},x_{2},y_{2})\in\mathrm{E}_{1}^{2}\times\mathrm{E}_{2}^{2}\) 成立。由此可推出空间 \(\operatorname{Sesq}_{\mathrm{G}_{1}\times\mathrm{G}_{2}}(\varrho,\varrho)\) 的维数为 1，这意味着表示 \(\varrho=\varrho_{1}\boxtimes\varrho_{2}\) 不可约（同上）。

回忆一下，U 表示模为 1 的复数所成的群。

推论 7. --- 设 \(\pi\) 是拓扑群 G 在希尔伯特空间 E 中的不可约酉表示。存在从 G 的中心 C 到 U 的连续同态 \(\chi\)，满足 \(\pi(z) = \chi(z) \cdot 1_{\mathrm{E}}\) 对所有 \(z \in \mathrm{C}\) 成立。特别地，若 G 是交换群，则 \(\dim(\mathrm{E}) = 1\)。

对于 \(z \in \mathrm{C}\)，映射 \(\pi(z)\) 属于 \(\mathrm{Hom}_{\mathrm{G}}(\pi, \pi)\)；因此它具有 \(\chi(z) \cdot 1_{\mathrm{E}}\) 的形式，其中 \(\chi(z)\) 是某个复数（命题 6）。由于 \(\pi(z)\) 是酉映射，有 \(|\chi(z)| = 1\)。此外，由于 \(\pi\) 是同态，映射 \(z \mapsto \chi(z)\) 是同态。取 E 中 \(v \neq 0\)；从 C 到 U 的映射 \(z \mapsto \chi(z)v = \pi(z)v\) 是连续的，因此同态 \(\chi\) 是连续的。

最后，若 G 是交换群，则 C = G，所以对 G 中所有 g，有 \(\pi(g) = \chi(g) \cdot 1_{\mathrm{E}}\)；于是 E 的每个一维子空间都是 \(\pi\) 的子表示，而由于 \(\pi\) 不可约，空间 E 必为一维。

定义 6. --- 设 π 是拓扑群 G 在希尔伯特空间 E 中的不可约酉表示。满足对 C 中所有 z 有 π(z) = χ(z) · \(1_{\mathrm{E}}\) 的从 G 的中心 C 到 U 的同态 χ，称为 π 的中心特征标。

注. --- 设 \(\pi\)（分别为 \(\varrho\)）是拓扑群 G（分别为 H）在希尔伯特空间 E（分别为 F）中的不可约酉表示。记 \(\chi\)（分别为 \(\eta\)）为 \(\pi\)（分别为 \(\varrho\)）的中心特征标。表示 \(\overline{\pi}\) 的中心特征标是 \(\overline{\chi}\)；而 G × H 的不可约酉表示 \(\pi \boxtimes \varrho\)（推论 6）的中心特征标，是 G × H 上的特征标 \(\chi \boxtimes \eta\)：\((g, h) \mapsto \chi(g)\eta(h)\)。

\subsection{半单性}

定义 7. --- 设 \(\varrho\) 是拓扑群 G 在希尔伯特空间 E 中的酉表示；称 \(\varrho\) 为半单的，当且仅当存在不可约子表示族 \((\mathrm{F}_i)_{i\in \mathrm{I}}\)，它们是 \(\varrho\) 的子表示，且 E 是子空间 \(\mathrm{F}_i\) 的希尔伯特直和。

有时也说半单酉表示具有离散分解，或称其可离散分解。

如果 \((\varrho_{i})_{i\in\mathrm{I}}\) 是拓扑群 G 的半单酉表示族，那么表示 \(\varrho_{i}\) 的希尔伯特直和是半单的。

命题 7. --- 设 G 是拓扑群，\(\varrho\) 是 G 在复希尔伯特空间 E 中的酉表示。表示 \(\varrho\) 半单，当且仅当 \(\varrho\) 的每个非零子表示都包含不可约子表示。

假设 \(\varrho\) 半单。设 \((\mathrm{F}_i)_{i\in \mathrm{I}}\) 是 \(\varrho\) 的不可约子表示族，使得 E 是子空间 \(\mathrm{F}_i\) 的希尔伯特直和。设 F 是 E 的一个非零子表示。存在 \(i\in \mathrm{I}\)，使得像为 \(\mathrm{F}_i\) 的正交投影在 F 上的限制非零。该正交投影于是通过限制到子空间定义了非零 G-态射 \(p_i\)，从 F 到 \(\mathrm{F}_i\)。由舒尔引理，该 G-态射 \(p_i\) 是满射（V, p. 388 的推论 5）。F 中 \(p_i\) 的核的正交补是 F 的一个与 \(\mathrm{F}_i\) 同构的子表示（V, p. 383 的命题 3），因而不可约。

证明逆命题。令 \(\mathcal{F}\) 为 E 的闭子空间 F 所成的集合，其中 F 在 G 下稳定，且 \(\varrho\) 在 F 中诱导的子表示不可约。令 \(\mathcal{O}\) 为所有满足 \(\mathcal{G}\) 是 \(\mathcal{F}\) 的子集且其成员两两正交的 \(\mathcal{G}\) 所成的集合。

集合 \(\mathcal{O}\) 具有有限特征（E, III, p. 34, 定义 2）。于是令 \(\mathcal{G}\) 为 \(\mathcal{O}\) 的一个极大元（E, III, p. 35, 定理 1）。令 E \(_1\) 为 E 中的子空间，它是 \(\mathcal{G}\) 中不可约子表示 F 的希尔伯特直和。若 E \(_1\) ≠ E，则 E \(_1\) 的正交补非零，而根据假设，G 在 E \(_1^{\circ}\) 中的子表示包含一个不可约子表示。其表示空间 F 与 \(\mathcal{G}\) 中的元素正交，于是 \(\mathcal{G} \cup \{F\} \in \mathcal{O}\)；这与 \(\mathcal{G}\) 的极大性矛盾。因此 E = E \(_1\)，证明完毕。

推论 1. --- 设 G 是拓扑群，\(\varrho\) 是 G 在复希尔伯特空间 E 中的半单酉表示。\(\varrho\) 的每个子表示（以及 \(\varrho\) 的每个商表示）都是半单的。

如果 \(\varrho_{1}\) 是 \(\varrho\) 的子表示，那么 \(\varrho_{1}\) 的每个非零子表示都包含不可约子表示（命题 7），故 \(\varrho_{1}\) 是半单的（同上）。

由 V, p. 383 的命题 3，\(\varrho\) 的每个商表示都与 \(\varrho\) 的一个子表示同构；因此它是半单的。

推论 2. --- 拓扑群 G 的每个有限维酉表示 \(\varrho\) 都是半单的。

这由 V, p. 379 的命题 7 和引理 3 得出。

推论 3. --- 设 \(\varrho_{1}\) 和 \(\varrho_{2}\) 是拓扑群 G 的有限维酉表示。表示 \(\varrho_{1}\) 与 \(\varrho_{2}\) 同构，当且仅当 \(\chi_{\varrho_{1}} = \chi_{\varrho_{2}}\)。

这由推论 2 以及 A, VIII, p. 389 的命题 1, b) 得出。

\subsection{酉表示的类}

引理 9. --- 设 E 是 K 上的希尔伯特空间，F ⊂ E 是 E 的稠密向量子空间。那么 E 的希尔伯特维数小于或等于 F 的维数。

若 E 的希尔伯特维数有限，则它等于 E 的维数。设 E 的希尔伯特维数无限，则子空间 F 的维数也是无限的。设 B 是 E 的标准正交基，B' 是 F 的基。对于每个 \(x \in B\)，存在 \(f(x) \in B'\)，使得 \(\langle x \mid f(x) \rangle \neq 0\)，否则将有 \(x \in F^{\circ} = \{0\}\)。于是定义映射 \(f: B \to B'\)。

对于每个 \(y \in B'\)，集合 \(f^{-1}(y)\) 包含于满足 \(x \in B\) 且 \(\langle x | y \rangle \neq 0\) 的集合中，因而是可数的（EVT, V, p. 21, 命题 4）。由 E, III, p. 50, 命题 4，得到 \(\text{Card(B)} = \text{Card}(f(B)) \leqslant \text{Card(B')}\)。

记 \(Is_{G}(\pi_{1},\pi_{2})\) 为关系

“G 是拓扑群，且 \(\pi_{1}\) 和 \(\pi_{2}\) 是

K 上希尔伯特空间中的酉表示，并且二者

同构。”关于 \(\pi_{1}\) 和 \(\pi_{2}\)，这是一个等价关系。对于 K 上希尔伯特空间中 G 的任意酉表示 \(\pi\)，记 \(\mathrm{cl}(\pi)\) 为 \(\pi\) 的等价类（参见 E, II, p. 47）；因此它是 G 的一个与 \(\pi\) 同构的酉表示；称 \(\mathrm{cl}(\pi)\) 为 \(\pi\) 的类。K 上希尔伯特空间中的酉表示 \(\pi_{1}\) 和 \(\pi_{2}\) 同构，当且仅当 \(\mathrm{cl}(\pi_{1}) = \mathrm{cl}(\pi_{2})\)。

设 G 是拓扑群，c 是一个基数。关系

“λ 是 G 在

希尔伯特维数 \(\leqslant\mathfrak{c}\) 的复希尔伯特空间中的酉表示的类”

是关于 \(\lambda\) 的集化关系（E, II, p. 3）。事实上，每个维数 \(\leqslant \mathfrak{c}\) 的 K 上希尔伯特空间都等距同构于 \(\ell^2 (\mathfrak{c})\) 的一个希尔伯特子空间（EVT, V, p. 23, cor. 2），于是断言由 E, II, p. 47 得出。

设 \(\pi\) 是 G 在希尔伯特空间 E 中的不可约酉表示。取 E 中的非零元素 \(x\)。由于 \(\pi\) 不可约，向量 \(x\) 是 \(\pi\) 的循环向量，这意味着 E 的希尔伯特维数 \(\leqslant\) Card(G)（将引理 9 应用于由 \(\pi(g)x\)（\(g \in G\)）所生成的稠密子空间）。因此，G 在 K 上希尔伯特空间中的不可约酉表示的各类，属于 G 在 K 上希尔伯特维数 \(\leqslant\) Card(G) 的希尔伯特空间中的酉表示的类所成的集合；所以它们构成一个集合。

定义 8. --- 我们用 \(\widehat{\mathbf{G}}\) 表示 \(\mathbf{G}\) 在复希尔伯特空间中的不可约酉表示的类所成的集合。称 \(\widehat{\mathbf{G}}\) 为 \(\mathbf{G}\) 的酉对偶。

因此，对于 G 在复希尔伯特空间中的每个不可约酉表示 \(\pi\)，都有 \(\operatorname{cl}(\pi) \in \widehat{\mathbf{G}}\)。

备注. --- 1) 假设 G 是交换群。由 V, p. 390 的推论 7，G 在复希尔伯特空间中的每个不可约酉表示都是一维的。如果 G 局部紧，则集合 \(\widehat{G}\) 可与 G 的酉特征标集合等同（II, p. 201 的定义 1），所以记号 \(\widehat{G}\) 与 II, p. 201 的定义 2 中引入的记号相容。

\begin{enumerate}
\def\labelenumi{\arabic{enumi})}
\setcounter{enumi}{1}
\item
  如果 G 是有限群，则集合 \(\widehat{\mathbf{G}}\) 与 \(\mathbf{C}[\mathrm{G}]\)-单模的类所成的集合之间存在双射（A, VIII, p. 47），后者在 A, VIII, p. 396 中也记作 \(\widehat{\mathbf{G}}\)。
\item
  对于 \(\pi \in \widehat{\mathbf{G}}\)，我们把 \(\overline{\pi}\) 与 \(\widehat{\mathbf{G}}\) 中 \(\pi\) 的共轭表示的类等同。
\item
  如果 \(\pi\) 是 G 在有限维复希尔伯特空间中的不可约酉表示，则其特征标 \(\chi_{\pi}\) 只取决于 \(\pi\) 的类。因此可以谈论 G 的有限维复不可约酉表示的特征标集合。
\end{enumerate}

\subsection{同型分量}

定义 9. --- 设 G 是拓扑群，\(\pi\) 是 G 的连续不可约表示。设 \(\varrho\) 是 G 在分离局部凸空间 E 中的连续表示。称 \(\pi\)-同型分量（对应 \(\varrho\)）为 \(\varrho\) 的所有与 \(\pi\) 同构的子表示的表示空间之和在 E 中的闭包。该子空间记为 \(\mathrm{M}_{\pi}(\varrho)\)。

空间 \(\mathrm{M}_{\pi}(\varrho)\) 是 E 的闭子空间，并定义 \(\varrho\) 的一个子表示。该空间只取决于 \(\pi\) 在 \(\widehat{G}\) 中的类。

命题 8. — 设 G 是拓扑群，\(\varrho\) 是 G 在复希尔伯特空间 E 中的酉表示。

\begin{enumerate}
\def\labelenumi{\alph{enumi})}
\item
  设 \(\pi_1\) 和 \(\pi_2\) 是 G 的非同构不可约酉表示。子空间 \(\mathrm{M}_{\pi_1}(\varrho)\) 和 \(\mathrm{M}_{\pi_2}(\varrho)\) 正交；
\item
  设 \(\varrho'\) 是 G 的酉表示，且 \(u\) 是从 \(\varrho\) 到 \(\varrho'\) 的 G-态射。对于 G 的每个不可约酉表示 \(\pi\)，有 \(u(\mathrm{M}_{\pi}(\varrho)) \subset \mathrm{M}_{\pi}(\varrho')\)。
\end{enumerate}

设 \(E_{1}\) 和 \(E_{2}\) 是 E 的子空间，分别定义与 \(\pi_{1}\) 和 \(\pi_{2}\) 同构的子表示。E 上像为 \(E_{2}\) 的正交投影通过限制到子空间定义了 \(\mathrm{Hom}_{\mathrm{G}}(\pi_{1},\pi_{2})\) 的一个元素；该元素为零（V, p. 387 的推论 2），这证明了 \(E_{2}\) 与 \(E_{1}\) 正交。断言 \(a\) 得证。

对于断言 \(b\)，注意到根据定义，\(\mathrm{M}_{\pi}(\varrho)\) 是 E 中由如下元素生成的空间的闭包：元素 \(x \in \mathrm{E}\) 属于闭子空间 \(\mathrm{F} \subset \mathrm{E}\)，该子空间在 \(\varrho\) 下稳定，且其子表示 \(\varrho_{\mathrm{F}}\) 是 \(\varrho\) 的子表示并与 \(\pi\) 同构。因此只需证明在这种情形下 \(u(x) \in \mathrm{M}_{\pi}(\varrho')\)。令 \(\mathrm{H} = u(\mathrm{F})\)；它是 \(\varrho'\) 的表示空间的闭子空间（V, p. 388 的推论 5），在 \(\varrho'\) 下稳定，而由限制到子空间得到的 u 是从 F 到 H 的满射 G-态射。如果 H 非零，则由 V, p. 387 的推论 2，该 G-态射是同构。因此 G 在 H 上的表示与 \(\pi\) 同构，从而 \(u(x)\) 属于 \(\mathrm{M}_{\pi}(\varrho')\)。

对于任意向量空间 H 及 H 的子空间族 \((\mathrm{H}_{i})_{i\in\mathrm{I}}\)，称这些空间 \((\mathrm{H}_{i})_{i\in\mathrm{I}}\) 为直和，如果族 \((\mathrm{H}_{i})_{i\in\mathrm{I}}\) 满足 A, II, p. 18 的命题 11 的等价条件。

命题 9. --- 设 G 是拓扑群，\(\varrho\) 是 G 在复希尔伯特空间 E 中的酉表示。设 \(\pi\) 是 G 在复希尔伯特空间 \(\mathrm{E}_{\pi}\) 中的不可约酉表示。记 \(v\) 为从 \(\mathrm{Hom}_{\mathrm{G}}(\pi, \varrho) \otimes \mathrm{E}_{\pi}\) 到 E 的线性映射，满足 \(v(u \otimes x) = u(x)\) 对所有 \((u, x) \in \mathrm{Hom}_{\mathrm{G}}(\pi, \varrho) \times \mathrm{E}_{\pi}\) 成立。

线性映射 \(v\) 是单射，其像是 \(\varrho\) 的所有与 \(\pi\) 同构的子表示的表示空间之和。特别地，\(v\) 的像在 \(M_{\pi}(\varrho)\) 中稠密。

V, p. 388 的推论 5 表明，v 的像是 \(\varrho\) 的所有与 \(\pi\) 同构的子表示的表示空间之和。

证明 \(v\) 是单射。设 \((u_i)_{i \in I}\) 是向量空间 Hom \(_G(\pi, \varrho)\) 的一组基。对于 \(i \in I\)，记 F \(_i\) 为 \(u_i\) 的像，记 \(\widetilde{u}_i \colon E_\pi \to F_i\) 为由 \(u_i\) 通过限制到子空间得到的 G-同构。

首先对 I 的有限子集 J 的基数作归纳，证明子空间 \((\mathrm{F}_{i})_{i\in\mathrm{J}}\) 是直和。

J 为空的情形是显然的。设 J 非空，并且所需性质对 I 中基数至多为 \(\text{Card(J)} - 1\) 的子集成立。

取 J 中的一个固定元素 \(j\)。归纳假设说明子空间 \(\mathrm{F}_{i}\)（\(i\in\mathrm{J}-\{j\}\)）构成直和。令 \(\mathrm{F}^{\prime}\) 为它们的和；现在只需证明 \(\mathrm{F}_{j}\cap\mathrm{F}^{\prime}=\{0\}\)。

假设 \(F_{j}\) 与 \(F'\) 的交不为 0；该交是 \(F_{j}\) 的一个子表示，而由于后者不可约，故有 \(F_{j} \cap F' = F_{j}\)，从而 \(F_{j} \subset F'\)。

对于 \(i \in \mathrm{J} - \{j\}\)，记 \(\operatorname{pr}_i: \mathrm{F}' \to \mathrm{F}_i\) 为投影，记 \(\iota_i: \mathrm{F}_i \to \mathrm{F}'\) 为包含映射。我们有规范同构

\[
\operatorname{Hom} _ {\mathrm{G}} \left(\mathrm{F} _ {j}, \mathrm{F} ^ {\prime}\right)\rightarrow \bigoplus_ {i \in \mathrm{J} - \{j \}} \operatorname{Hom} _ {\mathrm{G}} \left(\mathrm{F} _ {j}, \mathrm{F} _ {i}\right)
\]

（公式（1）, p. 377）使得 \(u \colon \mathrm{F}_j \to \mathrm{F}'\) 的像为族 \((\mathrm{pr}_i \circ u)_{i \in \mathrm{J} - \{j\}}\)（A, II, p. 13, cor. 1）。V, p. 387 的推论 2 表明，非零的 \(\widetilde{u}_i \circ \widetilde{u}_j^{-1}\) 是空间 \(\mathrm{Hom}_{\mathrm{G}}(\mathrm{F}_j, \mathrm{F}_i)\) 的一组基。

因此，G-态射族 \((\iota_{i} \circ \widetilde{u}_{i} \circ \widetilde{u}_{j}^{-1})_{i \in \mathrm{J} - \{j\}}\) 是 \(\mathrm{Hom}_{\mathrm{G}}(\mathrm{F}_{j}, \mathrm{F}')\) 的一组基。

记 \(\iota\) 为 \(F_{j}\) 到 \(F'\) 的包含映射；它是 \(\mathrm{Hom}_{\mathrm{G}}(\mathrm{F}_{j}, \mathrm{F}')\) 的一个元素，因而是各 G-态射 \(\iota_{i} \circ \widetilde{u}_{i} \circ \widetilde{u}_{j}^{-1}\)（\(i \in \mathrm{J} - \{j\}\)）的线性组合。于是 \(u_{j}\) 是各映射 \(u_{i}\)（\(i \neq j\)）的线性组合，这与族 \((u_{i})_{i \in \mathrm{I}}\) 的线性无关性矛盾。因此断言由归纳得到证明。

现在证明 v 是单射。设 w 是 \(\mathrm{Hom}_{\mathrm{G}}(\pi,\varrho)\otimes\mathrm{E}_{\pi}\) 的一个元素。存在唯一的有限支集族 \((x_{i})_{i\in\mathrm{I}}\)，其元素属于 \(E_{\pi}\)，使得

\[
w = \sum_ {i \in \mathrm{I}} u _ {i} \otimes x _ {i}
\]

（A, II, p. 62, cor. 1），于是有

\[
v (w) = \sum_ {i \in I} u _ {i} (x _ {i}).
\]

由前述内容，条件 \(v(w) = 0\) 因而蕴含 \(u_{i}(x_{i}) = 0\) 对所有 \(i \in I\) 成立，因而对所有 i 有 \(x_{i} = 0\)，这是因为 \(u_{i}\) 是单射，从而 w = 0。

命题 10. --- 设 G 是拓扑群。设 \(\pi\) 是 G 在复希尔伯特空间 E \(_{\pi}\) 中的不可约酉表示，\(\varrho\) 是 G 在复希尔伯特空间 E 中的酉表示。

存在 E 的闭不变子空间族 \((\mathrm{E}_{i})_{i\in\mathrm{I}}\)，使得 \(\varrho\) 在 \(E_{i}\) 中的子表示与 \(\pi\) 同构，对所有 \(i\in I\) 成立，并且 \(\mathrm{M}_{\pi}(\varrho)\) 是子空间 \(E_{i}\) 的希尔伯特直和。此外，满足这些性质的族 \((\mathrm{E}_{i})_{i\in\mathrm{I}}\) 的指标集 I 的基数与所选族无关。

证明满足所述条件的族存在。令 \(\mathcal{O}\) 为 \(\mathrm{Hom}_{\mathrm{G}}(\pi, \varrho)\setminus\{0\}\) 的子集 C 所成的集合，使得子空间 \(u(\mathrm{E}_{\pi})\)（\(u \in \mathrm{C}\)）两两正交。集合 \(\mathcal{O}\) 按包含关系排序。它非空，因为空集属于其中；它具有有限特征（E, III, p. 34, 定义 2），因为 C 属于 \(\mathcal{O}\) 当且仅当 C 的每个二元子集属于 \(\mathcal{O}\)。由 E, III, p. 35, 定理 1，\(\mathcal{O}\) 存在极大元 C。令 F 为由 \(u(\mathrm{E}_{\pi})\)（\(u \in \mathrm{C}\)）所成子空间生成的子空间的闭包；它是 \(u(\mathrm{E}_{\pi})\)（\(u \in C\)）所成空间的希尔伯特直和。我们将证明 \(\mathrm{F} = \mathrm{M}_{\pi}(\varrho)\)，从而证明族 \((u(\mathrm{E}_{\pi}))_{u \in \mathrm{C}}\) 具有所需性质。

按定义，\(u(\mathrm{E}_{\pi}) \subset \mathrm{M}_{\pi}(\varrho)\) 对所有 \(u \in \mathrm{C}\) 成立，故 \(\mathrm{F}\) 包含于 \(\mathrm{M}_{\pi}(\varrho)\)。为证明反向包含，只需证明：若 \(v\) 是从 \(\pi\) 到 \(\varrho\) 的 G-态射，则其像包含于 \(\mathrm{F}\)。令 \(p\) 为 \(\mathrm{E}\) 上像为 \(\mathrm{F}^{\circ}\) 的正交投影；由于 \(\mathrm{F}\) 是 \(\mathrm{E}\) 的子表示，它是 G-态射（V, p. 383 的命题 4）。G-态射 \(p \circ v\) 的像与 \(\mathrm{F}\) 正交；因此该像为零（否则 \(\mathrm{C} \cup \{p \circ v\} \in \mathcal{O}\)，这与 \(\mathrm{C}\) 的极大性矛盾），从而 \(v\) 的像包含于 \(\mathrm{F}\)。

现在设 \((\mathrm{E}_{i})_{i\in\mathrm{I}}\) 和 \((\mathrm{F}_{j})_{j\in\mathrm{J}}\) 是 E 的两两正交闭不变子空间族，使得 G 在 \(E_{i}\)（分别在 \(F_{j}\)）中的子表示与 \(\pi\) 同构，对所有 \(i\in I\)（分别对每个 \(j\in J\)）成立，并且 \(\mathrm{M}_{\pi}(\varrho)\) 既是族 \((\mathrm{E}_{i})_{i\in\mathrm{I}}\) 的希尔伯特直和，也是族 \((\mathrm{F}_{j})_{j\in\mathrm{J}}\) 的希尔伯特直和。

若 I 有限，则

\[
\dim \operatorname{Hom} _ {\mathrm{G}} (\pi , \mathrm{M} _ {\pi} (\varrho)) = \dim \operatorname{Hom} _ {\mathrm{G}} \Bigl (\mathrm{E} _ {\pi}, \bigoplus_ {i \in \mathrm{I}} \mathrm{E} _ {i} \Bigr) = \operatorname{Card} (\mathrm{I})
\]

（V, p. 377 的公式（1）和 V, p. 387 的推论 2）。于是对于 J 的每个有限子集 L，有

\[
\begin{array}{c} \operatorname{Card} (\mathrm{L}) = \dim \operatorname{Hom} _ {\mathrm{G}} \Bigl (\mathrm{E} _ {\pi}, \bigoplus_ {j \in \mathrm{L}} \mathrm{F} _ {j} \Bigr) \\ \leqslant \dim \operatorname{Hom} _ {\mathrm{G}} (\pi , \operatorname{M} _ {\pi} (\varrho)) = \operatorname{Card} (\mathrm{I}) \end{array}
\]

（同上）。这证明了 J 也是有限的，且 \(\text{Card(I)} = \text{Card(J)}\)，正如所要求的。同样，若 J 有限，则 I 有限且具有相同的基数。

现在假设 I 和 J 都是无限的。对于 \(j \in J\)，记 \(p_{j}\) 为 E 上像为 \(F_{j}\) 的正交投影。对于 \(i \in I\)，取 \(x_{i}\) 为 \(E_{i}\) 中的非零元素；它是 \(E_{i}\) 的循环向量。注意，由于 \(p_{j}\) 通过限制到子空间得到的是从 \(E_{i}\) 到 \(F_{j}\) 的 G-态射，空间 \(p_{j}(E_{i})\) 为零，当且仅当 \(p_{j}(x_{i}) = 0\)（事实上，空间 \(E_{i} \cap \text{Ker}(p_{j})\) 要么为零，要么等于 \(E_{i}\)）。

对于每个 \(j \in J\)，存在 \(f(j) \in I\)，使得 \(p_{j}(\mathrm{E}_{f(j)})\) 不为 0（否则，正交投影 \(p_{j}\) 在 \(\mathrm{M}_{\pi}(\varrho)\) 上为零）。于是定义了从 J 到 I 的映射 f。对于每个 \(i \in I\)，集合 \(f^{-1}(i)\) 是可数的。事实上，该集合包含于由 \(j \in J\) 中满足 \(p_{j}(\mathrm{E}_{i})\) 非零者组成的集合中，也就是满足 \(p_{j}(x_{i}) \neq 0\) 的那些 j。现在由（EVT, V, p. 20, cor. 2），有

\[
\sum_ {j \in \mathrm{J}} \| p _ {j} (x _ {i}) \| ^ {2} = \| x _ {i} \| ^ {2},
\]

因此，\(j \in J\) 中满足 \(p_{j}(x_{i}) \neq 0\) 的元素所成集合是可数的。由 E, III, p. 50, 命题 4，得到 \(\text{Card}(\text{J}) = \text{Card}(f(\text{J})) \leqslant \text{Card}(\text{I})\)。交换 I 和 J 的角色，得到 \(\text{Card}(\text{I}) = \text{Card}(\text{J})\)。

推论. --- 设 G 是拓扑群，\(\varrho\) 是 G 在复希尔伯特空间 E 中的酉表示。G 的 \(\pi\)-同型分量（其中 \(\pi \in \widehat{\mathbf{G}}\)）的希尔伯特直和是 \(\varrho\) 的最大半单子表示。

事实上，\(\pi\)-同型分量对每个 \(\pi \in \widehat{\mathbf{G}}\) 两两正交（命题 8，\(a\)），且每个都是半单的（命题 10，\(a\)），所以空间 \(\mathrm{M}_{\pi}(\varrho)\) 对 \(\pi \in \widehat{\mathbf{G}}\) 的希尔伯特直和 F 定义了 \(\varrho\) 的一个半单子表示。此外，\(\varrho\) 的每个不可约子表示都是 \(\varrho\) 的某个同型分量的子表示，故推论得证。

定义 10. --- 设 G 是拓扑群。设 \(\varrho\) 是 G 在复希尔伯特空间 E 中的酉表示，\(\pi\) 是 G 在复希尔伯特空间中的不可约酉表示。

称 \(\pi\) 在 \(\varrho\) 中的重数为指标集的基数，所针对的是 E 的在 G 下稳定的闭子空间族 \((\mathrm{E}_i)_{i\in \mathrm{I}}\)，其中 \(\varrho\) 在 \(\mathrm{E}_i\) 中诱导的子表示与 \(\pi\) 同构，对每个 \(i\in I\) 成立，并且 \(\mathrm{M}_{\pi}(\varrho)\) 是子空间 \(\mathrm{E}_i\) 的希尔伯特直和。

如果 \(\pi\) 在 \(\varrho\) 中的重数有限，则由 V, p. 377 的公式（1）和 V, p. 387 的推论 2，它等于空间 \(\mathrm{Hom}_{\mathrm{G}}(\pi, \varrho)\) 的维数（分别等于 \(\mathrm{Hom}_{\mathrm{G}}(\varrho, \pi)\) 的维数）。一般而言这并不总是成立。

注. --- 一个酉表示 \(\varrho\) 可能非零，但相对于 G 的所有不可约表示的 \(\varrho\) 的所有同型分量都为零（参见 V, p. 426, 注）。

\section{局部紧群的表示}

本段中，向量空间均在域 K = R 或 C 上。记 G 为赋予左哈尔测度 \(\mu\) 的局部紧群。记 \(\mathcal{L}^{p}(\mathrm{G}) = \mathcal{L}_{\mathbf{C}}^{p}(\mathrm{G}, \mu)\)，\(\mathrm{L}^{p}(\mathrm{G}) = \mathrm{L}_{\mathbf{C}}^{p}(\mathrm{G}, \mu)\)，其中 \(p \in [1, +\infty]\)。

\subsection{某些表示的连续性}

命题 1. — 设 H 是局部紧群。设 \(\varrho_{1}\)（分别为 \(\varrho_{2}\)）是 G（分别为 H）在分离局部凸 K-向量空间 E \(_{1}\)（分别为 E \(_{2}\)）中的连续表示。令 F 表示赋予紧收敛拓扑的空间 \(\mathcal{L}(\mathrm{E}_{1};\mathrm{E}_{2})\)。由 \(\varrho\) 通过 \(\varrho(g,h)u=\varrho_{2}(h)\circ u\circ\varrho_{1}(g^{-1})\)（\((g,h)\in\mathrm{G}\times\mathrm{H}\)）定义的 G×H 到 F 的表示是连续的。

记 \(\mathcal{L}_c(\mathrm{E}_1;\mathrm{E}_2)\) 为赋予紧收敛拓扑的空间 \(\mathcal{L}(\mathrm{E}_1;\mathrm{E}_2)\)。设 \(\mathfrak{F}_1\)（分别为 \(\mathfrak{F}_2,\mathfrak{F}_3\)）是趋于 \(e\) 的 G 中滤子（分别为趋于 0 的 \(\mathcal{L}_c(\mathrm{E}_1;\mathrm{E}_2)\) 中滤子、趋于 \(e\) 的 H 中滤子）。由于 G 和 H 局部紧，存在 \(\mathrm{C}\in \mathfrak{F}_1\) 和 \(\mathrm{D}\in \mathfrak{F}_3\)，它们是相对紧的。集合 \(\varrho_1(\mathrm{C}^{-1})\) 在 \(\mathcal{L}(\mathrm{E}_1)\) 中等度连续（参见 INT, VIII, p. 129, § 2, no. 1, 注 2, \(a'\)）。由 EVT, III, p. 33 的命题 9 和 EVT, III, p. 31 的命题 4，由 \((u,v)\mapsto u\circ v\) 定义的从 \(\mathcal{L}_c(\mathrm{E}_1;\mathrm{E}_2)\times \varrho_1(\mathrm{C}^{-1})\) 到 \(\mathcal{L}_c(\mathrm{E}_1;\mathrm{E}_2)\) 的映射是连续的。因此，滤基 \(\mathfrak{F}_2\circ \varrho_1(\mathfrak{F}_1^{-1})\) 在 \(\mathcal{L}_c(\mathrm{E}_1;\mathrm{E}_2)\) 中趋于 0。集合 \(\varrho_2(\mathrm{D})\) 在 \(\mathcal{L}(\mathrm{E}_2)\) 中等度连续（参见 INT, VIII, p. 129, § 2, n° 1, rem. 2），并且对每个 \(x\in \mathrm{E}_2\)，集合 \(\varrho_2(\mathrm{D})x\subset \mathrm{E}_2\) 是相对紧的。因此，赋予紧收敛拓扑的 \(\varrho_2(\mathrm{D})\) 在 \(\mathcal{L}(\mathrm{E}_2)\) 中相对紧（TG, X, p. 18, cor. 1）。由 \((u,v)\mapsto u\circ v\) 定义的从 \(\overline{\varrho_2(\mathrm{D})}\times \mathcal{L}_c(\mathrm{E}_1;\mathrm{E}_2)\) 到 \(\mathcal{L}_c(\mathrm{E}_1;\mathrm{E}_2)\) 的映射，由 EVT, III, p. 33 的命题 9 和 EVT, III, p. 31 的命题 4，是连续的。因此，滤基 \(\varrho (\mathfrak{F}_1\times \mathfrak{F}_3)(\mathfrak{F}_2)\) 在 \(\mathcal{L}_c(\mathrm{E}_1;\mathrm{E}_2)\) 中趋于 0。这就证明了断言。

推论. --- 设 \(\varrho\) 是 G 在分离局部凸 K-向量空间 E 上的连续表示。当 E' 赋予紧收敛拓扑时，逆步表示 \(\breve{\varrho}\) 是连续的。

取 \(H = \{e\}\)、\(\varrho_{1} = \varrho\)，并令 \(\varrho_{2}\) 为 K 上的平凡表示，即由命题得到。

注. --- 当 E' 赋予强拓扑时，逆步表示不一定连续（参见 INT, VIII, p. 191, § 2, 习题 3, d)）。可以证明，如果空间 E 是半自反的，则它是连续的（同上，c)）。

\subsection{表示向测度空间的延拓}

本节中假设 K = C。

设 \(\varrho\) 是 G 在分离局部凸拟完备空间 E 中的连续线性表示。记 \(\mathcal{M}_{c}(\mathrm{G})\) 为 G 上具有紧支撑且赋予紧收敛拓扑的测度空间；它是空间 \(\mathcal{C}(\mathrm{G})\) 的对偶。对于每个测度 \(\nu \in \mathcal{M}_{c}(\mathrm{G})\)，置

\[
\varrho (\nu) = \int_ {\mathrm{G}} \varrho (g) d \nu (g).
\]

这是 \(\mathcal{L}(\mathrm{E})\) 的一个元素。特别地，\(\varrho(\varepsilon_g) = \varrho(g)\) 对所有 \(g \in \mathrm{G}\) 成立。

根据 INT, VIII, p. 136, § 2, no. 6，映射 \(\nu \mapsto \varrho(\nu)\) 是从 \(\mathcal{M}_c(\mathrm{G})\) 到赋予紧收敛拓扑的空间 \(\mathcal{L}(\mathrm{E})\) 的连续线性映射。根据 INT, VIII, p. 145, § 3, no. 3, prop. 11，它是含单位代数的态射。

设 \(x \in \mathrm{E}\)。当 \(\mathcal{M}_c(\mathrm{G})\) 赋予紧收敛拓扑时，由 \(\nu \mapsto \varrho(\nu)x\) 定义的从 \(\mathcal{M}_c(\mathrm{G})\) 到 E 的映射是连续的（INT, VI, p. 27, § 1, no. 7, prop. 16 以及 EVT, III, p. 31, prop. 4）。

对于在 G 的某个紧子集之外为零的 \(f \in \mathcal{L}^{1}(\mathrm{G})\)，测度 \(f \cdot \mu\) 具有紧支撑；我们把 E 的自同态 \(\varrho^{\mu}(f)\)（也记为 \(\varrho(f)\)）记为 \(\varrho(f \cdot \mu)\)。此自同态只取决于其类 \(\widetilde{f}\)，即 \(f\) 在 \(\mathrm{L}^{1}(\mathrm{G})\) 中的类，我们也把它记为 \(\varrho^{\mu}(\widetilde{f})\) 或 \(\varrho(\widetilde{f})\)。

同样，设 \(\varrho\) 是 G 在巴拿赫空间 E 中的连续有界线性表示。对于每个有界测度 \(\nu \in \mathcal{M}^{1}(\mathrm{G})\)，置

\[
\varrho (\nu) = \int_ {\mathrm{G}} \varrho (g) d \nu (g).
\]

根据 INT, VIII, 同上，映射 \(\nu \mapsto \varrho(\nu)\) 是从 G 上有界复测度代数 \(\mathcal{M}^{1}(G)\) 到巴拿赫代数 \(\mathcal{L}(E)\) 的连续含单位巴拿赫代数态射。

令 \(\rho\) 为 G 上的函数 \(g\mapsto\|\varrho(g)\|\)；INT, VIII, p. 145 的命题 11 中记作 \(\mathcal{M}^{\rho}\) 的代数（其元素是测度 \(\nu\)，满足 \(\rho\in\mathcal{L}^{1}(\nu)\)）与巴拿赫代数 \(\mathcal{M}^{1}(G)\) 重合。事实上，置 \(M = \sup \rho\)。有 \(M > 0\)，因为 \(\rho(e)=1\)，且 \(M^{-1} \leqslant \rho \leqslant M\)，因为不等式 \(\|\varrho(e)\|\leqslant\|\varrho(g)\|\|\varrho(g^{-1})\|\) 对所有 \(g\in G\) 成立。因此，\(\rho\in\mathcal{L}^{1}(\nu)\) 当且仅当 \(\nu\) 是有界测度。

若 \(f \in \mathcal{L}^{1}(\mathrm{G})\)，自同态 \(\varrho^{\mu}(f)\)（或简记为 \(\varrho(f)\)）就是 \(\varrho(f \cdot \mu)\)；对于类 \(\tilde{f}\)，它是 \(f\) 在 \(\mathrm{L}^{1}(\mathrm{G})\) 中的类，同样如此。

引理 1. --- 设 \(\varrho\) 是 G 在希尔伯特空间 E 中的酉表示。映射 \(\nu \mapsto \varrho(\nu)\) 从 \(\mathcal{M}^{1}(\mathrm{G})\) 到 \(\mathcal{L}(\mathrm{E})\)，是对合巴拿赫代数的含单位态射。

由前述内容，只需证明态射 \(\nu \mapsto \varrho(\nu)\) 是对合的。设 \(\nu \in \mathcal{M}^{1}(\mathrm{G})\)。测度 \(\nu^{*}\) 是测度 \(\check{\nu}\) 的共轭测度（I, p. 99, 例 4）。根据共轭测度的定义（INT, III, p. 52, § 1, no. 5），计算得

\[
\begin{array}{c} \langle x \mid \varrho (\nu) y \rangle = \int_ {\mathrm{G}} \langle x \mid \varrho (g) y \rangle d \nu (g) \\ = \int_ {\mathrm{G}} \langle \varrho (g ^ {- 1}) x \mid y \rangle d \nu (g) = \int_ {\mathrm{G}} \langle \varrho (g) x \mid y \rangle d \check {\nu} (g) \\ = \overline {{\int_ {\mathrm{G}} \langle y \mid \varrho (g) x \rangle d \nu^ {*} (g)}} = \langle \varrho (\nu^ {*}) x \mid y \rangle \end{array}
\]

对 E 中所有 \(x\) 和 \(y\) 成立，从而 \(\varrho(\nu)^{*} = \varrho(\nu^{*})\)。

设 \(\varrho\) 是 G 在局部凸拟完备空间 E（分别在巴拿赫空间 E）中的连续线性表示（分别为连续有界表示）。如果 F 是 E 的闭子空间，并且子表示 \(\varrho_{\mathrm{F}}\) 由 \(\varrho\) 在 F 上诱导，那么对于每个测度 \(\nu \in \mathcal{M}_c(\mathrm{G})\)（分别为 \(\nu \in \mathcal{M}^1 (\mathrm{G})\)），线性映射 \(\varrho (\nu)\) 保持子空间 F 不变，并且在 F 上与 \(\varrho_{\mathrm{F}}(\nu)\) 重合。

反之，设 \(F \subset E\) 是闭子空间，并且在每个线性映射 \(\varrho(f)\) 下稳定，其中 \(f \in \mathcal{H}(G)\)。则 F 定义 \(\varrho\) 的一个子表示（INT, VIII, p. 139, § 2, no. 7, prop. 10）。

回忆（A, VIII, p. 388），函数 \(f \colon \mathbf{G} \to \mathbf{C}\) 称为中心函数，是指对所有 \(g\) 和 \(h\)（均属于 \(\mathbf{G}\)），有 \(f(gh) = f(hg)\)。这等价于说 \(f\) 在共轭作用下不变。

定义 1. --- 有界测度 \(\nu \in \mathcal{M}^{1}(\mathrm{G})\) 称为中心的，是指 \(\varepsilon_{g} * \nu = \nu * \varepsilon_{g}\) 对所有 \(g \in \mathrm{G}\) 成立。

如果 G 是单模群，且 \(f \in \mathcal{C}(\mathrm{G})\) 关于 \(\mu\) 可积，那么测度 \(f \cdot \mu\) 中心，当且仅当 f 是中心函数。

设 \(\varrho\) 是 G 在巴拿赫空间 E 中的连续有界表示（分别是 G 在拟完备分离局部凸空间 E 中的连续表示）。对于 G 上的每个有界中心测度 \(\nu\)（分别为 G 上的任意紧支撑中心测度 \(\nu\)），线性映射 \(\varrho(\nu)\) 是从 \(\varrho\) 到 \(\varrho\) 的 G-态射。事实上，对每个 \(g \in G\)，有

\[
\varrho (\nu) \varrho (g) = \varrho (\nu * \varepsilon_ {g}) = \varrho (\varepsilon_ {g} * \nu) = \varrho (g) \varrho (\nu).
\]

\subsection{半单性判据}

命题 2. --- 设 \(\varrho\) 是 G 在复希尔伯特空间 E 中的酉表示。设 \(\mathfrak{F}\) 是 \(\mathcal{M}_{c}(\mathrm{G})\) 上的滤基，并且对于紧收敛拓扑趋于测度 \(\varepsilon_{e}\)。假设存在 \(\mathrm{M} \in \mathfrak{F}\)，使得 \(\varrho(\nu)\) 对于每个 \(\nu \in \mathrm{M}\) 都是 E 的紧自同态。

那么酉表示 \(\varrho\) 是半单的，并且 G 的每个不可约酉表示在 \(\varrho\) 中的重数有限。

先证明一个引理。

引理 2. --- 设 \(\varrho_{1}\) 是 G 的一个非零酉表示，并且与 \(\varrho\) 的一个子表示同构。存在测度 \(\nu \in M\)，使得 \(\varrho_{1}(\check{\nu} * \nu)\) 是 \(\varrho_{1}\) 的表示空间上的非零紧埃尔米特自同态。特别地，存在非零实数 \(\lambda\)，使得 \(\varrho_{1}(\check{\nu} * \nu)\) 关于 \(\lambda\) 的特征空间非零。

不妨设 \(\varrho_{1}\) 是 \(\varrho\) 的子表示。令 E \(_{1}\) 为其表示空间。对于每个测度 \(\nu \in \mathrm{M}\)，自同态 \(\varrho_{1}(\nu)\) 由 \(\varrho(\nu)\) 限制到子空间得到。因此，根据假设及 III, p. 5 的命题 3，\(\varrho_{1}(\nu)\) 是紧的。

由于 \(\mathfrak{F}\) 趋于测度 \(\varepsilon_{e}\)，且收敛所在的空间是赋予紧收敛拓扑的 \(\mathcal{M}_{c}(\mathrm{G})\)，而空间 \(\mathrm{E}_{1}\) 非零，存在 \(\nu\in\mathrm{M}\)，使得 \(v=\varrho_{1}(\nu)\) 是 \(\mathrm{E}_{1}\) 的非零自同态（参见 V, p. 400, n°2）。自同态 \(u=\varrho_{1}(\check{\nu}*\nu)\) 等于 \(v^{*}\circ v\)（V, p. 401 的引理 1）；因此它非零（EVT, V, p. 39, prop. 2），是埃尔米特的且是紧的，因为 \(v\) 是紧的。

由于 \(u\) 是非零埃尔米特自同态，其谱不只含 0（I, p. 110 的例 1）。最后，由于 \(u\) 是紧的，每个 \(\lambda \in \mathrm{Sp}(u)\setminus\{0\}\) 都是 \(u\) 的特征值（III, p. 83 的命题 2）。

现在证明命题。

我们将使用 V, p. 391 的命题 7 来证明 \(\varrho\) 半单。设 \(\varrho_{1}\) 是 \(\varrho\) 的一个非零子表示，\(\mathrm{E}_{1} \subset \mathrm{E}\) 是其表示空间；我们必须证明 \(\varrho_{1}\) 包含一个不可约子表示。

取 \(\nu \in \mathrm{M}\)，使得 \(u = \varrho_{1}(\check{\nu}*\nu)\) 是 \(\varrho_{1}\) 的表示空间上的非零紧埃尔米特自同态，并取非零的 \(\lambda\)，使得 \(u\) 关于 \(\lambda\) 的特征空间 F 非零（引理 2）。空间 F 是有限维的（III, p. 90 的命题 5）。存在子表示 \(\varrho_{2}\)，它是 \(\varrho_{1}\) 在子空间 \(\mathrm{E}_{2}\)（\(\mathrm{E}_{1}\) 的子空间）上的表示，使得 \(\mathrm{E}_{2} \cap \mathrm{F}\) 非零且维数最小。取非零元素 \(x\)，它属于 \(\mathrm{E}_{2} \cap \mathrm{F}\)；令 \(\mathrm{E}_{3}\) 为 \(\varrho_{1}\) 中由 \(x\) 生成的子表示。于是 \(\mathrm{E}_{3} \subset \mathrm{E}_{2}\)，故 \(\mathrm{E}_{3} \cap \mathrm{F} = \mathrm{E}_{2} \cap \mathrm{F}\)，这是由 \(\mathrm{E}_{2} \cap \mathrm{F}\) 的维数极小性得到的。

证明表示 \(E_{3}\) 不可约。设 \(E_{4} \subset E_{3}\) 是在 G 下稳定的闭子空间。有 \(E_{2} \cap F = (E_{4} \cap F) \oplus (E_{4}^{\circ} \cap F)\)，其中 \(E_{4}^{\circ}\) 表示 \(E_{4}\) 在空间 \(E_{3}\) 中的正交补（事实上，若 \(y \in E_{2} \cap F\)，令 \(y_{4} \in E_{4}\) 为 y 在 \(E_{4}\) 上的正交投影；由于 \(E_{4}\) 和 \(E_{4}^{\circ}\) 在 u 下稳定，向量 \(u(y_{4})\) 是 \(u(y) = \lambda y\) 在 \(E_{4}\) 上的正交投影，从而 \(u(y_{4}) = \lambda y_{4}\)，这意味着 \(y_{4} \in E_{4} \cap F\)，断言得证）。于是由 \(E_{2} \cap F\) 的维数极小性，或者有 \(E_{4} \cap F = E_{2} \cap F\)，或者有 \(E_{4}^{\circ} \cap F = E_{2} \cap F\)。在第一种情形中，\(x \in E_{4}\)，从而 \(E_{4} = E_{3}\)；在第二种情形中，\(x \in E_{4}^{\circ}\)，从而 \(E_{4}^{\circ} = E_{3}\) 且 \(E_{4} = \{0\}\)。因此表示 \(E_{3}\) 不可约。

由 V, p. 391 的命题 7，得到表示 \(\varrho\) 是半单的。

设 \(\pi\) 是 G 的不可约酉表示，且其在 \(\varrho\) 中的重数非零；令 \(E_{\pi}\) 为其表示空间。存在测度 \(\nu \in M\) 和非零实数 \(\lambda\)，使得 \(u = \pi (\check{\nu} * \nu)\) 是 \(E_{\pi}\) 上的非零紧埃尔米特自同态，并且 \(u\) 关于 \(\lambda\) 的特征空间非零（引理 2）。令 F 为 \(v = \varrho (\check{\nu} * \nu)\) 关于 \(\lambda\) 的特征空间。对于每个子表示 \(E_1\)，它是 E 的与 \(\pi\) 同构的子表示；在 \(E_1\) 上由 \(v\) 限制到子空间得到的自同态与 \(u\) 等同，且 \(E_1\) 与 F 的交非零。因此，\(\pi\) 在 \(\varrho\) 中的重数小于空间 F 的维数，而该维数有限（III, p. 90 的命题 5）。

例. --- 假设 G 是单模群，例如实半单李群。设 \(\Gamma \subset G\) 是离散子群，使得商空间 \(X = \Gamma \backslash G\) 紧。群 G 通过右乘作用在 X 上。记 \(\beta\) 为 \(\Gamma\) 上的计数测度，并置 \(\widetilde{\mu} = \mu / \beta\)（INT, VII, p. 44, § 2, n° 2, 定义 1）；这是 X 上的有界 G-不变测度。

对于每个 \(f \in \mathcal{L}^2(\mathrm{X},\widetilde{\mu})\) 和每个 \(g \in \mathrm{G}\)，定义函数 \(\varrho(g)f \in \mathcal{L}^2(\mathrm{X},\widetilde{\mu})\)，使其满足 \(\varrho(g)f(x) = f(x \cdot g)\)。映射 \(\varrho(g)\) 是连续线性映射，通过取商在 \(\mathrm{L}^2(\mathrm{X},\widetilde{\mu})\) 上诱导出一个酉映射，仍记为 \(\varrho(g)\)。映射 \(\varrho\) 是 G 在 \(\mathrm{L}^2(\mathrm{X},\widetilde{\mu})\) 中的酉表示（INT, VII, p. 135, § 2, n° 5, prop. 8）。

取 \(\varphi \in \mathcal{K}(\mathrm{G})\)。对于 \(\mathrm{L}_1\) 和 \(\mathrm{L}_2\) 这两个 \(\mathrm{G}\) 的任意紧子集，\(\mathrm{T}\) 是 \(\Gamma\) 与紧集 \(\mathrm{L}_1\mathrm{Supp}(\varphi)\mathrm{L}_2^{-1}\) 的交，且该交是有限的。对于每个 \((g,h) \in \mathrm{L}_1 \times \mathrm{L}_2\)，级数 \(\sum_{\gamma \in \Gamma} \varphi(g^{-1}\gamma h)\) 与有限和 \(\sum_{\gamma \in \mathrm{T}} \varphi(g^{-1}\gamma h)\) 重合。由于 \(\mathrm{G}\) 局部紧，该级数的和 \(k_{\varphi}(g,h)\) 是 \(\mathrm{G} \times \mathrm{G}\) 上的连续函数。

取 \((g,h) \in G \times G\) 以及 \((\gamma, \eta) \in \Gamma \times \Gamma\)。有 \(k_{\varphi}(\gamma g, \eta h) = k_{\varphi}(g, h)\)，因此 \(k_{\varphi}\) 通过取商定义了 \(X \times X\) 上的连续函数，记为 \(\widetilde{k}_{\varphi}\)。又有 \(\widetilde{k}_{\varphi} \in \mathcal{L}^{2}(X \times X, \widetilde{\mu} \otimes \widetilde{\mu})\)，因为假设空间 X 是紧的。

记 \(\dot{x}\) 为 G 中元素 \(x\) 在规范投影下于 X 中的像。取 \(\varphi \in \mathcal{K}(\mathrm{G})\)。对于 \(f \in \mathcal{L}^2(\mathrm{X},\widetilde{\mu})\) 和 \(x \in \mathrm{G}\)，有

\[
\begin{array}{c} \varrho (\varphi) f (\dot {x}) = \int_ {\mathrm{G}} \varphi (g)   (\varrho (g) f) (\dot {x}) d \mu (g) = \int_ {\mathrm{G}} \varphi (g) f (\dot {x} \cdot g) d \mu (g) \\ = \int_ {\mathrm{G}} \varphi (x ^ {- 1} y) f (\dot {y}) d \mu (y) = \int_ {\Gamma \backslash \mathrm{G}} \Bigl (\sum_ {\gamma \in \Gamma} \varphi (x ^ {- 1} \gamma y) \Bigr) f (\dot {y}) d \widetilde {\mu} (\dot {y}) \end{array}
\]

（INT, VII, p. 46, § 2, no. 3, prop. 5）。由于 \(\widetilde{k}_{\varphi}\) 属于空间 \(\mathcal{L}^2 (\mathrm{X}\times \mathrm{X},\widetilde{\mu}\otimes \widetilde{\mu})\)，所以自同态 \(\varrho (\varphi)\) 在 \(\mathrm{L}^2 (\mathrm{X},\widetilde{\mu})\) 上与核为 \(\widetilde{k}_{\varphi}\) 的希尔伯特–施密特自同态重合；因此该自同态是紧的（III, p. 33 的推论 1）。

存在函数序列 \((\varphi_{n})_{n\in\mathbf{N}}\)，其元素属于 \(\mathcal{K}_{+}(\mathrm{G})\) 且积分为 1，使得 \(\varphi_{n}\cdot\mu\) 趋于测度 \(\varepsilon_{e}\)（在赋予紧收敛拓扑的 \(\mathcal{M}_{c}(\mathrm{G})\) 中）（INT, VIII, p. 139, § 2, no. 7, cor. 2）。因此命题 2 蕴含表示 \(\varrho\) 半单，且 G 的不可约酉表示在 \(\varrho\) 中的重数有限。

在 \(\varrho\) 中重数非零的 G 的不可约酉表示，称为群 G 的 \(\Gamma\)-自守表示。

\subsection{正则表示}

设 \(p\) 是满足 \(\geqslant 1\) 的实数，\(\mu'\) 是 \(G\) 上的右哈尔测度。对于每个 \(g \in G\) 及每个函数 \(f \in \mathcal{L}^p(G, \mu)\)，记 \(\gamma_G^{(p)}(g)f\) 为函数 \(x \mapsto f(g^{-1}x)\)（定义在 \(G\) 上）。映射 \(g \mapsto \gamma_G^{(p)}(g)\) 是 \(G\) 在 \(\mathcal{L}^p(G)\) 中的连续线性表示。通过取商，它诱导出 \(G\) 在 \(L^p(G)\) 中的连续等距线性表示（INT, VIII, p. 135, § 2, no. 5, prop. 8），记为 \(\gamma_G^{(p)}\)。

同样，对于每个 \(g \in G\) 及每个函数 \(f \in \mathcal{L}^{p}(G, \mu')\)，记 \(\boldsymbol{\delta}_{\mathrm{G}}^{(p)}(g)f\) 为 G 上的函数 \(x \mapsto f(xg)\)。映射 \(g \mapsto \boldsymbol{\delta}_{\mathrm{G}}^{(p)}(g)\) 是 G 在 \(\mathcal{L}^{p}(G, \mu')\) 中的连续线性表示。通过取商，它诱导出 G 在 \(\mathrm{L}^{p}(\mathrm{G}, \mu')\) 中的连续等距线性表示（参见 INT, VIII, p. 136, § 2, no. 5）。

称 \(\boldsymbol{\gamma}_{\mathrm{G}}^{(p)}\) 为 G 在 \(\mathcal{L}^{p}(\mathrm{G})\) 或 \(\mathrm{L}^{p}(\mathrm{G})\) 中的左正则表示，称 \(\boldsymbol{\delta}_{\mathrm{G}}^{(p)}\) 为 G 在 \(\mathcal{L}^{p}(\mathrm{G},\mu^{\prime})\) 或 \(\mathrm{L}^{p}(\mathrm{G},\mu^{\prime})\) 中的右正则表示。

引理 3. — 设 p 是大于等于 1 的实数。G 在 L \(^{p}\) (G,μ) 中的左（分别为右）正则表示（分别在 L \(^{p}\) (G,μ') 中）是忠实的。

更确切地，令 q 为 p 的共轭指数，令 g 为 G 中满足 \(g \neq e\) 的元素。

\begin{enumerate}
\def\labelenumi{\alph{enumi})}
\item
  存在函数 \(\varphi\in\mathcal{K}(\mathrm{G})\)，它在 \(\mathrm{L}^{q}(\mathrm{G},\mu)\) 中为正且非零，并满足 \(\langle\varphi,\boldsymbol{\gamma}_{\mathrm{G}}^{(p)}(g)\varphi\rangle=0\)；
\item
  存在函数 \(\varphi\in\mathcal{K}(\mathrm{G})\)，它在 \(\mathrm{L}^{q}(\mathrm{G},\mu')\) 中为正且非零，并满足 \(\langle\varphi,\boldsymbol{\delta}_{\mathrm{G}}^{(p)}(g)\varphi\rangle=0\)。
\end{enumerate}

考虑左正则表示 \(\boldsymbol{\gamma}_{\mathrm{G}}^{(p)}\) 的情形，右正则表示的情形类似。断言 \(a\) 蕴含 \(\boldsymbol{\gamma}_{\mathrm{G}}^{(p)}\) 是忠实的，因为若 \(g \neq e\) 是 G 中的元素，且 \(\varphi\) 如 \(a\) 所述，则有 \(\varphi \neq \boldsymbol{\gamma}_{\mathrm{G}}^{(p)}(g)\varphi\)，因为 \(\langle \varphi, \varphi \rangle = \int_{\mathrm{G}} \varphi^2 > 0\)。

因此证明 \(a\)。令 \(g \neq e\) 属于 G。令 C 是 \(e\) 的一个紧对称邻域，使得 \(g \notin \mathrm{C}^2\)，并令 \(\varphi \in \mathcal{K}(\mathrm{G})\) 是积分为 1 且支撑包含于 C 的连续正函数；函数 \(\varphi\) 在 \(\mathrm{L}^p(\mathrm{G}, \mu)\) 中非零。由于 \(\mathrm{C} \cap g^{-1}\mathrm{C} = \emptyset\)，有

\[
\langle \varphi , \boldsymbol {\gamma} _ {\mathrm{G}} ^ {(p)} (g) \varphi \rangle = \int_ {\mathrm{G}} \varphi (x) \varphi (g ^ {- 1} x) d \mu (x) = 0.
\]

设 \(\varrho\) 是 G 在巴拿赫空间 E 中的连续有界表示。对于每个 \(f\in \mathrm{L}^1 (\mathrm{G})\)（分别为 \(f^{\prime}\in \mathrm{L}^{1}(\mathrm{G},\mu^{\prime}))\)）以及每个 \(g\in \mathrm{G}\)，有

\[
\varrho (g) \varrho (f \cdot \mu) = \varrho (\varepsilon_ {g} * (f \cdot \mu)) = \varrho (\boldsymbol {\gamma} _ {\mathrm{G}} ^ {(1)} (g) f \cdot \mu),\tag{1}
\]

\[
\varrho (f ^ {\prime} \cdot \mu^ {\prime}) \varrho (g) = \varrho ((f ^ {\prime} \cdot \mu^ {\prime}) * \varepsilon_ {g}) = \varrho (\boldsymbol {\delta} _ {\mathrm{G}} ^ {(1)} (g ^ {- 1}) f ^ {\prime} \cdot \mu ^ {\prime}),\tag{2}
\]

（INT, VIII, p. 144, § 3, no. 2, 公式（5））。

若 G 是单模群，置 \(\mu' = \mu\)。G 在 \(\mathcal{L}^{p}(\mathrm{G})\)（分别在 \(\mathrm{L}^{p}(\mathrm{G})\)）中的双正则表示 \(\varrho_{\mathrm{G}}^{(p)}\)，是 \(G \times G\) 在 \(\mathcal{L}^{p}(\mathrm{G})\)（分别在 \(\mathrm{L}^{p}(\mathrm{G})\)）中的表示，满足

\[
\boldsymbol {\varrho} _ {\mathrm{G}} ^ {(p)} (g _ {1}, g _ {2}) = \boldsymbol {\gamma} _ {\mathrm{G}} ^ {(p)} (g _ {1}) \boldsymbol {\delta} _ {\mathrm{G}} ^ {(p)} (g _ {2}) = \boldsymbol {\delta} _ {\mathrm{G}} ^ {(p)} (g _ {2}) \boldsymbol {\gamma} _ {\mathrm{G}} ^ {(p)} (g _ {1}).
\]

这是连续线性表示（V, p. 377 的引理 1）。它满足

\[
(\boldsymbol {\varrho} _ {\mathrm{G}} ^ {(p)} (g _ {1}, g _ {2}) f) (x) = f (g _ {1} ^ {- 1} x g _ {2})
\]

对所有 \(f \in \mathcal{L}^p(\mathrm{G})\)、所有 \((g_1, g_2) \in \mathrm{G} \times \mathrm{G}\) 及每个 \(x \in \mathrm{G}\) 成立。

注. --- G 在 \(\mathrm{L}^{p}(\mathrm{G},\mu)\) 中的双正则表示不一定忠实；其核是 G 的中心在映射 \(g\mapsto(g,g)\) 下的像（V, p. 487 的习题 4）。

当 \(p = 2\) 时，G 的左正则表示 \(\gamma_{\mathrm{G}}^{(2)}\) 在复希尔伯特空间 \(\mathrm{L}^2 (\mathrm{G},\mu)\) 中是酉的，因为它是等距的。同样，右正则表示 \(\delta_{\mathrm{G}}^{(2)}\) 在 \(\mathrm{L}^2 (\mathrm{G},\mu')\) 中也是酉的。

我们简记 \(\gamma_{\mathrm{G}} = \gamma_{\mathrm{G}}^{(2)}\) 和 \(\delta_{\mathrm{G}} = \delta_{\mathrm{G}}^{(2)}\)，并将这些表示称为 G 的左、右正则表示。

若 G 是单模群，则双正则表示 \(\varrho_{\mathrm{G}}^{(2)}\) 是 \(G \times G\) 在 \(\mathrm{L}^{2}(\mathrm{G}, \mu)\) 中的表示，且是酉的。我们简记它为 \(\varrho_{G}\)。

引理 4. --- 设 \(p\) 是满足 \(\geqslant 1\) 的实数。对于 \(f \in \mathcal{L}^1(\mathrm{G})\)，线性映射 \(\gamma_{\mathrm{G}}^{(p)}(f)\) 与 \(\mathrm{L}^p(\mathrm{G})\) 上由 \(\varphi \mapsto f *^\mu \varphi\) 定义的自同态重合。

这由 INT, VIII, p. 157, § 4, n°2 的命题 6 得出，同时要注意 INT, VIII, p. 165 的公式（14）。

注. --- 回忆一下，对于 G 上的任意函数 \(f\)，在 G 上定义函数 \(\check{f}\)，使得 \(\check{f}(g) = f(g^{-1})\) 对所有 \(g \in \mathrm{G}\) 成立（INT, VII, p. 12, § 1, n° 1, 公式（12））。

可以检验，若 \(f \in \mathcal{L}^{1}(\mathrm{G})\)，线性映射 \(u = \boldsymbol{\delta}_{\mathrm{G}}^{(p)}(f)\) 满足关系式

\[
\widetilde {u (\varphi)} = f * \check {\varphi}
\]

对所有 \(\varphi\in\mathrm{L}^{p}(\mathrm{G},\mu^{\prime})\) 成立。
\subsection{等变函数}

在本小节中，我们取定 G 的一个闭子群 H。用 \(\varpi\) 表示 G 到 G/H 上的典范投射。

除 G 上的哈尔测度 \(\mu\) 外，我们再取定 H 上的一个左哈尔测度 \(\beta\)。

根据 INT, VII, 第 56 页, § 2, 第 5 目, 定理 2，存在连续函数 \(\kappa\colon\mathbf{G}\to\mathbf{R}_{+}^{*}\)，使得 \(\kappa(xh)=\Delta_{\mathrm{H}}(h)\Delta_{\mathrm{G}}(h)^{-1}\kappa(x)\) 对每个 \((x,h)\in\mathbf{G}\times\mathbf{H}\) 成立。我们取定这样一个函数 \(\kappa\)。用 \(\nu\) 表示 G/H 上的测度 \((\kappa\cdot\mu)/\beta\)；由出处同上，它是 G 下拟不变的非零正测度。它的支撑等于 G/H（INT,

VII, 第 10 页, § 1, 第 1 目）。根据 INT, VII, 第 43 页, § 2, 第 2 目, 命题 4，测度 \(\nu\) 是 G/H 上唯一的使得 G 上的测度 \(\nu^{\sharp}\) 等于 \(\kappa \cdot \mu\) 的测度。我们把空间 G/H 赋以测度 \(\nu\)（因而例如我们写 \(\mathcal{L}^{p}(\mathrm{G / H}) = \mathcal{L}^{p}(\mathrm{G / H},\nu)\)）。

 我们称 G 的子集 \(S \subset G\) 是模 H 零测的，如果 \(\varpi(S)\) 是 \(\nu\) -零测的。这一条件不依赖于 G 与 H 上哈尔测度的选取。它蕴含 S 是局部 \(\mu\) -零测的（INT, VII, 第 47 页, § 2, 第 3 目, 命题 6, a)）。我们称 G 的元素所具有的一个性质模 H 几乎处处为真，如果 G 中使该性质不成立的元素所成的集是模 H 零测的。

设 \(\pi\) 是 H 在希尔伯特空间 E 中的一个酉表示。用 \(\mathcal{F}_{\pi}(\mathrm{G})\) 表示 G 上取值于 E 的函数 \(f\) 所成的向量空间，其中 \(f(xh) = \pi(h)f(x)\) 对一切 \((x,h) \in \mathrm{G} \times \mathrm{H}\) 成立。我们称 \(\mathcal{F}_{\pi}(\mathrm{G})\) 的元素为 G 上的 \(\pi\) -等变函数。

G 在 C 上的平凡表示所对应的空间 \(\mathcal{F}_{1}(\mathrm{G})\) 与 G/H 上复值函数的空间 \(\mathcal{F}(\mathrm{G / H})\)，借助映射 \(f\mapsto f\circ \varpi\)（即 \(\mathcal{F}(\mathrm{G / H})\) 到 \(\mathcal{F}_1(\mathrm{G})\) 中的映射）等同。

对任意函数 \(f\)（属于 \(\mathcal{F}_{\pi}(\mathrm{G})\)），函数 \(\|f\|\) 属于 \(\mathcal{F}_{1}(\mathrm{G})\)，因为 \(\pi\) 是酉的。这使我们可以把 \(\|f\|\) 看作 \(\mathrm{G/H}\) 上的一个函数，例如我们将写作 \(\|f(x\mathrm{H})\|\)，以表示该函数在元素 \(x\mathrm{H}\)（\(\mathrm{G/H}\) 的元素）处的值。

设 \(f\) 是 \(\mathcal{F}_{\pi}(\mathrm{G})\) 中的一个函数；若 \(\|f\|\) 在 G/H 的一个紧子集外为零，则称它模 H 在一个紧集外为零。这等价于说，存在 G 的紧子集 K，使 \(f\) 在 \(\mathrm{K} \cdot \mathrm{H}\) 之外为零（TG, III, 第 33 页, 命题 10）。

用 \(\mathcal{K}_{\pi}(\mathrm{G})\) 表示 G 上属于 \(\mathcal{F}_{\pi}(\mathrm{G})\) 且模 H 具有紧支撑的连续函数所成的空间。

与 \(\mathcal{H}_{\pi}(\mathrm{G})\) 类似的一个空间，当 \(\pi\) 是 H 到 \(\mathbf{R}_{+}^{*}\) 中的连续同态时，已在 INT, VII, 第 39 页, § 2, 第 1 目 中定义过。

设 \(p \in [1, +\infty[\)。对每个 \(f \in \mathcal{F}_{\pi}(\mathrm{G})\)，记

\[
\mathrm{N} _ {p} (f) = \left(\int_ {\mathrm{G/H}} ^ {*} \| f \| ^ {p} d \nu\right) ^ {1 / p}.
\]

它是一个实数或 \(+\infty\)。用 \(\mathcal{F}_{\pi}^{p}(\mathrm{G},\nu)\)，或简作 \(\mathcal{F}_{\pi}^{p}(\mathrm{G})\)，表示由函数 \(f\in \mathcal{F}_{\pi}(\mathrm{G})\) 中使 \(\mathrm{N}_p(f)\)

为有限者所成的子空间。空间 \(\mathcal{F}_{\pi}^{p}(G)\) 赋以映射 \(N_{p}\) 后是一个半赋范空间。

空间 \(\mathcal{K}_{\pi}(\mathrm{G})\) 包含于 \(\mathcal{F}_{\pi}^{p}(\mathrm{G})\)。它在 \(\mathcal{F}_{\pi}^{p}(\mathrm{G})\) 中的闭包记作 \(\mathcal{L}_{\pi}^{p}(\mathrm{G},\nu)\)，或简作 \(\mathcal{L}_{\pi}^{p}(\mathrm{G})\)；我们称它为 G 上 \(\pi\) -等变函数中模 H 以指数 \(p^{e}\) 可积的那些函数所成的空间。

下面命题的各断言，当 \(\pi\) 是一维平凡表示时，是 INT, IV, 第 128 页, § 3, 第 3 目, 命题 6 及第 131 页, § 3, 第 4 目, 定理 3 的推论。

命题 3. --- a) 设 \((f_n)_{n \in \mathbf{N}}\) 是 \(\mathcal{F}_\pi^p(\mathrm{G})\) 中的一个序列，使得级数

\[
\sum_ {n = 0} ^ {+ \infty} \mathrm{N} _ {p} (f _ {n})
\]

收敛。于是以 \(f_n(g)\) 为通项的级数对位于一个模 H 零测集 T 之外的每个 \(g\) 绝对收敛。设 \(f\) 是 G 上取值于 E 的函数，它在 G−T 上等于该级数的和，而在 T 上等于 0。这时我们有 \(f \in \mathcal{F}_\pi^p(\mathrm{G})\)，且以 \(f_n\) 为通项的级数收敛于 \(f\)，这里收敛是在空间 \(\mathcal{F}_\pi^p(\mathrm{G})\) 中取的；

\begin{enumerate}
\def\labelenumi{\alph{enumi})}
\setcounter{enumi}{1}
\item
  设 \((f_n)\) 是 \(\mathcal{L}_{\pi}^{p}(\mathrm{G})\) 中收敛于一个函数 \(f\) 的序列。存在 \((f_{n_k})\)——它是 \((f_n)\) 的一个子列——使得 \(f_{n_k}(g)\) 收敛于 \(f(g)\)，对位于一个模 H 零测集之外的每个 \(g\) 成立；
\item
  半赋范空间 \(\mathcal{F}_{\pi}^{p}(\mathrm{G})\) 与 \(\mathcal{L}_{\pi}^{p}(\mathrm{G})\) 是完备的。
\end{enumerate}

在断言 \(a\) ) 中，说以 \(f_n\) 为通项的级数收敛于 \(f\)——这里是指在空间 \(\mathcal{F}_{\pi}^{p}(\mathrm{G})\) 中——意思是部分和的序列 \(f_0 + \cdots + f_n\) 收敛于 \(f\)，收敛是在 \(\mathcal{F}_{\pi}^{p}(\mathrm{G})\) 中取的。这时也称 \(f\) 是该级数的一个和。

我们来证明 \(a\) )。根据 INT, IV, 第 128 页, § 3, 第 3 目, 命题 6，存在一个 \(\nu\) -零测集 S ⊂ G/H，使得以 \(\|f_n(gH)\|\) 为通项的级数对 gH \(\notin\) . S 绝对收敛。此外，对于 gH \(\notin\) . S 等于该级数之和、对于 gH \(\in\) . S 为零的函数 h 满足 N \(_p(h)\) \textless{} + \(\infty\)。

集合 T = \(\overline{\varpi}(S)\) 是模 H 零测的。对每个 \(g \notin T\)，以 \(f_{n}(g)\) 为通项的级数在 E 中绝对收敛。令 \(f(g) = \sum f_{n}(g)\)（对 \(g \notin T\)），否则令 \(f(g) = 0\)。我们有 \(f \in \mathcal{F}_{\pi}(G)\)。注意 \(\|f(gH)\| \leqslant h(gH)\) 对所有 \(g \in G\) 成立，由此 \(N_{p}(f) \leqslant\)

\(N_{p}(h) < +\infty\)。于是我们有 \(f \in \mathcal{F}_{\pi}^{p}(G)\)。类似地，由此得出

\[
\mathrm{N} _ {p} \Big (f - \sum_ {n = 0} ^ {k} f _ {n} \Big) \leqslant \sum_ {n = k + 1} ^ {+ \infty} \mathrm{N} _ {p} (f _ {n})
\]

对所有 \(k \in \mathbf{N}\) 成立，因此级数 \(\sum f_n\) 收敛于 \(f\)，收敛是在 \(\mathcal{F}_{\pi}^{p}(\mathrm{G})\) 中取的。断言 \(a)\) 得证。

我们来证明 \(b\) )。序列 \((\|f_n - f\|)_{n \in \mathbf{N}}\) 在空间 \(\mathcal{L}^p(\mathrm{G}/\mathrm{H})\) 中收敛于 0。根据 INT, IV, 第 131 页, § 4, 第 3 目, 定理 3，存在一个子序列 \((\|f_{n_k} - f\|)_{k \in \mathbf{N}}\)，它 \(\nu\) -几乎处处收敛于 0。于是序列 \((f_{n_k}(g))_{k \in \mathbf{N}}\) 收敛于 \(f(g)\)，对位于一个模 H 零测集之外的每个 \(g\) 成立。

最后证明 \(c\) )。设 \((f_n)_{n \in \mathbf{N}}\) 是 \(\mathcal{F}_\pi^p(\mathrm{G})\) 中的一个柯西序列。存在严格递增的整数序列 \((n_k)_{k \in \mathbf{N}}\)，使得 \(\mathrm{N}_p(f_{n_{k+1}} - f_{n_k}) \leqslant 2^{-k}\) 对所有 \(k \in \mathbf{N}\) 成立。对每个 \(k \in \mathbf{N}\)，令 \(h_k = f_{n_{k+1}} - f_{n_k} \in \mathcal{F}_\pi^p(\mathrm{G})\)。由 \(a\) )，以 \(h_k\) 为通项的级数在 \(\mathcal{F}_\pi^p(\mathrm{G})\) 中收敛；记其和为 \(h\)。对每个 \(\ell \in \mathbf{N}\)，由此得出

\[
f _ {n _ {0}} + \sum_ {k = 0} ^ {\ell} h _ {k} = f _ {n _ {\ell + 1}},
\]

于是 \(f_{n_{0}} + h\) 是序列 \((f_{n})\) 的一个聚点。从而该序列收敛（TG, II, 第 14 页, 命题 5 的推论 2）。这样，空间 \(\mathcal{F}_{\pi}^{p}(\mathrm{G})\) 是完备的；又因空间 \(\mathcal{L}_{\pi}^{p}(\mathrm{G})\) 在 \(\mathcal{F}_{\pi}^{p}(\mathrm{G})\) 中是闭的，它也是完备的。

推论. --- 设 \((f_{n})_{n\in\mathbf{N}}\) 是 \(\mathcal{L}_{\pi}^{p}(\mathrm{G})\) 中的一个柯西序列，并设 \(f\in\mathcal{F}_{\pi}(\mathrm{G})\) 使得 \(f_{n}(g)\) 模 H 几乎处处收敛于 \(f(g)\)。则 \(f\in\mathcal{L}_{\pi}^{p}(\mathrm{G})\)，且 \((f_{n})_{n\in\mathbf{N}}\) 在 \(\mathcal{L}_{\pi}^{p}(\mathrm{G})\) 中收敛于 f。

事实上，函数 \(f\) 与序列 \((f_n)\) 在 \(\mathcal{L}_{\pi}^{p}(\mathrm{G})\) 中的极限模 H 几乎处处重合。

用 \(\mathrm{L}_{\pi}^{p}(\mathrm{G})\) 表示与半赋范空间 \(\mathcal{L}_{\pi}^{p}(\mathrm{G})\) 相伴随的分离赋范拓扑向量空间；它是一个巴拿赫空间。

设 \(f\) 是 \(\mathrm{G}\) 上取值于 \(\mathrm{E}\) 的一个函数。用 \(\mathrm{S}_{f}\) 表示由 \(x\in \mathrm{G}\) 中使函数 \(h\mapsto f(xh)\)（定义在 \(\mathrm{H}\) 上）不属于 \(\mathcal{L}_{\mathrm{E}}^{1}(\mathrm{H})\) 者所成的集。我们有 \(\mathrm{S}_f\cdot h = \mathrm{S}_f\) 对所有 \(h\in \mathrm{H}\) 成立。

设 \(x \in \mathrm{G} - \mathrm{S}_f\)。由于 \(\pi\) 是酉表示，函数 \(h \mapsto \pi(h)^* f(xh)\) 可测（它是 H 上的映射 \(h \mapsto (h, f(xh))\)——后者到 \(\mathrm{H} \times \mathrm{E}\) 中可测——与 H × E 到 E 中的连续映射 \((g,x)\mapsto\pi(g)^{*}x\) 的复合，参见 INT, IV, 第 174 页, § 5, 第 3 目, 定理 1），并且它在 H 上可积。

我们通过下式在 G 上定义一个函数 \(f^{\pi}\)：

\[
f ^ {\pi} (x) = \int_ {\mathrm{H}} \pi (h) ^ {*} f (x h) d \beta (h)\tag{3}
\]

其中 \(x \in \mathrm{G} - \mathrm{S}_f\)；而 \(f^\pi(x) = 0\)（当 \(x \in \mathrm{S}_f\) 时）。

命题 4. --- 设 \(f \in \mathcal{L}_{\mathrm{E}}^{1}(\mathrm{G})\)。

\begin{enumerate}
\def\labelenumi{\alph{enumi})}
\item
  集合 \(\mathrm{S}_f\) 模 H 零测，且 \(f^\pi \in \mathcal{F}_\pi(\mathrm{G})\)；
\item
  设 C 是 G 的一个紧子集。若 f 连续且支撑含于 C，则 \(S_{f}\) 为空，函数 \(f^{\pi}\) 属于 \(\mathcal{H}_{\pi}(\mathrm{G})\) 且其支撑含于 \(C \cdot H\)。
\end{enumerate}

\(a\) ) 的第一部分由 INT, VII, 第 57 页, § 2, 第 5 目, c) 得出。设 \(w \in H\)。若 \(x \in G - S_f\)，则 \(xw \in G - S_f\)，且

\[
\begin{array}{c} f ^ {\pi} (x w) = \int_ {\mathrm{H}} \pi (h) ^ {*} f (x w h) d \beta (h) \\ = \int_ {\mathrm{H}} \pi (w ^ {- 1} y) ^ {*} f (x y) d \beta (y) = \pi (w) f ^ {\pi} (x), \end{array}
\]

而若 \(x \in S_{f}\)，则 \(f^{\pi}(xw) = 0 = \pi(w)^{*}f^{\pi}(x)\)。因此函数 \(f^{\pi}\) 属于 \(\mathcal{F}_{\pi}(\mathrm{G})\)。

现在设 \(f\) 连续且支撑含于 C。于是对每个 \(x \in G\)，映射 \(h \mapsto f(xh)\) 属于 \(\mathcal{K}(H)\)，从而可积，这就证明了 \(S_f = \varnothing\)。

我们来证明 \(f^{\pi}\) 连续。设 \(x \in G\)。设 U 是 \(x\) 在 G 中一个相对紧的开邻域。

对每个 \(y \in \mathrm{U}\)，我们有

\[
\begin{array}{l} \| f ^ {\pi} (y) - f ^ {\pi} (x) \| \leqslant \int_ {\mathrm{H}} \| f (y h) - f (x h) \|   d \beta (h) \\ \qquad = \int_ {\mathrm{H} \cap (y ^ {- 1} \mathrm{C} \cup x ^ {- 1} \mathrm{C})} \| f (y h) - f (x h) \|   d \beta (h) \\ \qquad \leqslant \beta (\mathrm{H} \cap \mathrm{U} ^ {- 1} \mathrm{C}) \sup _ {h \in \mathrm{U} ^ {- 1} \mathrm{C}} \| f (y h) - f (x h) \|, \end{array}
\]

而 \(f^{\pi}\) 的连续性则由 f 在 G 上的一致连续性得出。

最后，若 \(x \in G\) 满足 \(f^{\pi}(x) \neq 0\)，则存在 \(h \in H\) 使 \(f(xh) \neq 0\)，于是 xh 属于 C，而 x 属于 C·H。由于 C·H 在 G 中闭，我们断定 \(f^{\pi}\) 的支撑含于 C·H。

 设 C 是 G 的一个紧子集。设 \(u \in \mathcal{K}_{+}(G)\) 是满足 \(u(x) > 0\)（对一切 \(x \in C\)）的函数。设 v 是 G 上由下式定义的函数

\[
v (x) = \int_ {\mathrm{H}} u (x h) d \beta (h)\tag{4}
\]

对所有 \(x \in G\) 成立；沿用前面的记号，我们有 \(v = u^{1}\)，它对应于 H 的一维平凡表示。函数 v 连续且为正；它属于 \(\mathcal{F}_{1}(G)\)，其支撑含于 \(\operatorname{Supp}(u) \cdot H\)，并且我们有

\[
\inf _ {x \in \mathrm{C} \cdot \mathrm{H}} v (x) > 0\tag{5}
\]

（INT, VII, 第 39--40 页, § 2, 第 1 目, 命题 1 及引理 1, a)）。

引理 5. --- 设 \(f \in \mathcal{F}_{\pi}(\mathrm{G})\) 是一个 \(\mu\) -可测且在 \(\mathrm{C} \cdot \mathrm{H}\) 之外为零的函数，使得函数 \(\|f\|\) \(\nu\) -可积（在 \(\mathrm{G}/\mathrm{H}\) 上）。设 s 是 \(\mathrm{G}\) 上取值于 E 的函数，满足

\[
 s (x) = \left\{ \begin{array}{l l} v (x) ^ {- 1} f (x) & \text{若 } x \in \mathrm{C} \cdot \mathrm{H} \\ 0 & \text{若 } x \in \mathrm{G} - \mathrm{C} \cdot \mathrm{H}. \end{array} \right.
\]

\begin{enumerate}
\def\labelenumi{\alph{enumi})}
\item
  函数 s 是 \(\mu\) -可测的；它属于 \(\mathcal{F}_{\pi}(\mathrm{G})\) 且在 C·H 之外为零；
\item
  函数 us 属于 \(\mathcal{L}^1 (\mathrm{G})\)，且模 H 几乎处处有 \((us)^{\pi} = f\)。
\end{enumerate}

函数 \(s\) 按定义在 C \(\cdot\) . H 之外为零；它是 \(\mu\) -可测的（INT, IV, 第 193 页, § 5, 第 10 目, 命题 16），因为函数 \(f\) \(\mu\) -可测且 \(v(x) > 0\) 对一切 \(x \in \mathrm{C} \cdot \mathrm{H}\) 成立。我们有 \(s \in \mathcal{F}_{\pi}(\mathrm{G})\)，因为 \(v \in \mathcal{F}_1(\mathrm{G})\)。

函数 \(f\) 是局部 \(\mu\) -可积的，因为 \(\|f\|\) 是 \(\nu\) -可积的函数（INT, VII, 第 47 页, § 2, 第 3 目, 命题 6, c)，注意测度 \(\nu^{\sharp} = \kappa \cdot \mu\) 等价于 \(\mu\)），从而由公式 (5)，函数 \(s\) 亦然。函数 \(us\) 可测且在 G 的一个紧子集之外为零；由于 \(u\) 有界，由此得出 \(us\) 是 \(\mu\) -可积的。

对每个 \(x\in \mathrm{G - S}_{us}\)，我们有

\[
\begin{array}{c} (u s) ^ {\pi} (x) = \int_ {\mathrm{H}} \pi (h) ^ {*} (u (x h) s (x h)) d \beta (h) \\ = \int_ {\mathrm{H}} u (x h) s (x) d \beta (h) = v (x) s (x), \end{array}
\]

因为 \(s(xh) = \pi(h)s(x)\)；最后一个断言由命题 4, a) 得出。

当 \(H = \{e\}\) 且 \(\pi\) 是一维时，下面的命题就是 INT, IV, 第 184 页, § 5, 第 6 目, 定理 5。

命题 5. --- 设 \(p \in [1, +\infty[\)。空间 \(\mathcal{L}_{\pi}^{p}(\mathrm{G})\) 是由满足如下条件的函数 \(f \in \mathcal{F}_{\pi}(\mathrm{G})\) 所成的空间：\(f\) \(\mu\) -可测，且函数 \(\|f\|\) 属于 \(\mathcal{L}^{p}(\mathrm{G}/\mathrm{H})\)。

设 \(f \in \mathcal{L}_{\pi}^{p}(G)\)。它是 \(\mathcal{L}_{\pi}^{p}(G)\) 中由 \(\mathcal{K}_{\pi}(G)\) 的元素组成的序列的极限。由命题 3, b) 及叶戈罗夫定理（INT, IV, 第 175 页, § 5, 第 4 目, 定理 2），函数 \(f\) 因此是 \(\mu\) -可测的；从而 G/H 上的函数 \(\|f\|\) 是 \(\nu\) -可测的（INT, VII, 第 47 页, § 2, 第 3 目, 命题 6, b)）。由于 \(N_{p}(f)\) 有限，函数 \(\|f\|\) 属于 \(\mathcal{L}^{p}(G/H)\)（INT, IV, 第 184 页, § 5, 第 6 目, 定理 5）。

 我们来证明逆命题。设 \(f\) 是一个 \(\mu\) -可测、属于 \(\mathcal{F}_{\pi}(\mathrm{G})\) 的函数，且满足 \(\|f\|\in\mathcal{L}^{p}(\mathrm{G}/\mathrm{H})\)。设 \(\varepsilon>0\)。我们来证明存在 \(\widetilde{f}\in\mathcal{K}_{\pi}(\mathrm{G})\)，使得

\[
\mathrm{N} _ {p} (f - \widetilde {f}) ^ {p} = \int_ {\mathrm{G/H}} ^ {*} \| f - \widetilde {f} \| ^ {p} d \nu \leqslant \varepsilon ,
\]

如此即完成证明。

 先设存在 G 的紧子集 C，使 \(f\) 在 \(\mathrm{C} \cdot \mathrm{H}\) 之外为零。设 \(u \in \mathcal{K}_{+}(\mathrm{G})\) 是满足 \(u(x) > 0\)（对一切 \(x \in \mathrm{C}\)）的函数，并按上式定义 \(v = u^{1}\)。用 \(\varphi\) 表示 \(u\) 的支撑的特征函数。

设 q 是 p 的共轭指数。若 p = 1，设 w 是 G 上等于 \(\sup_{x \in G} u(x)\) 的常值函数。若 p \textgreater{} 1，我们定义

\[
w (x) = \left(\int_ {\mathrm{H}} u (x h) ^ {q} d \beta (h)\right) ^ {1 / q}
\]

 对所有 \(x \in G\) 成立。在所有情形下，函数 w 都连续且为正；它属于 \(\mathcal{K}_{1}(G)\)（命题 4 应用于一维平凡表示及函数 \(u^{q}\)），从而在 G 上有界。我们令 \(W = \sup_{x \in G} w(x)\)。

 设 \(s\) 是 \(\mathbf{G}\) 上取值于 \(\mathbf{E}\) 的、把引理 5 应用于 \(f\) 时所得到的函数。令 \(g = \varphi s\)。函数 \(g\) 是 \(\mu\) -可测的，且满足 \(\| g \| \leqslant \| f \| / \inf_{x \in \mathrm{C} \cdot \mathrm{H}} v(x)\)。由于 \(g\) 在 \(u\) 的支撑之外为零且 \(\kappa\) 连续，我们有 \(g \in \mathcal{L}_{\mathrm{E}}^{p}(\mathrm{G}, \kappa \cdot \mu)\)。设 \(\widetilde{g}\) 是

\(\mathcal{K}_{\mathrm{E}}(\mathrm{G})\) 中的一个函数，使得

\[
\int_ {\mathrm{G}} ^ {*} \| g - \widetilde {g} \| ^ {p} \kappa d \mu \leqslant \frac {\varepsilon}{\mathrm{W} ^ {p}}.
\]

我们有 us = ug，从而模 H 几乎处处有 \(f = (us)^{\pi} = (ug)^{\pi}\)（引理 5, b)）。令 \(\widetilde{f} = (u\widetilde{g})^{\pi}\)；我们有 \(\widetilde{f} \in \mathcal{K}_{\pi}(\mathrm{G})\)（命题 4, b)）。对每个 \(x \in G - S_{ug}\)，我们得到

\[
\| (u g) ^ {\pi} (x) - (u \widetilde {g}) ^ {\pi} (x) \| \leqslant \int_ {\mathrm{H}} ^ {*} u (x h) \| g (x h) - \widetilde {g} (x h) \| d \beta (h)
\]

由此得

\[
\| (u g) ^ {\pi} (x) - (u \widetilde {g}) ^ {\pi} (x) \| ^ {p} \leqslant w (x) ^ {p} \int_ {\mathrm{H}} ^ {*} \| g (x h) - \widetilde {g} (x h) \| ^ {p} d \beta (h),
\]

这里在 \( p \) \textgreater{} 1 的情形用了赫尔德不等式。由于 \(S_{ug}\) 模 H 零测，由此得出

\[
\begin{array}{c} \int_ {\mathrm{G/H}} ^ {*} \| f - \widetilde {f} \| ^ {p} d \nu \leqslant \mathrm{W} ^ {p} \int_ {\mathrm{G/H}} ^ {*} \Big (\int_ {\mathrm{H}} ^ {*} \| g (x h) - \widetilde {g} (x h) \| d \beta (h) \Big) ^ {p} d \nu (x \mathrm{H}) \\ = \mathrm{W} ^ {p} \int_ {\mathrm{G}} ^ {*} \| g - \widetilde {g} \| ^ {p} \kappa d \mu \leqslant \varepsilon \end{array}
\]

这是根据 INT, VII, 第 46 页, § 2, 第 3 目, 命题 5, b)，它适用于 \(g - \widetilde{g}\) 在 G 的一个紧子集之外为零的情形。这就证明了当 \(f\) 模 H 在一个紧子集之外为零时所要求的性质。

现在考虑一般情形。由假设 \(\|f\|\in\mathcal{L}^{p}(G/H)\)，存在 G/H 的紧子集 L，使得

\[
\int_ {(\mathrm{G} / \mathrm{H}) - \mathrm{L}} ^ {*} \| f \| ^ {p} d \nu \leqslant \frac {\varepsilon}{2}
\]

（参见 INT, IV, 第 152 页, § 4, 第 6 目, 定理 4）。设 \(\varphi_{\mathrm{L}}\) 是 L 的特征函数，并令 \(f_{\mathrm{L}} = (\varphi_{\mathrm{L}} \circ \varpi)f\)。我们有

\[
\int_ {\mathrm{G} / \mathrm{H}} ^ {*} \| f - f _ {\mathrm{L}} \| ^ {p} d \nu = \int_ {(\mathrm{G} / \mathrm{H}) - \mathrm{L}} ^ {*} \| f \| ^ {p} d \nu \leqslant \frac {\varepsilon}{2}.
\]

函数 \(f_{\mathrm{L}}\) 是 \(\mu\) -可测的且模 H 在一个紧集之外为零。它属于 \(\mathcal{L}_{\pi}^{p}(\mathrm{G})\)，于是由前一种情形，存在 \(\widetilde{f} \in \mathcal{K}_{\pi}(\mathrm{G})\) 使得 \(\mathrm{N}_p(f_{\mathrm{L}} - \widetilde{f}) \leqslant (\frac{\varepsilon}{2})^{1/p}\)，从而

\[
\int_ {\mathrm{G/H}} ^ {*} \| f - \widetilde {f} \| ^ {p} d \nu \leqslant \varepsilon ,
\]

这正是所要证的。

 我们来考虑 \(p = 2\) 的情形。对 \(f_1\) 与 \(f_2\)（均属于 \(\mathcal{F}_{\pi}(\mathrm{G})\)），函数 \(x \mapsto \langle f_1(x) | f_2(x) \rangle\) 属于 \(\mathcal{F}_1(\mathrm{G})\)，因为表示 \(\pi\) 是酉的，从而通过取商定义了 \(\mathrm{G / H}\) 上的一个函数，我们和先前一样把它等同于 \(\langle f_1 | f_2 \rangle\)。我们有不等式 \(|\langle f_1 | f_2 \rangle| \leqslant \| f_1 \| \| f_2 \|\)（在 \(\mathcal{F}(\mathrm{G / H})\) 中）。

若 \(f_1\) 与 \(f_2\) 属于 \(\mathcal{L}_\pi^2(\mathrm{G})\)，则函数 \(\langle f_1 | f_2 \rangle\) 属于 \(\mathcal{L}_1^1(\mathrm{G})\)。特别地，它在 \(\mathrm{G/H}\) 上可积。映射

\[
(f _ {1}, f _ {2}) \mapsto \int_ {\mathrm{G} / \mathrm{H}} \langle f _ {1} | f _ {2} \rangle d \nu
\]

是 \(\mathcal{L}_{\pi}^{2}(\mathrm{G})\) 上的一个正埃尔米特形式；其相应的半范数是半范数 \(N_{2}\)。特别地，空间 \(\mathcal{L}_{\pi}^{2}(\mathrm{G})\) 是一个预希尔伯特空间，而 \(\mathrm{L}_{\pi}^{2}(\mathrm{G})\) 是与 \(\mathcal{L}_{\pi}^{2}(\mathrm{G})\) 相伴随的希尔伯特空间。

\subsection{诱导表示}

我们沿用前一日中关于 H 上的测度 \(\beta\) 与 G/H 上的测度 \(\nu\) 的记号和约定，以及函数 \(\kappa\colon\mathrm{G}\to\mathbf{R}_{+}^{*}\)。

存在一个连续函数 \(\eta\)，它定义在 \(\mathrm{G} \times \mathrm{G} / \mathrm{H}\) 上且取值于 \(\mathbf{R}_{+}^{*}\) 中，使得

\[
\eta (x, y \mathrm{H}) = \frac {\kappa (x y)}{\kappa (x)}
\]

对所有 \((x,y)\in\mathrm{G}\times\mathrm{G}\) 成立，且 \(\boldsymbol{\gamma}_{\mathrm{G}/\mathrm{H}}(x)\nu=(y\mapsto\eta(x^{-1},y))\cdot\nu\) 对 \(x\in G\) 成立（INT, VII, 第 56 页, § 2, 第 5 目, 定理 2, c)）。

设 \(p \in [1, +\infty[\)，并设 \(\pi\) 是 H 在复希尔伯特空间 E 中的一个酉表示。

引理 6. --- 设 \(f \in \mathcal{K}_{\overline{\pi}}(\mathrm{G})\)。对每个 \(g \in \mathrm{G}\)，函数

\[
\widetilde {f} \colon x \mapsto \eta (g ^ {- 1}, x \mathrm{H}) ^ {1 / p} f (g ^ {- 1} x)
\]

（它映 G 到 E 中）属于 \(\mathcal{K}_{\overline{\pi}}(\mathrm{G})\)，且满足 \(\mathrm{N}_p(\widetilde{f}) = \mathrm{N}_p(f)\)。

可以毫无困难地验证 \(\tilde{f} \in \mathcal{K}_{\overline{\pi}}(\mathrm{G})\)。由于

\[
\mathrm{N} _ {p} (\widetilde {f}) ^ {p} = \int_ {\mathrm{G/H}} ^ {*} \| \widetilde {f} \| ^ {p} d \nu = \int_ {\mathrm{G/H}} ^ {*} \gamma_ {\mathrm{G/H}} (g) (\| f \| ^ {p}) (y) \eta (g ^ {- 1}, y) d \nu (y)
\]

 以及 \((y\mapsto \eta (g^{-1},y))\cdot \nu = \pmb{\gamma}_{\mathrm{G / H}}(g)\nu\)，我们得到 \(\mathrm{N}_p(\widetilde{f}) = \mathrm{N}_p(f)\)。

 由这个引理可知，存在一个连续等距表示 \(\widetilde{\pi}\)，它从 G 映到 \(\mathrm{L}_{\overline{\pi}}^{p}(\mathrm{G})\) 中，使得对 \(f\in \mathcal{K}_{\pi}(\mathrm{G})\) 与 \(g\in \mathrm{G}\)，元素 \(\widetilde{\pi}(g)f\) 是上面定义的函数 \(\widetilde{f}\) 的类。若 \(p=2\)，则该表示是酉表示。

定义 2. --- 我们称如此定义的 G 在空间 \(L_{\overline{\pi}}^{2}(G)\) 中的酉表示为 G 的由 H 的表示 \(\pi\) 相对于 \(\kappa\) 所诱导的酉表示。把它记作 \(\operatorname{Ind}_{\mathrm{H}}^{\mathrm{G}}(\pi, \kappa)\)，或简作 \(\operatorname{Ind}_{\mathrm{H}}^{\mathrm{G}}(\pi)\)。

 注. --- 1) 设 \(\varrho\) 是 H 在复希尔伯特空间 F 中的一个酉表示，并设 \(u\colon\pi\to\varrho\) 是一个 H-态射。对每个函数 \(f\in\mathcal{K}_{\overline{\pi}}(\mathrm{G})\)，用 \(v(f)\) 表示 G 到 F 中的函数 \(g\mapsto u(f(g))\)；它属于 \(\mathcal{K}_{\overline{\varrho}}(\mathrm{G})\)，且满足 \(\mathrm{N}_{p}(v(f))\leqslant\|u\|\mathrm{N}_{p}(f)\)。由 \(\mathcal{K}_{\overline{\pi}}(\mathrm{G})\) 到 \(\mathcal{K}_{\overline{\varrho}}(\mathrm{G})\) 中的把 \(f\) 映为 \(v(f)\) 的线性映射因此延拓为 \(\mathrm{Ind}_{\mathrm{H}}^{\mathrm{G}}(\pi)\) 到 \(\mathrm{Ind}_{\mathrm{H}}^{\mathrm{G}}(\varrho)\) 中的一个连续 H-态射，记作 \(\mathrm{Ind}_{\mathrm{H}}^{\mathrm{G}}(u)\)。我们有 \(\mathrm{Ind}_{\mathrm{H}}^{\mathrm{G}}(1_{\pi})=1_{\mathrm{Ind}_{\mathrm{H}}^{\mathrm{G}}(\pi)}\)。设 \(\sigma\) 是 H 的一个酉表示，并设 \(v\colon\varrho\to\sigma\) 是一个 H-态射；我们有 \(\mathrm{Ind}_{\mathrm{H}}^{\mathrm{G}}(v\circ u)=\mathrm{Ind}_{\mathrm{H}}^{\mathrm{G}}(v)\circ\mathrm{Ind}_{\mathrm{H}}^{\mathrm{G}}(u)\)。

换言之，把 \(\pi\) 映为 \(\mathrm{Ind}_{\mathrm{H}}^{\mathrm{G}}(\pi)\)、把 \(u\) 映为 \(\mathrm{Ind}_{\mathrm{H}}^{\mathrm{G}}(u)\) 的构造是一个从 H 的酉表示的范畴到 G 的酉表示的范畴的函子（参见 CAT, I, § 2, 编纂中）。

\begin{enumerate}
\def\labelenumi{\arabic{enumi})}
\setcounter{enumi}{1}
\tightlist
\item
  设 \(\kappa^{\prime}\colon \mathrm{G}\to \mathbf{R}_{+}^{*}\)，使得
\end{enumerate}

\[
\frac {\kappa^ {\prime} (x h)}{\kappa^ {\prime} (x)} = \frac {\Delta_ {\mathrm{H}} (h)}{\Delta_ {\mathrm{G}} (h)} = \frac {\kappa (x h)}{\kappa (x)}
\]

 对所有 \((x,h)\in\mathrm{G}\times\mathrm{H}\) 成立。设 \(\nu^{\prime}\) 是 G/H 上的拟不变测度 \((\kappa^{\prime}\cdot\mu)/\beta\)。函数 \(\kappa^{\prime}\kappa^{-1}\) 通过取商定义了一个连续函数 \(\xi\colon G/H\to R_{+}^{*}\)，使得 \(\nu^{\prime}=\xi\cdot\nu\)。用 \(\alpha\) 表示 \(\mathcal{K}_{\pi}(G)\) 的由 \(f\mapsto(\kappa^{\prime}\kappa^{-1})^{1/p}f\) 定义的自同态；它满足

\[
\int_ {\mathrm{G} / \mathrm{H}} ^ {*} \| f \| ^ {p} d \nu^ {\prime} = \int_ {\mathrm{G} / \mathrm{H}} ^ {*} \| \alpha (f) \| ^ {p} d \nu
\]

并且它使我们可以把空间 \(\mathcal{L}_{\pi}^{p}(\mathrm{G},\nu^{\prime})\) 与 \(\mathcal{L}_{\pi}^{p}(\mathrm{G},\nu)\)、以及空间 \(\mathrm{L}_{\pi}^{p}(\mathrm{G},\nu^{\prime})\) 与 \(\mathrm{L}_{\pi}^{p}(\mathrm{G},\nu)\) 等同起来。此外，\(\alpha\) 诱导表示 \(\mathrm{Ind}_{\mathrm{H}}^{\mathrm{G}}(\pi,\kappa^{\prime})\) 与 \(\mathrm{Ind}_{\mathrm{H}}^{\mathrm{G}}(\pi,\kappa)\) 的一个等距同构。

\subsection{闭中心子群的情形}

在本目中，我们假设群 G 是单模的。

 考虑 G 的中心的一个闭子群 Z，并用 dz 表示 Z 上的一个哈尔测度。我们给商群 G/Z 赋以哈尔测度 \(\nu = \mu/dz\)。设 \(\chi\) 是 Z 的一个酉特征标。商群 G/Z 是单模的（INT, VII, 第 61 页, § 2, 第 7 目, 命题 11 的推论）。

引理 7. --- 我们有

\[
\mathrm{N} _ {1} (\check {f}) = \mathrm{N} _ {1} (f), \quad \langle f _ {1} | f _ {2} \rangle = \langle \check {f} _ {1} | \check {f} _ {2} \rangle .\tag{6}
\]

对所有 \(f \in \mathcal{F}_{\chi}(\mathrm{G})\) 以及一切 \(f_1\) 与 \(f_2\)（均属于 \(\mathcal{L}_{\chi}^{2}(\mathrm{G})\)）成立。

这些公式都是定义的直接推论。

 对每个 \(g \in \mathbf{G}\) 与每个 \(f \in \mathcal{F}_{\chi}(\mathbf{G})\)，函数 \(x \mapsto f(g^{-1}x)\) 与 \(x \mapsto f(xg)\) 都属于 \(\mathcal{F}_{\chi}(\mathbf{G})\)。分别把它们记作 \(\boldsymbol{\gamma}_{\mathrm{G},\chi}(g)f\) 与 \(\boldsymbol{\delta}_{\mathrm{G},\chi}(g)f\)。映射 \(\boldsymbol{\gamma}_{\mathrm{G},\chi}\) 与 \(\boldsymbol{\delta}_{\mathrm{G},\chi}\) 是 \(\mathbf{G}\) 在 \(\mathcal{F}_{\chi}(\mathbf{G})\) 中的线性表示。对每个 \(z \in \mathbf{Z}\)，我们有

\[
\boldsymbol {\gamma} _ {\mathrm{G}, \chi} (g z) = \overline {{\chi (z)}} \boldsymbol {\gamma} _ {\mathrm{G}, \chi} (g), \quad \boldsymbol {\delta} _ {\mathrm{G}, \chi} (g z) = \chi (z) \boldsymbol {\delta} _ {\mathrm{G}, \chi} (g).\tag{7}
\]

设 \(f\in \mathcal{F}_{\chi}(\mathrm{G})\) 且 \(g\in \mathrm{G}\)；我们有

\[
| \boldsymbol {\gamma} _ {\mathrm{G}, \chi} (g) f | = \boldsymbol {\gamma} _ {\mathrm{G} / \mathrm{Z}} (g \mathrm{Z}) | f |, \quad | \boldsymbol {\delta} _ {\mathrm{G}, \chi} (g) f | = \boldsymbol {\delta} _ {\mathrm{G} / \mathrm{Z}} (g \mathrm{Z}) | f |,\tag{8}
\]

其中出现在这些等式中的所有函数都等同于 G/Z 上的函数。

 子空间 \(\mathcal{K}_{\chi}(\mathrm{G})\)（\(\mathcal{F}_{\chi}(\mathrm{G})\) 的）在 \(\gamma_{\mathrm{G},\chi}\) 与 \(\delta_{\mathrm{G},\chi}\) 下稳定。设 \(p\) 是一个实数 \(\geqslant 1\)。公式 (8) 蕴含：表示 \(\gamma_{\mathrm{G},\chi}\) 与 \(\delta_{\mathrm{G},\chi}\) 在限制到 \(\mathcal{K}_{\chi}(\mathrm{G})\) 上后可延拓为 G 到 \(\mathcal{L}_{\chi}^{p}(\mathrm{G})\) 中的连续等距线性表示，分别记作 \(\gamma_{\mathrm{G},\chi}^{(p)}\) 与 \(\delta_{\mathrm{G},\chi}^{(p)}\)。通过取商，这些表示也在 \(\mathrm{L}_{\chi}^{p}(\mathrm{G})\) 中定义了 G 的等距表示，仍以同样的记号记之。

 表示 \(\gamma_{\mathrm{G},\chi}^{(2)}\) 与 \(\delta_{\mathrm{G},\chi}^{(2)}\)（\(\mathrm{L}_{\chi}^{2}(\mathrm{G})\) 中的）是酉表示，将分别简记为 \(\gamma_{\mathrm{G},\chi}\) 与 \(\delta_{\mathrm{G},\chi}\)——当不致与 \(\mathcal{F}_{\chi}(\mathrm{G})\) 中的表示混淆时。我们也用 \(\varrho_{\mathrm{G},\chi}\) 表示 \(\mathrm{G} \times \mathrm{G}\) 在 \(\mathcal{L}_{\chi}^{2}(\mathrm{G})\) 或 \(\mathrm{L}_{\chi}^{2}(\mathrm{G})\) 中的由下式定义的连续表示

\[
\pmb {\varrho} _ {\mathrm{G}, \chi} (g, h) = \pmb {\gamma} _ {\mathrm{G}, \chi} (g) \circ \pmb {\delta} _ {\mathrm{G}, \chi} (h) = \pmb {\delta} _ {\mathrm{G}, \chi} (h) \circ \pmb {\gamma} _ {\mathrm{G}, \chi} (g)
\]

对所有 \((g,h)\in \mathrm{G}\times \mathrm{G}\) 成立（参见 V, 第 377 页, 引理 1）。

\(\gamma_{\mathrm{G},\chi}\) 在 \(\mathrm{L}_{\chi}^{2}(\mathrm{G})\) 中的表示不是别的，正是诱导表示 \(\mathrm{Ind}_{\mathrm{Z}}^{\mathrm{G}}(\overline{\chi})\)。

当 \(Z = \{e\}\) 时，由于 G 单模，下面的引理可由 INT, VIII, 第 166 页, § 4, 第 5 目, 命题 12 得出。

引理 8. --- 设 \(f_1 \in \mathcal{K}_\chi(\mathrm{G})\) 且 \(f_2 \in \mathcal{L}_\chi^2(\mathrm{G})\)。

\begin{enumerate}
\def\labelenumi{\alph{enumi})}
\item
  函数 \(f\)——它定义在 \(\mathrm{G}\) 上，由 \(f(g) = \langle f_1 | \boldsymbol{\gamma}_{\mathrm{G},\chi}(g)f_2 \rangle\)（对所有 \(g \in \mathrm{G}\)）给出——属于 \(\mathcal{L}_{\chi}^{2}(\mathrm{G})\)，且满足 \(\mathrm{N}_2(f) \leqslant \mathrm{N}_1(f_1)\mathrm{N}_2(f_2)\)；
\item
  函数 \(f\)——它定义在 \(G\) 上，由 \(f(g) = \langle f_1 | \delta_{G,\chi}(g)f_2 \rangle\)（对所有 \(g \in G\)）给出——属于 \(\mathcal{L}_{\chi}^{2}(G)\)，且满足 \(N_2(f) \leqslant N_1(f_1)N_2(f_2)\)。
\end{enumerate}

我们来证明 \(a\) )，断言 \(b\) ) 的证明是类似的。函数 \(f\) 连续，因此 \(\mu\) -可测。对每个 \(z \in \mathbf{Z}\) 与每个 \(g \in \mathbf{G}\)，我们有

\[
f (g z) = \langle f _ {1} | \gamma_ {\mathrm{G}, \chi} (g z) f _ {2} \rangle = \overline {{\chi (z)}} \langle f _ {1} | \gamma_ {\mathrm{G}, \chi} (g) f _ {2} \rangle
\]

（公式 (7), 第 417 页），因此 \(f \in \mathcal{F}_{\overline{\chi}}(\mathrm{G})\)。V 第 413 页的命题 5 蕴含现在只需证明 \(\mathrm{N}_2(f) \leqslant \mathrm{N}_1(f_1)\mathrm{N}_2(f_2)\)。

先设 \(f_2\) 属于 \(\mathcal{K}_{\chi}(\mathrm{G})\)。对每个 \(g \in \mathrm{G}\)，按定义有

\[
f (g) = \int_ {\mathrm{G/Z}} \overline {{f _ {1}}} \gamma_ {\mathrm{G}, \chi} (g) f _ {2} d \nu ,
\]

其中函数 \(\overline{f_1} \gamma_{\mathrm{G},\chi}(g)f_2\) 等同于 G/Z 上的一个函数。

我们定义一个函数 \(f_3\)（在 \(G\) 上）：令 \(f_3(g) = 0\)（若 \(f_1(g) = 0\)），否则令 \(f_3(g) = f_1(g)|f_1(g)|^{-1/2}\)。函数 \(f_3\) 属于 \(\mathcal{F}_{\chi}(G)\) 且满足 \(f_1 = |f_1|^{1/2}f_3\)；它是 \(\mu\) -可测且模 \(Z\) 在一个紧集外为零的，因为 \(f_1\) 如此。由于 \(|f_1|^{1/2} \in \mathcal{K}_1(G)\)，由此得出

\[
\overline {{f _ {1}}} \gamma_ {\mathrm{G}, \chi} (g) f _ {2} = | f _ {1} | ^ {1 / 2} \overline {{f _ {3}}} \gamma_ {\mathrm{G}, \chi} (g) f _ {2},
\]

\[
| \overline {{f _ {3}}} \gamma_ {\mathrm{G}, \chi} (g) f _ {2} | = | f _ {3} | | \gamma_ {\mathrm{G}, \chi} (g) f _ {2} |
\]

对所有 \(g \in \mathbf{G}\) 成立。

 设 \(g \in G\)。由柯西--施瓦茨不等式，我们有

\[
\begin{array}{l} | f (g) | ^ {2} \leqslant \Big (\int_ {\mathrm{G/Z}} | f _ {1} | ^ {1 / 2} | f _ {3} | | \boldsymbol {\gamma} _ {\mathrm{G}, \chi} (g) f _ {2} | d \nu \Big) ^ {2} \\ \leqslant \Big (\int_ {\mathrm{G/Z}} | f _ {1} | d \nu \Big) \Big (\int_ {\mathrm{G/Z}} | f _ {3} | ^ {2} | \boldsymbol {\gamma} _ {\mathrm{G}, \chi} (g) f _ {2} | ^ {2} d \nu \Big). \end{array}
\]

由于有 \(|f_{3}|^{2}=|f_{1}|\)（在 \(\mathcal{K}(\mathrm{G}/\mathrm{Z})\) 中），我们通过在 G/Z 上积分得到

\[
\int_ {\mathrm{G} / \mathrm{Z}} | f (g) | ^ {2} d \nu (g) \leqslant \mathrm{N} _ {1} (f _ {1}) \int_ {\mathrm{G} / \mathrm{Z}} \left(\int_ {\mathrm{G} / \mathrm{Z}} | f _ {1} (x) | | f _ {2} (g ^ {- 1} x) | ^ {2} d \nu (x)\right) d \nu (g).
\]

由函数

\[
(g, x) \mapsto | f _ {1} (x) | | f _ {2} (g ^ {- 1} x) | ^ {2}
\]

通过取商所得到的 G/Z × G/Z 上的函数是 \((\nu \otimes \nu)\) 可测且具紧支撑的，从而是 \((\nu \otimes \nu)\) -缓增的（INT, V, 第 4 页, § 5, 第 2 目, 定义 2）。由 INT, V, 第 93 页, § 8, 第 3 目, 命题 7 得

\[
\begin{array}{l} \int_ {\mathrm{G} / \mathrm{Z}} \Bigl (\int_ {\mathrm{G} / \mathrm{Z}} | f _ {1} (x) |   | f _ {2} (g ^ {- 1} x) | ^ {2} d \nu (x) \Bigr) d \nu (g) \\ = \int_ {\mathrm{G} / \mathrm{Z}} | f _ {1} (x) | \Bigl (\int_ {\mathrm{G} / \mathrm{Z}} | f _ {2} (g ^ {- 1} x) | ^ {2} d \nu (g) \Bigr) d \nu (x) = \mathrm{N} _ {1} (f _ {1}) \mathrm{N} _ {2} (f _ {2}) ^ {2}. \end{array}
\]

于是我们有 \(\mathrm{N}_{2}(f)^{2}\leqslant\mathrm{N}_{1}(f_{1})^{2}\mathrm{N}_{2}(f_{2})^{2}\)，这就在此情形下建立了所要的性质。

 我们来考虑一般情形。用 \(u\) 表示 \(\mathcal{K}_{\chi}(\mathrm{G})\) 到 \(\mathcal{L}_{\overline{\chi}}^{2}(\mathrm{G})\) 中的把 \(f_{2}\) 映为 \(f\) 的线性映射。设 \(f_{2} \in \mathcal{L}_{\chi}^{2}(\mathrm{G})\)，并设 \((f_{2,n})_{n \in \mathbf{N}}\) 是 \(\mathcal{K}_{\chi}(\mathrm{G})\) 中的收敛于 \(f_{2}\)（在 \(\mathcal{L}_{\chi}^{2}(\mathrm{G})\) 中）的序列。令 \(f_{n} = u(f_{2,n})\)；序列 \((f_{n})_{n \in \mathbf{N}}\) 在 \(\mathcal{L}_{\overline{\chi}}^{2}(\mathrm{G})\) 中是柯西序列，因为前一种情形蕴含 \(\mathrm{N}_{2}(f_{n} - f_{m}) \leqslant \mathrm{N}_{1}(f_{1})\mathrm{N}_{2}(f_{2,n} - f_{2,m})\) 对一切 \(n\) 与 \(m\)（\(\mathbf{N}\) 中的）成立。设 \(f \in \mathcal{L}_{\overline{\chi}}^{2}(\mathrm{G})\) 使得 \((f_{n})\) 收敛于 \(f\)（V, 第 409 页, 命题 3, \(c\) )）。由于 \(\mathrm{N}_{2}(f_{n}) \leqslant \mathrm{N}_{1}(f_{1})\mathrm{N}_{2}(f_{2,n})\) 对所有 \(n \in \mathbf{N}\) 成立，由此得出 \(\mathrm{N}_{2}(f) \leqslant \mathrm{N}_{1}(f_{1})\mathrm{N}_{2}(f_{2})\)。

存在一个子序列 \((f_{n_{k}})_{k\in\mathbf{N}}\)，使得 \(f_{n_{k}}(g)\) 收敛于 \(f(g)\)，对所有位于 G 中模 Z 零测子集之外的 \(g\in G\) 成立（出处同上, b)）。但另一方面，对每个 \(g\in G\)，我们有

\[
f _ {n _ {k}} (g) = \left\langle f _ {1} \mid \gamma_ {\mathrm{G}, \chi} (g) f _ {2, n _ {k}} \right\rangle\rightarrow \left\langle f _ {1} \mid \gamma_ {\mathrm{G}, \chi} (g) f _ {2} \right\rangle .
\]

于是我们有 \(f(g) = \langle f_1 | \gamma_{\mathrm{G},\chi}(g)f_2 \rangle\)，对所有位于 G 中模 Z 零测子集之外的 \(g\) 成立。由于 \(f \in \mathcal{L}_{\overline{\chi}}^2(\mathrm{G})\) 且 \(\mathrm{N}_2(f) \leqslant \mathrm{N}_1(f_1)\mathrm{N}_2(f_2)\)，引理得证。

\subsection{平方可积表示}

在本目中，所考虑的希尔伯特空间都是复的。假设 G 是单模的，其中心记作 C。对 C 的每个闭子群 Z，用 \(\beta_{Z}\) 表示 Z 上的一个哈尔测度，并给 G/Z 赋以哈尔测度 \(\nu_{Z} = \mu/\beta_{Z}\)（INT, VII, 第 44 页, § 2, 第 2 目, 定义 1）。

设 \(\pi\) 是 G 在希尔伯特空间 E 中的一个不可约酉表示，\(\chi \in \widehat{\mathbf{C}}\) 是其中心特征标（V, 第 390 页, 定义 6）。对 E 中所有 \(x\) 与 \(y\)，用 \(f_{x,y}\) 表示矩阵系数 \(g \mapsto \langle x \mid \pi(g)y \rangle\)；它是 G 上取复值的连续函数。

设 E 中的 \(x\) 与 \(y\)。对 \(z \in \mathrm{C}\) 与 \(g \in \mathrm{G}\)，我们有

\[
f _ {x, y} (z g) = \langle x \mid \pi (z g) y \rangle = \langle x \mid \chi (z) \pi (g) y \rangle = \chi (z) f _ {x, y} (g),\tag{9}
\]

因此 \(f_{x,y}\) 属于空间 \(\mathcal{F}_{\chi}(\mathrm{G})\)（V, 第 408 页）。此外，对每对 \((g_1, h_1) \in \mathrm{G} \times \mathrm{G}\) 与每个 \(g \in \mathrm{G}\)，我们有

\[
f _ {\pi (g _ {1}) x, \pi (h _ {1}) y} (g) = \langle \pi (g _ {1}) x | \pi (g) \pi (h _ {1}) y \rangle = f _ {x, y} (g _ {1} ^ {- 1} g h _ {1}).\tag{10}
\]

关系式 (9) 为下面的定义提供了依据。

定义 3. --- 设 \(\pi\) 是 G 在希尔伯特空间 E 中的一个不可约酉表示。我们称 \(\pi\) 是模中心平方可积的，如果由函数 \(|f_{x,y}|\) 通过取商得到的 G/C 上的函数属于 \(\mathcal{L}^2 (\mathrm{G / C})\)，对所有 \((x,y)\in \mathrm{E}\times \mathrm{E}\) 成立。

这一条件不依赖于 G/C 上哈尔测度的选取。

若 \(\pi\) 的矩阵系数都属于 \(\mathcal{L}^2 (\mathrm{G})\)，则称 \(\pi\) 是平方可积的；G 的一个平方可积不可约酉表示的存在蕴含 G 的中心是紧的（V, 第 487 页, 习题 5）。

存在这样的群 G，它没有任何平方可积表示，即使模中心也没有（参见 V, 第 516 页, 习题 32）。

命题 6. --- 设 \(\pi\) 是 G 在希尔伯特空间 E 中的一个不可约酉表示。设 \(\chi\) 是 \(\pi\) 的中心特征标。则 \(\pi\) 模中心平方可积当且仅当存在 E 中非零元素 \(x_0\) 与 \(y_0\)，使得由 \(|f_{x_0,y_0}|\) 通过取商诱导的 G/C 上的函数属于 \(\mathcal{L}^2 (\mathrm{G / C})\)。

设存在 E 的非零元素 \(x_{0}\) 与 \(y_{0}\)，使得 \(|f_{x_{0},y_{0}}| \in \mathcal{L}^{2}(\mathrm{G/C})\)。只需证明这时 \(\pi\) 模中心平方可积。

 设 F 是由元素 \(x \in E\)——这里 \(|f_{x,y_{0}}|\) 属于 \(\mathcal{L}^{2}(G/C)\)——组成之集。这是 E 的一个向量子空间；它包含 \(x_{0}\)，从而非零。上面的关系式 (10) 蕴含 F 在 \(\pi\) 下稳定；由于表示 \(\pi\) 不可约，空间 F 在 E 中稠密。

 设 \(x \in \mathrm{F}\)。由于 \(f_{x,y_0}\) 属于 \(\mathcal{F}_{\chi}(\mathrm{G})\) 且 \(\mu\) -可测，又 \(|f_{x,y_0}| \in \mathcal{L}^2(\mathrm{G/C})\)，我们有 \(f_{x,y_0} \in \mathcal{L}_{\chi}^2(\mathrm{G})\)（V, 第 413 页的命题 5 应用于 Z = C）。用 \(u\) 表示 E 到 \(\mathrm{L}_{\chi}^{2}(\mathrm{G})\) 的部分算子，其定义域为 F，且把 \(x \in \mathrm{F}\) 映为 \(f_{x,y_0}\) 在 \(\mathrm{L}_{\chi}^{2}(\mathrm{G})\) 中的类。

 我们来证明部分算子 \(u\) 是闭的。设 \((x_{n}, u(x_{n}))_{n \in \mathbf{N}}\) 是 \(u\) 的图中的元素组成的序列，它在 \(\mathrm{E} \times \mathrm{L}_{\chi}^{2}(\mathrm{G})\) 中收敛。设 \(x\) 是 \((x_{n})\) 的极限，并设 \(f \in \mathcal{L}_{\chi}^{2}(\mathrm{G})\) 是一个函数，其类是序列 \((u(x_{n}))\) 的极限。

函数 \(u(x_{n})\) 是矩阵系数 \(f_{x_{n},y_{0}}\) 的类。对每个 \(g \in G\)，我们有

\[
f _ {x _ {n}, y _ {0}} (g) = \langle x _ {n} | \pi (g) y _ {0} \rangle \rightarrow \langle x | \pi (g) y _ {0} \rangle = f _ {x, y _ {0}} (g)
\]

当 \(n \to +\infty\)。因此我们有 \(f_{x,y_{0}} \in \mathcal{L}_{\chi}^{2}(\mathrm{G})\)，且模 C 几乎处处有 \(f = f_{x_{0},y}\)（V, 第 409 页, 命题 3 的推论）；这就是说 \(u(x)\) 是 f 在 \(\mathrm{L}_{\chi}^{2}(\mathrm{G})\) 中的类。于是我们证明了 u 是闭的。

u 的定义域 F 在 \(\pi\) 下稳定，且关系式 (10) 表明 u 满足 \(u \circ \pi(g) = \gamma_{\mathrm{G},\chi}(g) \circ u\)，对所有 \(g \in G\)。从而我们也有等式 \(u^{*} \circ \gamma_{\mathrm{G},\chi}(g) = \pi(g) \circ u^{*}\)，对所有 \(g \in G\)（V, 第 381 页, 引理 6），由此得 \((u^{*} \circ u) \circ \pi(g) = \pi(g) \circ (u^{*} \circ u)\)。而部分算子 \(u^{*} \circ u\) 是自伴的（IV, 第 241 页, 命题 12），特别地是闭的（IV, 第 236 页, 命题 7），于是 V, 第 387 页, 推论 1 蕴含 \(u^{*} \circ u\) 的定义域等于 E。更有，我们有 F = E，也就是说，函数 \(|f_{x,y_{0}}|\) 属于 \(\mathcal{L}^{2}(\mathrm{G}/\mathrm{C})\)，对所有 \(x \in E\) 成立。

设 \((x,y)\in\mathrm{E}\times\mathrm{E}\)。我们有 \(f_{x,y_{0}}\in\mathcal{L}_{\chi}^{2}(\mathrm{G})\)。把 \(\delta_{G,\chi}\) 代替 \(\gamma_{G,\chi}\)，用同样方法可证：\(y\in E\) 中使函数 \(|f_{x,y}|\) 属于 \(\mathcal{L}^{2}(\mathrm{G}/\mathrm{C})\) 的那种元素所成之集等于 E。命题得证。

在本节余下部分，我们取定 C 的一个使 C/Z 为紧的闭子群 Z。

引理 9. --- 设 \(\pi\) 是 G 在希尔伯特空间 E 中的一个模中心平方可积的不可约酉表示。设 \(\chi\) 是 \(\pi\) 的中心特征标到 Z 的限制。对 E 中所有 \(x\) 与 \(y\)，我们有 \(f_{x,y} \in \mathcal{L}_{\chi}^{2}(\mathrm{G})\)。

 函数 \(f_{x,y}\) 连续，因此 \(\mu\) -可测。由公式 (9)（第 420 页）我们有 \(f_{x,y} \in \mathcal{F}_{\chi}(\mathrm{G})\)。此外，由 INT, VII, 第 64 页, § 2, 第 8 目, 推论 1, c)，我们有 \(\mathrm{N}_{2}(f_{x,y}) < +\infty\)，因为 C/Z 紧且 \(|f_{x,y}| \in \mathcal{L}^{2}(\mathrm{G}/\mathrm{C})\)。于是断言由 V 第 413 页的命题 5 得出。

命题 7. --- 设 π 是 G 在希尔伯特空间 E 中的一个模中心平方可积的不可约酉表示。设 χ 是 π 的中心特征标到 Z 的限制。

存在一个实数 c \textgreater{} 0 与一个唯一的 \((\mathbf{G} \times \mathbf{G})\) -等距态射 w，它把酉表示 \(\overline{\pi} \boxtimes \pi\) 映到 \(\mathrm{L}_{\chi}^{2}(\mathrm{G})\) 中，使得对所有 \((x, y) \in \overline{\mathrm{E}} \times \mathrm{E}\)，元素 \(w(x \otimes y)\) 是 \(\mathrm{L}_{\chi}^{2}(\mathrm{G})\) 中函数 \(c^{1/2} f_{x,y}\) 的类。

 对每对 \((x,y)\in\mathrm{E}\times\mathrm{E}\)，我们有 \(f_{x,y}\in\mathcal{L}_{\chi}^{2}(\mathrm{G})\)（引理 9）。用 v 表示 \(\overline{E}\otimes E\) 到 \(\mathrm{L}_{\chi}^{2}(\mathrm{G})\) 中的唯一线性映射，使得 \(v(x\otimes y)\) 是 \(f_{x,y}\) 的类，对所有 \((x,y)\in\overline{\mathrm{E}}\times\mathrm{E}\) 成立。

我们在下面证明如下的引理。

引理 10. --- 存在一个实数 c \textgreater{} 0，使得线性映射 \(w = c^{1/2}v\) 是等距的。

假设该引理成立，注意公式 (10)（第 420 页）可以写成

\[
v (\overline {{\pi}} (g _ {1}) x \otimes \pi (h _ {1}) y) = \varrho_ {\mathrm{G}, \chi} (g _ {1}, h _ {1}) v (x \otimes y)\tag{11}
\]

 对所有 \((g_{1}, h_{1}) \in \mathbf{G} \times \mathbf{G}\) 与每对 \((x, y) \in \overline{\mathbf{E}} \times \mathbf{E}\) 成立。\(\overline{\mathbf{E}} \otimes \mathbf{E}\) 到 \(\mathrm{L}_{\chi}^{2}(\mathrm{G})\) 中的等距线性映射 w 可连续延拓（仍记作 w）到 \(\overline{E} \widehat{\otimes}_{2} E\) 上。由连续性与线性性，公式 (11) 蕴含 w 是一个 \((\mathrm{G} \times \mathrm{G})\) -态射——它从 \(\overline{\pi} \boxtimes \pi\) 映到 \(\mathrm{L}_{\chi}^{2}(\mathrm{G})\) 中——这就完成了命题的证明。

 我们来证明该引理。对所有 \(x \in E\)，用 \(u_{x}\) 表示线性映射 \(y \mapsto v(x \otimes y) = f_{x,y}\)（它从 E 映到 \(\mathrm{L}_{\chi}^{2}(\mathrm{G})\) 中）。我们有 \(u_{x} \in \mathrm{Hom}_{\mathrm{G}}(\pi, \gamma_{\mathrm{G},\chi})\)。由公式 (10)，根据 V 第 388 页的推论 5，存在一个实数 \(\lambda_{x} \geqslant 0\)，使得 \(\lambda_{x}u_{x}\) 是等距的。

设 E 中的 \(x\) 与 \(y\)。我们有

\[
\| f _ {x, y} \| ^ {2} = \int_ {\mathrm{G/Z}} \overline {{f}} _ {x, y} f _ {x, y} d \nu_ {\mathrm{Z}} = \int_ {\mathrm{G/Z}} \check {\overline {{f}}} _ {x, y} \check {f} _ {x, y} d \nu_ {\mathrm{Z}}
\]

因为 G/Z 是单模的（V 第 417 页, 引理 7）。注意到 \(\check{f}_{x,y} = \overline{f}_{y,x}\)，我们得到

\[
\lambda_ {x} \| y \| ^ {2} = \| f _ {x, y} \| ^ {2} = \int_ {\mathrm{G/Z}} f _ {y, x} \overline {{f}} _ {y, x} d \nu_ {\mathrm{Z}} = \| f _ {y, x} \| ^ {2} = \lambda_ {y} \| x \| ^ {2}.
\]

这表明正实数 \(\lambda_{x}/\|x\|^{2}\) 与 E 中非零元素 x 的选取无关。它是严格正的，因为对每个非零的 x，函数 \(f_{x,x}\) 连续且在 e 处取值 \(\|x\|^{2}>0\)，由此 \(\|u_{x}(x)\|=\|f_{x,x}\|>0\)。用 \(c^{-1}\) 记这个实数。

对每对 \((x,y)\in\mathrm{E}\times\mathrm{E}\)，由此得出

\[
\| v (x \otimes y) \| ^ {2} = \| f _ {x, y} \| ^ {2} = \lambda_ {x} \| y \| ^ {2} = c ^ {- 1} \| x \| ^ {2} \| y \| ^ {2} = c ^ {- 1} \| x \otimes y \| ^ {2}.
\]

利用 EVT, V, 第 29 页, 推论 1，我们推出线性映射 \(w = c^{1/2}v\)（它从 \(\overline{\mathrm{E}}\widehat{\otimes}_2\mathrm{E}\) 映到 \(\mathrm{L}_{\chi}^{2}(\mathrm{G})\) 中）是等距的，这正是所要证的。

 推论. --- 设 \(\pi\) 是 G 的一个不可约酉表示，\(\chi\) 是其中心特征标到 Z 的限制。表示 \(\pi\) 模中心平方可积当且仅当它同构于表示 \(\gamma_{\mathrm{G},\overline{\chi}}\)（从 G 到 \(\mathrm{L}_{\overline{\chi}}^{2}(\mathrm{G})\) 中）的一个子表示（相应地，同构于表示 \(\delta_{\mathrm{G},\chi}\)（从 G 到 \(\mathrm{L}_{\chi}^{2}(\mathrm{G})\) 中）的一个子表示）。

我们来证明关于 \(\gamma_{\mathrm{G},\overline{\chi}}\) 的断言，第二个断言的证明是类似的。

 先设表示 \(\gamma_{\mathrm{G},\overline{\chi}}\) 存在一个同构于 \(\pi\) 的子表示 E。由于空间 E 非零，存在函数 \(f_1 \in \mathcal{K}_{\overline{\chi}}(\mathrm{G})\)，其类 \(\widetilde{f}_1\) 不与 E 正交。我们有 \(f_1 \in \mathcal{L}_{\overline{\chi}}^1 (\mathrm{G})\)。用 \(\widetilde{f}_{1,\mathrm{E}}\) 表示 \(\widetilde{f}_1\) 在 E 上的正交投影；它是 E 的一个非零元素。再设非零的 \(\widetilde{f}_2 \in \mathrm{E}\)。映射 \(h: g \mapsto \langle \widetilde{f}_{1,\mathrm{E}} | \gamma_{\mathrm{G},\overline{\chi}}(g)\widetilde{f}_2 \rangle\) 是 \(\pi\) 的一个矩阵系数。对每个 \(g \in \mathrm{G}\)，我们有 \(h(g) = \langle \widetilde{f}_1 | \gamma_{\mathrm{G},\overline{\chi}}(g)\widetilde{f}_2 \rangle\)，因为 \(\widetilde{f}_1 - \widetilde{f}_{1,\mathrm{E}}\) 与 E 正交。由 V 第 418 页的引理 8，函数 \(h\) 属于 \(\mathcal{L}_{\chi}^{2}(\mathrm{G})\)，于是命题 6 蕴含 \(\pi\) 模中心平方可积。

反之，设 \(\pi\) 模中心平方可积。设 \(x_0 \in \mathrm{E}\) 是一个非零向量。由命题 7 与公式 (10)，映射 \(y \mapsto \check{f}_{x_0,y}\) 是 \(\pi\) 到 \(\gamma_{\mathrm{G},\overline{\chi}}\) 中的一个单射 G-态射。

定义 4. --- 设 Z 是 C 的一个使 C/Z 为紧的闭子群。设 π 是 G 的一个模中心平方可积的不可约酉表示。满足命题 7 之性质的唯一实数 c \textgreater{} 0 称为 π 相对于 Z 的形式次数。把它记作 \(c_{\mathrm{Z}}(\pi)\)。

形式次数依赖于 Z 上哈尔测度的选取。若把 Z 上的哈尔测度 \(\beta_{Z}\) 乘以一个实数 t \textgreater{} 0，则 G/Z 上的测度 \(\nu_{Z} = \mu/\beta_{Z}\) 被乘以 \(t^{-1}\)，而 \(\pi\) 的形式次数被乘以 t。

形式次数由下述性质所刻画：

命题 8（正交关系）. --- 设 π 是 G 在希尔伯特空间 E 中的一个模中心平方可积的不可约酉表示。则有

\[
c _ {\mathrm{Z}} (\pi) \int_ {\mathrm{G} / \mathrm{Z}} \overline {{\langle x | \pi (g) x ^ {\prime} \rangle}} \langle y | \pi (g) y ^ {\prime} \rangle d \nu_ {\mathrm{Z}} (g) = \overline {{\langle x | y \rangle}} \langle x ^ {\prime} | y ^ {\prime} \rangle
\]

对所有 \((x,y,x',y')\in \mathrm{E}^4\) 成立

用 \(w\) 表示命题 7 中的态射。我们有

\[
\int_ {\mathrm{G} / \mathrm{Z}} \overline {{\langle x | \pi (g) x ^ {\prime} \rangle}} \langle y | \pi (g) y ^ {\prime} \rangle d \nu_ {\mathrm{Z}} (g) = \langle f _ {x, x ^ {\prime}} | f _ {y, y ^ {\prime}} \rangle
\]

而根据出处同上，由此得出

\[
\begin{array}{c} \langle f _ {x, x ^ {\prime}} | f _ {y, y ^ {\prime}} \rangle = \frac {1}{c _ {\mathrm{Z}} (\pi)} \langle w (x \otimes x ^ {\prime}) | w (y \otimes y ^ {\prime}) \rangle \\ = \frac {1}{c _ {\mathrm{Z}} (\pi)} \langle x \otimes x ^ {\prime} | y \otimes y ^ {\prime} \rangle = \frac {1}{c _ {\mathrm{Z}} (\pi)} \langle x | y \rangle_ {\overline {{\mathrm{E}}}} \langle x ^ {\prime} | y ^ {\prime} \rangle_ {\mathrm{E}}, \end{array}
\]

由此即得结论。

除上述命题外，对于互不同构的平方可积不可约表示，我们还有下述关系。

命题 9. --- 设 \(\pi_1\) 与 \(\pi_2\) 是 G 在希尔伯特空间 E \(_1\) 与 E \(_2\) 中的两个互不同构的不可约酉表示。假设 \(\pi_1\) 与 \(\pi_2\) 都模中心平方可积，且它们的中心特征标到 Z 的限制相同。则有

\[
\int_ {\mathrm{G} / \mathrm{Z}} \overline {{\langle x \mid \pi_ {1} (g) x ^ {\prime} \rangle}} \langle y \mid \pi_ {2} (g) y ^ {\prime} \rangle d \nu_ {\mathrm{Z}} (g) = 0
\]

对所有 \((x,x^{\prime},y,y^{\prime})\in \mathrm{E}_{1}^{2}\times \mathrm{E}_{2}^{2}\) 成立

 对 \(i=1,2\)，用 \(w_{i}\) 表示关于表示 \(\pi_{i}\) 的命题 7 中的态射。由 V 第 384 页的引理 8 \(b\) ) 与 V 第 394 页的命题 8 的断言 \(b\) )，\(w_{i}\) 的像含于 \(\pi_{i}\) -同型分量（\(\delta_{\mathrm{G},\chi}\) 的）中。由出处同上的断言 \(a\) )，\(w_{1}\) 的像因此与 \(w_{2}\) 的像正交。于是得出 \(\langle w_{1}(x\otimes x^{\prime})|w_{2}(y\otimes y^{\prime})\rangle=0\)，对所有 \((x,x^{\prime},y,y^{\prime})\in\mathrm{E}_{1}^{2}\times\mathrm{E}_{2}^{2}\) 成立，这就是所要证的公式。

注. --- 命题 8 与命题 9 的正交关系推广了 A, VIII, 第 399 页中的正交关系（另见 V 第 457 页 § 2 中 G 为紧的情形）。

\subsection{交换群正则表示的子表示}

在本节中，G 是一个局部紧交换群，\(\mu\) 是 G 上的一个哈尔测度。用 \(\hat{G}\) 表示 G 的对偶群（II, 第 201 页, 定义 2），用 \(\hat{\mu}\) 表示 \(\mu\) 在 \(\hat{G}\) 上的对偶哈尔测度（II, 第 214 页, 定义 4）。可测性概念总是相对于 \(\mu\) 与 \(\hat{\mu}\) 而言。

我们打算确定左正则表示 \(\gamma_{G}\)（G 在 \(L^{2}(G,\mu)\) 中的）的所有子表示。由于 G 是交换的，我们还有 \(\delta_{\mathrm{G}}(g)=\gamma_{\mathrm{G}}(g^{-1})\)，故这些子表示也是右正则表示的子表示。

对 \(\widehat{\mathbf{G}}\) 的每个可测子集 M，用 \(\mathrm{E_M}\) 表示满足下列条件的 \(f \in \mathrm{L}^2 (\mathrm{G},\mu)\) 组成的集合：傅里叶变换 \(\mathcal{F}_{\mathrm{G}}(f)\) 在 M 上几乎处处为零（参见 II, 第 210 页, 第 3 目）。它是连续线性映射 \(f \mapsto \varphi_{\mathrm{M}}\mathcal{F}_{\mathrm{G}}(f)\)（从 \(\mathrm{L}^2 (\mathrm{G},\mu)\) 到 \(\mathrm{L}^2 (\widehat{\mathrm{G}},\widehat{\mu})\) 中）的核，其中 \(\varphi_{\mathrm{M}}\) 表示 M 的特征函数（参见 II, 第 215 页, 定理 1），因而是 \(\mathrm{L}^2 (\mathrm{G},\mu)\) 的一个闭子空间。

若 \(\widehat{G}\) 的可测子集 M 与 N 使得 \((\mathrm{M} \cup \mathrm{N}) - (\mathrm{M} \cap \mathrm{N})\) 是局部零测的，则称它们相差一个局部零测集意义下相等。等价地，可测子集 M 与 N 相差一个局部零测集意义下相等，当且仅当 M 与 N 的特征函数在 \(\mathrm{L}^{\infty}(\widehat{\mathrm{G}}, \widehat{\mu})\) 中相等。

命题 10. --- a) 设 M 是 \(\widehat{G}\) 的一个可测子集。空间 \(\mathrm{E}_{\mathrm{M}}\) 是表示 \(\gamma_{\mathrm{G}}\) 的一个子表示；

\begin{enumerate}
\def\labelenumi{\alph{enumi})}
\setcounter{enumi}{1}
\item
  设 M 与 N 是 \(\widehat{\mathbf{G}}\) 的可测子集。我们有 \(\mathrm{E_M} = \mathrm{E_N}\) 当且仅当 M 与 N 相差一个局部零测集意义下相等；
\item
  \(\gamma_{\mathrm{G}}\) 的任何子表示都形如 \(\mathrm{E}_{\mathrm{M}}\)，其中 M 是 \(\widehat{\mathrm{G}}\) 的一个可测子集。
\end{enumerate}

用 \(\eta\) 表示 G 到 \(\widehat{\mathbf{G}}\) 的典范映射（参见 II, 第 216 页, 注 1）。由于 \(\mathrm{E_M}\) 是闭的，断言 \(a\) ) 由以下公式得出

\[
\boldsymbol {\gamma} _ {\mathrm{G}} (x) (f) = \varepsilon_ {x} * f, \quad \mathcal {F} _ {\mathrm{G}} (\varepsilon_ {x} * f) = \eta (x) \mathcal {F} _ {\mathrm{G}} (f)
\]

其中 \(x\in \mathrm{G}\)，\(f\in \mathrm{L}^2 (\mathrm{G},\mu)\)。

设 M 与 N 是相差一个局部零测集意义下相等的可测子集。由于此时 \(\varphi_{M}\) 等于 \(\varphi_{N}\)（在 \(L^{\infty}(\widehat{G},\widehat{\mu})\) 中），这一条件蕴含 \(E_{M}=E_{N}\)。

反之，设 M 与 N 不是相差一个局部零测集意义下相等的。那么在交换 M 与 N 的角色之后，存在紧子集 K 使得集合 \(L = K \cap (M - (M \cap N))\) 不是零测的。设 \(\varphi_{\mathrm{L}} \in \mathrm{L}^{2}(\widehat{\mathrm{G}}, \widehat{\mu})\) 是 L 的特征函数所在的类，并令 \(f = \overline{\mathcal{F}}_{\widehat{\mathrm{G}}}(\varphi_{\mathrm{L}}) \in \mathrm{L}^{2}(\mathrm{G}, \mu)\)；我们有 \(\mathcal{F}_{\mathrm{G}}(f) = \varphi_{\mathrm{L}}\)（II, 第 220 页, 定理 2 的推论）。于是 f 属于 \(E_{N}\)，因为 \(L \cap N\) 是空集；但不属于 \(E_{M}\)，因为 \(\varphi_{\mathrm{M}}\mathcal{F}(f) = \varphi_{\mathrm{M}}\varphi_{\mathrm{L}} = \varphi_{\mathrm{L}}\)。因此 \(E_{M} \neq E_{N}\)，这就证明了 b)。

现在设 E 是 \(\gamma_{\mathrm{G}}\) 的一个子表示。设 \(p_{\mathrm{E}}\) 是 \(\mathrm{L}^2 (\mathrm{G},\mu)\) 中以 E 为像的正交投影算子，并令 \(q_{\mathrm{E}} = \mathcal{F}_{\mathrm{G}}\circ p_{\mathrm{E}}\circ \mathcal{F}_{\mathrm{G}}^{-1}\)。投影算子 \(p_{\mathrm{E}}\) 属于 \(\mathrm{Hom}_{\mathrm{G}}(\gamma_{\mathrm{G}},\gamma_{\mathrm{G}})\)（V, 第 383 页, 命题 4），故它与 \(\gamma_{\mathrm{G}}(f)\) 交换，对所有 \(f\in \mathrm{L}^{1}(\mathrm{G},\mu)\)（参见 V, 第 401 页）。这意味着它与自同态 \(\varphi \mapsto f*\varphi\) 交换，对所有 \(f\in \mathrm{L}^{1}(\mathrm{G},\mu)\)。于是由 V 第 407 页的引理 4，自同态 \(q_{\mathrm{E}}\)（\(\mathrm{L}^2 (\widehat{\mathrm{G}},\widehat{\mu})\) 的）与乘以 \(g\) 的乘法自同态交换，对每个满足下列条件的函数 \(g\in \mathcal{C}_0(\widehat{\mathrm{G}})\)：它属于 \(\mathrm{L}^1 (\mathrm{G},\mu)\) 的傅里叶变换在 \(\mathcal{C}_0(\widehat{\mathrm{G}})\) 中的像（II, 第 223 页, 命题 14）。由于傅里叶变换的像在 \(\mathcal{C}_0(\widehat{\mathrm{G}})\) 中稠密（II, 第 209 页, 命题 5 的推论），态射 \(g\mapsto m_g\) 的连续性蕴含 \(q_{\mathrm{E}}\) 与 \(m_g\) 交换，对每个函数 \(g\in \mathcal{C}_0(\widehat{\mathrm{G}})\)。

于是由 IV 第 188 页的命题 7 的推论，存在函数 \(\varphi \in \mathcal{L}^{\infty}(\widehat{\mathrm{G}},\widehat{\mu})\) 使得 \(q_{\mathrm{E}} = m_{\varphi}\)。由于 \(q_{\mathrm{E}} = q_{\mathrm{E}}^{2}\)，我们有 \(m_{\varphi} = m_{\varphi}^{2} = m_{\varphi^{2}}\)，从而 \(\varphi = \varphi^{2}\)（在 \(\mathrm{L}^{\infty}(\widehat{\mathrm{G}},\widehat{\mu})\) 中）（IV, 第 186 页, 命题 5）；这意味着 \(\varphi\) 在 \(\mathrm{L}^{\infty}(\widehat{\mathrm{G}},\widehat{\mu})\) 中等于 \(\widehat{\mathrm{G}}\) 的某个可测子集 N 的特征函数所在的类。令 \(\mathrm{M} = \widehat{\mathrm{G}} - \mathrm{N}\)。设 \(f \in \mathrm{L}^{2}(\mathrm{G},\mu)\)。我们有：\(f \in \mathrm{E}\) 当且仅当 \(p_{\mathrm{E}}(f) = f\)，当且仅当 \(\varphi \mathcal{F}_{\mathrm{G}}(f) = \mathcal{F}_{\mathrm{G}}(f)\)，而这等价于 \(f \in \mathrm{E}_{\mathrm{M}}\)。

注记. --- 设 \(\chi\) 是 G 的一个特征标。若 G 不是紧的，则 \(\chi\) -同型分量（G 在 L \(^{2}\) (G) 中的正则表示的）

是平凡的。若 G 是紧的，则 G 的正则表示的 \(\chi\) -同型分量的维数为 1，且函数 \(\chi \in \mathrm{L}^{2}(\mathrm{G})\) 是它的一个基。

\subsection{群 R 的酉表示}

在本节中，E 表示一个复希尔伯特空间。群 R 赋以勒贝格测度。

引理 11. --- 设 u 是 E 上的一个自伴部分算子。对 \(t \in \mathbf{R}\)，令 \(\varrho(t) = e^{itu} \in \mathcal{L}(\mathrm{E})\)。则映射 \(\varrho\) 是 \(\mathbf{R}\) 的一个酉表示。

算子 \(\varrho(t)\) 由普遍可测函数演算定义（IV, 第 272 页, 定义 5）；它是 E 的一个自同态，因为函数 \(x \mapsto e^{itx}\) 在 \(\mathrm{Sp}(u)\) 上有界（IV, 第 275 页, 命题 5, \(a\) )）。

由 IV 第 275 页的命题 5，我们有 \(\varrho(0) = 1_{\mathrm{E}}\)，\(\varrho(t)^* = \varrho(-t)\) 对所有 \(t \in \mathbf{R}\)，且 \(\varrho(t_1 + t_2) = \varrho(t_1)\varrho(t_2)\) 对所有 \(t_1\) 与 \(t_2\)（都属于 \(\mathbf{R}\)）。特别地，自同态 \(\varrho(t)\) 对每个 \(t \in \mathbf{R}\) 都是酉的。对每个 \(x \in \mathrm{E}\)，从 \(\mathbf{R}\) 到 E 的由 \(x \mapsto \varrho(t)x\) 定义的映射在 \(t = 0\) 处连续（IV, 第 276 页, 命题 6），故 \(\varrho\) 是 \(\mathbf{R}\) 在 E 中的一个酉表示（V, 第 380 页, 引理 4）。

引理 12. --- 设 \(\varrho\) 是 \(\mathbf{R}\) 在 E 中的一个酉表示。设 D 是 E 中由满足下列条件的元素 \(x\) 组成的集合：映射 \(\psi_x: t \mapsto \varrho(t)x\) 在 0 处可微。D 到 E 中的映射 \(u\) 由 \(x \mapsto i^{-1} \psi_x'(0)\) 给出，它定义了 E 上一个对称的部分算子。

设 \(f \in \mathcal{D}(\mathbf{R})\) 且 \(x \in \mathrm{E}\)。令 \(y = \varrho(f)x\)。对每个 \(h \in \mathbf{R}\)，我们有

\[
\begin{array}{r l} & {\psi_ {y} (h) = \varrho (h) \varrho (f) x = \varrho (f) \varrho (h) x} \\ & {\qquad = \int_ {\mathbf {R}} f (t) \varrho (t + h) x d t = \int_ {\mathbf {R}} f (t - h) \psi_ {x} (t) d t.} \end{array}
\]

从 \(R^{2}\) 到 C 的由 \((t,h)\mapsto f(t-h)\) 定义的函数是无穷次可微的；它对 h 的导数是由 \((t,h)\mapsto -f'(t-h)\) 定义的函数，后者被 \(\|f'\|_{\infty}\) 所控制。当 h=0 时，这个导数在 t 属于一个紧集之外时为零。由于对所有 t 有 \(\|\psi_{x}(t)\|=\|x\|\)，由 IV 第 197 页的命题 2 推出，函数 \(\psi_{y}\) 在 0 处可微且其导数由下式给出

\[
\psi_ {y} ^ {\prime} (0) = - \int_ {\mathbf {R}} f ^ {\prime} (t) \psi_ {x} (t) d t = - \varrho (f ^ {\prime}) x.
\]

因此 \(y \in D\)。由于空间 \(\mathcal{D}(\mathbf{R})\) 在 \(\mathrm{L}^{1}(\mathbf{R})\) 中稠密，由 INT, VIII, 第 139 页, § 2, 第 7 目, 命题 10 以及（IV, 第 202 页, 命题 4）推出，空间 D 在 E 中稠密。

设 \(x_{1}\) 与 \(x_{2}\) 属于 D。我们计算

\[
\begin{array}{c} \langle u (x _ {1}) \mid x _ {2} \rangle = \lim _ {h \to 0} \frac {i}{h} \langle (\varrho (h) - 1 _ {\mathrm{E}}) x _ {1} \mid x _ {2} \rangle \\ = \lim _ {h \to 0} \frac {i}{h} \langle x _ {1} \mid (\varrho (- h) - 1 _ {\mathrm{E}}) x _ {2} \rangle = \langle x _ {1} \mid u (x _ {2}) \rangle . \end{array}
\]

因此，部分算子 \(u\) 是对称的。

定义 5. --- 设 \(\varrho\) 是 \(\mathbf{R}\) 在 E 中的一个酉表示。引理 12 中定义的对称部分算子称为 \(\varrho\) 的无穷小生成元。

定理 1（斯通）. --- 映射 σ 把 R 在 E 中的每个酉表示 ρ 对应到其无穷小生成元 u，它定义了从 R 在 E 中的酉表示集合到 E 上自伴部分算子集合上的一个双射。逆双射 τ 把一个自伴部分算子对应到酉表示 \(t \mapsto e^{itu}\)。

我们先证明几个引理。

引理 13. --- 设 \(\varrho\) 是 \(\mathbf{R}\) 在 E 中的一个酉表示，\(u\) 是其无穷小生成元。

\begin{enumerate}
\def\labelenumi{\alph{enumi})}
\item
  u 的定义域是满足下列条件的 \(x \in E\) 组成的集合：映射 \(\psi_{x} \colon t \mapsto \varrho(t)x\) 在 R 上可微；
\item
  对每个 \(t \in \mathbf{R}\) 与每个 \(x \in \mathrm{dom}(u)\)，我们有 \(\varrho(t)x \in \mathrm{dom}(u)\) 且 \(\psi_x'(t) = i u(\varrho(t)x)\)；
\item
  部分算子 u 是本质自伴的。
\end{enumerate}

对每个 \(x \in \mathrm{E}\)，用 \(\psi_x\) 表示从 \(\mathbf{R}\) 到 \(\mathrm{E}\) 中的由 \(\psi_x(t) = \varrho(t)x\) 定义的映射。对每个 \(t \in \mathbf{R}\) 与每个 \(h \in \mathbf{R}\)，我们有

\[
\psi_ {x} (t + h) - \psi_ {x} (t) = \varrho (t) \big (\psi_ {x} (h) - \psi_ {x} (0) \big),
\]

这证明了 \(\operatorname{dom}(u)\) 是满足下列条件的元素 \(x \in E\) 组成的空间：\(\psi_{x}\) 在 R 上可微；并且确立了 \(\psi_{x}^{\prime}(t) = \varrho(t)\psi_{x}^{\prime}(0) = i\varrho(t)(u(x))\) 对所有 \(t \in R\)。

设 \(x \in E\) 且 \(t \in R\)。我们有 \(\psi_{\varrho(t)x}(s) = \psi_x(s + t)\)，对所有 \(s \in R\)。于是我们有 \(\varrho(t)x \in \text{dom}(u)\)，只要 \(x \in \text{dom}(u)\)，而且由 a) 有 \(u(\varrho(t)x) = i^{-1}\psi_x'(t) = \varrho(t)u(x)\)。由此得到断言 b)。

我们来证明 c)。设 \(x \in \text{Ker}(u^{*} - i1_{\text{E}})\)。我们来证明 x = 0。设 \(y \in \text{dom}(u)\)，并设 f 是 R 上由 \(f(t) = \langle \psi_{y}(t) | x \rangle\)（\(t \in R\)）定义的函数。函数 f 是有界的，因为 \(\| \psi_{y}(t) \| = \| y \|\) 对所有 \(t \in R\)；它在 R 上可微，且对每个 \(t \in R\)，断言 b) 蕴含

\[
\begin{array}{r l} f ^ {\prime} (t) = \langle \psi_ {y} ^ {\prime} (t) | x \rangle & = \langle i u (\varrho (t) y) | x \rangle \\ & = - i \langle \varrho (t) y | u ^ {*} (x) \rangle = - i \langle \psi_ {y} (t) | i x \rangle = f (t) \end{array}
\]

因为 \(u^{*}(x)=ix\)。因此我们有 \(f(t)=f(0)e^{t}\)，对所有 \(t\in R\)（FVR, IV, 第 27 页）。由于 f 有界，函数 f 为零，特别地我们有 \(\langle y|x\rangle=f(0)=0\)。由于空间 dom(u) 在 E 中稠密，我们得出 x=0。类似地，我们证明 \(\operatorname{Ker}(u^{*}+i1_{\mathrm{E}})\) 化为 0。由 IV 第 261 页的推论 3，部分算子 u 因此是本质自伴的。

引理 14. --- 设 u 是 E 上的一个自伴部分算子，\(\varrho(t) = e^{itu}\) 是由 u 定义的 R 的酉表示。\(\varrho\) 的无穷小生成元等于 u。

设 \(x \in \text{dom}(u)\)。令 \(\psi_{x}(t) = \varrho(t)x\)，对所有 \(t \in R\)。对每个非零实数 h，我们有

\[
\frac {1}{h} \left(\psi_ {x} (t + h) - \psi_ {x} (t)\right) = \left(\frac {1}{h} \left(e ^ {i h u} - 1 _ {\mathrm{E}}\right)\right) e ^ {i t u} x = \left(\frac {1}{h} \left(e ^ {i h u} - 1 _ {\mathrm{E}}\right)\right) \varrho (t) x.
\]

当 h 趋于 0 时，我们有

\[
\frac {1}{h} (e ^ {i h t} - 1) \rightarrow i t
\]

对所有 \(t \in R\) 成立。此外，

\[
\left| \frac {1}{h} (e ^ {i h t} - 1) \right| = | t | \left| \frac {1}{h} \int_ {0} ^ {h} e ^ {i t s} d s \right| \leqslant | t |.
\]

于是由 IV 第 276 页的命题 6 推出，函数 \(\psi_x\) 在 \(\mathbf{R}\) 上可微且满足 \(\psi_x'(t) = i u(\varrho(t)x)\)，对所有 \(t \in \mathbf{R}\)。因此，\(u\) 的定义域含于满足下列条件的 \(x \in E\) 组成的集合：\(\psi_x\) 在 0 处可微且此时有 \(\psi_x'(0) = i u(x)\)。按定义，这意味着 \(\varrho\) 的无穷小生成元是 \(u\) 的一个扩张。由于这两个算子都是对称的而 \(u\) 是自伴的（IV, 第 238 页, 注 5），它们因此相等。

引理 15. --- 设 \(\varrho\) 是 \(\mathbf{R}\) 在 E 中的一个酉表示。则无穷小生成元 \(u\)（\(\varrho\) 的）是自伴的，且 \(\varrho(t) = e^{itu}\) 对所有 \(t \in \mathbf{R}\)。

由引理 13(c)，部分算子 u 是本质自伴的。因此它的闭包 \(\overline{u}\) 是一个自伴算子。用 \(\pi\) 表示由 \(\pi(t) = e^{it\overline{u}}\) 定义的 R 的酉表示（引理 11）。

对每个 \(x \in \mathrm{E}\)，用 \(\psi_x\)（相应地 \(\widetilde{\psi}_x\)）表示从 \(\mathbf{R}\) 到 \(\mathrm{E}\) 中的由 \(\psi_x(t) = \varrho(t)x\)（相应地由 \(\widetilde{\psi}_x(t) = \pi(t)x\)）定义的映射。

由引理 13 的 a) 与 b) 两部分，空间 \(\mathrm{dom}(u)\) 是 E 中由满足下列条件的元素 \(x \in \mathrm{E}\) 组成的子空间：映射 \(\psi_x\) 在 \(\mathbf{R}\) 上可微，且对所有 \(x \in \mathrm{dom}(u)\) 与每个 \(t \in \mathbf{R}\)，我们有 \(\psi_x'(t) = i u (\psi_x(t))\)。同样，空间 \(\mathrm{dom}(\overline{u})\) 是 E 中由满足下列条件的元素 \(x \in \mathrm{E}\) 组成的子空间：映射 \(\widetilde{\psi}_x\) 在 \(\mathbf{R}\) 上可微，且对所有 \(x \in \mathrm{dom}(\overline{u})\) 与每个 \(t \in \mathbf{R}\)，我们有 \(\widetilde{\psi}_x'(t) = i \overline{u} (\widetilde{\psi}_x(t))\)。

设 \(x \in \text{dom}(u) \subset \text{dom}(\overline{u})\)。令 \(f = \psi_x - \widetilde{\psi}_x\)。这是从 \(\mathbf{R}\) 到 E 中的一个可微函数。对每个 \(t \in \mathbf{R}\)，由此得出

\[
f ^ {\prime} (t) = i u (\psi_ {x} (t)) - i \overline {{u}} (\widetilde {\psi} _ {x} (t)) = i \overline {{u}} (f (t))
\]

因为 \(\psi_x(t) \in \mathrm{dom}(u)\) 且 \(u \subset \overline{u}\)。令 \(g = \|f\|^2\)；它是从 \(\mathbf{R}\) 到 \(\mathbf{R}\) 中的一个可微映射且满足 \(g(0) = 0\)。对 \(t \in \mathbf{R}\)，由 FVR, I, 第 28 页, 命题 2 得到

\[
\begin{array}{r l} g ^ {\prime} (t) & = \langle f ^ {\prime} (t) | f (t) \rangle + \langle f (t) | f ^ {\prime} (t) \rangle \\ & = \langle i \overline {{u}} (f (t)) | f (t) \rangle + \langle f (t) | i \overline {{u}} (f (t)) \rangle = 0 \end{array}
\]

因为 \(\overline{u}\) 是自伴的。因此 f = 0，故 \(\varrho(t)x = \pi(t)x\) 对所有 \(t \in R\)。E 的连续自同态 \(\varrho(t)\) 与 \(\pi(t)\) 在 dom(u) 上一致，因而对所有 \(t \in R\) 相等。于是我们有 \(\pi = \varrho\)；由于 \(\overline{u}\) 是 \(\pi\) 的无穷小生成元（引理 14），我们有 \(\overline{u} = u\)，这就证明了 u 是自伴的。

现在我们可以证明定理 1 了。

映射 \(\sigma\) 与 \(\tau\) 是完全有定义的（分别是引理 15 与引理 11）。

设 \(\varrho\) 是 \(\mathbf{R}\) 的一个酉表示。用 \(u\) 表示它的无穷小生成元。关系式 \(\varrho(t) = e^{itu}\)（对所有 \(t \in \mathbf{R}\)）（引理 15）证明了 \(\tau \circ \sigma\) 是恒同映射。

设 u 是 E 上的一个自伴部分算子。引理 14 表明酉表示 \(t \mapsto e^{itu}\) 的无穷小生成元等于 u，因而 \(\sigma \circ \tau\) 是恒同映射。

设 u 是 E 上的一个自伴部分算子，\(\varrho(t) = e^{itu}\)（\(t \in R\)）是与之相伴的 R 在 E 中的酉表示。

设 \(x \in \text{dom}(u)\)。此时所满足的方程 \(\partial_{t} \varrho(t)x = i u(\varrho(t)x)\)（引理 13, b)）称为薛定谔方程。

推论. --- 设 \(\varrho\) 是 \(\mathbf{R}\) 在希尔伯特空间 E 中的一个酉表示。存在一个局部紧空间 X、X 上的一个正测度 \(\mu\) 以及 X 上的一个实值连续函数 \(g\)，使得 \(\varrho\) 同构于表示 \(\pi\)——\(\mathbf{R}\) 在 \(\mathrm{L}^2 (\mathrm{X},\mu)\) 中的——它由 \(\pi (t)f = e^{itg}f\) 定义（对所有 \(t\in \mathbf{R}\) 与每个 \(f\in \mathrm{L}^2 (\mathrm{X},\mu)\)）。

设 \(u\) 是 E 上使得 \(\varrho(t) = e^{itu}\)（对所有 \(t \in \mathbf{R}\)）的自伴算子（定理 1）。存在一个局部紧空间 X、X 上的一个正测度 \(\mu\)、一个等距同构 \(\theta\)（\(L^2(X, \mu)\) 到 E 上的） 以及 X 上的一个实值连续函数 \(g\)，使得 \(u = \theta \circ m_g \circ \theta^{-1}\)（IV, 第 266 页, 定理 1）。断言由公式 \(e^{itu} = \theta \circ e^{itm_g} \circ \theta^{-1} = \theta \circ m_{e^{itg}} \circ \theta^{-1}\)（IV, 第 269 页, 引理 4）得出。

注. --- 假设 u 是 E 的一个自同态。由 \(\varrho(t) = e^{itu}\) 定义的 R 在 E 中的酉表示此时满足不等式 \(\|\varrho(t) - 1_{\mathrm{E}}\| \leqslant |t||u\|\)，对所有 \(t \in R\)，且映射 \(\varrho\)（从 R 到巴拿赫空间 \(\mathcal{L}(\mathrm{E})\) 中的）是线性微分方程

\[
\frac {1}{i} \frac {d \varrho}{d t} = u \circ \varrho
\]

（参见 FVR, IV, 第 26 页, §6）的唯一解。

例. --- 设 \(\varrho\) 是 \(\mathbf{R}\) 在 \(\mathrm{L}^2 (\mathbf{R})\) 中的正则表示。\(\varrho\) 的无穷小生成元是以 \(\mathcal{D}(\mathbf{R})\) 为定义域、由 \(f\mapsto -if'\) 定义的微分算子的闭包。

\section{正定函数}

在本节中，除非另有说明，所考虑的所有向量空间，以及所有希尔伯特空间与代数，都是在 C 上的。

\subsection{普遍正核}

在本章中，X 是一个分离拓扑空间。

定理 1. --- 设 \(f \in \mathcal{C}(X \times X)\)。下列条件等价：

\begin{enumerate}
\def\labelenumi{(\roman{enumi})}
\tightlist
\item
  对 X 的每个紧子集 Y 及 Y 上的每个正测度 \(\mu\)，\(\mathrm{L}^{2}(\mathrm{Y},\mu)\) 的由核 \(f|(\mathrm{Y}\times\mathrm{Y})\)（III, 第 29 页, 定义 1）定义的自同态是正的，换句话说，我们有
\end{enumerate}

\[
\int_ {\mathrm{Y} \times \mathrm{Y}} \overline {{h (x)}} h (y) f (x, y) d (\mu \otimes \mu) (x, y) \geqslant 0
\]

对所有 \(h \in \mathcal{L}^{2}(\mathrm{Y}, \mu)\) 成立；

\begin{enumerate}
\def\labelenumi{(\roman{enumi})}
\setcounter{enumi}{1}
\tightlist
\item
  对每个整数 \(n \in N\)、X 中每个族 \((x_{i})_{0 \leqslant i \leqslant n}\) 以及复数的任意族 \((t_{i})_{0 \leqslant i \leqslant n}\)，我们有
\end{enumerate}

\[
\sum_ {i = 0} ^ {n} \sum_ {j = 0} ^ {n} \bar {t} _ {i} t _ {j} f (x _ {i}, x _ {j}) \geqslant 0;
\]

\begin{enumerate}
\def\labelenumi{(\roman{enumi})}
\setcounter{enumi}{2}
\item
  存在一个复希尔伯特空间 E 和一个具有完全系像的连续映射 \(g \colon X \to E\)，使得 \(f(x, y) = \langle g(x) | g(y) \rangle\) 对所有 \(x\) 与 \(y\)（\(X\) 中的）成立；
\item
  存在一个复希尔伯特空间 E 和一个连续映射 \(g: X \to E\)，使得对所有 X 中的 x 与 y 有 \(f(x, y) = \langle g(x) | g(y) \rangle\)。
\end{enumerate}

(iii) 蕴含 (iv) 是显然的，而通过考虑离散测度 \(\mu\)，即 \(\{0,\ldots,n\}\) 上的计数测度经映射 \(i\mapsto x_{i}\) 的像，以及满足下式的函数 h，可知 (i) 蕴含 (ii)

\[
h(x) = \sum_{\substack{0\leqslant i\leqslant n\\ x_{i} = x}}t_{i}.
\]

我们来证明 (iv) 蕴含 (i)。设存在一个复希尔伯特空间 E 和一个连续映射 \(g \colon \mathrm{X} \to \mathrm{E}\)，使得 \(f(x, y) = \langle g(x) \mid g(y) \rangle\) 对所有 \((x, y) \in \mathrm{X} \times \mathrm{X}\)。设 Y 是 X 的一个紧子集，\(\mu\) 是 Y 上的一个正测度，且 \(h \in \mathcal{L}^{2}(Y, \mu)\)。我们有

\[
\begin{array}{r l} & {\int_ {\mathrm{Y} \times \mathrm{Y}} \overline {{h (x)}} h (y) f (x, y) d (\mu \otimes \mu) (x, y)} \\ & {\qquad = \int_ {\mathrm{Y} \times \mathrm{Y}} \overline {{h (x)}} h (y) \langle g (x) | g (y) \rangle d (\mu \otimes \mu) (x, y)} \\ & {\qquad = \left\langle \int_ {\mathrm{Y}} h (x) g (x) d \mu (x) \Big | \int_ {\mathrm{Y}} h (y) g (y) d \mu (y) \right\rangle \geqslant 0} \end{array}
\]

（INT, V, 第 97 页, § 8, 第 4 目, 命题 9）。

最后我们来证明 (ii) 蕴含 (iii)。设 \(\widetilde{\mathrm{E}}\) 是 X 上具有限支撑的复测度组成的空间。对 \(\mu_{1}\) 与 \(\mu_{2}\)（\(\widetilde{\mathrm{E}}\) 中的），我们令

\[
\langle \mu_ {1} \mid \mu_ {2} \rangle = \int_ {\mathrm{X} \times \mathrm{X}} f (x, y) (\overline {{\mu}} _ {1} \otimes \mu_ {2}) (x, y).
\]

这样在 \(\widetilde{\mathrm{E}}\) 上定义的半双线性形式是一个正埃尔米特形式。事实上，设 \(\mu \in \widetilde{\mathrm{E}}\)；存在 X 中的有限族 \((x_{i})_{0\leqslant i\leqslant n}\) 及复数 \((t_i)_{0\leqslant i\leqslant n}\)，使得 \(\mu = \sum_{i=0}^{n} t_i \varepsilon_{x_i}\)。于是我们有

\[
\langle \mu | \mu \rangle = \sum_ {i = 0} ^ {n} \sum_ {j = 0} ^ {n} \bar {t} _ {i} t _ {j} f (x _ {i}, x _ {j}) \geqslant 0
\]

（由假设）。定义 \(\widetilde{g}\colon\mathbf{X}\to\widetilde{\mathbf{E}}\) 为 \(\widetilde{g}(x)=\varepsilon_{x}\)。映射 \(\widetilde{g}\) 的像生成 \(\widetilde{\mathbf{E}}\)。此外，对每个 \((x,y)\in\mathbf{X}\times\mathbf{X}\)，一方面有 \(f(x,y)=\langle\widetilde{g}(x)|\widetilde{g}(y)\rangle\)，另一方面有

\[
\| \widetilde {g} (x) - \widetilde {g} (y) \| ^ {2} = f (x, x) + f (y, y) - f (x, y) - f (y, x),
\]

这蕴含 \(\widetilde{g}\) 连续，因为 \(f\) 连续。

设 E 是 \(\widetilde{E}\) 的分离完备希尔伯特空间（EVT, V, 第 8 页, 命题 4 的推论），\(g\colon X \to E\) 是 \(\widetilde{g}\) 与典范映射 \(\widetilde{E} \to E\) 的复合。则映射 g 连续，它的像在 E 中是完全系，且 \(f(x,y) = \langle g(x) | g(y) \rangle\) 对所有 \((x,y) \in X \times X\)。

用来证明 (ii) 蕴含 (iii) 的方法称为盖尔范德–奈马克–塞加尔构造（GNS 构造）。

定义 1. --- 若函数 \(f \in \mathcal{C}(\mathrm{X} \times \mathrm{X})\) 满足定理 1 的等价条件，则称它是 X 上的一个普遍正核。

若 f 是 X 上的一个普遍正核，则满足出处同上条件 (iv) 的偶 \((E,g)\) 称为 f 的一个希尔伯特实现；若条件 (iii) 得到满足，则称它是一个循环希尔伯特实现。

用 \(\mathrm{Noy}_{+}(\mathrm{X})\) 表示 X 上普遍正核的集合。

设 X' 是一个分离拓扑空间，\(h: X \to X'\) 是一个连续映射。映射 \(f \mapsto f \circ (h, h)\)（\(\mathcal{C}(X' \times X')\) 到 \(\mathcal{C}(X \times X)\) 中的）通过过渡到子空间诱导出映射 \(\mathrm{Noy}_{+}(X') \to \mathrm{Noy}_{+}(X)\)。

命题 1. --- 设 f 是 X 上的一个普遍正核。设 \((\mathrm{E}_{1}, g_{1})\) 与 \((\mathrm{E}_{2}, g_{2})\) 是 f 的希尔伯特实现。假设 \((\mathrm{E}_{1}, g_{1})\) 是一个循环希尔伯特实现。则存在唯一的连续线性映射 \(u: E_{1} \to E_{2}\) 使得 \(g_{2} = u \circ g_{1}\)。这个映射是等距的。若 \((\mathrm{E}_{2}, g_{2})\) 也是循环的，则 u 是一个同构。

\(u\) 的唯一性由 \(g_{1}\) 的像在 \(\mathrm{E}_{1}\) 中是完全系这一事实得出。设 \(\mathrm{F} = \mathbf{C}^{(\mathrm{X})}\)，\((e_x)_{x\in \mathrm{X}}\) 是 \(\mathrm{F}\) 的典范基。对 \(j = 1\) 与 \(j = 2\)，用 \(u_{j}\) 表示从 \(\mathrm{F}\) 到 \(\mathrm{E}_{j}\) 中的由 \(u_{j}(e_{x}) = g_{j}(x)\) 确定的线性映射，并用 \(\mathrm{F}_{j}\) 表示它的像；空间 \(\mathrm{F}_{1}\) 在 \(\mathrm{E}_{1}\) 中稠密。设 \(t = \sum t_{x}e_{x}\) 是 \(\mathrm{F}\) 的一个元素。我们有

\[
\begin{array}{c} \| u _ {1} (t) \| ^ {2} = \sum_ {x, y} \bar {t} _ {x} t _ {y} \langle g _ {1} (x)   |   g _ {1} (y) \rangle = \sum_ {x, y} \bar {t} _ {x} t _ {y} f (x, y) \\ = \sum_ {x, y} \bar {t} _ {x} t _ {y} \langle g _ {2} (x)   |   g _ {2} (y) \rangle = \| u _ {2} (t) \| ^ {2}. \end{array}
\]

因此，存在从 \(F_{1}\) 到 \(F_{2}\) 中的线性等距映射 v 使得 \(u_{2}=v\circ u_{1}\)，特别地，\(g_{2}(x)=v(g_{1}(x))\) 对所有 \(x\in X\)。由于 \(g_{1}\) 的像在 \(E_{1}\) 中是完全系，这个映射延拓为从 \(E_{1}\) 到 \(E_{2}\) 中的等距线性映射 u，使得 \(g_{2}=u\circ g_{1}\)。

由 I 第 107 页的引理 8，\(u\) 的像在 \(\mathrm{E}_2\) 中是闭的。若 \((\mathrm{E}_2,g_2)\) 是一个循环希尔伯特实现，则 \(u\) 的像还在 \(\mathrm{E}_2\) 中稠密，因而 \(u\) 是一个同构。

命题 2. --- 集合 Noy \(_{+}\) (X) 是空间 \(\mathcal{C}(\text{X} \times \text{X})\) 中的一个自伴锥；它在乘法下稳定，且当给 \(\mathcal{C}(\text{X} \times \text{X})\) 赋以逐点收敛拓扑时它是闭的。

若 \(f \in \mathrm{Noy}_{+}(\mathrm{X})\)，则 \(tf \in \mathrm{Noy}_{+}(\mathrm{X})\) 对每个实数 \(t \geqslant 0\)，且 \(\overline{f} \in \mathrm{Noy}_{+}(\mathrm{X})\)，这是初等的。

若 \((\mathrm{E}_{1}, g_{1})\) 与 \((\mathrm{E}_{2}, g_{2})\) 分别是 X 上普遍正核 \(f_{1}\) 与 \(f_{2}\) 的希尔伯特实现，则偶 \((\mathrm{E}_{1} \oplus \mathrm{E}_{2}, g_{1} + g_{2})\)（相应地，偶 \((\mathrm{E}_{1} \widehat{\otimes}_{2} \mathrm{E}_{2}, g_{1} \otimes g_{2})\)）是 \(f_{1} + f_{2}\)（相应地 \(f_{1}f_{2}\)）的一个希尔伯特实现；因而它们是普遍正核。

\(\mathrm{Noy}_{+}(\mathrm{X})\) 的刻画 (ii)（定理 1）蕴含这个集合在赋以逐点收敛拓扑的 \(\mathcal{C}(\mathrm{X}\times\mathrm{X})\) 中是闭的。

\subsection{关于全纯函数演算的补充注记}

对 C 的每个子集 X，用 \(X^{*}\) 表示 X 在复共轭下的像。设 U 是 C 的一个开子集，\(g: U \to C\) 是一个全纯函数。于是函数 \(f^{*}: z \mapsto \overline{g(\overline{z})}\) 在 \(U^{*}\) 上有定义且全纯。映射 \(f \mapsto f^{*}\) 是从 \(\mathcal{O}(U)\) 到 \(\mathcal{O}(U^{*})\) 上的一个连续双射，使得 \((f_{1}f_{2})^{*} = f_{1}^{*}f_{2}^{*}\) 与 \((f_{1} + f_{2})^{*} = f_{1}^{*} + f_{2}^{*}\) 对 \(f_{1}\) 与 \(f_{2}\)（\(\mathcal{O}(U)\) 中的）成立。

设 C 是 C 的一个紧子集。映射 \(f \mapsto f^{*}\)（\(\mathcal{O}(U)\) 到 \(\mathcal{O}(U^{*})\) 中的）当 U 取遍 C 的包含 C 的开子集时，诱导出空间 \(\mathcal{O}(C)\) 到 \(\mathcal{O}(C^{*})\) 上的一个连续双射（I, 第 49 页, 第 1 目），它也记作 \(f \mapsto f^{*}\)，且满足 \((f_{1}f_{2})^{*} = f_{1}^{*}f_{2}^{*}\) 与 \((f_{1} + f_{2})^{*} = f_{1}^{*} + f_{2}^{*}\) 对 \(f_{1}\) 与 \(f_{2}\)（\(\mathcal{O}(C)\) 中的）。

命题 3. --- 设 A 是一个含单位元的对合巴拿赫代数。设 \(a \in \mathrm{A}\)。\(a^{*}\) 的谱是 \(\mathrm{Sp}_{\mathrm{A}}(a)^{*}\)，即 \(\mathrm{Sp}_{\mathrm{A}}(a)\) 在复共轭下的像。对每个 \(f \in \mathcal{O}(\mathrm{Sp}_{\mathrm{A}}(a))\)，我们有 \(f(a)^{*} = f^{*}(a^{*})\)。

第一个断言由 I 第 97 页得出。由前所述，映射 \(\varphi\)（\(\mathcal{O}(\mathrm{Sp}_{\mathrm{A}}(a))\) 到 A 中的）由 \(f \mapsto (f^{*}(a^{*}))^{*}\) 定义，它是 \(\mathcal{O}(\mathrm{Sp}_{\mathrm{A}}(a))\) 到 A 中的一个连续幺态射，使得在 \(\mathrm{Sp}_{\mathrm{A}}(a)\) 的一个邻域中 \(\mathbf{C}\) 的恒同函数的芽的像等于 \(a\)。因此，\(\varphi\) 就是全纯函数演算的映射 \(f \mapsto f(a)\)（I, 第 74 页, 定理 5）。

引理 1. --- 设 D 是 C 中的单位开圆盘。存在唯一的在 D 上定义的全纯函数 f，使得 \(f(z)^{2} = 1 - z\) 对所有 \(z \in D\)，且 \(f(0) = 1\)。我们有 \(f^{*} = f\)。

幂级数

\[
\sum_ {n = 0} ^ {+ \infty} \binom {1 / 2} {n} (- z) ^ {n}
\]

的收敛半径是 1，其和 \(f\) 定义了 D 上的一个全纯函数（VAR, R1, 第 27 页, 3.2.9），它在 0 处取值 1。它满足 \(f(x)=\sqrt{1-x}\) 对所有 \(x \in D \cap R\)（FVR, III, 第 19 页），因而 \(f(z)^{2} = 1 - z\) 对 \(z \in D\)，因为差 \(f(z)^{2} - (1 - z)\) 是一个各阶导数都在 0 处为零的全纯函数（VAR, R1, 第 27 页, 3.2.5）。按定义即可验证 \(f^{*} = f\)。

设 \(g\) 是 D 上的一个全纯函数，使得 \(g(z)^2 = 1 - z\) 对所有 \(z \in \mathrm{D}\)，且 \(g(0) = 1\)。函数 \(g\) 不取零值，于是 D 上的连续函数 \(f / g\) 取值于 \(\{-1, 1\}\)；由于 D 是连通的且 \(f(0) = g(0)\)，我们有 \(f = g\)。
\subsection{正线性泛函}

定义 2. --- 设 A 是对合代数。A 上的线性泛函 \(\lambda\) 称为正的，如果 \(\lambda(a^{*}a) \in \mathbf{R}_{+}\) 对所有 \(a \in \mathrm{A}\) 都成立。

若 A 是对合巴拿赫代数，我们以 \(A_{+}^{\prime}\) 表示 A 上连续正线性泛函的集合。

设 A 为对合巴拿赫代数。集合 \(A_{+}^{\prime}\) 是 A 上 C-线性泛函所在实向量空间中的尖凸锥。

引理 2. --- 设 A 为对合巴拿赫代数，且 \(\lambda\) 是 A 上的正线性泛函。

\begin{enumerate}
\def\labelenumi{\alph{enumi})}
\tightlist
\item
  对 A 中任意 \(a\) 和 \(b\)，有 \(\lambda(a^{*}b)=\overline{\lambda(b^{*}a)}\)，且
\end{enumerate}

\[
\left| \lambda \left(a ^ {*} b\right) \right| ^ {2} \leqslant \lambda \left(a ^ {*} a\right) \lambda \left(b ^ {*} b\right);
\]

\begin{enumerate}
\def\labelenumi{\alph{enumi})}
\setcounter{enumi}{1}
\tightlist
\item
  若 A 含单位元，则线性泛函 \(\lambda\) 连续，且其范数等于 \(\lambda(1)\)。
\end{enumerate}

映射 \((a,b)\mapsto\lambda(a^{*}b)\) 是 A 上的正埃尔米特形式；因此，对 A 中任意 a 和 b，它满足 \(\lambda(a^{*}b)=\overline{\lambda(b^{*}a)}\) 以及 \(|\lambda(a^{*}b)|^{2}\leqslant\lambda(a^{*}a)\lambda(b^{*}b)\)（EVT, V, p. 2, 注以及 p. 3, 命题 2）。

证明 \(b\)。设 \(a \in A\) 是范数 \(< 1\) 的埃尔米特元。\(a\) 的谱半径小于 \(\|a\|\)，所以 \(a\) 的谱包含于 C 的开单位圆盘 D 中（I, p. 24, 定理 1）。设 \(f\) 是 D 上满足 \(f(z)^{2} = 1 - z\) 且对所有 \(z \in D\) 成立的全纯函数（V, p. 435, 引理 1）。对元 \(a\) 和函数 \(f\) 应用全纯函数演算。元 \(b = f(a)\) 满足 \(b^{2} = 1 - a\)（I, p. 74, 定理 5）。由 V, p. 435, 命题 3，还有 \(f(a)^{*} = f^{*}(a^{*}) = f(a)\)，故 \(b\) 是埃尔米特元。于是 \(\lambda(1 - a) = \lambda(b^{*}b) \geqslant 0\)，从而 \(\lambda(a) \leqslant \lambda(1)\)。

现在设 \(b \in A\) 且范数 \(< 1\)。元 \(b^*b\) 是埃尔米特元且 \(\|b^*b\| < 1\)，因此在应用 \(a\) )（取 \(a = 1\)）的情形，得到

\[
| \lambda (b) | ^ {2} \leqslant \lambda (1) \lambda (b ^ {*} b) \leqslant \lambda (1) ^ {2},
\]

当 \(b = 1\) 时取等号。因此线性泛函 \(\lambda\) 连续，且其范数等于 \(\lambda(1)\)。断言 \(b\) ) 得证。

例。--- 设 X 是紧拓扑空间，且 A 是星代数 \(\mathcal{C}(\mathrm{X})\)。A 上的正线性泛函与 X 上的正测度相对应（INT, III, p. 52, § 1, n° 6, 定理 1）。

引理 3. --- 设 A 是具有逼近单位元的含单位元对合巴拿赫代数（I, p. 120, 定义 7）。设 λ 是 A 上的连续正线性泛函。

\begin{enumerate}
\def\labelenumi{\alph{enumi})}
\item
  对 A 中任意 a，有 \(\lambda(a^{*})=\overline{\lambda(a)}\) 以及 \(|\lambda(a)|^{2}\leqslant\|\lambda\|\lambda(a^{*}a)\)；
\item
  设 \(\widetilde{\mathrm{A}}\) 是通过向 A 添加单位元得到的对合代数，e 是其单位元。存在 \(\widetilde{\lambda}\) 在 \(\widetilde{\mathrm{A}}\) 上的连续正线性泛函，它延拓 \(\lambda\)，且满足 \(\widetilde{\lambda}(e) = \|\lambda\|\)；
\item
  对 A 中任意 \(a\) 和 \(b\)，有 \(|\lambda(b^{*}ab)|\leqslant\|a\|\lambda(b^{*}b)\)。
\end{enumerate}

证明 \(a\) )。设 \(\mathfrak{F}\) 是 A 的逼近单位元，且 \(a \in \mathrm{A}\)。利用引理 2、\(a\) ) 以及逼近单位元的定义，可得

\[
\lambda (a ^ {*}) = \lim _ {f, \mathfrak {F}} \lambda (f a ^ {*}) = \lim _ {f, \mathfrak {F}} \overline {{\lambda (a f ^ {*})}} = \lim _ {f, \mathfrak {F}} \overline {{\lambda ((f a ^ {*}) ^ {*})}} = \overline {{\lambda (a)}},
\]

因此（同上）

\[
| \lambda (a) | ^ {2} = \lim _ {f, \mathfrak {F}} | \lambda (f a) | ^ {2} \leqslant \lambda (a ^ {*} a) \operatorname * {limsup} _ {f, \mathfrak {F}} \lambda (f f ^ {*}) \leqslant \| \lambda \| \lambda (a ^ {*} a),
\]

因为 \(\mathfrak{F}\) 是 A 的单位球上的滤子。

证明 \(b\) )。不妨设 \(\lambda\) 非零，进而设 \(\| \lambda \| = 1\)。对 \(a \in \mathrm{A}\) 和 \(z \in \mathbf{C}\)，定义 \(\widetilde{\lambda}(a + z \cdot e) = \lambda(a) + z\)。映射 \(\widetilde{\lambda}\) 是 \(\widetilde{\mathrm{A}}\) 上延拓 \(\lambda\) 的连续线性泛函，并满足 \(\widetilde{\lambda}(e) = 1\)。它是正的：对任意 \(a \in \mathrm{A}\) 和 \(z \in \mathbf{C}\)，根据 \(a\) )，计算得

\[
\begin{array}{r l} \widetilde {\lambda} ((a + z \cdot e) ^ {*} (a + z \cdot e)) & = \lambda (a ^ {*} a) + \overline {{z}} \lambda (a) + z \lambda (a ^ {*}) + | z | ^ {2} \\ & = | z + \lambda (a) | ^ {2} + \lambda (a ^ {*} a) - | \lambda (a) | ^ {2} \geqslant 0. \end{array}
\]

最后，设 \(b\) 是 A 的一个元。映射 \(a \mapsto \widetilde{\lambda}(b^{*}ab)\) 是含单位元对合巴拿赫代数 \(\widetilde{A}\) 上的正线性泛函。因此它连续，且范数等于其在 \(e\) 处的值（引理 2，\(b\) )），而该值等于 \(\lambda(b^{*}b)\)，从而得到断言 \(c\) )。

命题 4. — 设 A 为对合巴拿赫代数。

\begin{enumerate}
\def\labelenumi{\alph{enumi})}
\item
  对每个希尔伯特空间 E、每个对合代数态射 \(\varphi\colon A\to\mathcal{L}(E)\) 以及每个向量 \(x\in E\)，在 A 上由 \(\lambda(a)=\langle x|\varphi(a)x\rangle\) 定义的线性泛函是连续正线性泛函；
\item
  设 \(\lambda \in \mathrm{A}_{+}^{\prime}\)。映射 \(f\) 从 \(\mathrm{A} \times \mathrm{A}\) 到 \(\mathbf{C}\)，由 \(f(a, b) = \lambda (a^{*}b)\) 对所有 \((a, b) \in \mathrm{A} \times \mathrm{A}\) 定义，是 \(\mathrm{A}\) 上的普遍正核。
\end{enumerate}

证明 \(a\) )。对每个 \(a \in \mathrm{A}\)，有

\[
\lambda (a ^ {*} a) = \langle x \mid \varphi (a ^ {*} a) x \rangle = \| \varphi (a) x \| ^ {2} \geqslant 0
\]

所以 \(\lambda\) 是 A 上的正线性泛函；由于 \(\varphi\) 连续，它也是连续的（I, p. 104, 命题 2）。

证明 \(b\) )。函数 \(f\) 连续。对每个 \(n \in \mathbf{N}\)、A 中的每个族 \((a_i)_{0 \leqslant i \leqslant n}\) 以及每个复数族 \((t_i)_{0 \leqslant i \leqslant n}\)，有

\[
\sum_ {i = 0} ^ {n} \sum_ {j = 0} ^ {n} \bar {t} _ {i} t _ {j} f (a _ {i}, a _ {j}) = \lambda \left(\left(\sum_ {i = 0} ^ {n} t _ {i} a _ {i}\right) ^ {*} \left(\sum_ {j = 0} ^ {n} t _ {j} a _ {j}\right)\right) \geqslant 0,
\]

从而得到结论（V, p. 433, 定义 1）。

定义 3. — 设 A 是对合巴拿赫代数，且 \(\lambda \in \mathrm{A}_{+}^{\prime}\) 是 A 上的连续正线性泛函。\(\lambda\) 的希尔伯特实现是一个三元组 (E, \(x, \varphi\) )，其中 E 是希尔伯特空间，\(x\) 是 E 的一个元，\(\varphi\) 是从 A 到 \(\mathcal{L}(\mathrm{E})\) 的对合代数态射，并且 \(\lambda(b^{*}a) = \langle \varphi(b)x | \varphi(a)x \rangle\) 对所有 \((a,b) \in \mathrm{A}^{2}\) 都成立。

如果元素 \(\varphi(a)x\)（其中 \(a\in\mathrm{A}\)）在 E 中构成全集，则称 (E, \(x,\varphi\) ) 是 \(\lambda\) 的循环希尔伯特空间实现。

命题 5. --- 设 A 是具有逼近单位元的对合巴拿赫代数。A 上的每个连续正线性泛函都有一个循环希尔伯特实现 (E, x, φ)。若 A 含单位元，则存在一个其中态射 φ 保单位元的此类实现。

由引理 3，b)，不妨设 A 是含单位元的代数。

设 \(\lambda\) 是 A 上的连续正线性泛函。设 (E, \(g\) ) 是 A 上的普遍正核 \(f\) 的循环希尔伯特实现，该核由 \(f(a,b) = \lambda(a^{*}b)\) 定义（命题 4 和定理 1）。映射 \(g\) 连续，其像在 E 中是全的，并且 \(\lambda(a^{*}b) = \langle g(a)|g(b)\rangle\) 对所有 \((a,b) \in \mathrm{A}^2\) 都成立。置 \(x = g(1) \in \mathrm{E}\)。

从 A 到 E 的映射 g 是线性的：事实上，对 \((a,b,c)\in\mathrm{A}^{3}\) 和 \((s,t)\in\mathbf{C}^{2}\)，有

\[
\begin{array}{c} \langle g (c) | g (s a + t b) \rangle = \lambda (c ^ {*} (s a + t b)) \\ = s \lambda (c ^ {*} a) + t \lambda (c ^ {*} b) = \langle g (c) | s g (a) + t g (b) \rangle , \end{array}
\]

由于 \(g\) 的像在 E 中是全的，故断言成立。

特别地，g 的像 F 是 E 的稠密向量子空间。g 的核是 A 的左理想：若 \(g(b) = 0\)，则对 A 中任意 a 和 c，有

\[
\langle g (a b) \mid g (c) \rangle = \lambda ((a b) ^ {*} c) = \lambda (b ^ {*} (a ^ {*} c)) = \langle g (b) \mid g (a ^ {*} c) \rangle = 0,
\]

由于 F 在 E 中稠密，所以 \(g(ab) = 0\)。

设 \(a \in \mathrm{A}\)。由于 \(g\) 的核是 \(\mathrm{A}\) 的左理想，存在一个线性映射 \(\widetilde{\varphi}(a) \colon \mathrm{F} \to \mathrm{F}\)，使得 \(\widetilde{\varphi}(a)(g(b)) = g(ab)\) 对所有 \(b \in \mathrm{A}\) 成立。此外，对每个 \(b \in \mathrm{A}\)，有

\[
\| \widetilde {\varphi} (a) (g (b)) \| ^ {2} = \| g (a b) \| ^ {2} = \lambda \left(b ^ {*} a ^ {*} a b\right) \leqslant \| a ^ {*} a \| \lambda \left(b ^ {*} b\right) \leqslant \| a \| ^ {2} \| g (b) \| ^ {2}
\]

（引理 3，c)），所以 \(\widetilde{\varphi}(a)\) 连续且范数 \(\leqslant \|a\|\)。因此存在唯一的自同态 \(\varphi(a) \in \mathcal{L}(\mathrm{E})\)，它通过对子空间的传递诱导出 \(\widetilde{\varphi}(a)\)。

设 \(a\) 和 \(b\) 属于 A。有

\[
(\varphi (a) \circ \varphi (b)) (g (c)) = g (a b c) = \varphi (a b) (g (c))
\]

对 A 中任意 c 都成立，因此由于 F 在 E 中稠密，\(\varphi(ab) = \varphi(a) \circ \varphi(b)\)。

映射 \(\varphi\) 从 A 到 \(\mathcal{L}(\mathrm{E})\) 是线性的：事实上，对任意 \((a,b)\in\mathrm{A}^{2}\) 和 \((s,t)\in\mathbf{C}^{2}\)，有

\[
\varphi (s a + t b) (g (c)) = g ((s a + t b) c) = (s \varphi (a) + t \varphi (b)) g (c)
\]

对 A 中任意 \(c \in A\) 都成立，从而由于 F 在 E 中稠密，\(\varphi(sa + tb) = s\varphi(a) + t\varphi(b)\)。由于 \(\|\varphi(b)\| \leqslant \|b\|\) 对所有 \(b \in A\) 成立，映射 \(\varphi\) 连续。又有 \(\varphi(1) = 1_{\mathrm{E}}\)，因为 \(\varphi(1)(g(a)) = g(a)\) 对所有 \(a \in A\) 成立，且 F 在 E 中稠密。

最后，设 a、b、c 属于 A。有

\[
\begin{array}{c} \langle g (c) \mid \varphi (a ^ {*}) (g (b)) \rangle = \langle g (c) \mid g (a ^ {*} b) \rangle = \lambda (c ^ {*} a ^ {*} b) \\ = \langle g (a c) \mid g (b) \rangle = \langle \varphi (a) (g (c)) \mid g (b) \rangle \end{array}
\]

由于 F 在 E 中稠密，从而 \(\varphi(a^{*})=\varphi(a)^{*}\)。

总之，映射 \(\varphi\) 是从 A 到 \(\mathcal{L}(\mathrm{E})\) 的连续对合代数态射，且 \(g(a) = \varphi(a)x\) 对所有 \(a \in \mathrm{A}\) 成立；因此三元组 \((\mathrm{E}, x, \varphi)\) 是 \(\lambda\) 的循环希尔伯特实现。

命题 6. --- 设 A 是对合巴拿赫代数，λ 是 A 上的连续正线性泛函，且 (E₁, x₁, φ₁) 和 (E₂, x₂, φ₂) 是 λ 的希尔伯特实现。设 (E₁, x₁, φ₁) 是循环希尔伯特实现。

\begin{enumerate}
\def\labelenumi{\alph{enumi})}
\item
  存在唯一的从 \(E_{1}\) 到 \(E_{2}\) 的连续线性映射 u，它是从 \(\varphi_{1}\) 到 \(\varphi_{2}\) 的表示态射（I, p. 11），并且满足 \(u(\varphi_{1}(a)x_{1}) = \varphi_{2}(a)x_{2}\) 对所有 \(a \in A\) 成立；
\item
  线性映射 u 是等距的；
\item
  若 \(\left(\mathrm{E}_{2},x_{2},\varphi_{2}\right)\) 是循环的，则 u 是同构；
\item
  若 \(\left(\mathrm{E}_2,x_2,\varphi_2\right)\) 是循环的且 A 含单位元，则 \(u(x_{1}) = x_{2}\)。
\end{enumerate}

对 \(j \in \{1,2\}\)，定义 \(\gamma_j \colon A \to E_j\)，使得 \(\gamma_j(a) = \varphi_j(a)x_j\) 对所有 \(a \in A\) 成立。按定义，\((E_j, \gamma_j)\) 是由普遍正核 \(f\) 在 \(A\) 上定义的希尔伯特实现，该核由 \((a,b) \mapsto \lambda(a^*b)\) 定义。希尔伯特实现 \((E_1, \gamma_1)\) 是循环的；由 V, p. 434, 命题 1，存在唯一的连续线性映射 \(u \colon E_1 \to E_2\)，使得 \(\gamma_2 = u \circ \gamma_1\)，且该映射是等距的。此外，若 \((E_2, x_2, \varphi_2)\) 也是循环的，则 \(u\) 是同构。要证明 \(a)\)、\(b)\) 和 \(c)\)，只需证明 \(u\) 是从 \(\varphi_1\) 到 \(\varphi_2\) 的表示态射。

设 \(a \in \mathrm{A}\)。对每个 \(b \in \mathrm{A}\)，有

\[
\begin{array}{c} (u \circ \varphi_ {1} (a)) (\gamma_ {1} (b)) = (u \circ \varphi_ {1} (a b)) x _ {1} = (u \circ \gamma_ {1}) (a b) \\ = \gamma_ {2} (a b) = \varphi_ {2} (a) (\gamma_ {2} (b)) = \varphi_ {2} (a) (u (\gamma_ {1} (b))), \end{array}
\]

因此，连续线性映射 \(u \circ \varphi_1(a)\) 和 \(\varphi_2(a) \circ u\) 在 \(E_1\) 中由 \(\gamma_1\) 的像生成的子空间上一致；由假设，该子空间在 \(E_1\) 中稠密，所以 \(u \circ \varphi_1(a) = \varphi_2(a) \circ u\)。

最后证明 \(d\) )。因此设 A 含单位元，且 \((\mathrm{E}_2, x_2, \varphi_2)\) 是循环的。存在循环希尔伯特空间实现 \((\mathrm{E}_3, x_3, \varphi_3)\) 的 \(\lambda\)，使得 \(\varphi_3\) 是保单位元的态射（命题 5）。设 \(j \in \{1, 2\}\)。将前述结果应用于 \((\mathrm{E}_j, x_j, \varphi_j)\) 和 \((\mathrm{E}_3, x_3, \varphi_3)\)，可知存在等距同构 \(\widetilde{u}_j\)，从 \(\mathrm{E}_j\) 到 \(\mathrm{E}_3\)，使得 \(\widetilde{u}_j \circ \varphi_j(a) = \varphi_3(a) \circ \widetilde{u}_j\) 对所有 \(a \in \mathrm{A}\) 成立。取 \(a = 1_{\mathrm{A}}\)，可知 \(\varphi_j\) 保单位元。于是 \(x_2 = \varphi_2(1_{\mathrm{A}})x_2 = u(\varphi_1(1_{\mathrm{A}})x_1) = u(x_1)\)。

注。--- 保持命题中的记号。可能 \(u(x_{1})\) 不同于 \(x_{2}\)（V, p. 493, 习题 3，b)）。不过，三元组 \((\mathrm{E}_{2}, u(x_{1}), \varphi_{2})\) 也是 \(\lambda\) 的希尔伯特实现。

\subsection{星代数的表示}

设 A 是星代数。以 \(A_{+}\) 表示 A 的正元组成的闭凸锥（I, p. 115, 定义 6）。A 上的线性泛函 \(\lambda\) 是正的，当且仅当 \(\lambda(\mathrm{A}_{+}) \subset \mathbf{R}_{+}\)（I, p. 118, 定理 2）。

命题 7. --- 设 A 是星代数。A 上的每个正线性泛函 \(\lambda\) 都是连续的。

先证明 \(\lambda\) 在 \(\mathrm{A}_{+}\) 与 A 的单位球的交集上有界。反之，存在序列 \((x_{n})_{n\geqslant 1}\) 属于 \(\mathrm{A}_{+}\)，满足 \(\| x_n\| \leqslant 1\) 且 \(\lambda (x_n)\geqslant n\)，对每个整数 \(n\geqslant 1\) 都成立。通项为 \(n^{-2}x_{n}\) 的级数收敛到 A 中的一个元 \(x\)（参见 TG, IV, p. 33, 例 3）。对每个整数 \(\mathrm{N}\geqslant 1\)，由于

\[
\sum_ {n = \mathrm{N} + 1} ^ {+ \infty} \frac {1}{n ^ {2}} x _ {n} \in \mathrm{A} _ {+},
\]

由（I, p. 116, 命题 14）可得

\[
\sum_ {n = 1} ^ {\mathrm{N}} \frac {1}{n} \leqslant \sum_ {n = 1} ^ {\mathrm{N}} \frac {1}{n ^ {2}} \lambda (x _ {n}) = \lambda (x) - \lambda \left(\sum_ {n = \mathrm{N} + 1} ^ {+ \infty} \frac {1}{n ^ {2}} x _ {n}\right) \leqslant \lambda (x),
\]

这与事实矛盾（TG, IV, p. 33, 例 4）。

由于 A 的单位球中的每个元都是 A 中至多四个范数 \(\leqslant 1\) 的正元的系数绝对值不超过 1 的线性组合（I, p. 96, 引理 2 及 I, p. 117, 公式 (4)），所以 \(\lambda\) 连续。

存在一些对合巴拿赫代数，它们具有不连续的正线性泛函（V, p. 493, 习题 3，a)）。

命题 8. --- 设 A 是星代数，a 是 A 的非零元。存在正线性泛函 \(\lambda \in \mathrm{A}_{+}^{\prime}\)，使得 \(\lambda(a^{*}a) > 0\)。

考虑 A 的埃尔米特元组成的实巴拿赫空间 \(A_{h}\)。集合 \(A_{+}\) 是 \(A_{h}\) 中的闭凸尖锥（I, p. 116, 命题 14）。元 \(a^{*}a\) 是正的（I, p. 118, 定理 2）且非零，故埃尔米特元 \(-a^{*}a\) 不是正的。由 EVT, II, p. 42, 推论 5，存在连续实线性泛函 \(\lambda\in A_{h}^{\prime}\)，使得 \(\lambda(-a^{*}a)<0\) 且 \(\lambda(\mathrm{A}_{+})\subset\mathbf{R}_{+}\)。线性泛函 \(\lambda\) 延拓为 A 上的埃尔米特 C-线性泛函（参见 I, p. 98），该泛函是正的，并具有所要求的性质。

定理 2（盖尔范德--奈马克）。--- 设 A 是星代数。存在希尔伯特空间 E 以及 \(\varphi\) 从 A 到 \(\mathcal{L}(\mathrm{E})\) 的等距对合代数态射。若 A 含单位元，则存在这样的保单位元态射。

对每个 \(b \in \mathrm{A} - \{0\}\)，取 A 上的连续正线性泛函 \(\lambda_b\)，使得 \(\lambda_b(b^*b) > 0\)（命题 8）。由于 A 具有逼近单位元（I, p. 121, 命题 18），存在希尔伯特实现 \((\mathrm{E}_b, x_b, \varphi_b)\) 的 \(\lambda_b\)（命题 5）。若 A 含单位元，可设 \(\varphi_b\) 保单位元（同上）。有 \(\| \varphi_b(b) \|^2 = \lambda_b(b^*b) \neq 0\)。

令 E 为各空间 \(E_{b}\)（其中 \(b\) 属于 \(A - \{0\}\)）的外希尔伯特直和。对每个 \(a \in A\)，以 \(\varphi(a) \in \mathcal{L}(E)\) 表示唯一的连续线性映射，其在每个 \(E_{b}\) 上的限制都与 \(\varphi_{b}(a)\) 一致，其中 \(b \in A - \{0\}\)。映射 \(a \mapsto \varphi(a)\) 是对合代数态射；它是单射，因而是等距的（I, p. 112, 命题 9）；若 A 含单位元，还满足 \(\varphi(1) = 1_{\mathrm{E}}\)。证明完毕。

*因此，以星代数为对象、以对合代数态射为态射的范畴，等价于以希尔伯特空间的自同态代数的闭对合子代数为对象、以对合代数态射为态射的范畴。*

\subsection{拓扑群上的正定函数}

在本节中，G 是拓扑群，其单位元记为 e。所考虑的希尔伯特空间都是复的。

定义 4. --- 连续函数 \(\varphi\in\mathcal{C}(\mathrm{G})\) 称为 G 上的正定函数，如果函数 \(f\) 由 \(f(g,h)=\varphi(g^{-1}h)\) 定义在 \(\mathrm{G}\times\mathrm{G}\) 上，且是 \(\mathrm{G}\) 上的普遍正核。

以 \(\operatorname{Pos}(\mathbf{G})\) 表示 G 上正定函数的集合，并以 \(\operatorname{Pos}_{1}(\mathbf{G})\)（分别以 \(\operatorname{Pos}_{\leqslant1}(\mathbf{G})\)）表示 \(\varphi\in\operatorname{Pos}(\mathbf{G})\) 中满足 \(\varphi(e)=1\)（分别满足 \(\varphi(e)\leqslant1\)）的函数所成的子集。

例。--- 设 \(\varrho\) 是 G 在希尔伯特空间 E 上的酉表示，且 \(x \in \mathrm{E}\)。设 \(\varphi\) 是由 \(\varphi(g) = \langle x \mid \varrho(g)x \rangle\) 对所有 \(g \in \mathrm{G}\) 定义的对角矩阵系数；函数 \(\varphi\) 连续。由此可得

\[
\varphi (g ^ {- 1} h) = \langle x | \varrho (g) ^ {*} \varrho (h) x \rangle = \langle \varrho (g) x | \varrho (h) x \rangle
\]

对所有 \((g,h)\in\mathrm{G}\times\mathrm{G}\) 都成立。因此 \(\varphi\in\mathrm{Pos}(\mathrm{G})\)（V, p. 432, 定理 1, (iv)）。\(\varphi\in\mathrm{Pos}_{1}(\mathrm{G})\) 当且仅当 \(\|x\|=1\)。

定义 5. --- 设 \(\varphi\) 是 G 上的正定函数。\(\varphi\) 的希尔伯特实现是一个对 \((\varrho, x)\)，其中 \(\varrho\) 是 G 在希尔伯特空间 E 上的酉表示，\(x \in \mathrm{E}\)，且 \(\varphi(g) = \langle x \mid \varrho(g)x \rangle\) 对所有 \(g \in \mathrm{G}\) 成立。

若 x 是 \(\varrho\) 的循环向量，则称其为 \(\varphi\) 的循环希尔伯特实现。

命题 9. --- 设 \(\varphi\in\mathcal{C}(\mathrm{G})\)。以下条件等价：

\begin{enumerate}
\def\labelenumi{(\roman{enumi})}
\item
  函数 \(\varphi\) 是 G 上的正定函数；
\item
  存在 \(\varphi\) 的循环希尔伯特实现；
\item
  存在 \(\varphi\) 的希尔伯特实现。
\end{enumerate}

条件 (ii) 蕴含条件 (iii)，而由上述例子，条件 (iii) 蕴含条件 (i)。

最后证明 (i) 蕴含 (ii)。设 (E, \(\gamma\) ) 是由 \(f(g,h)=\varphi(g^{-1}h)\) 定义的普遍正核（其中 \((g,h)\in\mathrm{G}\times\mathrm{G}\)）的循环希尔伯特实现。设 \(k\in\mathrm{G}\)。G 上的连续函数 \(\gamma_{k}:g\mapsto\gamma(kg)\) 满足

\[
\langle \gamma_ {k} (g) | \gamma_ {k} (h) \rangle = \langle \gamma (k g) | \gamma (k h) \rangle = f (k g, k h) = \varphi ((k h) ^ {- 1} k g) = f (g, h)
\]

对所有 \((g,h)\in\mathrm{G}\times\mathrm{G}\) 都成立，因此 \((\mathrm{E},\gamma_{k})\) 是 f 的希尔伯特实现。由 V, p. 434, 命题 1，存在唯一的酉算子 \(\varrho(k)\) 属于 \(\mathcal{L}(\mathrm{E})\)，使得 \(\gamma_{k}=\varrho(k)\circ\gamma\)。对每个 \(g\in G\) 以及每个 \((k_{1},k_{2})\in\mathrm{G}\times\mathrm{G}\)，可得

\[
\varrho (k _ {1} k _ {2}) (\gamma (g)) = \gamma (k _ {1} k _ {2} g) = \varrho (k _ {1}) (\varrho (k _ {2}) (\gamma (g))),
\]

从而 \(\varrho(k_1k_2) = \varrho(k_1)\varrho(k_2)\)，因为 \(\gamma\) 的像是 E 的全集。

设 \(k \in G\)。对每个 \(g \in G\)，有 \(\varrho(g)(\gamma(k)) = \gamma(gk)\)，因此由 \(g \mapsto \varrho(g)(\gamma(k))\) 定义的从 G 到 E 的映射连续。由于自同态 \(\varrho(g)\) 是酉的，故可推出 \(\varrho\) 是 G 在 E 上的酉表示（V, p. 380, 引理 4）。再置 \(x = \gamma(e)\)。向量集 \(\varrho(g)x = \varrho(g)(\gamma(e)) = \gamma(g)\) 在 E 中是全的，其中 \(g\) 遍历 G，故 x 是 \(\varrho\) 的循环向量。由于

\[
\langle x \mid \varrho (g) x \rangle = f (e, g) = \varphi (g)
\]

对所有 \(g \in G\) 都成立，所以对 \((\varrho, x)\) 是 \(\varphi\) 的循环希尔伯特实现。

命题 10. — 设 \(\varphi\) 是 G 上的正定函数，\((\varrho_{1}, x_{1})\) 和 \((\varrho_{2}, x_{2})\) 是 \(\varphi\) 的希尔伯特实现，且希尔伯特表示 \((\varrho_{1}, x_{1})\) 是循环的。存在唯一的等距 G-态射 \(u\)，从 \(\varrho_{1}\) 到 \(\varrho_{2}\)，使得 \(u(x_{1}) = x_{2}\)。若 \((\varrho_{2}, x_{2})\) 也是循环的，则 \(u\) 是同构。

对 \(1 \leqslant j \leqslant 2\)，令 \(\mathrm{E}_j\) 为 \(\varrho_j\) 的空间，并令 \(\gamma_j\) 为 \(\mathbf{G}\) 上由 \(\gamma_j(g) = \varrho_j(g)x_j\) 对所有 \(g \in \mathbf{G}\) 定义的函数。对 \((\mathrm{E}_1, \gamma_1)\) 和 \((\mathrm{E}_2, \gamma_2)\) 是正定函数 \((g, h) \mapsto \varphi(g^{-1}h)\) 的希尔伯特实现，且 \((\mathrm{E}_1, \gamma_1)\) 是循环希尔伯特实现。由 V, p. 434, 命题 1，存在唯一的等距线性映射 \(u: \mathrm{E}_1 \to \mathrm{E}_2\)，使得 \(\gamma_2 = u \circ \gamma_1\)。特别地，有 \(x_2 = \gamma_2(e) = u(\gamma_1(e)) = u(x_1)\)。此外，对任意 \(g\) 和 \(h\) 属于 \(\mathbf{G}\)，有

\[
\varrho_ {2} (g) u \left(\gamma_ {1} (h)\right) = \varrho_ {2} (g) \gamma_ {2} (h) = \gamma_ {2} (g h) = u \left(\gamma_ {1} (g h)\right) = u \left(\varrho_ {1} (g) \gamma_ {1} (h)\right),
\]

由于元素 \(\gamma_{1}(h)\)（其中 \(h \in G\)）所成的集合在 \(E_{1}\) 中是全的，这意味着 u 是 G-态射。根据同上，若 \((\varrho_{2}, x_{2})\) 也是循环的，则 u 是等距同构。

命题 11. — 设 \(\varphi \in \mathrm{Pos}_{1}(\mathrm{G})\)，且 \((\varrho, x)\) 是 \(\varphi\) 的循环希尔伯特实现。则 \(\varrho\) 是 G 的不可约表示，当且仅当 \(\varphi\) 是 \(\mathrm{Pos}_{1}(\mathrm{G})\) 的极点（EVT, II, p. 57, 定义 1）。

令 E 为 \(\varrho\) 的空间。先设 \(\varrho\) 不是不可约的。令 F 是 E 的一个在 \(\varrho\) 下稳定的、非零且不同于 E 的闭子空间。写作 \(x = x_1 + x_2\)，其中 \(x_1 \in \mathrm{F}\)，\(x_2 \in \mathrm{F}^\circ\)。于是 \(1 = \|x\|^2 = \|x_1\|^2 + \|x_2\|^2\)。由 \(x_1\) 生成的 E 的子表示包含于 F 中，所以 \(x_1 \neq x\)，因为 \(x\) 是 \(\varrho\) 的循环向量，从而 \(x_2 \neq 0\)。同理可验证 \(x_1 \neq 0\)。

对 \(j = 1\) 和 \(j = 2\)，令 \(\varphi_j\) 为 G 上满足

\[
\varphi_ {j} (g) = \frac {1}{\| x _ {j} \| ^ {2}} \langle x _ {j} | \varrho (g) x _ {j} \rangle
\]

对所有 \(g \in G\) 均成立。我们有 \(\varphi_{j} \in \mathrm{Pos}_{1}(G)\)（V, p. 443, 例）。由于 \(\varphi = \|x_{1}\|^{2}\varphi_{1} + \|x_{2}\|^{2}\varphi_{2}\)，只需验证 \(\varphi_{1} \neq \varphi_{2}\) 即可证明 \(\varphi\) 不是 \(\mathrm{Pos}_{1}(G)\) 的极点；为此只需证明 \(\varphi \neq \varphi_{1}\)。

反证。假设 \(\varphi = \varphi_1\)。由于 \(\langle x_1 | \varrho(g)x_2 \rangle = 0\) 对所有 \(g \in G\) 成立，应有

\[
\frac {1}{\| x _ {1} \| ^ {2}} \langle x _ {1} | \varrho (g) x \rangle = \frac {1}{\| x _ {1} \| ^ {2}} \langle x _ {1} | \varrho (g) x _ {1} \rangle = \varphi_ {1} (g) = \varphi (g) = \langle x | \varrho (g) x \rangle
\]

对所有 \(g \in G\) 都成立，从而 E 的向量子空间中的每个元素 y 都满足 \(\langle x_{1} | y \rangle = \langle \|x_{1}\|^{2}x | y \rangle\)，其中该子空间由元素 \(\varrho(g)x\)（其中 \(g \in G\)）生成；因此对于所有 \(y \in E\)，由于 x 是 \(\varrho\) 的循环向量，这都成立。这将意味着 \(x_{1} = \|x_{1}\|^{2}x\) 也是 \(\varrho\) 的循环向量，矛盾，故断言成立。

现在设 \(\varrho\) 不可约，并证明 \(\varphi\) 是 \(\mathrm{Pos}_1(\mathrm{G})\) 的极点。设 \(\varphi_1 \neq \varphi_2\) 是 \(\mathrm{Pos}_1(\mathrm{G})\) 中的元，且 \(t_1, t_2 \in [0,1]\) 满足 \(t_1 + t_2 = 1\) 和 \(\varphi = t_1\varphi_1 + t_2\varphi_2\)。对 \(j \in \{1,2\}\)，设 \((\varrho_j, x_j)\) 是 \(\varphi_j\) 的循环希尔伯特实现，并令 \(\mathrm{E}_j\) 为 \(\varrho_j\) 的空间。令 \(x_3 = t_1^{1/2}x_1 + t_2^{1/2}x_2\)。则 \((\varrho_1 \oplus \varrho_2, x_3)\) 是 \(\varphi\) 的希尔伯特实现。由于 \((\varrho, x)\) 是循环的，存在等距 G-态射 \(u: \mathrm{E} \to \mathrm{E}_1 \oplus \mathrm{E}_2\)，使得 \(u(x) = x_3\)（命题 10）。

令 \(j = 1\) 或 \(j = 2\)。由于 \(\varrho\) 不可约，存在 \(\lambda_j \geqslant 0\)，使得 G-态射 \(u_j = \mathrm{pr}_j \circ u\) 从 E 到 \(E_j\) 满足

\[
\langle u _ {j} (y) \mid u _ {j} (y ^ {\prime}) \rangle = \lambda_ {j} \langle y \mid y ^ {\prime} \rangle
\]

对 E 中任意 \(y\) 和 \(y'\) 都成立（若 \(u_j \neq 0\)，这是 V, p. 388, 推论 5；否则可取 \(\lambda_j = 0\)）。对每个 \(g \in G\)，可得

\[
\begin{array}{r} t _ {j} \varphi_ {j} (g) = \langle t _ {j} ^ {1 / 2} x _ {j} | \varrho_ {j} (g) (t _ {j} ^ {1 / 2} x _ {j}) \rangle = \langle u _ {j} (x) | \varrho_ {j} (g) (u _ {j} (x)) \rangle \\ = \langle u _ {j} (x) | u _ {j} (\varrho (g) x) \rangle = \lambda_ {j} \varphi (g). \end{array}
\]

由于 \(\varphi_{j}(e) = \varphi(e) = 1\)，可推出 \(\varphi_{j} = \varphi\) 当 \(t_{j} \neq 0\) 时成立。因此假设 \(\varphi_{1} \neq \varphi_{2}\) 蕴含 \(t_{1}\) 或 \(t_{2}\) 必为零，这就证明了 \(\varphi\) 是 \(\mathrm{Pos}_{1}(\mathrm{G})\) 的极点。

引理 4. --- 设 \(\varphi\in\mathrm{Pos}(\mathbf{G})\)。函数 \(\varphi\) 在 \(\mathbf{G}\) 上有界，且 \(\|\varphi\|_{\infty}=\varphi(e)\)。此外，有

\[
| \varphi (g ^ {- 1} h) - \varphi (h) | \leqslant \sqrt {2 \varphi (e) (\varphi (e) - \mathcal {R} (\varphi (g)))}\tag{1}
\]

对所有 \((g,h)\in \mathbf{G}\times \mathbf{G}\) 都成立。

设 \((\varrho, x)\) 是 \(\varphi\) 的希尔伯特实现。有 \(\varphi(e) = \|x\|^2\)。因此对每个 \(g \in G\)，由柯西–施瓦茨不等式可得 \(|\varphi(g)| = |\langle x | \varrho(g)x\rangle| \leqslant \varphi(e)\)。这就证明了第一个断言。

对每个 \((g,h)\in \mathrm{G}\times \mathrm{G}\)，有

\[
\varphi (g ^ {- 1} h) - \varphi (h) = \langle x | \varrho (g ^ {- 1} h) x \rangle - \langle x | \varrho (h) x \rangle = \langle \varrho (g) x - x | \varrho (h) x \rangle
\]

由于 \(\varrho\) 是酉的，从而

\[
| \varphi (g ^ {- 1} h) - \varphi (h) | \leqslant \| \varrho (g) x - x \| \| \varrho (h) x \| \leqslant \varphi (e) ^ {1 / 2} \| \varrho (g) x - x \|.
\]

注意到

\[
\| \varrho (g) x - x \| ^ {2} = 2 \| x \| ^ {2} - 2 \mathcal {R} (\langle x | \varrho (g) x \rangle) = 2 (\varphi (e) - \mathcal {R} (\varphi (g))).
\]

注。--- 集合 \(\operatorname{Pos}(G)\)（分别为 \(\operatorname{Pos}_{1}(G)\) 和 \(\operatorname{Pos}_{\leqslant1}(G)\)）是星代数 \(\mathcal{C}_{b}(G)\) 的凸自伴子集。它们在赋予逐点收敛拓扑的空间 \(\mathcal{C}_{b}(G)\) 中是闭的。集合 \(\operatorname{Pos}(G)\) 是实巴拿赫空间 \(\mathcal{C}_{b}(G)\) 中顶点为 0 的凸锥。

\subsection{局部紧群的酉对偶}

设 G 是局部紧拓扑群。给 G 赋予左哈尔测度 \(\mu\)。对 \(p \in [1, +\infty]\)，以 \(\mathcal{L}^{p}(\mathrm{G})\)（分别以 \(\mathrm{L}^{p}(\mathrm{G})\)）表示空间 \(\mathcal{L}_{\mathbf{C}}^{p}(\mathrm{G}, \mu)\)（分别表示空间 \(\mathrm{L}_{\mathbf{C}}^{p}(\mathrm{G}, \mu)\)）。将空间 \(\mathcal{C}_{b}(\mathrm{G})\) 与其在 \(\mathrm{L}^{\infty}(\mathrm{G})\) 中的像等同起来。

以 \(\Delta\) 表示 G 的模函数。回忆 \(L^1(G)\) 是对合巴拿赫代数，其对合由映射 \(f \mapsto f^*\) 通过取商诱导，其中 \(f^*(y) = \Delta^{-1}(y)\overline{f(y^{-1})}\) 对所有 \(f \in \mathcal{L}^1(G)\) 和 \(y \in G\) 成立（参见 I, p. 99, 例 4）。根据 INT, VIII, p. 172, §4, no. 7, 命题 20，巴拿赫代数 \(L^1(G)\) 具有逼近单位元。

设 \(\varrho\) 是 G 在复希尔伯特空间 E 上的酉表示。映射 \(f\mapsto\varrho(f)\) 是从 L \(^{1}\) (G) 到 \(\mathcal{L}\) (E) 的连续对合代数态射（V, p. 401, 引理 1）；它是 L \(^{1}\) (G) 在 E 上的非退化表示（INT, VIII, p. 139, § 2, no. 7, 命题 10, (i)）。

命题 12. --- 设 E 是复希尔伯特空间，\(\widetilde{\pi}\) 是从 L \(^{1}\) (G) 到 \(\mathcal{L}\) (E) 的对合代数态射。若表示 \(\widetilde{\pi}\) 非退化，则存在 G 在 E 上的唯一酉表示 \(\pi\)，使得 \(\widetilde{\pi}(f)=\pi(f)\) 对所有 \(f\in\mathrm{L}^{1}(\mathrm{G})\) 成立。

唯一性由 INT, VIII, p. 139, § 2, n° 7, 引理 4 的推论 3 得出。

证明 \(\pi\) 的存在性。由 I, p. 104, 命题 2，有 \(\| \widetilde{\pi}\| \leqslant 1\)。令 F 为 E 中由形如 \(\widetilde{\pi}(f)x\) 的向量生成的子空间，其中 \(f\) 属于 \(\mathrm{L}^1 (\mathrm{G})\)，\(x\) 属于 E；由于表示 \(\widetilde{\pi}\) 非退化，F 在 E 中稠密。

对 e 的每个紧邻域 V，令 \(\varphi_{V}\) 为支集包含于 V 且满足 \(\int\varphi_{V}\mu=1\) 的连续正函数。令 B 为 e 的紧邻域滤子的一个基。

设 \(g\) 属于 \(G\)，设 \(f\) 属于 \(L^1(G)\)。有 \((\varphi_V * \varepsilon_g) * f \to \varepsilon_g * f\) 在 \(L^1(G)\) 中沿 \(\mathfrak{B}\) 的截面滤子收敛（INT, VIII, p. 172, § 4, no. 7, 命题 20），所以 \(\widetilde{\pi}(\varphi_V * \varepsilon_g)\widetilde{\pi}(f)\) 收敛到 \(\widetilde{\pi}(\varepsilon_g * f)\)（在 \(\mathcal{L}(E)\) 中）。这意味着 \(\widetilde{\pi}(\varphi_V * \varepsilon_g)\) 沿 \(\mathfrak{B}\) 的截面滤子逐点收敛于映射 \(\pi(g)\)，从 \(F\) 到 \(F\)；由于 \(\| \widetilde{\pi}(\varphi_V * \varepsilon_g) \| \leqslant 1\) 对每个 \(V\)（它是 \(e\) 的紧邻域）成立，该映射由 EVT, III, p. 26, 推论 3 可知是连续的。因此存在唯一的 \(E\) 的自同态，它通过对子空间的传递诱导出 \(\pi(g)\)。仍以 \(\pi(g)\) 表示该自同态；有 \(\| \pi(g) \| \leqslant 1\)。

设 \(f \in \mathrm{L}^1(\mathrm{G})\)。由于 \(\widetilde{\pi}(\varphi_{\mathrm{V}} * \varepsilon_g) \widetilde{\pi}(f)\) 收敛到 \(\widetilde{\pi}(\varepsilon_g * f)\)，有 \(\pi(g) \widetilde{\pi}(f) = \widetilde{\pi}(\varepsilon_g * f)\)。

当 \(g = e\) 时，该关系表明 \(\pi(e)\) 在 \(F\) 上是恒等映射，因而 \(\pi(e) = 1_{E}\)。设 \(g_1\) 和 \(g_2\) 属于 \(G\)。有

\[
\pi (g _ {1}) \pi (g _ {2}) \widetilde {\pi} (f) = \pi (g _ {1}) \widetilde {\pi} (\varepsilon_ {g _ {2}} * f) = \widetilde {\pi} (\varepsilon_ {g _ {1} g _ {2}} * f) = \pi (g _ {1} g _ {2}) \widetilde {\pi} (f),
\]

所以 \(\pi(g_1)\pi(g_2) = \pi(g_1g_2)\) 在 F 上成立，因而在 E 上也成立。这证明了映射 \(g \mapsto \pi(g)\) 是 G 在 E 上的表示。由于 \(\|\pi(g)\| \leqslant 1\) 且 \(\|\pi(g)^{-1}\| = \|\pi(g^{-1})\| \leqslant 1\)，E 的自同态 \(\pi(g)\) 是等距的，因而是酉的（EVT, V, p. 40, 命题 3）。

设 \(f \in \mathrm{L}^{1}(\mathrm{G})\) 且 \(x \in E\)。映射 \(g \mapsto \varepsilon_{g} * f\) 从 G 到 \(\mathrm{L}^{1}(\mathrm{G})\) 连续（INT, VIII, p. 136, § 2, no. 5 及 p. 144, § 3, no. 2, 公式 (5)），且 \(\widetilde{\pi}\) 连续，因此映射 \(g \mapsto \widetilde{\pi}(\varepsilon_{g} * f)x = \pi(g)\widetilde{\pi}(f)x\) 在 \(g = e\) 处连续。由于 F 在 E 中稠密，表示 \(\pi\) 是酉表示（V, p. 380, 引理 4）。

设 \(f_1\) 和 \(f_2\) 属于 \(L^1(G)\)。根据 INT, VIII, p. 127, §1, no. 5, 命题 7，有关系

\[
f _ {1} * f _ {2} = \int_ {\mathrm{G}} f _ {1} (g) (\varepsilon_ {g} * f _ {2}) d \mu (g)
\]

在 L \(^{1}\) (G) 中，从而

\[
\begin{array}{l} \widetilde {\pi} (f _ {1}) \widetilde {\pi} (f _ {2}) = \widetilde {\pi} (f _ {1} * f _ {2}) = \int_ {\mathrm{G}} f _ {1} (g) \widetilde {\pi} (\varepsilon_ {g} * f _ {2}) d \mu (g) \\ \qquad = \int_ {\mathrm{G}} f _ {1} (g) \pi (g) \widetilde {\pi} (f _ {2}) d \mu (g) = \Big (\int_ {\mathrm{G}} f _ {1} (g) \pi (g) d \mu (g) \Big) \widetilde {\pi} (f _ {2}), \end{array}
\]

利用 INT, VI, p. 9, § 1, no. 1, 命题 1，可得

\[
\widetilde {\pi} (f) = \int_ {\mathrm{G}} f (g) \pi (g) d \mu (g) = \pi (f)
\]

对所有 \(f \in \mathrm{L}^{1}(\mathrm{G})\) 都成立。命题得证。

命题 13. --- 设 \(\varphi\in\mathrm{L}^{\infty}(\mathrm{G})\)。以 \(\lambda_{\varphi}\) 表示由 \(f\mapsto\langle f,\varphi\rangle\) 定义在 \(\mathrm{L}^{1}(\mathrm{G})\) 上的连续线性泛函。则 \(\varphi\) 是 G 上某个正定函数的等价类，当且仅当 \(\lambda_{\varphi}\) 是 \(\mathrm{L}^{1}(\mathrm{G})\) 上的连续正线性泛函。

设 \(\lambda_{\varphi}\) 是 \(L^{1}(G)\) 上的连续正线性泛函。令 \((E, x, \widetilde{\pi})\) 为 \(\lambda\) 的循环希尔伯特实现（V, p. 438, 命题 5）。

令 \(\pi\) 为 G 在 E 上的酉表示，使得 \(\pi(f) = \widetilde{\pi}(f)\) 对所有 \(f \in \mathrm{L}^1(\mathrm{G})\) 成立（命题 12）。于是 \(\lambda_{\varphi}(f) = \langle x | \pi(f)x \rangle\) 对所有 \(f \in \mathrm{L}^1(\mathrm{G})\) 成立。

由于

\[
\pi (f) = \int_ {\mathrm{G}} f (g) \pi (g) d \mu (g)
\]

对所有 \(f\in \mathrm{L}^1 (\mathrm{G})\) 都成立，从而

\[
\int_ {\mathrm{G}} f (g) \varphi (g) d \mu (g) = \lambda_ {\varphi} (f) = \langle x | \pi (f) x \rangle = \int_ {\mathrm{G}} f (g) \langle x | \pi (g) x \rangle d \mu (g)
\]

对所有 \(f \in \mathrm{L}^{1}(\mathrm{G})\) 都成立。因此，\(\varphi\) 是 \(\mathrm{L}^{\infty}(\mathrm{G})\) 中由 G 上的函数 \(g \mapsto \langle x | \pi(g)x \rangle\) 所定义的等价类，而该函数是 G 上的正定函数（V, p. 443, 命题 9）。

反之，设 \(\varphi\in\mathrm{Pos}(\mathrm{G})\)。令 \(f\in\mathcal{K}(\mathrm{G})\)。于是有

\[
\begin{array}{c} \lambda_ {\varphi} (f ^ {*} * f) = \int_ {G} \varphi (y) \Big (\int_ {G} \Delta (z) ^ {- 1} \overline {{f (z ^ {- 1})}} f (z ^ {- 1} y) d \mu (z) \Big) d \mu (y) \\ = \int_ {G} \varphi (y) \Big (\int_ {G} \overline {{f (z)}} f (z y) d \mu (z) \Big) d \mu (y) \end{array}
\]

G × G 上由 \((z, y) \mapsto \varphi(y)\overline{f(z)}f(zy)\) 定义的连续函数是有界的，并且

\[
\begin{array}{c} \int_ {\mathrm{G}} ^ {*} | \varphi (y) | \Big (\int_ {\mathrm{G}} ^ {*} | \overline {{f (z)}} f (z y) | d \mu (z) \Big) d \mu (y) \\ \leqslant \varphi (e) \int_ {\mathrm{G}} ^ {*} \int_ {\mathrm{G}} ^ {*} | f (z) | | f (z y) | d \mu (z) d \mu (y) \\ = \varphi (e) \Big (\int_ {\mathrm{G}} ^ {*} | f (z) | d \mu (z) \Big) ^ {2} \end{array}
\]

由 INT, V, p. 94, § 8, no. 3, 命题 8，利用勒贝格–富比尼定理（INT, V, p. 96, § 8, no. 4, 定理 1），可推出

\[
\begin{array}{c} \lambda_ {\varphi} (f ^ {*} * f) = \int_ {\mathrm{G}} \overline {{f (z)}} \Bigl (\int_ {\mathrm{G}} \varphi (y) f (z y) d \mu (y) \Bigr) d \mu (z) \\ = \int_ {\mathrm{G} \times \mathrm{G}} \overline {{f (z)}} f (y) \varphi (z ^ {- 1} y) d (\mu \otimes \mu) (z, y). \end{array}
\]

由 V, p. 432, 定理 1, (i)，可知 \(\lambda_{\varphi}(f^{*} * f) \geqslant 0\)，这是将其应用于 \(\mu\) 在 \(f\) 的支集上诱导的测度的结果（INT, IV, p. 186, § 5, n° 7, 定义 4）。由连续性可推出，\(\lambda_{\varphi}(f^{*} * f) \geqslant 0\) 对所有 \(f \in \mathrm{L}^1(\mathrm{G})\) 成立。

空间 \(\mathrm{L}^{\infty}(\mathrm{G})\) 和 \(\mathcal{C}_{b}(\mathrm{G})\) 赋予其巴拿赫空间拓扑，并以 \(f \mapsto \|f\|_{\infty}\) 表示范数。这里将 \(\mathrm{L}^{\infty}(\mathrm{G})\)、\(\mathcal{C}_{b}(\mathrm{G})\) 或 \(\operatorname{Pos}(\mathrm{G})\) 上由弱拓扑 \(\sigma(\mathrm{L}^{\infty}(\mathrm{G}), \mathrm{L}^{1}(\mathrm{G}))\) 诱导的拓扑称为弱拓扑。

由于 \(L^{\infty}(G)\) 与 \(L^{1}(G)\) 的对偶空间相等同（INT, V, p. 61, §5, no. 8, 定理 4），\(L^{\infty}(G)\) 的每个闭球对于弱拓扑都是紧的（EVT, III, p. 17, 推论 3）。

推论。--- 在赋予弱拓扑的 \(\mathrm{L}^{\infty}(\mathrm{G})\) 中，集合 \(\operatorname{Pos}(\mathrm{G})\) 是闭的，集合 \(\operatorname{Pos}_{\leqslant 1}(\mathrm{G})\) 是紧的。

第一个断言由命题得出，第二个断言则由 EVT, III, 同上得出。

一般而言，\(\mathrm{Pos}_{1}(\mathrm{G})\) 对弱拓扑并不紧，群 R 的例子说明了这一点（V, p. 497, 习题 10）。

命题 14. --- \(\mathrm{Pos}_{\leqslant 1}(\mathrm{G})\) 的极点集合等于 \(\mathrm{Pos}_1(\mathrm{G})\) 的极点集合与零函数的并。此外，\(\mathrm{Pos}_1(\mathrm{G})\) 的极点集合的闭凸包包含 \(\mathrm{Pos}_1(\mathrm{G})\)。

由命题 13 的推论，集合 \(\mathrm{Pos}_{\leqslant 1}(\mathrm{G})\) 是尖凸锥 \(\mathrm{Pos}(\mathrm{G})\) 的一个截面，该截面由规范 \(\varphi \mapsto \varphi(e)\) 定义（EVT, II, p. 61, 定义 3 及命题 4）。因此，它的极点是零函数以及元素 \(\varphi\)，它属于 \(\mathrm{Pos}_1(\mathrm{G})\) 且属于 \(\mathrm{Pos}(\mathrm{G})\) 的极端生成元（EVT, II, p. 62, 推论 1）。这些正是 \(\mathrm{Pos}_1(\mathrm{G})\) 的极点（EVT, II, p. 61, 命题 3）。第一个断言由此得出。

证明第二个断言。设 \(\varphi\in\mathrm{Pos}_{1}(\mathrm{G})\)。由于集合 \(\mathrm{Pos}_{\leqslant 1}(\mathrm{G})\) 对弱拓扑是紧的（命题 13 的推论），存在滤子 \(\mathfrak{F}\)，在 \(\mathrm{Pos}_{\leqslant 1}(\mathrm{G})\) 的极点集合的凸包上收敛于 \(\varphi\)（EVT, II, p. 59, 定理 1）。由于巴拿赫空间 \(\mathrm{L}^{\infty}(\mathrm{G})\) 的闭球对于弱拓扑是闭的，且 \(\|\varphi\|_{\infty}=1\)，有 \(\lim_{\psi,\mathfrak{F}}\|\psi\|_{\infty}=1\)（事实上，否则将存在实数 \(c<1\) 以及滤子 \(\mathfrak{G}\)，它在半径为 \(c\) 的 \(\mathrm{L}^{\infty}(\mathrm{G})\) 闭球上，且比 \(\mathfrak{F}\) 更细，这将意味着 \(\varphi\) 属于该闭球）。

根据对 \(\mathrm{Pos}_{\leqslant1}(\mathrm{G})\) 极点的描述，极点的凸包中的每个元 \(\psi\) 都是 G 上具有如下形式的 \(\mathrm{Pos}_{\leqslant1}(\mathrm{G})\) 上的正定函数

\[
\psi = \sum_ {i \in \mathrm{I}} t _ {i} \psi_ {i},
\]

其中 I 是有限集，\(\psi_{i}\) 是 \(\mathrm{Pos}_{1}(\mathrm{G})\) 的极点，对所有 \(i \in I\)，且 \(t_{i} \in [0,1]\)。有 \(\sum_{i} t_{i} = \psi(e) = \|\psi\|_{\infty} \leqslant 1\)（V, p. 446, 引理 4）。若 \(\psi \neq 0\)，函数

\[
\frac {\psi}{\| \psi \| _ {\infty}} = \sum_ {i \in \mathrm{I}} \frac {t _ {i}}{\psi (e)} \psi_ {i}
\]

因此属于 \(\mathrm{Pos}_{1}(\mathrm{G})\) 极点的凸包。由于 \(\lim_{\psi,\mathfrak{F}}\|\psi\|_{\infty}^{-1}\psi=\varphi\)，可知 \(\varphi\) 属于 \(\mathrm{Pos}_{1}(\mathrm{G})\) 极点的闭凸包。断言得证。

\subsection{不可约表示的存在性}

沿用上一节的记号。

定理 3（赖科夫）。--- Pos1(G) 上的弱拓扑与紧收敛拓扑重合。

先证明几个引理。

引理 5. --- 设 X 是局部紧拓扑空间，\(\nu\) 是 X 上的正测度。在 \(\mathcal{C}_{b}(X)\) 的每个有界子集上，由弱拓扑 \(\sigma(L^{\infty}(X), L^{1}(X))\) 诱导的拓扑比紧收敛拓扑粗。

设 B 是 \(\mathcal{C}_{b}(\mathrm{X})\) 的有界子集，且 \(M \in R_{+}\) 满足 \(\|\varphi\|_{\infty} \leqslant M\) 对所有 \(\varphi \in B\)。设 \(\psi \in L^{1}(\mathrm{X}, \nu)\) 且 \(\varepsilon > 0\)。取 X 的一个紧子集 K，使得

\[
\int_ {\mathrm{X-K}} | \psi | d \nu <   \varepsilon .
\]

于是对 B 中任意 \(\varphi_{1}\) 和 \(\varphi_{2}\)，有

\[
\begin{array}{r l} & {\langle \varphi_ {1} - \varphi_ {2}, \psi \rangle = \int_ {\mathrm{X}} \psi (\varphi_ {1} - \varphi_ {2}) d \nu} \\ & {\qquad = \int_ {\mathrm{K}} \psi (\varphi_ {1} - \varphi_ {2}) d \nu + \int_ {\mathrm{X-K}} \psi (\varphi_ {1} - \varphi_ {2}) d \nu ,} \end{array}
\]

从而

\[
| \langle \varphi_ {1} - \varphi_ {2}, \psi \rangle | \leqslant \sup _ {x \in \mathrm{K}} | \varphi_ {1} (x) - \varphi_ {2} (x) | \| \psi \| _ {1} + 2 \mathrm{M} \varepsilon ,
\]

引理得证。

引理 6. --- 设 \(\psi \in \mathrm{L}^{1}(\mathrm{G})\)。令 B 是巴拿赫空间 \(\mathrm{L}^{\infty}(\mathrm{G})\) 的有界子集。当 B 赋予弱拓扑而映射 \(\varphi \mapsto \psi * \varphi\) 的值域 \(\mathcal{C}_b(\mathrm{G})\) 赋予紧收敛拓扑时，该映射从 B 到 \(\mathcal{C}_b(\mathrm{G})\) 连续。

设 \(\varphi\in\mathrm{L}^{\infty}(\mathrm{G})\) 且 \(\eta\in\mathrm{L}^{1}(\mathrm{G})\)。函数 \(\Delta^{-1}\check{\eta}\) 属于 \(\mathrm{L}^{1}(\mathrm{G})\)，且 \(\langle\eta,\check{\psi}\rangle=\langle\Delta^{-1}\check{\eta},\psi\rangle\)（参见 INT, VII, p. 19, § 1, n°3, 公式 (22)）。因此，映射 \(\varphi\mapsto\check{\varphi}\) 是赋予弱拓扑的空间 \(\mathrm{L}^{\infty}(\mathrm{G})\) 上的自同构。

设 \(\varphi\in\mathrm{L}^{\infty}(\mathrm{G})\)。由 INT, VIII, p. 167, § 4, no. 5, 命题 14，函数 \(\psi*\varphi\) 属于 \(\mathcal{C}_{b}(\mathrm{G})\)，且满足

\[
\begin{array}{r l} & {(\psi * \varphi) (g) = \int_ {\mathrm{G}} \psi (h) \varphi (h ^ {- 1} g) d \mu (h)} \\ & {\qquad = \int_ {\mathrm{G}} \psi (g y) \check {\varphi} (y) d \mu (y) = \langle \check {\varphi}, \pmb {\gamma} _ {\mathrm{G}} ^ {(1)} (g ^ {- 1}) \psi \rangle} \end{array}
\]

对所有 \(g \in G\) 都成立。因此，线性映射 \(u: \varphi \mapsto \psi * \varphi\) 从赋予弱拓扑的空间 \(L^{\infty}(G)\) 到赋予逐点收敛拓扑的 \(\mathcal{C}_{b}(G)\) 是连续映射。

设 \(\varphi\in\mathrm{B}\) 且 \((g,h)\in\mathrm{G}\times\mathrm{G}\)。由上式，有

\[
| u (\varphi) (g) - u (\varphi) (h) | \leqslant \| \check {\varphi} \| _ {\infty} \| (\boldsymbol {\gamma} _ {\mathrm{G}} ^ {(1)} (g ^ {- 1}) - \boldsymbol {\gamma} _ {\mathrm{G}} ^ {(1)} (h ^ {- 1})) \psi \| _ {1}.
\]

由于 B 有界且 G 在 \(L^{1}(G)\) 上的左正则表示连续（V, p. 405, no. 4），这意味着 \(u(B)\) 是 \(\mathcal{C}_{b}(G)\) 的等度连续子集。由前述结果及 TG, X, p. 16, 定理 1，断言得证。

引理 7. --- 设 \(\psi \in \mathrm{L}^1 (\mathbf{G})\) 满足 \(\psi \geqslant 0\) 且 \(\int \psi = 1\)。令 p 表示 \(\mathcal{C}_b(\mathbf{G})\) 上由 \(p(\varphi) = |\langle \varphi ,\psi \rangle |\) 对所有 \(\varphi \in \mathcal{C}_b(\mathbf{G})\) 定义的半范数。对每个 \(\varphi \in \mathrm{Pos}_1(\mathbf{G})\)，有

\[
\| \psi * \varphi - \varphi \| _ {\infty} \leqslant \sqrt {2 p (1 - \varphi)}.
\]

设 \(\varphi\in\mathrm{Pos}_{1}(\mathrm{G})\) 且 \(x\in\mathrm{G}\)。根据 INT, VIII, p. 167, § 4, no. 5, 命题 14，得到

\[
\begin{array}{r l} | \psi * \varphi (x) - \varphi (x) | = & \left| \int_ {\mathrm{G}} (\varphi (y ^ {- 1} x) - \varphi (x)) \psi (y) d \mu (y) \right| \\ & \leqslant \int_ {\mathrm{G}} | \varphi (y ^ {- 1} x) - \varphi (x) | \psi (y) d \mu (y) \\ & \leqslant \int_ {\mathrm{G}} \sqrt {2 (1 - \mathcal {R}   \varphi (y))} \psi (y) d \mu (y) \end{array}
\]

应用 V, p. 446 的界 (1)。于是由柯西–施瓦茨不等式可得

\[
\begin{array}{c} \| \psi * \varphi - \varphi \| _ {\infty} \leqslant \sqrt {2} \Bigl (\int_ {\mathrm{G}} (1 - \mathcal {R} (\varphi)) \psi   d \mu \Bigr) ^ {1 / 2} \Bigl (\int_ {\mathrm{G}} \psi   d \mu \Bigr) ^ {1 / 2} \\ \leqslant \sqrt {2} \sqrt {p (1 - \varphi)}, \end{array}
\]

从而得到结论。

引理 8. --- 设 K 是拓扑域，E 和 F 是 K 上的拓扑向量空间。设 f 是从 E 到 F 的映射，X 是 E 的子集，且 \(x \in X\)。若对于 F 中 0 的每个邻域 W，都存在 E 中 x 的一个邻域 U 以及从 X 到 F 的连续映射 g，使得映射 f \textbar{} X 在 x 处连续，当 \((f - g)(\mathrm{U} \cap \mathrm{X}) \subset \mathrm{W}\)。

设 \(W_{0}\) 是 F 中 0 的一个邻域，\(V_{0}\) 是 F 中 0 的一个对称邻域，满足 \(V_{0} + V_{0} + V_{0} \subset W_{0}\)。设 \(U_{0}\) 是 E 中 x 的一个邻域，g 是从 X 到 F 的连续映射，满足 \((f - g)(\mathrm{U}_{0} \cap \mathrm{X}) \subset \mathrm{V}_{0}\)。

存在 E 中 x 的一个邻域 \(U_{1}\)，使得 \(g(\mathrm{U}_{1} \cap \mathrm{X}) \subset g(x) + \mathrm{V}_{0}\)。对每个 \(y \in U_{0} \cap U_{1} \cap X\)，有

\[
\begin{array}{c} f (y) - f (x) = (f (y) - g (y)) + (g (y) - g (x)) + (g (x) - f (x)) \in V_{0} + V_{0} + V_{0} \subset W_{0}, \end{array}
\]

所以 \(f(y) \in f(x) + W_{0}\)，从而得到结论。

现在证明定理 3。以 \(\mathrm{Pos}_{1}(\mathrm{G})_{f}\)（分别以 \(\mathrm{Pos}_{1}(\mathrm{G})_{c}\)）表示赋予弱拓扑（分别赋予紧收敛拓扑）的集合 \(\mathrm{Pos}_{1}(\mathrm{G})\)；同样，以 \(\mathcal{C}_{b}(\mathrm{G})_{f}\)（分别以 \(\mathcal{C}_{b}(\mathrm{G})_{c}\)）表示赋予弱拓扑（分别赋予紧收敛拓扑）的空间 \(\mathcal{C}_{b}(\mathrm{G})\)。

从 \(\mathrm{Pos}_{1}(\mathrm{G})_{c}\) 到 \(\mathrm{Pos}_{1}(\mathrm{G})_{f}\) 的恒等映射连续（引理 5）。反之，以 \(\iota\) 表示 \(\mathrm{Pos}_{1}(\mathrm{G})_{f}\) 到 \(\mathcal{C}_{b}(\mathrm{G})_{c}\) 的包含映射。证明 \(\iota\) 连续即可完成证明。

设 \(\varphi_{1} \in \mathrm{Pos}_{1}(\mathbf{G})\)。为验证 \(\iota\) 在 \(\varphi_{1}\) 处连续，应用引理 8。设 W 是 \(\mathcal{C}_{b}(\mathbf{G})_{c}\) 中 0 的一个邻域。只需找到线性映射 \(u\)，从 \(\mathcal{C}_{b}(\mathbf{G})\) 到 \(\mathcal{C}_{b}(\mathbf{G})\)，它通过对子空间的传递诱导出从 \(\mathrm{Pos}_{1}(\mathbf{G})_{f}\) 到 \(\mathcal{C}_{b}(\mathbf{G})_{c}\) 的连续映射，并找到 \(\mathrm{Pos}_{1}(\mathbf{G})_{f}\) 中 0 的一个邻域 U，使得 \((\iota - u)(\mathrm{U} \cap \mathrm{Pos}_{1}(\mathbf{G})) \subset \mathrm{W}\)。

存在 \(\varepsilon > 0\)，使得 W 包含满足 \(\varphi \in \mathcal{C}_b(\mathrm{G})\) 且 \(\| \varphi \|_{\infty} \leqslant \varepsilon\) 的函数所成的集合。由于 \(\varphi_1(e) = 1\)，存在 G 中 \(e\) 的一个紧邻域 V，使得

\[
\sup _ {x \in \mathrm{V}} | 1 - \varphi_ {1} (x) | \leqslant \frac {\varepsilon^ {2}}{4}.
\]

令 \(\varphi_{\mathrm{V}}\) 为 V 的特征函数，并置 \(\psi_{\mathrm{V}} = \mu (\mathrm{V})^{-1}\varphi_{\mathrm{V}}\)。线性映射 \(u\colon \varphi \mapsto \psi_{\mathrm{V}}*\varphi\) 从 \(\mathrm{Pos}_1(\mathrm{G})_f\) 到 \(\mathcal{C}_b(\mathrm{G})_c\) 连续（引理 6）。

以 \(q_{V}\) 表示由 \(\varphi\mapsto|\langle\varphi,\psi_{V}\rangle|\) 在 \(\mathcal{C}_{b}(G)\) 上定义的半范数；它对于弱拓扑连续。由于 \(\psi_{V}\) 在 V 外为零，可得 \(q_{\mathrm{V}}(1-\varphi_{1})\leqslant\varepsilon^{2}/4\)；因此在 \(\varphi_{1}\) 的 \(\mathcal{C}_{b}(\mathrm{G})_{f}\) 中存在一个邻域 U，使得 \(q_{\mathrm{V}}(1-\varphi)\leqslant\varepsilon^{2}/2\) 对所有 \(\varphi\in U\) 成立。

设 \(\varphi\in\mathrm{U}\cap\mathrm{Pos}_{1}(\mathrm{G})\)。由引理 7，有

\[
\| (\iota - u) (\varphi) \| _ {\infty} \leqslant \sqrt {2 q _ {\mathrm{V}} (1 - \varphi)} \leqslant \varepsilon
\]

所以 \((\iota - u)(\mathrm{U}) \subset \mathrm{W}\)。定理得证。

一般而言，\(\operatorname{Pos}_{\leqslant1}(\mathbf{G})\) 上的弱拓扑不与紧收敛拓扑重合，群 R 的例子说明了这一点。

回忆我们以 \(\widehat{G}\) 表示 \(G\) 的不可约酉表示的等价类集合（V, p. 393, 定义 8）。

定理 4（盖尔范德--赖科夫）。--- 对 G 中每个 \(x \neq e\)，存在 G 的不可约酉表示 \(\pi\)，使得 \(\pi(x)\) 不是 \(\pi\) 的空间上的恒等自同态。

设 \(x \in G\) 满足 \(\pi(x)\) 对每个 \(\pi \in \widehat{G}\) 都是恒等映射。则有 \(\varphi(x) = \varphi(e) = 1\) 对每个极点 \(\varphi\)（属于 \(\mathrm{Pos}_{1}(\mathrm{G})\)）（V, p. 444, 命题 11）；又由于命题 14，这对每个 \(\varphi \in \mathrm{Pos}_{1}(\mathrm{G})\) 都成立，因为在线性泛函 \(\varphi \mapsto \varphi(e)\) 在赋予弱拓扑的 \(\mathrm{Pos}_{1}(\mathrm{G})\) 上是连续的（定理 3）。但是若 \(x \neq e\)，则存在范数为 1 的 \(f \in L^{2}(\mathrm{G})\)，使得 \(\langle f | \boldsymbol{\gamma}_{\mathrm{G}}(x)f \rangle = 0\)（V, p. 406, 引理 3）。由于这个等式左端具有 \(\varphi(x)\) 的形式，其中 \(\varphi \in \mathrm{Pos}_{1}(\mathrm{G})\)，故矛盾。

若 G 不是局部紧的，则 G 的不可约酉表示未必能分离 G 的点。

推论。--- 对 G 中任意 \(g \neq h\)，存在不可约表示 \(\pi \in \widehat{G}\) 及 \(\pi\) 的一个矩阵系数 f，使得 \(f(g) \neq f(h)\)。特别地，\(\Upsilon(\mathrm{G})\) 是 \(\mathcal{C}_{b}(\mathrm{G})\) 的能分离 G 的点的子代数。

设 \(g \neq h\) 属于 G。存在不可约表示 \(\pi \in \widehat{G}\)，使得 \(\pi(g) \neq \pi(h)\)（定理 4），所以在 \(\pi\) 的空间中存在 x 和 y，使得 \(\langle x | \pi(g)y \rangle \neq \langle x | \pi(h)y \rangle\)。

\subsection{局部紧交换群上的正定函数}

在本节中，G 是局部紧交换群，\(\mu\) 表示 G 上的哈尔测度。以 \(\widehat{G}\) 表示 G 的对偶群，以 \(\widehat{\mu}\) 表示 \(\mu\) 的对偶哈尔测度（II, p. 214, 定义 4）。以 F 表示 G 上有界测度的巴拿赫空间 \(\mathcal{M}^{1}(G)\) 上的傅里叶变换。

由于 \(\mathcal{C}_0(\mathrm{G})\) 是 \(\mathcal{K}(\mathrm{G})\) 在 \(\mathcal{C}_b(\mathrm{G})\) 中的闭包，巴拿赫空间 \(\mathcal{M}^1 (\mathrm{G})\)（\(\mathcal{K}(\mathrm{G})\) 赋予 \(\mathcal{C}_b(\mathrm{G})\) 的拓扑的对偶）与 \(\mathcal{C}_0(\mathrm{G})\) 的对偶相等同（参见 INT, III, p. 56, § 1, no. 8）。在 \(\mathcal{M}^1 (\mathrm{G})\) 上赋予与此对偶性相联系的弱拓扑。

引理 9. --- 傅里叶变换是从赋予弱拓扑的 \(\mathcal{M}^{1}(\mathrm{G})\) 到 \(\mathcal{C}_{b}(\widehat{\mathrm{G}})\) 的连续映射，后者赋予由弱拓扑 \(\sigma(\mathrm{L}^{\infty}(\widehat{\mathrm{G}}),\mathrm{L}^{1}(\widehat{\mathrm{G}}))\) 诱导的拓扑。

设 \(\mathfrak{F}\) 是 \(\mathcal{M}^1 (\mathrm{G})\) 上弱收敛于 G 上有界测度 \(\nu\) 的滤子。对每个 \(\varphi \in \mathrm{L}^1 (\widehat{\mathrm{G}})\)，有 \(\mathcal{F}(\varphi)\in \mathcal{C}_0(\mathrm{G})\)（II, p. 209, 命题 4），从而

\[
\begin{array}{r l} & {\int_ {\widehat {\mathrm{G}}} \varphi \mathcal {F} (\nu) d \widehat {\mu} = \int_ {\mathrm{G}} \mathcal {F} (\varphi) d \nu} \\ & {\qquad = \lim _ {\eta , \mathfrak {F}} \int_ {\mathrm{G}} \mathcal {F} (\varphi) d \eta = \lim _ {\eta , \mathfrak {F}} \int_ {\widehat {\mathrm{G}}} \varphi \mathcal {F} (\eta) d \widehat {\mu},} \end{array}
\]

由傅里叶变换的转置性质（II, p. 221, 命题 13）。引理得证。

定理 5（博赫纳）。--- 连续函数 \(\varphi\) 在 \(\widehat{\mathbf{G}}\) 上且属于 \(\operatorname{Pos}(\widehat{\mathbf{G}})\)，当且仅当存在 \(\nu\) 在 \(\mathbf{G}\) 上的有界正测度，使得 \(\varphi = \mathcal{F}(\nu)\)。

设 \(\nu \in \mathcal{M}^{1}(\mathrm{G})\) 是正测度。其傅里叶变换连续。对每个有限族 \((x_{i})_{i\in \mathrm{I}}\)（属于 \(\widehat{\mathbf{G}}\)）以及复数的每个有限族 \((t_i)_{i\in \mathrm{I}}\)，可得

\[
\begin{array}{r} \sum_ {i \in \mathrm{I}} \sum_ {j \in \mathrm{I}} \bar {t} _ {i} t _ {j} \mathcal {F} (\nu) (x _ {i} ^ {- 1} x _ {j}) = \sum_ {i \in \mathrm{I}} \sum_ {j \in \mathrm{I}} \bar {t} _ {i} t _ {j} \int_ {\mathrm{G}} x _ {i} \overline {{x}} _ {j} d \nu \\ = \int_ {\mathrm{G}} \left| \sum_ {i \in \mathrm{I}} \bar {t} _ {i} x _ {i} \right| ^ {2} d \nu \geqslant 0, \end{array}
\]

由于 \(\nu\) 是正测度。因此，\(\nu\) 的傅里叶变换是 \(\hat{G}\) 上的正定函数（V, p. 432, 定理 1, (ii)）。

证明逆向蕴含。给集合 \(\operatorname{Pos}_{\leqslant1}(\widehat{\mathbf{G}})\) 赋予弱拓扑，该拓扑像前面一样由弱拓扑 \(\sigma(\mathrm{L}^{\infty}(\widehat{\mathrm{G}}),\mathrm{L}^{1}(\widehat{\mathrm{G}}))\) 诱导。集合 \(\operatorname{Pos}_{\leqslant1}(\widehat{\mathbf{G}})\) 是紧凸集（V, p. 448, 命题 13 的推论）。由 V, p. 450, 命题 14，V, p. 444, 命题 11 以及 V, p. 390, 推论 7，\(\operatorname{Pos}_{\leqslant1}(\widehat{\mathbf{G}})\) 的极点是零函数和 \(\widehat{G}\) 中的元素。

令 N 为 G 上质量 \(\leqslant 1\) 的正有界测度的集合；它在赋予弱拓扑的 \(\mathcal{M}^{1}(G)\) 中是紧的（EVT, III, p. 17, 推论 3）。由引理 9 及证明的第一部分，G 上的傅里叶变换通过对子空间的传递定义了从 \(\mathcal{N}\) 到 \(\mathrm{Pos}_{\leqslant 1}(\widehat{\mathrm{G}})\) 的连续映射；由齐次性，只需证明此映射是满射。\(\mathcal{F}(\mathcal{N})\) 是傅里叶变换下 \(\mathcal{N}\) 的像，是紧凸集；它包含零函数和 \(\widehat{\mathrm{G}}\) 中的元素（事实上，由 II, p. 220, 定理 2，这些元素都具有 \(\mathrm{ev}_x\colon \chi \mapsto \chi (x)\) 的形式，其中 \(x\) 是 G 中某个元，且 \(\mathrm{ev}_x = \mathcal{F}(\varepsilon_{x^{-1}})\)）。因此集合 \(\mathcal{F}(\mathcal{N})\) 包含 \(\mathrm{Pos}_{\leqslant 1}(\mathrm{G})\) 的极点，由克赖因–米尔曼定理（EVT, II, p. 59, 定理 1）可得 \(\mathcal{F}(\mathcal{N}) = \mathrm{Pos}_{\leqslant 1}(\mathrm{G})\)。证明完毕。

当 \(G = R^{k}\)，其中整数 \(k \in N\) 时，本定理对应于 INT, IX, p. 94, § 6, no. 12, 命题 11。

\section{紧群的表示}

在本段中，除非另有明确说明，所考虑的所有拓扑向量空间都是复的。

固定一个紧拓扑群 G，其单位元记为 e。群 G 是单模的（INT, VII, p. 20, § 1, n° 3, 命题 3 的推论）。以 μ 表示 G 上的归一化哈尔测度（即满足 μ(G) = 1 的测度），并且对 \(1 \leqslant p \leqslant +\infty\)，以 \(\mathcal{L}^{p}(G)\)（分别以 \(\mathrm{L}^{p}(\mathrm{G})\)）表示空间 \(\mathcal{L}_{\mathbf{C}}^{p}(\mathrm{G}, \mu)\)（分别表示空间 \(\mathrm{L}_{\mathbf{C}}^{p}(\mathrm{G}, \mu)\)）。卷积总是相对于测度 μ 取。设 \(p \in [1, +\infty]\)。将 \(\mathcal{C}(\mathrm{G})\) 与 \(\mathrm{L}^{p}(\mathrm{G})\) 的一个子空间等同，这是允许的，因为 μ 的支集等于 G。

对每个不可约酉表示 \(\pi \in \widehat{\mathbf{G}}\)，令 \(\mathrm{E}_{\pi}\) 为 \(\pi\) 的空间。

\subsection{有限维表示的半单性}

回忆（INT, VII, p. 71, § 3, n° 1, 引理 1）：对于 G 在希尔伯特空间 E 上的每个连续表示 \(\varrho\)，E 上存在非退化正埃尔米特形式 \(q\)，使得由 q 定义的 E 的拓扑向量空间结构与 E 的原结构相同，并且 \(\varrho\) 是 G 在赋予内积 q 的希尔伯特空间 E 上的酉表示。

命题 1. --- 设 \(\varrho\) 是 G 在分离拓扑向量空间 E 上的有限维线性表示。E 上存在内积 \(q\)，使得 \(\varrho\) 是赋予 \(q\) 的希尔伯特空间 E 上的酉表示。特别地，\(\varrho\) 是半单的。

由于 E 是有限维的，E 上存在希尔伯特空间结构。应用前述注记以及 V, p. 392, 推论 2 即得结果。

注。--- 稍后我们将看到（V, p. 466, 推论 4），G 的每个酉表示都是半单的。

推论。--- 设 \(\varrho_{1}\) 和 \(\varrho_{2}\) 是 G 的有限维表示。表示 \(\varrho_{1}\) 和 \(\varrho_{2}\) 同构，当且仅当它们的特征标相等。

这由命题及 V, p. 392, 推论 3 得出。

\subsection{不可约表示}

引理 1. --- G 的每个不可约酉表示都是平方可积的。

事实上，不可约酉表示的矩阵系数连续且有界，所以在 G 上平方可积，因为 G 是紧的。

命题 2. --- 设 π 是 G 在希尔伯特空间 E 上的不可约酉表示。E 的维数有限，且等于 π 的形式次数。

表示 \(\pi\) 是平方可积的（引理 1），所以形式次数 \(c_{\mu}(\pi)\) 相对于 \(\mu\) 有定义（V, p. 423, 定义 4）；它是严格为正的实数。设 \((e_i)_{i \in I}\) 是 E 中的有限标准正交族。设 \(x\) 是 E 中范数为 1 的元（这样的元存在，因为 E 非零）。对 \(i \in I\)，有

\[
c _ {\mu} (\pi) \int_ {\mathrm{G}} | \langle e _ {i} | \pi (g) x \rangle | ^ {2} d \mu (g) = 1
\]

（V, p. 424, 命题 8）。对 \(i \in I\) 将此公式求和，得到

\[
\begin{array}{r l} & {\mathrm{Card(I)} = c _ {\mu} (\pi) \int_ {\mathrm{G}} \sum_ {i \in \mathrm{I}} | \langle e _ {i} | \pi (g) x \rangle | ^ {2} d \mu (g)} \\ & {\qquad \leqslant c _ {\mu} (\pi) \int_ {\mathrm{G}} \| \pi (g) x \| ^ {2} d \mu (g) = c _ {\mu} (\pi)} \end{array}
\]

由贝塞尔不等式（EVT, V, p. 21, 命题 4）。因此 I 的基数上界为 \(c_{\mu}(\pi)\)。这意味着 E 的维数有限。

于是可将前述结果应用于 E 的标准正交基 \((e_{i})_{i\in\mathrm{I}}\)。由 EVT, V, p. 22, 命题 5，得到

\[
\begin{array}{r} \dim (\mathrm{E}) = c _ {\mu} (\pi) \int_ {\mathrm{G}} \sum_ {i \in \mathrm{I}} | \langle e _ {i} | \pi (g) x \rangle | ^ {2} d \mu (g) \\ = c _ {\mu} (\pi) \int_ {\mathrm{G}} \| \pi (g) x \| ^ {2} d \mu (g) = c _ {\mu} (\pi), \end{array}
\]

证明完毕。

特别地，因此可以谈论特征标 \(\chi_{\pi}\)，其中 \(\pi\) 是 G 的一个不可约酉表示。它是 G 上的连续函数。

推论 1. --- 设 \(\pi\) 是 G 在希尔伯特空间 E 上的不可约酉表示。\(\pi\) 的特征标是对合代数 \(\mathrm{L}^{1}(\mathrm{G})\) 的埃尔米特元。

设 \(x \in G\)。由于 \(G\) 是单模的，有 \(\chi_{\pi}^{*}(x) = \overline{\operatorname{Tr}(\pi(x^{-1}))}\)，又有 \(\chi_{\pi}^{*}(x) = \operatorname{Tr}(\pi(x))\)，这是由于 \(\pi\) 是酉表示。

推论 2. --- a) 设 \(\pi\) 是 G 在希尔伯特空间 E 上的不可约酉表示。有

\[
\int_ {\mathrm{G}} \overline {{\langle x \mid \pi (g) x ^ {\prime} \rangle}} \langle y \mid \pi (g) y ^ {\prime} \rangle d \mu (g) = \frac {1}{\dim (\mathrm{E})} \overline {{\langle x \mid y \rangle}} \langle x ^ {\prime} \mid y ^ {\prime} \rangle
\]

对所有 \((x,y,x',y')\in \mathrm{E}^4\) 都成立。

\begin{enumerate}
\def\labelenumi{\alph{enumi})}
\setcounter{enumi}{1}
\tightlist
\item
  设 \(\pi_{1}\)（分别为 \(\pi_{2}\)）是 G 在希尔伯特空间 \(E_{1}\)（分别为 \(E_{2}\)）上的不可约酉表示。若 \(\pi_{1}\) 和 \(\pi_{2}\) 不同构，则有
\end{enumerate}

\[
\int_ {G} \overline {{\langle x \mid \pi_ {1} (g) x ^ {\prime} \rangle}} \langle y \mid \pi_ {2} (g) y ^ {\prime} \rangle d \mu (g) = 0
\]

对所有 \((x,x^{\prime},y,y^{\prime})\in \mathrm{E}_{1}^{2}\times \mathrm{E}_{2}^{2}\) 都成立

这立即由 V, p. 424, 命题 8 和 V, p. 424, 命题 9 得出，同时考虑引理 1 及公式 \(c_{\mu}(\pi) = \dim(\mathrm{E})\)（命题 2）。

推论 3. --- 设 \(\pi_1\) 和 \(\pi_2\) 是 G 的不可约表示。在希尔伯特空间 \(\mathrm{L}^{2}(\mathrm{G})\) 中，有 \(\langle \chi_{\pi_1} | \chi_{\pi_2} \rangle = 1\)（当 \(\pi_1\) 和 \(\pi_2\) 同构时），否则有 \(\langle \chi_{\pi_1} | \chi_{\pi_2} \rangle = 0\)。换言之，类 \(\pi \in \widehat{\mathbf{G}}\) 的特征标所成的族是 \(\mathrm{L}^{2}(\mathrm{G})\) 中的标准正交族。

这由 V, p. 457, 命题 1 及推论 2 得出；注意，对于 G 的有限维酉表示的空间的每个标准正交基 \((e_{i})_{i\in\mathrm{I}}\)，以及表示 \(\pi\) 和每个 \(g\in G\)，都有公式

\[
\chi_ {\pi} (g) = \sum_ {i \in I} \langle e _ {i} | \pi (g) e _ {i} \rangle .
\]

推论 4. --- a) 设 \(\varrho\) 是 G 的连续有限维表示。有

\[
\chi_ {\varrho} = \sum_ {\pi \in \widehat {G}} \dim (\operatorname{Hom} _ {G} (\pi , \varrho)) \chi_ {\pi} = \sum_ {\pi \in \widehat {G}} \dim (\operatorname{Hom} _ {G} (\varrho , \pi)) \chi_ {\pi}.
\]

\begin{enumerate}
\def\labelenumi{\alph{enumi})}
\setcounter{enumi}{1}
\tightlist
\item
  设 \(\varrho_{1}\) 和 \(\varrho_{2}\) 是 G 的连续有限维表示。有
\end{enumerate}

\[
\left\langle \chi_ {\varrho_ {1}} \mid \chi_ {\varrho_ {2}} \right\rangle = \dim (\operatorname{Hom} _ {\mathrm{G}} (\varrho_ {1}, \varrho_ {2})) = \dim (\operatorname{Hom} _ {\mathrm{G}} (\varrho_ {2}, \varrho_ {1})).
\]

\begin{enumerate}
\def\labelenumi{\alph{enumi})}
\setcounter{enumi}{2}
\tightlist
\item
  G 的连续有限维表示 \(\varrho\) 不可约，当且仅当 \(\| \chi_{\varrho}\|^{2} = 1\)。
\end{enumerate}

由于 \(\varrho\) 是半单的，它同构于其 \(\pi\)-同型分量 \(\mathrm{M}_{\pi}(\varrho)\) 的直和，其中 \(\pi \in \widehat{\mathrm{G}}\)。以 \(m_{\pi}(\varrho)\) 表示 \(\pi\) 在 \(\varrho\) 中的重数（V, p. 398, 定义 10）；于是有

\[
\chi_ {\varrho} = \sum_ {\pi \in \widehat {G}} m _ {\pi} (\varrho) \chi_ {\pi}.
\]

因此，断言 a) 由以下事实得出：

\[
m _ {\pi} (\varrho) = \dim \operatorname{Hom} _ {\mathrm{G}} (\pi , \varrho) = \dim \operatorname{Hom} _ {\mathrm{G}} (\varrho , \pi)
\]

（V, p. 377, 公式 (1) 及 V, p. 387, 推论 2）。

由特征标的双线性和标准正交性，断言 \(a\) ) 蕴含

\[
\langle \chi_ {\varrho_ {1}} \mid \chi_ {\varrho_ {2}} \rangle = \sum_ {\pi \in \widehat {\mathrm{G}}} \dim (\operatorname{Hom} _ {\mathrm{G}} (\pi , \varrho_ {1})) \dim (\operatorname{Hom} _ {\mathrm{G}} (\pi , \varrho_ {2})),
\]

另一方面，有典范同构

\[
\begin{array}{c} \operatorname{Hom} _ {\mathrm{G}} (\varrho _ {1}, \varrho _ {2}) \to \bigoplus_ {(\pi _ {1}, \pi _ {2}) \in \widehat {\mathrm{G}} \times \widehat {\mathrm{G}}} \operatorname{Hom} _ {\mathrm{G}} (\mathrm{M} _ {\pi _ {1}} (\varrho _ {1}), \mathrm{M} _ {\pi _ {2}} (\varrho _ {2})) \\ \qquad \qquad \qquad \qquad \qquad \qquad \qquad \qquad \qquad \qquad \qquad \qquad \qquad \qquad \qquad \qquad \qquad \qquad \qquad \qquad \qquad \qquad \qquad \qquad \qquad \qquad \qquad \qquad \qquad \qquad \qquad \qquad \qquad \qquad \qquad \to \bigoplus_ {\pi \in \widehat {\mathrm{G}}} \operatorname{Hom} _ {\mathrm{G}} (\mathrm{M} _ {\pi} (\varrho _ {1}), \mathrm{M} _ {\pi} (\varrho _ {2})) \end{array}
\]

（V, p. 377, 公式 (1)），从而可得

\[
\dim (\operatorname{Hom} _ {\mathrm{G}} (\varrho_ {1}, \varrho_ {2})) = \sum_ {\pi \in \widehat {\mathrm{G}}} m _ {\pi} (\varrho_ {1}) m _ {\pi} (\varrho_ {2}),
\]

从而得到断言 \(b\)。断言 \(c\) 也由 \(a\) 和舒尔引理（V, p. 386, 命题 6）得出。

推论 5. --- 设 \(\pi_1\) 和 \(\pi_2\) 是 G 的不可约酉表示。则 \(\pi_1(\overline{\chi}_{\pi_2}) = 0\)；若 \(\pi_1\) 不同构于 \(\pi_2\)，而 \(\pi_1(\overline{\chi}_{\pi_1})\) 是乘以 \(1/\dim(\pi_1)\) 的映射。

\(\pi_2\) 的特征标是 G 上的连续中心函数，因此线性映射 \(u = \pi_1(\overline{\chi}_{\pi_2})\) 有定义且属于空间 \(\mathrm{Hom}_{\mathrm{G}}(\pi_1,\pi_1)\)。由舒尔引理（V, p. 386, 命题 6），它是一个倍乘映射，其迹为

\[
\begin{array}{c} \operatorname{Tr} (u) = \operatorname{Tr} \Bigl (\int_ {\mathrm{G}} \overline {{\chi_ {\pi_ {2}} (g)}}   \pi_ {1} (g) d \mu (g) \Bigr) \\ = \int_ {\mathrm{G}} \overline {{\chi_ {\pi_ {2}} (g)}} \chi_ {\pi_ {1}} (g) d \mu (g). \end{array}
\]

由正交关系，\(u\) 的迹为零；若 \(\pi_1\) 不同构于 \(\pi_2\)，否则等于 1。断言得证。

注。--- 当 G 有限时，舒尔正交关系（分别为特征标正交公式）与 A, VIII, p. 399 中的公式（分别与 A, VIII, p. 400, 命题 4 中的公式）一致。另一方面，有限群特征标的“第二正交关系”（A, VIII, p. 402, 公式 (32)）在 G 紧时没有完全对应的形式。

当 G 有限时，与 e 的共轭类对应的第二正交公式的特殊情形是公式

\[
\frac {1}{\mathrm{CardG}} \sum_ {\pi \in \widehat {\mathrm{G}}} \dim (\pi) \chi _ {\pi} (g) = \left\{ \begin{array}{l l} 1 & \text {if} g = e \\ 0 & \text {otherwise}, \end{array} \right.
\]

这也可解释为计算 G 的正则表示 \(\gamma_{G}\) 的特征标，并且等价于对从 G 到 C 的每个函数 f 都有公式 \(\mathrm{Tr}(\boldsymbol{\gamma}_{\mathrm{G}}(f)) = f(e)\)。

再设 G 是任意紧群。对每个函数 \(f \in \mathcal{C}(\mathrm{G})\)，\(\boldsymbol{\gamma}_{\mathrm{G}}(f)\) 是 \(\mathrm{L}^{2}(\mathrm{G})\) 的自同态，与自同态 \(\varphi \mapsto f * \varphi\) 重合，且具有有限迹（V, p. 407, 引理 4 及 III, p. 33, 推论 2）。由同上及 IV, p. 177, 定理 2，还有

\[
\operatorname{Tr} \left(\boldsymbol {\gamma} _ {\mathrm{G}} (f)\right) = \int_ {\mathrm{G}} f \left(x ^ {- 1} x\right) d \mu (x) = f (e).
\]

命题 3. --- 设 \(G = G_{1} \times G_{2}\)，其中 \(G_{1}\) 和 \(G_{2}\) 是紧群。从 \(\widehat{G}_{1} \times \widehat{G}_{2}\) 到 \(\widehat{G}\) 的映射 b 将 \((\pi_{1}, \pi_{2})\) 送到 \(\pi_{1} \boxtimes \pi_{2}\) 的类，且 b 是双射。

由 V, p. 389, 推论 6，映射 b 定义良好。

设 \(\pi_1\) 和 \(\pi_2\) 是 \(\widehat{\mathbf{G}}\) 的元素。设 \(\pi \in \widehat{\mathbf{G}}\)。\(\pi\)-同型分量是限制 \(\varpi\) 的 \(\pi_1 \boxtimes \pi_2\) 在 \(\mathrm{G}_1 \times \{e\}\) 上，其等于 \(\varpi\) 当 \(\pi_1 = \pi\) 时，否则为零（V, p. 384, 引理 8）。因此，酉表示 \(\pi_1 \boxtimes \pi_2\) 在同构意义下确定 \(\pi_1\)；同样，它确定 \(\pi_2\)，这证明 \(b\) 是单射。

证明 \(b\) 是满射。设 \(\pi\) 是 \(G_1 \times G_2\) 在希尔伯特空间 E 上的不可约酉表示。它是有限维的（命题 2）。设 \(\pi_1\) 是 \(G_1\) 在希尔伯特空间 \(E_1\) 上的不可约酉表示，使得 \(\pi\) 到 \(G_1 \times \{e\}\) 的限制包含一个与 \(\pi_1\) 同构的子表示（V, p. 379, 引理 3）。令 \(F = \text{Hom}_{G_1}(\pi_1, \pi)\)。这是非零有限维向量空间。对 \(h \in G_2\) 和 \(u \in F\)，令 \(\varrho(h)(u)\) 为从 \(E_1\) 到 E 的线性映射，由 \(x \mapsto \pi(e, h)(u(x))\) 定义。由于 \(G_1 \times \{e\}\) 与 \(\{e\} \times G_2\) 在 G 中交换，映射 \(\varrho(h)(u)\) 属于空间 F。映射 \(\varrho\) 是 \(G_2\) 在空间 F 上的线性表示；它是连续的。由于 F 非零，存在不可约子表示 \(\pi_2\) 的 \(\varrho\)（V, p. 379, 引理 3）。令 \(E_2\) 为 \(\pi_2\) 的空间。于是，线性映射 \(v: E_1 \otimes E_2 \to E\) 满足 \(x \otimes u \mapsto u(x)\) 对 \(x \in E_1\) 和 \(u \in E_2\) 成立，是从 \(\pi_1 \boxtimes \pi_2\) 到 \(\pi\) 的非零 G-态射；由于表示 \(\pi\) 和 \(\pi_1 \boxtimes \pi_2\) 都不可约且有限维，态射 \(v\) 是同构（V, p. 378, 引理 2）。

本命题可与 A, VIII, p. 208 的定理 1 比较。

\subsection{彼得--外尔定理}

回忆我们以 \(\Theta(\mathrm{G})\) 表示 G 的有限维酉表示的矩阵系数空间。空间 \(\Theta(\mathrm{G})\) 是 \(\mathcal{C}(\mathrm{G})\) 的一个在复共轭下稳定的含单位元子代数（V, p. 386, 命题 5）。

对 \(\pi \in \widehat{\mathrm{G}}\)，令 \(\varrho_{\mathrm{G}}(\pi)\) 为 \(\mathcal{C}(\mathrm{G})\) 中由 \(\pi\) 的矩阵系数生成的子空间。将其与 \(\mathrm{L}^2 (\mathrm{G})\) 的一个子空间等同。

空间 \(\Theta(G)\) 等于各空间 \(\varrho_{G}(\pi)\) 的和，其中 \(\pi\in\widehat{G}\)。事实上，\(\Theta(G)\) 的每个元都是 G 的不可约表示的矩阵系数之和，因为由 V, p. 457, 命题 1，G 的有限维表示是半单的。此外，该和是直和，因为由 V, p. 458, 推论 2，各空间 \(\varrho_{G}(\pi)\) 两两正交。

命题 4. — 空间 \(\Theta(\mathrm{G})\) 在 \(\mathcal{C}(\mathrm{G})\) 和 \(\mathrm{L}^2(\mathrm{G})\) 中都稠密。

含单位元子代数 \(\Theta(\mathrm{G})\) 的 \(\mathcal{C}(\mathrm{G})\) 在共轭下稳定。由 V, p. 457, 命题 2，它与子代数 \(\Upsilon(\mathrm{G})\) 重合，故能分离 G 的点（V, p. 454, 定理 4 的推论）。因此，由 TG, X, p. 40, 推论 2，它在 \(\mathcal{C}(\mathrm{G})\) 中稠密，从而在 \(\mathrm{L}^2(\mathrm{G})\) 中也稠密（参见 INT, IV, p. 155, §4, no. 7, 命题 13）。

回忆 G 的双正则酉表示是表示 \(\varrho_{G}\)，它从 G × G 到 \(\mathrm{L}^{2}(\mathrm{G})\)，满足 \((g_{1}, g_{2}) \mapsto \gamma_{\mathrm{G}}(g_{1}) \delta_{\mathrm{G}}(g_{2})\)。

命题 5. --- 设 \(\pi \in \widehat{\mathbf{G}}\)。空间 \(\varrho_{\mathrm{G}}(\pi)\) 是 \(\varrho_{\mathrm{G}}\) 的一个子表示，与 \(\overline{\pi} \boxtimes \pi\) 同构。特别地，它是 \(\mathrm{G} \times \mathrm{G}\) 的不可约表示。

由 V, p. 422, 命题 7（应用 \(Z = \{e\}\)），空间 \(\varrho_{G}(\pi)\) 是 \(\varrho_{G}\) 的子表示，且从 \(\overline{\mathrm{E}}_{\pi} \otimes \mathrm{E}_{\pi}\) 到 \(L^2(G)\)、将 \(x \otimes y\) 送到矩阵系数 \(g \mapsto \langle x | \pi(g)y \rangle\) 的线性映射，是一个 \((G \times G)\)-同构，从 \(\overline{\pi} \boxtimes \pi\) 到 \(\varrho_{G}(\pi)\)。因此子表示 \(\varrho_{G}(\pi)\) 是不可约的（V, p. 461, 命题 3）。

定理 1（彼得--外尔）。--- 设 G 是紧拓扑群。G 的双正则表示 \(\varrho_{\mathrm{G}}\) 是各子表示 \(\varrho_{\mathrm{G}}(\pi)\)（其中 \(\pi \in \widehat{\mathrm{G}}\)）的希尔伯特直和。

各空间 \(\varrho_{\mathrm{G}}(\pi)\) 两两正交；它们是 G 的双正则表示的不可约子表示（命题 5），并且两两不同构（V, p. 461, 命题 3）。

定理即由以下事实得出：空间 \(\Theta(\mathrm{G})\) 是各空间 \(\varrho_{\mathrm{G}}(\pi)\) 的和，其中 \(\pi\) 遍历 \(\widehat{\mathrm{G}}\)，且该和在 \(\mathrm{L}^2(\mathrm{G})\) 中稠密（命题 4）。

推论 1. --- 对每个 \(\pi\in\widehat{G}\)，子表示 \(\varrho_{\mathrm{G}}(\pi)\) 是 \((\overline{\pi}\boxtimes\pi)\)-同型分量的 \(\varrho_{\mathrm{G}}\)。

令 F 为 \((\overline{\pi} \boxtimes \pi)\)-同型分量的 \(\varrho_{G}\)。空间 F 包含 \(\varrho_{\mathrm{G}}(\pi)\)（命题 5）。此外，对每个 \(\tau \in \widehat{G}\)（不同于 \(\pi\)），F 与 \(\varrho_{\mathrm{G}}(\tau)\) 的交为零（V, p. 461, 命题 3）。应用定理可推出 \(\mathrm{F} = \varrho_{\mathrm{G}}(\pi)\)。

推论 2. --- 设 \(\pi_1\) 和 \(\pi_2\) 是不同构的 G 的不可约表示。\((\overline{\pi}_1 \boxtimes \pi_2)\)-同型分量的 \(\varrho_G\) 为零。

推论 3. --- G 的右（分别为左）正则表示同构于希尔伯特直和

\[
\bigoplus_ {\pi \in \widehat {\mathrm{G}}} \pi^ {\dim (\pi)}.
\]

由定理及命题 5，G 的双正则表示限制到子群 \(\mathrm{H} = \{e\} \times \mathrm{G}\) 后，同构于各表示 \(\overline{\pi} \boxtimes \pi\) 在 H 上的限制的希尔伯特直和，其中 \(\pi \in \widehat{\mathrm{G}}\)。这些表示同构于 \(\dim(\pi)\) 个 \(\pi\) 的直和（V, p. 384, 引理 8）。这给出右正则表示的结果，左正则表示的情形类似。

推论 4. --- 设 \(\pi\in\widehat{G}\)。\(\pi\)-同型分量的 \(\delta_{G}\)（分别为 \(\overline{\pi}\)-同型分量的 \(\gamma_{G}\)）等于 \(\varrho_{G}(\pi)\)。

论证与前一推论类似，只需将 \(\varrho_{\mathrm{G}}(\pi)\) 看作 \(\delta_{\mathrm{G}}\)（分别为 \(\gamma_{\mathrm{G}}\)）的子表示。

注。--- 1) 若 G 有限，这些陈述对应于 A, VIII, p. 398 的注记中的结果。

\begin{enumerate}
\def\labelenumi{\arabic{enumi})}
\setcounter{enumi}{1}
\item
  假设 G 是交换的。集合 \(\widehat{G}\) 与 G 的对偶群相等同（V, p. 393, 注记）。对每个 \(\chi \in \widehat{G}\)，空间 \(\varrho_{\mathrm{G}}(\chi)\) 是 \(\mathrm{L}^2 (\mathrm{G})\) 中由 \(\chi\) 生成的一维子空间。因此，G 的定理 1 等价于 II, p. 215 的定理 1 的推论。
\item
  若存在整数 \(n \geqslant 0\)，使得 G 是 \(\mathbf{GL}(\mathbf{C}^{n})\) 的紧子群，则 G 在 \(\mathbf{GL}(\mathbf{C}^{n})\) 中的恒等表示足以分离 G 的点，随后可以直接证明命题 4，再证明彼得–外尔定理，而无需援引盖尔范德–赖科夫定理。
\item
  彼得–外尔定理特别蕴含，从 G 到 \(\prod_{\pi\in\widehat{G}}\mathbf{U}(\mathrm{E}_{\pi})\) 的自然连续同态是单射。
\end{enumerate}

\subsection{矩阵系数与 G-有限函数}

命题 6. --- \(\mathrm{L}^{2}(\mathrm{G})\) 的下列子空间相等：

\begin{enumerate}
\def\labelenumi{\alph{enumi})}
\item
  空间 \(\Theta (\mathbf{G})\)
\item
  \(\varrho_{\mathrm{G}}(\pi)\) 的 \(\mathrm{L}^2 (\mathrm{G})\) 子空间的代数直和。
\item
  \(\gamma_{\mathrm{G}}\) 的 G-有限向量空间（参见 V, p. 376）；
\item
  \(\delta_{G}\) 的 G-有限向量空间；
\item
  \((\mathbf{G} \times \mathbf{G})\)-有限向量空间的 \(\varrho_{\mathrm{G}}\)。
\end{enumerate}

特别地，\(\gamma_{\mathrm{G}}\)、\(\delta_{G}\) 或 \(\varrho_{\mathrm{G}}\) 的每个 G-有限向量都属于 \(\mathcal{C}(\mathrm{G})\)。

以 \(F_{a}\)（分别以 \(F_{b}\)、\(F_{c}\)、\(F_{d}\)、\(F_{e}\)）表示由条件 a)（分别由 b)、c)、d)、e)）定义的空间。我们已经注意到 \(F_{a} = F_{b}\)。

有 \(F_{b} \subset F_{c}\)，因为 \(\varrho_{\mathrm{G}}(\pi)\) 是 \(\gamma_{G}\) 的有限维子表示，对所有 \(\pi \in \widehat{G}\) 成立。反之，设 f 是 \(\gamma_{G}\) 的 G-有限向量。子空间 \(E_{f}\) 由函数 \(\gamma_{\mathrm{G}}(g)f\)（其中 \(g \in G\)）生成，是 \(\gamma_{G}\) 的有限维子表示。它等于其 \(\pi\)-同型分量的直和，对所有 \(\pi \in \widehat{G}\)（V, p. 457, 命题 1）。根据 V, p. 463, 推论 4，这意味着 \(E_{f}\) 包含于各空间 \(\varrho_{\mathrm{G}}(\pi)\) 的和中，从而 \(F_{c} \subset F_{b}\)。

用类似论证可得 \(F_{b} = F_{d}\)；由于 \(\varrho_{\mathrm{G}}(\pi)\) 是 \(\varrho_{G}\) 的子表示，同样可证明 \(F_{b} = F_{e}\)。

推论。--- 设 \(E \subset L^{2}(G)\) 是定义 \(\gamma_{G}\)（分别为 \(\delta_{G}\)、\(\varrho_{G}\)）的一个子表示的有限维向量子空间。则 E 包含于 \(\mathcal{C}(G)\)。

事实上，E 的每个元都是表示 \(\gamma_{G}\) 的 G-有限向量。

\subsection{分离拟完备空间中的表示}

命题 7. --- 设 \(\varrho\) 是 G 在分离、局部凸且拟完备空间 E 上的连续表示。E 的有限维子表示之和在 E 中稠密。

设 \(x \in \mathrm{E}\)，U 是 \(x\) 的一个开邻域。测度 \(\nu \in \mathcal{M}(\mathrm{G})\) 满足 \(\varrho(\nu)x \in \mathrm{U}\) 的所成集合，对于紧收敛拓扑在 \(\mathcal{M}(\mathrm{G})\) 中是开的（参见 V, p. 400, no. 2）。它包含 \(\varepsilon_e\)；因此它包含形如 \(\nu = f_1 \cdot \mu\) 的测度，其中 \(f_1 \in \mathcal{C}(\mathrm{G})\)（INT, VIII, p. 171, § 4, n° 7, 命题 19）；于是它还包含形如 \(f_2 \cdot \mu\) 的测度，其中 \(f_2\) 是 G-有限函数（V, p. 464, 命题 6 及 V, p. 462, 命题 4）。

\(\gamma_{G}\) 中由 \(f_{2}\) 生成的子表示 F 是有限维的。令 \(\widetilde{F}\) 为线性映射 \(f \mapsto \varrho(f)x\) 从 F 到 E 的像。空间 \(\widetilde{F}\) 是有限维的，且包含属于 U 的元 \(\varrho(f_{2})x\)。由于 \(\varrho(g)\varrho(f) = \varrho(\gamma_{\mathrm{G}}(g)f)\) 对所有 \(g \in G\) 和每个 \(f \in F\) 成立（V, p. 406, 公式 (1)），空间 \(\widetilde{F}\) 是 \(\varrho\) 的一个与 U 相交的子表示。命题得证。

推论 1. --- 设 π 是 G 在局部凸分离拟完备空间 E 上的连续不可约表示。空间 E 是有限维的。

推论 2. --- 设 \(\varrho\) 是 G 在局部凸分离拟完备空间 E 上的连续表示，\(\pi\) 是 G 的连续不可约表示。

\begin{enumerate}
\def\labelenumi{\alph{enumi})}
\item
  从 E 到 E 的 G-态射 \(p_{\pi} = \dim(\pi)\varrho(\overline{\chi}_{\pi})\) 是 E 的连续投影，其像是 \(\pi\)-同型分量的 \(\varrho\)；
\item
  若 \(\varrho\) 是酉表示，则投影 \(p_{\pi}\) 是 E 的正交投影，其像为 \(\mathrm{M}_{\pi}(\varrho)\)。
\end{enumerate}

证明 \(a\)。线性映射 \(p_{\pi} = \dim(\pi)\varrho(\overline{\chi}_{\pi})\) 定义良好，因为 E 拟完备且 G 紧（V, p. 401）。它属于 \(\mathrm{Hom}_{\mathrm{G}}(\varrho,\varrho)\)，因为 \(\pi\) 的特征标是 G 上的连续中心函数。

设 \(\varpi\) 是 E 的有限维子表示，F 是其空间。映射 \(p_{\pi}\) 通过对子空间的传递在 F 上诱导自同态 \(\dim(\pi)\varpi(\overline{\chi}_{\pi})\)。因此，当 \(\varpi\) 不同构于 \(\pi\) 时此自同态为零，否则是 F 的恒等映射（事实上，当 \(\varpi\) 不可约时由 V, p. 460, 推论 5 得到一般情形，而一般情形由此推出，因为由 V, p. 457, 命题 1，\(\varpi\) 是半单的）。

最后，由于 E 的有限维子表示之和在 E 中稠密（命题 7），且 \(p_{\pi}\) 连续，我们得到 \(p_{\pi}\) 是一个投影，其像是 \(\pi\)-同型分量的 \(\varrho\)。若 E 是希尔伯特空间，则由 V, p. 458, 推论 1，投影 \(p_{\pi}\) 是埃尔米特的，因而是正交投影（I, p. 133, 引理 3, (ii)）。

推论 3. --- 设 \(\varrho\) 是 G 在局部凸分离拟完备空间 E 上的连续表示。自同态

\[
x \mapsto \int_ {\mathrm{G}} \varrho (g) x d \mu (g)
\]

是 E 的一个投影，其像是 E 中不变向量的空间 \(E^{G}\)。特别地，若 E 有限维，则

\[
\mathrm{dim} (\mathrm{E} ^ {\mathrm{G}}) = \int_ {\mathrm{G}} \chi_ {\varrho} (g) d \mu (g).
\]

第一个断言是前一推论在 \(\pi\) 为 G 的一维平凡表示时的特殊情形。第二个断言由此得出，因为 \(E^{G}\) 的维数是到 \(E^{G}\) 上的正交投影的迹，即

\[
\mathrm{dim} (\mathrm{E} ^ {\mathrm{G}}) = \mathrm{Tr} \Bigl (\int_ {\mathrm{G}} \varrho (g) d \mu (g) \Bigr) = \int_ {\mathrm{G}} \chi_ {\varrho} (g) d \mu (g).
\]

特别地，当 E 有限维时，E 中存在向量 \(x \neq 0\)，使得 \(\varrho(g)x = x\) 对所有 \(g \in G\) 成立，当且仅当

\[
\int_ {\mathrm{G}} \chi_ {\varrho} (g) d \mu (g) \neq 0.
\]

推论 4. --- 设 \(\varrho\) 是 G 在希尔伯特空间 E 上的酉表示。空间 E 是各 \(\pi\)-同型分量 \(\mathrm{M}_{\pi}(\varrho)\)（其中 \(\pi \in \widehat{\mathrm{G}}\)）的希尔伯特直和。特别地，表示 \(\varrho\) 是半单的。

对应于 G 的不同构不可约表示的 \(\varrho\) 的同型分量彼此正交（V, p. 394, 命题 8）。由于 G 的每个有限维酉表示都是半单的（V, p. 457, 命题 1），各同型分量 \(M_{\pi}(\varrho)\) 对 \(\pi \in \widehat{G}\) 的和在 E 中稠密（命题 7）。推论由此得出。
\subsection{特征标与共轭类}

令 \(\varrho\) 为 G 在有限维希尔伯特空间 E 中的酉表示。\(\varrho\) 的特征标满足 \(|\chi_{\varrho}| \leqslant \dim(E)\) 。令 \((e_i)_{i \in I}\) 为 E 的标准正交基；特征标 \(\chi_{\varrho}\) 是对角矩阵系数 \(g \mapsto \langle e_i | \varrho(g)e_i \rangle\)（属于 \(\varrho\)）之和。特别地，\(\varrho\) 的特征标是 G-有限函数；若 \(\varrho\) 不可约，则 \(\chi_{\varrho} \in \varrho_{\mathrm{G}}(\varrho)\) 。

有如下公式

\[
\chi_ {\varrho_ {1} \oplus \varrho_ {2}} = \chi_ {\varrho_ {1}} + \chi_ {\varrho_ {2}}, \quad \chi_ {\varrho_ {1} \otimes \varrho_ {2}} = \chi_ {\varrho_ {1}} \chi_ {\varrho_ {2}}, \quad \chi_ {\breve {\varrho}} = \chi_ {\overline {{\varrho}}} = \overline {{\chi}} _ {\varrho}
\]

（参见 A, VIII, 第 388--389 页）。

命题 8. --- 有

\begin{enumerate}
\def\labelenumi{(\arabic{enumi})}
\tightlist
\item
  \(\chi_{\pi}*\chi_{\sigma} = 0\) 对于 \(\pi, \sigma\) 均属于 \(\widehat{\mathrm{G}}\) 且 \(\pi \neq \sigma\) 的情形成立，
\end{enumerate}

\[
\chi_ {\pi} * \chi_ {\pi} = \frac {1}{\dim (\pi)} \chi_ {\pi} \quad\text{对于每个 }\pi\in\widehat {\mathrm{G}}.\tag{2}
\]

令 \(\pi\) 和 \(\sigma\) 为 G 的不可约表示。有

\[
\dim (\pi) \left(\chi_ {\pi} * \chi_ {\sigma}\right) = \dim (\pi) \boldsymbol {\gamma} _ {G} \left(\chi_ {\pi}\right) \chi_ {\sigma}
\]

（V, 第 407 页引理 4）。函数 \(\dim(\pi)\gamma_{\mathrm{G}}(\chi_{\pi})\chi_{\sigma}\) 是将 \(\chi_{\sigma}\) 投影到 \(\overline{\pi}\)-同型分量 \(\gamma_{\mathrm{G}}\) 上的正交投影（V, 第 465 页推论 2），即投影到 \(\varrho_{\mathrm{G}}(\pi)\) 上（V, 第 463 页推论 4）。由于 \(\chi_{\sigma}\) 属于 \(\varrho_{\mathrm{G}}(\sigma)\)，结论由 V, 第 462 页定理 1 得出。

引理 2. --- 等价关系“\(x \in G\) 且 \(y \in G\) 且 \(x\) 与 \(y\) 共轭”在 G 中的图像是闭集。

事实上，该图像是从紧空间 G × G 到自身、由 \((x, y) \mapsto (x, yxy^{-1})\) 定义的连续映射的像。

记 G \(^{\sharp}\) 为赋予商拓扑的 G 的共轭类空间；由引理 2 以及 TG, I, 第 78 页命题 8，它是紧空间。令 \(\varpi\) : G → G \(^{\sharp}\) 为典范投影。由从 \(\mathcal{C}(G^{\sharp})\) 到 \(\mathcal{C}(G)\) 的映射 \(f \mapsto f \circ \varpi\)，将星代数 \(\mathcal{C}(G^{\sharp})\) 识别为 \(\mathcal{C}(G)\) 中由连续中心函数组成的星子代数。

G \(^{\sharp}\) 上的测度与 G 上的中心测度相对应（V, 第 402 页定义 1）。

于是，在 \(\mathcal{C}(\mathrm{G}^{\sharp})\) 上由 \(f\mapsto\int_{\mathrm{G}}f\) 定义的线性泛函是 \(\mathrm{G}^{\sharp}\) 上质量为 1 的正测度，记为 \(\mu^{\sharp}\)。对于每个 \(p\in[1,+\infty]\)，记 \(\mathcal{L}^{p}(\mathrm{G}^{\sharp})\)（分别记 \(\mathrm{L}^{p}(\mathrm{G}^{\sharp})\)）为 \(\mathcal{L}_{\mathbf{C}}^{p}(\mathrm{G}^{\sharp},\mu^{\sharp})\)（分别为 \(\mathrm{L}_{\mathbf{C}}^{p}(\mathrm{G}^{\sharp},\mu^{\sharp})\)）。我们将 \(\mathrm{L}^{p}(\mathrm{G}^{\sharp})\) 识别为 \(\mathcal{C}(\mathrm{G}^{\sharp})\) 在 \(\mathrm{L}^{p}(\mathrm{G})\) 中的闭包；特别地，它是 \(\mathrm{L}^{p}(\mathrm{G})\) 的闭子空间。

我们还记 \(\Theta(G^{\sharp}) = \Theta(G) \cap \mathcal{C}(G^{\sharp})\) 为 G 的中心矩阵系数空间。它是 \(\mathcal{C}(G^{\sharp})\) 的含单位对合子代数。

当 G 为李群时，在 LIE, IX, 第 71 页中，空间 \(\Theta(G^{\sharp})\) 也记为 \(Z\Theta(G)\)。

引理 3. --- 令 \(\pi\) 为 G 的不可约酉表示。向量空间 \(\mathrm{L}^2 (\mathrm{G}^\sharp)\cap \varrho_{\mathrm{G}}(\pi)\) 是一维的，并由 \(\pi\) 的特征标生成。

考虑 G 的酉表示 \(\sigma\)，它作用在空间 \(\varrho_{\mathrm{G}}(\pi)\) 上，并由 \(\sigma(g) = \varrho_{\mathrm{G}}(g, g)\) 定义。由于 G × G 的酉表示 \(\varrho_{\mathrm{G}}(\pi)\) 同构于 \(\overline{\pi} \boxtimes \pi\)（V, 第 462 页命题 5），其特征标为

\[
\chi_ {\sigma} (g) = \chi_ {\overline {{\pi}} \boxtimes \pi} (g, g) = | \chi_ {\pi} (g) | ^ {2}
\]

对所有 \(g \in G\) 都成立。\(\mathrm{L}^2 (\mathrm{G}^\sharp) \cap \varrho_{\mathrm{G}}(\pi)\) 是该表示的不变元素子空间，且其维数等于

\[
\int_ {\mathrm{G}} \chi_ {\sigma} (g) d \mu (g) = \int_ {\mathrm{G}} | \chi_ {\pi} (g) | ^ {2} d \mu (g) = 1
\]

（V, 第 466 页推论 3 以及第 459 页推论 3）。由于 \(\pi\) 的特征标是 \(L^2(G^\sharp) \cap \varrho_G(\pi)\) 的非零元素，它就是该空间的一个基。

命题 9. --- G 的不可约酉表示的特征标族 \((\chi_{\pi})_{\pi \in \widehat{\mathbf{G}}}\) 是 \(\mathrm{L}^2 (\mathrm{G}^\sharp)\) 的标准正交基。

根据彼得–外尔定理（V, 第 462 页定理 1），不变元素构成闭子空间 \(L^{2}(G^{\sharp})\)，它是 \(L^{2}(G)\) 的子空间，由 G 的酉表示 \(g \mapsto \varrho_{\mathrm{G}}(g, g)\) 的不变元素构成；它是各子空间 \(L^{2}(G^{\sharp}) \cap \varrho_{\mathrm{G}}(\pi)\)（\(\pi \in \widehat{G}\)）的希尔伯特直和。因此命题由前一引理和 V, 第 459 页推论 3 得出。

推论 1. --- 向量空间 \(\Theta(\mathrm{G}^{\sharp})\) 在 \(\mathcal{C}(\mathrm{G}^{\sharp})\) 和 \(\mathrm{L}^{2}(\mathrm{G}^{\sharp})\) 中都是稠密的。

第二个断言由命题 9 得出。对于第一个断言，考虑 G 的线性表示 \(\varrho\)，它作用在巴拿赫空间 \(\mathcal{C}(\mathrm{G})\) 上，并由下式定义：

\[
\varrho (g) f (x) = f (g ^ {- 1} x g)
\]

对所有 \(f \in \mathcal{C}(\mathrm{G})\) 及每个 \(x \in \mathrm{G}\) 都成立。它是连续等距表示，而 \(\mathcal{C}(\mathrm{G}^{\sharp})\) 是该表示的不变元素空间。因此，连续投影 \(p = \varrho(1)\) 作用在 \(\mathcal{C}(\mathrm{G})\) 上，且其像为 \(\mathcal{C}(\mathrm{G}^{\sharp})\)（V, 第 466 页推论 3）；其范数满足 \(\leqslant 1\)。

令 \(f \in \mathcal{C}(\mathrm{G}^{\sharp})\) 且 \(\varepsilon > 0\)。存在 \(\widetilde{f} \in \Theta(\mathrm{G})\)，使得 \(\|f - \widetilde{f}\| \leqslant \varepsilon\)（V, 第 462 页命题 4），于是 \(\|f - p(\widetilde{f})\| = \|p(f - \widetilde{f})\| \leqslant \varepsilon\)；由于 \(p(\widetilde{f}) \in \mathcal{C}(\mathrm{G}^{\sharp})\)，我们得到 \(\Theta(\mathrm{G}^{\sharp})\) 在 \(\mathcal{C}(\mathrm{G}^{\sharp})\) 中稠密。

记 R(G) 为 G 在分离的有限维复向量空间中的连续表示所成的格罗滕迪克环（参见 LIE, IX, 第 70 页和 A, VIII, 第 182 页；将其应用于 G 在分离的有限维复向量空间中的连续表示这一加性类，并将其视为 \(\mathbf{C}[G]\)-模）。由于 G 在有限维分离拓扑向量空间中的每个连续线性表示都是半单的（V, 第 457 页命题 1），阿贝尔群 R(G) 是自由的，而不可约酉表示 \(\pi \in \widehat{\mathbf{G}}\) 的类构成一个基（A, VIII, 第 186 页命题 7）。

推论 2. --- a) G 的不可约表示的特征标族是 \(\Theta(G^{\sharp})\) 的一个基；

\begin{enumerate}
\def\labelenumi{\alph{enumi})}
\setcounter{enumi}{1}
\tightlist
\item
  映射 \(u: \mathrm{R}(\mathrm{G}) \otimes_{\mathbf{Z}} \mathbf{C} \to \Theta(\mathrm{G}^{\sharp})\)，满足 \(u(\pi \otimes 1) = \chi_{\pi}\)（对每个 \(\pi \in \widehat{\mathrm{G}}\)），是 \(\mathbf{C}\) 上代数的同构。
\end{enumerate}

第一个断言由命题 9 和推论 1 得出。

由于表示 \(\pi \in \widehat{\mathbf{G}}\) 的类构成自由 \(\mathbf{Z}\)-模 R(G) 的一个基，映射 \(u\) 定义良好。它是 \(\mathbf{C}\)-代数的态射，并由 \(a\) 可知它是同构。

推论 3. --- 令 H 为局部紧拓扑群，且 G 是 H 的紧子群。令 \(\varrho\) 为 H 在希尔伯特空间 E 中的酉表示。若 \(\pi\)-同型分量（在 \(\varrho\) 限制到 G 后）对所有 \(\pi \in \widehat{\mathbf{G}}\) 都是有限维的，则 H 的表示 \(\varrho\) 是半单的，并且 H 的每个不可约酉表示在 \(\varrho\) 中的重数都是有限的。

记 \(\sigma\) 为 \(\varrho\) 限制到 H 的紧子群 G 后的表示。令 \(\pi\) 为 G 的不可约表示。自同态 \(\sigma(\chi_{\pi}) \in \mathcal{L}(\mathrm{E})\) 是有限秩的，因为其像是 \(\overline{\pi}\)-同型分量，即 \(\sigma\) 中的相应子空间（V, 第 465 页推论 2），按假设它是有限维的。因此，自同态 \(\sigma(f)\) 对每个 \(f \in \Theta(\mathrm{G}^{\sharp})\) 都是有限秩的；对每个 \(f \in \mathcal{C}(\mathrm{G}^{\sharp})\)，它都是紧的（V, 第 468 页推论 1 以及 III, 第 4 页命题 2 的推论）。

记 \(j\) 为从 G 到 H 的典范嵌入。这是一个连续映射，并且对于 G 上的每个测度 \(\nu\) 都是 \(\nu\)-真映射（INT, V, 第 69 页，第 6 节第 1 号注 1）。令 \(\mathfrak{B}\) 为 G 中元素 \(e\) 的紧邻域滤子。对于每个 V \(\in\) \(\mathfrak{B}\)，G 上存在一个连续正中心函数 \(f_{\mathrm{V}}\)，其支撑包含于 V 且在 G 上的积分等于 1。令 \(\beta_{\mathrm{V}}\) 为像测度 \(j(f_{\mathrm{V}} \cdot \mu)\)；它是 H 上具有紧支撑的正测度，并满足 \(\varrho(\beta_{\mathrm{V}}) = \sigma(f_{\mathrm{V}})\)（INT, V, 第 71 页，第 6 节第 2 号定理 1）。根据上述结果，在 \(\mathcal{M}_{c}(\mathrm{H})\) 上，由 \(\mathfrak{B}\) 经映射 \(\mathrm{V} \mapsto \beta_{\mathrm{V}}\) 得到的滤子基满足 V, 第 402 页命题 2 的条件，断言得证。

推论 4（外尔等分布判据）

令 M 为 G 上的中心正测度集，满足 \(\nu(G)=1\) 对所有 \(\nu\in M\) 都成立。令 \(\mathfrak{F}\) 为 M 上的滤子。滤子 \(\mathfrak{F}\) 弱收敛到测度 \(\mu^{\sharp}\)，其中测度定义在 \(G^{\sharp}\) 上，其充要条件是：对每个非平凡不可约酉表示 \(\pi\)，都有

\[
\lim _ {\nu , \mathfrak {F}} \int_ {\mathrm{G}} \chi_ {\pi} (x) d \nu (x) = 0.
\]

由于 \(\nu (\mathrm{G}^{\sharp}) = 1 = \mu^{\sharp}(\mathrm{G}^{\sharp})\) 对 \(\nu \in \mathrm{M}\) 成立，上述假设意味着

\[
\lim _ {\nu , \mathfrak {F}} \int_ {\mathrm{G} ^ {\sharp}} \chi_ {\pi} (x) d \nu (x) = \int_ {\mathrm{G} ^ {\sharp}} \chi_ {\pi} (x) d \mu^ {\sharp} (x)
\]

对每个表示 \(\pi\in\widehat{\mathrm{G}}\) 都成立，因而它等价于如下条件

\[
\lim _ {\nu , \mathfrak {F}} \int_ {\mathrm{G} ^ {\sharp}} f (x) d \nu (x) = \int_ {\mathrm{G} ^ {\sharp}} f (x) d \mu^ {\sharp} (x)
\]

由于线性性，对每个函数 \(f \in \Theta(\mathrm{G}^{\sharp})\) 都成立。由于空间 \(\Theta(\mathrm{G}^{\sharp})\) 在 \(\mathcal{C}(\mathrm{G}^{\sharp})\) 中稠密（推论 1），所以该假设等价于滤子 \(\mathfrak{F}\) 收敛到 \(\mu^{\sharp}\)，并且在 \(\mathcal{M}(\mathrm{G}^{\sharp})\) 中收敛；该空间赋予由 \(\mathcal{C}(\mathrm{G})\) 上逐点收敛定义的拓扑。由于 G 紧，这一拓扑与弱收敛拓扑一致（参见 INT, III, 第 59 页，第 1 节第 9 号）。

\subsection{傅里叶余变换}

在集合 \(\widehat{\mathbf{G}}\) 上赋予离散拓扑。记 \(\mathrm{F}(\widehat{\mathbf{G}})\) 为各 \(\operatorname{End}(\mathrm{E}_{\pi})\) 的乘积代数，其中 \(\pi\) 遍历 \(\widehat{\mathbf{G}}\)，记 \(\mathrm{F}_b(\widehat{\mathbf{G}})\) 为各 \(\operatorname{End}(\mathrm{E}_{\pi})\) 的乘积星代数（I, 第 103 页例 5）；这就是由族 \((u_{\pi})_{\pi \in \widehat{\mathbf{G}}}\) 组成且满足 \(\sup \| u_{\pi} \| < +\infty\) 的集合。

\[
\pi \in \widehat {\mathbf {G}}
\]

记 \(\mathrm{F}_{0}(\widehat{\mathrm{G}})\) 为 \(\mathrm{F}_{b}(\widehat{\mathrm{G}})\) 的闭星子代数，它由族 \((u_{\pi})_{\pi\in\widehat{\mathrm{G}}}\) 组成，这些族满足 \(\|u_{\pi}\|\) 在无穷远处趋于 0。

令 \(\nu \in \mathcal{M}^1 (\mathrm{G})\)。对于每个 \(\pi \in \widehat{\mathrm{G}}\)，有 \(\| \pi (\nu)\| \leqslant \| \nu \|\)，所以族 \((\pi (\nu))_{\pi \in \widehat{\mathrm{G}}}\) 属于 \(\mathrm{F}_b(\widehat{\mathrm{G}})\)。

定义 1. --- 对每个测度 \(\nu \in \mathcal{M}^{1}(\mathrm{G})\)，将族 \((\pi (\nu))_{\pi \in \widehat{\mathrm{G}}}\) 视为 \(\mathrm{F}_b(\widehat{\mathrm{G}})\) 中的元素，记为 \(\overline{\mathcal{F}}_{\mathrm{G}}(\nu)\)。这样定义的从 \(\mathcal{M}^1 (\mathrm{G})\) 到 \(\mathrm{F}_b(\widehat{\mathrm{G}})\) 的映射称为 G 的傅里叶余变换，而 \(\overline{\mathcal{F}}_{\mathrm{G}}(\nu)\) 称为测度 \(\nu\) 的傅里叶余变换。

对于 \(f \in \mathrm{L}^1(\mathrm{G})\)，记 \(\overline{\mathcal{F}}_{\mathrm{G}}(f) = \overline{\mathcal{F}}_{\mathrm{G}}(f \cdot \mu)\)。

对于任意表示 \(\pi \in \widehat{\mathrm{G}}\)，映射 \(\nu \mapsto \pi(\nu)\) 是从 \(\mathcal{M}^{1}(\mathrm{G})\) 到 \(\operatorname{End}(\mathrm{E}_{\pi})\) 的含单位对合巴拿赫代数态射（V, 第 401 页引理 1）。因此，傅里叶余变换是从 \(\mathcal{M}^{1}(\mathrm{G})\) 到 \(\mathrm{F}_b(\widehat{\mathrm{G}})\) 的含单位对合巴拿赫代数态射。

\subsection{傅里叶变换}

沿用前一号的记号。

令 \(\pi \in \widehat{\mathbf{G}}\)。向量空间 \(\operatorname{End}(\mathrm{E}_{\pi})\) 赋予希尔伯特空间结构，其内积为

\[
\langle u _ {1} \mid u _ {2} \rangle = \dim (\pi) \operatorname{Tr} (u _ {1} ^ {*} u _ {2}) = \dim (\pi) \operatorname{Tr} (u _ {2} u _ {1} ^ {*})
\]

其中 \(u_{1}\)、\(u_{2}\) 属于 \(\operatorname{End}(\mathrm{E}_{\pi})\)（参见 EVT, V, 第 52 页定理 1）。

记 \(\|u\|_{2} = \sqrt{\langle u|u\rangle}\) 为 \(\operatorname{End}(\operatorname{E}_{\pi})\) 中元素 u 的范数，其中 \(\pi \in \widehat{G}\)。对于每个 \(g \in G\)，有 \(\|\pi(g)\|_{2} = \dim(\pi)\)，因为 \(\pi(g)\) 是酉的。

此处记作 \(\|u\|_{2}\) 的范数，与 EVT, V, 第 52 页定义的范数相差因子 \(\dim(\pi)\)；后者定义在 \(E_{\pi}\) 上的希尔伯特–施密特算子空间中。

引理 4. --- 令 \(\pi\) 为 G 的不可约表示。映射 \(f \mapsto \pi(f)\) 通过取子空间定义了从 \(\varrho_{\mathrm{G}}(\overline{\pi})\) 到 \(\operatorname{End}(\operatorname{E}_{\pi})\) 的等距同构。

置 \(\varepsilon_{\pi}(g,h)u = \pi(g)\circ u\circ\pi(h^{-1})\)，其中 \((g,h)\in\mathrm{G}\times\mathrm{G}\) 且 \(u\in\mathrm{End}(\mathrm{E}_{\pi})\)。映射 \(\varepsilon_{\pi}\) 是 \(\mathrm{G}\times\mathrm{G}\) 在 \(\mathrm{End}(\mathrm{E}_{\pi})\) 中的连续表示。它是酉的。事实上，由于 \(\pi\) 本身是酉的，有

\[
\begin{array}{c} \| \varepsilon_ {\pi} (g, h) u \| _ {2} ^ {2} = \dim (\pi) \operatorname{Tr} \Big ((\pi (g) u \pi (h ^ {- 1})) ^ {*} \pi (g) u \pi (h ^ {- 1}) \Big) \\ = \dim (\pi) \operatorname{Tr} (\pi (h) u ^ {*} u \pi (h) ^ {- 1}) = \dim (\pi) \operatorname{Tr} (u ^ {*} u) = \| u \| _ {2} ^ {2} \end{array}
\]

对所有 \((g,h)\in \mathrm{G}\times \mathrm{G}\) 及每个 \(u\in \operatorname {End}(\operatorname {E}_{\pi})\) 都成立。

映射 \(\Psi\)，定义为 \(f\mapsto\pi(f)\)，是一个 \((G\times G)\)-态射，从 \(\varrho_{G}(\overline{\pi})\) 到 \(\operatorname{End}(E_{\pi})\)，因为

\[
\pi (\varrho_ {\mathrm{G}} (g, h) f) = \int_ {\mathrm{G}} f (g ^ {- 1} x h) \pi (x) d \mu (x) = \pi (g) \pi (f) \pi (h ^ {- 1})
\]

对所有 \((g,h)\in\mathrm{G}\times\mathrm{G}\) 及每个 \(f\in\varrho_{\mathrm{G}}(\overline{\pi})\) 都成立。由于 \(\varrho_{\mathrm{G}}(\overline{\pi})\) 是不可约表示（V, 第 462 页定理 1，\(a\)），存在 \(\lambda\in\mathbf{C}^{*}\)，使得映射 \(\lambda\Psi\) 为零映射或等距映射（V, 第 388 页推论 5，\(a\)）。令 \(f=\dim(\pi)\overline{\chi}_{\pi}\in\varrho_{\mathrm{G}}(\overline{\pi})\)。得到 \(\pi(f)=1_{\mathrm{E}_{\pi}}\)（V, 第 465 页推论 2），从而 \(\|\pi(f)\|_{2}=\dim(\pi)=\|f\|\)（V, 第 459 页推论 3）。因此，映射 \(\Psi\) 是等距的。由于 \(\varrho_{\mathrm{G}}(\overline{\pi})\) 与 \(\operatorname{End}(\mathrm{E}_{\pi})\) 维数相同，\(\Psi\) 是等距同构。

记 \(\mathrm{F}^{2}(\widehat{\mathrm{G}})\) 为希尔伯特空间 \(\mathrm{End}(\mathrm{E}_{\pi})\) 的希尔伯特直和。元素 \(x \in \mathrm{F}^{2}(\widehat{\mathrm{G}})\) 的范数仍记为 \(\|x\|_{2}\)。

在 LIE, IX, 第 79 页中，该空间记为 \(\mathrm{L}^{2}(\widehat{\mathrm{G}})\)；这里我们不采用这一记号，以免与空间 \(\ell^{2}(\widehat{\mathrm{G}})\) 混淆。令 \((u_{\pi})_{\pi\in\widehat{\mathrm{G}}}\) 为 \(\mathrm{F}^{2}(\widehat{\mathrm{G}})\) 的元素。由于

\[
\sum_ {\pi \in \widehat {G}} \| u _ {\pi} \| _ {2} ^ {2} = \sum_ {\pi \in \widehat {G}} \dim (\pi) \operatorname{Tr} (u _ {\pi} ^ {*} u _ {\pi}) <   + \infty
\]

并且 \(\|u_{\pi}\|^{2}\leqslant\mathrm{Tr}(u_{\pi}^{*}u_{\pi})\)（参见 EVT, V, 第 52 页公式 (33)），在 \(\mathcal{L}(\mathrm{E}_{\pi})\) 中自同态 \(u_{\pi}\) 的范数在无穷远处趋于 0。因此可以将 \(\mathrm{F}^{2}(\widehat{\mathrm{G}})\) 识别为 \(\mathrm{F}_{0}(\widehat{\mathrm{G}})\) 的子空间。

令 \(\pi \in \widehat{\mathbf{G}}\) 且 \(u \in \mathrm{End}(\mathrm{E}_{\pi})\)。记 \(\mathcal{F}_{\pi}(u)\) 为 G 上由 \(\mathcal{F}_{\pi}(u)(g) = \langle \pi(g) | u \rangle = \dim(\pi) \mathrm{Tr}(\pi(g)^{*} u)\) 定义的函数（对所有 \(g \in \mathbf{G}\)）。它是 G 上的连续函数。若 G 是紧实李群，则函数 \(\mathcal{F}_{\pi}(u)\) 在 G 上解析（LIE, III, 第 8 节第 1 号定理 1）。

由于 G 是紧的，可以将 \(L^{2}(G)\) 识别为 \(L^{1}(G)\) 的子空间。

定理 2. --- G 的傅里叶余变换通过取子空间诱导出从 L²(G) 到 F²(Ĝ) 的等距同构。其逆 \(\mathcal{F}_{\mathrm{G}}\) 将 F²(Ĝ) 中的元素 \((u_{\pi})_{\pi \in \widehat{\mathrm{G}}}\) 对应于级数之和

\[
\sum_ {\pi \in \widehat {\mathrm{G}}} \mathcal {F} _ {\pi} (u _ {\pi}),
\]

该级数在 \(\mathrm{L}^2 (\mathrm{G})\) 中收敛。

根据彼得–外尔定理（V, 第 462 页定理 1），希尔伯特空间 \(L^{2}(G)\) 是各空间 \(\varrho_{\mathrm{G}}(\overline{\pi})\)（\(\pi \in \widehat{G}\)）的希尔伯特直和。对于每个 \(\pi \in \widehat{G}\)，线性映射 \(f \mapsto \pi(f)\) 从 \(\varrho_{\mathrm{G}}(\overline{\pi})\) 到 \(\operatorname{End}(\operatorname{E}_{\pi})\) 是等距同构（引理 4）。因此，傅里叶余变换在 \(L^{2}(G)\) 上的限制定义了从 \(L^{2}(G)\) 到 \(F^{2}(\widehat{G})\) 的等距同构。

令 \(f \in \mathrm{L}^2(\mathrm{G})\)。令 \(\pi \in \widehat{\mathrm{G}}\)，并令 \(f_{\overline{\pi}} \in \mathcal{C}(\mathrm{G})\) 为 \(f\) 在 \(\varrho_{\mathrm{G}}(\overline{\pi})\) 上的正交投影（V, 第 462 页定理 1）。由于该空间是 \(\overline{\pi}\)-同型分量，来自 \(\pmb{\delta}_{\mathrm{G}}\)（V, 第 463 页推论 4），由 V, 第 465 页推论 2，有 \(f_{\overline{\pi}} = \dim(\pi)\pmb{\delta}_{\mathrm{G}}(\chi_{\pi})(f)\)。于是对每个 \(x \in \mathrm{G}\)，有

\[
\begin{array}{r l} & {f _ {\overline {{\pi}}} (x) = \dim (\pi) \int_ {\mathrm{G}} \chi_ {\pi} (g) (\pmb {\delta} _ {\mathrm{G}} (g) f) (x) d \mu (g)} \\ & {\qquad = \dim (\pi) \int_ {\mathrm{G}} \chi_ {\pi} (g) f (x g) d \mu (g)} \\ & {\qquad = \dim (\pi) \int_ {\mathrm{G}} \chi_ {\pi} (x ^ {- 1} y) f (y) d \mu (y)} \\ & {\qquad = \dim (\pi) \mathrm{Tr} (\pi (x ^ {- 1}) \pi (f)) = \langle \pi (x) | \pi (f) \rangle ,} \end{array}
\]

即 \(f_{\overline{\pi}} = \mathcal{F}_{\pi}(\pi(f))\)。

由于 \(f\) 在 \(\mathrm{L}^2 (\mathrm{G})\) 中是族 \((f_{\overline{\pi}})\) 的和，根据彼得–外尔定理，得到

\[
f = \sum_ {\pi \in \widehat {\mathrm{G}}} \mathcal {F} _ {\pi} (\pi (f))
\]

其中级数在 \(L^{2}(G)\) 中收敛。这就证明了定理。

定义 2. --- \(\mathcal{F}_{\mathrm{G}}\) 是从 \(\mathrm{F}^2 (\widehat{\mathrm{G}})\) 到 \(\mathrm{L}^2 (\mathrm{G})\) 的等距同构，也是傅里叶余变换所诱导同构的逆；称其为 G 的傅里叶变换。称 \(\mathrm{F}^2 (\widehat{\mathrm{G}})\) 中元素的像为该元素的傅里叶变换。

注记。--- 1) \((u_{\pi})_{\pi \in \widehat{\mathrm{G}}} \in \mathrm{F}^2(\widehat{\mathrm{G}})\) 的傅里叶变换因此是下列级数在 \(\mathrm{L}^2(\mathrm{G})\) 中的等价类

\[
g \mapsto \sum_ {\pi \in \widehat {\mathrm{G}}} \langle \pi (g) | u _ {\pi} \rangle = \sum_ {\pi \in \widehat {\mathrm{G}}} \dim (\pi) \operatorname{Tr} (u _ {\pi} \circ \pi (g) ^ {- 1}).
\]

\begin{enumerate}
\def\labelenumi{\arabic{enumi})}
\setcounter{enumi}{1}
\tightlist
\item
  令 \(f \in \mathrm{L}^2(\mathrm{G})\)。于是有普朗歇雷尔公式
\end{enumerate}

\[
\| f \| ^ {2} = \| \overline {{\mathcal {F}}} _ {\mathrm{G}} (f) \| ^ {2} = \sum_ {\pi \in \widehat {\mathrm{G}}} \| \pi (f) \| ^ {2}.
\]

此外，根据定理 2，\(f\) 是 \(\mathrm{L}^2 (\mathrm{G})\) 中族 \((f_{\pi})_{\pi \in \widehat{\mathrm{G}}}\) 的和，其中

\[
\begin{array}{l} f _ {\pi} (x) = \mathcal {F} _ {\pi} (\pi (f)) (x) \\ \qquad = \langle \pi (x) | \pi (f) \rangle = \dim (\pi) \int_ {\mathrm{G}} f (g) \chi_ {\pi} (x ^ {- 1} g) d \mu (g) \end{array}
\]

对所有 \(x \in G\) 及每个 \(\pi \in \widehat{G}\) 都成立。

设 G 为交换群。由于 G 的不可约表示维数为 1（V, 第 390 页推论 7），且对偶群 \(\widehat{\mathbf{G}}\) 是离散的（II, 第 233 页命题 18），代数 \(\mathrm{F}_b(\widehat{\mathbf{G}})\) 和 \(\mathrm{F}_0(\widehat{\mathbf{G}})\) 分别识别为代数 \(\mathcal{C}_b(\widehat{\mathbf{G}})\)（在 \(\widehat{\mathbf{G}}\) 上的连续有界函数的代数）和代数 \(\mathcal{C}_0(\widehat{\mathbf{G}})\)（在 \(\widehat{\mathbf{G}}\) 上在无穷远处趋于 0 的连续函数的代数）。由于 G 是紧的，\(\widehat{\mathbf{G}}\) 上与 \(\mu\) 对偶的哈尔测度是计数测度 \(\widehat{\mu}\)（II, 第 233 页命题 18）。希尔伯特直和 \(\mathrm{F}^2 (\widehat{\mathrm{G}})\) 由各空间 \(\operatorname{End}(\mathrm{E}_{\pi})\)（\(\pi \in \widehat{\mathrm{G}}\)）组成，并识别为希尔伯特空间 \(\mathrm{L}^2 (\widehat{\mathrm{G}},\widehat{\mu})\)。

令 \(\nu \in \mathcal{M}^1 (\mathrm{G})\)。于是对每个酉特征标 \(\chi \in \widehat{\mathrm{G}}\)，有

\[
\chi (\nu) = \int_ {\mathrm{G}} \chi \nu
\]

这证明了上面定义的傅里叶余变换与 II, 第 206 页定义的 G 的傅里叶余变换一致。

每个 \(f \in \mathcal{L}^2(\mathrm{G})\) 都可积，而公式 \(\| \overline{\mathcal{F}}_{\mathrm{G}}(f)\| _2 = \| f\|\) 就是普朗歇雷尔公式（II, 第 215 页定理 1）。

命题 10. --- 傅里叶余变换是从 \(\mathcal{M}^{1}(\mathrm{G})\) 到 \(\mathrm{F}_{b}(\widehat{\mathrm{G}})\) 的单射对合巴拿赫代数态射，并将 \(\mathrm{L}^{1}(\mathrm{G})\) 映入 \(\mathrm{F}_{0}(\widehat{\mathrm{G}})\)。

由于 G 是紧的，空间 \(\mathcal{C}(\mathrm{G})\) 包含于 \(\mathcal{L}^{1}(\mathrm{G})\) 和 \(\mathcal{L}^{2}(\mathrm{G})\)，并且在二者中都稠密。

令 \(\nu \in \mathcal{M}^1 (\mathrm{G})\) 满足 \(\overline{\mathcal{F}}_{\mathrm{G}}(\nu) = 0\)。于是对 G 上的每个连续函数 \(f\)，都有 \(\overline{\mathcal{F}}_{\mathrm{G}}(\nu *f) = 0\)。现在 \(\nu *f\) 属于 \(\mathcal{C}(\mathrm{G})\)（INT, VIII, 第 152 页，第 4 节第 2 号命题 3）。由于根据定理 2，傅里叶余变换在 \(\mathrm{L}^2 (\mathrm{G})\) 上是单射，因而在空间 \(\mathcal{C}(\mathrm{G})\) 上也是单射，所以 \(\nu *f = 0\) 对 \(f\in \mathcal{C}(\mathrm{G})\) 成立。特别地，可得

\[
\int_ {\mathrm{G}} f (x) d \nu (x) = (\nu * \check {f}) (e) = 0
\]

对每个函数 \(f \in \mathcal{C}(\mathrm{G})\) 都成立（INT, VIII, 上引处），所以测度 \(\nu\) 为零。

傅里叶余变换下 \(\mathcal{C}(\mathrm{G})\) 的像包含于 \(\mathrm{F}^2 (\widehat{\mathrm{G}})\)，从而也包含于 \(\mathrm{F}_0(\widehat{\mathrm{G}})\)。由于 \(\mathrm{F}_0(\widehat{\mathrm{G}})\) 在 \(\mathrm{F}_b(\widehat{\mathrm{G}})\) 中是闭的，且 \(\mathcal{C}(\mathrm{G})\) 在 \(\mathrm{L}^1 (\mathrm{G})\) 中稠密，傅里叶余变换下 \(\mathrm{L}^1 (\mathrm{G})\) 的像也包含于 \(\mathrm{F}_0(\widehat{\mathrm{G}})\)。

命题 11. --- a) 令 \(u = (u_{\pi})_{\pi \in \widehat{\mathbf{G}}}\) 为 \(\mathrm{F}(\widehat{\mathbf{G}})\) 中的元素。若族 \((\mathcal{F}_{\pi}(u_{\pi}))_{\pi \in \widehat{\mathbf{G}}}\) 在 \(\mathcal{C}(\mathbf{G})\) 中一致可求和，则其和 \(f\) 是 \(\mathbf{G}\) 上的连续函数，且其傅里叶余变换为 \(u\)；

\begin{enumerate}
\def\labelenumi{\alph{enumi})}
\setcounter{enumi}{1}
\tightlist
\item
  令 \(f \in \mathrm{L}^1(\mathrm{G})\)。若族 \((\mathcal{F}_\pi(\pi(f)))_{\pi \in \widehat{\mathrm{G}}}\) 在 \(\mathcal{C}(\mathrm{G})\) 中一致可求和，则其和是 \(\mathrm{G}\) 上的连续函数，且它在 \(\mathrm{L}^1(\mathrm{G})\) 中的等价类等于 \(f\)。
\end{enumerate}

证明断言 \(a\)。由于 G 是紧的，\(f\) 是族 \((\mathcal{F}_{\pi}(u_{\pi}))\) 的和，属于 \(\mathrm{L}^2 (\mathrm{G})\)，且级数

\[
\sum_ {\pi \in \widehat {\mathrm{G}}} \mathcal {F} _ {\pi} (u _ {\pi})
\]

在 \(L^{2}(G)\) 中收敛于 f（INT, IV, 第 127 页，第 3 节第 3 号命题 4）。因此 \(\mathcal{F}_{\mathrm{G}}((u_{\pi})_{\pi}) = f\) 在 \(L^{2}(G)\) 中成立，从而 \((u_{\pi})_{\pi \in \widehat{\mathrm{G}}} = \overline{\mathcal{F}}_{\mathrm{G}}(f)\)（定理 2）。

证明断言 b)。可以将断言 a) 应用于族 \((\pi(f))_{\pi \in \widehat{\mathrm{G}}}\)。族 \((\mathcal{F}_{\pi}(\pi(f)))_{\pi \in \widehat{\mathrm{G}}}\) 的和 g 是连续函数，因而属于 \(L^{2}(G)\)，并满足 \(\overline{\mathcal{F}}_{\mathrm{G}}(g) = (\pi(f))_{\pi \in \widehat{\mathrm{G}}}\)。由于 \(\overline{\mathcal{F}}_{\mathrm{G}}\) 在 \(L^{2}(G)\) 上是单射，所以在 \(L^{2}(G)\) 中 f = g，因而在 \(L^{1}(G)\) 中也有 f = g。

将 \(\mathcal{C}(\widehat{\mathrm{G}})\) 识别为 \(\widehat{\mathrm{G}}\) 上复值函数的代数，它是代数 \(\mathrm{F}(\widehat{\mathrm{G}})\) 的中心：该映射将 \(f\colon \widehat{\mathrm{G}}\to \mathbf{C}\) 对应于中心元素 \((f(\pi)1_{\mathrm{E}_{\pi}})_{\pi \in \widehat{\mathrm{G}}}\)（属于 \(\mathrm{F}(\widehat{\mathrm{G}})\)）。记 \(\mathcal{M}^1 (\mathrm{G}^\sharp)\) 为 G 上中心有界测度的空间。它是巴拿赫代数 \(\mathcal{M}^1 (\mathrm{G})\) 的闭子空间。对于每个测度 \(\nu \in \mathcal{M}^1 (\mathrm{G}^\sharp)\) 和每个表示 \(\pi \in \widehat{\mathrm{G}}\)，有 \(\pi (\nu)\in \mathrm{End}_{\mathrm{G}}(\mathrm{E}_{\pi})\)，因而根据舒尔引理（V, 第 386 页命题 6），\(\pi (\nu)\) 是 \(\mathrm{E}_{\pi}\) 的恒等自同态的数量倍。傅里叶余变换在 \(\mathcal{M}^1 (\mathrm{G}^\sharp)\) 上的限制因而识别为从 \(\mathcal{M}^1 (\mathrm{G}^\sharp)\) 到 \(\mathcal{C}(\widehat{\mathrm{G}})\) 的含单位对合巴拿赫代数态射。该态射是单射；\(\mathrm{L}^1 (\mathrm{G}^\sharp)\) 的像包含于 \(\mathcal{C}_b(\widehat{\mathrm{G}})\)，而其在 \(\mathrm{L}^2 (\mathrm{G}^\sharp)\) 上的限制是到 \(\mathrm{L}^2 (\widehat{\mathrm{G}},\widehat{\mu})\) 的等距同构。

\subsection{弗罗贝尼乌斯--舒尔指标与拉森择一定理}

回忆（LIE, VIII, 第 7 节第 6 号定义 2）：若 G 在有限维复向量空间 E 中的不可约表示存在一个在 G 下不变的非零对称双线性形式（分别地，存在一个在 G 下不变的非零交错双线性形式），则称该表示为正交型（分别为辛型）。若 E 上不存在在 G 下不变的非零双线性形式，则称其为复型。

当 E 中的不可约表示为正交型（分别为辛型）时，E 上 G-不变对称（分别为交错）双线性形式的空间是一维的，并且每个非零 G-不变双线性形式都是分离的。

正交型表示有时称为实型表示，辛型表示有时称为四元数型表示（参见 LIE, IX, 附录 II）。

令 \(\pi\) 为 G 在复向量空间 E 中的不可约表示。上述三种情形由量

\[
\mathrm{FS} (\pi) = \int_ {\mathrm{G}} \chi_ {\pi} (g ^ {2}) d \mu (g),
\]

区分；称其为 π 的弗罗贝尼乌斯--舒尔指标。事实上有

\[
\operatorname{FS} (\pi) = \left\{ \begin{array}{l l} 1 & \text{若 }\pi\text{ 属于正交型，} \\ - 1 & \text{若 }\pi\text{ 属于辛型，} \\ 0 & \text{若 }\pi\text{ 属于复型。} \end{array} \right.
\]

（LIE, IX, 附录 2，第 103 页命题 3 及第 105 页命题 4）。

当 G 为李群时，可以利用 LIE, IX, 第 69 页命题 1 计算 \(\mathrm{FS}(\pi)\)。

定义 3. --- 令 \(\varrho\) 为 G 在希尔伯特空间 E 中的有限维酉表示。令 \(k\) 为正整数。称表示 \(\overline{\varrho}^{\otimes k} \otimes \varrho^{\otimes k}\) 的 G-不变元素子空间的维数为 \(2k\) 阶的 \(\varrho\) 的绝对矩，记为 \(\mathrm{M}_{2k}(\varrho)\)。

定理 3（拉森择一定理）。--- 设 G 是无限群。令 π 为 G 在维数 ≥ 2 的有限维希尔伯特空间 E 中的忠实酉表示。

\begin{enumerate}
\def\labelenumi{\alph{enumi})}
\item
  假设 G 的导出群是无限的。有 \(M_{4}(\pi) \geqslant 2\)，且等号成立当且仅当 G 包含 SU(E)；
\item
  假设 \(\dim(E) \geqslant 3\)，\(\pi\) 为正交型或辛型，且 \(\dim(E) \neq 4\)（在 \(\pi\) 为正交型时）。则 \(M_4(\pi) \geqslant 3\)，且等号成立当且仅当 G 的复化李代数等于 E 上某个分离的 G-不变双线性形式的李代数。*当且仅当 G 的单位分支 \(G_0\) 是这种双线性形式的自同构群的极大紧子群时，情形才成立。*
\end{enumerate}

证明中要用到下列引理。

引理 5. --- 令 \(\pi\) 为 G 在有限维希尔伯特空间 E 中的酉表示。令 \((\varrho_{i})_{i\in\mathrm{I}}\) 为 G 的非零酉表示族，\((n_{i})_{i\in\mathrm{I}}\) 为满足 \(\geqslant 1\) 的整数族，使得 G 的表示 \(\overline{\pi}\otimes\pi\)（分别地，表示 \(\pi\otimes\pi\)）同构于直和 \(\bigoplus_{i\in\mathrm{I}}\varrho_{i}^{n_{i}}\)。则有

\[
\mathrm{M} _ {4} (\pi) \geqslant \sum_ {i \in \mathrm{I}} n _ {i} ^ {2}
\]

等号成立当且仅当表示 \(\varrho_{i}\) 都不可约且两两不同构。

记 \(\chi_{i}\) 为 \(\varrho_{i}\) 的特征标，其中 \(i\in I\)。有

\[
| \chi_ {\pi} | ^ {2} = \chi_ {\overline {{\pi}} \otimes \pi} = \sum_ {i \in \mathrm{I}} n _ {i} \chi_ {i}, \quad (\text { resp. } \chi_ {\pi} ^ {2} = \chi_ {\pi \otimes \pi} = \sum_ {i \in \mathrm{I}} n _ {i} \chi_ {i}).
\]

根据定义和 V, 第 466 页推论 3，有

\[
\mathrm{M} _ {4} (\pi) = \int_ {\mathrm{G}} | \chi_ {\pi} | ^ {4} d \mu ,
\]

从而在两种情形下均有公式

\[
\mathrm{M} _ {4} (\pi) = \sum_ {i, j} n _ {i} n _ {j} \int_ {\mathrm{G}} \chi_ {i} \overline {{\chi}} _ {j} d \mu = \sum_ {i, j} n _ {i} n _ {j} \dim \mathrm{Hom} _ {\mathrm{G}} (\varrho_ {i}, \varrho_ {j})
\]

（V, 第 459 页推论 4, b)）。令 \(i \in I\)。由于表示 \(\varrho_{i}\) 非零，空间 \(\mathrm{Hom}_{\mathrm{G}}(\varrho_{i}, \varrho_{i})\) 非零，于是得到下界

\[
\mathrm{M} _ {4} (\pi) \geqslant \sum_ {i \in \mathrm{I}} n _ {i} ^ {2}.
\]

此外，由于 \(n_{i} \geqslant 1\)，等号成立当且仅当 \(\dim \operatorname{Hom}_{\mathrm{G}}(\varrho_{i}, \varrho_{j}) = 0\)（当 \(i \neq j\) 时），并且对每个 i 都有 \(\dim \operatorname{Hom}_{\mathrm{G}}(\varrho_{i}, \varrho_{i}) = 1\)。第二个条件成立当且仅当表示 \(\varrho_{i}\) 不可约（上引处）。此时根据舒尔引理（V, 第 387 页推论 2），第一个条件成立当且仅当表示 \(\varrho_{i}\) 两两不同构。

若 E 是交换环 A 上的模，则称 E 的对称平方（分别为外平方）为 A-模 \(S^{2}(E)\)（分别为 \(\wedge^{2}E\)）。

引理 6. --- 令 \(q\) 为维数 \(\geqslant 3\) 的有限维复向量空间 E 上的分离双线性形式。

\begin{enumerate}
\def\labelenumi{\alph{enumi})}
\item
  若 \(q\) 是对称的，则正交李代数 \(\mathfrak{so}(q)\) 的伴随表示同构于 \(\wedge^2\mathrm{E}\)，即自然表示 \(\mathfrak{so}(q) \to \mathfrak{gl}(\mathrm{E})\) 的外平方；
\item
  若 q 是交错的，则辛李代数 \(\mathfrak{sp}(q)\) 的伴随表示同构于 \(\mathsf{S}^2 (\mathrm{E})\)，即自然表示 \(\mathfrak{sp}(q)\to \mathfrak{gl}(\mathrm{E})\) 的对称平方。
\end{enumerate}

在 C 上赋予恒等对合自同构。于是双线性形式 q 是 \(\varepsilon\)-埃尔米特的（A, IX, 第 3 节第 1 号定义 1），其中情形 a) 中 \(\varepsilon = 1\)，情形 b) 中 \(\varepsilon = -1\)。对于每个 \(u \in \text{End}(\text{E})\)，记 \({}^{t}u\) 为关于 q 的 u 的伴随。于是有 \(q(x, u(y)) = q(^{t}u(x), y)\)，对所有 \((x, y) \in \text{E} \times \text{E}\) 成立。令 v 为 E \(\otimes\) E 的唯一自同构，使得 \(v(x \otimes y) = y \otimes x\) 对所有 \((x, y) \in \text{E}^{2}\) 成立。

存在唯一的线性映射 \(w\) 从 \(\mathrm{E} \otimes \mathrm{E}\) 到 \(\operatorname{End}(\mathrm{E})\)，使得 \(w(a \otimes b)\) 是由 \(x \mapsto q(a, x)b\) 定义的线性映射，对所有 \((a, b, x) \in \mathrm{E}^3\) 成立。映射 \(w\) 是同构。对于每个 \(s \in \mathrm{E} \otimes \mathrm{E}\)，有

\[
{ } ^ { t } ( w ( s ) ) = \varepsilon w ( v ( s ) ) .\tag{3}
\]

事实上，对于 \((x,y)\) 和 \((a,b)\) 在 \(\mathrm{E} \times \mathrm{E}\) 中的所有情形，可得

\[
q (x, w (a \otimes b) y) = q (a, y) q (x, b) = \varepsilon q (b, x) q (a, y) = q (\varepsilon w (b \otimes a) x, y).
\]

最后，记 \(x \mapsto x^*\) 为由 q 导出的从 E 到 \(\mathrm{E}'\) 的同构，即由 \(q\) 给出，满足 \(\langle x^*, y \rangle = q(x, y)\)，对所有 \((x, y) \in \mathrm{E}^2\) 成立。

记号确定后，令 H 为 q 的正交群（分别为辛群）。它是 \(\mathbf{GL}(\mathrm{E})\) 的李子群，其李代数 h 是由 \(u \in \operatorname{End}(\operatorname{E})\) 满足 \(^{t}u = -u\) 的 End(E) 子空间（LIE, III, 第 146 页推论 1）。由于形式 q 在 H 下不变，故有

\[
\langle (g a) ^ {*}, b \rangle = q (g a, b) = \langle a, g ^ {- 1} b \rangle
\]

对所有 \(g \in H\) 及每个 \((a, b) \in E \times E\) 都成立。

在 End(E) 上赋予 H 的伴随表示 Ad。线性映射 \(w\) 是 H 的表示之间的态射：事实上，对每个 \(g \in \mathrm{H}\) 及每个 \((a, b, x) \in \mathrm{E}^3\)，有

\[
\begin{array}{c} w (g   (a \otimes b)) (x) = \langle (g a) ^ {*}, x \rangle   g b = \langle a, g ^ {- 1} x \rangle   g b \\ = g   (\langle a, g ^ {- 1} x \rangle   b) = \mathrm{Ad} (g) (w (a \otimes b)) x, \end{array}
\]

由线性性即得结论。

由于 v 也是 H 的表示之间的态射，线性映射 \(\theta = w - \varepsilon w \circ v\) 也具有这一性质；因此它是 h-模之间的态射。

记 \(F \subset E \otimes E\) 为由元素 \(s \in E \otimes E\) 且满足 \(v(s) = -\varepsilon s\) 所成的子空间。若 \(\varepsilon = -1\)，则从 \(E \otimes E\) 到 E 的对称平方的典范映射在 F 上的限制是 C-向量空间同构（A, III, 第 72 页）；若 \(\varepsilon = 1\)，则从 \(E \otimes E\) 到外平方的典范映射在 F 上的限制是 C-向量空间同构（上引处，第 82 页）。

F 在 \(\theta\) 下的像包含于 \(\mathfrak{h}\) 中：事实上，对每个 \(s\in\mathrm{F}\)，公式 (3) 蕴含

\[
{ } ^ { t } \theta ( s ) = { } ^ { t } w ( s ) - \varepsilon { } ^ { t } w ( v ( s ) ) = \varepsilon w ( v ( s ) ) - w ( s ) = - \theta ( s ) .
\]

按 F 的定义，映射 \(\theta\) 在 F 上的限制与 \(2w\) 的限制重合，因而线性映射 \(\theta\) 在 F 上是单射。由于 \(\dim(F) = \dim(h)\)。由（A, III, 第 75 页定理 1 和第 87 页推论 1，以及 LIE, VIII, 第 192 页和第 201 页），可知 \(\theta\) 通过取子空间诱导出从 F 到 \(h\) 的 \(h\)-模同构，故引理得证。

引理 7. --- 令 \(H_{2}\) 为实李群，令 \(H_{1}\) 为 \(H_{2}\) 的闭子群。\(H_{1}\) 的复化李代数可识别为 \(H_{1}\) 在 \(H_{2}\) 的复化李代数上的伴随表示的一个子表示。

事实上，限制到 \(H_{1}\) 的 \(H_{2}\) 李代数上的伴随表示，正可识别为 \(H_{1}\) 的伴随表示。

引理 8. --- 令 E 为维数 n 的复希尔伯特空间。

\begin{enumerate}
\def\labelenumi{\alph{enumi})}
\item
  群 \(\mathbf{SU}(\mathrm{E})\) 和 \(\mathbf{U}(\mathrm{E})\) 是连通的；
\item
  SU(E) 的导出群等于 SU(E)；
\item
  U(E) 的导出子群等于 SU(E)。
\end{enumerate}

不妨设 \(n \geqslant 1\) 且 \(E = C^{n}\)。

令 A 为 \(\mathbf{U}(\mathrm{E})\) 中由对角矩阵组成的子群；它同胚于 \(U^{n}\)，因而是连通的（TG, VI, 第 11 页推论 2 以及 I, 第 83 页命题 8）。它与 \(\mathbf{SU}(\mathrm{E})\) 的交同胚于 \(U^{n-1}\)，所以也是连通的。

由 IV, 第 149 页定理 1，群 \(\mathbf{U}(\mathrm{E})\)（分别为 \(\mathbf{SU}(\mathrm{E})\)）是连通子集 \(g\mathrm{Ag}^{-1}\)（\(g \in \mathrm{G}\)）的并（分别是连通子集 \(g(\mathrm{A} \cap \mathbf{SU}(\mathrm{E}))g^{-1}\) 的并）；由于单位元属于这些集合中的每一个，空间 \(\mathbf{U}(\mathrm{E})\)（分别为 \(\mathbf{SU}(\mathrm{E})\)）是连通的（TG, I, 第 81 页命题 2）。

证明 \(b\)。当 \(n = 1\) 时断言成立，因为此时 \(\mathbf{SU}(\mathrm{E})\) 退化为单位元。因此设 \(n \geqslant 2\)。此时 \(\mathbf{SU}(\mathrm{E})\) 的李代数是单的（LIE, IX, 第 20 页，第 3 节第 4 号），所以 \(\mathbf{SU}(\mathrm{E})\) 的导出群在 \(\mathbf{SU}(\mathrm{E})\) 中具有有限指数。由于 \(\mathbf{SU}(\mathrm{E})\) 由 \(a\) 可知连通，故得到所需结果。

断言 \(c\) 由 \(b\) 得出，因为 \(\mathbf{U}(\mathrm{E})\) 的导出群包含于 \(\mathbf{SU}(\mathrm{E})\)。

现在证明定理 3。记 n 为希尔伯特空间 E 的维数。由于表示 \(\pi\) 忠实，不妨设 G 是 U(E) 的紧子群。群 G 在实李群 U(E) 中是闭的，因而是紧实李群（LIE, III, 第 8 节第 2 号定理 2）。根据假设，G 是无限的，故其维数 \(\geqslant 1\)。令 g 为其李代数；它非零。

证明 \(a\)。假设导出群 D(G) 是无限的。有 G-不变分解 \(\operatorname{End}(\mathrm{E}) = \mathbf{C}1_{\mathrm{E}} \oplus \mathfrak{sl}(\mathrm{E})\)。G 在 \(\mathbf{C}1_{\mathrm{E}}\) 和 \(\mathfrak{sl}(\mathrm{E})\) 上的表示不同构，因为 \(\dim \mathfrak{sl}(\mathrm{E}) \geqslant 2\)。由引理 5，\(\mathrm{M}_4(\pi) \geqslant 2\)，且等号成立当且仅当 G 在 \(\mathfrak{sl}(\mathrm{E})\) 中的表示不可约。现在，G 的导出群包含于 \(\mathbf{SL}(\mathrm{E})\)，所以 D(G) 的李代数的复化可识别为 \(\mathfrak{sl}(\mathrm{E})\) 的一个子空间，而后者是 G 在 \(\mathfrak{sl}(\mathrm{E})\) 上的表示的一个子表示（引理 7）。由于 D(G) 无限，该子表示非零。因此，\(\mathrm{M}_4(\pi) = 2\) 当且仅当 \((\mathcal{D}\mathfrak{g})_{(\mathbf{C})} = \mathfrak{sl}(\mathrm{E})\)。由于 D(G) 包含于 \(\mathbf{SU}(\mathrm{E})\)，且 \(\mathbf{SU}(\mathrm{E}) = \mathrm{D}(\mathbf{SU}(\mathrm{E}))\)（引理 8），这个条件等价于 \(\mathbf{SU}(\mathrm{E}) \subset \mathbf{G}\)。

为证明断言 \(b\)，沿用 LIE, VIII, 第 13 节第 2--4 号的记号。

设 \(\pi\) 为辛型且 \(\dim(\mathrm{E}) \geqslant 3\)。令 \(q\) 为 E 上非零的 G-不变交错双线性形式；它是分离的（LIE, IX, 第 103 页附录 2 命题 3），且 G 是辛群 \(\mathbf{Sp}(q)\) 的紧子群。记 \(q^{*} \in \wedge^{2}\mathrm{E}\) 为交错形式 \(q\) 所对应的元素。G 的复化李代数 \(\mathfrak{g}_{(\mathbf{C})}\) 包含于 \(\mathfrak{sp}(q)\)。现在，由于 \(\dim(\mathrm{E}) \geqslant 3\)，\(\mathfrak{sp}(q)\) 在 \(\mathrm{E} \otimes \mathrm{E}\) 中的表示具有如下直和分解

\[
\mathrm{E} \otimes \mathrm{E} = \mathsf {S} ^ {2} (\mathrm{E}) \oplus \wedge^ {2} \mathrm{E} = \mathsf {S} ^ {2} (\mathrm{E}) \oplus \mathrm{E} _ {2} \oplus \mathbf {C} q ^ {*},
\]

其中 \(E_{2}\) 是 \(\mathfrak{sp}(q)\) 的第二基本表示（LIE, VIII, 第 202 页，第 13 节第 3 号 (IV)）。表示 \(E_{2}\) 和 \(C q^{*}\) 是 \(\mathfrak{sp}(q)\) 的不可约表示（上引处），而表示 \(S^{2}(E)\) 非零。

因此由引理 5，\(\mathrm{M}_{4}(\pi) \geqslant 3\)，且等号成立当且仅当 G 在 \(\mathsf{S}^{2}(\mathrm{E})\)、\(E_{2}\) 和 \(Cq^{*}\) 中的表示都不可约且两两不同构。

设 \(M_{4}(\pi)=3\)。由于表示 \(S^{2}(E)\) 可识别为 \(\mathfrak{sp}(q)\) 的伴随表示（引理 6），它包含 G 的伴随表示的复化作为子表示（引理 7）。因此有 \(\mathfrak{g}_{(\mathbf{C})}=\mathfrak{sp}(q)\)。

反之，设 \(\mathfrak{g}_{(\mathbf{C})} = \mathfrak{sp}(q)\)。\(\mathfrak{sp}(q)\) 的伴随表示和表示 \(\mathrm{E}_2\)，后者属于 \(\mathfrak{sp}(q)\)，都不可约，其维数分别为 \(n(n + 1)/2\) 和 \(n(n - 1)/2 - 1\)（LIE, VIII, 第 202 页，第 13 节第 3 号 (IV)）；这两个维数不同且均 \(\geqslant 2\)，因为 \(n \geqslant 3\)，从而 \(\mathrm{M}_4(\pi) = 3\)。

最后，设 \(\pi\) 为实型，且 \(\dim(\mathrm{E}) \geqslant 3\) 以及 \(\dim(\mathrm{E}) \neq 4\)。令 \(q\) 为 \(\mathrm{E}\) 上的非零对称双线性形式，且在 G 下不变；它是分离的（LIE, IX, 第 103 页附录 2 命题 3），并且 \(\mathrm{G}\) 是正交群 \(\mathbf{O}(q)\) 的紧子群。记 \(q^{*} \in \mathbb{S}^{2}(\mathrm{E})\) 为与 \(q\) 对应的元素。

G 的复化李代数 \(\mathfrak{g}_{(\mathbf{C})}\) 包含于 \(\mathfrak{so}(q)\)。\(\mathfrak{so}(q)\) 在 E \(\otimes\) E 中的表示具有如下直和分解

\[
\mathrm{E} \otimes \mathrm{E} = \wedge^ {2} \mathrm{E} \oplus \mathsf {S} ^ {2} (\mathrm{E}) = \wedge^ {2} \mathrm{E} \oplus \mathsf {S} _ {0} ^ {2} (\mathrm{E}) \oplus \mathbf {C} q ^ {*},
\]

其中 \(\mathsf{S}_{0}^{2}(\mathbf{E})\) 是 \(q^{*}\) 在 \(\mathsf{S}^{2}(\mathbf{E})\) 中的正交补。这些表示的维数至少为 2，因为 \(\dim(\mathbf{E})\geqslant3\)。由引理 5，\(\mathrm{M}_{4}(\pi)\geqslant3\)，且等号成立当且仅当 G 在 \(\wedge^{2}\mathrm{E}\) 和 \(\mathsf{S}_{0}^{2}(\mathbf{E})\) 中的表示不可约且不同构。

设 \(\mathrm{M}_{4}(\pi)=3\)。由于表示 \(\wedge^{2}E\) 可识别为 \(\mathfrak{so}(q)\) 的伴随表示（引理 6），它包含 G 的伴随表示的复化作为子表示（引理 7）。因此条件 \(\mathrm{M}_{4}(\pi)=3\) 蕴含 \(\mathfrak{g}_{(\mathbf{C})}=\mathfrak{so}(q)\)。

反之，设 \(\mathfrak{g}_{(\mathbf{C})} = \mathfrak{so}(q)\)。\(\mathfrak{so}(q)\) 的伴随表示是不可约的，因为 \(\dim(\mathrm{E}) \neq 4\)（LIE, VIII, 第 13 节，第 193 页 (I) 及第 206 页 (I)），且维数为 \(n(n-1)/2\)。表示 \(\mathsf{S}_0^2(\mathrm{E})\) 是 \(\mathfrak{so}(q)\) 的表示，其最高权为 \(2\varpi_1\)。将其维数与由 Weyl 公式计算的 \(\mathfrak{so}(q)\) 的最高权为 \(2\varpi_1\) 的不可约表示的维数比较（参见 LIE, VIII, 第 9 节第 2 号定理 2 以及 LIE, VI, 图版 II 和 IV），可验证 \(\mathsf{S}_0^2(\mathrm{E})\) 是不可约的。因此得到 \(\mathrm{M}_4(\pi) = 3\)，当 \(\mathfrak{g}_{(\mathbf{C})} = \mathfrak{so}(q)\) 时。

注记。--- 1) “G 是无限的”这一条件对于定理 3 是必要的（V, 第 507 页习题 9）。

\begin{enumerate}
\def\labelenumi{\arabic{enumi})}
\setcounter{enumi}{1}
\item
  令 \(n \in \mathbf{N}\)。对于紧群的每个酉表示 \(\varrho\)，若它作用在维数为 \(n\) 的希尔伯特空间中，则有 \(\mathrm{M}_{4}(\varrho) \leqslant n^{2}\)；等号可能成立（V, 第 508 页习题 15）。
\item
  可以证明（参见 R. GURALNICK 和 P.H. TIEP，《Decomposition of small tensor powers and Larsen's conjecture》，Representation Theory 9 (2005), 138--208）：若假设 E 的维数 \(\geqslant 7\)，并将假设 \(M_{4}(\pi) = 2\)（分别为 \(M_{4}(\pi) = 3\)）替换为 \(M_{8}(\pi) = 24\)（分别为 \(M_{8}(\pi) = 105\)），则可以省略 G 为无限这一条件。
\end{enumerate}

\specialchapter{习题}

\exercisesection{}

\begin{enumerate}
\def\labelenumi{\arabic{enumi})}
\item
  给出一个群 G、一个域 K 以及 G 的一个忠实表示的例子，使得相应的 \(K[G]\)-模并不忠实。
\item
  令 \(\pi\) 为 G 在希尔伯特空间 E 上的酉表示，且 \(x \in E\) 的范数为 1，并满足
\end{enumerate}

\[
\sup _ {g \in \mathrm{G}} \| \pi (g) x - x \| \leqslant 1.
\]

存在非零的 G-不变向量 \(y \in E\)。

\begin{enumerate}
\def\labelenumi{\arabic{enumi})}
\setcounter{enumi}{2}
\tightlist
\item
  令 E 为希尔伯特空间。记 \(\mathbf{U}(\mathbf{E})\) 为 E 的酉自同态群，记 \(\mathbf{GL}(\mathbf{E})\) 为 E 的可逆自同态群。
\end{enumerate}

\begin{enumerate}
\def\labelenumi{\alph{enumi})}
\item
  群 \(\mathbf{U}(\mathrm{E})\) 赋予由 \(\mathcal{L}(\mathrm{E})\) 上逐点收敛拓扑诱导的拓扑后是拓扑群，记为 \(\mathbf{U}(\mathrm{E})_{s}\)；若 E 是无限维的，则 \(\mathbf{GL}(\mathrm{E})\) 赋予由逐点收敛拓扑诱导的拓扑后不是拓扑群。
\item
  \(\mathbf{U}(\mathrm{E})_{s}\) 的拓扑与由紧收敛拓扑诱导的拓扑一致。
\item
  若 E 是可数生成的，则群 \(\mathbf{U}(\mathrm{E})_{s}\) 可度量化且完备。
\item
  令 L 为由 \(u \in \mathcal{L}(\mathrm{E})\) 且满足 \(u^{*} = -u\) 的元素所成的空间；它是实李群 G = U(E) 的李代数，其中 G 赋予由 \(\mathcal{L}(\mathrm{E})\) 的巴拿赫空间拓扑诱导的拓扑（LIE, III, 第 146 页推论 2）。从 L 到 \(\mathbf{U}(\mathrm{E})_s\) 的指数映射是连续的。
\item
  若 E 是无限维的，则 \(\mathbf{U}(\mathrm{E})_{s}\) 不是型 \(\{L\}\) 的实李群。（考虑 \(\mathbf{U}(\mathrm{E})_{s}\) 中由 E 的某个固定标准正交基下的对角酉自同态组成的子群 K，并注意 K 是紧的。）
\item
  指数映射不是从 L 到 \(\mathbf{U}(\mathbf{E})_{s}\) 的局部同胚。（考虑它在某个固定标准正交基下的对角自同态上的限制。）
\end{enumerate}

\begin{enumerate}
\def\labelenumi{\arabic{enumi})}
\setcounter{enumi}{3}
\tightlist
\item
  令 G 为拓扑群。若存在一个线性泛函 \(\lambda\colon\mathcal{C}_{b}(\mathrm{G};\mathbf{R})\to\mathbf{R}\)，使得
\end{enumerate}

\begin{enumerate}
\def\labelenumi{(\roman{enumi})}
\item
  若 \(f \geqslant 0\)，则 \(\lambda(f) \geqslant 0\)。
\item
  有 \(\lambda(1)=1\)。
\item
  对每个 \(g \in \mathbf{G}\) 及每个 \(f \in \mathcal{C}_b(\mathbf{G};\mathbf{R})\)，有 \(\lambda(\boldsymbol{\gamma}(g)f) = \lambda(\boldsymbol{\delta}(g)f)\)，其中
\end{enumerate}

\[
\boldsymbol {\gamma} (g) f (x) = f \left(g ^ {- 1} x\right), \quad \boldsymbol {\delta} (g) f (x) = f (x g).
\]

若 G 是离散群，此定义与 EVT, IV, 第 73 页习题 4 中的定义一致。

\begin{enumerate}
\def\labelenumi{\alph{enumi})}
\tightlist
\item
  拓扑群 G 可和，当且仅当对于每个整数 \(n \in N\)、每个有限族 \((f_i)_{1 \leqslant i \leqslant n}\)（属于 \(\mathcal{C}_b(\mathbf{G}; \mathbf{R})\)）、G 中的两个族 \((g_i)_{1 \leqslant i \leqslant n}\) 和 \((h_i)_{1 \leqslant i \leqslant n}\) 以及每个实数 a，条件
\end{enumerate}

\[
\sum_ {i = 1} ^ {n} (f _ {i} (x) - f _ {i} (g _ {i} x h _ {i})) \geqslant a
\]

对所有 \(x \in G\) 成立都蕴含 \(a \leqslant 0\)。（为证明该条件充分，考虑形如

\[
x \mapsto \sum_ {i = 1} ^ {n} (f _ {i} (x) - f _ {i} (g _ {i} x h _ {i}))
\]

的函数所成的集合，并应用 Hahn--Banach 定理。）

\begin{enumerate}
\def\labelenumi{\alph{enumi})}
\setcounter{enumi}{1}
\item
  每个紧拓扑群都是可和的，每个阿贝尔拓扑群也都是可和的。
\item
  令 G 为非交换自由群。离散群 G 不是可和的。
\item
  对每个希尔伯特空间 E 及 G 在 E 中的每个连续有界表示，E 上都存在一个希尔伯特范数 q，使得 G 在赋予 q 的希尔伯特空间 E 中的表示是酉的。
\end{enumerate}

\begin{enumerate}
\def\labelenumi{\arabic{enumi})}
\setcounter{enumi}{4}
\tightlist
\item
  令 G 为拓扑群，H 为 G 的开子群。记 p 为典范投影 \(G \rightarrow G/H\)。
\end{enumerate}

\begin{enumerate}
\def\labelenumi{\alph{enumi})}
\item
  令 \(\sigma\colon\mathrm{G/H}\to\mathrm{G}\) 为满足 \(p\circ\sigma=\mathrm{Id}_{\mathrm{G/H}}\) 的映射。记 \(\beta\colon\mathrm{G}\times\mathrm{G/H}\to\mathrm{H}\) 为由 \(\beta(g,x)=\sigma(gx)^{-1}g\) 定义的映射。映射 \(\beta\) 是连续的。
\item
  令 \(\pi\) 为 H 在希尔伯特空间 E 中的有界表示。记 F 为希尔伯特空间 \(\ell^2 (\mathrm{G / H};\mathrm{E})\)。公式
\end{enumerate}

\[
\varrho (g) f (x) = \pi (\beta (g ^ {- 1}, x) ^ {- 1}) f (g ^ {- 1} x)
\]

定义了 G 在 F 中的连续有界表示。

\begin{enumerate}
\def\labelenumi{\alph{enumi})}
\setcounter{enumi}{2}
\tightlist
\item
  若 F 上存在一个希尔伯特范数 q，使得表示 \(\varrho\) 是 G 在 \((F,q)\) 中的酉表示，则 E 上存在一个希尔伯特范数 \(\widetilde{q}\)，使得 \(\pi\) 是 H 在 \((E,\widetilde{q})\) 中的酉表示。
\end{enumerate}

\begin{enumerate}
\def\labelenumi{\arabic{enumi})}
\setcounter{enumi}{5}
\tightlist
\item
  \begin{enumerate}
  \def\labelenumii{\alph{enumii})}
  \tightlist
  \item
    令 E 为复巴拿赫空间，令 \(u \in \mathcal{L}(\mathrm{E})\)。若存在 \(n \in N\)，使得 \(n \geqslant 1\) 且 \(u^{n} = 1_{E}\)，则 u 的谱包含于 n 次单位根群，并由 u 的特征值组成。b) 令 G 为交换拓扑群，且 G 的每个元素都有有限阶。G 在巴拿赫空间中的每个连续不可约表示维数均为 1，且取值于 \(U \subset C^{*}\)——模为 1 的复数所成的群。
  \end{enumerate}
\end{enumerate}

\begin{enumerate}
\def\labelenumi{\alph{enumi})}
\setcounter{enumi}{2}
\tightlist
\item
  令 G 为交换拓扑群，且 G 的有限阶元素集在 G 中稠密。G 在巴拿赫空间中的每个连续不可约表示维数均为 1，且取值于 \(U \subset C^{*}\)——模为 1 的复数所成的群。
\end{enumerate}

\begin{enumerate}
\def\labelenumi{\arabic{enumi})}
\setcounter{enumi}{6}
\tightlist
\item
  对每个整数 \(n \geqslant 1\)，记 \(s_{n}\) 为由下式定义的函数 \(Z \rightarrow R\)
\end{enumerate}

\[
s _ {n} (k) = \left\{ \begin{array}{l l} (1 - k) ^ {n (1 - k)} & \text{若 } k \leqslant - 1 \\ 1 & \text{若 } k = 0 \\ (k + 1) ^ {- (k + 1) / n} & \text{若 } k \geqslant 1. \end{array} \right.
\]

记 S 为函数 \(s_n\) 所成的集合。

\begin{enumerate}
\def\labelenumi{\alph{enumi})}
\item
  对每个整数 \(n \geqslant 1\) 及每个 \((k_{1}, k_{2}) \in \mathbf{Z}^{2}\)，有 \(s_{4n}(k_{1})s_{4n}(k_{2}) \geqslant s_{n}(k_{1} + k_{2})\)。
\item
  对于 \(f \in \mathbf{C}(\mathrm{T})\)，记 \((a_k(f))_{k \in \mathbf{Z}}\) 为 \(f\) 在 0 处的 Laurent 展开的系数。对于每个 \(s \in \mathrm{S}\) 和每个 \(f \in \mathbf{C}(\mathrm{T})\)，级数
\end{enumerate}

\[
\sum_ {k \in \mathbf {Z}} | a _ {k} (f) | s (k)
\]

收敛，并且集合

\[
\mathcal {V} (s, \varepsilon) = \{f \in \mathbf {C} (\mathrm{T}) | \sum_ {k \in \mathbf {Z}} | a _ {k} (f) | s (k) <   \varepsilon \}
\]

对 \(s \in S\) 和 \(\varepsilon > 0\) 定义的这些集合，构成 \(\mathbf{C}(T)\) 上某个局部凸拓扑的 0 的开邻域基本系。乘法 \(\mathbf{C}(T) \times \mathbf{C}(T) \to \mathbf{C}(T)\) 关于此拓扑是连续的。

\begin{enumerate}
\def\labelenumi{\alph{enumi})}
\setcounter{enumi}{2}
\tightlist
\item
  上一问题中定义的拓扑可度量化但不完备。
\end{enumerate}

\(d)\) 对于 \(f \in \mathbf{C}(\mathrm{T})^*\)，记 \(\varrho(f)\) 为映射 \(g \mapsto fg\)，从 \(\mathbf{C}(\mathrm{T})\) 到 \(\mathbf{C}(\mathrm{T})\)。映射 \(\varrho\) 是 \(\mathbf{C}(\mathrm{T})^*\) 在局部凸空间 \(\mathbf{C}(\mathrm{T})\) 中的连续表示。

\begin{enumerate}
\def\labelenumi{\alph{enumi})}
\setcounter{enumi}{4}
\tightlist
\item
  表示 \(\varrho\) 是不可约的。
\end{enumerate}

\exercisesection{}

在本节习题中，除非另有说明，G 表示赋有左哈尔测度的局部紧拓扑群，该测度记为 \(\mu\)。

\begin{enumerate}
\def\labelenumi{\arabic{enumi})}
\item
  令 \(\pi\) 为 G 在复希尔伯特空间 E 上的酉表示。表示 \(\pi\) 不可约，当且仅当 \(\mathcal{H}(\mathrm{G})\) 在 \(\mathcal{L}(\mathrm{E})\) 中的像关于逐点收敛拓扑是稠密的。（使用 IV, 第 326 页习题 24。）
\item
  令 H 为 G 的闭子群。令 \(\nu\) 为 G/H 上的拟不变测度（INT, VII, 第 56 页，第 2 节第 5 号定理 2），并令 \(\varrho\) 为与之关联的、取值于 \(\mathbf{R}_{+}^{*}\) 的 G 上函数。由下式定义的映射 \(\varrho(g)\)
\end{enumerate}

\[
\varrho (g) f (x) = \left(\frac {\varrho (g ^ {- 1} x)}{\varrho (x)}\right) ^ {1 / 2} f (g ^ {- 1} x)
\]

定义了 G 在 \(\mathcal{L}^{2}(\mathrm{G / H},\nu)\) 上的连续线性表示，并通过取商定义了 G 在 \(\mathrm{L}^2 (\mathrm{G / H},\nu)\) 上的酉表示，称为 \(\mathrm{G / H}\) 的拟正则表示。

\begin{enumerate}
\def\labelenumi{\arabic{enumi})}
\setcounter{enumi}{2}
\tightlist
\item
  令 G = R，令 H 为赋予离散拓扑并配备计数测度的群 R；该测度是哈尔测度。
\end{enumerate}

\begin{enumerate}
\def\labelenumi{\alph{enumi})}
\item
  对每个 \(t \in G\) 及每个 \(f \in L^{2}(H)\)，映射 \(\varrho(t)f: x \mapsto f(x + t)\) 属于 \(L^{2}(H)\)；映射 \(\varrho(t)\) 是 \(L^{2}(H)\) 上的酉映射。b) 对所有 \(f,g \in L^{2}(H)\)，由 \(t \mapsto \langle f | \varrho(t)g \rangle\) 定义的从 G 到 C 的映射是可测的。
\item
  令 \(f\) 和 \(g\) 属于 \(\mathrm{L}^2(\mathrm{H})\)。若存在 \(t \in \mathbf{R}\) 使得 \(\langle f | \varrho(t)g \rangle\) 非零，则映射 \(t \mapsto \langle f | \varrho(t)g \rangle\) 从 \(\mathbf{G}\) 到 \(\mathbf{C}\) 不连续。（特别地，E 可数生成这一假设不能从 INT, VIII, 第 171 页第 4 节第 7 号推论 2 中省略。）
\end{enumerate}

\(d)\) 存在 \(f \in \mathrm{L}^2(\mathrm{H})\)，使得 \(\mathrm{L}^2(\mathrm{H})\) 中由元素 \(\varrho(t)f\)（其中 \(t \in \mathbf{R}\)）生成的子空间在 \(\mathrm{L}^2(\mathrm{H})\) 中稠密。

\begin{enumerate}
\def\labelenumi{\arabic{enumi})}
\setcounter{enumi}{3}
\item
  令 \(p \geqslant 1\)。G 在 \(\mathrm{L}^{p}(\mathrm{G}, \mu)\) 上的双正则表示的核，是 G 的中心在同态 \(g \mapsto (g, g)\) 下的像。
\item
  令 G 为局部紧单模群。
\end{enumerate}

\begin{enumerate}
\def\labelenumi{\alph{enumi})}
\item
  若存在 G 的有限维平方可积表示，则 G 是紧的。
\item
  若存在 G 的平方可积表示，则 G 的中心是紧的。
\end{enumerate}

\begin{enumerate}
\def\labelenumi{\arabic{enumi})}
\setcounter{enumi}{5}
\tightlist
\item
  假设群 G 是单模的。记 \(\widehat{G}_{d} \subset \widehat{G}\) 为 G 的不可约酉平方可积表示的等价类集。对于 \(\pi \in \widehat{G}_{d}\)，记 \(p_{\pi}\) 为 \(L^{2}(G)\) 上的正交投影，其像是 G 的右正则表示的 \(\pi\)-同型分量。
\end{enumerate}

\(a)\) 令 \(f \in \mathcal{L}^2(\mathrm{G})\) 且 \(\pi \in \widehat{\mathrm{G}}_d\)。记 E 为 \(\pi\) 的表示空间。映射 \(g \mapsto f(g)\pi(g)\) 属于 \(\mathrm{L}^1(\mathrm{G};\mathcal{L}(\mathrm{E}))\)；自同态

\[
v = \int_ {\mathrm{G}} f (g) \pi (g) d \mu (g)
\]

是希尔伯特–施密特自同态，记为 \(v = \pi(f)\)。

\begin{enumerate}
\def\labelenumi{\alph{enumi})}
\setcounter{enumi}{1}
\item
  令 \(f \in \mathcal{L}^2(\mathrm{G})\)。有 \(\| \pi(\overline{f}) \|_2 = \| p_\pi(f) \|\)。
\item
  令 \(L_{d}^{2}(G)\) 为 \(L^{2}(G)\) 的闭子空间，它由 \(\pi\)-同型分量生成，这些分量属于 \(L^{2}(G)\)，且 \(\pi \in \widehat{G}_{d}\)。它是各正交投影 \(p_{\pi}\) 的像（对 \(\pi \in \widehat{G}_{d}\)）的希尔伯特直和。对于每个 \(f \in L_{d}^{2}(G)\)，有
\end{enumerate}

\[
\| f \| ^ {2} = \sum_ {\pi \in \widehat {\mathrm{G}} _ {d}} c (\pi) \| \pi (f) \| _ {2} ^ {2}
\]

其中 \(c(\pi)\) 是 \(\pi\) 相对于 \(\{e\}\) 的形式次数。

\begin{enumerate}
\def\labelenumi{\arabic{enumi})}
\setcounter{enumi}{6}
\tightlist
\item
  令 \(n \in N\)。确定 \(R^{n}\) 的酉自守表示（即对每个离散子群 \(\Gamma\)，它们都是 \(\Gamma\)-自守的 \(R^{n}\) 的表示）。
\end{enumerate}

8) 假设 G 是局部紧单模群。令 Γ 为 G 的离散子群，使商空间 X = Γ\textbackslash G 是紧的。令 μ 为 G 上的哈尔测度，β 为 Γ 上的计数测度；记 \(\widetilde{\mu} = \mu/\beta\)。

令 \(\varrho\) 为 V, 第 404 页例 3 中定义的 G 在 \(\mathrm{L}^{2}(\mathrm{X},\widetilde{\mu})\) 上的酉表示。

对于 \(\gamma\in\Gamma\)，记 \(\Gamma_{\gamma}\)（分别为 \(\mathbf{G}_{\gamma}\)）为 \(\gamma\) 在 \(\Gamma\)（分别在 \(\mathbf{G}\)）中的中心化子。\(a)\) 群 \(\mathbf{G}_{\gamma}\) 是单模的。

\begin{enumerate}
\def\labelenumi{\alph{enumi})}
\setcounter{enumi}{1}
\item
  令 \(\varphi\in\mathcal{K}(\mathrm{G})\)。对于每个 \(\Gamma\)-自守酉表示 \(\pi\)，自同态 \(\pi(\varphi)\) 都有有限迹。
\item
  有
\end{enumerate}

\[
\operatorname{Tr} (\varrho (\varphi)) = \sum_ {\pi \in \widehat {\mathrm{G}}} m _ {\pi} (\varrho) \operatorname{Tr} (\pi (\varphi))
\]

其中 \(m_{\pi}(\varrho)\) 是 \(\pi\) 在 \(\varrho\) 中的重数（“迹公式的谱侧”）。

\begin{enumerate}
\def\labelenumi{\alph{enumi})}
\setcounter{enumi}{3}
\tightlist
\item
  令 \(\Gamma^{\sharp}\) 为 \(\Gamma\) 的共轭类集。对于 \(\gamma \in \Gamma^{\sharp}\)，记 \(\beta_{\gamma}\) 为 \(\Gamma_{\gamma}\) 上的计数测度。有
\end{enumerate}

\[
\mathrm{Tr} (\varrho (\varphi)) = \sum_ {\gamma \in \Gamma^ {\sharp}} \int_ {\Gamma_ {\gamma} \backslash G} \varphi (x ^ {- 1} \gamma x) d (\mu / \beta_ {\gamma}) (x).
\]

\begin{enumerate}
\def\labelenumi{\alph{enumi})}
\setcounter{enumi}{4}
\tightlist
\item
  固定 \(\mu_{\gamma}\) 在 \(G_{\gamma}\) 上的哈尔测度，对所有 \(\gamma \in \Gamma\)。有
\end{enumerate}

\[
\operatorname{Tr} (\varrho (\varphi)) = \sum_ {\gamma \in \Gamma^ {\sharp}} \int_ {G _ {\gamma} \backslash G} \int_ {\Gamma_ {\gamma} \backslash G _ {\gamma}} \varphi (x ^ {- 1} y ^ {- 1} \gamma y x) d (\mu_ {\gamma} / \beta_ {\gamma}) (y) d (\mu / \mu_ {\gamma}) (x).
\]

\(f)\) 对每个 \(\gamma \in \Gamma\)，测度 \(\mu_{\gamma} / \beta_{\gamma}\) 在 \(\Gamma_{\gamma} \backslash G_{\gamma}\) 上是有限的。\(g)\) 有

\[
\mathrm{Tr} (\varrho (\varphi)) = \sum_ {\gamma \in \Gamma^ {\sharp}} (\mu_ {\gamma} / \beta_ {\gamma}) (\Gamma_ {\gamma} \backslash \mathrm{G} _ {\gamma}) \int_ {\mathrm{G} _ {\gamma} \backslash \mathrm{G}} \varphi (x ^ {- 1} \gamma x) d (\mu / \mu_ {\gamma}) (x)
\]

（“迹公式的几何侧”；将此表达式与 c) 中所得表达式相等联系起来的公式，称为 Γ-自守表示的“迹公式”。）

\begin{enumerate}
\def\labelenumi{\alph{enumi})}
\setcounter{enumi}{7}
\tightlist
\item
  置 G = R 且 \(\Gamma = Z\)。验证此时迹公式等价于泊松公式（参见 II, 第 239 页）。
\end{enumerate}

\begin{enumerate}
\def\labelenumi{\arabic{enumi})}
\setcounter{enumi}{8}
\item
  假设 G 是无紧因子的完美实半单李群。G 没有非平凡的有限维酉表示。
\item
  \begin{enumerate}
  \def\labelenumii{\alph{enumii})}
  \tightlist
  \item
    G 在 \(L^{2}(G)\) 上的左正则表示的每个矩阵系数都属于 \(\mathcal{C}_{0}(G)\)。
  \end{enumerate}
\end{enumerate}

\begin{enumerate}
\def\labelenumi{\alph{enumi})}
\setcounter{enumi}{1}
\tightlist
\item
  若 G 非紧，则 G 在 L²(G) 上的左正则表示不含非零有限维子表示。
\end{enumerate}

\begin{enumerate}
\def\labelenumi{\arabic{enumi})}
\setcounter{enumi}{10}
\tightlist
\item
  令 E 为希尔伯特空间。
\end{enumerate}

\begin{enumerate}
\def\labelenumi{\alph{enumi})}
\item
  令 \(\varrho\) 为群 \(\mathbf{R}\) 在 E 中的酉表示。令 D 为 E 的稠密子空间，在 \(\varrho\) 下稳定，并且对每个 \(x\in\mathrm{D}\)，映射 \(\psi_{x}:t\mapsto\varrho(t)x\) 在 R 上可微。令 u 为定义域为 D 且满足 \(u(x) = i^{-1}\psi_x'(0)\)（对所有 \(x \in D\)）的部分算子。则 u 本质自伴，且 \(\overline{u}\) 是 \(\varrho\) 的无穷小生成元。
\item
  令 u 为 E 上的自伴部分算子，D 为 \(\text{dom}(u)\) 的稠密子空间。假设对每个 \(t \in R\)，有 \(\exp(itu)(D) \subset D\)。则空间 D 在 \(E_{u}\) 中稠密。
\item
  令 u 为 E 上的自伴部分算子。令 F 为希尔伯特空间，v 为 F 上的自伴部分算子。部分算子 \(u \otimes 1_{F} + 1_{E} \otimes E\) 在 \(E \otimes F\) 上定义了 \(E \widehat{\otimes}_{2} F\) 上的一个本质自伴部分算子。
\end{enumerate}

\begin{enumerate}
\def\labelenumi{\arabic{enumi})}
\setcounter{enumi}{11}
\tightlist
\item
  令 G 为实李群，令 \(\varrho\) 为 G 在希尔伯特空间 E 中的酉表示。若元素 \(x \in \mathrm{E}\) 为 \(\mathrm{C}^{\infty}\) 类，亦即从 G 到 E 的映射 \(g \mapsto \varrho(g)x\) 为 \(\mathrm{C}^{\infty}\) 类，则称其为类向量。a) 令 \(\pi\) 为 G 在 F 中的酉表示，且 \(u: \mathrm{E} \to \mathrm{F}\) 为 G-态射。对于每个 \(x \in \mathrm{E}\) 的 \(\mathrm{C}^{\infty}\) 类向量，向量 \(u(x)\) 在 F 中也是 \(\mathrm{C}^{\infty}\) 类的。
\end{enumerate}

记 \(\mathrm{C}^{\infty}(\varrho)\) 为 E 中 \(C^{\infty}\) 类向量的空间。由问题 a)，每个从 \(\varrho\) 到 \(\pi\) 的 G-态射通过取子空间定义了从 \(\mathrm{C}^{\infty}(\varrho)\) 到 \(\mathrm{C}^{\infty}(\pi)\) 的线性映射。

\begin{enumerate}
\def\labelenumi{\alph{enumi})}
\setcounter{enumi}{1}
\tightlist
\item
  空间 \(\mathrm{C}^{\infty}(\varrho)\) 在 E 中稠密。
\end{enumerate}

\begin{enumerate}
\def\labelenumi{\arabic{enumi})}
\setcounter{enumi}{12}
\tightlist
\item
  令 \(\pi\) 为 G 在复希尔伯特空间 E 中的不可约酉表示。令 \(\varrho\) 为 G 在复希尔伯特空间 F 中的酉表示。假设存在一个从 E 到 F 的闭稠定部分算子 \(u\)，其定义域是 E 的 G-稳定子空间，并且满足 \(u \circ \pi(g) = \varrho(g) \circ u\) 对所有 \(g \in \mathbf{G}\) 成立。
\end{enumerate}

部分算子 \(u\) 属于 \(\mathcal{L}(\mathrm{E};\mathrm{F})\)，并且存在 \(\lambda > 0\)，使得 \(\lambda u\) 是等距的。

\begin{enumerate}
\def\labelenumi{\arabic{enumi})}
\setcounter{enumi}{13}
\tightlist
\item
  令 Z 为 G 的中心 C 的子群，且 C/Z 是紧的；令 \(\nu\) 为 G/Z 上的哈尔测度。令 \(\pi\) 为 G 在复希尔伯特空间 E 中的不可约酉表示。假设 \(\pi\) 模中心平方可积。令 \(u \in \mathcal{L}(\mathrm{E})\) 为核型自同态（IV, 第 169 页定义 4）。
\end{enumerate}

\begin{enumerate}
\def\labelenumi{\alph{enumi})}
\tightlist
\item
  对 E 中的每个 \(x\) 和 \(y\)，函数
\end{enumerate}

\[
g \mapsto \langle x | \pi (g) u \pi (g ^ {- 1}) y \rangle
\]

属于 \(\mathcal{F}(\mathrm{G / Z})\)。

\begin{enumerate}
\def\labelenumi{\alph{enumi})}
\setcounter{enumi}{1}
\tightlist
\item
  对 E 中的每个 \(x\) 和 \(y\)，有
\end{enumerate}

\[
\int_ {\mathrm{G} / \mathrm{Z}} \langle x \mid \pi (g) u \pi (g ^ {- 1}) y \rangle d \nu (g) = \frac {1}{c _ {\mathrm{Z}} (\pi)} \langle x \mid y \rangle \operatorname{Tr} (u).
\]

\begin{enumerate}
\def\labelenumi{\arabic{enumi})}
\setcounter{enumi}{14}
\tightlist
\item
  令 Z 为 G 的中心 C 的子群，且 C/Z 是紧的；令 \(\nu\) 为 G/Z 上的哈尔测度。令 \(\pi\) 为 G 在复希尔伯特空间 E 中的不可约酉表示。记 \(\chi\) 为 \(\pi\) 的中心特征标在 Z 上的限制。
\end{enumerate}

记 \(f_{x,y}\) 为矩阵系数 \(g \mapsto \langle x | \pi(g)y \rangle\)，其中 \((x, y) \in \mathrm{E}^{2}\)。令 \(p \geqslant 1\) 为实数，令 \(q\) 为 p 的共轭指数。

\begin{enumerate}
\def\labelenumi{\alph{enumi})}
\item
  令 \(y \in \mathrm{E}\)，使得矩阵系数 \(f_{x,y}\) 属于 \(\mathcal{L}_{\chi}^{p}(\mathrm{G})\)。映射 \(x \mapsto f_{x,y}\) 是从 \(\pi\) 到 \(\boldsymbol{\delta}_{\mathrm{G},\chi}^{(p)}\) 的连续 G-态射。
\item
  令 \(\varrho\) 为 G 在复希尔伯特空间 F 中的不可约酉表示，其中心特征标限制所得的 \(\chi\) 定义在 \(Z\) 上。假设 \(p > 1\)，且 \(\pi\) 的矩阵系数属于 \(\mathcal{L}_{\chi}^{p}(\mathrm{G})\)，\(\varrho\) 的矩阵系数属于 \(\mathcal{L}_{\chi}^{q}(\mathrm{G})\)。令 \(x_0 \in \mathrm{E}\) 且 \(y_0 \in \mathrm{F}\)。对于每个 \(x \in \mathrm{E}\)，存在从 F 到 C 的连续线性泛函 \(\ell_x\)，使得
\end{enumerate}

\[
\ell_ {x} (y) = \int_ {\mathrm{G/Z}} \overline {{\langle x _ {0} | \pi (g) x \rangle}} \langle y _ {0} | \varrho (g) y \rangle d \nu (g)
\]

对所有 \(y \in \mathbf{F}\) 都成立。

\begin{enumerate}
\def\labelenumi{\alph{enumi})}
\setcounter{enumi}{2}
\item
  存在从 E 到 F 的连续线性映射 u，使得有 \(\ell_{x}(y)=\langle u(x)\mid y\rangle\)，对所有 \((x,y)\in\mathrm{E}\times\mathrm{F}\) 成立。
\item
  有
\end{enumerate}

\[
\int_ {\mathrm{G} / \mathrm{Z}} \overline {{\langle x | \pi (g) x ^ {\prime} \rangle}} \langle y | \varrho (g) y ^ {\prime} \rangle d \nu (g) = 0
\]

对 \((x,x^{\prime},y,y^{\prime})\in \mathrm{E}^{2}\times \mathrm{F}^{2}\) 都成立。

\begin{enumerate}
\def\labelenumi{\alph{enumi})}
\setcounter{enumi}{4}
\tightlist
\item
  假设 p = 1。\(\pi\) 的矩阵系数与 \(\varrho\) 的矩阵系数正交。
\end{enumerate}

\begin{enumerate}
\def\labelenumi{\arabic{enumi})}
\setcounter{enumi}{15}
\tightlist
\item
  令 \(\mathbf{G} = \mathbf{S}\mathbf{L}(2, \mathbf{R})\)，令 H 为 C 中由满足 \(\mathcal{I}(z) > 0\) 的复数 z 组成的开子集。令 \(\mu\) 为 C 上的勒贝格测度。对于 G 中的每个元素
\end{enumerate}

\[
g = \left( \begin{array}{c c} a & b \\ c & d \end{array} \right)
\]

及每个 \(z \in \mathbf{H}\)，置 \(j(g, z) = cz + d\)。

\begin{enumerate}
\def\labelenumi{\alph{enumi})}
\tightlist
\item
  群 G 通过下式在 H 上连续地左作用
\end{enumerate}

\[
\left( \begin{array}{c c} a & b \\ c & d \end{array} \right) \cdot z = \frac {a z + b}{c z + d}.
\]

测度 \(\nu = y^{-2} \cdot \mu\) 是 G-不变的。

\begin{enumerate}
\def\labelenumi{\alph{enumi})}
\setcounter{enumi}{1}
\tightlist
\item
  令 n 为满足 \(n \geqslant 2\) 的整数。记 \(E_{n}\) 为全纯函数 \(f: H \to C\) 所成的向量空间，其中函数 \(z \mapsto |f(z)|^{2} \mathcal{I}(z)^{n}\) 属于 \(\mathcal{L}^{2}(\mathbf{H}, \nu)\)。空间 \(E_{n}\) 在赋予半双线性形式
\end{enumerate}

\[
\langle f \mid g \rangle = \int_ {\mathbf {H}} \overline {{f (z)}} g (z) \mathcal {I} (z) ^ {n} d \nu (z)
\]

后，是可数型且无限维的希尔伯特空间。

\begin{enumerate}
\def\labelenumi{\alph{enumi})}
\setcounter{enumi}{2}
\tightlist
\item
  对每个 \(g \in G\) 及每个 \(f \in E_{n}\)，置
\end{enumerate}

\[
\pi_ {n} (g) f (z) = j (g ^ {- 1}, z) ^ {- n} f (g ^ {- 1} \cdot z).
\]

映射 \(\pi_n\) 定义了 G 在 \(\mathrm{E}_n\) 中的酉表示。\(d)\) 由函数 \(f_n\) 定义为 \(f_n(z) = (z + i)^{-n}\)，它属于 \(\mathrm{E}_n\)；它满足 \(\pi_n(k)f_n = \chi_{-n}(k)f_n\) 对所有 \(k \in \mathbf{SO}(\mathbf{R}^2)\) 成立，其中 \(\chi_n\) 是满足下式的特征标

\[
\chi_ {n} \Big (\left( \begin{array}{c c} \cos (\theta) & \sin (\theta) \\ - \sin (\theta) & \cos (\theta) \end{array} \right) \Big) = e ^ {i n \theta}
\]

对所有 \(\theta\in\mathbf{R}\) 都成立。

\begin{enumerate}
\def\labelenumi{\alph{enumi})}
\setcounter{enumi}{4}
\item
  令 F 为 \(E_{n}\) 的子表示。\(E_{n}\) 的正交投影，其像是 \(\chi_{-n}\)-同型分量，属于将 \(\pi_{n}\) 限制到 \(\mathbf{SO}(\mathbf{R}^{2})\) 后的表示，由 \(f \mapsto -(2i)^{n}f(i)f_{n}\) 给出。（使用 V, 第 465 页推论 2 以及柯西积分公式，FVC，即将出版。）
\item
  表示 \(\pi_n\) 是不可约的。
\end{enumerate}

\(g)\) 矩阵系数 \(g \mapsto \langle f_n | \pi_n(g)f_n \rangle\) 非零且属于 \(\mathcal{L}^2(\mathrm{G})\)。

\begin{enumerate}
\def\labelenumi{\alph{enumi})}
\setcounter{enumi}{7}
\tightlist
\item
  表示 \(\pi_{n}\) 是平方可积的，并且存在常数 c > 0，使得 \(\pi_{n}\) 的形式次数为 \(c(n - 1)\)。
\end{enumerate}

\begin{enumerate}
\def\labelenumi{\arabic{enumi})}
\setcounter{enumi}{16}
\item
  令 G 为连通实李群。G 存在忠实有限维酉表示，当且仅当群 G 局部同构于紧群。
\item
  令 X 为赋有有界正测度 \(\mu\) 的局部紧空间。对于每个巴拿赫空间 F，记 S(X, \(\mu\) ; F) 为从 X 到 F 的 \(\mu\)-可测函数按 \(\mu\)-几乎处处相等关系所得的等价类空间，并赋予依测度收敛拓扑（INT, IV, 第 194 页，第 5 节第 11 号）。记 A 为 \(\mathcal{L}(\mathrm{L}_{\mathbf{C}}^{2}(\mathrm{X},\mu))\) 的交换子代数，它由函数 \(f \in \mathcal{L}_{\mathbf{C}}^{\infty}(\mathrm{X},\mu)\) 所对应的乘法自同态组成。
\end{enumerate}

记 \(e\colon \mathbf{R}\to \mathbf{C}\) 为同态 \(t\mapsto \exp (2i\pi t)\)。

令 E 为拓扑向量空间。

\begin{enumerate}
\def\labelenumi{\alph{enumi})}
\tightlist
\item
  令 \(f \in \mathrm{S}(\mathrm{X}, \mu; \mathbf{R})\)，使得 \(|f(x)| \geqslant 1\) 对 \(\mu\)-几乎所有 \(x \in \mathrm{X}\) 都成立。存在实数 \(t \in ]0, 1[\)，使得
\end{enumerate}

\[
\mu (\{x \in \mathrm{X} \mid t f (x) \in [ 1 / 4, 3 / 4 ] + \mathbf {Z} \}) \geqslant \frac {2}{5} \mu (\mathrm{X}).
\]

\begin{enumerate}
\def\labelenumi{\alph{enumi})}
\setcounter{enumi}{1}
\item
  令 u 为从 E 到 \(\mathrm{S}(\mathrm{X},\mu;\mathbf{R})\) 的线性映射。对于每个 \(x\in E\)，记 \(\varrho_{u}(x)\) 为 A 中由函数 \(e\circ u(x)\) 的乘法自同态给出的元素，该函数属于 \(\mathcal{L}_{\mathbf{C}}^{\infty}(\mathrm{X},\mu)\)。映射 \(\varrho_{u}\) 是从 G 到 \(\mathrm{L}_{\mathbf{C}}^{2}(\mathrm{X},\mu)\) 的连续表示，当且仅当映射 u 连续。（若 u 不连续，使用上一问题验证 \(\varrho\) 不连续。）
\item
  令 \(\varrho\) 为 E 在 \(\mathrm{L}^2 (\mathrm{X},\mu)\) 中的酉表示，且其像包含于 \(\mathcal{A}\)。存在唯一的连续线性映射 \(u\colon \mathrm{E}\to \mathrm{S}(\mathrm{X},\mu ;\mathbf{R})\)，使得 \(\varrho = \varrho_{u}\)。（先考虑 \(\mathrm{E} = \mathbf{R}\) 的情形，然后应用斯通定理。）
\end{enumerate}

\begin{enumerate}
\def\labelenumi{\arabic{enumi})}
\setcounter{enumi}{18}
\item
  令 G 为有限群，H 为 G 的子群。对于 H 在有限维复希尔伯特空间 E 中的每个酉表示 \(\pi\)，证明诱导酉表示 \(\operatorname{Ind}_{\mathrm{H}}^{\mathrm{G}}(\pi)\) 同构于 A, VIII, 第 392 页第 21 节第 3 号中定义的诱导表示。
\item
  令 G 为无穷远处可数的局部紧群，令 H 为 G 的闭子群。令 \(\pi\) 为 H 在希尔伯特空间 E 中的酉表示，令 \(\varrho\) 为由 \(\pi\) 诱导的 G 的酉表示。记 F 为 \(\varrho\) 的表示空间。
\end{enumerate}

在 G 上固定哈尔测度 \(\mu\)，在 H 上固定哈尔测度 \(\beta\)，并记 \(\kappa\colon\mathbf{G}\to\mathbf{R}_{+}^{*}\) 为满足 \(\kappa(xh)=\Delta_{\mathrm{H}}(h)\Delta_{\mathrm{G}}(h)^{-1}\kappa(x)\) 对 \((x,h)\in\mathbf{G}\times\mathrm{H}\) 成立，且满足 \(\kappa(e)=1\) 的连续函数。

\begin{enumerate}
\def\labelenumi{\alph{enumi})}
\item
  若 G/H 具有有界 G-不变测度，且存在非零的 H-不变向量 \(x \in E\)，则存在非零的 G-不变向量 \(y \in F\)。
\item
  反之，设 \(y \in F\) 是 \(\varrho\) 的非零 G-不变向量。证明存在非零函数 \(f \in L_{\overline{\pi}}^{2}\)，使得
\end{enumerate}

\[
f (x) = \left(\frac {\kappa (g ^ {- 1} x)}{\kappa (x)}\right) ^ {1 / 2} f (g ^ {- 1} x)
\]

对所有 \((g, x) \in \mathrm{G}^{2}\) 都成立。推出 \(\kappa = 1\)，从而 G/H 上的测度 \(\mu/\beta\) 有定义且有界，最后证明 \(f(e) \in \mathrm{E}\) 是非零 H-不变向量。
\exercisesection{}

\begin{enumerate}
\def\labelenumi{\arabic{enumi})}
\item
  设 G 为局部紧群。设 \(\pi\) 为 G 在复希尔伯特空间 E 中的酉表示，且 \(x \in \mathrm{E}\) 。我们有 \(\pi(g)x = x\) （对所有 \(g \in \mathrm{G}\) ），当且仅当正定函数 \(g \mapsto \langle x | \pi(g)x \rangle\) 为常数。
\item
  设 G 为局部紧群。对 G 上的每个正定函数 \(\varphi\) ，我们用 \((\pi_{\varphi}, x_{\varphi})\) 表示 \(\varphi\) 的一个循环希尔伯特表示。
\end{enumerate}

\begin{enumerate}
\def\labelenumi{\alph{enumi})}
\item
  设 \(\varphi_1, \varphi_2\) 与 \(\varphi_3\) 为 G 上的正定函数。若 \(\varphi_3 = \varphi_1 + \varphi_2\) ，则存在从 \(\pi_{\varphi_1}\) 到 \(\pi_{\varphi_3}\) 的单射 G-态射。
\item
  若 \(\varphi\) 是非零的常值正定函数，则 \(\pi_{\varphi}\) 是一维的平凡表示。
\item
  对 G 在希尔伯特空间 E 中的每个酉表示 \(\pi\) 以及每个 \(x \in \mathrm{E}\) ，存在从 \(\pi_{\varphi}\) 到 \(\pi\) 的单射 G-态射，其中 \(\varphi(g) = \langle x \mid \pi(g)x \rangle\) 。
\item
  设 H 为 G 的闭子群，\(\varphi\) 为 G 上的正定函数。用 \(\psi\) 表示 \(\varphi\) 在 H 上的限制；它是 H 上的正定函数，且存在从 \(\pi_{\psi}\) 到 \(\operatorname{Res}_{\mathrm{H}}^{\mathrm{H}}(\pi)\) 的单射 H-态射。
\end{enumerate}

\begin{enumerate}
\def\labelenumi{\arabic{enumi})}
\setcounter{enumi}{2}
\tightlist
\item
  \begin{enumerate}
  \def\labelenumii{\alph{enumii})}
  \tightlist
  \item
  给出一个对合巴拿赫代数的例子，它具有一个不连续的正线性泛函。
  \end{enumerate}
\end{enumerate}

\begin{enumerate}
\def\labelenumi{\alph{enumi})}
\setcounter{enumi}{1}
\tightlist
\item
  给出一个对合巴拿赫代数 A、A 上一个连续的正线性泛函 \(\lambda\) ，以及希尔伯特实现 \((\mathrm{E}_{1}, e_{1}, \varphi_{1})\) 与 \((\mathrm{E}_{2}, e_{2}, \varphi_{2})\) （均为 \(\lambda\) 的实现）的例子，使得 \((\mathrm{E}_{1}, e_{1}, \varphi_{1})\) 是循环希尔伯特实现，并且 V, p.~440 命题 6 中唯一的线性映射 u 不满足 \(u(e_{1}) = e_{2}\) 。
\end{enumerate}

\begin{enumerate}
\def\labelenumi{\arabic{enumi})}
\setcounter{enumi}{3}
\tightlist
\item
  设 A 与 B 为星代数。从 A 到 B 的线性映射 u 称为正的，如果 \(u(\mathrm{A}_{+}) \subset \mathrm{B}_{+}\) ；称为完全正的，如果对每个整数 \(k \in N\) ，映射 \(u_{k} = 1_{C^{k}} \otimes u: C^{k} \otimes A \to C^{k} \otimes B\) 是正的。
\end{enumerate}

\begin{enumerate}
\def\labelenumi{\alph{enumi})}
\tightlist
\item
  设 A 为含幺星代数，E 与 F 为复希尔伯特空间，v 为从 E 到 F 的线性映射，\(\varphi\) 为从 A 到 \(\mathcal{L}(\mathrm{F})\) 的对合代数的保幺态射，使得 \(\|\varphi\|^{2} = \|\varphi(1)\|\) 。映射 \(u: A \to \mathcal{L}(E)\) （由 \(u(a) = v^{*} \circ \varphi(a) \circ v\) 定义）是完全正的。
\end{enumerate}

本习题的目标是证明这一断言的逆命题（“斯廷斯普龙定理”）。设 A 为含幺星代数，E 为复希尔伯特空间，\(u \colon \mathrm{A} \to \mathcal{L}(\mathrm{E})\) 为完全正线性映射。

\begin{enumerate}
\def\labelenumi{\alph{enumi})}
\setcounter{enumi}{1}
\tightlist
\item
  设 \(F = A \otimes E\) 。F 上存在一个正的半双线性型 q，使得
\end{enumerate}

\[
q (a \otimes x, b \otimes y) = \langle x | \varphi (a ^ {*} b) y \rangle
\]

对所有 \((a,b)\in \mathrm{A}^2\) 与 \((x,y)\in \mathrm{E}^2\) 成立。

设 \(\widehat{\mathbf{F}}\) 是 \(\mathbf{F}\) 关于半双线性型 \(q\) 的分离完备希尔伯特空间。

\begin{enumerate}
\def\labelenumi{\alph{enumi})}
\setcounter{enumi}{2}
\tightlist
\item
  设 \(b \in A\) 。F 上使 \(a \otimes x \mapsto ba \otimes x\) （对所有 \(a \in A\) 与 \(x \in E\) ）成立的自同态，通过过渡到商诱导出 \(\varphi(b)\) （\(\widehat{F}\) 的一个自同态），满足 \(\|\varphi(b)\|\leqslant\|b\|\) ；映射 \(\varphi\colon A\to\mathcal{L}(F)\) 是对合代数的保幺态射。
\end{enumerate}

\begin{enumerate}
\def\labelenumi{\alph{enumi})}
\setcounter{enumi}{3}
\tightlist
\item
  设 \(v\) 为从 E 到 \(\widehat{\mathrm{F}}\) 的线性映射，使得 \(v(x)\) 是 \(1 \otimes x\) 所在的类（对所有 \(x \in \mathrm{E}\) ）。我们有 \(\|v\|^2 = \| \varphi(1)\|\) ，且 \(u(a) = v^* \circ \varphi(a) \circ v\) 对所有 \(a \in \mathrm{A}\) 成立。
\end{enumerate}

\begin{enumerate}
\def\labelenumi{\arabic{enumi})}
\setcounter{enumi}{4}
\tightlist
\item
  设 X 为分离拓扑空间。连续函数 \(f: X \times X \to R\) 称为条件负定型的，如果它满足下列条件
\end{enumerate}

\begin{enumerate}
\def\labelenumi{(\roman{enumi})}
\item
  对每个 \(x \in X\) ，有 \(f(x, x) = 0\) ；
\item
  对所有 \(x\) 与 \(y\) （属于 \(X\) ），我们有 \(f(x,y) = f(y,x)\) ；
\item
  对每个整数 \(n \in N\) 、X 中元素的每个族 \((x_{i})_{0 \leqslant i \leqslant n}\) 以及实数的每个族 \((t_{i})_{0 \leqslant i \leqslant n}\) （满足
\end{enumerate}

\[
\sum_ {i = 1} ^ {n} t _ {i} = 0,
\]

），我们有

\[
\sum_ {i = 0} ^ {n} \sum_ {j = 0} ^ {n} t _ {i} t _ {j} f (x _ {i}, x _ {j}) \leqslant 0.
\]

我们用 Noy \(_{-}\) (X) 表示 X × X 上条件负定型函数构成的集合。

\begin{enumerate}
\def\labelenumi{\alph{enumi})}
\item
  设 E 为实希尔伯特空间，\(\varphi\colon\mathrm{X}\to\mathrm{E}\) 为连续映射。函数 \(f(x,y)=\|\varphi(x)-\varphi(y)\|^{2}\) 是条件负定型的。
\item
  设 \(f\in\mathrm{Noy}_{-}(\mathrm{X})\) 且 \(x_{0}\in\mathrm{X}\) 。存在一个实希尔伯特空间 E 与一个连续映射 \(\varphi\colon\mathrm{X}\to\mathrm{E}\) ，使得
  \begin{enumerate}
  \def\labelenumii{\alph{enumii})}
  \tightlist
  \item
    对每个 \((x,y)\in \mathrm{X}\times \mathrm{X}\) ，有 \(f(x,y) = \| \varphi (x) - \varphi (y)\| ^2\) ；
  \item
    族 \((\varphi(x)-\varphi(x_{0}))_{x\in\mathrm{X}}\) 在 E 中是完全的。
  \end{enumerate}
\end{enumerate}

我们称 \((\mathrm{E},\varphi)\) 为 f 的一个循环希尔伯特表示。

\begin{enumerate}
\def\labelenumi{\alph{enumi})}
\setcounter{enumi}{2}
\item
  设 \(f \in \mathrm{Noy}_{-}(\mathrm{X})\) 且 \(x_0 \in \mathrm{X}\) 。设 \((\mathrm{E}, \varphi)\) 与 \((\mathrm{E}', \varphi')\) 为 \(f\) 的循环希尔伯特表示。存在唯一的仿射等距 \(u: \mathrm{E} \to \mathrm{E}'\) ，使得 \(\varphi' = u \circ \varphi\) 。
\item
  集合 \(\mathrm{Noy}_{-}(\mathrm{X})\) 是 \(\mathcal{C}(\mathrm{X}\times \mathrm{X};\mathbf{R})\) 中的一个凸锥。对每个普遍正核 \(g\in \mathrm{Noy}_{+}(\mathrm{X})\) （\(g\) 取实值），映射 \(f(x,y) = g(x,x) - g(x,y)\) 属于 \(\mathrm{Noy}_{-}(\mathrm{X})\) 。
\item
  设 \(\mathfrak{F}\) 为 \(\mathrm{Noy}_{-}(\mathrm{X})\) 上逐点收敛于 \(f\colon \mathrm{X}\times \mathrm{X}\to \mathbf{R}\) 的滤子。若 \(f\) 连续，则 \(f\in \mathrm{Noy}_{-}(\mathrm{X})\) 。
\item
  设 \(f: X \times X \to R\) 为满足上述条件 (i) 与 (ii) 的连续函数。设 \(x_{0} \in X\)。设 \(g: X \times X \to R\) 满足 \(g(x, y) = f(x, x_{0}) + f(y, x_{0}) - f(x, y)\) 。则 \(f \in \text{Noy}_{-}(X)\) 当且仅当 \(g \in \text{Noy}_{+}(X)\) 。
\item
  假设 X 非空。设 \(f: \mathrm{X} \times \mathrm{X} \to \mathbf{R}\) 为满足上述条件 (i) 与 (ii) 的连续函数。则 \(f \in \mathrm{Noy}_{-}(\mathrm{X})\) 当且仅当 \(\exp(-tf) \in \mathrm{Noy}_{+}(\mathrm{X})\) 对所有 \(t \in \mathbf{R}_{+}\) 成立（“舍恩贝格定理”）。（为证 \(\exp(-f)\) 在 \(f \in \mathrm{Noy}_{-}(\mathrm{X})\) 时是正定型的，固定 \(x_0 \in \mathrm{X}\) ；证明
\end{enumerate}

\[
g (x, y) = \exp (f (x, x _ {0}) + f (y, x _ {0}) - f (x, y))
\]

以及

\[
h (x, y) = \exp (- f (x, x _ {0})) \exp (- f (y, x _ {0}))
\]

都属于 Noy \(_{+}\) (X)。）

若 G 是拓扑群，我们称函数 \(\varphi \in \mathcal{C}(\mathrm{G})\) 是负定型的，如果函数 \(f\) 是 \(\mathrm{G} \times \mathrm{G}\) 上由 \(f(g,h) = \varphi(h^{-1}g)\) 定义的条件负定核。

\begin{enumerate}
\def\labelenumi{\arabic{enumi})}
\setcounter{enumi}{5}
\tightlist
\item
  设 G 为离散拓扑群。设 A 为 \(\mathcal{L}(\ell^{2}(\mathrm{G}))\) 中由左正则表示 \(\lambda\) 的像生成的对合子代数。这是一个含幺星代数。
\end{enumerate}

对每个正定型函数 \(\varphi\in\mathrm{Pos}(\mathrm{G})\) ，存在唯一的从 A 到 A 的完全正线性映射 \(u_{\varphi}\) （参见习题 4），使得 \(u_{\varphi}(\boldsymbol{\lambda}(g))=\varphi(g)\boldsymbol{\lambda}(g)\) 对所有 \(g\in\mathrm{G}\) 成立。（利用 \(\varphi\) 的一个循环希尔伯特表示。）

\begin{enumerate}
\def\labelenumi{\arabic{enumi})}
\setcounter{enumi}{6}
\tightlist
\item
  设 I 为有限集，\(G = F(I)\) 为在 I 上构造的自由群（A, I, p.~84）。设 A 为 \(\mathcal{L}(\ell^{2}(G))\) 中由左正则表示 \(\lambda\) 的像生成的对合子代数。
\end{enumerate}

\begin{enumerate}
\def\labelenumi{\alph{enumi})}
\item
  对每个 \(g \in G\) ，用 \(|g|\) 表示 g 的唯一既约分解的长度（A, I, p.~146, 习题 19）。映射 \(d(g, h) = |h^{-1}g|\) 是 G 上的一个距离。
\item
  设 \(\lambda \in \mathbf{R}_{+}^{*}\) 。G 上由 \(\psi_{\lambda}(g) = \exp(-\lambda|g|)\) 定义的函数是正定函数。（证明函数 \(g \mapsto |g|\) 是负定型函数，参见习题~5。）
\item
  对每个整数 \(k \in N\) ，设 \(G_{k}\) 为由 \(g \in G\) （满足 \(|g| = k\) ）组成的集合，并设 \(\varphi_{k}\) 为 \(G_{k}\) 的特征函数。设 f 与 g 为 G 上的函数，其支集分别含于 \(G_{k}\) 与 \(G_{l}\) 。对每个 \(m \in N\) ，我们有 \((f * g)\varphi_{m} = 0\) ，除非 \(k + l - m\) 为偶数且 \(|k - l| \leqslant m \leqslant k + l\) ；在这一情形，我们有
\end{enumerate}

\[
\left\| (f * g) \varphi_ {m} \right\| \leqslant \| f \| \| g \|.
\]

\begin{enumerate}
\def\labelenumi{\alph{enumi})}
\setcounter{enumi}{3}
\tightlist
\item
  设 f 为 G 上具有有限支集的函数。我们有
\end{enumerate}

\[
\| \boldsymbol {\lambda} (f) \| \leqslant \sum_ {n \in \mathbf {N}} (n + 1) \| f \varphi_ {n} \|,
\]

以及

\[
\left\| \boldsymbol {\lambda} (f) \right\| \leqslant 2 \left(\sum_ {g \in G} (1 + | g |) ^ {4} | f (g) | ^ {2}\right) ^ {1 / 2}.
\]

\begin{enumerate}
\def\labelenumi{\alph{enumi})}
\setcounter{enumi}{4}
\item
  设 \(\psi\) 为 G 上的函数，使得 \(M = \sup_{g \in G} |\psi(g)| (1 + |g|)^2\) 有限。存在唯一的从 A 到 A 的线性映射 \(u_\psi\) ，使得 \(u_\psi(\boldsymbol{\lambda}(g)) = \psi(g)\boldsymbol{\lambda}(g)\) 对所有 \(g \in G\) 成立。我们有 \(\|u_\psi\| \leqslant 2M\) 。
\item
  设 X 为 G 上由带有限支集的函数 \(\psi\) 组成的集合，其中 \(\|u_{\psi}\|\leqslant1\) 。存在 X 上的一个滤子 \(\mathfrak{F}\) ，使得
\end{enumerate}

\[
\lim _ {\psi , \mathfrak {F}} u _ {\psi} = 1 _ {\mathrm{A}}
\]

在赋予逐点收敛拓扑的空间 \(\mathcal{L}(\mathrm{A})\) 中。特别地，A 的底巴拿赫空间具有逼近性质。（考虑函数 \(\psi_{\lambda,n}(g) = (\varphi_0 + \dots + \varphi_n)\psi_\lambda\) ，其中 \(\lambda \in \mathbf{R}_+^*\) ，\(n \in \mathbf{N}\) ，并利用上一习题。）

\begin{enumerate}
\def\labelenumi{\arabic{enumi})}
\setcounter{enumi}{7}
\tightlist
\item
  我们把群 \(\mathbf{R}\) 与它的对偶群通过映射 \((x,y)\mapsto\) \(\exp (i\pi xy)\) 等同起来。设 \(\alpha \in \mathbf{R}_{+}\) 。
\end{enumerate}

\begin{enumerate}
\def\labelenumi{\alph{enumi})}
\item
  \(\mathbf{R}\) 上存在其傅里叶变换为函数 \(f_{\alpha}(x) = \exp(-|x|^{\alpha})\) 的正测度，当且仅当 \(\alpha \in ]0,2]\) 。在这一情形，我们把这一测度记作 \(\mu_{\alpha}\) ，它是唯一的；称之为参数 \(\alpha\) 的稳定测度。
\item
  测度 \(\mu_{\alpha}\) 的总质量为 1；R 上的恒等函数 f 关于 \(\mu_{\alpha}\) 可积，当且仅当 \(\alpha = 2\) 。
\end{enumerate}

\begin{enumerate}
\def\labelenumi{\arabic{enumi})}
\setcounter{enumi}{8}
\tightlist
\item
  本习题的目标是给出 V, p.~455 定理 5 的另一个证明。设 G 为局部紧交换拓扑群，带有哈尔测度 \(\mu\) 。我们用 \(\widehat{f}\) 表示 \(f \in \mathrm{L}^1(\mathrm{G})\) 的傅里叶变换，并用 \(\mathcal{F}(\mathrm{G}) \subset \mathcal{C}_0(\widehat{\mathrm{G}})\) 表示 \(\mathrm{L}^1(\mathrm{G})\) 上傅里叶变换的像。
\end{enumerate}

\begin{enumerate}
\def\labelenumi{\alph{enumi})}
\item
  空间 \(\mathcal{F}(\mathrm{G})\) 是 \(\mathcal{C}_0(\mathrm{G})\) 的稠密子空间。
\item
  对每个 \(f \in \mathrm{L}^{1}(\mathrm{G})\) ，f 在巴拿赫代数 \(\mathrm{L}^{1}(\mathrm{G})\) 中的谱半径等于 \(\|\widehat{f}\|_{\infty}\) 。
\item
  设 \(\varphi\in\mathrm{Pos}(\mathrm{G})\) 。它是 \(\mathrm{G}\) 上一致连续且有界的函数；对每个 \(f\in\mathrm{L}^{1}(\mathrm{G})\) ，我们有
\end{enumerate}

\[
\int_ {G} \int_ {G} f (x) \overline {{f (y)}} \varphi (x - y) d (\mu \otimes \mu) (x, y) \geqslant 0.
\]

\begin{enumerate}
\def\labelenumi{\alph{enumi})}
\setcounter{enumi}{3}
\tightlist
\item
  设 \(\lambda\) 为 \(\mathrm{L}^{1}(\mathrm{G})\) 上由 \(\varphi\) 确定的连续线性型，并令 \(b(f,g)=\lambda(f^{*}*g)\) （对所有 \((f,g)\in\mathrm{L}^{1}(\mathrm{G})\times\mathrm{L}^{1}(\mathrm{G})\) ）。映射 \(b\) 是 \(\mathrm{L}^{1}(\mathrm{G})\) 上的一个正埃尔米特型。
\end{enumerate}

\begin{enumerate}
\def\labelenumi{\alph{enumi})}
\setcounter{enumi}{4}
\tightlist
\item
  对每个 \(f \in \mathrm{L}^1(\mathrm{G})\) ，我们有 \(|\lambda(f)|^2 \leqslant b(f, f)\) 。
\end{enumerate}

\begin{enumerate}
\def\labelenumi{\alph{enumi})}
\setcounter{enumi}{5}
\tightlist
\item
  设 \(f \in \mathrm{L}^{1}(\mathrm{G})\) ，\(g = f^{*} * f \in \mathcal{C}_{b}(\mathrm{G})\) 。定义 \(g_{1} = g\) 与 \(g_{n+1} = g_{n} * g\) （对每个整数 \(n \geqslant 1\) ）。对每个整数 \(n \geqslant 1\) ，我们有 \(|\lambda(f)| \leqslant \|g_{n}\|_{1}^{2^{-n}}\) 。
\end{enumerate}

\begin{enumerate}
\def\labelenumi{\alph{enumi})}
\setcounter{enumi}{6}
\tightlist
\item
  对每个 \(f \in \mathrm{L}^{1}(\mathrm{G})\) ，我们有 \(|\lambda(f)| \leqslant \| \widehat{f} \|_{\infty}\) 。存在连续线性映射 \(\widehat{\lambda}\) （定义在 \(\mathcal{F}(\mathrm{G})\) 上），使得 \(\lambda(f) = \widehat{\lambda}(\widehat{f})\) 对所有 \(f \in \mathrm{L}^{1}(\mathrm{G})\) 成立。
\end{enumerate}

\begin{enumerate}
\def\labelenumi{\alph{enumi})}
\setcounter{enumi}{7}
\tightlist
\item
  存在有界测度 \(\nu\) （在 \(\widehat{G}\) 上），使得 \(\lambda(f) = \nu(\widehat{f})\) 对每个 \(f \in \mathrm{L}^{1}(\mathrm{G})\) 成立；测度 \(\nu\) 是正的，并且有
\end{enumerate}

\[
\varphi (x) = \int_ {\widehat {\mathrm{G}}} \langle \chi , x \rangle d \nu (\chi)
\]

对所有 \(x \in G\) 成立。

\begin{enumerate}
\def\labelenumi{\arabic{enumi})}
\setcounter{enumi}{9}
\tightlist
\item
  \begin{enumerate}
  \def\labelenumii{\alph{enumii})}
  \tightlist
  \item
    设 \(\mathrm{G} = \mathbf{R}\) 。集合 \(\operatorname{Pos}_1(\mathrm{G})\) 在弱拓扑下不是紧的。
  \end{enumerate}
\end{enumerate}

\begin{enumerate}
\def\labelenumi{\alph{enumi})}
\setcounter{enumi}{1}
\tightlist
\item
  设 G = R/Z。\(\operatorname{Pos}_{\leqslant1}(\mathbf{G})\) 上的弱拓扑与紧收敛拓扑不一致。
\end{enumerate}

\begin{enumerate}
\def\labelenumi{\arabic{enumi})}
\setcounter{enumi}{10}
\tightlist
\item
  设 G 为局部紧群。
\end{enumerate}

设 \(\varrho\) 与 \(\pi\) 为 G 的酉表示。我们称 \(\varrho\) 弱包含于 \(\pi\) ，并记作 \(\varrho \leqslant \pi\) ，如果 \(\varrho\) 的对角矩阵系数所成的集合包含于 \(\pi\) 的对角矩阵系数的有限和所成的集合在赋予紧收敛拓扑的空间 \(\mathcal{C}(\mathrm{G})\) 中的闭包内。

\begin{enumerate}
\def\labelenumi{\alph{enumi})}
\item
  关系“\(\varrho\) 与 \(\pi\) 是 G 的酉表示且 \(\varrho \leqslant \pi\)”是一个预序关系。若 \(\varrho'\) （相应地 \(\pi'\) ）同构于 \(\varrho\) （相应地 \(\pi\) ），则 \(\varrho \leqslant \pi\) 当且仅当 \(\varrho' \leqslant \pi'\) 。
\item
  若 \(\varrho\) 同构于 \(\pi\) 的一个子表示，则 \(\varrho \leqslant \pi\) 。若 I 与 J 是非空有限集，且若 \(\varrho \leqslant \pi\) ，则 \(\bigoplus_{i \in I} \varrho \leqslant \bigoplus_{j \in J} \pi\) 。
\item
  设 X 为 \(\varrho\) 的空间 E 中生成 \(\pi\) 的一个子集。我们有 \(\varrho \leqslant \pi\) ，当且仅当每个矩阵系数 \(g \mapsto \langle x | \varrho(g)x \rangle\) （带 \(x \in X\) ）都属于 \(\pi\) 的对角矩阵系数的有限和所成的集合在赋予紧收敛拓扑的空间 \(\mathcal{C}(\mathrm{G})\) 中的闭包。（证明 E 中满足这一性质的元素所成的集合 Y 是 E 的一个闭向量子空间。）
\item
  设 H 为 G 的闭子群。设 \(\pi\) 与 \(\varrho\) 为 G 的酉表示，使得 \(\varrho \leqslant \pi\) ；则我们有 \(\operatorname{Res}_{\mathrm{H}}^{\mathrm{G}}(\varrho) \leqslant \operatorname{Res}_{\mathrm{H}}^{\mathrm{G}}(\pi)\) 。
\item
  设 H 为 G 的闭子群，\(\nu\) 为 H 上的哈尔测度。设 \(\pi\) 与 \(\varrho\) 为 H 的酉表示，使得 \(\varrho \leqslant \pi\) ；则我们有 \(\operatorname{Ind}_{\mathrm{H}}^{\mathrm{G}}(\varrho) \leqslant \operatorname{Ind}_{\mathrm{H}}^{\mathrm{G}}(\pi)\) 。（令 \(\varrho' = \operatorname{Ind}_{\mathrm{H}}^{\mathrm{G}}(\varrho)\) ，\(\pi' = \operatorname{Ind}_{\mathrm{H}}^{\mathrm{G}}(\pi)\) ；利用 (c)，化归为逼近形如 \(g \mapsto \langle f | \varrho'(g)f \rangle\) 的矩阵系数，其中 f 形如
\end{enumerate}

\[
f (x) = \int_ {\mathrm{H}} f (x h) \varrho (h) y d \nu (y)
\]

（其中 \(\varphi\in\mathcal{K}(\mathrm{G})\) ，\(y\) 属于 \(\varrho\) 的空间。）

\begin{enumerate}
\def\labelenumi{\arabic{enumi})}
\setcounter{enumi}{11}
\item
  设 G = R，\(\pi\) 为 G 在 \(\mathrm{L}^{2}(\mathbf{R})\) 中的正则表示。证明 \(\chi \leqslant \pi\) 对每个特征标 \(\chi \in \widehat{G}\) 成立。
\item
  设 \(\varrho\) 为 G 在希尔伯特空间 E 中的不可约酉表示，\(\pi\) 为 G 在 F 中的酉表示。设 \(x \in E\) 为范数 1 的向量。假设 \(\varrho \leqslant \pi\) （参见习题~11）。
\end{enumerate}

我们用 \(A \subset \mathcal{C}(G)\) 表示 \(\pi\) 的规范化的对角矩阵系数所成的集合，用 \(\widetilde{A}\) 表示 A 在 \(L^{\infty}(G)\) 中关于弱拓扑的闭凸包。

设 \(x \in E\) 为范数 1 的向量，\(f: g \mapsto \langle x \mid \varrho(g)x \rangle\) 为相伴的对角矩阵系数。

\begin{enumerate}
\def\labelenumi{\alph{enumi})}
\item
  集合 \(\widetilde{\mathrm{A}}\) 是 \(\mathrm{L}^{\infty}(\mathrm{G})\) 中的一个紧凸集，而 \(f\) 是 \(\widetilde{\mathrm{A}}\) 的一个极点。
\item
  函数 f 属于 A 在 \(\widetilde{A}\) 中关于弱拓扑的闭包。
\item
  函数 f 属于 A 在赋予紧收敛拓扑的 \(\mathcal{C}(G)\) 中的闭包。
\item
  G 在 C 中的平凡表示弱包含于 \(\pi\) ，当且仅当对 G 的每个紧子集 C 与每个实数 \(\varepsilon > 0\) ，存在范数为 1 的 \(x \in F\) ，使得
\end{enumerate}

\[
\sup _ {g \in \mathrm{C}} \| \pi (g) x - x \| <   \varepsilon .
\]

\begin{enumerate}
\def\labelenumi{\arabic{enumi})}
\setcounter{enumi}{13}
\tightlist
\item
  设 G 为局部紧群。我们称 G 具有卡日丹性质 ]T[ \(^{(1)}\) ，如果 G 的酉表示 \(\pi\) 具有非零的平凡子表示 \(^{(2)}\) ，当且仅当 G 在 C 上的平凡表示 \(\mathbf{1}_{\text{G}}\) 弱包含于 \(\pi\) 。
\end{enumerate}

\begin{enumerate}
\def\labelenumi{\alph{enumi})}
\item
  设 G 为具有性质 ]T[ 的局部紧群。设 H 为 G 中由 G 的紧子集生成的开子群所成的集合。群 G 是诸子群 H ∈ H 的并。
\item
  设 \(\pi\) 为 G 的这样一个表示，它是 G 在 G/H（\(H \in H\) ）上的拟正则表示的希尔伯特和（参见 V, p. 486 的习题 2）。我们有 \(1_{G} \leqslant \pi\) 。
\item
  群 G 由 G 的一个紧子集生成。特别地，若 G 是离散的，则 G 由 G 的一个有限子集生成。
\end{enumerate}

\begin{enumerate}
\def\labelenumi{\arabic{enumi})}
\setcounter{enumi}{14}
\tightlist
\item
  设 G 为局部紧群，C 为 G 的子集，\(\varepsilon > 0\) 。我们称 \((\mathrm{C}, \varepsilon)\) 是 G 的一个卡日丹对，如果对 G 在希尔伯特空间 E 中的每个酉表示 \(\pi\) ，存在向量 \(x \in E\)
\end{enumerate}

使得

\[
\sup _ {g \in \mathrm{C}} \| \pi (g) x - x \| <   \varepsilon \| x \|,
\]

  便蕴含 \(\pi\) 有一个非零的平凡子表示。我们称 C 为卡日丹集，如果存在 \(\varepsilon > 0\) 使得 \((C, \varepsilon)\) 是一个卡日丹对。

\begin{enumerate}
\def\labelenumi{\alph{enumi})}
\tightlist
\item
  设 \((\mathrm{C}, \varepsilon)\) 为 G 的一个卡日丹对。设 \(\pi\) 为 G 在希尔伯特空间 E 中的酉表示。设 \(x \in E\) 非零且 \(\delta > 0\) ，使得
\end{enumerate}

\[
\sup _ {g \in \mathrm{C}} \| \pi (g) x - x \| <   \delta \varepsilon \| x \|.
\]

则我们有 \(\|p(x)-x\|\leqslant\delta\|x\|\) ，其中 p 是 E 的正交投影算子，其像为 G 不变向量所成的空间。

\begin{enumerate}
\def\labelenumi{\alph{enumi})}
\setcounter{enumi}{1}
\item
  局部紧群 G 具有性质 ]T[，当且仅当存在一个卡日丹对 (C,ε)。（为证性质 ]T[ 蕴含这一断言，用反证法，考虑诸表示 π \(_{C,\varepsilon}\) 的希尔伯特和表示，这些表示由 G 的紧子集 C 与 ε > 0 参数化，使得 π \(_{C,\varepsilon}\) 不含非平凡子表示，但存在范数为 1 的 x \(_{C,\varepsilon}\) 满足 sup \(_{g\in\mathbf{C}}\|\pi(g)x_{\mathrm{C},\varepsilon}-x_{\mathrm{C},\varepsilon}\|<\varepsilon.\)
\item
  设 G 为具有性质 ]T[ 的局部紧群。G 的每个生成 G 的紧子集 C 都是卡日丹集。
\item
  设 G 为具有性质 ]T[ 的局部紧群，C 为 G 的一个紧子集，其内部非空且是卡日丹集。则 C 生成 G。特别地，若 G 是离散的，则卡日丹集恰好是 G 的有限生成子集。（由 C 生成的 G 的子群 H 是开的；考虑 G 在 G/H 上的拟正则表示 \(\pi\) ，并证明 H 的特征函数 \(\varphi\) （H \(\in\) G/H）在 \(\pi\) 的不变向量空间上的投影等于 \(\varphi\) 。）
\end{enumerate}

\begin{enumerate}
\def\labelenumi{\arabic{enumi})}
\setcounter{enumi}{15}
\item
  设 G 为在无穷远处可数的局部紧群，H 为 G 的闭子群。若 G 具有性质 ]T[，且若 G/H 容许一个有界的 G 不变测度，则 H 具有性质 ]T[。（利用 V, p.~492 的习题 20。）
\item
  设 \(n \geqslant 2\) 为整数。我们考虑实李群 \(\mathbf{G} = \mathbf{R}^n \rtimes \mathbf{SL}(n, \mathbf{R})\) ，并用 H 表示 \(\mathbf{R}^n \times \{e\}\) （\(\mathbf{G}\) 的子群），它与 \(\mathbf{R}^n\) 等同。
\end{enumerate}

设 \(\pi\) 为 G 在复希尔伯特空间 E 中的酉表示。假设对 G 的每个紧子集 C 与每个 \(\varepsilon > 0\) ，存在非零的 \(x \in E\) ，使得

\[
\sup _ {g \in \mathrm{C}} \| \pi (g) x - x \| <   \varepsilon \| x \|.
\]

\begin{enumerate}
\def\labelenumi{\alph{enumi})}
\item
  存在 E 中由范数为 1 的向量构成的序列 \((x_{m})\) ，使得函数 \(\varphi_{m}(g)=\langle x_{m}|\pi(g)x_{m}\rangle\) 在赋予紧收敛拓扑的 \(\mathcal{C}(\mathrm{G})\) 中收敛于常值函数 1。（考虑一个其并等于 G 的紧子集序列 \((\mathrm{C}_{m})_{m\in\mathbf{N}}\) ，并取 \(x_{m}\) 使其对 \((\mathrm{C}_{m},1/(m+1))\) 满足上述条件。）
\item
  对每个 \(m \in N\) ，存在总质量为 1 的正测度 \(\mu_{m}\) （在 \(R^{n}\) 上），其傅里叶变换与 \(\varphi_{m}\) 在 H 上的限制等同。
\item
  假设 \(\mu_{m}(\{0\}) = 0\) 对所有 \(m \in N\) 成立。这时把 \(\mu_{m}\) 与 \(R^{n} - \{0\}\) 上总质量为 1 的正测度等同起来，并用 \(\nu_{m}\) 表示 \(\mu_{m}\) 在射影空间 \(\mathbf{P}_{n-1}(\mathbf{R})\) 上的像测度。存在子序列 \((\nu_{n_{k}})_{k \in \mathbf{N}}\) ，它狭收敛于总质量为 1 的正测度 \(\nu\) （在 \(\mathbf{P}_{n-1}(\mathbf{R})\) 上）；测度 \(\nu\) 是 \(\mathbf{SL}(n, \mathbf{R})\) -不变的。由此导出矛盾。
\item
  设 \(m \in N\) 满足 \(\mu_{m}(\{0\}) > 0\) 。设 \(\psi\) 为 \(\varphi_{m}\) 在 H 上的限制。存在 c > 0，使得 \(\psi = c + \widetilde{\psi}\) ，其中 \(\widetilde{\psi}\) 是 \(R^{n}\) 上的正定型函数。
\item
  表示 \(\pi\) 含有 H 不变的向量。（利用习题 2。）我们称 (G,H) 具有相对性质 ]T[。
\end{enumerate}

\begin{enumerate}
\def\labelenumi{\arabic{enumi})}
\setcounter{enumi}{17}
\tightlist
\item
  设 \(\mathrm{G} = \mathbf{SL}(2, \mathbf{R})\) 。用 N 与 P 表示子群
\end{enumerate}

\[
\mathrm{N} = \Big \{\left( \begin{array}{c c} 1 & t \\ 0 & 1 \end{array} \right) | t \in \mathbf {R} \Big \}, \qquad \mathrm{P} = \Big \{\left( \begin{array}{c c} a & t \\ 0 & a ^ {- 1} \end{array} \right) | a \in \mathbf {R} _ {*}, t \in \mathbf {R} \Big \}
\]

（G 的子群）。设 \(\pi\) 为 G 在复希尔伯特空间 E 中的酉表示。假设存在 \(x \in \mathrm{E}\) ，使得 \(\pi(n)x = x\) 对所有 \(n \in \mathrm{N}\) 成立。

\begin{enumerate}
\def\labelenumi{\alph{enumi})}
\tightlist
\item
  齐性空间 \(\mathbf{SL}(2, \mathbf{R}) / \mathrm{N}\) 与 \(\mathbf{R}^2 - \{0\}\) 通过映射 \(g \mapsto ge_1\) 等同，其中 \(e_1\) 是 \(\mathbf{R}^2\) 的典范基的第一个向量；齐性空间 \(\mathbf{SL}(2, \mathbf{R}) / \mathrm{P}\) 与子空间 \((\mathbf{R} - \{0\}) \times \{0\}\) （\(\mathbf{SL}(2, \mathbf{R}) / \mathrm{N}\) 的）等同。
\item
  G 上的每个连续函数 \(f\) ，若满足 \(f(ngn') = f(g)\) （对所有 \((n, n') \in \mathrm{N}^2\) ），则满足 \(f(pgp') = f(g)\) （对所有 \((p, p') \in \mathrm{P}^2\) ）。（把 \(f\) 解释为连续函数 \(\varphi\) （在 \(\mathbf{R}^2 - \{0\}\) 上、在 N 的自然作用下不变），并验证它在 \(\mathbf{SL}(2, \mathbf{R}) / \mathrm{P}\) 上是常数。）
\end{enumerate}

\begin{enumerate}
\def\labelenumi{\alph{enumi})}
\setcounter{enumi}{2}
\tightlist
\item
  有 \(\pi(g)x = x\) 对所有 \(g \in G\) 成立。
\end{enumerate}

\begin{enumerate}
\def\labelenumi{\arabic{enumi})}
\setcounter{enumi}{18}
\tightlist
\item
  设 \(n \geqslant 2\) 为整数。考虑实李群 \(\mathrm{G} = \mathbf{SL}(n, \mathbf{R})\) 。用 N 与 H 表示 G 中元素分别为分块形式
\end{enumerate}

\[
\left( \begin{array}{c c} \operatorname{Id} _ {n - 1} & x \\ 0 & 1 \end{array} \right), \quad (x \in \mathbf {R} ^ {n - 1}), \qquad \left( \begin{array}{c c} g & 0 \\ 0 & 1 \end{array} \right), \quad (g \in \mathbf {S L} (n - 1, \mathbf {R}))
\]

的闭子群。群 N 与 \(R^{n-1}\) 等同，群 H 与 \(\mathbf{SL}(n-1,\mathbf{R})\) 等同。

\begin{enumerate}
\def\labelenumi{\alph{enumi})}
\item
  设 \(\pi\) 为 G 在复希尔伯特空间 E 中的酉表示。E 的元素 \(x\) 在 G 下不变，当且仅当它在 N 下不变。
\item
  若 \(n \geqslant 3\) ，则群 G 具有性质 ]T[。（把上一问题的结果与上一习题结合起来。）
\item
  若 \(n \geqslant 3\) ，则群 \(\widetilde{G} = R^{n} \rtimes G\) 具有性质 ]T[。（设 \(\pi\) 为满足 \(1_{\widetilde{G}} \leqslant \pi\) 的酉表示；由上一问题，存在非零向量 \(x \neq 0\) ，它被 \(\widetilde{G}\) 的子群 G 保持不变；证明相应的矩阵系数是常数。）
\item
  若 \(n \geqslant 3\) ，则群 \(\mathbf{SL}(n, \mathbf{Z})\) 具有性质 ]T[。（利用 INT, VII, p.~116 的习题 7。）
\item
  设 \(n \geqslant 3\) 。群 \(\mathbf{SL}(n, \mathbf{Z})\) 是有限生成的。设 H 为 \(\mathbf{SL}(n, \mathbf{Z})\) 的有限指标正规子群所成的族；它是无限的。对 \(\mathbf{SL}(n, \mathbf{Z})\) 的每个有限生成子集 S，凯莱图族 \((\mathcal{G}(\mathbf{SL}(n, \mathbf{Z})/\mathrm{H}, \mathrm{SH}/\mathrm{H}))_{\mathrm{H} \in \mathcal{H}}\) 是一个扩展图族（参见 TA, 2, p.~221, 习题 7 及 IV, p.~323 的习题 16）。
\end{enumerate}

\begin{enumerate}
\def\labelenumi{\arabic{enumi})}
\setcounter{enumi}{19}
\tightlist
\item
  设 G 为局部紧群，\(\mu\) 为 G 上的左哈尔测度。我们用 \(\Delta\) 表示 G 的模。对 G 上的任意函数 f，我们将记 \(\widetilde{f} = \Delta f^{*}\) 。
\end{enumerate}

G 上的测度 \(\nu\) 称为正定型的，如果 \(\langle \nu, \tilde{f} * f \rangle \geqslant 0\) 对每个函数 \(f \in \mathcal{K}(\mathrm{G})\) 成立。G 上的局部可积函数 \(\varphi\) 称为正定型的，如果 \(\langle \varphi, f^* * f \rangle \geqslant 0\) 对每个函数 \(f \in \mathcal{K}(\mathrm{G})\) 成立。

\begin{enumerate}
\def\labelenumi{\alph{enumi})}
\item
  假设 \(\nu\) 有界。测度 \(\nu\) 是正定型的，当且仅当 \(\gamma(\nu) \geqslant 0\) 。
\item
  设 \(\varphi\in\mathcal{K}(\mathrm{G})\) 。函数 \(\varphi\) 作为局部可积函数是正定型的，当且仅当它按照 V, p.~442 的定义 4 是正定型的。
\item
  设 \(\nu\) 为 G 上的一个正定型测度。公式
\end{enumerate}

\[
(f \mid g) _ {\nu} = \langle \nu , \widetilde {f} * g \rangle
\]

在 \(\mathcal{K}(\mathrm{G})\) 上定义了一个正埃尔米特型。我们用 \(\mathrm{E}_{\nu}\) 表示赋予了该埃尔米特型的 \(\mathcal{K}(\mathrm{G})\) 的分离完备希尔伯特空间。

\begin{enumerate}
\def\labelenumi{\alph{enumi})}
\setcounter{enumi}{3}
\tightlist
\item
  对 \(f \in \mathcal{K}(\mathrm{G})\) ，函数 \(\sigma(g)f \colon g \mapsto (\varepsilon_g * f)\Delta^{1/2}\) 属于 \(\mathcal{K}(\mathrm{G})\) ，且 \(\| \sigma(g)f \|_{\nu} = \|f\|_{\nu}\) 。映射 \(g \mapsto \sigma(g)\) 延拓为 G 在 \(\mathrm{E}_{\nu}\) 中的一个酉表示。
\end{enumerate}

\begin{enumerate}
\def\labelenumi{\alph{enumi})}
\setcounter{enumi}{4}
\tightlist
\item
  设 U 为 G 中 e 的开邻域，\(\nu\) 为 G 上的正定型测度，使得存在实数 \(c \geqslant 0\) 满足
\end{enumerate}

\[
\langle \nu , f * \widetilde {f} \rangle \leqslant c \left(\int_ {\mathrm{G}} f d \mu\right) ^ {2}
\]

  其中 \(f \in \mathcal{K}(\mathrm{G})\) 为正的且在 U 外为零。存在 G 上的连续正定型函数 \(\varphi\) ，使得 \(\nu = \varphi \Delta^{-1/2} \cdot \mu\) 。（设 \(j\) 为从 \(\mathcal{K}(\mathrm{G})\) 到 \(\mathrm{E}_{\nu}\) 的典范态射；若 \(\mathfrak{B}\) 是 \(e\) 的邻域滤子的一组基（其元素都含于 U），而 \(f_{\mathrm{V}} \in \mathcal{K}_{+}(\mathrm{G})\) 是支集含于 \(V \in B\) 、积分为 1 的函数，证明 \(j(f_{\mathrm{V}})\) 沿 B 狭收敛于元素 \(y_{0}\) （属于 \(E_{\nu}\) ），使得

\[
(j (f) \mid y _ {0}) = \int_ {\mathrm{G}} f ^ {*} d \mu
\]

对 \(f \in \mathcal{K}(\mathrm{G})\) 成立；令 \(\varphi(g) = (y_0 \mid \sigma(g)y_0)\) ，验证 \(\varphi\) 具有所需的性质。

\begin{enumerate}
\def\labelenumi{\alph{enumi})}
\setcounter{enumi}{5}
\tightlist
\item
  设 \(\varphi_{1}\) 为 G 上的连续正定函数，\(\varphi_{2}\) 为局部可积的正定型函数。若 \(\varphi_{1}-\varphi_{2}\) 是正定型的，则 \(\varphi_{2}\) 局部几乎处处等于一个连续的正定型函数。（验证可以应用上一问题。）
\end{enumerate}

\begin{enumerate}
\def\labelenumi{\arabic{enumi})}
\setcounter{enumi}{20}
\tightlist
\item
  我们沿用上一习题的记号。
\end{enumerate}

\begin{enumerate}
\def\labelenumi{\alph{enumi})}
\item
  设 \(f \in \mathcal{L}^{2}(G)\) 作为局部可积函数是正定型的。\(L^{2}(G)\) 上以 \(\mathcal{K}(G)\) 为定义域、由 \(\varphi \mapsto \gamma(\varphi)f\) 定义的部分算子是正的。我们用 \(\varrho(f)\) 表示 f 的弗里德里希斯延拓（IV, p.~295 的例 8）；它是 \(L^{2}(G)\) 上的一个正自伴部分算子。
\item
  设 \(f \in \mathcal{L}^2(\mathrm{G})\) 。我们称 \(f\) 是清醒的，如果存在实数 \(c \geqslant 0\) ，使得 \(\| \boldsymbol{\gamma}(\varphi)f \| \leqslant c \| \varphi \|\) 对所有 \(\varphi \in \mathcal{K}(\mathrm{G}) \subset \mathcal{L}^2(\mathrm{G})\) 成立。这时我们用 \(\varrho(f)\) 表示 \(\mathrm{L}^2(\mathrm{G})\) 中延拓 \(\varphi \mapsto \boldsymbol{\gamma}(\varphi)f\) 的自同态。若 \(f\) 是正定型的，这一定义与上一个定义一致，并且这一条件此时等价于说 \(\varrho(f)\) 是正的。
\item
  设 \(f_{1}\) 与 \(f_{2}\) 为 \(\mathcal{L}^{2}(G)\) 中的清醒函数。若 \(\varrho(f_{1}) = \varrho(f_{2})\) ，则 \(f_{1}\) 与 \(f_{2}\) 在 \(L^{2}(G)\) 中相等。
\end{enumerate}

\begin{enumerate}
\def\labelenumi{\alph{enumi})}
\setcounter{enumi}{3}
\tightlist
\item
  设 \(f \in \mathcal{L}^2(\mathrm{G})\) 为正定型的。对所有 \(g \in \mathrm{G}\) ，我们有 \(\varrho(f) \circ \boldsymbol{\gamma}(g) = \boldsymbol{\gamma}(g) \circ \varrho(f)\) 。此外，对每个 \(h \in \operatorname{dom}(\varrho(f))\) ，我们有 \(\varrho(f)h = h * f\) 。
\end{enumerate}

\begin{enumerate}
\def\labelenumi{\alph{enumi})}
\setcounter{enumi}{4}
\tightlist
\item
  设 \(f_{1}\) 与 \(f_{2}\) 为 \(\mathcal{L}^{2}(\mathrm{G})\) 中的清醒元素，使得 \(\varrho(f_{1})\) 与 \(\varrho(f_{2})\) 交换。函数 \(f_{1}*f_{2}\) 是 G 上的连续正定函数；它是清醒的，并且我们有
\end{enumerate}

\[
\varrho (f _ {1} * f _ {2}) = \varrho (f _ {1}) \varrho (f _ {2}), \quad \langle f _ {1} | f _ {2} \rangle \geqslant 0.
\]

\begin{enumerate}
\def\labelenumi{\alph{enumi})}
\setcounter{enumi}{5}
\tightlist
\item
  在上一问题的记号与假设下，若 \(f_{2} - f_{1}\) 是正定型的，则
\end{enumerate}

\[
\left\| f _ {2} - f _ {1} \right\| ^ {2} \leqslant \left\| f _ {2} \right\| ^ {2} - \left\| f _ {1} \right\| ^ {2}.
\]

\begin{enumerate}
\def\labelenumi{\alph{enumi})}
\setcounter{enumi}{6}
\tightlist
\item
  设 \((f_n)_{n\in \mathbf{N}}\) 为 \(\mathcal{L}^2 (\mathrm{G})\) 中两两交换的清醒正定型元素构成的序列，且 \(f_{n + 1} - f_n\) 对所有 \(n\in \mathbf{N}\) 是正定型的。若序列 \((\| f\| _n)\) 有界，则序列 \((f_n)\) 在 \(\mathcal{L}^2 (\mathrm{G})\) 中收敛。
\end{enumerate}

\begin{enumerate}
\def\labelenumi{\arabic{enumi})}
\setcounter{enumi}{21}
\tightlist
\item
  我们沿用上一习题的记号。
\end{enumerate}

设 \(f\) 为属于 \(\mathcal{L}^2 (\mathrm{G})\) 的连续正定型函数。

\begin{enumerate}
\def\labelenumi{\alph{enumi})}
\item
  存在 [0,1] 上多项式函数的递增序列 \((p_{n})_{n\in\mathbf{N}}\) ，它在 [0,1] 上一致收敛于函数 \(t\mapsto\sqrt{t}\) 。
\item
  假设 \(f\) 是清醒的，且 \(0 \leqslant f \leqslant 1\) （在 \(\mathcal{L}(\mathrm{L}^2(\mathrm{G}))\) 中）。存在序列 \((h_n)\) （在 \(\mathcal{L}^2(\mathrm{G})\) 中），满足下列条件：
\end{enumerate}

\begin{enumerate}
\def\labelenumi{(\roman{enumi})}
\item
  每个 \(h_n\) 都是清醒的且是正定型的；
\item
  函数 \(h_{n+1} - h_n\) 对所有 \(n \in \mathbf{N}\) 是正定型的；
\item
  自同态 \(\varrho(h_n)\) 两两交换；
\item
  有 \(\varrho(h_n)^2 \leqslant \varrho(f)\) （对所有 \(n\) ）。
\end{enumerate}

\begin{enumerate}
\def\labelenumi{\alph{enumi})}
\setcounter{enumi}{2}
\tightlist
\item
  序列 \((h_n)_{n\in \mathbf{N}}\) 在 \(\mathcal{L}^2 (\mathrm{G})\) 中收敛；若 \(h\) 是 \((h_n)\) 的一个极限，则 \(h\) 是正定型且清醒的，并满足 \(\varrho (h)^2 = \varrho (f)\) ；我们有 \(f = h*h\) 。
\end{enumerate}

\begin{enumerate}
\def\labelenumi{\alph{enumi})}
\setcounter{enumi}{3}
\tightlist
\item
  我们不再假设 \(f\) 是清醒的。对 \(t \in \mathbf{R}\) ，用 \(p_t\) 表示部分自伴算子 \(\varrho(f)\) 的由 \([- \infty, t]\) 定义的谱投影。我们有 \(p_t \circ \boldsymbol{\gamma}(g) = \boldsymbol{\gamma}(g) \circ p_t\) 对所有 \(t \in \mathbf{R}\) 与 \(g \in G\) 成立。
\end{enumerate}

\begin{enumerate}
\def\labelenumi{\alph{enumi})}
\setcounter{enumi}{4}
\tightlist
\item
  对 \(t \in \mathbf{R}\) ，令 \(f_t = p_t(f)\) ；它是 \(\mathrm{L}^2(\mathrm{G})\) 中某个连续正定的清醒函数（属于 \(\mathcal{L}^2(\mathrm{G})\) ）的类。（验证 \(\gamma(\varphi)f_t = \varrho(f)p_t(g)\) 对 \(g \in \mathcal{H}(\mathrm{G})\) 成立，由此推出 \(\varrho(f_t) = \varrho(f) \circ p_t\) ，再应用习题 20 的结果。）
\end{enumerate}

\begin{enumerate}
\def\labelenumi{\alph{enumi})}
\setcounter{enumi}{5}
\tightlist
\item
  存在 \(h \in \mathcal{L}^2(\mathrm{G})\) ，使得 \(f = h * h\) 。（考虑 \(f_n = p_n(f)\) （\(n \in \mathbf{N}\) ），并应用部分 \(c\) 。）
\end{enumerate}

\begin{enumerate}
\def\labelenumi{\arabic{enumi})}
\setcounter{enumi}{22}
\tightlist
\item
  我们记 \(G = R_{+}^{*}\) ；它是一个局部紧交换群，且测度 \(\mu = dx/x\) 是 G 上的哈尔测度。群 G 与群 iR 通过映射 \((x, it) \mapsto x^{it}\) 处于对偶；\(\mu\) 的对偶测度是 iR 上的勒贝格测度（参见 II, p. 236 的推论 4）。我们用 F 表示 G 的傅里叶变换，用 \(\overline{F}\) 表示 \(\widehat{G}\) 的傅里叶余变换。
\end{enumerate}

我们用 X 表示 G 上的恒等函数。

\begin{enumerate}
\def\labelenumi{\alph{enumi})}
\item
  设 \(\nu\) 为 G 上的正测度。我们用 \(\mathrm{I}_{\nu}\) 表示这样的 \(\sigma \in \mathbf{R}\) 所成的集合：函数 \(\mathrm{X}^{\sigma}\) 是 \(\nu\) -可积的；用 \(\mathrm{B}_{\nu}\) 表示这样的 \(s \in \mathbf{R}\) 所成的集合：\(\mathcal{R}(s) \in \mathrm{I}_{\nu}\) 。集合 \(\mathrm{I}_{\nu}\) 是 \(\mathbf{R}\) 的一个区间，而函数 \(s \mapsto \nu(\mathrm{X}^s)\) 在 \(\mathrm{B}_{\nu}\) 的内部有定义且全纯。
\item
  设 \(f\) 为 G 上的局部可积函数。我们记 \(\mathrm{I}_{f} = \mathrm{I}_{f\cdot \mu}\) 与 \(\mathrm{B}_f = \mathrm{B}_{f\cdot \mu}\) 。对每个 \(s \in \mathrm{B}_f\) ，我们令
\end{enumerate}

\[
\widehat {f} (s) = \int_ {\mathrm{G}} f (x) x ^ {s} d \mu (x).
\]

若存在 \(\delta > 0\) 与 \(c \geqslant 0\) ，使得 \(|\widehat{f}(s)| \leqslant c(1 + |s|)^{-1 - \delta}\) 对 \(s \in \mathrm{B}_f\) 成立，则对所有 \(\sigma\) （属于 \(\mathrm{B}_f\) 的内部），我们有

\[
f (x) = \frac {1}{2 \pi} \int_ {\mathbf {R}} \widehat {f} (\sigma + i t) x ^ {- \sigma - i t} d t.
\]

\begin{enumerate}
\def\labelenumi{\alph{enumi})}
\setcounter{enumi}{2}
\tightlist
\item
  设 \(f \in \mathcal{L}^{2}(\mathrm{G})\) 。假设 0 属于 \(I_{f}\) 的内部。则我们有 \(\overline{\mathcal{F}} f(it) = \widehat{f}(it)\) 对几乎所有 \(t \in \mathbf{R}\) 成立。
\end{enumerate}

\begin{enumerate}
\def\labelenumi{\alph{enumi})}
\setcounter{enumi}{3}
\item
  设 \(f\) 与 \(g\) 为 \(\mathbf{G}\) 上的可测函数。设 \(\sigma\) 为实数，使得 \(\mathrm{X}^{-2\sigma}f\in \mathrm{L}^2 (\mathrm{G})\) 且 \(\mathrm{X}^{2\sigma}g\in \mathrm{L}^2 (\mathrm{G})\) 。则我们有

\[
\int_ {\mathrm{G}} f g d \mu = \frac {1}{2 \pi} \int_ {\mathbf {R}} \widehat {f} (- \sigma - i t) \widehat {g} (\sigma + i t) d t.
\]

\end{enumerate}

\begin{enumerate}
\def\labelenumi{\alph{enumi})}
\setcounter{enumi}{4}
\tightlist
\item
  设 \(f\) 与 \(g\) 为 \(G\) 上的可测函数。假设
\end{enumerate}

\begin{enumerate}
\def\labelenumi{(\roman{enumi})}
\item
  存在 \(c \in \mathbf{R}\) ，使得 \(|f(x)| \leqslant cx / (1 + x)\) 对所有 \(x \in \mathbf{G}\) 成立；
\item
  有 ]0,1[ ⊂ I \(_{g}\) ，且对每个满足 0 < σ < 1 的 σ，函数 X \(^{\sigma}\) g 有界；
\item
  存在 \(c' \in \mathbf{R}\) ，使得 \(|\widehat{g}(s)| \leqslant c'(1 + |s|)^{-2}\) 对所有 \(s \in \mathbf{C}\) （满足 \(\mathcal{R}(s) \in ]0, 1[\) ）成立。
\end{enumerate}

则我们有

\[
\int_ {\mathrm{G}} f g d \mu = \frac {1}{2 \pi} \int_ {\mathbf {R}} \widehat {f} (- \sigma - i t) \widehat {g} (\sigma + i t) d t
\]

对所有 \(\sigma\in\mathbf{R}\) （满足 \(0<\sigma<1\) ）成立。

\begin{enumerate}
\def\labelenumi{\alph{enumi})}
\setcounter{enumi}{5}
\tightlist
\item
  设 \(\sigma \in \mathbf{R}\) ，\(f \in \mathcal{C}(\sigma + i\mathbf{R})\) 。G 上存在正测度 \(\nu\) ，使得 \(\sigma \in \mathrm{I}_{\nu}\) 且 \(\widehat{\nu}(\sigma + it) = f(\sigma + it)\) 对所有 \(t \in \mathbf{R}\) 成立，当且仅当对每个整数 \(n \geqslant 1\) 、复数的每个族 \((t_i)_{0 \leqslant i \leqslant n}\) 与任意族 \((z_i)_{0 \leqslant i \leqslant n}\) （\(\mathcal{R}(z_i) = \sigma\) 对所有 \(i\) ），我们有
\end{enumerate}

\[
\sum_ {i = 0} ^ {n} \sum_ {j = 0} ^ {n} \bar {t} _ {i} t _ {j} f \left(\frac {1}{2} (\overline {{z}} _ {i} + z _ {j})\right) \geqslant 0.
\]

\exercisesection{}

\begin{enumerate}
\def\labelenumi{\arabic{enumi})}
\tightlist
\item
  设 G 为紧拓扑群，B 为 G 的不可约酉表示的特征标所成的族。
\end{enumerate}

\begin{enumerate}
\def\labelenumi{\alph{enumi})}
\item
  偶 \((\mathrm{R}(\mathrm{G}),\mathcal{B})\) 是一个自然的含幺基代数（参见 III, p.~130 的习题 6）。
\item
  假设 G 有限。代数 R(G) 这时是有限秩的；计算 \(p_f(\pi)\) （对每个 \(\pi \in \widehat{G}\) ），并辨识与上述引文问题 j) 中定义的元素 r₀ 相对应的 G 的表示。
\end{enumerate}

\begin{enumerate}
\def\labelenumi{\arabic{enumi})}
\setcounter{enumi}{1}
\tightlist
\item
  设 \(n \geqslant 1\) ，\(p \in [1, +\infty]\) 。我们给 \(\mathbf{R}^n\) 赋予范数
\end{enumerate}

\[
\| x \| = \left(\sum_ {i = 1} ^ {n} | x _ {i} | ^ {p}\right) ^ {1 / p} \text {   if   } p \neq + \infty ,
\]

而 \(\|x\| = \sup_{i}|x_{i}|\) （若 \(p = +\infty\) ）。设 G 为赋有该范数的度量空间 \(R^{n}\) 的等距群。

\begin{enumerate}
\def\labelenumi{\alph{enumi})}
\item
  设 \(\sigma\) 为 \(\{1,\ldots,n\}\) 的一个置换，\(\varepsilon\in\{-1,1\}^{n}\) 。设 \((e_{i})\) 为 \(\mathbf{R}^{n}\) 的典范基。\(\mathbf{R}^{n}\) 中使 \(e_{i}\mapsto\varepsilon_{i}e_{\sigma(i)}\) 的唯一自同态属于 G。这种形式的等距所成的集合是 G 的一个子群 W；它是 \(\mathfrak{S}_{n}\) 被 \((\mathbf{Z}/2\mathbf{Z})^{n}\) 的一个扩张。
\item
  群 G 是紧的。
\item
  群 G 在 \(R^{n}\) 上的作用是不可约的；群 G 含于二次型 \(\sum x_{i}^{2}\) 的正交群。
\item
  若 \(p \neq 2\) ，则 G = W。
\end{enumerate}

\begin{enumerate}
\def\labelenumi{\arabic{enumi})}
\setcounter{enumi}{2}
\tightlist
\item
  我们回顾，称群 \(\Gamma\) 是剩余有限的，如果 \(\Gamma\) 的有限指标子群的交约化为 \(\{e\}\) （A, I, p.~129, 习题 5, b)）。
\end{enumerate}

对群 \(\Gamma\) ，我们用 \(\widehat{\Gamma}\) 表示射影极限 \(\widehat{\Gamma} = \varprojlim \Gamma / H\) ，其中 \(H\) 取遍 \(\Gamma\) 的按包含排序的有限指标正规子群所成的集合。我们称 \(\widehat{\Gamma}\) 是 \(\Gamma\) 的射有限完备化。

\begin{enumerate}
\def\labelenumi{\alph{enumi})}
\item
  群 \(\Gamma\) 是剩余有限的，当且仅当对每个异于单位元的 \(g \in \Gamma\) ，存在有限群 F 与同态 \(\varphi: \Gamma \to F\) ，使得 \(\varphi(g) \neq e\) 。
\item
  证明 \(\Gamma\) 是剩余有限的，当且仅当从 \(\Gamma\) 到 \(\widehat{\Gamma}\) 的典范同态是单射。
\item
  设 K 为域。\(\mathbf{GL}(n,\mathrm{K})\) 的每个有限生成子群（参见 A, I, p.~145, 习题 16）都是剩余有限的（“马尔采夫定理”）。（证明存在整数 \(m\geqslant 0\) 、一个等于 \(\mathbf{Z}\) 或某个有限域的环 A，以及该子群到 \(\mathbf{GL}(m,\mathrm{R})\) 的一个单射同态，其中 R 形如 \(\mathrm{A}[\mathrm{X}_1,\dots ,\mathrm{X}_r,1 / f]\) 。）
\item
  设 \(\Gamma\) 为有限生成的离散群。假设存在从 \(\Gamma\) 到某个紧拓扑群的单射同态。对每个异于单位元的 \(\gamma \in \Gamma\) ，存在 \(\Gamma\) 的一个有限维不可约酉表示，使得 \(\gamma\) 在 \(\pi\) 下的像不是恒等映射。
\item
  有限生成群是剩余有限的，当且仅当它同构于某个紧群的一个子群。
\end{enumerate}

¶ 4) 设 G 为群，A 与 B 为 G 的子群。设 \(\varphi\colon A\to B\) 是一个同构。

设 \(r'\) 为自由群 F(G) 中的一个族，使得 G 同构于群 F(G, \(r'\) ) （A, I, p.~87, n°6）。我们称 G 依 \(\varphi\) 而得的 HNN 扩张为群 \(G^{\varphi} = F(X, r)\) ，其中 X 是集合 G、A、B 与一个元素 t 之和，而 r 是 \(r'\) 与关系子 \(t^{-1}at = \varphi(a)\) （对所有 \(a \in A\) ）的和族。

我们用 \(j\colon G\to G^{\varphi}\) 表示唯一的同态，使得 \(j(g)\) 是 \(g\) 在 \(G^{\varphi}\) 中所代表的元素。

\begin{enumerate}
\def\labelenumi{\alph{enumi})}
\tightlist
\item
  设 \(k \geqslant 1\) 为整数，\((g_{0}, \varepsilon_{1}, g_{1}, \ldots, \varepsilon_{k}, g_{k})\) 为一个有限族，使得对每个 i 有 \(g_{i} \in G\) 且 \(\varepsilon_{i} \in \{-1, 1\}\) 。假设不存在满足 \(1 \leqslant j \leqslant k - 1\) 的整数 j 使得
\end{enumerate}

\[
(\varepsilon_ {j}, \varepsilon_ {j + 1}) = (- 1, 1), \quad g _ {i} \in A
\]

或

\[
(\varepsilon_ {j}, \varepsilon_ {j + 1}) = (1, - 1), \quad g _ {i} \in \mathrm{B}.
\]

则元素

\[
j (g _ {0}) t ^ {\varepsilon_ {1}} j (g _ {1}) \dots t ^ {\varepsilon_ {k}} j (g _ {k})
\]

（G \(^{φ}\) 的这个元素）不是单位元（“布里顿引理”。）

\begin{enumerate}
\def\labelenumi{\alph{enumi})}
\setcounter{enumi}{1}
\item
  同态 j 是单射。
\item
  设 \(a\) 与 \(b\) 为不同的符号。设 G 为群 \(\mathrm{F}(\{a,b\}, a^{-1}b^{2}ab^{-3})\) 。记 \(b_{1} = a^{-1}ba\) 。设 H 为 G 的一个有限指标正规子群，\(\pi: \mathrm{G} \to \mathrm{G}/\mathrm{H}\) 为典范投影。有 \([b_{1}, b] \in \mathrm{Ker}(\pi)\) 。
\item
  群 G 不容许忠实的有限维表示。（证明 G 不是剩余有限的。）
\end{enumerate}

\begin{enumerate}
\def\labelenumi{\arabic{enumi})}
\setcounter{enumi}{4}
\tightlist
\item
  设 G 为紧群，\(\varrho\) 为 G 的一个忠实的有限维酉表示。设 \(\pi \in \widehat{G}\) 。
\end{enumerate}

\begin{enumerate}
\def\labelenumi{\alph{enumi})}
\tightlist
\item
  存在整数 a 与 b，使得 \(\pi\) 同构于 \(\varrho^{\otimes a} \otimes \overline{\varrho}^{\otimes b}\) 的一个子表示。（例如可以利用魏尔斯特拉斯–斯通定理。）
\end{enumerate}

我们用 \(w(\pi)\) 表示 \(a + b\) 当 \((a, b)\) 取遍满足这一性质的整数组所成的集合时的最小值。

\begin{enumerate}
\def\labelenumi{\alph{enumi})}
\setcounter{enumi}{1}
\tightlist
\item
  设 \(\Lambda\) 为 \(G^{0}\) 的权格，q 为 \(\Lambda\otimes \mathbf{R}\) 上的一个范数。存在实数 C 与 D，使得对每个表示 \(\pi\in\widehat{G}\) （其在 \(G^{0}\) 上的限制具有最高权 \((\lambda_{1},\ldots,\lambda_{k})\) ），我们有
\end{enumerate}

\[
w (\pi) \leqslant \mathrm{C} \sup (q (\lambda_ {1}), \dots , q (\lambda_ {k})) + \mathrm{D}.
\]

（先考虑连通半单群的情形，并对 G 的适当选取的有限表示族应用 a)。）

\begin{enumerate}
\def\labelenumi{\arabic{enumi})}
\setcounter{enumi}{5}
\tightlist
\item
  设 G 为紧群，\(\varrho\) 为 G 的一个不可约表示。存在中心函数 \(f \in \Theta(G^{\sharp})\) ，使得
\end{enumerate}

\[
f = \sum_ {\pi \in \widehat {\mathrm{G}}} c (\pi) \chi_ {\pi}
\]

其中整数 \(c(\pi)\) 满足

\begin{enumerate}
\def\labelenumi{(\roman{enumi})}
\item
  有 \(c(\pi) \geqslant 0\) （对所有 \(\pi \in \widehat{G}\) ）；
\item
  \(c(1) \leqslant c(\varrho)\) ，且除 \(\varrho\) 是一维的并满足 \(\varrho^{2} = 1\) 的可能情形外，该不等式是严格的；
\item
  \(\mathcal{R}(f) \geqslant 0.\)
\end{enumerate}

¶ 7) 对每个整数 \(n \geqslant 1\) ，我们用 \(\mu_n\) 表示 \(\mathbf{C}\) 上群 \(\mathbf{U}(\mathbf{C}^n) \subset \mathrm{GL}(\mathbf{C}^n)\) 的规范化哈尔测度在迹映射 \(g \mapsto \mathrm{Tr}(g)\) 下的像；它是总质量为 1 的紧支集测度。

\begin{enumerate}
\def\labelenumi{\alph{enumi})}
\tightlist
\item
  设 \(a\) 与 \(b\) 为自然数，满足 \(a + b \leqslant n\) 。我们有
\end{enumerate}

\[
\int_ {\mathbf {C}} z ^ {a} \overline {{z}} ^ {b} d \mu_ {n} (z) = \left\{ \begin{array}{l l} 0 & \text{若 } a \neq b \\ a! & \text{若 } a = b. \end{array} \right.
\]

\begin{enumerate}
\def\labelenumi{\alph{enumi})}
\setcounter{enumi}{1}
\tightlist
\item
  测度序列 \((\mu_{n})_{n\geqslant1}\) 狭收敛于实向量空间 C 上协方差为 \(\mathrm{Q}(z)=|z|^{2}\) 的高斯测度。
\end{enumerate}

\begin{enumerate}
\def\labelenumi{\arabic{enumi})}
\setcounter{enumi}{7}
\item
  设 G 为紧群，\(\varrho_{1}\) 与 \(\varrho_{2}\) 为 G 的不可约酉表示。设 \(n \geqslant 2\) 为整数。若 \(S^{n}\varrho_{1}\) 同构于 \(S^{n}\varrho_{2}\) ，则存在 G 的一个有限指标子群 H，使得 \(\varrho_{1}\) 在 H 上的限制同构于 \(\varrho_{2}\) 在 H 上的限制。
\item
  设 \(k\) 为含 5 个元素的有限域，\(\mathrm{G} = \mathbf{S}\mathbf{L}(2,k)\) 。
\end{enumerate}

\begin{enumerate}
\def\labelenumi{\alph{enumi})}
\tightlist
\item
  群 G 的共轭类由 e、-e 以及七个元素
\end{enumerate}

\[
g _ {3} = \left( \begin{array}{c c} 1 & 1 \\ 0 & 1 \end{array} \right), \qquad g _ {4} = \left( \begin{array}{c c} 1 & 2 \\ 0 & 1 \end{array} \right), \qquad g _ {5} = \left( \begin{array}{c c} - 1 & 1 \\ 0 & - 1 \end{array} \right), \qquad g _ {6} = \left( \begin{array}{c c} - 1 & 2 \\ 0 & - 1 \end{array} \right)
\]

\[
g _ {7} = \left( \begin{array}{c c} 2 & 0 \\ 0 & - 2 \end{array} \right), \qquad g _ {8} = \left( \begin{array}{c c} 0 & 1 \\ - 1 & 1 \end{array} \right), \qquad g _ {9} = \left( \begin{array}{c c} 0 & 1 \\ - 1 & - 1 \end{array} \right).
\]

代表。确定每个共轭类的基数。

\begin{enumerate}
\def\labelenumi{\alph{enumi})}
\setcounter{enumi}{1}
\tightlist
\item
  设 \(\alpha = \frac{1}{2}(1 + \sqrt{5})\) 。存在 G 的一个二维忠实不可约酉表示 \(\pi\) ，其特征标 \(\chi\) 满足 \(\chi(-e) = -2\) 以及
\end{enumerate}

\[
\begin{array}{c} \chi (g _ {4}) = \alpha , \quad \chi (g _ {6}) = - \alpha , \quad \chi (g _ {3}) = - \alpha - 1, \quad \chi (g _ {5}) = \alpha + 1 \\ \chi (g _ {7}) = 0, \quad \chi (g _ {8}) = 1, \quad \chi (g _ {9}) = - 1. \end{array}
\]

（可以考虑 \(\mathbf{GL}(2,k)\) 的四维不可约表示在 G 上的限制，参见 A, VIII, p.~422, 习题 27。）

\begin{enumerate}
\def\labelenumi{\alph{enumi})}
\setcounter{enumi}{2}
\tightlist
\item
  有 \(\mathrm{M}_4(\pi) = 2\) 。
\end{enumerate}

\begin{enumerate}
\def\labelenumi{\arabic{enumi})}
\setcounter{enumi}{9}
\tightlist
\item
  \begin{enumerate}
  \def\labelenumii{\alph{enumii})}
  \tightlist
  \item
    存在一个正交型的不可约表示 \(\pi\) （交错群 \(\mathfrak{A}_5\) 的），使得 \(\mathrm{M}_4(\pi) = 3\) 。
  \end{enumerate}
\end{enumerate}

\begin{enumerate}
\def\labelenumi{\alph{enumi})}
\setcounter{enumi}{1}
\tightlist
\item
  设 R 为 \(E_{6}\) 型（相应地 \(E_{7}\) 型、\(E_{8}\) 型）的根系，W 为 R 的外尔群。设 \(\pi\) 为 W 的几何表示（LIE, V, § 4, 第 3 目）。它是不可约的、正交型的，且满足 \(M_{4}(\pi)=3\) 。
\end{enumerate}

（可以利用诸如 GAP 或 MAGMA 的软件来确定 W 的共轭类及其在 \(\pi\) 下的像的列表。）

\begin{enumerate}
\def\labelenumi{\arabic{enumi})}
\setcounter{enumi}{10}
\item
  设 G 为紧群，\(\pi\) 为 G 在希尔伯特空间 E 中的有限维酉表示。若 \(\mathrm{M}_{4}(\pi) \leqslant 5\) ，则表示 \(\pi\) 是不可约的。
\item
  设 G 为连通紧半单李群，其复化李代数同构于单李代数 \(e_{6}\) 。
\end{enumerate}

\begin{enumerate}
\def\labelenumi{\alph{enumi})}
\item
  存在 G 的一个 27 维忠实不可约表示 \(\varrho\) ；这样的表示是复型的。
\item
  表示 \(\varrho\) 的四阶矩等于 3。
\end{enumerate}

\begin{enumerate}
\def\labelenumi{\arabic{enumi})}
\setcounter{enumi}{12}
\tightlist
\item
  设 G 为连通紧半单李群，g 为其李代数。假设存在 G 的一个 27 维忠实不可约表示 \(\varrho\) 。我们假设 \(\varrho\) 的四阶矩等于 3，且表示 \(\varrho\) 是复型的。
\end{enumerate}

\begin{enumerate}
\def\labelenumi{\alph{enumi})}
\item
  李代数 g 是单的。
\item
  李代数 g 同构于 \(\mathfrak{e}_{6}\) 。
\end{enumerate}

\begin{enumerate}
\def\labelenumi{\arabic{enumi})}
\setcounter{enumi}{13}
\item
  设 G 为连通紧半单李群。假设存在 G 的一个 27 维忠实不可约表示 \(\varrho\) 。G 的复化李代数同构于 \(\mathfrak{e}_6\) ，当且仅当表示 \(\operatorname{End}(\varrho)\) 中存在一个 78 维的不可约子表示。
\item
  设 \(q\) 为奇素数的幂，F 为含 \(q\) 个元素的有限域。考虑由 F 上系数的 3 阶上三角幺幂矩阵组成的海森伯群 H \(_3\) (F)。对 \((a, b, c) \in \mathrm{F}^3\) ，我们记
\end{enumerate}

\[
\sigma (a, b, c) = \left( \begin{array}{c c c} 1 & a & c \\ 0 & 1 & b \\ 0 & 0 & 1 \end{array} \right) \in \mathrm{H} _ {3} (\mathrm{F}).
\]

因此映射 \(\sigma\) 是双射。

\begin{enumerate}
\def\labelenumi{\alph{enumi})}
\item
  \(\mathrm{H}_{3}(\mathrm{F})\) 的中心 C 通过把 \(\sigma(a,b,c)\) 映到 c 的群同态与 F 等同。
\item
  存在 \(q^{2}\) 个 \(\mathrm{H}_{3}(\mathrm{F})\) 的一维酉特征标；这些特征标在 C 上是平凡的。
\item
  设 X 为 F 中由元素 \(\sigma(0,b,c)\) （\((b,c)\in\mathrm{F}^{2}\) ）组成的指标为 q 的子群。子群 X 与 \(F^{2}\) 等同。设 \(\psi_{1}\) 与 \(\psi_{2}\) 为 F 的酉特征标，令 \(\psi(\sigma(0,b,c))=\psi_{1}(b)\psi_{2}(c)\) 。酉表示 \(\varrho\) （\(\mathrm{H}_{3}(\mathrm{F})\) 的、由 X 的特征标 \(\psi\) 诱导）是不可约的，当且仅当 \(\psi_{2}\) 非平凡；它在同构意义下只依赖于 \(\psi_{2}\) ，其中心特征标等于 \(\psi_{2}\) 。\(\varrho\) 的特征标的支集等于 C。
\item
  在前面两个问题中构造的不可约表示就是 \(\mathrm{H}_{3}(\mathrm{F})\) 的全部不可约酉表示。
\end{enumerate}

\begin{enumerate}
\def\labelenumi{\alph{enumi})}
\setcounter{enumi}{4}
\tightlist
\item
  设 \(\varrho\) 为 \(\mathrm{H}_{3}(\mathrm{F})\) 的 q 维不可约酉表示之一。表示 \(\varrho\) 是忠实的，且满足 \(\mathrm{M}_{4}(\varrho)=q^{2}\) 。
\end{enumerate}

\begin{enumerate}
\def\labelenumi{\arabic{enumi})}
\setcounter{enumi}{15}
\tightlist
\item
  设 \(r \geqslant 1\) 为整数，G 为 \(\mathbf{U}(\mathbf{C}^{r})\) 的连通紧子群，它在 \(C^{r}\) 上不可约地作用。我们用 D 表示 \(\mathbf{U}(\mathbf{C}^{r})\) 中由对角矩阵组成的子群。
\end{enumerate}

\begin{enumerate}
\def\labelenumi{\alph{enumi})}
\tightlist
\item
  设 A 为 D 的一个含于 G 在 \(\mathbf{U}(\mathbf{C}^{r})\) 中的正规化子内的子群。设 S 为 \(\widehat{A}\) 中由对角系数 \(a = (a_{i,j}) \mapsto a_{i,i}\) （\(1 \leqslant i \leqslant r\) ）组成的子集。假设 S 含有 r 个元素，且方程
\end{enumerate}

\[
\chi_ {1} \chi_ {2} ^ {- 1} = \chi_ {3} \chi_ {4} ^ {- 1}
\]

（其中 \(\chi_{i} \in S\) ）只有形如 \((\chi_{1}, \chi_{1}, \chi_{3}, \chi_{3})\) 与 \((\chi_{1}, \chi_{2}, \chi_{1}, \chi_{2})\) 的解（这时称 \(S\) 是一个 Sidon 集）。

若 G 是半单的，证明 \(\mathbf{G} = \mathbf{S}\mathbf{U}(\mathbf{C}^{r})\) 。（考虑 A 在 \(\operatorname{End}(\mathbf{C}^{r})\) 上的表示，并由假设推出 G 的复化李代数在 D 的共轭下稳定。）

\begin{enumerate}
\def\labelenumi{\alph{enumi})}
\setcounter{enumi}{1}
\tightlist
\item
  群 G 含有 \(\mathbf{SU}(\mathbf{C}^{r})\) ，当且仅当存在元素 \(g \in G\) ，使得 g 的特征值互不相同并在 \(C^{*}\) 中构成一个 Sidon 集。
\end{enumerate}

\begin{enumerate}
\def\labelenumi{\arabic{enumi})}
\setcounter{enumi}{16}
\tightlist
\item
  设 G 为紧群，H 为 G 的闭子群。设 \(\varrho\) 为 H 在有限维希尔伯特空间 E 中的酉表示。
\end{enumerate}

\begin{enumerate}
\def\labelenumi{\alph{enumi})}
\tightlist
\item
  设 \(\pi \in \widehat{\mathbf{G}}\) 。空间 \(\mathrm{Hom}_{\mathrm{G}}(\pi, \mathrm{Ind}_{\mathrm{H}}^{\mathrm{G}}(\varrho))\) 是有限维的。

\end{enumerate}

\begin{enumerate}
\def\labelenumi{\alph{enumi})}
\setcounter{enumi}{1}
\item
  对每个 \(u \in \mathrm{Hom}_{\mathrm{G}}(\pi, \mathrm{Ind}_{\mathrm{H}}^{\mathrm{G}}(\varrho))\) ，\(u\) 的像含于 \(\mathcal{C}(\mathrm{G};\mathrm{E})\) 在 \(\mathrm{L}^2 (\mathrm{G};\mathrm{E})\) 中的像。
\item
  存在线性映射 \(\Phi\) ，它从空间 \(\mathrm{Hom}_{\mathrm{G}}(\pi, \mathrm{Ind}_{\mathrm{H}}^{\mathrm{G}}(\varrho))\) 到空间 \(\mathrm{Hom}_{\mathrm{H}}(\mathrm{Res}_{\mathrm{H}}^{\mathrm{G}}(\pi), \varrho)\) ，使得对每个 \(u \in \mathrm{Hom}_{\mathrm{G}}(\pi, \mathrm{Ind}_{\mathrm{H}}^{\mathrm{G}}(\varrho))\) 与每个 \(x \in \mathrm{E}_{\pi}\) ，\(\Phi(u)(x)\) 是连续函数在 \(e\) 处的值，该连续函数表示 \(u(x)\) 。
\item
  线性映射 \(\Phi\) 是同构。
\item
  假设 \(\varrho\) 不可约。\(\pi\) 在 \(\mathrm{Ind}_{\mathrm{H}}^{\mathrm{G}}(\varrho)\) 中的重数等于 \(\varrho\) 在 \(\mathrm{Res}_{\mathrm{H}}^{\mathrm{G}}(\pi)\) 中的重数（弗罗贝尼乌斯互反公式，参见 A, VIII, p. 392）。
\end{enumerate}

\begin{enumerate}
\def\labelenumi{\arabic{enumi})}
\setcounter{enumi}{17}
\item
  假设 G 是紧的。设 \(\varrho\) 与 \(\varrho'\) 为 G 的酉表示。表示 \(\varrho\) 弱包含于 \(\varrho'\) （参见 V, p. 497 的习题 11），当且仅当对 G 的每个不可约酉表示 \(\pi\) ，条件 \(m_{\pi}(\varrho) \geqslant 1\) 蕴含 \(m_{\pi}(\varrho') \geqslant 1\) 。（化归到 \(\varrho'\) 不可约的情形；对每个规范化的对角矩阵系数 \(f\) （属于 \(\varrho\) ），应用 V, p.~498 的习题 13 证明存在 \(\varrho'\) 的一个与 \(f\) 不正交的对角矩阵系数，再用正交关系得出结论。）
\item
  设 G 为赋有规范化哈尔测度 \(\mu\) 的紧群，\(f \in \mathcal{C}(\mathrm{G})\) 为 G 上的连续函数。存在不可约表示 \(\pi \in \widehat{G}\) ，使得 \(f = \dim(\pi)^{-1}\chi_{\pi}\) ，当且仅当
\end{enumerate}

\[
\int_ {\mathrm{G}} f (x g x ^ {- 1} h) d \mu (x) = f (g) f (h)
\]

对所有 \((g,h)\in \mathrm{G}^2\) 成立

\begin{enumerate}
\def\labelenumi{\arabic{enumi})}
\setcounter{enumi}{19}
\tightlist
\item
  设 H 为赋有左哈尔测度 \(\nu\) 的局部紧拓扑群。设 G 为 H 的紧子群。设 \(\varrho\) 为 H 在希尔伯特空间 E 中的不可约酉表示。
\end{enumerate}

我们假设 \(\pi\) 在 \(\varrho\) 到 G 的限制中的重数对每个不可约表示 \(\pi \in \widehat{G}\) 都是有限的。我们把 G 上的函数与 H 上紧支集的测度等同起来。

\begin{enumerate}
\def\labelenumi{\alph{enumi})}
\item
  对每个 \(\pi \in \widehat{\mathrm{G}}\) ，E 的自同态 \(\varrho(f * \chi_{\pi})\) 是有限秩的。
\item
  对每个 \(f \in \mathrm{L}^{1}(\mathrm{H})\) ，自同态 \(\varrho(f)\) 是紧的。
\end{enumerate}

\begin{enumerate}
\def\labelenumi{\arabic{enumi})}
\setcounter{enumi}{20}
\tightlist
\item
  设 \(n \geqslant 1\) 为整数。
\end{enumerate}

\begin{enumerate}
\def\labelenumi{\alph{enumi})}
\tightlist
\item
  对每个整数 \(r > n^{2}\) 与每个族 \((a_{1}, \ldots, a_{r})\) （\(\mathcal{L}(\mathbf{C}^{n})\) 的元素），我们有
\end{enumerate}

\[
\sum_ {\sigma \in \mathfrak {S} _ {r}} \varepsilon (\sigma) a _ {\sigma (1)} \dots a _ {\sigma (r)} = 0.\tag{4}
\]

\begin{enumerate}
\def\labelenumi{\alph{enumi})}
\setcounter{enumi}{1}
\tightlist
\item
  我们用 \(r(n)\) 表示最小的整数 \(r \geqslant 1\) ，使公式 (4) 对每个族 \((a_i)_{1 \leqslant i \leqslant r}\) 成立。我们有 \(r(n) \geqslant r(n-1) + 2\) 。（考虑一族 \((a_1, \ldots, a_{r(n-1)-1})\) （属于 \(\mathcal{L}(\mathbf{C}^{n-1})\) ），使得记作 b 的上述表达式非零；把 \(\mathcal{L}(\mathbf{C}^n)\) 的元素解释为矩阵，选取 i 与 j 使 \(b_{i,j}\) 非零；把诸 \(a_i\) 与 b 视为 \(\mathcal{L}(\mathbf{C}^n)\) 的元素（通过添加一行一列零），然后定义矩阵 \(a_{r(n-1)}\) 与 \(a_{r(n-1)+1}\) ，使得 \((a_1, \ldots, a_{r(n-1)+1})\) 的表达式 (4) 非零。）
\end{enumerate}

\begin{enumerate}
\def\labelenumi{\arabic{enumi})}
\setcounter{enumi}{21}
\tightlist
\item
  设 \(r \geqslant 1\) 为整数，\(\mathrm{H} \subset \mathbf{GL}(r, \mathbf{C})\) 为闭子群。
\end{enumerate}

\begin{enumerate}
\def\labelenumi{\alph{enumi})}
\item
  H 上形如 \(g \mapsto \langle \varrho(g)x, \lambda \rangle\) 的函数所成的集合（其中 \(\varrho\) 是 H 在有限维复向量空间 E 中的连续线性表示，\(x \in E\) ，\(\lambda \in E'\) ）是 \(\mathcal{C}(\mathrm{H})\) 的稠密对合子代数。
\item
  对任意非零的函数 \(f \in \mathcal{K}(\mathrm{H})\) ，存在 H 在有限维复向量空间中的一个连续线性表示 \(\varrho\) ，使得 \(\varrho(f)\) 非零。
\item
  假设 H 是连通半单实李群。对每个非零的 \(f \in \mathcal{K}(H)\) ，存在 H 在有限维复向量空间中的不可约连续线性表示 \(\varrho\) ，使得 \(\varrho(f)\) 非零。
\end{enumerate}

\begin{enumerate}
\def\labelenumi{\arabic{enumi})}
\setcounter{enumi}{22}
\tightlist
\item
  我们沿用上一习题的假设与记号，并假设 H 是连通半单实李群。
\end{enumerate}

设 G 为 H 的紧子群。假设存在 H 的连通可解闭子群 N，使得 H = GN。（对每个半单实李群，当 G 是 H 的极大紧子群时这一条件满足；例如当 G = SL(n, R) 时，参见 INT, VII, p.~91, § 3, n° 3, 命题 7。）

设 \(\pi \in \widehat{\mathbf{G}}\) ，\(e_{\pi} = \dim (\pi)\overline{\chi}_{\pi}\) 。设 \(\varrho\) 为上一习题最后一问给出的 H 的不可约表示；我们用 E 表示 \(\varrho\) 的空间。

\begin{enumerate}
\def\labelenumi{\alph{enumi})}
\item
  \(\varrho\) 到 G 的限制存在循环向量。（利用李定理，LIE, III, p.~241, 推论，找出非零的 \(x \in E\) ，使 C \(x\) 在 N 的作用下不变。）\(\pi\) 在 \(\varrho\) 到 G 的限制中的重数至多等于 \(\dim(\pi)\) 。
\item
  设 \(r\) 为习题 22 中定义的整数 \(r(\dim (\pi)^2)\) ；恒等式
\end{enumerate}

\[
\sum_ {\sigma \in \mathfrak {S} _ {r}} \varepsilon (\sigma) a _ {\sigma (1)} \dots a _ {\sigma (r)} = 0
\]

对每个元素族 \(a_{i}=e_{\pi}*g_{i}*e_{\pi}\) （属于 H 上紧支集测度的代数 \(\mathcal{M}_{c}(\mathrm{H})\) ，其中 \(g_{i}\in\mathcal{K}(\mathrm{H})\) （对 \(1\leqslant i\leqslant r\) ））成立。

\begin{enumerate}
\def\labelenumi{\alph{enumi})}
\setcounter{enumi}{2}
\tightlist
\item
  设 \(\varpi\) 为 H 在希尔伯特空间 F 中的不可约酉表示。设 \(F_{\pi}\) 为 \(\pi\)-同型分量（在 \(\varpi\) 到 G 的限制中），\(p\) 为 F 的以 \(F_{\pi}\) 为像的正交投影算子。对每个族 \((u_i)_{1 \leqslant i \leqslant r}\) （属于 \(\mathcal{L}(F_{\pi})\) ），我们有
\end{enumerate}

\[
\sum_ {\sigma \in \mathfrak {S} _ {r}} \varepsilon (\sigma) u _ {\sigma (1)} \dots u _ {\sigma (r)} = 0
\]

  （利用 V, p.~486 的习题 1）。

\begin{enumerate}
\def\labelenumi{\alph{enumi})}
\setcounter{enumi}{3}
\tightlist
\item
  \(\pi\) 在 \(\varpi\) 到 G 的限制中的重数至多等于 \(\dim(\pi)\) 。（利用 V, p.~510 的习题 22 的结果。）
\end{enumerate}

\begin{enumerate}
\def\labelenumi{\arabic{enumi})}
\setcounter{enumi}{23}
\tightlist
\item
  称拓扑群 G 是极小概周期的，如果从 G 到紧群的每个连续同态都是平凡的。
\end{enumerate}

\begin{enumerate}
\def\labelenumi{\alph{enumi})}
\item
  拓扑群 G 是极小概周期的，当且仅当 G 不容许非平凡的有限维酉表示。若 G 是交换的，则 G 是极小概周期的，当且仅当 G 的每个连续酉特征标都是平凡的。
\item
  设 G 为拓扑群，\(g \in G\) 满足：对每个整数 \(k \geqslant 1\) ，存在整数 \(m \geqslant 1\) ，使得 \(g^{km}\) 共轭于 g。则有 \(\varrho(g) = 1_{\mathrm{E}}\) 对 G 在希尔伯特空间 E 中的每个有限维酉表示 \(\varrho\) 成立。
\item
  设 k 为 C 的子域。赋有离散拓扑的群 \(\mathbf{SL}(2,k)\) 是极小概周期的。（验证上一问题适用于 \(\begin{pmatrix}1 & z \\ 0 & 1\end{pmatrix}\) （对所有 \(z \in k\) ）。）
\end{enumerate}

\begin{enumerate}
\def\labelenumi{\arabic{enumi})}
\setcounter{enumi}{24}
\tightlist
\item
  设 G 为赋有左哈尔测度 \(\nu\) 的局部紧群，\(K \subset G\) 为赋有规范化哈尔测度 \(\mu\) 的 G 的紧子群。
\end{enumerate}

我们用 \(\mathcal{H}(\mathrm{G},\mathrm{K})\) 表示由函数 \(\varphi \in \mathcal{H}(\mathrm{G})\) （满足 \(\varphi (k_1gk_2) = \varphi (g)\) ，对所有 \(g\in \mathrm{G}\) 与每个 \((k_{1},k_{2})\in \mathrm{K}\times \mathrm{K}\) ）所成的空间。

空间 \(\mathcal{H}(\mathrm{G},\mathrm{K})\) 是代数 \(\mathcal{M}_c(\mathrm{G})\) 的一个含幺子代数；我们称 \((\mathrm{G},\mathrm{K})\) 为盖尔范德对，如果代数 \(\mathcal{H}(\mathrm{G},\mathrm{K})\) 是交换的。\(a)\) 映射 \(u^{\mathrm{K}}\) 把 \(f\) 对应到函数 \(f^{\mathrm{K}}\) （由

\[
f ^ {\mathrm{K}} (g) = \int_ {\mathrm{K} \times \mathrm{K}} f (k _ {1} g k _ {2}) d (\mu \otimes \mu) (k _ {1}, k _ {2})
\]

所定义）；它是 \(\mathcal{K}(\mathrm{G})\) 中的一个投影算子，其像为 \(\mathcal{H}(\mathrm{G},\mathrm{K})\) 。

\begin{enumerate}
\def\labelenumi{\alph{enumi})}
\setcounter{enumi}{1}
\tightlist
\item
  设 \(f\in \mathcal{K}(\mathrm{G})\) 。我们有
\end{enumerate}

\[
\int_ {\mathrm{G}} f (g) d \nu (g) = \int_ {\mathrm{G}} f ^ {\mathrm{K}} (g) d \nu (g)
\]

并且 \((f\varphi)^{\mathrm{K}} = f^{\mathrm{K}}\varphi\) 对每个函数 \(\varphi \in \mathcal{H}(\mathrm{G},\mathrm{K})\) 成立。

\begin{enumerate}
\def\labelenumi{\alph{enumi})}
\setcounter{enumi}{2}
\item
  设 \(C \subset G\) 为紧子集。存在函数 \(\varphi \in \mathcal{H}(G, K)\) ，使得 \(\varphi(x) = 1\) 对所有 \(x \in C\) 成立。
\item
  设 (G,K) 为盖尔范德对。群 G 是幺模的。（注意 G 的模属于 \(\mathcal{H}(\mathrm{G},\mathrm{K})\) ，并把 \(\varphi \in \mathcal{H}(\mathrm{G},\mathrm{K})\) 的积分解释为适当选取的卷积在 \(e\) 处的值。）
\item
  假设存在连续对合 \(\theta\colon\mathrm{G}\to\mathrm{G}\) ，使得 \(\mathrm{K}\theta(x)\mathrm{K}=\mathrm{K}x^{-1}\mathrm{K}\) 对所有 \(x\in\mathrm{G}\) 成立。则 (G,K) 是盖尔范德对。（验证映射 \(f\mapsto\nu(f\circ\theta)\) （在 \(\mathcal{K}(\mathrm{G})\) 上）是 G 上的一个左哈尔测度，进而它等于 \(\underline{\nu}\) ；推出 \((f*g)\circ\theta=(f\circ\theta)*(g\circ\theta)\) ，并与公式 \(f*g=\check{g}*f\) 比较。）
\item
  若 G 是紧的，则 \((\mathbf{G} \times \mathbf{G}, \Delta)\) 是盖尔范德对，其中 \(\Delta\) 是 \(G \times G\) 的对角子群，即同态 \(x \mapsto (x, x)\) （从 G 到 \(G \times G\) ）的像。
\end{enumerate}

\begin{enumerate}
\def\labelenumi{\arabic{enumi})}
\setcounter{enumi}{25}
\tightlist
\item
  我们继续沿用上一习题的记号。我们用 \(\varphi_{\mathrm{K}}\) 表示 K 的特征函数。
\end{enumerate}

\begin{enumerate}
\def\labelenumi{\alph{enumi})}
\item
  设 \(\pi\) 为 G 在复希尔伯特空间 E 中的酉表示。\(\mathcal{M}_c(\mathrm{G})\) 在 E 中的表示通过过渡到子空间，定义了代数 \(\mathcal{H}(\mathrm{G},\mathrm{K})\) 在 E 的 K 不变向量所成空间 \(\mathrm{E}^{\mathrm{K}}\) 中的表示。
\item
  若 E 容许循环向量 x，使得 \(x \in E^{K}\) ，则 \(\pi\) 是不可约的。
\item
  假设 \(\pi\) 不可约。若空间 \(E^{K}\) 不约化为 0，则 \(\mathcal{H}(G,K)\) 在 \(E^{K}\) 中的表示是不可约的。
\end{enumerate}

\begin{enumerate}
\def\labelenumi{\alph{enumi})}
\setcounter{enumi}{3}
\tightlist
\item
  设 \(\varphi\) 为 G 上的正定函数，\((\varrho, x)\) 为 \(\varphi\) 的一个循环希尔伯特实现。我们有 \(\varphi \in \mathcal{H}(\mathrm{G}, \mathrm{K})\) ，当且仅当 \(x \in \mathrm{E}^{\mathrm{K}}\) 。
\end{enumerate}

\begin{enumerate}
\def\labelenumi{\arabic{enumi})}
\setcounter{enumi}{26}
\tightlist
\item
  我们继续沿用上一习题的记号。假设 (G, K) 是盖尔范德对。
\end{enumerate}

我们称 \(\varphi\in\mathcal{H}(\mathrm{G},\mathrm{K})\) 是 K-球函数，如果映射

\[
f \mapsto \int_ {\mathrm{G}} f \check {\varphi} d \nu
\]

是代数 \(\mathcal{H}(\mathrm{G},\mathrm{K})\) 的一个非零特征标。

\begin{enumerate}
\def\labelenumi{\alph{enumi})}
\tightlist
\item
  设 \(\varphi \in \mathcal{H}(\mathrm{G},\mathrm{K})\) 。证明下列条件等价：
\end{enumerate}

\begin{enumerate}
\def\labelenumi{(\roman{enumi})}
\item
  函数 \(\varphi\) 是 K-球函数；
\item
  对所有 \(x\) 与 \(y\) （属于 \(\mathrm{G}\) ），我们有
\end{enumerate}

\[
\int_ {\mathrm{K}} \varphi (x k y) d \mu (k) = \varphi (x) \varphi (y);
\]

\begin{enumerate}
\def\labelenumi{(\roman{enumi})}
\setcounter{enumi}{2}
\tightlist
\item
  有 \(\varphi(e)=1\) ，且对每个函数 \(\psi\in\mathcal{H}(\mathrm{G},\mathrm{K})\) ，存在 \(c\in\mathbf{C}\) 使得 \(\psi*\varphi=c\varphi\) 。
\end{enumerate}

\begin{enumerate}
\def\labelenumi{\alph{enumi})}
\setcounter{enumi}{1}
\item
  我们用 \(\mathcal{H}^1 (\mathrm{G},\mathrm{K})\) 表示由 \(f\in \mathrm{L}^1 (\mathrm{G})\) （满足 \(\varrho_{\mathrm{G}}(k_1,k_2)f = f\) ，对所有 \((k_{1},k_{2})\in \mathrm{K}\times \mathrm{K}\) ）所成的空间。它是对合巴拿赫代数 \(\mathrm{L}^1 (\mathrm{G})\) 的一个巴拿赫子代数；它是交换的。
\item
  代数 \(\mathcal{H}(\mathrm{G},\mathrm{K})\) 与 \(\mathcal{H}^1 (\mathrm{G},\mathrm{K})\) 的一个稠密子代数等同。\\
\item
  把 \(\varphi \in \mathcal{H}(\mathrm{G},\mathrm{K})\) 对应到线性型
\end{enumerate}

\[
f \mapsto \int_ {\mathrm{G}} f \check {\varphi} d \nu
\]

的映射（在 \(\mathcal{H}^{1}(\mathrm{G},\mathrm{K})\) 上）定义了从 G 上有界 K-球函数所成的空间到 \(\mathcal{H}^{1}(\mathrm{G},\mathrm{K})\) 的非零特征标所成空间上的一个双射。（注意 \(\mathcal{H}^{1}(\mathrm{G},\mathrm{K})\) 的每个特征标都是连续的且范数 \(\leqslant 1\) ，参见 I, p.~104 的命题 2。）

\begin{enumerate}
\def\labelenumi{\arabic{enumi})}
\setcounter{enumi}{27}
\tightlist
\item
  设 G 为局部紧群，K 为 G 的紧子群。
\end{enumerate}

\begin{enumerate}
\def\labelenumi{\alph{enumi})}
\item
  偶 (G, K) 是盖尔范德对，当且仅当对 G 在复希尔伯特空间 E 中的每个酉表示 \(\pi\) ，\(E^{K}\) 的维数至多为 1。（为证该条件是必要的，证明交换对合巴拿赫代数在希尔伯特空间中的不可约表示是一维的；为证它是充分的，利用 V, p.~454 的定理 4 的推论验证代数 \(\mathcal{H}(G,K)\) 的特征标分离其点。）
\item
  假设 (G, K) 是盖尔范德对。设 \(\varphi \in \mathcal{H}(\mathrm{G},\mathrm{K})\) 。假设 \(\varphi\) 是正定函数且 \(\varphi(e) = 1\) 。证明 \(\varphi\) 是 K-球函数，当且仅当 \(\varphi\) 有一个希尔伯特表示 \((\varrho, x)\) ，使得 \(\varrho\) 不可约。
\item
  假设 \((\mathrm{G},\mathrm{K})\) 是盖尔范德对。G 上正定 K-球函数所成的集合与由不可约表示 \(\pi\in\widehat{G}\) （满足 \(E_{\pi}^{K}\) 非零）所成的集合双射对应。
\end{enumerate}

\begin{enumerate}
\def\labelenumi{\arabic{enumi})}
\setcounter{enumi}{28}
\tightlist
\item
  设 G 为紧群，K 为 G 的闭子群。假设 (G, K) 是盖尔范德对。我们用 \(\nu\) （相应地 \(\mu\) ）表示 G（相应地 K）上的规范化哈尔测度。
\end{enumerate}

设 \(\varphi\) 为 G 上的一个 K-球函数。

\begin{enumerate}
\def\labelenumi{\alph{enumi})}
\item
  我们用 \(E_{\varphi}\) 表示 \(\mathcal{C}(G)\) 中形如 \(\alpha * \varphi\) 的函数所成的子空间，其中 \(\alpha\) 是 G 上有限支集的测度。对每个 \(\psi \in \mathcal{H}(G, K)\) ，自同态 \(v_{\psi}: f \mapsto f * \psi\) （\(\mathcal{C}(G)\) 的）是紧的，且 \(E_{\varphi}\) 含于 \(v_{\psi}\) 的一个特征空间。
\item
  空间 \(E_{\varphi}\) 是有限维的。它与正则表示 \(\gamma_{G}\) 的 G-有限向量所成空间的一个含 \(\varphi\) 的子空间等同。
\item
  \(E_{\varphi}\) 的 K 不变向量所成的空间是一维的，且 G 在 \(E_{\varphi}\) 中的表示是不可约的。
\end{enumerate}

\begin{enumerate}
\def\labelenumi{\alph{enumi})}
\setcounter{enumi}{3}
\tightlist
\item
  设 \((f_i)_{i\in I}\) 为 \(\mathrm{E}_{\varphi}\) 的规范正交基，\(k\colon \mathbf{G}\times \mathbf{G}\to \mathbf{C}\) 为由

\[
k (x, y) = \sum_ {i \in \mathrm{I}} \overline {{f _ {i} (x)}} f _ {i} (y).
\]

定义的函数。有

\[
f (x) = \int_ {\mathrm{G}} k (x, y) f (y) d \nu (y)
\]

对所有 \(f \in \mathrm{E}_{\varphi}\) 与 \(x \in \mathrm{G}\) 成立。

\end{enumerate}

\begin{enumerate}
\def\labelenumi{\alph{enumi})}
\setcounter{enumi}{4}
\item
  有 \(k(x,y)=\dim(\mathrm{E}_{\varphi})\varphi(y^{-1}x)\) 对所有 \((x,y)\in\mathrm{G}\times\mathrm{G}\) 成立。
\item
  函数 \(\varphi\) 是正定型的。
\item
  设 \(\varphi\) 为 G 上的 K-球函数，\((\pi, x)\) 为 \(\varphi\) 的一个希尔伯特实现。表示 \(\pi\) 是不可约的；我们有
\end{enumerate}

\[
\varphi (x) = \int_ {\mathrm{K}} \chi_ {\pi} (x ^ {- 1} k) d \mu (k)
\]

对所有 \(x \in \mathrm{G}\) 成立，且 \(\int_{\mathrm{G}} |\varphi|^2 d\nu = \dim(\pi)^{-1}\) 。

\begin{enumerate}
\def\labelenumi{\arabic{enumi})}
\setcounter{enumi}{29}
\tightlist
\item
  设 \(n \in N\) 。用 G 表示仿射群 \(\mathrm{A}(\mathbf{R}^{n})\) （空间 \(R^{n}\) 的）中由元素 \(g \in \mathrm{A}(\mathbf{R}^{n})\) （其相伴线性映射 \(g_{0}\) 属于 \(\mathbf{SO}(\mathbf{R}^{n})\) ）所成的子群；它是一个实李群（参见 A, II, p. 131 与 LIE, III, p. 102, 例）。我们把 G 的元素记作 \(g = (g_{0}, a)\) ，其中 \(g_{0} \in \mathbf{SO}(\mathbf{R}^{n})\) ，\(a \in R^{n}\) 。
\end{enumerate}

\begin{enumerate}
\def\labelenumi{\alph{enumi})}
\item
  偶 \((\mathbf{G},\mathbf{SO}(\mathbf{R}^n))\) 是盖尔范德对。
\item
  设 \(\varphi \in \mathcal{H}(\mathrm{G},\mathbf{SO}(\mathbf{R}^n))\) 。用 \(f\colon \mathbf{R}^n\to \mathbf{C}\) 表示由 \(f(r) = \varphi (1_{\mathbf{R}^N},g(e_1))\) 定义的函数。映射 \(u\colon \varphi \mapsto f\) 是从 \(\mathcal{H}(\mathrm{G},\mathrm{K})\) 到 \(\mathcal{K}(\mathbf{R}^n)\) 的一个单射复代数态射；其像是 \(\mathcal{K}(\mathbf{R}^n)\) 中由函数 \(f\) （满足 \(f(g_0(x)) = f(x)\) ，对所有 \(g_0\in \mathbf{SO}(\mathbf{R}^n)\) 与每个 \(x\in \mathbf{R}^n\) ）组成的子空间。
\item
  我们用 \(\Delta\) 表示 \(\mathbf{R}^n\) 上的拉普拉斯算子。函数 \(\varphi \in \mathcal{K}(\mathrm{G})\) 是球函数，当且仅当函数 \(f = u(\varphi)\) 是 \(\mathrm{C}^{\infty}\) 类的、满足 \(f(0) = 1\) ，且存在 \(\lambda \in \mathbf{C}\) 使得 \(\Delta f = \lambda f\) 。（利用 \(\Delta(f \circ g_0) = (\Delta f) \circ g_0\) 对所有 \(g_0 \in \mathbf{SO}(\mathbf{R}^n)\) 成立这一事实。）
\end{enumerate}

\begin{enumerate}
\def\labelenumi{\arabic{enumi})}
\setcounter{enumi}{30}
\tightlist
\item
  设 \(n\) 为整数 \(\geqslant 1\) 。我们用 N 表示 \(\mathbf{R}^n\) 上的欧几里得范数，用 \(\lambda_n\) 表示 \(\mathbf{R}^n\) 上的勒贝格测度。我们用 \(\omega_n\) 表示球面 \(\mathbf{S}_n \subset \mathbf{R}^{n+1}\) 上唯一的、在正交群 \(\mathbf{O}(n+1, \mathbf{R})\) 的作用下不变的、总质量为 1 的正测度（参见 INT, VIII, p.~116, 习题 8）。
\end{enumerate}

我们把 \(R^{n}\) 与射影空间 \(\mathbf{P}_{n}(\mathbf{R})\) 中射影超平面 H 的余集等同起来（参见 TG, VI, p.~15, 命题 5），并把 \(R^{n}\) 上的测度解释为 \(\mathbf{P}_{n}(\mathbf{R})\) 上支集含于 H 的余集中的测度。

对每个 \(g \in \mathbf{GL}(n+1, \mathbf{R})\) ，我们用 \(\widetilde{g}\) 表示由 g 定义的 \(\mathbf{P}_{n}(\mathbf{R})\) 的同胚。

\begin{enumerate}
\def\labelenumi{\alph{enumi})}
\tightlist
\item
  正测度 \(c_{n}\) （在 \(\mathbf{R}^n\) 上）由
\end{enumerate}

\[
c _ {n} = \pi^ {- (n + 1) / 2} \Gamma \Big (\frac {n + 1}{2} \Big) (1 + \mathrm{N} ^ {2}) ^ {- (n + 1) / 2} \lambda_ {n}
\]

定义的（“柯西测度”）是有界的且总质量为 1。若 \(n = 1\) ，它就是参数 1 的稳定测度（参见习题 V, p.~496）。对每个射影超平面 \(\mathrm{H}'\) （\(\mathbf{P}_n(\mathbf{R})\) 的），我们有 \(c_n(\mathrm{H}') = 0\) 。

\begin{enumerate}
\def\labelenumi{\alph{enumi})}
\setcounter{enumi}{1}
\item
  从 \(\mathbf{R}^{n+1} - \{0\}\) 到 \(\mathbf{P}_n(\mathbf{R})\) 的典范映射通过过渡到子空间，定义了连续满射 \(\sigma: \mathbf{S}_n - (\mathbf{R}^n \times \{0\}) \to \mathbf{R}^n\) 。我们有 \(\sigma(\omega_n) = c_n\) 。
\item
  设 \(g \in \mathbf{GL}(n+1, \mathbf{R})\) 。存在 \(a \in \mathbf{GL}(n, \mathbf{R})\) 与 \(b \in R^n\) ，使得像测度 \(\widetilde{g}(c_n)\) 等于 \(c_n\) 在仿射映射 \(x \mapsto a(x) + b\) 下的像。
\end{enumerate}

在下面，我们打算证明这一断言的一个逆命题。

\begin{enumerate}
\def\labelenumi{\alph{enumi})}
\setcounter{enumi}{3}
\tightlist
\item
  设 \(\mu\) 为 \(\mathbf{P}_{n}(\mathbf{R})\) 上的有界正测度，使得每个射影超平面都是 \(\mu\) -可忽略的。由 \(g \in \mathbf{GL}(n+1, \mathbf{R})\) （满足 \(g(\mu) = \mu\) 且 \(|\det(g)| = 1\) ）所成的集合是 \(\mathbf{GL}(n+1, \mathbf{R})\) 的一个紧子群。

\end{enumerate}

\begin{enumerate}
\def\labelenumi{\alph{enumi})}
\setcounter{enumi}{4}
\item
  设 \(\mu\) 为 \(\mathbf{P}_n(\mathbf{R})\) 上总质量为 1 的有界正测度，使得对每个 \(g \in \mathbf{GL}(n + 1, \mathbf{R})\) ，存在仿射映射 \(f\) （\(\mathbf{R}^n\) 的）满足 \(g(\mu) = f(\mu)\) 。每个射影超平面都是 \(\mu\) -可忽略的。
\item
  设 \(K_{\mu}\) 为由 \(g \in \mathbf{GL}(n+1, \mathbf{R})\) （满足 \(g(\mu) = \mu\) ）所成的子群。存在 \(h \in \mathbf{GL}(n+1, \mathbf{R})\) ，使得 \(K_{\mu} \subset h \mathbf{O}(n+1, \mathbf{R}) h^{-1}\) 。
\item
  设 \(K = h^{-1}K_{\mu}h\) 。我们有 \(\mathbf{O}(n + 1, \mathbf{R}) = \mathbf{G} \cdot \mathbf{K}\) ，其中 G 是 \(\mathbf{O}(n + 1, \mathbf{R})\) 中由矩阵
\end{enumerate}

\[
\left( \begin{array}{c c} x & 0 \\ 0 & 1 \end{array} \right)
\]

（\(x\in \mathbf{O}(n,\mathbf{R})\) ）组成的子群

\begin{enumerate}
\def\labelenumi{\alph{enumi})}
\setcounter{enumi}{7}
\tightlist
\item
  测度 \(h^{-1}(\mu)\) 与 \(c_{n}\) 重合。（注意群 K 传递地作用于 \(S_{n}\) ，且像测度 \(\sigma(h^{-1}(\mu))\) 是 \(S_{n}\) 上的一个 K 不变测度。）
\end{enumerate}

\begin{enumerate}
\def\labelenumi{\arabic{enumi})}
\setcounter{enumi}{31}
\tightlist
\item
  设 G 为局部紧群。假设存在 G 的一个既紧又开的子群 H。我们用 \(\mu\) 表示 G 上满足 \(\mu(\mathrm{H}) = 1\) 的唯一的左哈尔测度，用 \(\varphi\) 表示 H 的特征函数。
\end{enumerate}

\begin{enumerate}
\def\labelenumi{\alph{enumi})}
\item
  设 \(\pi\) 为 G 在复希尔伯特空间 E 中的一个平方可积表示。证明自同态 \(\pi(\varphi)\) 是希尔伯特–施密特自同态。
\item
  空间 \(E^{H}\) 是有限维的，且满足
\end{enumerate}

\[
\dim (\mathrm{E} ^ {\mathrm{H}}) \leqslant \frac {1}{c (\pi)}
\]

其中 \(c(\pi)\) 是 \(\pi\) 相对于 \(\{e\}\) 的形式次数。（利用 V, p.~487 的习题 6。）

\begin{enumerate}
\def\labelenumi{\alph{enumi})}
\setcounter{enumi}{2}
\tightlist
\item
  离散无限群没有平方可积表示；若 \(n \geqslant 2\) 为整数，则群 \(\mathbf{SL}(n, \mathbf{Z})\) 没有模中心的平方可积表示。
\end{enumerate}

\specialchapter{历史注记}

（第 I 至 V 章）

特征值与傅里叶级数的概念在 1830 年已为人熟知，并贯穿十九世纪的许多分析领域（参见~ÉHM, pp.~114, 260, 275）。但谱理论与调和分析——按我们在这本书中所理解的意义——直到大约一个世纪之后才开始具有其现在的形式；这与它们在 1930 年前后同群表示研究的汇合、随后在 1940 年前后同赋范代数研究的汇合是同时发生的。

我们在这里仅限于勾画它们从 1906 年（希尔伯特在其积分方程工作中引入“连续谱”概念之时）到 1945–1950 年间演化的主要脉络；在此期间，这里所陈述的理论基本上获得了保持至今的形式。

\histitem{连续谱的发现}

在关于拓扑向量空间的注记中（参见 ÉHM, pp. 262–263），我们描述了 I. 弗雷德霍姆如何在 1900 年前后——继承 C. 诺伊曼、V. 沃尔泰拉与 H. 庞加莱的工作——开始考虑方程

\[
u (x) - \lambda \int_ {\mathrm{I}} \mathrm{K} (x, y) u (y) d y = f (x)\tag{1}
\]

其中未知函数为 \(u\colon I\to\mathbf{C}\) ，这里区间 \(I\subset\mathbf{R}\) 是紧的，且核 \(K\) 在 \(I\times I\) 上连续。他巧妙地使用“无穷行列式”——其基础是庞加莱与 H. 冯·科赫的思想——得到了对复变量 \(\lambda\) 的依赖为亚纯的解族，进而证明了关于解的存在的著名“择一定理”（参见 III, p.~81, 定理 5）。但弗雷德霍姆随即把他的结果特殊化到 \(\lambda = -1\) 的情形，而没有利用方程 (1) 是一个特征值问题 \(^{(1)}\) 这一事实。

当核 K 对称时，我们还看到希尔伯特在其关于积分方程的第一篇论文 [31] 中如何认识到弗雷德霍姆问题与求二次型的“主轴”之间的亲缘关系。从核 K 的有限“截段”（离散化）的这种约化出发，他通过取极限得到公式

\[
\int_ {\mathrm{I}} \mathrm{K} (s, t) x (s) x (t) d t = \sum_ {n = 1} ^ {\infty} \frac {1}{\lambda_ {n}} \int_ {\mathrm{I}} \varphi_ {n} (s) x (s) d s,\tag{2}
\]

其中 \(\lambda_{n}\) 是核 K 的特征值，\(\varphi_{n}\) 构成相伴特征向量 \(^{(2)}\) 的规范正交系，并且只要 x 在 I 上平方可积等式就成立（参见~ÉHM, p.~264）。

希尔伯特面临的主要障碍涉及从有限和过渡到公式 (2) 的积分。他从这一公式推出了特征值的存在性——不过，他受有限维中确定特征值的经典方法的启发（参见~IV, p.~153 第 4 目），为后者给出了一个与取极限无关的变分刻画。E. 施密特在 1905 年将证明，借鉴 H. 施瓦茨的工作，如何不经过弗雷德霍姆与希尔伯特的方法所依赖的“无穷行列式”而得到 (2)（参见~ÉHM, p.~265）。

但在 1906 年，希尔伯特 [32] 已转向约化二次型

\[
\mathrm{B} (x, x) = \sum_ {p, q = 0} ^ {\infty} b _ {p q} x _ {p} x _ {q}, \qquad x = (x _ {p}) _ {p \in \mathbf {N}}
\]

这一更一般的问题，二次型定义在 \(E = \ell_{\mathbf{C}}^{2}(\mathbf{N})\) 上。他心里显然想着：通过过渡到 E 的规范正交基下的系数，问题 (1) 可化为求“无穷线性方程组”

\[
x _ {p} + \sum_ {q = 0} ^ {\infty} b _ {p q} x _ {q} = f _ {p} \quad (p \in \mathbf {N}).
\]

的解。然而他注意到，弗雷德霍姆问题所处理的是非常特殊的二次型：\(\mathrm{B}(x_{n}, x_{n})\) 趋于 \(\mathrm{B}(x, x)\) （只要 \(x_{n}\) 弱收敛于 x）；他称之为“全连续的”（vollstetig）。

希尔伯特强调它们与 E 上有界的二次型之间的本质差别，并决定着手对后者作更一般的约化。他的方法依旧是从“有限截段”

\[
(x _ {0}, \ldots , x _ {n}) \mapsto \sum_ {p, q = 0} ^ {n} b _ {p q} x _ {p} x _ {q}
\]

的约化出发，然后让 n 趋于无穷。为完成这一取极限的过程，他利用了 T. 斯蒂尔杰斯在其关于连分数的工作 [71] 中引入的积分思想。他证明每个有界二次型经过变量的一个正交变换后都可以写成

\[
\sum_ {p \in {\bf N}} \frac {1}{\lambda_ {p}} x _ {p} ^ {2} + \int_ {s \in {\bf R}} \frac {1}{s} d \sigma (s, x).\tag{3}
\]

的形式。在这个公式中，\(\lambda_{p}\) 是 B 的特征值，\(x\mapsto\sigma(s,x)\) （对固定的 \(s\in\mathbf{R}\) ）是 E 上正定的分离二次型，而 \(s\mapsto\sigma(s,x)\) \(^{(3)}\) （对固定的 \(x\in\mathrm{E}\) ）是连续的递增函数，它在 \(-\infty\) 处趋于 0，其极限值 \(\|x\|^2\) 在 \(+\infty\) 处达到（斯蒂尔杰斯记号参见 ÉHM, p. 284）。

用 S 表示 \(\overline{R}\) 中这样的点所成的集合：不存在使 \(s \mapsto \sigma(s, x)\) 对所有 x 都保持常数的邻域；于是 (3) 中的积分自然可以在 S 上取。希尔伯特把集合 S 称为“连续谱”（Streckenspektrum），把 B 的特征值所成的集合称为“点谱”（Punktspektrum），把两者的并称为“谱”（Spektrum）。这个来自光学术语的词，曾在 1897 年被 W. 维尔丁格 [91] 用于带周期系数的微分方程的研究。

希尔伯特很快便从他的“谱”方法中为经典分析的若干问题汲取了深刻的革新，这些问题涉及常微分方程（特别是 Sturm–Liouville 型）、偏微分方程或复变量函数的研究。我们在这里不复述二十世纪头二十五年间随之而来的丰富成果，而请读者参看 E. 海林格与 O. 托普利茨的综述 [29]。不过我们仍要提到其中一些关于积分方程的工作，在它们之中可以看到抽象理论后来发展的萌芽。

对希尔伯特及其学生而言——他们忠实于线性代数中的德国传统——模型始终是主轴定理；因此所研究的总是双线性型或二次型的谱，特别是在“希尔伯特”空间 \(E = \ell_{\mathbf{C}}^{2}(\mathbf{N})\) 上。可以把 E 上任何连续双线性型 B 对应于由 \(\langle Ax, y \rangle = B(x, y)\) 刻画的 E 的连续自同态 A，这一事实为希尔伯特学派所知，尤其是在 E. 施密特、M. 弗雷歇与 F. 里斯 1907–1908 年的工作之后。但正是 F. 里斯在 1913 年一本令人赞叹的书 [59] 中指出了在此背景下把自同态置于首位的益处。

里斯给出了 \(\mathcal{L}(\mathrm{E})\) 的元素 A 的谱的现代定义 [59, p.~139]，指出它是一个有界闭集，并证明预解式 \(\lambda \mapsto (\lambda 1_{\mathrm{E}} - \mathrm{A})^{-1}\) 在 A 的谱的余集上是全纯的。尤其重要的是，他赋予 \(\mathcal{L}(\mathrm{E})\) 的代数结构以中心地位：借鉴沃尔泰拉的工作，他对对称自同态发展了泛函演算的第一个版本，使他能够简单地重新得到希尔伯特与施密特的约化结果。若 A 是 \(\mathcal{L}(\mathrm{E})\) 中的对称元素，里斯说明如何定义 \(f(\mathrm{A})\) ——对每个函数 \(f\) ，它在 \(\mathbf{R}\) 上（下或上）半连续、取值于 \(\mathbf{R}\) ——并指出 \(f \mapsto f(\mathrm{A})\) 就是我们今天所说的代数态射 [59, pp.~129--130]。把这一思想应用于区间 \([\xi, +\infty[\) 的特征函数，就对每个 \(\xi\) 得到 E 的一个对称自同态 \(\mathrm{A}_{\xi}\) ；借助斯蒂尔杰斯积分，里斯随后证明了希尔伯特关系式 (3) 的一个版本：

\[
\langle \mathrm{A} x, y \rangle = \int_ {\mathbf {R}} \xi d \langle \mathrm{A} _ {\xi} x, y \rangle
\]

对 \(x, y \in \mathrm{E}\) 成立，他甚至把这个等式写成

\[
\mathrm{A} = \int_ {\mathbf {R}} \xi d \mathrm{A} _ {\xi}.
\]

于是 A 的谱是这样的值 \(\xi_{0}\) 所成的集合：\(A_{\xi}\) 在 \(\xi_{0}\) 的任何邻域内都不保持常数；而点谱是 \(\xi \mapsto A_{\xi}\) 的不连续点所成的集合。

在以理论的应用作结之前不久，里斯指出 [59, p.~146]，当自同态 A 来自希尔伯特的某个“全连续”双线性型时，其结果具有特别简单的形式：谱约化为特征值所成之集（可能还须加上 0），特征值排成一个趋于 0 的序列，且每个非零特征值的重数是有限的（III, p.~90, 命题 5）。

\histitem{紧算子}

不到三年之后，里斯在其关于泛函方程的论文 [60] 中大大强化了这些评注；该文以历经一个多世纪仍未减损的清晰度，发展了巴拿赫空间上紧算子谱理论的要点。这篇文本宣称的目标是，通过考虑赋有一致收敛拓扑的空间 \(E = \mathcal{C}(I)\) ，对紧区间 I 上的方程 (1) 重新表述弗雷德霍姆理论。但正如我们在早先的一条注记中已经指出的（ÉHM, p.~268），里斯清楚地意识到他的结果适用于任何巴拿赫空间——这一概念要到下一个十年才被引入。

受弗雷歇关于一般拓扑的思想启发，里斯这时考虑 E 的把每个有界序列映为相对紧序列的自同态。与希尔伯特一样，他仍称其为“全连续的”，并指出在 \(\mathrm{E} = \ell_{\mathbf{C}}^{2}(\mathbf{N})\) 的情形，这一新概念等价于他在 1913 年的书中 \(^{(4)}\) 用过的全连续性概念。

里斯研究形如 \(B = 1_{E} - A\) 的自同态，其中 A 是 E 的全连续自同态，并从“局部紧的赋范向量空间必是有限维的”这一——似乎正是为该场合而发现的——事实导出它们的所有谱性质。他容易地证明 B 的核是有限维的，且它的像是闭的、余维数有限。如同 E. 韦尔 [87] 在有限维中所做的那样，他接着考虑迭代 \(B^{k}\) 。他证明它们的核所成的序列与像所成的序列都是平稳的，然后引入我们今天所谓的 B 的零空间与余零空间，并高效地证明 E 是它们的拓扑直和。把这一观察应用于自同态 \(B_{\lambda} = 1_{E} - \lambda A\) （\(\lambda \in C\) ），他便毫不困难地推出巴拿赫空间上全连续自同态的谱性质，而且这几乎是最终的形式。

关于这些算子的谱的主要结果中，似乎只有那些需要对函数空间概念、特别是对偶概念有更成熟认识的结果才超出了里斯的范围。例如，1920 年代末，T. 希尔德布兰特 [33] 与 J. 绍德尔 [63] 将（在此时已充分确立的巴拿赫空间框架内）证明全连续自同态的伴随自同态是全连续的。

此外，里斯的全连续性概念仍然依赖于序列的使用；巴拿赫 [4] 在其 1932 年的书中摆脱了这一点，转而考虑把每个有界集映为相对紧集的映射。“紧”线性映射这一术语要到二十世纪下半叶才逐渐确立，大概是随着 E. 希尔关于算子半群的书 [34] 而流行开来。

1930 年代提出的若干问题长期悬而未决。在这方面，我们注意到 1973 年对“逼近问题”的解决 [18]，巴拿赫与 S. 马祖尔曾在 1932 与 1938 年使这一问题闻名（见 III, p.~16 的注记 6 与 III, p.~112 的习题 25）。

与此同时，洛莫诺索夫 [40] 证明了：在局部凸空间中，与一个非零紧自同态交换的自同态存在非平凡的不变闭子空间（III, p.~13 的推论 2）。这一结果是对任意连续自同态是否存在此类子空间的问题（“不变子空间问题”）最重大的进展之一。该问题在可分希尔伯特空间的连续自同态情形至今仍未解决；P. 恩夫洛 [19] 在 1987 年构造了第一个不具有非平凡不变闭子空间的巴拿赫空间自同态的例子。
\histitem{弗雷德霍姆指标与扰动}

指标这一概念于1920年出现在弗里茨·诺特关于积分方程的工作中。在1904年国际数学家大会的演讲中，希尔伯特考察了黎曼在势理论中提出的一个问题：给定平面内一个有界开域 \(\Omega\) ，其边界为光滑简单闭曲线 \(\Gamma\) ，以及三个函数 \(a\) 、\(b\) 、\(f\) ，它们定义在 \(\Gamma\) 上，求一个函数 \(z: \overline{\Omega} \to \mathbf{C}\) ，它在 \(\Omega\) 上全纯、在 \(\overline{\Omega}\) 上连续，并满足

\[
a \mathcal {R} (z) + b \mathcal {I} (z) = f \quad \text { on } \Gamma .
\]

将 \(\Gamma\) 用实数 \(s \in [0, 2\pi]\) 参数化后，希尔伯特证明该问题归结为求解关于 \(\varphi = \mathcal{R}(z)\) 的一个“具有奇异核”的积分方程，即

\[
a (s) \varphi (s) + \int_ {0} ^ {2 \pi} \mathrm{K} (s, t) \varphi (t) d t = f (s).\tag{4}
\]

这里，核 \(\mathrm{K}(s, t)\) 对某个连续函数 A 具有形式 \(b(s) \cot\left(\frac{t-s}{2}\right) + \mathrm{A}(s, t)\) 。因此它沿对角线奇异，上述积分须按柯西主值的意义理解。

诺特 {[}53{]} 指出，与满足弗雷德霍姆择一定理的积分方程不同，对于上述形式的核 K 和非零右端项 \(f\) ，(4) 可能有非平凡解族。他证明解空间由整数

\[
n = \frac {1}{2 \pi} \int_ {\Gamma} d (\log (a - i b)).\tag{5}
\]

若 \(n < 0\) ，则不存在非平凡解；若 \(n \geqslant 0\) ，方程(4)有一个依赖于 2n 个参数的解族。诺特使用“指标”（Index）一词称呼整数 n，并在其中认出一个绕数；在复变函数研究中，这一概念自十九世纪末以来已属经典。

后来的工作继续研究需要扩展弗雷德霍姆理论的方程实例。不过，在弗雷德霍姆映射的概念被系统提炼出来之前，还要再过二十多年；与此同时，对紧算子的扰动的研究也日趋成熟。

1909年，H. 外尔 {[}81{]} 证明，若定义在 \(\ell_{\mathbf{C}}^{2}(\mathbf{N})\) 上的两个有界二次型之差是完全连续的，则它们的本质谱相同（参见 III，第89页，定理2）。此外，里斯在1916年的回忆录中（当然没有使用现代语言）注意到，巴拿赫空间 E 上的紧算子构成 \(\mathcal{L}(\mathrm{E})\) 的双边理想。然而，似乎直到20世纪40年代才有人沿着这两条评注所指示的道路探索。1941年，J. 卡尔金受 J. 冯·诺伊曼关于算子代数的工作启发（我们很快将有机会讨论，见第537页），指出在希尔伯特空间框架下研究模此理想的同余关系的意义 {[}10{]}。他说明了如今以自己名字命名的代数的结构（参见 III，第75页）如何使人能够简单地恢复外尔的结果。

与此同时，大约1942年，J. 迪厄多内 {[}15{]} 对研究赋范空间之间的连续线性映射 \(u\) 的扰动产生兴趣，所考虑的情形是扰动按算子范数的意义“很小”。他只研究 \(u\) 是严格态射且其核有限维或有限余维的情形；他证明 \(u\) 的任意“微小扰动”也具有同样性质，并说明对于弗雷德霍姆映射（参见 III，第40页的 \(cf. n^{\circ} 2\) ），指标局部恒定（III，第58页定理1），III，第55页§4呈现的若干结果也如此。

这两种思想在1950年前后汇合，此时 F. 阿特金森 {[}3{]}、I. 戈赫伯格 {[}26{]}、S. 米赫林 {[}46{]} 和 B. 尤德 {[}93{]} 研究紧算子的扰动。他们明确指出了它与诺特指标的联系，并提出弗雷德霍姆映射 {[}3，第8页{]} 和里斯映射 {[}93，第5节{]} 的概念。随后十年中，人们对后者展开了系统研究；本书第三章呈现的结果在这一时期大致取得了今天的形式。

这些理论通常在巴拿赫空间框架中展开；然而，各种应用促使人们将其推广到更广泛的拓扑向量空间类。因此，弗雷歇空间的情形于1954年自然出现在 H. 嘉当和 J-P. 塞尔的工作中 {[}13{]}，该工作证明紧复解析流形取值于相干层的上同调是有限维的，其证明调用了 L. 施瓦茨证明的变形结果 {[}66{]}（参见 III，第73页定理2）。

在公式(5)的情形，若注意到右端（绕数）在同伦下不变，指标在变形下的不变性便得到充分说明。这样的指标拓扑表达式很快因微分几何中产生的弗雷德霍姆应用而闻名。若 D 是紧 \(C^{\infty}\) 流形 X 上的微分算子（参见 VAR，§14），且 D 是“椭圆的”（参见{[}2，第1节{]}），则将 D 延拓到适当的 Sobolev 空间定义一个弗雷德霍姆映射。寻求一个利用 X 的“特征类”给出其指标的拓扑表达式的公式——这一研究由 I. 盖尔范德约1960年的工作等所提出 {[}23，第I卷，第65页{]}——将在1963年导向 M. 阿蒂亚和 I. 辛格的“指标定理” {[}2{]}。后者将引发分析、拓扑和微分几何交汇处的非凡发展。我们无法在此描述他们开辟的广阔领域；这些领域至今仍极其富有活力，并且不断从弗雷德霍姆理论汲取营养。

\histitem{部分算子与谱定理}

一般的希尔伯特空间谱理论在很大程度上是为解决量子力学的数学基础问题而创立的，尤其是在冯·诺伊曼的推动下。

关于量子理论直至 W. 海森堡那篇引入量子力学的奠基性论文 {[}28{]}的发展详述，参见 M. 贾默的工作 {[}35{]}。1924年夏发表该文与1929年发表文章 {[}50{]} 相距不到五年；在后者中，冯·诺伊曼以极为清晰严谨的方式呈现了希尔伯特空间上部分（常称“无界”）埃尔米特算子的全部基本结果。

通过给出希尔伯特空间的第一种公理化表述 {[}49{]}，冯·诺伊曼完成了一个酝酿已久的思想（参见 ÉHM，第267页），并阐明了施密特、弗雷歇、菲舍尔和里斯在1910年前已熟知的序列空间或函数空间之间的各种同构。这种抽象程度的提升使冯·诺伊曼在概念上取得重大进展。

部分算子的思想便是如此；它似乎仍以未成形的形态盘旋在希尔伯特之后重新研究弗雷德霍姆方程的几部著作之上（见上文第1号）。希尔伯特注意到，许多 Sturm--Liouville 型微分方程，特别是在无界区间情形，可化为弗雷德霍姆方程；但是要应用积分方法，必须放宽对核施加的条件，并离开弗雷德霍姆、希尔伯特和施密特的结果适用的框架。多年来，人们研究了越来越多的奇异核，直到开始考虑这样的方程：其左端只有在未知量属于 L²(I) 的子空间时才有意义。在这方面必须提到 E. 希尔布 {[}30{]} 1908年的工作、H. 外尔关于 Sturm--Liouville 理论的1910年工作 {[}82{]}，以及最重要的 T. 卡勒曼 {[}11{]} 1923年的工作，后者在非常一般的框架中展开。

通过统一这些问题，冯·诺伊曼使这些经典工作呈现全新面貌。他尤其强调对称算子与自伴算子之间的根本区别，并用抽象“边界条件”来解释对称算子的自伴扩张，从而推广了在 \(R^{n}\) 的有界开集上拉普拉斯算子的狄利克雷问题和诺伊曼问题中出现的众所周知的边界条件。

与此同时，冯·诺伊曼 {[}51{]} 提出了正规算子 \(^{(5)}\) 的概念，将其作为谱理论的自然框架，并强调由算子及其伴随生成的交换代数的作用。

依托这一精确的数学形式主义——它还揭示了海森堡与薛定谔两种观点的等价性——以及对实验步骤和结果的哥本哈根诠释，非相对论量子力学从此达到一种此后几乎未有变化的精致程度；尽管这套物理理论的结果惊人、甚至看似悖论，其有效性自那以后仍不断得到实验确认，而且精度非凡。

从形式观点看，谱理论的基础由冯·诺伊曼奠定。与此同时，在后者首次提出这些思想之后，M. 斯通于1928年至1930年间进行了类似研究，并获得密切相关的结果。斯通于1932年发表的已知结果的清晰完整表述 {[}72{]}，对希尔伯特谱理论的传播产生了重要影响。

然而，表述上的其他改进也对这些结果的传播很重要。事实上，冯·诺伊曼给出的谱定理（参见 IV，第266页定理1）长期以来被认为非常困难，主要是因为它是在向量值测度框架中表述的。

后来，P. 哈尔莫斯 {[}27{]} 强调通过 \(L^{2}\) 空间上的乘法算子解释谱定理，从而揭示了它在有界算子情形下本质上初等的一面。最近，S. 沃罗诺维奇（在希尔伯特模上的正规算子这一稍有不同的 C*-代数框架中 {[}92{]}）通过引入

“有界化”（参见 IV，第265页第2号），使得可以像处理自伴算子一样简单地处理正规部分算子。

此外，若 D 是 \(R^{n}\) 的开集 U 上的形式对称标量微分算子，则研究 D 在 \(L^{2}(U)\) 上的自伴扩张自然与研究 U 上的广义函数相联系（拉普拉斯算子的情形参见 IV，第242页第6号）。关于后者的历史，参见 J. 迪厄多内的著作 {[}16，第VII章{]}——这是一段曲折的历史，紧接在我们刚才述及的时期之后展开，其中 S. 索伯列夫于1936年 {[}69{]} 和 L. 施瓦茨十年后的决定性贡献值得一提 {[}65{]}。

从纯数学观点看，谱理论与黎曼几何以一种特别富有成果的方式相互作用。因此，M. 卡茨于1966年重新讨论 S. 博赫纳提出的问题，给出了一种引人注目且极其生动的表述 {[}36{]}：确定一个黎曼流形的哪些几何性质可以从拉普拉斯算子的谱性质推出。物理学家 A. 舒斯特于1882年已经预见到这个问题 {[}64{]}：要弄清振动系统发出的不同音调，在某些特殊情形下可能可以解决，也可能不行；但是若要解决逆问题，根据钟能够发出的声音来确定钟的形状，即便是最有技巧的数学家也会束手无策。（“确定振动系统发出的谐波，在某些特殊情形下可能可以解决，也可能不行；但如果要解决逆问题、从钟能够发出的声音确定其形状，最有技巧的数学家也会感到困惑。”）。

\histitem{谐波分析与群论的联系}

我们在《积分学》一书的历史注记中（参见 ÉHM，第275页）描述过，与热方程有关的问题如何引导傅里叶提出把任意函数表示为三角函数之和这一革命性思想。

我们无法在此追溯整个十九世纪傅里叶级数和积分的发展，也无法追溯该理论提出的问题对分析演变的深刻影响——其影响甚至颠覆了实数的观念，并促成集合论的诞生（参见 ÉHM，第42页）。但为了理解本书所阐述思想的形式，有必要说明谐波分析在1925年以后如何同群论和希尔伯特谱理论发生联系。

傅里叶思想与 R/Z 和 R 的群结构之间联系的认识，以及将这些思想推广到其他群，来得相当晚。

在有限群及其特征标理论的早期阶段，自然可以辨认出一些如今看来包含有限傅里叶分析雏形的结果。例如，Z/qZ 的可逆元群的对偶，以及表达一个元素的特征函数为特征标线性组合的正交关系，已隐约出现在1837年狄利克雷的文章 {[}17{]} 中；该文证明当 a 与 q 互素时，存在无穷多个素数 \(p \equiv a \mod q\) 。他借助高斯的思想；高斯在研究具有整数系数及固定判别式的二元二次型时引入了“特征标”一词（参见 {[}22，第230节{]}）。

但高斯和狄利克雷都没有使用群的语言。戴德金在1879年讲述这些工作时，才认识到它们潜在的作用，并为交换群定义了特征标的概念。H. 韦伯大大扩充了戴德金的评述：先是1882年 {[}77{]}，然后是1886年 {[}78，第112页{]}；他在那里引入有限阿贝尔群 A 的对偶群，并注意到后者同构于 A。他当时的目的是构造具有给定伽罗瓦群的域 Q 的阿贝尔扩张，并未考虑这类群上的谐波分析问题。

在非交换群情形，有限群表示理论的基础在20世纪初前后，继 G. 弗罗贝尼乌斯、W. 伯恩赛德和 I. 舒尔的工作之后得到坚实奠定（参见 ÉHM，第154页）。1905年以后，特征标的正交关系在理论组织中发挥关键作用已很清楚。但它们与谐波分析的亲缘关系仍隐藏不显。

我们应归功于赫尔曼·外尔将这些领域联系起来：他构想了紧群表示的一般理论。在1925年至1927年间发表的几部重要著作中，他认识到这类理论与傅里叶分析及希尔伯特谱理论之间的联系，并为许多后来的发现奠定基础。

是舒尔的一封信把外尔引向这条道路的，缘于二人对不变量理论的共同兴趣。我们在关于 Haar 测度的注记（ÉHM，第289–291页）中描述过，A. 胡尔维茨早在1897年就已构想利用“不变积分”来确定在正交群下不变的 \(R^{n}\) 上多项式。舒尔似乎直到1924年才知道胡尔维茨的结果；那时他突然理解到，不变积分可用于把自己为有限群建立的特征标理论和正交关系推广到群 \(\mathbf{O}(n)\)。外尔立即看出，舒尔的方法与埃利·嘉当的工作相结合，使得可以构造紧半单李群的不可约表示。外尔将嘉当和舒尔的思想融为一体的三篇论文发表于1925年 {[}83{]}，其中已经包含 LIE, IX, §6–7 的基本结果；参见 ÉHM，第328–330页。

但在次年一篇与学生 F. 彼得合写且同样著名的文章 {[}54{]} 中，外尔才揭示紧群表示论、傅里叶分析和希尔伯特谱理论之间的联系。这篇不足二十页的文章孕育了许多未来的发展。

彼得和外尔摆脱了对所研究群的结构所作的任何代数假设，只假定 G 是一个配备不变测度的紧拓扑群。对连通李群而言，这种测度的存在是显然的；不久之后 A. Haar 将在一般情形建立它，部分动机正来自此处呈现的结果（参见 ÉHM，第291页）。

他们研究的出发点是矩阵系数的正交性。舒尔注意到，对每个不可约表示等价类选取一个代表 \(\pi \in \widehat{\mathrm{G}}\) 、\(\pi\) 的空间 E 上一个 G-不变内积以及 E 的一个标准正交基，就得到一族在 \(\mathrm{L}^2 (\mathrm{G})\) 中正交归一的矩阵系数。彼得和外尔计划证明这族函数是完备的。

在有限群情形，弗罗贝尼乌斯曾通过研究群代数 C{[}G{]} 证明了相应结果。彼得和外尔接着考虑 G 上连续函数空间 \(\mathcal{C}(\mathrm{G})\) ，并为其赋予卷积乘法。这似乎是首次将卷积乘法作为抽象运算、明确地同群结构联系起来（参见 ÉHM，第295页）。此外，彼得和外尔把 \(\mathcal{C}(\mathrm{G})\) 的一个元素称为 Gruppenzahl（“群数”），并以 xy 表示两个 Gruppenzahlen 的（卷积）乘积。他们还为 \(\mathcal{C}(\mathrm{G})\) 配备对合 \(f \mapsto \widetilde{f}\) ，其中 \(\widetilde{f}(g) = \overline{f(g^{-1})}\) ，对于 \(g \in G\) 。

在引入对合代数 \(\mathcal{C}(\mathrm{G})\) 后，彼得和外尔的第一个观察是，每个 G 的酉表示 \(\pi\) 都给出一个表示 \(f \mapsto \pi(f)\) ，其代数为 \(\mathcal{C}(\mathrm{G})\) （参见 V，第400页）。他们明确给出了 \(\mathcal{C}(\mathrm{G})\) 的埃尔米特元素概念，也（隐含地）给出了正元素概念，为应用希尔伯特空间技巧打开道路。

彼得和外尔不再犹豫，称算子 \(\pi(f)\) 为 \(f\) 的“傅里叶系数”，并通过指出舒尔正交关系蕴含贝塞尔不等式，继续与傅里叶理论作平行比较

\[
\sum_ {\pi \in \widehat {\mathrm{G}}} \dim (\pi) \operatorname{Tr} (\pi (f) \pi (f) ^ {*}) \leqslant (f * \widetilde {f}) (e) = \| f \| _ {2} ^ {2}.\tag{6}
\]

希尔伯特--施密特谱理论随后使他们能够证明这是等式（“普朗谢雷尔公式”）。对 \(\varphi\) （\(\mathcal{C}(\mathrm{G})\) 中的每个元素），彼得和外尔关联连续核 \(K_{\varphi}\colon G\times G\to C\) ，其定义为 \(\mathrm{K}_{\varphi}(x,y)=\varphi(xy^{-1})\) 。随后他们观察到，如果 \(\varphi\) 形如 \(f*\widetilde{f}\) 且 \(f\in\mathcal{C}(\mathrm{G})\) ， \(^{(6)}\) 则 \(K_{\varphi}\) 是希尔伯特--施密特理论意义下的正埃尔米特核。施密特算法的一个变体 \(^{(7)}\) 使他们可以构造核 \(K_{\varphi}\) 的特征函数标准正交基，继而将(6)中等式的证明归结为：由 \(K_{\varphi}\) 定义的算子的迹等于其特征值之和。

当然，只有表示 \(\pi\) 满足 \(\pi(f) \neq 0\)才可能出现在(6)中；\(K_{\varphi}\) 的谱理论表明，只要 \(f\) 不恒为零，就存在这样的表示。彼得和外尔很容易由此推出，不可约表示分离 G 的点——这是盖尔范德--赖科夫定理的一个特例，稍后将讨论该定理。

在弗罗贝尼乌斯的著名工作中，与(6)类似的等式曾使人能够证明 G 的正则表示包含所有不可约表示。但这涉及将该等式用于 C{[}G{]} 的单位元 f；此时显然有 \(\pi(f) \neq 0\) 对所有 \(\pi\) 都成立 。依彼得和外尔之见，卷积没有单位元解释了他们证明中的许多困难；然而，他们借用了傅里叶级数理论中一个注定具有巨大影响的思想，指出可将(6)应用于一列构成卷积逼近单位的函数，从而推出弗罗贝尼乌斯结果的类似物（参见 I，第120页第10号）。

同样地，从(6)产生的等式出发，其极化形式

\[
\sum_ {\pi \in \widehat {\mathrm{G}}} \dim (\pi) \operatorname{Tr} (\pi (x) \pi (y)) = (x * y) (e)\tag{7}
\]

通过令 \(y\) 为卷积逼近单位，可得到 \(\mathcal{C}(\mathrm{G})\) 的每个元素都能由不可约表示的矩阵系数的线性组合一致逼近。这些方法刚刚使外尔能对 H. 玻尔的“几乎周期函数”理论的基本结果给出新证明 {[}84{]}；我们很快将回到该理论。彼得和外尔指出，其中可以看出对非紧群表示理论的第一个分析应用，这里该群是 \(\mathbf{R}\) 的平移群。

年轻的量子理论很快给外尔提供了重新回到 R 的表示的机会。1927年底，他注意到（无论有限与否的）群和表示对于阐明该理论基础的意义 {[}85{]}。次年，他在一本轰动性的著作 {[}86{]} 中再次讨论这一问题。

在1927年的文章中，外尔注意到，若 u 是希尔伯特空间上的连续自伴算子，则 \(t \mapsto e^{itu}\) 是 R 的酉表示。但存在并非这种形式的 R 的酉表示；重新采用冯·诺伊曼关于用对称部分算子表示可观测物理量的思想，外尔提出每个 R 的酉表示都形如 \(t \mapsto e^{itu}\) ，其中 u 是希尔伯特空间上的部分自伴算子。这一观点使他能够把海森堡和薛定谔的“正则关系”（冯·诺伊曼已作详细研究）的分析，归结为同一希尔伯特空间上两个 R 的酉表示之间的联系。斯通 {[}73{]} 在1932年、沿着他对自伴算子谱性质的研究，证明了外尔所期待的结果（V，第428页，定理1）。

\histitem{局部紧交换群}

到20世纪30年代初，谐波分析与群表示之间的互动已经确立。接下来的十年中，二者相互推动，互动大为加强，分别促成对方迅速发展。

这种互动的一个极其肥沃的土壤是几乎周期函数的研究；H. 玻尔于1924年在实变量函数的情形中引入了它们 {[}9{]}。根据他的一个基本定理，这些函数正是 \(\mathbf{R}\) 上可由 \(x \mapsto e^{i\lambda x}\) 、\(\lambda \in \mathbf{R}\) 的线性组合一致逼近的有界连续函数（参见 II，第292页，习题54）。博赫纳 {[}6{]} 于1926年证明，\(f \in \mathcal{C}_b(\mathbf{R})\) 是几乎周期的，当且仅当其平移 \(x \mapsto f(x - a)\) 对于 \(a \in \mathbf{R}\) 在一致收敛拓扑下构成 \(\mathcal{C}_b(\mathbf{R})\) 的相对紧子集。博赫纳和冯·诺伊曼注意到，若 G 是拓扑群，这一刻画显然在 G 上（右侧或左侧）提供了几乎周期函数的概念。20世纪30年代初，在冯·诺伊曼 {[}52{]} 及彼得和外尔的方法推动下，这一理论有望迅速发展。

另一个富有成果的领域出现在 N. 维纳工作的边缘。他在1932年一篇著名论文 {[}88{]} 中把代数方法引入所谓塔伯尔问题的研究——也就是在已知某些加权平均行为的情况下，研究函数（或序列）沿滤子的渐近行为。他的方法受到 R. 施密特的思想启发。

维纳工作中的一个辅助结果（{[}88，引理 IIe{]}）尤其令人印象深刻并激发了诸多发展：若 \(f\) 是一个傅里叶级数绝对收敛且 \(f\) 不为零的周期函数，则 \(1/f\) 的傅里叶级数绝对收敛（参见 I，第38页，例）。早在1932年，R. 佩利和 N. 维纳 {[}89{]} 就勾勒出这一收敛结果、离散交换群的对偶群概念与几乎周期函数理论之间的联系。

在这一分析背景下，代数拓扑促使 L. 庞特里亚金为局部紧阿贝尔群的特征标研究提供决定性推动。他早在1932年就宣布，特征标群概念的意义在于把欧氏空间紧子集的同调群与其补集的同调群置于对偶关系中 {[}55{]}。在1934年发表的论文 {[}58{]} 中，他证明离散群的特征标群是紧的，随后在这一框架中建立对偶定理。虽然庞特里亚金使用了彼得和外尔的结果以及刚刚在一般情形建立存在性的哈尔测度，但他的主要目标是拓扑应用 {[}57{]}，并未明确关注谐波分析 \(^{(8)}\)：这些对偶定理是通过细致研究紧阿贝尔群的结构证明的。

庞特里亚金的工作立即引起极大兴趣。应冯·诺伊曼的建议，E. 范坎彭次年将对偶性推广到所有局部紧交换群 {[}37{]}。不久之后，他把新不变分析应用到几乎周期函数的研究，给出一个引人注目的结果 {[}38{]}；A. 外尔在同一时期独立发现了该结果 {[}79{]}：若 G 是局部紧阿贝尔群，可以得到一个紧群 \(\widetilde{G}\)：先给对偶群 \(\widehat{G}\) 赋予离散拓扑，再把由此得到的离散群的紧对偶群定义为 \(\widetilde{G}\)。群 G 随后典范地嵌入 \(\widetilde{G}\)（参见 II，第292页，习题54），对 \(\widetilde{G}\) 上连续函数应用彼得--外尔理论的简单结果，就能得到当时关于 G 上几乎周期函数的所有已知结果。

与此同时，在这一时期，经典分析中早已熟知的正定函数出现在 R 的酉表示研究中。J. 默瑟在1910年以前不久、结合积分方程理论研究了普遍正核 {[}45{]}；而其序列类比则在同一时期引出大量著作，尤其与矩问题有关（参见 IV，第359页，习题21）。与此同时，M. 马蒂亚斯于1923年已系统研究 R 上的正定函数 {[}43{]}；在博赫纳手中，它们成为傅里叶分析和概率论的基本工具之一 {[}7{]}。1933年，博赫纳 {[}8{]} 与里斯 {[}61{]} 独立注意到，R 的酉表示的每个对角矩阵系数都是 R 上的正定函数。这一观察使他们能够给出斯通定理的更简单证明，并将对一般理论的发展产生决定性影响（参见第7号）。

这一时期的结果由 A. 外尔整理在一本1937年前完成、1940年出版的著作中 {[}80{]}。在那里，他详细阐述彼得--外尔、庞特里亚金和范坎彭的结果。但在交换群情形，庞特里亚金和范坎彭强调群的结构和对偶定理，而外尔系统地发展了谐波分析，引入傅里叶变换并证明普朗谢雷尔公式与博赫纳定理（V，第455页，定理5）。他尤其强调卷积（称为“合成乘积”）的作用，以及正定函数的作用；后者如今定义在任意群上，并同表示的对角矩阵系数联系起来，至少在有限维表示情形如此。

因此在交换情形达到的这种一般性，稍晚些时候将为数论开辟新的视野。早在1936年，为向类域论引入代数方法，C. 谢瓦莱 {[}14{]} 引入了数域的 idèle 群；该群后来会作为 adèle 环的可逆元群出现（参见 AC，VII，第221--222页及 ÉHM，第143页）。1950年，J. 泰特 {[}75{]} 将证明，对 adèle 和 idèle 群的谐波分析使得可以容易地恢复 Hecke L-函数的函数方程，预示了把 adèle 群上的谐波分析同自守形式研究联系起来的非凡发展。

可以说，到20世纪30年代末，不变谐波分析的主要结果已经在紧群和局部紧交换群上建立。在交换情形，外尔的证明仍像庞特里亚金和范坎彭的证明一样，依赖于对群结构的精细了解，粗略说即 II 第2节第244页的分类结果。下一个十年借助盖尔范德学派的新方法将说明如何摆脱这一依赖（参见 H. 嘉当和 R. 戈德芒 {[}12{]}）。

\histitem{算子代数}

我们已经看到，里斯在1913年阐释希尔伯特学派工作时，赋予希尔伯特空间自同态代数 \(\mathcal{L}(\mathrm{E})\) 的作用，也赋予全纯函数演算和谱的抽象概念以作用。在1916年的回忆录中，他似乎已经预示赋范代数是谱理论的自然框架；经过巴拿赫及其学派在20世纪20年代

的工作，这一思想占据中心地位也就不足为奇了。M. 长友 {[}48{]} 于1936年引入巴拿赫代数概念，而 S. 马祖尔 {[}44{]} 于1938年证明，唯一作为域的 C 上赋范代数就是 C 本身（I，第26页，推论2）。

大约在20世纪30年代中期，与谱理论相关，抽象考察希尔伯特空间上的算子代数的机会不断增多。F. 默里和 J. 冯·诺伊曼在一系列深刻且有影响的工作 {[}47{]} 中，系统研究了等于自身双交换子的 \(\mathcal{L}(\mathrm{E})\) 的对合子代数（“冯·诺伊曼代数”）。另一方面，S. 斯廷于1936年提出引入“算子代数”的公理化概念，并深入研究相关结构 {[}70{]}。此外，斯通 {[}74{]} 受谱投影研究启发，于1937年研究其所有元素均为幂等元的含单位代数（“布尔代数”：参见 ÉHM，第146页）。他注意到，若 X 是局部紧拓扑空间，则开集的特征函数和处处稠密集补集的特征函数在 \(\mathcal{C}(\mathrm{X})\) 中生成一个布尔代数 A。随后他证明 A 的理想自然对应于空间 X 的闭子集，从而引入了一个后来特别富有成果的思想。

从1939年起，盖尔范德的观点为这些相对分散的结果打开了新视野。他的工作似乎部分受到维纳思想的推动，特别是其关于绝对收敛傅里叶级数的定理，以及彼得和外尔的方法。1939至1946年间，他与 M. 奈马克、D. 拉伊科夫和 G. 希洛夫合作，奠定了交换巴拿赫代数理论和星代数理论的基础（参见 {[}23{]}，第I卷，第169--400页）。特别地，他给出了维纳定理的经典证明——由 P. 莱维 {[}39{]} 推广——本书中可见该证明（I，第38页，例）。

在盖尔范德的方法中，与交换巴拿赫代数 A 相关的本质对象是 A 的极大理想空间 X。事实上，通过盖尔范德--马祖尔定理（I，第26页，推论2），A 的每个元素都在 X 上定义一个函数。众所周知，这一观点后来会在代数几何的发展中变得十分重要（参见 ÉHM，第146--148页）。这一方法本身受到了上文提到的斯通关于布尔代数的工作的启发。

似乎是 L. 卢米斯在1953年出版的课程 {[}41，第53页{]} 中首次强调 A 的特征标空间来阐述盖尔范德理论，并按本书所理解的意义定义盖尔范德变换。这一观点也许最清楚的优点是，特征标空间的拓扑性质直接来自巴拿赫空间对偶单位球在弱拓扑下的紧性。

从最早的工作起，盖尔范德就通过柯西积分在巴拿赫代数中构造了一元全纯函数演算，并由此推出：当极大理想空间不连通时，对合巴拿赫代数允许按环的乘积作非平凡分解。任意巴拿赫代数的情形直到1953年才得到处理，那时希洛夫在多元情形发展出一种函数演算。{[}68{]}。

至于星代数，盖尔范德和奈马克早在1942年就已证明它们可以实现为希尔伯特空间上的算子代数（参见 V，第442页定理2），它们很快成为大量深刻研究的对象；推动力来自群表示的应用、默里和冯·诺伊曼的工作以及量子力学的数学形式化。特别值得一提的是 I. 西格尔的贡献；他是系统研究星代数表示的先驱之一，这项研究与局部紧群的酉对偶研究相联系，我们将在几行后讨论后者。

尽管盖尔范德、奈马克、拉伊科夫和希洛夫在1946年以前已经奠定 C*-代数理论的基本基础，某些技术细节后来才得到改进。例如，R. 阿伦斯 {[}1{]} 说明了如何避免调用“希洛夫边界”（参见 I，第171页，习题26）的存在性，以证明盖尔范德变换是一个对合态射。此外，盖尔范德--奈马克关于含单位 C*-代数的理论起初把 \(1 + x^{*}x\) 对所有 \(x\) 可逆这一事实作为公理加入。然而他们猜想这个公理并非必要；这一断言的证明虽然初等却很精细，直到1952年前后才在 M. 福卡米亚 {[}21{]} 的工作、经 I. 卡普兰斯基 {[}62{]} 接续之后给出。

\histitem{局部紧群的表示}

直到1939年，拓扑群表示的研究局限于紧情形和交换情形。另一方面，物理学家在量子力学中发现了许多理由，需要研究不可约酉表示，以寻找能够模拟基本粒子状态的希尔伯特空间。P. 狄拉克在 E. 马约拉纳工作的基础上，于1936年前后指出，在这一视角下，洛伦兹群 SO(3,1) 和庞加莱群 SO(3,1) × \(R^{4}\) 的意义。他们的方法仍与当时数学家的方法相去甚远。但1939年，E. 维格纳 {[}90{]} 以默里和冯·诺伊曼的工作为基础，对庞加莱群的不可约酉表示作出分类，打开了一道缺口。

我们似乎应将研究局部紧群表示的一般路径归功于盖尔范德。他在1942年写的一篇回忆录中（参见 {[}23{]}，第II卷，第3--17页）与拉伊科夫注意到，这一理论之所以可能，是因为酉表示与正定函数之间存在联系。他们证明每个正定函数都有希尔伯特实现（V，第432页定理1），并由此推出不可约酉表示分离 G 的点（V，第454页定理4）。对合代数 L \(^{1}\) (G) 上的正线性泛函发挥了基本作用：盖尔范德与拉伊科夫揭示了它们同正定函数相联系的关系（参见 V，第448页命题13）。对合代数的包络 C* 代数在 {[}24，第48节{]} 中已经（未命名）出现，不过当时还没有直接提出把代数 Stell(G) 与 G 联系起来（参见 I 第125页定义9）。

1945年以后，酉表示的研究突飞猛进。与 C*-代数的互动以及对量子理论基础的分析很快带来有价值的一般结果：约1947年，西格尔把各种 \(^{(9)}\) 星代数同一个局部紧群联系起来 {[}67{]}，并编织起它们的表示与所研究群的表示之间的联系。不久后，G. 麦基提炼出诱导表示的概念 {[}42{]}，它将在酉表示的具体结果中无处不在。

与此同时，对揭示其多样性与深度的例子的研究大大推动了该理论。盖尔范德和奈马克研究 SL(2, C) 这一例子，于1947年分离出诱导表示概念的一个特殊情形，随后证明它是研究该群表示的关键 {[}23，第II卷，第41--124页{]}。同年，V. 巴格曼研究群 SL(2, R)，发现一个可数的平方可积表示族，称为“离散系列”，以及它们的矩阵系数之间的正交关系（参见 V，第424页命题8） {[}5{]}。正如戈德芒立即认识到的 {[}25{]}，Bargmann 正交关系由所有局部紧群的平方可积表示满足。

很快，对特殊群类的研究又一次、并且长时期内优先于抽象理论。特别是，哈里什-钱德拉将着手一般研究约化李群的表示——这项史诗般的任务，以及朗兰兹将为数论预见的卓越后果，我们都无法在此描述。

\specialchapter{参考文献}

{[}1{]} R. ARENS - "关于盖尔范德和奈马克的一个定理", Proc. Natl. Acad. Sci. USA 32 (1946), p. 237-239, Zbl 0060.27101.

{[}2{]} M. F. ATIYAH and I. M. SINGER -- "紧流形上椭圆算子的指标", Bull. Amer. Math. Soc. 69 (1963), p. 422--433, Zbl 0118.31203; Michael Atiyah 文集, vol. 3, Clarendon Press, 1988, p. 13--24.

{[}3{]} F. V. ATKINSON -- "赋范空间中线性方程的正规可解性", Mat. Sb., Nov. Ser. 28 (1951), p. 3--14, Zbl 0042.12001（俄文）。

{[}4{]} S. BANACH - 线性运算理论, Warsaw, 1932, Zbl 0005.20901.

{[}5{]} V. BARGMANN - "洛伦兹群的不可约酉表示", Ann. of Math. (2) 48 (1947), p. 568-640, Zbl 0045.38801.

{[}6{]} S. BOCHNER -- "几乎周期函数理论的贡献。II：多变量函数", Math. Ann. 96 (1926), p. 383--409, JfM 52.0265.01; 论文集, vol. 2, Amer. Math. Soc., 1992, p. 79--106.

{[}7{]} \_\_\_\_ 傅里叶积分讲义, Leipzig, 1932, Zbl 0006.11001.

{[}8{]} \_\_\_\_ "酉算子线性族的谱分解", Sitzungsber. Preuß. Akad. Wiss., Phys.-Math. Kl. (1933), p. 371--376, Zbl 0007.16603; 论文集, vol. 1, Amer. Math. Soc., 1992, p. 579--586.

{[}9{]} H. BOHR - "关于几乎周期函数的理论。I. 傅里叶级数理论的推广", Acta Math. 45 (1924), p. 29-127, JfM 50.0196.01.

{[}10{]} J. W. CALKIN - "希尔伯特空间有界算子环中的双边理想与同余关系", Ann. of Math. (2) 42 (1941), p. 839-873, Zbl 0063.00692.

{[}11{]} T. CARLEMAN - 关于具有实对称核的奇异积分方程, Uppsala, 1923, JfM 49.0272.01.

{[}12{]} H. CARTAN and R. GODEMENT -- "局部紧阿贝尔群中的对偶理论与谐波分析", Ann. Sci. Éc. Norm. Supér. (3) 64 (1947), p. 79--99, Zbl 0033.18801; Henri Cartan 文集, vol. III, Springer, 1979, p. 1203--1223.

{[}13{]} H. CARTAN and J-P. SERRE -- "关于紧解析流形的有限性定理", C. R. Acad. Sci., Paris 237 (1953), p. 128--130, Zbl 0050.17701; Henri Cartan 文集, vol. II, Springer, 1979, p. 684--686.

{[}14{]} C. CHEVALLEY - "无限扩张的类域论推广", J. Math. Pures Appl. (9) 15 (1936), p. 359-371, Zbl 0015.15101.

{[}15{]} J. DIEUDONNÉ -- "关于赋范空间的同态", Bull. Sci. Math. (2) 67 (1943), p. 72--84, Zbl 0028.23301; 数学选集, Hermann, 1981, p. 268--281.

{[}16{]} J. DIEUDONNÉ - 泛函分析史, North-Holland Publ. Co., 1981, Zbl 0478.46001.

{[}17{]} P. L. DIRICHLET -- "证明每个首项和公差互素的无限算术级数包含无穷多个素数的定理", Abh. Preuß. Akad. Wiss. 8 (1837), p. 45--81; 文集, vol. 1, Reimer (Berlin), 1889, p. 315--356.

{[}18{]} P. ENFLO - "巴拿赫空间逼近问题的反例", Acta Math. 130 (1973), p. 309-317, Zbl 0267.46012.

{[}19{]} \_\_\_\_ "关于巴拿赫空间不变子空间问题", Acta Math. 158 (1987), no. 3-4, p. 213--313, Zbl 0663.47003.

{[}20{]} J. M. G. FELL - "C*代数的对偶空间", Trans. Amer. Math. Soc. 94 (1960), p. 365-403, Zbl 0090.32803.

{[}21{]} M. FUKAMIYA - "关于盖尔范德和奈马克的一个定理以及 B*代数", Kumamoto J. Sci., Math. 1 (1952), p. 17-22, Zbl 0049.08605.

{[}22{]} C. F. GAUSS - Disquisitiones Arithmeticae, Göttingen, 1801.

{[}23{]} I. M. GELFAND - 论文集, 3 vol., Springer, 1987.

{[}24{]} I. M. GELFAND and M. A. NAIMARK -- "带对合的赋范环及其表示", Izv. Akad. Nauk SSSR, Ser. Mat. 12 (1948), p. 445--480, Zbl 0031.03403; I. M. Gelfand 论文集, vol. I, Springer, 1987, p. 364--400（俄文）。

{[}25{]} R. GODEMENT -- "关于 V. Bargmann 正交关系。II. 一般证明", C. R. Acad. Sci. Paris 225 (1947), p. 657--659, Zbl 0029.19907.

{[}26{]} I. C. GOHBERG -- "关于赋范空间中的线性方程", Dokl. Akad. Nauk SSSR, n. Ser. 76 (1951), p. 477--480, Zbl 0042.11903（俄文）。

{[}27{]} P. R. HALMOS - "谱定理说了什么？", Amer. Math. Monthly 70 (1963), p. 241-247, Zbl 0132.35606; 选集, vol. II, Springer, 1983, p. 59-66.

{[}28{]} W. HEISENBERG -- "关于量子理论形式规则在反常塞曼效应问题中的一种修改", Z. Phys. 26 (1924), p. 291--307, JfM 51.0740.07; 文集：科学原始论文, vol. 1, Springer, 1985, p. 306--325.

{[}29{]} E. HELLINGER and O. TOEPLITZ -- "积分方程与无穷多个未知量的方程", Enzyklopädie der mathematischen Wissenschaften, art. II C 13, Leipzig (Teubner), 1927, JfM 53.0350.01.

{[}30{]} E. HILB - "关于任意函数的积分表示", Math. Ann. 66 (1909), p. 1-66, JfM 39.0380.03.

{[}31{]} D. HILBERT - "线性积分方程一般理论的基础。第一篇通信。", Nachr. Ges. Wiss. Göttingen, Math.-Phys. Kl. (1904), p. 49-91, JfM 35.0378.02.

{[}32{]} \_\_\_\_ "线性积分方程一般理论的基础。第四篇通信。", Nachr. Ges. Wiss. Göttingen, Math.-Phys. Kl. (1906), p. 157--227, JfM 37.0351.03.

{[}33{]} T. H. HILDEBRANDT - "关于完全连续线性变换", Acta Math. 51 (1928), p. 311-318, JfM 54.0427.03.

{[}34{]} E. HILLE - 泛函分析与半群, Colloq. Publ., vol. 31, Amer. Math. Soc., 1948, Zbl 0033.06501.

{[}35{]} M. JAMMER - 量子力学的概念发展, McGraw-Hill, 1966.

{[}36{]} M. KAC - "能听出鼓的形状吗？", Am. Math. Mon. 73 (1966), p. 1-23, Zbl 0139.05603.

{[}37{]} E. R. VAN KAMPEN - "局部双紧阿贝尔群及其特征标群", Ann. of Math. (2) 36 (1935), p. 448-463, Zbl 0011.39401.

{[}38{]} \_\_\_\_ "几乎周期函数与紧群", Ann. of Math. (2) 37 (1936), p. 78--91, Zbl 0014.01702.

{[}39{]} P. LÉVY - "关于傅里叶级数的绝对收敛", Compos. Math. 1 (1934), p. 1-14, Zbl 0008.31201.

{[}40{]} V. I. LOMONOSOV - "与完全连续算子交换的算子族的不变子空间", Funct. Anal. Appl. 7 (1974), p. 213-214, Zbl 0293.47003.

{[}41{]} L. H. LOOMIS - 抽象谐波分析导论, Van Nostrand Co., 1953, Zbl 0052.11701.

{[}42{]} G. W. MACKEY - "局部紧群表示的不可约性。I", Proc. Nat. Acad. Sci. U.S.A. 35 (1949), p. 537-545, Zbl 0035.06901.

{[}43{]} M. MATHIAS - "关于正傅里叶积分", Math. Z. 16 (1923), p. 103--125, JfM 49.0207.03.

{[}44{]} S. MAZUR - "关于线性环", C. R. Acad. Sci., Paris 207 (1938), p. 1025-1027, Zbl 0020.20101.

{[}45{]} J. MERCER - "正型和负型函数及其与积分方程理论的联系", Philos. Trans. R. Soc. Lond., Ser. A 209 (1909), p. 415-446, JfM 40.0408.02.

{[}46{]} S. G. MIKHLIN - "关于希尔伯特空间中线性方程的求解", Dokl. Akad. Nauk SSSR (N. S.) 57 (1947), p. 11-12, Zbl 0029.05101（俄文）。

{[}47{]} F. J. MURRAY and J. VON NEUMANN - "关于算子环", Ann. of Math. (2) 37 (1936), p. 116-229, Zbl 0014.16101; John von Neumann 文集, vol. III, Pergamon Press, 1961, p. 6-120.

{[}48{]} M. NAGUMO - "线性度量环中的一些解析研究", Jpn. J. Math. 13 (1936), p. 61-80, Zbl 3024132.

{[}49{]} J. VON NEUMANN -- "量子力学的数学基础", Göttingen Nachrichten 1927 (1927), p. 1--57, JfM 53.0848.03; 文集, vol. I, Pergamon Press, 1961, p. 151--207.

{[}50{]} \_\_\_\_ "埃尔米特函数算子的一般特征值理论", Math. Ann. 102 (1929), p. 49--131, JfM 55.0824.02; 文集, vol. II, Pergamon Press, 1961, p. 3--85.

{[}51{]} \_\_\_\_ "关于函数运算代数与正规算子理论", Math. Ann. 102 (1929), p. 370--427, JfM 55.0825.02; 文集, vol. II, Pergamon Press, 1961, p. 86--143.

{[}52{]} \_\_\_\_ "群上的几乎周期函数。I", Trans. Amer. Math. Soc. 36 (1934), p. 445--492, Zbl 0009.34902; 文集, vol. II, Pergamon Press, 1961, p. 454--501.

{[}53{]} F. NOETHER -- "关于一类奇异积分方程", Math. Ann. 82 (1920), p. 42--63, JfM 47.0369.02.

{[}54{]} F. PETER and H. WEYL -- "闭连续群的本原表示的完备性", Math. Ann. 97 (1927), p. 737--755, JfM 53.0387.02; Hermann Weyl Gesammelte Abhandlungen, Bd. III, Springer, 1968, p. 58--75.

{[}55{]} L. PONTRYAGIN -- "闭集的一般对偶定理", Verhandlungen Kongreß Zürich, 2, p. 195-197, 1932, JfM 58.0646.08.

{[}56{]} \_\_\_\_ "几乎周期函数与位置分析", C. R. Acad. Sci., Paris 196 (1933), p. 1201--1203, Zbl 0006.42801.

{[}57{]} \_\_\_\_ "闭集的拓扑对偶一般定理", Ann. of Math. (2) 35 (1934), no. 4, p. 904--914, Zbl 0010.18004.

{[}58{]} \_\_\_\_ "拓扑交换群理论", Ann. of Math. (2) 35 (1934), p. 361--388, Zbl 0009.15601.

{[}59{]} F. RIESZ - 无穷多个未知量的线性方程组, Gauthier-Villars, 1913, JfM 44.0401.01.

{[}60{]} \_\_\_\_ "关于线性泛函方程", Acta Math. 41 (1916), p. 71--98, JfM 46.0635.01 ; 全集, t. II, Gauthier-Villars, 1960, p. 1053--1080.

{[}61{]} \_\_\_\_ "关于斯通和博赫纳定理", Acta Litt. Sci. Szeged 6 (1933), p. 184--198, Zbl 0007.30903; 全集, t. II, Gauthier-Villars, 1960, p. 1135--1149.

{[}62{]} J. SCHATZ - "对 M. Fukamiya 所著《关于盖尔范德和奈马克的一个定理以及 B*代数》一文的评论", Math. Reviews 14 (1952), p. 884.

{[}63{]} J. SCHAUDER - "关于线性的完全连续泛函运算", Studia Math. 2 (1930), p. 183-196, JfM 56.0354.01.

{[}64{]} A. SCHUSTER -- "谱的起源", Report on the fifty-second meeting of the British association for the advancement of Science, J. Murray (London), 1882, p. 120--143.

{[}65{]} L. SCHWARTZ - "函数、微分、傅里叶变换概念的推广及其数学和物理应用", Ann. Univ. Grenoble, Sect. Sci. Math. Phys., n. Ser. 21 (1946), p. 57-74, Zbl 0060.27504; 科学全集, t. I, Soc. Math. France, 2011, p. 61-88.

{[}66{]} \_\_\_\_ "同态与完全连续映射", C. R. Acad. Sci., Paris 236 (1953), p. 2472--2473, Zbl 0050.33301 ; 科学全集, t. I, Soc. Math. France, 2011, p. 345--346.

{[}67{]} I. E. SEGAL - "局部紧群的群代数", Trans. Amer. Math. Soc. 61 (1947), p. 69-105, Zbl 0032.02901.

{[}68{]} G. E. SHILOV -- "关于交换赋范环分解为理想直和", Mat. Sb., Nov. Ser. 32 (1953), p. 353--364, Zbl 0052.34203（俄文）。

{[}69{]} S. SOBOLEV - "求解正规双曲型线性方程柯西问题的新方法", Rec. Math. Moscou, n. Ser. 1 (1936), p. 39-71, Zbl 0014.05902.

{[}70{]} S. W. P. STEEN - "算子理论导论：I. 实算子。模", Proc. Lond. Math. Soc. (2) 41 (1936), p. 361-392, Zbl 0014.16201.

{[}71{]} T.-J. STIELTJES - "连分数研究", Ann. Fac. Sci. Toulouse Sci. Math. Sci. Phys. 8 (1894), no. 4, p. J1-J122, JfM 25.0326.01.

{[}72{]} M. H. STONE - 希尔伯特空间中的线性变换及其在分析中的应用, Colloq. Publ., Am. Math. Soc., vol. 15, Amer. Math. Soc., 1932, Zbl 0005.40003.

{[}73{]} \_\_\_\_ "关于希尔伯特空间中的单参数酉群", Ann. of Math. (2) 33 (1932), p. 643--648, Zbl 0005.16403.

{[}74{]} \_\_\_\_ "布尔环理论在一般拓扑中的应用", Trans. Am. Math. Soc. 41 (1937), p. 375--481, Zbl 0017.13502.

{[}75{]} J. T. TATE — “数域中的傅里叶分析与黑克 ζ 函数”, 学位论文, Princeton, 1950; 文集, part I, Amer. Math. Soc., 2016, p. 1–43, Zbl 1407.01030.

{[}76{]} O. TOEPLITZ — “费耶尔定理的代数类似物”, Math. Z. 2 (1918), p. 187–197, JfM 46.0157.02.

{[}77{]} H. WEBER — “证明：每个真正本原二次型都能表示无穷多个素数的定理。”, Math. Ann. 20 (1882), p. 301–330, JfM 14.0141.01.

{[}78{]} \_\_\_\_ “阿贝尔数域理论”, Acta Math. 9 (1887), p. 105–130, JfM 18.0055.04.

{[}79{]} A. WEIL — “关于几乎周期函数”, C. R. Acad. Sci., Paris 200 (1935), p. 38–40, Zbl 0010.35303; 科学全集, vol. I, Springer, 1979, p. 101–103.

{[}80{]} \_\_\_\_ 拓扑群上的积分及其应用, Actualités scientifiques et industrielles, vol. 869, Hermann, 1940, Zbl 0063.08195.

{[}81{]} H. WEYL — “关于差为完全连续的有界二次型”, Rend. Circ. Mat. Palermo 27 (1909), p. 373–392, JfM 40.0395.01; Gesammelte Abhandlungen, Bd. I, Springer, 1968, p. 175–194.

{[}82{]} \_\_\_\_ “关于带奇点的常微分方程及任意函数的相应展开”, Math. Ann. 68 (1910), p. 220–269, JfM 41.0343.01; Gesammelte Abhandlungen, Bd. I, Springer, 1968, p. 248–297.

{[}83{]} \_\_\_\_ “通过线性变换表示连续半单群的理论。I、II、III，补编”, Math. Z. 23 (1925), p. 271–309, JfM 51.0319.01; Gesammelte Abhandlungen, Bd. II, Springer, 1968, p. 543–647.

{[}84{]} \_\_\_\_ “积分方程与几乎周期函数”, Math. Ann. 97 (1926), p. 338–356, JfM 52.0260.02; Gesammelte Abhandlungen, Bd. III, Springer, 1968, p. 38–57.

{[}85{]} \_\_\_\_ “量子力学与群论”, Z. Phys. 46 (1927), p. 1–46, JfM 53.0848.02; Gesammelte Abhandlungen, Bd. III, Springer, 1968, p. 90–135.

{[}86{]} \_\_\_\_ 群论与量子力学, S. Hirzel (Leipzig), 1928, Zbl 0537.22023.

{[}87{]} E. WEYR — “关于双线性型理论”, Monatsh. Math. Phys. 1 (1890), p. 163–236, JfM 22.0141.02.

{[}88{]} N. WIENER — “塔伯尔定理”, Ann. of Math. (2) 33 (1932), p. 1–100, Zbl 0004.05905; 论文选集, MIT Press, 1964, p. 261–360.

{[}89{]} N. WIENER and R. E. A. C. PALEY — “无限阿贝尔群的特征标的解析性质”, Verhandlungen Kongreß Zürich, 2, p. 95, 1932, JfM 58.0445.02.

{[}90{]} E. WIGNER — “关于非齐次洛伦兹群的酉表示”, Ann. of Math. (2) 40 (1939), p. 149–204, Zbl 0020.29601; 文集, part A, Springer, 1993, p. 334–389.

{[}91{]} W. WIRTINGER — “关于黎曼双曲型微分方程积分方法及其在振动问题中应用的贡献”, Math. Ann. 48 (1897), p. 365–389, JfM 27.0281.01.

{[}92{]} S. L. WORONOWICZ — “与 C*代数相关的无界元与非紧量子群”, Comm. Math. Phys. 136 (1991), no. 2, p. 399–432, Zbl 0743.46080.

{[}93{]} B. YOOD — “在线性变换加上完全连续变换后保持的性质”, Duke Math. J. 18 (1951), p. 599–612, Zbl 0043.11901.

\specialchapter{记号索引}

\(\mathcal{L}^{\mathrm{c}}(\mathrm{E};\mathrm{F})\) (紧线性映射) 2 \(\mathrm{E}_{(\mathbf{C})}\) (复化向量空间) 2 \(u_{(\mathbf{C})}\) (复化线性映射) 3 \(\mathcal{L}^{\mathrm{f}}(\mathrm{E};\mathrm{F})\) (有限秩线性映射) 4 \(\mathcal{C}_0(\mathrm{X};\mathrm{K}), \mathcal{C}_0(\mathrm{X})\) (在无穷远处趋于0的连续函数) 20 \(\mathcal{N}^{p,q}(\mathrm{X} \times \mathrm{Y})\) (核空间) 27 \(u_k\) (带核的积分算子) 29 \(k_\circ\)（映射 \(y \mapsto k_y\) ，其中 \(k_y(x) = k(x,y)\) ）30 \(\mathcal{L}^{p',q'}(\mathrm{X},\mathrm{Y},\mu,\nu)\) (核空间) 31 \(u \equiv v\) (模 \(\mathcal{L}^{\mathrm{f}}(\mathrm{E};\mathrm{F})\) 的同余) 40 \(\mathcal{F}(\mathrm{E};\mathrm{F})\) (弗雷德霍姆映射) 41\\
ind(u) (弗雷德霍姆映射的指标) 43 \(\mathcal{I}(\mathrm{E};\mathrm{F})\) (E 到 F 的同构) 57 \(\mathcal{M}\mathcal{D}(\mathrm{E};\mathrm{F})\) (单射直态射) 57 \(\mathcal{E}\mathcal{D}(\mathrm{E};\mathrm{F})\) (满射直态射) 57 \(\mathcal{R}(\mathrm{E})\) (里斯自同态) 60\\
((u)) (u 的余范数) 61 \(\mathcal{M}(\mathrm{E};\mathrm{F})\) (严格单射态射) 67 \(\mathcal{Q}\mathcal{M}(\mathrm{E};\mathrm{F})\) (核有限维的严格态射) 67 \(\mathcal{E}(\mathrm{E};\mathrm{F})\) (连续满射态射) 70 \(\mathcal{Q}\mathcal{E}(\mathrm{E};\mathrm{F})\) (像具有有限余维的态射) 70 \(\mathcal{C}alk(\mathrm{E})\) (巴拿赫空间 E 的卡尔金代数) 75 \(\mathrm{N}_{\lambda}\) (幂零空间) 78 \(\mathrm{I}_{\lambda}\) (余幂零空间) 78 \(e_{\lambda}(u)\) (谱投影) 82 \(m_{\lambda}(u)\) (谱重数) 82 \(\mathrm{Sp}_{\mathrm{s}}(u)\) (谱的敏感点) 83 \(\mathrm{Sp}_{\mathrm{e}}(u)\) (本质谱) 83 \(\mathrm{Sp}_{n}(u)\) (谱中使得 \(u - \lambda 1_{\mathrm{E}}\) 具有指标 n 的元素) 86 \(\mathrm{Sp}_{\omega}(u)\) (\(\mathrm{Sp}_{n}(u)\) 的并集的补集) 86 \(\mathcal{D}_{\mathrm{B}}(\mathrm{E})\) (基 B 中对角的自同态) 146 \(\mathrm{I}_{\mathrm{E}}\) (F ≠ E 的有限维子空间的维数) 151 \(\lambda_n(u)\) (u 的特征值递减序列) 152

\(r_{\mathrm{F}}(u), \mathrm{R}_{\mathrm{F}}(u)\) （瑞利商）\ldots.. 153 \(\mathcal{F}_n\) （维数 \(n\) 的子空间）\ldots.. 153 \((\alpha_n(u))_{n \in \mathrm{J}}\) （奇异值序列）\ldots.. 156 \((\alpha_n(u))_{n \in \mathrm{I_E}}\) （奇异值的扩展序列）\ldots.. 156 \(\mathcal{L}_1(\mathrm{E})\) （有限迹自同态）\ldots.. 164 \(\|u\|_1\) （\(\mathcal{L}_1(\mathrm{E};\mathrm{F})\) 上的范数）\ldots.. 167 \(\mathcal{L}_1(\mathrm{E};\mathrm{F})\) （核型线性映射）\ldots.. 169 \(\mathcal{L}_{\mathrm{u}}(\mathrm{X})\) （普遍可测函数）\ldots.. 183 \(\mathcal{L}_{\mathrm{u}}^{\infty}(\mathrm{X})\) （有界普遍可测函数）\ldots.. 183 \(m_g, \widetilde{m}_g\) （乘法自同态）\ldots.. 186 \(\mathcal{G}_{\mathrm{A}}\) （盖尔范德变换）\ldots.. 190 \(\mathcal{H}_{\mathrm{A}}\) （盖尔范德变换的逆）\ldots.. 190 \(\mathcal{F}\) （傅里叶变换）\ldots.. 196 \(\overline{\mathcal{F}}\) （余傅里叶变换）\ldots.. 196 \(\mathcal{D}(\mathrm{U})\) （U 上的测试函数）\ldots.. 200 \(\mathcal{D}_{\mathbf{R}}(\mathrm{U})\) （实值测试函数）\ldots.. 201 \(\mathcal{D}'(\mathrm{U})\) （U 上的广义函数）\ldots.. 203 \(\partial^\alpha f\) （广义函数的偏导数）\ldots.. 208 \(\mathcal{S}(\mathbf{R}^n)\) （\(\mathbf{R}^n\) 上的施瓦茨函数）\ldots.. 209 \(\mathcal{S}'(\mathbf{R}^n)\) （\(\mathbf{R}^n\) 上的缓增广义函数）\ldots.. 215 \(\mathcal{M}^t(\mathbf{R}^n)\) （缓增测度空间）\ldots.. 216 \(W^{k,p}(U)\) （指标 \(k\) 、指数 \(p\) 的索伯列夫空间）\ldots.. 222 \(W_0^{k,p}(U)\) （\(\mathcal{D}(U)\) 在 \(W^{k,p}(U)\) 中的闭包）\ldots.. 222 \(H^k(U)\) （指标 \(k\) 、指数2的索伯列夫空间）\ldots.. 222 \(H_0^k(U)\) （\(\mathcal{D}(U)\) 在 \(H^k(U)\) 中的闭包）\ldots.. 222 \(\mathcal{P}(\mathrm{E};\mathrm{F})\) （从 E 到 F 的部分算子）\ldots.. 224 \(\mathcal{P}(\mathrm{E})\) （E 上的部分算子）\ldots.. 224 dom( \(u\) ) （\(u\) 的定义域）\ldots.. 224 \(1_{\mathrm{D}}\) （定义域为 D 的恒等部分算子）\ldots.. 225 \(u \circ v\) （部分算子 \(u\) 与 \(v\) 的复合）\ldots.. 225 \(u \subset v\) （\(v\) 是 \(u\) 的扩张）\ldots.. 226 \((u_i)_{i \in I}\) （乘积部分算子）\ldots.. 226 \(u + v\) （部分算子的和）\(u\)、\(v\)\ldots.. 226 \(\overline{u}\) （\(u\) 的闭包）\ldots.. 228 \(E_u\) （由 \(u\) 的定义域定义的赋范空间）\ldots.. 230 \(\|x\|_u\) （空间 \(E_u\) 上的范数）\ldots.. 230 \((x | y)_u\) （定义希尔伯特空间 \(E_u\) 的内积）\ldots.. 230 \(m_g, \widetilde{m}_g\) （乘法部分算子）\ldots.. 231 \(u^*\) （\(u\) 的伴随）\ldots.. 235 \(\mathcal{A}(\mathrm{E})\) （自伴部分算子）\ldots.. 238 \(^t\mathrm{D}\) （转置微分算子）\ldots.. 242 \(H_D\) （微分算子 \(D_+\) 的定义域）\ldots.. 242 \(D_-\) （定义域为 \(\mathcal{D}(U)\) 的微分算子）\ldots.. 242 \(D_+\) （定义域为 \(H_D\) 的微分算子）\ldots.. 242 \(R(u,\lambda)\) （\(u\) 的预解式）\ldots.. 244 \(\mathrm{PSp}_{\varepsilon}(u)\) （\(\varepsilon\)-伪谱）\ldots.. 250 \(b(u)\) （\(u\) 的有界化）\ldots.. 262 \(\Omega(\mathrm{E})\) （使 \(v\) 正且单射的 \(1_{\mathrm{E}} - v^*v\) 的集合）\ldots.. 264 \(B(v)\) （部分算子 \(v\circ((1_{\mathrm{E}} - v^*v)^{1/2})^{-1}\) ）\ldots.. 264

\(\beta (z) = z / \sqrt{1 + |z|^2}\) 266 \(D_f\) （\(f(u)\) 的定义域）271 \(f(u)\) （普遍可测函数演算）272 \(p_A\) （由 A 定义的谱投影）277 \(u^\alpha\) （函数演算）289 \(|u|\) （\(u\) 的绝对值）290 \(q_u\) （与 \(u\) 关联的正部分形式）291 \(E_q\) （预希尔伯特空间）292 \((x|y)_q\) （\(E_q\) 上的内积）292 \(\| x \| _q\) （\(E_q\) 上的范数）292 \(u_q\) （表示 \(q\) 的部分算子）292 \(\mathrm{Sp}_{\mathrm{s}}(u)\) （\(u\) 的敏感谱）297 \(\mathrm{Sp}_{\mathrm{e}}(u)\) （\(u\) 的本质谱）297 \(\mathrm{Sp_h}(u)\) （谱的上部）298 \(\varrho_{e}\) （本质谱的下界）298 \(\mathrm{Sp_b}(u)\) （谱的下部）298 \(\mathcal{F}_n\) （维数 \(n\) 的子空间）299 \(\mathcal{F}_n^u\) （\(\mathrm{dom}(u)\) 中维数 \(n\) 的子空间）299 \(\widetilde{r}_{\mathrm{F}}(u), \widetilde{\mathrm{R}}_{\mathrm{F}}(u)\) （瑞利商）300 \(\|v\|_u\) （关于 \(v\) 的 \(u\) 的相对范数）306\\
vp （柯西主值）337 \(\chi_\varrho\) （表示的特征标）374 \(\mathrm{Res}_{H}^{G}(\varrho)\) （\(\varrho\) 到 H 的限制）375 \(\mathrm{Hom}_{G}(\varrho_1,\varrho_2)\) （G-态射空间）375 \(\varrho^n = n \varrho\) （表示 \(\varrho\) 的副本之和）376 \(\check{\varrho}\) （逆步表示）377 \(\mathrm{Sesq_G}(\varrho, \pi), \mathrm{Sesq_G(E,F)}\) （G-不变半双线性形式）380 \(\overline{\varrho}\) （共轭表示）382 \(\varrho_1 \boxtimes \varrho_2\) （外张量积）385 \(\varrho_1 \otimes \varrho_2\) （张量积）385 \(\varrho^{\otimes n}\) （\(\varrho\) 的张量幂）385 \(\Upsilon(G)\) （矩阵系数）386 \(\Theta(G)\) （有限维矩阵系数）386 \(\operatorname{cl}(\pi)\) （酉表示的类）392 \(\widehat{G}\) （酉对偶）393 \(M_\pi (\varrho)\) （同型分量）394 \(\mathcal{M}_c(G)\) （紧支撑测度）400 \(\mathcal{M}^1(G)\) （有界测度）401 \(\gamma_{G}^{(p)}\) （左正则表示）405 \(\delta_{G}^{(p)}\) （右正则表示）405 \(\varrho_{G}^{(p)}\) （双正则表示）406 \(\gamma_{G}\) （\(L^2\) 中的左正则表示）407 \(\delta_{G}\) （\(L^2\) 中的右正则表示）407 \(\varrho_{G}\) （\(L^2\) 中的双正则表示）407 \(\mathcal{F}_{\pi}(G)\) （π-等变函数）408 \(\mathcal{K}_{\pi}(G)\) （在）\(\mathcal{F}_{\pi}(G)\) 中模 H 具有紧支集的连续函数）\ldots{} 408 \(N_p(f)\) （在）\(\mathcal{F}_π(G)\) 上的半范数）\ldots{} 408 \(\mathcal{L}_{\pi}^{p}(G)\) （模 H 对幂 \(p^e\) 可积的 π-等变函数）\ldots{} 409

\(L_{\pi}^{p}(G)\) （与 \(\mathcal{L}_{\pi}^{p}(G)\) 关联的巴拿赫空间）\ldots.. 410 \(\mathcal{L}_{\pi}^{2}(G)\) （模 H 平方可积的 \(\pi\)-等变函数）\ldots.. 415 \(L_{\pi}^{2}(G)\) （与 \(\mathcal{L}_{\pi}^{2}(G)\) 关联的希尔伯特空间）\ldots.. 415 \(\operatorname{Ind}_{H}^{G}(\pi)\) （由 H 的表示 π 诱导的 G 表示）\ldots.. 416 \(\gamma_{G,\chi}\) （在 \(\mathcal{F}_{\chi}(G)\) 或 \(L_{\chi}^{2}(G)\) 上的左正则表示）\ldots.. 417 \(\delta_{G,\chi}\) （在 \(\mathcal{F}_{\chi}(G)\) 或 \(L_{\chi}^{2}(G)\) 上的右正则表示）\ldots.. 417 \(\gamma_{G,\chi}^{(p)}\) （\(L_{\chi}^{p}(G)\) 中的左正则表示）\ldots.. 417 \(\delta_{G,\chi}^{(p)}\) （\(L_{\chi}^{p}(G)\) 中的右正则表示）\ldots.. 417 \(\gamma_{G,\chi}\) （\(L_{\chi}^{2}(G)\) 中的酉左正则表示）\ldots.. 417 \(\delta_{G,\chi}\) （\(L_{\chi}^{2}(G)\) 中的酉右正则表示）\ldots.. 417 \(\varrho_{G,\chi}\) （\(L_{\chi}^{2}(G)\) 中的酉双正则表示）\ldots.. 417 \(c_{Z}(\pi)\) （形式次数）\ldots.. 423 \(\mathrm{Noy}_{+}(\mathrm{X})\) （普遍正核）\ldots.. 434 \(A_{+}^{\prime}\) （连续正线性泛函）\ldots.. 436 \(\mathrm{Pos}(G)\) （正定函数）\ldots.. 443 \(\mathrm{Pos}_{1}(G)\) （满足 \(\varphi(e)=1\) 的正定函数）\ldots.. 443 \(\varrho_{G}(\pi)\) （\((\overline{\pi}\boxtimes\pi)\)-同型分量的 \(\varrho_{G}\)）\ldots.. 462 \(\Theta(G^{\sharp})\)，\(Z\Theta(G)\) （中心矩阵系数）\ldots.. 467 \(R(G)\) （格罗滕迪克环）\ldots.. 469 \(\mathrm{F}(\widehat{G})\) （End(Eπ) 的乘积代数）\ldots.. 470 \(\mathrm{F}_{b}(\widehat{G})\) （有界族的子代数）\ldots.. 470 \(\|u\|_{\mathrm{F}_{b}}\) （\(\mathrm{F}_{b}(\widehat{G})\) 中的范数）\ldots.. 470 \(\mathrm{F}_{0}(\widehat{G})\) （趋于0的族的子代数）\ldots.. 470 \(\overline{\mathcal{F}}_{G}(\nu)\) （ν 的傅里叶余变换）\ldots.. 471 \(\overline{\mathcal{F}}_{G}(f)\) （f 的傅里叶余变换）\ldots.. 471 \(\|u\|_{2}\) （End(Eπ) 中的范数）\ldots.. 471 \(\mathrm{F}^{2}(\widehat{G})\) （End(Eπ) 的希尔伯特直和）\ldots.. 472 \(\mathcal{F}_{G}(f)\) （f 的傅里叶变换）\ldots.. 473 \(\mathcal{C}(\widehat{G})\) （\(\widehat{G}\) 上的函数）\ldots.. 475 \(M^{1}(G^{\sharp})\) （中心有界测度）\ldots.. 475 \(\mathrm{FS}(\pi)\) （弗罗贝尼乌斯--舒尔指标）\ldots.. 476 \(M_{2k}(\varrho)\) （阶为 2k 的绝对矩）\ldots.. 476 \(S^{2}(E)\) （A-模的对称平方）\ldots.. 478 \(\wedge^{2}E\) （A-模的外平方）\ldots.. 478

\specialchapter{术语索引}

A 卡尔金代数，75 自然的，130 弗雷德霍姆择一定理，79 拉森的，476 弱混合映射，
322
紧线性映射，2 希尔伯特--施密特，23 有限秩的，4 有限迹的，164 由核定义的，25 核型的，169 严格的，2 逼近空间的---，14 逼近性质，14 自伴部分等距---，238 本质部分等距---，238 自守（表示---），405

B 基底（自然代数---），131 底部（谱的---部分），298 双交换子（---定理），326 双正则（表示---），406 博赫纳（---定理），455

有界（表示---），374 有界化，262 单位球，1 布里顿（---引理），506 C 函数演算，185 普遍可测的，272 卡尔金（---代数），75 特征标，374 中心的，390 亏指标偶的---，256 ---的亏指标，256 外平方，478 对称的，478 柯西（---测度），515 柯西--柯瓦列夫斯卡娅（---定理），343 中心特征标---，390 中心测度---，402 矩阵系数，386 有限维的，386 对角的，386 部分算子的核心，231 紧（线性映射---），2 完全连续，104，107 同型分量，394 边界条件，258 余幂零空间，45

共轭（表示---），382 余范数，61 切格常数，323 盖尔范德--奈马克--塞加尔构造，433 逆步（表示---），377 广义函数与函数的卷积，339 对应，224 余傅里叶变换，196，471 分布的，219 余傅里叶变换，471 盖尔范德对，512 卡日丹的，499 亏指标偶，256 外尔等分布判据，470 循环希尔伯特实现---，433，438，443 向量---，191

D 极分解，291 形式次数，423 广义函数的偏导数，208 对角自同态---，24 基中的对角自同态---，146 基中的对角部分算子---，234 直态射---，56 直向量子空间---，55 广义函数，203 紧支撑的，336 缓增的，215 部分算子的定义域，224 酉对偶，393

E 弗雷德霍姆自同态，41 乘法的，186 里斯的，46

有限迹的，23，164 对角的，24 基 B 中对角的，146 卡日丹集，499 预解式，244 薛定谔方程，431 积分的，80 等变（函数---），408 逼近空间，14 施瓦茨空间，209 索伯列夫空间，222 本质的（谱---），297 自伴扩张，255 弗里德里希斯的，295 几乎解析的，281 部分算子的扩张，226 F 层，204 部分算子的闭包，228 向量丛，36 忠实（表示---），374 多项式增长函数，214 中心的，402 勒夫纳的，326 施瓦茨的，209 负定型的，495 正定型的，442 插值的，327 代表的，386 球的，513 测试的，200 部分埃尔米特形式，291 正线性泛函，436 部分正的，291 关联的部分正的，291 闭的部分正的，292 不变半双线性的，380 形式的（---次数），423 形式对称（---微分算子），243 埃尔弗--舍斯特兰德公式，284 莱布尼茨公式，196，209

普朗谢雷尔的，474 斯通的，353 迹公式，488 弗雷德霍姆择一定理的---，79 ---的映射，41 ---的自同态，41 ---的映射的指标，43 弗里德里希斯（---扩张），295 弗罗贝尼乌斯--舒尔（---指标），476 G 盖尔范德--奈马克（---定理），442 盖尔范德--奈马克--塞加尔（---构造），433 盖尔范德--赖科夫（---定理），454 无穷小生成元，428 部分算子的图像，224 H 汉堡（---矩问题），359 高（谱的---部分），298 埃尔弗--舍斯特兰德（---公式），284 海林格--特普利茨（---定理），240 希尔伯特实现---，433，438，443 循环希尔伯特实现---，433，438，443 希尔--吉田（---定理），346 I 恒等（对合自构---），478 μ-本质像，181 弗罗贝尼乌斯--舒尔指标，476 指标，43 亏指标的，256 加藤不等式，344 不可约矩阵---，132 不可约表示---，378 表示的同构，375 同型的（分量---），394 K 加藤

---的不等式，344 加藤--雷利希（---定理），306 ---的卡日丹对，499 ---的卡日丹集，499 ---的卡日丹性质 {]}T{[}，498 凯斯滕（---定理），323 克赖因--鲁特曼（---定理），94

L 拉普拉斯算子，243 狄利克雷的，310 拉森（---择一定理），476 莱布尼茨（---公式），196 舒尔引理，12，386 卢伊（---定理），343 勒夫纳（---定理），330 外尔定律，311 洛莫诺索夫（---定理），13

M 不可约矩阵，132 本原的，138 可约的，132 极大–极小，154，300 默瑟（---定理），177 中心测度，402 柯西的，515 谱的，191，268 缓增的，216 极小–极大，154，300 下有界（部分算子---），238 绝对矩，476 测度的矩，359 ---的矩问题（汉堡的），359 ---的矩问题（斯蒂尔切斯的），364 穆尔--彭罗斯（---伪逆），320 表示的态射，375 直的，56 ---的乘法自同态，186 ---的乘法部分算子，231 重数，245 ρ 中的 π 的，398

谱的，82 有限谱的，82

N 自然（代数---），130 负的（负定型函数---），495 模 H 零测（集合---），
408
幂零空间，45，48 非退化（表示---），
373
正规（部分算子---），265 相对范数，306 部分算子的核，225 普遍正的，433 核型（线性映射---），
169
\lettergroup{0}

带指标的算子，41 带核的，25 紧的，2 弗雷德霍姆的，41 希尔伯特--施密特的，23 马尔可夫的，322 乘法的，186 有限秩的，4 里斯的，46 斯图姆--刘维尔的，350 有限迹的，23 沃尔泰拉的，118 微分的，235 标量微分的，235 转置微分的，242 带核的积分的，29

部分算子，224 定义域稠密的，227 具有紧预解式的，309 伴随的，235 自伴的，238 双射的，225 ---的核心，231 乘法的，231 ---的定义域，224 本质自伴的，238 ---的扩张，226 可闭的，227 闭的，227

---的闭包，228 ---的图像，224 在---下的像，225 在---下的原像，225 单射的，225 正规的，265 ---的核，225 正的，238 乘积的，226 逆的，226 ---的约化，226 相对有界的，306 表示正形式的，
292
---的预解式，244 和，226 ---的谱，244 满射的，225 对称的，238 正交性（---的关系），424，
458
\lettergroup{P}

相差局部零测集意义下相等的集合，425 佩龙--弗罗贝尼乌斯（---定理），

彼得--外尔（---定理），462 普朗谢雷尔（---公式），474 谱的敏感点，83 第一类正交多项式，361 第二类的，363 ---型的正定函数，442 部分算子---，238 正的（线性泛函---），436 本原的（矩阵---），138 极大–极小原理，154，300 极小–极大原理，154，300 酉表示的张量积，385 酉表示的外张量积，385 谱投影，54，277 逼近性质，14 卡日丹的 {]}T{[}，498

模 H 几乎处处为真，408 伪谱，250 Q 拟逆，40 典范的，48 拟同构，41 R 希尔伯特实现，433，438，443 循环希尔伯特实现，433，438，443 可约（矩阵---），132 约化，226 正则表示---右，405 表示---左，405 正交关系，424，458 Γ-自守表示，405 双正则的，406 有界的，374 共轭的，382 逆步的，377 平方可积的，420 模中心平方可积的，420 正交型的，476 四元数型的，476 实型的，476 辛型的，476 忠实的，374 诱导的，416 不可约的，378 ---的表示同构，375 线性的，374 连续线性的，374 ---的态射，375 非退化的，373 酉表示的张量积，385 张量幂，385 拟正则的，486 商的，376 右正则的，405 左正则的，405 ---对子群的限制，375 ---的副本之和，376

---的希尔伯特直和，酉，384 子---，375 平凡的，374 酉的，379 预解式，244 具有紧---的部分算子，309 表示的限制，375 S 绍德尔（---定理），9 薛定谔（---方程），431 舒尔（---引理），12，386 施瓦茨空间的---，209 ---函数，209 连续半群，345 半单（酉表示---），391 敏感的（谱的---点），83，297 斯梅尔（---定理），124 索伯列夫（---空间），222 分部求和，160 适应于 u 的子空间，153，299 坐标的，132 不变的，12 直向量，55 子表示，375 由一个子集生成的，376 谱测度---，268 ---的重数，82 有限重数---，82 ---投影，277 ---定理，194，266 谱，244 复的，54 连续的，245 本质的，83，297 ---的敏感点，83 剩余的，245 敏感的，297 稳定的，86 斯蒂尔切斯（---矩问题），364 斯通的---公式，353

---定理，428 特征值序列，152 广义函数的支撑，336 符号，127 对称（部分算子---），238

\lettergroup{T}

博赫纳定理，455 博苏克--乌拉姆定理，64 柯西--柯瓦列夫斯卡娅定理，343 盖尔范德--奈马克定理，442 盖尔范德--赖科夫定理，454 海林格--特普利茨定理，240 希尔--吉田定理，346 加藤--雷利希定理，306 凯斯滕定理，323 克赖因、克拉索塞利斯基和米尔曼定理，64 克赖因--鲁特曼定理，94 卢伊定理，343 勒夫纳定理，330 洛莫诺索夫定理，13 默瑟定理，177 佩龙--弗罗贝尼乌斯定理，101 彼得--外尔定理，462 椭圆正则性定理，341 雷利希--孔德拉绍夫定理，335 绍德尔定理，9 斯梅尔定理，124 斯通定理，428 冯·诺伊曼定理，258 双交换子定理，326 施瓦茨核定理，342 谱定理，194，266 特普利茨（---算子），127 弱拓扑，6 C \(^{r}\)-紧收敛的，34 弱的，449 T-预解式，285 迹

有限迹映射，164 迹公式，488

傅里叶变换，196，473 盖尔范德变换，190

傅里叶变换，473 T-预解式（拓扑---），285 平凡（表示---），374

U

自然含单位代数

酉对偶---，393 ---表示的张量积，385 表示---，379 半单表示---，391 表示---的希尔伯特直和，384

逼近单位元

普遍可测函数演算---，272 函数---，183

V

柯西主值，337 特征值，245 特征值（---族），146，234
234
奇异值，151，156 循环向量，191，376 G-有限，376 沃尔泰拉（---算子），118 冯·诺伊曼的---定理，258 ---的双交换子定理，
326
\lettergroup{W}

外尔等分布判据---，470 ---不等式，162 ---定律，311

\specialchapter{自伴性判据}

这里汇总正文中陈述的、可用来证明某些部分算子自伴的判据。

固定一个复希尔伯特空间 E 以及 E 上的部分算子 u。在下列每一种情形中，该算子都是自伴的：

\begin{enumerate}
\def\labelenumi{(\roman{enumi})}
\item
  存在一个单射且为埃尔米特的 \(v \in \mathcal{L}(\mathrm{E})\) ，使得 \(u = v^{-1}\) 。（IV，第239页命题9）；
\item
  算子 \(u\) 是对称的，且 \(\mathrm{dom}(u) = \mathrm{E}\) （IV，第240页命题10）；
\item
  算子 \(u\) 是对称的，并诱导从 \(\mathrm{dom}(u)\) 到 E 的双射（IV，第240页命题10的推论）；
\item
  存在 \(\lambda \in \mathbf{C}\) ，使得 \(u + \lambda 1_{\mathrm{E}}\) 和 \(u + \overline{\lambda} 1_{\mathrm{E}}\) 均为满射（IV，第240页命题11）；
\item
  存在一个定义域稠密的闭部分算子 \(v\) ，使得 \(u = v^{*} \circ v\) （IV，第241页命题12）；
\item
  算子 \(u\) 是对称且闭的，并且其亏指标为 \((0,0)\) （IV，第257页命题26的推论），也就是说，\(\mathrm{Ker}(u^{*} - i) = \mathrm{Ker}(u^{*} + i) = \{0\}\) ；
\item
  空间 E 具有 \(\mathrm{L}^{2}(\mathrm{X},\mu)\) 的形式，其中 X 是局部紧拓扑空间，\(\mu\) 是 X 上的测度，而 u 是乘法算子，即乘以一个在 X 上 \(\mu\) -可测且局部地 \(\mu\) -几乎处处取实值的函数（命题23的推论）；
\item
  存在 E 上的正规部分算子 v，以及一个取实值的普遍可测函数 \(f \in \mathcal{L}_{\mathrm{u}}(\mathrm{Sp}(v))\) ，使得 \(u = f(v)\) （IV，第273页推论1）；
\item
  存在 E 上的部分埃尔米特形式 q，使 u 是表示 q 的部分算子（IV，第293页命题14）；
\item
  有 \(u = v + w\) ，其中 v 是自伴部分算子，w 是关于 v 相对有界且满足 \(\|w\|_{v} < 1\) 的对称部分算子（IV，第306页定理5）；
\item
  算子 \(u\) 是 E 上的酉表示 \(\varrho\) （其群为 \(\mathbf{R}\)）的无穷小生成元（V，第428页定理1）。
\end{enumerate}

\end{document}
